% 本文件由 LaTeX 排版 Agent 根据有效 translation.json 自动生成。
% author=Augustine of Hippo; work_id=1201; title=天主之城

\chapter{天主之城}

% structure: p0001 source=h1 target=omitted reason=page-title-already-rendered-by-parent

\Emph{第一卷} 奥古斯丁谴责异教徒，他们将世界的灾难，尤其是最近罗马被哥特人洗劫，归咎于基督教及其禁止崇拜诸神。他谈到生活的福祸，这些当时与以往一样，同样临到善人和恶人身上。最后，他斥责那些不知羞耻的人，他们责难基督徒，说他们的妇女被士兵玷污。\par

\Emph{第二卷} 在这卷中奥古斯丁回顾了在基督之前，且遍行假神崇拜之时，罗马人所遭受的那些灾难；并证明，罗马人非但未得诸神庇佑免于不幸，反而被他们用那唯一，或至少是最大的灾难——道德堕落与灵魂的邪恶——所压倒。\par

\Emph{第三卷} 正如在前一卷中奥古斯丁证明了道德和灵性上的灾难，在本卷中他证明了外在与身体上的祸患：自建城以来，罗马人便不断遭受这类祸患；即便在假神独受崇拜、基督尚未降来之前，他们也不曾从这样的灾祸中得到解脱。\par

\Emph{第四卷} 本卷证明，罗马帝国的疆域辽阔与长久延续，不应归功于\Person{朱庇特}或异教徒的诸神，因为人们相信连最微小的事物和最卑下的功能都单独托付给了他们；而应归功于唯一的真天主，幸福的主宰，世间的王国都由他的权能和裁断得以建立和维持。\par

\Emph{第五卷} 奥古斯丁首先讨论命定论，为要驳斥那些倾向于把罗马帝国的强权和扩张归诸命运的人；如前卷所示，这绝不能归功于假神。随后，他证明天主的预知与我们的自由意志并无矛盾。接着，他论及古罗马人的风尚，并说明罗马人自身的德性在何种意义上，以及在天主的旨意中，使得他们的版图得以扩张，尽管他们并不敬拜他。最后，他阐述了基督徒皇帝真正的幸福是什么。\par

\Emph{第六卷} 此前的论证是针对那些认为应为现世利益而崇拜诸神的人；现在则转向那些相信应为永生而崇拜诸神的人。奥古斯丁用随后五卷来驳斥后一种信念，并首先表明，就连异教神学中最受推崇的作家\Person{瓦罗}本人对诸神的看法也是何等低劣。奥古斯丁采纳瓦罗将这种神学区分为三类——神话的、自然的和民间的——并随即证明，神话的和民间的神学均无助于来生的幸福。\par

\Emph{第七卷} 本卷表明，永生不能藉着敬拜\Person{雅努斯}、\Person{朱庇特}、\Person{萨图恩}以及民间神学中其他\Quote{拣选的神明}而获得。\par

\Emph{第八卷} 奥古斯丁现在进入第三类神学，即自然神学，并探讨崇拜自然神学的诸神是否有助于获得来世的真福。他选择与柏拉图学派讨论这个问题，因为柏拉图体系在哲学中\Quote{\OriginalTerm{显居首位}{facile princeps}}，且与基督教真理最为接近。在论证中，他首先驳斥\Person{阿普列乌斯}，以及所有主张应崇拜魔鬼作为神明与人类之间的信使和中保的人；证明人绝不可能通过魔鬼与善神和好，因为魔鬼是罪恶的奴隶，他们喜好并赞助那些善良智慧之人憎恶并谴责的事物——亵渎神明的诗人虚构、戏剧表演和巫术。\par

\Emph{第九卷} 在前一卷中已表明必须弃绝魔鬼的敬拜，因为他们以千百种方式表明自己是邪恶的灵；本卷中奥古斯丁回应那些主张魔鬼中有区别的人，即有些是恶的，另一些是善的；在驳倒这一区分之后，他证明，为人类提供永恒真福的职分不属于任何魔鬼，而只属于基督。\par

\Emph{第十卷} 在这卷中奥古斯丁教导说，善良的天使只希望，以祭献所呈奉的、称为\Quote{\OriginalTerm{敬拜}{latreia}}的神圣尊荣，归于他们所事奉的天主。然后他转而驳斥\Person{波斐留斯}关于灵魂净化和解脱的原则与方式。\par

\Emph{第十一卷} 这里开始本著作的第二部分，论及两座城——地上之城和天上之城——的起源、历史和归宿。首先，奥古斯丁在本卷中说明这两座城最初是如何由善天使和恶天使的分离而形成的；并借此机会讨论了世界的创造，正如\WorkTitle{圣经}中\BibleBook{}{创世纪}的开篇所描述的那样。\par

\Emph{第十二卷} 奥古斯丁首先就天使提出两个问题：亦即，为何有些天使里有善意，而另一些里却有恶意？以及，善者的福乐和恶者的苦难的原因是什么？然后他论及人的创造，并教导说，人并非从永恒而来，乃是受造的，且唯独由天主所造。\par

\Emph{第十三卷} 本卷教导说，死亡是惩罚，并且起源于亚当的罪。\par

\Emph{第十四卷} 奥古斯丁再次论及第一个人的罪，并教导说，它是人的肉欲生活和邪恶情感的原因。他特别证明，伴随情欲的羞耻是那种悖逆的公正惩罚，并探询如果人没有犯罪，将如何能够没有情欲而繁衍后代。\par

\Emph{第十五卷} 在前四卷中，奥斯定论述了两座城的起源，即地上之城和天上之城；随后在接下来的四卷中，他阐述它们的成长与进展；为此，他阐释了圣史中与此主题有关的主要段落。在这第十五卷中，他通过解说 \BibleBook{创世}{纪} 中从加音与亚伯尔直到洪水的事迹，开启了其著作的这一部分。\par

\Emph{第十六卷} 本卷前部，从第一章至第十二章，根据圣经展示了从诺厄到亚巴郎、地上之城与天上之城的进展；后部则以仅天上之城从亚巴郎到以色列诸王的进展为主题。\par

\Emph{第十七卷} 本卷追溯了从撒慕尔到达味、直至基督的列王与先知时期的天主之城的历史；并诠释了记载于 \BibleBook{列王}{纪}、\BibleBook{圣咏}{集} 和 \WorkTitle{撒罗满的著作} 中的有关基督和教会的预言。\par

\Emph{第十八卷} 奥斯定追溯了地上之城与天上之城从亚巴郎直到世界终末并行发展的轨迹；并提及有关基督的神谕，既有西彼拉们所说的，也有那些在罗马建城后写作的神圣先知——欧瑟亚、亚毛斯、依撒意亚、米该亚及他们的继承者——的预言。\par

\Emph{第十九卷} 本卷讨论了两座城——地上之城与天上之城——的终向。奥斯定回顾了哲学家们关于至善的观点，以及他们力图在此生为自己寻求幸福的徒劳努力；在驳斥这些观点的同时，他借此机会指出，属于天上之城或基督子民的和平与幸福，无论是在现时还是将来，究竟为何。\par

\Emph{第二十卷} 论最后的审判，以及 \WorkTitle{旧约} 与 \WorkTitle{新约} 中有关它的宣告。\par

\Emph{第二十一卷} 论魔鬼之城所注定的终局，即受永罚者的永罚；以及不信对此提出的论点。\par

\Emph{第二十二卷} 本卷论述天主之城的终向，即圣人们的永福；确立并解释了身体复活的信德；全书结束时则展示，圣人们身披不朽的精神的身体，将如何享用其福乐。\par

\SourceRecord{Translated by Marcus Dods.；Nicene and Post-Nicene Fathers, First Series；Vol. 2.；Edited by Philip Schaff.；Buffalo, NY: Christian Literature Publishing Co.,；1887.；<http://www.newadvent.org/fathers/1201.htm>.}

% structure: p0001 source=h1 target=section reason=page-title

\section{\WorkTitle{天主之城（第一册）}}

奥古斯丁谴责异教徒，他们将世间的灾难，尤其是最近哥特人对\Place{罗马}的洗劫，归咎于基督教及其禁止崇拜诸神的规定。他谈到生活中的福祉与灾祸，这些福祸无论何时都同样临于善人与恶人。最后，他斥责那些基督徒妇女被士兵强暴却反过来羞辱基督徒的人厚颜无耻。\par

% structure: p0003 source=h2 target=subsection reason=relative-heading-rank; base=h2

\subsection{序言：解释他从事这项工作的计划。}

光荣的天主之城是我在这部著作中的主题，这是你，我最亲爱的儿子\Person{马尔切利努斯}，提议的，也是我因承诺而应归于你的。我承担为它辩护，以反对那些宁愿信奉自己的神明而不信奉此城缔造者的人——这座城的光荣无与伦比，无论我们看它如何在时日短促的流转中仍凭\OriginalTerm{信德}{faith}而生活，并在不虔敬者中作旅居的过客，还是看它将如何安居于永恒居所的稳固安宁中，如今它正耐心等候，期待\Quote{公义将归于审判}，并凭借自身的卓越，获得最终的胜利与完美的和平。这是一项伟大而艰辛的工作；但天主是我的助佑。我深知需要何等的能力才能说服骄傲的人，\OriginalTerm{谦卑}{humility}之德何等伟大，它并非凭人的傲慢，而是赖天主的\OriginalTerm{恩宠}{grace}，将我们提拔到高于这变幻场景中摇摆不定的一切世俗尊严之上。因为我们所谈论的这座城的君王与缔造者，已在圣经中向祂的子民宣布了天主律法的这句话：\Quote{天主拒绝骄傲人，却赐恩宠予谦卑人。}但这也是天主的特权，骄傲心灵的膨胀野心也妄加觊觎，极其喜欢将此列为其属性之一：\par

\Quote{对被贬抑的灵魂显示怜悯，\newline{}压碎骄傲的子孙。}\par

因此，根据我们承担的这项工作的计划要求，并随机会的出现，我们也必须谈论地上之城，虽然它是万国的女主人，却受其统治欲所统治。\par

% structure: p0007 source=h2 target=subsection reason=relative-heading-rank; base=h2

\subsection{第1章 基督名字的敌人，蛮族因基督之故在洗劫该城时赦免了他们。}

因为属于这地上之城的，是我必须为其辩护天主的城所要反对的敌人。他们中有许多人的确从亵渎的谬误中回头，成了此城相当可信赖的公民；但也有许多人对它燃满仇恨，对其救主显明的恩惠如此忘恩负义，竟忘了他们如今若没有在逃避敌人刀剑时，在他们现在自夸得以保全性命的神圣场所找到避难之所，他们现在就不可能说出半句贬低它的话。难道那些因蛮族尊重基督而被赦免的罗马人，如今竟成了敌基督之名的吗？殉道者的\OriginalTerm{圣髑匣}{reliquaries}与宗徒们的教堂为此作证；因为在城市遭洗劫时，这些地方是向所有投奔者——无论是基督徒还是异教徒——敞开的避难所。嗜血的敌人一直肆虐到他们的门槛；在那里，他那杀人的暴怒才承认了界限。敌人中若有怜悯之心的人，便将他们手下留情的人带到那里，以免可能遭遇不那么仁慈的人。事实上，当那些在其他各处都显得冷酷的杀人者来到这些场所，这些禁止战争在其他各处所许可的暴行的地方时，他们那狂暴的屠杀之欲受到了约束，他们捉拿俘虏的渴望也熄灭了。就这样，众多的人逃脱了，如今他们却指责基督宗教，将降临于其城的灾祸归咎于基督；但他们自己生命的保全——这得归功于蛮族对基督持有的尊重——他们不归因于我们的基督，却归于自己的好运。如果他们有正确的判断，就应反把敌人加诸的严酷与苦难归于\OriginalTerm{天主眷顾}{divine providence}，那眷顾惯于藉惩戒改革人们堕落的品行，也以类似的磨难操练义人和堪受赞美的人——或是当他们在考验中经过之后，将他们转移到更好的世界，或是为更深远的目的让他们仍留在地上。同时，他们也当把这归因于这些基督教时代的精神：这些嗜血的蛮族，违背战争惯例，因着基督的缘故赦免了他们，无论这怜悯是普遍显于各处，还是特别显于那些奉献给基督名字的场所，其中最大的被选作避难所，好让那广施怜悯的渴望得以充分满足，使广大群众可在其中找到庇护。因此，他们理应感谢天主，并以真诚宣认投奔祂的名下，好使他们能逃脱永火的惩罚——这些口是心非的人当初假托此名，才逃脱了眼前的毁灭。因为你们看到那些傲慢无耻地侮辱基督仆人的人中间，有许多人当初若没有假装是基督的仆人，便不能逃脱那毁灭与屠杀。然而如今，他们怀着忘恩负义的骄傲与最邪恶的疯狂，冒着在永恒黑暗中受罚的危险，执拗地反对那个他们当初为贪享此生短暂光明而欺诈地赖以保护自己的名字。\par

% structure: p0009 source=h2 target=subsection reason=relative-heading-rank; base=h2

\subsection{第2章 胜利者因所奉神明之故而赦免被征服者，这完全违背战争惯例。}

有无数战争的史书，既包括\Place{罗马}建城以前，也包括它兴起并扩张版图以后；且读这些史书，找出一个例子，在其中，一座城被异族攻陷时，胜利者赦免了那些被发现逃到他们神明庙宇中寻求庇护的人；或某个蛮族将军下令，凡在这座或那座神殿中被发现的人，谁也不得杀害。难道\Person{埃涅阿斯}没有见过\par

\Quote{垂死的\Person{普里阿摩斯}在圣坛前，\newline{}玷污他亲自祝圣的祭坛？}\par

\Quote{\Person{狄俄墨得斯}与\Person{尤利西斯}，\newline{}拖着血红的双手，杀死卫兵，\newline{}将你命运的神像从庙中拖出，\newline{}触摸她贞洁的发辫，并用鲜血染污\newline{}她所佩戴的童贞花冠？}\par

接下来说的也非事实，\newline{}自此国运转向，\newline{}希腊日渐式微。\par

此后他们以火与剑征服并毁灭了\Place{特洛伊}；此后当\Person{普里阿摩斯}逃往祭坛时被他们斩首。\Place{特洛伊}的灭亡也不是因为失去了\Person{密涅瓦}。因为\Person{密涅瓦}自己先失去了什么，竟让她自己灭亡？也许是她的守卫？无疑；正是她的守卫。因为一旦守卫被杀，她就能被偷走。事实上，并非那像保护了人，而是人保护了那像。她连自己的守卫都无法保护，又怎能被呼求来保卫城与公民呢？\par

% structure: p0015 source=h2 target=subsection reason=relative-heading-rank; base=h2

\subsection{第三章 —— \Place{罗马}人并未表现出他们通常的睿智，因为他们竟信赖那些无能保卫\Place{特洛伊}的神祇会施恩于他们。}

罗马人竟然乐于将他们的城托付给这些神祇保护！啊，何等可怜的错误！当我们如此谈论他们的神祇时，他们向我们发怒，然而，他们非但不向自己的作家发怒，反而花钱学习他们所说的；事实上，这些作家的教师还被认为应得公帑薪酬和别的荣誉。这里有\Person{维吉尔}，男孩子们诵读他，为使这位大诗人，这位最著名且最受推崇的诗人，能以其思想浸润他们童贞的心灵，并不易被遗忘，正如\Person{贺拉斯}那句名言所说：\par

\Quote{新桶长留初味。}\newline{}那么，我且说，在这位\Person{维吉尔}笔下，\Person{尤诺}被表现为与特洛伊人为敌，并挑唆风王\Person{埃俄洛斯}攻击他们，她说：\newline{}我所憎恨的种族正破海航行，\newline{}将\Place{特洛伊}迁往\Place{意大利}，\newline{}带着被征服的家神。…\par

明智的人应该把罗马的防御托付给这些被征服的神祇吗？但有人会说，这只是\Person{尤诺}的话，她像一个愤怒的女人，不知道自己在说什么。那么，那位常常被冠以\Quote{虔敬}之名的\Person{埃涅阿斯}本人又说了什么呢？他不是说：\par

看哪！\Person{庞托斯}从死亡中飞逃脱身，\newline{}高处的\Person{阿波罗}祭司，\newline{}怀抱着那被征服的神祇，双手发颤，\newline{}急促地寻求庇护？\newline{}那么，岂非显而易见：那些神祇（他毫不犹豫地称之为\Quote{被征服的}）更像是被托付给\Person{埃涅阿斯}，而非他托付于神祇？因为有人对他说：\newline{}你祖国将她家龛中的神祇，\newline{}交你照管？\par

那么，如果\Person{维吉尔}说这些神祇竟是如此，并且被征服了，而且被征服时若无一人保护便无法脱逃，那么，以为罗马曾明智地托付给这些守护者，若非失去它们罗马便不会被攻克，这该是何等的疯狂！实际上，崇拜被征服的神祇，将之当作保护者与斗士，这不过是在崇拜不祥的徵兆，而非良善的神明。难道不是以下的想法更为明智：并非因为神祇先灭亡了，罗马才绝不会陷入如此大灾；而是若非罗马尽可能长久地保全它们，它们早已灭亡了？因为，有谁稍加思索还看不出，这是一个多么愚蠢的假设：它们在战败的守护者之下不可能被征服，它们灭亡仅仅是因为失去了守护的神祇——事实的惟一原因是，它们选定了注定要灭亡的神祇作为保护者。因此，诗人们吟咏这些关于被征服之神祇的诗句时，并非有意虚构谎言，而是像诚实的人一样，说出了真理由他们强索出的实情。不过，这一点我们将在另一处更合适的地方加以仔细而充分的讨论。眼下，我将就我所能，简要说明我对这些忘恩负义之人的看法：他们因自己邪恶的行径理当遭受灾祸，却亵渎地归咎于\Person{基督}；而因\Person{基督}之名，他们的恶行竟获幸免的事，他们甚至不屑注意；他们在疯狂亵渎的傲慢中，用自己的嘴唇攻击祂的圣名——那正是他们为活命曾虚情假意求告的同一圣名。在那些为\Person{基督}而祝圣的地方，凡为祂的缘故，敌人不曾伤害他们；在那里，他们曾勒住舌头，以保安全并受庇护；然而，一旦离开这些避难所，他们就放纵舌头，向祂投掷充满仇恨的诅咒。\par

% structure: p0021 source=h2 target=subsection reason=relative-heading-rank; base=h2

\subsection{第四章 —— \Place{特洛伊}\Person{尤诺}的避难所并未从\Place{希腊}人手中救下任何人；而宗徒们的教会则保护了所有逃往其中的蛮族。}

\Place{特洛伊}本身，\Place{罗马}民族之母，正如我说过的，虽则\Place{希腊}人崇拜同样的神祇，却未能在其诸神的圣所中保护自己的公民免遭\Place{希腊}人的火与剑。不仅如此，而且\par

\Person{菲尼克斯}与\Person{尤利西斯}被派定，\newline{}在\Person{尤诺}寝旁的空庭中，\newline{}看守掠物；\newline{}自焚燃的殿宇抢夺而来，\newline{}\Place{伊里昂}的巨量财宝藏于此处，\newline{}丰盛的祭坛、巨金的碗，\newline{}与掳来的衣服，胡乱卷作\newline{}杂乱的一堆；\newline{}同时，童男与主妇，惊惶失措，\newline{}排成长列，站于近旁。\par

换句话说，选定这个奉献给如此伟大女神的地方，并非为了不让任何人被掳走，而是要让所有俘虏都被关押其中。现在请比较一下这个\Quote{\OriginalTerm{避难所}{asylum}}——它既非寻常神祇的避难所，也非等闲小神的避难所，而是\Person{朱庇特}亲姐妹兼妻子、众神之后的避难所——与为纪念宗徒们而建的教堂。前者收聚了从燃烧的神庙中抢救出来、从诸神手中夺来的战利品，不是为了将其归还被征服者，而是分配给征服者；而在后者，凡是属于教堂的物件，即使是在别处寻回，也都要以最虔敬的态度和最高的敬意送归原处。在前者，自由丧失；在后者，自由得以保全。前者充斥严酷的奴役，后者则严拒奴役。人们被赶入那座庙宇，成为敌人的财物，任凭他们摆布；而人们被他们的敌人怀着恻隐之心引进这些教堂，好能重获自由。总之，温和的希腊人将\Person{朱诺}的神庙用以满足自己的贪婪和傲气；而这些\OriginalTerm{基督}{Christ}的教堂，甚至为野蛮的蛮族选作施行谦卑与慈悲的合宜场所。然而，或许希腊人在那场胜利中毕竟没有破坏那些他们与特洛伊人共同崇拜的神明之庙，也未敢将逃到那里的悲惨被征服特洛伊人加以杀戮或掳掠？或许\Person{维吉尔}以诗人的笔法描绘了从未真正发生过的事？但毫无疑问，他所描绘的乃是敌人在劫掠城市时的惯常做法。\par

% structure: p0025 source=h2 target=subsection reason=relative-heading-rank; base=h2

\subsection{\Person{凯撒}关于敌人劫掠城市时普遍惯例的陈述}

关于这一惯例，就连\Person{凯撒}本人也为我们提供了确凿的证词；因为他在元老院就阴谋者发表意见时，说道（正如极为忠实的历史学家\Person{撒路斯提乌斯}所记）\Quote{处女和男孩遭到玷污，婴孩从父母怀中夺走，主妇听凭征服者任意处置，庙宇和房屋被劫掠，杀戮与纵火遍地；总之，到处充斥着武器、尸体、鲜血和哀号。}如果他在这里没有提到庙宇，我们或许会以为敌人惯于饶过诸神的住所。而罗马的庙宇曾面临的这类灾难，并非来自外敌，而是来自\Person{喀提林}及其同伙，即罗马最高贵的元老和公民。但人们或许会说，这些人是堕落之辈，是祖国之逆子。\par

% structure: p0027 source=h2 target=subsection reason=relative-heading-rank; base=h2

\subsection{即便罗马人攻取城市时，也未曾饶过躲藏在庙宇中的被征服者}

那么，我们的论证又何必提及那些彼此交战的众多民族，他们从未在对方的神庙中饶过被征服者呢？让我们来看看罗马人自身的做法吧；我再说一遍，让我们回顾并审视罗马人，他们最引以为豪的是 \Quote{饶恕被征服者，并挫败骄傲者}，并且他们宁可 \Quote{宽恕伤害，而不加以报复}；在他们为了扩张疆土而攻陷、占领、摧毁的那么多伟大城市中，请告诉我们，他们曾习惯于豁免哪些庙宇，使逃入其中者得以自由？抑或他们确实做过此事，而史家却对此事隐而不提？难道能相信，那些极为热切地搜求可赞美事迹之人，竟会遗漏他们自己眼中最卓著的虔诚证据吗？\Person{马尔库斯·马凯路斯}，一位卓越的罗马人，他攻占了装饰极为华美的城市\Place{叙拉古}，据称他曾为其即将降临的毁灭而哀叹，甚至在造成流血之前就已为之挥泪。他还采取措施保全敌人的贞洁。因为在下令攻城之前，他颁布了敕令，禁止侵犯任何自由人。然而，该城仍按战争惯例遭到劫掠；我们也从未读到，即便由这样一位贞洁而温和的统帅下令，仍规定不得伤害逃入某座庙宇者。这一点肯定不会被漏记，因为他的泪水和他保全贞洁的敕令都未能被略而不提。\Person{费边}，\Place{塔伦图姆}城的征服者，因未将神像作为战利品而受到赞扬。当他的秘书问他，如何处置大量缴获的众神雕像时，他以一句玩笑掩饰了自己的克制。他问那些神像是何种模样；当人们报告说不仅有许多巨大的神像，而且有些还持有武器时，他说：\Quote{哦，让我们把愤怒的众神留给塔伦图姆人吧。} 可见，既然罗马史家不会默然不提一位将军的哭泣和另一位将军的玩笑，不会不提这位的贞洁怜悯和那位的风趣克制，那么，倘若他们曾为了尊敬敌方某位神明，展示出某种特定的宽仁，即在任何庙宇中禁止杀戮或掳掠，又怎会被漏记呢？\par

% structure: p0029 source=h2 target=subsection reason=relative-heading-rank; base=h2

\subsection{罗马劫掠中所发生的残忍行为符合战争惯例，而仁慈之举则源于\OriginalTerm{基督}{Christ}之名的感化力}

\Place{罗马}在最近这场灾难中所遭受的一切劫掠——一切杀戮、抢夺、焚烧和苦难——都是战争习俗的结果。然而新奇的是，野蛮的蛮族却表现得如此温和：他们挑选并保留最大的教堂，用以收容那些得到宽恕的民众；教堂里面没有一人被杀，没有一人被强行拖走；许多俘虏被发了善心的敌人领往那里获得释放；没有一人被无情的敌人带去做奴隶。凡看不出这事当归功于基督之名与基督徒精神的人，是瞎眼的；凡看出这事却不赞颂的人，是忘恩负义的；凡阻止别人赞颂此事的人，是疯狂的。但凡明智之人，绝不会将这种宽仁归于蛮族。他们那残暴嗜血的心，被那许久以前藉著先知说\BibleQuote{我必用杖惩罚他们的过犯，用鞭责罚他们的罪孽；只是我必不将我的慈爱全然收回}的那位所震慑、遏制并奇妙地节制。\par

% structure: p0031 source=h2 target=subsection reason=relative-heading-rank; base=h2

\subsection{第八章 论不分善恶而临于善人与恶人的利与害}

可能有人说：\Quote{那么，为何这神圣的慈悯竟也临及那不虔敬和忘恩负义的人呢？}答案是，这正是那位天天\BibleQuote{他使太阳上升，光照恶人，也光照善人；降雨给义人，也给不义的人}\BibleRef{玛 5:45}的天主的慈爱。因为虽然这些人中有的因思想此事而悔罪改过，但正如宗徒所说的，有些人\BibleQuote{竟轻视天主丰厚的慈爱、宽容与忍耐，因执拗不悔改的心，为自己积蓄愤怒，以待愤怒和天主正义审判显现的日子，那时，他要照每人的行为予以报应}\BibleRef{罗 2:4}，然而天主的忍耐仍在邀请恶人悔改，正如天主的鞭笞教育善人忍耐一样。同样，天主的仁慈也怀抱善人，为要爱护他们；而天主的严厉也阻止恶人，为要惩罚他们。依天主的上智安排，似乎定意在来世为义人预备善事，这是不义的人不得享受的；也为恶人预备恶事，这是善人不受折磨的。至于今生的福乐与灾祸，天主却愿意两者共有；好使我们不至于过分贪求那些连恶人也同样享受的东西，也不至于以不适当的畏惧退缩于那些连善人也常遭遇的灾祸。\par

而且，那些我们称为逆境和顺境的事件，其作用也大不相同。因为善人不因世间的顺境而趾高气扬，也不因世间的逆境而灰心丧志；恶人却因世福而腐败堕落，因世祸而自感受罚。然而，即便在现今对世事的分施中，天主也常明白昭示他的干预。因为，若每一项罪过如今都遭到显明的惩罚，似乎就没有什么留给最后的审判了；相反，若任何罪过如今都不受明显的神性惩罚，人们便会断定根本没有天主的眷顾。至于今生的福利，也是如此：天主若不仁慈地显明赏赐某些祈求的人，我们便会说这些福利并不掌握在他手中；他若赏赐给所有求的人，我们便会以为侍奉他只得到这样的报酬；这样侍奉，不使我们成为虔诚，反使我们成为贪婪而不知足。因此，虽然善人也如恶人一样受苦，但我们不能因他们所遭受的相同，就以为他们之间没有差别。因为即使在相同的苦难中，受苦者也仍有不同；虽遭受同样的痛苦，德行与恶行却并非同一回事。因为，同样的火使金子熠熠生辉，却使糠秕冒烟；同样的连枷下，禾秸被打得粉碎，谷粒却被洁净；同样的压榨下，酒糟与油不相混合，同样，同样的苦难暴力证明、炼净、澄清善人，却定罪、毁灭、铲除恶人。所以，在同一苦难中，恶人憎恨天主并亵渎，善人却祈祷并赞颂。造成如此巨大的差别，不在于遭受什么祸患，而在于遭受祸患的是什么样的人。因为，同样的搅动使污泥发出恶臭，而使香膏散发芬芳。\par

% structure: p0034 source=h2 target=subsection reason=relative-heading-rank; base=h2

\subsection{第九章 论对善人与恶人一同施予纠正的原因}

那么，在那个灾难频仍的时期，基督徒究竟受了什么苦，以至于凡正确而忠诚地思考以下情形的人，都能从中获益呢？首先，他们必须谦卑地反省那些激怒\Person{天主}、使世界充满如此可怕灾难的罪过；因为尽管他们远非邪恶、不道德和不敬虔之人的极端恶行，但他们也不该认为自己如此清洁无瑕，以至于善到连这些现世的灾祸都不配受。因为每一个人，无论生活多么堪称赞许，在有些事情上仍会屈服于\OriginalTerm{肉欲}{lust of the flesh}。他虽然未陷入穷凶极恶的罪行、肆无忌惮的堕落和可憎的亵渎，却仍会犯一些罪，或偶尔为之，或越是不以为意的罪，犯得越频繁。不过，暂且不说这一点，我们又能轻易找到一个人，对那些人——正是因那可憎的骄傲、奢华、贪婪，以及可咒的不义与不虔，\Person{天主}如今按祂的预言惩罚大地——持有恰当而公道的评价呢？谁又按着我们应该与他们相处的方式，与他们一同生活呢？因为我们常常邪恶地自蔽双目，不肯寻找教导劝诫他们的机会，有时甚至不肯责备指责他们，这或是因为我们畏难退缩，或是因为羞于得罪他们，或是因为害怕失去良好的友谊，唯恐妨碍我们的晋升，或在某种世俗之事上损害我们——这事或是因我们贪婪的性情想要获得，或是因我们软弱怕失去。因此，虽然恶人的行为令善人厌恶，所以善人不会与他们同陷于那来生等待着这类人的\OriginalTerm{永罚}{damnation}；但他们因为惧怕而姑息恶人那该受永罚的罪过，所以即使他们自己的罪轻微且属\OriginalTerm{小罪}{venial}，他们在今生也理当与恶人一同受鞭笞，尽管在永恒中他们完全免于惩罚。这是公正的：当\Person{天主}让他们与恶人一同受苦时，他们便觉得此生苦涩；正因贪爱此生的甜美，他们才不愿以苦言对待那些罪人。\par

若有人因寻求更合适的时机，或因担心自己的斥责会使他们变得更坏，或其他软弱的人可能因此灰心，不再努力度良善而虔诚的生活，甚至被逐出\OriginalTerm{信德}{faith}，而不去责备与挑剔那些作恶的人；这种人的疏忽似乎并非出于\OriginalTerm{贪婪}{covetousness}，而是出于\OriginalTerm{爱德}{charity}的考量。但应受指责的是，那些自己远离恶人的行为，且生活作风截然不同的人，却姑息别人身上的那些他们本该加以斥责并使之戒除的过错；他们之所以姑息，是因为怕得罪人，唯恐损害自己在那些善人可以清白而合法享用的事物上的利益——尽管他们享用起来，比那些在此世为旅客、并宣认盼望\OriginalTerm{天乡}{heavenly country}的人所应有的态度更为贪婪。因为不仅那些享受婚姻生活、有儿女（或渴望有儿女）、拥有房产和家业的软弱弟兄——宗徒在\OriginalTerm{教会}{churches}中向他们讲话，警告并教导他们应如何生活：妻子待丈夫，丈夫待妻子，儿女待父母，父母待儿女，仆人待主人，主人待仆人——不仅这些软弱的弟兄因许多现世短暂的事物而乐于获得、苦于失去，并因这些事物而不敢得罪那些其污秽邪恶的生活令他们大为不悦的人；就连那些生活水准更高、未受婚姻生活的缠累、只用清贫衣食的人，也常常顾念自己的安全和美名，因惧怕恶人的诡计和强暴而不去指责他们。而且，他们虽然不怕到被引去做同样恶事的地步，也不惧怕任何威胁或强暴；可是，对于他们自己不肯参与的恶事，他们却常常不肯去指责，而指责或许正能阻止这些恶事的实施。他们不加干涉，是因为担心如果指责达不到良好效果，自己的安全或名誉可能受损或毁灭；并不是因为他们看到自己的保全和美名是必要的，以便能影响那些需要他们教导的人，而是因为他们懦弱地贪图别人的奉承与尊敬，惧怕众人的评断，以及肉身的痛苦或死亡；也就是说，他们的不干预是出于自私，而非出于\OriginalTerm{爱德}{love}。\par

因此，我认为，当天主乐意以\OriginalTerm{暂世的惩罚}{temporal punishments}降罚一社会的放荡风俗时，善人与恶人一同受惩，其主要原因在此。他们同受惩罚，并非因为他们过着同样败坏的生活，而是因为善人与恶人，虽程度不同，都贪恋此\OriginalTerm{现世的生命}{present life}；他们本该轻看这生命，好让恶人因善人的榜样受劝诫而归正，可以把握\OriginalTerm{永生}{life eternal}。恶人若不愿与善人同为寻求永生之伴，就应如仇人般受爱，并耐心对待。因为只要他们还活着，是否改过迁善尚不可知。那些自私的人，比那藉先知对他们说：\BibleQuote{他死于自己的罪恶，我却要从警卫员手中追讨他的血债。}\BibleRef{则 33:6} 因为教会中设立了人民的警卫员或监督者，为使他们不姑息地谴责罪恶。我们所谈的这种罪，纵使某人不是警卫员，但他在今生各种关系所接触之人的行为中，见到许多应受谴责的事，却因怕得罪人、怕失去那些虽可合法贪恋、他却过分热衷的世俗福分而视若无睹，这人也不能无罪。最后，另有原因使善人遭受\OriginalTerm{暂世灾祸}{temporal calamities}——即约伯的案例所展示的：为使人灵受试探，并显明它以何等虔诚信赖的刚毅、何等无私的爱，坚附于天主。\par

% structure: p0038 source=h2 target=subsection reason=relative-heading-rank; base=h2

\subsection{第10章——圣人失去暂世财物一无所失}

这些是必须牢记的考量，以便回答这个问题：是否有任何恶事临到忠信和虔敬的人，而不能转化为益处呢？难道我们要说这个问题毫无必要，而宗徒说：\BibleQuote{我们知道，凡爱天主的人，一切事都协助他们获得益处。}\BibleRef{罗 8:28} 是徒托空言吗？\par

他们失去了一切所有。失去了\OriginalTerm{信德}{faith}吗？失去了\OriginalTerm{虔敬}{godliness}吗？失去了那藏于内心的隐密之人，在天主眼中极其宝贵的财物吗？\BibleRef{伯前 3:4} 他们失去了这些吗？因为这些才是基督徒的财富，富有的宗徒曾对他们说：\BibleQuote{虔敬与知足是大利，因为我们没有带什么到世界上，同样也不能带走什么。只要有吃有穿，就当知足。至于那些想望致富的人，却陷于诱惑，堕入罗网和许多不智及有害的欲望中，这欲望叫人沉溺于败坏和灭亡中。因为贪爱钱财乃万恶的根源；有些人曾因贪求钱财而离弃了信德，使自己受了许多刺心的痛苦。}\par

因此，那些在罗马遭劫掠时丧尽了尘世财产的人，如果他们如同那位外表贫穷、内心却丰富的宗徒所教导的那样，拥有财物犹如没有拥有——也就是说，他们用世物，如同不用——便能以虽受重试炼、却未被压倒的\Person{约伯}的话说：\BibleQuote{我赤身脱离母胎，也要赤身归去；主赐的，主收回；主所喜悦的，就这样发生了：愿主的名受赞美！}\BibleRef{约 1:21} 约伯如同一位良善的仆人，将自己主人的旨意视为他的伟大财富，因服从这旨意，他的灵魂得以富有；在有生之年失去那些他不久后死亡时必须抛下的财物，也并不使他忧伤。但对于那些较为软弱的灵魂，尽管不能说他们爱尘世财产胜过基督，却仍以某种过度的依恋执着于它们，他们通过失去这些财物所感受到的痛苦，才发现自己往日对它们的贪爱是多么地犯罪。因为他们的忧伤是自找的；正如上文所引宗徒的话：\BibleQuote{使自己受了许多刺心的痛苦。} 他们长久轻忽这些言语的劝告，如今从经历中学到教训，这正是好事。因为当宗徒说：\BibleQuote{那些想望致富的人，却陷于诱惑}等等时，他所谴责的并非拥有财富本身，而是贪图财富。因为在另一处他说：\BibleQuote{至于今世的富人，你要劝告他们，不要心高气傲，也不要寄望于无常的财富，惟独寄望于那将万物丰富地供给我们享用的天主；又要劝他们行善，在善工上致富，甘心施舍，乐意通财，为自己积蓄良好的根基，以备将来，能享受那真正的生命。}\BibleRef{弟前 6:17} 那些如此使用自己财产的人，虽遭受轻微损失，却因巨大的收益而得到安慰；他们因慷慨施舍而稳妥储存起来的财富，带给了他们更多喜乐，远胜过因焦虑与自私的囤积而完全失去的财物所带来的忧伤。因为在地上，唯有那些他们羞于从世上带走的，才会朽坏。我们的主这样命令说：\BibleQuote{你们不要在地上为自己积蓄财宝，因为在地上有虫蛀，有锈蚀，在地上也有贼挖洞偷窃；但该在天上为自己积蓄财宝，因为那里没有虫蛀，没有锈蚀，那里也没有贼挖洞偷窃。因为你的财宝在那里，你的心也必在那里。}\BibleRef{玛 6:19} 凡听从了这一命令的人，在磨难时期便证实了他们不轻看这位最可信赖的导师、最忠实有力之财宝守护者的劝告，是多么明智。因为若许多人因财宝藏在敌人恰巧找不到的地方而欢欣，那么那些顺从他们的天主的劝告，携财宝逃入敌人绝不可能攻到的堡垒中的人，其喜乐是何等更为确实！因此，我们那位自愿抛弃巨大财富、虽在圣德上极为富有却成为赤贫的\Place{诺拉}主教\Person{保利诺}，当蛮族洗劫\Place{诺拉}并将他俘掳时，据他后来告诉我，他心中默默祈祷说：\Quote{主啊，求祢不要让我为金银焦虑，因为祢知道我的全部财宝在哪里。} 因为他的全部财宝，就在那位曾亲自预告这些灾祸将在世上发生的天主教导他隐藏和储存之处。因此，那些听从了主的警告，得知应在何处、如何积蓄财宝的人，在蛮族入侵时，就连他们的尘世财物也未丧失；而那些如今正因未服从主而懊悔的人，即便未能从预防损失的智慧中，至少也从随之而来的经历中，学到了正确使用世间财物之道。\par

然而，一些善良的基督徒曾遭受酷刑，为要强迫他们向敌人交出自己的财产。他们确实既不能交出、也不能丧失那使他们成为善人的善。然而，如果他们宁愿受酷刑，也不愿交出\OriginalTerm{不义之财}{mammona iniquitatis}，那么我要说，他们并非善人。其实，倒该让他们想起，若为了金钱便忍受如此重苦，那么为了基督，如果有必要，他们更当承受一切折磨；这样他们才能受教，去爱那以永恒真福丰富所有为祂受苦者的主，而不是爱金银；为了金银而受苦是可怜的，不论他们用撒谎保全了它，还是用说真话而失去了它。因为在这种酷刑之下，没有人因宣认基督而失去基督，没有人保存了财产，除非否认财产的存在。因此，可能那教导他们应寄心于一种不能丧失之财富的酷刑，远比那些毫无裨益、只令焦虑的主人不安和痛苦的财产更有用。不过，我们也被提醒，有些人并无财产可交出，却仍受了酷刑，因为他们如此说时，未获相信。然而，这些人或许也曾贪慕财富，并非以圣善的顺服甘愿贫穷；对于这类人，需要显明，不仅实际的拥有，连对财富的贪欲也应得如此剧烈的痛苦。即使他们确实没有任何隐秘的金银积蓄，因为他们生活在更好生命的希望中——我确实不知是否有这样的人因猜疑他拥有财富而受酷刑；但若有，那么当他受审问时，他宣认了\OriginalTerm{神贫}{sancta paupertas}，便无疑宣认了基督。虽然几乎不能指望蛮族会相信他，但任何宣认神贫的人，绝不可能在不受天上赏报的情况下遭受酷刑。\par

此外，人们说漫长的饥荒击倒了许多基督徒。但这一苦难，信友们也以虔敬的忍耐将其转为善用。因为那些被饥荒彻底夺去生命的人，如同被一种温和的疾病解救一般，脱离了此生的种种不幸；而那些仅受饥饿困扰的人，则学会了更节俭地生活，并习惯了更长时间的守斋。\par

% structure: p0044 source=h2 target=subsection reason=relative-heading-rank; base=h2

\subsection{第11章 论此生的终结，无关其长久迟延与否}

但有人又说，许多基督徒惨遭屠杀，以种种骇人听闻的残忍方式被处死。诚然，若此事难以忍受，那这确实是一切生于斯世者共同的命运。至少我确信：从未有人死去，若非注定终有一死。生命的终点使最长与最短的寿命等同。因为对于均已止息的两样事物，既无优劣，亦无大小之分。生命以何种死亡终结又有何紧要？既然人一旦死去，便无需再度经受同样的试炼。在日常生活中，每个人都仿佛时刻面临无数种死亡的威胁，只要命运未定，我便要问：与其活在所有死亡的恐惧中，不如遭受一次死亡而了结，岂不更好？我并非不知那种怯懦的恐惧，它驱使我们宁可长久活在惧怕众多死亡之下，也不愿一死了之，从而逃离一切；然而，肉体的软弱与怯懦的退缩是一回事，灵魂经过深思熟虑的合理信念则完全是另一回事。那作为善生终结的死亡，不应被视为恶事；因为死亡之所以成为恶，只因其后随之而来的报应。因此，注定要死的人们无需关心将死于何种死亡，而应关心死亡将把他们引向何处。基督徒既深知，那位被狗舔疮的虔敬乞丐之死，远胜于那身穿紫袍细麻衣的邪恶富人之死，那么，这些恐怖的死亡又怎能为善生而死者带来伤害呢？\par

% structure: p0046 source=h2 target=subsection reason=relative-heading-rank; base=h2

\subsection{第12章 论死者的埋葬：拒绝给予基督徒埋葬并不伤害他们}

此外，有人提醒我们，在当日那种大屠杀中，连尸体都无法得到埋葬。但虔敬的信赖并不因如此不祥的境遇而惊惶；因为信徒铭记着已获保证：他们连一根头发也不会失落，因此，即使他们被野兽吞噬，他们蒙福的复活也不会因此受阻。真理绝不会说：\BibleQuote{那杀身体不能杀灵魂的，不要怕他们}\BibleRef{玛 10:28}，假如仇敌对被杀害者身体所做的任何事，竟能损害来世的生命。或者，有人会持如此荒谬的立场，争辩说：那些杀身体的人，不该在死前害怕他们、怕他们杀身体，而该在死后害怕他们、怕他们剥夺埋葬吗？若果如此，那么基督所说的话便是虚假的了：\BibleQuote{你们不要害怕那些杀害肉身，而后不能更有所为的人}\BibleRef{路 12:4}；因为看起来他们确实能对死尸造成重大伤害。我们绝不认为真理会如此虚假。经上说那些杀身体的人\Quote{能有所为}，是因为致死的一击仍有感觉，身体尚有知觉；但此后，他们便再无能为，因为已死的身体已无知觉。因此，确实有许多基督徒的尸体未得埋葬，横陈于地；但无人能将他们与天堂隔绝，也无人能将他们与那充满祂临在的大地隔绝，祂知晓将从何处复活祂所造的。的确，圣咏中说：\BibleQuote{将你仆人的尸首，给天空的飞鸟作食物；将你圣者的尸体，给地上的野兽作食物。在耶路撒冷周围流他们的血，如水一般，无人埋葬。}但说这话，更是为彰显行事者的残忍，而非受苦者的悲惨。在人眼中，这似乎是严酷而可悲的命运，然而\BibleQuote{主的圣者们的去世，在主的眼中十分珍贵}。因此，所有为死者举行的最后职务与礼仪，精心的葬礼安排，墓穴的装备，以及出殡的盛况，与其说是死者的安慰，不如说是生者的慰藉。倘若一场奢华的葬礼能为恶人带来任何益处，那么鄙陋的埋葬，或根本没有埋葬，便可能伤害虔敬的人。他的众多仆役为身穿紫袍的\OriginalTerm{财主}{Dives}操办了人眼中华丽的葬礼；但在天主眼中，那位疮痍满身的乞丐从天使手中所得的葬礼，却更为隆重；天使们并没有将他抬往大理石墓穴，而是将他高举到亚巴郎的怀抱中。\par

我着手捍卫天主之城，所要反对的那些人，对此嗤之以鼻。然而，就连他们自己的哲学家也轻视精心的埋葬；常有整个军队为尘世的家国而战而亡，从不费心探究自己是否会被弃尸沙场，或成为野兽的食物。对于这种对丧葬的崇高轻蔑，有诗句说得好：\Quote{无墓者，以天为穹顶。} 他们对基督徒未得埋葬的尸体，岂不更不该加以嘲弄？因为基督徒已蒙许诺，血肉之躯必将恢复，身体亦将重新塑造，其所有肢体不仅从大地，也会从其他元素最隐秘的深处汇聚而来——那些死人的尸体曾隐藏其间！\par

% structure: p0049 source=h2 target=subsection reason=relative-heading-rank; base=h2

\subsection{第13章 埋葬圣人遗体的理由}

但死者的遗体却不可因此受到轻贱而弃置不葬；尤其不可这样对待义人与忠信者的身体，因为圣神曾使用它们，作为行一切善工的工具与器皿。倘若父亲的衣袍、戒指或任何用过之物，因儿女对他的爱而成为珍贵，那么对于我们所爱之人的遗体，我们岂不更该照料？——他们穿戴这身体，比任何衣物都更为贴身和亲密！因为身体并非外在的饰物或辅助，而是人之本性的组成部分。因此，古时的义人虔诚地尽好安葬的本分，为亡者预备坟墓，举行殡葬礼；他们自己生前也嘱咐子孙有关下葬的事，有时甚至指示要将遗体迁往心爱的地方。多俾亚因着天使的见证而受赞扬，说他是因埋葬死者而悦乐天主。\BibleRef{多 12:12} 我们的主自己，虽然第三日要复活，但对那虔诚的妇人将珍贵的香液倾倒在祂的肢体上，为祂的安葬而作此美事，也予以称扬，并令人传颂。\BibleRef{玛 26:10} 福音也称赞那些小心翼翼从十字架上取下祂的身体，敬爱地用细麻布裹好，并妥为安葬的人。\BibleRef{若 19:38} 这些事例当然并不证明遗骸有任何感觉；但它们表明，天主的眷顾甚至及于死者的身体，这类虔诚的本分因培植对复活的信德而蒙祂悦纳。我们也可从中汲取有益的教训：如果天主连慈爱之心施于无知觉的死者身上的任何善工都不遗忘，那么祂必定更会酬报我们施予生者的爱德。至于圣祖们有关自己遗体的安葬与迁葬所说的话，他们意在让人从先知性的意义上去理解；关于这些，我们在此不必多言，前面所述已经足够。不过，如果维持活人生存所必需的物品，如食物与衣服，即使匮乏虽痛苦且熬炼人，却不能摧毁善人的坚毅和德行的忍耐，也不能从他们灵魂中根除虔敬，反而使之更为丰盛，那么，缺少葬礼及其他惯例中的遗体照料，又岂能使那些已在有福者隐秘居所中安息的亡者成为不幸的呢！因此，在罗马及其他城镇被洗劫时，虽然基督徒的遗体被剥夺了这些最后的本分，这既非生者的过失，因为他们无法提供；对死者也非一种灾祸，因为他们不会有失去之感。\par

% structure: p0051 source=h2 target=subsection reason=relative-heading-rank; base=h2

\subsection{十四、论圣人们的被掳，以及天主的安慰从未弃绝他们。}

但有人说，许多基督徒甚至被掳走。这诚然是极其可悲的命运，倘若他们能被掳往一个找不到自己天主的地方的话。然而对于这灾难，圣经也提供了极大的安慰。三位青年\EditorialNote{达 3}曾被掳；\Person{达尼尔}曾被掳；其他先知亦然：而安慰人的天主并未舍弃他们。同样，祂也没有在蛮族（尽管还是人）手中放弃自己的子民——祂是那位没有将先知遗弃在怪物腹中的主。\EditorialNote{纳 1} 不过，这些事在我们与之辩论的人看来，与其说被相信，不如说被嘲笑；尽管他们相信自己书中所读到的，说\Place{梅廷纳}的著名琴师\Person{亚里翁}被抛入海中时，是被海豚的背接住并载到岸上的。然而我们关于先知\Person{约纳}的故事，远为更难以置信——越神奇便越难以置信，而越展示大能便越神奇。\par

% structure: p0053 source=h2 target=subsection reason=relative-heading-rank; base=h2

\subsection{十五、论\Person{雷古鲁斯}，他为我们立了为宗教缘故甘愿忍受被掳的榜样；然而这对身为诸神敬拜者的他并无益处。}

但在他们自己的著名人物中，他们有一个非常崇高的榜样，为顺从宗教顾虑而自愿忍受被俘。罗马将军\Person{马尔库斯·阿蒂利乌斯·雷古鲁斯}落在迦太基人手中成为俘虏。但迦太基人更急于与罗马人交换俘虏，而非扣留他们，于是派遣雷古鲁斯作为特使，和他们的使者一同去谈判这一交换，但先用誓言约束他，如果他不实现他们的愿望，就必须返回\Place{迦太基}。他去后，说服元老院采取了相反的方针，因为他认为交换俘虏对\OriginalTerm{罗马共和国}{Roman republic}无益。在他如此施加影响之后，罗马人并未强迫他返回敌人那里；但他自愿履行了他所起的誓言。但迦太基人用一种精巧、残忍而可怕的酷刑将他处死。他们把他关在一个狭窄的箱子里，迫使他站立其中，箱内四周布满锐利的钉子，使他无法倚靠任何部分而不感到剧痛；就这样，他们用剥夺睡眠的方式将他处死。确实，他们公正地赞扬了那种超越如此可怕命运的德性。然而，他所凭以起誓的\OriginalTerm{神祇}{gods}，正是现在被认为因禁止其崇拜而向人类施加这些现今灾祸的\OriginalTerm{神祇}{gods}。但如果这些\OriginalTerm{神祇}{gods}，人们为了在此生获得幸福而特别崇拜，对于一位遵守誓约的人，尚且愿意或容许这样的惩罚加在他身上，那么他们对发假誓者盛怒之下，又会施加何等更残酷的惩罚呢？但是，我何不从这推理中得出双重结论呢？雷古鲁斯对\OriginalTerm{神祇}{gods}确实如此敬畏，以致为了誓约，他不肯留在本国，也不肯去别处，而是毫不犹豫地回到他最凶恶的敌人那里去。如果他认为这样做对此生有利，那他肯定大错特错了，因为这引向了他生命的可怕终结。事实上，他以自己的榜样教导说，\OriginalTerm{神祇}{gods}并不保障崇拜者的现世幸福；因为他自己，一位崇拜\OriginalTerm{神祇}{gods}的人，既在战斗中被击败而成了俘虏，又因不肯违反他所凭以起誓的\OriginalTerm{神祇}{gods}之誓言，而遭受了一种新的、前所未闻的、过于恐怖的刑罚，被折磨至死。如果假设\OriginalTerm{神祇}{gods}的崇拜者会在来生得享幸福，那么，他们又为何毁谤\OriginalTerm{基督教}{Christianity}的影响呢？他们为何断言，这城遭受灾难是因为它不再崇拜自己的\OriginalTerm{神祇}{gods}了呢？尽管它可能像雷古鲁斯一样殷勤地崇拜它们，却仍可能像雷古鲁斯一样不幸。或者说，难道有人竟会如此令人惊异盲目，在明显的事实面前，大胆妄图争辩说，虽然一个人崇拜\OriginalTerm{神祇}{gods}也可能不幸，但整个城市却不会如此？这就是说，他们的\OriginalTerm{神祇}{gods}的能力更适合保护群体而非个人——仿佛群体不是由个人组成的。\par

但如果他们声称，\Person{马尔库斯·雷古鲁斯}即使在囚禁并忍受这些身体折磨时，仍可能享有有德灵魂的\OriginalTerm{真福}{blessedness}，那么，就让他们承认，一座城市也可因\OriginalTerm{真德性}{true virtue}而获得\OriginalTerm{真福}{blessedness}。因为群体和个体的\OriginalTerm{真福}{blessedness}都源于同一源头；因为群体无非是个体的和谐集合。因此，我暂时不关心讨论雷古鲁斯拥有何种德性；单凭他那极为崇高的榜样，就足以迫使他们承认，崇拜\OriginalTerm{神祇}{gods}不是为了身体的安逸或外在的利益；因为他宁可失去所有这类东西，而不愿冒犯他凭以起誓的\OriginalTerm{神祇}{gods}。但是，对于那些以拥有这样一位公民为荣，却害怕拥有一座像他一样的城市的人，我们该作何感想呢？如果他们不害怕这一点，那就让他们承认，即使他们像他那样勤勉地崇拜\OriginalTerm{神祇}{gods}，类似雷古鲁斯所遭遇的灾祸也可能临到一个城市；并让他们不再把自己的不幸归咎于\OriginalTerm{基督教}{Christianity}。然而，此时我们所关注的是那些被俘的\OriginalTerm{基督徒}{Christians}，那些借这次灾祸，以不亚于无礼的轻率来辱骂我们极有益健康的宗教的人，应该斟酌此事，并闭口无言；因为如果一位最谨守崇拜\OriginalTerm{神祇}{gods}的人，为了遵守他对\OriginalTerm{神祇}{gods}的誓言，竟失去祖国而无望找到另一个，落入敌人手中，遭受长久而精巧的酷刑而死，这对他们的\OriginalTerm{神祇}{gods}并不构成谴责，那么，更不应该把信靠其权能者的被俘归咎于基督徒之名，因为他们怀着对\OriginalTerm{天乡}{heavenly country}的坚定期盼，知道即使在自己的家中，他们也是旅客。\par

% structure: p0056 source=h2 target=subsection reason=relative-heading-rank; base=h2

\subsection{第十六章——论\OriginalTerm{祝圣贞女}{Consecrated Virgins}及其他基督徒贞女在俘虏中所遭受的侵犯，她们的意志并未同意；并论此事是否玷污了她们的\OriginalTerm{灵魂}{Souls}。}

但他们自以为对基督教提出了决定性的指控，因为他们加剧了俘虏的恐惧，声称不仅已婚妇女和未婚少女，甚至献身守贞的贞女也遭到侵犯。但事实上，就这一点而言，陷入困境的并非基督徒的\OriginalTerm{信德}{faith}、\OriginalTerm{虔敬}{piety}，甚至也不是\OriginalTerm{贞洁}{chastity}的美德；唯一的困难在于如何以一种同时满足廉耻和理性的方式处理这一主题。在讨论中，我们不会急于答复控告我们的人，而是要安慰我们的朋友。因此，首先应当确立一个不可动摇的立场：使生活成为善的德行，其宝座在灵魂之中，并从那里统治身体的肢体；身体因意志的圣洁而成为圣洁的；只要意志保持坚定不动摇，别人对身体所做的或加于身体之上的事，对于忍受者来说便没有过错，只要他无法不犯罪地逃脱。但是，不仅痛苦可能被加诸于他人的身体，而且淫欲也可能在他人的身体上得到满足。每当后一种情况发生时，羞耻便会侵袭一个全然纯洁的灵魂——因廉耻尚未离去——之所以羞耻，是因为担心那若没有某种感官快感便无法忍受的行为，会被认为也伴随着意志的某种赞同。\par

% structure: p0058 source=h2 target=subsection reason=relative-heading-rank; base=h2

\subsection{第十七章 出于对惩罚或耻辱的恐惧而自杀}

因此，即使这些贞女中的一些人为了避免这样的羞辱而自尽，凡有恻隐之心的人谁会不予以宽恕呢？至于那些不愿结束自己生命的人，以免似乎以自己的一项罪来逃避他人的罪，若有人将此举指控为极大的邪恶，那他自己也难逃愚昧之咎。因为，如果擅自执法，甚至处死一个未经公开判决的罪犯都是不合法的，那么自杀者肯定是\OriginalTerm{杀人罪}{homicide}，而且在他自己的死亡上罪孽更重，因为正是他使自己被处死的那个过失，他更为清白无辜。难道我们不是公正地憎恶\Person{犹达斯}的行为吗？真理本身岂不也判定，他因上吊自尽，非但没有补赎，反而加重了他那极其邪恶的背叛之罪，因为他因那导致死亡的悲伤而对天主的仁慈绝望，没有给自己留下任何挽回的\OriginalTerm{忏悔}{penitence}的余地？一个未犯有任何该当此罚之罪行的人，更应避免对自己施加暴力！因为犹达斯在自杀时，杀的是一个恶人；但他离世时，不仅背负着基督的死亡，也背负着他自己的死亡：因为他虽因自己的罪行而自杀，但他的自杀却是另一个罪行。那么，一个没有做过任何恶事的人，为何要对自己作恶，通过杀死自己——一个无辜者——来逃避他人的罪行之举，并在自己身上犯下自己的罪，以免他人的罪在他身上发生呢？\par

% structure: p0060 source=h2 target=subsection reason=relative-heading-rank; base=h2

\subsection{第十八章 论他人淫欲可能对身体施加的暴力，而心灵保持不受侵犯}

然而，是否有人担心，即便他人的淫欲也可能会玷污受害者呢？若它是他人的，便不会玷污；若它玷污了，那便不是他人的，而是受害者也共同参与了。然而，既然纯洁是灵魂的一种德行，并以\OriginalTerm{刚毅}{fortitude}为其同伴之德，宁可忍受一切凶恶也不赞同邪恶；而且，没有人，无论多么高尚和纯洁，能一直支配自己的身体，他只能控制自己意志的同意或拒绝。那么，一个心智健全的人怎能设想，如果他的身体被强行抓住并用以满足他人的淫欲，他就会因此丧失纯洁呢？因为，如果纯洁可以如此被摧毁，那么纯洁定然不是灵魂的德行；它也不能被列为那些使生活变得美好的善物之中，而只能是身体的善物，与强壮、美丽、完好无损的健康，总之，所有那些即使减少也丝毫不削弱我们生活良善与正直的善物属于同一范畴。但如果纯洁不比这些更好，为何要为保全它而危及身体呢？另一方面，如果它属于灵魂，那么，即使在身体受到侵犯时，它也不会丧失。更有甚者，神圣\OriginalTerm{节德}{continence}的德行，当它抵抗肉欲的不洁时，甚至圣化身体，因此，当这种节德保持不屈时，身体的\OriginalTerm{圣洁}{sanctity}也得以保全，因为圣善地使用身体的意愿依然存在，而且，就身体本身而言，能力也依然存在。\par

因为身体的圣洁不在于其肢体的完整，也不在于它们免受任何触摸；因为身体会遭遇各种意外，使其受暴力侵犯和伤害，而施行救治的外科医生所进行的操作，常令旁观者作呕。接生婆，姑且说，在试图确定一个女孩的\OriginalTerm{童贞}{virginity}时（无论是出于恶意、意外，还是技艺不精），破坏了它：我想，没有人会如此愚蠢，以至于相信，因这一器官完整性的破坏，那贞女便丧失了其身体的圣洁。因此，只要灵魂保持这种坚定意志——它甚至圣化身体，那么他人淫欲所施加的暴力，便丝毫不会影响这身体的圣洁，这种圣洁因自身持守的节德而完好无损。假设一个贞女违背了她向天主所发的誓言，前去与她的引诱者相会，意图顺从他：我们能否说，在她前往之时，她仍拥有身体的圣洁，而那时她已经丧失并摧毁了圣化身体的灵魂的圣洁？我们绝不能如此误用词语。我们宁愿得出这样的结论：即使身体受到侵犯，灵魂的圣洁依然存留时，身体的圣洁便不会丧失；同样，当灵魂的圣洁被侵犯时，身体的圣洁也就丧失了，即使身体本身保持完整。因此，一个因他人的罪、且未经她本人任何同意而被侵犯的妇女，没有任何理由自杀；更不用说，为了避免这样的侵犯而自杀：因为在那种情况下，她实施确定的\OriginalTerm{杀人罪}{homicide}，以阻止一项不确定的、且不属于她自己的罪行。\par

% structure: p0063 source=h2 target=subsection reason=relative-heading-rank; base=h2

\subsection{第十九章 \Person{路克雷霞}因遭受暴行而自杀}

因此，这就是我们的立场，而且似乎足够清楚。我们主张，当一个妇女在被侵犯时，若她的灵魂对罪恶未予同意，而保持不可侵犯的贞洁，那么罪过就不在她，而在于侵犯她的人。但是，那些我们不但要保卫其灵魂，还要保卫这些受辱的基督徒俘虏的圣洁身体的人——他们或许敢于反驳我们的立场吗？然而，所有人都知道他们多么高声赞扬\Person{路克雷霞}的贞洁，那位古\Place{罗马}的高贵主妇。当王子\Person{塔克文}之子侵犯了她的身体时，她将这位年轻浪子的恶行告知了她的丈夫\Person{科拉提努斯}，以及她的亲戚\Person{布鲁图斯}，这两位都是地位崇高、充满勇气的人，并让他们发誓报仇。然后，她心如刀绞，无法承受这耻辱，便结束了自己的生命。我们该称她什么呢？淫妇，还是贞妇？这毫无疑问。一位演说家论及这悲惨事件时，说得既贴切又真实：\Quote{这里有一个奇迹：有两个人，但犯奸淫的只有一个。}这话极为有力而真实。因为这位演说家看到了两具身体的结合中，一方的肮脏情欲和另一方的贞洁意愿，注重的不是肢体的接触，而是灵魂的巨大差异，所以他说：\Quote{有两个人，但犯奸淫的只有一个。}\par

但是，为何她并未参与罪行，却要承受两者中更重的惩罚呢？因为那奸淫者只与他的父亲一起被流放；她却遭受了极刑。如果那并非她所愿的玷污不算是污秽，那么她这位贞妇受惩罚也就不算是正义了。我向你们呼吁，\Place{罗马}的法律和法官们。即使在犯下滔天大罪之后，你们也不允许未经审判就将罪犯处死。那么，如果有人将这案件带到你们的法庭上，并向你们证明，有一位妇女不仅未经审判，而且贞洁无辜，却被杀害了，难道你们不会对那谋杀者施以相应的严厉惩罚吗？这罪是\Person{路克雷霞}犯下的；那位备受称赞和颂扬的\Person{路克雷霞}，杀死了无辜的、贞洁的、受辱的\Person{路克雷霞}。你们宣判吧。但如果你们不能宣判，因为没有可惩罚之人，那你们为何还要如此过度地颂扬那位杀死无辜贞妇的人呢？毫无疑问，你们会发现，在阴间的法官面前为她辩护是不可能的，如果那里的法官真如你们的诗人所乐于描绘的那样的话；因为她就在那些\par

\Quote{他们清白无辜，却自寻死路，\newline{}只因厌恶光明，\newline{}在疯狂中抛弃了生命。\newline{}若是她想与其他人一同返回，}\par

\Quote{命运阻断了归途：围绕着他们的居所，\newline{}那迟缓而可憎的冥河之水缓缓流淌，\newline{}并以九重锁链将他们束缚。}\par

或许她不在那里，因为她自尽时自知有罪，而非无罪？只有她自己知道原因；但如果她是在那行为中感受到了快感，对\Person{塞克斯图斯}有了某种同意，尽管他是如此暴力地侵犯她，然后她感到懊悔，认为只有死亡才能赎罪呢？即使如此，她仍然不该自杀，如果她能通过她的假神们完成有效的忏悔的话。然而，如果是这种情况，如果并非两人，而是只有一人犯奸淫的说法是错误的；如果事实是两人都参与了，一人通过公开的侵犯，另一人通过秘密的同意，那么她杀死的就不是一个无辜的妇女；因此，她那博学的辩护者们就可以主张，她并不属于下界居民中的那一类\Quote{清白无辜却自寻死路}的人。但\Person{路克雷霞}的这个案例陷入了两难：如果你减轻她的杀人罪，你就确认了她的奸淫罪；如果你为她洗脱奸淫罪，你就加重了她的杀人罪；而且没有出路，当人们问道：如果她是淫妇，为何要赞扬她？如果是贞妇，为何要杀死她？\par

然而，我们的目的是反驳那些无法理解真正圣洁是什么，并因此侮辱我们受辱的基督徒妇女的人，就这点而言，在这位高贵的\Place{罗马}主妇的例子中，人们为了赞扬她而说\Quote{有两个人，但犯奸淫的只有一人}，这就足够了。因为人们确信\Person{路克雷霞}能够超越任何对奸淫抱以同意的念头。因此，既然她因遭受暴行而自杀，而那暴行中她并无罪责，那么很显然，她这一行为并非出于对贞洁的热爱，而是出于耻辱的重压。她感到羞耻的是，如此肮脏的罪行竟在她身上发生，尽管她并未参与；这位血管里流着\Place{罗马}人爱荣耀之血的贵妇，被一种骄傲的恐惧所攫住，担心如果她继续活下去，人们会以为她自愿容忍了对她所犯的恶行。她无法向人展示自己的良心，但她断定，自己施加的惩罚可以证明她的心态；一想到她对别人对她施加的污秽暴行的忍耐，会被解释为与那人同谋，她便羞愤难当。基督徒妇女们遭受了同样的暴行，却存活了下来，她们的决定却不是这样的。她们拒绝为了他人的罪而惩罚自己，以免在那些她们并无参与的罪行之外，再增添自己的罪行。如果她们因为羞耻而杀人，正如敌人的情欲驱使她们犯奸淫一样，那么她们就会这样做。她们在自己的灵魂内，在自己良心的见证下，享有贞洁的荣耀。在天主眼中，她们也被视为纯洁，这便使她们满足；她们别无所求：有机会行善便足够了，她们拒绝逃避世人的怀疑所带来的困苦，以免因此偏离天主的法律。\par

% structure: p0070 source=h2 target=subsection reason=relative-heading-rank; base=h2

\subsection{第二十章 基督徒在任何情况下都无权自杀}

此事并非毫无意义：在神圣正典书卷的任何段落中，都找不到天主的诫命或准许，让我们可以为获得永生之乐，或为了躲避、摆脱任何事物而夺取自己的性命。不，这法律若正确解读，甚至禁止自杀，因为经上说：\BibleQuote{不可杀人。} 这一点尤其可从\Quote{你的近人}一词的省略得到证明，因为在禁止作假见证时加上了这几个字：\BibleQuote{不可对你的近人作假见证。} 也不可因为人只对自己作假见证，就认为他没有违犯这条诫命。因为对近人的爱，是以对自己的爱来规度的，正如经上所载：\BibleQuote{应爱你的近人，如你自己。} 那么，若一个人对自己说假话，其作假见证的罪并不比作证伤害近人来得轻；虽然禁止作假见证的诫命只提到了近人，而不用心理解的人，可能以为人可被允许作害己的假见证。那么更有理由明白，人不可以杀自己，因为在\BibleQuote{不可杀人}这条诫命中，既没有附加任何限制，也没有对任何人、尤其对受命者本人，作出任何例外！因此，有些人试图将这命令扩展至兽类和牲口，好像它禁止我们剥夺任何受造物的生命。但若如此，为何不也扩展到植物，以及一切扎根于地、由地滋养之物呢？因为这类受造物虽然没有感觉，但人却也说它们有生命，因而能死；所以，若对其施加暴力，便能杀死它们。同样，宗徒论及此类种子时也说：\BibleQuote{你所撒的种子，若不先死了，决不能生出来；} \WorkTitle{圣咏}中也说：\BibleQuote{祂以冰雹打死了他们的葡萄树。} 那么，我们是否必须把摘一朵花也视为违犯\BibleQuote{不可杀人}这条诫命呢？我们岂能如此疯狂地认同\OriginalTerm{摩尼教徒}{Manichæans}的愚昧错谬？那么，且放下这些呓语，如果当我们说\BibleQuote{不可杀人}时，我们并未将此理解为针对植物，因为它们没有感觉；也不理解为针对无理性的飞禽、游鱼、走兽、爬虫，因为它们因缺乏理性而与我们区隔，并因此按造物主公正的安排而被置于我们权下，供我们为自身用途而杀或养；那么，这诫命就只剩下被理解为专对人而言。这诫命便是：\BibleQuote{不可杀人}；因此，既不可杀他人，也不可杀自己，因为凡杀己者，所杀的仍然不外是人。\par

% structure: p0072 source=h2 target=subsection reason=relative-heading-rank; base=h2

\subsection{第21章 在何种情况下处死他人可不负杀人罪责。}

然而，天主的权威对祂自己的那不可杀人的法律，定下了一些例外。这些例外分为两类：或由一项普遍的法律所证成，或由赐予某人一段时期的特别使命所证成。而在后一种情况中，被授权者，只是运用权柄者手中的剑，他本人对所施行的死亡不负责任。因此，凡服从天主的命令或遵照祂的法律而作战的人，都亲自代表了公共正义或治国的智慧，并以此身份处死恶人；这样的人丝毫没有违反\BibleQuote{不可杀人}这条诫命。\Person{亚巴郎}实在不仅被视为没有残忍之罪，更因其虔敬而备受称赞，因为他准备服从天主而非自己的私情，杀献亲子。至于\Person{依弗大}杀了自己的女儿，是否应视为奉命于天主的命令，这个问题也颇有理由。因为他曾许愿说，当他战胜归回时，无论谁先从家门出来迎接他，他就要把谁献为全燔祭；而他的女儿出来迎接了他。同样，\Person{三松}倾覆房子，与敌人同归于尽；他得以成义，只基于这一点：那藉他行奇事的圣神，曾暗中指示他如此行事。那么，除了这两类情况——或由一项普遍适用的公正法律，或由天主本身、一切正义之源的特别指示所证成——之外，凡杀人者，无论杀己还是杀他，都陷于谋杀之罪。\par

% structure: p0074 source=h2 target=subsection reason=relative-heading-rank; base=h2

\subsection{第22章 自杀绝不能出于豪迈之德。}

但那些亲手自杀的人，也许因其灵魂之伟大而让人敬佩，然而其判断之健全却不能得到称许。不过，你若更仔细地审视此事，便很难称之为灵魂之伟大；这种伟大竟促使人自杀，而不愿忍受命运的一些艰辛，或他未卷入的罪过。这岂不更是心灵软弱的证明，无法忍受身体的奴役之苦或庸众的愚见吗？难道不应当说，那面对而非逃避人生诸恶，且相较于良心的光明与纯洁，对世人的判断，尤其是常笼罩在错误迷雾中的庸众之见，视若等闲的，才是更伟大的心灵吗？因此，若自杀被视为气度恢宏之举，那么再没有谁比\Person{克勒奥姆布若图斯}在气度恢宏上更胜一筹了；据传，他读了\Person{柏拉图}论灵魂不朽的著作后，从墙上跳下，由此离开了这一生，前往他所相信的更美好的生命。因为他既未受灾难逼迫，也未遭到任何真假难辨而他无法安然度过的指控；总之，没有任何动机，唯独气度恢宏驱使他寻求死亡，摆脱今生这甜蜜的羁留。然而，他所读的\Person{柏拉图}本人会告诉他，这行动虽气度恢宏，却并不正当；因为那洞悉灵魂不朽的同一明哲，也辨明以自杀寻求不朽，应受禁止而非鼓励；否则\Person{柏拉图}必定早已亲自实行，或至少推崇自杀了。\par

再者，据说许多人自杀是为了防止敌人下手。但我们不是在问是否有人这样做过，而是问是否应这样做。健全的判断甚至应优先于榜样，而且榜样若与理性的声音和谐，方为可取；但并非所有榜样，而仅那些以虔诚著称、堪当效法的榜样才如此。论到自杀，我们无法引证圣祖、先知或宗徒的榜样；尽管我们的主耶稣基督在告诫门徒受迫害时应从这城逃到那城时，大可利用那机会劝他们亲手自杀，以此逃避迫害者。但祂既没有这样做，也没有提出这种离世的方式，尽管祂是在对祂自己的朋友们说话，祂曾应许为他们预备永久的居所，这就表明，从那些\Quote{忘记天主的外邦}所取的这类榜样，并不允许敬拜独一真天主的人效法。\par

% structure: p0077 source=h2 target=subsection reason=relative-heading-rank; base=h2

\subsection{我们对\Person{加图}自杀一例的看法：他因无法忍受\Person{凯撒}的胜利而自尽。}

除了已充分谈论过的\Person{卢克雷提亚}，鼓吹自杀的人很难找到任何其他可援引的先例，除非是\Person{加图}之例，他在\Place{乌提卡}自杀。人们援引他的榜样，并非因为只有他这样做，而是因为他被视作博学而优秀的人，以致可以振振有词地主张，他所做的在以前和现在都是善行。但对于他的这一行为，我所能说的只是，他自己的朋友们，那些和他同样开明的人，明智地劝阻了他，因而断定他这行为出自软弱而非坚强的精神，并非出自预杜耻辱的荣誉感，而是出自畏避艰辛的懦弱。事实上，\Person{加图}借他给自己所钟爱的儿子的劝告，定了自己的罪。因为如果在\Person{凯撒}治下生活是一种耻辱，为何作父亲的却要催促儿子去承受这种耻辱，鼓励他完全信赖\Person{凯撒}的宽宏大量呢？为何他不说服儿子与自己同死呢？如果\Person{托尔夸图斯}因儿子违背命令与敌军交战且战胜而被处死而备受称赞，那么战败的\Person{加图}，为何却饶过了自己战败的儿子，尽管他不饶恕自己呢？违背命令却获胜，比起违背时人的荣誉观念而屈从于胜利者，哪一样更为可耻呢？因此，\Person{加图}不可能认为在\Person{凯撒}治下生活是可耻的；因为假若他这样认为，父亲的剑就会使儿子摆脱这耻辱了。真相是，他的儿子，他既希望也盼望会得到\Person{凯撒}饶恕的儿子，其受疼爱之情，并不比他对\Person{凯撒}因宽恕其子而获荣耀的嫉妒更强烈（正如传说\Person{凯撒}本人也这样说）；或者，若嫉妒一词过重，就让我们说他为此荣耀竟落于\Person{凯撒}而\Emph{感到羞愧}。\par

% structure: p0079 source=h2 target=subsection reason=relative-heading-rank; base=h2

\subsection{在\Person{雷古卢斯}胜过\Person{加图}的那种美德上，基督徒尤为卓越。}

我们的对手反感我们宁取圣\Person{约伯}、而非\Person{加图}；约伯宁愿忍受肉身的可怕灾祸，也不肯借自杀脱离一切痛苦；他们也反感其他圣人——据我们权威且可信的典籍所载，这些圣人宁愿忍受被掳和敌人的压迫，也不肯自杀。然而他们自己的典籍却许可我们宁取\Person{马尔库斯·雷古鲁斯}、而非\Person{马尔库斯·加图}。因为\Person{加图}从未征服过\Person{凯撒}；当他被\Person{凯撒}征服时，不屑于屈从，为逃避这种屈从而自杀身亡。与此相反，\Person{雷古鲁斯}先前曾征服\Place{迦太基}人，并作为罗马军队的统帅为\Place{罗马}共和国赢得了一场没有公民可以为之悲叹、连敌人也不得不钦佩的胜利；然而后来，轮到他被他们击败时，他宁愿作他们的俘虏，也不愿借自杀逃脱他们的辖制。在\Place{迦太基}人的统治下他默默忍耐，对\Place{罗马}人的爱却始终不渝：他既未从被征服者那里夺去自己被俘虏的躯体，也未从未被征服者那里夺去自己不屈的精神。他所以没有自杀也并非出于贪生。这一点既明显又充分地表明：他因承诺与誓言，毫不犹豫地回到同一批敌人那里——他在元老院里的言辞对他们的触怒甚至比战场上的武器更为深重。既然他如此蔑视生命，宁愿以激怒的敌人可能设计的任何折磨来结束自己的生命，也不肯亲手了结，他就再清楚不过地宣示了，他认为自杀是多么大的罪行。在\Place{罗马}人所有著名而杰出的公民当中，没有比这个人更值得夸耀的了：他既未因顺境而败坏——赢得这般胜利后，他依然极其贫穷；也未因逆境而折断——他毫不畏惧地返回，面对最为悲惨的结局。然而，如果那些最勇敢、最著名的英雄，他们仅是为保卫一个地上的祖国而战，而且他们虽侍奉的是假神，却向假神献上真诚的敬拜，并谨守向他们所发的誓言；如果这些人，按照战争的惯例和权利，可以杀败被征服的敌人，却在自己被敌人征服时，仍不敢结束自己的性命；如果他们虽毫不惧怕死亡，却宁愿忍受为奴，也不肯自杀；那么，基督徒，真天主的敬拜者，向往天上家园的人，若在天主的眷顾下暂时被交于敌人手中，为要考验他们或惩戒他们，岂不更当远远避开这种行为吗！当然，陷入这屈辱境地的基督徒决不会被至高者遗弃，因为至高者为了他们的缘故自谦自卑。他们也不该忘记，他们不受任何战争法则或军事命令的约束，甚至连被征服的敌人都不应杀掉；如果一个人不可以杀死已经得罪过他、或者可能将要得罪他的敌人，那又是什么人如此疯狂，竟主张人可以因敌人已经得罪或将要得罪他而自杀呢？\par

% structure: p0081 source=h2 target=subsection reason=relative-heading-rank; base=h2

\subsection{第二十五章 我们不可试图以罪来防止罪。}

可是，有人告诉我们，有理由担心：当身体屈服于敌人的情欲时，感官隐秘的快乐会诱惑灵魂赞同犯罪，因此必须采取措施以免这一灾难性后果。那么，自杀岂不正是防止敌人的罪、也防止基督徒受到诱惑而犯罪之恰当方式吗？首先，凡受天主及其智慧引导、而非受肉身贪欲引导的灵魂，决不会赞同由他人情欲在自身肉体中引起的欲望。而且，无论如何，如果如真理所明言的，自杀是可憎且应受永罚的恶行是真实的，那么谁是那样的傻瓜，竟会说：\Quote{让我们现在就犯罪吧，以免将来可能犯罪吧；让我们现在就去杀人，以免日后犯奸淫吧}？如果我们如此被邪恶控制，以致纯全之无罪已无可能，我们充其量只能在诸罪之间作一选择，那么，一个将来且不确定的奸淫之罪，岂不胜过当前确凿的谋杀之罪吗？犯下一个忏悔可以治愈的恶行，岂不胜过触犯一桩不容治愈之痛悔的罪行吗？我说这话是为了那些害怕受诱而赞同施暴者情欲、认为应当对自己施以暴力从而防止——不是别人的罪，而是自己的罪——的男人或女人。然而，一个信靠天主、安息在对祂援助的盼望中的基督徒，其心思决不会——我再说决不会——对呈现出来的肉身快乐，做出可耻的赞同，无论这些快乐如何呈现。况且，如果那仍居住在我们有死之肢体中、不顾我们意志而随从自己律法的\OriginalTerm{贪欲的悖逆}{lustful disobedience}，它在反叛之人的身体中的活动，与它在睡眠之人身体中的活动一样，都是无可指摘的。\par

% structure: p0083 source=h2 target=subsection reason=relative-heading-rank; base=h2

\subsection{第二十六章 在某些特殊情形下，圣人的榜样是不可效法的。}

但是，他们说，在教难时期，有些圣洁妇女为躲避那些以强暴威胁她们的人，便投河自尽，明知自己会溺死；以这种方式死去后，她们在天主教会中被尊为殉道者。对于这类人，我不敢妄加论断。我无法断定，教会是否曾蒙天主赐予某种神圣权威，且由可信的证据证实，如此敬礼她们的纪念：也许确有此事。也许她们并非出于人的判断而误入歧途，而是受天主的智慧驱使，实行了自尽的行为。我们知道\Person{三松}的事便如此。当天主命人行一事，并以明白的证据指明这是祂的命令时，谁敢斥责这服从为罪恶？谁会控告如此虔诚的顺命？然而，并非每个人都能像\Person{亚巴郎}献子那样，因这行为是可赞美的而自视为正当。士兵遵照合法授予他权力的长官的命令杀人，国家的任何法律都不控告他犯谋杀罪；不，若他没有执行命令杀了那人，他反会被控告背叛国家、藐视法律。但如果他凭自己的权威、出于自己的冲动行事，在此情况下，他便犯下流人血的罪行。因此，他因无命令而行这事受到了惩罚，而这事他若受命而不行亦将受罚。如果将军的命令能产生如此大的区别，天主的命令难道毫无分别？那么，明知自杀是不合法的人，若领受了天主的命令——这命令我们不可忽视——却可以自尽。只是他必须十分确定这神圣的命令已指示明白。至于我们，我们仅能在良心秘密向我们揭示的范围内，成为这些秘密的知情人，并仅在此范围内作出判断：\BibleQuote{因为除了人内里的心神外，有谁能知道那人的事呢？}\BibleRef{格前 2:11}但这一点我们断言、我们坚持、我们以一切方式宣布为正当：无人应自己加给自己自愿的死亡，因为这是逃脱现世的患难，而投入永世的祸患；无人应为他人之罪而如此行，因为这是为避免一个不能玷污他的罪过，而招致自身的重罪；无人应为自己过去的罪而如此行，因为他更需要此生，好能借着悔改治愈这些罪；无人应为获得死后所期望的更美好生命而结束此世，因为凡亲手结束自己生命的人，死后并无更美好的生命。\par

% structure: p0085 source=h2 target=subsection reason=relative-heading-rank; base=h2

\subsection{第27章 是否应为避免犯罪而寻求自愿死亡。}

还有一种自杀的理由，我先前曾提及，且被认作颇为合理——即为防止自己因快乐的诱惑或痛苦的暴力而陷入罪恶。如果这理由是正当的，那么我们就应当敦促人，在甫经重生的圣洗后，获赦一切罪恶之后，立即毁灭自己。那时，一切过往的罪业尽消，正是逃避一切未来之罪的最佳时刻。倘若借自杀即可合法地达成此逃离，那为何不特别在此时实行呢？为何任何领受圣洗的人会阻止自己下手自尽？为何一个已脱离此生危险的人，在他可如此轻易地摆脱一切患难，且经上写道：\BibleQuote{喜爱危险的人，必陷于其中}\BibleRef{德 3:27}，却还要重新将自己暴露于危险之中？他为何热爱，或至少面对，如此众多、严重的危险，而滞留于此本可合法离去的人生？然而，是否有人如此盲昧颠倒，远离真理，竟认为：人因害怕一位主人（他的主人）的压迫而引人犯罪，就当结束自己的生命；却仍当活着，将自己暴露于这世界的每时每刻的诱惑中，包括一个主人的压迫所带来的一切邪恶，以及此生无法回避的无数其他苦毒？那么，我们又有何理由耗费时间，借劝勉来激励已领受圣洗者或守童贞的贞洁、或寡妇的节德、或婚姻的忠贞，既然我们拥有远为简单便捷的方法——劝那些刚领受圣洗的人，即结束自己的生命，从而纯洁妥善地归到他们的主那里去？若有人以为此等劝告可行，我不说他是愚人，乃是疯子。如此，他有何面目对人说：\Quote{你自杀吧，免得你生活在一个行为如野蛮人那般淫乱的主人之下时，在细微的罪上再添重罪。}如果他不能不邪恶地说：\Quote{你既已洗净一切罪恶，现在就自杀吧，免得你生活在这个以其不洁的欢乐诱惑人、以可怕的残忍折磨人、以错谬和恐怖克胜人的世界时，再次陷入相似甚至更重的罪过。}那他又怎能如此说呢？如此说既是邪恶的，自杀便是邪恶的。因为若有什么自杀的正当理由，这理由便是了。既然连这也不算正当，就没有任何正当理由了。\par

% structure: p0087 source=h2 target=subsection reason=relative-heading-rank; base=h2

\subsection{第28章 天主以何等裁断，容许仇敌在节德基督徒身上放纵欲情。}

所以，你们这些基督的忠实仆人，虽然你们的贞洁被敌人当作玩物，但不要让你们的生活成为重担。你们有伟大而真实的安慰，只要你们保持无愧的良心，并知道自己未曾同意那些被允许向你们施行罪恶暴行的人的罪。如果你们要问为何允准此事，这确实是世界创造者兼统治者的深奥的\OriginalTerm{眷顾}{providence}；且\BibleQuote{祂的判断何其难测，祂的踪迹何其难寻。}\BibleRef{罗 11:33}然而，你们当忠实地审问自己的灵魂，是否因自己的正直、\OriginalTerm{节制}{continence}和贞洁而过分自高；是否过于贪求这些德行所获得的人的称赞，以致嫉妒那些拥有这些德行的人。就我而言，我不知道你们的内心，所以我不加指责；我甚至听不到你们审问自己的内心时，它们所作的答复。然而，如果你们的内心答复说，情况正如我所设想的可能那样，那么不要惊奇，你们失去了那能赢得人称赞的事物，却保留了那不能向人展示的事物。如果你们没有同意犯罪，那是因为天主在祂的恩宠上加上援助，使恩宠不致失落，且因为羞于见人取代了人的荣耀，使荣耀不致被贪爱。但在这两方面，即使你们中胆怯的人也有安慰，一方面因经验而受认可，另一方面因磨练而受惩戒；一方面得称为义，另一方面得蒙纠正。至于那些在审问时，内心答复说：她们从未因童贞、寡居或婚姻贞洁的德行而骄傲，反而俯就卑微的人，以颤栗之情喜乐于天主的这些恩赐；她们从未嫉妒任何人同样卓越的圣德与纯洁，反而超越人的赞美——人的赞美往往因所赞美的德行稀罕而愈益丰厚——她们反倒愿意自己的人数增多，而不愿因人数稀少而个个显眼——我说，即使是这样的忠信妇女，也不可抱怨天主允准野蛮人如此粗暴地凌辱她们；也不可让自己相信，当天主允准那些无人能犯而不受罚的罪行时，竟忽略了她们的品格。因为有些极其可耻邪恶的欲望，目前因天主隐秘的审判而被放任，以待公开的最终审判。此外，那些因自己有德的贞洁而不自觉有任何过分骄傲的基督教妇女——她们因此而无罪地遭受了掳掠者的暴行——可能仍隐伏着某种软弱，若没有遭受到城陷时所受的羞辱，这软弱或许会使她们陷入骄傲与轻慢的态度。因此，正如有些人被死亡移去，免得邪恶改变他们的性情，这些妇女遭受凌辱，也是免得顺境败坏她们的端庄。所以，无论是那些因仍然身为童贞而已经骄傲的妇女，还是那些若非遭遇敌人的暴力就可能骄傲的妇女，她们都没有失去贞洁，反而获得了谦卑；前者得免于已萌的骄傲，后者得免于即将滋长的骄傲。\par

我们必须进一步注意到，那些受苦者中有些人可能曾以为，节制是身体的善，只要身体不受侵犯就存在，却不明白身体和灵魂的纯洁，系于由天主恩宠所坚固的意志的坚定，且不能强行从不愿的人夺走。他们现在大概已从此错误中解脱出来。因为当他们反思自己如何凭良心事奉天主，当他们再次坚信祂绝不会抛弃那些如此事奉祂并呼求祂援助的人，当他们毫无怀疑地思考贞洁多么中悦于祂时，他们就不得不作出结论：祂绝不可能允许这些灾难临到祂的圣人，假如这些灾难能毁灭祂亲自赐予并乐于在他们身上看到的圣德的话。\par

% structure: p0090 source=h2 target=subsection reason=relative-heading-rank; base=h2

\subsection{第29章——基督的仆人应当如何回答那些指责基督没有从敌人的暴怒中拯救他们的不信者？}

因此，整个至高至真的天主家庭，有自己独有的安慰——这安慰不会欺骗人，且内含比地上摇摇欲坠、终将败坏的事物所能提供的更可靠的希望。他们不会拒绝这现世生命的训练，在其中他们为永生而受教导；也不会哀叹其中的经历，因为他们像旅客一般使用世上的美物，不为所绊，而世上的患难则或考验他们，或造就他们。至于那些在他们受试炼时侮辱他们，当患难临到时说\Quote{你们的天主在哪里？}的人，我们可以反问他们，当他们的神明遭受他们正是为了逃避这些灾难而崇拜、或主张应该崇拜的那些灾难时，他们的神明在哪里？因为基督的家庭已有现成的答复：我们的天主无所不在，全然遍在；不受任何地方的限制。祂能不知不觉地临在，也能不移动而离去；当祂使我们遭受逆境时，或是为考验我们的完善，或是为纠正我们的缺失；而作为我们耐心忍受今世苦难的回报，祂为我们保留了永生的赏报。但你们是谁，竟使我们屈尊与你们谈论你们自己的神明，更不用说我们的天主——那位\BibleQuote{唯应敬畏上主，超越诸神，因为万民的神都属虚无，唯上主造了诸天。}的？\par

% structure: p0092 source=h2 target=subsection reason=relative-heading-rank; base=h2

\subsection{第30章——那些抱怨基督教的人，其实是想要毫无约束地生活在可耻的奢侈中。}

如果那位著名的\Person{西庇阿·纳西卡}如今还在世——他曾是你们的\OriginalTerm{大司祭}{pontiff}，在布匿战争引发的恐慌中，元老院一致推选他为最佳公民，以接待弗里吉亚女神——那么他定会抑制你们这种无耻行径，尽管你们或许几乎不敢正视这样一个人物。因为你们在灾祸中为何抱怨基督信仰，岂不是因为你们渴望无拘无束地享受奢侈的放纵，过一种堕落放荡的生活，而不受任何不安或灾难的干扰？的确，你们对和平、繁荣和富足的渴望，并非出于任何正当使用这些恩赐的目的，即以节制、清醒、自律和虔敬来使用；你们的目的毋宁是在无穷无尽的愚昧享乐中纵情狂欢，从而由繁荣中滋生出一种道德的瘟疫，其灾难性将比最凶残的敌人更甚千倍。正是这种灾难，正是你们的\OriginalTerm{首席大司祭}{chief pontiff}、整个元老院公认的最佳人物西庇阿所担心的，当他拒绝同意毁灭罗马的对手\Place{迦太基}，并反对主张将其毁灭的\Person{加图}时。他惧怕安逸——那种软弱心灵的敌人——并洞悉一种健康的畏惧将是公民合适的监护者。他并没有错；事态的发展证明了他所言之明智。因为当\Place{迦太基}被毁灭，\Place{罗马}共和国解除了它巨大的焦虑之源后，随着事物的繁荣，一系列灾难性的恶果立即接踵而至。首先是和谐被削弱，进而被激烈血腥的暴乱所摧毁；接着，由于一连串祸因，内战随之而来，带来了屠杀、流血、无法无天和残酷的剥夺与掠夺，以至于那些在德行尚存的年代只预期敌人伤害的罗马人，如今德行丧失，反而在同胞手下遭受更为惨烈的暴行。统治欲，连同其他恶习，在罗马人中间比任何其他民族更为炽烈；一旦它占据了少数强有力者，便将其他人——精疲力竭、困顿不堪的人——都置于其轭下。\par

% structure: p0094 source=h2 target=subsection reason=relative-heading-rank; base=h2

\subsection{统治欲在罗马人中如何逐步增长。}

这种激情一旦寓居于骄傲的灵魂，要到何种地步才会止息？它只会通过一连串的推进，直至达到王座。而要获得这样的推进，除了肆无忌惮的野心之外别无他法。而肆无忌惮的野心除非在一个因贪婪和奢侈而腐败的民族中，才能有所作为。此外，一个民族因繁荣而变得贪婪和奢侈；这正是那位极其明智的人\Person{纳西卡}在反对摧毁\Place{罗马}之敌最伟大、最强大、最富裕的城市时所力图避免的。他认为如此则恐惧将成为欲望的约束，欲望受约束便不会在奢侈中放纵，奢侈被阻止则贪婪便告终结；这些恶习一旦被逐，德行便会兴盛，为国家带来巨大的利益；而自由——德行的适当伴侣——将不受束缚地常在。出于类似的理由，并受到同样深思熟虑的爱国情怀所激励，那位你们的\OriginalTerm{首席大司祭}{chief pontiff}——我仍是指那位被全元老院一致裁定为\Place{罗马}最佳人物的人——对元老院在剧院周围建造环形座位的提议泼了冷水，并在一场极有分量的演说中警告他们不要容许\Place{希腊}的奢华风气腐蚀\Place{罗马}的阳刚之气，劝说他们不要屈服于外来放纵的消沉柔弱之风。他的话如此具有权威和力量，以致元老院受到感动，甚至禁止使用那些迄今习惯带进剧院供公民临时使用的长凳。这样一个人，倘若敢于反对他视为神明者的权威，他将会何等急切地想要将戏剧表演本身从\Place{罗马}驱逐出去！因为他不知道他们是恶毒的魔鬼；或者，即便他知道，他也认为应当安抚他们，而非鄙视他们。因为当时尚未向外邦人启示那属天的教义，这教义能以信德净化他们的心灵，以谦卑的虔敬改变他们的自然倾向，使他们从侍奉骄傲的魔鬼转而寻求天上，乃至超越诸天之上的事物。\par

% structure: p0096 source=h2 target=subsection reason=relative-heading-rank; base=h2

\subsection{戏剧表演的设立。}

那么，你们这些对此一无所知的人，应当知道；你们这些假装不知的人，应当被提醒：当你们抱怨那位将你们从这些统治下解救出来的祂时，要知道那些戏剧表演——无耻愚昧和放纵的展览——是在\Place{罗马}设立的，并非出于人们的邪恶欲望，而是由你们的神明所指定。你们本可以为\Person{西庇阿}而非为这样的神明奉献神圣的尊荣，那还更情有可原。神明们并不如他们的\OriginalTerm{大司祭}{pontiff}那样有道德。但现在请留心听，如果你们的心智尚未因深陷谬误的麻醉而失去辨别清醒真理的能力。神明们命令举行戏剧表演以向他们致敬，为要阻止一场肉身的瘟疫；他们的大司祭却禁止建造剧院，以防止一场道德的瘟疫。因此，如果你们还有足够的精神明悟，能选择灵魂胜于身体，那就选择你们所要敬拜的吧。此外，虽然瘟疫止住了，但这并非因为舞台剧的纵情疯狂占据了一个迄今只习惯于竞技表演的好战民族；而是这些狡猾邪恶的精灵，预见到瘟疫不久便会如期消退，便抓住机会感染他们崇拜者的心灵——而非身体——以一种远为严重的疾病。这种瘟疫令神明们大为喜悦，因为它以如此浓重的黑暗蒙蔽了人们的心智，以如此污秽的丑陋羞辱了他们，以至于甚至就在最近（后代能够相信吗？），一些从那\Place{罗马}浩劫中逃出并在\Place{迦太基}找到避难所的人，竟如此深染此疾，以至于日复一日，他们似乎争相在剧院中疯狂地追逐演员。\par

% structure: p0098 source=h2 target=subsection reason=relative-heading-rank; base=h2

\subsection{罗马的覆灭并未纠正罗马人的恶习。}

哦，昏迷的人哪！这是何等盲目，直可谓疯狂，竟占据了你们？我们听闻，连东方民族都为你的覆亡而哀哭，远方强大的邦国也把你的倾覆视作公众的灾难而悲痛，你们自己却如何仍挤进剧场，涌入满座；总之，比以往任何时候都更疯狂地演戏？这便是那污浊的瘟疫斑点，便是\Person{西庇阿}曾力图保全你们免于沦丧的美德与荣誉的毁坏，因此他禁止建造剧场；这便是他渴望你们仍能畏惧敌人的原因，因为他目睹繁荣会多么轻易地败坏并毁灭你们。他不认为那种城墙矗立而道德沦丧的共和国是繁荣的。可是，心怀恶念的魔鬼的诱惑在你们身上比明智之人的防备更有力。因此，你们自己所行的伤害，你们并不许归咎于自身；然而你们所遭受的伤害，却归咎于基督宗教。你们被好运败坏，又未受逆境磨炼，所以在恢复平安稳固的国家时，你们所渴望的并非共和国的安宁，而是自身邪恶奢靡的不受惩处。\Person{西庇阿}愿你们受敌人逼迫，好使你们不致沉溺于奢靡之风；但你们竟如此自暴自弃，即使被敌人击垮，你们的奢靡也未被抑制。你们错失了灾难中的教益；你们变得最为可悲，却仍旧最为放荡。\par

% structure: p0100 source=h2 target=subsection reason=relative-heading-rank; base=h2

\subsection{第34章 论天主仁慈减轻城邑毁灭}

你们至今仍活着，乃因天主存心宽仁，使你们得以警醒悔改，革新生活。正是祂容许你们这些忘恩之人逃脱敌人刀剑，或藉称自己为祂仆人，或藉避难於殉道者之圣地。\par

相传\Person{罗慕路斯}与\Person{雷穆斯}为增加所建城市人口，开设一庇护所，使人人可得避难与免罪——此乃预表近日为尊崇基督所发生之事。罗马毁灭者效法建城者之先例。然建城者之所以行此事者，为增其公民之数；而毁灭者之所以行此事者，为使敌人之数不至减损；故建城者所为不足为荣也。\par

% structure: p0103 source=h2 target=subsection reason=relative-heading-rank; base=h2

\subsection{第35章 论藏于恶人中的教会之子和教会内假基督徒}

让主基督的赎民家庭和君王基督的旅居之城以这些及类似的答复（若有更充分、更恰切的答复）回答他们的敌人。但愿此城记念，在敌人中隐伏着那些预定将为其同胞的人，这样她才不会认为忍受敌人所加诸的苦待是徒劳的，直到他们成为信仰的宣认者。同样，天主之城，当其旅居此世时，在自己的共融中，以圣事与之联结的，也有一些人将不在圣者的位分上永居。这些人中，有些现尚未被认出；另一些人则公开表明，且大胆地与我们的敌人一同抱怨天主，尽管佩戴着祂的圣事标记。今日你们可见到这些人与我们一道挤满教堂，明日则与不敬虔者拥挤戏院。然而，即使对于这样的人，我们也有更少的理由对他们的重归绝望，若在最公开的敌人中如今也有某些人，连他们自己也不知晓，却预定将成为我们的朋友。事实上，这双城在此世互相纠缠混杂，直至末日的审判实施其分离。现在我要在天主的助佑下，开始论述这双城的兴起、进展与结局；我写这些，是为光荣天主之城，为使它与另一城相比，能发出更璀璨的光辉。\par

% structure: p0105 source=h2 target=subsection reason=relative-heading-rank; base=h2

\subsection{第36章 在后继讨论中将要处理哪些主题}

但我仍有一些话，为反驳那些将罗马共和国的灾难归于我们宗教的人，因为我们的宗教禁止向邪神祭献。为此目的，我必须记述该城及其所属行省在禁止祭献之前所遭受的全部或足以说明问题的灾难；因为若非我们的宗教在彼时早已照耀他们并禁止其祭献，这些灾难他们无疑都会归咎于我们。然后我必须接着阐明，那位掌管一切王国的真天主，愿意赐予他们何等的社会福祉以使其帝国渐增。我必须说明祂为何如此行，以及他们的假神非但丝毫未助益他们，反而以诡诈欺骗大为加害。最后，我必须面对那些，在这点上被无可辩驳的明证说服折服后，仍力图坚持他们敬拜那些神，不是为人生现今的利益，而是为死后可享的福乐。倘若我没有弄错，这将是本著最艰难的部分，也值得最高华的论证；因为我们必须与那些哲学家交锋，不是普通的哲学乌合之众，而是最著名的哲士，他们在许多观点上与我们一致，例如关于灵魂不死，真天主创造世界，并藉其天主眷顾管理所造一切。然而，既然他们在其他观点上与我们分歧，我们便不应畏缩揭露其错误，以便仰赖天主赐予的才能，驳倒恶人的反驳，从而维护天主之城、真宗教及敬拜天主，唯独这敬拜结合了真实永恒幸福的许诺。那么，就让我们在此结束，以便在新一卷书中进入这些主题。\par

\SourceRecord{Translated by Marcus Dods.；Nicene and Post-Nicene Fathers, First Series；Vol. 2.；Edited by Philip Schaff.；Buffalo, NY: Christian Literature Publishing Co.,；1887.；<http://www.newadvent.org/fathers/120101.htm>.}

% structure: p0001 source=h1 target=section reason=page-title

\section{\WorkTitle{天主之城}（第二卷）}

在这本书中，奥古斯丁回顾了那些罗马人在基督的时代之前，且在伪神的崇拜被普遍遵行时，所遭受的灾难；并证明，罗马人非但没有因他们的神而得免于不幸，反而被他们陷入那唯一的，或至少是最大的灾难中——即道德的腐化与灵魂的恶习。\par

% structure: p0003 source=h2 target=subsection reason=relative-heading-rank; base=h2

\subsection{必须对答复对手的必要性加以限制}

如果人类脆弱的理智不妄图抵抗那清晰明白的真理，而是服膺于有益的教义，如同服膺于医治疾病的良药，直到借着信德和虔敬，从天主获得那为医治它所需的恩宠，那么，那些拥有正确观念，并用恰当言语表达的人，便无需长篇大论，来反驳空洞臆测的错误了。然而，这种理智的软弱如今比以往更普遍、更贻害无穷，以致即使真理已经被证明得如同人能向人证明的那样充分，他们仍将自己不合理性的幻想视为真理本身，这或是由于他们极大的盲目，使他们看不清那明明摆在他们眼前的道理，或是由于他们固执己见的顽梗，使他们无法承认自己所见之事的效力。因此，我们时常不得不对那些已然清晰的问题，作更详尽的说明，以便，不妨说，不仅将它们呈现在眼前，甚至呈现于触感之前，好叫那些闭眼不顾的人也能感觉得到。但是，如果我们坚持一个原则，即我们必须总去答复那些反驳我们的人，那么我们的讨论将会有怎样的结局，我们的言论又能有何界限呢？因为那些或是不能理解我们的论证，或是在反驳的习惯中变得如此刚硬，以致虽能理解却仍不肯屈服的人，他们反驳我们，并且，正如经上所载，\Quote{说顶撞的话}，且是不可救药地骄矜。如果我决意在他们每次厚颜漠视我们的论证时，便驳斥他们的异议，或是在他们每次能以任何方式反驳我们的陈述时，便加以回应，你就会看到，我们承担的是何等无休无止、徒劳无益而又令人痛苦的工作了。因此，我的儿子\Person{马凯利努斯}啊，我不愿意我的著作被你，或是任何其他那些我为甘愿并以基督徒爱德而奉献此书的人，这样来评判——如果你们总是要求，对于你们所听见的针对读到的每一异议，都要予以答复的话；因为这样，你们就会变得如同宗徒所说的那些无知妇女，她们\BibleQuote{永远学习，而总达不到认识真理的地步。}\BibleRef{弟后 3:7}\par

% structure: p0005 source=h2 target=subsection reason=relative-heading-rank; base=h2

\subsection{第一卷内容概述}

在前一卷书中，我已开始谈论\WorkTitle{天主之城}，并已决意，仰赖天主的助佑，将全书奉献于此主题。我最初努力，便是要答复那些人，他们将使世界疮痍满目的战争，尤其是新近蛮族蹂躏\Place{罗马}的事件，归咎于基督的宗教，因为基督宗教禁止祭献魔鬼。我指出，他们更应将此事归功于基督，因为正是因祂的圣名，蛮族违背一切惯例与战争法，将最大的教堂开放为避难所；在许多事例中，他们对基督表现出如此的敬畏，以致不仅是基督的真正仆人，就连那些因恐惧而假扮为仆人的人，也都免于遭受按战争惯例本可合法加予的一切苦难。由此产生了这样的问题：为什么那邪恶而负恩的人，竟被允许分享这些利益；又为什么战争的艰苦与灾难，竟同样落在虔诚与不虔敬的人身上。在对这重大问题作出适切充分的答复时，我花费了相当可观的篇幅，这一方面是为了消除许多人的焦虑，当他们观察到天主的祝福，以及日常人间的祸福，无差别地临于善人与恶人时，他们便感到不安；但更重要的是，我要为那些遭受敌人凌辱的圣洁的妇女们，献上慰藉。那些凌辱虽刺激了她们的端庄，却未污损她们的贞洁；我愿她们不致因生命而羞惭，尽管她们并无过错可羞。随后，我还简要地驳斥了那些以最无耻的放纵，侮辱那些遭受灾难的可怜基督徒，尤其是侮辱那些心碎哀恸、谦卑自抑的、圣洁的妇女们的人；这些家伙自身，却是最堕落不堪、毫无丈夫气概的败类，完全退化，不配称为真正的罗马人。真罗马人的赫赫事功，在历史上被大量记载，并到处受人称颂，而他们的后代，却是他们荣誉的最大仇敌。实在说来，罗马这座由古英雄的勤劳所建立并壮大的城，当它的围墙仍矗立时，就已被他们的后代更可耻地毁灭了，比现今城墙被拆毁更甚。因为在这次毁灭中，倒塌的是石头和梁木；而在那些败类所导致的毁灭中，倒塌的却是城邑的道德屏障与华饰，而非砖石的壁垒；他们心中燃烧的情欲，较那焚烧他们房屋的火焰，更具毁灭性。我就这样结束了我的第一卷。如今我继续谈论，自这座城及其所属各行省立城以来，它们所遭受的那些灾难；如果他们在那早期的时代，反对他们那些虚假欺骗之神的福音教义，也像现在这样广博自由地被宣扬，那么他们同样会将所有这些灾难，完全归咎于基督宗教。\par

% structure: p0007 source=h2 target=subsection reason=relative-heading-rank; base=h2

\subsection{我们只需阅读历史，便可见到在基督宗教开始与诸神崇拜竞争之前，罗马人所遭受的灾难}

但请记得，在叙述这些事时，我仍需要应对无知之辈；他们无知到了竟产生出那句常言：\Quote{干旱与基督教形影不离。} 他们当中确实有一些受过良好教育并对历史有鉴赏力的人，我所谈论之事物在他们观察之下是明摆着的；但是，为了煽动无知的群众敌视我们，他们假装对这类事件无知，尽其所能使俗众相信：那些在特定地方及特定时期普遍加于人类的灾难，乃是基督教的后果，而基督教正日益广传，且拥有令他们自己的诸神黯然失色的荣名与光辉。那么，让他们同我们一起回想，在基督降生成人之前，在祂的名号以他们徒然嫉妒的荣耀传扬于万民之前，罗马的繁荣曾被如何多样且屡屡发生的灾难所摧残。让他们在这件事上为他们的神辩护吧，如果他们能够的话；因为他们坚持说自己崇拜诸神，就是为了免受这些灾难；他们如今只要遭受丝毫灾难，便归咎于我们。因为，在我将要讲述的灾难中，这些神为何在基督之名的宣讲冒犯到他们并终止其祭献之前，听任这些灾难落在自己的崇拜者身上呢？\par

% structure: p0009 source=h2 target=subsection reason=relative-heading-rank; base=h2

\subsection{诸神的崇拜者从未从他们那里得到任何健康的道德诫命，且在举行其敬拜时，种种污秽之事皆被实行。}

首先，我们要问，为什么他们的神不采取步骤改善其崇拜者的道德。真天主不理睬那些不寻求祂帮助的人，这是公正的；但那些神，其崇拜者如今正抱怨被禁止崇拜的神，为什么没有颁布任何本可引导其虔信者走向德性生活的律法呢？确实，人类关心诸神的敬拜，诸神反过来也关心人类的行为，这本是公正的。但是，他们回答说，人步入歧途是出于自己的意志。谁否认这点呢？但这丝毫不减损这些人类守护神的责任，他们本应以明白的言辞颁布美好生活的律法，而非向他们的崇拜者隐匿。他们理应派遣先知，去教导并谴责违犯这些律法的人，并公开宣告那等待着作恶者的惩罚，以及行善者可期盼的赏报。可曾有哪座神殿的墙壁回荡过这样的警告之声？我年轻的时候，有时会去参加那些亵渎神圣的娱乐和表演；我看见司祭们宗教狂热中发狂，听见唱诗班的声音；我喜爱那些为尊崇诸位神祇、女神，以及童贞的\Person{刻莱斯蒂斯}和众神之母\Person{贝雷辛提亚}而举行的可耻赛会。在她被祝圣的净化节日里，在她卧榻前唱出的作品，其歌词是如此淫秽龌龊——我甚至不说那是众神之母，而是任何元老或正派人的母亲都无法入耳的——更确切地说，是如此不洁，以至于连那些污言演员自己的母亲，也无法成为听众之一。因为对父母的自然尊敬，是一种连最堕落者也无法忽视的纽带。因此，这些演员为尊崇众神之母而做出的淫秽动作和污言秽语，当着广大男女观众的面，他们羞于在自家母亲面前排演。我设想，那些因好奇而从各处聚集的人群，其受到冒犯的廉耻，定会因羞愧而四散。如果这些就是神圣礼仪，那什么是亵渎？如果这是净化，那什么是污秽？这个庆典被称为\OriginalTerm{桌节}{Tables}，仿佛在设宴，让不洁的魔鬼可以在此找到合适的享受。因为不难看出，喜欢这类淫秽的必然是何种邪灵，除非有人被这些假托神名的那恶灵所蒙蔽，或是不信其存在，或是过着一种宁愿讨好畏惧它们而非真天主的生活。\par

% structure: p0011 source=h2 target=subsection reason=relative-heading-rank; base=h2

\subsection{论为尊崇众神之母而行的淫秽之事。}

在本事上，我更愿请那些并非乐于这些丑行、反而致力终止它们的人来作我的判断同席，而是那个被元老院选为最堪当亲手接受那邪神\OriginalTerm{西布莉}{Cybele}像并引入城中的\Person{西庇阿·纳西卡}。他会告诉我们，他是否骄傲于见到自己的母亲受国家如此尊崇，以致被判予神圣的荣耀；正如希腊人、罗马人及其他民族，将他们曾受惠者判予神圣的荣耀，并相信他们凡间的恩人因此成为不死，列于神位之中。可以肯定，若可能，他必愿自己的母亲享有这样的幸福。但如果我们继续问他，在向她表敬的种种仪式中，他是否愿举行如此可耻的祭礼，他不会立刻扬声说，他宁愿他的母亲僵死，也不愿她作为一个女神，聆听这些猥亵之词吗？一位品行如此严谨，以致运用其罗马元老的影响力阻止在该城建造一座专奉阳刚德行的剧院的人，难道会希望他的母亲以一位罗马主妇当年听闻便会颊染羞赧的言词，被当作女神来求情吗？他岂能相信，一位可敬妇女的端庄，会因被尊为神而彻底转变，以至于她竟任凭人以如此粗鄙无礼的言语来呼求、颂扬；若她尚在人间时听到此类话语，而不仅不掩耳速离，反倒听进去，她的亲属、她的丈夫、她的子女都会为她羞惭？因此，众神之母正是那种最放荡者也不耻有之为母的品格，她欲迷惑罗马人的心灵，便要求他们最优秀的公民为她服役，不是要以有益的劝言来更成全他的德行，而是要以她的诡诈来缠累他，如同经上论到的那种女人：\BibleQuote{淫妇猎取宝贵的生命。}\BibleRef{箴 6:26}她的意图是以一份表面上的神圣证词，去夸扬这位心志高远之人的卓越，好让他倚仗自己德行的卓异，而不再致力于真正的虔诚与宗教；缺少了这些，天赋的才智不论何等辉耀，终将化为骄傲，归于虚幻。因为，那女神若在她的神圣节庆中要求最优秀者听闻都会羞愧难当的丑事，那么她指定要最优秀的人，不是出于诡诈的用意还能是什么呢？\par

% structure: p0013 source=h2 target=subsection reason=relative-heading-rank; base=h2

\subsection{第六章：论异教诸神从未教诲圣洁的生活。}

就是由于这个缘故，那些神明全然忽视了崇拜他们的城市与民族的生活及道德，未曾投放任何可畏的禁令，以阻止他们彻底败坏，并保护他们免于那些可怕又可憎的邪恶——这些邪恶打击的不是庄稼与葡萄的收成，不是房屋与财产，也不是隶属于灵魂的身体，而是灵魂本身，那辖管全人的心神。若有此类禁令，就请出示，请证实。他们会告诉我们，在那些获授宗教秘仪的人中，已教诲了纯洁与正直，并对\Emph{精英}耳语了隐秘的德行劝勉；但这不过是虚夸。他们应为我们指出或列举，曾有哪一处场地，被祝圣为聚会之所，在那里，人们非但不听闻优伶猥亵的歌曲与放荡的表演，也不举行那些最污秽无耻的\OriginalTerm{富伽利亚节}{Fugalia}（名称起得好，因为它驱逐了端庄与正感），反而奉诸神之名被命令约束贪婪、遏制淫乱、战胜野心；简言之，人们能够在那里学到\Person{佩尔西乌斯}猛烈鞭策他们的那所学校所教导的，他说：\Quote{学习吧，你们这些被弃的受造之物，探究万物的因由；我们是什么，我们为何而生；我们生活顺利的法则是什么；我们凭何技巧可绕过终点而不致触礁；我们应对财富设下何种限度，我们可合法渴求什么，污秽的利益有何用处；我们应为国家与家庭奉献多少；简言之，学习天主命你何所是，祂定你身居何位。}让他们给我们举出那些惯常由神明传达这类训诲、且崇拜他们的百姓惯常前往聆听的处所吧，正如我们能指出，凡基督教信仰所及之地，为此目的而建造的教堂一样。\par

% structure: p0015 source=h2 target=subsection reason=relative-heading-rank; base=h2

\subsection{第七章：论哲学家的劝言何以不能产生道德效果：因其缺乏神圣教导的权威，且人天生的恶倾向驱使人效法神明的榜样，而不是服从人的规诫。}

但也许他们会提醒我们哲学家学派及其辩论？首先，这些不属于\Place{罗马}，而属于\Place{希腊}；即使我们承认它们现在是罗马的，因为\Place{希腊}本身已成为罗马的一个行省，哲学家的教导仍不是诸神的诫命，而是人的发现；他们受自身思辨能力的驱使，努力探究自然的隐秘法则，伦理学中的是与非，以及辩证法中何者符合逻辑规则，何者不合逻辑与错误。其中一些人，蒙天主的助佑，取得了伟大的发现；但当他们放任自流时，便被人性的软弱所出卖，陷入谬误。这是由\OriginalTerm{天主眷顾}{divine providence}所安排的，为抑制他们的骄傲，并以他们为例指明，只有谦卑才能通往至高之境。关于这一点，若真理的主天主允许，我们将在适当的地方再作详述。但若哲学家们有某些发现足以引导人们走向美德与幸福，给他们授予神圣的荣耀岂非更公正？阅读\Person{柏拉图}的著作于一座\Quote{柏拉图殿}中，岂不比置身魔鬼的庙宇，目睹\Person{库柏勒}的祭司自阉，不男不女者受祝圣，疯狂的信徒自割，以及其他这些神的仪式所规定的残忍或可耻、或者可耻地残忍或残忍地可耻的礼节，更符合一切美德的情操？让青年们聆听诸神的律法公开宣讲，而非徒然地赞美祖先的习俗与法律，岂非更适宜的教育，也更可能激励他们向善？的确，所有罗马诸神的崇拜者，一旦被\Person{佩尔西乌斯}所谓的\Quote{情欲的灼人毒药}所占据，便宁愿观看\Person{朱庇特}的所作所为，也不愿听\Person{柏拉图}的教诲或\Person{卡托}的指责。因此，\Person{泰伦斯}剧中那个放荡的年轻人，当他看到墙上的一幅壁画描绘了寓言中\Person{朱庇特}化作金雨降入\Person{达那厄}怀中的情景，便以此作为自己放纵的权威先例，并夸耀自己是神的模仿者。\Quote{那是什么神？}他说，\Quote{那位用他的雷霆震撼至高殿宇的神。而我，与他相比一个可怜的生物，还要犹豫吗？不；我干了，而且全心全意。}\par

% structure: p0017 source=h2 target=subsection reason=relative-heading-rank; base=h2

\subsection{第8章：戏剧表演公开展示诸神可耻行径，却取悦他们而非冒犯}

但是，有人会插话说，这些都是诗人的寓言，并非诸神自己的启示。那么，我无意在戏剧表演和神秘仪式的淫秽之间作出裁断；我只说一点，且历史为我作证：那些以诗人虚构为主要吸引力的表演，并非因\Place{罗马}人无知的热忱而被引入诸神的节庆，而是诸神自己下达了最紧急的命令，甚至强求\Place{罗马}人向他们奉上这些庄严的庆典和仪式。我在前一卷已提及此事，并指出戏剧表演最初是在\Place{罗马}因一场瘟疫，由大司祭授权而举行的。有谁不更倾向于采纳那些具有神圣认可的戏剧中所展现的榜样，来规范自己的生活，而非那些仅凭人的权威所书写和颁布的戒律呢？如果诗人们将\Person{朱庇特}描述为奸淫者，是对他的歪曲，那么贞洁的诸神理应愤怒地报复如此邪恶的虚构，而不是鼓励传播这虚构的表演。这些戏剧中，最无害的是喜剧和悲剧，即诗人专为舞台所写的剧作；它们虽常处理不洁的题材，却没有许多其他表演特有的污言秽语；正是这些戏剧，长辈们强迫子弟阅读并学习，作为所谓自由而绅士教育的一部分。\par

% structure: p0019 source=h2 target=subsection reason=relative-heading-rank; base=h2

\subsection{第9章：希腊人依顺诸神所允许的诗的放纵，被古罗马人所限制}

古罗马人对此事的看法，由 \Person{西塞罗} 在其著作 \WorkTitle{论共和国} 中予以证实；在这部作品中，对话者之一 \Person{西庇阿} 说：\Quote{除非社会风俗事先认可了同样的淫荡，否则观众绝不可能容忍喜剧中的淫荡。} 而在早期，希腊人在放纵中保留了一定的克制，并制定法律，规定喜剧若想批评任何人，必须指名道姓。因此在 \Person{西塞罗} 的同一著作中，\Person{西庇阿} 还说：\Quote{它没有诽谤过谁？不，它没有苛责过谁？它放过谁了？就算它攻击那些乱政者和结党营私之辈，危害国家的人——如 \Person{克里昂}、\Person{克利奥丰}、\Person{希佩波罗斯}。这尚可容忍，虽然由公共监察官来谴责这些人，比由诗人来嘲讽他们更为得体；但是，\Person{伯里克利} 以最高尊严在战争与和平中主持国家大局，却用粗俗的诗句玷污他的名声，这种做法与诗人的身份极不相称，就如同我们的 \Person{普劳图斯} 或 \Person{奈维乌斯} 将 \Person{普布利乌斯·西庇阿} 和 \Person{格奈乌斯·西庇阿} 搬上喜剧舞台，或者 \Person{凯基利乌斯} 讽刺 \Person{加图} 一样。} 紧接着他继续说：\Quote{虽然我们的 \WorkTitle{十二铜表法} 只对极少数罪行规定了死刑，但这少数之中就有这么一条：若有人吟唱诽谤性讽刺诗，或撰写旨在给他人带来恶名或耻辱的讽刺作品。这是明智的规定。因为我们的生命应当由长官的裁决和明达的正义来判断，而不是由诗人轻浮的幻想来评判；我们也不应被迫聆听诽谤，除非我们有权回复，并在适当的法庭前为自己辩护。} 我从 \Person{西塞罗} 的 \WorkTitle{论共和国} 第四卷中援引了这些，认为这是明智的；我逐字引用了原文，但为便于理解而省略了一些词，并稍作调换。这段摘录显然与我试图解释的问题密切相关。\Person{西塞罗} 又做了一些评论，并总结说，古罗马人不允许在舞台上对任何活人进行赞美或谴责。但希腊人，如我所说，虽然不那么道德，但在允许这种罗马人所禁止的放纵时，却更符合逻辑；因为他们看到，他们的神明赞同并欣赏低俗喜剧中那些不仅针对人、甚至针对他们自身的粗鄙言辞；而且，无论归给他们的那些卑劣行为是诗人的虚构，还是他们真实的恶行被纪念并在剧院中演出，都是如此。但愿观众认为他们只配嘲笑，而不配效仿！显见，当神明们都不介意自己的名誉受到玷污时，要饶过领袖人物和普通公民的好名声，那便是极度的傲慢了。\par

% structure: p0021 source=h2 target=subsection reason=relative-heading-rank; base=h2

\subsection{魔鬼任凭人们将虚假或真实的罪行归于他们，意在害人}

人们为这种做法辩解称，关于神明的故事并非真实，而是虚假和纯粹的虚构，但这只会使事情变得更糟，如果我们根据我们的宗教所教导的道德来判断；并且考虑到魔鬼的恶意，他们还能对人类施行什么更狡猾、更诡诈的计谋呢？当诽谤指向一位正直而有益于国家的领袖时，其可谴责的程度不是与它的不实和毫无根据成正比吗？那么，当神明成为如此邪恶、如此蛮横的不义行为的对象时，又有什么惩罚足够呢？但那些被这些人尊为神明的魔鬼，却乐意于甚至将那些他们本无罪过的恶行都归于自己，只要他们能用这些观念的网罗缠住人的心灵，并吸引他们与自己一同走向那预定的惩罚：无论是这些魔鬼因喜爱人类的愚昧而让人当神明来崇拜的那些人确实做过此类事情，而魔鬼则通过千般恶毒、欺骗的伎俩取代他们，从而接受崇拜；或者，即使这些确实是人类的罪行，这些邪恶的灵也乐于让它们归于更高的存在，以便似乎从天界本身传达出充分的许可，让人去犯下可耻的罪恶。因此，希腊人看到自己所侍奉的神明的品性，便认为诗人确实不应克制在舞台上揭露人类的恶行，这不是因为他们希望在这方面与自己的神明相似，就是因为他们害怕，若要求自身拥有比他们声称神明所拥有的更洁白无瑕的名声，可能会激怒神明。\par

% structure: p0023 source=h2 target=subsection reason=relative-heading-rank; base=h2

\subsection{希腊人允许伶人担任公职，理由是讨神喜悦的人不应被同伴轻看}

希腊人同样以其理性精神，将相当显赫的公民荣誉授予这些戏剧的演员。在上文提及的 \WorkTitle{De Republica} 一书中，提到一位口才极佳的雅典人 \Person{Aeschines}，他年轻时曾是悲剧演员，后来成为政治家；雅典人还曾一再派遣另一位悲剧演员 \Person{Aristodemus} 作为前往 \Person{Philip} 的全权大使。因为他们认为，那些在戏剧演出中担任主要演员的人，既然看到戏剧令众 \OriginalTerm{神}{gods} 如此喜悦，就不宜将他们定为无耻之人而加以谴责。毫无疑问，这是希腊人的不道德行为，但同样毫无疑问的是，他们的行为与其众神的品性相符；因为众神允许甚至要求诗人与演员撕碎他们的神圣名声，既然如此，他们又怎能保护公民的行为免遭言论攻击？既然他们已经确知，他们所崇拜的众神喜爱这些戏剧，他们又怎能鄙弃那些在剧场中表演这些戏剧的人？不仅如此，他们怎能不授予他们最高的公民荣誉？他们怎能一方面尊敬那些为他们向众神献上可悦纳祭品的司铎，另一方面却将那些代表人民向众神提供其所要求的喜乐或荣耀的演员打得声名狼藉——依照司铎们的说法，若未得到这些喜乐与荣耀，众神便会发怒。饱学的 \Person{Labeo} 在此问题上堪称权威，他认为，敬拜善恶 \OriginalTerm{善恶神明}{good and evil deities} 的区别应体现为不同的敬拜仪式；邪恶的神应以血腥祭品与哀恸之礼来平息，而善良的神则当以欢乐愉快的庆祝来敬拜，例如（如他自己所说）\Emph{e.g.}，戏剧、节庆与宴会。这一切，我们将在天主的助佑下日后讨论。目前，回归本题：无论是向所有神祇不分善恶地呈献一切祭品——仿佛众神皆为善类（而设想有恶神存在本就不合宜；然而异教徒的这些神全是恶者，因为他们并非神，而是 \OriginalTerm{魔鬼}{evil spirits}）——还是按 \Person{Labeo} 的看法，向不同的神祇呈献不同的祭品，希腊人都同样有理由，既尊敬献祭的司铎，也尊敬演戏的演员，以免背上冒犯所有众神的罪名——如果戏剧令所有神祇悦乐；或者更糟，只令那些善良的神祇悦乐。\par

% structure: p0025 source=h2 target=subsection reason=relative-heading-rank; base=h2

\subsection{第12章 罗马人拒绝给予诗人针对人的同样许可，却允许他们针对众神，由此显示出他们对自己比对众神更为细腻的敏感。}

然而，正如 \Person{Scipio} 在同一场讨论中所夸耀的那样，罗马人拒绝让他们的行为与良好声名遭受诗人们的攻击与诽谤，甚至规定若有任何人胆敢撰写此类诗句，即构成死罪。就他们自身而言，这诚然是极为可敬的举动；但就对待众 \OriginalTerm{神}{gods} 而言，却是傲慢而不虔敬的：因为他们明知，众神不仅容忍，而且乐于受到诗人侮辱性言辞的攻击，而他们自己却不容忍受同样的对待；他们的宗教礼仪规定为可悦纳于神的，他们的法律却禁止为有害于自身者。那么，\Person{Scipio}，你怎能赞扬罗马人拒绝给予诗人这种许可，使任何公民不受诽谤，而你却知道众神并不在这种保护之列？难道你认为你们的 \Place{元老院} 比 \Place{Capitol} 更值得敬重？难道在你眼中，整个罗马城比满天的神祇更有价值，以至于你禁止诗人对任何公民口出恶言，却允许他们不受惩罚地对众神任意诋毁，而无需 \OriginalTerm{元老}{senator}、\OriginalTerm{监察官}{censor}、\OriginalTerm{元首}{prince} 或 \OriginalTerm{大祭司}{pontiff} 的干预？诚然，\Person{Plautus} 或 \Person{Nævus} 攻击 \Person{Publius} 与 \Person{Cneius Scipio} 是不可容忍的，\Person{Cæcilius} 讽刺 \Person{Cato} 是难以忍受的；然而，你们的 \Person{Terence} 以至尊 \Person{Jove} 的邪恶榜样煽动青年的淫欲，却被认为是完全合宜的。\par

% structure: p0027 source=h2 target=subsection reason=relative-heading-rank; base=h2

\subsection{第13章 罗马人应当理解，凡期望在放荡的娱乐中受人敬拜的神，不配获得神圣的尊荣。}

然而，\Person{西庇阿}若还在世，或许会如此回答：\Quote{我们怎能对神明亲自祝圣的事物施加惩罚呢？因为那些上演、表演这类言行的戏剧娱乐，乃是神明引入罗马社会的，他们命人将这些娱乐奉献并演出，以尊崇自己。}但这岂不正是最明显的证据，证明他们并非真神，也根本不配从共和国领受神圣的荣耀吗？假如他们曾要求，为了尊崇他们，需拿罗马公民来取笑，那么每一个罗马人都会对这一可憎的提议表示愤慨。那么，我要问，当他们提议用自己的罪行作为颂扬自己的素材时，他们怎能被认为配受敬拜呢？这种伎俩岂不暴露了他们，并证明他们乃是可憎的魔鬼吗？因此，罗马人虽然迷信到侍奉那些毫不掩饰渴望在淫荡戏剧中受敬拜的为神明，却仍充分顾念他们祖传的尊严与德性，以致拒绝将希腊人给予演员的那种奖赏给予演员。关于这一点，\Person{西塞罗}记录了\Person{西庇阿}的证言：\Quote{他们[罗马人]视喜剧及一切戏剧表演为可耻之事，因此不仅禁止演员担任普通公民可得的官职与荣誉，还下令监察官将其名字打上烙印，并从部落名册中删除。}这真是一项卓越的法令，也是罗马睿智的又一明证；但我但愿他们的明智能更彻底、更一贯。因为，当我听说，若有任何罗马公民选择舞台为业，他不仅断绝了自己一切可嘉的职业，甚至被逐出自己的部落时，我不禁高呼：这才是真正的罗马精神，这才配得上一个珍视自身声誉的邦国。但接着有人打断了我的狂喜，质问道：演员被剥夺一切荣誉，而戏剧却被算作对神明应有的尊荣，这一致性何在？因为，在很长一段时间里，罗马的道德未经戏剧表演污染；倘若它们是为了迎合公民的口味而被采用，那么它们便会与风气的松弛携手引入。但事实是，正是神明要求上演它们以取悦自己。那么，敬拜这神的演员反被逐出，这有什么公义可言？你们怎能一面敬拜那强要戏剧的神，一面又羞辱那演剧的人呢？这就是希腊人与罗马人所卷入的争论。希腊人认为，他们给予演员尊荣是正当的，因为他们敬拜那些要求戏剧的神明；罗马人则不同，他们不容许演员以其名字玷污自己的平民部落，更不用说元老院阶层了。整个讨论可归结为以下三段论。希腊人提供了大前提：若应敬拜这样的神明，则这样的人也必定应受尊崇。罗马人补充了小前提：但这样的人绝不应受尊崇。基督徒得出结论：因此，这样的神明绝不应受敬拜。\par

% structure: p0029 source=h2 target=subsection reason=relative-heading-rank; base=h2

\subsection{\Person{柏拉图}将诗人逐出有序之城，胜于那些渴望以戏剧受尊崇的神明。}

我们还需探究，为何那些编写戏剧的诗人——按\WorkTitle{十二铜表法}，他们被禁止损害公民的名誉——却比演员更受尊重，尽管他们可耻地诽谤了神明的品格？那些诗意的、辱神的作品，其演员被打上烙印，而其作者却受尊崇，这难道正当吗？在此，我们难道不应将桂冠授予一位希腊人，也就是\Person{柏拉图}吗？他在构建他的\WorkTitle{理想国}时，主张应将诗人作为国家的敌人逐出城邦。他无法容忍神明被诗人的虚构所玷污，也无法容忍公民的心智因诗人的虚构而堕落、麻木。现在，请将你所见于\Person{柏拉图}的人性——他将诗人逐出城邦，以免公民受害——与你所见于这些神明身上的神性作一比较：这些神明为了自身的尊荣，强要上演戏剧。\Person{柏拉图}曾力图——尽管未能成功——说服那些轻浮、淫逸的希腊人，连写作这类戏剧也当戒绝；而这些神明却运用其权威，从庄重、清醒的罗马人那里强行索取了这些戏剧的表演。而且，他们还不满足于仅仅上演，更将这些戏剧奉献给自己、祝圣给自己，为自身的荣耀而隆重举行。那么，一个国家将神圣的荣耀授予哪一方更为适宜——是授予禁止这些邪恶放荡戏剧的\Person{柏拉图}，还是授予那些乐于蒙蔽人眼、使之无法领悟\Person{柏拉图}力图教诲之真理的魔鬼呢？\par

\Person{柏拉图}这位哲学家被\Person{拉贝奥}提升到半神的地位，从而与如\Person{赫丘利}、\Person{罗慕路斯}之流平起平坐。\Person{拉贝奥}将半神置于英雄之上，但两者他都算作神明。然而我毫不怀疑，他认为这个被他视为半神的人，不仅比英雄更值得尊重，甚至比众神本身也更值得尊重。罗马人的法律与\Person{柏拉图}的学说有相似之处：后者对诗的虚构一概谴责，而前者则限制讽刺的自由，至少当讽刺的对象是人时。\Person{柏拉图}甚至不容许诗人居住在他的城邦里；罗马的法律禁止演员登记为公民；若非害怕得罪那些曾要求演员服务的众神，他们很可能已将演员彻底驱逐了。因此显而易见，罗马人不能、也不会合理地期望从他们的神那里领受规范其行为的法律，因为他们自己制定的法律远远超越诸神的道德，并使后者蒙羞。众神要求以舞台剧来荣耀自己；罗马人则将演员排除于一切公民荣誉之外；前者命令通过戏剧表演来颂扬他们自身的丑行；后者命令任何诗人不得玷污任何公民的名誉。然而那位半神\Person{柏拉图}抗拒了这等神祇的欲望，并向罗马人指出了他们自己的天才所未完成之事；因为他完全将诗人排除在他的理想国之外，无论他们是创作不顾真理的虚构，还是以神圣行为的伪装为可怜人树立最坏的榜样。至于我们，确实将\Person{柏拉图}既不算作神，也不算作半神；我们甚至不会将他与天主的任何一位圣天使相比，不会与传达真理的先知相比，不会与基督的任何宗徒或殉道者相比，甚至不会与任何一位忠信的基督徒相比。我们持此看法的理由，在天主的助佑下，我们将在适当之处予以说明。然而，既然他们希望他成为半神，我们认为他确实更有资格获得那个品级，且在各方面都更为优越，即便不能与\Person{赫丘利}和\Person{罗慕路斯}相比（尽管没有任何史家讲述过、也没有任何诗人传唱过他杀害兄弟或犯下任何罪行），但肯定胜过普里阿普斯、犬头神或热病女神——这些神明部分是从外邦人那里引进的，部分则是通过国内祭礼祝圣的。那么，这样的神祇又怎能被期望去颁布良善而有益的法律，无论是为预防道德与社会的恶行，还是为革除已滋生的恶行？——这些神祇利用自己的影响力，甚至去播种和滋养放荡，他们规定，那些或真或假归于他们的行为，应通过戏剧表演向民众公布，从而以看似神圣嘉许的气息，无端地煽起人类欲望的火焰。\Person{西塞罗}论及诗人时，徒然地以此言反对这种状况：\Quote{当诗人赢得了犹如无误裁判般端坐的民众的鼓掌喝彩时，是何等的黑暗蒙蔽了心灵，何等的恐惧侵袭，何等的激情使之燃烧！}\par

% structure: p0032 source=h2 target=subsection reason=relative-heading-rank; base=h2

\subsection{是虚荣而非理性造成了某些罗马神祇。}

然而，难道不是显而易见，是虚荣而非理性主宰了对某些假神的选择吗？这位被他们视为半神的\Person{柏拉图}，倾其全部辩才以保护人免于最危险的属灵灾祸，却甚至被认为不配拥有一座小神龛；而\Person{罗慕路斯}，因为他们可以称其为自己人，便受到比许多神祇更高的尊敬，尽管他们的隐秘教义只能赋予他半神的品级。他们为他指派了一位弗拉门祭司，也就是说，一位在其宗教中极受尊崇的祭司（又以圆锥形法冠为标志），以致只为三位神祇设立了弗拉门——为朱庇特设立的是朱庇特弗拉门，为马尔斯设立的是马尔斯弗拉门，为\Person{罗慕路斯}设立的是奎里努斯弗拉门（因为当同胞的热情把\Person{罗慕路斯}置于众神之列时，他们给了他奎里努斯这个新名）。因此，通过这一荣誉，\Person{罗慕路斯}被置于朱庇特的兄弟涅普顿和普路托之上，甚至置于他们的父亲萨图恩本人之上。他们指派同等的祭司去侍奉他如同侍奉朱庇特；而在给予马尔斯（相传为\Person{罗慕路斯}之父）同等荣誉时，这难道不是为了\Person{罗慕路斯}的缘故，而非为尊崇马尔斯吗？\par

% structure: p0034 source=h2 target=subsection reason=relative-heading-rank; base=h2

\subsection{若神祇真有任何对正义的顾念，罗马人本应从他们领受良法，而不是不得不从其他民族借用。}

再者，如果罗马人能够从他们的神祇那里领受生活准则，他们就不会在罗马建城若干年后从雅典人那里借用\Person{梭伦}的法律；而且他们并未原封不动地保留借来的法律，而是努力改进和修订它们。尽管\Person{莱克格斯}假装受阿波罗授权为拉栖第梦人制定法律，明智的罗马人却不愿相信此事，也没有被诱导从\Place{斯巴达}借用法律。据说继\Person{罗慕路斯}为王位的\Person{努马·庞皮利乌斯}制定了一些法律，然而这些法律不足以管理城邦事务。在这些法规中有许多涉及宗教仪式，然而据传他甚至没有从神祇那里领受这些。因此，对于道德恶行，即生活和行为的恶——这些恶如此强大，据最智慧的异教徒称，它们摧毁国家而其城池却完好无损——，他们的神祇并未为保护其崇拜者免受这些恶而作出丝毫规定，相反，却特别着力于增加它们，正如我们先前所努力证明的那样。\par

% structure: p0036 source=h2 target=subsection reason=relative-heading-rank; base=h2

\subsection{论萨宾妇女的劫掠及罗马最辉煌时代所犯的其他罪行。}

但是，我们也许可以从\Person{萨卢斯特}的话中找到罗马诸神如此忽视罗马人的原因，他说：\Quote{在罗马人中间，公正与德行的盛行，与其说是由于法律的力量，不如说是由于天性。}我猜想，这种天生的公正和善良的品性，应当归因于抢劫萨宾妇女的行为。的确，还有什么比这更公正、更有德行呢？——用武力抢走那些陌生而作客的姑娘，她们被表演的借口诱骗而来，没有父母的同意，每个人凭自己的能力带走。如果萨宾人在罗马人求亲时拒绝是不对的，那么罗马人在被拒绝后强行抢走她们，岂不是更大的错误？罗马人更应当因为他们最初求亲时邻族拒绝将女儿嫁给他们而宣战，而不是在偷走她们之后因对方索回而战。起初就应该宣战；那时\Person{马尔斯}本应帮助他好战的儿子，用武力报复因拒绝婚姻而造成的伤害，同时赢得他所渴望的女人。一个胜利者凭借\Quote{\OriginalTerm{战争的权利}{right of war}}，夺走那些毫无权利而被拒绝的处女，或许还有些表面上的理由；但并没有\Quote{\OriginalTerm{和平的权利}{right of peace}}，能使他有权带走那些未给予自己的人，并对义愤填膺的父母发动不义的战争。这个暴力行为倒是有一件幸事：虽然它通过竞技场的比赛而被纪念，但这并没有使它成为\Place{罗马}城或罗马国家的一个先例。如果有人要指责这行为的后果，倒不如说罗马人尽管\Person{罗穆卢斯}犯下这等恶行，却还将他封为神；因为不能责备他们把这行为变成抢劫妇女的任何先例。\par

再者，我想，由于这种天生的公正和德行，在国王\Person{塔奎尼乌斯}被驱逐后——他的儿子曾强暴了\Person{卢克丽霞}——执政官\Person{尤尼乌斯·布鲁图斯}迫使\Person{卢基乌斯·塔奎尼乌斯·科拉提努斯}，也就是卢克丽霞的丈夫、他自己的同僚，一个善良无辜的人，辞去职务并流放，唯一的罪名是他属于塔奎尼乌斯的家族和血统。这一不义之举得到了人民的赞同，或至少默许，而正是这些人民将科拉提努斯和布鲁图斯两人提升到执政官职位。这种公正和德行的另一个例子，可见于他们对待\Person{马库斯·卡米卢斯}的方式。这位杰出的人物，在迅速征服了当时罗马最可畏的敌人维爱人之后——这场战争持续了十年，罗马军队由于指挥不力而遭受了常见的灾难——在恢复了罗马的安全之后，那时\Place{罗马}已开始为自己的安危战栗，在攻占了敌人最富庶的城市之后，却遭到那些嫉妒他成功的人的恶意指控，以及平民保民官的傲慢攻击；他见自己保全了城邦却得不到丝毫感激，且必定会被定罪，于是流亡出走，甚至缺席被判处一万阿斯的罚金。然而不久之后，他忘恩负义的祖国又不得不为了抵抗高卢人而寻求他的保护。然而我现在无法一一列举罗马所遭受的全部可耻而不义的行为，当时贵族试图压服平民，平民抗拒他们的侵犯，而双方的支持者与其说是出于任何公正或德行的考虑，不如说是受制于对胜利的渴望。\par

% structure: p0039 source=h2 target=subsection reason=relative-heading-rank; base=h2

\subsection{第十八章。——\Person{萨卢斯特}的\WorkTitle{历史}所揭示的罗马人的生活，无论是处于焦虑的困境还是安逸的放松中。}

因此，我将暂且停下，引述\Person{撒路斯提乌斯}本人的证词；他对罗马人的赞扬之词（即\Quote{公平和美德在他们当中盛行，借助法律的力量并不比借助天性更多}）引发了这番讨论。他指的是驱逐国王之后的那段时期，在此期间，这座城市在极短的时间里变得伟大。然而，同一位作者在他历史著作的第一卷，也就是开篇之处承认，即使在那时，当政府从国王转到执政官手中只过去了极短的时间，权势更大的人就开始行事不义，导致人民背离贵族，并在城中引发了其他混乱。因为\Person{撒路斯提乌斯}陈述说，罗马人在第二次和第三次布匿战争之间比任何其他时候都享有更大的和谐和更纯洁的社会状态，其原因不是他们对良好秩序的热爱，而是害怕他们与\Place{迦太基}的和平可能被破坏（正如我们提过的，\Person{纳西卡}反对摧毁\Place{迦太基}时也考虑到了这一点，因为他认为恐惧会抑制邪恶，并维护健康的生活方式），然后他接着说：\Quote{然而，在\Place{迦太基}被摧毁之后，不和、贪婪、野心，以及其他通常由繁荣滋生的恶习，比以往任何时候都更加滋长了。}如果它们\Quote{滋长了}，而且是\Quote{比以往任何时候都更加}，那么它们早就出现了，并且一直在滋长。因此，\Person{撒路斯提乌斯}为他所说的话补充了理由。他说，\Quote{因为，权势者的压迫手段，以及随之而来的平民脱离贵族，以及其他内部纷争，从一开始就存在，而事务的公平和适度正义的管理，只持续了驱逐国王后的一段短时间，那时这座城市正忙于严重的\Place{托斯卡纳}战争和\Person{塔尔昆}的复仇。}你看到，即使在驱逐国王后的那段短暂时期，他承认，恐惧是那个公平和良好秩序时期的原因。事实上，他们害怕\Person{塔尔昆}对他们发动的战争，那时他已被逐出王位和城市，并与\Place{托斯卡纳}人结盟。请继续看他说的话：\Quote{此后，贵族把人民当作奴隶对待，命令他们像国王曾经做的那样遭受鞭打或斩首，把他们从自己的土地上赶走，对那些无产可失的人施行残酷的暴政。人民被这些压迫措施压垮，尤其是被过高的高利贷压垮，并且被迫为持续的战争贡献钱财和亲身服役，最终拿起武器，撤离到\Place{阿文廷山}和\Place{圣山}，从而为自己争取到了保民官和保护性法律。但是，只有第二次布匿战争才结束了双方的不和与争斗。}你看到，即使在驱逐国王之后区区几年，罗马人是什么样的人；他说的就是这些人，\Quote{公平和美德在他们当中盛行，借助法律的力量并不比借助天性更多}。\par

现在，如果这些就是罗马共和国展现最美好、最优秀的时日，那么对于随后的时代，我们该说什么或作何感想呢？用同一位历史学家\Person{撒路斯提乌斯}的话来说，\Quote{它从当初那个美好而富有美德的城市，一点一点地变化，最终变得彻底邪恶而放荡}？正如他所提到的，这是在\Place{迦太基}被摧毁之后。\Person{撒路斯提乌斯}对这一时期的简短总结和概述，可以在他自己的历史中读到，其中他展示了由繁荣助长的放纵风气最终如何导致内战。他说：\Quote{从这时起，原始的风尚，不再像迄今为止那样经历不知不觉的渐变，而是如洪流般被一扫而空：年轻人被奢侈和贪婪败坏到如此地步，以至于可以公正地说，没有父亲能有一个儿子，既能守住自己的祖产，又能不染指他人的财产。}\Person{撒路斯提乌斯}又详述了\Person{苏拉}的种种恶行以及共和国的普遍败坏状况；其他作家也有类似的论述，尽管语言不那么引人注目。\par

然而，我想你现在看到，或者至少任何留心的人都能看到，在我们天上的君王来临之前，那城已陷入了何等邪恶的深渊。因为这些事情不仅发生在\Person{基督}开始教导之前，而且发生在他由\Person{童贞玛利亚}出生之前。如果，他们不敢将那些前代的可怕灾难归咎于他们的神祇——这些灾难在\Place{迦太基}被毁前还较可忍受，之后却变得不堪忍受且可怕至极，尽管正是这些神祇用他们恶毒的诡计，将那些恶念灌输到人心中，由此蔓生出四面八方如此可怕的罪恶——那么，他们为何将这些当下的灾祸归咎于\Person{基督}呢？他教导赐予生命的真理，禁止我们敬拜虚假骗人的神祇，并以他神圣的权威憎恶并谴责那些邪恶有害的人欲，逐渐将他的子民从这被罪恶败坏、正走向毁灭的世界中分别出来，为要使他们成为一座永恒的城，其荣耀不建立在虚浮的欢呼之上，而建立在真理的审判之上。\par

% structure: p0043 source=h2 target=subsection reason=relative-heading-rank; base=h2

\subsection{论\Person{基督}废除对诸神的敬拜之前，罗马共和国中滋长的败坏}

这便是那个\OriginalTerm{罗马共和国}{res publica Romana}，\Quote{它从昔日美好而有德的城市，逐渐蜕变为彻底邪恶而放荡。}第一个这么说的并不是我，而是他们自己的作家；我们曾付费向他们学习，而他们早在基督降临之前就写下了这些。你们看，在基督降临之前，在\Place{迦太基}陷落之后，\Quote{古老的风尚，不像以往那样经历着不易察觉的渐变，而是如被洪流般席卷而去；年轻人又是如何因奢侈与贪婪而堕落的。}现在，让他们反过来向我们宣读，他们的神明为罗马人民颁布的、旨在反对奢侈与贪婪的任何法律吧。但愿他们对贞洁与端庄缄口不言，而不曾要求人民进行那些下流可耻的行为，并用他们所谓的神性，给予这些行为有害的庇护。让他们读一读我们在\WorkTitle{先知书}、\WorkTitle{福音}、\WorkTitle{宗徒大事录}或\WorkTitle{书信}中的诫命吧；让他们细察那大量反对贪婪与奢侈的规诫，这些规诫在各处为此而聚集的会众中诵读，它们传入耳中，不是以哲学讨论那种游移不定的声音，而是以天主自己的神谕从云端发出的雷霆之声。然而，他们并不将奢侈与贪婪，以及那些冷酷放荡的习尚——这些甚至在基督降临前就已使共和国成为彻底邪恶而败坏——归咎于他们的神明；相反，无论他们的骄傲与柔弱在这末后的日子里招致了何种苦难，他们都疯狂地归咎于我们的宗教。倘若世上的君王及其全体臣民，倘若世上的一切王侯与审判官，倘若青年男女、老人与少年，各个年龄段，以及男女两性；倘若\Person{洗者若翰}所劝导的税吏与士兵，全都一同听从并遵守基督宗教关于公义与德行的规诫，那么，共和国必将以自己的福乐装饰全地，并在永生中达至君王荣耀的顶峰。但因这人听从，那人却嘲笑，而且大多数人都迷恋恶习的温言媚语，而不喜爱德行的有益严苛，所以，基督的百姓，无论境遇如何——无论是君王、王侯、审判官、士兵、外省人，贫者或富者，为奴的或自主的，男性或女性——都受命忍受这个邪恶放荡的尘世共和国，好能因这坚忍，在那至圣至尊的天使议会和天国的共和国中，为自己赢得显赫的地位；在那里，天主的旨意即是法律。\par

% structure: p0045 source=h2 target=subsection reason=relative-heading-rank; base=h2

\subsection{第20章——论那些攻击基督宗教者所真正喜欢的幸福与生活}

可是，这些神明的崇拜者与仰慕者，却乐于模仿他们那些可耻的恶行，丝毫不在意共和国是否少几分堕落与放纵。他们说：\Quote{只要共和国保持不败，只要它繁荣富庶；让它因胜利而辉煌，或更妙，安享太平；这关我们什么事呢？我们所关心的就是：每个人都能增加财富，以应付日常的挥霍，并使权势者能为了自己的目的压服弱者。让穷人为了生计去巴结富人，好在其庇护下享受懒散苟安；让富人将穷人当作随从，来助长其傲慢。让人民鼓掌喝彩的，不是那些保护他们利益的人，而是那些供他们享乐的人。不要求任何严峻的义务，不禁止任何污秽的行径。让君王们不以臣民的正义，而以臣民的奴性，来衡量自己的繁荣。让各行省忠于君王，不是将其视为道德导师，而是视其为主宰自己财产的主子和享乐的供应者；不是出于由衷的敬畏，而是出于曲意逢迎的奴性恐惧。让法律关注人对他人财产的伤害，甚于对自己人格的伤害。若是某人骚扰邻舍，或侵害其财产、家人或人身，便予以制裁；但自己家内的事务，以及自愿加入者的相处，则任凭各人任意而行，不受惩罚。要大量供应公娼，任何想使用的人皆可享用，尤其对那些无力在私宅中蓄养一个的人。要兴建最大、最富丽的房舍；在其中预备最为奢华的宴席，任何乐意的人，不论白昼黑夜，都可在此赌博、宴饮、呕吐、挥霍。到处都要听到舞者的窸窣声，戏院里的高声淫笑；最残酷与最淫逸的享乐，要接连不断，维持一场永无止息的亢奋。倘若这样的幸福有谁厌恶，便将他烙上公敌的印记；若有谁企图改变或终结它，便让他噤声、遭放逐、被消灭。要将那些能为百姓获致并保守此种状况的神明，算为真神。任凭他们按照自己的心愿受敬拜；任凭他们从或与自己的崇拜者要求任何游戏；只要他们能确保这样的福乐不会被仇敌、瘟疫或任何灾祸危及。}一个理智健全的人，岂会将一个这样的共和国——我不说与罗马帝国相比——而与\Person{萨达那帕鲁斯}的宫殿相提并论呢？那位古代的君王如此沉溺于享乐，甚至让人在他的墓上铭刻，说他死后所拥有的，只是他生前凭嗜欲所吞吃与消耗的东西而已。若这些人有这样一个君王，他虽然放纵自己，却对他们不加严厉约束，那么他们比古罗马人对\Person{罗慕路斯}所作过的，还要更热切地为他献上一座神庙和一位祭司吧。\par

% structure: p0047 source=h2 target=subsection reason=relative-heading-rank; base=h2

\subsection{第21章——\Person{西塞罗}对罗马共和国的看法}

但如果我们的对手不在乎罗马共和国被腐败的行为玷污得多么肮脏与可耻，只要它还能凝聚并继续存在，而且他们因而对\Person{撒路斯提乌斯}关于它\Quote{全然邪恶与堕落}状况的见证不屑一顾，那么他们又将如何看待\Person{西塞罗}的陈述，即即便在他的时代，它已经完全消亡，根本不再有任何罗马共和国存在？他介绍\Person{西庇阿}（即那位摧毁了\Place{迦太基}的西庇阿）在讨论共和国，当时已经预感到它将被\Person{撒路斯提乌斯}所描述的那种腐败迅速毁灭。事实上，在讨论进行时，根据\Person{撒路斯提乌斯}的说法，\Person{格拉古}兄弟之一，即最初煽动叛乱的主要人物，已经被处死。他的死确在同一本书中被提到。现在，\Person{西庇阿}在第二卷的结尾说：\Quote{正如由里拉琴、笛子和人声产生的不同声音之间必须保持某种和谐，受过训练的耳朵无法忍受听到紊乱或不协调的声音，但通过调节即使彼此非常不同的声音，也能产生完全而绝对的协和；同样，当理性被允许调节国家的各种因素时，就能从上层、下层和中层阶级如同各种声音一样获得完美的协和；音乐家在歌唱中所称的和谐，在国家事务中就是协和，这是任何共和国最牢固的纽带和最好的保障，在正义已消亡的地方，无论如何巧思也无法维持。}然后，当他稍微更充分地发挥，并更详尽地阐明了它的存在对国家的益处以及它的缺失对国家的毁灭性影响之后，讨论中的一位在场者\Person{皮卢斯}插话要求更彻底地审视这个问题，并要求自由地讨论正义的主题，以便探究当时日渐流传的那句格言有什么道理，即\Quote{没有不义，共和国便无法治理。}\Person{西庇阿}表示愿意讨论和审视这句格言，并认为它是没有根据的，而且在讨论共和国的问题上不可能取得进展，除非确立不仅这句格言——即\Quote{没有不义，共和国便无法治理}——是错误的，而且确立这样一个真理：没有最绝对的正义，共和国便无法治理。这个问题的讨论被推迟到第二天，在第三卷中以极大的热情进行。因为\Person{皮卢斯}主动为\Quote{没有不义，共和国便无法治理}的立场辩护，同时特别努力澄清自己并不真正赞同这种观点。他极力为不义反对正义的主张辩护，并试图用巧妙的理由和例证证明前者对共和国有利，后者无用。然后，应众人请求，\Person{莱利乌斯}试图为正义辩护，并竭尽全力证明，对于国家而言，没有什么比不义更为有害；没有正义，共和国便既不能治理，甚至无法继续存在。\par

当这个问题被讨论到大家满意之后，\Person{西庇阿}回到原来的讨论主线，并赞许地重述了他自己对共和国的简短定义，即它是人民的利益。\Quote{人民}，他将其定义为不是所有的人群或乱民，而是由于共同承认法律和共同利益而联合起来的人群。然后他指出定义在辩论中的用处；并且从他自己的这些定义中，他得出结论：一个共和国，或\Quote{人民的利益}，只有在它被良好而正义地治理时才存在，不论是由一个君主、一个贵族政体还是由全体人民治理。但当君主是不义的，或者如希腊人所说，是一个暴君；或贵族是不义的，并形成一个派别；或者人民本身是不义的，并且像\Person{西庇阿}因没有更好的名称而称呼他们的那样，自身成为暴君，那么共和国不仅被玷污（正如前一天所证明的），而且根据这些定义进行合法推论，它便完全不复存在。因为当暴君派系性地对国家施加统治时，它就不可能是人民的利益；如果人民是不义的，它也不再是人民，因为它不再符合人民的定义——\Quote{由于共同承认法律和共同利益而联合起来的人群。}\par

因此，当罗马共和国成为\Person{撒路斯提乌斯}所描述的那种样子时，它便不是如他所说的\Quote{极度邪恶放荡}，而是根本不再存在了，如果我们承认那些最优秀的代言人在论共和国时所持的辩论理由的话。\Person{西塞罗}本人，不是以\Person{西庇阿}或任何其他人的口吻说话，而是抒发己见，在第五卷的开头，引用了诗人\Person{恩尼乌斯}的一句诗后，说了以下的话。诗中写道：\Quote{罗马的严正道德及其公民是她的保障。} \Quote{这首诗，}西塞罗说，\Quote{在我看来，这句诗仿佛具有神谕般言简意赅的真理。因为若没有共同体的道德，公民将一无所用；若没有杰出的人物，平民的道德亦不足以建立，更遑论长久地维持如此伟大、疆域辽阔且公义昭著的共和国了。因此，在我们这个时代之前，祖传的习俗塑造了我们最卓越的人物，而他们又反过来保留了父辈的习俗和制度。然而，我们的时代接受了那个已开始衰老的另一时代所留下的共和国\OriginalTerm{杰作}{chef-d'oeuvre}，不仅忽略了恢复其原有的色彩，甚至未曾费心去保留哪怕是大致的轮廓和最突出的特征。因为诗人所称的罗马的保障，那种原始道德还剩下什么呢？它已如此过时并被遗忘，以致于人们不但不实行它，甚至根本不知道它。至于公民，我又能说什么呢？道德因伟人的匮乏而败坏；对于这种匮乏，我们不仅必须找出原因，更必须如同被控死罪的罪犯那样为罪责负责。因为正是由于我们自己的恶行，而非任何不幸，我们才只保留了共和国的空名，而早已丧失了其实质。}\par

这是\Person{西塞罗}的坦言，远在\Person{阿非利加努斯}去世之后（他在其著作\WorkTitle{De Republica}中把阿非利加努斯作为对话者引入），但仍是在基督来临之前。然而，如果他哀叹的这些灾难是在基督教传播并开始得势之后才被悲悼的，那么我们的敌手中还有谁不会认为这些灾难应归咎于基督徒呢？那么，他们的神祇为何当时不采取措施来防止那个共和国的衰败和灭亡呢？对它的灭亡，早在基督降生成人之前很久，西塞罗就唱着如此悲切的挽歌。其仰慕者需要探究，即便在那些古人和淳风的时代，真正的正义是否曾在其中盛行；或许，用西塞罗偶然的话来说，那时它也只是彩色涂绘而非活生生的现实？不过，若天主愿意，我们将在别处考虑这一点。因为我打算在适当的地方指出——根据西塞罗本人，借西庇阿之口简要阐述的共和国与人民的定义，以及他本人和参与那场辩论者提供的诸多见证——罗马从来不是一个共和国，因为真正的正义从未在其中占据过位置。但若采取较为可行的共和国定义，我承认存在过某种共和国的形式，并且古罗马人对它的治理无疑比其现代代表要好得多。但事实是，真正的正义只存在于那位以\Person{基督}为创建者和统治者的共和国中，倘若有人愿意称此为共和国的话；而我们确实不能否认它是人民的福祉。但若这个名字，因在其他语境中已为人所熟知，被认为与我们惯常的说法格格不入，那么我们至少可以说，在这座城中存在着真正的正义；圣经论及这座城说：\BibleQuote{论及你，天主之城，确有光荣的事！}\par

% structure: p0052 source=h2 target=subsection reason=relative-heading-rank; base=h2

\subsection{第二十二章 罗马的神祇从未采取任何措施以防止共和国因道德败坏而毁灭。}

但与当前问题相关的是这一点：无论我们的对手说共和国曾经或现在是多么令人钦佩，但根据他们自己最博学的作家的见证，可以肯定的是，早在基督来临之前，它就已经变得极度邪恶放荡，实际上已不复存在，而是毁灭于伤风败俗。为了阻止这一点，这些守护神无疑应当向那些在众多庙宇中、以如此多样的祭司和祭献、数不清的各式礼仪、众多节庆典礼和无数盛大竞技来崇拜他们的人民，颁布道德训诫和生活准则。但在这所有一切中，魔鬼只关心自己的利益，丝毫不关心崇拜者如何生活，甚至反而费心促使他们过一种放荡的生活，只要他们向它们奉献这些荣誉，并对它们心存畏惧。如果有人否认这一点，就让他拿出、指出、读出神祇为反对暴乱而颁布的法律——正是这些法律\Person{格拉古兄弟}在造成全面混乱时加以违反；或是\Person{马略}、\Person{秦纳}和\Person{卡尔波}在将国家拖入内战时所打破的法律——那些内战起因极其不义无理，过程残酷，结局更为残酷；或是\Person{苏拉}所蔑视的法律，其生平、品行和作为，如\Person{撒路斯提乌斯}和其他史家所述，为全人类所憎恶。谁会否认，在那个时候，共和国已经消亡了？\par

也许他们会大胆地为神祇辩护，说神祇是因公民的放荡而离弃了这城，正如\Person{维吉尔}的诗句所言：\par

从每座庙宇，每处圣所，\newline{}那些使此邦神圣者，都已离去。\par

但首先，若果如此，他们便不能抱怨基督宗教，仿佛它得罪了神祇，使他们抛弃了\Place{罗马}，因为罗马的伤风败俗早已将一群小神，犹如苍蝇般，从城内的祭台上驱走了。然而，在原始道德败坏之前很久，当\Place{罗马}被高卢人攻占并焚毁时，这一群神灵又在哪里呢？也许他们那时在场，但睡着了吗？因为那时，除\Place{卡匹托尔山}外，全城都落入敌人手中；而这座山也本要被攻占，若非——警觉的鹅唤醒了沉睡的众神！这便促成了鹅的节庆，\Place{罗马}在此节庆中几乎堕落到崇拜鸟兽的埃及人的迷信水平。然而，对于这些由敌军或某种灾祸造成的外来祸患，且它们主要触及身体而非灵魂，我暂且不论。目前我所谈的是道德的衰败，它起初几乎不知不觉地褪去了自身的辉煌色彩，随后便完全消失，如被洪流冲走一般，使共和国陷入如此灾难性的毁灭，以至于尽管房屋宫墙依旧矗立，主要作家们仍毫不讳言共和国已经毁灭。那么，神祇们\Quote{离开每一座庙宇，每一处圣殿}，并舍弃此城任其毁灭，这便是一种公义之举了，倘若此城一向轻蔑他们教导公义与道德生活的法律。可是请问，这是怎样的神祇呢？他们不愿与敬拜他们的人民同住，对这些人败坏的生活，他们却从未做任何事加以匡正。\par

% structure: p0057 source=h2 target=subsection reason=relative-heading-rank; base=h2

\subsection{第23章。—— 此生变幻不系于魔鬼之恩仇，而系于真天主之旨意。}

然而，此外，神祇们不是在纵容人欲的满足，而不是加以威严的约束吗？因为\Person{马略}，一个出身微贱、白手起家的人，残酷地挑起并指挥内战，却得到他们如此有力的援助，乃至七次担任执政官，并在第七任执政官职上寿终正寝，逃脱了随后立即掌权的\Person{苏拉}之手。那么，他们为何不也帮助他，以约束他勿犯如此多的暴行呢？因为若说神祇们在他成功中没有插手，这无异于承认一个绝不平凡的事：一个人即使自己的神祇不眷顾他，也能获得此生为人热切渴望的福乐；正如马略虽遭神祇敌视，却满载命运的恩赐，享有健康、权柄、财富、荣誉、高位、长寿；而另一方面，像\Person{雷古卢斯}那样的人，尽管神祇是他的朋友，却饱受囚禁、奴役、贫困、不眠、痛苦与残酷死亡的折磨。承认这一点，便是概括地坦白：神祇无用，其敬拜多余。如果神祇教导人民的，是那些与灵魂之美德、与能获死后赏报的生活正直全然相反的事；如果在暂时的、转瞬即逝的福祉方面，他们既不伤害所恨之人，也不福佑所爱之人，那么为何要敬拜他们，为何要如此热切地呼求崇拜？为何人在艰难可悲的紧急关头，怨叹仿佛神祇怒而退隐？为何基督宗教因他们而遭受最可耻的诽谤？如果在今世事务上他们有行善或作恶的能力，为何他们支持罗马公民中最坏的\Person{马略}，却抛弃了最好的\Person{雷古卢斯}？这岂不证明他们自己最为不义邪恶吗？即便有人设想，正因如此才更该敬畏、敬拜他们，这也是错误的；因为我们并未读到\Person{雷古卢斯}敬拜他们不如\Person{马略}殷勤。也看不出，既然神祇据说更偏爱马略而非雷古卢斯，就该选择邪恶的生活。因为\Person{梅特卢斯}，罗马人中最受尊敬者，有五个儿子担任执政官，他在此世也事业兴旺；而\Person{喀提林}，最邪恶的人，陷入贫困，并在他自己罪恶招来的战争中惨败，生而悲惨，死也悲惨。真实而确实的福乐，惟属敬拜那位惟独能赐予福乐的天主之人。\par

由此可见，当共和国因放荡的风俗而正被毁灭时，其神灵并未通过引导或纠正风俗来阻止其毁灭，反倒加剧了既存的道德沦丧与败坏，从而加速了它的毁灭。他们不必佯装其善性因城邦的邪恶而感到震惊，故而愤怒地撤离了。因为他们确实在那里，这是确凿无疑的；他们被察觉，被证实：他们既不能打破沉默以引导他人，也不能保持沉默以隐藏自己。此外，我不详细论述下面的事实：\Place{明图尔纳}的居民怜悯\Person{马略}，将他托付给丛林中的女神\OriginalTerm{玛里卡}{Marica}，好让她使他在一切事上顺利；他从当时所处的绝望深渊中立即安然返回\Place{罗马}，并以一支无情军队的无情首领的身份进入了这城；若有人想要知道他的胜利何等血腥，何等不像一个公民，而且比任何外敌都更加残酷无情，就让他们去读历史吧。然而，正如我所说，我不详细论述这一点；我既不把\Person{马略}的血腥福乐归于我所不知的明图尔纳女神\OriginalTerm{玛里卡}{Marica}，而是归之于天主的隐秘\OriginalTerm{眷顾}{providence}，为要堵住我们对手的口，并使那些不受激情驱使，而以审慎思考事件的人能脱离谬误。即使邪魔在这些事上拥有某种能力，他们所拥有的也只是全能者的隐秘定断所分配给他们的能力，为要让我们不过于看重尘世的繁荣，因为我们看到它甚至时常被赐予像\Person{马略}这样的恶人；另一方面，也让我们不把它视为恶事，因为我们看到许多善良而敬虔的、敬拜唯一真天主的人，尽管有邪魔作祟，却格外成功；最后，要让我们不以为为了尘世的福乐或灾祸而应当讨好或惧怕这些不洁之灵：因为正如地上的恶人不能为所欲为，这些邪魔也同样不能，他们只能在天主的定断所允许的范围内行事，而祂的判断全可理解，无人能公正地指责。\par

% structure: p0060 source=h2 target=subsection reason=relative-heading-rank; base=h2

\subsection{第24章。——\Person{苏拉}的作为，其中邪魔夸耀他们得到了他的帮助。}

显然可确定的是，\Person{苏拉}——其统治是如此的残酷，以致相比之下，人们竟怀念起他所要报复的前一状态——当他首次向\Place{罗马}进军，要与\Person{马略}交战时，献祭时所得的\OriginalTerm{鸟占}{auspices}如此吉利，据\Person{李维}记述，\OriginalTerm{占卜官}{augur}\Person{波司图米乌斯}甚至表示，若\Person{苏拉}不在众神帮助下达成所图，他甘愿失掉自己的头颅。你们看，众神并未离开\Quote{每个殿宇和圣龛}，因为他们仍在预言这些事的结果，然而却未采取任何步骤来纠正\Person{苏拉}本人。他们的预兆应许他大获成功，但没有一句威吓能制服他的邪恶情欲。后来，他在\Place{亚细亚}进行对\Person{米特里达梯}的战争时，\OriginalTerm{朱庇特}{Jupiter}通过\Person{路奇乌斯·提提乌斯}向他传达了一个信息，大意是他将战胜\Person{米特里达梯}；事情果然如此发生了。再后来，当他正计划返回\Place{罗马}，要以公民的鲜血为他本人及朋友所受的伤害复仇时，\OriginalTerm{朱庇特}{Jupiter}又通过第六军团的一名士兵向他传达了第二个信息，大意是正是他曾预言了战胜\Person{米特里达梯}的胜利，而如今他应许赐予他从敌人手中收复共和国的能力，尽管将伴有巨大的流血。\Person{苏拉}立即询问那士兵有什么形象显现给他；在得到答复后，他认出那形象正是\OriginalTerm{朱庇特}{Jupiter}从前用来向他传达关于战胜\Person{米特里达梯}的确证时所运用的同一形象。那么，众神在这一事上如何能自辩，他们费心去预言这些虚幻的成功，却疏于纠正\Person{苏拉}，并阻止他挑起一场如此可悲而凶残的内战——这内战不仅玷污了共和国，更是将其毁灭了？真相是，正如我常常说的，如\WorkTitle{圣经}告知我们的，也如事实本身充分表明的：人们发现，邪魔只追求他们自己的目的，那就是要让人把他们当作神灵来尊崇和敬拜，并诱导人向他们献上那种将他们与其罪行相连的敬拜，从而使人与他们同处于一种共同的邪恶以及天主的审判之下。\par

此后，当\Person{苏拉}来到\Place{塔伦图姆}，在那里献祭时，他看见祭牲肝脏的头上有一个金冠的形像。于是同一位占卜者\Person{波斯图米乌斯}将此解释为预表一场显著的胜利，并命只应由他一人食用内脏。不久之后，一个名叫\Person{卢基乌斯·庞提乌斯}的人的奴隶喊叫道：\Quote{我是贝娄娜的信使；胜利属于你，苏拉！}接着他又说\Place{卡匹托尔}将被焚毁。他一说出这预言，就离开了军营，但第二天回来时比先前更加激动，喊道：\Quote{卡匹托尔着火了！}而它确实着火了。这一点，一个魔鬼既容易预知，也容易迅速宣布。但请注意，与我们主题相关的是，那些亵渎\OriginalTerm{救主}{Saviour}——祂把信徒的意志从魔鬼的统治下解救出来——的人，愿意生活在怎样的神祇之下。那人在预言的神魂超拔中喊出：\Quote{胜利属于你，苏拉！}而且为了证明他是藉着一种神灵说话，他还预言了一件不久就将发生的事，这件事确实在一个离这个说话的神灵所在之处很远的地方发生了。但他从未喊出：\Quote{苏拉，停止你的恶行吧！}——这些恶行是在\Place{罗马}由那位得胜者犯下的，而牛犊肝脏上的金冠曾被作为他胜利的神圣证据显示给他。如果这样的征兆通常是由正义之神而非邪恶的魔鬼所赐，那么他所求问的内脏肯定更早就会向\Person{苏拉}预示那将临到这座城和他自己头上的残酷灾难。因为那场胜利对于他权势的高升，远不如对于他的野心那样致命；因为藉着它，他变得欲壑难填，又因成功而变得如此傲慢和鲁莽，以至于可以说他给自己造成了道德上的毁灭，甚于给敌人造成身体上的毁灭。但是这些真正可悲可叹的灾难，神祇们并没有预先给他任何提示，无论是通过内脏、鸟占、梦兆还是预言。因为他们害怕他改过，更甚于害怕他失败。是的，他们煞费苦心地使这位自己同胞的光荣征服者被他自身可耻的恶行所征服和俘虏，从而更加顺服地成为魔鬼本身的奴隶。\par

% structure: p0063 source=h2 target=subsection reason=relative-heading-rank; base=h2

\subsection{恶神如何通过赋予其榜样以近乎神性的权威，有力地怂恿人行恶。}

如今，谁不由此明白——除非他宁愿效仿这样的神祇，而不愿藉着\OriginalTerm{天主的恩宠}{divine grace}离开他们的交谊——谁看不出这些恶神是何等急切地力图通过他们的榜样，赋予罪恶以一种近乎神性的权威呢？以下事实岂不证明了这一点：他们曾在\Place{坎帕尼亚}一片广阔的平原上被人看见，彼此演练后来在那里发生于\Place{罗马}军队之间、造成大量流血的战役吗？因为起初有人听见巨大的撞击声，后来许多人报告说他们看见两支军队一连数日在那里交战。当这场战斗停止时，他们发现地面全都留下了人和马的脚印，正如一场大战所会留下的一样。那么，如果这些神祇确实在彼此交战，人类的内战就得到了充分的开脱；然而，顺便请注意，这样好斗的神祇必定是非常邪恶或非常可悲的。然而，如果那只是一场假战，他们这样做又是何用意呢？难道不是要使\Place{罗马}人的内战看来并非恶行，而是对神祇的效仿吗？因为内战已经开始了；而在此之前，一些可悲的战役和可憎的屠杀已经发生。许多人已经被一位士兵的故事所触动：他在剥取被杀敌人的战利品时，从被剥去衣甲的尸体上认出了自己的兄弟，于是对内战发出深深的诅咒后，当场在兄弟的尸体旁自杀了。为了掩饰这类悲剧的惨痛，并在这可怖的战争中煽起愈益高涨的热情，这些被尊奉和崇拜为神祇的邪恶魔鬼，便想出这个计谋，让自己以内战的状态显现出来，好使\Place{罗马}人不因对同胞的愧疚而退缩不战，反使人类的罪行得以借神性的榜样而开脱。这些恶神也以同样狡诈的手段，命令为他们设立并奉献我已提及的戏剧表演。在这些表演中，戏剧的诗歌作品和情节将种种恶行归之于神祇，以致人人都可放心效仿他们，无论是相信神祇确实做过这些事，还是虽不相信，却看出他们极其渴望自己被描绘成做过这些事。而且，为了不让人以为诗人在描绘神祇彼此交战时，是在诽谤他们，把不配的行为归给他们，众神祇自己，为完成这骗局，竟亲自在人们眼前展现他们自己的战斗，不仅借着剧场中的表演，还以亲身在真正的战场上显现来证实诗人的作品。\par

我们不得不提出这些事实，因为那些作者竟毫不顾忌地声称并写道，罗马共和国在公民腐化堕落的风俗中早已败坏，并在我们的主耶稣基督降临之前便已不复存在。如今，他们并不将这种败坏归咎于自己的神明，却把我们基督在世间的恶事归咎于他，而这些恶事并不能毁灭善人，无论是生是死。他们这样做，尽管我们的基督颁布了如此多的诫命，教导美德，抑制恶行；而他们自己的神明对于侍奉他们的共和国，非但未作出任何努力以保存它，用这样的诫命阻止其毁灭，反而以他们有害的榜样败坏其道德，加速了它的灭亡。我想，如今恐怕没有人敢说，共和国当时之所以败亡，是因为神明们\Quote{离开了每一座神庙，每一处圣所}，仿佛他们是美德之友，因人类的恶行而愤怒。不，有太多来自内脏占卜、鸟占、预言的事物，他们据此夸耀自己能预知未来，掌控战争的命运——这一切都证明他们当时在场。倘若他们真的不在场，罗马人在这些内战中，绝不至于因其自身的激情而放纵，反倒因这些神明的煽动而更加疯狂。\par

% structure: p0066 source=h2 target=subsection reason=relative-heading-rank; base=h2

\subsection{魔鬼暗地里传授某些隐晦的道德训导，而公开的庆典却教导一切邪恶。}

鉴于事实如此——鉴于神明们那些污秽残忍的行为，可耻罪恶的事迹，无论是真实的还是虚构的，都是应他们自己的要求而公之于众，并被祝圣，作为神圣且定期的庆典，以显扬他们；鉴于他们誓言要对那些拒绝将这一切展示在众人眼前的人施行报复，以便这些行为被树立为值得效仿的榜样——为什么这些魔鬼，既然他们乐于沉浸在这种淫秽之中，承认自己是不洁的灵体；既然他们喜欢自己那些真实的或想象的恶行与不义，并要求放肆之人、强迫端庄之人庆祝这些放荡行为，从而宣称自己是罪恶和淫荡生活的煽动者——我问，为什么他们却被描绘成，在圣所的秘密中，给少数自己的选民传授一些良好的道德训诫呢？倘若如此，这恰好进一步证明了这些有害邪灵的恶毒诡计。因为正直与贞洁的影响力如此之大，以至于所有人，或几乎所有人，都为对这些美德的赞扬所感动；没有任何人因恶习而堕落至如此地步，以至于心中不残留丝毫荣誉感。因此，正如圣经所言，若非魔鬼有时将自己\BibleQuote{化作光明的天使}\BibleRef{格后 11:14}，他便无法实现其欺骗的意图。于是，在公开场合，厚颜无耻的污秽以喧嚣的叫嚷充满民众的耳朵；在私下里，伪装的贞洁以几乎听不见的耳语对少数人低语：可耻的事有敞开的舞台，而可赞美的事却落下帷幕：优雅掩面，丢脸招摇：邪恶的行为吸引满座观众，美德的言辞却难觅听众，仿佛贞洁应感羞愧，淫乱反当夸耀。这种混乱除了在魔鬼的庙宇中外，还能在哪里盛行呢？除了在欺诈的巢穴中外，还能在哪里呢？因为秘密的训诫是作为对少数有德者的安抚而给予的；邪恶的榜样则公开展示，以鼓动无数的恶人。\par

那些接受\OriginalTerm{切莱斯蒂斯}{Cœlestis}秘仪的人，是在何时何地接受了良好的训导，我们无从知晓。我们所确知的是，在她的圣龛前——龛中安放着她的像——在从四面八方聚集而来、拥挤站立的人群中，我们曾作为全神贯注的观众，观看正在进行的表演；我们随意转目，只见一边是妓女们的奢华展示，另一边则是这位童贞女神；我们看到这位童贞女神既受人祈祷朝拜，也接受淫秽的仪式。在那里，我们看不到面带羞容的哑剧演员，也没有过分拘谨的女伶；淫秽仪式所要求的一切，都得到了充分的满足。我们清楚地看到了什么使这位童贞女神喜悦，而那些目睹景象的已婚妇女，从庙里回家时，已变得更加老练。的确，有些较为审慎的妇女，别过脸去不看表演者下流的动作，却通过偷偷一瞥学会了邪恶的伎俩。因为，她们因对男性应有的端庄而不敢大胆观看那些不雅的手势；但她们更不敢以贞洁的心去谴责她们所朝拜的女神的神圣仪式。然而，这种放荡——若在自己家中施行，也只能在暗处为之——却在庙宇中作为公开的课程来实践；如果人心中还存有任何羞耻感，那么这种羞耻感便会用来惊叹：人无法肆意犯下的恶行，竟成为神明宗教教义的一部分，而省略这种表演竟会招致神明的愤怒。那是什么样的灵体，能通过隐秘的灵感煽动人的败坏，驱使人们犯奸淫，并以成熟的邪恶为食呢？它若不是那个以这种宗教仪式为乐、在庙宇中设立魔鬼的偶像、喜爱在表演中观看邪恶影像的灵体，又能是谁呢？它在暗中低语一些义言，以欺骗少数善人；又在公开场合散布放纵的邀请，以俘获无数恶人。\par

% structure: p0069 source=h2 target=subsection reason=relative-heading-rank; base=h2

\subsection{罗马人为取悦神明而奉献的那些戏剧中的淫秽，极大地促成了公共秩序的倾覆。}

\Person{西塞罗}，一位举足轻重的人物，也算是一位哲学家，在即将出任\OriginalTerm{市政官}{edile}时，希望公民们明白，在他担任此官职的诸多职责中，有一项是必须举办赛会以媚悦\Person{福罗拉}女神。而这些赛会越是淫秽，便越被视为虔诚的媚悦之举。在另一处，当他已是执政官，国家处于极大的危险中时，他说赛会已连续举行了十天，凡是能平息众神的事，一件也没有遗漏 ；仿佛以节制激怒神明，远不如以放荡来安抚他们；或者说，以正直的生活激起他们的仇恨，远不如以如此不堪的淫秽来平息。因为无论那些威胁国家的人是何等残酷凶猛——神祇正是为他们的缘故而受到讨好——，他们所能造成的伤害，也比不上那些以最肮脏的恶行赢得的诸神联盟来得更大。为消除威胁人们肉身的危险，神明竟以这种方式得到安抚，而这方式却将美德从他们的灵魂中驱走；诸神在尚未先攻陷并洗劫公民的道德之前，并不肯充当城垛的守卫者来抵御外敌。对这种神祇的媚悦之举——一种如此放荡、不洁、无耻、邪恶、污秽的媚悦之举，其表演者被罗马人与生俱来且值得赞美的德行所排斥，不得担任公职，被从部落中除名，被视为玷污且声名狼藉——\Emph{我说}，这种媚悦之举如此肮脏、可憎，与一切宗教情感相去甚远；有关诸神恶行的这些虚构而诱人堕落的传说，这些他们或是以可耻邪恶的方式真实犯下，或是以更可耻邪恶的方式捏造出来的丑行，整个城市的人都通过表演者的言语和姿势公开学习。他们眼见诸神乐于此类恶行，便相信诸神不仅希望将这些东西表演给他们看，更希望他们自己加以模仿。至于他们所说的那种良善正直的教诲，却是在如此秘密的情况下，只对极少数人传授（如果当真传授过的话），以至于他们似乎更怕它泄露出去，而不是担心无人践行。\par

% structure: p0071 source=h2 target=subsection reason=relative-heading-rank; base=h2

\subsection{基督宗教是赋予健康的宗教。}

因此，那些人不过是自暴自弃、忘恩负义的可怜虫，深深而牢固地受着那邪灵的束缚，他们抱怨并咕哝着说，人们藉着基督之名，从这些不洁之灵的地狱般的奴役中，并从分担他们惩罚的处境中，被拯救出来，脱离了那充满毒害的不虔敬的黑夜，而进入了最有益健康的虔诚的光明之中。唯有如此的人才会抱怨说，群众涌向教会及其贞洁的崇拜活动，在那里，男女之间有着合宜的隔离；在那里，他们学习如何度过此世生活，以赢取来世的永福；在那里，\WorkTitle{圣经}与正义的教导在高台上向众人宣讲，使那些遵行圣言的人听见而得\OriginalTerm{救恩}{salvation}，那些不遵行的人听见而受\OriginalTerm{审判}{judgment}。即便有些人带着嘲弄这类训诫的态度进来，他们的傲慢要么被突然的转变所熄灭，要么因恐惧或羞耻而受到约束。因为在那里，没有任何污秽邪恶的行为被展示出来供人观看或模仿；所推荐的，是真天主的诫命；所叙述的，是祂的奇迹；所赞美的，是祂的恩赐；所求的，是祂的恩惠。\par

% structure: p0073 source=h2 target=subsection reason=relative-heading-rank; base=h2

\subsection{规劝罗马人弃绝异教。}

可敬的罗马种族——你们的\Person{斯凯沃拉}和\Person{西庇阿}、\Person{雷古鲁斯}和\Person{法布里修斯}等伟人的后裔啊，这才是配得上你们心志的宗教。要切慕这样的宗教，要把它与魔鬼的污秽虚妄和诡诈邪恶分别开来。如果你们本性中有任何卓越的德行，唯有藉着真正的\OriginalTerm{虔诚}{piety}才能使其净化与成全，而不虔敬则使其毁伤并受罚。现在选择你们所要追随的吧，这样你们的荣耀就不在于自己，而在于\OriginalTerm{真天主}{true God}，在祂内没有错谬。因为世俗的荣耀你们已经有份；但出于\OriginalTerm{天主隐秘的眷顾}{secret providence of God}，真正的宗教并未曾提供给你们的祖先选择。醒来吧，现在已是白昼；正如你们已经借着某些人而醒来——我们以这些人的完美德行和为\OriginalTerm{真信仰}{true faith}所受的苦难而自豪：因为他们四面与敌对势力争战，并以英勇的死亡战胜了所有敌人，用自己的血为我们赢得了这属于我们的国度；我们邀请你们来到这国度，并劝勉你们加入这个城市的子民行列，这城市自身也拥有真正赦罪的特有圣所。不要听从你们那些堕落的子孙，他们诽谤\Person{基督}和基督徒，并将这些灾难的时代归咎于他们，尽管他们所渴望的时代是那种他们能逍遥法外、享受罪恶而不受惩罚的时代，远胜于和平的生活。即便对她的地上国度，罗马也从未有过这样的野心。现在，要把握那天上的国度，它易于赢得，在那国度里你们将真实地、永远地为王。因为在那里你们将找不到\OriginalTerm{维斯塔圣火}{vestal fire}，找不到\OriginalTerm{卡皮托利山石}{Capitoline stone}，唯有唯一的真天主。\par

此处既无时限，亦无终极可命；\newline{}惟有无穷无疆之王权赐予。\par

因此，不要再追求伪诈的神明了；务要弃绝他们，轻视他们，冲入真正的自由中。他们并非神明，乃是凶恶的灵，你们的永恒真福将是他们的严厉惩罚。按血统论，你们认作起源的\Person{朱诺}，对特洛伊人夺取\Place{罗马}城堡，还没有像你们仍视为神明的这些魔鬼对人类永远居所的嫉妒那样刻毒。你们自己也曾毫不犹豫地对他们作出了判决：你们用赛会安抚他们，却把演戏的人视为无耻之徒。那么，请容许我们为你们的自由辩护，抵抗那些把庆祝自己的羞耻和丑行当作轭加在你们颈上的不洁之灵。你们已将那些演绎神明罪行的人逐出荣耀的职位；恳求真天主吧，求祂从你们中间除去那些以罪行自娱的神明——若这些罪行真是他们所为，那便极其可耻；若只是假造，便是最恶毒的捏造。你们主动将所有演员和优伶放逐出公民的行列，这很好。更完全地醒悟吧：天主的尊威不能靠玷污人的尊严来平息。那么，你们怎能相信那些喜好如此淫荡戏剧的神明属于天上圣洁的权势呢？因为那些演戏的人，连你们自己也拒绝接纳他们进入\Place{罗马}公民之列，哪怕是最低等级的公民。那座天上的城，远比\Place{罗马}更为荣耀，在那里，你们拥有真理作为胜利，圣洁作为尊严，幸福作为和平，永生作为生命。它更不允许这样的神明加入其团体，如果你们羞于让这样的人加入你们的团体。因此，若你们想要进入那有福的城，就要远避魔鬼的团体。以可耻之事讨好的人，不配高尚心灵的人崇拜。所以，愿这些神明借着基督宗教的洁净而从你们的敬拜中被抹去，正如那些人被监察官的标记从你们的公民权中抹去一样。\par

然而，论到肉身的利益——这是恶人唯一渴望享受的福分，以及肉身的苦难——这是他们唯一畏惧遭受的——我们将在下一卷中指出，魔鬼并没有人们以为的那种权能；而且即使他们拥有，我们倒更应因此轻视这些利益，而不该为了这些利益去崇拜那些神明，并且因崇拜他们而错失他们嫉妒我们得享的这些利益。但是，他们甚至没有那些为了现世利益而崇拜他们的人所归于他们的能力，这一点，我说，我将在下一卷中证明；所以，就让我们在此结束目前的论述。\par

\SourceRecord{Translated by Marcus Dods.；Nicene and Post-Nicene Fathers, First Series；Vol. 2.；Edited by Philip Schaff.；Buffalo, NY: Christian Literature Publishing Co.,；1887.；<http://www.newadvent.org/fathers/120102.htm>.}

% structure: p0001 source=h1 target=section reason=page-title

\section{\WorkTitle{天主之城（第三卷）}}

如前卷中\Person{奥古斯丁}已证明关于道德及灵性方面的灾祸，本卷他则证明关于外在及肉身的灾难：自\Place{罗马}建城以来，罗马人便不断遭受此类灾难；甚至当伪神在\Person{基督}来临前无可匹敌地受崇拜时，他们也没有使罗马人免于这些灾祸。\par

% structure: p0003 source=h2 target=subsection reason=relative-heading-rank; base=h2

\subsection{第一章 论唯独恶人所惧怕、而世界在诸神受崇拜时仍不断遭受的灾祸}

关于道德与灵性的恶，是首当远避的，我想，已说得足够多，足以表明那些伪神并未采取任何措施阻止其崇拜者被这类灾祸所淹没，反而加重了他们的毁灭。看来我现在必须讲述那些唯独异教徒所惧怕的灾祸——饥荒、瘟疫、战争、劫掠、被俘、屠杀，以及第一卷中已列举过的类似灾祸。因为恶人只将那不使人变恶的事视为恶；他们称赞善事却不以为耻，却在所称赞的善事中依然为恶。他们拥有破败的房屋比拥有败坏的生活更令其忧伤，仿佛人的最大好处是自身之外一切皆好。然而，就连异教徒唯独惧怕的这些灾祸，他们的诸神也未加阻挡，即使在这些神被最无限制地崇拜时也是如此。因为在我们救主来临前的各个时代和各个地方，人类遭受了无数、有时是难以令人置信的灾祸；那时，除了希伯来民族以及天主最隐秘、最公义的审判认为堪受恩宠的个别人士外，这世界所崇拜的不正是那些神吗？但为避免冗长，我将对任何其他民族所遭受的重大灾难保持沉默，只讲述\Place{罗马}和\Place{罗马帝国}所发生的事；我所说的\Place{罗马帝国}，是指罗马本城，以及那些在\Person{基督}来临之前，已通过结盟或征服而成为国家机体成员的那些土地。\par

% structure: p0005 source=h2 target=subsection reason=relative-heading-rank; base=h2

\subsection{第二章 论希腊人与罗马人共同崇拜的诸神，是否理当容许\Place{伊利昂}的毁灭}

首先，那么，为什么\Place{特洛伊}或\Place{伊利昂}，罗马民族的摇篮（因为我切不可忽略或掩饰我在第一卷中提及的），尽管它敬拜与希腊人相同的神，却被希腊人征服、攻陷和毁灭呢？有人回答说，\Person{普里阿摩斯}为其父亲\Person{拉俄墨冬}的伪誓付出了代价。那么，的确，\Person{拉俄墨冬}曾雇佣\Person{阿波罗}与\Person{尼普顿}为他做工。因为故事说他应许给他们工钱，随后却毁约。我感到惊讶的是，那位著名的占卜者\Person{阿波罗}竟为如此浩大的工程劳苦，却从不怀疑\Person{拉俄墨冬}会欺骗他，不付工钱。还有\Person{尼普顿}，他的叔叔，\Person{朱庇特}的兄弟，海中之王，竟然不知将要发生的事，这实在不相宜。因为\Person{荷马}（他生活在\Place{罗马}建立之前并著书）描绘他预言了\Person{埃涅阿斯}后裔的伟大事业，而\Person{埃涅阿斯}事实上建立了\Place{罗马}。正如\Person{荷马}所说，\Person{尼普顿}也曾在云雾中从\Person{阿喀琉斯}的愤怒中救出\Person{埃涅阿斯}，尽管（按照\Person{维吉尔}）\par

\Quote{他全心要摧毁\newline{}他自己所造、背誓的\Place{特洛伊}}\par

那么，像\Person{阿波罗}和\Person{尼普顿}这样伟大的神，竟对那将使他们的工钱被诈取之事毫不知晓，白白为\Place{特洛伊}建造城墙，只为换得忘恩负义之人的感谢。人们或许有些怀疑，相信这些家伙是神，是否比欺骗这样的神更为罪孽。就连\Person{荷马}本人也未全然相信这故事，因为他一方面将\Person{尼普顿}描绘成特洛伊人的敌人，另一方面又介绍\Person{阿波罗}为他们的保卫者，尽管故事暗示二者都被那欺诈行为所冒犯。因此，如果他们相信他们的神话，就让他们羞于崇拜这类神；如果他们不信这些神话，就不要再提什么\Quote{特洛伊的伪誓；}或者，让他们解释，诸神如何憎恶特洛伊人的伪誓，却喜爱罗马人的伪誓。因为甚至在如此庞大而腐败的\Place{罗马}城中，\Person{喀提林}的阴谋如何能提供如此众多的人，其手与舌靠伪誓和内乱为生？除了伪誓，还有什么能败坏如此多元老院议员所宣布的判决？还有什么能败坏人民的投票和他们在一切案件中的裁决？因为看来，古代的起誓惯例即使在极度腐败中仍被保留，不是为了通过宗教敬畏约束邪恶，而是为了藉增添伪誓之罪来完成罪恶的总额。\par

% structure: p0009 source=h2 target=subsection reason=relative-heading-rank; base=h2

\subsection{第三章 论诸神不会因\Person{帕里斯}的奸淫而被冒犯，因为这在神自己中间如此普遍}

因此，没有理由声称神明（他们说那帝国靠他们而立，尽管已证明他们被希腊人打败过）因特洛伊人的背誓而愤怒。同样，另一些人为神明辩护说，他们不再保护特洛伊并非出于\Person{帕里斯}通奸的愤慨。因为神明的习性是诱使和教导人作恶，而非报复恶行。\Person{撒路斯提乌斯}说：\Quote{据我所闻，\Place{罗马}城最早是由特洛伊人建造和居住的，他们在\Person{埃涅阿斯}的带领下逃离故土，四处漂泊，没有固定居所。} 因此，如果神明认为应当惩罚\Person{帕里斯}的通奸，那么首当其冲的应是罗马人，或者至少罗马人也该受罚；因为那通奸正是由\Person{埃涅阿斯}的母亲促成的。神明怎会憎恶\Person{帕里斯}的罪行，却对自己姐妹\Person{维纳斯}同样的行为视而不见呢？（姑且不提其他）\Person{维纳斯}与\Person{安喀塞斯}通奸，因而成了\Person{埃涅阿斯}的母亲。难道是因为前者\Person{墨涅拉俄斯}受到了伤害，而后者\Person{伏尔甘}默许了罪行？我看，神明们对自己的妻子极少吃醋，甚至毫不介意让她们与凡人分享。但也许有人会怀疑我在轻蔑地嘲弄神话，而没有以应有的严肃态度处理如此重大的问题。那么，让我们假定\Person{埃涅阿斯}不是\Person{维纳斯}的儿子。我愿承认；但\Person{罗慕路斯}更不可能是\Person{马尔斯}的儿子了吧？为什么一个可以而另一个就不行？或者说，神与女子交合是合法的，人与女神交合就是非法的？这是一个苛刻，甚至难以置信的条件：\Person{马尔斯}依\Person{维纳斯}的律法所获许的，\Person{维纳斯}自己却不得自己律法的许可。然而，这两个例子都有罗马的权威支持；因为近代的\Person{凯撒}深信自己是\Person{维纳斯}的后裔，不亚于古代\Person{罗慕路斯}相信自己是\Person{马尔斯}的儿子。\par

% structure: p0011 source=h2 target=subsection reason=relative-heading-rank; base=h2

\subsection{第四章. 瓦罗认为，人假装自己是神明的后代是有益的。}

有人会说，但这一切你相信吗？我当然不信。因为就连\Person{瓦罗}那样博学的\OriginalTerm{外教人}{heathen}，也几乎承认这些故事是虚假的，尽管他并没有大胆而肯定地这样说。但他主张，让勇敢的人相信自己是神明的后代——即便是虚假的——对国家是有益的；因为这样一来，人心珍视这种神圣出身的信念，就会更勇敢地投身于伟大的事业，并更坚毅地去实现它们，从而凭借这信念本身获得更大的成功。你看到，\Person{瓦罗}的这个意见为虚谎敞开了多么广阔的园地，我尽可能用自己的话表达了它；而且不难理解，在一个认为对公民有益、甚至在关于神明的事上也可以说谎的社群中，许多宗教和神圣传说都应是捏造的。\par

% structure: p0013 source=h2 target=subsection reason=relative-heading-rank; base=h2

\subsection{第五章. 神明既未对\Person{罗慕路斯}之母的通奸表示愤慨，却惩罚\Person{帕里斯}的通奸，这是难以置信的。}

至于\Person{维纳斯}是否为凡人父亲\Person{安喀塞斯}生下\Person{埃涅阿斯}，或\Person{马尔斯}是否使\Person{努米托尔}的女儿怀上\Person{罗慕路斯}，我们且悬而不论。因为我们自己的\WorkTitle{圣经}提示了一个极为相似的问题：\OriginalTerm{堕落的天使}{fallen angels}是否与人的女子交合，使大地当时充满了巨人，就是非常高大强壮的人。因此，目前我将讨论限定于这一两难境地：如果他们书中所载关于\Person{埃涅阿斯}之母和\Person{罗慕路斯}之父的事属实，神明怎能因人犯下他们自己犯下却毫无不悦的通奸而恼怒人呢？若属虚假，则即使在这种情况下，神明也不会因人真的犯下通奸而愤怒，因为他们甚至喜爱人们将这些事虚假地归之于他们。此外，如果否认\Person{马尔斯}的通奸，好让\Person{维纳斯}也摆脱指责，那么\Person{罗慕路斯}的母亲就失去了神明诱惑的借口而被置于不理。因为\Person{西尔维娅}是\OriginalTerm{维斯塔贞女}{vestal priestess}，神明理应为这种亵渎行为对罗马人施以比惩罚\Person{帕里斯}的通奸更严厉的报复。因为就连罗马人自己在远古时代，一旦发现维斯塔贞女有通奸行为，便会将她活埋，而普通妇女即使受罚，也从不因此罪行而被处死；这样，他们更热切地维护他们视为神圣的庙宇的纯洁，胜过维护婚床的纯洁。\par

% structure: p0015 source=h2 target=subsection reason=relative-heading-rank; base=h2

\subsection{第六章. 神明对\Person{罗慕路斯}杀害兄弟的罪行未加惩罚}

我另举一例：如果人的罪恶如此激怒那些神明，以致他们抛弃\Place{特洛伊}，任其遭受刀兵之灾，以惩罚\Person{帕里斯}的罪行，那么\Person{罗慕路斯}杀兄弟的行为，应当比希腊丈夫的诱骗更激使神明们反对罗马人：一个初生的城邦中的杀兄弟之罪，理当比一个已繁荣的城邦中的奸淫之罪更激起他们的愤怒。对我们所讨论的问题而言，无论\Person{罗慕路斯}是下令杀害他的兄弟，还是亲手杀了他，都没有区别；这桩罪行，许多人厚颜无耻地否认，许多人出于羞愧怀疑它，许多人因悲伤而掩饰它。我们不必停下来审查衡量史家在这一主题上的证言。所有人都一致同意，\Person{罗慕路斯}的兄弟是被杀死的，不是被敌人，不是被外人。如果正是\Person{罗慕路斯}命令或实施了这一罪行，\Person{罗慕路斯}比\Person{帕里斯}更堪称是罗马人的首领，而非特洛伊人的首领；那么为何抢夺他人之妻的那人招来神明对特洛伊人的愤怒，而夺取兄弟性命的那人却获得了同一些神明的护佑？如果，另一方面，该罪行并非由\Person{罗慕路斯}的手或意志所为，那么整个城邦都应为此负责，因为它没有予以惩治，从而犯下了不是杀兄弟，而是弑亲之罪，这更恶劣。因为兄弟二人都是该城的建立者，其中一人被奸邪所阻，未能成为统治者。就我所见，那么，无法将任何邪恶归咎于\Place{特洛伊}，以证明神明抛弃它、任其毁灭是正当的；也无法将任何善事归于\Place{罗马}，以解释神明为何赐予它繁荣；除非真相是：他们因战败而逃离\Place{特洛伊}，转投\Place{罗马}，在那里施行他们特有的欺骗。尽管如此，他们在\Place{特洛伊}仍保有一席之地，以便欺骗那些重新定居在这些土地上的未来居民；而在\Place{罗马}，通过更广泛地施展他们邪恶的伎俩，他们因更丰厚的荣耀而得意洋洋。\par

% structure: p0017 source=h2 target=subsection reason=relative-heading-rank; base=h2

\subsection{\Place{伊利昂}被\Person{马略}的副将\Person{芬布里亚}所毁}

而且我们肯定会问，可怜的\Place{伊利昂}做了什么错事，以至于在\Place{罗马}内战的初起之时，它竟遭受\Person{芬布里亚}——\Person{马略}党人中最十足的恶棍——之手，遭遇一场比希腊人的洗劫更为猛烈和残酷的毁灭。因为当希腊人攻下它时，许多人逃走了，许多未逃走的人也被允许活着，尽管沦为俘虏。但\Person{芬布里亚}从一开始就下令不留一个活口，并将该城连同其全部居民一起付之一炬。\Place{伊利昂}就是这样获得了报偿，不是来自那些因邪恶行径而被触怒的希腊人，而是来自那些建立于其废墟之上的罗马人；而双方同样崇拜的神明，却完全无所作为，或者更准确地说，是无能为力。那么，难道这也真实吗，即在这个时刻，在\Place{特洛伊}修复了希腊战火造成的损毁之后，所有曾支撑其王位的神明，\Quote{都抛弃了每一座圣殿，每一处神龛？}\par

但若如此，我问其缘由；因为在我看来，神明的行为与市民的行为一样应受谴责，而市民的行为却值得称赞。因为后者紧闭城门抗拒\Person{芬布里亚}，为\Person{苏拉}保全该城，因此被激怒的将领焚烧殆尽。到此为止，\Person{苏拉}的事业是两者中更可敬的；因为至此为止他用兵是为了恢复共和，而他的良好意图尚未遭遇任何挫败。那么，特洛伊人还能做什么更好的事吗？还能做什么更光荣的、对\Place{罗马}更忠诚的，或更配得上其亲缘关系的事，胜过为罗马人中更好的一方保全他们的城市，并对一个弑亲叛国者关上城门？让那些神明的辩护者来考虑一下，这种行为给\Place{特洛伊}带来了怎样的毁灭。神明抛弃了犯淫乱的人民，将\Place{特洛伊}交给希腊人的战火，以便从它的灰烬中兴起一个更贞洁的\Place{罗马}。但他们为何第二次又抛弃这同一座城，它现在已与\Place{罗马}联盟，并非向它高贵的女儿开战，而是对\Place{罗马}最正义的派系保持着最坚定、最虔诚的忠诚？为何他们将它交付毁灭，不是希腊英雄之手，而是最卑劣的罗马人之手？或者，如果神明并不偏爱\Person{苏拉}的事业——不幸的特洛伊人正是为这一事业而坚守此城——那么他们为何自己又预言并许诺\Person{苏拉}如此的成功呢？我们该称他们为幸运者的奉承者，而非不幸者的救助者吗？那么，\Place{特洛伊}被毁灭，并非因为神明抛弃了它。因为那些总是警醒着行骗的\OriginalTerm{魔鬼}{demons}，已尽了他们之所能。因为，据\Person{李维}告诉我们，当所有雕像被推倒并与城一同焚烧时，据说只有\Person{弥涅耳瓦}的像被发现完好无损地矗立在庙宇废墟之中；这并非为使人在赞美中说，\Quote{诸神使此国神圣，} 而是为使人在辩护中不能说，他们 \Quote{已离开每一座圣殿，每一处神龛：} 因为这一奇迹被允许给他们，不是为证明他们强大，而是为证实他们确实在场。\par

% structure: p0020 source=h2 target=subsection reason=relative-heading-rank; base=h2

\subsection{是否应将\Place{罗马}托付给特洛伊的众神}

那么，将\Place{罗马}托付给已在\Place{特洛伊}的失陷中显明其软弱的特洛伊众神，这又是什么智慧呢？难道有人会说，当\Person{芬布里亚}攻打\Place{特洛伊}时，众神已经定居在\Place{罗马}了吗？那么，\Person{弥涅耳瓦}的像何以仍能矗立呢？此外，如果\Person{芬布里亚}毁灭\Place{特洛伊}时他们在\Place{罗马}，或许当\Place{罗马}本身被高卢人攻占并放火焚烧时，他们就在\Place{特洛伊}呢。但由于他们听觉极敏锐，行动极迅速，听闻鹅的嘎嘎叫声便迅速赶来，至少保卫了\Place{卡皮托尔}，虽然要保卫城中的其余部分，他们却迟迟未得到警告。\par

% structure: p0022 source=h2 target=subsection reason=relative-heading-rank; base=h2

\subsection{可否相信\Person{努马}统治时期的和平是由众神带来的}

人们也相信，\Person{罗慕路斯}的继任者\Person{努玛·庞皮利乌斯}在其整个统治期间享有和平，并关闭了\Person{雅努斯}神庙的大门——这些门通常在战时保持敞开——这乃是凭众神的助佑。且他被认为因在罗马人中设立了许多宗教仪式而得到了这样的酬报。诚然，那位国王若能以其闲暇从事有益之事，克服有害的好奇，以真实的虔诚寻求真天主，那么这如此罕见的闲暇本当令我们致贺。然而事实是，众神并非他闲暇的创造者；但或许他们若见他更为忙碌，便不会如此欺骗他。因为他们发现他愈是清闲，便愈是占据他的心神。\Person{瓦罗}向我们告知了他的一切努力，以及他为使这些神祇与他本人及城市联结所施用的计策；若天主愿意，我将在适当之处讨论这些问题。同时，既然我们正谈论众神所赐的恩惠，我愿承认和平是一项大恩；但这是真天主的恩惠，正如太阳、雨水以及其他生命之所需，常常也赐予不知感恩和邪恶之人。然而，倘若这巨大的恩赐是由他们的众神赐予罗马和庞皮利乌斯的，那么为何后来在更为堪当的时期，他们却从未将和平赐予罗马帝国呢？难道神圣仪式在首次设立时，比后来持续举行时更为有效吗？可是在努玛的时代，在他将这些仪式增添入礼仪之前，它们并不存在；而后来，它们已得以举行并保存下来，原可由此产生恩惠。那么，为何努玛统治的那四十三年——或如其他一些人更偏爱的，三十九年——在无间断的和平中度过，而此后，当崇拜已确立，那些被祈求的众神本身被认作城市的守护者和主保时，我们却难以找到整个时期中——从建城至\Person{奥古斯都}统治——有那么一年——即第一次布匿战争结束后的那一年——在那一年，罗马人竟能关闭战门，这成了一个奇迹？\par

% structure: p0024 source=h2 target=subsection reason=relative-heading-rank; base=h2

\subsection{第十章 罗马帝国是否应当通过如此狂烈的连绵战事扩张，它本可效法\Person{努玛}的和平道路而享有安宁与安全。}

他们是否回答说，若非持续不断且从不间断的战争，罗马帝国绝无可能如此辽阔，如此荣耀？这真是个恰当的论据！为何一个王国必须遭受纷扰，才能成为伟大？在人类身体这个小小世界里，拥有适中的身材并享有健康，岂不胜过通过非自然的折磨而达到巨人的庞大尺寸，且达到之后没有安宁，反而因肢体庞大而更加痛苦？倘若\Person{撒路斯提乌斯}所描绘的那些时代得以持续，会有什么恶果产生？岂非一切善果都会产生？他写道：\Quote{起初，国王们（因为这是世上帝国的最初称号）意见分歧：一部分修养心灵，另一部分注重身体；那时人们的生活没有贪欲，每个人都满足于自己的所有！}那么，为了罗马的繁荣，难道就必须让\Person{维吉尔}所谴责的那种事态继之而来：\par

末了，一个更卑劣的时代悄然来临，\newline{}伴随着战争不可驯服的狂暴，\newline{}以及贪婪的利欲之心？\par

然而，罗马人显然有一个看似有理的辩护，来承担并从事如此灾难性的战争——那就是，敌人的压迫迫使他们抵抗，因此他们不得不战斗，并非出于任何对人间赞美的贪婪，而是出于保护生命与自由的必需。好吧，姑且不论。以下便是\Person{撒路斯提乌斯}对此事的叙述：\Quote{因为当他们的国家，因法律、制度、领土而变得丰足，看似极尽繁荣且足够强大时，按照人性常例，富足产生了嫉妒。于是，邻近的国王和邦国拿起武器攻击他们。少数盟友提供了援助；其余的则因恐惧而远离危险。但罗马人，无论是在家还是战中，都警惕、积极，做好准备，互相鼓励，迎击敌人——以武力保卫他们的自由、国家、父母。后来，当他们以勇敢击退了危险，便向盟友和朋友提供援助，并通过施恩而非受惠来争得同盟。}这正是以体面的方式建树罗马的伟大。但是，在\Person{努玛}统治时期，我想知道，那长久的和平是在抵御邪恶邻人的入侵中维持的，还是因为这些入侵停止才得以维持和平？因为如果那时罗马也备受战争侵扰，却并未以武力回应武力，那么她当时用来平息敌人——既不靠战争征服他们，也不靠冲锋陷阵威吓他们——的方法，她本可始终使用，并和平统治，紧闭\Person{雅努斯}的大门。而若她无力做到这点，那么罗马享有的和平就不是出于她众神的意思，而是出于她周边邻人的意思，且仅仅持续到他们不愿以战争挑衅她的时候；除非这些可悲的神祇竟敢将不在他们权能之内，而在他者意志之中的事物，当作他们的恩惠卖给某个人。的确，这些魔鬼，在其被容许的范围内，能以其特有的邪恶恐吓或煽动恶人的心灵。但是，如果它们总是拥有这种能力，而且如果没有任何更隐秘、更高的力量对抗它们的努力，那么它们便将至高无上地主宰和平或战争的胜利——这些几乎总是由于某种人的情绪而发生，并且常常违背众神的意愿，正如不仅一些几乎不透露或暗示任何真实成分的虚假传说，甚至罗马历史本身所证明的那样。\par

% structure: p0028 source=h2 target=subsection reason=relative-heading-rank; base=h2

\subsection{第十一章 \Place{库迈}的\Person{阿波罗}雕像，其眼泪据信预示了\Person{希腊人}的灾难，这位神明无力援救。}

这正是神明软弱无力的又一个明证，见于有关\Place{库迈}的\OriginalTerm{阿波罗}{Apollo}的故事；据说在与\OriginalTerm{亚该亚人}{Achæans}及\Person{阿里斯托尼库斯}王的战争中，他的神像哭泣了四天。当占卜官们因这一异兆而惊恐不安，决定将神像投入海中时，库迈的老人们出面干预，并说在对安条克和佩尔修斯的战争中，同一神像也曾出现类似的异象；而元老院颁布法令，向阿波罗献上祭品，因为后来的事件证明对\OriginalTerm{罗马人}{Romans}有利。于是召来了据说技艺更为高超的占卜师，他们宣告阿波罗神像的哭泣对罗马人乃是吉兆，因为库迈是一个希腊殖民地，而阿波罗是在悲叹（并借此预示）即将降临到他故土\Place{希腊}——他正是从那里被迎来——的忧伤和灾祸。不久后即有消息传来，国王阿里斯托尼库斯战败被俘——这一失败显然违背了阿波罗的意旨；而阿波罗正是通过其大理石像流泪表明了这一点。这向我们表明，尽管诗人们的诗句是神话性质的，却也并非全无真理，而是以颇为贴切的笔法描绘了魔鬼的行事方式。因为在\Person{维吉尔}笔下，\OriginalTerm{狄安娜}{Diana}悲悼\Person{卡米拉}，而\OriginalTerm{赫拉克勒斯}{Hercules}则为注定死去的\Person{帕拉斯}哭泣。这或许就是为什么，\Person{努马·庞皮利乌斯}在享受长期和平时，既不知道也不探究这和平从何而来，闲暇时便开始思量应将罗马的安全与管理托付给哪些神明；他从未梦想过那位真实、全能、至高的天主眷顾着尘世之事，而只记得\Person{埃涅阿斯}带到\Place{意大利}来的特洛伊诸神，既未能保全特洛伊王国，也未能保全埃涅阿斯本人建立的\Place{拉维尼乌姆}王国，于是便认定他必须另寻一些神明，作为逃亡者的保护者和弱者的救助者，并将它们添加进那些较早的神祇之列——这些神祇或是随\Person{罗慕路斯}一同来到\Place{罗马}，或是当初在\Place{阿尔巴}毁灭时被接纳的。\par

% structure: p0030 source=h2 target=subsection reason=relative-heading-rank; base=h2

\subsection{罗马人给\Person{努马}引入的神明增添了大量的神，而他们的数量毫无帮助。}

然而，尽管庞皮利乌斯引入了如此丰富的礼仪，\Place{罗马}却并未感到满足。因为当时\OriginalTerm{朱庇特}{Jupiter}自己还没有主殿——是国王\Person{塔奎尼乌斯}建造了\Place{卡皮托利山}神殿。而\OriginalTerm{埃斯库拉庇乌斯}{Æsculapius}则离开\Place{埃皮达鲁斯}来到罗马，以便在这座首屈一指的城市里，有一片更广阔的天地施展他那高超的医术。还有\OriginalTerm{众神之母}{Mother of the gods}，也不知是从\Place{佩西努斯}的何处而来；因为她的儿子在卡皮托利山巅主事，她自己却隐匿于无名之中，这未免不太体面。但倘若她是所有神明的母亲，那她就不仅跟随她的某些子女来到了罗马，也让另一些子女追随她而来。我确实感到好奇，她可会是\OriginalTerm{狗头神}{Cynocephalus}的母亲——那位是在很久以后从\Place{埃及}而来的。至于\OriginalTerm{热病女神}{Fever}是否也是她的后代，这该由她的孙子埃斯库拉庇乌斯去裁断。但不论她出身如何，我相信那些外来的神明总不至于胆敢称一位身为罗马公民的女神为出身卑贱。谁能数得清那些受托守护罗马的神明呢？本地的与外来的，天堂的、大地的、地狱的、海洋的、泉水的、河流的；正如\Person{瓦罗}所说，还有确定的神与不确定的神，男神与女神：因为就像在动物之中一样，在各类神明中也存在着这些区分。那么，罗马既然享有如此众多神明的庇护，按理说必能免于某些重大而可怖的灾祸，我在这里只能略举数端。因为她借着祭台升腾的浓烟，如同点燃烽火一般，召引了一大群神明来保护自己，为他们设立并维持殿宇、祭台、祭献和祭司，由此便得罪了那位真实而至高的天主，原本这一切礼仪只应合法地献给他。实际上，当她神明不多时反倒更加昌盛；但随着她日渐强大，便以为应当拥有更多的神明，正如大船需要更多的船员来操控一般。我猜想，她是对那较少的神明，——她曾在其庇护下度过相对幸福的岁月——能否捍卫她的伟大感到绝望。因为即便在诸王治下（我已谈过的努马除外），若非存在着何等邪恶的纷争，又怎会导致罗慕路斯兄弟之死呢！\par

% structure: p0032 source=h2 target=subsection reason=relative-heading-rank; base=h2

\subsection{罗马人以什么权利或协议获得了他们的第一批妻子。}

为何无论是\OriginalTerm{朱诺}{Juno}，她与她的丈夫\OriginalTerm{朱庇特}{Jupiter}那时尚且眷顾着\par

\Quote{罗马的子孙，穿袍的族类}\par

还是\OriginalTerm{维纳斯}{Venus}自己，都不肯帮助她所眷爱的\Person{埃涅阿斯}的子孙，以某种正当而公道的方式寻得妻子呢？由于缺乏这一途径，\OriginalTerm{罗马人}{Romans}便陷入了不得不劫夺妻子，随后又与众岳丈刀兵相向的可悲境地；结果那些不幸的妇女，尚未从丈夫强加给她们的冤屈中恢复过来，便已收到了以她们父亲鲜血为代价的嫁妆。\Quote{但罗马人战胜了他们的邻邦。}不错，但双方又遭受了何等的创伤，亲族与邻人之间又发生了何等惨烈的杀戮！\Person{凯撒}与\Person{庞培}之间的战争，不过是区区一位岳父与一位女婿的争斗；而在开战之前，凯撒的女儿、庞培的妻子，即已香消玉殒。然而，\Person{卢卡努斯}又怀着何等深切而公义的哀伤呼喊道：\Quote{我歌咏那场比内战更惨烈的战事，它在\Place{埃马提亚}的平原上进行，罪行竟因胜利而得到正当化！}\par

罗马人当时是这样征战的：他们双手染着岳父们的鲜血，硬是从岳父们的怀抱中夺走那些可怜的姑娘——这些姑娘被夺走后，不敢为被杀的父母哭泣，怕触怒胜利的丈夫；甚至在战事正酣时，她们虽口中祈祷，却不知该为谁祈求。这样的婚姻，当然不是\Person{维纳斯}为罗马人准备的，而是\Person{柏洛娜}；或许当\Person{朱诺}在帮助他们时，那位地狱的复仇女神\Person{阿勒克托}就比当初\Person{朱诺}的祈祷煽动她攻击\Person{埃涅阿斯}时，更有机会伤害他们了。被俘的\Person{安德罗玛刻}比这些罗马新娘幸福。她虽身为奴隶，但嫁给\Person{皮洛斯}以后，不再有特洛伊人死在他手里；而罗马人却在战斗中亲手杀死他们所爱抚的新娘的亲生父亲。\Person{安德罗玛刻}是胜利者的俘虏，只能哀悼自己人民的死亡，却不必恐惧。萨宾妇女们，她们的亲人还在对抗中，当丈夫出战时她们为父亲的性命担忧，丈夫归来时又为父亲的死而哀哭，而她们的悲伤与恐惧，都无法自由表露。因为丈夫的胜利意味着同乡、亲戚、兄弟、父亲的毁灭，这给她们带来的是虔诚的痛苦，或是残忍的狂喜。此外，由于战争胜负无常，她们之中有的丈夫死在父亲刀下，有的则在自相残杀中同时失去丈夫和父亲。罗马人绝非毫无损伤地脱身；他们被逼退回城墙之内，紧闭城门防守；后来城门被诈开，敌人涌入城内，\Place{罗马广场}本身竟成了岳父与女婿们仇恨而凶猛的战场。那些劫持者确实彻底溃败，四散逃回家中，给他们原本可耻又可叹的胜利再添新的耻辱。就在此时，\Person{罗慕洛}对公民们的英勇不再抱希望，便祈求\Person{朱庇特}让他们稳住阵脚；由此这位神得到了\OriginalTerm{阻挡者}{Stator}的称号。然而即使如此，祸患仍可能不会终结，若非那些被劫的妇女们披头散发冲出来，扑倒在父母面前，不是用胜利的武器，而是用儿女的恳求，平息了他们正义的愤怒。于是，原本连自己的亲兄弟都不能容忍作为同僚的\Person{罗慕洛}，却被迫接纳\Place{萨宾}人的王\Person{提图斯·塔提乌斯}与他共治。然而，连自己的孪生兄弟的同伴关系都嫌恶的人，又能忍受一个外人多久呢？因此，\Person{提图斯·塔提乌斯}被杀后，\Person{罗慕洛}独揽王权，以便成为更大的神。请看，这就是激发了不自然战争的婚姻权利。这就是罗马人所谓的血缘、亲属、联姻、宗教的联盟。这就是受众神如此丰厚保护的那座城的生活。你也可看到，就这个话题还能说出许多严厉的话；但我们的旨趣将我们引开，需要我们谈论别的事情。\par

% structure: p0037 source=h2 target=subsection reason=relative-heading-rank; base=h2

\subsection{十四、罗马人对\Place{阿尔巴}人发动的邪恶战争，以及权力欲所赢得的胜利}

但在\Person{努马}统治之后，在其他国王时期，当\Place{阿尔巴}人被挑动战争时，产生了不仅对他们自己，也对罗马人同样悲惨的后果，又发生了什么呢？\Person{努马}的长期和平变得令人厌倦；罗马与\Place{阿尔巴}的军队以怎样的无尽屠杀和双方国家的损毁来终结了它！因为\Place{阿尔巴}是由\Person{埃涅阿斯}之子\Person{阿斯卡尼俄斯}所建，且比\Place{特洛伊}本身更确切地是罗马的母亲，却被罗马王\Person{图卢斯·霍斯提利乌斯}挑动应战，在冲突中双方既给予也遭受如此损害，以致最终双方都厌倦了争斗。于是他们想出一个办法，让战争由双方各出三胞胎兄弟的决斗来定胜负：罗马一方由\Person{霍拉提乌斯}三兄弟出战，\Place{阿尔巴}一方由\Person{库里阿提乌斯}三兄弟出战。\Person{霍拉提乌斯}兄弟中有两人被\Person{库里阿提乌斯}兄弟击败并杀死；但仅存的一名\Person{霍拉提乌斯}却杀死了\Person{库里阿提乌斯}三兄弟。这样，罗马最终获胜，但牺牲如此之大，唯有唯一生还者返抵家园。双方的损失归谁呢？\Person{埃涅阿斯}的后裔，\Person{阿斯卡尼俄斯}的子孙，\Person{维纳斯}的后代，\Person{朱庇特}的孙辈，他们的悲伤又归谁呢？因为这场战争，也是\Quote{比内战更恶劣}的战争，交战国竟是母女。除了这三胞胎兄弟的决斗外，又添了一桩残暴而可怕的悲剧。因为两族原先友好（彼此有亲缘关系，又是近邻），\Person{霍拉提乌斯}的一个姊妹早已许配给一位\Person{库里阿提乌斯}兄弟；当她看见自己的兄弟身上披着未婚夫的战利品时，禁不住泪水夺眶，竟被她愤怒的兄弟亲手杀死。在我看来，这个少女似乎比整个罗马民族更为人道。我不认为她理当受到指责，因为她为那位已与之缔结婚约的男子而悲泣，或者，也许她是为兄弟竟杀死了曾将妹妹许配给他的人而悲伤。我们为何称赞（在\Person{维吉尔}的作品中）\Person{埃涅阿斯}面对自己亲手杀死的敌人所流之泪呢？为何\Person{马尔克卢斯}在即将毁灭\Place{叙拉古}城之前，忆及其宏伟和正午般的辉煌，且念及万物有共同命运时，竟为之垂泪呢？我以人类的名义要求，如果人们因征服敌人后为之洒泪而受到称赞，那么，一位纤弱少女为被同胞兄弟杀害的恋人所流的泪，就不应视为罪恶。当那位少女因兄弟之手杀死未婚夫而哭泣之时，罗马却正为这等蹂躏加诸于其母邦，并且以耗费自己与\Place{阿尔巴}人共同的鲜血为代价而购得的胜利，欢天喜地。\par

为什么要向我提出\Quote{荣耀}和\Quote{胜利}这些空洞的名称和言词呢？撕下那疯狂幻想的伪装，看看赤裸裸的行为吧：赤裸裸地衡量它们，赤裸裸地判断它们。让对\Place{阿尔巴}的指控，如同当初\Place{特洛伊}因通奸受指控一样。但找不到这样的指控，或任何类似的指控：战争之所以燃起，只为……\par

在慵懒的耳中响起\newline{}\Person{图卢斯}与胜利的欢呼。\par

这种永不停歇的野心之恶习，正是那场自相残杀的内战的唯一动机——这恶习\Person{撒路斯提乌斯}在记述中曾予以谴责；因为当他简短而由衷地赞扬了那无贪欲、人人满足于自己所有的原始时代后，接着写道：\Quote{但自从亚细亚的\Person{居鲁士}，以及希腊的拉刻代蒙人和雅典人，开始征服城邦与民族，并将统治欲望视为战争的充足理由，且认为最大的光荣就在于最大的帝国；} 等等，现在无需引述。这种统治欲望以可怕的灾祸困扰并吞噬人类。正是出于这欲望，罗马在征服\Place{阿尔巴}并夸耀其罪行时，竟称之为光荣。因为，如我们的圣经所说：\BibleQuote{恶人夸耀自己心中的欲望，并祝福贪婪的人，却是主所憎恶的。}\BibleRef{咏 10:3} 那么，除去这些虚假的面具，这些骗人的粉饰，好叫事物能被真实地看见与细察。不要有人对我说，这个或那个人是\Quote{伟大}的人，因为他战胜了某某。角斗士战斗并获胜，这种野蛮行径也得到赞扬；但我想，承受任何懒惰的后果，也强过追求凭这类武器赢得的荣耀。而如果两名角斗士进入竞技场相斗，一位是父亲，另一位是他的儿子，谁能忍受这样的场面？谁不会对此感到憎恶？那么，一个女儿城邦向它的母亲发动的战争，又怎能算是光荣的呢？或者，其差别就在于：战场不是竞技场，广阔的平原上填满的不是两名角斗士，而是两个民族众多精英的尸骸；观看那竞赛的不是圆形剧场，而是整个世界，并为当时在世的人及其后代提供了一场亵渎的奇观，只要这名声流传下去？\par

然而那些守护罗马帝国、并如同此类竞赛的戏剧观众般的神祇们，仍不满足，直到\Person{贺拉提乌斯}兄弟的姊妹被自己兄弟的剑添作罗马方的第三个牺牲品，以致罗马虽赢得胜利，却有许多死亡要哀悼。之后，作为胜利的果实，\Place{阿尔巴}被摧毁，尽管正是在那里，特洛亚的神祇们在\Place{伊利翁}被希腊人劫掠后，并在他们离开\Place{拉维尼乌姆}——\Person{埃涅阿斯}曾在那块流放之地上建立王国——之后，建立了第三个避难所。但或许阿尔巴被摧毁，正是因为神祇们也同样离开了它，按他们惯常的方式，正如\Person{维吉尔}所说：\par

\Quote{从每一座庙宇，每一处圣所离去，\newline{}是那些使这国度神圣的神祇。}\par

确实离开了，现在又从他们的第三个避难所离开，好使罗马在献身于他们之后显得更为明智，因为他们已离弃了另外三座城市。\Place{阿尔巴}，其王\Person{阿穆利乌斯}曾放逐其兄弟，令他们不悦；罗马，其王\Person{罗慕路斯}曾杀害其兄弟，却令他们喜悦。但在阿尔巴被摧毁之前，据说其居民与罗马居民融合，使两城成为一城。好吧，就算如此，事实仍然是，\Person{阿斯卡尼乌斯}之城，特洛亚神祇的第三个退避之所，被其女儿城邦摧毁了。此外，为了拼凑这战争残渣的可怜混合，双方都流了大量鲜血。我又怎能详述继后朝代中屡次重演的相同战争，尽管它们看似已由辉煌的胜利所终结；以及那些一次又一次以大规模屠杀告终，却又一次又一次被订立和平与条约者的后代重新点燃的战争呢？对于这灾难的历史，我们有一个不小的证据：此后没有一位国王关闭过战争之门；因此，纵有所有守护神祇，也无一人能在和平中统治。\par

% structure: p0045 source=h2 target=subsection reason=relative-heading-rank; base=h2

\subsection{罗马诸王的生活与死亡情形}

那么，这些国王自身的结局又如何呢？关于\Person{罗穆路斯}，一则谄媚的传说告诉我们，他被\OriginalTerm{被接升天}{assumed}。但一些罗马史家记述，他因残暴被元老院撕成碎片，并有一人名\Person{尤利乌斯·普罗库鲁斯}被收买，放出风声说\Person{罗穆路斯}曾向他显现，并通过他命令罗马人民敬拜他为神；于是，那些开始对元老院的行为感到愤恨的民众，就这样被安抚平静下来。因为当时也发生了一次\OriginalTerm{日食}{eclipse}；无知无识的民众将此归因于\Person{罗穆路斯}的神力，他们不知道这是由太阳运行的确定法则所致：虽然，这太阳的悲伤本可被视为\Person{罗穆路斯}已被杀的证明，而这太阳光亮的隐失正表明了罪行；正如，事实上，当主因犹太人的残忍与不虔而被钉十字架时，\BibleRef{玛 27:45} 所记载的。因为可充分证明，后者的太阳变暗并非出于天体的自然规律，因为那时正是犹太人的逾越节，这节日只在满月时举行，而自然的日食只发生在月亮的下弦之时。\Person{西塞罗}也足够清楚地表明，\Person{罗穆路斯}的\OriginalTerm{神化}{apotheosis}是想象而非真实。当他在\WorkTitle{论共和国}中，借\Person{西庇阿}之口称赞他时，他说：\Quote{他获得了如此声望，以至于当他在一次日食中突然消失时，就被认为已被接入诸神的行列——没有最高美德声望的凡人，是不可能得到这种推想的。}从\Quote{他突然消失}这些字眼，我们应明白他是被暴风雨或凶手袭击的暴力所神秘地除掉了。因为其他作家不仅提到日食，也提到一阵突然的风暴，这风暴无疑或是给犯罪提供了机会，或是本身结束了\Person{罗穆路斯}的性命。至于\Person{图路斯·霍斯提利乌斯}，他是罗马的第三任国王，本人被闪电击毙。\Person{西塞罗}在同一书中说，\Quote{人们并没有认为他因这种死法而被神化，可能是由于罗马人不愿将他们对\Person{罗穆路斯}确信或信以为真的擢升弄得庸俗化，以免因随意赋予给任何人而使其遭到轻蔑。}在他的一篇抨击文章中，他也直言不讳地说：\Quote{这座城的奠基者\Person{罗穆路斯}，我们通过善意颂扬他的功绩，将他提升到了不朽和神明的地位；}这表明他的神化并非真实，只是出于对他的美德的客套而传扬和称呼的。在对话录\WorkTitle{霍滕修斯}中，当他谈及有规律的日食时，也说它们\Quote{产生同样的黑暗，正如笼罩\Person{罗穆路斯}之死的黑暗那样，而这死亡是发生在一次日食期间。}在此你看到，他毫不规避谈论他的\Quote{死}，因为\Person{西塞罗}更像是一位推理者，而非颂扬者。\par

罗马的其余国王，除了自然死亡的\Person{努马·庞皮利乌斯}和\Person{安库斯·马尔基乌斯}，都遭遇了多么可怕的结局啊！\Person{图卢斯·霍斯提利乌斯}，\Place{阿尔巴}的征服者和毁灭者，正如我所说，他自己连同全家都被雷电烧死。\Person{普里斯库斯·塔奎尼乌斯}被其前任的儿子们杀死。\Person{塞尔维乌斯·图利乌斯}被其女婿\Quote{高傲者}\Person{塔奎尼乌斯}可耻地谋杀，后者继承了王位。就罗马最优秀国王遭此恶毒弑亲而言，这并未将那些神明从祭坛和神庙驱走——这些神明据说因\Person{帕里斯}的通奸而动怒，以同样的方式对待可怜的\Place{特洛伊}，将其抛弃给希腊人的火与剑。不仅如此，那犯下谋杀罪的塔奎尼乌斯反而被允许继承其岳父之位。这个臭名昭著的弑亲者，在靠谋杀获得的统治期间，被允许在许多胜利的战争中凯旋，并利用战利品建造了\Place{卡比托利山}；而神明们没有离开，反而停留、帮助，并容忍他们的王\Person{朱庇特}在一个由弑亲者建造的、金碧辉煌的\Place{卡比托利山}上主持和统治。因为他并非在清白时建造了\Place{卡比托利山}，后来因罪行而被放逐；相反，他是靠着逆伦的罪行才赢得了建造\Place{卡比托利山}的统治权。后来他被罗马人放逐、被禁止入城时，却不是因为自己的，而是因为儿子在\Person{卢克蕾蒂亚}事件中的恶行——一个不仅他不知情，而且他不在场时犯下的罪行。因为那时他正在攻打\Place{阿尔代亚}，为罗马作战；我们无法说要是他知道了儿子的罪行会怎么做。尽管如此，尽管他的意见未被征询，也未被查明，人民剥夺了他的王权；当他带兵返回罗马时，军队被接纳，但他本人却被排拒，被自己的军队抛弃，城门当着他的面关闭。然而，在他向邻国求助，并以多次带来灾难却未成功的军事行动折磨罗马人之后，又被他最依赖的盟友离弃，对重获王国不再抱希望时，据说他在罗马城\Place{图斯库卢姆}隐居静修了十四年，他与妻子在那里一起变老，最终以比他岳父——那位死于女婿之手的人——令人满意得多的方式结束了生命；若传闻属实，他自己的女儿也是同谋。罗马人管这个塔奎尼乌斯，既不叫\Quote{残忍者}，也不叫\Quote{无耻者}，而是叫\Quote{高傲者}；这或许是他们自己的高傲对他在暴政中的盛气凌人感到愤恨。他们对他谋杀自己最优秀的国王、自己的岳父一事是如此不以为意，竟然选他当自己的国王。我不知道他们对如此大罪之人的如此慷慨酬报，是不是更为罪恶。然而，从未听到过神明离开祭坛的说法；或许有人会为这些神明辩护说，他们留在罗马是为了惩罚罗马人，而不是为了帮助他们、让他们受益，用空洞的胜利诱惑他们，用惨烈的战争消耗他们。这就是罗马人在备受赞誉的王政时期的生活，这一时期一直延续到第243年\Quote{高傲者}\Person{塔奎尼乌斯}被驱逐；在此期间，所有那些以如此多的鲜血和灾难换来的胜利，仅仅将罗马的统治从城市向外推进了差不多二十英里；这块地盘无论如何都无法与任何一个微不足道的\Place{格图利亚}小国相比拟。\par

% structure: p0048 source=h2 target=subsection reason=relative-heading-rank; base=h2

\subsection{第16章—— 论最初的罗马执政官，其中一位将另一位从罗马逐出，不久后自己在罗马死于一名受伤敌人之手，这样结束了一段种种逆伦谋杀的生涯。}

插入这个历史阶段，我们也应加上\Person{撒路斯提乌斯}所说的那一段，即国家在正义和节制中治理，而同时面临着对\Person{塔奎尼乌斯}和与\Place{埃特鲁里亚}战争的恐惧。因为只要\Place{埃特鲁里亚}人支持塔奎尼乌斯恢复王位的努力，罗马就为艰难的战事所折磨。因此他说，国家是在恐惧的压力下，而非在公正的影响下，才被正义和节制地治理。在这段极短的时期内，首次设立执政官的那一年是何等灾难性啊，当时王权被废除了！他们没有完成任期。因为\Person{尤尼乌斯·布鲁图}废黜了他的同僚\Person{卢基乌斯·塔奎尼乌斯·柯拉提努斯}，把他逐出罗马；不久后他自己也在战场上阵亡，他杀死敌人的同时被杀，此前他曾处死了自己的儿子和内兄内弟们，因为他发觉他们正阴谋恢复塔奎尼乌斯。正是这件事，\Person{维吉尔}在记录时不寒而栗，即使他看似在赞扬它；因为当他说：\par

\Quote{他呼唤自己叛逆的子孙，\newline{}为受威胁的自由而流血，}\newline{}他随即感叹道，\par

\Quote{不幸的父亲啊！不论后世\newline{}如何评判这事；}\newline{}也就是说，让后世随意评判这事吧，让他们赞美颂扬那杀死自己儿子的父亲，但他是悲惨的。然后他补充道，仿佛在安慰这样一个不幸的人：\par

\Quote{对国家的爱将压倒一切，\newline{}以及对赞美永不熄灭的渴望。}\par

在\Person{布鲁图}的悲剧结局中——他杀了自己的儿子们，虽然杀了敌人、\Person{塔奎尼乌斯}的儿子，却不能活过对方，反而被老塔奎尼乌斯所存活——他同僚\Person{柯拉提努斯}的清白难道不是得到了证实吗？\Person{柯拉提努斯}虽然是个好公民，却遭受了与塔奎尼乌斯本人同样的惩罚，当那位暴君被放逐时。据说布鲁图本人是塔奎尼乌斯的亲属。但柯拉提努斯不幸不仅带有塔奎尼乌斯的血统，还带有他的姓氏。那么，改变他的姓氏，而不是他的国家，才是适当的惩罚：简化他的名字，去掉那个词，只称为\Person{L. 柯拉提努斯}。但他没有被强迫丧失他可以无损害地丧失的东西，反而被剥夺了首任执政官的荣誉，被驱逐出他所爱的土地。难道这就是布鲁图的荣耀——这样一种对共和国既可憎又无益的不义吗？难道正是\Quote{对国家的爱，以及对赞美永不熄灭的渴望}驱使他这样做吗？\par

当\Person{暴君塔克文}被驱逐后，\Person{卢克雷提娅}的丈夫\Person{L. 塔克文·科拉提努斯}与\Person{布鲁图斯}一同被选为执政官。人民行事何等公正，他们看重公民的品格胜于其姓氏！布鲁图斯行事何等不公，竟剥夺了其新官职同僚的荣誉和国家，如果他厌恶其姓氏，本可以只剥夺其姓氏！这就是政府\Quote{以正义与节制治理}时所发生的灾殃和祸患。此外，接替布鲁图斯的\Person{卢克雷提乌斯}，在同一年结束前也因病去世。因此，接替科拉提努斯的\Person{P. 瓦莱里乌斯}，以及填补因卢克雷提乌斯之死所留空缺的\Person{M. 霍拉提乌斯}，共同结束了这多灾多难、犹如丧礼的一年，该年共有五位执政官。正是在这一年，罗马共和国开启了执政官这一新的荣誉与官职。\par

% structure: p0055 source=h2 target=subsection reason=relative-heading-rank; base=h2

\subsection{第十七章.—— 论执政官职设立后困扰罗马共和国的灾难，以及罗马诸神的袖手旁观}

之后，当恐惧逐渐减退——不是因为战争停止，而是因为它们不那么激烈了——那段以\Quote{以正义与节制治理}事物的时期便走到了尽头，随之而来的事态，\Person{撒路斯提乌斯}曾这样简短地描述：\Quote{于是贵族开始像榨取奴隶一样压迫平民，像国王曾做过的那样判处他们死刑或鞭笞，将他们赶出自己的田地，对那些一无所有者施行暴政。平民被这些压迫措施压垮，尤其是受高利贷盘剥，又被迫为连绵的战争既出钱又出力，最终拿起武器，退守到\Place{阿文廷山}和\Place{圣山}，从而为自己争得了保民官和保护性法律。但只有第二次布匿战争才使双方停止了纷争与冲突。}但我为何要花费时间写这些事，或者让别人花费时间读它们呢？让\Person{撒路斯提乌斯}的简要概括足以表明，在那段漫长时期，直至第二次布匿战争，共和国是何等凄惨——它对外为无休止的战争所困扰，内部又被公民间的争吵与不和撕裂。因此他们所夸耀的那些胜利，并非幸福者实在的喜乐，而是可怜人空洞的慰藉，是对骚动不安者的诱人煽动，使之制造一场又一场的灾难。但愿善良审慎的罗马人不要因我们这样说而生气；其实我们既无须恳求他们的饶恕，也无须谴责他们的愤怒，因为我们知道他们不会心怀愤恨。因为我们说的话并不比他们自己的作家更严厉，而且远不如他们那样精心和惊人；然而他们却勤读这些作家，并强迫子女学习。但是，那些生气的人，如果我像\Person{撒路斯提乌斯}那样说话，他们会怎么对我呢？\Quote{频繁的暴民骚乱、叛乱，直至最终的内战，变得司空见惯，而少数几位民众所依附的领袖，则以寻求元老院与平民利益为体面的借口，谋取最高权力；公民被评判为好或坏，不考虑他们对共和国的忠诚（因为所有人都同样腐败）；但那些富有而危险的有势力者，却被视为好公民，因为他们维护了现存秩序。}现今，如果那些史家认为，正当的言论自由要求他们不应对自己城邦的缺陷保持沉默——他们在许多地方曾高声赞誉这城邦，却不知道那另一座真实的城，在其中公民身份享有永久的尊荣——那么，我们又当如何呢？我们的自由应当大得多，因为我们对天主的望德更美好、更确实，而他们却将这个时代的灾难归咎于我们的基督，为要使那些学识较少、信心较弱的人远离那唯一能享受永恒和真福生命的城。我们控告他们的神明，并不比他们自己阅读和传播的那些作家所说的更为可怕。事实上，我们所说的一切都取自他们，而且还有更多更恶劣的事是我们无法说的。\par

那么，当罗马人因欺骗的诡计被诱去事奉他们，并遭受这些灾难时，那些为了此世微薄而虚幻的繁荣理应受到敬拜的神在哪里呢？当执政官\Person{瓦莱里乌斯}在保卫被流亡者和奴隶纵火的\Place{卡比托利欧山}时被杀时，他们在哪里？他本人更能保卫\Person{朱庇特}的神殿，而他前来救援的那群神明及其至高无上的大王却无法保卫他。当这座城被不停的骚乱耗尽，在某种平静中等待派往\Place{雅典}去借取法律的使节返回，却遭可怕的饥荒和瘟疫蹂躏时，他们在哪里？当人民再次被饥荒困扰，首次设立了一位市场长官时；当\Person{斯普里乌斯·梅利乌斯}，随着饥荒加剧，将粮食分发给饥饿的民众，却被指控觊觎王权，并在该市场长官的提议下，依据年迈的独裁官\Person{卢基乌斯·昆克提乌斯}的授权，由骑兵统领\Person{昆图斯·塞尔维利乌斯}处死——这一事件引发了一场严重而危险的骚乱时，他们在哪里？当那场非常严重的瘟疫降临\Place{罗马}时，因着这场瘟疫，人民在徒劳地向那些无能的神发出漫长、厌烦且无用的祈求之后，想到了举行前所未有的\OriginalTerm{神明宴}{Lectisternia}——也就是说，他们为敬拜神而摆放床榻，这个神圣仪式，或更确切地说是亵渎仪式，因此得名——他们在哪里？当罗马军队在连续十年的败绩中，在\Place{维爱}人那里屡遭重大损失，若非\Person{弗里乌斯·卡米卢斯}的救援便已全军覆没，而卡米卢斯后来却被忘恩负义的祖国放逐时，他们在哪里？当\Place{高卢}人攻占、洗劫、焚烧并蹂躏罗马时，他们在哪里？当那场令人难忘的瘟疫造成如此毁灭，连\Person{弗里乌斯·卡米卢斯}也死于其中时——他先是保卫了忘恩负义的共和国免受维爱人侵害，后又从高卢人手中拯救了她——他们在哪里？不仅如此，在这场瘟疫期间，他们引入了一种新的舞台戏剧瘟疫，它传播了更致命的感染，不是针对身体，而是针对罗马人的道德；他们在哪里？当另一场可怕的瘟疫降临这座城时——我指的是被归咎于数量惊人的罗马贵妇的下毒事件，她们的品行感染了比任何瘟疫都更致命的疾病——他们在哪里？或者当带兵的两名执政官在\Place{卡夫丁峡谷}被\Place{萨莫奈}人围困，被迫签订一项可耻的条约，六百名罗马骑士被扣为人质；同时，部队在放下武器、被剥去一切后，每人只着一件衣服，被迫从轭门下通过时，他们在哪里？或者当一场严重的瘟疫期间，闪电击中罗马营地并杀死了许多人时，他们在哪里？或者当罗马被另一场不堪忍受的瘟疫所迫，派人前往\Place{埃皮达鲁斯}去请医神\OriginalTerm{埃斯库拉庇乌斯}{Æsculapius}时——因为\Person{朱庇特}年轻时频繁的通奸，或许没有给这位长期统治\Place{卡比托利欧山}的众王之王留下任何空闲去研究医学——他们在哪里？或者有一次，\Place{卢卡尼亚}人、\Place{布鲁提}人、\Place{萨莫奈}人、\Place{托斯卡纳}人和\Place{塞农}高卢人联合对抗罗马，先杀了她的使节，然后击溃了一支由指挥官率领的军队，除了指挥官和七位将官外，还斩杀了一万三千人时，他们在哪里？或者当人民在罗马严重而持续不断的骚乱之后，最终洗劫了城市并撤退到\Place{贾尼科洛}山时；事态如此严重，以致\Emph{霍滕修斯}被任命为独裁官——这一官职他们只在极端紧急时才采用；而他在带回人民后，仍在任内便去世了——这在任何独裁官史无前例，对如今已有\OriginalTerm{埃斯库拉庇乌斯}{Æsculapius}在内的那些神来说，是一种羞辱时，他们在哪里？\par

在那段时期，各地战争频仍，因兵员匮乏，他们甚至征召\Emph{\OriginalTerm{无产者}{proletarii}}入伍；这名称之由来，乃是因为他们过于贫穷，无力自备武器，却有闲暇生儿育女。\Person{皮洛士}，希腊国王，当时声名远播，应\Place{塔伦托}人之邀，兴兵对抗罗马。正是对他，\Person{阿波罗}在被问及此行成败时，以几分戏谑口吻说出了一道如此模棱两可的神谕，以致无论何种结局发生，这位神明都可被算为灵验。因为他这样措辞神谕：无论\Person{皮洛士}被罗马人征服，还是罗马人被\Person{皮洛士}征服，这位预言神都能安稳地等待结果。随后，双方军队遭受了何等可怕的屠杀！然而\Person{皮洛士}成了胜者，本可据此宣称\Person{阿波罗}是真先知——依他所理解的神谕——若非罗马人在下一场战斗中变成了征服者。正当这些惨烈的战争进行之际，妇女们当中爆发了一场可怕的疫病。因为孕妇在分娩前便死去了。我想，\Person{厄斯库拉庇乌斯}会以此事为自己开脱，说他所司职的是首席医生，而非助产士。牲畜也同样死亡；以致人们相信，整个动物种群注定要灭绝。至于那个令人难忘的冬天，天气冷得不可思议，\Place{罗马广场}的积雪厚得骇人，整整四十天未化，\Place{台伯河}也结了冰；我对此又该说什么呢？这类事情若发生在我们的时代，我们必会听到仇敌何等的指控！还有另一场大瘟疫，肆虐良久，吞噬无数生灵；我又该说什么呢？尽管服用了\Person{厄斯库拉庇乌斯}的各种药剂，疫情在第二年反而更趋严重，直到最后人们求助于\WorkTitle{西彼拉占语书}——这种神谕，正如\Person{西塞罗}在其\WorkTitle{论占卜}中所说，其意涵全系于解读者，他们凭自己的能力或意愿作出可疑的推测。此次，瘟疫的起因被归咎于众多神殿已被用作私人住宅。就这样，\Person{厄斯库拉庇乌斯}暂时免于可耻的疏忽失职或医术不精的指控。可是，为何会有那么多人被容许毫无妨碍地占据圣所呢？难道不是因为人们长期向如此众多的神明徒然祈求，以致圣所渐渐失去了崇拜者，就这样空置下来，才可毫无冒犯地至少改作俗用吗？而那些神殿，当时为遏止瘟疫曾被费力地辨识并修复，随后又被弃置，再度用于同样的俗用了。若非它们如此湮没无闻，也就无法被指为\Person{瓦罗}博学多闻的证据，即他在论圣地的著作中引述了那么多人所未知的神殿。与此同时，修复神殿并未使瘟疫得治，反倒只为诸神提供了一个动听的借口。\par

% structure: p0059 source=h2 target=subsection reason=relative-heading-rank; base=h2

\subsection{第18章 罗马人在布匿战争中遭受的灾难，并未因诸神的护佑而减轻。}

而在布匿战争中，当胜利在两个王国之间长久难决，当两个强国竭尽全力、倾尽资源彼此相争时，多少较小的王国被粉碎，多少繁荣昌盛的大城被摧毁，多少邦国遭倾覆而灭亡，多少远近地区的土地变为荒芜！胜利者的一方多少次又成为战败者！多少生灵，无论从军者还是平民，惨遭屠戮！又有多少庞大的舰队在交战中受损，或因各种海难沉没！我们若试图数算或列举这些灾难，我们便会成为历史编纂者。那段时期，罗马惊恐万分，转向各种徒劳可笑的手段。据\WorkTitle{西彼拉占语书}的权威，\OriginalTerm{百年节赛会}{secular games}被重新设立，这个赛会早在一百年前便已首创，却在较为宁泰的岁月中被遗忘了。奉献给冥府诸神的赛会也由大祭司们重新恢复；因为它们同样在较为宁泰的岁月中被废弃了。这也难怪；因为当这些赛会恢复时，大量垂死之人使整个地狱为所得之丰而欢欣，并恣意作乐：诚然，那些凶猛的战事、灾难深重的争斗、血腥的胜利——时而偏向一方，时而偏向另一方——虽然对人类而言是至惨之祸，却为魔鬼们提供了绝佳的娱乐和丰盛的宴席。然而，在第一次布匿战争中，再没有比那场罗马人战败、\Person{雷古卢斯}被俘的灾难更为惨重的了。我们已在前面两卷中提过此人，他无疑是一位伟人，曾击败并征服迦太基人，本可以结束第一次布匿战争，若非他对赞誉和荣耀的过度渴求促使他强加给疲惫不堪的迦太基人以超出他们承受能力的苛刻条件。如果这位人物出人意料的被俘、不光彩的囚禁、他对誓言的忠诚，以及他极为残酷的死亡，都不能令诸神脸红，那么诸神的脸皮确实是铜铁般厚而毫无血色的。\par

\Place{台伯河}发生了罕见的洪水，几乎摧毁了城市的所有低洼地区；一些建筑物被洪水的暴力冲走，而另一些则被洪水退后仍积存的水浸泡而腐烂。这场灾祸之后又发生了一场破坏性更大的火灾，因为大火烧毁了\Place{古罗马广场}周围一些较高的建筑，甚至连它自己的庙宇——\Person{维斯塔}神庙也未能幸免；在这庙里，被选出来承受这份荣誉、更确切地说这份惩罚的贞女们，通过不断添加新的燃料，仿佛在赋予火焰永恒的生命。但在我们所谈论的那时，庙中的火并不满足于被保持不灭：它猛烈燃烧起来。当那些贞女们被火势吓得无法拯救那些\OriginalTerm{致命的神像}{fatal images}时——这些神像先前在三个城市被供奉，已给它们带去了毁灭——\OriginalTerm{祭司}{priest}\Person{麦特鲁斯}不顾自身安危，冲了进去，抢救出了那些\OriginalTerm{圣物}{sacred things}，虽然在此过程中自己被烧得半熟。因为要么那火甚至不认识他，要么就是\Person{火女神}本人在场——一位在场的女神并不会逃离火焰，假如她真在场的话。但你们看，一个人能给予\Person{维斯塔}的帮助，竟比\Person{维斯塔}能给予他的还要多。如果这些神都不能使自身免于火灾，那么对于他们号称守护的国家，他们又怎能帮助其抵御火灾或洪水呢？事实已经表明，他们是无用的。如果我们的对手主张，他们供奉偶像并非为了保障现世的福泽，而是作为永恒事物的象征；那么这样一来，尽管这些象征，如同所有物质可见之物，可能会消亡，但并不会损害它们为之被祝圣的事物；至于偶像本身，则可以为它们原先所服务的同一目的重新制造。若是如此，我们的这些异议便属徒然。然而，他们怀着可悲的盲目，认为通过可朽坏的神祇的干预，国家的现世福祉与世俗繁荣就能免于朽坏。因此，当人们提醒他们，即使神祇们与他们同在的时候，这种福祉与繁荣也遭受挫败时，他们虽无法为自己辩护，却羞于改变所持的看法。\par

% structure: p0062 source=h2 target=subsection reason=relative-heading-rank; base=h2

\subsection{第二次布匿战争的灾难，它使交战国力耗尽。}

至于第二次布匿战争，若要详述它给交战双方带来的灾难，那将是冗长乏味的；这场战争如此旷日持久且胜负不定，以至于（甚至连那些著书立说并非旨在叙述战争，而是颂扬\Place{罗马}统治的作家们也承认）最终获胜的一方与其说是征服者，倒不如说更像是被征服者。因为，当\Person{汉尼拔}从\Place{西班牙}倾巢而出，翻越\Place{比利牛斯山}，横扫\Place{高卢}，突破\Place{阿尔卑斯山}，并在其整个推进过程中一路劫掠征服、壮大实力，如同洪流般涌入\Place{意大利}时，随后的战事何等血腥，交战何等持续不断！\Place{罗马}人遭受了何等频繁的失败！多少城镇倒向了敌人，又有多少被攻占征服！那些战斗何等可怖，而\Place{罗马}人的败绩又曾多少次为\Person{汉尼拔}的军队增添荣耀！对于\Place{坎尼}那场极其惨重的败仗，我又能说什么呢？在那里，就连残忍如\Person{汉尼拔}者，也因杀戮他最痛恨的敌人过甚而感到餍足，竟下令饶恕他们。他从这片战场往\Place{迦太基}送去了三\OriginalTerm{蒲式耳}{bushel}的金指环，以此表示\Place{罗马}的显贵阶层当天折损之多，以至于用容量而非个数来想见此数更为容易；而对于那些没有指环区别身份、横尸遍野的普通士卒——其人数之众与其地位之低微成正比——则只能凭推测而无法准确报告了。事实上，此役过后士兵奇缺，\Place{罗马}人竟以特赦为许诺征募罪犯，以自由为饵征召奴隶，用这些声名狼藉之辈并非补给军队，而是创编了一支军队。但这些奴隶——或者赋予他们全部称谓的话，这些被征召为\Place{罗马共和国}作战的释奴——缺少武器。于是他们从庙宇中取出武器，仿佛\Place{罗马}人在对自己的神说：\Quote{放下你们长久以来徒然持有的那些武装吧，或许我们的奴隶能有效地使用你们这些神无力运用的东西。}那时，国库也已枯竭到无力支付军饷，私人资财被用于公共用途；个人捐输如此慷慨，以至于除却每人所佩戴的金指环和\OriginalTerm{布拉}{bulla}——那区区表明身份的标帜——外，没有一位元老，更不用说其他阶层和部落的人了，为自己保留任何黄金。但如果在我们这时代，他们落到这般贫困的地步，又有谁能忍受得住他们的责难呢？即便是现在，当为了不必要的享乐而花费在伶人身上的金钱，比那时发给军团的军饷还多时，他们的责难就已几乎令人无法容忍了。\par

% structure: p0064 source=h2 target=subsection reason=relative-heading-rank; base=h2

\subsection{\Place{萨贡托}人的毁灭，他们因效忠\Place{罗马}而灭亡，却未得到\Place{罗马}神祇的任何援助。}

但在第二次布匿战争的所有灾难中，没有比\Place{萨贡托}人的遭遇更令人悲痛、更易引起激烈怨言的了。这座\Place{西班牙}城市，与\Place{罗马}极为友好，却因忠于罗马人民而遭毁灭。当\Person{汉尼拔}背弃与罗马人的条约时，他寻找挑动战争的借口，于是猛烈进攻\Place{萨贡托}。消息传到\Place{罗马}，使者被派往\Person{汉尼拔}处，敦促他解围；当此劝告未被理会时，他们便前往\Place{迦太基}，抗议条约被破坏，随后未能达到目的而返回\Place{罗马}。同时围攻仍持续；在第八或第九个月，这座富庶却不幸的城市，虽为其本国及\Place{罗马}所珍爱，终被攻陷，并遭受了令人不忍卒读、更无法详述的暴行。然而，因这事与本文主旨直接相关，我将简要提及。首先，饥荒耗尽了\Place{萨贡托}人，以至于有人甚至吃人尸：至少史籍如此记载。随后，当他们精疲力竭，为至少免于落入\Person{汉尼拔}之手的羞辱，便公开堆起一座巨型火葬柴堆，纵身投入火焰，同时又用刀剑杀死自己的儿女与自身。这些神祇——这些放荡贪食者，垂涎于丰盛祭品，口中吐出虚假预言——在如此情形下就无能为力吗？难道他们不能为保护这座与罗马人民紧密结盟的城市而干预，或阻止它因忠于他们自己曾居间促成的联盟而灭亡吗？\Place{萨贡托}忠实地遵守了在这些神祇面前所订、并以誓言庄严缔结的条约，却被一个背誓者围攻、攻陷并毁灭。若后来当\Person{汉尼拔}兵临罗马城下时，是神祇们以雷电和风暴恐吓他，将他驱离，那么我要问，为何他们不事先如此干预？因我敢断言，以风暴示警，与其说是为保卫罗马人自己——他们为自己的事业而战，且有充足资源对抗\Person{汉尼拔}——不如说更体面地用于保卫罗马的盟友，后者因不愿背弃对罗马人的信义而身处险境，且自身没有资源。因此，若这些神祇真是罗马繁荣与荣耀的守护者，他们本应保护那荣耀免于\Place{萨贡托}灾难的污点；而相信罗马在\Person{汉尼拔}手中得以幸存，是出于那些无法拯救\Place{萨贡托}城因其忠于与罗马联盟而灭亡的神祇保佑，是多么愚蠢。若\Place{萨贡托}的居民曾是基督徒，并因基督信仰而如此受难（虽然，基督徒当然不会用火与剑对待自身），他们乃是怀着那源于对基督的信德的望德而受难——那不是暂世赏报，而是无尽永恒的幸福。那么，这些神祇的辩护者和卫道士，在因\Place{萨贡托}人的血而受责时，将何以自辩？因他们正是为此目的——确保此生短暂无常的繁荣——而公然崇拜并呼求这些神祇。难道除了在\Person{雷古卢斯}之死一事上提出的辩词外，还能说什么吗？虽然两事有所不同，一为个人，一为整个群体，然而毁灭的起因都是信守诺言。正是这使\Person{雷古卢斯}甘愿回到敌人那里，也正是这使\Place{萨贡托}人不愿叛乱投敌。那么，守信触怒了神祇吗？抑或可能不仅个人，甚至整个群体在神祇垂青之下仍会遭难灭亡？让我们的对手二者择一。如果一方面，那些神祇因守信而愤怒，就让他们招募背誓者为崇拜者。如果另一方面，在神祇垂青之下，人与国家仍可遭受巨大可怕的灾祸，并最终灭亡，那么他们的敬拜并不产生幸福果实。因此，那些认为因敬拜被废除而遭遇不幸的人，当放下他们的愤怒；因为即便神祇不仅与他们同在，且垂青于他们，他们仍可能被留下哀叹不幸的命运，甚或如同\Person{雷古卢斯}和\Place{萨贡托}人一般，遭受残酷折磨，最终悲惨灭亡。\par

% structure: p0066 source=h2 target=subsection reason=relative-heading-rank; base=h2

\subsection{\Place{罗马}对其解放者\Person{西庇阿}的忘恩负义，以及\Person{撒路斯提乌斯}称为最好时期的\Place{罗马}风尚}

我略过许多事情，为的是不超出我为自己设定的著作篇幅。我来到第二次布匿战争与最后一次布匿战争之间的时期，根据\Person{撒路斯提乌斯}，\Place{罗马}人在这段时期内生活在最伟大的德行与和谐之中。在这德行与和谐的时期，伟大的\Person{大西庇阿}，\Place{罗马}和\Place{意大利}的解放者，他以惊人的才能结束了第二次布匿战争——那场可怕、毁灭性、危险的冲突——他击败了\Person{汉尼拔}，征服了\Place{迦太基}，据说他的一生都奉献给了神明，在神庙中被珍视——这位西庇阿，在取得如此胜利之后，却不得不屈服于敌人的指控，离开了他以勇气拯救并解放了的祖国，在\Place{利特努姆}镇度过余生。他是如此不在意从流放中被召回，据说甚至吩咐过，连他的遗骸都不要留在那忘恩负义的祖国中。也正是在那时，代执政官\Person{格奈乌斯·曼利乌斯}在征服\Place{加拉太}人之后，将比一切敌军更具破坏力的\Place{亚细亚}奢侈风气引入了\Place{罗马}。正是在那时，铁制床架和昂贵的地毯首次被使用；也正是在那时，女歌手被允许出现在宴会上，其他放荡的恶习也被引入。但目前我打算讲的，不是人自愿犯下的恶，而是他们身不由己遭受的恶。所以我提到西庇阿的事例，他屈服于敌人，并在救国的流放中死去，是与当前的讨论相关的；因为这是他从那些\Place{罗马}神明那里得到的回报，他曾在\Person{汉尼拔}面前保住了他们的神庙，而人们崇拜神明仅仅是为了获得现世的幸福。然而，正如我们所见，\Person{撒路斯提乌斯}宣称\Place{罗马}的德行在那时从未如此美好，因此我判断提及那时引入的亚细亚奢靡是对的，以便可以看出，只有当那个时期与其它时期相比较时，他所说的才是正确的——在那些时期，道德肯定更败坏，内乱更激烈。因为在那时——我指的是第二次和第三次布匿战争之间——那臭名昭著的\WorkTitle{沃科尼亚法}被通过，该法禁止任何人立妇女，即便是独生女，为继承人；我无法想象还有什么法律比这更不公正。的确，在这两次布匿战争之间的间隔，\Place{罗马}的苦难在一定程度上较少。对外，他们的军队被战争消耗，却也被胜利所安慰；而国内没有像其他时候那样的骚乱。但当最后一次布匿战争以\Place{罗马}的竞争对手彻底毁灭而告终，它很快被\Person{小西庇阿}——他因此为自己赢得了阿非利加努斯的称号——所征服，那时\Place{罗马}共和国被如此众多的灾祸所淹没，这些灾祸源于繁荣和安全所引发的腐败道德，以至于\Place{迦太基}的突然倾覆，看起来比它长期的敌对行为更加严重地伤害了\Place{罗马}。在随后的整个时期，直到\Person{凯撒·奥古斯都}的时代，他看起来完全剥夺了\Place{罗马}人的自由——一种在他们自己看来已不再光荣，而是充满纷争和危险，并已变得极为颓废和衰败的自由——他再次将一切置于一位君主的意志之下，仿佛为共和国奄奄一息的暮年注入了新的生命，开启了一个新的\Emph{政制}——在这一整个时期，我说，在种种不同的场合遭受了许多军事灾难，所有这些我都略过不提。尤其是那极不名誉的\Place{努曼提亚}条约，它被耻辱所玷污；因为据说那神圣的鸡从围栏里飞了出来，从而预示着对执政官\Person{曼奇努斯}的灾难；就好像，在那些年里，小小的\Place{努曼提亚}城抵抗\Place{罗马}围城军队，成为共和国的恐惧之时，其他将领们都是在不祥的预兆下向它进军的。\par

% structure: p0068 source=h2 target=subsection reason=relative-heading-rank; base=h2

\subsection{二十二、——\Person{米特里达梯}的敕令，命令杀死在\Place{亚细亚}的所有\Place{罗马}公民。}

我说，这些事我略过不谈；但我绝不能对\Place{亚细亚}国王\Person{米特里达梯}的命令保持沉默，这命令说在一天之内，所有在\Place{亚细亚}任何地方居住的\Place{罗马}公民（那里有大量的人在从事他们的私人业务）都应被处死：而这命令被执行了。当时呈现了多么悲惨的景象啊，当每个人无论身在何处——在田野或路上，在城镇，在自己家中，或是在街上，在市场或庙宇，在床上或餐桌——都突然遭到背叛性的谋杀！想想垂死者的呻吟，旁观者的眼泪，甚至刽子手自己的眼泪。因为强迫这些受害者宿主不仅要在自己家中目睹这些可憎的屠杀，甚至还要亲手执行的，是何等残酷的必然：要突然改变脸色，从友谊的温和亲切，在和平之中着手战争之事；并且，我该说，给予并承受创伤，被杀者身体被刺穿，杀人者灵魂被刺伤！那么，所有这些被谋杀的人，是否都轻视了占卜呢？当他们离开家园，踏上那致命的旅程时，难道没有公共的或家中的神明可以求问吗？如果他们没有，那么我们的对手就没有理由在这一点上抱怨这些基督徒的时代，因为很久以前\Place{罗马}人就把占卜视为无用了。另一方面，如果他们确实咨询了预兆，让我们告诉他们，即使在当时这类事情未被禁止，反而被人的法律——如果不是神圣的法律——所许可的情况下，他们从中得到了什么好处。\par

% structure: p0070 source=h2 target=subsection reason=relative-heading-rank; base=h2

\subsection{二十三、——困扰\Place{罗马}共和国的内部灾难，以及随之而来、攫住所有家畜的不祥疯狂。}

但是，我们还是尽可能简明地提一下那些更令人烦恼的灾难吧，因为它们就发生在身边；我指的是那些被误称为\Quote{内部}的纷争，因为这类纷争毁灭的正是内部利益。那些骚动如今已演变为城内的战争，鲜血横流，派系之间彼此激愤，不再凭借争吵与言辞交锋，而是动用武力与刀兵。\Place{罗马}人的鲜血流成了怎样的血海啊！同盟战争、奴隶战争、内战，又在\Place{意大利}造成了多少荒残与破坏！在\Place{拉丁}人对\Place{罗马}发动同盟战争之前，所有为人所用的牲畜——狗、马、驴、牛，以及其余凡是服从于人的动物——全都突然变得狂野起来，忘记了家养的驯良，离弃棚圈四处游荡，生人乃至它们自己的主人都不敢毫无危险地就近。如果这算是一个凶兆，那么无论是不是凶兆，本身就是一场严重灾难的瘟疫，它所预示的祸患又该是何等严重！要是这件事发生在我们的时代，异教徒对我们的暴怒，一定远远超过当年他们的牲畜对他们自己的暴怒。\par

% structure: p0072 source=h2 target=subsection reason=relative-heading-rank; base=h2

\subsection{第24章。—— 论\Person{格拉古兄弟}的骚乱引发的内乱。}

内战起源于\Person{格拉古兄弟}因土地法而煽动的骚乱；因为他们想要将贵族非法占有的土地分给人民。然而，改革一个长期存在的弊端，乃是一件充满危险的行动，或者更确切地说，正如结果所证明的，是毁灭性的。因为伴随老\Person{格拉古}之死的是怎样的灾难啊！不久之后，当弟弟遭遇同样命运时，又发生了怎样的屠杀啊！因为贵族和平民都不分青红皂白地被屠杀；而且这不是通过法律权威与程序，而是由暴民和武装暴乱者执行的。在小\Person{格拉古}死后，执政官\Person{卢基乌斯·奥皮米乌斯}，他曾在城内与他交战，并击败和杀死了他本人及其同伙，还屠杀了许多公民，之后又对其他人进行了司法调查，据报告，被处死的人多达三千人。由此可以推断，在暴乱冲突中究竟有多少人倒下，因为就连司法调查的结果也是如此血腥。杀死\Person{格拉古}的刺客将他的头颅卖给执政官，以换取与头颅等重的黄金，这是事先约好的。在这场屠杀中，曾任执政官的\Person{马尔库斯·富尔维乌斯}和他的所有子女也都被处死。\par

% structure: p0074 source=h2 target=subsection reason=relative-heading-rank; base=h2

\subsection{第25章。—— 论根据元老院法令在这些骚乱与屠杀现场建造的\OriginalTerm{孔科耳狄亚}{Concordia}庙。}

元老院的确通过了一道冠冕堂皇的法令，就在那场灾难性的暴动发生之处，也就是各个等级众多公民丧命的地方，建起了一座\OriginalTerm{孔科耳狄亚}{Concordia}庙。我猜想，这是要让\Person{格拉古兄弟}受惩的纪念物刺痛鼓吹者的眼睛，影响他们的记忆。但这岂不就是要嘲弄神明，给这样一位女神建庙，如果她真在这座城里，就不会让自己被如此的分裂撕碎？还是说，孔科耳狄亚要为那场流血负责，因为她抛弃了公民的心思，因此被关进了那座神庙？因为如果他们有任何顾及一致性的想法，为什么不在那个地方建一座\OriginalTerm{狄斯科耳狄亚}{Discordia}神庙？或者，有什么理由说孔科耳狄亚是女神，而\OriginalTerm{狄斯科耳狄亚}{Discordia}却不是女神？\Person{拉贝奥}的区分在这里适用吗？他曾将一位视为善神，另一位视为恶神——这种区分似乎仅仅是因为他注意到在\Place{罗马}城，既有一座\OriginalTerm{费布里斯}{Febris}神庙，也有一座\OriginalTerm{萨卢斯}{Salus}神庙。但按照同样的道理，狄斯科耳狄亚也应当和孔科耳狄亚一样被尊为神。\Place{罗马}人冒了很大的风险，去招惹如此邪恶的一位女神，而且忘记了\Place{特洛伊}的毁灭就是因她发怒而造成的。因为，她因未被邀请与其他神明一同[参加\Person{佩琉斯}与\Person{忒提斯}的婚礼]而愤慨，便投下金苹果，在三位女神中制造了不和，从而引发了天上的争斗、\Person{维纳斯}的胜利、\Person{海伦}的抢掠以及\Place{特洛伊}的毁灭。因此，如果她也许因为\Place{罗马}人认为她配不上在他们城里的众神中拥有一座神庙而生气，从而用如此的骚乱扰乱了城邦，那么当她看到她敌手的神庙就建在那场屠杀的现场，换句话说，建在她亲手造就的灾难现场时，她又会激起多么狂烈的愤怒啊！那些博学多才的人对于我们嘲笑这些蠢事感到愤怒；然而，他们自己既是善神也是恶神的崇拜者，无法逃脱关于孔科耳狄亚与狄斯科耳狄亚的困境：他们要么忽略了这些女神的崇拜，却偏爱古老神庙中供奉的费布里斯和战神，要么他们崇拜了她们，但最终孔科耳狄亚抛弃了他们，而狄斯科耳狄亚则狂暴地将他们掷入内战。\par

% structure: p0076 source=h2 target=subsection reason=relative-heading-rank; base=h2

\subsection{第26章。—— 论\OriginalTerm{孔科耳狄亚}{Concordia}庙建成后接踵而至的各种战争。}

但他们以为，在演说家们目光所及之处修建和谐神殿，作为对\Person{格拉古兄弟}所受惩罚和死亡的纪念，就是为制止内乱设置了有效障碍。其效果如何，由随后发生的更为可悲的战争便可见一斑。因为在此之后，演说家们非但不以\Person{格拉古兄弟}为前车之鉴，反而力图超越他们的计划；就像保民官\Person{卢基乌斯·萨图尔尼努斯}、\OriginalTerm{裁判官}{praetor}\Person{盖乌斯·塞尔维利乌斯}，以及稍后的\Person{马尔库斯·德鲁苏斯}那样，所有这些人都煽动起骚乱，最初引发了流血，随后又引发了同盟战争，\Place{意大利}在这场战争中遭受重创，沦落到一片荒凉、破败不堪的可悲境地。接着是奴隶战争和内战；在这些战争中，发生了多少战斗，流了多少血，以至于\Place{意大利}诸民族——他们构成\Place{罗马帝国}的主力——几乎全被征服，犹如野蛮人一般！然后，甚至历史学家自己也发现难以解释：这场奴隶战争是如何由极少数人（肯定不到七十名角斗士）发动起来的，又有多少凶悍残忍之徒依附于他们，这群人击败了多少罗马将领，以及他们如何摧毁了许多地区与城市。那还不是唯一的一次奴隶战争：\Place{马其顿}行省，随后是\Place{西西里}和沿海地区，也被成群的奴隶洗劫一空。谁能充分描述，要么是海盗最初犯下的可怕暴行，要么是他们随后对\Place{罗马}发动的战争？\par

% structure: p0078 source=h2 target=subsection reason=relative-heading-rank; base=h2

\subsection{第27章 \Person{马里乌斯}与\Person{苏拉}之间的内战}

但当\Person{马里乌斯}，身上沾满同胞的鲜血——党派狂怒将这些同胞当作了祭品——自己也被击败并逐出城市时，它几乎还没来得及舒畅地呼吸，就用\Person{西塞罗}的话说：\Quote{\Person{秦纳}和\Person{马里乌斯}一同返回并占领了它。那时，的确，国家的首要人物被处死，它的光辉熄灭了。\Person{苏拉}后来为这场残酷的胜利进行了报复；但我们无需细说有多少生命损失，给共和国带来了多大的毁灭。} 因为对于这场报复——它比它所惩罚的罪行若不受惩罚地发生，更具毁灭性——\Person{卢坎}说：\Quote{疗法过度，且与疾病太过相似。有罪者灭亡了，但只有有罪者幸存下来：随后，私人的仇恨和愤怒，不受法律约束，得以肆意发泄。} 在那场\Person{马里乌斯}与\Person{苏拉}之间的战争中，除了战场上阵亡的人，城市也到处是尸体，在街道、广场、市场、剧院和神庙里；以至于很难计算，胜利者在胜利之前与之后，是为了成为胜利者而杀人更多，还是因为成了胜利者而杀人更多。\Person{马里乌斯}一取胜并从流放中归来，除了到处进行的屠杀，执政官\Person{屋大维}的头颅被悬挂在\OriginalTerm{演讲台}{rostrum}上；\Person{凯撒}和\Person{菲姆布里亚}在各自家中被暗杀；\Person{克拉苏}父子二人当着彼此的面被杀害；\Person{贝比乌斯}和\Person{努米托里乌斯}被钩子拖行，开膛破肚；\Person{卡图卢斯}饮下毒药，才逃出敌人之手；朱庇特的\OriginalTerm{祭司}{flamen}\Person{麦鲁拉}割开血管，用自己的血向他的神行灌奠礼。而且，凡\Person{马里乌斯}未曾伸手回礼的人，立刻就被砍倒在他面前。\par

% structure: p0080 source=h2 target=subsection reason=relative-heading-rank; base=h2

\subsection{第28章 \Person{苏拉}的胜利——\Person{马里乌斯}暴行的复仇者}

随后是\Person{苏拉}的胜利，他被称为\Person{马里乌斯}暴行的复仇者。然而，他的胜利不仅以大量流血为代价；当敌对行动结束后，敌意依然幸存，随后的和平与战争一样血腥。老\Person{马里乌斯}造成的屠杀尚未过去，同属该派的小\Person{马里乌斯}和\Person{卡尔博}又增添了更大的暴行。因为当\Person{苏拉}逼近，他们不仅对胜利绝望，对生命本身也绝望了，便不分敌友地大肆屠杀。而且，他们不满足于使\Place{罗马}的每个角落沾满鲜血，还包围了元老院，将元老们从\OriginalTerm{元老院会堂}{curia}像从监狱中一样带去处死。\OriginalTerm{大祭司}{pontiff}\Person{穆基乌斯·斯凯沃拉}在维斯塔祭坛前被杀，他曾紧抱祭坛，因为在\Place{罗马}没有比维斯塔神庙更神圣的地方了；他的血几乎熄灭了由贞女们持续照管的不灭之火。随后，\Person{苏拉}以胜利者身份进入城市，此前他在\Place{公共庄园}屠杀了七千人，这些人并非死于战斗，而是因投降后手无寸铁，被一道命令处死；即使在战争的狂怒已经熄灭之后，和平本身的狂怒也是如此猛烈。此外，在整个城市，\Person{苏拉}的每一个党徒都随意杀人，以至于死亡人数无法计算，直到有人向\Person{苏拉}建议，应允许一些人活下来，免得统治者没有臣民。随后，这种疯狂的、不加区别的杀人许可才被遏制，当公布了\OriginalTerm{剥夺名单}{proscription}，许多人如释重负，尽管这份名单上列出了两千名最高等级的元老和骑士的死刑令。如此庞大的数目诚然令人悲哀，但令人欣慰的是，它定下了界限；死于被杀戮者数量的悲哀，也不如其余人得以保全的喜悦那么大。然而，即便是这种铁石心肠的安全感，也不能不为那些注定要死的人所遭受的极刑而悲叹。因为，有人被刽子手赤手空拳撕成碎片；人们对待活人，比野兽撕扯丢弃的尸体还要残忍。另一个人被挖掉眼睛，四肢被一块块地切割，被迫在这种折磨下活了很长一段时间，或者说，死了很长一段时间。一些著名的城市像农庄一样被拍卖；有一个城市全体居民被判处屠杀，就像一个罪犯被判处死刑一样。这些事都发生在战争结束后的和平时期，不是为了更快地获得胜利，而是为了在获得胜利后，不让人轻视它。和平与战争在残忍上竞争，并且超过了它：战争推翻武装的敌军，和平却杀戮手无寸铁的人。战争让被攻击者有自由，若有可能便反击；和平赐予幸存者的不是生命，而是无法抵抗的死亡。\par

% structure: p0082 source=h2 target=subsection reason=relative-heading-rank; base=h2

\subsection{第29章——罗马在哥特与高卢入侵时所经历的灾难与内战制造者所造成的灾难之比较。}

外族的狂怒、野蛮人的残暴，怎能与公民对公民的这场胜利相比？对于罗马而言，哪一种情况更为灾难深重、更为丑恶、更为痛苦：是晚近的哥特人与古代的高卢人入侵，还是\Person{马略}与\Person{苏拉}及其党羽对同属一个躯体的人们所施的残忍？高卢人确实屠杀了他们在城内除唯一得到防守的\Place{卡比托利山}之外任何地方发现的所有元老；但他们至少将生命卖给了在卡比托利山上的人，尽管如果不能攻克那里他们本可将那些人饿死。哥特人则饶恕了如此之多的元老，以至于他们杀掉任何一个都令人更加惊讶。但\Person{苏拉}，在\Person{马略}尚在世时，就作为征服者盘踞在连高卢人都未曾侵犯的卡比托利山上，并从那里发出死刑令；而当\Person{马略}逃跑之后——尽管他注定要以更凶猛更嗜血的面目回来——\Person{苏拉}从卡比托利山上甚至发布了元老院的法令，屠杀许多公民并没收其财产。随后，当\Person{苏拉}离开之后，\Person{马略}派的人又有什么圣物或什么人他们不曾侵犯或饶恕，他们甚至对\Person{穆基乌斯}也毫不留情，此人乃公民、元老、\OriginalTerm{大司祭}{pontifex}，当时正抱着那座据说寄寓着罗马命运的祭台发出可怜的拥抱？\Person{苏拉}那份最终的剥夺令，尚不提无数其他的屠杀，所消灭的元老比哥特人能够劫掠的还要多。\par

% structure: p0084 source=h2 target=subsection reason=relative-heading-rank; base=h2

\subsection{第30章——论基督来临之前接踵而至且极为严酷频繁的战争之间的关联。}

那么，他们是凭着何等厚颜无耻、何等自信、何等狂妄、何等愚蠢，乃至疯狂，竟拒绝将这些灾祸归于自己的神明，却将现今的归于我们的基督！这些血腥的内战，据他们自己的史家承认，比任何对外战争都更可悲，声明说那不仅是灾难性的，而且对共和国绝对具有毁灭性；它们在基督来临之前很久就已开始，且彼此催生；如此，一串非正义的因由便从\Person{马略}与\Person{苏拉}的战争，引向了\Person{塞尔托里乌斯}与\Person{喀提林}的战争——前者被公敌宣告，后者由苏拉提拔；由此又引向\Person{雷必达}与\Person{卡图卢斯}的战争，其中一人想废除\Person{苏拉}的法案，另一人则要捍卫；又由此引向\Person{庞培}与\Person{凯撒}的战争，\Person{庞培}曾是\Person{苏拉}的党徒，他在权力上与\Person{苏拉}相等甚至超越了他，而\Person{凯撒}则谴责\Person{庞培}的权力，因为那不是他自己的，但在击败并杀死\Person{庞培}之后，他却超越了这权力。内战之链从他那里延伸到第二\Person{凯撒}，即后来被称为\Person{奥古斯都}的那位，在其统治时期基督诞生。因为连\Person{奥古斯都}本人也发动了许多内战；在这些战争中，许多顶尖人物丧生，其中包括那位共和国的巧手操纵者，\Person{西塞罗}。\Person{盖乌斯〔尤利乌斯〕凯撒}在征服\Person{庞培}后，虽以宽仁运用胜利，并赐予敌方党人生命与荣誉，却因图谋王权而遭疑，在\OriginalTerm{库里亚}{curia}内被一群旨在捍卫共和国自由的贵族元老密谋刺杀。其后，他的权力为\Person{安东尼}所觊觎，此人品格迥异，染满并败坏于各等邪恶，而\Person{西塞罗}则以捍卫共和国自由的同一理由激烈反对他。就在这个节骨眼，另一位\Person{凯撒}，即\Person{盖乌斯}的养子，也就是后来如我所说以\Person{奥古斯都}之名为人所知的那位，已经作为一个天资非凡的青年\Emph{\OriginalTerm{初次登场}{début}}了。这位年轻的\Person{凯撒}受\Person{西塞罗}青睐，为的是让他的势力能抗衡\Person{安东尼}的势力；因为\Person{西塞罗}希望\Person{凯撒}能推翻并摧毁\Person{安东尼}的权力，建立一个自由的国度——他是如此盲目，对未来毫无觉察：因为其所培植其晋升与势力的那个青年，竟然任由\Person{西塞罗}被杀，以此作为与\Person{安东尼}结盟的封印，并且将共和国的自由本身——他曾为捍卫它发表过如此之多的演说——置于自己的统治之下。\par

% structure: p0086 source=h2 target=subsection reason=relative-heading-rank; base=h2

\subsection{第31章——将目前的灾难归于基督和禁止多神崇拜乃厚颜无耻之举，因为即便在众神崇拜之时，人民也遭遇了此类灾祸。}

对于那些不因\Person{基督}的巨大恩惠而感恩的人，让他们为这些沉重的灾难责备他们自己的神吧。因为当这些灾难发生时，众神的祭台确实在熊熊燃烧，那里升起了\Quote{萨巴焚香与新鲜的花环；}的混合香气，祭司们荣耀加身，庙宇金碧辉煌；殿堂中充斥着祭祀、赛会和神圣的迷狂；而同时公民的血正被如此肆意地流淌，不仅在遥远的地方，更是在众神的祭台之间。\Person{西塞罗}不愿在庙宇中寻求庇护，因为\Person{穆基乌斯}曾在那里徒劳地寻求。但是，那些最不可宽恕地诽谤这个基督宗教时代的人，正是那些要么自己逃到特别奉献给\Person{基督}的地方寻求避难，要么被蛮族带到那里才得以安全的人。简而言之，不再重述我引用的许多例证，也不再添加那些冗长枚举的例证，我坚信这一点，而任何公正的判断都会欣然承认：如果人类在布匿战争之前就接受了基督宗教，而基督宗教引入后随之而来的是这些战争带给\Place{欧洲}和\Place{阿非利加}的同样毁灭性的灾难，那么如今控告我们的人当中，没有一个不会把这些灾难归咎于我们的宗教。如果基督宗教在\Place{高卢}人入侵之前，或在毁灭\Place{罗马}的洪水和火灾之前，或在一切事件中最灾难性的内战之前，就被接受和传播，他们的指控会是多么令人难以忍受，至少对\Place{罗马}人而言！而其他那些性质如此怪异以致被算作异象的灾难，如果它们发生在基督宗教时代以后，除了基督徒，还会归罪于谁呢？我不谈那些令人惊奇更甚于伤害的事情——牛说话，未出生的婴儿在母腹中发出话语，蛇飞翔，母鸡和女人变性；以及其他类似的异象，无论真假，并非记载于他们的想象作品，而是记载于他们的历史著作中，这些并不伤害人，只使人惊异。但是，当天空降下泥土，降下白垩，降下石头——不是冰雹，而是真正的石头——这无疑会造成严重破坏。我们在他们的书中读到，\Place{埃特纳}的火焰从山顶倾泻到附近的海岸，使海水沸腾，以致岩石烧毁，船上的沥青开始融化——这种景象令人难以置信地惊异，但同时也同样有害。他们还记载，另一次，同样猛烈的热力使\Place{西西里}充满灰烬，\Place{卡塔那}城的房屋因此毁坏并被掩埋——这一灾难感动了\Place{罗马}人同情他们，并免除那年的贡赋。我们也可以读到，当时已成为\Place{罗马}行省的\Place{阿非利加}，遭受了数不胜数的蝗虫侵袭，这些蝗虫吞食了树木的果实和叶子后，被驱入海中，形成巨大无边的乌云；以致它们淹死后被冲上岸时，空气被污染，引发了严重的瘟疫，据说仅在\Person{马西尼萨}王国就有八十万人死亡，邻近地区死亡人数更多。在\Place{乌提卡}，他们确信，当时驻守的三万士兵中，仅有十人幸存。可是，假如这些灾难发生在现在，哪一件不会被那些轻率指责我们、而我们不得不回应的人归咎于基督宗教呢？然而，他们崇拜自己的神是为了避免同类的较小灾祸，却从不反思，从前崇拜那些神的人并没有被保佑免于这些严重的灾难。\par

\SourceRecord{Translated by Marcus Dods.；Nicene and Post-Nicene Fathers, First Series；Vol. 2.；Edited by Philip Schaff.；Buffalo, NY: Christian Literature Publishing Co.,；1887.；<http://www.newadvent.org/fathers/120103.htm>.}

% structure: p0001 source=h1 target=section reason=page-title

\section{\WorkTitle{天主之城}（卷四）}

在本卷中，将证明罗马帝国的幅员辽阔与长久延续，不应归因于\OriginalTerm{朱庇特}{Jove}或\OriginalTerm{异教众神}{gods of the heathen}——这些神灵各自被认为仅仅掌管少数事物，甚至是最卑微的功能——而应归于\OriginalTerm{唯一真天主}{one true God}，即\OriginalTerm{真福之源}{author of felicity}；地上的王国皆是由祂的权能与审判所建立和维持的。\par

% structure: p0003 source=h2 target=subsection reason=relative-heading-rank; base=h2

\subsection{论第一卷中所讨论之事}

既然我已开始谈论\WorkTitle{天主之城}，我认为首先必要的是答复它的仇敌；这些人热切追求尘世之乐，渴求转瞬即逝的事物，却将他们所遭遇的一切痛苦——这痛苦实是天主出自怜悯的劝诫，而非严厉的惩罚——归咎于基督宗教，即唯一带来救恩且真实的宗教。在他们中间也有无知群氓，他们在有学识之人的权威煽动下，更加痛恨我们；他们因缺乏经验而以为，自己时代发生的异常之事，在以往的世代中从未发生过；而他们的这种看法，甚至被那些明知虚假却佯装不知的人所肯定，以便显得有正当理由抱怨我们。因此，有必要从他们的作者所记录和出版的历史书籍中，证明事实远非他们所想；同时教导他们：他们曾公开崇拜、或仍在暗中崇拜的\OriginalTerm{假神}{false gods}，其实是\OriginalTerm{最不洁的邪灵}{most unclean spirits}和\OriginalTerm{最邪恶诡诈的恶魔}{most malignant and deceitful demons}，甚至到了如此地步：他们乐于在自己的节庆中，为尊崇自己而庆祝那些罪恶——无论真实还是虚构，都属于他们；因此，只要似乎有神圣的权威可供效法，人类的无能便无法被阻止去犯下可憎的罪行。这些事我们并非凭空推测，而是部分源于近代记忆，因为我们亲眼见过这样的庆典，以及那些神明；部分源于那些人的著作，他们将这些事记录下来留给后世，并非为了谴责，而是为了尊崇他们自己的神明。因此，他们中最博学、最具权威的\Person{瓦罗}，在分列\OriginalTerm{人事}{things human}和\OriginalTerm{神事}{things divine}的著作中，依据各自的尊严，将某些归于人事，某些归于神事，却将\OriginalTerm{戏剧表演}{scenic plays}全然不归入人事，而归入神事；然而，如果国家中有善良正直的人，那么即便在人事中也本不该容许戏剧表演。他这样做并非出于己见，而是因为他在\Place{罗马}出生并受教育时，发现戏剧表演已被列为神事。正如我们在第一卷末尾简要说明了此后打算讨论的内容，并在随后的两卷中处理了其中一部分，现在我们看到读者期待我们继续探讨什么了。\par

% structure: p0005 source=h2 target=subsection reason=relative-heading-rank; base=h2

\subsection{论第二、三卷中所含内容}

那么，我们曾应许，要对那些将罗马共和国的灾难归咎于我们宗教的人说些话，并且我们要列举那些灾难，就是这座城市或其帝国所属各行省，在其祭献被禁止之前所遭受的，凡是我们能记起或认为足够的、又多又大的灾难；若我们的宗教在那时就已经照耀他们，或者以这种方式禁止了他们亵渎的仪式，所有这些灾难无疑都会被归咎于我们。这些事，我们认为已在第二卷和第三卷中充分处理了，第二卷论述了道德的恶，唯有这些或主要该被算作是恶；第三卷论述了只有愚人才惧怕遭受的恶——即身体或外在事物的恶——这些恶大部分善人也会遭受。但至于那些让他们自己变恶的恶，他们不仅不忍受，反而乐于接受。就那个城市及其帝国，我所述的恶是何等少啊！甚至连\Person{凯撒·奥古斯都}时代之前的恶都未罗列完全。若我选择列举并详述那些恶，不是人彼此加害的恶，如战争的蹂躏与毁灭，而是发生在地上事物中，来自世界本身元素的恶，又会怎样？对此类恶，\Person{阿普列尤斯}在其所著的\WorkTitle{论世界}一书的一段中简略地谈到了，他说地上的一切事物都易遭受变迁、倾覆和毁灭。因为，用他自己的话说：因过度地震，大地崩裂，城市及其居民被彻底毁灭；因突降的暴雨，整片地区被冲刷而去；那些先前曾是大陆的地方，因新奇涌来的波浪而变为岛屿，而另一些地方则因海水退去而成为人的脚可行走之地；因风暴，城市被倾覆；火焰从云中闪出，致使东方的地区被燃烧而毁灭；在西方的海岸，因洪水和泛滥而导致类似的毁灭。同样，往昔从\Place{埃特纳}的高耸火山口，由天主点燃的火河如洪流般涌下陡坡。若我想从历史中，只要可能，收集这些及类似事例，那么在基督的名字尚未打倒他们那些如此虚妄且有害于真正救恩的偶像之名的时代发生的事，我何时才能述说完呢？我曾应许，我也要指出，他们习俗中哪些，以及由于什么原因，真天主——所有王国都在他权下——曾俯就施恩，使他们的帝国得以扩张；以及他们所认为的那些神不但对他们毫无助益，反而因欺骗和诱惑他们而害了他们；所以，在我看来，我现在必须谈论这些事，主要是关于罗马帝国的扩张。因为我已说了不少，尤其是在第二卷中，关于那些被他们当作神来崇拜的魔鬼，通过其有害的欺骗，将许多邪恶引入了他们的风俗。但在已经完成的三卷中，只要显得合适，我们便已阐明，天主借着基督之名——蛮族也超越战争惯例对其致以极大敬意——给予了善人和恶人多少救助，正如经上记载：\BibleQuote{因为他使太阳上升，光照恶人，也光照善人；降雨给义人，也给不义的人。} \BibleRef{玛 5:45}\par

% structure: p0007 source=h2 target=subsection reason=relative-heading-rank; base=h2

\subsection{第三章：仅凭战争取得的辽阔帝国，是否应算作智者或幸福者的好事？}

现在，让我们来看看他们怎敢将罗马帝国的广袤与长久归功于那些他们自认为通过可鄙的娱乐仪式和可鄙之人的伺候而虔诚崇拜的神明。不过，我想先稍加探询：他们有什么理由、什么智慧，竟以帝国的广大辽阔为荣，却无法指出那些经常在黑暗的恐惧和残忍的欲望中，在战争杀戮和鲜血——无论流自内战还是外战，终究是人血——中打滚的人们的幸福；因此，他们的喜乐可比作玻璃的脆弱光彩，时刻惶恐它猝然破碎。为了更清楚地洞察这一点，我们切莫因空洞的夸耀而迷失，或因听闻人民、王国、行省这些响亮的名称而削弱注意力。试假设有两个人：因为每一个人，就如语言中的一个字母，可以说是城市或王国的一个元素，不论这王国在地球上如何广泛占据。假定这两人中一个贫穷，或更恰当说，境况中等；另一个极富。但富人满怀恐惧，于不满足中憔悴，因贪婪而焦灼，从未安枕，始终不安，在仇敌的不断争斗中喘息，借着这些苦难，他的家产固然惊人地增加，但这些增加也积累起极苦的忧虑。而那一个拥有中等财富的人，则以小而紧凑的产业为满足，对家人极其亲爱，与亲族、邻人和朋友共享最甜蜜的和平，在虔诚上笃信宗教，心地仁慈，身体健康，生活节俭，品行端正，良心坦然。我不知道可有人会如此愚蠢，竟至于在选择时犹豫不决。因此，如在这两人之间，同样在两个家庭、两个民族、两个王国中，这个安宁的试金石同样有效；我们若警醒而不带偏见地运用它，便很容易看出何处仅存幸福的外表，何处寓有真正的福乐。因此，如果真天主受到崇拜，并以纯正的礼仪和真实的德行得到事奉，那么善人长治广土是有益的。这益并非对他们自己如此有益，而是对他们所统御的人有益。因为就他们自身而言，他们的虔诚和正直——这是天主的宏大恩赐——足以赐予他们真正的福乐，使他们在此生善度生活，并在来世承受永生。因此，在这个世界上，善人的统治有益于人类事务，而非有益于他们自身。但恶人的统治主要伤害他们自己，因为这些统治者放纵更大的邪恶，毁灭自己的灵魂；而受他们奴役的人，惟有因自身的罪孽才受伤害。对于义人，不义的统治者加在他们身上的一切恶事，不是罪的刑罚，而是德行的考验。因此善人虽为奴隶，却是自由的；而恶人即便为王，也仍是奴隶，不只作一个人的奴隶，而是——更为可悲的——作他所染恶习的奴隶；多少恶习，多少主子。论及这些恶习时，圣经说：\BibleQuote{因为人被谁制胜，就是谁的奴隶。}\BibleRef{伯后 2:19}\par

% structure: p0009 source=h2 target=subsection reason=relative-heading-rank; base=h2

\subsection{没有正义的王国如何与强盗相似。}

因此，一旦正义被拿走，王国岂不就是大的强盗团夥？因为强盗团夥本身，不就是小王国吗？团夥由人组成，由首领的权威统辖，因盟约协议而结合，抢夺的赃物按商定的法律分配。倘若因接纳堕落的亡命之徒，这一恶行膨胀至于占据地方、设立巢穴、占领城池、征服人民，它便更堂而皇之地取王国之名，而这并非由于贪欲被消除，而是由于不受惩罚的状态扩展了，才使其实质昭然若揭。确实，一位被擒的海盗答复\Person{亚历山大大帝}的话，既巧妙又真实。当那位国王问这人，他抱着什么意图以敌对姿态占据海面时，那人以大胆的傲气答道：\Quote{和你夺取全地一样；只不过我因为使用一艘小船，就被称为强盗，而你因为拥有一支强大的舰队，就被尊为皇帝。}\par

% structure: p0011 source=h2 target=subsection reason=relative-heading-rank; base=h2

\subsection{论逃亡角斗士的权势如何与王权无异。}

因此，我不打算停下来探究\Person{罗穆卢斯}聚集了什么样的人，既然他对他们多有考量——如何通过将他们从原先的生活带入他的城邦共同体，使他们不再思想应得的惩罚，对惩罚的恐惧曾驱使他们做出更大的恶事；从而从此以后他们可以成为社会中更和平的成员。但我要说：罗马帝国通过征服许多民族已经变得强大，成为普遍畏惧的对象，然而它自己也大为惊恐，并费了很大力气才避免了一场灾难性的倾覆，因为仅仅是一小撮\Place{坎帕尼亚}的\OriginalTerm{角斗士}{gladiators}，从竞技中逃脱，便招募了一支大军，任命了三位将军，并且极为广泛而残酷地蹂躏了\Place{意大利}。让他们说吧，是哪位神帮助了这些人，以至于他们从一伙微不足道、受人蔑视的强盗发展成了一个王国，甚至为拥有如此大军和堡垒的罗马人所畏惧。或者他们是否因为这些人没有长久持续而否认有神的帮助？就好像任何人的生命能够长久似的。在这种情况下，神也不帮助任何人掌权，因为人人都很快死去；统治权也不能被视为一种益处，因为它在每个人身上，从而在所有一个个的人身上，都像蒸气一样短暂消散。因为对那些在\Person{罗穆卢斯}治下崇拜诸神而早已死去的人而言，罗马帝国在他们死后变得如此强大又有何意义？他们正在地下的权能面前为他们的案件辩护。无论那些案件是好是坏，都无关我们当前的问题。所有那些在短暂的生命中匆匆掠过帝国官职舞台的人都是这样，尽管官职本身历经漫长岁月，由不断死去的继任者填补。然而，如果连那些只持续最短时间的益处都应归功于神的帮助，那么这些角斗士的确得到了不小的帮助，他们挣脱了奴役的锁链，逃走并召集了一支庞大而极强大的军队，服从首领的意志与命令，令罗马威严甚为恐惧，他们未被数位罗马将军征服，夺取了许多地方，并赢得了许多胜利，享受他们想要的任何快乐，做他们欲望所驱之事，直到最终以极大困难被征服为止，生活得崇高而掌权。但让我们来看更大的事件。\par

% structure: p0013 source=h2 target=subsection reason=relative-heading-rank; base=h2

\subsection{第6章 论\Person{尼诺斯}的贪婪——他第一个对邻国发动战争以更广泛地统治}

\Person{优斯提努斯}用拉丁文撰写了希腊或更确切地说是外邦的历史，他像他所追随的\Person{特洛古斯·庞培}那样简明扼要，其著作开头如此写道：\Quote{在万民与各族事务的初期，政府掌握在国王手中，他们被提升到这崇高地位，不是靠讨好民众，而是因善人知道他们的节制。人民不受任何法律的约束；君王的决定代替了法律。习惯是守护而非扩张帝国的边界；王国被限制在各统治者故土的范围内。\Place{亚述}人的国王\Person{尼诺斯}，首先通过新的\OriginalTerm{征服欲}{new lust of empire}，改变了各民族古老且如同祖传的习俗。他首先对邻国发动战争，并完全征服了直到\Place{利比亚}边界的各民族，当时他们还不谙抵抗。}稍后他又说：\Quote{\Person{尼诺斯}通过持续占有巩固了他所获权威的伟大。他制服了最近的邻邦，便以增强的兵力继续攻击其他民族，并使每一次新的胜利成为下一次胜利的工具，征服了整个\Place{东方}的民族。}现在，不论他或\Person{特洛古斯}一般在多大程度上忠于事实——因为其他更可靠的作者表明他们有时说谎——但在其他作家中一致认为，\Place{亚述}王国在\Person{尼诺斯}王的扩张下幅员广阔。而且它持续了如此之久，以至于罗马帝国尚未达到同样的年龄；因为正如那些编写编年史的人所写，这个王国从\Person{尼诺斯}开始统治的第一年起，延续了一千二百四十年，直到它被转移给\Place{米底}人。但是，对邻国开战，由此进而攻击他国，并仅仅出于\OriginalTerm{统治欲}{lust of dominion}去粉碎和征服那些未曾加害于你的人，这该称作什么，如果不是\OriginalTerm{大盗}{great robbery}呢？\par

% structure: p0015 source=h2 target=subsection reason=relative-heading-rank; base=h2

\subsection{第7章 尘世王国在其兴衰中是否得到诸神的帮助或遗弃}

如果这个王国在没有神明的帮助之下就如此伟大而长久，为什么要把\Place{罗马}帝国的广阔疆土和漫长延续归因于\Place{罗马}神明呢？因为不论造成前者那种情况的原因是什么，后者也是同样原因。但如果他们争辩说，其他王国的繁荣也应归因于神明的帮助，我要问，是哪些神明呢？因为\Person{尼努斯}所征服的其他民族，那时并不崇拜别的神明。或者，如果\Place{亚述}人拥有自己的神明，他们可以说是更熟练的工匠，在建立和维持帝国方面更为高手——那么，这些神明是否已经死了，因为他们自己也失去了帝国呢？或者，他们是否因为被欺骗了报酬，或者被许以更大的报酬，而宁愿转而投靠\Place{玛待}人，随后又从\Place{玛待}人投靠\Place{波斯}人呢？因为\Person{居鲁士}邀请了它们，并许给它们更有利的东西？的确，这个民族，自从\Place{马其顿}的\Person{亚历山大}的王国以来——这个王国时间虽短，疆域却广——一直保有自己的帝国，直到今天在\Place{东方}还占据着不小的领土。如果是这样，那么这些神明要么是不忠的，抛弃自己的属民而投靠敌人——这是连凡人\Person{卡米卢斯}也不曾做的：他作为胜利者和最敌对国家的征服者，尽管感到\Place{罗马}，那个他为之付出甚多的\Place{罗马}，忘恩负义，却后来忘记伤害而记念祖国，再次从\Place{高卢}人手中解放了它；要么它们就不像神明所应有的那样强大，因为它们能被人的技能或力量所胜。或者，如果神明之间互相争战，它们不是被人所胜，而是某些属于特定城邦的神明被其他神明所胜，那么它们之间就有纷争，各为自己的部分而坚持。因此，一个城邦不该崇拜自己的神明，而宁可崇拜那些帮助自己崇拜者的其他神明。最后，不论神明这一方在改换阵营、逃跑、迁移或战斗失败方面的情况如何，当这些王国因战争的巨大破坏而丧失和转移时，\Person{基督}之名尚未在这些土地上被宣扬。因为，如果在超过一千二百年之后，当王国从\Place{亚述}人手中被夺走时，基督宗教就已经在那里宣讲另一个永恒王国，并制止了对虚假神明的亵渎崇拜，那么那个民族的愚人除了会说，这个保存了如此之久的王国之所以丧失，没有别的原因，只是因为它们抛弃了自己的宗教而接受了基督宗教之外，还能说什么呢？让那些我们所说的、可能曾发出这样愚昧言语的人，看看他们自己类似的言论，并感到脸红吧——如果他们还有任何羞耻感的话——因为他们也发出了类似的抱怨；尽管\Place{罗马}帝国现在是遭受苦难而非转变——这在\Person{基督}之名被听闻之前的其他时代也曾遭遇过，并在遭受苦难之后得以恢复——即使在今日，也不应对此绝望。因为谁知晓天主在这件事上的旨意呢？\par

% structure: p0017 source=h2 target=subsection reason=relative-heading-rank; base=h2

\subsection{\Place{罗马}人能认为哪位神明掌管他们帝国的增长与维持呢？因为他们相信，即使每一项小事都几乎不能托付给单独一位神明。}

接下来，如果诸位愿意，让我们询问，在罗马人所崇拜的如此众多的神明中，他们特别相信哪一位或哪些神明扩展并保守了这个帝国。当然，对于这项如此卓越、充满至高尊严的事业，他们绝不敢将任何部分归于女神\Person{克洛阿西娜}；或归于\Person{沃卢皮娅}，她的名字来自淫乐；或归于\Person{利本蒂娜}，她的名字来自情欲；或归于\Person{瓦提卡努斯}，他掌管婴儿的叫喊；或归于\Person{库尼娜}，她掌管婴儿的摇篮。然而，怎么可能在本书的这一部分中历数所有男神和女神的名号呢？他们几乎无法将这些名号容纳于大部头的书卷中，并给这些神祇分配各自对单一事物的特殊职能。他们甚至没有想过将土地的管理托付给任何一位神明：而是将农田托付给\Person{鲁西娜}；将山脊托付给\Person{尤加提努斯}；他们设立女神\Person{科拉蒂娜}管辖丘陵；设立\Person{瓦洛尼亚}管辖山谷。他们甚至找不到一位如此胜任的\Person{塞格提亚}，可以让她同时照管他们所有的谷物作物；而是在种子谷粒仍在地下时，让女神\Person{塞娅}掌管它；然后，当它长出地面并形成秸秆时，他们让女神\Person{塞格提亚}掌管它；当谷物被收获并储存后，他们让女神\Person{图提利娜}掌管它，以便安全保存。谁不认为女神\Person{塞格提亚}足以照料立着的庄稼，直到它从最初的绿叶变为干穗呢？然而，对于热爱众多神明的人们来说，她仍不足够，为的是让可怜的灵魂，轻看唯一真天主的贞洁拥抱，而向一群魔鬼行淫。因此，他们让\Person{普罗塞皮娜}掌管发芽的种子；让神\Person{诺多图斯}掌管茎秆的节与瘤；让女神\Person{沃隆蒂娜}掌管包裹麦穗的叶鞘；当叶鞘张开使穗子长出时，这被归于女神\Person{帕特拉娜}；当茎秆与新穗齐平时，因为古人用术语\Emph{\OriginalTerm{齐平}{hostire}}描述这一均等，这被归于女神\Person{霍斯提利娜}；当谷粒开花时，献给女神\Person{芙罗拉}；当充满浆液时，献给神\Person{拉克图努斯}；当成熟时，献给女神\Person{玛图塔}；当庄稼被连根拔起（\OriginalTerm{连根拔起}{runcated}——也就是从土中移除时——献给女神\Person{伦西娜}。我还没有把它们全部列举出来，因为我已对此感到厌烦，尽管他们并不感到羞耻。我只说了这很少的一些，为的是让人明白：他们绝不敢说是他们的神祇建立、扩展并保守了罗马帝国，因为这些神祇各自的职能被分配得如此之细，以至于没有一位被赋予全面的监管。因此，\Person{塞格提亚}何时能照料帝国呢？她连谷物和树木都不被允许照料。\Person{库尼娜}何时能思考战争呢？她的监管范围被限制在婴儿的摇篮之内。\Person{诺多图斯}何时能在战斗中提供帮助呢？他甚至连麦穗的叶鞘也管不着，只管茎秆的节瘤。每个人都为自己的家门设立一个看门人，因为他是人，就足够胜任；但这些人却设立了三位神：\Person{福尔库鲁斯}管门，\Person{卡耳德亚}管铰链，\Person{利门提努斯}管门槛。因此，\Person{福尔库鲁斯}无法同时照管铰链和门槛。\par

% structure: p0019 source=h2 target=subsection reason=relative-heading-rank; base=h2

\subsection{是否应将罗马帝国的辽阔疆域和长久存续归功于\Person{朱庇特}，其崇拜者相信他是主神。}

因此，暂且忽略或绕过那群渺小的神明，我们应当探究那些大神所起的作用，正是通过他们，罗马变得如此伟大，以至统治如此众多的民族如此之久。毫无疑问，这是\Person{朱庇特}的作为。因为他们主张他是所有男神和女神之王，正如他的权杖和高高山丘上的\Place{卡皮托利山}所表明的那样。关于那位神，他们传诵着一句最贴切的话语，尽管出自诗人之口：\Quote{万物皆充满\Person{朱庇特}。}\Person{瓦罗}相信，这位神虽然被以别的名字称呼，却也受到那些只崇拜唯一、无偶像的天主的人的崇拜。但如果是这样，为什么在罗马（以及其他民族中）他受到如此糟糕的对待，竟为他造了像呢？——这件事令\Person{瓦罗}本人非常不悦，以至于尽管他被如此大城的邪恶习俗所压倒，却毫不犹豫地既说又写，那些为人民设立偶像的人，既夺去了敬畏，又加添了错误。\par

% structure: p0021 source=h2 target=subsection reason=relative-heading-rank; base=h2

\subsection{那些为世界不同部分设立不同神明的人遵循了怎样的意见。}

为什么\OriginalTerm{朱诺}{Juno}又与他结合为妻，她既被称为\Quote{姐姐和同轭者}？他们说，因为我们在\OriginalTerm{以太}{ether}中有\OriginalTerm{朱庇特}{Jove}，在空气中有\OriginalTerm{朱诺}{Juno}；这两种元素结合在一起，一个在上，一个在下。那么，如果\OriginalTerm{朱诺}{Juno}也充满一部分，他就不是\Quote{万物充满了\OriginalTerm{朱庇特}{Jove}}所说的那一位了。难道他们各充满一个，这对夫妇的两个神都在两种元素中，并且同时都在每一种元素中吗？那么，为什么\OriginalTerm{以太}{ether}被分配给\OriginalTerm{朱庇特}{Jove}，空气被分配给\OriginalTerm{朱诺}{Juno}呢？此外，这两个本来应该足够了。为什么海洋被分配给\OriginalTerm{尼普顿}{Neptune}，大地被分配给\OriginalTerm{普路托}{Pluto}？而且，为了使这二位也不缺少配偶，\OriginalTerm{萨拉西亚}{Salacia}被配给\OriginalTerm{尼普顿}{Neptune}，\OriginalTerm{普洛塞庇娜}{Proserpine}被配给\OriginalTerm{普路托}{Pluto}。因为他们说，正如\OriginalTerm{朱诺}{Juno}拥有天界较低的部分（即空气），\OriginalTerm{萨拉西亚}{Salacia}拥有海洋较低的部分，\OriginalTerm{普洛塞庇娜}{Proserpine}拥有大地的较低部分。他们试图拼凑这些寓言，却找不到出路。因为如果真是这样，他们古代的圣贤就会主张世界的基本元素是三个，而不是四个，以便每种元素都有一对神明。现在，他们明确肯定\OriginalTerm{以太}{ether}是一回事，空气是另一回事。然而水，无论是上层的还是下层的水，终究是水。不管它多么不同，难道能不同到不再是水吗？至于下层的土地，无论被冠以何种神名，它又能是别的什么，不就是土地吗？看哪，既然整个物质世界由这四种或三种元素构成，那么\OriginalTerm{密涅瓦}{Minerva}该在哪里呢？她该拥有什么，充满什么呢？因为她与这两位一同被安置在\Place{卡比托利山}上，尽管她并不是他们婚姻的后代。或者，如果他们说她拥有\OriginalTerm{以太}{ether}的上层部分——因此诗人们虚构说她从\OriginalTerm{朱庇特}{Jove}的头颅中出生——那么她为什么不被视为众神的女王，因为她比\OriginalTerm{朱庇特}{Jove}更高呢？难道是因为把女儿置于父亲之前不合宜？那么，为什么对\OriginalTerm{朱庇特}{Jove}自己与\OriginalTerm{萨图恩}{Saturn}之间不遵守那条正义的原则呢？难道是因为\OriginalTerm{萨图恩}{Saturn}被征服了吗？他们打过仗吗？他们说，绝非如此；那是老妇人的寓言。看，我们不应相信寓言，而必须对众神抱有更可敬的看法！那么，他们为什么不分给\OriginalTerm{朱庇特}{Jove}的父亲一个位置，即便不是更高，至少也应是同等荣耀的位置呢？他们说，因为\OriginalTerm{萨图恩}{Saturn}是时间的长度。所以崇拜\OriginalTerm{萨图恩}{Saturn}的人是在崇拜时间；并且暗示众神之王\OriginalTerm{朱庇特}{Jove}是从时间中诞生的。因为如果说\OriginalTerm{朱庇特}{Jove}和\OriginalTerm{朱诺}{Juno}是从时间中诞生的，何有可耻之处呢？如果他是天，她是地，因为天和地都是被造的，因此不是永恒的。因为他们的博学智者也在他们的书中这么写。维吉尔的那句话也不是取自诗人的虚构，而是取自哲学家的书，\par

\Quote{于是全能之父\OriginalTerm{以太}{Ether}，以丰沛的雨水降落\newline{}进入他配偶欢悦的怀中，使之肥沃，}\par

——也就是说，进入\OriginalTerm{忒勒斯}{Tellus}或大地的怀中。然而，在此处他们又认为存在一些区别，并且认为在土地本身之中，\OriginalTerm{泰拉}{Terra}是一回事，\OriginalTerm{忒勒斯}{Tellus}是另一回事，\OriginalTerm{忒卢摩}{Tellumo}又是另一回事。他们把这些都当作神明，各以自己的名字称呼，各以独特的职司区分，各以自己的祭坛和仪式受崇拜。这同一片土地，他们还称为诸神之母，以至于根据他们的圣书而非诗书，诗人虚构反更可容忍；\OriginalTerm{朱诺}{Juno}不仅是\OriginalTerm{朱庇特}{Jove}的姐姐和妻子，还是他的母亲。他们以\OriginalTerm{刻瑞斯}{Ceres}之名崇拜同一片土地，也以\OriginalTerm{维斯塔}{Vesta}之名崇拜它；然而他们又更经常地断言\OriginalTerm{维斯塔}{Vesta}无非是火，属于炉灶，没有它城市便不能存在；因此通常让处女侍奉她，因为正如处女不能生产，火也不产任何东西——但这一切无稽之谈都应被那位由童贞女所生者彻底废除和扑灭。因为谁能容忍，他们给予火如此多的尊荣乃至贞洁，却有时甚至把\OriginalTerm{维斯塔}{Vesta}称为\OriginalTerm{维纳斯}{Venus}也不脸红，以至于尊荣的贞洁在她的女侍中消失？因为如果\OriginalTerm{维斯塔}{Vesta}就是\OriginalTerm{维纳斯}{Venus}，那么处女通过禁绝淫欲来恰当地侍奉她，如何可行？难道有两个\OriginalTerm{维纳斯}{Venus}，一个是处女，另一个不是少女？或者更确切地说，有三个\OriginalTerm{维纳斯}{Venus}：一个是童贞女神，也称作\OriginalTerm{维斯塔}{Vesta}；一个是人妻女神；还有一个是娼妓女神？\OriginalTerm{腓尼基人}{Phenicians}也曾向后者献祭，让他们女儿在出嫁前卖淫。其中哪一位是\OriginalTerm{伏尔甘}{Vulcan}的妻子呢？肯定不是那位处女，因为她有丈夫。但愿不是那娼妓，以免我们看似冤枉了\OriginalTerm{朱诺}{Juno}的儿子和\OriginalTerm{密涅瓦}{Minerva}的同工。因此应当理解她属于已婚者；但我们不愿他们效法她与\OriginalTerm{玛尔斯}{Mars}所作的事。他们说：\Quote{你又在扯回寓言了。}这算哪门子正义？——对我们诉说他们神明的这类事便发怒，而对他们自己在剧场里极乐意观看他们神明的罪行却不发怒？而且——若非有充分证据简直难以置信——这些表现他们神明罪行的戏剧表演，正是为了尊崇同样这些神明而设立的。\par

% structure: p0025 source=h2 target=subsection reason=relative-heading-rank; base=h2

\subsection{论异教博士们所辩护的众多神明其实为同一位\OriginalTerm{朱庇特}{Jove}}

因此，让他们在自然哲学的推理和争论中随心所欲地断言多少都行。时而让\OriginalTerm{朱庇特}{Jupiter}成为这有形世界的灵魂，充满并推动那由四种或他们乐意的任意多种元素构成并紧密组合的整体；时而让他把其中各部分让与他的姐姐和兄弟们：时而让他成为\OriginalTerm{以太}{ether}，从上方拥抱伸展在下的空气\OriginalTerm{朱诺}{Juno}；又有时，让他成为整个天穹连同空气，以雨水和种子使大地肥沃，视之为其妻子，同时也为其母亲（因为在神圣存在中这并不卑鄙）；然后再一次（免得必须全数罗列），就让这位独一的神，很多人认为最卓越的诗人曾如此描述过他，\par

\Quote{因为天主遍及万物，\newline{}遍及所有陆地，海域，以及诸天的深处，}\par

——让他在以太中是\Person{朱庇特}；在空气中是\Person{朱诺}；在海中是\Person{尼普顿}；在海的深处是\Person{萨拉西亚}；在陆地中是\Person{普路托}；在地的深处是\Person{普洛塞庇娜}；在家中的灶台上是\Person{维斯塔}；在工匠的熔炉中是\Person{伏尔甘}；在星辰间是\Person{索尔}、\Person{卢娜}和众星；在预言中是\Person{阿波罗}；在商业中是\Person{墨丘利}；在作为起始者的\Person{雅努斯}中；在作为终结者的\Person{特米努斯}中；在时间中是\Person{萨图恩}；在战争中是\Person{玛尔斯}与\Person{贝罗娜}；在葡萄园中是\Person{利柏耳}；在麦田中是\Person{刻瑞斯}；在森林中是\Person{狄安娜}；在学术中是\Person{密涅瓦}。最后，让他在那群所谓平民神祇中也是他：让他在\Person{利柏耳}的名义下掌管男子的种子，在\Person{利贝拉}的名义下掌管女人的种子：让他成为\Person{迪斯皮特}，那位将新生儿带到日光之下的神；让他成为\Person{梅娜}女神，她被指派管理女人的经期；让他成为\Person{卢基娜}，她是妇女临产时所呼求的女神；让他帮助那些正在诞生的人，把他们从大地的怀抱中抱起，故被称为\Person{奥普斯}；让他使哭泣的婴儿张开嘴，故被称为\Person{瓦提卡努斯}神；让他把婴儿从地上抬起，故被称为\Person{勒瓦娜}女神；让他守护摇篮，故被称为\Person{库尼娜}女神：让他不是别的，正是那些歌唱新生儿命运的、被称为\Person{卡耳门忒斯}的女神们中的那位：让他掌管偶然事件，故被称为\Person{福尔图娜}；在\Person{鲁弥娜}女神中，让他给婴儿哺乳，因为古人称乳房为\OriginalTerm{乳房}{ruma}；在\Person{波提娜}女神中，让他管理饮水；在\Person{埃杜卡}女神中，让他供应食物：因婴儿的恐惧，称他为\Person{帕文提亚}；因来临的盼望，\Person{维尼利亚}；因快感，\Person{沃卢皮亚}；因行动，\Person{阿革诺耳}；因那些刺激人投入许多行动的激发物，他被称为\Person{斯提穆拉}女神：让他成为\Person{斯特瑞尼亚}女神，因为使人奋发；\Person{努梅里亚}，教人计数；\Person{卡墨娜}，教人歌唱：让他既是为提供商议的\Person{孔苏斯}神，也是为启发思考的\Person{森提亚}女神：让他成为\Person{尤文塔斯}女神，她在脱下童年长袍后，掌管青年时期的开端：让他成为\Person{福尔图娜·巴尔巴塔}，她赋予成年人胡须，而他们并不愿尊崇这位神；以致这位神，无论是什么，至少应是一位男性神，或者因\OriginalTerm{胡子}{barba}而称\Person{巴尔巴图斯}，如同因\OriginalTerm{节}{nodus}而称\Person{诺多图斯}；或者，当然不叫\Person{福尔图娜}，但因他有胡须，可称\Person{福尔图尼乌斯}：让他，在\Person{尤加提努斯}神中，使婚姻结合；当童贞新娘的腰带解开时，他被呼求为\Person{维尔吉尼恩西斯}女神：让他成为\Person{穆图努斯}或\Person{图特努斯}，即在希腊人中被称为\Person{普里阿普斯}的神。如果他们对此不感羞耻，那么所有这些我列举的神，以及我未列举的其它神（因为我不认为有必要全列），让所有这些男神和女神都成为那独一的\Person{朱庇特}，无论是如某些人所愿，所有这些神都是他的部分，或是他的能力——正如那些乐于认为他是世界灵魂的人所想的那样，这也是他们那些最伟大的博士们中多数人的意见。如果事果然如此（我暂且不探究这会多么邪恶），那么倘若他们用一个更精要的简洁方式去崇拜独一神，又会损失什么呢？因为他若自己受崇拜，他的哪一部分会被藐视呢？但如果他们担心他的各部分因被忽略或轻看而发怒，那么事情就并非如他们所愿，即这整体仿佛一个活物的生命，把众神一起包含在其中，好像它们是它的能力、肢体或部分；反之，如果一部分比另一部分更容易被触怒、平息或激发，那么每一部分都各有自己与其它部分分离的生命。但如果有人说，全部合起来——即整个\Person{朱庇特}本身——若他的部分不逐一而细致地受崇拜，他也会被冒犯，这种说法是愚蠢的。当然，若那位单独拥有它们全部的神受崇拜，就没有任何部分会被忽略。因为，且不提其它无数的事，当他们说所有星辰都是\Person{朱庇特}的部分，且都有生命，有理性灵魂，因而毫无疑问是神时，难道他们看不见有多少星辰他们并不崇拜，不为多少星辰建庙或设祭坛，而实际上他们只为其中极少部分立坛献祭吗？因此，如果那些未个别受崇拜的部分感到不悦，难道他们不怕只在安抚少数几颗时，整个上天都已被激怒吗？然而，如果他们因为崇拜\Person{朱庇特}而崇拜他身上的所有星辰，用同样的简便方式，他们就可以单单在他里面向所有星辰祈求。因为这样一来，没有人会不悦，因为在他独一里面，全部都被祈求了。没有人会被藐视，而在只向某些神献上崇拜时，反倒有极大理由冒犯那被逾越的更大数目；尤其当丑陋赤裸地伸展的\Person{普里阿普斯}，比起那些从高处居所闪耀的星辰更受偏爱的时候。\par

% structure: p0029 source=h2 target=subsection reason=relative-heading-rank; base=h2

\subsection{论那些认为天主是世界灵魂、世界是天主身体的人的意见。}

有理性的、乃至任何种类的人，难道不该被激励去省察此意见的本质吗？因为并不需要出色才干即可——放下争竞的欲望，便能观察：如果天主是世界灵魂，而世界对祂——那位灵魂——而言如同身体，祂必是一个由灵魂和身体组成的活物，并且同一个天主管如一个自然的子宫，将万物包含在自己之内；这样，一切生物的生命和灵魂，都按各自诞生的方式，从祂那赋予整个团块生命的灵魂中获取，所以没有任何事物不是天主的一部分。如果这果真如此，谁看不出随之而来的是何等不虔敬与不敬畏的后果？比如，任何人无论践踏什么，他必践踏了天主的一部分；在杀死任何活物时，天主的一部分必遭屠戮。然而，我不愿说出所有可能被思想者想到、却又无法不带着不敬而言之的事情。\par

% structure: p0031 source=h2 target=subsection reason=relative-heading-rank; base=h2

\subsection{论那些断言只有理性动物才是独一天主之部分的人。}

但若他们争辩说，只有理性的动物，如人，才是天主的部分，若整个世界即是天主，我实在看不出他们如何能将禽兽排除在天主的部分之外。然而，何须为此争执？就理性的动物本身——即人类而言——还有什么比一个男孩受鞭笞时，竟以为天主的一部分也在受鞭笞，更可悲的信念呢？而且，除非彻底癫狂，谁能容忍这样的想法：天主的部分竟能变得好色、不义、不虔敬，且全然该被诅咒？简言之，若那些冒犯者正是祂自身的一部分，天主为何又要向不敬拜祂的人发怒呢？因此，他们不得不说，所有神明各有其生命，各自为己而生，没有一个神明是任何其他神明的一部分；然而，所有神明都应受敬拜——至少是所有能被认识与敬拜的；因为神明如此之多，全数敬拜实不可能。而在这一切之中，我认为，既然\Person{朱庇特}以君王之尊统御，他们便以为罗马帝国是由他建立并扩张的。因为若此事并非他成就，他们又能相信哪一位神明堪当如此大业呢？既然众神都须各司其职，各务其事，且不可相互侵扰。如此说来，人间的国度能否藉着众神之王得以扩展与增长呢？\par

% structure: p0033 source=h2 target=subsection reason=relative-heading-rank; base=h2

\subsection{国度的扩张不宜归因于\Person{朱庇特}；因为，若如他们所言，\Person{维多利亚}是一位女神，则她独自便足以成就此事。}

在此，我首先发问：为何连王国本身也不能是一位神明呢？若\Person{维多利亚}是一位女神，王国为何不能同样是呢？再者，若\Person{维多利亚}女神眷顾施恩，且总与那些她愿其得胜的人同在，此事又与\Person{朱庇特}何干呢？若有这位女神眷顾施恩，纵使\Person{朱庇特}闲散无事，哪个邦族能不被征服，哪个王国能不降服呢？然而，或许善人厌恶以极不义之恶作战，借主动挑衅和平无欺的邻邦来扩张王国？若他们确有此感，我完全赞同并称许他们。\par

% structure: p0035 source=h2 target=subsection reason=relative-heading-rank; base=h2

\subsection{善人是否适宜渴望更广阔的统治。}

那么，让他们审问：善人若因帝国扩张而欢欣，是否全然合宜呢？因为，正是那些与之为敌者的不义，助长了王国的壮大；倘若邻邦的和平与正义从未因任何错谬而招致战争，王国本会始终狭小。人事若更幸福如此，所有王国都将保持弱小，并享受睦邻之乐；那样，世上就会有极多的民族王国，就如同一座城中有极多的居民之家。因此，对恶人而言，发动战争、扩张王国以征服诸族，似是福乐，对善人而言，却是不得已之事。但由于让不义者统治较正义者更为恶劣，故那也被称为福乐，也非全无道理。但无疑，拥有一个和睦相处的善邻，远胜于征服一个需要通过战争解决的恶邻。你若期望那为你所恨或所惧之人落入可被你战胜的境地，这便是恶愿。因此，若罗马人通过正义、而非不虔不义的战争获得了如此庞大的帝国，他们岂不也应将那外邦的不义当作女神来敬拜吗？因为我们看到，这外邦不义对帝国的扩张贡献甚大，它使得外邦人变得如此不义，以至于他们成了可以发动正义战争的对象，帝国从而得以扩张。再者，若畏惧、惊惧与疟疾都配成为罗马之神，为何不义——至少是外邦的不义——不能也成一位女神呢？因此，通过这两者——即外邦的不义与\Person{维多利亚}女神，因为不义煽动战争之因，而\Person{维多利亚}则使战争完满结束——帝国得以扩张，纵使\Person{朱庇特}袖手旁观。\Person{朱庇特}在此又有何能参与呢？既然那些或许本可视为其恩赐的事物，却都被视为神明，被称为神，被当作神来敬拜，并为各自职分而受祈求。倘若他本身可被称为\Quote{帝国}，正如她被称为\Quote{胜利}一般，他或许也能在此事上有一席之地。或者，若帝国是\Person{朱庇特}的恩赐，\Person{维多利亚}不亦可视为他的恩赐吗？事实上，倘若他真的是那位万王之王、万主之主\BibleRef{默 19:16}，而非\Place{卡比托利山}上的那块石头，那么\Person{维多利亚}定会被如此看待。\par

% structure: p0037 source=h2 target=subsection reason=relative-heading-rank; base=h2

\subsection{罗马人既为万事万物及心灵一切活动分设众神，却选择将静息神庙建于城门之外，其原因何在。}

然而我极感诧异的是：他们将各项事物，以及（几乎）所有心灵的动态，都分派给不同的神明；他们呼求行动女神\OriginalTerm{阿革诺里亚}{Agenoria}，以激励人去行动；呼求刺激女神\OriginalTerm{斯提穆拉}{Stimula}，以煽动异常行为；呼求迟缓女神\OriginalTerm{穆尔基亚}{Murcia}，她不会使人过度亢奋，反而如\Person{庞波尼乌斯}所言，让人变得\OriginalTerm{迟钝的}{murcid}——即过于懒散怠惰；又呼求奋力女神\OriginalTerm{斯特瑞努亚}{Strenua}，以使人奋发努力；他们向这所有的男女神明献上庄严公开的敬礼，却不情愿公开承认那位他们称为\OriginalTerm{静息女神奎厄斯}{Quies}的女神——因为她使人安静——反而将她的庙宇建在\Place{科林城门}之外。这究竟是一颗不安之心的征兆，还是意在暗示：凡执意崇拜那一大群——当然，不是神明，而是魔鬼——的人，无法安居于静息之中？那位真正的医生召唤说：\BibleQuote{你们跟我学吧！因为我是良善心谦的：这样你们必要找得你们灵魂的安息。}\par

% structure: p0039 source=h2 target=subsection reason=relative-heading-rank; base=h2

\subsection{若至高权力属\Person{朱庇特}，\Person{维多利亚}是否亦当受崇拜？}

或者，他们也许会说，\Person{朱庇特}派遣女神\Person{维多利亚}，而她似乎服从众神之王的命令，来到他差遣她前往的人那里，并在他们那一方驻扎下来？这话说得确实不错，但不是对他们凭自己想象虚构为众神之王的\Person{朱庇特}说的，而是针对那位真正永恒的君王；因为他派遣的，不是\Person{维多利亚}（她并非一位位格者），而是他的天使，使他愿意的人得胜；他的谋略可能深藏不露，但绝不会不义。因为如果\Person{维多利亚}是女神，那么凯旋为什么不也是神，并作为丈夫、兄弟或儿子与\Person{维多利亚}相结合呢？事实上，他们关于诸神臆想了如此多的事，如果诗人们也编造了类似的故事，而我们加以讨论，他们必会回答说，那不过是诗人们可笑的虚构，不应归于真实的神明。然而，他们自己却并不觉得可笑，当他们不是在读诗人的作品，而是在庙宇中敬拜如此愚蠢的妄想时。因此，他们应当只恳求\Person{朱庇特}一人，只向他祈求一切。因为如果\Person{维多利亚}是女神，且服从作为她君王的\Person{朱庇特}，那么无论他差遣她去哪里，她都不敢抗拒而按自己的意愿行事，而不是按他的意愿。\par

% structure: p0041 source=h2 target=subsection reason=relative-heading-rank; base=h2

\subsection{第十八章 —— 那些以\OriginalTerm{幸福女神}{Felicity}与\OriginalTerm{命运女神}{Fortune}为女神者，凭什么理由区分了她们。}

此外，关于\OriginalTerm{幸福女神}{Felicity}也是一位女神的看法，我们还能说什么呢？她已获得一座庙宇；她配得一座祭台；人向她献上合宜的敬拜礼仪。那么，只应当敬拜她一位。因为哪里有她临在，什么美善能缺席呢？然而，一个人想要什么，竟以为\OriginalTerm{命运女神}{Fortune}也是一位女神，并敬拜她呢？难道幸福是一回事，命运是另一回事吗？命运的确可能既坏又好；但幸福若可能是坏的，就不再是幸福了。我们当然应当认为，所有性别的神（若他们也有性别）都仅是善的。\Person{柏拉图}如此说；其他哲学家如此说；所有可敬的共和国与民族的统治者也都如此说。那么，命运女神怎么会时善时恶呢？是否当她为恶时，她就不是女神，而是突然间变成了一位恶毒的恶魔？那样的话，有多少命运女神呢？有多少幸运的人，即享有好运的人，就有多少命运女神。然而，必然也有许多同时不幸的人，那么同一个命运女神，怎么可能同时既恶又善——对这些人恶，对那些人善呢？那作为女神的她，是否总是善的呢？那么她就是幸福本身了。既然如此，为何又给她两个名字呢？这倒还可以容忍；因为一件事物有两个名字，本是寻常的事。但为何有不同的庙宇、不同的祭台、不同的礼仪呢？他们说，是有理由的：因为幸福女神是善人凭着先前的功劳而获得的；而命运女神，所谓的好命运，无需任何功德的考验，就偶然地临到善人和恶人，因此她也得名于\OriginalTerm{幸运}{fortune}。那么，她既无分别地临到善人与恶人，又如何是善的呢？她如此盲目，随意奔向任何人，以致多半越过敬拜她的人，反而依附那些轻慢她的人，那么为何还要敬拜她呢？或者，倘若敬拜她的人能得些好处，被她看见并喜爱，那么她就追随功德，而非偶然来临了。这样一来，命运的定义何在？那认为她从偶然事件得名的看法又如何自圆其说？因为如果她真是命运，那么敬拜她便毫无益处。但如果她分辨她的敬拜者，以便施恩于他们，她就不是命运了。或者，难道\Person{朱庇特}也随己意差遣她吗？那么，就只应敬拜\Person{朱庇特}一人；因为命运不能抗拒他的命令，当他吩咐她、随意差遣她的时候。或者，至少让那些不愿有功德以邀来幸福女神的恶人去敬拜她吧。\par

% structure: p0043 source=h2 target=subsection reason=relative-heading-rank; base=h2

\subsection{第十九章 —— 论\OriginalTerm{妇女命运女神}{Fortuna Muliebris}}

对于他们称之为\OriginalTerm{福图娜}{Fortuna}的这位假想神明，他们竟赋予她如此大的能力，以致流传下这样一个传说：她的神像，由罗马的主妇们献上，并称为\OriginalTerm{妇女命运女神}{Fortuna Muliebris}，竟然开口说过话，并且不止一次地说，\Quote{主妇们的敬礼令她喜悦。}此事若真，也不值得我们惊奇。因为对于恶毒的魔鬼而言，欺骗并非难事；而他们更应留心其诡计与奸诈，因为说话的是那位偶然来临的女神，而非那位为酬报功德而来的女神。因为\OriginalTerm{福图娜}{Fortuna}多嘴多舌，而\OriginalTerm{费利西塔斯}{Felicitas}却沉默不语；这究竟是何缘故，岂不正是为了让人们不在乎正当生活，而与\OriginalTerm{福图娜}{Fortuna}结为朋友，因为她无需任何善工便能使人幸运吗？的确，如果\OriginalTerm{福图娜}{Fortuna}说话，她至少不该用女声，而应用男声；免得那些奉献此神像的人以为，如此重大的奇迹，竟是由女性的饶舌所成就的。\par

% structure: p0045 source=h2 target=subsection reason=relative-heading-rank; base=h2

\subsection{第二十章 —— 论异教徒以庙宇与神圣仪式尊崇的美德与信德，却忽略了其他同样应受敬拜的善德——倘若神性本当归于这些品质的话。}

他们还将\OriginalTerm{德性}{Virtue}立为女神，的确，如果它能成为女神，就比许多神更可取。而今，因为它并非女神，而是天主的恩赐，所以该当藉着祈祷向祂求得\Emph{它}——唯独祂能赐予这恩赐，于是所有虚假的神祇便销声匿迹了。但为何\OriginalTerm{信德}{Faith}被奉为女神，她自己又为何享有庙宇和祭坛？因为凡明智地接纳她的，就把自己变成了她的居所。但他们如何知道信德是什么？它的首要和最大的功能，是使人相信真天主。但为何德性还不够？它不也包括信德吗？因为他们认为该把德性分为四个部分——\OriginalTerm{智德}{prudence}、\OriginalTerm{义德}{justice}、\OriginalTerm{勇德}{fortitude}和\OriginalTerm{节德}{temperance}——而每一部分又各有其德性，信德便属於义德的范畴，并在我们当中那些明白以下这句话的人中居於首位：\BibleQuote{义人必因他的信德而生活。}\BibleRef{哈 2:4} 但如果信德是女神，我便纳闷这些热衷於众多神祇的人，为何竟亏待了那麽多其他女神，把她们忽略过去，他们原本同样能为她们奉献庙宇和祭坛。为何节德不配成为女神？有些罗马君主正因为节德而赢得了不小的荣耀。最后，为何勇德不是女神？它曾帮助\Person{穆修斯}将右手伸入火焰；帮助\Person{库尔提乌斯}为了祖国纵身跃入裂开的大地；帮助\Person{老德西乌斯}和\Person{小德西乌斯}舍身献祭以保全军队？——尽管我们可以质疑这些人是否具有\Emph{真正的}勇德，如果这与我们当前的讨论有关的话。为何智德和智慧竟未能在众神中占有一席之地？难道是因它们全都在德性这个总名之下受人敬拜？那麽他们也可如此敬拜真天主，因为其他所有神祇都被认为只是祂的部分。但在德性这一个名称下，已经包含了信德和\OriginalTerm{贞洁}{chastity}，而它们却独自在各自的庙宇中拥有了祭坛。\par

% structure: p0047 source=h2 target=subsection reason=relative-heading-rank; base=h2

\subsection{他们虽不理解这些是天主的恩赐，至少也应以\OriginalTerm{德性}{Virtue}和\OriginalTerm{幸福}{Felicity}为足。}

这些东西，不是真理而是虚妄造作了女神。因为它们是真实天主的恩赐，它们本身并非女神。然而，哪里有了德性和幸福，还寻求别的什么呢？对于一个德性和幸福都不能令其满足的人，还有什么能满足他呢？因为毫无疑问，德性包含了一切我们需要去行的事，幸福包含了一切我们所渴望的事。那么，如果人们崇拜\Person{朱庇特}是为了让他赐予这两样东西——因为若帝国的广大和长久算是一种善好，它就属于这同一幸福——为何不明白这两样并非女神，而是天主的恩赐呢？但如果人们判定她们是女神，那么至少那一大群其他神灵就不应当再去寻求了。因为，且让他们考虑他们幻想分配给众多男女神灵的一切职司，若有可能，让他们找出任何一个神灵能够赐予一个有德性、有幸福的人什么事物。当德性本身已经拥有一切时，还需向\Person{墨丘利}或\Person{密涅瓦}寻求什么教导呢？古人确实将德性界说为善生与正当生活的\OriginalTerm{技艺}{ars}本身。因此，因德性在希腊语中称为\OriginalTerm{德性}{\GreekText{ἀρετη}}，人便认为拉丁人由此衍生出\Emph{技艺}一词。然而，若德性非聪明人不能得，又何需父神\Person{卡提乌斯}使人谨慎即敏锐呢？幸福女神岂不能赐予此吗？因为生来聪明归属于幸福。因此，虽然幸福女神不能被尚未出生者敬拜，以至与之交友而赐予他聪明，她却能将这恩惠赐予敬拜她的父母，使聪明的儿女生于他们。分娩中的妇女何需呼求\Person{卢基娜}，若幸福女神在场，她们不仅会顺利分娩，更会得好子女？当儿女出生时，何需把他们托付给\Person{俄普斯}女神；在他们初啼时，托付给\Person{瓦提卡努斯}神；当他们躺在摇篮中时，托付给\Person{库尼娜}女神；吮乳时，托付给\Person{里米娜}女神；站立时，托付给\Person{斯塔提利努斯}神；来去时，托付给\Person{阿德奥娜}女神去，托付给\Person{阿贝奥娜}女神归；为使孩子心智良好，托付给\Person{门斯}女神；为使他们愿欲善美之物，托付给\Person{沃鲁姆努斯}神和\Person{沃鲁姆娜}女神；为缔结良缘，托付给婚姻诸神；为得最丰盛的果实，托付给乡村诸神，尤其托付给\Person{弗鲁克特斯卡}女神本身；为能善于作战，托付给\Person{马尔斯}和\Person{贝洛娜}；为得胜利，托付给\Person{维多利亚}女神；为受尊荣，托付给\Person{霍诺尔}神；为有丰足金钱，托付给\Person{佩库尼亚}女神；为有铜币银币，托付给\Person{埃斯库拉努斯}神及其儿子\Person{阿根提努斯}？因为他们把\Person{埃斯库拉努斯}认定为\Person{阿根提努斯}的父亲，原因是铜币的使用早于银币。然而我奇怪\Person{阿根提努斯}竟没有生出\Person{奥里努斯}，因金币后来也出现了。倘若他们能以他为神，他们必爱\Person{奥里努斯}胜于其父\Person{阿根提努斯}及其祖\Person{埃斯库拉努斯}，正如他们爱\Person{朱庇特}胜于\Person{萨图恩}。因此，为了这些或属于灵魂、或属于身体、或属于外在环境的恩惠，何须敬拜并呼求如此众多神灵呢——这些神祇我尚未悉数提及，他们自己也未能涵盖人类所有福益，逐一细分，零散琐碎，而单单一位幸福女神便能无比轻易、简括地赐予这一切？也不应需求任何其他神灵，无论是为赐予善物，还是为消除邪恶。因为，对于疲倦之人，何须呼求\Person{费松尼亚}女神；为驱除敌人，何须呼求\Person{佩洛尼亚}女神；对于病患，既已有医生\Person{阿波罗}或\Person{埃斯库拉庇乌斯}，或者若有大危，二者同求？既不必祈求\Person{斯皮尼恩西斯}神来除去田间荆棘，也毋须哀求\Person{鲁比戈}女神以免霉病——只要幸福女神独自临在守护，要么恶事根本不会兴，要么轻易便被驱散。最后，既然我们论的是这两位女神——\OriginalTerm{德性女神}{Virtue}与\OriginalTerm{幸福女神}{Felicity}，若幸福是德性的报酬，她便不是女神，而是天主的恩赐。但若她真是女神，为何不能说她也赐予德性本身呢？因为获得德性就是莫大的幸福。\par

% structure: p0049 source=h2 target=subsection reason=relative-heading-rank; base=h2

\subsection{论瓦罗自诩传授给罗马人的、关于应敬拜诸神之知识}

那么，瓦罗自夸他赐给同胞极大的恩惠，因为他不但列举了罗马人应当敬拜的诸神，而且还指出每一位神分别掌管什么，这算什么呢？他说：\Quote{正如知道一个医生的姓名和相貌，却不知道他是医生，这毫无益处；同样，明知\Person{埃斯库拉庇乌斯}是神，却不知他能赐予健康之恩，因而不知为何应当祈求他，也毫无益处。}他又以另一比喻肯定这一点，说：\Quote{一个人如果不知道谁是铁匠，谁是面包师，谁是织工，从谁可以求取用具，可找谁当帮手，可找谁当领袖，可找谁当教师，不但不能善生，连存活都不可能；}他断言：\Quote{这样一来，就不会有人怀疑有关诸神的知识是有益的，只要人能知道每一位神对于任何事物具有何种力量、禀赋或能力。因为由此我们便能知道，为何事应当呼求并祈求哪位神；免得像许多人所惯做的那样，向\Person{利伯尔}求水，向\Person{林菲斯}们求酒。}的确非常有用！谁不会感谢此人呢，假如他能揭示真实的事理，能教导人应当敬拜那唯一真实的天主，一切美善都是从他而来？\par

% structure: p0051 source=h2 target=subsection reason=relative-heading-rank; base=h2

\subsection{论幸福女神——罗马人虽崇奉众多神灵，却长期未以神圣尊荣敬拜她，尽管单是她便足以代替一切}

可是，如果他们的书籍和礼仪是真实的，且\Quote{幸福}是一位女神，那么为何不指定她为唯一应受敬拜的神呢？既然她能赐予一切，并一举使人幸福。因为除了为使自己幸福，还有谁会为了别的缘由而希求任何事物呢？为何要让卢库卢斯在如此晚近的时期，且在那么多罗马统治者之后，才为如此伟大的女神奉献殿宇呢？罗穆卢斯本人如此热望建立一座幸运的城市，为何不先于其他一切神明，为这女神建造殿宇呢？他为何还向别的神明祈求任何事物，因为若她与他同在，他便一无所缺。因为倘若这女神未曾垂顾他，就连他自己，也不会先成为国王，而后如他们所想，成为神明。因此，为何他为罗马人指定的神明是雅努斯、朱庇特、玛尔斯、皮库斯、法乌努斯、提伯里努斯、赫丘利，以及其他神明，如果还有更多的话？为何提图斯·塔提乌斯增加了萨图恩、俄普斯、太阳、月亮、伏尔甘、光，以及他所增添的别的，其中甚至有女神科洛阿奇娜，却忽略了幸福女神？为何努马任命了如此多男神女神，却缺了这一位？或许是因为在如此众多神明中他看不见她？当然，国王霍斯提利乌斯若认识或可能敬拜过这女神，就不会引进新的神明\Quote{恐惧}与\Quote{畏惧}来求取悦纳。因为在幸福面前，恐惧与畏惧就会消失——我不说被悦纳，而是被驱散。其次，我要问，罗马帝国远在任何人事奉幸福女神之前，便已极度扩张，这如何可能？那么，帝国岂非伟大有余，而幸福不足？因为那里若无真正的虔敬，何来真正的幸福？因为虔敬乃是对真天主的真实敬拜，而非拜多少假神就有多少魔鬼。然而，甚至后来，当幸福已被纳入神祇之列时，内战的极大不幸随之而来。或许幸福女神义愤难平，既因她被邀请得如此之晚，又因邀请她并非为尊崇，反似羞辱，因为与她同受敬拜的，还有普里阿普斯、科洛阿奇娜、恐惧、畏惧、疟疾，以及其他并非当受敬拜的神明，而是敬拜者本身的罪恶？最后，若他们认为将如此伟大的女神与最不配的乌合之众一同敬拜仍为美事，那么至少为何不以较其余神明更尊崇的方式敬拜她呢？因为幸福女神既未被安置于\OriginalTerm{协商诸神}{Consentes}——他们声称这些神祇得入朱庇特的会议——也未列于他们所谓的\OriginalTerm{特选神}{Select}之中，这岂非难以容忍？本当为她建造一座殿宇，无论在地势之高耸，还是在风格之庄严上，都出类拔萃。的确，为何不建造比朱庇特本人的更胜一筹的殿宇呢？因为除了幸福，谁曾将王国赐予朱庇特？我姑且假定他统治时是幸福的。然而，幸福确实比王国更为珍贵。因为无人怀疑，很容易找到害怕成为国王的人；却找不到不愿幸福的人。因此，如果认为可以通过鸟占或任何其他方式求问神明，他们应当就此事求问诸神本身，看他们是否愿意给幸福让位。倘若恰巧某一地点已为其他神明的殿宇和祭台所占据，而该处本可为幸福建造一座更宏大更崇高的殿宇，那么就连朱庇特本人也该让位，好让幸福反得占据卡比托利山巅的至高点。因为除了一位——这是不可能的——愿自己不幸福的人之外，无人会抗拒幸福。当然，若有人求问朱庇特，他绝不会做出玛尔斯、特耳米努斯和尤文塔斯那三位神祇所做的事，他们断然拒绝给他们的尊长和君王让位。因为，据他们的书籍记载，当国王塔奎纽斯想要建造卡比托利神殿时，察觉到他看来最配得且最合适的地点已被其他神明占据；他不敢做任何违背他们心意的事，相信他们会心甘情愿给一位如此伟大、是他们自己的主的神明让位——因为在卡比托利奠基时，那里已有许多神祇——他便通过鸟占求问他们是否愿意给朱庇特让位；结果除了我已指名的玛尔斯、特耳米努斯和尤文塔斯之外，全都愿意从那里迁走。因此卡比托利神殿的建造方式，使这三位神祇也能留在其内，但其表征如此晦暗，甚至最博学的人也不易知晓。那么，朱庇特本人当然绝不会轻蔑幸福，如同他自身被特耳米努斯、玛尔斯和尤文塔斯所轻蔑那样。然而，就连那些不愿意给朱庇特让位的神祇，也定会给幸福让位，因为是幸福使朱庇特作了他们的君王。或者，若他们不让位，他们这样做并非出于轻视她，而是因为他们宁愿在幸福的殿宇中隐而不显，也不愿在自己的地方无幸而显赫。\par

因此，\OriginalTerm{幸福女神}{Felicitas}既被安置在最大最高的地方，公民自当晓得应从何处寻求一切美善愿望的成全。于是，受本性自身的劝导，舍弃其他多余的神祇，单拜幸福女神一位，只向她祈祷，她的神殿独为渴望幸福之公民所常临——而无人不愿幸福；这样，那向一切神明寻求的幸福，便只向她本身寻求了。因为谁愿从任何神明领受幸福以外的其他事物，或那他认为引向幸福之物呢？因此，如果幸福女神有能力与她所喜悦的人同在（若她是女神，便有这能力），那么，竟然绕过她本身可求得的，反去求其他神，该是何等愚拙！所以，他们理当以场所的尊荣，将这位女神敬奉在其他众神之上。因为，我们读他们自己的作家时看到，古代罗马人对某位\Person{苏曼努斯}，将夜间雷霆归于他，所献的尊敬，竟超过\Person{朱庇特}，而白昼的雷霆则被认为属于朱庇特。但是，等到一座著名而宏伟的神殿为朱庇特建成后，由于建筑的尊严，群众蜂拥而至，以致如今几乎找不到一个记得曾读到苏曼努斯名字的人，甚至听都没再听人提起。然而，如果幸福不是女神，因为，如同真理所示，它乃是天主的恩赐，那么，就当寻求那位有能力赐予幸福的天主，抛弃那群有害的假神——虚妄无知之辈追逐这些假神，将天主的恩赐造作偶像，并以骄傲意志的顽梗，得罪了恩赐这些恩赐的主。因为，敬拜幸福如女神而离弃赐幸福的天主者，不能免于不幸；正如舔食画中面包而不向持有真面包的人购买者，不能免于饥饿。\par

% structure: p0054 source=h2 target=subsection reason=relative-heading-rank; base=h2

\subsection{第24章 异教徒用来为他们在众神中敬拜那些天主的恩赐本身所作辩护的种种理由。}

然而，我们可以思考他们的理由。他们说，难道可以相信我们的祖先愚钝到竟不知道这些事物是神明的恩赐，而非神明吗？但他们既知道，若非有某位神明白白赐予，无人能得着这些事物，便以他们认为由这些神明所赐之事物的名称，来称呼那些他们未找出名字的神明；为此目的，有时稍稍改变词语，例如，从\Emph{bellum}（战争）造出\Person{贝娄娜}，而非\Emph{bellum}；从\Emph{cunae}（摇篮）造出\Person{库尼娜}，而非\Emph{cunae}；从\Emph{seges}（禾谷）造出\Person{塞格提亚}，而非\Emph{seges}；从\Emph{pomum}（苹果）造出\Person{波摩娜}，而非\Emph{pomum}；从\Emph{bos}（牛）造出\Person{布波娜}，而非\Emph{bos}。有时，又对词语不作变动，就如事物本身的名称那样，于是赐金钱的女神便称为\Person{佩库尼亚}，而金钱并不被视作女神本身：同样，赐勇德的称为\Person{维耳图斯}；赐尊荣的称为\Person{霍诺耳}；赐和合的称为\Person{孔科尔迪亚}；赐得胜的称为\Person{维多利亚}。他们说，所以，称幸福为女神时，所指的非被赐予的事物本身，而是那位赐予幸福的神明。\par

% structure: p0056 source=h2 target=subsection reason=relative-heading-rank; base=h2

\subsection{第25章 论唯独应敬拜的独一天主，祂的名虽无人知晓，却被认作幸福的赐予者。}

我们既已得到这样的理由，或将更容易，如我们所愿，说服那些心肠尚未过于刚硬的人。因为，若说人性软弱已领悟到，幸福非由某位神明则不能赐予；若那些敬拜众神、且奉\Person{朱庇特}为其首领的人也已领悟到这一点；若他们不知那赐予幸福者的名，便同意以他们认为祂所赐之物的名来称呼祂——那么，结论便是，他们认为即使朱庇特，他们已在敬拜的，也不能赐予幸福，而那位被认为当以幸福本身之名来敬拜的，才确实能赐予幸福。我完全肯定这种说法：他们相信幸福是由一位他们不认识的某位天主所赐予的。因此，当寻求祂，当敬拜祂，便足够了。让那无数邪魔的队列被摒弃吧，让祂的恩赐满足世人的这位天主，也满足每一个人吧。因为，我说，对于领受幸福本身尚嫌不足的人，敬拜赐幸福的天主也不会令他满足。但是，谁若以此便感满足（而人不应再有更多所求），就当侍奉那独一的天主，幸福的赐予者。这位天主并非他们称为朱庇特的那位。因为，若他们承认朱庇特是幸福的赐予者，他们就不会在幸福本身之名下，去寻求另一位能赐幸福的神或女神；他们也不能容忍朱庇特自身以如此可耻的属性被敬拜。因为，据说他玷污他人之妻，他是一位美少年的无耻情人和强夺者。\par

% structure: p0058 source=h2 target=subsection reason=relative-heading-rank; base=h2

\subsection{第26章 论戏剧演出，即众神强令其敬拜者举行的那类庆祝。}

\Quote{但是，}\Person{西塞罗}说，\Quote{\Person{荷马}捏造了这些事，将人间的事转移到了众神身上：我宁愿将属神的事转移给我们。}这位诗人将如此罪行归诸众神，理所当然地冒犯了这位严肃的人。那么，为什么在这些罪行被惯常谈论、表演、展示以荣耀众神的戏剧，却被最有学问的人算作属神的事呢？\Person{西塞罗}应当谴责的，不是诗人的虚构，而是古人的习俗。他们岂不会回答说：\Quote{我们做了什么？}众神自己曾高声要求在尊崇他们的庆典中上演这些戏剧，曾苛刻地强求，曾威胁若不举行就要降灾，对怠慢严加报复，并对这种怠慢的补偿表示欢悦。在众神的德能奇事中，流传着这样一件事。一位名叫\Person{提图斯·拉丁尼乌斯}的\Place{罗马}乡下人在梦中得到启示，要去到元老院，并告诉他们重新开始\Place{罗马}的竞赛，因为在竞赛开始的第一天，一名被判死刑的罪犯当着人民的面被带去处决，这一悲伤事件扰乱了希望在竞赛中寻求乐趣的众神。当这位得到指示的农夫在次日不敢将其告知元老院后，在下一夜，同样的启示以更严厉的方式再次出现：他因疏忽而失去了儿子。第三夜，他受到警告，如果他仍然拒绝服从，更重的惩罚即将降临。即使如此，他仍不敢服从，结果患上了恶性可怕的疾病。但随后，在朋友的建议下，他向地方官报告了情况，并被用担架抬到了元老院；在叙述了梦境之后，他立刻恢复了体力，健步离开。元老院对如此伟大的奇迹感到惊异，下令以四倍的费用重新举办竞赛。哪一个明理之人看不出，人们被邪恶的魔鬼所劫持，除了藉着我主\Person{耶稣基督}，天主的\OriginalTerm{恩宠}{grace}，无人能从\OriginalTerm{魔鬼}{devil}的统治下得到释放，被迫向这样的神明展示那些，如果审慎考虑，他们理应谴责为可耻的戏剧呢？正是在这些戏剧中，诗人们歌颂众神的罪行，而这些戏剧正是由于众神的强制，由元老院的法令而重新设立的。在这些戏剧中，最无耻的演员将\Person{朱庇特}赞颂为贞洁的败坏者，并以此取乐。如果那是虚构，他本该被激怒；但若他对自己罪行的再现感到愉悦，即使那仅仅是神话，那么，每逢他受到崇拜，除了魔鬼之外，谁还会受到侍奉呢？难道这样一位比任何不悦此类事的\Place{罗马}人都更加卑劣的神，能够建立、扩张并保存\Place{罗马}帝国吗？一位受到如此\OriginalTerm{不幸}{infelicitously}崇拜的神，能够赐予\OriginalTerm{幸福}{felicity}吗？且若他未得到如此崇拜，还会更加\OriginalTerm{不幸}{infelicitously}地被激怒。\par

% structure: p0060 source=h2 target=subsection reason=relative-heading-rank; base=h2

\subsection{第27章 论\Person{斯凯沃拉}\OriginalTerm{司祭}{pontiff}所谈论的三类神祇}

据记载，博学的\OriginalTerm{大司祭}{pontiff}\Person{斯凯沃拉}曾区别三类神——一类由诗人引入，一类由哲学家，一类由政治家。他宣称第一类神是轻浮的，因为诗人编造了许多有关诸神的不堪之事；第二类不适合国家，因为它包含某些多余的事，以及某些若为人民所知便会造成损害的事。那些多余的事无关紧要，因为精通法律的律师们常说：\Quote{多余的事并无害处。}然而，那些摆在民众面前会造成损害的是什么呢？他说：\Quote{这些——\Person{赫丘利}、\Person{埃斯库拉庇乌斯}、\Person{卡斯托尔}与\Person{波路克斯}并不是神；因为有学问的人宣称，他们不过是人，且屈服于凡人的共同命运。}还有别的吗？\Quote{国家并没有神的真实形象；因为真天主既无性别，也无年龄，亦无确定的肢体。}大司祭不愿让百姓知道这些事，因为他并不认为它们是虚假的。因此，他认为在宗教事务上欺骗国家是有益的；\Person{瓦罗}本人在其论神圣事物的书中甚至毫不犹豫地这样说。绝妙的宗教！软弱者本可逃往其中寻求救助，因为他需要被拯救；而当他寻求使他得救的真理时，竟认为欺骗他是有益的。确实，在这些同样的书中，\Person{斯凯沃拉}并没有沉默他不接受诗人那类神的原因——即，\Quote{因为他们如此丑化诸神，以致他们甚至不能与好人相比，他们让一位神犯偷窃，另一位犯奸淫；或又说或做什么其他卑劣愚蠢的事；如三位女神争夺美貌的奖赏，而被\Person{维纳斯}击败的那两位毁灭了\Place{特洛伊}；\Person{朱庇特}变成公牛或天鹅以便与某位交媾；一位女神嫁给凡人，\Person{萨图恩}吞吃自己的子女；总之，没有任何可以想象得出的，无论是奇迹还是恶行，是那里找不到的，却又远非神的本质。}大司祭\Person{斯凯沃拉}啊，你若能够，就废除这些戏剧吧；教导百姓不要向不朽的神明献上这种尊崇，在其中，他们若是喜欢，可以欣赏神明的罪行，而且若他们愿意，甚至可以模仿，只要他们能够。但如果百姓这样回答你：你这位大司祭啊，是你把这些事带进我们中间的；那么你就去求问那些神明本身，你是受谁的唆使而吩咐这些事的，好叫他们不要命令人们向他们献上这些事。因为如果这些事是坏的，因此对于大多数神明完全不可信，那么那些不受惩罚地被捏造出来的神明所受的冤枉就更大了。但他们不会听你，他们是魔鬼，他们教导邪恶的事，他们以卑鄙的事为乐；他们不仅不认为这些被捏造在他们身上是冤枉，而且若不在他们规定的节日演出这些事，就是他们完全无法忍受的冤枉。但现在，你若为此事求告\Person{朱庇特}抵挡他们，主要是因为他的罪行更常在戏剧中被演出，难道事实不是：尽管你称他为统治和管理这整个世界的\Person{朱庇特}神，但正是你对他行了最大的冤枉，因为你认为他应当与那些神一同受敬拜，并尊他为他们的王？\par

% structure: p0062 source=h2 target=subsection reason=relative-heading-rank; base=h2

\subsection{敬拜诸神是否有助于罗马人获得并拓展帝国}

因此，这些神明，他们以如此尊崇被安抚，或者更确切地说，因这些尊崇而遭责难（因为以虚假之言说他们而令他们喜悦，这罪过比即使能说真事还要大），是根本不可能曾拓展并保存罗马帝国的。因为如果他们真能做到，他们宁愿将如此盛大的恩赐赠予希腊人，希腊人在这种神圣事物上——即戏剧——更加尊崇和相匹配地敬拜了他们，尽管他们并未使自己免受诗人的那些诽谤，任由诗人肆意撕裂神明，给予他们许可去任意辱骂任何人，并且不认为戏剧演员下贱，反而认为他们甚至配得特殊的荣誉。但正如罗马人能够拥有金币，尽管他们并不敬拜神\Person{奥里努斯}，同样他们也能有银币和铜币，却既不敬拜\Person{阿尔根提努斯}也不敬拜他的父亲\Person{埃斯库拉努斯}；其余的一切也是如此，详细叙述将令人生厌。因此，结论便是：他们无论如何也不可能获得如此疆域，如果真天主不愿意的话；而且，如果这些虚假且为数众多的神明不为人知或被藐视，而唯独祂以真诚的信德和德性被认识、受敬拜，那么他们在此世会拥有一个更好的王国，不论其疆域多大，且无论他们有没有在此世的王国，日后都将承受一个永恒的王国。\par

% structure: p0064 source=h2 target=subsection reason=relative-heading-rank; base=h2

\subsection{论那被认为预示罗马帝国力量与稳固的占卜之虚假}

他们宣称最为美好的那种\OriginalTerm{占卜}{augury}，就是我刚才提到的，即\OriginalTerm{玛尔斯}{Mars}、\OriginalTerm{特尔米努斯}{Terminus}和\OriginalTerm{尤文塔斯}{Juventas}甚至不愿给众神之王\OriginalTerm{朱庇特}{Jove}让位，这是什么样的占卜呢？因为他们说，这预兆意味着：奉献给玛尔斯的民族——即罗马人——不会向任何人放弃他们曾经占据的地方；同样，因特尔米努斯神的缘故，没有人能够扰乱罗马的边界；而且，因尤文塔斯女神，罗马的青年不会向任何人屈服。因而，让他们看看，若这些预兆将他列为对手，且不向他屈服本是光荣的，他们又如何能视他为众神之王，且是他自己王国的赐予者呢？然而，若这些事是真实的，他们便完全不必害怕。因为他们不愿承认，那些不愿向朱庇特屈服的神祇已向基督屈服。因为，耶稣基督未曾改变帝国的边界，便已证明自己能够将他们，不但从他们的庙宇，并且从他们敬拜者的心中赶出去。然而，在基督降生成人之前，甚至在我们从他们的书中引述的这些事能够被记录下来之前，但在\Person{塔尔奎纽斯}王统治下那个\OriginalTerm{预兆}{auspice}出现之后，罗马军队已数次被击溃或败逃，这便暴露了那个预兆的虚妄，他们从尤文塔斯女神未给朱庇特让位这一事实得出此预兆；而奉献给玛尔斯的民族，在罗马城中，被入侵的\Place{高卢}人践踏得胜；帝国的边界，由于许多城市倒向\Person{汉尼拔}，已被压缩至狭小的范围。这样，预兆的美好之处便化为乌有，仅剩下对朱庇特的顽抗，不是出于神祇，而是出于魔鬼。因为不屈服是一回事，而回到曾屈服之处则是另一回事。此外，后来在东方地区，罗马帝国的边界因\Person{哈德良}的意愿而改变；因为他将那三个富饶的行省，\Place{亚美尼亚}、\Place{美索不达米亚}和\Place{亚述}，让给了\Place{波斯}帝国。如此，那位按照这些书籍是罗马边境守护神的特尔米努斯，且因那最美好的预兆而未给朱庇特让位的，似乎更畏惧人间的君王哈德良，而非众神之王。上述行省后来虽曾收复，几乎在我们的记忆中，边疆又退缩了，那时\Person{尤利安}，沉溺于他们神祇的\OriginalTerm{谕示}{oracles}，以无度的冒失下令焚烧运粮船。军队因而断绝了供应，他本人不久也被敌人杀死，军团陷入困境，因失去统帅而惊惶失措，被逼至绝境，若非以和约规定了帝国的边界，使其维持至今，则无人能够逃脱；虽未如哈德良的割让那般损失惨重，却仍付出了不小的牺牲。那么，特尔米努斯神不向朱庇特屈服的占卜，是虚妄的，因为他屈服于哈德良的旨意，也屈服于尤利安的鲁莽，以及\Person{约维尼安}的窘迫。较为明智且严肃的罗马人已看出这些，但无力反对国家的习俗，这习俗要他们遵守魔鬼的礼仪；因为即使他们自己，虽觉这些事虚妄，却以为应当将那唯独归于天主的宗教敬拜，献给那在独一真天主的治理下所建立的万有之本性，正如宗徒所说：\BibleQuote{去崇拜事奉受造之物，以代替造物主；他是永远可赞美的。}\BibleRef{罗 1:25}这位真天主的助佑是必要的，为要派遣圣善而真正虔诚的人，他们愿为真宗教而死，好能从活人中除去虚假的宗教。\par

% structure: p0066 source=h2 target=subsection reason=relative-heading-rank; base=h2

\subsection{第三十章——连崇拜者自己也承认他们对外邦神祇所怀的想法。}

\Person{西塞罗}这位\OriginalTerm{占卜官}{augur}嘲笑\OriginalTerm{占卜}{auguries}，并责备人依据乌鸦和寒鸦的叫声来规正人生。但有人会说，一位主张万事万物皆不确定的\OriginalTerm{学园派}{Academic}哲学家，在这些事上没有任何权威性。他在其\WorkTitle{De Natura Deorum}的第二卷中引入了\Person{卢基利乌斯·巴尔布斯}，后者在指明迷信源于物理和哲学真理之后，对设立神像和虚构概念表示愤慨，如此说道：\Quote{你难道没有看到，从真实而有益的自然发现中，理性会被引入虚构和想象的神明之中吗？这便产生了错误的意见、纷乱的错谬，以及近乎老妇般的迷信。因为神明们的相貌、年龄、衣着和饰物都已为我们所熟悉；他们的谱系、婚姻、亲族关系以及所有关于他们的事，都被降格到与人间软弱相似的地步。他们甚至被描绘成有着纷扰的心灵；因为我们有关于神明的情欲、忧虑和怒气的记载。而且，正如神话所述，神明们也并非没有战争和战斗。这不只是在\Person{荷马}史诗中那样，双方各有神明保护两支对立的军队，他们甚至为了自己而发动战争，比如与\Person{提坦}族或\Person{巨人}族的战争。这类事情，无论说还是信，都极其荒谬：它们完全是轻浮而又毫无根据的。}看哪，这正是那些维护列国神明的人所承认的。随后，他接着说，有些事属于迷信，有些则属于宗教，他认为依照\OriginalTerm{斯多亚派}{Stoics}的观点来教导宗教是好的。\Quote{因为不仅是哲学家们，}他说，\Quote{我们的祖先也已经区分了迷信和宗教。}他说，\Quote{那些整天祈祷、献祭，以求子女比他们活得长久的人，就被称为迷信者。}谁看不出来，他因为害怕公众的偏见，试图赞扬古人的宗教，并且希望将它与迷信分开，却找不到办法呢？假若那些整天祈祷献祭的人被古人称为迷信者，那么那些设立了他所责备的各样年龄、独特衣着的诸神形象，并编造了神明谱系、婚姻及亲族关系的人，是否也被同样称呼呢？因此，当这些事情被斥为迷信时，他便将设立并敬拜这些形象的古人牵连到这一过错之中。不仅如此，他也牵连了自己：无论他如何运用口才努力想脱身并获得自由，他还是不得不敬拜这些形象；而且他也不敢在面向民众的演讲中悄声说出他在此讨论中所清楚宣明的道理。所以，我们基督徒应当感谢主我们的天主——不是感谢天地，如那位作者所主张的，而是感谢那造了天地的祂；因为这些迷信，那位\Person{巴尔布斯}像饶舌者一样都几乎不敢指责，而祂藉着\Person{基督}至深的谦卑、宗徒们的宣讲、殉道者为真理而死并与真理同活的信德，不仅从虔诚者的心中，甚至在迷信者的庙宇中，通过他们自愿的服事，将其推翻了。\par

% structure: p0068 source=h2 target=subsection reason=relative-heading-rank; base=h2

\subsection{瓦罗的意见：他虽谴责民间信仰，却认为敬拜应限于一位神，尽管他未能发现真天主。}

\Person{瓦罗}他自己怎么说呢？我们发现他将戏剧列入神圣事物之中，尽管不是出于他自己的判断，我们为此感到惋惜。他在许多段落中，像一位虔敬者那样劝人敬拜神明时，难道不是在这样做的同时承认，他自己并不相信他所叙述的由罗马国家所设立的这些事吗？以至于他毫不犹豫地断言，如果他建立一座新城邦，他可以依照自然的法则更好地列举诸神及其名称。但是，由于生于一个已经古老的民族，他说，他发现自己不得不接受诸神的传统名称和别号，以及与之相关的传说，而他探究并发表这些细节的目的，是为了引导人民敬拜神明，而不是轻视它们。通过这些话语，这位极为敏锐的人充分地表明，他没有把所有事都公之于众，因为那些事不仅他自己会藐视，甚至对一般民众来说也会显得可鄙，除非被默然略过。我本可能被认为是在推测这些事，但他在另一段中，在谈到宗教仪式时，曾公开说过：有许多真实的事，不仅对一般民众知道它们毫无益处，而且，让民众以另一种方式，即使是错误的方式去想，反而更为适宜；因此，希腊人将宗教仪式和奥秘隐藏在沉默与围墙之内。在这一点上，他无疑表达了那些所谓的智者用以统治城邦和民族的政策。然而，这种狡诈的手段极大地取悦了邪恶的魔鬼，他们同样掌控着欺骗者和受骗者，而除了藉着我们的主耶稣基督而来的天主的恩宠，没有任何事物能将人从他们的暴政下解救出来。\par

那位极其敏锐而博学的作者还说，在他看来，只有那些相信天主是世界的灵魂，并以设计和理性治理世界的人，才领悟了天主是什么。由此看来，尽管他未能达到真理——因为真天主不是灵魂，而是灵魂的创造者和作者——但若他得以摆脱习俗的偏见，他本可以承认并劝告他人，应当敬拜那位以设计和理性治理世界的独一天主；这样，就这一问题，与他辩论的只剩下这一点：他称天主为灵魂，而非灵魂的创造者。他还说，古罗马人曾长达一百七十多年敬拜无偶像的神明。\Quote{倘若这种习俗}，他说，\Quote{能延续至今，神明便会受到更纯洁的敬拜。}为支持这一观点，他援引犹太民族作为证据之一；并在那段话的末尾毫不犹豫地说，那些首先为人民祝圣偶像的人，既夺去了同胞的宗教敬畏，又加增了谬误，因为他们明智地认为，当神明以笨拙的偶像形式展示时，便容易遭到轻蔑。但他并没有说他们传播了谬误，而是说他们加增了谬误，因此他希望人们明白，在没有偶像时便已存在谬误。因此，当他说只有相信天主是治理世界的灵魂的人才领悟天主是什么，并认为宗教礼仪若无偶像会更为纯洁地遵守时，谁能看不出他离真理已多么近？因为倘若他能反对如此根深蒂固的谬误，他必定会提出自己的意见：既应敬拜独一天主，又应无偶像地敬拜；当他如此近地发现了真理，或许可以轻易地被提醒灵魂的可变性，从而领悟真天主是创造灵魂本身的不可改变的本性。既然事实如此，那么这类人在其著作中对多神论所倾泻的讥讽，与其说是出于试图说服他人，不如说是被天主隐秘的旨意所迫而承认它们。因此，我们若从这些著作中引用任何证言，便是为了驳斥那些不愿思考倾流至圣之血的独一祭献，以及所赐圣神的恩赐，能从何等巨大而邪恶的魔鬼权下将我们解救出来的人。\par

% structure: p0071 source=h2 target=subsection reason=relative-heading-rank; base=h2

\subsection{万民的领袖出于何种利益，希望虚假宗教在其所统治的人民中间延续。}

\Person{瓦罗}还谈到神明的谱系，说人民更倾向于诗人而不取法自然哲学家；因此他们的祖先——即古罗马人——相信神明的性别和谱系，并为他们安排婚配。这显然没有别的原因，无非是那些精明而智慧的人，在宗教事务上欺骗人民，并且在这件事上不仅敬拜魔鬼，还效法魔鬼，因为魔鬼最大的欲望就是欺骗。因为正如魔鬼只能占有所欺骗的人，那些居王者之位的人，虽然并不公义，却像魔鬼一样，假借宗教之名说服人民，将他们明知虚假的事当作真实接受；这样，便仿佛将他们更牢固地捆绑在公民社会中，好能以同样方式占有他们为臣民。然而，软弱无知的人，谁能逃脱国家君主与魔鬼的双重欺骗呢？\par

% structure: p0073 source=h2 target=subsection reason=relative-heading-rank; base=h2

\subsection{一切君王与国家的时期，皆由真天主的裁断和大能所定。}

因此，那位\OriginalTerm{真福}{felicity}的作者和赐予者天主，因为唯独祂是真天主，亲自将世间王国赐予善人和恶人。祂并非轻率而行，仿佛出于偶然——因为祂是天主，不是命运——而是依照事物和时期的秩序，这秩序对我们隐藏，但祂全然知晓；然而祂并不像服侍时间秩序那样受制于它，而是以主宰的身份治理并安排它。真福，祂只赐予善人。一个人无论是臣民还是君王，并无分别；他同样可以拥有或不拥有它。真福将在那再无君王与臣民的生命中圆满实现。因此，世间王国由祂赐予善人和恶人；免得祂的敬拜者在仍处于极其软弱的心灵引领之下时，将这些恩赐当作什么大事来贪求。这就是\BibleBook{}{旧约}的奥秘，其中隐藏着\BibleBook{}{新约}：即便在那些世间恩惠也被许诺的时代，属灵的人已经领悟了——尽管尚未公开宣布——这些世间之物所象征的永恒，以及在哪项天主的恩赐中才能找到真正的真福。\par

% structure: p0075 source=h2 target=subsection reason=relative-heading-rank; base=h2

\subsection{论犹太人的王国，它由那位独一真天主建立，并在他们持守真宗教时受其保全。}

因此，为使人知道这些属世的好东西——那些无法想象更好事物的人所渴求的一切——乃是在唯一天主的能力之内，而非在罗马人从前认为值得敬拜的众多假神手中，祂使自己的子民在埃及从极少数人繁衍增多，并用奇妙的征兆将他们从那里解救出来。当他们的后嗣难以置信地增多时，他们的妇女并未呼求\OriginalTerm{路喀娜}{Lucina}；那民族如此惊人地增长，是天主亲自解救，亲自拯救他们脱离迫害他们、想要杀死他们所有婴儿的埃及人之手。没有\OriginalTerm{鲁弥娜}{Rumina}，他们照样吮乳；没有\OriginalTerm{库尼娜}{Cunina}，他们照样安睡在摇篮中；没有\OriginalTerm{埃杜卡}{Educa}与\OriginalTerm{波提娜}{Potina}，他们照样吃喝；没有所有那些幼稚的神明，他们照样受教育；没有婚姻之神，他们照样结婚；没有敬拜\OriginalTerm{普里阿普斯}{Priapus}，他们照样有夫妻的结合；没有呼求\OriginalTerm{涅普顿}{Neptune}，分开的海水仍为他们开出一条通路，又用合拢的波浪淹没了追赶他们的仇敌。当他们从天上领受玛纳时，并未祝圣任何\OriginalTerm{曼尼亚}{Mannia}女神；口渴时被击打的磐石为他们涌出水来，他们也没有朝拜\OriginalTerm{宁芙}{Nymphs}与\OriginalTerm{林芙}{Lymphs}。没有\OriginalTerm{玛尔斯}{Mars}与\OriginalTerm{贝娄娜}{Bellona}的疯狂仪式，他们照样作战；而他们诚然并非没有胜利而得胜，但胜利不被视为一位女神，而是他们天主的恩赐。无\OriginalTerm{塞杰提亚}{Segetia}而有收成；无\OriginalTerm{布波娜}{Bubona}而有牛；无\OriginalTerm{梅洛娜}{Mellona}而有蜜；无\OriginalTerm{波摩娜}{Pomona}而有苹果。总而言之，罗马人认为必须向如此一大群假神祈求的一切，他们从唯一天主那里领受得更为幸福。倘若他们不以不虔敬的好奇心犯罪得罪祂——这好奇心像魔术一样诱惑他们，引他们归向陌生的神明和偶像，最终导致他们杀害基督——他们的国度本会存留至今，并且，若不在疆域上更辽阔，也定比罗马的国度更幸福。如今他们分散在几乎万国万族之中，乃是出于那唯一天主的眷顾；以便当假神的像、祭台、圣林和庙宇在各处被推翻，其祭献被禁止时，能从他们自己的经书中显明，这些事如何早已由他们的先知预言了；免得当这些事在我们书中被人读到的时候，或许会显得是我们捏造的。但现在，将后续内容留待下一卷，我们必须在此为本卷的冗长设一个界限。\par

\SourceRecord{Translated by Marcus Dods.；Nicene and Post-Nicene Fathers, First Series；Vol. 2.；Edited by Philip Schaff.；Buffalo, NY: Christian Literature Publishing Co.,；1887.；<http://www.newadvent.org/fathers/120104.htm>.}

% structure: p0001 source=h1 target=section reason=page-title

\section{\WorkTitle{天主之城（卷五）}}

\Person{奥古斯丁}首先讨论\OriginalTerm{命运}{fate}学说，为要驳斥那些倾向于将\Place{罗马帝国}的强盛与扩张归因于命运的人——如前一卷所示，这不能归于伪神。之后，他证明天主的\OriginalTerm{预知}{prescience}与我们的\OriginalTerm{自由意志}{free will}并不矛盾。随后他谈及古罗马人的\OriginalTerm{德行}{virtue}，阐明在何种意义上\Place{罗马帝国}的扩张是由于罗马人自身的德行，又在多大程度上是由于\OriginalTerm{天主的计划}{counsel of God}，尽管他们并未敬拜天主。最后，他解释了何谓基督徒皇帝的\OriginalTerm{真幸福}{true happiness}。\par

% structure: p0003 source=h2 target=subsection reason=relative-heading-rank; base=h2

\subsection{序言。}

既然确定的乃是：完全达成我们所渴望的一切即为\OriginalTerm{幸福}{felicity}，而这幸福并非女神，乃是天主的恩赐，因此人只能敬拜那位能使他们幸福的天主——若\OriginalTerm{幸福女神}{Felicity}本身便是一位女神，她自然应成为唯一的敬拜对象——既然，我说，这一点已确定，现在让我们继续探讨：为什么天主——祂虽能将那些不善、因而也不幸福之人所能享有的美好恩赐连同其他一切事物赐予人——却认为应当将如此广大且长久的统治权授予\Place{罗马帝国}；因为此非由他们所敬拜的众多伪神所造成，我们既已提出，也将在适当机会提出相当的证明。\par

% structure: p0005 source=h2 target=subsection reason=relative-heading-rank; base=h2

\subsection{第一章：\Place{罗马帝国}及一切王国的原因既非偶然，也不在于星辰的位置。}

\Place{罗马帝国}强大的原因，按照那些将事物称为\OriginalTerm{偶然的}{fortuitous}和\OriginalTerm{宿命的}{fatal}之人的判断或意见，既非偶然，也非宿命，他们称那些无原因或无出于可理解秩序之原因的事物为\OriginalTerm{偶然的}{fortuitous}，称那些不依赖天主和人的意志、由某种\OriginalTerm{秩序}{order}的必然性而发生的事物为\OriginalTerm{宿命的}{fatal}。简言之，人的王国是由\OriginalTerm{天主眷顾}{divine providence}所建立的。如果有人认为它们的存在是由于\OriginalTerm{命运}{fate}，因为他将天主的意志或大能本身称为命运，那就让他保留他的意见，但纠正他的用语。因为他为什么不一开始就说出来呢，当有人问他：\Quote{你所理解的\OriginalTerm{命运}{fate}是什么？}时，他随后将要说的。因为当人们听到那个词时，根据语言的通常用法，他们简单地理解为，在某人出生或受孕时可能存在的特定\OriginalTerm{星辰的位置}{position of the stars}的力量，有些人将这种力量完全与天主的意志分开，而另一些人则断言这也依赖于那意志。但是，那些认为星辰不依赖天主的意志而决定我们将要做什么，或将拥有什么好处，或将遭受什么灾祸的人，所有人都应拒绝倾听他们，不仅那些持有真宗教的人，而且那些愿意敬拜任何神明，甚至是伪神的人，也应如此。因为这种意见的实质无非是说，没有任何神明应当被敬拜或祈求。然而，我们当前的讨论并非旨在反对这些人，而是反对那些为保卫他们所认为的神明而反对基督宗教的人。然而，那些使星辰的位置依赖于天主的意志，并在某种程度上决定每个人的性格，以及将有何种善恶降临到他身上的人，如果他们认为这些星辰从天主的至高权能那里领受了那种力量，以便按照自己的意志决定这些事情，那么他们就极大地伤害了天界——在其最辉煌的元老院，也可以说是最灿烂的会议厅中，他们竟设想邪恶的行为被决定施行——这样的行为，如果任何地上的国家决定施行，它将被全人类的判决定罪推翻。那么，对于人的行为，当被赋予天界的必然性时，那位既是星辰也是人类之主的天主，对这些行为还有什么判断可言？或者，如果他们不说星辰——尽管它们确实从至高天主那里领受了某种力量——按照自己的判断决定这些事情，而只是说天主的命令通过它们在运用和实施这些必然性时被工具性地履行，那么我们是否就要对天主抱有那种我们甚至认为对星辰的意志都不应有的看法呢？但是，如果说星辰与其说是致使这些事情发生，不如说是预示它们，因此那\OriginalTerm{星辰的位置}{position of the stars}如同一种预言未来而非致使未来的语言——因为这曾是学识不凡之人的意见——那么，\OriginalTerm{占星家}{mathematicians}们通常并不这样说，例如，他们说\Person{玛尔斯}在某某位置\Quote{预示着}杀人，而不是\Quote{导致}杀人。但是，尽管如此，即便我们承认他们说话不得体，而我们应当接受哲学家们在预测他们认为从星辰位置中发现的事物时所使用的那种恰当的说话方式，那么，他们何以从来不能解释，为什么在双胞胎的生活、行为、遭遇的事件、职业、技艺、荣誉以及其他与人类生活相关的事情中，甚至在他们的死亡本身，常常有如此巨大的差异，以至于就这些事情而言，许多完全的陌生人比他们彼此之间更加相像，尽管他们在出生时只相隔极短的时间，但在受孕时却是由同一交合行为在同一瞬间生成的？\par

% structure: p0007 source=h2 target=subsection reason=relative-heading-rank; base=h2

\subsection{第二章：论双胞胎健康状况的差异。}

\Person{西塞罗}说，著名的医生\Person{希波克拉底}曾留下文字记载，他怀疑某对兄弟是双胞胎，因为他们同时染病，且疾病的进程、危象及消退时间在两人身上完全相同。斯多葛派的\Person{波西多纽斯}酷爱占星术，他解释此事时假设他们是在同一星座下受孕和出生的。在这个问题上，医生的推测远比占星术士的推测值得接受，也更近乎可信，因为，根据父母在交媾时的身体状况，胎儿的初始元素会受到影响，从而他们从同一母体获得生长和发育直至出生所需的一切，可能带着相似的体质降生。此后，在同一家庭中受养育，饮食相同，所处空气、地域、水质也相同——据医学证明，这些对身体健康状况好坏有极大影响——并且他们习惯于相同的操练，他们的体质会非常相似，以至于他们会在相同时间、因相同原因患上同样的疾病。然而，若要将他们出生或受孕时星辰的特定位置引为同时患病的原因，那便是极度的狂妄自大，因为在同一片天空下、同一地区，在同一时刻受孕和出生的事物种类繁多，状况各异，遭遇千差万别。但我们知道，双胞胎不仅行为方式不同，游历之地迥异，而且会患不同种类的疾病；对此，依我之见，希波克拉底会给出最简单的原因，即由于饮食和操练的差异——这并非源于体质，而是出于心志的倾向——他们可能在健康方面彼此有别。再者，波西多纽斯或任何主张星辰具有宿命影响的人，若不想在无知者心中灌输他们不明白的事，将极难就此自圆其说。至于他们企图从双胞胎出生间隔的极短时间中推衍出的那一点，即根据天穹中安放出生时刻标记之处，他们称之为\OriginalTerm{星位}{horoscope}的那一点，要么与双胞胎在性情、行为、习惯和运势上的差异极不相称，要么就与双胞胎地位高低——两者在这方面完全相同——相较而言极不相称；他们总将双胞胎间最大差异的原因归结于其中一人的诞生时刻；故此，若一人紧随另一人之后出生，以致星位毫无变化，我就要求两人在一切方面完全相似，而这在任何双胞胎中都不可能找到。但若第二人出生迟缓，致使星位改变，我则要求他们的父母不同，而这在双胞胎中是不可能的。\par

% structure: p0009 source=h2 target=subsection reason=relative-heading-rank; base=h2

\subsection{论占星家\Person{尼吉迪乌斯}在双胞胎出生问题上从陶轮引出的论证}

因此，那个关于陶轮的著名杜撰故事被提出来是毫无用处的，据说\Person{尼吉迪乌斯}在这问题上困惑不解时作了回答，并因此被称为\OriginalTerm{陶工}{Figulus}。他全力旋转陶轮，用墨水标记，以极快的速度两次击打，使得击打点似乎落在完全相同的部位上。当旋转停止后，发现他所作的标记在轮缘上相距不小。他说，同样，考虑到天球旋转的极大速度，即使双胞胎出生的间隔如同我给这轮子两次击打的间隔那么短，这一短暂的时间间隔也相当于天球上极远的距离。因此，他说，凡在双胞胎的习惯和运势上所能看到的任何差异，便由此而来。这个论证比那由陶轮转动而成的器皿更脆弱。因为若天穹中有如此大的意涵，无法通过星座观察来领会，以致在双胞胎中，遗产可能归于一人而不归另一人，那么，为什么对于非双胞胎的他人，他们竟敢在观察了其星座后，宣告那些属于无人能领会的秘密之事，并归因于各人出生的精确时刻呢？而今，若那些关于非双胞胎的出生时刻所作的预测，要以它们是根据对天穹更广阔空间的观察为依据来辩护，而分隔双胞胎出生、对应于天界微小部分的那些极短时刻，则关联着那些占星家通常不会被问及的琐事——因为谁会就何时该坐、何时该外出散步、何时及该食何物而去询问他们呢？——那么，当我们能指陈双胞胎在习惯、行为和命运上如此多的差异时，我们这样说又如何能自圆其说呢？\par

% structure: p0011 source=h2 target=subsection reason=relative-heading-rank; base=h2

\subsection{论性格行为极不相同的双胞胎\Person{厄撒乌}与\Person{雅各伯}}

\OriginalTerm{古圣祖}{ancient fathers}时代，就显赫人物而言，有一对孪生兄弟出生，他们出生时间极为接近，以致于先出生的那一位抓住了后出生者的脚跟。他们的生活和品行差异如此之大，行为举止如此不同，父母对他们的爱又如此有分别，以至于他们之间的对比甚至产生了相互敌意的反感。当我们说他们彼此如此不相像，以至于一人行走时另一人坐着，一人睡着时另一人醒着——这些差异被归因于那些无法被记录人出生时刻星位的人所察觉的微小空间片段，以便据此向\OriginalTerm{占星家}{mathematicians}求问——我们是否意指这些？这对孪生子中的一个长时间做雇工；另一个则从未服侍他人。其中一个得母亲喜爱；另一个则不然。其中一个失去了在他们民族中极为看重的尊荣；另一个则得到了它。至于他们的妻子、儿女和财产，我们又当怎么说呢？在所有这一切上，他们是多么不同啊！因此，如果这些事情与孪生子出生之间逝去的那些微小时间间隔有关，而不是归因于星座，那么为什么在别人的事例中，他们却从观察星座来预言这些事情呢？另一方面，如果这些事之所以被预言，是因为它们与可以观察和记录的时间间隔有关，而非与微不可察的时刻有关，那么在这件事上，\OriginalTerm{窑轮}{potters wheel}又有什么用处呢？无非是转动那些有着\OriginalTerm{泥土的心}{hearts of clay}的人，好使他们无法察觉\OriginalTerm{占星家}{mathematicians}言论的空洞罢了。\par

% structure: p0013 source=h2 target=subsection reason=relative-heading-rank; base=h2

\subsection{占星家如何被证实传授一门虚妄的学问}

\Person{希波克拉底}凭借医学的洞察力，怀疑一对病人是双胞胎，因为他观察到他们的疾病在同一时间发展至危象，又在同一时间消退——这些人不正充分驳斥了那些想要将实由体质相似所致之事归因于星辰影响的人吗？因为他们为何两人都患上同一种疾病，并且在同一时间，而不是按照出生的先后顺序一个接一个地生病？（因为他们当然不可能在同一时间出生。）或者，如果他们在不同时间出生这一事实绝不必然意味着他们必定在不同时间生病，那么他们为何坚持认为出生时间的差异是他们在其他方面差异的原因呢？为何他们能凭出生时间不同而在不同时间前往异域旅行、在不同时间结婚、在不同时间生育子女，并做许多其他事情，却因同样的理由不能在不同于对方的时间生病呢？因为，如果出生时刻的差异改变了星位，并导致所有其他方面的不相似，那么为何属于他们受孕的同时性却留在了他们的疾病发作中？或者，如果健康方面的命运牵涉于受孕时间，而其他事物的命运据称系于出生时间，那么他们就不应从察看出生的星象来预言任何关于健康的事情，如果受孕的时辰没有一并提供以便察考其星象的话。然而，如果他们说不察考受孕的时刻也能预言疾病的发作，因为这些由出生的时刻指示，那么他们又怎能根据两人中某一位的出生星位，告知他何时将生病，既然他的另一位有着不同出生星位的兄弟也必然在同一时间生病呢？再者，我要问，如果双胞胎出生的时间间隔如此之大，以致因着不同的星位，以及一切被赋予巨大影响的基本点之差，使得他们的星座产生差异，甚至从这种变化中得出命运的不同，那么既然他们不可能在不同时间受孕，这又如何可能呢？或者，如果在同一时刻受孕的两个人在出生方面能有不同的命运，那么为何在同一时刻出生的两个人不能在生与死方面有不同命运呢？因为如果两人受孕的那同一时刻并未阻止一人先于另一人出生，那么如果两人在同一时刻出生，又有什么能阻止他们在同一时刻死亡呢？如果同时受孕允许双胞胎在\Emph{胎中}受到不同的影响，为何同时出生不允许任何两个人在\Emph{世上}拥有不同的运数呢？这样一来，这门技艺的一切虚构——更确切地说是骗术——便会被一扫而空。这是多么奇怪的事：两个在同一时间，不，在同一时刻，在同一星位下受孕的孩子，会有不同的命运，将他们带至不同的出生时辰；而两个由不同母亲所生、在同一时刻、在同一星位下出生的孩子，却不能拥有不同的命运，以必然地引领他们走向不同的生与死的方式？他们尚在受孕时还没有命运，只因唯有出生才能拥有命运吗？那么，当他们说若能找出受孕的时辰，这些占星家便能预言许多事，这到底是什么意思呢？由此也产生了那个被一些人反复讲述的故事：某位贤哲挑选了一个时辰与他的妻子同房，以确保生出杰出的儿子。出自这一见解的，还有大占星家兼哲学家\Person{波西多尼乌斯}关于那对同时患病的双胞胎的答复，即：\Quote{这事之所以在他们身上发生，是因为他们是在同一时间受孕，并且在同一时间\Emph{出生}的。}他无疑加上了\Quote{受孕}一词，免得有人对他说他们不可能在同一时间出生，因为他知道无论如何他们必定是在同一时间受孕的；他如此是要表明，他不将两人同样且同时患上疾病的事实归于他们体质的相似作为近因，而是认为即使在健康相似这一点上，他们也是由星体的联系而连结在一起的。因此，如果受孕的时间对命运相似有如此大的作用，这些命运就不应因出生的情况而改变；或者，如果说双胞胎的命运因出生时间不同而被改变，那么我们为何不反而理解为：它们早已被改变，为的是让他们在不同时间出生呢？那么，当出生的次序能够改变他们在受孕时已有的命运时，难道不是生活在世上的人的意志改变了出生时的命运吗？\par

% structure: p0015 source=h2 target=subsection reason=relative-heading-rank; base=h2

\subsection{论不同性别的双胞胎。}

然而，即使在双胞胎的成孕中——这当然对双方同时发生——常常发生一个受孕为男，另一个为女。我认识一对不同性别的双胞胎。两人都还活着，正值青春；而且虽然他们在身体上，就性别差异所允许的程度而言，彼此相似，但在整个生活的范围和目标上却大不相同（考虑到男女生活之间必然存在的那些差异）——一个担任伯爵官职，几乎常年离家随军在国外服役，另一个从未离开过自己国家的土地，或她本乡的地区。更有甚者——而且如果星辰的命运值得相信，这就更难以置信了，不过如果我们考虑到人的意志和天主的白白恩赐，这就不足为奇了——他结了婚；她则是一位圣洁的童贞女：他生养了众多子女；她甚至从未结婚。然而，\OriginalTerm{星位}{horoscope}的效能岂不是极大吗？我想我已经说得足够表明那点的荒谬性了。但是，那些占星家说，不论星位在其他方面的效能如何，它对于出生无疑是有意义的。但为什么对于成孕也是那样呢？成孕无疑是以一次交合的行为发生的。而且，自然的效力如此之大，以致妇女一旦受孕，就不再会受孕。或者，他们也许在出生时改变了，或男变成女，或女变成男，因为星位的不同？然而，说某些\OriginalTerm{恒星}{sidereal}的影响力能够仅仅在身体上造成差异——比如，我们看到每年的季节随着太阳的接近与远离而循环，某些东西的大小因月亮的盈亏而增大或减小，例如海胆、牡蛎和海洋奇妙的潮汐——这虽然并非完全不合情理，但这并不意味\Emph{人的意志}就必须屈从于星辰的位置。不过，那些占星家当他们想要把我们的行为也束缚于\OriginalTerm{星座}{constellations}时，只是促使我们去探究，即使在这些身体方面，这些变化是否不可归因于星辰以外的其他原因。因为，有什么比性别更与身体密切相关呢？然而，在相同的星辰位置下，不同性别的双胞胎仍可被成孕。因此，还能断言或相信什么更大的荒谬，胜过说：成孕时对双方相同的星辰位置，竟不能使这个孩子与和她共有同一星座的兄弟具有不同性别，而他们出生时刻的星辰位置却能造成她与他之间在婚姻与圣洁童贞之间的巨大距离呢？\par

% structure: p0017 source=h2 target=subsection reason=relative-heading-rank; base=h2

\subsection{论结婚、栽种或播种的择日}

如今，有人会提出这样的说法吗：人们为某些特定行为选择特定日子，因而为他们的行为造成了某些新的宿命？比如，按照这种学说，某人本来注定了不该有一个显赫的儿子，反而该有一个可鄙的儿子，于是，因他是个有学问的人，便选择了一个时辰与妻子同房。这样，他便造成了一个先前没有的宿命，而从他自造的宿命中，有些事开始成了定数，而这定数并不包含在他出生时辰的宿命中。啊，多么愚蠢的荒谬！人们挑选日子结婚；而我认为，这是因为若不挑选日子，婚礼可能落在不吉利的日子，从而变得不幸。那么，星辰在出生时辰所命定的又成了什么呢？一个人能凭自己的选择改变那已为他决定的事，而他自己在选择日子时所决定的事，却不能被另一种力量所改变吗？因此，如果只有人，而不是天下万物，都受星辰影响，那么为什么人们还要选择某些日子作为种植葡萄、树木或播种谷物的吉日，另一些日子作为驯养牲口、或让雄畜与雌畜交配以使母牛母马受孕的吉日，以及其他类似的事呢？如果说某些被选定的日子会对此类事产生影响，是因为星辰排列根据时间的不同而支配着一切地上的物体，无论有无生命，那么请想一想，在同一瞬间出生、发生或起源的生物何其无穷无尽，而其结局却如此千差万别，这足以使任何小孩子都明白，那些关于日子的观察是荒唐可笑的。因为谁竟会疯狂到敢于断言，所有的树木、所有的草木、所有的走兽、蛇虫、飞鸟、鱼类和蠕虫，都会各有其出生或开始的时刻呢？然而，人们为考验占星家的技艺，惯常把他们带到哑口动物的生辰星座前，他们为了这一发现，在家中仔细记录动物的出生星座；人们推崇那些根据星座观察而断言这表示的是一头野兽而非一个人的诞生的占星家，胜过其他一切占星家。他们甚至敢说那是什么样的野兽，是产毛的兽，还是适于负重、适于耕田或适于守家的兽；因为人们也曾拿狗的宿命来考验占星家，而他们对此的回答则引来求问者的一片赞叹之声。他们这样欺骗人，使人以为当一个人出生时，所有其他生物的出生都中止了，以至于在他降生的同一时刻，在同一片天空下，甚至连一只苍蝇也不会诞生。如果这一点被承认适用于苍蝇，那么推论就不能停止，而必须从苍蝇上升到骆驼和大象。他们也不愿注意这样的事实：当选择了一个日子来播种田地时，那么多谷粒同时落入土中，同时萌芽，同时长出，同时成熟，同时结穗；然而，在所有那些同龄且可说是\Emph{\OriginalTerm{同胚生}{congerminal}}的穗子中，有些会被霉病毁掉，有些被飞鸟啄食，有些被人拔除。他们看到这些生物结局如此不同，怎能说它们各有不同的星座呢？他们是否会承认，为这类事选择日子是愚蠢的，并断言这些事不属星宿的定命范畴，而单单将人置于星辰影响之下——在这个世界上，唯有天主赐予了人自由意志？考虑到这一切，我们就有充分的理由相信，当占星家给出许多神奇的回答时，这不应归因于他们根据某种其实并不存在的技艺来标示和观察星宫图，而应归因于某些并非善类的神灵的隐秘启示，这些神灵忙于将关于星辰宿命影响的错误而有害的见解灌输进人的心灵，并在其中加以巩固。\par

% structure: p0019 source=h2 target=subsection reason=relative-heading-rank; base=h2

\subsection{第8章——论那些不以星辰方位而以取决于天主意志的因果连锁称为命运者。}

然而，对于这样一些人，他们所称的命运，并非指任何受造物受孕、出生或开始存在时星辰的排列，而是指那整个的因果连锁与环节，正是这因果连锁使每一事物成为现状，对于他们，我无需在纯属言辞的争议上费力争辩，因为他们将所谓的秩序和因果连锁归于至高天主的意志和权能，人们最正当、最真切地相信，这位天主在万事发生之前便预先知道一切，并且不会留下任何未安排的事；一切权能都出自祂，虽然并非所有的意志都出自祂。那么，下面这几行诗句——如果我没记错的话，作者是\Person{阿奈乌斯·塞内卡}——便证明了，他们所谓的命运，主要就是至高天主的意志，祂的权能不可抗拒地遍及万物：——\par

\Quote{至高的父，穹苍的统治者，\newline{}随祢喜悦引领我往何处；我必\newline{}迅即顺从，不加迟延，\newline{}看哪！我在这里。我速来遵行祢至尊的旨意；\newline{}倘若祢的命令拂逆我的心愿，我仍将\newline{}呻吟着随从祢，并以满心的\newline{}抗拒之苦去完成所派定的工作，\newline{}这乃出于邪恶；若为善者，\newline{}我本应虽艰难却仍以乐意的美德\newline{}承担并完成。\newline{}命运引领顺从的人前行；\newline{}但不顺从的人，则被拖着走。}\par

十分明显，在这最后一行诗里，他称那\Quote{命运}为他先前称为\Quote{至高父的意志}者，他说他随时预备顺从，以求被引领，而不是被拖曳，因为——正如\Quote{命运引领顺从的人前行，但不顺从的人，则被拖着走。}\par

以下\Person{西塞罗}译成拉丁文的\Person{荷马}诗句也支持这种见解：——\par

如此便是人之心灵，犹如那光明，\newline{}由\Person{朱庇特}亲自倾泻\newline{}荣耀普照丰饶的大地。\par

\Person{西塞罗}并不希望在此类问题上信赖诗人的观点；因为他说，当\OriginalTerm{斯多葛派}{Stoics}在主张\OriginalTerm{命运}{fate}的力量时，惯于引用这些来自\Person{荷马}的诗句，他并非在论述那位诗人的见解，而是在论述那些哲学家的见解，因为他们借这些在与命运相关的争论中所引用的诗句，最清楚地表明了他们所认定的命运是什么，因为他们把那位至高神称作\Person{朱庇特}，他们说，整个命运的链条都系于他。\par

% structure: p0026 source=h2 target=subsection reason=relative-heading-rank; base=h2

\subsection{第九章 论天主的预知与人的自由意志——驳\Person{西塞罗}的定义}

\Person{西塞罗}着手批驳\OriginalTerm{斯多葛派}{Stoics}的方式表明，他认为除非先摧毁了\OriginalTerm{占卜}{divination}，否则在论证中无法对他们取得任何成就。他试图通过否认任何有关未来事物的知识来达成此事，并竭力主张，无论在天主或人里面，都不存在这种知识，也没有对事件的预告。这样，他既否认了天主的\OriginalTerm{预知}{foreknowledge}，又试图用一些空洞的论证，以及对自己提出的某些极易被驳倒的\OriginalTerm{神谕}{oracles}，来推翻一切\OriginalTerm{预言}{prophecy}，即使是其清晰超乎光明的预言（尽管就连这些神谕他也并未驳倒）。\par

然而，在批驳那些\OriginalTerm{占星家}{mathematicians}的这些推测时，他的论证是得胜的，因为这些推测确实自我驳斥、自我摧毁。尽管如此，声称星宿具有命运影响的人，远较那些否认未来事件预知的人更可容忍。因为，既承认天主存在，同时又否认祂对未来事物有预知，这乃是最明显的愚妄。\Person{西塞罗}本人也看到了这一点，因此他试图主张圣经上的话所含的教义：\BibleQuote{愚妄的人心中说：没有天主。} 然而，他并未以自己的身份这样做，因为他看到这种观点多么令人厌恶和反感；因此，他在其\WorkTitle{论神性}一书中，让\Person{科塔}就此问题与\OriginalTerm{斯多葛派}{Stoics}争论，他宁愿将自己的意见给予\Person{卢齐利乌斯·巴尔布斯}，让他为斯多葛派的立场辩护，而不给科塔，因为科塔主张没有任何神性存在。不过，在他的\WorkTitle{论占卜}一书中，他亲自公开地反对未来事物预知的教义。但他这样做，似乎只是为了不承认\OriginalTerm{命运}{fate}的教义，从而破坏\OriginalTerm{自由意志}{free will}。因为他认为，一旦承认了对未来事物的预知，命运就必然随之而来，无法否认。\par

但是，任凭哲学家们的这些纠缠不清的辩论和争论如何进行吧，我们为了承认那位至高、真实的天主本身，便要承认祂的意志、至高权能和\OriginalTerm{预知}{prescience}。我们也不必害怕，免得我们不是因意志而做我们凭意志所做的事，因为祂，那预知绝对无误的天主，早已预知我们会做这事。这正是\Person{西塞罗}所畏惧的，所以他反对预知。\OriginalTerm{斯多葛派}{Stoics}也主张，并非万事都由\OriginalTerm{必然}{necessity}发生，尽管他们坚持万事都按\OriginalTerm{命运}{fate}发生。那么，\Person{西塞罗}对未来事物的预知到底害怕什么呢？无疑就是——如果一切未来之事都已被预知，它们就将按被预知的次序发生；如果它们按这次序发生，那么在天主那里就存在着已被预知的事物的某种秩序；如果存在着事物的某种秩序，那么也就存在着原因的某种秩序，因为没有一件事能发生，除非有某种\OriginalTerm{动力因}{efficient cause}在先。但如果存在着一个原因秩序，万事都按此秩序发生，那么，他说，所发生的一切就都是按命运发生的。但若是这样，那就没有什么在我们的权能之内，也就没有意志的自由；我们如果承认这点，他说，人类生活的整个秩序便被颠覆了。法律便成为无用的了。责备、称赞、规劝、劝诫都白费了；赏善罚恶的公正安排也荡然无存了。为了避免这些不光彩、荒谬且有害于人类的结果成为事实，\Person{西塞罗}选择拒绝未来事物的预知，并将虔敬的心灵逼到这两个选项之间，必须在两者中择一：要么有什么在我们权能之内，要么有预知——二者不能同时为真；但若肯定某一，便因此否定了另一。因此，他像一个真正伟大而睿智的人，且极为谨慎而灵巧地为人类好处着想的人那样，在这两者之中选择了意志的自由，并为了确立它而否认了未来事物的预知；就这样，他欲使人自由，却使他们成为亵渎者。然而，虔敬的心灵则两者都选择，两者都承认，并以虔敬的信德持守此二者。但如何可能呢？\Person{西塞罗}说：因为既已承认对未来的预知，就会有一连串的推论，最终导向这个结论：没有任何事能取决于我们的自由意志。进而言之，如果有什么事取决于我们的意志，我们就必须按照同样的推理步骤倒退回去，直到得出这样的结论：没有对未来的预知。因为我们按如下次序倒退回去：——如果有自由意志，万事就不都按命运发生；如果万事不都按命运发生，就不存在原因的某种秩序；如果不存在原因的某种秩序，那么在天主那里也不存在已被预知的事物的某种秩序——因为除非有动力因在先，事物不可能发生——，但是，如果不存在天主已预知的、固定而确定的原因秩序，万事便不能说都按祂所预知的那样发生。进而，如果万事并非完全按祂预知的那样发生，那么，他说，在天主内就没有对未来事件的任何\OriginalTerm{预知}{foreknowledge}。\par

面对理性那种亵渎而不虔敬的妄为，我们既要断言天主在万事发生之前便已知道一切，也要断言：我们凭自由意志所做的一切，都是我们明知且感觉到只因我们愿意才由我们完成的。但是，我们并不说万事皆因\OriginalTerm{命运}{fatum}而发生；相反，我们肯定没有一事是因命运而发生的。因为我们证明，\Quote{命运}这一名称，在谈论命运之人惯常的用法中，是指每个人受孕或出生时刻星辰的位置，那是一个毫无意义的词语；占星术本身即是虚妄。然而，那将最高效能归于天主意志的因果秩序，我们既不否认，也不以命运之名来称呼它；除非我们或许可将命运理解为从\OriginalTerm{说话}{fari}派生而来的\Quote{所说的事}；因为我们无法否认圣经上记载着：\BibleQuote{天主说过一次，我确实听到这两件事：威能属于天主；吾主，慈爱也非你莫属，因为你要照各人的行为予以赏报。} \BibleRef{咏 62:12-13} 而\Quote{祂说过一次}这句话，应理解为\Quote{\Emph{不可动摇地}}，也就是不可改变地祂说了，因为祂不可改变地知道一切将要发生的事，以及祂自己将做的一切事。因此，若是\Quote{命运}一词尚未被理解为另一种意义——我不愿人心不自觉地滑入那种意义——我们本可以在其源自\Quote{说话}一词的意义上使用它。但是，并不能由此推论说，尽管为天主而言，一切原因有某种确定的秩序，就必然没有任何事依赖于我们自由意志的行使；因为我们的意志本身就包含在那为天主所确知的因果秩序之中，并被祂的\OriginalTerm{预知}{praescientia}所涵盖；因为人的意志也是人的行为的原因；而预知一切事物原因的那一位，肯定不会在这些原因中忽略我们的意志。因为即使\Person{西塞罗}本人所做出的让步，也足以在这一论证中驳倒他自己。因为他虽说没有任何事是无因发生的，但并非所有原因都是命定的，存在着偶然因、自然因和自愿因，这能对他有什么帮助呢？他承认无论何事发生，必先有一原因，这就足够了。因为我们说，那些被称为偶然的原因，并非只是无因之事的虚名，而是潜伏的；我们或将其归于真天主的意志，或归于某种神灵的意志。至于自然因，我们绝不将它们与那位万有自然的创造者和塑造者的意志分离。至于自愿因，它们可归于天主、或天使、或人、或任何种类的动物——如果那些缺乏理性的动物按其本性寻求或躲避各种事物的本能行动，也可称为意志的话。当我论及天使的意志时，我指的是善天使的意志，我们称他们为天主的天使；或是恶天使的意志，我们称他们为魔鬼的天使，或恶魔。同样，论及人的意志，我指的是善人或恶人的意志。由此我们得出结论：一切发生之事，除非自愿因，即那属于\Quote{生命之灵}本性的原因，便没有动力因。因为风或气虽被称为\Quote{灵}，但由于它是形体，便不是生命之灵。因此，那赋予万物生命、是每一形体及每一受造之灵的创造者的生命之灵，正是天主自己，非受造的灵。在祂至高的意志中寓居着那作用于一切受造之灵的意志的能力：扶助善者，审判恶者，统御万有，将能力赐予某些，而不赐予另些。因为祂既是一切本性的创造者，也就是一切能力的赋予者，但并非一切意志的赋予者；因为邪恶的意志并非来自祂，它们与本性相悖，而本性却来自祂。至于形体，它们更服从于意志：有些服从于我们的意志，我指的是所有有死的活物的意志，但人之意志的支配甚于兽类。然而，一切形体尤其服从于天主的意志，一切意志也都服从于祂，因为它们除了祂所赋予的能力外，别无能力。因此，那创造而不被造的事物原因，就是天主；而其他一切原因，既创造也被造。一切受造之灵即是如此，尤其是理性的灵。因此，物质原因更可谓被造而非创造，不应列入动力因，因为它们只能借灵之意志而有所作为。那么，既然一个为天主的预知所确知的因果秩序，如何能迫使不存在任何依赖我们意志的事呢？因为我们的意志本身在因果秩序中占据极重要的位置。因此，\Person{西塞罗}与那些称此因果秩序为命定的，或更确切地将此秩序本身命名为命运的人争论；对此我们深恶痛绝，尤其因为这个词语，人们已习惯将其理解为不实之事。但是，他否认一切原因的秩序为最确定、并对天主的预知完全清晰，我们比斯多葛派更憎恶他的这种观点。因为他要么否认天主存在——事实上，他以假托角色在其著作\WorkTitle{论神性}中竭力如此——要么，若他承认天主存在，却否认祂预知未来之事，这不就是\BibleQuote{愚妄的人心中说：没有天主}吗？\BibleRef{咏 14:1} 因为不预知一切未来之事的，便不是天主。因此，我们的意志也具有天主所意愿并预知其当有的能力；故此，它们所拥有的任何能力，都在极为确定的界限之内；它们所将行的一切，都最为确实地将要发生，因为那位预知无误的，已预知它们将有能力行此事，并必将行此事。因此，若我定要将命运之名用于任何事物，我宁愿说命运属于双方中较弱的一方，意志属于较强的一方，即支配对方的一方，而不是说我们意志的自由被斯多葛派以他们特有的非常规用词称为\Emph{命运}的那种因果秩序所排除。\par

% structure: p0031 source=h2 target=subsection reason=relative-heading-rank; base=h2

\subsection{第十章——我们的意志是否受必然性支配。}

因此，我们也不必害怕那种\OriginalTerm{必然性}{necessity}；正是出于对这种必然性的畏惧，\OriginalTerm{斯多葛派}{Stoics}努力对事物的原因作出区分，以便能将某些事物从必然性的支配下解救出来，同时让另一些事物受其支配。在他们不愿使之受必然性支配的事物中，他们列入了我们的\OriginalTerm{意志}{will}，因为他们知道，如果意志受必然性支配，它便不会是自由的。因为如果我们将那不在我们权能之内、即使我们不愿意也仍能产生其效果者称为必然性——例如死亡的必然性——那么显而易见，我们赖以正直或邪恶地生活的意志，并不处于这样的必然性之下；因为我们做许多事，如果我们不愿意，我们肯定不会去做。这在意志行为本身中尤其真实——因为如果我们愿意，意志便存在；如果我们不愿意，它便不存在——因为如果我们不愿意，我们就不会愿意。但是，如果我们这样定义必然性：根据它，我们说任何事物必然具有某种性质，或以某种方式被做成，那么我不知道为什么我们还要害怕这种必然性会夺走我们\OriginalTerm{意志的自由}{the freedom of our will}。因为如果我们说天主必然永远活着、预知万事，我们并没有将天主的生命或天主的\OriginalTerm{预知}{foreknowledge}置于必然性之下；同样，当我们说他不能死、不能犯错时，他的权能也并未减少——因为这对他而言是如此不可能，以至于如果对他而言是可能的，他的权能便会较小。但毫无疑问，他确实被称为\OriginalTerm{全能的}{omnipotent}，尽管他既不能死，也不能犯错。因为被称为全能，是由于他做他所愿意的事，而不是由于他承受他所不愿的事；因为如果这种事临到他，他便绝不是全能的了。因此，正因为他是全能的，他不能做某些事。同样地，当我们说，当我们愿意时，我们必然通过\OriginalTerm{自由选择}{free choice}而愿意，我们既肯定了不容置疑的真理，也并没有因此将我们的意志置于一种摧毁自由的必然性之下。因此，我们的意志\Emph{作为意志而存在}，并且亲自做我们因愿意而做的一切事，这些事如果我们不愿意，就不会被做。但当一个人因另一人的意志而遭受他不愿意的事时，即使在这种情况下，意志仍保持其本质的有效性——我们指的不是施加痛苦者的意志，因为我们将之归于天主的权能。因为如果一个意志仅仅存在，却不能做它愿意的事，它就会被一个更强有力的意志所压倒。若非先有意志存在，即不是另一方的意志，而是那愿意却无法实现所愿者的意志，这事也不会发生。因此，凡一个人违背自己意志所受的痛苦，他不应归因于人的意志、或天使的意志、或任何受造之灵的意志，而应归因于赐予意志权能的那一位的意志。所以，并非因为天主预知了在我们意志权能之内的事，因此在我们的意志权能之内就什么都没有了。因为那预知这事的那位，并非预知虚无。再者，如果那位预知在我们意志权能之内之事者，并非预知虚无，而是预知了某事，那么毫无疑问，即使他预知了，在我们意志权能之内仍有某事。因此，我们绝不被迫，要么在保留天主的预知时，取消意志的自由；要么在保留意志的自由时，否认他预知未来的事，这是不虔敬的。但我们两者都接受。我们忠信而真诚地承认两者：前者，为使我们能正确地相信；后者，为使我们能正确地生活。因为凡对天主没有正确信仰的人，便生活邪恶。因此，我们断不可为了维持我们的自由，而否认那位藉其帮助我们已是或将获自由者的预知。所以，法律的制定、谴责、劝勉、赞扬和责备的运用都不是徒然的；因为他也预知了这些，并且它们大有功效，正如他所预知的那样功效显著。同样，祈祷也有功效，为求得他所预知将要赐给祈祷者的事物；并且公义地，善行已注定有赏报，罪恶有惩罚。因为一个人犯罪，并非因为天主预知他会犯罪。不，无可怀疑的是，当一个人犯罪时，正是他自己在犯罪，因为那预知无误者，预知犯罪的不是\OriginalTerm{命运}{fate}、\OriginalTerm{运气}{fortune}或其他什么东西，而是那人自己会犯罪，而如果他不愿意，他便不会犯罪。但如果他将不愿意犯罪，甚至这也为天主所预知。\par

% structure: p0033 source=h2 target=subsection reason=relative-heading-rank; base=h2

\subsection{第十一章——论天主在统摄万物的法则中的普遍眷顾。}

因此，至高而真天主，偕同祂的圣言与圣神（三者而成为一体），唯一全能的天主，是一切灵魂与一切形体的创造者与形成者；凡藉真理而非虚幻而享有真福的人，都是因着祂的恩赐而获此幸福。祂使人类成为由灵魂与身体构成的理性动物；当人类犯罪之后，祂既未容其不受惩罚而逍遥法外，也未将其遗弃而毫无仁慈。祂将存在普遍赐予善者与恶者，与石头同具；将生命滋养之能普遍赐予万物，与树木同具；将感官生命普遍赐予，与禽兽同具；将理智生命仅赐予天使。一切形式、一切种类、一切秩序都来自祂；尺度、数目、重量都来自祂；自然界中凡存在的事物，无论其类别与价值如何，都来自祂；形式的种子和种子的形式，以及种子与形式的运动，也都来自祂。祂还赋予肉体以起源、美丽、健康、繁殖能力、肢体的配置，以及各部分间有益健康的和谐协调；祂还赐予无理性的灵魂以记忆、感觉与欲望，但对于理性的灵魂，则除此之外还赐予了理智与意志。祂未曾使天地、天使与人类，不用说；甚至连最微小、最卑贱的动物的内脏、飞鸟的羽毛、植物的小花、树上的叶子，也没有一处没有各部分之间的和谐与某种相互的和平——这样一位天主，我们绝不能相信祂会将人类的一切王国、统治与奴役，置于祂的天主眷顾的法则之外。\par

% structure: p0035 source=h2 target=subsection reason=relative-heading-rank; base=h2

\subsection{第十二章：古代罗马人以何品德，虽未朝拜真天主，却堪受真天主扩展其帝国}

因此，让我们继续探讨罗马人的哪些品德，得到了真天主的协助（在祂手中掌握着地上万国），以使帝国得以扩展，以及祂这样做的原因。并且，为了把这个问题讨论得更清楚，我们写了前几卷，指出他们所认为应该用那些琐碎而愚蠢的礼仪去崇拜的神明，对于此事毫无作用；我们也已完成了本卷前文的论述，驳斥了命运论，免得可能已经有人不再将罗马帝国的扩张与维持归于敬拜那些神明，却可能将其归于某种命运，而非至高天主的强大旨意。所以，古代和早期的罗马人，尽管他们的历史记载显示，与所有其他民族一样，除希伯来人外，都敬拜假神，不向天主而向魔鬼献祭牺牲，然而，他们的历史家却这样称赞他们：他们是\Quote{渴求赞美，挥霍财富，渴望伟大的光荣，满足于微薄的产业。}他们对光荣有着最为炽烈的爱：为光荣他们愿生，为光荣他们不避死。对这一件事的热烈激情，压抑了其他一切欲望。最终，他们自己的国家，因为觉得臣服于人似乎不够光荣，而统治、号令他人才是光荣，因此他们首先热诚地渴望自由，继而又渴望成为霸主。正因如此，他们不甘忍受君王的统治，而将政权交到两位任期一年的官员手中，称他们为执政官，而非国王或主宰。但君王的排场似乎不符合治理者\OriginalTerm{治理者}{regentis}的行政管理，也不符合谋议者（即为公共利益着想者）\OriginalTerm{谋议者}{consulentis}的仁慈，反倒更适合主宰者\OriginalTerm{主宰者}{dominantis}的傲慢。因此，\Person{Tarquin}王被放逐之后，设立了执政官制度，正如前面提到的同一位作者在赞扬罗马人时所说的那样：\Quote{国家在获得自由之后，以惊人的速度成长起来，因为对光荣的渴望占据了它。}那么，正是这种对赞美的热切和对光荣的渴望，成就了那些按人间的判断而言，无疑是值得称赞和辉煌的许多奇功伟业。同一位\Person{Sallust}称赞他同时代的伟人\Person{Marcus Cato}和\Person{Caius Cæsar}，说很长一段时间国家没有德行出众的人，但在他的记忆中却有过这两个德行卓著、志趣各异的人物。在称颂凯撒的话语中，他特别提到，凯撒渴望有一个庞大的帝国、一支军队和一场新的战争，好让自己有施展才华与德行的舞台。因此，那些英雄气概的人常常祈祷，愿\Person{Bellona}拨动悲惨的民族投入战争，用她那血腥的鞭子抽打他们，使他们躁动不安，以便有机会展示他们的勇武。这，正是那种对赞美的渴求和对光荣的渴望所造成的结果。所以，首先由于对自由的热爱，随后又出于对统治权的热爱，并因为对赞美和光荣的渴望，他们成就了许多伟业；他们最卓越的诗人证实了这些动机都在激励着他们：\par

\Person{Porsenna}傲气冲天，\newline{}命\Place{Rome}为\Person{Tarquin}打开城门；\newline{}以武力围困城池，\newline{}\Person{Æneas}的子孙坚守阵地，以求胜利。\par

那时他们最大的抱负，或是英勇战死，或是自由地活着；但获得自由之后，对光荣的渴望如许强烈地占据了他们，以至于仅自由还不够，还必须追求统治权；而这些大志，正是同一位诗人借\Person{Jupiter}之口所说的：\par

\Person{Juno}自己，尽管她狂怒的警报\newline{}曾使海洋、大地和天空都仓皇备战，\newline{}如今也会化阴郁的怒容为微笑，\newline{}与我争相热情为\newline{}罗马的子孙——那穿长袍的民族——加冕。\newline{}这是我的旨意。终有一天，\newline{}当罗马的伟业世代延续，\newline{}古老的\Person{Assaracus}的子孙\newline{}将摆脱密尔弥冬人，\newline{}统治\Place{Phthia}与\Place{Mycenæ}，\newline{}并迫使\Place{Argos}俯首称臣。\par

这些事，实际上\Person{维吉尔}让\Person{朱庇特}预言为未来的事，而事实上他只是在自己心中回顾那些已经发生的事，并将其视为当前的现实。但我提及它们，是为了表明罗马人除了自由之外，如此高度重视统治，以至于统治在他们最赞颂的事物中占有一席之地。因此，那位诗人偏爱罗马人特有的技艺，即统治与命令、征服战胜列国的技艺，胜过其他民族的技艺，他说道：\par

\Quote{他人或许以更幸运的妙手，\newline{}从青铜或石头中唤醒面容，\newline{}为疑案辩护，描绘天穹，\newline{}道出行星何时升落；\newline{}但罗马人，你的使命是统治\newline{}万邦四境；\newline{}这才是你的才能：将和平的法则\newline{}强加于降敌，\newline{}怜悯卑微者，\newline{}并粉碎骄傲的子孙。}\par

他们越是少耽于逸乐，少在贪求积聚财富中消磨身心，并通过向不幸的公民勒索财富、将其挥霍于卑贱的优伶而败坏道德，便越是熟练地运用这些技艺。因此，当\Person{撒路斯提乌斯}著文、\Person{维吉尔}咏唱这些事时，大量涌现的那些卑劣之徒，并非通过这些技艺追求荣誉和光荣，而是通过奸诈和欺骗。所以同一位作者说：\Quote{起初搅动人心的，主要是\OriginalTerm{野心}{ambition}而非贪婪；然而，这一恶习更接近德行。因为光荣、荣誉和权力，善人与卑劣者同样渴望；但前者，}他说，\Quote{循正道奋力追求它们，而后者，因对\OriginalTerm{美善的技艺}{good arts}一无所知，便以诡诈欺骗来获取它们。}所谓借着美善的技艺寻求光荣、荣誉和权力，即借着德行，而非阴险的诡计来寻求；因为善人与卑劣者同样渴望这些，但善人力图循正道获取它们。那正道即是德行，他沿着它奋力前进，如同朝向占有的目标——即光荣、荣誉和权力。这种情感深植于罗马人的心灵，甚至从他们神明的庙宇也可看出；因为他们在极近处修建了德行女神和荣誉女神的庙宇，将天主的恩赐当作神明来敬拜。由此我们可以理解，他们当中的善人认为德行的目的是什么，以及他们最终将其归于何处，即荣誉；至于恶人，他们虽渴望荣誉，却毫无德行，并企图通过诡诈欺骗来占有它。\Person{加图}受到了更高级别的赞誉，因为有人论及他说：\Quote{他越不追求光荣，光荣越发跟随着他。}我们说更高级别的赞誉；因为罗马人所热烈渴望的光荣，是人们对人良好评价的判断。因此，德行更为优越，它只满足于自己良心的见证，而不需要人的任何评判。所以宗徒说：\BibleQuote{因为我们有这样的光荣：就是我们良心的见证。} \BibleRef{格后 1:12} 在另一处他又说：\BibleQuote{但各人应当察验自己的行为；这样，他的光荣就只在自己，而不在别人了。} \BibleRef{迦 6:4} 因此，他们为自己所渴望，而善人又试图借着美善的技艺获取的光荣、荣誉和权力，不应由德行去追求，而应由它们去追求德行。因为没有真正的德行，除非它指向那目的，即人的至高、终极的善。所以，就连加图所寻求的荣誉，他也不该去求，而国家本应因他的德行，不待他请求便授与他。\par

然而，在当时的两位伟大罗马人中，\Person{加图}是德行最接近德行真义的一位。因此，让我们参考加图自己的意见，来发现他对当时及从前国家状况形成了怎样的判断。\Quote{我不认为，}他说，\Quote{我们的祖先是通过武力使共和国由小变大的。若果真如此，我们今天的共和国本应远比他们那时繁荣，因为我们的盟友和公民人数要多得多；此外，我们拥有的武器和马匹也比他们那时丰富得多。但使他们伟大的不是这些，而我们根本没有这些：国内的勤勉，对外的公正统治，议事时自由的心灵，既不沉迷于罪恶，也不沉迷于情欲。与此相反，我们有的是奢侈和贪婪，国家贫困，公民富有；我们赞美财富，追求懒惰；善与恶之间不加区分；德行的一切奖赏都被阴谋手段攫取。毫不奇怪，当每个人都只为自己谋利益，当你在家是享乐的奴隶，在公共事务上是金钱和偏私的奴隶时，毫不奇怪，那没有防卫的共和国会遭受攻击。}\par

凡听见\Person{加图}或\Person{撒路斯提乌斯}这些话的人，大概会以为，对古代罗马人的这种赞扬适用于他们全体，或至少适用于他们中的大多数。事实并非如此；否则，\Person{加图}本人所写、且我在本作品第二卷中引述的那些事，就不会是真的了。他在那段话说，即便从国家初创时起，更有权势者就犯下了不义之行，导致平民与父老分离，此外还有其他内部纷争；唯一存在公正温和治理的时期，是在诸王被逐之后，且仅限于他们仍有理由畏惧\Person{塔克文}，并因他而对\Place{埃特鲁里亚}进行那场艰苦战争之时；但此后，父老们就像对待奴隶一样压迫平民，像国王们那样鞭打他们，把他们逐出土地，并排除其他所有人，将政府独自掌握在自己手中。而当父老们想要统治、平民们不愿服役之时，是第二次布匿战争结束了这些纷争；因为巨大的恐惧再次开始压迫他们不安的心灵，以另一种更大的焦虑将他们从那些纷扰中拉回，并带回公民和睦。然而，那时所成就的伟大事业，乃是通过少数按他们自己的方式算为良善之人的治理而完成的。正是由于这少数良善之人的智慧和远见，先使共和国能够忍受并减轻这些罪恶，它才日益强大。这位历史学家还断言，当他阅读和听闻罗马人民在和平与战争中、在陆地和海洋上许多辉煌的成就时，他想要理解究竟是什么特别支撑着这些伟业。因为他知道，罗马人常常以少量兵力与敌人庞大的军团作战；他也知道，他们以微薄资源与富有的国王进行战争。他说，经过深思熟虑，他认为很明显，少数公民卓越的德行成就了一切，这也解释了贫穷如何战胜财富，寡少如何战胜多数。但是，他补充道，在国家被奢侈和怠惰败坏之后，共和国仍能凭借自身的伟大，承受其官员和将领的恶习。因此，即使\Person{加图}的赞美也只适用于少数人；因为只有少数人拥有那种德行，引导人们通过真正的途径——即通过德行本身——追求光荣、尊荣和权力。\Person{加图}所说的那种在家勤劳，乃是出于渴望充实公库，即使结果是家中贫穷；所以，当他论到道德腐败所致的弊端时，他倒转过来说：\Quote{国家贫穷，家中富足。}\par

% structure: p0045 source=h2 target=subsection reason=relative-heading-rank; base=h2

\subsection{第十三章 论对赞美的爱，它虽是一项恶习，却被视为德行，因为它能遏制更大的恶习。}

因此，当东方各国长久以来声名显赫之后，\Person{天主}乐意兴起一个西方的帝国，它在时间上虽晚，但在疆域和伟大上却更为辉煌。而为了克服其他国家中存在的严重罪恶，祂特意将帝国交给这样一些人：他们为了荣誉、赞美和光荣，为国家悉心谋划，在国家的光荣中寻求自己的光荣，并毫不犹豫地将国家的安全置于自己的安全之上；他们为了这一个恶习，即对赞美的爱，而压制了对财富的贪欲和许多其他恶习。最正确的认识是承认对赞美的爱本身也是恶习；诗人\Person{贺拉斯}的观察也未忽略这一点，他说：\par

\Quote{你因野心而膨胀？听从劝告：\newline{}那本书若读三遍，便会让你松快。}\newline{}同一位诗人在一首抒情歌中，出于抑制统治欲的愿望，这样说：\par

\Quote{统治雄心壮志的精神，你便拥有\newline{}一个比连接遥远的\Place{加的斯}与\Place{利比亚}更广阔的王国，\newline{}如此}\par

\Quote{使\Place{迦太基}人都在服役之下。}\par

然而，那些抑制卑劣情欲的人，不是靠着由虔诚信仰获得的\Person{圣神}的能力，也不是靠着对可理解之美的爱，而是靠着对人间赞美的渴望——或者至少靠着这种对赞美的爱而更好地抑制它们——他们确实还不是圣洁的，只是较少卑劣而已。连\Person{西塞罗}也无法隐瞒这一事实；因为在他所写的\WorkTitle{论共和国}同一书中，谈及国家领袖的教育时，他说领袖应当从光荣中得到滋养，并进而说他们的祖先出于对光荣的追求，行了许多奇妙辉煌的事。因此，他们不但不抵制这种恶习，反而认为应该激发并点燃它，以为这会对共和国有利。但即使在他的哲学著作中，\Person{西塞罗}也没有掩饰这种有毒的观点，反而比白昼更清楚地承认了它。因为当论及那些着眼于真正的善、而非虚妄地追求人间赞美的研究时，他提出了以下普遍而一般的陈述：\par

\Quote{荣誉滋养技艺，光荣激励人们从事研究；而那些普遍不被看重的追求则总是被人忽视。}\par

% structure: p0052 source=h2 target=subsection reason=relative-heading-rank; base=h2

\subsection{第十四章 论根除对人间赞美的爱，因为义人的一切光荣都在于\Person{天主}。}

因此，毫无疑问，抵抗这贪欲远胜过屈服于它，因为人越洁净，不受这玷污，便越肖似天主；并且，虽然这恶习不能从心中完全根除——因为它甚至不断诱惑那些在德行上大有进步的人——但无论如何，对光荣的贪求总要被对仁义的喜爱所超越，如此，倘若在任何地方见到 \Quote{被忽略、普遍遭摈弃的事物}，如果它们是好的，是正当的，那么即使是对世人赞美的贪爱，也会羞愧地屈服于对真理的爱。因为这恶习与虔诚的信德如此为敌，如果心中贪爱光荣胜过敬畏或爱慕天主，以致主说：\BibleQuote{你们既然彼此寻求光荣，而不寻求从独一天主来的光荣，你们怎么能信呢？} \BibleRef{若 5:44} 同样，关于那些信了祂，却不敢公开承认祂的人，圣史说：\BibleQuote{他们贪爱世人的赞美，胜于天主的赞美；} \BibleRef{若 12:43} 对于这些，圣宗徒们却没有如此行；他们在那些地方宣讲基督之名时，这名不仅遭到摈弃，因而被轻忽——正如\Person{西塞罗}所说：\Quote{那些普遍遭摈弃的事物常被轻忽}——而且甚至被人极度憎恶，他们却谨守从那位既是善师、也是心灵的医师那里所听到的：\BibleQuote{凡在人前否认我的，我在我天上的父前，并在天主的使者前，也必否认他，} \BibleRef{玛 10:33} 在詈骂、侮辱、最惨烈的迫害和残酷的刑罚中，他们并没有因世人愤慨的喧嚷而不敢宣讲人类的救恩。当他们实践并宣讲神圣之事，度神圣的生活，如同攻克了铁石心肠，并将仁义的平安注入其中时，在基督的教会中便有极大的光荣随他们而来；但他们并不安于此，以之为德行的终向，反而将这光荣归于天主的光荣——正是靠天主的恩宠，他们才得以成为这样的人物——并力求用这同一火焰，点燃那些他们为其益处着想的人的心灵，使他们爱慕那位能使他们成为像他们自己一样的天主。因为他们的师傅曾教导他们，不要为了世人的光荣而追求善行，说：\BibleQuote{你们小心，不可在人前行你们的仁义，为叫他们看见；不然，你们在天之父前，就没有赏报了。} \BibleRef{玛 6:1} 但祂又怕他们领会错了，因怕讨人喜欢而隐藏自己的善，以致减少效益，于是指明他们应当为什么目的而显扬善行，说：\BibleQuote{你们的光也当在人前照耀，好使他们看见你们的善行，光荣你们在天之父。} \BibleRef{玛 5:16} 请注意，不是\BibleQuote{为叫他们看见你们}，即让他们的目光都投向你们——因为你们自己原是一无所是——而是\BibleQuote{为叫他们光荣你们在天之父}，使他们定睛注视祂，从而变得像你们一样。殉道者们追随了这些教导，他们无论在真正的德行——因为在于真虔诚——还是人数众多上，都远胜\Person{斯凯沃拉}、\Person{库尔提乌斯}和\Person{德西乌斯}之流。然而，那些罗马人是在一个属世之城，他们为城邦承担一切职务的目的，乃在于它的安全，以及一个王国：不是在天上的，而是在地上的——不是在永生的领域，而是在死亡与承替的领域，死人不断为垂死之人所接替——他们除了爱光荣，借以希望死后还能活在赞美者的口中外，又能爱什么呢？\par

% structure: p0054 source=h2 target=subsection reason=relative-heading-rank; base=h2

\subsection{第十五章 论天主赐给罗马人德行的现世赏报。}

现在，因此，对于那些天主不曾预定赐予永生、与祂的圣天使同在祂自己的天上之城的人——而真实的虔诚，唯独向真天主献上那希腊人称为 \OriginalTerm{钦崇}{\GreekText{λατρεία}} 的敬礼，才能引领人进入那团体——假若天主连那最卓越帝国的世间光荣也不赐给他们，那么对于他们借以追求如此伟大光荣的良好技艺——即他们的德行——便未给予任何赏报了。因为对于那些貌似行善，以求获得世人光荣的人，主也说：\BibleQuote{我实在告诉你们，他们已经获得了他们的赏报。} \BibleRef{玛 6:2} 同样，这些人为了共和国，轻看自己的私事；为了国库，抵制贪婪；以自由的精神为国家的益处筹谋；既不沾染法律所定的罪行，也不放纵情欲。他们借着这一切行为，如同沿着真道，奋力追求尊荣、权力和光荣；他们受几乎万民的尊敬；他们把帝国的法律加于许多民族；直至今日，无论文学还是历史中，他们仍在几乎万民中享有光荣。他们没有理由抱怨至高真天主的公义——\BibleQuote{他们已经获得了他们的赏报。}\par

% structure: p0056 source=h2 target=subsection reason=relative-heading-rank; base=h2

\subsection{第十六章 论天上之城的圣洁公民所得的赏报，罗马人的德行榜样对他们有益。}

但圣人的赏报却远不相同，他们甚至在此世为那座为爱恋此世者所憎恶的\OriginalTerm{天主之城}{City of God}忍受了凌辱。这座城是永恒的。在那里没有人出生，因为没有人死亡。那里有真正而圆满的真福——不是一位女神，而是天主的恩赐。我们从那里领受了信德的凭证，在此流徙途中，我们为它的美丽而叹息。那里没有太阳既照善人也照恶人，而是\OriginalTerm{义德的太阳}{Sun of Righteousness}只护佑善人。那里无需付出极大辛劳，通过在国内忍受匮乏来充实公共府库，因为那里有真理的公库。因此，罗马帝国的权势和荣耀之所以如此显著地扩展，不仅是为了酬报罗马城的公民，也是为了那永恒之城的公民，在他们的尘世旅居期间，勤奋而清醒地默观这些榜样，看出如果地上之城因其公民为了人间的荣耀而备受爱戴，那么他们为了永生，该是何等地爱慕天上的国度。\par

% structure: p0058 source=h2 target=subsection reason=relative-heading-rank; base=h2

\subsection{第十七章——罗马人发动战争得到什么利益，他们为被征服者的福祉贡献了多少}

因为就这短短几日便耗尽并终结的凡人的生命而言，一个人生活在谁的统治之下，只要统治者不强迫他行不虔不义之事，又有何关系呢？罗马人对于那些被他们征服并强加其法律的民族，除了在战争中伴随着大屠杀之外，又有何伤害呢？如果当初是征得各民族同意而为之，那会成功得多，但也就不会有征服的荣耀了，因为罗马人自己也未能免于他们加诸他人的法律。倘若这事在没有\OriginalTerm{马尔斯与贝娄娜}{Mars and Bellona}的情况下成就，以致没有胜利的位置，无人战斗便无人征服，那么罗马人与其他民族的境况岂不相同，特别是如果当初就即刻做成后来以最人道、最受欢迎的方式成就的事，即所有隶属罗马帝国的人都获得罗马公民权，且如果那原本仅为少数人的特权成为所有人的特权，唯一条件是那些没有自己土地的卑微阶级应由公共开支供养——这种供养税，由他们作为共和国的一员，出于衷心同意，交到贤明管理者手中，会比从被征服者那里强行榨取时优雅得多。因为我看不出，一些人征服而另一些人被征服，对人的安全、良好道德，当然还有尊严，有何益处，除非它带给他们那种最疯狂的人间荣耀的虚夸，那些对此炽烈渴求、热衷于战争的人，正\Quote{已经获得了他们的赏报}。他们的土地难道不纳贡吗？他们有什么学习的特权，是别人不得享有的吗？在其他国家，难道没有许多元老连罗马都没见过吗？除去外表，所有人终究不过是人。但即使这时代的乖谬允许所有更优秀者受到比他人更高的尊荣，人间的荣耀也不应被看重，因为它不过是无分量的轻烟。然而，即使在这些事上，我们也要善用天主的仁善。让我们思量，他们为了人间的荣耀，轻视了多少伟业，忍受了多少苦难，克制了多少欲望，他们如此德行，仿佛理应得到那荣耀作为赏报；愿这有助于我们抑制骄傲，以致既然那应许我们统治的\OriginalTerm{天主之城}{City of God}远胜此城，就如天远超过地，永生远胜暂世的欢乐，坚实的荣耀远胜空虚的赞美，或天使的团体远胜凡人的团体，那造了日月之主的荣耀远胜日月光辉，那么这伟大国度的公民，若为获得它而做了些许善工或忍受了些许邪恶，便不该自以为成就了什么大事，因为那些人为了这已然获得的地上之城，做出了如此伟业，忍受了如此苦难。尤其要思量这一切，因为那聚集公民进入天上国度的赦罪之恩，有其某种影子般的相似处可见于\OriginalTerm{罗穆卢斯的避难所}{asylum of Romulus}，那里逃脱各种罪罚的人聚集起来，形成那建邦立国的群众。\par

% structure: p0060 source=h2 target=subsection reason=relative-heading-rank; base=h2

\subsection{第十八章——基督徒若曾为永恒天乡做了何事，便绝不应自夸，因为罗马人为人间荣耀与地上之城做出了如此伟业}

那么，倘若为了这座地上之城，\Person{布鲁图斯}竟能处死自己的儿子——这是天城不强迫任何人作出的牺牲——那永恒而天上的城鄙视这世间一切欢乐，无论多么愉悦，又算得什么伟大的事呢？但无疑，处死自己的儿子，比为天国所要求做的事更难，即把那些原以为是为儿女积攒储存的东西分施给穷人，或者，若有任何逼迫我们如此行的试探临到，为了信仰和义德，就舍弃它们。因为使我们或我们儿女幸福的，并不是地上的财富；这些财富要么在我们有生之年失去，要么在我们死后为不知是谁、甚至可能是我们所不愿的人拥有。唯有天主使人幸福，祂是心灵的真财富。但至于\Person{布鲁图斯}，就连那位歌颂他的诗人也证实，杀死自己的儿子是他不幸的缘由，因为他说，\par

唤来他自己的叛逆骨肉\newline{}为受威胁的自由而流血。\newline{}不幸的父亲！无论\newline{}后世如何评判这行为。\newline{}但在接下来的诗句中，他安慰他的不幸，说：\newline{}\Quote{对国家的爱压倒一切。}\par

正是这两个因素，即自由和对人的赞美的贪求，迫使罗马人去成就可敬的事业。因此，如果会死之人为了自由，以及必朽者所追求的属人的赞美，父亲可以杀死儿子，那么我们为了真正的自由——这自由使我们脱离罪恶、死亡和魔鬼的辖制——不是为了贪求人的赞美，而是为了热心逃避（世俗），不是逃避\Person{塔克文王}，而是逃避魔鬼及其首领，我们就算不杀死自己的儿子，而把基督的穷人视为我们的儿子，这又算什么大事呢？如果另一位罗马首领，绰号\Person{托尔夸图斯}的，杀死了自己的儿子，不是因为儿子反对国家，而是因为儿子被敌人挑战，出于年轻冲动，虽为国而战，却违背了作为统帅的父亲所下的命令；他这样做，尽管他的儿子获胜了，因为担心藐视权威的榜样比杀敌的荣耀为害更大——我说，如果\Person{托尔夸图斯}这样做了，那么，为天上国度的法律，轻看一切世间的美物，而这些美物远不如儿子那么被爱的人，又有什么可夸耀的呢？如果\Person{富里乌斯·卡米卢斯}，虽曾被嫉妒他的人定罪，尽管他曾为国人卸下了最苦毒仇敌\Place{维爱人}的轭，后来又从\Place{高卢人}手中解救了忘恩负义的祖国，只因没有别处可以给他更好的机会过荣耀的生活——如果\Person{卡米卢斯}这样做了，那么，那人在教会中受到属血肉的仇敌极其严重且羞辱的伤害，却没有投靠异端敌人，也没有自己发起异端反对她，反而竭尽所能保护她免受异端者最有害的邪恶，既然没有另一个教会，我且不说在其中可以过荣耀的生活，而是可以获得永生，他有什么大功值得被高举呢？如果\Person{穆奇乌斯}，为了使与正在以极惨重战争压迫罗马人的\Person{波尔塞纳王}和好，当他没能杀死\Person{波尔塞纳}，而是误杀了别人时，他伸出右手放在火热的祭坛上，说有许多像他这样的人阴谋要害他，以致\Person{波尔塞纳}被他的大胆和想到这类人的阴谋吓坏了，立刻收回了所有战争计划，并缔结和约——我说，如果\Person{穆奇乌斯}这样做了，那么有人为了天国，可能将一只手，甚至全身投入火焰，并且不是出于自愿，而是受人迫害，谁又能说他对天国有功呢？如果\Person{库尔提乌斯}，策马全身武装跳入深渊，遵从他们神祇的神谕，神命令罗马人把拥有的最好的东西投入那深渊，他们只能理解为，既然他们在人和武器上优越，神命令将一个武装的人头朝下投入那毁灭中——如果他这样做了，那么一个人因同样的死而死，但不是自己主动跳入深渊，而是死于某个信仰仇敌之手，尤其当他从主——也是他国度的王——领受了更确定的神谕：\BibleQuote{你们不要害怕那杀害肉身，而不能杀害灵魂的}\BibleRef{玛 10:28}，我们岂能说这人为永恒之城做了什么大事呢？如果\Person{德基乌斯}家族献身死亡，用誓言祝圣自己，仿佛倒下去，以自己的血平息神怒，解救罗马军队——如果他们这样做了，那么圣洁的殉道者们，当他们为了在那有永生和幸福的天国里的一份，甚至流血致命时，不仅爱为他们流血的弟兄，而且遵照命令，也爱那使他们流血的仇敌，他们在爱的信德和信德的爱中彼此争胜，也不可自高，以为自己做了什么可夸的事。如果\Person{马尔库斯·普尔维路斯}，在献殿于\Person{朱庇特}、\Person{朱诺}和\Person{密涅瓦}时，有人故意告诉他儿子死了的假消息，想使他激动离去，这样献殿的荣耀就归于他的同僚；他如此无动于衷地接受了这个消息，竟下令将儿子抛尸不葬，爱荣耀之心战胜了丧亲之痛——如果他这样做了，那么对于传扬福音的人——天上之城的公民藉福音从各种错谬中被拯救出来，从各种流离中被聚集起来，主曾对他说，当他为埋葬父亲焦虑时：\BibleQuote{跟随我！任凭死人去埋葬他们的死人！}\BibleRef{玛 8:22}——谁能说他做了什么大事呢？\Person{雷古卢斯}，为了不违背誓言，甚至对最残忍的敌人，他从罗马返回他们那里，因为（据说当罗马人想留下他时他回答说）\Quote{在做了非洲人的奴隶后，他不能再有罗马荣耀公民的尊严}，而\Place{迦太基人}因为他曾在元老院中反对他们，用极刑将他处死。如果\Person{雷古卢斯}这样做了，那么为了对那国度的信实——信德本身引领人到达它的真福——还有什么痛苦是不该轻看的呢？或者，一个人如果为了对主应尽的信实，忍受了像\Person{雷古卢斯}为对敌人守信而遭受的同样痛苦，他又算偿还了主所赏赐的一切吗？再者，一个基督徒，为能在今生旅途的朝圣中，更轻松地走上通往那拥有真正财富的国度、即天主本身的道路，而选择了自愿的贫穷，当他听说或读到\Person{卢基乌斯·瓦莱里乌斯}，在担任执政官时去世，穷到需要人民捐钱办丧事；或听到\Person{昆提乌斯·辛辛纳图斯}，只有四亩地，亲手耕种，从犁旁被召立为独裁官——比执政官更尊贵——在战胜敌人赢得大荣耀后，却宁愿继续贫穷，他又怎敢以此自夸呢？那么，一个没有被任何现世赏报的许诺打动而放弃与那永恒天国的连结的人，当他听说\Person{法布里奇乌斯}没有被\Place{伊庇鲁斯}王\Person{皮洛士}的大礼打动而放弃罗马城，\Person{皮洛士}答应给他王国四分之一，他却宁愿贫穷地以平民身份留在那里，他怎能自夸做了大事呢？因为，当他们的共和国——即人民的利益、国家的利益、共同的利益——最为繁荣富足时，他们自己在家里却如此贫穷，以至于其中一位曾两次担任执政官的人，因为被发现有十磅银器，而被监察官从那穷人组成的元老院中开除——既然那些以凯旋充盈国库的人尚且如此贫穷，那么所有基督徒，他们以远为高尚的目的将自己的财富公用，即（按\WorkTitle{宗徒大事录}所记）他们按每人的需要分配，并且没有人说什么是自己的，一切都归公用\BibleRef{宗 2:45}——难道他们不应明白自己不可自夸吗？因为他们这样做是为获得天使的团体，而那些人几乎做了同样的事，为的是保存罗马人的荣耀。\par

这些事，以及罗马历史中所出现的类似事迹，若非罗马帝国因辉煌的胜利而扩张壮大，又怎能如此广为人知，并被传扬得如此赫赫有名呢？因此，通过这个如此辽阔、延续如此长久，又因如此伟大人物的德行而显耀辉煌的帝国，他们所追求的赏报便给予了他们热切的渴望，同时也为我们立下了榜样，其中包含着必要的劝诫，好使我们若看到自己并未为了最光荣的天主之城而持守那些德行——这些德行，无论如何，与那些他们为了地上之城的荣耀而持守的德行相似——便感到羞惭；并且，即使我们意识到自己持守了这些德行，也不可骄傲自大，因为正如宗徒所说：\BibleQuote{现时的苦楚，与将来在我们身上要显示的光荣，是不能较量的。}\BibleRef{罗 8:18}然而，就人间和短暂的荣耀而言，这些古代罗马人的一生被认为足够有价值。因此，在那隐藏于旧约而显明于新约的真理之光中，我们也看到，即：崇敬天主，并不是为了天主的上智安排不加以区分地赐予善人和恶人的地上和暂时的利益，而是为了永生、永恒的恩赐以及天上之城本身的团体——在这真理之光中，我们看到犹太人被极其公义地交给罗马人，作为他们荣耀的战利品；因为我们看到，这些罗马人，依靠地上的荣耀，并以他们那样的德行去寻求获得它，却征服了那些因其极大的堕落而杀害并拒绝了那赋予真正光荣及永恒之城者的人。\par

% structure: p0065 source=h2 target=subsection reason=relative-heading-rank; base=h2

\subsection{论真荣耀与支配欲的区别}

人间的荣耀之欲与支配之欲之间确有区别；因为，尽管对人间的荣耀过分喜爱的人也会极其倾向于热切地追求支配，但那些渴望甚至人间的赞美所给予的真荣耀的人，却努力不使那些正确判断他们的人不悦。因为有许多良好的品德，虽然并不被许多人拥有，却有许多人能够正确地判断它们；借着这些良好的品德，那些人奋力追求荣耀、尊荣和支配，\Person{撒路斯提乌斯}论到这些人说：\Quote{而他们却是循着真正的道路奋力前行。}\par

然而，谁若没有那使人畏惧不悦于判断其行为者的荣耀之欲，却渴望支配和权力，便往往寻求以最公开的罪行来获取他所爱的。因此，寻求荣耀的人，或是循着真正的道路，或是借着欺骗和诡计，奋力获得它，渴望在不善时显得善。所以，对有德行的人来说，轻看荣耀是一种大德；因为对荣耀的轻看为天主所见，却不显于人的判断。因为无论任何人在人眼前做什么，以显示自己是个轻看荣耀的人，倘若别人怀疑他这样做是为了博取更大的赞美——即更大的荣耀——他便无法向那些怀疑他的人证明事实并非如他们所怀疑的那样。然而，轻看赞美者之判断的人，也轻看怀疑者之轻率。如果他是真正良善的，他确实不轻视他们的救恩；因为那从圣神领受其德行的人，其义德是如此之大，以致他爱他的仇敌，甚至如此爱他们，希望那些憎恨和毁谤他的人能转向义德，成为他的同伴，不是在一个地上的国度，而是在天上的国度。但对于赞美他的人，尽管他不看重他们的赞美，却不轻视他们的爱；他也不回避他们的赞美，以免失去他们的爱。因此，他竭力使他们的赞美指向那位，凡在他身上真正值得赞美的，都是从祂那里领受的。但那轻看荣耀却贪求支配的人，在残忍和奢侈的恶习上胜过禽兽。某些罗马人正是如此，他们虽不贪图尊荣，却渴求支配；而历史证明，这样的人为数不少。然而，第一个登上这一恶习的顶峰，仿佛占据其卫城的人，是\Person{尼禄·凯撒}；因为他的奢侈如此之大，使人以为他没有什么男性可畏之处，而他的残忍又如此之甚，若非已知相反，无人会想到他性格中有任何柔弱之处。然而，即使是赐予这样的统治者，权力和支配也不例外于至高天主的上智安排，当祂判断人事的境况堪受如此之主时。对此，天主的言语清楚明了；因为天主的智慧如此说：\BibleQuote{藉着我，君王执政，暴君统辖大地。}\BibleRef{箴 8:15}然而，为了不让人以为\Quote{暴君}一词不是指邪恶不虔的君主，而是按照此词的古义指勇士，正如\Person{维吉尔}所说：\par

\Quote{要知道，当君主向君主致敬却不握手时，\newline{}盟约便不能成立。}\par

在另一处，论到天主，圣经极其明确地说：\BibleQuote{祂使伪善者作王，是因百姓的邪恶。}\BibleRef{约 34:30}所以，尽管我已尽我所能，说明了那位独一真而公义的天主，为何帮助那些依照某种地上之国的标准还算良善的罗马人，去获取如此伟大帝国的荣耀，但或许仍有一个更隐秘的原因，为天主所知甚于我们所知，取决于人类功绩的差异。至少，在一切真正虔诚的人中间，大家都同意：没有真正的虔诚——即对真天主的真正崇拜——就没有真正的德行；而那种是人的称赞之奴隶的，并非真正的德行。尽管如此，那些并非永恒之城——在圣经中被称为天主之城——的公民，若他们哪怕拥有那种德行，也比丝毫不具备时，对地上之城更有用。然而，对人类事务而言，再没有比这更幸运的事了：即藉天主的仁慈，那些具备了真正生活虔诚的人，若有治理人民的才能，便也应拥有权柄。但这样的人，不管他们在此生拥有多么伟大的德行，都唯独将它归功于天主的恩宠，是祂将这德行赐予了他们——愿欲、相信、寻求。同时，他们也明白自己在义德的完美上离那圣天使的团体中的完美有多远，他们正努力使自己适合那团体。然而，那种没有真正虔诚、是人的荣耀之奴隶的德行，不论受到何等赞美和吹嘘，都根本无法与圣徒们德行的微弱开端相比拟，因为圣徒们的希望寄托于真天主的恩宠和仁慈。\par

% structure: p0070 source=h2 target=subsection reason=relative-heading-rank; base=h2

\subsection{第二十章 德行服务于人的荣耀与服务于肉身的快乐同样可耻}

哲学家们——那些将人之善的目的置于德行本身，为要羞辱另一些哲学家；后者虽认可德行，却都以肉身的快乐为目的来衡量一切德行，并认为快乐应为其自身而追求，而德行则因快乐之故才需要——习惯于描绘这样一幅文字图画：快乐如同一位奢靡的女王坐在王座上，一切德行都作为奴隶臣服于她，察其颜色，好遵行她的吩咐。她命令明智时刻警醒，以发现快乐如何能施行统治并安然无恙。她命令正义尽力施予所能施的益处，以获取那些为肉身快乐所必需的友谊；谁也不可得罪，免得因违犯法律，快乐无法安全生活。她命令勇毅顽强地护卫女主人——即快乐——如果身体遭受任何不至于死的痛苦，便藉回忆往日的欢愉来缓解当前痛苦的尖锐。她命令节制甚至对最喜爱的食物也只取一定分量，免得因过度食用，任何东西扰乱身体健康而有害，这样快乐——\OriginalTerm{伊壁鸠鲁派}{Epicureans}认为主要在于身体健康的快乐——便遭到严重冒犯。因此，德行带着其一切荣耀的尊严，将成为快乐的奴隶，如同服事某个专横又声名狼藉的女人。\par

我们的哲学家们说，没有比这幅图画更可耻、更畸形的东西了，好人的眼睛更不能忍受它。他们说出了真理。但我认为，即便勾勒一幅德行作为人的荣耀的奴隶的图画，这幅画也并非足够得体；因为，那荣耀虽不是奢靡的女人，却也自大，内中充满许多虚浮。因此，描绘德行服事这荣耀，以致明智一无所筹，正义一无所分，节制无所节制，除非为取悦于人、服事虚荣，这是与德行的坚实稳固不相称的。他们也不能使自己免于如此卑劣的指控，当他们以轻视荣耀者自居，不顾他人判断，自以为智，且自我取悦时。因为他们的德行——如果那还能算德行的话——不过是以另一种方式屈从于人的称赞；因为那寻求取悦自己的人，仍在寻求取悦人。但那位怀着对天主的真正虔诚——他所爱、所信、所仰望的天主——的人，更专注于那些使他令自己不悦的事，而非那些令他喜悦的事（如果有任何这样的事的话），或者不如说，令他喜悦的并非他自己，而是真理；他能如今取悦于真理，他将这能力不归功于别的，只归功于那位他怕得罪的天主的仁慈，为他内得到医治的部分献上感谢，并为他内尚未痊愈的部分倾倒祈祷。\par

% structure: p0073 source=h2 target=subsection reason=relative-heading-rank; base=h2

\subsection{第二十一章 罗马的统治权是由那位一切权柄皆出于祂，且以其眷顾治理万物的天主所赐予的}

这些事既如此，我们不把赐予万国和帝国的权柄归于什么别的，只归于真天主，祂将天堂国度的幸福只赐给虔诚的人，却在今世将王权赐给虔诚者和不虔诚者，全凭祂所喜悦的，而祂的圣意常是公义的。因为我们虽就祂治理的法则略有所说，是照祂乐意向我们显明的程度，但探究人心中的隐秘，并藉明确的审查断定各王国的功过，对我们来说实为太重，远超我们的力量。所以，那独一真天主，从不以公义的审判和援助离弃人类，在祂愿意的时候，将王国赐给罗马人，且随祂意定其大小，就像祂对\Place{亚述}人，甚至对\Place{波斯}人所行的一样——据他们自己的典籍证明，波斯人只崇拜两个神，一善一恶——至于希伯来民族，我已如前所述略表必要的话，他们有国时期，惟独崇拜真天主。故同一天主，既将丰收赐给波斯人，尽管他们不崇拜女神\Person{塞革提亚}，又将地上的其他恩泽赐给他们，尽管他们不崇拜罗马人所设想的众多神明，那些神明各自掌管某个特定事物，甚至多个神明共掌一事——我说，同一天主将统治权赐给波斯人，尽管他们不崇拜那些罗马人自以为帝国仰赖的神明。这对个人也是如此，如同对国家一样。那位赐权柄与\Person{马略}的，也将权柄赐给\Person{盖乌斯·凯撒}；那位赐给\Person{奥古斯都}的，也赐给\Person{尼禄}；那位赐给最仁慈的皇帝——\Person{韦斯帕芗}父子的，也赐给残忍的\Person{图密善}；最后，无需逐一列举，那位将权柄赐给基督徒\Person{君士坦丁}的，也赐给背教者\Person{尤利安}，其天赋的才思被一种亵渎而可恶的好奇所欺骗，并受权欲的刺激。正是由于他沉溺于好奇，求问虚妄的神谕，自信胜利在望，竟烧毁了装载军队必需粮草的船只，因此以狂热的激情投身于轻率大胆的冒险，不久自食其卤莽的恶果而被杀，使其军旅在敌境中无粮可继，陷入绝境，若非更改罗马帝国边界，违背我在上卷所言界神\Person{特耳弥努斯}的征兆，决不能脱险；因为界神\Person{特耳弥努斯}虽未屈于\Person{朱庇特}，却向急需让了步。显然，这些事是由独一天主按己意统治掌管；祂的动机虽隐秘，岂因此就不公义吗？\par

% structure: p0075 source=h2 target=subsection reason=relative-heading-rank; base=h2

\subsection{二十二、战争的持续与结局取决于天主的旨意。}

因此，战争的持续长短也由祂按祂公义的旨意、喜悦和仁慈，为磨难或安慰人类，随时而适地决定，以至于有时较长，有时较短。海盗战争与第三次布匿战争以难以置信的速度结束。还有角斗士逃亡之战，虽然许多罗马将军和执政官都曾被击败，意大利惨遭蹂躏摧残，但仍在三年内结束，在其延续期间，它本身就是许多的终结。\Place{派森提}人、\Place{马尔西}人及\Place{佩利尼}人，并非远方部族，而是\Place{意大利}境内之民，在长期极忠顺地服属罗马轭下之后，试图昂首争取独立，尽管那时许多民族已屈服于罗马武力，\Place{迦太基}也已覆灭。在此意大利战争中，罗马人屡遭败北，两位执政官阵亡，之外还有贵族元老；但此灾祸并未迁延很久，第五年就告结束。然而第二次布匿战争，延亘十八年，给共和国带来极大的灾难和损伤，耗尽并几乎摧垮罗马人的力量；因为两场战役中约有七万罗马人阵亡。第一次布匿战争是在进行了二十三年之后告终。米特里达梯战争进行了四十年之久。为免有人以为在早期备受赞颂的时代，罗马人更为英勇能干，能迅速结束战争，我们可以指出\Place{萨莫奈}战争曾拖延将近五十年；在此战争中，罗马人惨败至被迫轭下受辱。但因他们不是为正义而爱光荣，反似为光荣而爱正义，他们竟破坏了已订立的和约与盟约。我提起这些事，因为有许多人昧于过去之事，也有一些明知而佯作不知，若在基督教时代见任何战争稍逾所期，便立刻对我们的宗教发烈怒和狂言，高呼若非如此，他们早已按古礼求告诸神，靠着昔日罗马人的英勇，在\Person{玛尔斯}和\Person{贝娄娜}的助佑下，迅速了结所有大战，这场战争也会很快结束。所以，让读过历史的人回想，古代罗马人曾发动了多少次持续甚久、胜败无常、导致悲惨杀戮的战争——正如世间一般常情，大地如汹涌的深海，常被风暴摇撼，这些灾祸的风暴程度不一——让他们有时承认他们不愿承认的事，不要疯狂地攻击天主，自取沉沦，又欺骗无知的人。\par

% structure: p0077 source=h2 target=subsection reason=relative-heading-rank; base=h2

\subsection{二十三、论\Place{哥特}王\Person{拉达盖苏斯}，本为魔鬼崇拜者，一日之间与其全部大军同遭败亡的战争。}

他们却并不以感恩之心提及天主近日在我们记忆中奇妙而仁慈的作为，反而尽其所能，企图将其埋没于普遍的遗忘之中。但若我们对这些事保持沉默，我们也同样忘恩负义。当哥特王\Person{拉达盖苏斯}率领庞大而野蛮的军队，在离城极近之处扎营，已逼近罗马人时，他竟在一天之内如此迅速而彻底地被击败，以至于罗马人无一受伤，更无死亡，而他的军队却有超过十万人被歼灭，他本人及其子嗣被俘，随即被处死，受到了应有的惩罚。因为，倘若如此不虔敬之人，率领如此强大而不虔敬的军队，攻入城内，他会放过谁呢？他会尊重哪个殉道者的坟墓？他对哪个人会表现出对天主的敬畏？他会不流谁的血？他会愿意保全谁的贞洁不受侵犯？那时，他们对他们的神祇的赞颂会有多么喧嚣！他们会何等傲慢地夸耀说，拉达盖苏斯之所以能征服，之所以能成就如此伟业，是因为他通过每日的祭献安抚并赢得了诸神——而基督宗教却不允许罗马人这样做！因为当他逼近至尊威仪一颔首便使他覆灭之地时，他的名声四处传播，我们当时在\Place{迦太基}便听说，异教徒们相信、宣布并夸耀，由于他每日据说向诸神献祭，因而获得了友好诸神的帮助和护佑，他一定不会被那些没有向罗马诸神举行此类祭献、甚至不允许任何人献祭的人所征服。如今，这些可怜之人却不因天主的莫大仁慈而感谢祂；天主决心惩罚人类的败坏，这败坏远比蛮族的败坏更应受严惩，却以如此温和的方式抑制祂的义怒：首先，使哥特王以奇妙的方式被征服，以免荣归于他所祈求的魔鬼，从而动摇软弱者的信心；然后，在\Place{罗马}城将被攻陷时，使其由那些蛮族攻取，这些蛮族一反以往所有战争的惯例，出于对基督宗教的敬畏，保护那些逃往圣地避难的人，并且如此反对魔鬼本身及不虔敬的祭献仪式，以至于他们似乎与魔鬼进行的战争远比与人类的战争更为可怕。如此，万物的真主宰既仁慈地鞭责了罗马人，又借着对魔鬼崇拜者的奇妙击败，表明那些祭献即使为了现世的安全也并非必要；这样，那些不顽梗坚持、而是谨慎思考的人，便不会因现时的迫切需要而离弃真宗教，反会在对永生最坚定的期望中更紧地持守它。\par

% structure: p0079 source=h2 target=subsection reason=relative-heading-rank; base=h2

\subsection{基督徒皇帝的幸福是什么，以及它在多大程度上是真实的幸福。}

因为我们并不是说某些基督徒皇帝之所以幸福，是因为他们统治长久，或是安详而终，留下儿子继承帝位，或是征服了国家的敌人，或是能够预防并镇压那些对他们怀有敌意的公民的叛乱。这些以及其他可悲生命中的恩赐或安慰，连某些魔鬼的崇拜者也有幸得到，他们并不属于天主的国，而这些基督徒皇帝却属于；这应归因于天主的仁慈，祂不愿让信靠祂的人把这些东西当作最高的善。但我们却说，他们是幸福的，如果他们秉公治理；如果在那些以极高尊敬赞誉他们、以过分谦卑奉承他们的人当中，他们不骄傲自大，而是记得自己不过是人；如果他们运用自己的权力，使之成为天主威仪的婢女，用以尽可能地扩展对祂的敬拜；如果他们敬畏、爱慕、钦崇天主；如果他们爱那个他们不畏惧其中有同伴的国度，胜过爱自己的国度；如果他们缓于惩罚，乐于宽恕；如果他们施行惩罚是为了统治和保卫国家的必要，而不是为了满足个人的仇恨；如果他们宽恕罪过，不是为了让不义不受惩罚，而是期望犯罪者能改过自新；如果他们以仁慈的温和与善行的慷慨来补偿不得不作出的任何严厉判决；如果他们的奢侈受到极大的遏制，正如它可能不受遏制一样；如果他们宁肯治理败坏的欲望，而不愿统治任何民族；如果他们做这一切，不是出于对虚空荣耀的贪求，而是出于对永恒真福的爱慕，并且不忘为他们自己的罪过，向真天主，他们的天主，献上谦卑、痛悔和祈祷的祭献。我们说，这样的基督徒皇帝，在现世是因着\OriginalTerm{望德}{hope}而幸福，并且注定在等待之事来临时，在现实的享乐中也是如此。\par

% structure: p0081 source=h2 target=subsection reason=relative-heading-rank; base=h2

\subsection{论天主赐予基督徒皇帝\Person{君士坦丁}的繁荣。}

因为美善的天主，为了不让人们以为，凡是相信应当敬拜祂以求永生的人，不能不认为，除非崇拜魔鬼——以为这些神灵对世福有极大的能力——否则无人能达到如此崇高的地位和地上的王权，因此，祂将如此满盈的世福赐给了\Person{君士坦丁}皇帝，他并不崇拜魔鬼，而是崇拜真天主本身，这样的福分甚至无人敢于奢望。祂还赐给了他建立一座城的荣耀，这座城作为\Place{罗马帝国}的伴侣，宛如\Place{罗马}的亲女儿，却没有任何魔鬼的庙宇或像。他作为独一的皇帝统治了很长时期，独自拥有并捍卫了整个罗马世界。在指挥和进行战争时，他战无不胜；在推翻暴君时，他最为成功。他享尽天年，因疾病和年老而逝世，留下他的儿子们继承帝国。然而，为防止任何皇帝为了求得\Person{君士坦丁}那样的福乐而成为基督徒，而每一个人都应当为永生而成为基督徒，天主取走了\Person{约维安}，远远早于\Person{尤利安}，并允许\Person{格拉提安}被暴君之剑所杀。但在他的情况下，灾难的减轻远胜于伟大的\Person{庞培}，因为\Person{庞培}无法被他留下作为内战继承人的\Person{加图}所复仇。但\Person{格拉提安}虽无需虔诚的心灵所求的那种安慰，却被他与之共治帝国的\Person{狄奥多西}所复仇，尽管他有一个小弟弟，但他更渴望忠实的联盟，而不是广大的权力。\par

% structure: p0083 source=h2 target=subsection reason=relative-heading-rank; base=h2

\subsection{第26章 论\Person{狄奥多西·奥古斯都}的信德与虔诚}

因此，\Person{狄奥多西}不仅在\Person{格拉齐安}生前始终对他尽忠，而且在他死后，如同真正的基督徒那样，将他的小弟弟\Person{瓦伦提尼安}保于自己的保护之下，奉为共治皇帝，当时\Person{瓦伦提尼安}已被其父的凶手\Person{马克西穆斯}驱逐。他以慈父般的爱守护他，尽管当时\Person{瓦伦提尼安}毫无资源，如果\Person{狄奥多西}怀有扩张帝国的欲望，而非造福他人的雄心，他完全可以毫无困难地除掉他。因此，当他收养了这个孩子，并为他保全了皇帝的尊严，以仁爱和善意安慰他，这给他带来的喜乐要远远大得多。后来，当那种成功使\Person{马克西穆斯}变得可怕时，\Person{狄奥多西}在重重焦虑之中，并未被引向听从某种渎神而非法的好奇心的建议，而是派人前往\Person{若望}那里——这位\Person{若望}居住于\Place{埃及}的沙漠中，因为他得知这位天主的仆人（其名声远扬）被赋予了预言的恩赐——并从\Person{若望}那里获得了必胜的保证。他立即以最深切的同情和敬意，诛杀了暴君\Person{马克西穆斯}，并将被驱逐的幼童\Person{瓦伦提尼安}恢复了他应有的帝国之份。不久之后，\Person{瓦伦提尼安}被秘密刺杀，或因其他阴谋或意外身亡，\Person{狄奥多西}再次从那位先知那里得到答复，并全心信赖，便起兵攻打暴君\Person{尤金尼乌斯}——此人被非法拥立继承那位皇帝——他击败了\Person{尤金尼乌斯}极为强大的军队，凭借祈祷多于刀剑。一些参与此战的士兵向我报告说，他们投掷的所有投射物都被一阵狂风从手中夺走，这阵风从\Person{狄奥多西}的军队方向吹向敌军；风不仅使掷向敌军的标枪以更快的速度射出，还将敌军自己投出的标枪转回来，刺入他们自己的身体。因此，诗人\Person{克劳狄安}，尽管与基督之名无份，却在他对\Person{狄奥多西}的赞歌中说：\Quote{哦，君主，你太蒙神所爱，为你埃俄罗斯从洞穴中倾出全副武装的风暴；为你天空也交战，所有的风都同心听从你的号角。}但胜利者，正如他所信并预言的，推倒了\Person{朱庇特}的雕像——这些雕像似乎是经由我不清楚的某种仪式，被奉献用于反对他，并竖立在\Place{阿尔卑斯山}上。而这些雕像的金制雷霆，他愉快而和蔼地赠给了他的报信使者，当时他们（趁此欢庆之机）开玩笑说，若能遭这样的雷霆击打，他们将是极幸福的。对于他自己敌人的儿子们，他们的父亲并非直接由他的命令所杀，而是死于战争的激烈，他们逃到教会寻求庇护，虽然他们还不是基督徒，他却急切地利用这个机会，引导他们归向基督信仰，并以基督徒的爱德对待他们。他并没有剥夺他们的财产，反而除了允许他们保留财产外，还赐予他们额外的荣耀。战争结束后，他不容许私人的仇恨影响对任何人的处置。他不像\Person{秦纳}、\Person{马略}、\Person{苏拉}及其同类那样，这些人甚至在内战已经结束时还不愿结束内战，与其说他们希望在内战结束时伤害任何人，不如说他们更因内战的发生而悲伤。在这一切事件中，自他执政之初，他就从未停止以最公义慈悲的法律，帮助遭受患难的教会对抗不虔敬者，而持异端的\Person{瓦伦斯}，因偏袒亚略派，曾大大压迫了教会。确实，他因成为这教会的一员而喜乐，远胜于作地上的君王。他下令到处推倒外邦人的偶像，深明即使世间的恩赐也不在魔鬼的权下，而在真天主的权下。还有什么比他信仰上的谦卑更为可叹呢？当他受某些近臣的催逼，为得撒洛尼人所犯的重罪施行报复时——他曾应主教们的祈祷，许诺要宽恕这罪行——他受到了教会纪律的约束，如此做了补赎，以至于他那至尊俯伏的景象，使代他求情的人民流泪，更甚于他们自觉冒犯时对盛怒的恐惧。这些以及其他类似的善工，若要细说就太长了，他带着它们离开了这短暂的世界，在此世间，人类最高的尊贵与崇高不过是云烟。这些善工的奖赏是永恒的幸福，天主是赐予者，但只赐予那些真诚虔敬的人。但此生所有其他的福祉与特权，如世界本身、光、空气、大地、水、果实，以及人的灵魂本身、他的身体、感官、心灵、生命，天主都同样慷慨赐予善人和恶人。而在这些福祉中，也应算上对帝国的拥有，其疆域大小，天主依其眷顾治理之所需，在不同时期加以调节。由此，我看出，我们现在必须回答那些人——他们已被最明显的证据驳倒并定罪，这些证据表明，为获取这些世间之物——而愚昧人所渴望的无非这些——众多假神毫无用处；他们却企图声称，敬拜诸神不是为了现世的利益，而是为了死后将来的生命。至于那些为了此世的友谊而甘愿崇拜虚幻，并不因自己被留在幼稚的见识中而悲伤的人，我想在这五卷书中已经充分回答了他们；当我发表了前三卷，它们开始传到许多人手中时，我听说某些人正准备写一篇回答来反驳它们，以书面形式。接着又有人告诉我，他们已经写好了回答，但在等待一个时机，可以毫无危险地将其发表。我愿劝告这些人，不要渴望那些对他们毫无益处的事物；因为当一个人只是不愿沉默时，他很容易自以为已经答复了论证。因为还有什么比虚妄更多言呢？即使它愿意，能比真理喊得更响，尽管如此，它并不比真理更有力。但让人勤勉思考我们所说的这一切，倘若他们不怀党派精神地判断，或许能清楚看出，这些事可能被他们那极其无耻的饶舌以及如同讥讽模仿的轻浮所撼动，而非连根拔除，那么，就让他们收起自己的荒谬，选择受智者纠正，强于受愚人赞美。因为，若他们等待时机，不是为了自由地讲真理，而是为了得到辱骂的许可，难道\Person{图利乌斯}论及某人的话不会应验在他们身上吗？\Quote{哦，可怜的人啊！谁给了他犯罪的自由？}因此，无论何人，若因有辱骂的许可而自以为幸福，那么若根本不许他如此，他将会幸福得多；因为他完全可以一直放弃空洞的夸耀，以自由咨询的方式与那些他所反对其观点的人争辩，并以符合其身份的方式——尊敬地、庄重地、坦诚地——倾听一切能由那些他以友善辩论方式咨询的人所提出的论证。\par

\SourceRecord{Translated by Marcus Dods.；Nicene and Post-Nicene Fathers, First Series；Vol. 2.；Edited by Philip Schaff.；Buffalo, NY: Christian Literature Publishing Co.,；1887.；<http://www.newadvent.org/fathers/120105.htm>.}

% structure: p0001 source=h1 target=section reason=page-title

\section{\WorkTitle{天主之城}（卷六）}

此前论证一直针对那些相信为获得现世利益而应崇拜神祇的人，如今则转向那些相信为获得永生而应崇拜神祇的人。奥古斯丁用以下五卷书来驳斥这后一种信念，并首先指出，即便是最受推崇的异教神学作家\Person{瓦罗}本人，对神祇所持的观点也何其卑微。奥古斯丁采用瓦罗把神学划分为神话的、自然的和公民的这三种类型的方法，并随即证明，无论神话神学还是公民神学，都丝毫不能有助于未来生命的幸福。\par

% structure: p0003 source=h2 target=subsection reason=relative-heading-rank; base=h2

\subsection{前言}

我认为，在前五卷书中，我已对那些认为应为了现世生命和世俗事务的利益，以希腊人称为\OriginalTerm{崇拜}{\GreekText{λατρεία}}的仪式与事奉——这原是只应归于唯一真天主的——去崇拜众多伪神的人，作了充分的驳斥；而基督信仰的真理却指明，那些伪神不过是无用之偶像，或污秽的邪灵与有害的魔鬼，或至多只是受造之物，决非造物主。谁不知道，面对极度的愚钝与固执，即便这五卷书，甚或不计其数的书籍，都无济于事，因为人们竟把在真理力量面前绝不屈服视为虚荣的荣耀——这其实是在毁灭那受此可憎恶习所辖制的人？因为，尽管医者不懈努力以求治愈，但疾病未被征服，这并非他的过失，而是由于病人已无可救药。然而，那些仔细衡量所读内容的人，在明白并思考之后，若没有，或没有过多地怀有那种属于长期怀抱之错误的固执，便会更倾向于认为，在已完成的五卷书中，我们所做的已超过了问题本身的要求，而不是所作讨论少于所需。他们也必定不会怀疑，一切憎恨——即无知者因今生的灾难和地上事物的毁灭与变迁而试图加于基督宗教的憎恨，而有识之士却不只佯装不见，反而违背自己的良心加以煽动，因为他们被一种疯狂的\OriginalTerm{亵渎}{impiety}所控制——他们，我说，必定不会怀疑，这种憎恨既无正当的思考与理性，又充满了最轻率的鲁莽和最有害的敌意。\par

% structure: p0005 source=h2 target=subsection reason=relative-heading-rank; base=h2

\subsection{第一章 论那些主张崇拜神祇不是为了现世，而是为了永世利益的人}

如今，依照既定的次序，接下来应当驳斥并教导那些人——他们主张，民族的诸神，即基督信仰的真理所摧毁的，应当受崇拜，不是为了今生，而是为了死后那生命——因此，我最好从\BibleBook{}{圣咏}的真实神谕开始我的辩论：\BibleQuote{凡以上主天主为希望，不理会虚幻和欺诈愚妄的人，真是有福。}然而，在一切虚幻和欺诈愚妄之中，那些哲学家还应得到更大的宽容，因为他们曾拒斥民众的那些意见和错误；因为民众给神祇设立偶像，或者编造许多关于那些他们称之为不死神明的虚假和不配之事，或者相信那些已被编造的事，并在相信之后，将它们混合到崇拜和神圣仪式之中。\par

与那些人，他们虽不公开表白自己的信念，却在论及此事时通过咕咕哝哝的不赞同表明他们不认可那些事，与他们讨论以下问题也许并非很不妥：为了死后将有之生命，我们应当崇拜的，不是那创造一切灵性与有形受造物的唯一天主，而是那些许多的神明——按照这些哲学家中某些人的看法，他们由那唯一天主所造，并由祂安置于各自崇高的领域，因而被认为比其他一切神明更卓越、更尊贵——吗？然而，谁又会声称，必须肯定并主张：我在第四卷中提到的某些神明，他们被一一分派了细小事务的职责，却能赐予永生？但那些极有技巧与极敏锐的人，他们夸耀自己曾为人类的大益而写作，教导因何原由应崇拜每一位神，以及应向每一位祈求什么，以免出现最可耻的荒谬——如同丑角为逗乐而常表演的那样，向\OriginalTerm{利贝尔}{Liber}求水，向\OriginalTerm{林菲斯}{Lymphs}求酒——这些人是否会确实地告诉任何恳求不朽神明的人说：当他向林菲斯求酒，而她们回答他\Quote{我们有水，向利贝尔求酒吧，}他便可以正当地说\Quote{如果你们没有酒，至少给我永生？}还有什么比这荒谬更怪异？这些林菲斯——因为她们通常很容易被逗笑——岂不哈哈大笑（除非她们像魔鬼一样企图欺骗），并回答恳求者说：\Quote{人啊，你以为我们掌管\Emph{\OriginalTerm{生命}{vitam}}，而你们却听说我们连\Emph{\OriginalTerm{葡萄藤}{vitem}}都没有？} 因此，向那些被宣称如此主管此一极可悲而短暂生命之各个细小关切，以及任何有益于维持扶持此生命之事的诸神，去寻求并盼望永生，实乃最厚颜无耻之愚妄；以至若有任何在一位神明照管与权能之下的事物，却向另一位去求，便如此不协调、如此荒谬，极像丑角的插科打诨，——当此事由知晓自己在做什么的丑角行之，在剧场里当得人笑；但当由无知而行的愚人行之，在世间便更当受人耻笑。据此，就各个城邦所立之诸神而言，学者们已巧妙地创制并传留下记忆：关于每一特定事物，当向何神或女神祈求——例如，当向利贝尔求什么，向林菲斯求什么，向\OriginalTerm{伏尔甘}{Vulcan}求什么，以及其余一切，其中有些我已在第四卷提及，有些我认为宜予省略。再者，若向\OriginalTerm{刻瑞斯}{Ceres}求酒、向利贝尔求饼、向伏尔甘求水、向林菲斯求火是错误，那么倘若有人向其中任何一位祈求永生，又该被视为何等更大的荒谬呢？\par

因此，当我们曾探究当相信何种神明能赋予人地上王国时，在讨论一切之后，已表明：即使地上王国也绝非由那些众多虚假神祇中的任何一位所建立，——若此，那么相信永生——即那无疑且无比地胜于一切地上王国的永生——竟能由这些神中的任何一位赐予任何人，岂非最疯狂的不虔？因为，我们之所以认为那样的神甚至不能赐予地上王国，并非由于他们极其伟大崇高，而地上王国是微小卑贱的，以他们如此之高的尊威，竟不屑垂顾；而是因为，尽管任何人出于对人类脆弱的体恤，可以轻看地上王国那即将倾颓的顶巅，但这些神所呈现的面貌，却显得甚至不配受托赐予并保全这些国度；因此，如果在那一大群神中——无论属于所谓\OriginalTerm{庶民之神}{plebeian gods}还是\OriginalTerm{贵族之神}{noble gods}——没有一位适合将必死的王国赐予必死之人（如我们在讨论此事的前两卷中所教导的），那么他岂不更不可能使必死之人成为不朽？\par

此外，如果按照我们与之争论的那些人的意见，应当崇拜神明，不是为了现世的生命，而是为了死后将有的生命，那么，他们当然不应该为了那些个别事物而崇拜神明；这些事物被分配和指派给这样一些神祇的权能——不是根据任何理性的真理法则，而只是凭空虚的猜测——那些主张为了获得这必死生命中的一切所欲之物而必须崇拜神明的人，就相信他们应当如此崇拜。对于这些人，我在前五卷中已经尽我所能地充分辩驳过了。既然如此，如果那些崇拜\OriginalTerm{青春女神犹文塔斯}{Juventas}的人确实在其青春年纪上表现出显著的活力，而轻视她的人要么在青年时代便死亡，要么在此期间变得迟钝，如同老年的麻痹；如果\OriginalTerm{胡子福耳图娜}{Fortuna Barbata}确实能使崇拜她的人面颊上长出的胡须比所有其他人更俊美、更优雅，而我们看到轻视她的人要么完全没有胡须，要么胡须长得难看；即使如此，我们也可以十分正确地说，这些神明各自拥有在其职司范围内受到一定限制的能力，因此，既不应向不能赐予胡须的犹文塔斯寻求永生，也不应在今生之后从胡子福耳图娜那里期待任何美善之物，因为她即使在此生也没有能力赐予那生长胡须的年纪本身。但如今，当甚至连他们以为受其权能支配的那些事情本身，也并不需要崇拜他们时——因为许多崇拜女神犹文塔斯的人在那个时候根本不强健，而许多不崇拜她的人却拥有青春的力量；同样，许多向胡子福耳图娜祈求的人或者完全长不出任何胡须，甚至连难看的胡须也长不出，尽管那些为了求得胡须而崇拜她的人反被有胡须的轻视者所讥笑——人心真是如此愚蠢，竟会相信对于这些现世短暂、转瞬即逝的恩赐——据说每一样都由这些神明中的一位掌管——其崇拜既是徒然又可笑的，却能对于永生结出果实来？而且，他们能赐予永生这一点，甚至连那些为了受到愚蠢百姓的崇拜而将现世事务细微划分并分配给他们，免得他们闲坐无事的人，也不曾如此宣称过；因为他们认为神明数目极其繁多。\par

% structure: p0010 source=h2 target=subsection reason=relative-heading-rank; base=h2

\subsection{关于\Person{瓦罗}对异邦众神所持的看法，我们当作何想——他既已表明这些神祇的种类与礼仪如此不堪，以致他若对其完全保持沉默，反倒是对他们更为敬重了。}

谁曾比\Person{玛尔库斯·瓦罗}更仔细地研究这些事？谁曾以更高的学识发现它们？谁曾更专注地思考它们？谁曾更敏锐地辨别它们？谁曾更勤勉、更充分地论述过它们？——他虽在口才方面不那么讨人喜欢，却是如此充满教益与智慧，以致在我们所谓\Quote{世俗}而他们称为\Quote{自由}的全部学术中，他之于事物的研习者，正如\Person{西塞罗}之于言辞的爱好者一样，能给人以教导。甚至\Person{西塞罗}本人也为他作证，在他的\WorkTitle{学院派}书中说，那场他与\Person{玛尔库斯·瓦罗}进行的辩论，他称赞道：\Quote{一个毫无疑问在所有的人中最为敏锐的人，并且无可置疑地，最有学问的人。}他没有说最雄辩或最流畅的人，因为他确实极缺乏这种才能，而只说\Quote{在所有的人中最为敏锐。}在那些书中——即\WorkTitle{学院派}——他主张一切事物都当存疑，却又这样评价他：\Quote{无可置疑地最有学问。}事实上，他对这一点如此确信，竟至将他惯于在一切事上运用的那种怀疑弃置一旁，仿佛当他即将为学院派的怀疑辩护时，就这一件事而言，他忘记了自己是学院派。但在第一书里，当他盛赞同一位\Person{瓦罗}的文学著作时，他说：\Quote{我们如同异乡人在自己的城邦中流浪徘徊，你的书籍仿佛将我们带回家中，好使我们最终明白自己是谁，身在何处。你向我们揭示了祖国的年代、时令的更替、圣物和祭司的法律；你向我们揭示了家庭的和公共的纪律；你为我们指明了宗教典礼的适当场所，并告知我们关于圣地的事。你向我们展示了所有神事和人事的名称、种类、职能和原因。}\par

这位如此卓越出色、博学多识的人，正如\Person{特伦提安}在一首极为优雅的诗中简练地评价他说：\par

\Quote{\Person{瓦罗}，一位无所不知的人，}\par

谁读了如此之多，以致我们惊异他何时有暇写作，又写了如此之多，以致我们几乎难信有人能全部读完——我指的是这个人，如此伟大的天才，如此博学——若他曾是他笔下所谓\OriginalTerm{神事}{divine things}的反对者和摧毁者，并曾说过这些事属于迷信而非宗教，那么即使那样，他或许也不会写出这么多可笑、可鄙、可憎的东西。然而，当他如此崇拜这些神祇，如此为其崇拜辩护，甚至在他那同一部著作中说，他担心这些神祇会消亡，不是由于敌人的攻击，而是由于公民的疏忽，而他正将它们从这种耻辱中解救出来，并通过这类书籍将它们保存并传承于善人的记忆中，其热忱远比据说\Person{梅特鲁斯}从火焰中救出\Person{维斯塔}的圣物、\Person{埃涅阿斯}从\Place{特洛伊}大火中救出家神更为有益；然而，他却又将这样的东西留给后世阅读，而这些东西智者和愚者都理当判定为不宜阅读，也与宗教的真理最为敌对；对此，除了认为这位极其敏锐博学之人——却未被\OriginalTerm{圣神}{Holy Spirit}所释放——被本国的习俗与法律所压倒，因无法对那些影响他的事物缄默，遂在赞扬宗教的借口下谈论它们之外，我们还能作何想呢？\par

% structure: p0015 source=h2 target=subsection reason=relative-heading-rank; base=h2

\subsection{第三章 \Person{瓦罗}为其所撰\WorkTitle{人事与神事古事纪}一书的编排}

他写了四十一卷\WorkTitle{古事纪}。他把这些分为\OriginalTerm{人事}{human things}和\OriginalTerm{神事}{divine things}。二十五卷专论人事，十六卷专论神事；其划分计划如下：即为人事的四个部分各分配六卷。因为他关注的问题是：谁执行，在何处执行，何时执行，执行什么。因此，前六卷论及人，次六卷论及地点，再次六卷论及时间，最后的六卷论及事物。不过四六得二十四。但他在卷首另置一卷独立的著作，综论这所有四个方面。\par

在神事方面，就那些为神而执行的事而言，他保持了同样的次序。因为\OriginalTerm{神圣之事}{sacred things}是由人在特定的地点和时间执行的。这四项我在前面提及的内容，他包含在十二卷书中，每项分配三卷。因为他写了前三卷论及人，随后三卷论及地点，接着三卷论及时间，最后三卷论及\OriginalTerm{祭礼}{sacred rites}——极其细致地区分了谁应执行、在何处执行、何时执行、执行什么。但是，由于还有必要说明——而且这正是人们特别期待的——应对谁执行这些祭礼，于是他写了最后三卷，论述神祇本身；这样五个三卷就成了十五卷。不过，如前所述，总共有十六卷。因为他还在这一部分的开头另置一卷独特的书，作为此后一切内容的导论；完成之后，他进而将前三卷按与人相关的\OriginalTerm{五重划分}{five-fold distribution}加以细分：第一卷论\OriginalTerm{大司祭}{high priests}，第二卷论\OriginalTerm{占卜官}{augurs}，第三卷论主持\OriginalTerm{神圣典礼}{sacred ceremonies}的\OriginalTerm{十五人司祭团}{fifteen men presiding over the sacred ceremonies}。第二部分三卷论地点，其一论\OriginalTerm{小圣堂}{chapels}，其二论\OriginalTerm{神庙}{temples}，其三论\OriginalTerm{圣地}{religious places}。随后的三卷论时间，即节庆日，他将其分为：一卷论\OriginalTerm{节日}{holidays}，另一卷论\OriginalTerm{竞技游戏}{circus games}，第三卷论\OriginalTerm{戏剧}{scenic plays}。第四部分三卷论\OriginalTerm{神圣事物}{sacred things}，一卷专论\OriginalTerm{祝圣}{consecrations}，另一卷专论\OriginalTerm{私人祭礼}{private sacred rites}，最后一卷专论\OriginalTerm{公共祭礼}{public sacred rites}。最后剩下的三卷，则仿佛追随这一辉煌行列，专论所有这些祭礼所为之目的，即神祇本身。第一卷论\OriginalTerm{确定的神}{certain gods}，第二卷论\OriginalTerm{不确定的神}{uncertain gods}，第三卷即最后一卷论\OriginalTerm{主要和特选的神}{chief and select gods}。\par

% structure: p0018 source=h2 target=subsection reason=relative-heading-rank; base=h2

\subsection{第四章 从\Person{瓦罗}的论述可知，众神的崇拜者视\OriginalTerm{人事}{human things}比\OriginalTerm{神事}{divine things}更为古老。}

在这整个极其美妙、极其精微的划分与区别的体系中，任何内心不固执而与自己为敌的人，都可以从我们已经说过的以及接下来将要说的内容里，轻而易举地看出，寻求并盼望永生，甚至奢望永生，都是徒劳的。因为这一切制度要么出自人类，要么出自魔鬼——不是他们所谓的善灵，更明白地说，是出自那些\OriginalTerm{不洁的灵}{unclean spirits}、无疑是\OriginalTerm{邪恶的灵}{malign spirits}的，他们以惊人的狡诈和隐秘，将有害的见解灌输给不敬神者的思想，有时还公然呈现于他们的认知中，使得人的心灵日益愚蠢，无法适应并持守那不变而永恒的真理，并企图以他们所能运用的种种\OriginalTerm{欺诈的证明}{fallacious attestation}来巩固这些见解。这位\Person{瓦罗}自己就证实，他先写人事，后写神事，是因为城邦先存在，然后这些事才由城邦设立。然而真宗教并非由任何属地的城邦所设立，相反，它建立了\OriginalTerm{天主之城}{celestial city}。这城乃是受\OriginalTerm{真天主}{true God}的默启与教导，真天主将\OriginalTerm{永生}{eternal life}赐予真正的崇拜者。\par

以下是\Person{瓦罗}给出的理由，他承认自己先写了关于人事，后写了关于神事，因为这些神事是由人设立的：—— \Quote{正如画家先于所绘的板，工匠先于建筑物，国家也先于那些由国家所设立的事物。} 但他却说，如果他是在写关于诸神的全部本性，他就会先写关于诸神，后写关于人——似乎他真的是在写关于诸神本性的某部分而不是全部；或者，似乎诸神本性的某部分（尽管不是全部）确实不应该被置于人的本性之前。那么，在最后三卷书中，当他详细解释那些确定的神、不确定的神和特选的神时，他如何似乎没有忽略诸神本性的任何部分？那么，他为什么又说，\Quote{如果我们写的是诸神的全部本性，我们就会先完成神事，然后再涉及人事呢？} 因为他写的是要么关于诸神的全部本性，要么关于其某部分，要么根本不涉及任何部分。如果是关于全部本性，它当然应该置于人事之前；如果是关于某部分，它为什么不根据事物的本性而先于人事呢？难道诸神的一部分不比全人类更应优先吗？但如果认为将神圣事物的一部分置于所有人事之前太过分，那么至少这一部分应当优先于\Place{罗马}人。因为他写关于人事的书，不是针对整个世界，而只是针对\Place{罗马}；他说，在写作顺序上，他把这些书恰当地置于关于神事的书之前，就像画家先于所绘的板，或工匠先于建筑物，最公开地承认，这些神事就像图画或建筑物一样，是由人设立的。只剩下第三种假设，即他应被理解为根本没有写关于任何神的本性，但他不愿公开这样说，而是留给聪明人去推断；因为当他说\Quote{不是\Emph{全部}}时，通常理解为\Quote{\Emph{部分}}，但\Emph{也可以}理解为\Emph{毫无所有}，因为毫无所有既不是全部也不是部分。事实上，正如他自己所说，如果他写的是关于诸神的全部本性，那么在写作顺序上，它本应出现在人事之前。然而，正如真理所表明的，即使\Person{瓦罗}沉默不言，神的本性也应该优先于罗马的事物，即使它不是\Emph{全部}，而只是\Emph{部分}。但它却被恰当地放在后面，因此它是\Emph{毫无所有}。因此，他的这种安排，不是由于他想把人事置于神事之前，而是由于他不愿让虚假的事物优先于真实的事物。因为在关于人事的部分，他遵循了事物的历史；但在关于他们所谓的神圣事物部分，他除了对虚妄事物的纯粹猜测外，还遵循了什么呢？无疑，这就是他以微妙的方式想要表明的；他不仅在写神事时将其放在人事之后，甚至还给出了这样做的理由；因为如果他隐瞒了这一点，或许有人会以某种方式为他的做法辩护，也有人会以另一种方式辩护。但在他给出的那个理由中，他没有留下任何让人随意猜测的余地，并充分证明了他是把人置于人的制度之上，而不是把人的本性置于神的本性之上。因此他承认，在写关于神事之书时，他所写的不是属于本性的真理，而是属于错误的不实；他在别处更公开地表达了这一点（正如我在第四卷中提到的），他说，如果他自己创立一座新城，他会按照本性的次序来写；但由于他找到的只是一座旧城，他不得不遵循它的习俗。\par

% structure: p0021 source=h2 target=subsection reason=relative-heading-rank; base=h2

\subsection{第5章 论瓦罗所述三种神学，即神话神学、自然神学与公民神学。}

对于他的这个主张，即神学有三种，也就是关于诸神的叙述，其中一种称为神话神学，另一种称为自然神学，第三种称为公民神学，我们该说什么呢？如果拉丁语用法允许，我们应当称他放在首位的那种神学为\Emph{传说的}，但让我们称它为\Emph{虚构的}；因为\Emph{神话的}一词源自希腊文\OriginalTerm{寓言}{\GreekText{μῦθος}}，意为寓言；第二种应称为\Emph{自然的}，现在的语言习惯允许这样称呼；第三种他本人用拉丁文命名为\Emph{公民的}。然后他说，\Quote{他们把诗人们主要使用的那种称为\Emph{神话的}；把哲学家们使用的那种称为\Emph{自然的}；把民众使用的那种称为\Emph{公民的}。至于我提到过的第一种，}他说，\Quote{其中有许多虚构，与不死者的尊严和本性相违背。因为我们在其中发现，有一位神从头部生出，另一位从大腿生出，还有一位从血滴中生出；而且，我们还发现诸神偷盗、奸淫、服侍人；总之，在这种神学中，各种事情都被归于诸神，这些事不仅可能发生在任何人身上，甚至可能发生在最可鄙的人身上。}\ 的确，在他能够、敢做、认为可以不受惩罚的情况下，他毫不含糊地表明了，虚构的寓言对诸神的本性造成了多么大的伤害；因为他不是在谈论自然神学，也不是公民神学，而是神话神学，他认为可以自由地加以指责。\par

让我们现在看看他关于第二种神学说了什么。他说：\Quote{我所解释的第二种神学，} \Quote{是哲学家们留下许多著作讨论的；他们探讨诸如以下问题：有哪些神，他们在哪里，他们的种类和特性如何，他们自何时开始存在，或者是否永恒存在；他们是火，如\Person{赫拉克利特}所信；或是数，如\Person{毕达哥拉斯}；或是原子，如\Person{伊壁鸠鲁}所说；以及诸如此类，人们在学校院墙内比在市集广场上更容易听到的问题。} 他对这种他们称为\OriginalTerm{自然神学}{physical theology}且属于哲学家的神学，毫无指摘，只是他叙述了他们彼此间的争论，由此产生了众多歧异的派别。然而，他把这种神学从市集广场，即从民众中移开，却关进了学校。可是，那第一种，最虚伪、最卑下的神学，他却未从公民中移除。哦，民众的宗教之耳，其中甚至包括罗马人的耳朵，竟不能忍受哲学家们对诸神的争论！但当诗人吟唱、优伶搬演那些贬损不朽者尊严与本性的事情——不仅可能降临于某人，就连最卑鄙之人也难逃——他们不仅忍受，还乐于倾听。不仅如此，他们甚至认为这些事情能讨神明喜悦，并借此获得神明的垂青。\par

但有人或许会说，让我们区分这两种神学：\OriginalTerm{神话神学}{mythical theology}与\OriginalTerm{自然神学}{physical theology}——即寓言神学与自然神学——与我们正在谈论的\OriginalTerm{城邦神学}{civil theology}。他预见到这一点，自己已区分了它们。现在让我们看看他如何解释城邦神学本身。我确实明白为何应该将寓言神学区分出来，因为它虚假、卑下、不配。然而，要想将自然神学与城邦神学区分开来，这无异于承认城邦神学本身是虚假的。因为如果自然神学是真实的，它有什么过错而应被排除？而如果这所谓的城邦神学不是自然的，它又有何价值而应被接纳？事实上，这正是他先写人事、后写神事的原因；因为在神事上，他遵循的不是本性，而是人的制度。让我们看看他的城邦神学。\Quote{第三种，} 他说，\Quote{是城邦公民，特别是司铎，应当知道并执行的。从中得知应合宜地敬拜哪个神，应合宜地举行哪些神圣仪式和祭献。} 让我们继续看下文。\Quote{第一种神学，} 他说，\Quote{特别适合剧院，第二种适合宇宙，第三种适合城邦。} 谁看不出他把桂冠给了哪一种？显然给了第二种，即他前文所说的哲学家的神学。因为他证实这种神学关乎宇宙，而他们认为宇宙无比卓越。但是那两种神学，第一种和第三种——即剧院神学和城邦神学——他是把它们区分了，还是统一了？因为尽管我们看到城邦在宇宙之中，但我们并不因此便认为城邦中的事物都关乎宇宙。因为有可能在城邦中，根据错误的意见而被崇拜和信仰的事物，在宇宙中或宇宙之外并不存在。然而，剧院岂不在城邦之中？剧院难道不是国家所建？国家建立剧院，目的何在？难道不是为了上演戏剧？戏剧演出又属于哪一类事？岂不正是属于那些神圣事物——瓦罗这些著作以如此高超才华所论及的神事吗？\par

% structure: p0025 source=h2 target=subsection reason=relative-heading-rank; base=h2

\subsection{论\OriginalTerm{神话神学}{mythic theology}，即\OriginalTerm{寓言神学}{fabulous theology}，与\OriginalTerm{城邦神学}{civil theology}：驳\Person{瓦罗}}

哦，\Person{马尔库斯·瓦罗}！你最敏锐，无疑最博学，但仍不过是人，不是天主——你并未被\OriginalTerm{天主的神}{Spirit of God}提升以看见并宣告神圣事物。你确实看出，神圣事物应与人的琐碎和谎言分开；但你却害怕冒犯民众至为腐败的意见，以及他们在公共迷信中的习俗。这些意见和习俗，当你从各方面考虑时，你自己也觉察到，而且你的全部著作大声宣告，它们与神的本性格格不入，甚至与人类心灵脆弱所设想的、存在于这个世界元素中的那些神也格格不入。最卓越的人类才能在此能做什么？广博的学问，在这困惑中又能于你何益？你渴望敬拜\OriginalTerm{自然神}{natural gods}，却被迫敬拜\OriginalTerm{城邦神}{civil gods}。你已发现有些神是\OriginalTerm{神话神}{fabulous gods}，你对这些神尽情倾吐自己的想法；而不论你愿与否，你也连带溅湿了城邦神。诚然，你说寓言神适应剧院，自然神适应宇宙，城邦神适应城邦；然而，宇宙是天主的工程，城邦和剧院却是人的工程；而且，在剧院中遭取笑的神，正是在庙宇中受崇拜的神；你举行赛会以荣耀的神，也正是你向其献祭的神。如果你曾说有些神是自然的，另一些为人所立，那么你对这些事的判断将会多么更加自由、更加精妙？而那些被人设立的神，诗人的著作提供一套说法，司铎的记述提供另一套说法——这两者却因共犯虚假而如此亲近，以致都为魔鬼所喜悦，而魔鬼正是与真理之道为敌的。\par

因此，暂且搁下他们所称为自然神学的那种神学，因为稍后将加以讨论；我们要问：真的有人满足于从那些诗歌的、戏剧的、舞台上的诸神那里寻求永生的希望吗？绝对不可！愿真天主使我们远离如此疯狂而亵圣的愚妄！什么？难道永生要从那些喜好这些表演，并因这些重现了它们自身罪行的表演而得取悦的诸神那里求取吗？在我看来，没有人会堕落到如此轻率狂怒的不虔之境。所以，无论是通过神话神学还是城邦神学，都没有人能获得永生。因为前者通过虚构播种了关于诸神的卑鄙之事，后者则通过珍惜它们而收获；前者散播谎言，后者收集谎言；前者以虚假的罪行对待神圣之事，后者则将这些罪行构成的戏剧纳入神圣之事中；前者在人的歌谣中宣扬关于诸神的不敬虚构，后者则将这些虚构分别为诸神自己的节庆而祝圣；前者歌唱诸神的恶行与罪愆，后者则喜爱它们；前者或创作或虚构，后者则或证实真实之事，或乐于虚假之事。两者都是卑鄙的，两者都该受谴责。然而，那戏剧性的神学公开教导可憎之事，而那城邦的神学却以那可憎之事装饰自己。永生难道要从这些污染了短暂今生的事物中期待吗？恶人的团体如果渗透进我们的情感，并赢得我们的赞同，难道不会污染我们的生命吗？那么，那些以其自身的罪行而被崇拜的魔鬼团体，难道不污染生命吗？——如果罪行是真实的，那么魔鬼是何等邪恶！如果罪行是虚假的，那么崇拜是何等邪恶！\par

当我们说这些话时，或许对于那些对这些事极其无知的人来说，会以为唯有那些在诗人的歌谣中传唱、在舞台上表演的关于诸神之事，才不配神圣的威严，荒唐可笑，更是可憎得不宜举行庆典；而那些并非由演员，而是由祭司们举行的神圣仪式，却是纯洁无瑕、毫无不当之处的。如果真是这样，那么就绝不会有人想到该用这些剧场中的可憎之举来向诸神致敬，诸神也绝不会命令人们为它们举行这些表演。然而，人们在剧场里表演这些事却毫不羞愧，因为类似的事在神庙里也在进行。总之，当前面提及的那位作者试图将城邦神学从神话神学和自然神学中区分出来，作为第三种独特的神学时，他希望人们理解，城邦神学其实是调和了二者，而非与任何一方截然分离。因为他曾说，诗人所写的东西，低于民众应遵循的标准；而哲学家所谈的，却又超出了民众宜于深究的范围。他说：\Quote{它们如此不同，以至于双方仍有不少内容被纳入了城邦神学的范畴；因此我们将指出城邦神学与神话神学的共通之处，尽管它本应与哲学家的神学联系更为紧密。}由此，城邦神学与神话神学并非完全脱节。尽管如此，在另一处，论及诸神的谱系时，他说民众更倾向于诗人而非自然神学家。因为在那里他说的是应当做的事，在另一处说的则是实际发生的事。他说后者是为了实用而写作，而诗人则是为了娱乐。因此，诗人的作品中那些民众不应效仿的事，就是诸神的罪行；而这些罪行却娱乐了民众和诸神。因为他说，诗人们是为了娱乐而写作，而非为了实用；然而他们所写的那些事情，恰是诸神所欲，民众所表演。\par

% structure: p0029 source=h2 target=subsection reason=relative-heading-rank; base=h2

\subsection{第7章 论神话神学与城邦神学的相似与一致}

因此，那种神话的、演剧的、舞台的、充满了一切卑鄙与不合宜的神学，被纳入了公民神学；那总体上理应被判定为应受谴责与摈弃的神学，其中的一部分却被宣布为值得培植和遵行——根本不是我所试图证明的那种不协调的部分，并非一个与整体身体格格不入、被不当附加和悬挂于其上的部分，而是一个完全和谐并与其余部分最协调一致的部分，犹如同一身体上的一个肢体。因为，那些神祇的形像、形态、年龄、性别、特征表明了什么？倘若诗人们笔下的\Person{朱庇特}有胡须而\Person{墨丘利}无须，祭司们岂非也如此？\Person{普里阿普斯}，祭司的比演员的更不淫秽吗？他接受崇拜者朝拜的形态，是否不同于他在舞台上取悦观众的姿态？在神像矗立的庙宇中，\Person{萨图恩}老年而\Person{阿波罗}年轻，与演员面具所表现时岂不一样？为何掌管门的\Person{福库鲁斯}与掌管门槛和门楣的\Person{利门提努斯}是男神，而介于他们之间、掌管铰链的\Person{卡耳德亚}却是女神？那些在神祇之事的典籍中记载的事情，难道不是庄重的诗人们认为不配写入诗行吗？剧院的\Person{狄安娜}佩带武器，而城中的\Person{狄安娜}仅是一位贞女？舞台上的\Person{阿波罗}是竖琴手，而\Place{德尔斐}的\Person{阿波罗}对此技艺一无所知？但与更可耻的事情相比，这些还算体面。那些将\Person{朱庇特}的乳母安置在\Place{卡匹托尔山}的人，他们自己对\Person{朱庇特}作何想法？他们岂非为\Person{欧赫墨罗斯}作证？\Person{欧赫墨罗斯}并非以寓言家的饶舌，而是以一位勤勉考察此事的史学家的严肃，撰述说所有这类神明都曾是凡人、终有一死。而那些任命\Person{埃普洛涅斯}为\Person{朱庇特}餐桌的寄生者的人，除了模仿神圣仪式之外，他们还能指望什么？因为如若某个喜剧演员说，\Person{朱庇特}的餐桌上有寄生者为他效劳，他定然显得意在博人一笑。\Person{瓦罗}说过这话——不是在嘲讽时，而是在颂扬诸神时说过的。他所写的神祇之事而非人事的书籍证明他写过这些话——不是在叙述舞台竞赛之处，而是在解释\Place{卡匹托尔}法律之时。总而言之，他被征服了，并且承认，由于他们按人的形象造了神，所以他们相信神也喜欢人的快乐。\par

因为，那些邪恶精灵也并非如此疏于己务，以至于竟不将有害的见解转化为戏谑，以此在人们心中巩固这些见解。由此也产生了那个关于\Person{赫丘利}守庙人的故事，故事说：他闲来无事，以掷骰子作消遣，交替地用一只手替\Person{赫丘利}掷，另一只手替自己掷，约定若他赢了，便用庙里的钱为自己备办晚餐，并雇一名情妇；但若\Person{赫丘利}赢了，他就要自费为\Person{赫丘利}提供同样的享乐。然后，当他被自己、仿佛是被\Person{赫丘利}击败时，他便把欠\Person{赫丘利}的晚餐，连同那最高贵的娼妓\Person{拉伦提娜}，一并献给了这位神明。而她呢，在庙里睡着之后，梦见\Person{赫丘利}与她交合，并对她说，她将在离开庙宇时遇到的第一个青年那里得到酬报，她应当相信这笔报酬是\Person{赫丘利}付给她的。于是，她出门时遇到的第一个青年便是富有的\Person{塔鲁提乌斯}，他长期供养她，临终时立她为继承人。她既得了极丰厚的财产，为免显得对神明的酬报忘恩负义，便转而立罗马人民为继承人，认为这最蒙众神悦纳；而她消失后，遗嘱被找到了。他们说，因着这立功的行为，她获得了神明的荣耀。\par

如果这些事情是诗人虚构并由模仿者表演的，无疑会被认为属于\OriginalTerm{神话神学}{fabulous theology}，并且认为应排除在\OriginalTerm{公民神学}{civil theology}的尊严之外。但是，当这些可耻的事情——并非出自诗人，而是民众；并非出自模仿者，而是神圣事物；并非出自剧场，而是庙宇，即非属神话神学，而属公民神学——被如此伟大的作者记述下来时，演员以戏剧艺术表现众神如此巨大的卑劣，并非徒劳；但祭司们试图通过所谓神圣的祭礼来表现他们那并不存在的高尚品格，则肯定是徒劳的。有\Person{朱诺}的祭礼，在她所爱的\Place{萨摩斯岛}上举行，她在那里嫁给了\Person{朱庇特}。有\Person{刻瑞斯}的祭礼，其中寻找被\Person{普路托}劫走的\Person{普洛塞庇娜}。有\Person{维纳斯}的祭礼，其中哀悼她心爱的\Person{阿多尼斯}，这位俊美的青年被野猪的牙齿杀死。有\Person{众神之母}的祭礼，其中她所爱的俊美青年\Person{阿提斯}因女人的嫉妒而被她阉割，遭到同样不幸、被称为伽利的人们哀悼他。既然这些事情比舞台上的所有可憎之事更加丑恶，他们为何竭力要将诗人关于众神的神话虚构与公民神学（他们希望公民神学属于城邦）分开，仿佛将可鄙下贱之物与高贵可敬之物分开呢？因此，更有理由感谢舞台演员，他们顾惜人们的眼睛，没有用戏剧表演将庙宇墙壁内隐藏的一切都暴露出来。当那些被公之于众的祭礼如此可憎时，那些隐藏在黑暗中的祭礼又能有什么好呢？当然，他们自己也亲眼看到了他们通过那些被阉割且柔弱的人们秘密进行的事情。然而他们未能遮掩这些人悲惨而可耻地被削弱和败坏的事实。让他们尽力说服那些能相信的人，通过这样的人，他们从事的是神圣事务吧，而他们无法否认，那些与他们神圣事物为伍、寄居其间的人正是如此。我们不知道他们做了些什么，但我们知道他们是通过什么人做的；因为我们知道在舞台上所表演的事情，即使在妓女的合唱队里，也从未出现过一个被阉割或柔弱的人。然而，这些事情仍由卑贱无耻的角色表演；的确，它们本不该由品格高尚的人表演。那么，那些祭礼又是什么呢？神圣性竟为了履行它们而选择了这样的人，甚至连舞台的淫秽都未曾接纳！\par

% structure: p0033 source=h2 target=subsection reason=relative-heading-rank; base=h2

\subsection{论外教教师为他们的神祇所力图提供的那类由自然解释构成的诠释}

但是，他们说，所有这些事情都有某些物理的，即自然的解释，展现其自然的含义；仿佛我们在这场辩论中是在寻求物理学，而不是神学，神学不是关于自然的，而是关于天主的。因为虽然那位是真正的天主者是天主，不是由于意见，而是由于本性，但并非整个自然都是天主；因为确实有人、兽、树、石的自然本性，其中没有一个是天主。因为，如果当询问到\Person{众神之母}的时候，整个解释体系的出发点实际上是众神之母就是大地，为什么我们还要进一步探寻？为什么我们还要考察其余一切？有什么能更明显地支持那些说所有这些神都是凡人的人呢？因为他们是\Quote{地生的}，意思是大地是他们的母亲。但在真正的神学中，大地是天主的工程，而不是母亲。但是，无论他们的祭礼可以怎样解释，以及它们与事物的本性有何关联，人变为柔弱无男子气概，这既不符合本性，反倒违反本性。这种病态、这种罪行、这种可憎之事，在他们的神圣事物中占有一席之地，即使堕落之人受到酷刑也几乎不会被逼承认自己犯有此事。再者，如果这些被证明比舞台上的可憎之事更加污秽的祭礼，因它们有自己的解释，并借此表明它们象征了事物的本性，而得以宽宥和辩护，那么为什么诗人的事物不以同样的方式获得宽宥和辩护呢？因为甚至对这些事物，也有许多人以类似的方式解释过，以至于就连他们所说的最怪异、最恐怖的事——即\Person{萨图恩}吞吃自己的子女——一些人也解释为：时间（由\Person{萨图恩}这个名字所表示）吞噬它所产生的一切；或者，如\Person{瓦罗}所认为的，萨图恩属于那些落回它们所从出的大地的种子。于是，有人这样解释，有人那样解释。对这门神学的其余一切，也是如此。\par

然而，它仍被称为\OriginalTerm{神话神学}{fabulous theology}，并连同其所有相关解释一同受到谴责、摈弃和排斥。不仅被哲学家们的\OriginalTerm{自然神学}{natural theology}所谴责，也被我们正谈论的这种据说关乎城邦与人民的\OriginalTerm{公民神学}{civil theology}判定为应予摈弃，因为它捏造了关于神明的卑鄙无稽之事。关于此事，我知道其中的秘密：那些写作这些著作的最敏锐且博学之士明白，这两种神学---即神话神学与公民神学---都当受摈弃；然而他们敢于摈弃前者，却不敢摈弃后者；他们将前者公诸谴责，而后者则被指出与前极其相似；这样做并非为了让人择取后者而放弃前者，而是为了让人明白后者同样应当一并受摈弃。于是，在那些不敢谴责公民神学的人毫无危险的情况下，两者皆被轻蔑，他们所谓的自然神学便可在素质更好的心灵中找到位置；因为公民神学与神话神学同为神话，也同为公民。凡睿智地细察二者之虚妄与秽行者，必发现它们皆是神话；而凡留心注意在公民神祇的节庆与城邦神圣礼仪中属于神话神学的戏剧表演者，将发现它们皆是公民。既然如此，赐予永生的能力又怎能归给那些神明---其偶像与神圣礼仪在其形象、年岁、性别、性格、婚配、谱系、礼仪各方面都证实它们与公开受摈弃的神话神明最为相似？在这一切事上，它们或被视为凡人，其神圣礼仪与节庆乃是按各人平生或临终之状而设立，\OriginalTerm{魔鬼}{demons}暗示且证实这种谬误，或实为最污秽的灵，利用各种时机潜入人心以行欺骗？\par

% structure: p0036 source=h2 target=subsection reason=relative-heading-rank; base=h2

\subsection{论神明的特殊职司}

至于那些神明被如此卑微琐细地分配的职司，以致人们认为必须根据各自特殊职能向每一位祈求---关于此点我们已经谈论不少，尽管尚未尽述---这些岂不是更符合模仿小丑的滑稽戏而非神圣威仪吗？如果有人为他的婴儿雇用两名保姆，一名只供食物，另一名只供饮料，正如这些人为此目的使用两位女神---\OriginalTerm{厄杜卡}{Educa}和\OriginalTerm{波提娜}{Potina}，那他必定显得愚蠢，并在家中行那配得上丑角的事。他们声称\OriginalTerm{利柏尔}{Liber}得名于\Quote{解放}，因为通过他，男性在交合时因排出精液而获得解放。他们还说\OriginalTerm{利柏拉}{Libera}（他们认为与\OriginalTerm{维纳斯}{Venus}相同）对女性执行同样的功能，因为他们说女性也排出精液；为此缘故，同一部位---男性与女性的---被安置在神庙中，男性的归利柏尔，女性的归利柏拉。此外他们还增加了归属利柏尔的女侍，以及激发情欲的酒。因此，酒神节以极度的疯狂来庆祝，关于此事，\Person{瓦罗}本人也承认，除非心神极度激越，酒神的信徒们不会做出这等举动。然而，这些事后来令较为清醒的元老院反感，遂下令禁止。在此，他们或许终于认识到，被视为神明的邪秽之灵对人心拥有何等大能。当然，此类事情不应在剧院中演出；因为那里是表演，而非发狂，尽管以喜好这类表演的神明为乐，本身就与发狂极其相似。\par

然而，他在虔敬者与迷信者之间所作的那种区分是何等荒谬？他说神祇为迷信者所畏惧，却被虔敬者如同父母般尊敬，而非当作敌人畏惧；又说所有神祇都如此良善，以致他们更愿意饶恕不虔敬者，而不愿伤害一个无辜者。然而他却告诉我们，有三位神祇被派作产妇的守护者，以防神祇\OriginalTerm{西尔瓦努斯}{Silvanus}闯入骚扰她；为表明这些守护者的临在，夜间有三个人绕行房屋，先用斧头击打门槛，接着用杵，第三次用扫帚扫刷，这样，在显示了这些农耕象征物之后，神祇西尔瓦努斯便被阻止进入，因为树木非斧不能砍伐或修剪，谷物非杵不能研碎，亦非扫帚不能聚拢。据此三种事物，三位神祇得名：\OriginalTerm{因特尔齐多纳}{Intercidona}，得名于斧头造成的切口；\OriginalTerm{皮鲁姆努斯}{Pilumnus}，得名于杵；\OriginalTerm{狄维拉}{Diverra}，得名于扫帚---藉着这些守护神祇，产妇得以免受神祇西尔瓦努斯的侵害。如此，善意神祇的守护竟不能抵御一个恶作剧神祇的恶意，除非他们三对一，并似乎以农耕的对立象征物与之战斗，而那位林中居民，既粗野、恐怖且未开化。这就是神祇的无辜吗？这就是他们的和谐吗？这些便是城邦的救治之神，比剧院中遭人嘲笑的事物更为可笑？\par

当男女结合时，\OriginalTerm{朱加提努斯}{Jugatinus}神主管。好吧，这一点可以容忍。但已婚妇女必须被带回家：\OriginalTerm{多米杜库斯}{Domiducus}神也被祈求。为了让她在家中，\OriginalTerm{多米提乌斯}{Domitius}神被引进。为了让她与丈夫同住，\OriginalTerm{曼图尔娜}{Manturnæ}女神被使用。还需要什么？让人间的羞怯得以保留。让血肉之欲继续下去吧，尊重那羞耻的秘密。为什么卧室里挤满了一群神祇，就连伴郎都已离去？而且，挤满这些神祇，并非为了因有他们临在而更加顾及贞洁，而是为了借助他们的帮助，那本性上软弱的女子，因处境的新奇而战栗，可以更容易地献出她的童贞。因为那里有\OriginalTerm{维尔吉尼恩西斯}{Virginiensis}女神，有神父\OriginalTerm{苏比古斯}{Subigus}，有神母\OriginalTerm{普雷玛}{Prema}，有\OriginalTerm{佩尔通达}{Pertunda}女神，还有\OriginalTerm{维纳斯}{Venus}和\OriginalTerm{普里阿普斯}{Priapus}。这算什么？如果一个人在做这件事时绝对需要众神的帮助，难道一个神或一个女神还不够吗？单是维纳斯还不够吗？据说她正是由此得名，因为没有她的力量，女人并不会停止成为处女。如果人心中有羞耻，而神祇却没有，那么，当新婚夫妇相信如此众多男男女女的神祇都在场，忙于此事时，他们是否会因羞耻而大为困窘，以致男人兴致减退，女人更不情愿呢？而且，诚然，如果维尔吉尼恩西斯女神在场是为了松开处女的腰带，如果苏比古斯神在场是为了让处女被置于男人之下，如果普雷玛女神在场是为了让她被置于男人之下后，被压住不动，那么佩尔通达女神在那里做什么？让她脸红；让她出去。让丈夫自己来做些事吧。除了他自己以外，任何人来做那使她得名的事，都是可耻的。但也许她得到容忍，因为据说她是一位女神，而非男神。因为如果她被视为男性，并被称为佩尔通达斯，那么丈夫为了保护妻子的贞洁，会要求更多的帮助来对抗他，胜过刚生产的女人对抗\OriginalTerm{西尔瓦努斯}{Silvanus}。但我为什么说这个呢，就连普里阿普斯也在那里，一个过度的男性，根据主妇们最尊贵、最虔诚的习俗，新婚新娘被命令坐在他那巨大而极为丑陋的器官上？\par

让他们继续下去，让他们尽其所能地狡辩，试图区分公民神学与神话神学，区分城市与剧场，区分神殿与舞台，区分祭司的圣物与诗人的歌曲，如同区分尊贵与卑贱、真实与虚妄、严肃与轻浮、庄重与滑稽、值得追求与应予摒弃的事物，我们明白他们所做的。他们知道，那种戏剧与神话神学依附于公民神学，并从诗人的诗歌中如同从镜子里反射回公民神学；因此，他们不敢谴责那种神学本身，便将它暴露出来，更随意地抨击和指责它的形象，为的是让那些领会他们用意的人憎恶这形象所反映的那张面孔——然而，众神自己，仿佛从同一面镜中看到自己，却如此喜爱它，以至于从两者中更能看清他们是谁、是什么。由此，他们也以可怕的命令迫使崇拜者将神话神学的污秽奉献给他们，列入他们的庆典，并视为神圣之事；这样，他们既更明显地表明自己是最不洁的精灵，也使那种被拒斥和摈弃的戏剧神学成为这种似乎被拣选和认可的公民神学的一个部分，以至于尽管整体是可耻、虚假、包含虚构神祇的，一部分见于祭司的文献，另一部分见于诗人的歌曲。它是否还有其他部分，那是另一个问题。目前，我认为，基于\Person{瓦罗}的划分，我已充分表明，公民神学和戏剧神学同属一种公民神学。因此，由于两者同样可耻、荒谬、不堪、虚假，虔诚的人绝不要指望从任何一种中获得永生。\par

最后，就连\Person{瓦罗}自己，在他对众神的叙述和列举中，也是从人的受孕那一刻开始的。他以\OriginalTerm{雅努斯}{Janus}为起点，开始列举那些掌管人类的神祇，一直延续到衰老之人死亡，并以\OriginalTerm{奈尼亚}{Nænia}女神作为终结，这位女神在老人的葬礼上被咏唱。此后，他开始叙述其他神祇，其管辖范围不是人本身，而是人的所有物，如食物、衣物以及此生所需的一切；对于所有这些神祇，他解释了他们各自的专职为何，以及应当为何事祈求他们。然而，尽管他如此谨细周全地勤奋探究，他却既未证明，甚至也未提及任何一位神的名称，能从祂求得永生——而这永生正是我们成为基督徒的唯一目标。那么，有谁会如此愚钝，竟看不出此人通过如此勤勉地展示和阐述公民神学，通过揭示它与那神话的、可耻的、不堪的神学的相似性，并且通过教导说那种神话神学也是公民神学的一部分，乃是在努力为那种唯独属于哲学家的自然神学在人们心中取得位置；他用如此巧妙的手法，指责神话神学，却不敢公开指责公民神学，仅仅通过展示它来表明其可指责之处；这样，当这两者都被有正确理解力的人所摒弃时，就只剩下自然神学可供选择了？但关于这一点，我们将在适当的地方，依靠真天主的帮助，更细致地加以讨论。\par

% structure: p0042 source=h2 target=subsection reason=relative-heading-rank; base=h2

\subsection{第十章。——论\Person{塞内加}的放肆，他比\Person{瓦罗}更猛烈地抨击公民神学。}

那个人所缺乏的那种自由——正是由于这种缺乏，他才不敢像公开抨击戏剧神学本身那样去抨击与戏剧神学极为相似的城邦神学——虽然并未完全，却已部分地为\Person{阿奈乌斯·塞内加}所享有，我们有证据表明他在我们\OriginalTerm{宗徒}{apostles}的时代活跃。我说他部分地拥有，因为他是在写作中拥有，而非在生活中。因为在他那本反对迷信的书中，他比\Person{瓦罗}更详尽、更猛烈地抨击了公民和城市神学，而非戏剧和寓言神学。当他谈到偶像时，他说：\Quote{他们用最没有价值、毫无生气的材料供奉神圣不可侵犯的不朽者的偶像。他们把这些偶像造成人、野兽和鱼的形象，有些还做成混合性别和异类身体。他们称之为神祇，而这些偶像如果获得生命，突然与人们照面，则会被视为怪物。}然后，过了一会儿，在赞美自然神学时，他阐述了一些哲学家的观点，然后对自己提出一个问题，说：\Quote{这里有人说，我能相信天和地是神，有些在月亮之上，有些在月亮之下吗？我要提出\Person{柏拉图}或逍遥学派的\Person{斯特拉托}吗？他们一个认为神没有身体，另一个认为神没有心智。}对此他回答说：\Quote{说实在的，在你看来，\Person{提图斯·塔提乌斯}、\Person{罗慕路斯}或\Person{图卢斯·霍斯提利乌斯}的幻梦，又有什么更真实呢？\Person{塔提乌斯}宣称女神\Person{克罗阿西娜}的神性；\Person{罗慕路斯}宣称\Person{皮库斯}和\Person{提贝里努斯}的神性；\Person{图卢斯·霍斯提利乌斯}宣称\Person{帕沃尔}和\Person{帕洛尔}的神性，这是人类最可憎的情感，一个是恐惧时心灵的激荡，另一个是身体的反应，实际上不是疾病，而是变色。}你宁愿相信这些是神祇，并将它们迎进天堂吗？但是，对于那些残酷可耻的仪式本身，他又是以何等自由来进行写作的啊！他说：\Quote{一个人阉割自己，另一个人割自己的手臂。那些在愤怒时需要威吓才能使之畏惧的诸神，当人们以这种方式谋求他们的好感时，他们从何找到畏惧的余地？然而，希望以这种方式受崇拜的神，根本不配受崇拜。心灵一旦被扰乱而脱离本位，就走火入魔至极，以至于人们用以取悦诸神的方式，就连最野蛮、以传说中极残酷闻名的暴君发泄怒火时也未曾用过。暴君曾撕裂人的四肢，却从未命令任何人撕裂自己的。为了满足王者的淫欲，有人被阉割；但从没有人因其主人的命令而对自己痛下毒手、自残形体。他们在神庙中自杀；他们以自己的伤口和血恳求。若有人得闲去看看他们所行和所遭受的，就会发现如此多不合体面人士、不配自由人、不像神智清醒者的行径，以至于无人不怀疑他们是疯了，假如只有少数人这样疯，那倒罢了；但如今这许多疯傻之人，反倒成了他们心智正常的辩护。}\par

他接着叙述了那些惯常在\Place{卡比托利欧山}进行的事情，并极其无所畏惧地坚称，那些事情只能令人相信若非出自嘻闹之人，便是疯人之手。因为他曾以嘲笑的口吻谈起\OriginalTerm{埃及宗教仪式}{Egyptian sacred rites}中哀悼丧亡的\Person{奥西里斯}，而一旦找到，又因他的重现而大喜过望，因为丧失与寻得都是假装的，然而那些毫无丧亡、也无所寻得的人由此引发的悲伤与喜乐，却是真实的。我这样说，他谈完这些后，说：\Quote{不过这种疯癫还有定时。一年疯一次，尚可容忍。走进\Place{卡比托利欧山}吧。有人向神祇传达神的命令；有人向\Person{朱庇特}报时；有人是\OriginalTerm{执法吏}{lictor}；有人是\OriginalTerm{涂油者}{anointer}，仅凭手臂动作模仿涂油。有些妇女为\Person{朱诺}和\Person{密涅瓦}梳理头发，她们站得远远的，不仅远离神像，甚至远离神庙。这些人以美发师的方式动着手指。有些妇女持着镜子。有些人呼唤诸神来诉讼中帮助自己。有些人把文件高持给神祇，向他们解释案情。一位博学杰出的\OriginalTerm{喜剧演员}{comedian}，年老衰迈，却每天在\Place{卡比托利欧山}上演\OriginalTerm{滑稽戏}{mimic}，仿佛诸神会乐见人们已经不再关心的东西。各类为不朽诸神工作的工匠，都闲居在那里。}过了片刻，他又说：\Quote{然而，这些人虽然为极无聊的目的献身于神，却并非为了任何可憎或无耻的目的。在\Place{卡比托利欧山}，有某些女子自以为得到\Person{朱庇特}的宠爱；即便面对那位——如果你相信诗人所言——最为善妒的\Person{朱诺}的目光，她们也毫不惧怕。}\par

这种自由，\Person{瓦罗}并未享有。他似乎只谴责诗性神学。对于公民神学，尽管他将其撕得粉碎，却不敢大胆抨击。但若我们关注真理，就会发现举行这些事的庙宇远比表演它们的剧场更加败坏。因此，对于公民神学的这些神圣礼仪，\Person{塞内卡}主张，智者最好的做法是在行为上佯装尊敬，心中却毫无敬意。\Quote{所有这些事，}他说，\Quote{一位智者将依照法律的命令而遵行，并非因为它们是取悦神祇的。}稍后他又说，\Quote{此外，我们还为神祇联姻，甚至违背自然，因为我们让兄弟姐妹结合？我们将\Person{贝娄娜}嫁给\Person{玛尔斯}，将\Person{维纳斯}嫁给\Person{伏尔甘}，将\Person{萨拉喀亚}嫁给\Person{奈普顿}。有些神祇我们任其未婚，仿佛没有匹配的对象，这实在多此一举，尤其当有某些未婚的女神，如\Person{波普洛尼亚}、\Person{富尔戈拉}或女神\Person{鲁米娜}，我不奇怪她们缺少追求者。所有这些可鄙的神祇群体，乃历代迷信所积聚，我们应当，}他说，\Quote{以这样的方式崇拜，同时始终牢记，这种崇拜更多属于习俗而非真实。}因此，公民神学所制定的那些法律与习俗，既非取悦神祇，亦非关乎真实。然而，这位被哲学变得仿佛自由了的人，因身为罗马人民的显赫元老，竟然崇拜他所谴责的、践行他所定罪的、敬拜他所指责的；因为，哲学确曾教给他一件大事——即在世间不陷于迷信，而是为了城市的法律和人的习俗，当一名演员，不是在舞台上，而是在庙宇中——这种行为更该受谴责，因为他如此欺诈地表演，以至于民众以为他是在真诚地行事。然而，舞台上的演员宁愿通过表演戏剧取悦观众，也不愿用虚假外表欺骗他们。\par

% structure: p0046 source=h2 target=subsection reason=relative-heading-rank; base=h2

\subsection{第11章 \Person{塞内卡}对犹太人的看法}

在公民神学的其他迷信中，\Person{塞内卡}也挑剔犹太人的圣物，尤其是安息日，断言他们守这些第七日是无益的，因藉着闲散，他们失去了生命的约七分之一，而且许多需要即刻关注的事也受到损害。然而，对于已经与犹太人极其敌对的基督徒，他却不敢提及，无论褒贬，唯恐若赞美他们，便是违背祖国的古老习俗，或者，若指责他们，可能是违背自己的意愿。\par

当他论及那些犹太人时，他说：\Quote{与此同时，那个最可咒骂的民族的风俗已取得如此力量，如今竟在各地被接受，被征服者已将法律加给了征服者。}他以这些话表达惊讶；且因不知\OriginalTerm{天主眷顾}{providence of God}正引他如此说，便用直白的话补充了意见，借此表明他对那些神圣制度之意义的看法：\Quote{因为，}他说，\Quote{然而，他们知道其礼仪的缘由，而民众大部分却不知为何举行他们的礼仪。}至于犹太人的庆节，它们究竟为何、在多大程度上由神圣权威所设立，而后又在适当时候由同一权威从天主之民身上移除——这永生的奥秘已启示给这子民——我们已在别处论及，尤其在驳斥摩尼教徒时，并打算在本著作中更合适的地方再作阐述。\par

% structure: p0049 source=h2 target=subsection reason=relative-heading-rank; base=h2

\subsection{第12章 万民神祇的虚妄一经揭露，便无可置疑，他们既不能于此生之事提供援助，亦不能将永生赐予任何人}

现在，既然有三种神学，希腊人分别称之为神话神学、自然神学和政治神学，在拉丁语中可称为神话神学、自然神学和公民神学；既然从神话神学（即使是众多伪神的崇拜者自己也极自由地谴责它），也从公民神学（该神学被证实是神话神学的一部分，甚至比之更糟），都不能指望从这些神学中获得永生——如果有人认为本书所言尚不足够，让他再添加上前几卷，尤其是第四卷中，关于\OriginalTerm{天主}{God}作为幸福赐予者的众多且各种论述。\par

因为如果\OriginalTerm{福}{Felicity}是一位女神，人们若不向福献身，又该向谁献身呢？然而，既然它并非女神，而是天主的恩赐，那么我们若不将自己献与那赐予幸福的天主，又该献与谁呢？我们虔诚地爱慕永生，在其中有着真正圆满的福。但我想，根据前文所述，无人应怀疑，那些以如此可耻的方式被敬拜的神明中，没有一个赐予幸福；而且，若不以那种方式敬拜他们，他们便更可耻地发怒，从而承认自己是最污秽的魔鬼。再者，不能赐予幸福者，怎能赐予永生？因为我们所说的永生，是指那有无穷幸福的生命。因为如果灵魂处于永罚之中——那些\OriginalTerm{不洁的魔鬼}{unclean spirits}也要在其中受折磨——那与其说是永生，不如说是永死。因为没有比死而不死的死亡更大更坏的死亡了。但由于灵魂按其本性，被造得不死不灭，不可能没有某种生命，因此它的极端死亡就是在永罚中与天主的生命隔绝。所以，只有赐予真正幸福的那一位，才赐予永生，也就是无尽幸福的生命。既然那些\OriginalTerm{公民神学}{civil theology}所敬拜的神明已被证明不能赐予这种幸福，那么，正如我们在前五卷中已指出的，他们便不应为那些现世和地上的事物受敬拜，更不应为了死后的永生而受敬拜——这正是我们在此卷中特别要证明的，其它各卷也对此有所贡献。但是，既然\OriginalTerm{积习}{inveterate habit}的力量根深蒂固，如果有人以为我论证得尚不够充分，不足以表明这种\OriginalTerm{公民神学}{civil theology}应当被弃绝和回避，那么请他注意另一卷书，它将在天主的助佑下接续本卷。\par

\SourceRecord{Translated by Marcus Dods.；Nicene and Post-Nicene Fathers, First Series；Vol. 2.；Edited by Philip Schaff.；Buffalo, NY: Christian Literature Publishing Co.,；1887.；<http://www.newadvent.org/fathers/120106.htm>.}

% structure: p0001 source=h1 target=section reason=page-title

\section{\WorkTitle{天主之城}（卷七）}

本书表明，通过对\Person{雅努斯}、\Person{朱庇特}、\Person{萨图恩}以及公民神学中其他\Quote{精选的神}的崇拜，并不能获得永生。\par

% structure: p0003 source=h2 target=subsection reason=relative-heading-rank; base=h2

\subsection{序言}

那些理解力较快较强的人——对他们来说，前几卷书已足够，甚至绰绰有余以实现预期目的——应当耐心平和地容忍我，在我试图以超常的勤奋根除那些败坏而古老的、与虔诚的真理相悖的观念时，这些观念因人类长久的错误，业已深深植根于未蒙光照的心灵；同时，我也按我微薄的力量，与那位真天主的恩宠合作，他能够成就此事，而我的工作也依赖他的帮助；并且，为了他人的缘故，他们不应认为自己已经不再需要的东西是多余的。当真正的、真正神圣的天主被推荐给人，作为他们应当寻求和崇拜的对象时——不是为了短暂必死生命的转瞬即逝的烟雾，而是为了永生，唯有永生才是真福，尽管我们也从他获得现世脆弱生活所需的帮助——这是一件极为重大的事。\par

% structure: p0005 source=h2 target=subsection reason=relative-heading-rank; base=h2

\subsection{第一章：既然公民神学中明显找不到神性，我们是否应相信神性可在精选的神中找到？}

如果说有人，我最后完成的第六卷尚未使他相信，这神性或所谓的神格——我们的作者也不犹豫地使用这个词，以更准确地翻译希腊人所称的\OriginalTerm{神性}{\GreekText{θεότης}}——我说，如果说有人，第六卷尚未使他相信，这神性或神格在那种他们称为公民神学、并由\Person{马库斯·瓦罗}在十六卷书中阐述的神学中是找不到的——也就是说，通过崇拜那些国家已确立为崇拜对象的神，是不可能获得永生真福的，并且是以那种形式崇拜——也许，当他读完本卷时，对于澄清这个问题，他将不再有任何进一步的要求。因为可能有人会认为，至少那些精选的、主要的神——瓦罗将他们编在最后一卷中，而我们尚未充分讨论——是应当为了真福生命（无非是永生）而崇拜的。关于此事，我不说\Person{特土良}所说的、也许机智多于真实的话：\Quote{如果神像洋葱一样被挑选，那么其馀的当然被视为坏而丢弃。}我不说这话，因为我看到，即便在精选的神之中，也有一些因某种更大更卓越的职事而被选出：正如在战争中，当新兵被选出后，又从其中选出一些来执行更重大的军事任务；在教会中，当有人被选为主教时，其馀的人并没有被弃绝，因为所有好基督徒都当之无愧地被称为选民；在建造房屋时，角石被选出，而其他用于结构其他部分的石头并未被弃绝；葡萄被选出食用，而留下酿酒的并未被弃绝。无需列举许多例证，因为此事显而易见。因此，从众多神中挑选某些神，并不能提供任何正当理由，让人轻看写这个主题的人、敬拜神的人，甚或神本身。我们更应寻求知道这些神是哪些，以及他们看来是为何目的被选出的。\par

% structure: p0007 source=h2 target=subsection reason=relative-heading-rank; base=h2

\subsection{第二章：哪些是精选的神，以及他们是否被认为免于普通神的职分。}

以下神，瓦罗确认为精选的神，专注于这一主题用了一卷书：\Person{雅努斯}、\Person{朱庇特}、\Person{萨图恩}、\Person{革尼乌斯}、\Person{墨丘利}、\Person{阿波罗}、\Person{马尔斯}、\Person{伏尔甘}、\Person{尼普顿}、\Person{索尔}、\Person{俄耳库斯}、父\Person{利贝尔}、\Person{忒卢斯}、\Person{刻瑞斯}、\Person{朱诺}、\Person{卢娜}、\Person{狄安娜}、\Person{密涅瓦}、\Person{维纳斯}、\Person{维斯塔}；在这二十位神中，十二位是男神，八位是女神。这些神被称为精选，是因为他们在世界中管理的职分更高，还是因为他们更为民众所知、受到更多崇拜？若是因为他们在世间执行的更大工作，我们就不该在那一群仿佛平民的神祇——他们被指派管理琐碎细微之事——中找到他们。因为，首先，从胎儿受孕开始，所有被细致分配给许多神祇的工作便开始了，\Person{雅努斯}本人为种子的接受开辟道路；那里还有\Person{萨图恩}，是关乎种子本身；有\Person{利贝尔}，他以种子的流溢释放男性；有\Person{利贝拉}，他们也会让她等同于\Person{维纳斯}，她将同样的益处赋予女性，即她也因种子的流溢而得到释放——所有这些都属于那些被称为精选的神。但也有女神\Person{梅娜}，她掌管月经；尽管是\Person{朱庇特}的女儿，却并不高贵。而同一位作者在他关于精选的诸神的书中，将月经这一领域也分配给了\Person{朱诺}本人，她甚至是精选的诸神中的皇后；而在这里，作为\Person{朱诺·卢西娜}，她与同一位\Person{梅娜}——她的继女——一起掌管同样的血液。也有两位极其晦暗的神，\Person{维图姆努斯}和\Person{森提努斯}——一个赋予胎儿生命，另一个赋予感觉；而且，尽管他们如此卑微，他们赐予的却远多于所有那些高贵而精选的神所赐予的。因为，的确，没有生命和感觉，妇女胎中怀有的整个胎儿，不过是极其卑贱无用的东西，不比泥浆和尘土好多少？\par

% structure: p0009 source=h2 target=subsection reason=relative-heading-rank; base=h2

\subsection{第三章：当许多低等神被赋予管理更崇高职事时，为何无法给出挑选某些神的理由。}

那么，究竟是何原因，竟驱使如此众多的特选神祇去掌管这些极其微小的职能，而在这上，他们却被那些虽鲜为人知且湮没无闻、却赋予了生命与感觉如此慷慨礼物的\Person{维图姆努斯}和\Person{森提努斯}所超越？因为特选的\Person{雅努斯}赐予入口，仿佛为种子开了一扇门；特选的\Person{萨图恩}赐予种子本身；特选的\Person{利贝尔}赐予男子同一种子的射出；而\Person{利贝拉}，即\Person{刻瑞斯}或\Person{维纳斯}，则将这同样赐予女子；特选的\Person{朱诺}（并非独自，而是与\Person{朱庇特}之女\Person{梅娜}一起）赐予月经，以利所孕之物的生长；而那默默无闻、地位卑微的\Person{维图姆努斯}赐予生命，默默无闻、地位卑微的\Person{森提努斯}赐予感觉——这最后两样东西，比其他那些要卓越得多，正如它们自身又逊于理智与智力。因为那些具有理智与领悟能力的事物，要优于那些没有智力与理智、如牲畜一般仅仅活着与感觉着的事物；同样，那些被赋予了生命与感觉的事物，也理应优先于那些既无生命又无感觉的事物。因此，赐予生命的\Person{维图姆努斯}，以及赐予感觉的\Person{森提努斯}，本应被归入特选神祇之列，而不应是接纳种子的\Person{雅努斯}、赐予或撒播种子的\Person{萨图恩}，以及促使并释放种子的\Person{利贝尔}和\Person{利贝拉}；这种子本身并不值得一顾，除非它达到了生命与感觉。然而，这些特选的恩赐却并非由特选神祇赐予，而是由某些籍籍无名、且就其尊严而言遭到忽视的神祇所赐。但如果有人回答说，\Person{雅努斯}掌管一切开端，因此将开辟受孕之路的职能分配给他并非没有道理；而\Person{萨图恩}掌管一切种子，因此撒播那使人得以生成的种子不能排除在他的职掌之外；\Person{利贝尔}和\Person{利贝拉}有权掌管一切种子的射出，因此他们也管辖那些与人类生育相关的种子；\Person{朱诺}掌管一切净化与分娩，因此她也有责任照管女人的净化与人的诞生——如果他们这样回答，那么就让他们在有关\Person{维图姆努斯}和\Person{森提努斯}的问题上找到答案吧，看他们是否愿意让这两位也同样掌管一切有生命和有感觉的事物。倘若他们承认这一点，那就让他们看看，他们即将把这两位置于何等崇高的地位。因为源于种子是在地上、属于地上的事，而生活与感觉却被认为是甚至诸星辰神祇也拥有的特性。但如果他们说，只有那些在肉身中得生命并由感官支持的事物，才被分配给\Person{森提努斯}，那么那位使一切事物生活与感觉的\Emph{天主}，为什么不也赐予肉身生命与感觉，在他那遍及万有的运作中，也将这份恩赐赋予胎儿呢？那么，\Person{维图姆努斯}和\Person{森提努斯}又有何用呢？可是，如果说这些仿佛是最边缘、最低等的事，是由那位统摄生命与感觉的\Emph{天主}，交给了这些神祇，如同交给仆役一般，那么这些特选神祇难道竟如此缺乏仆役，以至于找不到任何神祇可委派这些事务，却要凭藉他们全部的高位——他们似乎就因这高位而配得特选——而被迫与那些卑微的神祇一同履行其职吗？\Person{朱诺}是特选的神祇之后，是\Person{朱庇特}的姊妹与妻子；然而她却成了引导男童的\Person{伊特尔杜卡}，并与一对最卑微的女神——\Person{阿贝奥娜}和\Person{阿德奥娜}——一起履行此职。在那里他们还安置了女神\Person{梅娜}，她赐予男童良好的心智，而她并未被列入特选神祇之列；就好像能赐予人的还有什么比良好的心智更伟大的呢。但\Person{朱诺}之所以被列入特选，是因为她是\Person{伊特尔杜卡}和\Person{多米杜卡}（引领人在旅途上前行、并引领他重返家园的女神）；就好像，若人心智不良，出行一次并被引领回家，又有什么益处呢。然而，那位赐予此等恩赐的女神，却未被那些挑选者列入特选神祇，尽管她本当优先于\Person{密涅瓦}，因为在这细致的职能分配中，他们将男童的记忆力分配给了\Person{密涅瓦}。因为谁还会怀疑，拥有一副良好的心智，远比拥有再强大的记忆力要好得多呢？因为凡拥有良好心智的人，没有一个是坏的；而一些极坏的人却拥有令人赞叹的记忆力，他们越是无法忘记自己所想的恶事，就越是恶劣。然而，\Person{密涅瓦}却位列特选神祇之中，而女神\Person{梅娜}却被一群无足轻重的神祇所淹没。关于\Person{维耳图斯}，我又该说什么呢？关于\Person{斐丽西塔斯}，我又该说什么呢？——关于她们，我已在第四卷中多有论述；尽管人们将她们视为女神，却未认为宜给她们在特选神祇中安排一席之地，而在这些神祇中，他们却给了\Person{马尔斯}和\Person{奥尔库斯}一个位置，前者是死亡的制造者，后者是亡灵的接收者。\par

因此，既然我们看到，即便是那些被选中的神本身，也像元老院同人民协作一样，在其他神明当中共同完成那些被细致分配给众多神祇的细微工作；并且我们发现，某些未被认为配得上入选的神所管理的事务，反而比那些被称为\Quote{入选}的神所管理的更要伟大和美好；那么，我们就只能推测，他们被称为\Quote{入选}且\Quote{首要}的神，并非因为他们在世界中担任更高贵的职分，而是因为他们恰巧更为民众所知。就连\Person{瓦罗}本人也说过，某些父神和母神也难逃默默无闻的命运，正如人类中的某些人一样。因此，如果说\Person{福乐女神}或许不该被列入入选之神，因为他们并非凭功德，而是凭机遇才获得这等高贵地位，那么至少\Person{命运女神}应该位列其中，甚至应该排在他们之前；因为据说那位女神将所得的恩赐分配给每个人，并非依据任何理性的安排，而是全凭机遇的摆布。她本该在入选之神中占据至高的位置，因为正是在他们中间，她才特别彰显了自己所拥有的力量。我们看到，他们被选上，并非因某种杰出的美德或理性的幸福，而是由于\Person{命运女神}那种随机的力量——这些神的崇拜者以为她正是施展了这般力量。那位最雄辩的\Person{撒路斯提乌斯}或许正是想到了这些神，才说：\Quote{然而，实际上，命运主宰一切；它使万物闻名或湮没无闻，乃是出于任性而非依据真理。}他们实在找不出理由来解释，为何\Person{维纳斯}得以声名显赫，而\Person{美德女神}却默默无闻，毕竟两者的神性都曾受到他们庄严的认可，两者的功德更是无法相提并论。再说，如果她（\Person{维纳斯}）因备受追求而配得崇高地位——毕竟追求\Person{维纳斯}的人比追求\Person{美德女神}的人更多——那么，为何\Person{密涅瓦}备受尊崇，而\Person{金钱女神}却湮没无闻，虽然在整个人类当中，贪婪吸引的人远多于技艺？而且，即便在那些精通技艺的人里，你也很难找到不以谋取金钱利益为目的而从事本行的人；况且，任何事物，若它是其他事物的目的，其价值永远高于那为其他事物而存在的事物。那么，如果对神的这种拣选是基于愚昧大众的判断，为何\Person{金钱女神}没有被置于\Person{密涅瓦}之前，既然许许多多的工匠都是为了金钱？但如果这区分是由少数智慧之人作出的，为何\Person{美德女神}被置于\Person{维纳斯}之前，因为理性远远偏爱前者？无论如何，正如我已经说过的，\Person{命运女神}本身——按照那些赋予她最大影响力的人的说法，她使万物闻名或湮没无闻，乃是出于任性而非依据真理——既然她甚至能够对众神施加如此大的权能，凭着自己任性的判断，想使谁闻名便使谁闻名，想使谁湮没便使谁湮没；那么，\Person{命运女神}本人在入选之神中自当占据优越地位，因为即使对众神，她也拥有如此优越的权能。抑或我们应该设想，她之所以不在入选之列，原因仅仅在于，就连\Person{命运女神}自己也遭遇了厄运？那么，她是对自己不利了，因为她在使别人显贵的同时，自身却一直默默无闻。\par

% structure: p0012 source=h2 target=subsection reason=relative-heading-rank; base=h2

\subsection{四、— 那些名字不与丑闻相连的低等神祇，所受的对待优于那些丑行被传扬的入选之神。}

然而，任何一个热切寻求名望和声誉的人，可能会祝贺那些入选之神，称他们为幸运者，若非他看出他们被选上，更多地是给他们带来伤害而非荣耀。因为那群低等神祇，凭着自己的卑微和默默无闻，得以免于被丑闻吞没。的确，当我们看到他们仅仅依据人类意见的虚构，按照分配给各自的特有工作而被区分开来时，我们会发笑，就像是那些承包一小部分公共税收的人，或是银匠街上的工匠，在那里，一件器皿为了能完美出厂，要经过许多人之手，而它本来可以由一个完美的工匠一次完成。然而，之所以认为需要许多工匠联合的技巧，唯一的原因在于，让每一部分工艺由专门的工匠学习，这样能够迅速而容易地完成，总比要求每个人都精通一门工艺的所有部分要好，而那只能缓慢而艰难地达到。尽管如此，在非入选之神中，几乎找不到一个因任何罪行而给自己带来丑闻的，而在入选之神中，却几乎没有谁不曾烙上可耻的臭名。后者降格去从事前者的卑微工作，而前者却未曾高攀到后者的崇高罪行。关于\Person{雅努斯}，我一时想不起他有什么可耻的事；或许他就是这样一位，生活得比其余神祇更清白，远离过失和罪行。在\Person{萨图恩}逃亡时，他仁慈地接纳并款待了他；他与自己的客人分享王国，因此他们各自拥有一座城，一个是\Place{雅尼库鲁姆}，另一个是\Place{萨图尼亚}。然而，那些在崇拜神明方面寻求各种不体面之事的人，却以畸形怪诞的形象来丑化他——他们发现他的生活比其他神更少可耻——将其塑造成有时是双面，有时又仿佛是双重的四面。难道他们是想，既然大多数入选之神因犯下可耻罪行而失去廉耻，那么他那更大的清白，就该用更多的面孔来标示吗？\par

% structure: p0014 source=h2 target=subsection reason=relative-heading-rank; base=h2

\subsection{五、— 论异教徒更为隐秘的教义，以及关于自然解释的论述}

但让我们听听他们自己的物理学解释，他们试图以此粉饰最可悲错误的低劣，仿佛带着更深刻教义的表面。\Person{瓦罗}首先强烈推崇这些解释，说古人发明了神的形象、标志和装饰，是为了让那些参加秘仪的人用肉眼看见它们时，能用心灵的眼睛看见世界灵魂及其部分，也就是真神；而且那些以人形制作神像的人所意图的意义似乎是：凡人的心灵（在人体中的）与不朽的心灵非常相似，正如器皿可以被放置来代表神，例如，酒器可以放在\Person{利伯尔}的神庙里代表酒，通过容器来表示所盛之物。因此，通过具有人形的形象，就表示了理性的灵魂，因为人形可以说是器皿，他们归给神或众神的本性通常包含在其中。这就是那位最有学问的人深入探究并揭示出来的教义奥秘。但是，哦，你这最敏锐的人，你难道在这些奥秘中丧失了那种审慎吗？那种审慎使你形成清醒的看法：那些最初为人民设立神像的人剥夺了公民的敬畏并添加了错误，而且古代罗马人更纯洁地尊敬无像的神。正是由于考虑这些，你才有勇气对后来的罗马人说这些话。因为如果最古老的罗马人也崇拜偶像，恐怕你会因恐惧而默默压制所有关于设立偶像的愚蠢（尽管如此却是真实的）观点，并会更高声、更喋喋不休地颂扬那些由这些空虚有害的虚构构成的奥秘教义。你的灵魂，如此博学聪明（为此我为你深感悲痛），永远无法通过这些奥秘达到它的天主；即那位灵魂由祂而非与祂一起被造的天主，灵魂不是祂的一部分，而是祂的作品——那位不是万物灵魂，而是创造每一个灵魂的天主；每一个灵魂只有在祂的光中才是有福的，如果它不因忘恩负义而辜负祂的恩宠。\par

但本书随后内容将说明这些奥秘的本质以及应当赋予它们何种价值。同时，这位最有学问的人坦白他的看法：世界灵魂及其部分是真神，由此我们看出他的神学（即他很重视的那种自然神学）在其完整性中甚至能延伸到理性灵魂的本性。因为在这本（关于精选神的）书中，他通过预先说明自然神学讲了一些内容；我们将会看到他在那本书中是否能够通过物理学的解释，将公民神学——即他最后论及精选神时所说的——归结为这种自然神学。现在，如果他能够做到这一点，那么整个神学就是自然的；既然如此，又有什么必要将公民神学与自然神学如此仔细地区分开来呢？但如果这种区分是真实的区分，那么，既然连他极为称赞的这种自然神学都不是真实的（因为它虽然触及了灵魂，却没有触及那创造灵魂的真天主），那种主要关注形体之物的公民神学又是多么可鄙和虚假，正如其解释本身将表明的，这些解释是他们如此勤奋地寻找并阐明的，我不得不提及其中一些！\par

% structure: p0017 source=h2 target=subsection reason=relative-heading-rank; base=h2

\subsection{论\Person{瓦罗}的观点：天主是世界灵魂，然而世界在它的不同部分有许多灵魂，其本性是神圣的。}

同一个\Person{瓦罗}仍然在预先说明中说他认为天主是世界灵魂（希腊人称之为\OriginalTerm{宇宙}{\GreekText{κόσμος}}），并且这个世界本身就是天主；但正如一个智慧的人虽然由身体和心灵构成，却因心灵而被称为智慧的，同样，世界因心灵而被称为天主，尽管它由心灵和身体构成。这里他似乎在某种程度上至少承认了一位天主；但为了引入更多的神，他补充说世界分为两部分：天和地，它们每部分又分为二：天分为以太和空气，地分为水和陆地。在这四部分中，以太最高，空气第二，水第三，陆地最低。他说，所有这些部分都充满灵魂；以太和空气中的灵魂是不死的，水和陆地上的灵魂是有死的。从天的最高处到月球的轨道有灵魂，即星辰和行星；这些不仅被理解为神，而且被看到是这样。在月球轨道与云风区域开始之间，有空中的灵魂；但这些是用心灵看见的，而不是用眼睛，它们被称为英雄、\Person{拉瑞斯}和\Person{格尼乌斯}。这就是在预先说明中简要阐述的自然神学，它不仅让\Person{瓦罗}满意，也让其他许多哲学家满意。在我完成了关于公民神学——就它涉及精选神而言——尚需论述的内容后，必须借助天主的帮助更仔细地讨论它。\par

% structure: p0019 source=h2 target=subsection reason=relative-heading-rank; base=h2

\subsection{将\Person{雅努斯}和\Person{特米努斯}区分为两位不同的神是否合理。}

那么，\Person{瓦罗} 以之开端的 \OriginalTerm{雅努斯}{Janus} 究竟是谁？他就是世界。这倒是一个极简明确而不含糊的回答。既然如此，为何他们又说事物之始属于他，而事物之终则属于另一位被称作 \OriginalTerm{特尔米努斯}{Terminus} 的神呢？因为他们说有两月献给这两神，与始终相关——一月献给 \OriginalTerm{雅努斯}{Janus}，二月献给 \OriginalTerm{特尔米努斯}{Terminus}——这还不包括从三月起至十二月止的其他十个月。他们说这就是为什么在二月庆祝 \OriginalTerm{特尔米那利亚节}{Terminalia}，而在同一个月也举行那被他们称为 \OriginalTerm{净化礼}{Februum} 的神圣净化仪式，月名亦由此而来。那么，事物之始既属于即 \OriginalTerm{雅努斯}{Janus} 的世界，难道事物之终就不属于这世界，却另立一神司掌吗？他们难道不承认，他们所说始于这个世界的一切，也在这个世界内终结吗？既然在他的雕像上给了他两张脸，在施为上却只给他一半权能，这是何等愚妄！若说 \OriginalTerm{雅努斯}{Janus} 与 \OriginalTerm{特尔米努斯}{Terminus} 即是同一个，其一面脸指向开始，另一面指向终结，这岂不优雅得多？因为行事之人理应两面皆顾。每一举措若不回顾开端，也不前瞻终局，如何能成？因此前瞻之意向必与回顾之记忆相连。若已忘怀所开始者是什么，又如何知道如何完成？但如果他们以为，真福生命始于此世，而圆满于世界之外，并因此将开始之能仅归于 \OriginalTerm{雅努斯}{Janus}，即归于世界，那么他们本应更加推重 \OriginalTerm{特尔米努斯}{Terminus}，而不应将他排斥于选神之列。即便现在，当这两神代表时间性事物的始终之时，也本应给予 \OriginalTerm{特尔米努斯}{Terminus} 以更多尊荣。因为更大的喜乐在于事物完成之时；然而起头之事未至完成，总令人忧心忡忡；而那开始者所深切渴慕、心之所向、期盼、想望的正是此终局；若非达于完成，无人在其起始时便真地欢喜。\par

% structure: p0021 source=h2 target=subsection reason=relative-heading-rank; base=h2

\subsection{第八章 崇拜 \OriginalTerm{雅努斯}{Janus} 者为何造其双面之像，而有时又使其有四面的原因}

不过，现在且让他们拿出那尊双面像的解释来。因为他们说它有两张脸，一前一后，乃因我们张开的嘴巴似乎与世界相似：希腊人称上颚为 \OriginalTerm{天空}{\GreekText{οὐρανός}}，而据他说，某些拉丁诗人也曾称天为 \OriginalTerm{上颚}{palatum}；他们又说，从这张开的口中，朝着牙齿的方向有一出路，朝着食道的方向有一入路。瞧瞧，为了一个希腊词或一个与上颚有关的诗词用语，世界被低贬到了何等地步！若为唾液的缘故而崇拜这位神祇，那么在 \OriginalTerm{上颚}{palatum} 之天穹下，唾液的确有两个敞开的门户——一个可以吐出一部分，另一个可以咽下一部分。此外，还有什么比这更荒谬的呢：不在世界本身中寻找两个互相对着的门户——世界可以由此纳入或排出事物——却从我们的喉咙和食道，去找那与世界毫无相似之处的东西，来为 \OriginalTerm{雅努斯}{Janus} 编造一个世界的形象，只因据说世界与 \OriginalTerm{上颚}{palatum} 相似，而 \OriginalTerm{雅努斯}{Janus} 与那上颚却毫无相像？然而，当他们把他做成四面，称之为双重的 \OriginalTerm{雅努斯}{Janus} 时，他们解释说这是指世界的四方，仿佛世界像 \OriginalTerm{雅努斯}{Janus} 那样，透过四张脸向外观看。再者，若 \OriginalTerm{雅努斯}{Janus} 就是世界，而世界由四方构成，那么双面的 \OriginalTerm{雅努斯}{Janus} 像便是错的。若它是真的——因为整个世界有时只用东和西来理解——那么，当南北也被提及时，会有人称世界为双重吗？正如他们称有四面时的 \OriginalTerm{雅努斯}{Janus} 为双重一样？他们根本无法像解释双面 \OriginalTerm{雅努斯}{Janus} 那样，为世界解释出入的四个门户，在双面像那里，他们好歹在人的口中找到了某种与他相应的东西；除非 \OriginalTerm{尼普顿}{Neptune} 来给他们帮忙，递给他们一条鱼：那条鱼除了嘴和食道之外，还有鳃孔，两侧各一。然而，纵有这许多门户，除非是那听见真理说 \BibleQuote{我是门。} \BibleRef{若 10:9} 的灵魂，没有灵魂能逃脱这种虚妄。\par

% structure: p0023 source=h2 target=subsection reason=relative-heading-rank; base=h2

\subsection{第九章 论 \OriginalTerm{朱庇特}{Jupiter} 之能，并 \OriginalTerm{朱庇特}{Jupiter} 与 \OriginalTerm{雅努斯}{Janus} 之比较}

但他们也指明，他们希望人如何理解 \OriginalTerm{尤维}{Jove}（亦称 \OriginalTerm{朱庇特}{Jupiter}）。他们说，他是那位有能力掌管万事万物因何而成之神。\Person{维吉尔} 那句极为尊贵的诗就为此神之伟大作证：\par

\Quote{知晓事物之因者何其有福。}\par

但为何 \OriginalTerm{雅努斯}{Janus} 被置于他之先？让那位极为敏锐、博学多识的人回答我们这个问题。他说：\Quote{因为，\OriginalTerm{雅努斯}{Janus} 管辖最初之物，\OriginalTerm{朱庇特}{Jupiter} 则管辖至高之物。因此，\OriginalTerm{朱庇特}{Jupiter} 被公认为万有之王，实至名归；因为至高之物优于最初之物：虽然最初之物在时间上居先，至高之物却在尊严上超绝。}\par

如果人们将已完成之事的第一部分与最高部分区分开来，这种说法原本是可以正确理解的；例如，启程是行动的第一部分，到达则是最高部分。开始学习是已开始之事的第一部分，掌握知识则是最高部分。所有事物都是如此：开端在先，目的在最高。然而，关于此问题，已结合\Person{雅努斯}与\Person{忒耳弥努斯}讨论过了。但归因于\Person{朱庇特}的那些原因乃是产生结果者，而非被产生的事物；它们不可能在时间上被已造成或已完成的事物，或被此类事物的开端所阻止；因为产生者总是先于被产生之物。因此，尽管已造成或已完成之事的开端属于\Person{雅努斯}，它们却并不先于他们所归因于\Person{朱庇特}的动力因。因为正如若无动力因在先，便无任何事物发生，同样，若无动力因，也无任何事物开始发生。的确，如果人们称这位神为\Person{朱庇特}，一切受造本性与一切自然事物的所有原因都在他的权能之下，却以如此的侮辱和臭名昭著的指控来崇拜他，他们所犯的亵渎之罪，比完全否认任何神存在更为可憎。因此，对他们来说，更好的做法是用\Person{朱庇特}之名去称呼另一位神——一位配得上卑劣和罪恶之荣誉的神；以某种虚幻之物代替\Person{朱庇特}（据说\Person{萨图恩}曾得到一块石头代替儿子吞下），让他们以此作为亵渎的对象，而不要说成\Emph{那}位神既打雷又犯奸淫——统治整个世界，又放纵自己犯下如此多的淫荡行为——掌握着自然以及一切自然事物的最高原因，自身的原因却毫无良善。\par

接下来，我要问，如果\Person{雅努斯}是世界，那么他们还能在众神中为这位\Person{朱庇特}找到什么位置呢？因为\Person{瓦罗}将真神定义为世界的灵魂及其各个部分。因此，根据他们的说法，凡不属于这一定义的，就肯定不是真神。那么，他们会说\Person{朱庇特}是世界的灵魂，而\Person{雅努斯}是世界的身体——即这个可见的世界吗？如果他们这样说，他们就不可能断言\Person{雅努斯}是神。因为即便按照他们的说法，世界的身体也不是神，而是世界的灵魂及其各个部分才是。因此，\Person{瓦罗}看到这一点后说，他认为\OriginalTerm{天主}{God}是世界的灵魂，而这个世界本身就是\OriginalTerm{天主}{God}；但正如一个智慧之人虽由灵魂和身体构成，却因灵魂而被称为智慧，同样，世界虽由灵魂和身体构成，却因灵魂而被称为\OriginalTerm{天主}{God}。所以，世界的身体单独不是\OriginalTerm{天主}{God}，或只是灵魂是\OriginalTerm{天主}{God}，或灵魂和身体一起是\OriginalTerm{天主}{God}，但如此成为\OriginalTerm{天主}{God}不是凭借身体，而是凭借灵魂。因此，如果\Person{雅努斯}是世界，且\Person{雅努斯}是一位神，那么为了能让\Person{朱庇特}也是一位神，他们是否会认为\Person{朱庇特}是\Person{雅努斯}的某一部分呢？因为他们更倾向于将普遍存在归于\Person{朱庇特}；因此有句话说：\Quote{一切事物都充满了朱庇特。} 因此，他们也必须认为\Person{朱庇特}是世界，这样他才能成为神，并且是众神之王，从而统治其他众神——按他们的说法，这些神是他的各个部分。同样，还是那位\Person{瓦罗}，在他专门撰写的那本关于崇拜诸神的著作中，对\Person{瓦勒里乌斯·索拉努斯}的某些诗句作了这样的解释。这些诗句如下：\par

全能的朱庇特，君王、万物与众神之祖，\newline{}兼为诸神之母，一神而万有。\par

但\Person{瓦罗}在同一著作中解释这些诗句时说，正如雄性排射种子，雌性接受种子，同样，\Person{朱庇特}——他们相信他就是世界——既从他自身排射所有种子，又将它们接纳入自身。他说，正因为这个原因，\Person{索拉努斯}才写道：\Quote{朱庇特，祖先与母亲；} 并且同样有理由说，一与万是同一的。因为世界是一，而在那\Quote{一}之中包含着万有。\par

% structure: p0031 source=h2 target=subsection reason=relative-heading-rank; base=h2

\subsection{第十章 雅努斯与朱庇特的区分是否恰当}

既然\Person{雅努斯}是世界，\Person{朱庇特}也是世界，为何\Person{雅努斯}与\Person{朱庇特}是两位神，而世界却只有一个？为什么他们有不同的神庙、不同的祭坛、不同的仪式、不同的形象呢？如果这是因为开端的本性是一种，原因的本性是另一种，而一种获得了\Person{雅努斯}之名，另一种获得了\Person{朱庇特}之名；那么，如果一个人拥有两种不同的权柄或两种技艺，是否就会因为权柄或技艺的本性不同而被说成是两个法官或两个工匠呢？对同一位神而言也是如此：如果他拥有开端的能力和原因的能力，难道就必须认为他是两位神，仅仅因为开端和原因是两件事吗？但如果他们认为这样是对的，那么也让他们断言，\Person{朱庇特}有多少别名，他就有多少位神，因为他有许多能力；而赋予他这些别名的那些事物本就是多样且众多的。我且举出其中的几个。\par

% structure: p0033 source=h2 target=subsection reason=relative-heading-rank; base=h2

\subsection{第十一章 论朱庇特的别名，这些别名并非指向许多神，而是指向同一位神}

他们称他为\OriginalTerm{胜利者}{Victor}、\OriginalTerm{不可战胜者}{Invictus}、\OriginalTerm{援助者}{Opitulus}、\OriginalTerm{推动者}{Impulsor}、\OriginalTerm{稳定者}{Stator}、\OriginalTerm{百足者}{Centumpeda}、\OriginalTerm{仰卧者}{Supinalis}、\OriginalTerm{梁木}{Tigillus}、\OriginalTerm{滋养者}{Almus}、\OriginalTerm{乳哺者}{Ruminus}，以及其他名号，不胜枚举。但这些附加的名号是因着不同的理由与能力而赋予同一位神的，却并未因这众多理由而迫使他成为同样多的神。他们给予他这些名号，是因为他征服万物；因为他不被任何事物征服；因为他援助贫困者；因为他有能力推动、止息、建立、掷翻；因为他如梁木般维系并支撑世界；因为他滋养万物；因为他如同乳头，哺育动物。在此，我们觉察到既有大事，也有小事；然而，据称这一切皆由同一位神成就。我认为，他们设想一个世界为两位神——\OriginalTerm{朱庇特}{Jupiter}与\OriginalTerm{雅努斯}{Janus}——的原因与根源，比起维系世界与为动物哺乳这两件事来，彼此之间更为相近；然而，因着这两件在性质与尊严上都相去甚远的工作，却未必要有两位神；同一位朱庇特，因这一件被称为梁木，因那一件被称为乳哺者。我不愿说，为吮乳动物哺乳本应更符合\OriginalTerm{朱诺}{Juno}而非朱庇特，尤其是在这项工作上有女神\OriginalTerm{卢弥娜}{Rumina}相助并侍奉她；因为我认为，人们可以回应说，朱诺本身无非就是朱庇特，依据\Person{瓦莱里乌斯·索拉努斯}的诗句，其中有言：\par

全能的\OriginalTerm{朱庇特}{Jove}，诸王、万物与众神的祖先，\newline{}亦是众神的母亲，等等。\par

那么，为何他被称为乳哺者，而倘若他们更仔细探究，或许会发现他亦是那位女神卢弥娜呢？\par

若在一穗谷粒上由一位神管理关节处，另一位神管理外壳，这被认为与神的尊严不相称；那么，一件最低等的事——为动物哺乳以滋养它们——竟由两位神照管，其中之一是朱庇特本人，万物之王，他并不与自己的妻子同做此事，而是与某位卑微的\OriginalTerm{卢弥娜}{Rumina}（除非或许他本人就是卢弥娜，为雄性时是乳哺者，为雌性时是卢弥娜），这岂非更与神的尊严不相称！我原本一定会说他们不愿将阴性名号加诸朱庇特，若非在这些诗句中他被称作\Quote{祖先与母亲}，若非在他其他的别号中读到\OriginalTerm{金钱}{Pecunia}一名，而我们在那些小神中发现了一位女神，正如我在第四卷中已提及的。但既然男女都有\OriginalTerm{金钱}{pecuniam}，他何以不既被称为\OriginalTerm{佩库尼乌斯}{Pecunius}，又被称为\OriginalTerm{佩库尼亚}{Pecunia}呢？那是他们的事。\par

% structure: p0038 source=h2 target=subsection reason=relative-heading-rank; base=h2

\subsection{第12章——朱庇特亦被称为金钱。}

他们为这名称所作的解释何等雅致！他们说：\Quote{他也被称为\OriginalTerm{金钱}{Pecunia}，因为万物都属于他。}哦，对一个神祇之名的解释何其堂皇！正是：那拥有万物的，却最卑微、最侮辱地被称作金钱。与天地所包含的万有相比，人在金钱的名义下所拥有的一切总和又算什么呢？这个名号，无疑是贪婪加给朱庇特的，好让贪爱金钱的人自以为所爱的不是一位普通的神，而是万物之王本身。然而，若他被称为财富，那便大为不同了。因为财富是一回事，金钱是另一回事。我们称智者、义人、善人为富有，他们或毫无金钱，或仅有极少。他们因拥有德行而更为真正富有，因为凭藉德行，他们对于肉体所需之物，即使仅有眼前所有，也知足常乐。但我们称贪婪者为贫穷，他们永在渴求，永不知足。他们或许拥有极多金钱；但无论那金钱如何丰裕，他们仍难免匮乏。我们称天主本身为富有的，这理所当然；然而并非在金钱上，而是在全能上。因此，拥有大量金钱的人被称为富人，但若他们贪婪，内心却是贫穷的。同样，没有金钱的人被称为穷人，但若他们有智慧，内心却是富足的。\par

那么，智者当如何思想这种神学呢？在这种神学中，众神之王竟获得了一个\Quote{智者从不欲求}之物的名号。因为，若这套哲学关于永生有任何健全的教导，他们更应当称那位统管世界之神，不是金钱，而是智慧，爱智慧能涤除贪财——即爱金钱——的污秽！\par

% structure: p0041 source=h2 target=subsection reason=relative-heading-rank; base=h2

\subsection{第13章——当撒图尔努斯是何者、革尼乌斯是何者被阐明时，得出的结论乃是二者皆为朱庇特。}

然而，何必再多谈这位朱庇特呢？或许所有其他神祇都可与他等同；这样一来，当他成为一切时，多神存在的意见便仅是一种空洞无实的意见。倘若他不同的部分和能力被视为如此多的神，或者他们以为弥漫万物的心智原理，因这可见世界的整体所包含的不同部分，以及自然界纷繁的治理，而获取了众多神祇的名号，那么他们就全都归向于他。因为，撒图尔努斯又是什么呢？他说：\Quote{主要的神祇之一，他掌管所有的播种。}\Person{瓦莱里乌斯·索拉努斯}的诗篇的解释岂不教导说，朱庇特即是世界，他从自身发出一切种子，并收归于自身？\par

因此，掌管所有播种的正是他。革尼乌斯是什么呢？\Quote{他是被设立掌管万物，并拥有生育万物之能力的神。}除了世界，他们还会相信谁拥有这能力呢？对于世界，曾有言：\par

\Quote{全能的朱庇特，祖先与母亲？}\par

他在另一个地方说，\OriginalTerm{革尼乌斯}{Genius}是每一个人的理性灵魂，因而分别存在于各人身上，而世界相对应的灵魂则是天主，这恰恰又回到了同一件事——也就是说，世界的灵魂本身应当被视为仿佛是普世的革尼乌斯。因此，这便是他称之为\OriginalTerm{朱庇特}{Jupiter}者。因为如果每一个革尼乌斯都是一位神，每个人的灵魂又是一个革尼乌斯，那么每个人的灵魂便都是一位神了。然而，假如连这些神学家们自己也因这极端荒谬而不得不对此畏缩，那么剩下的便是，他们以特殊且卓越的区分将那一个革尼乌斯称为神，他们称其为世界的灵魂，也就是朱庇特。\par

% structure: p0046 source=h2 target=subsection reason=relative-heading-rank; base=h2

\subsection{第十四章 \OriginalTerm{墨丘利}{Mercury}与\OriginalTerm{玛尔斯}{Mars}的职司}

然而，他们未能找出如何将墨丘利和玛尔斯归属于世界的哪一部分，以及归属于元素中天主的工程；因此，他们至少将这两位神置于人类的工作之上，使他们成为说话和从事战争的助手。如今，墨丘利如果同时也拥有众神言语的权能，他也就辖管众神之王本身了——倘若朱庇特从他那里领受了说话的能力，并且按照朱庇特所乐意准许他的那样说话——这当然是荒谬的；但如果人们认为归属于他的，仅仅是对人类言语的权能，那么我们说，朱庇特竟肯俯就，不但给婴孩，也给野兽哺乳——他因此得名\OriginalTerm{鲁米努斯}{Ruminus}——却不愿让那使我们胜过野兽的言语之照管归属于他，这实在难以置信。这样一来，言语本身既属于朱庇特，又是墨丘利。然而，如果言语本身被说成就是墨丘利，正如那些以解释的方式论及他的话所表明的那样——因为据说他被称为墨丘利，意即\InnerQuote{在中间奔驰者}，因为言语在人与人之间奔驰：他们还说，希腊人称他为\OriginalTerm{赫耳墨斯}{\GreekText{Ἑρμῆς}}。\par

\SourceRecord{Translated by Marcus Dods.；Nicene and Post-Nicene Fathers, First Series；Vol. 2.；Edited by Philip Schaff.；Buffalo, NY: Christian Literature Publishing Co.,；1887.；<http://www.newadvent.org/fathers/120107.htm>.}

% structure: p0001 source=h1 target=section reason=page-title

\section{\WorkTitle{天主之城}（卷八）}

\Person{奥斯定}现在讨论第三种神学，即自然神学，并探讨敬礼自然神学中的神明是否有助于获得来世的真福。他更愿意与柏拉图主义者讨论此问题，因为柏拉图的体系在哲学中享有\OriginalTerm{最卓越}{facile princeps}的地位，且与基督宗教的真理最为接近。在展开论证时，他首先驳斥\Person{阿普列乌斯}，以及所有主张应敬礼魔鬼，以之作为神与人之间的信使和中介的人；并证明，那些本身就是邪恶奴隶的魔鬼，不可能使人与良善的神和解；魔鬼所喜爱并支持的，正是善良明智之人所痛恨并谴责的——诗人的亵渎虚构、戏剧表演和魔术技艺。\par

% structure: p0003 source=h2 target=subsection reason=relative-heading-rank; base=h2

\subsection{第1章：自然神学的问题应同那些追求更卓越智慧的哲学家们讨论。}

我们须以远比前几卷书中解答和阐述问题时所需的更强烈的专注，来探讨目前这个问题；因为我们要与之讨论他们所谓的自然神学的，不是普通人，而是哲学家。因为这种神学既不同于神话的，即戏剧的；也不同于城邦的，即城市的：前者展示诸神的罪行，后者显露他们的罪恶欲望，这表明与其说他们是神，不如说是邪恶的魔鬼。我们说，要就这种神学与之进行讨论的，是哲学家——这个名称，若译成拉丁语，意指那些宣称热爱智慧的人。然而，如果智慧是天主，而天主创造万物，正如圣经权威和真理所证实的\BibleRef{智 7:24}，那么哲学家就是热爱天主的人。但是，由于这个名称所指的事物，并非存在于所有以此名号自傲的人身上——因为并非所有称为哲学家的人都是真智慧的爱好者——我们必须从我们通过阅读得以了解其意见的人中，选出一些我们尚且配与之讨论此问题的人。因为我在这部作品中并未试图驳斥哲学家们的所有虚妄意见，而只针对那些与\OriginalTerm{神学}{theologia}有关者；按我们对此希腊词的理解，神学意谓对神性本质的记述或解释。而且，我也并未试图驳斥所有哲学家的一切虚妄神学意见，而只针对这样一些人：他们虽然同意存在某种神性，并且这种神性关心人类事务，但却否认敬礼唯一不变的天主足以在死亡之后和现时获得真福生活；并且主张，为获得那种生活，应当敬礼许多神明，这些神明确由那唯一的天主所创造，并安置于各自的领域。甚至比\Person{瓦罗}，这些人也更接近真理；因为瓦罗毫无困难地将自然神学整体扩展至世界和世界灵魂，而他们却承认天主超越一切具有灵魂本性的事物，不仅是这个常被称为天地的可见世界的造物主，也是一切灵魂的造物主，并且祂通过使理性灵魂——人类灵魂即属此类——分沾祂自己不变的无体之光，而赐予真福。任何对这些事略有了解的人，无不知晓其名取自其师\Person{柏拉图}的柏拉图主义哲学家。关于这位\Person{柏拉图}，我将就目前问题所需，简要陈述，并先提及在时间上先于他而在同一文学领域中的前辈。\par

% structure: p0005 source=h2 target=subsection reason=relative-heading-rank; base=h2

\subsection{第2章：关于两个哲学学派，即意大利学派和爱奥尼亚学派，及其创始人。}

就希腊文献而言，其语言地位比任何其他民族的语言都更显赫，历史提到两个哲学学派：一个称为意大利学派，起源于\Place{意大利}从前称为\Place{大希腊}的地区；另一个称为伊奥尼亚学派，起源于那些至今仍称为\Place{希腊}的地区。意大利学派的创始人是\Person{萨摩斯的毕达哥拉斯}，据说\Quote{\OriginalTerm{哲学}{philosophy}}一词也起源于他。因为从前那些因生活方式值得称赞而看似出类拔萃的人被称为\OriginalTerm{贤者}{sage}，\Person{毕达哥拉斯}在被问及他自称什么时，回答说他是\OriginalTerm{哲学家}{philosopher}，即智慧的学生或爱好者；因为他认为自称贤者是最傲慢的。伊奥尼亚学派的创始人是\Person{米利都的泰勒斯}，他是被称作\Quote{\OriginalTerm{七贤}{seven sages}}的七人之一，其中六人以他们的生活方式和某些生活格言而闻名。\Person{泰勒斯}则以探究万物的本性而著称；为了使自己的学派后继有人，他将自己的论述写成文字。然而，使他特别突出的是他通过天文计算甚至能预测日食和月食的能力。不过，他认为水是万物的本原，世界的一切元素、世界本身以及其中所生的一切，最终都由此构成。然而，对这当我们观察世界时显得如此奇妙的一切工作，他没有加上任何\OriginalTerm{神圣理智}{divine mind}的性质。他的学生\Person{阿那克西曼德}继承了他，对万物的本性持有不同意见；因为他并不像\Person{泰勒斯}那样认为万物源自一个本原，即水，而是认为每一事物都源自其自身的本原。他相信这些事物的本原在数量上是无限的，并认为它们产生了无数个世界，以及其中所产生的一切事物。他还认为，这些世界处于永恒的交替解体与再生过程中，每个世界持续的时间长短依情况而定；而且，在产生这一切事物活动的过程中，他并不比\Person{泰勒斯}更多地归因于神圣理智。\Person{阿那克西曼德}留下他的弟子\Person{阿那克西美尼}作为继承人，后者将万物的所有原因归因于一种\OriginalTerm{无限的气}{infinite air}。他既不否认也不忽视诸神的存在，但远非相信气是由他们造成的，相反，他坚持认为他们由气产生。然而，他的弟子\Person{阿那克萨戈拉}认识到，神圣理智是我们所见一切事物的产生原因，并说所有不同种类的事物，依照各自的方式和种类，是由一种由\OriginalTerm{同质微粒}{homogeneous particles}组成的无限质料产生的，但通过神圣理智的作用。\Person{阿那克西美尼}的另一位弟子\Person{第欧根尼}说，某种气是万物的原初实体，万物由此产生，但这种气具有\OriginalTerm{神圣理性}{divine reason}，没有它，就什么也不能从其中产生。\Person{阿那克萨戈拉}的弟子\Person{阿凯劳斯}继承了他，他也认为万物由同质微粒构成，每一具体事物皆由此形成，但那些微粒被神圣理智所渗透，神圣理智永恒地推动所有永恒的物体，即那些微粒，使它们交替结合与分离。据说\Person{柏拉图}的老师\Person{苏格拉底}是\Person{阿凯劳斯}的弟子；正因为\Person{柏拉图}，我才对整个这些学派的历史做了这一简短概述。\par

% structure: p0007 source=h2 target=subsection reason=relative-heading-rank; base=h2

\subsection{第三章 苏格拉底的哲学}

据说\Person{苏格拉底}是第一个将哲学的全部努力指向纠正和规范风俗的人，在他之前的所有人都将最大努力耗费在探究物理的、即自然的现象上。然而，在我看来，无法确切发现\Person{苏格拉底}这样做，是因为他厌倦了晦涩而不确定的事物，因此希望将心思转向发现某种清晰而确定的事物，这是获得幸福生活所必需的——这似乎是所有哲学家的劳苦、警醒和勤勉所指向的那个伟大目标——还是（如一些更偏袒他的人所推想的）他这样做，是因为他不愿意让被尘世欲望玷污的心灵试图向上提升自己，去接触神圣事物。因为他看到，他们是在寻求事物的原因——他相信这些原因最终都可归结为唯一真实至高的天主的意志——因此，他认为只有洁净的心灵才能领悟这些原因；所以，一切努力都应用于通过良好道德净化生活，好使心灵摆脱私欲的压抑重负，能以本有的活力向上升腾，直达永恒事物，并能以洁净的理智默观那无形无像、永不改变的光明之性体，一切受造本性的原因都存在于其中。然而，显而易见的是，他以令人惊叹的优雅风格和论证，以及最为尖锐而委婉的机智，搜寻并追索愚昧人的愚蠢，这些人自以为知道这或那——有时承认自己的无知，有时假装不知，甚至在他似乎倾注全部心力的那些道德问题上也是如此。由此产生了对他的敌意，最终导致他被诬告，并被判处死刑。然而，后来，那个曾公开判处他的雅典城邦，又公开哀悼他——民众的愤慨如此强烈地转向控告他的人，以至于其中一人死于群众的暴力，另一人仅以自愿的永久流放才逃过类似的惩罚。\par

因此，\Person{苏格拉底}在生时与死后皆卓著非凡，其哲学留下了众多门徒，他们竞相在探讨那些关乎\OriginalTerm{至善}{summum bonum}的道德议题上取得精通，一旦拥有此至善，便能使人得享真福；而在苏格拉底的论辩中，他提出各类问题、作出断言、继而将其推翻，因此并未清楚显明他究竟以何为至善，于是每一个人便从这些论辩中撷取最中己意者，且各人将至善置于自身看来所构成者。那称为至善者，乃是人一达到便得享真福的事物。然而，苏格拉底的门徒们对此至善所持的意见竟分歧如此之大，以至于（对于同一师门的追随者，此事几乎难以置信）有人将至善置于享乐，如\Person{亚里斯提卜}，有人则置于美德，如\Person{安提斯泰尼}。实际上，逐一详述不同门徒的各样见解，难免令人厌烦。\par

% structure: p0010 source=h2 target=subsection reason=relative-heading-rank; base=h2

\subsection{第四章 论\Person{苏格拉底}弟子中最杰出的\Person{柏拉图}及其三重划分的哲学}

然而，在\Person{苏格拉底}的门徒中，\Person{柏拉图}光辉远胜其余，毫不冤枉地使他们全都黯然失色。他生为\Place{雅典}人，出身名门，天赋才智远超同门，且拥有惊人禀赋。但他认为自身以及苏格拉底的训练远不足以将哲学臻于完善，便尽可能广泛游历，凡以培育任何可掌握之学问而闻名的地方，他必前往。于是他向\Place{埃及}人学习了他们视为重要的一切所持所教；并从埃及转赴\Place{意大利}那些因毕达哥拉斯学派而声名洋溢的地区，在最卓越的教师门下，极轻易地掌握了当时全意大利盛行的哲学。由于他对自己的老师苏格拉底怀有特殊的爱，便让苏格拉底作他一切对话录的发言者，将他从别人所学，或靠自己强大理智努力所得，悉数借其口说出，甚至以苏格拉底风格的优雅文雅，润饰其伦理辩论。智慧之学在于行动与默观，因此一部分可称为行动的，另一部分为默观的——行动部分关乎生活行为，即道德的规范；默观部分在于探究自然的原因及纯粹的真理——据说苏格拉底长于此学的行动部分，而\Person{毕达哥拉斯}则更注重默观部分，他以伟大理智的全部力量致力于此。至于\Person{柏拉图}，则被誉为将两部分结合为一，从而使哲学得以完满。随后他将哲学分为三部分：第一部分为伦理的，主要关注行动；第二部分为自然的，其对象是默观；第三部分为理性的，辨别真伪。虽然这最后一部分对行动与默观均属必需，然而默观特别要求探究真理本性的职分。因此，这一三重划分并不违背那使智慧之学由行动与默观构成之区分。至于柏拉图对每一部分所持的观点——即他认何者为一切行动的目的、一切自然的原因、一切理智的光明——讨论起来过于冗长，我们不可妄加断言。因为，正如柏拉图喜好并常效法其师苏格拉底众所周知的方法，即掩饰自己的知识或意见，所以要清楚发现他对各种事物究竟如何想望，并不比发现苏格拉底的真实意见更容易。尽管如此，我们必须将他在著作中所表达的某些观点纳入我们的作品，不论是他自己说出，还是转述他人所言并似乎表示赞同的意见——这些观点有时有利于真宗教，那正是我们的信德所采纳并捍卫的，有时却与之相悖，例如关于存在一位天主还是多位神明的问题，因为这关系到死后将要到来的真正有福的生命。因为那些被誉为最佳追随柏拉图的人——理当被置于外邦所有其他哲学家之上——并且据说在理解他方面显示出最敏锐洞察力的人，或许对天主持有这样的观念：承认在祂内可找到存在的原因、理解的最终理由，以及规范整个生活的目的。这三者之中，第一者被认为属于哲学的自然部分，第二者属于理性部分，第三者属于伦理部分。因为，若人受造，正是要藉着他内最卓越的部分，达至那超越万有者——即那唯一真实且绝对善的天主，若无祂，则无一物存在，无一教训有教益，无一操练有裨益——那么当寻求那位在祂内万物于我们稳固者，当发现那位在祂内一切真理于我们确凿者，当爱慕那位在祂内一切于我们正直者。\par

% structure: p0012 source=h2 target=subsection reason=relative-heading-rank; base=h2

\subsection{第五章 论在神学问题上我们尤其应与\Person{柏拉图主义者}辩论，他们的观点优于所有其他哲学家}

那么，如果\Person{柏拉图}将智慧之人定义为那位效法、认识、爱慕这位\OriginalTerm{天主}{God}，并藉着与祂在其自身的真福中相通而蒙受祝福的人，又何必与其他哲学家讨论呢？显然，没有谁比\OriginalTerm{柏拉图主义者}{Platonists}更接近我们。因此，让那以众神的罪行取悦人心的\OriginalTerm{神话神学}{fabulous theology}向他们让步；也让那\OriginalTerm{城邦神学}{civil theology}让步，在其中，不洁的\OriginalTerm{魔鬼}{demons}假借众神之名，诱使那些沉溺于尘世欢愉的地上之民，它们渴望藉着人类的谬误而受尊崇，用不洁的欲望充满敬拜者的心灵，煽动他们将这些魔鬼的罪行表演作为崇拜礼仪的一部分，而魔鬼自己却在观赏这些表演的观众身上找到最令它们愉悦的景象——在这种神学中，圣殿中凡尊贵的，都因与剧场的淫秽混合而被玷污；剧场中凡卑贱的，都藉着圣殿的可憎之事而得到辩护。这些哲学家还必须让位于\Person{瓦罗}的解释，他在解释中称那些神圣礼仪涉及天地以及可朽之物的种子和运作；因为，首先，那些礼仪并不具有他想要人们相信其所附有的意义，因此，他如此解释它们的尝试并不伴随真理；而且，即使它们具有这种意义，那些在自然界等级中低于\OriginalTerm{理性灵魂}{rational soul}之物，仍不该被理性灵魂当作其天主来崇拜；灵魂也不该将那些被\OriginalTerm{真天主}{true God}赋予其优越地位之物，摆在它本身之上视为神。对于与神圣礼仪有关的那些著作，也必须持同样看法；\Person{努马·庞皮利乌斯}特意将它们与他本人一同埋葬，以期隐藏，后来它们被犁翻出，奉元老院之命烧毁。而且，出于对\Person{努马}的充分尊重，我们要提到，与这些著作地位相同的，还有\Person{马其顿的亚历山大}写给他母亲的信，信中所言是由一位\Place{埃及}的大司祭\Person{利奥}传达给他的。在这封信中，不仅\Person{皮库斯}和\Person{法乌努斯}、\Person{埃涅阿斯}和\Person{罗慕路斯}，甚或\Person{赫拉克勒斯}，以及由\Person{塞墨勒}所生的\Person{利柏尔}，还有\Person{廷达瑞俄斯}的双生子，或其他任何被神化的凡人，就连那些主要的神本身——\Person{西塞罗}在其\WorkTitle{图斯库卢姆论辩}中曾未提其名而暗示的诸位，如\Person{朱庇特}、\Person{朱诺}、\Person{萨图恩}、\Person{伏尔甘}、\Person{维斯塔}，以及\Person{瓦罗}试图将其等同于世界之部分或\OriginalTerm{元素}{elements}的许多其他神——都被揭示为凡人。如我们所说，此事与\Person{努马}之事有相似之处；因为那位司祭因泄漏了一项奥秘而惧怕，恳切请求\Person{亚历山大}下令他的母亲烧掉那封传达这些消息的信。因此，让这两种神学，即神话神学与城邦神学，向那些\OriginalTerm{柏拉图主义哲学家}{Platonic philosophers}让步吧；他们承认\OriginalTerm{真天主}{true God}是万物的创造者、真理之光的泉源、一切真福的慷慨赐予者。不仅如此，那些因心灵受制于身体而臆想\OriginalTerm{万物本原}{principles of all things}是\OriginalTerm{物质的}{material}的哲学家，也必须向这些对如此伟大的天主作出伟大承认的人让步。例如\Person{泰勒斯}，他主张万物的\OriginalTerm{第一本原}{first principle}是水；\Person{阿那克西美尼}认为是气；\OriginalTerm{斯多葛派}{Stoics}认为是火；\Person{伊壁鸠鲁}断言它由\OriginalTerm{原子}{atoms}，也就是微粒构成；还有许多无需一一列举的人，他们认为物体——无论单一或复合，有生命或无生命，但终究是物体——是万物之因和本原。因为他们中的一些人——例如\OriginalTerm{伊壁鸠鲁派}{Epicureans}——相信有生命之物能从无生命之物中产生；另一些人则认为，无论有无生命，万物都出自一个有生命的本原，然而，由于万物都是物质的，它们出自一个物质的本原。因为\OriginalTerm{斯多葛派}{Stoics}认为，火，即构成这个可见世界的四种\OriginalTerm{物质元素}{material elements}之一，是有生命、有理智的，是这世界及其所包含的万物的创造者——它实际上就是天主。这些人以及与他们类似的人，只能推断出他们那受感官奴役的心灵所虚妄地提示给他们的事物。然而，他们自身内部有着他们看不到的某物：即使没有亲眼看见，而只是在思考它们时，他们也能在内心向自己呈现他们在外部见过的事物。但这种思维中的呈现已不再是一个物体，而只是物体的相似物；而那种藉以看到这一物体相似物的心灵能力，既不是物体，也不是物体的相似物；而且，判断那呈现是美或丑的能力，无疑高于被判断的对象。这原则就是人的\OriginalTerm{理智}{understanding}，即\OriginalTerm{理性灵魂}{rational soul}；它当然不是物体，因为它所观看和判断的物体相似物本身就不是物体。灵魂既不是土、水、气、火——我们看见这个世界就是由这四种称作\OriginalTerm{四元素}{four elements}的物体构成的。如果灵魂不是物体，它的创造者天主又怎能是物体呢？那么，正如我们说过的，所有这些哲学家，连同那些虽耻于宣称天主是物体，却又认为我们的灵魂与天主有相同本性的人，都当向柏拉图主义者让步。他们并未因灵魂的巨大可变性——将这种属性归于神圣本性是亵渎——而感到不安，而是说，是身体改变了灵魂，因为灵魂本身是不可改变的。他们倒不如说：\Quote{肉体为某物体所伤，因为它本身刀枪不入。}简言之，不可改变者，什么都不能改变它，因此，能被身体改变的东西，便不能恰当地说是不变的。\par

% structure: p0014 source=h2 target=subsection reason=relative-heading-rank; base=h2

\subsection{第六章 论柏拉图主义者在被称为\OriginalTerm{自然哲学}{physical}的那部分哲学中的含义}

这些哲学家，我们看到他们并非不配地在声名与荣耀中被高举于众人之上，他们已看出，没有物质形体是天主，因此，他们在寻求天主时超越了一切形体。他们已看出，凡可变者皆非至高天主，因此，他们在寻求至高者时超越了一切灵魂和所有可变的神体。他们也看出，在每一个可变事物中，那使其成为其所是的形式，无论其样式或本性如何，只能藉着那位真正\Emph{存在}者而\Emph{存在}，因为祂是不变的。因此，无论我们观察世界的整个形体，它的形状、性质与有序的运动，以及其中的一切形体；或是观察一切生命，无论是那滋养并维持的生命，如树木的生命，还是那除此之外还有感觉的生命，如野兽的生命；或是那更兼具所有这些与理智的生命，如人的生命；或是那不需营养的支持，而仅维持、感觉、理解的生命，如天使的生命——所有这一切只能藉着那位绝对\Emph{存在}者而\Emph{存在}。因为对祂来说，\Emph{存在}并非一码事，活着又是一码事，好像祂能\Emph{存在}却不活着；活着也不是一码事，理解又是一码事，好像祂能活着却不理解；理解也不是一码事，享真福又是一码事，好像祂能理解却不享真福。而对祂来说，活着、理解、享真福，就是\Emph{存在}。他们从这种不变性与纯一性中领悟到，万物必定是由祂所造，而祂自己不可能被任何所造。因为他们曾考量，凡\Emph{存在}者或是形体，或是生命，而生命优于形体，且形体的本性是可感的，生命的本性是可理解的。因此，他们偏爱可理解的本性胜过可感的。我们所说的可感事物，是指那些能由形体的视觉和触觉所感知的事物；可理解的事物，则是指那些能由心灵的洞察所理解的事物。因为凡形体的美，无论是形体的静态，如形状，还是其动态，如音乐，没有不是由心灵判断的。但若心灵本身没有一种更超越的形式——没有体积，没有语音，没有空间和时间——这事就绝无可能。可是，就在这些事上，若心灵不是可变的，一个人就不可能就感性形式做出比另一个人更精准的判断。聪明人比迟钝者判断得更准，娴熟者比生疏者判断得更准，练达者比未练者判断得更准；同一个人在获得经验后，也比之前判断得更准。但凡能有程度高低者，就是可变的；因此，深入思考这些事的有才之士得出结论：那第一形式，不可能存在于那些形式可变的事物中。所以，他们既已看到形体与心灵在形式上可以有美或不美的程度差别，且倘若没有形式，便不可能存在，便看到一定有一种存在，其中有第一形式，不变，因此不容许有比较的程度；他们极正确地相信，那就是万物的第一本原，它被非造，而万物皆由它而造。因此，天主已将自己所能被人认识的事显明给他们了，因为祂那看不见的事，可从受造物中领悟而看见；祂的永恒德能与天主性，也是一样，一切有形可见并短暂的事物都是由祂所造。\BibleRef{罗 1:19} 关于他们称为物理学，即自然神学的那部分，我们已经说得够多了。\par

% structure: p0016 source=h2 target=subsection reason=relative-heading-rank; base=h2

\subsection{柏拉图主义者在逻辑学，即理性哲学方面，究竟应被认作何等超越其余哲学家}

此外，就他们称为逻辑学，即理性哲学的那部分教义而言，我们断不可将他们与那些人相提并论：后者将辨别真理的能力归于身体感官，认为我们所学的一切，都须用这些不可靠且易错的规则来衡量。这些人就是伊壁鸠鲁派，以及所有同一学派的人。斯多葛派也是如此，他们将他们如此热爱的论辩技艺——他们称之为辩证法——也归诸身体感官，并声称心灵从感官构想出\OriginalTerm{概念}{\GreekText{ἒννοιαι}}，用以通过定义来说明事物。由此发展出他们学问和教育的整个体系与关联。关于这点，我常感到惊奇：他们怎能说唯有智者才是美的？因为他们是用哪一种身体感官感知那种美，又是用哪双肉身的眼睛看见智慧的秀雅形态呢？然而，我们公义地将之置于一切人之上的那些人，已将心灵所构思的事物与感官所感知的事物区分开来，既不从感官拿走任何其力所能及的东西，也不将任何超出其能力的东西归给它。而我们的悟性之光，我们藉以学习一切的光明，他们断言就是那位创造万物的天主本身。\par

% structure: p0018 source=h2 target=subsection reason=relative-heading-rank; base=h2

\subsection{柏拉图主义者在道德哲学方面亦居首位}

哲学余下的部分便是伦理学，即希腊人所谓\OriginalTerm{伦理}{\GreekText{ἠθική}}，其中讨论的是关于至善的问题——那使我们不再寻求其他，就能获得真福的善，只要我们使自己的一切行为都指向它，并且不是为其他事物而寻求它，而是为它本身而寻求它。因此它被称为目的，因为我们为了它而愿意其他事物，但它本身只为本身而被寻求。所以，这真福善，据某些人说，是从身体来到人身上的；据另一些人说，是从心灵来的；又据另一些人说，是从二者同来的。因为他们看到人本身是由灵魂和身体组成的；因此他们相信，他们的幸福必定从这三者中的某一方，或从二者一同而来，这幸福在于某种终极的善，这善能使他们享福，他们可以把自己的一切行为都归向它，而不需要再有任何其他事物可以将这善本身再归向它。这就是为什么那些增添了第三种善的人——他们称之为外在善，如尊荣、名誉、财富等——并不把它们视为终极善的一部分，即不是为了它们本身而寻求的，而是为了其他事物才寻求的，他们断言，这类善对善人是好的，对恶人是坏的。因此，无论他们是从心灵还是从身体，或是从二者同来寻求人的善，他们仍然只以为这善必须从人那里寻求。但是，从身体寻求善的，是从人的较低部分寻求；从心灵寻求善的，是从较高部分寻求；从二者寻求善的，是从整个人寻求。因此，不论他们是从某一部分，还是从整个人寻求善，他们仍然只是从人那里寻求；这三种不同的见解，并没有只产生三个意见纷歧的哲学学派，而是产生了许多学派。因为不同的哲学家，无论是关于身体的善、心灵的善，还是二者相容的善，都持有不同的意见。所以，让所有这些人都让位给那些哲学家吧，他们不认为人是通过享受身体，或通过享受心灵，而是通过享受天主而成为有福的——然而，这种享受祂，不像心灵享受身体或享受它自身那样，也不像一个朋友享受另一个朋友那样，而是像眼睛享受光明一样，如果我们可以在这类事物之间作某种比拟的话。至于这比拟的性质是什么，若天主助我，我将在别处尽力说明。目前只需提到，柏拉图确定终极的善就是按德行生活，并断言只有认识和效法天主的人才能获得美德——而这认识和效法就是福乐的唯一原因。因此，他毫不犹豫地认为，哲学思考就是爱天主，而天主的本性是无形的。由此必然得出，智慧的学生，即哲学家，将在开始享受天主时成为有福的。因为虽然享有自己所爱的并不必然有福（许多人因爱了不该爱的事物而痛苦，享有它们时更为痛苦），然而，不享有自己所爱的，没有人是有福的。因为即使那些爱不该爱事物的人，也不认为自己在仅仅爱时就是有福的，而是在享有它们时才认为是有福的。那么，谁，除了最为痛苦的人，会否认享有自己所爱的、并且爱的是真实而至高的善的人，是有福的呢？按照柏拉图，真实而至高的善就是天主，因此他将热爱天主的人称为哲学家；因为哲学的目的在于获得幸福的生活，而热爱天主的人在享受天主中就是有福的。\par

% structure: p0020 source=h2 target=subsection reason=relative-heading-rank; base=h2

\subsection{论最接近基督信仰的哲学}

因此，无论哪些哲学家对至高天主持有如此看法：认为祂既是万物的创造者，是认识事物的光明，也是行事所应趋向的善；我们在祂内拥有自然的第一原理、教义的真理和生命的幸福——无论将这些哲学家称为柏拉图主义者更合适，还是给他们的学派另起他名；无论只有伊奥尼亚学派的杰出人物，如柏拉图本人及正确理解他的人持此看法，还是我们也将意大利学派包括在内，因为毕达哥拉斯和毕达哥拉斯学派及所有持类似观点的人；最后，无论我们是否也包括各民族中被视为智者与哲学家，并确曾见到并教导此理的人，无论他们是亚特兰蒂斯人、利比亚人、埃及人、印度人、波斯人、迦勒底人、西徐亚人、高卢人、西班牙人或其他民族的人——我们都认为他们优于其他所有哲学家，并承认他们最接近我们。\par

% structure: p0022 source=h2 target=subsection reason=relative-heading-rank; base=h2

\subsection{基督宗教的卓越超越哲学家的一切学问}

因为，尽管一个只受过教会文献教导的基督徒可能连柏拉图主义者的名字都不知道，甚至不知道曾经存在过两个操希腊语的哲学学派，即伊奥尼亚学派和意大利学派，但相对于人事，他并非如此闭塞，以至于不知道哲学家宣称研究、甚至拥有智慧。然而，对于按照这个世界的元素、而非按照创造世界本身的天主进行哲学思考的人，他是有所提防的；因为他受到宗徒训令的警告，并且忠实地听从了所说的话：\BibleQuote{你们要小心，免得有人以哲学和虚伪的欺诈，按照世俗的元素迷惑你们。}\BibleRef{哥 2:8} 然后，为了不让他认为所有哲学家都是如此行的，他听到同一位宗徒论到他们中的某些人说：\BibleQuote{因为认识天主已为他们显示出来，因为天主已给他们显示了。事实上，自从天主创世以来，祂那看不见的美善，即祂永远的大能和天主性，都可凭祂所造的万物，辨认洞察出来。}\BibleRef{罗 1:19} 并且，当他向雅典人讲话时，在说了一件关于天主的、很少人能够理解的大能之事后，他说：\BibleQuote{我们生活、行动、存在，都在祂内，正如你们有的诗人也说过的。}\BibleRef{宗 17:28} 他也清楚地知道，即使在他们的错误中，也要提防这些哲学家。因为，宗徒说了\Quote{天主已借着所造之物，将自己不可见的属性显示给他们，使他们能以理智看见}之后，接着又说，他们没有恰当地崇拜天主自己，因为他们将只应归于天主的崇敬，也给予了不该给予的其他事物——\BibleQuote{因为他们虽然认识了天主，却没有以天主的光荣来荣耀祂，也没有感谢祂，反而在思想上成了徒然，他们冥顽不灵的心陷入了黑暗。他们自负为智者，反成了愚蠢的，将不可腐朽的天主的光荣，变为可朽坏的人、飞禽、走兽和爬虫的偶像的样式；}\BibleRef{罗 1:21}——宗徒在这里要我们理解，他指的是那些以智慧之名自夸的罗马人、希腊人和埃及人；但关于这一点，我们以后要同他们辩论。然而，关于他们与我们一致之处，我们看重他们超过其他所有人——即关于那位唯一的天主，宇宙的创造者，祂不仅超越一切形体，是无形的，而且超越一切灵魂，是不朽的——祂是我们的本原、我们的光明、我们的善。而且，虽然那位基督徒不知道他们的著作，在争论中不使用他没有学过的词语——不称哲学中研究自然的那部分为\OriginalTerm{自然哲学}{natural}（拉丁术语）或\OriginalTerm{物理学}{physical}（希腊术语），不称处理如何发现真理问题的那部分为\OriginalTerm{理性}{rational}或\OriginalTerm{逻辑学}{logical}，不称论及道德、指示如何寻求善和躲避恶的那部分为\OriginalTerm{道德}{moral}或\OriginalTerm{伦理学}{ethical}——但他并非因此不知道，我们拥有那种按天主的肖像受造的本性，那种使我们认识天主和我们自己的教义，以及那种通过依附于祂而使我们蒙福的恩宠，都是来自那位唯一真实且至高至善的天主。因此，这就是我们看重这些人超过其他所有人的原因，因为其他哲学家在寻求事物的原因、努力发现正确的学习和生活方式时，耗尽了他们的心智和精力；而这些人因认识了天主，已经找到了宇宙由以构成的原因所在之处，发现了真理由以发现的光明，以及幸福由以畅饮的泉源。因此，所有对天主怀有这些思想的哲学家，无论是柏拉图主义者还是其他人，都与我们一致。但是我们认为，与柏拉图主义者辩护我们的主张更为合宜，因为他们的著作更为知名。因为希腊人——他们的语言在外邦人的语言中占据最高的地位——对这些著作大加赞扬；而拉丁人，被这些著作的卓越或名声所吸引，比研究其他著作更热心地研究它们，并通过将它们翻译成我们的语言，使之获得了更大的声誉和名望。\par

% structure: p0024 source=h2 target=subsection reason=relative-heading-rank; base=h2

\subsection{柏拉图如何得以如此接近基督信仰的知识。}

与我们一起分享基督恩宠的某些人，当他们听说并读到\Person{柏拉图}对天主持有某些观念，并在其中辨认出与我们的宗教真理有相当大的一致性时，感到惊奇。有些人由此得出结论，认为他前往\Place{埃及}时曾听到先知\Person{耶肋米亚}的讲论，或是在同一地区旅行时，读过先知的经卷；这一观点，我本人在一些著作中也曾表达过。然而，根据年代史所载的精确日期推算，\Person{柏拉图}大约出生于\Person{耶肋米亚}说预言的时代之后百年，而他活到八十一岁，故从他去世到\Place{埃及}王\Person{托勒密}请人从\Place{犹太}将希伯来先知书送往自己的地方，并委托七十位精通希腊语的希伯来人翻译并收藏，其间大约相隔七十年。因此，\Person{柏拉图}在\Place{埃及}之行中，既不可能见到早已过世的\Person{耶肋米亚}，也不可能阅读这些尚未译成他通晓的希腊语的经卷，除非——诚然如此——由于他求知欲极强，也曾像研读\Place{埃及}著作那样，通过一名通译来研究这些作品；当然，他并未将其译出（即便\Person{托勒密}得以翻译，也是以盛大恩惠为报偿才获得的便利，尽管威逼其王权似乎也是足够的动机），而是通过交谈，尽可能多地了解其内容。支持这一假设的，是《\BibleBook{创}{世纪}》开头几节：\BibleQuote{在起初天主创造了天地。大地还是混沌空虚，深渊上还是一团黑暗，天主的神在水面上运行。}\BibleRef{创 1:1} 因为在《\WorkTitle{蒂迈欧篇}》中，论及世界的形成时，他说天主首先将土与火结合；由此可见，他把火置于天中。这一观点与\BibleQuote{在起初天主创造了天地}这句话有某种相似之处。接着，\Person{柏拉图}谈到水与气这两种中介元素，借它们将土与火这两个极端彼此联合；从这一点来看，他大概是这样理解\BibleQuote{天主的神在水面上运行}这句话的；由于没有充分注意圣经对天主之神的称谓，他便可能以为那处经文说的是四种元素，因为气也被称为\Quote{\OriginalTerm{灵}{spirit}}。至于\Person{柏拉图}所说的哲学家是热爱天主的人，在这神圣经卷中再没有比这更显明的了。然而，在这方面最引人注目、也最让我几乎倾向于相信\Person{柏拉图}并非不知道这些经卷的，便是当天主的话藉天使传给圣\Person{梅瑟}时，针对他所问的回答；当时他问那位命令他去解救希伯来人民出离\Place{埃及}的天主叫什么名字，得到的回答是：\BibleQuote{我是自有者；你将要这样对以色列子民说：\InnerQuote{\Emph{自有者}打发我到你们这里来}}\BibleRef{出 3:14} 这仿佛是说，与那位真正\Emph{是}者（因他永恒不变）相比，那些受造而可变的事物则\Emph{不是}——这一真理正是\Person{柏拉图}热切持守并极力推崇的。我不知道在\Person{柏拉图}之前的人的著作中，哪里还能找到这一见解，除非是在那一卷书中，其中记载道：\BibleQuote{我是自有者；你将要这样对以色列子民说：\InnerQuote{\Emph{自有者}打发我到你们这里来}}。\par

% structure: p0026 source=h2 target=subsection reason=relative-heading-rank; base=h2

\subsection{柏拉图学派虽论及唯一真天主，却仍认为应向众多神祇举行敬礼。}

然而，我们无须确定他是从何源头学到这些事——无论来自他以前的古人著作，还是更可能出自宗徒的话：\BibleQuote{因为那关于天主而可认识的，已显示给他们了，天主已将它显示给他们了。因为自创世以来，他那看不见的美善，即他永远的德能和天主性，都可凭他所造的万物，辨认洞察出来。}\BibleRef{罗 1:20} 无论他从何处得到这一知识，我认为，我已充分说明，我选择\Person{柏拉图}派哲学家作为讨论的伙伴并非不恰当；因为我们刚才所讨论的问题正是自然神学的问题——即，为了死后的幸福，敬礼应向一位天主还是向多位神祇举行。我特别选择他们，是因为他们关于那创造天地的唯一天主的较为正确的思想，使他们在哲学家中享有盛名。这在后世的判断中使他们凌驾于所有其他人之上，以致于尽管\Person{柏拉图}的门徒\Person{亚里士多德}，一位才能卓越的人，口才虽不如\Person{柏拉图}，但在那方面却远胜许多人，创立了逍遥学派——如此称呼，因为他们在辩论时习惯于来回走动——并且，他甚至在他老师还在世时，就以其盛名吸引了许多门徒入其门下；\Person{柏拉图}去世后，他那个被称为\OriginalTerm{学园}{Academy}的学校由他妹妹的儿子\Person{斯彪西波}和他钟爱的门徒\Person{色诺克拉底}接续，他们及后学都因此学校的名称而被称为学园派；然而，那些选择追随\Person{柏拉图}的最杰出的后来哲学家们，却不愿被称为逍遥学派或学园派，而宁愿被称为\Person{柏拉图}学派。其中著名者有希腊人\Person{普罗提诺}、\Person{杨布里科斯}和\Person{波尔菲里}，以及\Place{阿非利加}人\Person{阿普列乌斯}，他精通希腊和拉丁两种语言。然而，所有这些人，以及同派的其他人士，连同\Person{柏拉图}本人，都认为应向许多神祇举行敬礼。\par

% structure: p0028 source=h2 target=subsection reason=relative-heading-rank; base=h2

\subsection{第十三章 论\Person{柏拉图}认为众神全善且为德行之友}

因此，虽然他们在许多其他重要方面与我们不同，但关于我刚刚指出的这一点重要差异，因为问题关键在此，我首先要问他们：他们认为应当向哪些神祇举行祭礼——是向善的，还是向恶的，抑或既向善的又向恶的？然而，我们持有\Person{柏拉图}的看法，他宣称所有神祇都是善的，没有一个神祇是恶的。由此可以推出，祭礼应当向善的举行，因为那样才是向神祇举行的；因为如果他们不是善的，他们也就不是神祇。现在，如果事实如此（关于神祇，我们还能相信别的什么呢？），这无疑就推翻了这样一种意见：即我们应当通过祭礼来讨好凶神，以免他们伤害我们，而应当呼求善神，以求他们助佑我们。因为根本就不存在凶神，正如他们所说，应当向善神献上这种礼仪应有的尊崇。那么，那些喜爱戏剧表演、甚至要求将之纳为神圣事物、并为其荣耀而演出的神祇，究竟又是什么品格呢？这些神祇的能力证明了他们的存在，但他们对这类事物的喜好却证明了他们是恶的。因为众所周知\Person{柏拉图}对戏剧表演的意见。他认为诗人自己，因为他们创作了如此不配神祇的威严和良善的诗歌，就应当被逐出城邦。那么，那些为了这些戏剧表演而与\Person{柏拉图}本人相争的神祇，又是什么品格呢？他不容神祇因虚假的罪行而受诽谤；众神祇却命令用这些罪行来庆祝他们自己的荣耀。\par

最终，当他们命令开始这些表演时，他们不仅要求低下的事，还行了残忍的事：从\Person{提图斯·拉提尼乌斯}那里夺走了他的儿子，并因为他不服从命令而降下疾病，直到他满足了他们的要求才消除疾病。然而，\Person{柏拉图}虽然认为他们邪恶，却不认为应当惧怕他们；他怀着最坚定的信念，毫不犹豫地要将诗人们一切亵圣的愚妄从一个秩序井然的国家中清除出去，而这些正是那些神祇所喜爱的，因为他们自身就是不洁的。但\Person{拉贝奥}却将这位\Person{柏拉图}（如我在第二卷中已经提到的）归入半神之列。\Person{拉贝奥}认为，恶神应当用血腥的祭品和与之相伴的斋戒来讨好，善神则应用表演和一切与欢悦相关的事物来讨好。那么，作为半神的\Person{柏拉图}，为何如此坚决地敢于剥夺那些乐趣——不是剥夺半神们的，而是剥夺众神的，更何况是那些善神的——因为他认为这些乐趣是可鄙的呢？而且，那些神祇本身无疑也驳斥了\Person{拉贝奥}的意见，因为在\Person{拉提尼乌斯}的事上，他们显示自己不仅是放荡嬉戏的，还是残忍可怕的。因此，让柏拉图学派的人向我们解释这些吧，因为他们遵照他们老师的意见，认为所有神祇都是善良、尊贵的，并与智慧者的德行为友，认为对任何神祇持有别种想法都是不虔的。他们说：\Quote{我们会解释的}。那么就让我们仔细听他们说。\par

% structure: p0031 source=h2 target=subsection reason=relative-heading-rank; base=h2

\subsection{第十四章 论那些认为理性灵魂有三类——属天神、属空中神怪与属地之人——的意见}

他们说，凡具有理性灵魂的动物分为三类：即神、人、神怪。众神占据最高的区域，人占据最低的，神怪占据中间的区域。因为众神的居所是天，人的是地，神怪的是空中。他们区域地位的不同，相应他们本性的尊贵也不同；因此众神优于人和神怪。人则被置于众神和神怪之下，无论就他们所居区域的秩序，还是就他们功绩的差异而言，都是如此。因此，处于中间位置的神怪，他们既低于众神，因为他们居住在较低的区域，又高于人，因为他们居住在较高的区域。因为他们与众神同有身体的不死性，却与人同有心灵的情感。他们说，正因如此，他们喜欢剧场的猥亵和诗人的虚构，也就不足为怪了，因为他们也受制于人的情感，而众神却远离这些，与之毫无关涉。由此我们得出结论：\Person{柏拉图}所剥夺的戏剧之乐，并非剥夺那些全是善良崇高的众神的，而是剥夺神怪的，他通过对诗人们的虚构加以谴责和禁止而这样做。\par

许多人都写过这些事：其中有\Place{马道拉}的柏拉图主义者\Person{阿普列乌斯}，他专门就这个主题写了一整部作品，题为\Emph{\WorkTitle{论苏格拉底的神}}。他在那里讨论并解释，侍奉\Person{苏格拉底}的那位究竟是哪一类神明，那是一种密友精灵，据说正是这位密友提醒他放弃任何不会带来益处的事。他十分明确地断言，并以大量篇幅证明，那不是一位神，而是一个\OriginalTerm{魔鬼}{demon}；他还极其勤勉地讨论了\Person{柏拉图}关于众神的崇高地位、人类的卑微地位，以及魔鬼的中间地位的看法。既然情况如此，\Person{柏拉图}怎么敢于把剧场的一切乐趣，即便不是从他所去除一切人间接触的众神那里，但肯定是从魔鬼那里夺走，藉着把诗人逐出国度呢？显然，他是想藉此劝诫人类的灵魂，虽然仍被困在这些必死的肢体中，要鄙视魔鬼那些可耻的命令，憎恶他们的不洁，反而选择德行的光辉。然而，如果\Person{柏拉图}在回应和禁止这些事上显出了德行，那么魔鬼去命令这些事肯定就是可耻的了。因此，要么\Person{阿普列乌斯}错了，\Person{苏格拉底}的密友精灵并不属于这类神祇，要么\Person{柏拉图}持有自相矛盾的意见，时而敬拜魔鬼，时而又把魔鬼所喜爱的事情从一个治理良好的国度中清除出去，要么\Person{苏格拉底}就不该因与那个魔鬼的友谊而受到祝贺，\Person{阿普列乌斯}竟对此感到如此羞涩，以至于将自己的书题为\Emph{\WorkTitle{论苏格拉底的神}}，而根据他讨论的主旨——他在其中如此勤勉且长篇大论地区分神与魔鬼——他本该题作\Emph{\WorkTitle{论苏格拉底的魔鬼}}，而非\Emph{\WorkTitle{论苏格拉底的神}}。但他宁可将这一点放在讨论本身之中，而不是放在书的标题里。因为，藉着那照亮人类社会的纯正教义，所有或几乎所有的人，都对魔鬼这个名字如此惊惧，以至于任何人在阅读\Person{阿普列乌斯}那篇阐述魔鬼尊严的论文之前，只要先读到\Emph{\WorkTitle{论苏格拉底的魔鬼}}这个标题，就肯定会以为作者不是个心智正常的人。然而，即使是\Person{阿普列乌斯}，在魔鬼身上又找到了什么可称赞的呢，除了身体的精微与力量，以及更高的居住场所？因为当他概括地谈论他们的举止时，他没有说什么好的，倒说了很多坏的。最终，没有人在读了那本书之后会感到惊讶，他们竟渴望甚至让舞台上的猥亵之事也成为神圣事物，或者，他们既然希望被视为神，就乐于享受诸神的罪行，或者，所有那些神圣的典礼——其猥亵惹人发笑，其可耻的残忍引起惊骇——都符合他们的情欲。\par

% structure: p0034 source=h2 target=subsection reason=relative-heading-rank; base=h2

\subsection{魔鬼并不因有气体般的身体或居所优越就比人更好。}

因此，凡真正虔诚、服从真\OriginalTerm{天主}{God}的人，不要以为魔鬼因为拥有更好的身体就比人更好。否则，他就得把许多野兽置于自己之上了，因为这些野兽在感官的敏锐、行动的轻松与迅速、力量以及长久不衰的体力方面都胜过我们。有哪一个人能在视力上比得上鹰或秃鹫呢？谁能在嗅觉的敏锐上比得上狗呢？谁能在奔跑的速度上比得上野兔、牡鹿和一切飞鸟呢？谁能在力量上比得上狮子或大象呢？谁能在寿命的长度上比得上那些据说会连同皮肤一起脱去衰老、重返青春的蛇呢？但是，我们既然因拥有理性和理解力而胜过这一切，也就更该藉着过上良善和有德的生活而胜过魔鬼。因为\OriginalTerm{天主眷顾}{divine providence}给了他们比我们更好的身体，正是为了让我们胜过他们的那个长处，能因此而受到褒扬，被认为远比身体更值得关心，并且我们要学会，与良善的生活相比——在这方面我们胜过他们——魔鬼身体上的卓越应当受到蔑视，我们知道，我们也必将拥有身体的不死——不是那种遭受永罚折磨的不死，而是那种伴随灵魂洁净而来的不死。\par

但是，至于地理位置高低，若因\OriginalTerm{魔鬼}{demon}住在空中而我们在地上，就认为因此他们应被置于我们之上，这全然可笑；因为按照这个逻辑，我们会把一切鸟都置于自己之上。然而鸟在飞累了或需补充食物时，会回到地上休息或觅食，据说魔鬼却不这么做。那么，他们是否因此要说鸟比我们优越，魔鬼又比鸟优越呢？如果这样想是疯狂的，那就没有理由认为魔鬼仅因居住在更高元素中就有权要求我们的宗教崇拜。然而实际上，空中的鸟不但未被置于我们这些住在地上的人之上，反而因我们内的理性灵魂的尊严而臣服于我们；同样，魔鬼尽管属气，也并不因空气高于土地而比我们这属地的好，恰恰相反，人应被置于魔鬼之上，因为他们的绝望与虔诚人的望德无法相提并论。甚至\Person{柏拉图}的那一法则——他据此排列四元素，在两极端元素，即最具活动性的火与不动摇的土之间，插入两个中间元素，气与水，气比水高多少，火比气高多少，水也比土高多少——这一法则，我说，已充分告诫我们，不要依照元素的等级来评判有生之物的功德。\Person{阿普列尤斯}自己也说，人同其他动物一样是陆生动物，但却远胜于水生动物，尽管柏拉图将水本身置于土之上。他藉此要我们明白，当论及动物的功德时，不可遵循相同的次序，尽管这看来在形体等级中是正确的；因为似乎更高阶的灵魂可能居于较低级的身体，较低阶的灵魂也可能居于较高级的身体。\par

% structure: p0037 source=h2 target=subsection reason=relative-heading-rank; base=h2

\subsection{第16章——柏拉图主义者\Person{阿普列尤斯}对魔鬼的品行与作为的看法}

同一位\Person{阿普列尤斯}，在论及魔鬼的品行时，说他们与人类一样受心灵纷扰；他们因伤害而被激怒，因侍奉和礼物而得安抚，以尊荣为乐，喜爱各种祭礼，若有疏忽便被触怒。此外，他还说，占卜者、预言者、先知们的占卜，梦的启示，都依赖于他们；魔术家的奇迹也出自他们。但在给他们下简短定义时，他说：\Quote{魔鬼具有\OriginalTerm{动物本性}{animal nature}，\OriginalTerm{灵魂受动}{passive in soul}，\OriginalTerm{心灵有理性}{rational in mind}，\OriginalTerm{身体属气}{aerial in body}，\OriginalTerm{时间上永恒}{eternal in time}。}\Quote{这五项中，前三项是他们与我们共有的，第四项是他们独有的，第五项是他们与\OriginalTerm{神}{gods}共有的。}但我看到，他们与\OriginalTerm{神}{gods}共有前两项中的两项，而这两项他们也与我们共有。因为他说神也是动物；且当他为各类存在分配各自元素时，将我们归于其他在地上生活知觉的陆生动物。因此，如果魔鬼在类别上是动物，这不仅与我们和人也与神和野兽共有；如果他们在心灵上有理性，这他们与神和人共有；如果他们在时间上永恒，这他们仅与神共有；如果他们在灵魂上受动，这他们仅与人共有；如果他们在身体上属气，这仅是他们独自的。所以，他们有动物本性，这不算什么，因为野兽亦然；在心灵上有理性，他们不高于我们，因为我们亦然；至于在时间上永恒，若他们并非有福，又有何益？因为暂时的幸福胜于永恒的悲惨。再者，就灵魂受动而言，他们如何在这方面优于我们？因为我们也如此，但若非我们悲惨，我们就不会如此。同样，就身体属气而言，这有何价值？因为无论何种灵魂都高于任何身体。因此，本应由灵魂发出的宗教崇拜，绝不归于那低于灵魂的事物。此外，如果他在那些所谓属于魔鬼的事物中，列举美德、智慧、幸福，并断言他们与神永远共有这些，他一定会赋予他们极大值得渴求和珍视的东西。即便如此，我们也不应因此像敬拜天主那样敬拜他们，反而应该敬拜那位我们深知他们从其获得这些的天主。然而，他们实在多么不配享有神圣的尊荣——那些属气的动物，仅有理性以便能感受痛苦，有感受性以便实际痛苦，永恒以便其痛苦永无终结！\par

% structure: p0039 source=h2 target=subsection reason=relative-heading-rank; base=h2

\subsection{第17章——人是否应该崇拜那些他们必须摆脱其邪恶的魔鬼？}

所以，姑且略去其他，只专注于他所说魔鬼与我们相同之处，我们且问：如果四大元素都充满其特有的生物，火与气充满不朽者，水与土充满有死者，那为什么魔鬼的灵魂却为情欲的狂风暴雨所扰动呢？——因为希腊字\OriginalTerm{情欲}{\GreekText{παθος}}的意思是纷扰，所以他选择称魔鬼为\Quote{灵魂有情欲的}，因为情欲一词源自\OriginalTerm{情欲}{\GreekText{πάθος}}，意指一种违反理性的心神激动。那么，这些为禽兽所无的东西，为什么会在魔鬼心中呢？因为倘若禽兽显出类似的事，那并不是纷扰，因为禽兽没有理性，谈不上违反。而对人而言，这些纷扰的原因在于愚昧或可怜，因为我们尚未享有那终将赐予我们的圆满智慧所带来的真福，届时我们将脱离现世的可朽生命。但他们却说，众神却无此纷扰，因为他们不仅永恒，而且享真福；因为他们虽也有同样的理性灵魂，却至洁无瑕，绝无污染。所以，若众神因享真福而非可怜之物，故无纷扰；禽兽因既不能享真福亦不能受苦，故无纷扰；那么魔鬼便如人一样，因其并非真福而是可怜之物，故亦受制于纷扰。因此，我们若出于任何宗教情感而听命于魔鬼，是何等的愚昧，甚至疯狂！因为真宗教的本分正是将我们从那使我们肖似魔鬼的堕落中解救出来！就连\Person{阿普莱乌斯}本人，虽然百般善待魔鬼，且认为他们配受神圣的尊崇，也迫不得已承认他们会动怒；而真宗教命令我们不要动怒，反要抵抗愤怒。魔鬼因礼物而被收买；而真宗教命令我们不可因所收的礼物而偏待任何人。魔鬼因尊荣而受奉承；而真宗教命令我们绝不可为此等事所动。魔鬼憎恨某些人，喜爱另一些人，并非出于明智而冷静的判断，而是出于他所谓的\Quote{灵魂有情欲}；而真宗教命令我们甚至要爱仇敌。最后，真宗教命令我们抛弃内心所有的不安、精神上的烦乱，以及灵魂中一切的风波与狂浪，而\Person{阿普莱乌斯}却断言魔鬼的灵魂中无时不在翻腾不已。那么，若非出于愚昧和可怜的错谬，你为何要去崇拜一个你生活上绝不欲效仿的对象呢？你既不愿效法他，又为何要给予他宗教上的敬礼呢？因为宗教的最高本分正是效法你所崇拜的那一位！\par

% structure: p0041 source=h2 target=subsection reason=relative-heading-rank; base=h2

\subsection{第十八章 何种宗教教人应用魔鬼的代求，以求获得善神的恩宠}

因此，\Person{阿普莱乌斯}及与他持相同想法的人，将魔鬼置于空中，即苍穹与大地之间，以使其能向众神传达人的祈祷，向人传达众神的答复，这种赋予魔鬼尊荣的做法，是徒劳无益的：因为据说\Person{柏拉图}主张，没有神与人直接交往。相信这些事的人认为，人与神、神与人直接交往不合时宜，而由魔鬼与双方交往，向神呈上人的请愿，向人转达神的允准，则是合宜的；以致一个贞洁、与邪术罪行无涉的人，必须借助那些喜爱此等罪行的魔鬼作为保护人，以使众神俯听，尽管他正因不喜爱这些罪行，本应更得神的乐意垂听。魔鬼喜爱舞台上那些贞洁所不喜爱的可憎之事。他们喜爱术士邪术中的\Quote{\Emph{千般害人之术}}，那是无辜者所不喜爱的。然而，贞洁与无辜者若想从众神那里求得什么，绝不能凭己功获得，除非其仇敌充当他们的中介。\Person{阿普莱乌斯}不必试图为诗人的虚构和舞台上的嬉笑辩护。如果人的廉耻竟能如此背弃自己，不仅喜爱可耻之事，甚至还以为这些能悦神，那么我们可以引用他们自己的最高权威和导师\Person{柏拉图}的话来反驳。\par

% structure: p0043 source=h2 target=subsection reason=relative-heading-rank; base=h2

\subsection{第十九章 论邪术的不虔敬，此术依赖恶神的援助}

再者，对于邪术，有些极端不幸且极端不虔敬的人竟以夸耀为乐，难道公众舆论本身不足以作证吗？如果邪术是应受崇拜的神灵所作，为什么法律要如此严厉地惩罚它们？难道可以说，是基督徒制定了惩罚邪术的那些法律吗？最杰出的诗人若非因为这些邪术无疑有害人类，又怎会说：\par

\Quote{我指天发誓，以你可爱的生命为誓，\newline{}我极不情愿提起这些武器，\newline{}去迎接未来的战斗，\newline{}拿起邪术的剑与盾。} \newline{}他在另一处论及邪术时又说：\newline{} \Quote{我曾见他把生长的庄稼移到别处，}\par

这指的是这样一个事实：按照这种流毒有害、应予诅咒的教导所传授的法术，一处田地的果实据说能被转移到另一处田地。\Person{西塞罗}不是告诉我们吗？在\WorkTitle{十二铜表法}——即罗马人最古老的法律——之中，有一条成文的法律，规定了对施行此事者应施加的惩罚。最后，\Person{阿普列乌斯}本人不是在基督徒法官面前被控行邪术吗？倘若他当时知道这些法术是神圣而虔敬的，符合神圣大能的作为，他就不仅应该承认，更应该公开宣扬它们，反而谴责那些禁止这些事并宣布它们应受惩罚的法律；其实这些事本该备受钦佩与尊崇。因为若这样做，他要么就能说服法官接受他自己的意见，要么，倘若法官偏袒不公的法律，即使他赞扬推崇这些事，仍判他死刑，那么魔鬼就会赐予他的灵魂应得的奖赏，因为他为了宣扬并显扬其神圣作为，不惧怕丧失自己的生命。正如我们的殉道者，当那宗教被指控为他们的罪行时——他们深知凭此宗教，他们得以安全且在永恒中至为荣耀——他们不选择否认它以逃避现世的惩罚，反而借着宣认、自承、传扬它，借以忠诚与刚毅为它忍耐一切，并借虔诚的平静为它而死，羞辱了那禁止该宗教的法律，并使其被撤回。然而，这位柏拉图派哲学家留存了一篇极为详尽且雄辩的演说，他在其中为自己辩解，反驳施行这些法术的指控，声称自己与这些法术全然无关，并只希望借着否认那些不可能清白施行的事，来表明自己的无辜。但那些术士——他认为他们理应受罚——所行的一切奇迹，都是按魔鬼的教导并借其力量施行的。那么，他为何认为那些魔鬼应受尊敬呢？因为他断言，魔鬼对于向神明呈上我们的祈祷是必要的，而如果我们希望我们的祈祷达于真正的天主，则必须远避魔鬼的作为。再者，我请问，按他所想，魔鬼向善神呈上的是人的哪种祈祷？若是邪术的祈祷，善神必不接纳；若是合法的祈祷，善神也不会通过这类存在来接受。然而，如果一个悔罪的罪人倾出祈祷，尤其当他犯了某种巫术之罪时，他是否借着那些魔鬼的转求而获赦免呢？正是因他们的煽动和协助，他才陷入他所哀哭的罪中。或者，魔鬼们自己为了为悔罪者赢得宽恕，首先因他们欺骗了人而自己悔过？关于魔鬼，从未有人如此说过；因为若果如此，他们就不会胆敢为自己寻求神圣的尊崇了。因为一个渴望借悔改获得赦罪之恩宠的人，怎会如此呢？——这种可憎的骄傲不可能与堪得赦免的谦卑共存。\par

% structure: p0047 source=h2 target=subsection reason=relative-heading-rank; base=h2

\subsection{第二十章 我们是否应相信善神更愿与魔鬼交往，而非与人交往}

然而，有什么紧急且最为迫切的原因，迫使魔鬼在神明与人之间作中介，好让他们献上人的祈祷，并带回神明的答复呢？若真如此，试问，那原因是什么，那极大的必要又是什么？因为他们说，没有神明与人交往。\newline{}何等可赞美的天主之圣洁！祂不与哀求的人交往，却与傲慢的魔鬼交往！\newline{}祂不与悔罪的人交往，却与欺骗人的魔鬼交往！\newline{}祂不与投奔神性的人交往，却与假装神性的魔鬼交往！\newline{}祂不与寻求宽恕的人交往，却与劝人作恶的魔鬼交往！\newline{}祂不与那些借哲学著作将诗人逐出井然有序之国的人交往，却与那些请求国家君主和司祭演出诗人嘲弄之戏剧的魔鬼交往！\newline{}祂不与禁止将罪行归于神明的人交往，却与乐于虚构地表现神明罪行的魔鬼交往！\newline{}祂不与以公正法律惩处术士罪行的人交往，却与教导并施行邪术的魔鬼交往！\newline{}祂不与远避模仿魔鬼的人交往，却与埋伏着要欺骗人的魔鬼交往！\par

% structure: p0049 source=h2 target=subsection reason=relative-heading-rank; base=h2

\subsection{第二十一章 神明是否以魔鬼为使者与通事，以及他们是否自愿受其欺骗，抑或浑然不知}

但毫无疑问，这里正显现出这种荒谬之极大必要性，实在有辱诸神：据称关切人类事务的\OriginalTerm{以太之神}{ethereal gods}，若没有\OriginalTerm{空中魔鬼}{aerial demons}前来通风报信，竟不会知道地上的人们在做什么，因为以太高悬，远离大地，而空气则与以太及大地都相毗邻。何等妙绝的智慧！这些人对诸神的想法，还能是什么别的呢？他们说诸神都是至善的，关切人类事务，以免显得不配受人崇拜，然而另一方面，由于\OriginalTerm{元素}{elements}的间隔，他们却对地上的事一无所知。正是由于这个缘故，他们设想魔鬼是必需的中介，诸神借此可以了解人类事务，并在必要时借助他们救助人类；而魔鬼本人正因为这种职分，被认为配受崇拜。倘若如此，那么魔鬼因身体的邻近而比人更容易为这些善神所知，胜过人因心灵的美善而被知晓。可悲的必然性，或者莫如说可憎且空洞的谬误，我且不将虚空归咎于神性吧！因为如果诸神能不受身体障碍地运用心神看见我们的心灵，他们就不需要魔鬼作中介，把我们的心意传给他们；但是如果以太之神凭借身体觉察到心灵的有形标志，如面容、语言、动作，并由此了解魔鬼告诉他们的事情，那么他们也可能被魔鬼的谎言所欺骗。再者，如果诸神的神性不会为魔鬼所欺骗，那么他们也就不可能不知道我们的行为。但我愿这些人告诉我：魔鬼是否已将诗人虚构的、关于诸神之罪恶的故事触怒了\Person{柏拉图}一事告知诸神，而隐瞒了他们自己从这些故事中得到的乐趣？或者他们两件事都隐瞒了，宁愿诸神对整个事情一无所知？或者两件都告知了，即\Person{柏拉图}对诸神的虔敬审慎，以及他们自己那有损诸神的贪欲？抑或他们隐瞒了\Person{柏拉图}的主张——按\Person{柏拉图}之见，他不愿诸神因诗人不敬的放纵、用虚假指控而蒙受丑化——而他们却不耻且无畏地张扬自己的邪恶，正是这邪恶使他们喜爱戏剧演出，而演出中歌颂的正是诸神的丑行。任凭他们在这四种可能性中选择哪一种吧，让他们想想，无论选哪一种，都会迫使他们把诸神想得何等邪恶。因为他们若选第一种，就必须承认，善神不可能与善的\Person{柏拉图}同住，尽管\Person{柏拉图}力图禁绝有害于诸神的事，而他们却与欢乐于伤害诸神的恶鬼同住；其理由是，他们以为善神只能通过恶鬼的中介才能认识远方的善人，而恶鬼他们却因临近而能认识。他们若选第二种，说魔鬼将这两件事都隐瞒了，以致诸神对于\Person{柏拉图}最虔敬的法则和魔鬼亵渎的快乐都全然无知，那么，在这种情况下，诸神通过中介魔鬼，能对人间的任何事有裨益地知道些什么呢？他们竟然不知道那些因善人虔诚而确立的、为维护善神尊严以对抗恶鬼贪欲的事。但他们若选第三种，回答说这些中介魔鬼不但传达了禁止加害诸神的\Person{柏拉图}意见，也传达了他们自己在加害诸神上的喜好，我倒要问，这样的传达岂非更是一种侮辱？现在，诸神既听到并知道这两件事，不但容许那些心怀恶意、渴望并做出违背诸神尊严与\Person{柏拉图}宗教信仰之事的恶鬼亲近，而且通过身边这些邪恶的魔鬼，将好东西送给远方的善良\Person{柏拉图}；因为他们身处元素连锁序列中的这样一个位置，只能接触到控告他们的人，却不能接触到为他们辩护的人——他们知道双方的真相，却不能改变空气与大地的比重。还剩下第四种假设；但这比其余几种更糟糕。因为谁能容忍这样的说法：魔鬼将使诸神蒙羞的、诗人们的诽谤虚构，以及剧场里丢脸的戏谑，连同他们自己对这一切极为炽烈的贪欲和最为甜美的乐趣，都让诸神知晓了，却对他们隐瞒了\Person{柏拉图}以哲学家的严肃所给出的意见——所有这些都应从一个治理良好的\OriginalTerm{共和国}{republic}中清除出去；结果，善良的诸神如今竟被迫通过这样的信使，去知道最邪恶之存在（也就是信使本人）的恶行，却不能获知哲学家们的善举，尽管前者是在损害诸神自身，而后者却是在荣耀诸神？\par

% structure: p0051 source=h2 target=subsection reason=relative-heading-rank; base=h2

\subsection{尽管\Person{阿普列乌斯}的意见如此，我们仍应拒绝魔鬼崇拜。}

因此，这四种选择中，一个也不可采纳；因为我们不敢对神明推想如此不体面的事，而采纳其中任何一种都会导致我们如此设想。因此，我们绝不可相信\Person{阿普列乌斯}和相同学派的其他哲学家的观点，即魔鬼在神明与人之间充当信使和传译者，将我们的祈求从我们这里带给神明，又将神明的帮助带回给我们。相反，我们必须相信他们是一些极想害人的精灵，完全与正义无缘，骄傲自大，嫉妒得面色苍白，诡计多端；他们确实居住在这空气中，如同在监牢里，这符合他们自己的品性，因为他们从高天之上被抛下，被判罚住在这元素中，作为其不可挽回的叛逆的公正报应。然而，尽管空气位于大地和水域之上，他们并不因此在功绩上比人优越，反倒人——虽然就肉身而言并不超越他们，却因心灵的虔敬而远远胜过他们，因为他们选择了真天主作为自己的帮助者。然而，他们对许多显然不配分享真宗教的人实施暴虐统治，如同对待被征服的俘虏——他们用奇妙而虚假的征兆，无论是奇事还是预言，说服这些人中的大多数，使他们相信自己是神圣的。不过，对于那些更加留意和勤谨地思考他们恶行的人，他们未能说服他们相信他们是神，于是就假扮自己是神明与人之间的信使。事实上，有些人甚至认为这后一种尊荣也不应归于他们，他们不相信他们是神，因为他们看出他们是邪恶的，而在他们看来，神明全都是善的。然而，他们不敢说魔鬼完全不配任何神圣尊荣，因为害怕触怒民众，民众出于根深蒂固的迷信，通过举行许多仪式和建造许多庙宇来侍奉魔鬼。\par

% structure: p0053 source=h2 target=subsection reason=relative-heading-rank; base=h2

\subsection{第 23 章——\Person{赫尔墨斯·特里斯墨吉斯忒斯}对偶像崇拜的看法，以及他从何得知埃及迷信将被废除。}

埃及人\Person{赫尔墨斯}，人称\Person{特里斯墨吉斯忒斯}，对这些魔鬼持有不同的看法。\Person{阿普列乌斯}虽然否认他们是神，但他说他们居于神明与人类之间的中位，似乎是人在神与人之间所需的中介时，并未区分对他们应有的敬拜和对上界的神明应有的宗教崇敬。然而，这位埃及人说，有些神明是由至高的天主造的，有些则是由人造的。任何听到我这样说的人，无疑会以为这指的是偶像，因为偶像是人手所造；但他断言，可见可触摸的偶像只不过是神明们的身体，并且有某些灵居于其中，这些灵是应邀来居住在里面的，它们有能力加害于人，或满足那些向它们奉献神圣尊崇和侍奉的人们的愿望。因此，他宣称，借由某种技艺将这些不可见的灵与可见的、有形的物结合，使之仿佛成为有生命的身体，专献给且交付给那些居住其中的灵——这就是造神，并补充说人类已经领受了这种伟大而奇妙的权能。我将引用这位埃及人的话，如被翻译成我们的语言所述：\Quote{既然我们着手论述人与神之间的关系和交往，\Person{埃斯库拉庇乌斯}啊，你要知道人的能力和力量。正如主和父，或者说那至高的，亦即\OriginalTerm{天主}{Deus}，是天界神明的创造者；同样，人是庙宇中神明的创造者，这些神明乐意居住在人附近。}稍后他又说：\Quote{因此，人类永远不忘自己的本性和起源，坚持不懈地模仿神性；正如主和父造了永生的神明，使他们相似自己，人类也照自己面貌的肖像塑造了自己的神明。}当这位\Person{埃斯库拉庇乌斯}——他特别对他说话——回答他说：\Quote{你是指雕像吗，\Person{特里斯墨吉斯忒斯}？}---\Quote{是的，雕像，}他回答说，\Quote{尽管你们不信，\Person{埃斯库拉庇乌斯}啊---那些雕像，是有生命、充满感觉和灵性的，能行如此伟大奇妙的事；那些雕像能预知未来，并通过抽签、预言、梦境和许多其他方式预告未来；它们降疾病于人，又再治疗好他们，照着各人的功过赐予他们欢乐或忧伤。难道你不知道，\Person{埃斯库拉庇乌斯}啊，埃及是天界的影像，或者更真确地说，是那里一切被命定和完成的事物的转移和下降，它实际上，如果可以这样说，就是全世界的圣殿？然而，正如明智之人应该预先知道一切，你不该不知道这一点：有一个时候将要来到，那时将显明，埃及人以虔诚的心和最严谨的勤勉侍奉神明，全都是徒然的，他们所有神圣的敬拜都将化为乌有，并被发现是枉然的。}\par

\Person{赫尔墨斯}接着长篇大论地阐述了这段话的内容，他似乎预言了现今的时代，在这时代，基督宗教以其超凡的真理和圣洁，正以猛烈的力量和自由推翻一切虚妄的谎言，好使真救主的恩宠能将人从那由人所造的神明中解救出来，并使他们臣服于那造人的天主。但当\Person{赫尔墨斯}预言这些事时，他却以这些魔鬼嘲弄之事的朋友身份发言，并未明示\Person{基督}的名字。相反，他悲叹这一切的废除，仿佛它已发生——正是藉着恪守这些事，\Place{埃及}曾维持着天堂的模样——他以一种悲哀的预言为基督教作了见证。关于这样的人，宗徒曾说：\BibleQuote{他们虽然认识了天主，却没有以他为天主而予以光荣或感谢，反而在思想上成了空洞，他们冥顽不灵的心陷入了黑暗；他们自负为智者，反而成了愚妄，将不可朽坏的天主的光荣，改归于可朽坏的人的形象}，\BibleRef{罗 1:21}等等，因为整段经文太长，不便全部引用。因为\Person{赫尔墨斯}就那唯一真天主、世界的创造者，说了许多合乎真理的陈述。我不知道他怎会被那\BibleQuote{冥顽不灵的心}弄得如此糊涂，竟至失足表达出这样的愿望：但愿人们永远臣服于那些他承认是人所造的神明，并哀叹它们将来的清除；仿佛再也没有什么比人类被自己双手所造之物奴役更可悲的了，因为人若敬拜自己双手的制品，那么人更易失去人性，而非他双手的制品能因他的敬拜而成为神。因为更容易发生的是，人虽拥有尊荣的地位，却因缺乏理智而沦为与走兽相似，而非人的制品变得优于天主的化工——那按祂自己肖像所造的，即人本身。因此，人若选择了自己所造之物而舍弃造他的天主，那么他离弃造他的主，正是理所当然的。\par

这些空虚、欺骗、有害、亵圣的事，让\Place{埃及}人\Person{赫尔墨斯}感到悲伤，因为他知道它们被清除的时刻即将来临。但他的悲伤表露得如此厚颜无耻，正如他的知识获得得如此轻率无知；因为并非圣神将这些事启示给他，如同祂启示给圣先知们一样——那些圣先知预见这些事，便欢欣地说：\BibleQuote{如果人能制造神，看，它们并不是神}，\BibleRef{耶 16:20}；在另一处又说：\BibleQuote{到那一天——上主说——我要将偶像的名号从大地上除去，不再被人记起}，\BibleRef{匝 13:2}。圣先知\Person{依撒意亚}更明确地论\Place{埃及}预言此事说：\BibleQuote{看，上主乘着轻快的云彩进入埃及，埃及的偶像在他面前战栗，埃及人的心灵在身内消溶}，\BibleRef{依 19:1}，以及其它类似的话。与先知同列的，还有那些因知道自己所盼望的已经实现而欢欣的人——比如\Person{西默盎}、\Person{亚纳}，他们立即认出了诞生的\Person{耶稣}；或\Person{依撒伯尔}，她在圣神内认出了成孕的祂；或\Person{伯多禄}，他因圣父的启示说：\BibleQuote{你是基督，永生天主之子}，\BibleRef{玛 16:16}。但对这位\Place{埃及}人，那些魔鬼指出了他们自己毁灭的时刻；当主在肉身内存世时，他们也战栗地说：\BibleQuote{时期还没有到，你就来这里毁灭我们吗？}，\BibleRef{玛 8:29}。所谓期前毁灭，或是指他们预料将临、却未想到会如此突然发生的毁灭，或仅指他们因被揭露而遭轻视的那种毁灭。这确实算是期前的毁灭，即在审判的时刻之前——到那时，他们将与所有参与其邪恶的人一同受永罚，正如真宗教所宣布的；这宗教既不会错误，也不会引人错误；因为它不像\Person{赫尔墨斯}那样，他随各种学说的风飘来飘去，将真事与假事混杂，悲叹一种宗教即将灭亡，而后来却承认那宗教是谬误。\par

% structure: p0057 source=h2 target=subsection reason=relative-heading-rank; base=h2

\subsection{第二十四章 赫尔墨斯如何公开承认其祖先的错误，却仍为之悲痛}

在很长的间歇之后，\Person{赫耳墨斯}再次回到人造之神这个主题，如下说道：\Quote{但关于这个主题已经谈够了。让我们回到人和理性，那是神圣的恩赐，人因此被称为理性动物。因为关于人所说的一切，虽然奇妙，却不及关于理性所说的奇妙。因为人发现神圣的本性，并制造它，超越了所有奇妙之事的奇妙。因此，\InnerQuote{由于我们的祖先在对众神的认识上大错特错了，由于不信和疏于对他们崇拜和事奉的注意，他们发明了这种造神术}；这种术一旦发明，他们便将从宇宙本性借来的适宜能力与之结合，因为不能制造灵魂，便招来魔鬼或天使的灵魂，将其与这些圣像和神圣奥秘结合，以便通过这些灵魂，偶像得以有能力对人行善或作恶。}我不知道是否就连魔鬼本身，即便通过诅咒，也能像他在这话里所承认的那样承认：\Quote{因为我们的祖先在对众神的认识上大错特错了，由于不信和疏于对他们崇拜和事奉的注意，他们发明了造神术。}他说这种导致他们发现造神术的错误程度是适度的吗，还是他仅满足于说\Quote{他们错了？}不；他必须加上\Quote{大错特错，}并说\Quote{\Emph{他们大错特错了}}。因此，正是他们祖先的这种大错和不信，他们不留意对众神的崇拜和事奉，才是造神术的起源。然而这位智者却在为这术在将来某个时候的毁灭而悲伤，仿佛它是一种神圣的宗教。他岂不是确实被神圣的影响所迫，一方面揭示他祖先过去的错误，另一方面又被魔鬼的影响所迫，哀悼魔鬼未来的惩罚吗？因为如果他们的祖先因在对众神的认识上大错特错，由于不信和心灵上疏于对他们的崇拜和事奉，而发明了造神术，那么，当真理纠正错误，信德反驳不信，皈依矫正疏离时，所有这令人憎恶的、与神圣宗教相悖的术所成就的，被那宗教除去，又有什么奇怪呢？\par

因为如果他只是说了，没有提及原因，他的祖先发现了造神术，那么，如果我们对正确和虔敬之事有所顾及，我们有责任思考并看出，如果他们未曾偏离真理，如果他们相信配得上天主的事，如果他们曾留意天主的崇拜和事奉，他们就永远不可能达到此种术艺。然而，如果只有我们说，这种术艺的原因在于人的大错和不信，以及心灵偏离并不信从神圣宗教的疏离，那么抗拒真理之人的厚颜在某种程度上还可以容忍；但是，当那个在所有事情中最赞赏人获得行使的这种能力，并为将有那个时候，那时所有人造的众神的虚构甚至将被法律命令除去而悲伤的人——当这个人尽管如此，仍然承认并解释了导致发现这种术艺的原因，说他们的祖先因大错和不信，因不留意众神的崇拜和事奉，而发明了这种造神术——我们该说什么，或者更确切地说，我们该做什么，除了尽力感谢主我们的天主，因为祂以与导致这些事建立的原因相反的原因将它们除去了？因为错误的盛行所建立的，真理之道将它除去；不信所建立的，信德将它除去；对天主的崇拜和事奉的疏离所建立的，皈依唯一真实而圣洁的天主将它除去。这情形不但在埃及——\Person{赫耳墨斯}中的魔鬼之灵只为这个国家悲叹——而且在整个大地，全地都向主歌唱新歌，正如真正神圣的、真正预言的圣经所预言的，其中记载说：\BibleQuote{向主歌唱新歌，全地都要向主歌唱。}因为这篇圣咏的标题是，\Quote{圣殿在被掳后建成之时。}因为一座圣殿正在全地为主建造，也就是天主之城，即圣教会，在此之后，即在魔鬼曾掳去那些人之后，他们藉着对天主的信德，成为这圣殿中的活石。因为人虽造了神，但并不因此，制造者不被它们所掳，当他藉着崇拜它们而被引入与它们交往——不是与愚钝偶像的交往，而是与狡猾魔鬼的交往；因为偶像算什么，不过正如同一圣经中所描述的，\BibleQuote{它们有眼，却看不见，}并且，虽经巧艺塑造，却仍没有生命和感觉？但是，不洁之神，藉着那邪恶的术与这些同样的偶像结合，可怜地掳去了崇拜者的灵魂，将他们拖入与自己的交往中。因此，宗徒说：\BibleQuote{我们知道偶像算不得什么，但外邦人所祭献的，是祭献给魔鬼，而不是祭献给天主；我不愿意你们与魔鬼交往。}\BibleRef{格前 10:19}所以，在这次被掳之后，其中人受到邪恶魔鬼的囚禁，天主的圣殿正在整个大地上建造；因此那篇圣咏的标题如此说，其中写道：\BibleQuote{向主歌唱新歌；全地都要向主歌唱。向主歌唱，赞美祂的名；日复一日，宣扬祂的救恩。在列邦中述说祂的光荣，在万民间述说祂的奇妙作为。因为主伟大，应受赞美：祂超越所有神祇，可敬可畏。因为列邦的所有神祇都是魔鬼：但主创造了诸天。}\par

因此，那位因偶像崇拜将被废除、魔鬼对崇拜它们之人的统治将被终结的时代来临而悲伤的人，在魔鬼的影响下，希望那种俘虏状态永远持续，而\WorkTitle{诗篇}则庆祝这种状态的终结，即在全地上建造主的殿宇。\Person{赫耳墨斯}忧心地预言了这些事，先知则欢欣地预言了；因为那位藉古代先知歌唱这些事的圣神是得胜的，甚至\Person{赫耳墨斯}本人也被神奇地迫使承认：那些他但愿不要被废除、因预见其废除而悲伤的事，并非由明智、忠信和虔敬的人设立的，而是由犯错、不信、厌恶崇拜和事奉诸神的人设立的。虽然他称它们为神，然而，当他说它们是由我们绝不应成为的那种人所造时，他便表明——无论他愿意与否——它们不应被那些与这些塑像者不相似的人，即明智、忠信和虔敬的人所崇拜，同时也显明那些制造它们的人竟陷于崇拜那些并非神的为神。因为先知的话是真实的：\BibleQuote{人岂能为自己制造神？其实那不是神。}\BibleRef{耶 16:20}因此，\Person{赫耳墨斯}称这些被那种崇拜者承认、由那种人制造的诸神为\Quote{人做之神}，也就是说，它们乃魔鬼，藉着某种不知为何的技艺，被自己私欲的锁链束缚在偶像上。然而，他并不同意\Person{柏拉图}学派的\Person{阿普列尤斯}的那个观点——我们已经指出其不合宜与荒谬之处——即认为它们是天主所造的诸神与同一天主所造的人之间的解释者和转求者，将人的祈祷呈于天主，并将天主应答这些祈祷而赐予的恩惠带给人。因为，若是相信人之所造的诸神，比人自己——他们乃同一天主所造——在天主所造的诸神面前更有影响力，那真是愚不可及。也要考虑，魔鬼通过不虔的技艺被人束缚于偶像，从而被造成神，但这只对那种人而言是神，并非对所有人。那么，除了那错误、不信、厌恶真天主的人，无人会制造的神，算是什么神呢？再者，若那些在庙宇中受崇拜的魔鬼，是借着某种怪异的技艺被引入偶像——即它们可见的像——之中的，而引进它们的人正是那些偏离、厌恶崇拜和事奉诸神，却以此技艺造神的人。若说，这些魔鬼既不是人与诸神之间的中保或解释者，一方面由于它们自身最邪恶卑劣的品行，另一方面也因为人虽错误、不信、厌恶崇拜和事奉诸神，但无疑仍比他们自己所召来的魔鬼更好。那么，剩下的便是肯定：它们所拥有的力量是作为魔鬼的力量，通过赐予虚假的恩惠而加害——欺骗使伤害更深——或是公然不加掩饰地向人行恶。然而，除非它们得到深奥隐秘的天主眷顾的允许，否则它们无法做出此类恶事，而且只能限于被允许的程度。而当它们被允许时，并非因为它们处于人与诸神之间的中间位置，并借着与诸神的友谊而对人拥有大能；因为这些魔鬼绝不可能与居于神圣天乡的善神为友——我们指的乃是圣天使与具有理性的受造物，无论是座天使、宰制天使、率领天使还是掌权天使——它们在性情与品格上与后者的差距，犹如邪恶远离美德，凶恶远离善良。\par

% structure: p0061 source=h2 target=subsection reason=relative-heading-rank; base=h2

\subsection{第25章——论圣天使与人可能共有的事}

因此，我们绝不能试图通过臆想的魔鬼的中保作用，来获利于诸神的善意或恩惠，或者更准确地说，获利于善天使的善意或恩惠，而是要通过在拥有善意上相似于他们，借此我们与他们同在，与他们共处，并与他们一起崇拜同一位天主，尽管我们不能用肉身的眼睛看见他们。但我们与他们之间的距离，不在于空间，而在于生活的功绩，这是由于我们的意志与他们可悲地不相像，以及我们品格的软弱所致；因为我们生活在地上，处于肉身生命的状态，这一事实本身并不阻止我们与他们相通。只有当我们在心灵的不洁中思念地上的事时，才会受到阻碍。但在现今，当我们正在被治愈、为使我们最终能如他们一样时，我们借着信德得以亲近他们，只要借着他们的帮助，我们相信那位是他们的真福的，也是我们的真福。\par

% structure: p0063 source=h2 target=subsection reason=relative-heading-rank; base=h2

\subsection{第26章——异教徒的一切宗教都涉及死人}

有件事确实值得注意：这位埃及人在表达他的悲伤，说将来\Place{埃及}要失去那些东西——他承认那些是由错误、不信、厌恶事奉神圣宗教的人所发明的——时，在其他话中说道：\Quote{那时，那最神圣的神龛和殿宇之地，将布满坟墓和死人，}好像如果这些东西不被拿走，人就不会死似的！好像尸体能被埋葬在不是土地的地方似的！好像随着时间推移，坟墓的数量不必与死人数目成正比增加似的！但是，那些心地乖谬、反对我们的人，以为他所悲伤的是我们殉道者的纪念物将要取代他们的殿宇和神龛，于是，他们便有理由认为，异教徒在殿宇中崇拜的是神明，而我们在坟墓中崇拜的却是死人。这些不虔敬的人如此盲目，就像被山绊倒一样，却看不见摆在眼前的事实——在所有异教徒的文献中，几乎找不到任何一位神明，不是曾经为人，死后被奉以神明的尊荣。我不想赘述\Person{瓦罗}说过的话，他说所有死人都被他们视为神——\OriginalTerm{曼涅斯}{Manes}——并且他通过那些几乎为所有死者举行的神圣仪式来证明这一点，其中他提到了葬礼竞技，认为这是神性的最高证明，因为竞技通常只为神明举行。我们正在讨论的\Person{赫尔墨斯}本人，在同一本书中，仿佛预言未来，悲伤地说：\Quote{那时，那最神圣的神龛和殿宇之地，将布满坟墓和死人，}他证明埃及的神明都是死人。因为他先说，他们的祖先在对神明的认识上大错特错，不信也不注意事奉敬拜神明，却发明了制造神明的技艺；这技艺发明后，他们将宇宙本性中固有的某种相宜的能力与之结合，又将这能力与这技艺混合，于是召来魔鬼或天使的灵魂（因为他们不能制造灵魂），并使这些灵魂进入或与圣像和神圣奥迹结合，好通过这些灵魂使神像有能力对人行善或作恶；——说完这些，他接下来举例证明说：\Quote{你的祖父，啊，\Person{埃斯库拉庇乌斯}，医学的首先发现者，在\Place{利比亚}的一座山上，靠近鳄鱼之岸，有人为他修建了一座殿宇，殿中安放着他那属土的人，即他的身体——因为他那更优秀的部分，或者更确切地说，整个的他，如果整个的人都存在于理性的生命中，已经回到天上去了——如今仍以他的神性，赐给病弱的人一切帮助，就如从前他惯于借医术所赐的一样。}因此他说，一个死人在他葬身之处被当作神明来崇拜。他以谎话欺骗人，因为那人\Quote{回到天上去了。}然后他又加上：\Quote{我的祖父\Person{赫尔墨斯}，我因他而得名，住在以他的名字命名的地区，他不是帮助并保护从各处来到他面前的凡人吗？}因为这位老\Person{赫尔墨斯}，也就是\Person{墨丘利}，据他说是他的祖父，据说葬在\Place{赫耳墨波利斯}，即以他名字命名的城市；所以这里有两个神明，他断言他们都曾是人，即\Person{埃斯库拉庇乌斯}和\Person{墨丘利}。关于\Person{埃斯库拉庇乌斯}，希腊人和拉丁人的看法相同；但关于\Person{墨丘利}，许多人并不认为他从前曾是有死之人，尽管\Person{赫尔墨斯}作证说\Person{墨丘利}是他的祖父。但这两个同一名字的是两个不同的人吗？我不愿多争辩他们是同一人还是不同的人。只要知道\Person{赫尔墨斯}所说的这个\Person{墨丘利}，正如\Person{埃斯库拉庇乌斯}一样，根据同一位\Person{特里斯墨吉斯提斯}——在他的同胞中备受尊崇，且是\Person{墨丘利}本人的孙子——的证言，是一位曾经是人的神明，这就足够了。\par

\Person{赫尔墨斯}接着说：\Quote{但我们知道\Person{伊西斯}，\Person{奥西里斯}的妻子，当她喜悦时赐下多少好处，发怒时又能造成多大的阻挠吗？}然后，为了说明通过这种技艺产生了由人造的神明，他接着说：\Quote{因为属地的和世俗的神明容易发怒，它们是由人用两种本性造成并组合而成的；}由此我们明白，他相信魔鬼原是死人的灵魂，正如他所说，它们借着某样由那些大错特错、不信、无宗教的人所发明的技艺，被引入神像，因为制造这些神明的人不能制造灵魂。所以，当他说\Quote{两种本性}时，指的是灵魂和身体——魔鬼是灵魂，神像是身体。那么，那悲伤的抱怨，说\Place{埃及}，最神圣的神龛和殿宇之地，将要布满坟墓和死人，又当如何呢？确实，那欺诈的灵，\Person{赫尔墨斯}就是受它的默示说了这些话，也不得不通过他承认，那块土地早已充满了坟墓和死人，就是他们所当作神明来崇拜的。但那正是魔鬼的悲伤借他的口表达出来，它们因即将在殉道者的坟墓那里承受的惩罚而哀伤。因为在许多这样的地方，它们受折磨，被迫承认，并被逐出它们所附的人体。\par

% structure: p0066 source=h2 target=subsection reason=relative-heading-rank; base=h2

\subsection{二十七、论基督徒向殉道者表示的尊荣的性质}

然而，我们并不为这些殉道者建造庙宇，也不为他们设立司铎、礼仪和\OriginalTerm{祭献}{sacrifices}；因为他们不是我们的神，而是他们的天主就是我们的天主。诚然，我们尊敬他们的\OriginalTerm{圣髑匣}{reliquaries}，视之为天主圣者的纪念物，他们曾为真理奋斗，直至身体死亡，为使真宗教得以宣明，虚假伪造的宗教得以揭露。若在他们之前有人曾认为这些宗教确实是虚假伪造的，他们也不敢表达其信念。但有谁曾听见信友的司铎，站在为尊崇敬拜天主而建于某位\OriginalTerm{殉道者}{martyr}圣髑之上的祭台前，在祈祷中说：\Quote{我向你呈上祭献，\Person{伯多禄}啊，或\Person{保禄}啊，或\Person{西彼廉}啊？} 因为在他们的墓前，祭献是向天主呈上的——就是那位使他们既成为人又成为殉道者，并与圣天使同享天上尊荣的天主；我们之所以如此尊敬他们的纪念，是因为藉此我们既能因他们的胜利而感谢真天主，又能通过重新记念他们，激励自己效法他们，追求获得同样的冠冕与棕榈枝，并呼求他们所呼求的同一天主助佑。所以，虔敬者可能在殉道者的处所献上的任何敬礼，都只是向他们纪念表示的敬意，而非如同对神明一般向亡者献上的神圣礼仪或祭献。甚至那些带食物去那里的人——其实更好的基督徒并不这样做，且在世上大部分地方根本无人如此行——这样做是为了使食物通过殉道者的\OriginalTerm{功劳}{merits}，并因殉道者之主的名而被圣化，他们先呈上食物并献上祈祷，然后取走食用，或部分施予穷乏之人。但凡明了基督徒那独一无二的祭献——就是在那些地方所举行的祭献——的人，也便知道这些祭献并非向殉道者呈上的。所以，我们既不以属神的敬礼，也不以属人的罪恶——那些他们用来敬拜其神明的方式——来敬礼我们的殉道者；我们既不会向他们献祭，也不会将神明的罪恶转变为他们的神圣礼仪。因为谁若愿意且有能力的，可读一读\WorkTitle{亚历山大致其母奥林匹亚斯书}，信中他讲述了司铎\Person{莱昂}所启示他的事；读过的人要回想起信的内容，便会看到那些可憎之事如何流传下来，不是通过诗人的作品，而是通过埃及人的奥秘著作，关乎女神\Person{伊西斯}——\Person{奥西里斯}之妻，以及两者的父母，所有这些人物，据这些著作称，都是王族。据说\Person{伊西斯}在向父母献祭时，发现了一片大麦，她带了些麦穗给她的王夫，以及他的谋士\Person{墨丘利}，因此人们将她等同于\Person{刻瑞斯}。读过此信的人或可从中看到那些人是什么品性，他们死后竟如同神明般被设立了神圣礼仪，以及他们的哪些行为为这些礼仪提供了契机。他们绝不敢在任何方面把那些人——尽管他们认为他们是神——与我们神圣的殉道者相比，尽管我们不认为殉道者是神。因为我们不为我们的殉道者设立司铎和呈献祭献，如同他们为他们的亡者所做的那样，因为那是不相称、不当且不合法的，这些唯独归于天主；如此，我们也不以他们自己的罪恶，或以那些在可耻戏剧中庆祝神明的罪恶来取悦他们；这些罪恶要么是他们在身为凡人时真实犯下的，要么——如果他们从来不是凡人——是为了取悦有害的\OriginalTerm{魔鬼}{demons}而捏造的虚构罪恶。\Person{苏格拉底}的神，如果他真有一位神，绝不会属于此类魔鬼。然而，或许那些希图在制造神明之术上出类拔萃的人，把这样的神强加给了一个对其术全然陌生、毫无关联的人。我们还需要多说什么呢？但凡稍有智慧的人，都不会以为为了死后的真福生命，就当敬拜魔鬼。但或许他们会说，所有的神都是好的，而魔鬼中有些是坏的，有些是好的，应当敬拜的是好的那些，好能藉着他们达至永远的真福生命。我们将在下一卷书中致力于审视这种见解。\par

\SourceRecord{Translated by Marcus Dods.；Nicene and Post-Nicene Fathers, First Series；Vol. 2.；Edited by Philip Schaff.；Buffalo, NY: Christian Literature Publishing Co.,；1887.；<http://www.newadvent.org/fathers/120108.htm>.}

% structure: p0001 source=h1 target=section reason=page-title

\section{\WorkTitle{天主之城}（第九卷）}

在前一卷中已表明，必须弃绝魔鬼的敬拜，因为它们以千百种方式宣称自己是邪恶的精灵，\Person{奥思定}在本卷中反驳那些声称魔鬼中有区别的人，即有些是恶的，有些是善的；并且，在驳斥了这一区分之后，他证明了，提供人类永福的职分不属于任何魔鬼，而只属于\Person{基督}。\par

% structure: p0003 source=h2 target=subsection reason=relative-heading-rank; base=h2

\subsection{第一章——讨论已到达之处，以及尚待处理之事。}

有些人提出，既有善神也有恶神；但另一些人，对神祇怀有更深的敬意，将如此多的尊崇和赞美归于它们，以致排除了任何神祇可能是邪恶的假设。然而，那些主张既有恶神也有善神的人，将魔鬼归在\Quote{神祇}的名下，有时——尽管更罕见——也称神祇为魔鬼；因此他们承认，被他们奉为众神之王和首领的\Person{朱庇特}，被\Person{荷马}称为魔鬼。另一方面，那些主张所有神祇都是善的、且远远优于那些被公正地称为善人的人，则因魔鬼的行为——这些行为他们既不能否认，也不能归咎于他们所肯定其善性的神祇——而感到有必要区分神祇与魔鬼；因此，每当他们在不可见之灵显现其能力的行为或情感中发现任何冒犯之处，他们便相信这是出于魔鬼，而非神祇。同时，他们相信，既然没有神祇能与人类直接交往，这些魔鬼便处于中保的地位，带着祈祷上升，带着恩赐返回。这是柏拉图学派的观点，他们是其哲学家中最能干、最受尊敬的，因此我们选择与他们辩论这个问题——敬拜众多神祇对于获得来生的永福是否有所帮助。这就是为什么在前一卷中，我们探究了魔鬼——它们以善良智慧之人所厌恶和咒骂的事物为乐，以诗人们所写的并非关于人、而是关于神祇本身的亵渎和不道德的故事为乐，以魔法妖术的邪恶与罪恶的暴力为乐——如何能被看作比人类更亲近神祇、与神祇更友好，并能在善人与善神之间担任中保；而我们已证明这是绝对不可能的。\par

% structure: p0005 source=h2 target=subsection reason=relative-heading-rank; base=h2

\subsection{第二章——在低于神祇的魔鬼中，是否有一些善灵，在其守护下，人的灵魂能达到真福。}

因此，按照前卷结尾处的承诺，本卷应包含一场讨论，不是讨论神祇之间的差异——按照柏拉图学派，所有神祇都是善的——也不是讨论神祇与魔鬼之间的差异——前者被他们以巨大的间隔与人类隔开，而后者被置于神祇与人类之间——而是讨论魔鬼自身的差异，因为他们确实作出了这种区分。我们将就与我们的主题相关的部分加以讨论。人们普遍而惯常地相信，有些魔鬼是恶的，另一些是善的；这一观点，无论是柏拉图学派的还是任何其他派别的，都绝不应被沉默地忽略，以免有人认为他应当敬拜善的魔鬼，以便通过它们的中保作用而被神祇接纳——他相信所有神祇都是善的——并在死后与神祇同住；而他实际上会因此陷入邪恶精灵的网罗，远远偏离真天主，因为唯独与祂同在、且在祂内，人的灵魂，即理性的和理智的灵魂，才得享真福。\par

% structure: p0007 source=h2 target=subsection reason=relative-heading-rank; base=h2

\subsection{第三章——\Person{阿普列乌斯}如何论述魔鬼，他虽不否认它们具有理性，却未将德性归于它们。}

那么，善的魔鬼与恶的魔鬼之间有何区别呢？因为柏拉图学派的\Person{阿普列乌斯}，在一篇关于这整个主题的论文中，关于它们的空气般的身体谈了很多，却对灵性的德性只字未提——如果它们是善的，它们必然被赋予了这些德性。因此，他对于能给予它们幸福的事物未置一词；但他却给出了它们不幸的证明，承认它们的理智——凭此它们位列理性存在——不仅没有以德性浸润和巩固，从而能抵抗一切非理性的情欲，反而以某种方式被暴风雨般的情感所激动，因此与愚人的心灵处于同一水平。他自己的话如下：\Quote{正是这类魔鬼，诗人们在未犯严重错误的情况下，想像神祇憎恨和喜爱人群中的某些人，使一些人兴旺尊贵，而敌对和困扰另一些人。因此，怜悯、愤慨、悲伤、快乐，每一种人类的情感，魔鬼都经历着，带着同样的心灵扰动，以及同样的情感与思想起伏。这些骚动和风暴使它们远离了天上神祇的宁静。} 在这些话中，难道还有疑问吗？他所说的被情欲搅动如暴风雨般的海洋的，并非它们灵性本质的较低部分，而是魔鬼赖以位列理性存在的那个心灵本身。因此，它们甚至不能与智慧之人相比，智慧之人以不受扰乱的心灵抗拒这些在此生中不可避免的扰动——人类软弱永远无法幸免——并且不让自己赞同或做出任何可能偏离智慧和正直法则之事。它们在品性上，虽然不在身体外貌上，与邪恶愚昧的人相似。我甚至可以说它们更坏，因为它们在邪恶中变老，且惩罚无法纠正。它们的心灵，正如\Person{阿普列乌斯}所言，如同被风暴抛掷的海洋，在其灵魂中没有任何真理或德性的支撑点，能让它们抵抗其动荡和败坏的情感。\par

% structure: p0009 source=h2 target=subsection reason=relative-heading-rank; base=h2

\subsection{第四章——逍遥学派和斯多葛学派关于心灵情感的观点。}

在哲学家们中间，关于这些心灵的激情，希腊人称之为\OriginalTerm{激情}{\GreekText{παθη}}，而我们自己的一些作家，如\Person{西塞罗}，称它们为\Quote{扰动}，\Quote{感受}，或者为了更准确地翻译希腊词，称之为\Quote{激情}。有人认为，即使智者也会受到这些扰动的影响，但受到理性的节制和约束，理性为它们设立法则，将其限制在必要的范围内。这是柏拉图主义者和亚里士多德主义者的观点；因为\Person{亚里士多德}是\Person{柏拉图}的门徒，也是逍遥学派的创始人。但另一些人，如斯多葛派，则认为智者不会受到这些扰动的影响。但\Person{西塞罗}在其著作\WorkTitle{论目的}中指出，斯多葛派与柏拉图主义者和逍遥学派在此问题上的分歧，是言辞上的而非实质上的；因为斯多葛派拒绝将\Quote{善}这个词用于外在的和身体上的好处，他们认为唯一的善是德性，即良好生活的技艺，而这仅存在于心灵之中。其他哲学家则使用通常的习惯用语，毫不踌躇地称这些事物为善，尽管与引导我们生活的德性相比，它们是渺小的，微不足道的。因此，显然，无论这些外在事物被称作善还是好处，双方对它们的评价是相同的，而在这件事上，斯多葛派只是满足于一种新的措辞而已。那么，在我看来，在这个问题上，即智者究竟是受心灵激情的影响，还是完全不受其扰，争论也是言辞的而非实质的；因为我认为，如果考察的是实质而非单纯的言辞，斯多葛派实际上与柏拉图主义者和逍遥学派持完全相同的见解。为了简短起见，我略去可能引用的其他证据，只讲一个我认为具有决定性的证据。\Person{奥卢斯·革利乌斯}，一位学识渊博、文辞优雅的人，在他的著作\WorkTitle{阿提卡之夜}中记述，他曾与一位著名的斯多葛派哲学家一同航行；他详细而饶有兴味地记述，我则简略叙述如下：当船只遭遇猛烈风暴，颠簸危险时，那位哲学家吓得脸色苍白。船上的人注意到了这一情形，他们虽然自身也面临死亡威胁，却好奇地想看看，一位哲学家是否也会像常人一样惊慌失措。风暴过后，待到安全无虞，可以重新交谈时，船上一位富裕奢华的亚细亚乘客便开始嘲笑那位哲学家，奚落他竟因恐惧而面容失色，而他本人面对即将到来的毁灭却无动于衷。但那位哲学家借用了苏格拉底派哲学家\Person{阿里斯提波斯}的回答——\Person{阿里斯提波斯}发现自己被同样品性的人嘲弄时，曾如此答道：\Quote{你无需为一个放荡的败家子的灵魂担忧，而我却有理由为\Person{阿里斯提波斯}的灵魂感到惊慌。}这样打发那位富人之后，\Person{奥卢斯·革利乌斯}为了求知而非刁难，询问那位哲学家他恐惧的原因是什么？那位哲学家很乐意教导一位如此热心求知的人，立即从行囊中取出一卷斯多葛派哲学家\Person{爱比克泰德}的书，书中阐述的学说与斯多葛学派创始人\Person{芝诺}和\Person{克吕西波}的见解正相吻合。\Person{奥卢斯·革利乌斯}说，他在那本书中读到，斯多葛派主张，外界事物在灵魂中造成某些印象，他们称之为\OriginalTerm{印象}{phantasiæ}，而灵魂并无能力决定是否或何时受到这些印象的侵袭。当这些印象由令人惊恐和畏惧的事物引起时，它们必然触动甚至智者的灵魂，以至于片刻之间他会因恐惧而战栗，或因忧伤而沮丧，因为这些印象发生在理性和自制的运作之前；但这并不意味着心灵接受了这些不良印象，或赞同、同意它们。因为他们认为，这种同意取决于人自己；智者的心灵与愚人的心灵区别在于：愚人的心灵屈服于这些激情并同意它们，而智者的心灵，尽管无法避免为其所侵，却仍以不动摇的坚定保持着对依据理性应当追求或避免的那些事物的真实而稳固的确信。关于\Person{奥卢斯·革利乌斯}记述他在\Person{爱比克泰德}的书中读到的斯多葛派的观点和学说，我已尽力转述，也许未能如他那般选词精当，但更为简短，且我认为更为明晰。如果这是真的，那么在关于心灵的激情和扰动的问题上，斯多葛派与其他哲学家的观点便没有区别，或几乎没有区别，因为双方都一致主张，智者的心灵和理性不会屈从于它们。或许，斯多葛派这样断言的意思是，智者所具有的智慧不受任何错误的蒙蔽，不被任何污点所玷染，但在保持智慧不受扰动的条件下，他仍会受到今生所谓的善与恶（或如他们更愿称呼的，好处与坏处）所造成的印象的影响。因为无需多言，倘若那位哲学家对他认为即将失去的东西——生命和身体的安全——毫不在意，他就不会因危险而如此恐惧，乃至面色苍白暴露了他的恐惧。然而，他可能经受这种心灵扰乱，却仍然保持坚定的确信：生命和身体的安全——即风暴的暴力所威胁要摧毁的——并非那种能使其拥有者成为善人的善，如同拥有义德那样。但就他们坚持我们不应称之为善而应称之为好处而言，他们是在言辞上争执而忽略了实质。因为，称它们为善还是好处，有什么不同呢？既然斯多葛主义者和逍遥主义者同样担心失去它们，他们虽然称呼不同，但评价相同。双方都向我们保证，倘若因面临丧失这些善或好处的威胁而被唆使去做某些不道德或犯罪的事，他们宁愿失去那些维护身体舒适和安全的事物，也不愿去做违犯义德的事。因此，一个坚定树立了此等决心的心灵，绝不容许任何扰乱在与理性相悖的情况下得逞，即使它们攻击灵魂的较软弱部分；不仅如此，它还统治它们，在拒绝同意并抵抗它们的同时，施行德性的统御。\Person{维吉尔}在描述\Person{埃涅阿斯}时便赋予他这样的品质，他说，\par

他屹立不动，不为眼泪所动，\newline{} 最温柔的话语，他亦不怜悯地倾听。\par

% structure: p0012 source=h2 target=subsection reason=relative-heading-rank; base=h2

\subsection{侵袭基督徒灵魂的偏情，并不诱使他们作恶，而是锻炼他们的德性。}

目前我们无需对圣经教义——即基督徒知识的纲要——中关于这些\OriginalTerm{偏情}{passions}的教导作详尽周密的阐述。它将心灵本身服从于天主，使祂能治理并扶助它，又将偏情置于心灵之下，好节制、约束并引导它们趋于正当的用途。在我们的伦理学中，我们并不太追究一个虔诚的灵魂是否发怒，而追究他为何发怒；不是他是否悲伤，而是他悲伤的原因；不是他是否恐惧，而是他恐惧的对象。因为我认为，任何思维正确的人，都不会责备那为求改正他人而对作恶者所生的愤怒，不会责备那意在救助受苦者的悲伤，不会责备那唯恐身陷险境的人毁灭。\OriginalTerm{斯多葛派}{Stoics}确实惯于谴责怜悯。然而，我们之前提到的那个斯多葛派，如果他被怜悯所动而去救助一个同类，比他因恐惧船难而烦乱，要体面得多！\Person{西塞罗}称赞\Person{恺撒}的话，要美好、更人道且更契合虔诚情感得多。他说：\Quote{在你的诸德性中，没有比你的怜悯更值得赞赏和悦人心意的了。}怜悯是什么？岂不就是对他人不幸的同感，这同感促使我们若可能就去帮助他么？这种情感服从于理性，当怜悯的施行不违反正义时，例如救济贫苦者、宽恕悔罪者。\Person{西塞罗}，这位善于辞令的人，毫不犹疑地称此为一种德性，而斯多葛派却恬不知耻地将其归入恶习，尽管如那位杰出的斯多葛派\Person{爱比克泰德}的著作中所引述的\Person{芝诺}与\Person{克律西波斯}——这一学派的创立者——的意见，他们承认这类偏情会侵扰智者的灵魂，而他们本希望智者脱离一切恶习。因此可推知，在他们看来，这些偏情本身并不被视为恶习，因为它们侵袭智者，却并未迫使他作出违反理性与德性的行为；所以，\OriginalTerm{逍遥学派}{Peripatetics}或\OriginalTerm{柏拉图主义者}{Platonists}与斯多葛派的意见实为一体。但是，正如\Person{西塞罗}所言，这些可悲的希腊人热中于争辩过于寻求真理，单纯的语词论争正是他们的祸根。然而，我们或可合理地发问：即便我们追随德性时，我们之情动于这些偏情，是否属于此生的软弱一部分？因为圣善的天使们，在他们惩罚那由天主的永恒律法判定受罚者时，感受不到愤怒；在救助遭受苦难者时，没有怜悯；在护佑身处险境者时，亦无恐惧；然而通常的语言也将这些心理情绪归之于他们，因为他们虽无我们的软弱，他们的行动却类似于那些使我们被这些情绪所触发的行动；因此，圣经中甚至说天主发怒，却毫无内心的纷扰。因为这个词语是指祂复仇的效果，而非指扰乱心神的情绪。\par

% structure: p0014 source=h2 target=subsection reason=relative-heading-rank; base=h2

\subsection{论据\Person{阿普列乌斯}所说，那些被认为在神祇与人之间作中介的魔鬼为偏情所激动的道理。}

我们暂且搁下关于圣天使的问题，来考察\OriginalTerm{柏拉图主义者}{Platonists}的意见：那些中介于神祇与人之间的\OriginalTerm{魔鬼}{demons}，为\OriginalTerm{偏情}{passions}所激动。因为，如果他们的心即便暴露在偏情的侵袭下，仍能保持自由并超越其上，\Person{阿普列乌斯}就不能说，他们的心犹如被暴风掀起的大海，为偏情所颠簸。那么，他们的心——即他们灵魂的上部，使之成为理性存在者，且若这心确实存在，就应当辖制和约束灵魂下部的狂暴偏情——我说，他们的这颗心，据这位柏拉图主义者所说，正被偏情的狂风摧荡。因此，魔鬼的心便臣服于恐惧、愤怒、贪欲以及一切类似情绪之下。那么，在他们之中，哪一部分是自由且装备着智慧的，使之能讨神祇喜悦，并作为引导人们走向纯净生活的合适向导呢？既然他们最高的部分，既沦为偏情的奴隶并臣服于邪恶，只会使他们更处心积虑地欺骗和诱惑，这与他们拥有的欲望的精神力量和强弱程度成正比。\par

% structure: p0016 source=h2 target=subsection reason=relative-heading-rank; base=h2

\subsection{柏拉图主义者主张，诗人错谬地描述神祇遭受派系纷争之苦，而魔鬼——而非神祇——才易受此影响。}

但若有人说，诗人并非全无真实地声称魔鬼们——不是所有魔鬼，而只是那些邪恶的——会强烈地爱或恨某些人（因为\Person{阿普莱乌斯}正是就它们说，它们被强烈的情感潮流所驱使），那么，当\Person{阿普莱乌斯}在同一上下文中，以其\OriginalTerm{气状的身体}{aerial bodies}将所有的魔鬼，而不仅仅是邪恶的，描述为介于神与人之间的中介时，我们如何能接受这种解释呢？按照他的说法，诗人的虚构在于将魔鬼变成神，赋予它们神的名字，并把它们作为盟友或敌人指派给个别的人，运用这种诗意的特许，尽管他们声称神在特性上与魔鬼大不相同，且因其天上的居所和充满的\OriginalTerm{真福}{beatitude}而远高于它们。我认为，这就是诗人的虚构：说这些不是神的是神，并且，在神的名号之下，它们为它们所爱或恨的人，以强烈的党派情感彼此争斗。\Person{阿普莱乌斯}说，这离真理不远，因为尽管它们被错误地以神的名字称呼，它们仍按其本身的特性被描述为魔鬼。他说，\Person{荷马}的\Person{密涅瓦}便属于这一类，\Quote{她介入希腊人的队伍，以阻止\Person{阿喀琉斯}。} 因为他认为这个\Person{密涅瓦}是诗意的虚构；他认为\Person{密涅瓦}是一位女神，并将她列入那些他相信居住在崇高的以太境界、远离与人交往、全都美善而蒙福的神之中。但是，有一个魔鬼支持希腊人而反对特洛伊人，正如另一位，同一诗人在\Person{维纳斯}或\Person{马尔斯}（这些神高居天上的住所，超然于尘世之上）的名义下所提到的，是特洛伊人的盟友和希腊人的仇敌；而这些魔鬼为自己所爱的对抗自己所恨的——在这一切上，他承认诗人所说非常接近真理。因为诗人关于这些存在所说的话，\Person{阿普莱乌斯}同样赋予了它们那扰乱人类的、猛烈而狂暴的激情，因此它们能够产生爱和恨，这些情感并非基于公义，而是出于派系之心，就像竞赛或狩猎的观众会偏爱或偏恶一样。这位\OriginalTerm{柏拉图主义者}{Platonist}似乎极其担心，诗人的虚构会被认为属于那些神，而非属于顶着它们名字的魔鬼。\par

% structure: p0018 source=h2 target=subsection reason=relative-heading-rank; base=h2

\subsection{第八章——阿普莱乌斯如何界定居住在天上的神、占据空中的魔鬼和住在地上的人}

\Person{阿普莱乌斯}对魔鬼的定义——当然包括所有的魔鬼——是：它们在\OriginalTerm{本性}{nature}上是\OriginalTerm{生物}{animals}，灵魂受\OriginalTerm{情感}{passion}支配，心灵有\OriginalTerm{理性}{reasonable}，身体为\OriginalTerm{气状}{aerial}，存在是\OriginalTerm{永恒的}{eternal}。然而，在这五种特质中，他丝毫未提及任何唯独属于善人而非恶人的品质。因为当\Person{阿普莱乌斯}首先谈及天上者，然后扩展其描述，以包括那些远居于下界之地上的存在，从而在描述了理性存在的两个极端之后，得以进而谈论中介的魔鬼时，他说：\Quote{因此，人，被赋予理性和言语的官能，灵魂不死，肢体却会朽坏；心神软弱而烦乱，身体迟钝而腐朽，性格各不相同，无知却彼此相似，顽固于自己的大胆，固执于自己的希望，劳苦徒然无益，命运日渐衰微；他们的族类永存，自己却个个消亡，一代代人子孙充盈，他们的生命迅疾，智慧却迟缓，死亡突然，生命充满哀号——这便是住在地上的人。} 他在历数那属于大多数人的众多品质时，当他谈到他们的智慧迟缓时，岂不是忘记了那少数人的特性？倘若删去这一点，那他如此精心刻画的这幅人类画像便有了缺陷。而当他在赞美神的卓越时，他声言他们的卓越正在于他以为人需靠智慧才能达致的那真福。因此，如果他希望我们相信部分魔鬼是善的，他就应该在他的描述中加入一些东西，使我们能看出它们与神一样享有一份真福，或与人一样具有某种智慧。然而，事实上，他没有提及任何可以使善者区别于恶者的良善品质。因为尽管他出于害怕冒犯不是魔鬼本身，而是他正为其写作的敬拜它们的人，而克制没有充分详述它们的邪恶，但他已向有洞察力的读者充分表明了他对它们的看法；因为仅在身体永恒这一条上，他将它们比作神——他宣称所有的神都是善而蒙福的，并且完全没有他自己所谓的魔鬼的那些狂暴的情感；至于灵魂，他相当明确地断言，它们与人相似而非与神相似，且这种相似不在于拥有智慧（智慧连人也能够获得），而在于情感的搅扰，这搅扰支配着愚昧和邪恶的人，但被善人和智者如此管理，以至于他们宁愿不容其进入内心，而非去战胜它。因为假如他有意让人理解魔鬼不仅仅在身体上、也在灵魂上在永恒方面类似神，那他肯定也会让人分享这一特权，因为作为\OriginalTerm{柏拉图主义者}{Platonist}，他自然必定认为人的灵魂是永恒的。因此，当描述这众生类别时，他说他们的灵魂不死，肢体却会朽坏。于是，结论是：既然人因其身体可朽而不能与神共享永恒，那么魔鬼则因其身体不朽而与神共享永恒。\par

% structure: p0020 source=h2 target=subsection reason=relative-heading-rank; base=h2

\subsection{第九章——魔鬼的转求能否为人赢得天上众神的友谊}

那么，人怎能希望通过这样的中介者\OriginalTerm{中介者}{mediators}获得神明的友谊呢？这些中介者与人一样，在每个生物更优越的部分——即灵魂——上有缺陷，而只在较低劣的部分——即身体——上与神明相似。因为一个生物或动物由灵魂和身体构成，这两部分中灵魂无疑是更优越的；即使它邪恶而软弱，也显然比最健康、最强壮的身体更好，因为其本性的更高卓越不会因邪恶的污染而降低到身体的层次，正如金子，即便失去光泽，也比最纯的银或铅更珍贵。然而，这些中介者，通过他们的中介，人事与神事得以和谐，却拥有与神明共有的永恒身体，与人共有的邪恶灵魂——仿佛这些魔鬼\OriginalTerm{魔鬼}{demons}用来联合神与人的宗教是身体性的，而非精神性的事。那么，是怎样的邪恶或惩罚使这些虚伪狡诈的中介者仿佛被倒悬起来，以致他们较低的部分——身体——与在上的神明相连，而较高的部分——灵魂——却被束缚于在下的人间；以服役的部分与天上的神明结合，而以统治的部分与地上的居民一同受难？因为身体是仆人，正如\Person{撒路斯提乌斯}所说：\Quote{我们以灵魂统治，以身体服从；}并补充道：\Quote{前者我们与神明共有，后者我们与野兽共有。}他在这里说的是人；人也像野兽一样有必死的身体。这些魔鬼，我们的哲学界朋友为我们提供作为与神明之间的中介者，确实可以说，灵魂和身体，前者我们与神明共有，后者我们与人共有；但如我所言，他们仿佛被倒悬捆绑，将仆人——身体——与神明共有，将主人——灵魂——与悲惨的人共有——他们较低的部分被高举，较高的部分被贬抑。因此，若有人认为，他们不像地上的动物那样因死亡而灵魂与身体分离，所以在永恒性上与神明相似，那么他们的身体不应被视为永恒胜利的战车，而应被视为永恒惩罚的锁链。\par

% structure: p0022 source=h2 target=subsection reason=relative-heading-rank; base=h2

\subsection{十、--- 根据\Person{普罗提诺}，身体必死的人比身体永恒的魔鬼更不悲惨。}

\Person{普罗提诺}，其记忆犹新，享有比柏拉图任何其他门徒更好地理解柏拉图的声誉。在谈论人的灵魂时，他说：\Quote{父\OriginalTerm{父}{Father}出于怜悯使他们的锁链成为必死的；}也就是说，他认为，由于父的慈悲，人有了必死的身体，就不致永远被困于此生的苦难中。但这种慈悲，魔鬼\OriginalTerm{魔鬼}{demons}被认为不配获得，他们得到了一个受情感支配的灵魂，并结合了一个不像人那样必死、而是永恒的身体。因为他们本应比人更为有福，倘若他们像人一样有必死的身体，又像神明一样有蒙福的灵魂。他们也应当与人不相上下，倘若他们结合了悲惨的灵魂，至少像人一样接受了必死的身体，这样死亡或许能使他们摆脱烦恼——只要他们至少达到若干程度的虔诚。但事实是，他们不仅不比人更为有福，与人一样有着悲惨的灵魂，而且更为悲惨，因为他们永远被束缚于身体；因为他并未让我们推断通过智慧和虔诚的某个进展他们能成为神明，而是明确地说他们永远是魔鬼。\par

% structure: p0024 source=h2 target=subsection reason=relative-heading-rank; base=h2

\subsection{十一、--- 论柏拉图主义者的意见：人的灵魂脱离身体后成为魔鬼。}

他说，的确，人的灵魂是魔鬼\OriginalTerm{魔鬼}{demons}，且人若良善，就变为\Emph{\OriginalTerm{拉列斯}{Lares}}，若邪恶，就变为\Emph{\OriginalTerm{勒穆瑞斯}{Lemures}}或\Emph{\OriginalTerm{拉尔瓦}{Larvæ}}，若不确定其善恶，则变为\Emph{\OriginalTerm{曼内斯}{Manes}}。谁不能一眼看出这只是一个将人卷入道德毁灭的漩涡呢？因为，无论人曾多么邪恶，如果他们认为自己死后会变成拉尔瓦或神圣的曼内斯，他们越喜欢加害他人，就会变得越坏；因为拉尔瓦是由恶人变成的有害魔鬼，这些人必会以为死后他们会受到祭祀和神明般的尊崇，好让他们能施加伤害。但这个问题我们不能深究。他还说，受祝福的人在希腊文中被称为\OriginalTerm{尤戴蒙}{\GreekText{εὐδαίμονες}}，因为他们是好的灵魂，即好的魔鬼，以此证实他关于人的灵魂是魔鬼的意见。\par

% structure: p0026 source=h2 target=subsection reason=relative-heading-rank; base=h2

\subsection{十二、--- 论柏拉图主义者区分人性和魔鬼本性的三种相反特质。}

但我们现在所谈论的，乃是他所描述的那些正当介于\OriginalTerm{神}{gods}与人之间的存在；就其本性而言是生物，心灵是有理性的，灵魂受情欲支配，身体是气性的，寿命是永恒的。当他区分了安置在至高天的神与安置在大地上的人，不仅根据位置，也根据其本性不等的尊荣时，他这样总结道：\Quote{这里你有两种生物：神，因居所的崇高、生命的永恒、本性的完美而远远有别于人；因为他们的居所间隔如此之大，以致彼此间无法有亲密的交流；而且一种的生命是永恒而无衰的，另一种的生命却是易逝而不稳固的；并且神的灵是在福乐中被高举，人的灵却沉于苦难之中。}在这里，我发现在存在的两极，即最高者和最低者身上，被赋予了三种对立的品质。因为，在提到我们应赞叹神的三项品质之后，他又以别的话重述了这三项，作为人的缺陷的陪衬。这三项品质是：\Quote{居所的崇高、生命的永恒、本性的完美。}他再次提及这些，以便凸显它们与人类状况的对比。正如他提到了\Quote{居所的崇高}，他说，\Quote{他们的居所间隔如此之大；}正如他提到了\Quote{生命的永恒}，他说，\Quote{神圣的生命是永恒而无衰的，人的生命却是易逝而不稳固的；}正如他提到了\Quote{本性的完美}，他说，\Quote{神的灵是在福乐中被高举，人的灵却沉于苦难之中。}因此，他将这三者归与神：崇高、永恒、福乐；而将相反者归与人：居所的卑下、有死、苦难。\par

% structure: p0028 source=h2 target=subsection reason=relative-heading-rank; base=h2

\subsection{\OriginalTerm{魔鬼}{demons}如何能在\OriginalTerm{神}{gods}与人之间担任中介，如果它们与两者毫无共同之处，既不似神那样有福，也不像人那样悲惨。}

现在，如果我们试图在这些对立之间寻找\OriginalTerm{魔鬼}{demons}所占据的中间状态，那么关于它们的位置就毫无疑问了；因为在最高和最低位之间，有一个正确被视为并称为中间位置的地方。剩下的另外两种品质，我们必须更仔细地考察，以便看出它们是完全与魔鬼无关，还是以怎样的方式赋予它们而不损害其中间地位。我们可以放弃它们与魔鬼无关的想法。因为我们不能说，魔鬼作为理性的生物，既不是有福也不是悲惨的，如同我们说那些没有感觉和理性的野兽和植物那样，也不能像我们说中间位置既不是最高也不是最低那样。魔鬼既然是有理性的，就必定要么悲惨，要么有福。同样，我们不能说它们既不是有死的也不是不朽的；因为一切有生命之物要么永远活着，要么以死终结生命。此外，我们的作者说过魔鬼是永恒的。那么，我们还能设想什么呢？无非是这些中间的存在在上述两种品质中，一项与\OriginalTerm{神}{gods}相似，另一项与人相似。因为如果它们从上面接受了两者，或者从下面接受了两者，它们就不再是中间的，而要么上升到上面的神那里，要么下降到下面的人那里。所以，既然已经证明它们必然拥有这两项品质，如果它们各从一方接受一项，它们就将保持中间位置。因此，既然它们不能从下面得到永恒，因为那里没有永恒可得，它们就必须从上面得到；于是，它们别无选择，只能通过从人那里接受悲惨来完成它们的中间地位。\par

那么，根据\OriginalTerm{柏拉图主义者}{Platonists}的说法，居于最高位置的\OriginalTerm{神}{gods}享有永恒的福乐，或谓有福的永恒；居于最低位置的人则享有有死的悲惨，或谓悲惨的有死性；而居于中间的\OriginalTerm{魔鬼}{demons}则享有悲惨的永恒，或谓永恒的悲惨。至于\OriginalTerm{\Person{阿普列乌斯}}{Apuleius}在给魔鬼的定义中所包括的那五项，他并未如他所承诺的那样，表明魔鬼是中间的。因为其中三项——即它们的本性是生物，心灵是有理性的，灵魂受情欲支配——他说它们与人共有；一项，即它们的永恒，与神共有；还有一项专属于它们自己，即它们的\OriginalTerm{气性}{aerial}身体。那么，当它们有三项与最低者共有，只有一项与最高者共有时，它们如何是中间的呢？谁看不出，当它们倾向于并被压低至最低极端时，其中间位置就被放弃了，且其倾向越甚，放弃越远？然而，也许我们应当因它们气性身体这一特性而接受它们为中间者，正如两极各有其对应的身体——神有\OriginalTerm{以太}{ethereal}身体，人有\OriginalTerm{地上}{terrestrial}身体——并且因为它们与人有共有的两种品质，也与神共有，即它们的生物本性和理性心灵。因为\Person{阿普列乌斯}本人在谈到神和人时说：\Quote{你们拥有两种生物本性。}而柏拉图主义者习惯于将理性心灵归于神。剩下两种品质，它们受情欲支配的倾向，以及它们的永恒——前者它们与人共有，后者与神共有；这样，它们既不飞升至最高极端，也不被压至最低极端，而是完美地平衡在它们的中间位置。然而，这恰恰是构成魔鬼的永恒悲惨或悲惨永恒的情况。因为凡是说它们的灵魂受情欲支配的人，也会说它们是悲惨的，若不是为敬拜它们的人感到羞耻的话。此外，由于世界并非由偶然的巧合统治，而是如柏拉图主义者自己所承认的，由至高天主的\OriginalTerm{天主眷顾}{providence}统治，魔鬼的悲惨就不会是永恒的，除非它们的邪恶是极大的。\par

那么，如果\Emph{有福者}被正确地称为\OriginalTerm{尤德蒙}{eudemons}，介于神与人之间的魔鬼就不是尤德蒙。那么，那些善的魔鬼居于何处？他们高于人而低于神，既援助人，又服事神。因为如果他们是善的且永恒的，他们无疑是有福的。但永恒的福乐摧毁了他们中间的特性，使他们与神相似，与人相远。因此，柏拉图主义者将徒劳地努力证明，善的魔鬼如果既不死不灭又有福，如何能被正确地认为居于不死不灭且有福的神与有死且悲惨的人之间。因为如果他们与神共同拥有不死不灭和有福，而不与人共有这两者中的任何一个（人既悲惨又有死），他们岂不是远离人而与神联合，而非介于两者之间吗？他们若是在一种特性上与一方相同，在另一种特性上与另一方相同，他们就是居间的，正如人是介于天使与野兽之间的某种中项——野兽是无理性且有死的动物，天使是有理性且不死不灭的，而人低于天使、高于野兽，与后者共有有死，与前者共有理性，因此是有理性且有死的动物。因此，当我们寻求有福的不死不灭者与悲惨的有死者之间的中项时，我们应当找到一种存在，要么是有死且有福的，要么是不死不灭且悲惨的。\par

% structure: p0032 source=h2 target=subsection reason=relative-heading-rank; base=h2

\subsection{第14章——人虽是有死的，能否享有真福？}

人间有一大问题：人是否能是有死且有福的？有些人，对人的境况持较谦卑的看法，否认人在继续这有死的生命时能享有真福；另一些人则鄙弃这种想法，大胆主张，即使有死，人也可藉由获得智慧而成为有福的。但如果这是事实，为什么这些智者不成为悲惨的有死者与有福的不死不灭者之间的中介呢？因为他们与后者共有有福，与前者共有有死。当然，如果他们有福，他们就不嫉妒任何人（因为有什么比嫉妒更悲惨呢？），而是竭尽全力帮助悲惨的有死者获得真福，使他们死后能成为不死不灭的，并与有福且不死不灭的天使为伴。\par

% structure: p0034 source=h2 target=subsection reason=relative-heading-rank; base=h2

\subsection{第15章——论人而天主者\Person{耶稣基督}，天主与人之间的中保。}

但是，如果（正如更可能和可信的那样）只要人是有死的，就必然也是悲惨的，那么我们必须寻求一位中保，他不仅是人，而且也是天主，好藉着他那有福的有死性，将人从有死的悲惨中领到有福的不死不灭里。在这位中保身上，需要两件事：他成为有死的，且不持续有死。他确实成为了有死的，但这并非使\OriginalTerm{圣言}{Word}的神性软弱，而是摄取了肉体的软弱。他也没有在肉体中持续有死，而是从死者中复活了它；因为他中保工作的果实正是，那些他为之成为中保并完成救赎的人，不应永远停留在肉体的死亡中。因此，我们与天主之间的中保，理当兼具短暂的有死与永恒的福乐，好藉着那短暂者而与有死之人相似，并能将他们从有死中转移到那永恒者那里。所以，善的天使不能在悲惨的有死者和有福的不死不灭者之间作中介，因为他们自己也是有福且不死不灭的；但恶的天使却能作中介，因为他们如同后者一样不死不灭，又如同前者一样悲惨。与这些恶天使相对立的是善的中保，他，与他们的不死不灭和悲惨相反，选择了暂时有死，且能永远保持有福。正是这样，他以自己死亡的谦卑和自己福乐的慈祥，摧毁了那些骄傲的不死不灭者和有害的可怜虫，并阻止了他们以不死不灭的夸口引诱人走向悲惨，这些人的心他藉着信德加以洁净，并由此将他们从这些恶者的不洁统治下解救出来。\par

人，既是有死且可怜的，又远离不死和真福者，应当选择何种中介，以便能与不死和真福相连？魔鬼的不死，对人或许有些吸引，却是可怜的；基督的有死，可能冒犯人，却已不复存在。在前者中，有对永恒痛苦的恐惧；在后者中，死亡既不可能永恒，便不再可畏，而永恒的福乐，则必须爱慕。因为那不死而可怜的中保置身中间，阻止我们走向有福的不死，因为那阻碍此路的，即痛苦，仍存留在他内；但那有死而有福的中保则置身中间，为的是经过有死后，能使有死者成为不死的（祂以自己的复活显示了这一能力），并使祂们脱离可怜，提升到真福的团体中，而祂自己从未离开过这团体。因此，有一位邪恶的中保，分离朋友；又有一位善的中保，使仇敌和好。那些施行分离者为数众多，因为真福的群众之所以有福，仅由于他们分沾唯一的天主；邪恶的天使被剥夺了这分沾，便沦为可怜，并置身中间，阻碍而非帮助人获得这真福，且以其众多数目阻止我们抵达那唯一的真福之善，为获得它，我们不需要许多中保，而只需要一位中保，即那非受造的天主\OriginalTerm{圣言}{Word}，万物是藉着祂造成的，我们因分沾祂而有福。我说祂是中保，并非因为祂是圣言，因为作为圣言，祂是至高的有福和至高的不死，因而远离可怜的凡人；但祂作为人，才是中保，因为祂借着祂的人性向我们显示：为获得那真福与真福之善，我们无需寻求其他中保，引导我们逐步攀登，而是那位真福而有福的天主，自己成了我们人性的分沾者，赐予我们容易的途径，去分沾祂的天主性。因为祂在解救我们脱离有死和可怜时，并不引我们到不死而有福的天使那里，使我们借着分沾他们的本性而成为不死和有福的，而是直接引我们到那天主圣三，天使们正是通过分沾天主圣三而获得真福。因此，当祂选择取了奴仆的形体，成为低于天使的，以便作我们的中保时，祂仍以天主的形体，超越天使之上——祂自己同时是地上的生命之道和天上的生命本身。\par

% structure: p0037 source=h2 target=subsection reason=relative-heading-rank; base=h2

\subsection{柏拉图主义者断定天上神祇拒绝接触属地事物并与世人交往，因此需要魔鬼代为转求，这是否合理。}

那位柏拉图主义者所宣称的、\Person{柏拉图}曾表示的观点是不真实的：\Quote{没有一位神祇与世人交往。}他说，这正是他们高贵的首要证据，即他们永不被与人的接触所玷污。因此，他承认魔鬼是被玷污的；从而推出，他们无法洁净那些玷污他们的人，于是众人都同成为不洁，魔鬼因与人交往，人因敬拜魔鬼。或者，如果他们声称魔鬼并不因与人交往和打交道而受玷污，那么魔鬼就比神祇更好了，因为神祇若这样做，便会被玷污。因为我们被告知，这便是神祇的荣耀，他们如此崇高，以致任何世人的交往都不能玷污他们。他确实断言，那至高者天主，万物的创造者，我们称之为真天主的，\Person{柏拉图}曾论及祂是唯一的天主，人类言语的贫乏甚至无法粗略地描述祂；而且，即便是智者，当他们的精神力量尽可能摆脱与肉身的纠葛时，对祂本性的洞察也仅如照亮黑暗的一道闪电。那么，如果这位至高者天主，真正超乎万物之上，却在智者脱离肉身时，以一种可理解的、不可言喻的临在造访他们的心灵，尽管这只是偶尔发生，如同瞬息闪电划过黑暗，其他神祇又何故如此高远，不与世人接触，仿佛会被玷污似的？仿佛只需举目仰望那些赐予大地必需光明的天体，便足以驳斥这种说法。如果星宿——按照他的说法，它们是可见的神祇——在我们注视它们时并未被玷污，那么魔鬼在人们相当接近看见他们时，也不会被玷污。但或许玷污神祇的是人声，而非目光；因此魔鬼被派作中介，将人的言语传达给神祇，而神祇因害怕玷污而远离？至于其他感官，我又能说什么呢？就嗅觉而言，无论是亲身在场的魔鬼，还是神祇——即使他们在场并吸入活人的气息——都不会被玷污，只要他们不被祭献中尸体的恶臭所污染。至于味觉，他们无需修复肉身的衰败，因而不必向人求食。而触觉则完全在他们自己的能力范围内。因为虽然接触这一名称似乎特别关涉触觉，然而神祇若有此意，尽可与人交融，看见也被看见，听见也被听见；哪里需要触碰呢？因为人如果蒙恩得见神祇或善的魔鬼，或与之交谈，便不敢奢求触碰；倘若出于过分好奇竟想触碰，又怎能不经神祇或魔鬼同意而遂其所愿？毕竟，人连一只麻雀也无法触碰，除非将它关进笼子。\par

那么，没有什么能阻止神明以形体与人交往，能看见也被看见，能说话也能听见。如果如我所说，魔鬼这样与人交往，并不受玷污，而神明若是这样做就会受玷污，那么魔鬼就比神明更不容易受玷污了。要是连魔鬼也会被玷污，他们又如何能帮助人获得死后的真福呢？这些中保自己已受玷污，非但不能洁净人、将他们洁净地呈献给未受玷污的神明，反而自身污秽，又怎能为人提供帮助？如果他们不能给人带来这种益处，他们友好的中介又有何用？难道其结果是，人不能进入神明之中，反倒是人和魔鬼共同处于玷污的状态，从而被排除在真福之外？也许有人会说，魔鬼像海绵之类的东西，在洁净朋友的过程中，自己会变得更加污秽，而别人则随之清洁。但若如此，那么因害怕玷污而避免与人接触或交往的神明，却与远为污秽的魔鬼交往。又或者，神明若不玷污自己就不能洁净人，却能不玷污自己而洁净被人类接触所玷污的魔鬼？谁会相信这样的蠢事，除非魔鬼用诡计欺骗了他？如果看见与被看见就是玷污，而\Person{阿普列乌斯}亲自称为可见的、\Quote{世界的璀璨光明}的神明以及其他星辰，都被人看见了，难道我们能相信魔鬼——除非它们愿意，否则人是看不见的——就更不容易受玷污吗？如果只有看见而非被看见才会玷污，那么他们就必须否认他们视为世界璀璨光明的这些神明，在将自己的光芒洒向大地时会看见人。他们的光芒照在各种污秽上并未受玷污，难道我们就该认为，如果神明与人交往，甚至如果需要接触才能帮助他们，就会受到玷污吗？因为大地与日月光辉之间有接触，而这并没有玷污光。\par

% structure: p0040 source=h2 target=subsection reason=relative-heading-rank; base=h2

\subsection{为获得真福生命——即分享至高之善——人所需要的中保不是魔鬼，而是唯独基督。}

我极为惊讶，那些宣称一切物质和可感事物都完全低于精神与理智事物的博学之士，竟在谈及真福生命时提及身体接触。难道他们忘记了\Person{普罗提诺}的那句名言？——\Quote{我们必须飞向所爱的故乡。那里是父，那里是我们的一切。什么船队或飞行能载我们到那里？我们的路，就是成为相似天主。}那么，如果一个人越相似天主，就越接近祂，那么与天主的距离，无非就是与祂的不相似。人的灵魂越是渴求短暂易变的事物，就越不相似于那无形、不变、永恒的实体。既然下面这些可朽不洁的事物，无法与上面不朽的纯洁交往，就确实需要一位中保来排除这一困难；但不是那种中保：它借着一具不朽的身体而类似最高层次的存在，却又因灵魂有病而类似最低层次，这只能使它嫉妒我们得医治，而非帮助我们获愈。我们需要一位中保，祂借着身体的必朽性与我们下面的人结合，同时又能借着祂精神的永恒正义，赐给我们真正属天的帮助，以洁净并释放我们；祂虽在世上，却始终是属天的。愿永不受玷污的天主绝对不怕因祂所取的人身，或因祂以人形生活在人们中间而受玷污。因为，即使祂的降生成人没有告诉我们别的，单单这两个有益的真理就足够了：真正的天主性不会因肉身受玷污，魔鬼并不因为它们没有肉身就应被视为比我们更优越。因此，正如圣经所说，这位就是\BibleQuote{在天主与人之间的中保，就是人耶稣基督，}\BibleRef{弟前 2:5}祂的天主性使祂与圣父同等，祂的人性使祂相似我们，至于这两者，此处不宜尽我所能详述。\par

% structure: p0042 source=h2 target=subsection reason=relative-heading-rank; base=h2

\subsection{欺诈的魔鬼虽承诺通过其代祷引领人归向天主，实则意图使他们偏离真理之路。}

至于魔鬼，这些虚假和欺诈的中保，虽然他们灵性的不洁常显露出其悲惨与邪恶，却借着其气的身体的轻盈和所居之地的性质，设法使我们偏离，阻碍我们的灵性进展；他们既不帮助我们接近天主，反倒阻拦我们抵达祂。因为即使是在那错误而误导人的物质道路上——义德不循此路——因为我们上升归向天主，不是靠身体的上升，而是靠与祂无形的或精神的相似——我说的是魔鬼的朋友们按各种元素的重量安排的这条物质道路：气的魔鬼被置于以太神明与属地的人之间，他们便设想神明享有此特权，即由这空间间隔而不受人类接触的玷污。这样，他们就以为魔鬼被人玷污，而非人被魔鬼洁净，且神明自身若非因其空间优越便也会受玷污。有谁会如此可悲，竟期望通过这样一条人玷污、魔鬼被玷污、神明也可受玷污的道路来获得净化呢？谁不愿选择那条我们借以逃脱魔鬼的玷污，并由永不受玷污的天主涤除污秽，从而与未受玷污的天使为伴的道路呢？\par

% structure: p0044 source=h2 target=subsection reason=relative-heading-rank; base=h2

\subsection{即便在其崇拜者中，\Quote{魔鬼}之名也从未有过好意义。}

然而，在这些我姑且称之为\OriginalTerm{拜魔者}{demonolators}的人中，连同\Person{拉贝奥}在内，有人声称，他们称为\OriginalTerm{魔鬼}{demons}的，却被其他人称为\OriginalTerm{天使}{angels}。如果我不想显得仅仅在争辩词句，就必须论及良善的天使。柏拉图学派也并不否认其存在，却宁愿称他们为好的\OriginalTerm{魔鬼}{demons}。但我们基督徒遵循\OriginalTerm{圣经}{Scripture}，从而知道有些天使是善的，有些是恶的，却从未在圣经中读到有好的魔鬼；无论何处出现此词或任何同类词语，总是指邪灵。这种惯用法已如此普遍，甚至在被称作异教徒、主张魔鬼当与神明一同受敬拜的人中间，也几乎没有一个人，无论怎样博学，敢于以称赞的口吻对他的奴隶说\Quote{你附了魔鬼}；也没有人会怀疑，若有人对他这样说，他定会视之为咒诅。这样，当我们使用\OriginalTerm{魔鬼}{demon}一词触犯人，而几乎每个人都将坏的意义与之相连时，我们为什么还要让自己非得费心去解释，而不干脆使用\OriginalTerm{天使}{angel}一词来轻松避免这种必要呢？\par

% structure: p0046 source=h2 target=subsection reason=relative-heading-rank; base=h2

\subsection{论那种使魔鬼自高自大的知识。}

然而，这名称的起源本身就暗示了值得注意的事，如果我们将其与\OriginalTerm{圣经}{divine books}加以比较。他们被称为\OriginalTerm{魔鬼}{demons}，来自一个希腊词，意思是知识。\OriginalTerm{宗徒}{apostle}在\OriginalTerm{圣神}{Holy Spirit}感动下说：\BibleQuote{知识使人自高自大，但爱德建立人。}\BibleRef{格前 8:1}这只能理解为，没有爱德，知识毫无益处，反而使人膨胀，以虚浮的风气自高。魔鬼们拥有知识却没有爱德，因此他们如此自高或骄傲，以至于贪图那些他们明知应归于真\OriginalTerm{天主}{God}的属神尊崇和宗教敬礼，并力所能及地向所有受他们影响的人索取这些。为了对抗魔鬼的这股骄傲——人类因之受其奴役，作为应得的惩罚——\OriginalTerm{天主}{God}以仆人的形象显现，施展了祂谦卑的莫大影响；但人却在骄傲上效法魔鬼，却不效法他们的知识，并为污秽所膨胀，未能认识祂。\par

% structure: p0048 source=h2 target=subsection reason=relative-heading-rank; base=h2

\subsection{主乐意在何种程度上让魔鬼认识自己。}

魔鬼们自己也清楚地认识天主的这次显现，以致他们对那带着肉身软弱的主说：\BibleQuote{纳匝肋人耶稣！我们与祢有什么相干？祢来提前毁灭我们吗？}\BibleRef{谷 1:24}从这些话显然可见，他们大有知识，却毫无爱德。他们惧怕祂惩罚的权能，却不爱慕祂的义德。祂按自己的意愿让他们认识，只让他们认识所必需的。但祂显示自己，并不是像向圣天使那样——圣天使认识祂是\OriginalTerm{圣言}{Word of God}，分享祂的永恒并因此而欢欣——而是为了恐吓那些魔鬼，因祂即将从他们的暴政下解救那些被预定进入祂王国和祂荣耀的人；这荣耀是永恒真实且真正永恒的。因此，祂不是借着那永生和不变之光——这光照耀虔诚的人，他们的灵魂因信祂而洁净——而是借祂权能的一些现世效果及祂神秘临在的迹象，使自己为魔鬼认识；这些迹象，恶灵的官能比人类的软弱更容易察觉。然而，当祂认为应逐渐抑制这些迹象并退入更深隐晦时，魔鬼的君主便怀疑祂是否是\OriginalTerm{基督}{Christ}，并企图通过试探祂来弄清此事；祂也容许自己被试探，为使人性成为我们效法的榜样。但那次试探之后，如\OriginalTerm{圣经}{Scripture}所载，有善良圣洁的天使来服事祂\BibleRef{玛 4:3}，这些天使令不洁之灵战栗；从此祂越来越清楚地显明祂的伟大，以致即使祂肉身的软弱看似可鄙，也无魔鬼敢抗拒祂的权威。\par

% structure: p0050 source=h2 target=subsection reason=relative-heading-rank; base=h2

\subsection{圣天使的知识与魔鬼的知识之间的区别。}

善的天使们因此轻视那一切关于物质的、暂时的事物知识——魔鬼们以此沾沾自喜——这并不是因为他们对此无知，而是因为\OriginalTerm{爱天主}{love of God}，那使他们圣化的爱，为他们极其宝贵；并且，相比于那不仅非物质，且永恒不变、不可言喻的美——他们被对它的圣爱所点燃——他们轻视一切在它之下的，以及一切不是它的，好使他们在自己内的一切美善中，享有那美善之源本身。因此，甚至对于那些暂时可变的事物，他们也有更为确切的知识，因为他们是在天主的\OriginalTerm{圣言}{Word of God}中，即在创造世界的圣言中，静观它们的理与因——那些因，凭此，一个事物得以确立，另一个被拒绝，而万物得以安排。但魔鬼们不在天主的\OriginalTerm{智慧}{Wisdom of God}中看到这些永恒的、仿佛枢纽性的暂时事物的原因，只是由于对隐藏于我们的征兆有更大熟悉，故比人能预见更大一部分将来。有时，他们预言的也是他们自己的意图。最后，魔鬼常被欺骗，而天使却永不会被欺骗。因为，借助暂时而可变的事物，推测时间中可能发生的变化，并靠自己的意志和能力改变这些事物——这在一定程度上被允许给魔鬼——这是一回事；而在永存于祂智慧中的永恒不变的天主法律中预见时代之变化，并借着分享祂的\OriginalTerm{圣神}{Spirit}而认识天主的意志，这最无误且最强大的万有之因，则是另一回事；而这是凭着公正的裁决赐给圣天使的。因此，他们不仅永恒，还是有福的。而他们所享有福乐的美善，即是创造他们的\OriginalTerm{天主}{God}。因为他们无止尽地享有对祂的观视与分沾。\par

% structure: p0052 source=h2 target=subsection reason=relative-heading-rank; base=h2

\subsection{第二十三章。——外邦人的诸神被妄称为神，然而\WorkTitle{圣经}也以此名称呼圣天使和义人。}

如果\OriginalTerm{柏拉图主义者}{Platonists}宁愿称这些天使为神而不称\OriginalTerm{魔鬼}{demons}，并且将他们与他们的创始者和宗师\Person{柏拉图}所主张由至高天主创造的不朽者同列，他们尽可如此，因为我不会在词语上徒费精力。因为他们若说这些存在是不朽的，但由至高天主所创造，其真福在于依附其创造者而非凭自己的力量，那么无论他们叫这些存在何名，他们所说的就是我们所说的。而这一点是所有或最优秀的柏拉图主义者的意见，可从他们的著作察知。至于名称本身，若他们认为合适，称这样真福而不朽的受造物为神，这也不至于在我们之间引起严重的争论，因为在我们自己的\WorkTitle{圣经}中我们读到：\BibleQuote{主，众神之主，说了话}；又说：\BibleQuote{请众感谢众神之神}；又说：\BibleQuote{他是超越众神的伟大君王}。论到所谓\BibleQuote{他超越众神，可敬可畏}，便立即加上理由，因为接着说：\BibleQuote{因为外邦人的众神全是虚无，主却创造了诸天}。他说\BibleQuote{超越众神}，却加上\BibleQuote{外邦人的}，也就是说，超越一切外邦人所视为神者，换言之，即魔鬼。他应受他们以那种畏惧而敬畏，就是他们向主喊说：\BibleQuote{你竟来毁灭我们吗？}但论到\BibleQuote{众神之主}，却不可理解为魔鬼之神；我们断不敢说，\BibleQuote{超越众神的伟大君王}意谓\BibleQuote{超越一切魔鬼的伟大君王}。但同一\WorkTitle{圣经}也称属于天主子民的人为\BibleQuote{神}：\BibleQuote{我曾说过：你们是神，都是至高者的子民。}因此，当天主被称为众神之主时，这应理解为这些神；同样，当他被称为超越众神的伟大君王时，亦应如此理解。\par

然而，有人或许会说，若人因属于天主的子民而被称为神，且天主借着人和天使向他们发言，那么那些不朽者，既已享有人们通过敬拜天主所寻求达至的真福，岂不更配得这名称吗？对此我们有何回答呢？无非是：在\WorkTitle{圣经}中，人比我们希望在复活时与之同等的不朽真福之灵更明白地被称为神，并非没有理由，因为恐怕不信的软弱，被这些存在的卓越所压倒，竟擅自将其中某个立为神。对人的情形而言，却无需防范此事。再者，天主子民中人当更明白地被称为神，这是正当的，为要使他们确信并证实，那位被称为众神之主的天主是他们的天主；因为，虽然那些居住在天上的不朽真福之灵被称为神，他们却未被称作众神之神，也就是说，未被称为组成天主子民之人——即那些受有\BibleQuote{我曾说过：你们是神，都是至高者的子民}之人——的神。因此宗徒说过：\BibleQuote{因为虽然有称为神的，或在天上，或在地上，就如那许多\InnerQuote{神}和许多\InnerQuote{主}，可是为我们只有一个天主，就是圣父，万物都出于他，而我们也归于他；也只有一个主，就是耶稣基督，万物借他而有，我们也借他而有。} \BibleRef{格前 8:5}\par

因此，我们无需费力争执名称，因为事实如此明显，不容丝毫怀疑。我们所说，受派遣向人宣告天主旨意的天使属于真福不朽者的等级，并不使柏拉图主义者满意，因为他们相信，这种职务并非由他们称为神者——换言之，并非由真福的不朽者——执行，而是由魔鬼执行，他们不敢肯定魔鬼是真福的，只肯定其不朽；或者，即便他们将某些魔鬼列入真福不朽者之中，也只作为善的魔鬼，而非作为居住在高天之上、远离一切人类接触的神。然而，尽管这似乎只是名词上的争辩，但魔鬼之名如此可憎，以至于我们无论如何不能忍受将其用于圣天使。因此，现在让我们在确信中结束本书：无论我们如何称呼这些不朽而真福的灵体——他们仍只是受造物——他们并不充当中保，将悲惨的凡人引入永恒的真福，因为他们与凡人之间隔着一道双重的鸿沟。至于那些充当中保的其他灵体，就他们与在上者同享不朽、与在下者同受悲惨（因他们的邪恶受罚而理应悲惨）而言，不能赐予我们真福，反而嫉妒我们拥有他们自身所被排除在外的真福。因此，魔鬼的朋友们提不出任何重大理由，说明我们为何应当将这些魔鬼当作援助者来敬拜，而不应当将其视为我们利益的叛徒而远避。至于那些善的灵体，因而是不仅不朽且真福的，他们以为我们应加之以神的称号，并为了承受来生而奉献敬拜和祭献，我们将依靠天主的助佑，在下一书中尽力表明：这些灵体——无论你以何名称呼他们，赋予他们何种本性——都切愿惟独对创造他们的天主，并因天主自我通传于他们而获得真福的那一位，奉行宗教敬拜。\par

\SourceRecord{Translated by Marcus Dods.；Nicene and Post-Nicene Fathers, First Series；Vol. 2.；Edited by Philip Schaff.；Buffalo, NY: Christian Literature Publishing Co.,；1887.；<http://www.newadvent.org/fathers/120109.htm>.}

% structure: p0001 source=h1 target=section reason=page-title

\section{\WorkTitle{天主之城}（卷十）}

\Person{奥斯定}在本卷中教导说，善的天使愿意只有天主——他们自己所侍奉的那位——接受那通过祭献而呈上的神圣尊荣，即所谓的\OriginalTerm{钦崇}{\GreekText{λατρεία}}。他接着反驳\Person{波斐留斯}，讨论灵魂的洁净与得救的原则和途径。\par

% structure: p0003 source=h2 target=subsection reason=relative-heading-rank; base=h2

\subsection{第一章 柏拉图主义者自己认定，唯天主能赐予天使或人幸福，但仍需探究：究竟那些他们教导我们为了获得幸福而应敬拜的神灵，愿意我们将祭献献给他们自己，还是唯独献给唯一的天主。}

凡运用理智的人，无不认定所有人都渴望幸福。但谁是幸福的，或如何成为幸福，这些问题因人类理解力的软弱引发了无休止而激烈的争论，哲学家们为此徒耗精力与闲暇。要引述并讨论他们各色各样的观点，既冗长又无必要。读者或许记得我们在第八卷中，挑选了一些哲学家，以便讨论有关来世幸福生活的问题：我们能否通过向唯一真天主、万神的创造者献上神圣敬礼，或通过敬拜多神而达到幸福？因此，读者不会期望我们在此重述同样的论证，尤其因为即便他已忘记，重读一遍便可恢复记忆。我们之所以挑选柏拉图主义者——他们公认为哲学家中最卓越的——是因为他们有智慧看出，人的灵魂，虽为不朽而具理性或理智，除非分享那造生灵魂与世界之天主的光，便不能幸福；也看出众人所渴望的幸福人生，除非以纯洁、神圣的爱依附于那唯一的至善、不变的天主，否则无法达到。然而，就连这些哲学家，无论是迁就民众的愚昧无知，还是如宗徒所说：\BibleQuote{他们冥顽不灵，心思荒谬}\BibleRef{罗 1:21}，却也认为或容人以为应敬拜多神，以致他们中间有些人以为，就连魔鬼也应通过敬礼和祭献而受到神圣尊崇（这一错误我先前已驳斥）。现在，我们必须依靠天主的助佑，探明那些居住在天界，属于统领、宰制、异能等品级的不朽而蒙福的神灵——柏拉图主义者称之为神祇，有些则称为善的精灵，或像我们一样称为天使——他们对我们宗教敬礼和虔敬的看法如何；也就是说，更明白地讲，这些天使究竟愿意我们将祭献和敬拜，以及我们自身和财物的奉献，献给他们自己，还是只献给天主——他们和我们的天主。\par

因为这是应归于\OriginalTerm{神性}{Divinity}的敬拜，或更准确地说，应归于\OriginalTerm{天主性}{Deity}；而要表达这种敬拜，由于我找不到任何足够精确的拉丁词，必要时我将采用希腊词\OriginalTerm{钦崇}{\GreekText{λατρεία}}。在圣经中，这个词被译为\Quote{服务}。但那应归于人的服务，即宗徒所写奴仆应服从自己主人\BibleRef{弗 6:5}所指的那种服务，在希腊文中通常用另一个词表示；至于通过敬拜唯独献给天主的服务，在那些由神圣神谕写作的人笔下，总是或几乎总是称为\OriginalTerm{钦崇}{\GreekText{λατρεία}}。这不能简单地称为\Quote{\OriginalTerm{敬礼}{cultus}}，因为那样就显得并非天主所独享；因为同一个词也被用于我们对逝者或生者所怀的尊敬。由此也衍生出\Quote{农业}、\Quote{殖民者}等词。外教人称他们的神为\Quote{\OriginalTerm{天堂居民}{cœlicolæ}}，并非因为他们敬拜天，而是因为他们住在其中，宛如殖民——不是指我们所说的那些依附本土、在主人支配下耕种土地的殖民者，而是如那位拉丁语言大师所说：\Quote{有一座古城居住着泰尔移民}。他称他们为移民，并非因为他们耕种土地，而是因为他们居住在那城中。同样，从大城分出去的城市也称为殖民地。因此，虽然事实上，若在特殊意义上使用该词，\Quote{敬礼}只能献给天主，然而，由于该词也被用于其他事物，对天主专有的敬礼在拉丁语中无法仅用这个词来表达。\par

\Quote{\OriginalTerm{宗教}{religion}}一词似乎更明确地表达只应归于天主的崇拜，因此拉丁译者用这个词来代表\OriginalTerm{宗教}{\GreekText{θρησκεία}}；然而，由于不仅未受教育者，就连最有学识的人也使用宗教一词来表达人间的纽带、关系和亲缘，若在讨论对天主的崇拜时使用这个词，就不可避免地会引入歧义，因为我们无法说宗教无非就是对天主的崇拜，而不至于违背将此词用于表示社会关系遵守的通常用法。同样，\Quote{\OriginalTerm{虔敬}{piety}}，或如希腊人所说，\OriginalTerm{虔敬}{\GreekText{εὐσέβεια}}，通常被理解为对天主崇拜的恰当名称。然而这个词也被用于对父母的孝顺。普通民众也将其用于爱德工作，我想，这是因为天主命令人履行这些工作，并宣称祂喜悦这些工作，以代替或优先于祭献。由于这种用法，天主本身甚至也被称为虔敬的，而希腊人却从不在这个意义上使用\OriginalTerm{虔敬的}{\GreekText{εὐσεβεῖν}}，尽管他们的普通民众也将\OriginalTerm{虔敬}{\GreekText{εὐσέβεια}}用于爱德工作。因此，在圣经的某些章节中，他们试图通过不使用更一般的词\OriginalTerm{虔敬}{\GreekText{εὐσέβεια}}，而使用\OriginalTerm{对天主的崇拜}{\GreekText{θεοσέβεια}}——其字面意思就是崇拜天主——来保持区别。而另一方面，我们却无法用一个词来表达这两种观念中的任何一种。那么，这在希腊语中称为\OriginalTerm{敬礼}{\GreekText{λατρεία}}，在拉丁语中称为\Quote{\OriginalTerm{服务}{servitus}}，但却是只应归于天主的侍奉；这在希腊语中称为\OriginalTerm{宗教}{\GreekText{θρησκεία}}，在拉丁语中称为\Quote{\OriginalTerm{宗教}{religio}}，但却是仅仅使我们与天主相连的宗教；他们称之为\OriginalTerm{对天主的崇拜}{\GreekText{θεοσέβεια}}，而我们无法用一个词来表达，只能称之为对天主的崇拜——我们说，这崇拜只属于那位真天主，祂使敬拜祂的人成为神。因此，无论这些天上不朽而蒙福的居民是谁，如果他们不爱我们，不希望我们得享真福，我们就不应崇拜他们；如果他们爱我们并愿我们幸福，他们不可能希望我们经由与他们自身所享受的不同的方式而获得幸福——因为他们怎能希望我们的真福出自一个源头，而他们的出自另一个源头呢？\par

% structure: p0007 source=h2 target=subsection reason=relative-heading-rank; base=h2

\subsection{第二章 柏拉图主义者普罗提诺关于自上而来的光照的观点}

然而，在这些更值得尊敬的哲学家面前，我们于此问题上并无争议。因为他们认识到，并以各种形式在他们的著作中充分表达，这些神灵与我们有着同一的幸福源泉——某种可理知的光，这光是他们的天主，与他们自身有别，并光照他们，使他们被光渗透，并通过分享天主而享有完美的幸福。普罗提诺在评论柏拉图时，反复而有力地断言，即便是他们所相信的世界灵魂，其真福也并非源自与我们不同的其他源泉，而是源自那与它不同并创造了它的光，藉着这光的可理知照耀，它在可理知的事物中享有光明。他还将那些属灵之物比作广阔而显著的天空中的天体，仿佛天主是太阳，灵魂是月亮；因为他们认为月亮从太阳获得光明。因此，那位伟大的柏拉图主义者说，理性的灵魂，或更确切地说，理智的灵魂——他将居于天上的蒙福不朽者的灵魂也归入此类——除了天主，世界的创造者和灵魂本身，没有凌驾于它之上的本性，并且这些天上的神灵所享有的真福生命和真理之光，其源泉与我们相同。这与福音书中我们读到的话一致：\BibleQuote{曾有一人，由天主派遣来的，名叫若翰。这人来，是为作证，为给光作证，为使众人藉他而信。他不是那光，只是为给那光作证。那普照每人的真光，正在进入这世界。}\BibleRef{若 1:6} 这个区分充分证明，像若翰那样的理性或理智的灵魂，不能自身成为光，而需要从另一个、那真光那里接受光照。这位若翰在作证时自己承认说：\Quote{从他的满盈中，我们都领受了。}\par

% structure: p0009 source=h2 target=subsection reason=relative-heading-rank; base=h2

\subsection{第三章 柏拉图主义者虽对宇宙的创造者有所认识，却因敬拜天使——无论是善的或恶的——而误解了对天主的真正崇拜}

事情既是这样，如果柏拉图派，或那些与他们思想一致的人，认识天主，却光荣祂为天主并感谢祂，如果他们并未陷入虚妄的思想，也不制造或随从流行的谬误，他们就必定会承认：那有福的不朽者，若不钦崇那既是他们的、也是我们的唯一天主，万神之神，便不能保持幸福；而我们可悲的凡人，若不如此，也无法达到幸福的境界。对祂，我们应当奉献希腊语称为\OriginalTerm{钦崇}{\GreekText{λατρεία}}的敬礼，无论是外在的，还是内在的；因为我们都是祂的圣殿，无论是我们每个人，还是我们全体，因为祂屈尊就卑，居住于每一个体内，也居住于整个和谐的身体内，祂在整个身体内并不比在每一个体内更大，因为祂既不扩张，也不分割。我们的心向祂高举时，就是祂的祭台；那位为我们代祷的司祭，是祂的独生子；我们为祂的真理奋斗，直至流血，便是向祂祭献流血的牺牲；当我们满怀圣善热诚的爱，来到祂面前时，便是向祂奉献最甘饴的馨香；我们将自己以及祂在我们内的恩赐，全然奉献并交托给祂；我们在隆重的庆节和指定的日子里，向祂祝圣我们对祂恩惠的纪念，免得因岁月流逝而让忘恩的遗忘潜入；我们在心之祭台上，献上由灼热爱火点燃的谦卑和赞颂的祭献。我们得以看见祂，按祂能被人看见的程度；我们得以依附祂，正是为了使我们从一切罪恶和偏情的污秽中洗净，并因祂的名而被祝圣。因为祂是我们幸福的泉源，祂是我们一切渴望的终向。依附着祂，或者更确切地说，再依附于祂——因为我们曾脱离并失落了祂——我说，再依附于祂，我们借着爱趋向于祂，为能在祂内安息，并藉达到这目的而寻得我们的真福。因为我们的善，哲学家们曾为之激烈争论，无非是与天主结合。可以说，正是借着心灵上拥抱祂，理性灵魂才被真正的德行所充满并孕育。我们受命要以全心、全灵、全力爱慕这善。这善，我们应当被那些爱我们的人引向这善，也应当引领我们所爱的人向这善。如此便满全了这两条诫命，全部法律和先知书都系于其上：\BibleQuote{你应当全心、全灵、全意，爱你的主，你的天主；}\BibleQuote{你应当爱近人如你自己。}\BibleRef{玛 22:37}为使人在自爱中具有智慧，就为他指定了一个目的，让他将自己的一切行为都指向这目的，好使他获得幸福。因为凡爱自己的人，所希求的无非是此。而摆在他面前的这目的便是\Quote{亲近天主}。因此，当一位怀有这智慧的自爱的人，受命去爱近人如自己时，所命令的无非是他要尽力劝勉近人去爱天主而已。这就是对天主的钦崇，这就是真宗教，这就是纯正的虔敬，这便是只应归于天主的敬礼。因此，任何不朽的异能者，无论具有何等德行，若爱我们如同爱自己，就必定希望我们藉顺从那一位——他本人也是因顺从而获得幸福——来寻获我们的幸福。他若不朝拜天主，便是可悲的，因为离开了天主；他若朝拜天主，便不可能愿意自己取代天主而受人朝拜。反之，这些大能者由衷赞许天主的谕令，经上记载说：\BibleQuote{那向任何神祭献，而不唯独向主祭献的，必被全然灭绝。}\BibleRef{出 22:20}\par

% structure: p0011 source=h2 target=subsection reason=relative-heading-rank; base=h2

\subsection{祭献只应归于真天主。}

然而，暂且放下其它敬拜天主的宗教仪式不谈，肯定没有人敢说祭献可归于除天主外的任何一位。确实，神圣敬礼中的许多部分，常被滥用来向人表示尊敬，或是出于过分的谦卑，或是出于有害的谄媚；然而，即便如此，那些受人如此敬拜、尊敬，甚至朝拜的人，仍被视为不过是人；有谁曾想过要献祭，除非是对那位他知道、设想或假想为神的呢？至于祭献是何等古老的天主敬礼部分，那两兄弟——\Person{加音}与\Person{亚伯尔}——已充分显示，天主拒绝了长子的祭献，而悦纳了幼子的祭献。\par

% structure: p0013 source=h2 target=subsection reason=relative-heading-rank; base=h2

\subsection{天主所不需要的祭献，但为表明祂所需要的事而愿人遵行的祭献。}

有谁会如此愚昧，竟以为献给天主的祭品是祂所需，是为祂自己的某种用途呢？圣经在许多地方都驳斥了这种想法。为避免冗长，仅引圣咏中的一句简短的话便足矣：\BibleQuote{我对主说：祢是我的天主，因为祢不需要我的好处。}因此我们必须相信，天主不仅不需要牛羊或任何其它世上的物质之物，甚至也不需要人的正义；并且，无论向天主献上怎样正当的敬礼，受益的不是祂，而是人。因为没有人会说自己因饮水而给了水泉什么好处，或因看见而给了光明什么好处。古代教会献上动物祭品这一事实，今天的天主子民读到却不仿效，这无非证明了：那些祭品象征着我们为了亲近天主并激励近人同样行事而做的事。因此，祭献，是不可见祭献的有形圣事或神圣标记。因此，圣咏中的那位忏悔者，或可能就是圣咏作者本人，恳求天主怜悯他的罪过，说道：\BibleQuote{祢若喜爱祭献，我必向祢献上；祢并不喜悦全燔祭。天主的祭献就是痛悔的心；一颗痛悔和谦卑的心，天主绝不轻看。}请注意，他正是在表达天主拒绝祭献的话语中，表明天主确实要求祭献。祂不渴望被宰杀野兽的祭献，而是渴望一颗痛悔之心的祭献。因此，他说天主所不喜悦的祭献，乃是天主所喜悦之祭献的象征。天主并不是按照愚昧人所想的那样要求这些祭献，即以它们来取悦自己。因为，若祂不愿让祂所要求的祭献——\Emph{例如}，因忏悔的忧伤而痛悔谦卑的心——由那些祂曾被以为因悦己而喜悦的祭献来象征，那么旧约的法律就绝不会命人献上那些祭献；而这些祭献注定要在适当的时机被取代，以免人们认为那些祭献本身，而非它们所象征的事物，才是天主所喜悦的或我们所呈上的蒙悦纳之物。因此，在另一首圣咏的另一段中，他说：\BibleQuote{若我饥饿，我必不告诉你；因为世界和其中所充满的都是我的。难道我吃牛犊的肉，或喝山羊的血吗？}仿佛天主在说：就算这些是我所需的，我也绝不会向你索求我手中已有的东西。然后他接着提到这些象征的意义：\BibleQuote{你们要向天主奉献赞颂之祭，向至高者还你们的誓愿。你在患难之日呼求我，我必拯救你，你也要光荣我。}同样，另一先知书中说：\BibleQuote{我应带着什么到主面前，叩拜至高者天主呢？我应带着全燔祭，带着一岁的牛犊到祂面前吗？主岂喜悦万千的公羊，或万道溪流的油吗？为了我的过犯，我应献上我的长子，为我灵魂的罪献上我亲生的儿子吗？人啊，已指示你何为善，主要求你的是什么：无非就是履行正义，爱好仁慈，谦卑地与你的天主同行。}\BibleRef{米 6:6}在这位先知的话语中，这两件事被清晰地区分和阐明：天主为祭献本身并不要求这些祭献，祂所要求的是这些祭献所象征的祭献。在\BibleBook{希伯来}{书}中说道：\BibleQuote{行善和奉献，不可忘记；因为这样的祭献是天主所喜悦的。}\BibleRef{希 13:16}因此，当经上写道：\BibleQuote{我喜欢仁爱胜过祭献，}\BibleRef{欧 6:6}这无非意味着一种祭献优于另一种祭献；因为通常所说的祭献只是真正祭献的象征。然而，仁爱乃是真正的祭献，因此正如我刚才所引的，经上说：\BibleQuote{这样的祭献是天主所喜悦的。}因此，我们在关于会幕或圣殿礼仪中祭献的事上读到的一切神圣规定，都应指向对天主和对近人的爱。因为经上写道：\BibleQuote{全部法律和先知书，都系于这两条诫命。}\BibleRef{玛 22:40}\par

% structure: p0015 source=h2 target=subsection reason=relative-heading-rank; base=h2

\subsection{六、真而全的祭献}

因此，真正的祭献是指任何使我们得以在圣洁的共融中与天主结合，并指向那使我们能真正蒙福的至高的善和终向的工作。所以，即使是我们向人所显示的仁慈，若不是为天主的缘故，也不是祭献。因为，虽然祭献是由人奉献的，但祭献本身是神圣的事物，正如那些称之为\OriginalTerm{祭献}{sacrifice}的人所要指明的。同样，人自己，因天主的名而被祝圣，并奉献于天主，就成为一种祭献，只要他死于世俗而活于天主。因为这是人对自己所施仁慈的一部分；正如经上记载：\BibleQuote{你要怜悯你的灵魂，悦乐天主。}\BibleRef{德 30:24} 我们的身体，当我们以节制磨炼它时，若我们为天主的缘故，按所应当的那样去做，以免将我们的肢体交给罪恶，作为不义的利器，反而作为正义的利器献给天主，也是一样祭献。\BibleRef{罗 6:13} 为劝勉我们作这祭献，宗徒说：\BibleQuote{所以弟兄们！我以天主的仁慈请求你们，献上你们的身体当作生活、圣洁和悦乐天主的祭品：这才是你们合理的敬礼。}\BibleRef{罗 12:1} 那么，如果身体，这灵魂用以作为仆人或工具的较低之物，当它被正当使用并指向天主时，便是一个祭献，那么当灵魂将自己献于天主，为被祂的爱火所燃烧，领受祂的美善，成为祂所悦纳的，舍弃世俗欲望的形状，并按永恒可爱的肖像重新塑造时，灵魂本身岂不更成为祭献吗？宗徒接着又说：\BibleQuote{你们不可与此世同化，反而应以更新的心思变化自己，为使你们能辨别什么是天主的旨意，什么是善事，什么是悦乐天主的事，什么是成全的事。}\BibleRef{罗 12:2} 既然，真正的祭献是指为自己或别人所作的、指向天主的仁慈工作，而仁慈的工作无非是为解救困苦或赐予幸福；并且，离了那经上所说的\BibleQuote{亲近天主乃是我的美好}的美善，就没有真正的幸福；因此，整个被救赎的城邑，即圣人的会众或团体，就借着那位大司祭被献于天主，作为我们的祭献；祂曾以奴仆的形体，在祂的苦难中为我们将自己献于天主，好使我们成为这光荣之首的肢体。因为祂所献上的正是这个形体，祂也是在这个形体内被献上，因为依此形体，祂是\OriginalTerm{中保}{Mediator}、是我们的司祭、是祭品。因此，宗徒在劝勉我们献上身体作为生活、圣洁和悦乐天主的祭品，作为我们合理的敬礼，不要与此世同化，却要以更新的心思变化自己，好能辨别那善的、悦乐天主的、和成全的天主旨意，即我们自己的真正祭献之后，说：\BibleQuote{我因所赐给我的恩宠，告诉你们每一位：不可把自己估计得太高，而过了份；但应按照天主所分与各人的信德尺度，估计得适中。就如我们在一个身体上有许多肢体，但每个肢体，都有不同的作用；同样，我们众人在基督内，也都是一个身体，彼此之间，每个都是肢体。按我们各人所受的恩宠，各有不同的恩赐。}\BibleRef{罗 12:3} 这就是基督徒的祭献：我们虽多，但在基督内却是一个身体。这也是教会在\OriginalTerm{祭台上的圣事}{sacrament of the altar}中不断举行的祭献，这圣事是信友们所熟知的；在这圣事中，教会教导说，她自己在向天主呈上的奉献中被奉献。\par

% structure: p0017 source=h2 target=subsection reason=relative-heading-rank; base=h2

\subsection{第7章 论圣天使的爱——这爱促使他们愿意我们只敬拜独一真天主，而非他们自己}

这些有福且不死的神体，居住在\Place{天上的居所}，享受他们造物主丰满的相通，坚定于祂的永恒，安固于祂的真理，因祂的恩宠而成圣，他们既然满怀慈悲与温柔地垂顾我们这些可悲的凡人，愿意我们成为不死而幸福的，就不愿我们向他们自己献祭，而只愿我们向那位献祭——他们深知自己同我们一样，是献给祂的祭品。因为我们和他们一同组成那座唯一的天主之城，圣咏中论及这城说：\BibleQuote{天主之城啊，人们论到你，曾说了许多光荣的事。} 此城的人类部分在下土旅居，天使部分则从上界援助。因为从这座天上的城——在那里，天主的旨意就是可理解的、永不改变的法律；从这天上议会——因为他们正议论有关我们的事——有那\WorkTitle{圣经}，借着众天使的职分，降到我们中间，其上记载：\BibleQuote{祭其他神祇，而不只祭上主的，应被完全毁灭。}\BibleRef{出 22:20} 这部\WorkTitle{圣经}，这法律，这些规诫，已由如此众多的\Emph{奇迹}所证实，足以清楚显明：这些愿意我们像他们一样的不死且幸福的神体，究竟要我们向谁献祭。\par

% structure: p0019 source=h2 target=subsection reason=relative-heading-rank; base=h2

\subsection{第8章 论天主借天使的职分所施行的奇迹，以证实祂的应许，坚定虔诚信者的信德}

若要一一细数所有古代奇迹，未免显得冗长，这些奇迹是为了证明天主的许诺，即祂在数千年前向亚巴郎所作的许诺——地上万民都要因他的后裔蒙受祝福。\BibleRef{创 18:18} 因为谁能不惊奇，亚巴郎那不能生育的妻子，竟在连多产妇女也无法生育的年纪生了一个儿子？或者，当亚巴郎祭献时，有火焰从天降下，从那些剖开的肉块中间经过？又或者，他殷勤款待的那些人形的天使，既重申了天主关于后裔的许诺，又预言索多玛将被从天降下的火焚毁？\EditorialNote{创 18} 还有他的侄子罗特，在火正要降下时，被天使从索多玛救出，而他的妻子在逃走时回头一看，立刻变成了一根盐柱，如同一座神圣的警示，告诫我们，凡蒙救者，不应留恋他所离开的东西？何等惊人还有梅瑟所行的奇事，为解救天主的子民脱离埃及的奴役，当时压迫那子民的法郎，即埃及王，其术士虽蒙允许行了一些奇事，却为使他们更彻底地败北！他们凭借恶神或恶魔所擅长的法术和咒语行了这些事；然而梅瑟，既有正义在握，力量便远胜于他们，又得天使之助，便奉创造天地的主之名，轻易胜过了他们。事实上，那些术士在第三灾时便无能为力了；而梅瑟则藉着委托给他的奇迹，给那地方降下了十灾，直到法郎和埃及人的硬心终于屈服，放走了天主的子民。但他们很快又后悔，想要追赶离去的希伯来人，希伯来人已走过干涸的海床，而追赶者却被回流的海水淹没吞噬。我还要怎样述说在旷野中引领子民时，天主大能频繁而惊人的彰显呢？——比如那不能饮用的苦水，因天主的命令投入一根木头便失去了苦味，解了人渴？又如那自天降下解除饥饿的\OriginalTerm{玛纳}{manna}，若有人多收过定量便生虫发臭，但在安息日前一天，虽可收双倍（安息日不可收取），却仍保持新鲜？又如他们渴望肉食，似乎无法供应如此众多人口时，飞鸟布满营地，使食欲变为饱饫？又如何有敌人迎面阻拦，以兵器挡住去路，却被击败，希伯来人却无一伤亡，当时梅瑟伸开双手形成\OriginalTerm{十字形}{form of a cross}祈祷？又如何在天主子民中兴起叛乱之人，脱离天主所立的团体，被大地活活吞没，作为无形惩罚的有形象征？又如何击打磐石出水，足够全营饮用？又如何因罪孽的正义惩罚，民受致命火蛇咬伤，但一仰望那被高举的铜蛇便得痊愈，如此不仅受苦的百姓得愈，更以这死亡消灭死亡的方式，将十字架死亡的预象置于他们面前？正是这条铜蛇为纪念此事而被保存，后来却被误入歧途的百姓当作偶像朝拜，最后被虔敬敬畏天主的君王\Person{希则克雅}毁掉，这是他极大的功德。\par

% structure: p0021 source=h2 target=subsection reason=relative-heading-rank; base=h2

\subsection{第9章 论与恶魔崇拜有关的非法技艺，柏拉图主义者\Person{波菲利}采纳一些而摒弃另一些。}

这些奇迹以及许多同类奇迹，一一细述未免冗长，它们之施行，乃是为了将崇拜唯一真天主的信仰推荐于人，并禁止崇拜众多假神。此外，它们乃是凭借单纯的信仰和对天主的虔信而行的，并非靠着干涉不可见世界所获得的罪恶法术和符咒；那种技艺，有人称之为\OriginalTerm{魔法}{magic}，或用更可憎的名字称为\OriginalTerm{招魂术}{necromancy}，或用较体面的名称称为\OriginalTerm{通神术}{theurgy}。因为他们试图区分那些被百姓称为魔法的术士，这些人行招魂之术，沉溺于非法技艺而遭谴责，与那些因施行通神术而似乎值得称赞的人，然而真相是，这两类人都是那些他们奉天使之名所呼求之恶魔的骗人仪式的奴隶。\par

因为连 \Person{波斐留斯} 也答应，借着 \OriginalTerm{通神术}{theurgy} 的帮助，灵魂可以得到某种净化，尽管他说时有些犹豫和羞愧，并且否认这种法术能使任何人回到天主那里；因此你可以看出他的意见在哲学主张和一种他觉得狂妄亵圣的法术之间摇摆不定。因为他一方面警告我们要避免它，视之为骗人的、法律禁止的、对行此法者危险的；然后，又像是为了顺应它的鼓吹者，他宣称它对于洁净灵魂的某一个部分是有用的，确实，不是理智的部分，理智部分是用来认识没有可感形象的可理知之物的真理，而是 \OriginalTerm{精神部分}{spiritual part}，它是认知物质事物影像的。他说，这精神部分，借着某些 \OriginalTerm{通神祝圣}{theurgic consecrations}，或如他们所称呼的，\OriginalTerm{密迹}{mysteries} 的帮助，得以准备并适宜于与神灵和天使交往，并得见众神。然而，他承认，这些通神秘迹不能将那种使理智之魂能看见自己的天主并认知真实存在之物的纯净性赋予理智之魂。从这一承认，我们可以推断这些是何种的神，以及通神祝圣所赋予的看见他们的那种景象如何，如果凭借它不能看见真实存在之物的话。他进一步说，理智之魂，或如他偏爱称呼的，理性之魂，即使精神部分未经通神术的洁净，也能进入诸天，而且这种法术不能足够净化精神部分，使之进入不死与永恒。因此，尽管他把天使与魔鬼区分开来，声音后者的居所在空气中，而前者住于 \OriginalTerm{以太}{ether} 和 \OriginalTerm{九天}{empyrean}，尽管他建议我们培植某些魔鬼的友谊，使它们能在我们死后帮助我们，并至少将我们稍微提离地面——因为他承认我们必须经由另一条路才能达到天使们的天上团体——但同时他又明确警告我们要避免魔鬼的团体，说，死后赔补其罪的灵魂，痛恨那曾缠绕它的魔鬼的崇拜。而就通神术本身而言，尽管他推荐它是使天使和魔鬼和解的，他却不能否认，它涉及的那些神灵，或者自己嫉妒灵魂的洁净，或者为嫉妒灵魂洁净者的伎俩服务。他通过某个加色丁人的口抱怨这件事：他说，\Quote{在 \Place{加色丁} 的一个善人抱怨，他尽最大努力净化自己灵魂的努力被挫败，因为另一个人在这些事上有势力，嫉妒他的洁净，向那些神灵祈祷，用咒语束缚它们，不让它们听从他的请求。因此，} \Person{波斐留斯} 补充说，\Quote{一人所束缚的，另一人不能解开。} 由此他下结论说，通神术是一种不仅在神人之间成就善事也成就恶事的诡计；而且众神也有情欲，也会被 \Person{阿普列乌斯} 归于魔鬼和人的情感所搅扰和动荡不安，但他凭借与 \Person{柏拉图} 共认的居处的崇高，保护了众神免于这些。\par

% structure: p0024 source=h2 target=subsection reason=relative-heading-rank; base=h2

\subsection{第十章 论那应许靠呼求魔鬼而给灵魂带来虚幻净化的\OriginalTerm{通神术}{theurgy}}

可是在这里我们有另一位、也比 \Person{阿普列乌斯} 博学得多的柏拉图主义者，即 \Person{波斐留斯}，他断言，靠着某种通神术，连众神自身也会遭受情欲和动荡；因为借着咒语，他们被束缚惊吓，以致不能赋予人灵魂的洁净——他们是如此被那发号施令行邪恶的人所惊吓，以至于不能借同样的通神术摆脱那种恐惧，去实现向他们祈祷之人的正义命令，或者行他所求的善事。谁看不出这一切都是欺骗人的魔鬼的虚构呢，除非他是它们可怜的奴才，隔绝于真正解放者的恩宠？因为如果那位 \Place{加色丁} 人所交涉的是良善之神，那么一个寻求洁净自己灵魂的心怀善意的人，当然会比一个试图阻止他的心怀恶意的人对神更有影响力。或者，如果众神是公义的，认为此人不配得他所寻求的净化，那么无论如何他们也不该被嫉妒者所惊吓，也不该像 \Person{波斐留斯} 所肯认的那样，因一位更强神祇的恐惧而受阻，而应该凭自己自由的判断单纯地拒绝施恩。令人惊讶的是，那位心善的 \Place{加色丁} 人，他想用通神仪式来净化自己的灵魂，竟然找不到一位更高的神祇，这位神祇能够进一步吓唬那些受了惊吓的神，强迫他们赐恩，或者安抚他们的恐惧，从而使他们能够不受胁迫而行善——即便假定善良的通神术士自己没有仪式可以清除他为了净化自己灵魂而呼求的那些神的恐惧污点。为什么居然有一位神有能力吓唬低级的神，却没有一位神有能力解除他们的恐惧呢？难道找得到一位听从嫉妒者、吓唬众神不让做善事的神，却找不到一位听从心怀善意者、消除众神恐惧以便他们向他行善的神吗？哦，卓越的通神术！哦，可钦佩的灵魂净化！——在这通神术中，不洁嫉妒的暴力比纯洁神圣的恳求更有力量。不如我们憎恶并远离这类邪恶之灵的欺骗，聆听纯正的教训吧。至于那些用亵圣的仪式来操办污秽净化的人，他们据他说在受戒状态中看见（纵然我们可以质疑这种异象）天使或神某种奇妙可爱的显现，这就是宗徒提到 \BibleQuote{撒殚化形为光明的天使}\BibleRef{格后 11:14}时所指的。因为这些都是那长久渴想用许多假神的欺骗敬拜来缠住悲哀的灵魂，使他们偏离真正天主的敬拜的精灵的骗人异象；只有凭借这真正天主，灵魂才得以洁净与治愈；而且他，正如论及 \Person{普罗特斯} 所说的，\Quote{变成一切形状}，无论是在作对头攻击我们，或是装成朋友的样子，都是一样有害的。\par

% structure: p0026 source=h2 target=subsection reason=relative-heading-rank; base=h2

\subsection{第十一章 波菲利致阿内博的信——在信中他询问魔鬼之间的差别。}

波菲利在致埃及人阿内博的信中采取了一种较好的论调，在信中，他装成一个向他请教的人，从而揭露并驳斥这些亵渎神圣的技艺。在那封信中，他确实摒弃了所有的\OriginalTerm{魔鬼}{daemones}，主张它们如此愚蠢，竟会被祭品的烟雾所吸引，因此不住在以太中，而是住在月亮下面的大气中，甚至就在月亮里面。然而，他还没有胆量把一切欺骗、恶意和愚蠢的行为都归于所有的魔鬼，尽管这些行为正当地激起了他的愤慨。因为，虽然他承认魔鬼作为一个种类是愚蠢的，但他却迁就流俗的观念，称其中一些为善意的魔鬼。他感到惊讶的是，祭献不仅能左右诸神，还能强迫他们去做人们所希望的事；他感到困惑不解的是，如果神与魔鬼的区别在于其\OriginalTerm{无形体性}{incorporeitas}，那太阳、月亮以及其他可见的天体——他并不怀疑它们是有形体的——又怎能被视为神；此外，如果它们是神，为什么有些被称为有益的，而另一些被称为有害的，并且它们既是有形体的，又如何能被列入无形体的神之中。他进一步询问，并仍以疑惑的口吻说，\OriginalTerm{占卜者}{divinatores}和行奇迹者究竟是具有异常强大灵魂的人，还是这种行事的能力是由外来的神灵传达的。他倾向于后一种观点，理由是，他们正是利用石头和药草来对人施加咒语，打开关闭的门，并施行类似的奇事。因此，他说，有人认为存在一种生灵，其特性就是听从人类——一种欺诈的、诡计多端的、能变换种种形象，模仿诸神、魔鬼和死人的生灵——正是这类生灵制造出这一切看似或善或恶的事情，但对于真正善的东西，他们从不帮助我们，事实上完全不了解，因为他们使作恶变得容易，却给那些热切追求美德的人设置障碍；他们充满骄傲与鲁莽，喜欢祭品的香气，迷恋谄媚之言。对于这些以及其他关于这类欺诈、恶意的神灵的特性——他们进入人们的心灵，在睡梦和清醒时迷惑他们的感官——他并非当作自己确信之事来描述，而只是带着如此多的猜疑和疑虑，以致把它们说成是世人通常接受的意见。我们应该同情这位伟大的哲学家，他在了解和自信地攻击整个魔鬼家族时所遇到的困难，这件事任何一个信基督的老妇都能毫不犹豫地描述，并且极其痛恨。然而，也许他是不敢得罪他写信的对象阿内博——他本人正是这些神秘仪式的最杰出的赞助者——或者其他那些惊叹于这些魔术表演如同神明工作、并与诸神的崇拜紧密相关的人。\par

然而，他继续以此为题，并仍以探究者的身份，提到了一些事情，凡头脑清醒之人，都只会将这些事归咎于那些心怀恶意且擅于欺骗的势力。他问道：为什么在祈求了较善的一类灵体之后，却命令较恶的一类去满足人们邪恶的欲望？为什么它们不理睬一个刚刚离开女人怀抱的人，而它们自己却毫无顾忌地引诱人们犯下乱伦与奸淫之罪？为什么它们的司祭被命令禁绝肉食，以免受肉体浊气的污染，而它们自己却被祭品的烟气及其他浊气所吸引？为什么受祕者被禁止触碰尸体，而它们的种种神秘仪式却几乎全凭尸体来举行？为什么一个沉溺于种种恶习的人，会不向\OriginalTerm{魔鬼}{demon}或亡魂，而是向日月或某一天体发出威胁，以想象的恐怖来恐吓它们，以便从它们那里强索真实的恩惠？——因为他威胁说要捣毁天空，以及其他这类不可能做到的事——以便那些神明，像幼稚的孩童一般，被这些虚妄而荒谬的威胁所惊吓，从而执行所命的任何事。\Person{波菲利}进一步转述，有一个叫\Person{查雷蒙}的人，深谙这些神圣的、或不如说渎神的奥迹，他曾写道，著名的\Place{埃及}的\Person{伊希斯}与其丈夫\Person{奥西里斯}的奥迹，对神明们具有极大效力，能迫使它们执行所命之事，当那位使用咒语的人威胁要揭露或废除这些奥迹，并以威胁的声音喊道，若不听从他的命令，他就要散掉\Person{奥西里斯}的肢体时，便会产生这种效力。\Person{波菲利}有理由感到惊奇，一个人竟向神明——并非微不足道的神明，而是天上的，且闪耀着星芒的神明——发出如此疯狂而空洞的威胁，而这类威胁竟能产生实效，以不可抗拒的力量束缚它们，并使它们惊恐不安，从而满足他的愿望。他并非没有理由地，以探究这些可惊可怪之事的缘由的探究者身份，让人可以理解，这些事是由那一类灵体所为，他此前曾像是引述他人的见解一样描述过它们——这类灵体，如他所说，并非因其本性，而是因其自身的腐化而进行欺骗，它们假扮神明和亡人，但并非他所说的魔鬼，因为实际上它们就是魔鬼。至于他的那种想法，即藉由草药、石头、动物，以及某些咒语、声响、画符——有时是凭空想象的，有时则是仿照天体运动的——人在地上创造出能产生各种效果的力量，这一切不过是这些魔鬼对臣服于它们的人所施的骗术，为的是讥笑愚弄那些受骗者，供自己取乐。那么，要么\Person{波菲利}在其疑问与探究中是真诚的，他提及这些事情，是为了证明并确定无疑地表明，这些事并非那些帮助我们获得生命的势力所为，而是欺骗性的魔鬼干的；要么，若以更善意的眼光看待这位哲学家，他对那位耽于这些错误并引以为傲的\Place{埃及}人采用了这种方法，以免因采取教师的姿态而冒犯他，或因公开攻击者的争辩而扰乱其心神，而是凭借扮演探究者的角色，怀着急于求学的谦逊态度，将他的注意力引向这些事情，并显明它们是当受何等藐视与摈弃。在这封信的结尾部分，他请求\Person{阿内博}告诉他，\Place{埃及}的智慧指明何者为通往真福的道路。至于那些与神明相交、只是为了寻找逃亡的奴隶、获取财产、撮合婚姻之类的事而烦扰神明的人，他宣称，他们自诩睿智全然是虚妄的。他补充说，纵使承认这些神明在其他方面的宣示是真的，它们就真福所作的启示，却如此不明智、令人不满，以致于它们不可能是神明，也不可能是\OriginalTerm{善神}{good demons}，而要么是那个被称为欺诈者的灵，要么纯粹是想象的虚构。\par

% structure: p0029 source=h2 target=subsection reason=relative-heading-rank; base=h2

\subsection{第十二章 论真天主藉圣天使之职所行的奇迹}

由于通过这些技艺所行的奇迹远远超出人的能力，我们除了相信那些看似神奇且属神、同时却不属于对唯一天主的崇拜——正如柏拉图主义者们充分证实的，一切真福全在于与这位天主的结合——的预言和行为，是邪恶神灵的戏弄，他们以此试图诱惑并阻碍真正虔诚的人之外，还能作何选择呢？另一方面，我们不得不相信，一切奇迹，无论是藉由天使或其他方式所行，只要它们是在推崇那唯一真福所系的唯一天主的崇拜和宗教，它们就是由那些以真实而虔诚的方式爱我们的人所行，或藉由他们，天主亲自在他们内工作。因为我们不能听信那些主张无形的天主不行可见奇迹的人；因为他们自己也相信天主创造了世界，而他们当然不会否认世界是可见的。这世上发生的任何奇事，当然都不及这整个世界本身——我指的是天地及其中的一切——更为神奇，而这一切确是天主所造。然而，正如造物主自身对人而言是隐藏而不可理解的，创造的方式也是如此。因此，尽管这可见世界的恒常奇迹因其常在我们眼前而很少被重视，但当我们激发自己去默观它时，它却是一个比最罕见、最闻所未闻的奇事更大的奇迹。因为人本身，就是一个比任何藉他而行的奇迹更大的奇迹。所以，那创造了可见天地的天主，并不鄙弃在天地间行可见的奇迹，以此唤醒那沉溺于可见事物中的灵魂去崇拜他自己——那无形者。但这些奇迹的地点和时间，取决于他那不变的意志，在他的意志中，未来的事物如同已经成就了一样被安排。因为他在推动时间中的事物时，自身并不在时间中运动；他不以一种方式认识将要发生的事，又以另一种方式认识已经发生的事；他聆听祈祷者的方式，也与他预见谁将祈祷的方式无异。因为，即使当他的天使们聆听我们时，也是他亲自在他们内聆听我们，如同在他那并非人手所造的真正殿宇中，即在他那些圣人身上；而他的回应，虽在时间内完成，却是由他永恒的任命所安排的。\par

% structure: p0031 source=h2 target=subsection reason=relative-heading-rank; base=h2

\subsection{第十三章 论无形的天主，他常使自己成为可见的，并非按其本体，而是按观看者所能承受的。}

我们也不必奇怪，天主虽然是无形的，却常常有形地显现给圣祖们。因为，正如那传达心中默念的声音并非思想本身，同样，天主按其本性虽是无形的，却借以成为可见的那形式，也并非天主本身。然而，正如那思想本身在声音中被听到一样，在那形式下被人看见的，也正是他本身；圣祖们也认识到，虽然那有形的身体并非天主，他们却看见了无形的天主。因为，尽管\Person{梅瑟}曾与天主交谈，但他仍说：\BibleQuote{如果我在你眼前蒙恩，求你将你自己显示给我，使我看见并认识你。} \BibleRef{出 33:13}并且，既然那并非赐给一个人或少数开明的人，而是赐给整个众多民族的法律，应当伴有令人敬畏的标记，因此，藉着天使的服役，在山上，当法律通过一个人赐给民众时，大奇迹便在众人眼前施行，大众都目睹了那可怕的景象。因为以色列人相信\Person{梅瑟}，并不像\Place{拉刻代孟}人相信他们的\Person{吕库古}那样，以为他是从\Person{朱庇特}或\Person{阿波罗}那里领受了所赐给他们的法律。因为当那命令崇拜唯一天主的法律赐给民众时，神奇的标记和地震，就是那属神的智慧判定为足够的，在众目睽睽之下发生了，好使他们知道，正是那造物主，能够如此利用受造物来颁行他的法律。\par

% structure: p0033 source=h2 target=subsection reason=relative-heading-rank; base=h2

\subsection{第十四章 唯一天主应当受崇拜，不仅是为了永恒的真福，也连及现世的繁荣，因为万事都由天主眷顾所统御。}

人类的教养，由天主的子民所代表，如同个人的教养一样，经历了某些时代，或可谓世纪，以便能逐渐从世俗之事提升至属天之事，从可见之物提升至不可见之物。这一目的被如此清晰地维持着，以至于即使在承诺现世赏报的时期，唯一天主仍被呈现为崇拜的对象，以免人在涉及此短暂生命的现世福祉时，承认真正造物主及灵魂之主以外的任何其他对象。因为，谁若否认天使或人所能赐予我们的一切，都在独一全能者的手中，便是一个疯子。柏拉图主义者\Person{普罗提诺}论述天主眷顾时，从花叶之美证明了，从那其美不可见且不可言的至高天主，天主眷顾直降至下界的这些世事；他论证说，所有这些脆弱且易逝的事物，若不由那其不可见且不变的美丽恒常渗透万物的他所塑造，便不可能具有如此精致而细巧的美。主耶稣也证实了这一点，他说：\BibleQuote{你们观察一下田间的百合花怎样生长：它们既不劳作，也不纺织；可是我告诉你们：连\Person{撒罗满}在他极盛的荣华时所披戴的，也不如这些花中的一朵。田地里的野草今天还在，明天就投在炉中，天主尚且这样装饰，信德薄弱的人哪，何况你们呢！} \BibleRef{玛 6:28} 因此，最好的是，那仍懦弱地贪求世俗之物的人灵，应被培养得惯于向唯一天主寻求，即使这些微小的现世恩惠，以及此短暂生命的尘世需求——这些在与永恒的真福相比时是可鄙的——如此，即便是对这些事物的欲望，也不致使其偏离对那位我们藉着轻视和弃绝这类事物而归向的天主的崇拜。\par

% structure: p0035 source=h2 target=subsection reason=relative-heading-rank; base=h2

\subsection{第十五章 论圣天使的服役，他们藉此成就天主眷顾。}

因此，就如我已说过的，正如我们在\BibleBook{宗徒}{大事录}中所读到的\BibleRef{宗 7:53}那法律是天主眷顾乐意藉天使的安排颁赐命令人只敬拜唯一天主的法律。但在他们中间，天主本身的位格曾以可见的方式显现，当然，不是以其固有的本体——这本对必死的眼目常是看不见的——而是借着受造物顺从造物主所提供的无误标记。祂也使用了人类言语的词句，逐一发出音节，然而按照祂的本性，祂并非以肉体的方式说话，而是以属灵的方式；不是对着感官，而是对着心灵；不是用时态中的言词，而是，如果可以这么说，以永恒的方式，既没有开始说话，也没有结束。祂所说的，不是被身体的耳朵，而是被祂的使者的心灵之耳准确地听见，他们因享有祂不变的真理而享有不朽的真福；而他们以某种不可言喻的方式所领受的指示，他们在可感可见的世界里毫不迟延、毫无困难地执行。这法律乃是依照世界的时代而颁赐的，起初包括属地的应许，如我已说过的，但这些应许却象征属天的；这些永恒的真福虽有许多人参与了其可见标记的庆典，却很少有人理解。尽管如此，那法律的话和可见的礼仪却一致地命令人敬拜唯一的天主——不是一群神明中的一位，而是那位创造了天和地，以及一切与祂自己不同的灵魂与神体的造物者。祂创造了；其他一切都是受造的；并且，无论是存在还是善存，万有都需要那位创造它们的祂。\par

% structure: p0037 source=h2 target=subsection reason=relative-heading-rank; base=h2

\subsection{第十六章 关于永生之路，我们应相信那些要求我们向他们献上神圣钦崇的天使，还是那些教导我们将神圣侍奉归于天主而非他们自己的天使？}

那么，关乎真福与永生这件事，我们究竟应当相信哪些天使呢？——是那些愿意人用宗教礼仪和节庆来敬拜他们，并要求人向他们献祭的天使；还是那些宣称这一切敬拜只应归于唯一天主、造物主，并教导我们以真正的虔诚向祂奉献，而他们自己已因享见祂而获真福，并应许我们也将如此的天使呢？因为那享见天主的异象，其美善如此伟大，如此无限令人渴慕，以至于\Person{普罗提诺}毫不犹豫地说，那拥有其他一切福乐极其丰盛，却没有这一异象的人，是极其不幸的。因此，既然有些天使行奇迹为诱使我们敬拜这位天主，而另一些则行奇迹为诱使我们敬拜他们自己；既然前者禁止我们敬拜后者，而后者却不敢禁止我们敬拜天主，那么，我们究竟该听从哪一类呢？让柏拉图学派来回答，或者任何哲学家，或者\OriginalTerm{神术士}{theurgists}，或更确切地说，\OriginalTerm{\Emph{佩里乌术士}}{periurgists}——因为这个名字对从事这类法术的人来说恰如其分。总之，让所有人回答——只要在他们内还残存着他们受造时作为理性受造物所拥有的那自然知觉的些许火花——我说，让他们告诉我们，我们应当向那些命令我们向他们献祭的神明或天使献祭，还是应当向那位由禁止我们敬拜他们自己或那些其他者的天使所命令我们敬拜的唯一者献祭？如果没有哪一方行过奇迹，而只是发布了命令，一方命令人们向他们自己献祭，另一方则禁止如此，并命令我们向天主献祭，那么，一虔诚的心灵必不难分辨，哪一道命令出自骄傲的傲慢，哪一道出自真正的宗教。我更进一步说。如果只有那些为自己要求祭献的天使行了奇迹，而禁止此事并命令只向唯一天主献祭的天使，却认为宜完全舍弃使用可见的奇迹，那么，凡不仅用眼，更运用理性的人，都当选择后者的权威。但是，既然天主为将祂真理的谕旨推荐给我们，便借着这些不朽的使者——他们宣扬的是祂的尊威，而非自己的骄傲——行了无比伟大、确定且显著的奇迹，好使虔诚者中的软弱之辈，不致被那些要求我们向他们献祭，并企图以令人惊骇的感官证据来说服我们的假宗教所诱离真理；那么，当他发现真理有比虚假更惊人的证据来作前导时，谁还会如此丧尽理智，竟不去选择并跟随真理呢？\par

至于历史记载归因于\Person{异教神祇}的奇迹——我不是指那些由于某些未知的物理原因间歇发生的、由\Emph{天主眷顾}所安排与指定的异象，例如怪异的诞生和异常的气象现象，无论它们仅是令人惊骇还是也造成伤害，并且据说通过与\Emph{魔鬼}的交往及其最狡诈的诡计而引起或消除——而是指那些显然是由它们的力量和能力所行的异象，例如：\Person{埃涅阿斯}在逃亡中从特洛伊携带的\Emph{家神}能够移动位置；\Person{塔尔奎尼乌斯}用剃刀切割磨刀石；\Emph{埃皮达鲁斯蛇}在\Person{阿斯克勒庇俄斯}前往\Place{罗马}的途中依附于他作为伴侣；载有\Emph{弗里吉亚母神}像的船只，由众多人力和牛群都无法移动，却被一个柔弱的女子用腰带系住船只而拉动，以此证明她的贞洁；一位\Emph{维斯塔贞女}的童贞受到质疑，她从\Place{台伯河}中打捞出满满一筛子的水而滴水不漏，从而消除了嫌疑：诸如此类的事件，无论在伟大还是德能上，都绝无法与我们所读到的在天主子民中所行的奇迹相比。更不可将那些异教国家的法律所禁止和惩罚的奇事——我指的是\Emph{魔法与通神术}的奇事，其中大部分仅仅是加于感官的幻象，例如招引月亮，正如\Person{卢坎}所说，\Quote{使它对植物施加更强的力量？}——相提并论。即使其中有些似乎与虔诚者所行的奇迹相等，但行奇迹的目的将二者区分开来，并表明我们的奇迹无可比拟地更为卓越。因为那些奇迹推崇对众多神祇的敬拜，它们越是要求敬拜，就越不值得敬拜；而我们的这些奇迹推崇对唯一天主的敬拜，祂通过\WorkTitle{圣经}的见证以及\Emph{祭献}的最终废除，证明自己不需要这类奉物。因此，若有天使为自己要求\Emph{祭献}，我们必须更喜爱那些不为自己，而是为万有的创造者——他们侍奉的天主——要求\Emph{祭献}的天使。因为他们以此证明他们多么真诚地爱我们，因为他们希望通过\Emph{祭献}使我们顺服，不是归向他们自己，而是归向那位他们在瞻仰中自身就享有真福的天主，并引领我们归向那位他们自己从未偏离过的天主。另一方面，若有天使希望我们向多位神祇献祭，确实不是向他们自己，而是向他们所奉的神祇，那么我们在此情况下也必须更喜爱那些属于万神之神的天主的天使，他们命令我们敬拜祂，以排除我们敬拜任何其他神祇。然而，更有可能的是——正如他们的骄傲与欺诈所显示——他们既非善天使，也非善神的使者，而是邪恶的\Emph{魔鬼}，他们希望\Emph{祭献}不是献给唯一而至高的天主，而是献给他们自己。那么，除了那位善天使所侍奉的唯一天主，我们还能选择什么更好的保护呢？那些善天使命令我们不要向他们自己献祭，而要向那位我们应以自身为祭献的天主献祭？\par

% structure: p0040 source=h2 target=subsection reason=relative-heading-rank; base=h2

\subsection{论约柜及天主用以证实律法与许诺的奇迹标记}

因此，由天使安排所颁布的天主法律，命令唯独万神中唯一天主接受神圣崇拜，而排除其他一切神明，这法律存放在约柜中，即称为\OriginalTerm{约柜}{ark of the testimony}的柜子。这个名称充分表明，并非那些礼仪中所崇拜的天主被局限和幽禁在那个地方——虽然他的回应连同感官可察觉的标记从那里发出——而是他的旨意从那个宝座上传扬出来。法律本身也刻在石版上，正如我所说，存放在约柜中；在旷野逗留期间，司铎们恭敬地抬着它，连同帐幕，这帐幕同样称为\OriginalTerm{约柜的帐幕}{tabernacle of the testimony}；那时还有一个伴随的标记，白天如云彩，夜间如火柱；云彩移动，营地就迁移；云彩停驻，营地就扎营。除了这些标记和发自约柜所在处的声音外，还有别的神奇证据证实法律。因为当约柜抬过\Place{约旦河}，进入预许之地时，河水上流停住，下流奔流，以致为约柜和人民都露出干地可以过去。然后，当他们抬着约柜绕他们遇到的第一座仇敌多神教城七次时，城墙突然倒塌，虽然无人手攻击，也无破城槌撞击。后来，当他们定居在预许之地，竟因自己的罪而被敌人掳去了约柜，掳掠者得意地将约柜放在他们所敬神明的庙里，关在那里，但次日开启庙门，发现他们常向之祈祷的偶像倒在地上，可耻地粉碎了。随后，他们自己因前兆而惊恐，并受到更可耻的惩罚，便将约柜归还给他们夺来的人。归还的方式如何？他们将约柜放在一辆车上，套上刚被夺去牛犊的母牛，任她们自择道路，期望借此显示天主的旨意；那些母牛无人驱赶或引导，径直走上通往希伯来人的道路，不顾自己牛犊的哀鸣，就这样将约柜归还给敬拜它的人。对天主而言，这些奇迹是小事，但为恐吓并给人健全的教训，它们却具有极大的功效。因为如果哲学家，尤其如\OriginalTerm{柏拉图主义者}{Platonists}，正如我刚才提到的，被公正地认为比其他人更有智慧，因为他们教导说，甚至这些属地和微小的事物也受\OriginalTerm{天主眷顾}{Divine Providence}的掌管——他们从无数可观察到的美丽中推断出这一点，这些美丽不仅存在于动物的身体中，也见于植物和草丛——那么，在预言的时间发生的这些事情，岂不更明显地证明了天主的临在吗？而在这些事情中，那禁止向任何天上的、地上的或阴间的神明奉献祭献，却命令只向天主奉献祭献的宗教得到了嘉许；天主因他爱我们和我们爱他而独赐福给我们；他通过安排那些祭献的预定时间，并预言它们将由一位更好的司铎转为更好的祭献，从而证明了他并不贪求这些祭献，而是借着它们使人联想到其他更实质的祝福。而这一切，并非要使他本人因这些恭敬而受光荣，而是要激励我们朝拜和依附他，被他的爱火所燃，这乃是为了我们的益处，而非为他的益处？\par

% structure: p0042 source=h2 target=subsection reason=relative-heading-rank; base=h2

\subsection{第18章 驳斥那些否认应相信教会书籍中关于天主子民所受奇迹教育的人。}

难道有人会说这些奇迹是假的，从未发生过，有关它们的记载都是谎言吗？凡这样说的人，并断言在这类事上没有任何记录可信，也可以说没有神明关心人类事务。因为他们诱使人崇拜他们，惟独通过奇迹的工作，这些在异教史书中有记载，神明以此炫耀自己的能力，而非提供真正的服务。这就是为何我们在这正在写作第十卷的著作中，没有着手驳斥那些否认有神能存在，或争辩神能并不干涉人类事务的人，而是驳斥那些宁愿相信自己的神明，而不信我们天主的人——我们天主是圣善而最荣耀之城的建立者，他们未知他也是这可见又变幻的世界不可见的、不变的奠基者，是真福生活的最真实赐予者，这真福生活不在于受造之物，而在于他自身。因为他最可信赖的先知这样说道：\BibleQuote{亲近天主对我是多么的美好}\BibleRef{咏 73:28}。哲学家们追问，那为我们一切职责所应趋向的目的和至善究竟是什么？圣咏作者并没有说：拥有巨大财富，或穿戴皇袍、紫袍、权杖和冠冕，对我多么美好；也没有像某些哲学家那样厚颜地说：享受感官快乐对我多么美好；也没有像他们中较优秀者似乎所说的：我的善是我的灵性力量；而是说：\BibleQuote{亲近天主对我是多么的美好}\BibleRef{咏 73:28}。这是他由天主那里学得的，圣洁的天使伴随奇迹的见证，将他呈现为唯一应受崇拜的对象。因此他自己也成了天主的祭品，天主的灵性之爱点燃了他，他渴望投身于他无可言喻且无形体的拥抱中。再者，如果多神崇拜者（无论他们幻想自己的神是哪一类）相信他们民族史书、魔法书或更可敬的\OriginalTerm{通神术}{theurgy}书籍中所载的奇迹，是由这些神明所行，他们有什么理由拒绝相信记载于这些书卷中的奇迹——我们对这些书卷应予以更大的信赖，正如这些书卷教导我们应唯独向他奉献祭献的那位更伟大一样？\par

% structure: p0044 source=h2 target=subsection reason=relative-heading-rank; base=h2

\subsection{第19章 论真宗教教导向唯一真实无形的天主献有形祭献之合理性。}

关于那些认为这些可见的祭献适宜奉献给其他神祇，而不可见的祭献——即心神的纯洁和意志的圣善之恩宠——应作为更伟大、更美善的祭献，奉献给那位比一切更伟大更美善的不可见的天主，他们必定忘记，这些可见的祭献乃是不可见之物的标记，正如我们发出的言语是事物的标记。因此，正如我们在祈祷或赞颂中，将可理解的言语献给那位，我们心中对祂呈上我们正表达的情感，同样我们应明白，在祭献中，我们只有对那位我们应在心中将自己作为不可见的祭品献给祂的天主，才奉献可见的祭献。正是在这时，天使们，以及一切因善良和虔诚而强大的上等权能，才会喜悦地注视我们，与我们同乐，并竭尽全力协助我们。然而，如果我们向它们献上这样的敬拜，它们会拒绝；当它们受命来到人间而显现于感官时，它们更会明确禁止。这类例子在圣经中可见。有些人妄想应该以朝拜或祭献，将唯独归于天主的尊荣同样给予天使，但天使们自己制止了他们，并命令他们将这尊荣只归于他们所知的应得的天主。在这件事上，圣天使为天主的圣人们所效仿。因为\Person{保禄}和\Person{巴尔纳伯}，在\Place{吕考尼雅}施行了治愈的奇迹后，被当作神明，吕考尼雅人想要向他们献祭，而他们却谦卑而虔诚地拒绝了这份尊荣，并向众人宣讲了他们应当信仰的天主。那些欺诈而骄傲的灵体，之所以要求敬拜，仅仅是因为他们知道这本应归于真天主。正如\Person{波菲利}所言，且有些人所想象的那样，他们喜悦的并不是祭牲的气味，而是神性的尊荣。事实上，他们随处可得充裕的气味，若想要更多，也能自己供给。但那些妄自尊大的灵体，喜悦的不是尸体的烟气，而是他们所欺骗和奴役的祈求者的灵魂，他们阻碍这灵魂亲近天主，阻止他向天主奉献自己作为祭献，却引诱他向他神献祭。\par

% structure: p0046 source=h2 target=subsection reason=relative-heading-rank; base=h2

\subsection{第二十章　论那由天主与人之间的中保所完成的至高而真实的祭献}

因此，那位真正的\OriginalTerm{中保}{Mediator}，就祂取了奴仆的形象，成为天主与人之间的中保——即人\Person{耶稣基督}——而言：虽然在天主的形象中，祂与父一同接受祭献，与父同为唯一天主，然而在奴仆的形象中，祂宁愿成为祭品，也不接受祭献，以免有人借此以为应向任何受造物奉献祭献。如此，祂既是奉献祭品的司祭，也是被奉献的祭品。祂愿意在教会的祭献中有一项每日的标记，这教会是祂的身体，学习借着祂奉献自己。这真实的祭献，昔日圣人们众多的古代祭献乃是它的种种标记；它被如此多样地预表，正如一件事物用多种不同的言辞来表达，使我们在反复谈论时不觉厌倦。在这至高而真实的祭献面前，一切虚假的祭献都让路了。\par

% structure: p0048 source=h2 target=subsection reason=relative-heading-rank; base=h2

\subsection{第二十一章　论那授予魔鬼以试探并光荣圣人的权力，圣人不靠取悦空中的神灵，而是靠居于天主内得胜}

授予魔鬼的权力，在指定和安排得当的时期，使他们得以表达对天主之城的敌意，鼓动受他们支配的人反对它，不仅能从那些自愿献祭者接受祭献，还通过暴烈的迫害向不愿者强索祭献——这种权力非但无害，甚至有益于教会，因为它完成了殉道者的数目；天主之城视这些殉道者为更加光耀、更受尊敬的公民，因为他们与不虔之罪相争，甚至流血。如果教会的常用语允许，我们或许能更优雅地称这些人为我们的英雄。因为据说这个名字源于\Person{朱诺}，她在希腊语中称为\Person{赫拉}，因此，根据希腊神话，她的一个儿子被称为\Person{赫罗斯}。这些神话神秘地表示，\Person{朱诺}是空气的女主人，人们认为空气中有魔鬼和英雄居住，所谓英雄，是指那些功绩卓越的亡者的灵魂。但我们称我们的殉道者为英雄，却是出于截然相反的理由——我先前说过，假如教会用语允许的话——不是因为他们与魔鬼一同住在空中，而是因为他们战胜了这些魔鬼或空中势力，其中包括\Person{朱诺}本人；不论她究竟是谁，诗人们通常把她描绘成敌视德行、嫉妒那些显著向往天界的人，这倒并非不适合。然而，\Person{维吉尔}不幸地退让，屈服于她；因为尽管他让她说，\Quote{我被\Person{埃涅阿斯}战胜了，}但\Person{赫勒努斯}却给埃涅阿斯本人这样的虔诚劝告：\par

\Quote{向\Person{朱诺}许愿：用礼物和祈祷\newline{}压服她那女王般的灵魂。}\par

根据这种看法，\Person{波菲利}——不过他所表达的，与其说是他自己的观点，不如说是别人的——也说，除非先取悦那邪恶的神灵，否则善良的神或守护神不会来到人那里，这意味着邪恶的神明比善良的神明更有力量；因为在邪恶的神明尚未被安抚并让路之前，善良的神明无法给予任何帮助；如果邪恶的神明加以阻挠，善良的神明也无能为力；而邪恶的神明却能加害于人，善良的神明却无法阻止。这不是真正的、真正神圣的宗教的方式；我们的殉道者并不是这样战胜\Person{朱诺}——即那些嫉妒虔诚者德行的空中势力。我们的英雄，如果我们能这样称呼他们的话，战胜\Person{赫拉}，靠的不是祈求的礼物，而是神圣的德行。正如\Person{西庇阿}，他凭勇力征服了\Place{阿非利加}，因此被称为阿非利加努斯，这比他用礼物讨好敌人而赢得他们的怜悯更为适宜。\par

% structure: p0052 source=h2 target=subsection reason=relative-heading-rank; base=h2

\subsection{第二十二章　圣人对抗魔鬼的能力以及心灵的真正净化从何而来}

借着真正的虔敬，属于\Person{天主}的人驱逐那敌挡虔敬的空中的敌对势力；不是用安抚，而是用驱魔的方式；他们战胜仇敌的一切诱惑，不是向仇敌祈祷，而是向自己的\Person{天主}祈求，以对抗它。因为魔鬼只能征服那些与罪联盟的人；因此，它被那取了人性且无罪者的名字所征服；这取了人性者自己既是\OriginalTerm{司铎}{Priest}又是\OriginalTerm{祭献}{Sacrifice}，为使\OriginalTerm{罪过的赦免}{remission of sins}成就，也就是借着\OriginalTerm{中保}{Mediator}，即\Person{天主}与人之间的中保，人\Person{耶稣基督}，透过祂我们与\Person{天主}和好，罪污得以洁净。因为唯独罪使人远离\Person{天主}，在今生我们得洁净不是靠自己的德行，而是靠\Person{天主}的怜悯；是靠祂的宽恕，不是靠我们自己的能力。因为无论我们自称有何德行，都是祂的美善赐予的。我们在肉身之时，若不是在获得赦免中生活，直到脱下这肉身，我们就会过分自恃。这就是为什么借着\OriginalTerm{中保}{Mediator}，赐给了我们这\OriginalTerm{恩宠}{grace}：我们这被罪恶的肉身所污染的人，得以借着\OriginalTerm{相似于罪恶的肉身}{likeness of sinful flesh}而被洁净。靠着\Person{天主}的这恩宠——祂在其中向我们显示了祂极大的怜悯——我们在今生受\OriginalTerm{信德}{faith}引导，并在今生之后，借着对\OriginalTerm{不变真理}{immutable truth}的直观，被导向最完全的成全。\par

% structure: p0054 source=h2 target=subsection reason=relative-heading-rank; base=h2

\subsection{论柏拉图主义者所主张的灵魂净化之原理。}

甚至\Person{波菲利}也断言，神谕启示我们，向日月献祭无法使人得洁净；这暗示我们，向任何神祇献祭都不能使人得洁净。因为倘若被视为天体诸神之首的日月之祭都不能洁净人，那还有什么奥秘能洁净人呢？他在同一处还说，唯有\Quote{原理}才能洁净人，以免有人根据他说向日月献祭不能洁净人，便以为向其他神祇中的某一位献祭或许可以。而我们知道，作为柏拉图主义者，他所说的\Quote{原理}意指什么。因为他论及圣父和圣子，称圣子（用希腊文写作）为圣父的\OriginalTerm{理智或心灵}{intellect or mind}；但对圣神，他或者只字不提，或者未曾明言，因为我不明白他所说的那位于这两者之间的另一位是指什么。因为，如果像\Person{普罗提诺}在讨论\OriginalTerm{三个主要本体}{three principal substances}时那样，他意欲以此第三位指\OriginalTerm{自然的灵魂}{soul of nature}，那么他肯定不会将它置于这两者之间，即圣父与圣子之间。因为\Person{普罗提诺}将自然的灵魂置于圣父的理智之后，而\Person{波菲利}则令其居中，不置于后，而置于其他两者之间。无疑，他是按自己的见识，或自以为合宜而说的；但我们却断言，圣神不仅是圣父的，也不仅是圣子的，而是两者的。因为哲学家们随心所欲地说话，在最艰深的问题上，他们也不怕冒犯虔敬者的耳朵；而我们却必须按照一定的规则说话，以免言语的自由导致对我们所谈论之事务本身的见解不敬。\par

% structure: p0056 source=h2 target=subsection reason=relative-heading-rank; base=h2

\subsection{论唯一真实的原理，唯有它能净化并更新人性。}

因此，当我们论及天主时，并不肯定两个或三个\OriginalTerm{本原}{principles}，正如我们不能任意肯定两个或三个神；虽说论到每一位：圣父、圣子、圣神，我们宣认每一位都是天主：但我们不说，如\OriginalTerm{萨贝里异端者}{Sabellian heretics}所说，父与子相同，圣神与父和子相同；而是说，父是子的父，子是父的子，而发自父和子的圣神，既不是父也不是子。所以，人仅由一个本原得到净化，这话说得确实不错，尽管\OriginalTerm{柏拉图主义者}{Platonists}使用\Emph{本原}的复数是个错误。但\Person{波菲利}处于那些嫉妒势力的支配下，他既羞于又害怕摆脱其影响，拒绝承认基督就是那\OriginalTerm{本原}{Principle}，我们藉其\OriginalTerm{降生成人}{Incarnation}而得到净化。实际上，他轻视了基督，正因为祂所取的血肉本身，为给我们洁净而奉献祭献——这是一个大奥迹，波菲利的高傲所不能理解的，那位真实而仁慈的救主却以自己的谦卑贬抑了这骄傲，通过祂所取的死亡，向必死之人彰显自己，而恶毒欺诈的\OriginalTerm{中保}{Mediator}却以不具此死亡为傲，并向可怜人许诺，作为不死者的恩赐，提供骗人的助佑。因此，那位善良真实的中保指出，恶的是罪，而不是肉躯的实体或本性；因为肉躯连同人的灵魂，可以无罪地被取、被保留、被置于死地，并藉复活转变为更好的。祂也指出，死亡本身，虽是罪的惩罚，祂为了我们却无罪地承受了，我们不可因罪而逃避死亡，反而如机会允许，应为义德而忍受它。因为祂能藉着自己的死亡而赎罪，是因祂死了，却不是因自己的罪。但波菲利没有认出祂就是本原，否则他就会认出祂是净化者。本原不是基督的肉躯或人的灵魂，而是那创造万物的圣言。所以，肉躯不是凭自己的德能净化，而是凭那摄取它的圣言的德能，当\BibleQuote{圣言成了血肉，寄居在我们中间。}\BibleRef{若 1:14}。因为，当祂以神秘方式谈及吃祂的肉时，那些不明白的人大感震惊而离开，说：\Quote{这话生硬，谁能听得下去呢？}祂回答那些留下的说：\BibleQuote{使人生活的是圣神，肉一无所用。}\BibleRef{若 6:60}。因此，本原既取了人的灵魂和肉躯，就净化信者的灵魂和肉躯。所以，当犹太人问祂是谁时，祂回答说，祂就是\Emph{本原}。而我们这些属血气的、脆弱的、容易犯罪的、陷于无知黑暗的人，若非由祂经由我们曾是及我们不曾是的来净化、医治，便根本无法理解。因为我们曾是人，却不曾是义人；而在祂的降生成人中，有一人性，但却是正义、无罪的。这就是那伸向堕落跌倒之人的手，那中介；这就是那\BibleQuote{由天使所颁布的}后裔，法律也是借着天使服役而赐下，命令崇拜唯一的天主，并应许这位中保将要来临。\par

% structure: p0058 source=h2 target=subsection reason=relative-heading-rank; base=h2

\subsection{在法律之下和之前的诸圣，都因信德成义，归于基督降生成人的奥迹。}

正是借着对这奥迹的信德，以及虔敬的生活，即使古时的圣者们也能获得洁净；无论是在法律被颁布给希伯来人之前（因为那时天主和天使们已作为导师临在），还是在法律之下的时期——虽然属灵之事的应许以预像出现，看似属血肉的，因此称为旧约。因为在那个时代有先知们生活，由他们，如同由天使一样，宣布了同一应许；他们当中就有一位，关于人的终向与至善，我刚才引用了他的伟大而神圣的感悟：\BibleQuote{亲近天主是我的福乐。} 在这首圣咏中，旧约与新约的区别被清楚宣告。因为圣咏作者说，当他看见那不敬虔的人充分享受属血肉和属地的应许时，他的脚几乎滑跌，他的步履几乎跌倒；并且当他看到那些藐视天主的人，在他指望从天主手中得到的繁荣中反倒加增，他觉得自己似乎是徒然事奉了天主。他也说，他渴望理解何以如此，在探究这事时，他劳而无功，直到他进入天主的圣所，明白了那些他错误地视为有福之人的结局。于是他便明白，他们正是被他们所夸耀的事倾覆，正如他所说；他们因自己的不义而被消灭丧亡；那整个属世繁荣的架构，就如人从梦中醒来，突然发现自己失去了在睡梦中所想象的一切欢乐。且如同在这尘世或属世之城中，他们自视为伟大，他说：\BibleQuote{上主，在祢的城中，祢必使他们的影像归于虚无。} 他也指出，对他而言，只向那位掌管万有的独一真天主寻求甚至地上的福分，也是何等有益，因为他说：\BibleQuote{我在祢面前竟如牲畜，但我常与祢同在。} 他说\Quote{如牲畜}，意思是自己愚昧无知。因为我原应向祢寻求那些不敬虔的人不能像我一样享受的事物，而非那些我看见他们享用丰足的事物，并因此认为自己是徒然事奉祢，因为他们拒不事奉祢，却得到了我所没有的。然而，\BibleQuote{我常与祢同在}，因为即使在我渴求这类事物时，我也没有向别的神明祈求。因此他继续说：\BibleQuote{祢握住我的右手，以祢的训言引导我，随后以荣耀接我上升；} 仿佛一切属地的好处都是左手边的祝福，尽管当他看见恶人享受这些时，他的脚几乎滑跌。他说：\BibleQuote{在天上，除祢以外，我还能有谁？在地上，除祢以外，我别无所求。} 他责怪自己，且对自己不满，实属应当；因为他虽然在天上拥有如此宽广的产业（他后来明白了），却仍在地上向他的天主寻求一种短暂易逝的幸福——我们可以说，一种泥淖般的幸福。他说：\BibleQuote{我的肉体和我的心神虽已衰竭，我心之天主啊！} 幸福的衰竭，从下界的事物转向上界的事物！因此在另一首圣咏中他说：\BibleQuote{我的灵魂渴望上主的庭院，甚至衰竭。} 然而，虽然他提到了自己的心和肉体两者都已衰竭，他并没有说\Quote{我心灵和肉体的天主}，而是说\BibleQuote{我心灵的天主}；因为借着心灵，肉体得以洁净。因此，主说：\BibleQuote{先洁净杯盘的里面，好叫它的外面也成为洁净的。}\BibleRef{玛 23:26} 然后他说，天主自己——不是从他那里获得的任何事物，而是他本身——就是他的产业。\BibleQuote{我心的天主，我的产业，直到永远。} 在人类所选择的众多对象中，唯有天主使他满足。\BibleQuote{看，凡远离祢的，必要灭亡；凡对祢不忠而背叛祢的，祢必要消除，}——也就是说，那些委身于众多神明的人。然后便是那节似乎全篇圣咏都在为其预备的经文：\BibleQuote{亲近天主是我的福乐，}——不是远离，不是与众多神明行淫。然后，当在我们内一切有待救赎的都得到救赎时，这种与天主的合一便将臻于圆满。但至于目前，我们必须如他接着所说的，\BibleQuote{将我们的望德寄托于天主。} 宗徒说：\BibleQuote{因为我们得救在于望德；所见的望德不是望德，因为谁还希望所见的事物呢？但我们若希望那未看见的，就必坚忍等待。}\BibleRef{罗 8:24} 那么，既然我们现今在此望德中立定，让我们实行圣咏作者进一步所指示的，并按我们的尺度成为天主的天使或使者，宣扬他的旨意，赞美他的光荣和他的恩宠。因为当他说了\BibleQuote{将我的望德寄托于天主}之后，他接着说：\BibleQuote{好使我在\Place{熙雍}女郎的城门，宣扬祢的一切颂赞。} 这就是最光荣的天主之城；这就是认识并敬拜唯一真天主的城：她为圣天使们所称扬，他们邀请我们加入他们的团契，渴望我们在这城中成为他们的同乡；因为他们不愿我们把他们当神明来敬拜，而是希望我们与他们一同敬拜他们的、也是我们的天主；不愿我们向他们祭献，而是与他们一起成为对天主的祭献。因此，任何人只要放下恶意的固执，思考这些事，就必确信：这些有福且不朽的灵体，他们并不嫉妒我们（他们若嫉妒，便不是有福的了），反而爱我们，渴望我们像他们自己一样有福；当我们与他们一同敬拜唯一的真天主——圣父、圣子和圣神——时，他们比我们向他们自己祭献和敬拜，更喜悦地看顾我们，更大力地帮助我们。\par

% structure: p0060 source=h2 target=subsection reason=relative-heading-rank; base=h2

\subsection{第二十六章 论\Person{波斐留斯}在宣认真天主与崇拜邪魔之间摇摆不定的软弱}

我不知道为什么会这样，但在我看来，\Person{波斐利}为他的朋友\OriginalTerm{通神者}{theurgists}感到羞愧；因为他已知晓我所引述的一切，却未坦诚地谴责多神崇拜。事实上，他说有些天使降临世间，向通神者启示神圣真理，而另一些天使则在世上宣报属于圣父的事物，即祂的高深。那么，我们能否相信，那些职责是宣报圣父旨意的天使，愿意我们臣服于除他们所宣报其旨意的那一位以外的任何存在呢？因此，就连这位\OriginalTerm{柏拉图主义者}{Platonist}本人也明智地指出，我们应当效法他们，而非呼求他们。这样，我们就不必担心因不向唯一真天主的这些不朽而蒙福的臣属献祭而得罪他们；因为他们知道，祭献只应归于唯一真天主，他们正是通过效忠于祂而获得真福，因此他们不愿人们以象征祭献奥秘的仪式或实质将他们当作天主来祭献。这样的傲慢属于骄傲而可怜的恶魔，他们的性情与那些服从天主者的虔敬截然相反，而他们的真福在于依附天主。而且，为了使我们也能达至这一真福，他们以真诚的善意适宜地帮助我们，不篡取对我们的任何统治权，而是向我们宣报那位我们共同臣服的统治者。那么，哲学家啊，你为什么仍然不敢直言反对那些敌视真德与真天主之恩赐的势力呢？你已经区分了宣报天主旨意的天使与那些被某种我不知道的方术召下、降临于通神者的天使。为什么你仍将宣示神圣真理的尊荣归于后者？如果他们不宣报圣父的旨意，他们能发出什么神圣启示呢？这些不正是被一个忌妒者的咒语所束缚的恶灵吗？——他们被缚，以致无法将灵魂的纯洁赐予他人，而且正如你所说，他们无法被一位渴望纯洁的好人从这些束缚中释放，并重获对自己行为的控制权。你仍然怀疑这些是邪恶的恶魔吗？或者，你或许佯装不知，以免得罪那些以秘密仪式引诱你，并将这些疯狂而有害的鬼蜮伎俩当作厚礼教给你的通神者？你竟敢将这些嫉妒的势力——或不如称之为害虫吧，如你自己所承认的，他们更配称为奴隶而非君主——高举到空中，甚至高举到天堂；难道你不感到羞耻，竟将他们置于你的星辰诸神之列，从而轻蔑星辰本身？\par

% structure: p0062 source=h2 target=subsection reason=relative-heading-rank; base=h2

\subsection{第二十七章 \Person{波斐利}的不虔敬，甚至比\Person{阿普莱乌斯}的错误更为恶劣。}

你的同道柏拉图主义者\Person{阿普列乌斯}的错误是何等更可容忍、更符合人的情感啊！因为他只将人类情感的疾病和风暴归因于居住在月下等级的魔鬼，而且关于他所尊崇的神明，他作出这一承认也仿佛出于勉强；但对于那些居住在天上的、更高的神明，无论可见的，如太阳、月亮和其他星辰，因其光辉而显著，或不可见的，却为他所信奉的，他则竭尽所能使之远离这些扰乱的丝毫污染。因此，你们不是从柏拉图，而是从你们的\Place{迦勒底}教师们那里学会了把人性的罪恶提升到世界的以太境界和最高天，乃至天上的穹苍，好让你们的\OriginalTerm{通神术士}{theurgists}能够从你们的神明那里获得神圣启示；然而，你们却通过自己的理智生活，使自己超越这些神圣启示，因为这种生活无需这些\OriginalTerm{通神净化}{theurgic purifications}，认为这些对哲学家并非必需。可是，作为对你们教师们的酬谢，你们却将这些法术推荐给其他人，那些不是哲学家的人，可以被说服去使用你们自己所承认、对你们这些能够达到更高事物的人毫无用处的东西；这样一来，那些不能利用哲学美德的人——因为哲学对大众而言过于艰难——便可能受你们的煽动，去求助于通神术士，通过他们得以净化，不是净化灵魂的理智部分，而是其精神部分。而既然不适合哲学的人构成了人类中无可比拟的大多数，那么被迫去求教于你们这些隐秘而非法的教师，而非去柏拉图学园的人就会更多。因为这些最不洁的魔鬼，冒充为天上的神明，你们已成了他们的传令官和使者，他们应许说，那些通过\OriginalTerm{通神术}{theurgy}在灵魂的精神部分得以净化的人，固然不会回到圣父那里，却将居住在天上的神明中间，在气层之上。然而，基督来救拔脱离魔鬼暴政的群众，不听信这类幻想。因为在祂里面，他们有最慈悲的净化，心灵、精神和身体共同参与其中。因为，为要治愈整个人脱离罪的瘟疫，祂取了无罪而又完整的人性。若你们曾认识祂，若你们曾将自己交托给祂以求医治，而不是靠你们自己那脆弱软弱的人德，或靠那些有害而好管闲事的法术，那该多好！祂必不会欺骗你们；因为你们自己的神谕，按你们自己的说法，也承认祂是圣洁而不朽的。那位最著名的诗人也论到祂，确实是以诗的方式，因为他将这些话放在另一个人物口中，但若归之于基督，便是真实的了，他说：\Quote{在你的护佑下，若有任何罪行的痕迹存留，都将被抹去，大地亦将从其永久的恐惧中得释放。}诗人以此表明，由于此生所固有的软弱，即使德行和义德上取得了最大的进步，仍会给罪行——若不是罪行本身，至少是罪行的痕迹——留下存在的余地，而这些痕迹唯有这诗中所言的那位救主才能消除。因为\Person{维吉尔}在\WorkTitle{第四牧歌}的几乎最后一行告诉我们，他并不是凭自己的幻想说出这话，他说：\Quote{库迈女预言家所预言的末世如今已来临；}由此可见，这分明是由\Person{库迈女预言家}所口授的。但那些通神术士，或者更确切地说，那些冒充神明外貌和形体的魔鬼，借虚假的景象和虚幻形式的骗人伪冒来玷污人的精神，而非净化它。那些自身精神不洁的，怎能洁净人的精神呢？倘若他们不是不洁的，就不会被一个嫉妒者的咒语所束缚，也不会害怕或吝惜赐给他们所应许的空洞惠赐。但就我们的目的而言，你们承认理智灵魂，即我们的心灵，不能通过通神术而成义，这便足够了；并且你们也认为，即使灵魂的精神部分或较低部分也不能因此行为而成为永恒和不朽的，尽管你们坚称它可以借此得到净化。然而，基督应许了永生；因此世人纷纷归向祂，这令你们大为愤怒，也令你们大为震惊和困惑。你们勉强承认通神术引人迷途，以其无知而愚蠢的教导欺骗了无数人，并且承认通过祈祷和祭献求助于天使和掌权者是最明显的错误，但这又有什么用呢？你们同时为了免于在这些法术上徒劳无功的指责，竟指点人们去求通神术士，好让那些不按照理智灵魂的准则生活的人，可以通过他们来净化自己的精神灵魂，这又有什么益处呢？\par

% structure: p0064 source=h2 target=subsection reason=relative-heading-rank; base=h2

\subsection{\Person{波菲利}何以如此盲目，竟不认识真智慧——\Person{基督}}

因此，你把人引入最明显的错误。然而，你自称热爱德行与智慧，却做出如此大的伤害而不以为耻。如果你在这宣称上真诚而忠信，你就会承认基督，天主的德能与天主的智慧，就不会因虚妄学识的骄傲而背离祂有益的谦卑。然而你承认，灵魂的灵性部分可以借着贞洁的德行得到净化，无需借助你耗费时间学习的那类神智术和密仪。你甚至有时说，这些密仪在死后并不提升灵魂，因此，此生结束后，它们似乎对你所谓的灵性部分也毫无益处；但你却利用每个机会重提这些法术，就我所见，目的无非是显得像个熟练的神智术士，取悦那些对非法之术好奇的人，或是激发他人同样的好奇。然而，我们完全称赞你说这法术当受畏惧，既由于法律的禁令，也由于施行它本身所涉及的危险。但愿至少在这点上，那些可怜的崇奉者能听从你，好使他们能从完全沉迷中退出，甚或根本不受其沾染！你的确说过，无知及由它产生的无数恶习，无法被任何密仪除去，而只能由\OriginalTerm{圣父的理智}{\GreekText{πατρικὸς} \GreekText{νοῦς}}，即意识到圣父旨意的圣父之心灵或理智除去。但你不相信基督就是这理智；因为你因祂由女人所取的身体和十字架的耻辱而鄙视祂；因为你的高超智慧藐视这类低微可鄙的事物，而飞向更崇高的领域。但祂应验了圣先知们关于祂的真实预言：\BibleQuote{我要摧毁智者的智慧，废除贤者的聪明。} \BibleRef{依 29:14}因为祂所摧毁和废弃的，不是祂在他们内的恩赐，而是他们据为己有却不属于祂的东西。因此，宗徒引用了先知的这见证后，又说：\BibleQuote{智者在哪里？经师在哪里？这世代的诡辩者又在哪里？天主岂不是让这世界的智慧变成了愚妄吗？因为世人没有凭自己的智慧认识天主，天主就乐意以所传的愚妄道理拯救那些相信的人。的确，犹太人要求的是神迹，希腊人寻求的是智慧，而我们所宣讲的却是被钉在十字架上的基督：这为犹太人固然是绊脚石，为外邦人是愚妄，但为那些蒙召的，不拘是犹太人或希腊人，基督却是天主的德能和天主的智慧：因为天主的愚妄总比人明智，天主的懦弱也总比人坚强。} \BibleRef{格前 1:19}这在那些自命有智慧和强有力的人看来，被藐视为软弱愚蠢的事；然而这正是医治弱者的恩宠，他们不骄傲地夸耀自己的幸福，反而谦卑地承认自己的真实可怜。\par

% structure: p0066 source=h2 target=subsection reason=relative-heading-rank; base=h2

\subsection{论我主耶稣基督的降生成人，柏拉图学派因不虔敬而羞于承认。}

你宣扬圣父和祂的圣子，你称圣子为圣父的理智或心灵，并在两者之间安置第三位，我们猜想你意指圣神，并以你自己的方式称这三位为神。在这一点上，尽管你的用语不准确，但你确实在某种程度上，如同透过一层帕子，看到了我们应当勉力追求的；但不变的天主圣子的降生成人，我们借此得救，并能够达到我们所信的、或部分理解的事物，这正是你拒绝承认的。你仿佛远远地，以模糊的视线，看到了我们应当居住的家乡；但你不认识通向它的道路。然而你相信恩宠，因为你说只有少数人因着睿智而蒙恩到达天主。因为你没有说：\Quote{少数人认为合适或愿意}，而是说：\Quote{这是恩赐给少数人的}，明确承认天主的恩宠，而非人的能力。你更明确地使用这词，当你按照\Person{柏拉图}的观点，毫不怀疑地认为，人在今生绝不能获得完美的智慧，但所缺少的一切，会在来生由天主的眷顾和恩宠补足给那些度理智生活的人。啊，如果你承认了天主在主耶稣基督内的恩宠，以及祂那取了人的灵魂和肉身的降生成人，你便似乎成了恩宠最明亮的榜样！但我在做什么呢？我知道对死人说话是无用的，至少对你而言是无用的，但或许对于那些高度尊重你、因热爱智慧或好奇那些你不该学习的法术而爱你的人，并非徒劳无益；我以你的名义向这些人说话。天主的恩宠不可能以更优雅的方式举荐给我们了：天主的独生子，自身保持永恒不变，却取了人性，借着人性的中介，赐给我们祂爱的希望，使我们能从人的境况，来到如此遥远的祂那里——从有死的到不死的；从可变的到不变的；从不义者到义者；从可怜者到真福者。既然祂已赐给我们渴望真福和不死的自然本能，祂自己继续是真福的；却取了有死的身体，借着忍受我们所恐惧的，教导我们轻视它，好使祂将我们切望的赐给我们。\par

为使你们顺服此真理，所需的是\Emph{谦卑}，而要使你们做到这一点，则极为困难。对于像你们这样习惯于思辨、本应倾向于相信此事的人来说，有什么是不可信的？——我且说，有人宣称天主取了人的灵魂和身体，这其中有什么是不可信的？你们自己赋予\OriginalTerm{智性灵魂}{intellectual soul}——那毕竟就是人的灵魂——如此卓越的地位，以至于主张它能与你们所信为天主子的\OriginalTerm{圣父}{Father}的理智成为\OriginalTerm{同性同体}{consubstantial}。那么，如果祂为了众人的\OriginalTerm{救恩}{salvation}，以一种不可言喻且独一无二的方式取了一个灵魂，这又有什么不可思议呢？而且，我们的本性自身证明，除非身体与灵魂结合，否则人便不完全。当然，若非这是万事中最寻常的，这会更难令人相信；因为我们更容易相信精神与精神的结合，或用你们的话说，\OriginalTerm{无形的}{incorporeal}与无形的结合，即使一个是属人的，另一个是神圣的，一个是可变的，另一个是不可变的，也胜过相信\OriginalTerm{有形的}{corporeal}与无形的结合。但或许，那史无前例的由童贞女所生之身体令你们震惊？然而，这远非难事，反而应有助于你们接受我们的宗教，因为一个神奇的人物是以神奇的方式诞生的。或者，你们难以接受如下事实，即祂的身体被交于死亡之后，经由\OriginalTerm{复活}{resurrection}转变为一种更高层次的身体，不再是会死的，而是\OriginalTerm{不朽的}{incorruptible}，祂便带着它升入了高天？或许你们拒不相信，因为你们记得\Person{波斐留}，正是在我大量引用的那些书中，即论述灵魂回归的那些书中，屡屡教导说，必须脱离各种身体，好使灵魂能与天主同在福乐中。但在此处，你们不应追随波斐留，反倒应当纠正他，尤其因为你们和他一样相信那些有关这有形可见的世界及其庞然物质构架的灵魂的不可思议之事。因为，作为\Person{柏拉图}的学者，你们主张世界是一个有生命的存在，且是极为幸福的有生命的存在，你们还希望它也是永恒的。那么，若为灵魂的幸福，必须抛弃身体，世界又怎能永不脱离身体却又不失其幸福呢？太阳与其他星辰，你们不仅承认它们是物体，对此所有明眼人都一致同意，而且，为了遵循你们所认为的更深洞见，你们宣称它们是极有福的生命体，且连同他们的身体一起是永恒的。那么，当基督教信仰向你们提出时，你们为何忘记或佯装不知你们常常讨论或教导的事？你们为什么拒绝成为基督徒，理由竟是你们持有你们自己实际上在推翻的观点？这岂不是因为\Person{基督}在\Emph{谦卑}中来临，而你们\Emph{骄傲}吗？圣徒们复活后的身体的精确性质，有时会在那些最精通基督教圣经的人中间引起讨论；然而在我们当中，丝毫不会怀疑它们将是永恒的，且其性质以基督复活后的身体为范例。但无论它们的性质如何，既然我们主张它们绝对是不朽的、不死的，且不会为灵魂的默观——借此灵魂固定于天主——带来任何阻碍，而你们又说，在诸天界中，恒福者的身体是永恒的，那么，你们为何主张，为了福乐，必须脱离一切身体呢？你们为何这样寻找一个貌似合理的理由来逃避基督教信仰呢？若不是因为，我再重复一遍，基督是谦卑的，而你们是骄傲的。你们羞于被纠正吗？这是骄傲者的恶习。确实，对博学之士而言，从柏拉图的学园转至成为基督的门徒，是一种降格；基督曾藉祂的\OriginalTerm{圣神}{Holy Spirit}使一个渔夫思想并说出：\BibleQuote{在起初已有圣言，圣言与天主同在，圣言就是天主。圣言在起初就与天主同在。万物是藉着祂而造成的；凡受造的，没有一样不是由祂而造成的。在祂内有生命，这生命是人的光。光在黑暗中照耀，黑暗决不能胜过祂。} \BibleRef{若 1:1} 年高的圣人\Person{辛普利齐亚努斯}，后来成为\Place{米兰}的主教，曾告诉我，有位柏拉图主义者惯常说，这篇题为\WorkTitle{若望福音}的圣福音开头部分，应当以金字书写，并悬挂在各处教会中最显眼的地方。但骄傲者不屑以天主为师，因为\BibleQuote{圣言成了血肉，寄居在我们中间。} \BibleRef{若 1:14} 因此，对这些可怜人而言，他们有病还不够，还要夸耀其病，且羞于能治愈他们的良药。这样做，他们得到的不是高升，而是更灾难性的堕落。\par

% structure: p0069 source=h2 target=subsection reason=relative-heading-rank; base=h2

\subsection{第三十章 \Person{波斐留}对柏拉图主义的修正与改变}

倘若认为更动\Person{柏拉图}所涉及的任何事都是不恰当的，那么\Person{波菲利}自己又为何进行修订，而且为数不少呢？因为确凿无疑的是，柏拉图曾写道：人的灵魂死后会回到畜类的身体里。\Person{波菲利}的老师\Person{普罗提诺}也持这种意见；然而波菲利却公正地拒绝了它。他主张，人的灵魂固然回到人的身体，但不是回到它们离开的那个身体，而是进入另一些新的身体。他回避了另一种意见，以免一个转生为骡子的女人可能将自己的儿子驮在背上。然而，他却未回避另一种理论，这种理论承认：一位母亲可能转生为少女，并嫁给自己的儿子。那由神圣而真实的众天使所教导、由受天主圣神感动的先知所宣告、并由那位被祂的先驱预报者所预言为即将来临的救主以及祂所派遣且以福音充满全世界的宗徒们所宣讲的信经，是多么更为可敬啊！我再说，灵魂一次而永远地回到自己原本的身体，较之反复进入不同的身体，这信仰又是多么更为可敬！尽管如此，如我所说，波菲利确实对这种意见作出了相当大的改进，至少他主张人的灵魂只能转入人的身体，并且毫不迟疑地拆毁了柏拉图想将灵魂投入其中的那些畜类牢狱。他还说，天主将灵魂放入世界，是为了让它认识物质的邪恶，然后回归圣父，并永远摆脱与物质的污染性接触。虽然这里有些想法并不恰当（因为灵魂之被赋予身体，更是为了行善；因为它若不实际作恶，便不会学到恶），但他仍纠正了其他柏拉图主义者的意见，而且是在一个绝非不重要的观点上，因为他宣称，那已被涤除一切邪恶、蒙接纳到圣父面前的灵魂，将永不再遭受此生的种种苦害。借着这种意见，他彻底推翻了柏拉图所心爱的那个信条：如同活人变成死人，死人也要变成活人；并且他摧毁了\Person{维吉尔}似乎从柏拉图那里采纳的观念：那些被遣送到\Place{乐土}（对有福者之乐的诗意称呼）的净化了的灵魂，被召唤到\Place{勒忒河}，即遗忘之河，以遗忘过去，\par

好让他们再次返回尘寰，\newline{}不记生前种种，\newline{}带着盲目的欲求，\newline{}渴望回到血肉之身。\par

这种看法未获波菲利的认可，而且极其公正；因为相信灵魂竟渴望从那若非确信其恒久便不能真正有福的生命中返回，回到此生，回到可朽身体的污染中，仿佛完全净化的结果竟只是让污秽变得可欲，这确是愚蠢的。因为如果完全的净化导致遗忘所有邪恶，而对邪恶的遗忘又产生对身体的渴望——灵魂在这身体中将再度与邪恶纠缠——那么，至高的幸福将成为不幸的原因，智慧的圆满将成为愚昧的原因，最纯净的洁净将成为污秽的原因。而且，无论灵魂的福乐持续多久，如果为了有福就必须受骗，那么这福乐就不可能建立在真理之上。因为灵魂若不脱离恐惧，便不能有福。但是，要脱离恐惧，它就必须处于一种虚假的印象之下，以为它将永远有福——这印象是\Emph{虚假的}，因为它注定在某个时候也会受苦。那么，其欢乐建立在虚谎之上的灵魂，又怎能因真理而喜乐呢？波菲利看到了这一点，因此说：净化后的灵魂回归圣父，从此永远不再被邪恶的污染性接触所缠扰。因此，某些柏拉图主义者所主张的、认为有一种必然的循环带动灵魂离开并又将其带回同样情境的意见，是错误的。然而，即便这种看法是真的，知道它又有什么益处呢？难道柏拉图主义者敢于声称他们优于我们，只是因为我们在此生中不知道那些事——而他们自己在那另一个更美好的生命中，当他们在纯洁和智慧上达到圆满时，也注定会不知道这些事，并且他们若要有福，就必须不知道这些事？如果说这种话是极其荒谬愚蠢的，那么无疑，我们应当优先采纳波菲利的意见，而非那种主张灵魂在持续交替的幸福与痛苦中循环的观点。倘若这是正确的，那么这里便有一位柏拉图主义者正在修正柏拉图，这里便有一个人看到了柏拉图未曾看到的，而且他毫不退缩地纠正如此卓越的一位导师，宁取真理而舍柏拉图。\par

% structure: p0073 source=h2 target=subsection reason=relative-heading-rank; base=h2

\subsection{反驳柏拉图主义者用以主张人的灵魂与天主同为永恒的那些论据}

那么，我们为何不更相信那些人的才能无法测透的事理中的神圣启示呢？我们为何不信神圣的断言：灵魂并非与天主\OriginalTerm{同永恒}{co-eternal}，而是受造的，并且曾有一时不存在呢？因为柏拉图主义者似乎为自己找到了充足的理由来拒绝这项教义，他们断言凡非一直存在者便不能是永恒的。然而，柏拉图在论及世界以及其中由至高者所造的神祇时，极为明确地说它们有开始，却无终结，而是依造物者的至高旨意永远持续。但为了解释这话，柏拉图主义者发现他所说的开始并非时间的开始，而是原因的开始。他们说：\Quote{犹如一只脚从永恒就在尘土中，那底下便总有一个脚印；无人会怀疑这脚印是由脚的压力造成的，也不会怀疑虽然脚印由脚造成，二者却无先后之别；因此，}他们说，\Quote{世界和其中被造的神祇也一直存在，它们的造物者也一直存在，然而它们却是被造的。}那么，如果灵魂一直存在，我们是否要说它的悲惨也一直存在呢？因为倘若灵魂内有什么不是自永恒而有，却在时间中开始的，为何不可能灵魂本身虽先前不存在，却在时间中开始存在呢？同样，它的\OriginalTerm{真福}{blessedness}——正如他所承认的，在灵魂经历诸恶之后将更为稳固，且确实无尽——无疑在时间中有开始，却要永远存在，尽管先前并无此物。因此，那想要证明凡无始者外无一物可无尽的整个论证便站不住脚了。因为在此我们见到灵魂的真福有开始，却无终结。所以，让人之无能屈服于天主的权威；让我们从永福的\OriginalTerm{神体}{spirits}那里接受关乎真宗教的信仰，这些神体不为自己寻求那他们知道只当归于他们的天主、也是我们的天主的尊荣，并命令我们只向祂献祭——正如我已多次说过且必须再说的，我们和他们应当一同成为祭献，藉着那位在祂所取的人性中，并依此人性祂愿意成为我们的司铎，将自己奉献于死亡，为我们作了祭献。\par

% structure: p0075 source=h2 target=subsection reason=relative-heading-rank; base=h2

\subsection{论灵魂得救的普遍之道——\Person{波斐理}未能寻得，因其寻求不正，惟有基督的恩宠将其敞开。}

这就是那拥有拯救灵魂之\OriginalTerm{普世道路}{universal way}的宗教；因为除了这条道路，无人能得救。这是一种\OriginalTerm{君王之路}{royal way}，唯有它导向一个不像一切现世尊荣那样动摇、而是坚立于永恒根基之上的国度。当\Person{波菲利}于其著作\WorkTitle{灵魂的回归}第一书结尾处说，无论是来自至真的哲学，还是来自\Place{印度}人的观念与修行，或是来自\Place{加色丁}人的推理，或是来自任何来源，都未曾有一套教授体系，能提供拯救灵魂的普世道路，且历史文献的阅读也未使他认识那条道路时，他明确承认存在这样一条道路，但他尚未认识它。他如此辛勤学到的一切，关于灵魂得救之事，无一令他满意；凡他在别人看来（即便不是在自己看来）似乎知道并信服的，也无一令他满足。因为他觉察到，在如此重大的事上，仍缺少一种有权柄的、正宜于遵循的权威。当他说他未从任何至真的哲学中学到一套拥有灵魂得救之普世道路的体系时，在我看来，这足够清楚地表明，他所师从的哲学要么不是至真的，要么并未包含这样一条道路。那不具备这条道路的哲学，怎能是至真的呢？因为拯救灵魂的普世道路，无非是所有灵魂都藉它普遍得救，且若无它便无一灵魂能得救的那条道路。并且，当他补充说道：\Quote{或是来自\Place{印度}人的观念与修行，或是来自\Place{加色丁}人的推理，或是来自任何来源，}他以最清晰的语言宣告，他从印度人或加色丁人那里所学到的，并不包含这条拯救灵魂的普世道路；然而他却不免声称，正是从加色丁人那里，他获得了那些他如此频繁提及的\OriginalTerm{神谕}{divine oracles}。那么，他所说的这条拯救灵魂的普世道路究竟指什么？它尚未被任何至真的哲学，或被那些拥有洞悉神事之卓越见识的民族（因他们更沉溺于对天使的好奇及幻想性的学问与崇拜中）的教义体系所显明。他承认自己不知的这条普世道路，若非那不属于任何民族专有、却是为万民共有、且由天主所赐的道路，又能是什么呢？\Person{波菲利}，一位才能非凡的人，并不怀疑这样一条道路存在；因为他相信，天主的上智安排不可能使人类缺乏这条拯救灵魂的普世道路。因为他并未说这条道路不存在，而是说这一伟大的恩赐与援助尚未被发现，且未曾为他所知。这不足为奇；因为\Person{波菲利}生活在一个时代，那时这条拯救灵魂的普世道路——换言之，基督宗教——正遭受拜偶像者、魔鬼崇拜者以及世俗统治者的迫害，为使殉道者或真理的见证人数目得以圆满并被祝圣，并藉他们证明：我们必须为神圣信德的缘故，并为真理的举荐，忍受一切肉身的苦难。\Person{波菲利}目睹这些迫害，便断定这条道路注定迅速消亡，因此它并非拯救灵魂的普世道路，却未看见那使他动摇并阻止他成为基督徒的这件事本身，反而有助于我们宗教的确立及更有效力的举荐。\par

因此，这就是灵魂获得拯救的\OriginalTerm{普世之路}{universal way}，是天主慈悲为普世万民所赐予的道路。已经认识或将要认识此路的民族，都不应问：为何如此快？或为何如此迟？因为派遣此路的那一位之计划，非人力所能测透。\Person{波尔菲里}感觉到了这一点，因此他只说，天主的这份恩赐尚未被他接受，也尚未达到他的知识。因为尽管情况如此，他并未因此断定这道路本身不存在。我说，这就是信者得救的普世之路，忠信的\Person{亚巴郎}曾得到天主的保证，\BibleQuote{地上万民都要因你的后裔蒙受祝福。}\BibleRef{创 22:18}他本是\Place{加色丁}人；但为了使他能接受这些伟大许诺，并从他繁衍出一个后裔，\BibleQuote{是藉天使，经过一位\OriginalTerm{中保}{Mediator}之手而设立的，}\BibleRef{迦 3:19}在这中保身上，便可找到那条为拯救灵魂而向万民敞开的普世之路；为此，他被命令离开自己的故乡、亲族和父家。于是他自己首先脱离了加色丁人的迷信，藉着服从，敬拜了独一真天主，忠信地信赖了天主的许诺。这就是普世之路，神圣的预言论及此路时说：\BibleQuote{愿天主怜悯我们，降福我们，以自己的慈容光照我们；为使世人认识你的道路，万民得知你的救恩。}因此，许久以后，我们的救主从亚巴郎的后裔中取了肉躯，论及自己说：\BibleQuote{我是道路、真理、生命。}\BibleRef{若 14:6}这就是普世之路，很久以前就预言过：\BibleQuote{到了末日，上主的圣殿山必要矗立在群山之上，超乎一切山岳，万民都要向它涌来。将有许多民族前去，说：来，我们攀登上主的圣山，往雅各伯天主的殿里去！祂必指示我们祂的道路，教给我们循行祂的途径，因为法律将出自熙雍，上主的话将出自耶路撒冷。}\BibleRef{依 2:2}因此，这条道路不属某个民族，而属于万民。上主的法律和圣言并未停留在\Place{熙雍}和\Place{耶路撒冷}，而是从那里发出，传遍普世。因此，中保本人在复活后，对惊惶的宗徒们说：\BibleQuote{我还与你们同在的时候，曾对你们说过：梅瑟法律、先知书和圣咏上，关于我所记载的一切，都必须应验。耶稣遂开启他们的明悟，叫他们理解经书；又对他们说：\InnerQuote{经上曾这样记载：基督必须受苦，第三天从死者中复活；并且必须从耶路撒冷开始，因祂的名向万邦宣讲悔改，以得罪之赦。}}\BibleRef{路 24:44}这就是灵魂得救的普世之路，昔日，圣天使和圣先知们在少数蒙受天主恩宠的人中，尽可能地揭示了它，尤其在希伯来民族中。这个民族的国体，通过会幕、圣殿、司祭职和祭献，就像是被祝圣，用以预表和预告那个将要从万民中聚集而成的\WorkTitle{天主之城}。有时藉清晰的言辞，有时藉许多隐晦的预像，这道路被宣示出来；但后来，中保亲自取了肉躯，祂的圣宗徒们也来了，他们揭示了新约的恩宠如何更公开地说明那隐晦暗示给前代的事，这符合人类各时代的关系，也是天主凭其智慧乐意安排的；天主又以奇迹异事为他们作证，其中一些我已提及过。因为不仅出现了天使的异像，听见了天上使者的言语，而且天主的人，以纯朴虔敬的言语武装起来，从人的身体和感官中驱逐不洁之神，治愈残疾和疾病；地上的走兽、海中的生物、空中的飞鸟、无灵之物、自然元素、星辰，都服从他们出自天主的命令；地狱的势力在他们面前退避，死者复活。我不提救主本人特有的奇迹，尤其是祂的诞生和复活；在诞生中，祂只造成了童贞女生育的奥秘；而在复活中，祂提供了一个所有人在最后都要经历的复活的实例。这道路净化整个的人，使凡人每一部分都为不死不灭做好准备。因为，为防止我们为\Person{波尔菲里}所说的理智部分寻求一种净化，为他所说的神魂部分寻求另一种净化，又为肉体本身寻求另一种净化，我们最强大、最真实的净化者和救主，承担了整个人性。若不藉着这条道路——它曾在许诺时期和宣布许诺实现时期常与世人同在——没有人曾得救，没有人现在得救，将来也没有人能得救。\par

关于\Person{波菲利}声称\OriginalTerm{灵魂得救的普遍途径}{universal way of the soul's deliverance}还没有通过任何他所知的历史而为他所知晓，我要问，有什么历史比那以权威的声音占据全世界的历史更值得注意呢？或者，有什么比那既以同样的清晰叙述过去的事件，又预言未来，并且其中尚未应验的预言因那些已经应验的而使我们不得不相信的历史更可信呢？因为无论\Person{波菲利}还是任何\Person{柏拉图}主义者，都不会藐视\OriginalTerm{占卜和预言}{divination and prediction}，即使是关乎今生和尘世的事物，尽管他们公正地藐视\OriginalTerm{普通的占卜}{ordinary soothsaying}和与\OriginalTerm{邪术}{magical arts}相关的占卜。他们否认这些是伟人们的预言，或认为它们重要，他们是对的；因为这些预言要么是基于对\OriginalTerm{次要原因}{subsidiary causes}的先见，如同对专业的眼光来说，通过某些预兆症状就能预见疾病的大部分过程，要么是\OriginalTerm{不洁的魔鬼}{unclean demons}预言他们决定要做的事，以便他们用权威的表象影响恶人的思想和欲望，并诱使人类的软弱去模仿他们不洁的行为。行走在普遍途径上的圣人并不关心预言这类事情，认为它们重要，尽管为了举荐信仰，他们知道甚至常常预言那些不能由人的观察所发觉、也不易由经验证实的事情。但他们还预言了其他真正重要和神圣的事件，只要他们得以知道天主的旨意。因为\OriginalTerm{基督的降生成人}{incarnation of Christ}，以及在祂内并奉祂的名所成就的一切重要奇迹；人的\OriginalTerm{悔改}{repentance}和意志归向天主；罪过的赦免，\OriginalTerm{成义的恩宠}{grace of righteousness}，虔诚者的\OriginalTerm{信德}{faith}，以及世界各地相信\OriginalTerm{真天主性}{true divinity}的众多群众；偶像崇拜和魔鬼崇拜的倾覆，以及信友通过试探的考验；坚忍者的净化，以及他们脱离一切邪恶；\OriginalTerm{审判的日子}{day of judgment}，\OriginalTerm{死者复活}{resurrection of the dead}，\OriginalTerm{恶人的团体}{community of the ungodly}的\OriginalTerm{永罚}{eternal damnation}，以及最荣耀的\OriginalTerm{天主之城}{city of God}的\OriginalTerm{永恒王国}{eternal kingdom}，这城在\OriginalTerm{享见天主}{vision of God}中永享福乐——这些事都在这途径的圣经中被预言和应许了；我们看到其中许多已经应验，以至于我们公正而虔诚地相信其余的也必将发生。至于那些不相信、因此也不明白这就是按照圣经的真实预言和陈述，直通享见天主并与祂永恒共融的途径的人，他们可以攻击我们的立场，但他们无法攻破它。\par

因此，在这十卷书中，虽然我敢说没有满足某些人的期望，但靠真天主和主的助佑，我通过驳斥那些不虔敬之人的异议——他们喜爱自己的神胜过\OriginalTerm{圣城的建立者}{Founder of the holy city}，也就是我们着手讨论的那位——来满足了某些人的愿望。在这十卷书中，前五卷是针对那些认为我们应该为着\OriginalTerm{今生的福乐}{blessings of this life}而敬拜诸神的人，后五卷是针对那些认为我们应该为着\OriginalTerm{死后的生命}{life after death}而敬拜诸神的人。现在，为了履行我在第一卷中所作的承诺，我将依靠天主的助佑，继续讲述我认为关于\OriginalTerm{两座城}{two cities}的\OriginalTerm{起源、历史和应得的结局}{origin, history, and deserved ends}所须说的一切；正如已经指出的，这两座城在现世中交错混杂在一起。\par

\SourceRecord{Translated by Marcus Dods.；Nicene and Post-Nicene Fathers, First Series；Vol. 2.；Edited by Philip Schaff.；Buffalo, NY: Christian Literature Publishing Co.,；1887.；<http://www.newadvent.org/fathers/120110.htm>.}

% structure: p0001 source=h1 target=section reason=page-title

\section{天主之城（卷十一）}

这里开始本作品第二部分，论述两座城——地上之城与天上之城——的起源、历史与归宿。首先，奥古斯丁在本卷中说明两城如何最初由善恶天使的分离而形成；并借此机会讨论世界的创造，正如圣经在\BibleBook{创世}{纪}卷首所描述的。\par

% structure: p0003 source=h2 target=subsection reason=relative-heading-rank; base=h2

\subsection{第一章——论本部分，其中我们开始解释两城的起源与终结。}

我们所说的天主之城，正是那圣经所为之作证的；这圣经以其神圣权威超越万国一切著作，并感化了各种心智，这并非出于偶然的智性活动，而显然是由于明确的眷顾安排。因为经上记载着：\BibleQuote{天主之城啊，有荣耀的事乃指着你说的。}又在一首圣咏中我们读到：\BibleQuote{主本为大，在我们天主的城中，在祂的圣山上，当大受赞美。祂使全地欢喜。}稍后，在同一篇圣咏中：\BibleQuote{我们在万军之主的城中，在我们天主的城中，所听见的，正如我们所看见的。天主必坚立这城，直到永远。}在另一篇中：\BibleQuote{有一道河，这河的分汊，使天主的城欢喜；这城就是至高者帐幕的圣所。天主在其中，城必不动摇。}从这些及类似的见证（若一一引述则过于冗长），我们得知有一座天主之城，其建造者在我们心中激起了渴慕成为其中公民的爱。对这圣城的建造者，地上之城的公民却偏爱他们自己的神祇，不认识祂乃是万神的天主——不是虚假之神的天主，\Emph{即}不虔敬而骄傲之神的天主，这些神被剥夺了祂不变而慷慨分施的真光，因而沦入一种贫乏的能力，贪婪地攫取私有的特权，向受欺骗的属民寻求神圣尊崇——乃是虔敬圣洁之神的天主，这些神更乐于服从独一者，而非使众多者服从自己，并且宁愿敬拜天主，而不愿被当作天主敬拜。对于这城的仇敌，我们在前十卷书中，已按我们的能力并赖我们的主与君王的助佑，作了答复。如今，认清了所期待于我的，谨记我的许诺，也仰赖同样的扶助，我将尽力论述两城（即地上之城与天上之城）的起源、进展与应得的归宿；这两城，如我们所说，在现今的世界里相互混杂，仿佛彼此纠缠在一起。首先，我要说明这两城的基础最初如何奠基在天使们之间产生的分歧上。\par

% structure: p0005 source=h2 target=subsection reason=relative-heading-rank; base=h2

\subsection{第二章——论对天主的认识，除经由天主与人之间的\OriginalTerm{中保}{Mediator}，亦是人而耶稣基督外，无人能臻至。}

人若默观了整个受造界，无论有形体的与无形体的，并辨识了其可变性之后，能超越其上，凭借心灵不断的飞升，达至天主不变的本体，并在那默观的高峰，从天主亲自得知唯独祂造了非天主本质的一切——这是一件伟大而极稀有的事。因为天主与人说话，并不是藉着某种可听的受造物在耳边喧响，以致大气振动，连接发声者与闻声者；也不是藉着有身体假象的神体，一如我们在梦境或类似状态中所见的；即使在这种情况，祂也仿佛对身体之耳发言，因为祂是藉着身体的假象说话，且有真实空间间隔的表象——因为异象正是有形物体的逼真再现。天主说话并非藉着这些，而是藉着真理本身，只要有人准备好以心灵而非以身体来聆听。因为祂是对人内优于其余一切、且唯独天主自己胜过它的那一部分说话。既然人最恰当的理解（或若不能理解，至少\Emph{相信}）是按照天主的肖像受造，无疑正是人用以超越与野兽共有的低级部分的那一部分，使他更靠近至高者。可是心灵本身，虽在本性上能理性和理智，却被愚钝而根深柢固的恶习所困，不仅不能乐于并安居于祂不变的光中，甚至不能忍受它，除非逐步获得医治、更新，堪当承受此等福乐，为此心灵必须先接受信德的浸润，从而得以净化。又为能在信德中更坚定地向真理迈进，真理本身，天主、圣子，摄取人性而不失其天主性，奠定并建立了这信德，好能经由一位\OriginalTerm{天主而人者}{God-man}，为人开辟一条通向人之天主的道路。因为这正是天主与人之间的\OriginalTerm{中保}{Mediator}，亦是人而耶稣基督。正因为祂是人，所以祂是中保，又是道路。若在行路者与他要去的地方之间有一条道路，便有希望抵达；倘若没有道路，或他不知道路何在，知道应去何处又有何益？如今，唯一确然无误、远离一切错谬的道路，便是那同一位既是天主又是人者，天主是我们的终向，人则是我们的道路。\par

% structure: p0007 source=h2 target=subsection reason=relative-heading-rank; base=h2

\subsection{第三章——论由\OriginalTerm{圣神}{Divine Spirit}所著之\OriginalTerm{正典圣经}{Canonical Scriptures}的权威。}

这位\OriginalTerm{中保}{Mediator}，先借着先知，后亲口发言，再后又借着宗徒，说出了他认为足够的话，此外还产生了被称为\OriginalTerm{正典}{canonical}的圣经，它具有至高无上的权威，我们在一切不应无知、却又无法靠自己得知的事上，都对此圣经予以认可。因为如果我们通过自己内外感官的见证来认识眼前的对象，那么对于远离我们感官的对象，我们就需要别人提供他们的见证，因为我们不能凭自己的感官认知它们，于是我们便相信那些曾或正亲身经历这些对象的人。因此，就如对于未见过的可见之物，我们相信见过的人（对于一切可感的事物也是如此），同样，对于那些由心灵和理性所察觉的事物，\Emph{即}，那些远离我们内在感官的事物，我们理当相信那些曾在无形的真光中看见它们、或常久默观它们的人。\par

% structure: p0009 source=h2 target=subsection reason=relative-heading-rank; base=h2

\subsection{第四章 世界既非无始，亦非由天主新的定旨所造，祂并非后来才愿欲先前所不愿欲者。}

在一切可见之物中，世界是最伟大的；在一切不可见之物中，最伟大的是天主。我们看到，世界是存在的；我们相信，天主是存在的。天主创造了世界，我们最可靠的是从天主自己那里相信。可是我们曾在何处听闻祂呢？最明确的莫过于在圣经中，祂的先知说：\BibleQuote{在起初天主创造了天地。}\BibleRef{创 1:1} 先知是否在天主创造天地时在场呢？不；但是天主的智慧，那创造万有的智慧，却在场\BibleRef{箴 8:27}，智慧渗入圣善的灵魂，使他们成为天主的朋友和祂的先知，并无声地告诉他们祂的工程。他们也由天主的天使施教，这些天使常瞻仰圣父的面容\BibleRef{玛 18:10}，并向适当的人宣示祂的旨意。这些先知中的一位曾说并写道：\BibleQuote{在起初天主创造了天地。}他是如此适合作天主的见证人，以致将这一切启示给他的同一圣神，也使他能在如此久远之前就预言我们的信仰也将到来。\par

然而，天主为何要在那时选择创造天地，而在此之前祂并未创造呢？若提出此问题的人想表明世界是永恒的、无始的，因而并非由天主所造，那他们便受了大骗，在不可救药的邪恶疯狂中胡言乱语。因为即便先知的声音沉默，世界本身却以其井然有序的变化运行、以及一切可见之物的美好面貌，自身作证，既证明它受造，也证明它只能由天主所造——天主的伟大与美善，是不可言喻、不可见的。至于那些虽然承认世界由天主所造，却将之归因于非时间性的、仅为受造性的开端的人，他们想以此在某种难以理解的方式上，说明世界始终是一个被造的世界而存在；他们这样主张，似乎是为天主辩护，免得祂被指控为任意匆忙，或是突然以全新的想法创造世界，或是随意改变祂的旨意，尽管祂是永恒不变的。但我不明白这种假定在其他方面——尤其关于灵魂方面——如何能成立。因为如果他们主张灵魂与天主同样永恒，那么他们将完全无法解释，为何灵魂会遭受新的不幸，而这不幸在以往无尽的永恒中并未存在过。因为倘若他们说，灵魂的福与祸不停交替，那么他们必须进一步说，这种交替将永远持续；由此产生荒谬的结果：灵魂虽被称为有福，却并非真福，因为它预见到自己的不幸与屈辱。然而，如果灵魂并未预见，且自以为将永不受辱、永不离苦，而常享有福乐，那么它之所以有福，乃是由于受了欺骗；再没有比此更愚昧的说辞了。如果他们又设想，在以往永世中，灵魂的不幸与其福乐交替，但如今，一旦灵魂获得释放，便永不再归于不幸，那么他们仍然认为灵魂以往从未真正有福，而今终于开始享受一种新而不确定的幸福；就是说，他们必须承认，灵魂遭受了某种新事，且是重要而显著的事，这是在以往整个永世中从未发生过的。如果他们否认天主的永恒计划包含灵魂的这种新经验，那就是否认天主是其福乐的作者，这是无可言喻的不敬虔。反之，如果他们声称灵魂未来的福乐源于天主的新定旨，那么他们又怎能说明天主不受他们所不悦的那种可变性之控诉呢？此外，如果他们承认灵魂是在时间中受造，但绝不随时间而消亡——它如同数字，有始无终——既然如此，灵魂既曾经历不幸并从中得救，便永不再重返不幸；他们必定会承认，这一切的发生绝不违背天主的永恒旨意。那么，就让他们同样相信世界吧：世界亦可在时间中受造，而天主在创造时，并未改变祂永恒的谋划。\par

% structure: p0012 source=h2 target=subsection reason=relative-heading-rank; base=h2

\subsection{第五章 我们不当寻求理解世界之前的无限时间，亦不当寻求理解无限空间。}

接下来，我们必须看看，对于同意天主是世界的创造者，却对其创造的时间有疑难的那些人，能作何答复；同时，对于我们就创造的地点可能提出的疑难，他们又能作何答复。 因为，正如他们追问世界为何在那个时候而非更早被造，我们也可以问，世界为何恰恰在此处而非别处被造。 如果他们想象世界之前有无限的时间，其间天主不可能无所作为，同样，他们也可以设想世界之外有无限的空间领域；如果有人说，全能者不能停手不作工，那么，他们岂不必然要采用\Person{伊壁鸠鲁}那无数世界的梦想吗？只不过有这样的分别：他主张那些世界是由原子的偶然运动形成和毁灭的，而他们却坚持那些世界是由天主之手造成的，倘若他们认为，在围绕世界向各方无限伸展的无边无际的空间里，天主不能安息，并且他们假设祂所造的世界也不能毁灭。 因为这里的问题是与那些和我们一样相信天主是精神体、除祂以外一切存在之创造者的人有关的。至于其他的人，与他们争论宗教问题乃是屈尊俯就，因为他们只在那些敬拜众多神明的人中间获得了一点名声，而在其他哲学家中显得突出，仅因他们虽离真理尚远，却比其余人更接近真理而已。 那么，这些人既然不在任何地点限制、约束或分割\OriginalTerm{天主的性体}{divine substance}，而是按天主所配得的，承认它完全地、虽以精神的方式临在于各处，他们或许会说，这性体并不临在于世界之外如此广大的空间中，而只占据其中一点——与世界之外的无限相比，微乎其微的一点——即世界所在的这一处吗？我想他们不至于荒唐至此。 他们既然主张只有一个世界，体积固然广大，却是有限的，有其确定的位置，并且是由天主的工作造成的，那么，让他们对于世界之前无限时间里天主的安息，给出与他们对世界之外无限空间中天主的安息同样的解释吧。 正如我们不能因为世界被安置在它所占的那一地点而非别处，就推论说这是出于偶然而不是由于天主的理据——尽管人的理性无法理解为何如此安置，而且所选的地点并无任何功绩能使它优先于其他无限的地点——同样，我们也不应以为，天主在那时刻而非更早创造世界是出于偶然，尽管在无限的过去里，先前的时间一直在流逝，而且没有任何差异使某一时刻能优先于另一时刻被选中。 但是，如果他们声称，人设想无限空间是徒劳的，因为世界之外别无空间，我们便答说，根据同样道理，设想天主在世界之前的安息时间也是徒劳的，因为世界之前并无时间。\par

% structure: p0014 source=h2 target=subsection reason=relative-heading-rank; base=h2

\subsection{第六章. — 世界与时间同有一始，非一先一後}

因为如果永恒与时间的正确区分在于：时间没有某种运动和变迁便不存在，而在永恒中毫无变化，那么谁看不到，除非有某种受造物被造，能以某种运动产生变化——此种运动和变化的各部分不能同时存在，而是彼此相继——并且，在这些较短或较长的延续间隔中，时间便开始？ 因此，天主，在其永恒中毫无变化，是时间的创造者和安排者；我看不出，如何能说祂是在一段时间过去之后才创造世界，除非说在世界之前已经有一些受造物，藉其运动时间才得以流逝。 如果神圣而无误的圣经说：\BibleQuote{在起初天主创造了天地}，为的是使人理解祂之前没有造任何东西——因为如果祂在别的之前造了什么，那东西更应说是\Quote{在起初}受造的——那么，毫无疑问，世界不是在时间内，而是与时间同时被造的。 因为在时间内被造的东西，是在某段时间之后和之前被造的——在已过的时间之后，在未来的时间之前。但那时不可能有已过的时间，因为没有受造物，其运动可用来度量持续。 然而，世界是与时间同时被造的，如果在世界的创造中，变化与运动均被造了，这从起初六日或七日的次序看来是明显的。 因为在这些日子里，有晚上，有早晨，被计算着，直到第六日，天主那时所造的一切都完成了，在第七日，天主的安息被神秘而崇高地标志出来。 这些日子究竟是什么样的日子，我们极难、或许不可能构想，更遑论言说了！\par

% structure: p0016 source=h2 target=subsection reason=relative-heading-rank; base=h2

\subsection{第七章. — 论未有太阳之时，已有早晚的最初几日的性质}

我们确实看到，我们平常的日子只有太阳落山才有傍晚，只有太阳升起才有早晨；但起初的三天却是在没有太阳的情况下度过的，因为据记载，太阳是在第四天造成的。首先，光由天主的话造成，我们读到，天主将光与黑暗分开，称光为昼，称黑暗为夜；但那是何种光，又借着怎样的周期性运行形成了傍晚与早晨，却非我们感官所能企及；我们既不能理解那光如何，却必须毫不迟疑地相信。因为它要么是某种物质的光，或是从世界高处发出，远非我们视力所及，或是从后来太阳被点燃的那地方发出；要么以光之名被象征的是圣城，由圣天使和有福的灵所组成，关于这城，宗徒说：\BibleQuote{那在天上的耶路撒冷是我们永远的母；}\BibleRef{迦 4:26} 在另一处又说：\BibleQuote{你们都是光明之子，白昼之子；我们不属于黑夜，也不属于黑暗。}\BibleRef{得前 5:5} 然而，在某种方式上，我们也可恰当地说这日子有早晚。因为受造物的知识，相比于造物主的知识，不过犹如微明；因此，当受造物被引向对造物主的赞美与爱时，它便破晓而成为早晨；而当受造物不因爱恋受造之物而离弃造物主时，黑夜便绝不降临。总之，圣经依次记述那些日子时，从不提及夜这个词。它从不说\BibleQuote{有夜晚}，而是说\BibleQuote{过了晚上，过了早晨，这是第一天。}第二天和其余日子也是如此。的确，受造物本身被审视时的知识，比起在天主的智慧中——即在其被造的艺术中——观看它们时，可以说是更为暗淡。因此，以傍晚比喻比黑夜更为贴切；然而，正如我所说的，当受造物回归对造物主的赞美和爱时，早晨便会回归。当它如此行于对自身的认知时，这便是第一天；当行于对穹苍的认知时，即对那被称为在水上之水和下之水之间的天空的认知时，是第二天；当行于对大地、海洋和地上生出的一切的认知时，是第三天；当行于对大小光体及一切星辰的认知时，是第四天；当行于对所有在水中游动和在空中飞翔的动物的认知时，是第五天；当行于对所有生活在地上的一切动物及人自身的认知时，是第六天。\par

% structure: p0018 source=h2 target=subsection reason=relative-heading-rank; base=h2

\subsection{第八章。—— 对于天主在六日创造之后，在第七天安息，我们应如何理解。}

圣经说天主在第七天停止了一切工作，安息了，并圣化了这一天；我们不可幼稚地理解此事，好像工作对天主来说是一种劳苦似的——祂\BibleQuote{一发命，事便成就}，藉着属灵的、永恒的，而非可听闻且短暂的话而说的。但天主的安息，指的是那些在天主内安息的人所享有的安息，正如房屋的喜乐意指那房屋中欢欣之人的喜乐，虽然带来喜乐的并非房屋，而是别的事物。那么，如果房屋本身因着自己的美丽而使居住者喜乐，此等用语岂不更加明了！在此情形下，我们不仅以那种以容器代表所容之物的修辞手法称其为喜乐（例如我们说\Quote{剧场喝彩}，\Quote{草地低鸣}，意指剧场中的人喝彩，草地上的牛低鸣），还以那种将原因当作结果的修辞手法（例如说一封信是喜乐的，因为它使读信的人欢喜）。因此，圣经叙述最为恰当地指出天主安息了，意指那些在祂内、且蒙祂使之安息的人得享安息。这预言的叙述也向其所言说、并为其而写的人应许：他们在天主通过他们并在他们内所做的那些善工之后，若已藉着信德在此生努力亲近了天主，就要在祂内享受永恒的安息。此事曾通过安息日法律所规定的安息，预表给天主的古时子民；对此，我将在适当之处更详细地加以讨论。\par

% structure: p0020 source=h2 target=subsection reason=relative-heading-rank; base=h2

\subsection{第九章。—— 关于天使的创造，圣经教导我们应相信什么。}

目前，我既已着手讨论圣城的起源，首先是圣天使，他们构成这城的一大部分，且确实是最有福的一部分，因为他们从未被放逐；我要在天主的助佑下，并尽可能相称地，致力解释与此相关的圣经。圣经论及世界的创造时，并未明确说明天使是否或何时被造；但若提及他们，则是隐含于\Quote{天}这一名称下，正如经上所说：\BibleQuote{在起初天主创造了天地，}或者更可能是隐含于\Quote{光}这一名称，关于光我们稍后再说。但若说天使完全被遗漏，我不敢相信，因为经上记载天主在第七天安息，歇了祂所作的工；而这部书本身也以\BibleQuote{在起初天主创造了天地，}为开头，故在天与地之前，天主似乎没有造什么。那么，祂既以天地开始——而地本身，如圣经接着所说，起初混沌空虚，光尚未造成，黑暗笼罩渊面（也就是笼罩着天地间一种模糊的混沌，因为无光之处，必然黑暗）——之后，当六天内所记载完成的一切皆被创造和安排妥当，天使又怎会被遗漏，仿佛他们不在安息日所停工的天主之工中呢？然而，虽然天使系天主所造这一事实在此并未省略，但确实未显式提及；不过，圣经其他地方以极明确的方式予以肯定。因为在三青年火窑赞歌中如此说：\BibleQuote{主的一切化工，请赞美主；}而在其后逐一列举的这些化工中，也提到了天使。圣咏中也说：\BibleQuote{请你们从天上赞美主，请你们在高处赞美祂！祂的众天使，请赞美祂，祂的众军旅，请赞美祂！太阳和月亮，请赞美祂，闪烁的星辰，请赞美祂！天上的天和天上的水，请赞美祂！愿它们赞美主的名，因祂一命令，它们便受造。}此处，天使最明确地，且凭天主的权威，被说成是由天主所造，因为在论及其他天上事物时，也论到他们说：\BibleQuote{祂一命令，他们便受造。}那么，谁敢说天使是在六日创造之后才被造的呢？若有人如此愚昧，另一具有同等权威的圣经便驳斥了他的愚昧，其中天主说：\BibleQuote{当星辰被造时，众天使曾同声赞颂我。}\BibleRef{约 38:7}因此，天使存在于星辰之前；而星辰是在第四天被造的。那么，我们能说他们是在第三天被造的吗？绝非如此；因为我们知道那一天造了什么。地与水分开，各元素各具其形，地上生长出一切植物。那么，是在第二天吗？也不是；因为在那一天，穹苍形成于上下水之间，并被称为\Quote{天}，而星辰于第四天被造在这穹苍中。那么，毫无疑问，如果天使被包括在这六天的天主之工中，他们就是那被称为\Quote{昼}的光，圣经强调其统一性，称那日不是\Quote{第一日}，而是\Quote{一日}。因为第二日、第三日及其余的并非别的日子；而是那同一\Quote{一日}重复出现，以凑足六或七之数，好使人既认识天主的工程，也认识祂的安息。因为当天主说：\BibleQuote{有光！就有了光}时，若我们有理据将这光理解为天使的受造，那么无疑，他们受造是为分享那永光，这光就是天主不变之智慧，万物藉它而造，我们称之为天主独生子；这样，他们被创造他们的光所光照，自己便成为光，并因分受那不变的光和日子——即圣言，他们及万物皆藉祂而造——而被称为\Quote{昼}。\BibleQuote{那普照每人的真光，正在进入这世界}\BibleRef{若 1:9}——这光也光照每个纯洁的天使，使他成为光，但不是在自己内，而是在天主内；若有天使背离了这光，他便成为不洁，如同一切所谓的不洁之神，不再是在主内的光，而是本身成为黑暗，因为被剥夺了分受永光。因为恶并无实在的本性；乃是善的缺失获得了\Quote{恶}之名。\par

% structure: p0022 source=h2 target=subsection reason=relative-heading-rank; base=h2

\subsection{十、论纯一不变的天主圣三，圣父、圣子、圣神，唯一的天主，在祂内实体与属性同一。}

因此，存在一种唯独单纯，因而唯独不变的善，这就是天主。所有其他事物都是藉着这个善被创造的，但它们不是单纯的，所以不是不变的。我说\Quote{受造}，是指被造，而非受生。因为那由单纯善所生的，就如它一样单纯，并且与它相同。我们称这两者为父与子；两者连同圣神，是一个天主；而这位圣神，在圣经中，可以说，特别被赋予\Quote{圣}的称号。祂与父和子有别，因为祂既不是父，也不是子。我说\Quote{有别}，而非\Quote{另一事物}，因为祂与他们同等，是那单纯的善，不变且同永恒。这个天主圣三是一个天主；并不因为是天主圣三而不那么单纯。因为我们不是说善的本质之所以单纯，是因为唯独父拥有它，或唯独子，或唯独圣神；我们也不像\OriginalTerm{撒伯里乌派异端者}{Sabellian heretics}那样，说天主圣三只是名义上的，没有真实的位格区分；而是说它是单纯的，因为它是它所拥有的，除了各个位格彼此之间的关系之外。因为就这关系而言，父有子，但自己不是子；子有父，但自己不是父，这是真实的。但就祂自己本身，不考虑与他者的关系，每个位格都是祂所拥有的；这样，祂在自己内是生活的，因为祂有生命，并且祂自己就是祂所拥有的生命。\par

正是由于这个缘故，所以天主圣三的本性被称为单纯的，因为它没有任何会失落的东西，而且它不是像杯子与液体、物体与其颜色、空气与其光或热、心灵与其智慧那样，本身是一回事，其所有物是另一回事。因为这些东西中，没有一个就是它所拥有的：杯子不是液体，物体不是颜色，空气不是光与热，心灵不是智慧。因此，它们可以失去它们所拥有的，并转变或改变成其他性质和状态，以致杯子可以倒空所盛的液体，物体可以褪色，空气可以变暗，心灵可以变得愚钝。但在复活中所应许给圣人的不朽身体，固然不能失去其不朽的性质，但身体的实体与不朽的性质并不是一回事。因为不朽的性质完整地存在于每一个部分中，在一个部分里不会更大，在另一个部分里不会更小；因为没有哪个部分比另一个部分更不朽。身体本身，整体确实大于部分；它的一个部分较大，另一个部分较小，然而较大的部分并不比较小的部分更不朽。因此，并非在每个部分中都是一个完整身体的身体，是一回事；而完全充满各处的不朽性，是另一回事——因为不朽身体的每一个部分，不管在别的方面与其他部分如何不同，都是同样不朽的。\Emph{例如}，手并不因为大于手指就更不朽；所以，虽然手指与手大小不同，它们的不朽性却是相等的。因此，虽然不朽性不可分离地属于不朽的身体，但身体的实体是一回事，不朽的性质又是另一回事。所以身体并不是它所拥有的东西。同样地，灵魂本身，即使它永远明智（就如它得救赎后将永恒如此一样），也将是通过分有那不变的智慧而成为明智的，而它本身并不是那智慧；因为，尽管空气永不失去散布在其中的光，它并不因此就是光。我并不是说灵魂就是空气，像某些无法设想精神性实体的人所曾想象的那样；而是说，尽管两者大不相同，却有一种相似性，这使我们适宜地说，那非物质灵魂为天主单纯智慧的非物质光明所照亮，就如物质的空气被物质的光所照耀一样；而且，正如空气一旦失去这光就变暗（因为物质的黑暗无非就是缺乏光的空气），同样，灵魂失去了智慧之光就变暗。\par

据此，那些本质和真正神圣的事物被称为单纯的，因为在它们中性质与实体是同一的，并且因为它们在自己内，无需外来补充，就是神圣的、智慧的、有福的。的确，在圣经中，\OriginalTerm{智慧之神}{Spirit of wisdom}被称为\Quote{多方面的}\BibleRef{智 7:22}，因为它包含许多事物；但它所包含的，它也就是，它既是单一的，也就是这一切。因为智慧并非多个，而是一个，在其中蕴藏着不可数算、无限的知识宝藏，其中包含着由它所创造的可视、可变事物的所有不可见、不变的\OriginalTerm{理}{reasons}。因为天主创造任何事物都非无知；即使一个人间的工匠也不能说是如此。但若祂知道祂所创造的一切，祂便只创造祂所已知的事物。由此得出一个十分惊人却真实的结论：这世界除非存在，否则不能为我们所知；但除非它为天主所知，否则便不能存在。\par

% structure: p0026 source=h2 target=subsection reason=relative-heading-rank; base=h2

\subsection{十一、堕落的天使是否分享了圣天使自受造以来一直享有的真福？}

既然事实如此，我们称为天使的那些神体，从来不曾是黑暗，而是在受造之时就立即成为光明；然而他们受造并非只是为了以任何方式存在和生活，而是蒙受光照，为能明智而享真福地生活。其中有一些背离了这光，未能获得这明智而真福的生命；这生命当然是永恒的，而且伴有对其永恒性的确切把握；但他们仍具有理性生命，尽管因愚昧而昏暗不明，而且他们即使想要也无法失去这生命。但谁能断定他们在堕落之前享有那智慧到了何等程度呢？我们又怎能说，他们同等分享了这智慧，以致像那些藉此智慧而真正并完全享有真福、安息于永恒真福的确切把握中的天使一样呢？因为如果他们同等地分享了这真知识，那么恶天使就会与善天使一样永远享有真福了，因为他们同等地期盼着它。因为一个生命不论多长，如果注定要终结，就不能真正称为永恒的；它之所以称为生命，是因它活着，而称为永恒，是因它没有终结。因此，虽然并非所有永恒之物都因此享有真福（因为地狱之火也是永恒的），然而如果除了永生，没有生命能真正且完全地享有真福，那么这些天使的生命便不是享有真福的，因为他们的生命注定要结束，因此不是永恒的，无论他们是否知道这一点。或是恐惧，或是无知，使他们无法享有真福。而且即使他们的无知不至于大到在他们内产生完全虚假的期望，却使他们摇摆不定，无法确定他们的善是永恒的还是会有一天终止，那么这种对如此伟大命运的怀疑，便与我们相信圣善天使所享有的圆满\OriginalTerm{真福}{blessedness}不相容。因为我们并不把\Quote{真福}一词的应用范围限制得那样狭窄，只用于天主；尽管无疑地，祂如此真正地享有真福，以至于不能有更大的真福；而与祂的真福相比，天使的真福又算什么呢？即使按他们的容量，他们已完全享有真福。\par

% structure: p0028 source=h2 target=subsection reason=relative-heading-rank; base=h2

\subsection{第12章——尚未蒙受神圣赏报的义人之真福与乐园中原祖父母之真福的比较}

天使并非理性与理智的受造物中，我们称为有福的仅有成员。因为谁敢否认乐园中的原祖父母在犯罪之前是有福的，尽管他们不确定这真福会持续多久，以及是否会永恒（假使他们没有犯罪，这真福就会是永恒的）——我说，谁愿这样做呢？因为即使在现今，我们看到那些度正义而圣善的生活、怀有不朽的盼望、没有锥心的良心悔恨、反而易于获得天主对自己现时软弱之罪的宽赦的人，我们称之为有福的，并非不合宜。这些人，虽然确信如果他们坚持到底就必得赏报，却不确定自己是否会坚持到底。因为谁能知道自己在修持和增长恩宠中会坚持到底呢？除非他从天主那里得到某种启示的保证；天主在祂正义而隐秘的审断中，不欺骗任何人，却将此事告知少数人。因此，就当前的安慰而言，乐园中的元祖比在此不安定状态中的任何义人都更为有福；但就未来福善的希望而言，任何不仅推测、而且确实知道自己将永享至高天主、与天使为伴、远离灾祸的人——这样的人，无论遭受何等肉身的痛苦，都比那位甚至在乐园极大的真福中却对自己的命运毫无把握的人，更为有福。\par

% structure: p0030 source=h2 target=subsection reason=relative-heading-rank; base=h2

\subsection{第13章——是否所有天使都受造于同一共同的\OriginalTerm{真福}{felicity}状态，以致堕落者并不知道他们会堕落，而站立者在堕落者败亡之后获得了自己坚持到底的保证}

从这一切，任何人都容易想到，一个有理智的存在者所切愿的、作为其正当对象的真福，是由这两件事的结合而生的：即不间断地享有那不变的善，也就是天主；并且脱离一切疑惑，确知自己将永恒地住在这同一享有中。我们虔诚地相信，光明的天使确实是如此；但论到那些因自己的过失而失落了那光明的堕落天使，即使在它们犯罪以前，也并不享有这种真福，这是理性命我们得出的结论。然而，若是它们在堕落以前确曾度过一段生命，我们便必须承认它们曾有过某种真福，尽管不是带有预知的那种。或者，如果难以相信：当天使们受造时，有些天使在受造时就对于自己的坚守或堕落一无所知，而另一些天使却对自己的永恒幸福确知无疑——如果难以相信它们从起初并非一律平等，直到如今那些恶天使凭其自由意志从善的光明堕落为止，那么，更难相信的是，圣善的天使现在对于自身的永恒真福尚不确定，并且对自己所知道的，还不如我们据圣经所能推断的那么多。因为，有哪个天主教徒不知道，善天使中决不会再产生新的魔鬼，正如他知道现今这个魔鬼决不会再回到善天使的团体中呢？因为福音中的真理向圣徒和信者应许说，他们将与天主的天使同等；并且也向他们应许了，他们将\BibleQuote{进入永生。}\BibleRef{玛 25:46}但若是我们确知自己永不会从永恒的真福中堕落，而他们却不确知，那么我们便不是与它们同等，而是高过它们了。然而，真理既绝不会欺骗人，我们将来又要与它们同等，所以它们必定确知自己的真福。而恶天使既不可能对此确知，因为它们的真福注定要终结，那么随而来的结论就是：要么天使原不平等，要么，若原是平等的，那么善天使是在其他天使灭亡之后，才确知其永恒真福的；除非，或许有人会说，主论魔鬼的话——\BibleQuote{他从起初就是杀人者，不居于真理之中，}\BibleRef{若 8:44}——应当这样理解：他不仅从人类之初（那时人受造，他能用诡计杀害人）就是杀人者，而且还从他受造之时起就不居于真理中，因而从不曾与圣天使一同有福；反而拒绝服从他的造物主，高傲地洋洋自得，仿佛拥有一种私有的主权，并如此自欺欺人。因为全能者的统治是无可逃避的；凡不愿虔诚地按事物真相顺服的，便高傲地假想，并以一种并不存在的事物状态来欺骗自己；这样，圣宗徒若望所说的话就变得可理解了：\BibleQuote{魔鬼从起初就犯罪，}\BibleRef{若一 3:8}就是说，从他受造之时起，他便拒绝了正义，而这正义，惟有那虔诚顺服天主的意志才能享有。凡采纳这种见解的人，至少与那些异端者，即摩尼教徒，以及任何其它有害的教派不相一致；这些教派可能设想魔鬼从一个相反的恶本原获取了他固有的本性。这些人大为错误所愚弄，以致虽然他们同我们一样承认福音的权威，却没有注意到，主并没有说：\BibleQuote{魔鬼本性是与真理无缘的}，而是说：\BibleQuote{魔鬼不居于真理中}，主以此使我们明白，他是从真理中堕落的；他若居于其中，就会成为真理的分享者，并与圣天使一同永居于真福中。\par

% structure: p0032 source=h2 target=subsection reason=relative-heading-rank; base=h2

\subsection{解释论魔鬼所说，\BibleQuote{他不居于真理中，因为真理不在他内。}的意思。}

此外，似乎我们正在询问为何魔鬼不居于真理中，主便附上理由说：\BibleQuote{因为真理不在他内。} 其实，他若居于真理中，真理便在他内。但这说法是不寻常的。因为，依照经文所说，\BibleQuote{他不居于真理中，因为真理不在他内，} 看来好像真理不在他内是他不居于真理中的原因；而实际上，他不居于真理中反倒是真理不在他内的原因。在圣咏中也可找到同样的语式：\BibleQuote{我向你呼号，因为你俯听了我，天主啊！}，我们本应期望它这样说：\Quote{天主啊，你俯听了我，因为我向你呼号了。} 但当他先说\Quote{我呼号}之后，仿佛有人寻求这事的证据，他便借着天主俯听他的效果，证明他祈祷的恳切有效；好像说：我祈祷的证据，就是你俯听了我。\par

% structure: p0034 source=h2 target=subsection reason=relative-heading-rank; base=h2

\subsection{我们如何理解\BibleQuote{魔鬼从起初就犯罪}这句话。}

\Person{若望}论及魔鬼的话：\BibleQuote{魔鬼从起初就犯罪}\BibleRef{若一 3:8}，那些以为此言意指魔鬼受造即具罪性的人，乃是误解了；因为罪若属自然，便根本不成其为罪。他们又如何应对先知们的明证——或\Person{依撒意亚}以巴比伦王预表魔鬼时所说的：\BibleQuote{晨星！朝霞之子！你怎么从天坠下？}\BibleRef{依 14:12}，或\Person{厄则克耳}所说的：\BibleQuote{你曾在伊甸乐园，天主的园中；各种宝石是你的披覆}\BibleRef{则 28:13}，此处意味着他曾经毫无罪愆；因为稍后更明确地说：\BibleQuote{你当初在行径上原是完美的？}如果这些经文不能作他解，则我们也必须从\BibleQuote{他没有站立在真理中}这一句理解到他曾处于真理中，却未持守。至于\BibleQuote{魔鬼从起初就犯罪}这句，不可视为他从受造之始便犯罪，而是从他犯罪之始，当他因骄傲而开始犯罪之时。在\WorkTitle{约伯传}中也有一段论及魔鬼：\BibleQuote{这是天主的化工之始，造了它只为供天使戏弄}\BibleRef{约 40:19}，此语与\WorkTitle{圣咏}相符，那里说：\BibleQuote{盘踞在那里的大龙，是你造了来戏弄的。}\BibleRef{咏 104:26}但这些经文不可令我们设想魔鬼起初被造便是为供天使戏弄，而是他在犯罪之后方受此惩罚。因此，他的开端乃是天主的化工；因为任何本性，即使是兽类中最微末、最低下、最卑贱的，无不是那一位的工程，由祂发出了一切尺度、一切形式、一切秩序，离了这些，任何事物皆无法计划或构思。那么，此种天使本性，其尊贵超越祂所造的其他万有，岂不更是至高者的化工！\par

% structure: p0036 source=h2 target=subsection reason=relative-heading-rank; base=h2

\subsection{论受造物的等级与差别，按其效用或按其本性的自然层次}

因为，在那些存在而不属于造物主天主本质的事物中，有生命的高于无生命的；具有生育能力甚至欲望能力的，高于欠缺此能力的。而在有生命者中，有感觉的高于无感觉的，正如动物高于树木。在有感觉者中，有理智的高于无理智的——\Emph{例如}，人高于牲畜。在有理智者中，不死不灭的，如天使，高于有死的，如人。这些是按自然秩序的等级；但按每人在一物中所见的效用，则有种种不同的价值标准，以致我们有时偏爱某些无感觉之物甚于某些有感觉者。这种偏爱如此之深，以至于若有可能，我们会将后者从自然中全然除去，或由于不知其在自然界中的地位，或虽知之，却为自己的方便而牺牲之。\Emph{例如}，谁不愿家中有面包而非老鼠，有金子而非跳蚤？但这不足为奇，因为即使在人本身（其本性确实拥有最高尊严）的估价中，一匹马往往贵于一个奴隶，一颗宝石贵于一个婢女。因此，观想自然者的理性所引发的判断，与贫困者的急需或纵欲者的贪求所驱使的判断迥然不同；因为前者考量一物在受造物等级中自身有何价值，而急需则考量它如何满足需求；理性寻求理智之光所判定的真实，而快乐则寻求令肉体感官愉悦的刺激。然而在理性本性中，意志与爱的分量，可说，如此重大，以至于虽然在自然秩序中天使高于人，但按公义的尺度，善人却比恶天使更有价值。\par

% structure: p0038 source=h2 target=subsection reason=relative-heading-rank; base=h2

\subsection{邪恶的瑕疵并非本性，而是相反本性，其根源不在造物主，而在意志}

因此，我们须按本性，而非按魔鬼的邪恶，来理解此言：\BibleQuote{这是天主化工之始}；因为毫无疑问，邪恶只能是那原先未受败坏的本性才会有的瑕疵或恶习。恶习又与本性如此相反，以至于不能不损伤它。所以，背离天主若非发生在一个本应依附天主而存的本性上，便不成为恶习。因此，连邪恶的意志也是本性美善的有力证据。但天主，祂既是美善本性的至善创造者，也是邪恶意志的最公义统治者；如此，当它们恶用美善本性时，祂甚至能善用邪恶的意志。因此，祂使魔鬼（因天主的创造而为善，因己意志而为恶）从其高位被贬抑，并成为天使们的笑柄——就是说，祂使魔鬼的试探转而造福那些他企图加害的人。又因天主创造他时，绝不会不知他将来的邪恶，且预见到将由他的恶中引出的善，所以\WorkTitle{圣咏}说：\BibleQuote{你所造的这大龙，是要在其中戏弄的}\BibleRef{咏 104:26}，好叫我们看到，天主虽因自己的美善将他造得善良，却早已预知并安排在他变为邪恶之后要如何利用他。\par

% structure: p0040 source=h2 target=subsection reason=relative-heading-rank; base=h2

\subsection{论宇宙之美，因天主的安排，因对立而更显辉煌}

因为天主绝不会创造任何——我不说天使，就是人——如果祂预知他未来的邪恶，除非祂同样知道如何将此人用于善事，以此装饰时代的进程，仿佛一首以对照加以点缀的精美诗篇。因为所谓的对照属于语言最优雅的修辞之一。在拉丁文中，它们可称为\Quote{对置}，或更精确地说，\Quote{反置}；但这个词在我们中间并不常用，尽管拉丁文，甚至所有民族的语言，都利用相同的修辞手法。在\BibleBook{致格林多人}{后书}中，宗徒\Person{保禄}也优雅地运用了对照，在那个地方他说：\BibleQuote{以正义的武器在右在左，历经光荣和凌辱，恶名和美名；像是迷惑人的，却是真诚的；像是不为人所知的，却是人所共知的；像是待死的，看，我们却活着；像是受惩罚的，却没有被杀；像是忧苦的，却常常喜乐；像是贫穷的，却使许多人富足；像是一无所有的，却无所不有。} \BibleRef{格后 6:7} 因此，这些相反之事物的对照赋予语言以美，同样，此世行程之美也是通过相反事物的对照而成就的，仿佛是事物本身而非言辞的雄辩所安排的。这在\BibleBook{德训}{篇}中十分明确地如此陈述：\BibleQuote{善与恶相对，生与死相对；罪人与虔敬人也如此。你看至高者的一切化工，无不两两相对，一一对立。} \BibleRef{德 33:15}\par

% structure: p0042 source=h2 target=subsection reason=relative-heading-rank; base=h2

\subsection{第19章．—— 关于经句\BibleQuote{天主将光与黑暗分离}我们似乎应如何理解}

因此，虽然天主圣言的晦涩确实有这样的好处，即它引发许多关于真理的意见开始和讨论，每位读者都从中看出某些新的含义，然而，对于一段晦涩经文所说的含义，应当要么由明显事实的见证加以证实，要么在其他较少歧义的文本中加以断言。这种晦涩是有益的，无论作者的原意经过许多其他解释的讨论最终达到，还是尽管原意仍然隐藏，但通过晦涩的讨论而引出其他真理。在我看来，这与天主的工作并不矛盾，如果我们理解天使是在那最初的光被造时受造的，并且当经上说\BibleQuote{天主将光与黑暗分离；天主称光为昼，称黑暗为夜}时，在圣天使与不洁天使之间就有了分离。因为唯有祂能作出这种区分，祂也能在他们堕落之前预知他们会堕落，并预知他们将被剥夺真理之光，停留在骄傲的黑暗中。因为，就我们所熟悉的白天和黑夜而言，祂命令那些我们的感官所能看见的天体将光与黑暗分开。祂说：\BibleQuote{在天空中要有光体，以分别昼夜；}随后不久祂说：\BibleQuote{天主造了两个大光体：较大的控制白天，较小的控制黑夜，并造了星宿。天主将星宿摆放在天空，照耀大地，控制昼夜，分别明暗。} \BibleRef{创 1:14} 但是，在那光，即天使的神圣团体，藉着真理的照耀而灵性焕发，与那对立的黑暗，即那些背离正义之光的、天使灵性状况的污秽败坏之间，唯有祂自己能分开，因为他们的邪恶（不在于本性，而在于意志）虽在将来，对祂却不可能隐藏或不确定。\par

% structure: p0044 source=h2 target=subsection reason=relative-heading-rank; base=h2

\subsection{第20章．—— 关于紧接着光暗分离之后的经句\BibleQuote{天主见光好}}

于是，我们不可略过这段圣经而不注意到，当天主说\BibleQuote{要有光，就有了光}之后，便立即加上了\BibleQuote{天主见光好}。在记述祂将光与黑暗分开、称光为昼、称黑暗为夜之后，并没有这样的语句，好像祂的认可也加于如此黑暗，如同加于光一样。因为当黑暗并非遭非难的对象时，例如天体将此黑暗与我们眼所能见的光分开时，经上才在记录分离之后，而非之前，插入\BibleQuote{天主见光好}的语句。经文如此说：\BibleQuote{天主将星宿摆放在天空中，照耀大地，控制昼夜，分别明暗：天主见光好。}因为祂赞许两者，因为两者都是无罪的。但当天主说了\BibleQuote{要有光，就有了光；天主见光好}之后，圣经接着记述道：\BibleQuote{将光与黑暗分开。天主称光为昼，称黑暗为夜。}此处却未附上\BibleQuote{天主见光好}的语句，以免两者都被称为好，而其中之一却是恶的，不是由于本性，而是由于自己的过失。因此，在这种情况下，只有光得到了造物主的赞许，而天使的黑暗，虽然被预定，却未获赞许。\par

% structure: p0046 source=h2 target=subsection reason=relative-heading-rank; base=h2

\subsection{第21章．—— 论天主永恒不变的知识与意志，祂所造的一切在永恒计划中及实际结果中都使祂喜悦}

因为那不断重复的经文\BibleQuote{天主看了认为好}，除表示这作品在构思中已蒙悦纳——那就是天主的\OriginalTerm{智慧}{wisdom}——之外，还能有什么别的意思呢？毫无疑问，天主并不是在作品实际完成之后才知道它是好的；相反，若非祂预先已知，就什么也不会造成。因此，祂看见那在造成之前若非祂已看见就绝不会造成的是好的，这显然不是祂在发现，而是在教导它是好的。\Person{柏拉图}确实曾大胆地说，宇宙完成时，天主仿佛欢欣雀跃。而\Person{柏拉图}也并非愚笨至此，以为借此表示天主因创造的新奇而变得更加幸福；他不过是要表明，已经完成的作品正得到其创造者的悦纳，就如尚在构思中时一样。这并不是说天主的知识有多种，以不同的方式认识尚未存在的事物、现在存在的事物、以及已经过去的事物。因为祂不像我们那样前瞻将来、注视现在、回顾过去；而是以一种全然不同、且与我们思维方式相距极远的方式行事。因为祂并不通过思想的推移而由此及彼，而是以绝对的不变性注视一切；以致那些在时间内出现的事物，未来的确实尚未存在，现在的现在正存在，过去的已不再存在；但是这一切，都被祂在稳定而永恒的临在中一览无余。祂也不是以眼睛或心灵等不同方式观看，因为祂并非由心灵和身体组成；祂现在的知识也不异于祂曾有的或将有的知识，因为那些过去、现在、将来的时间变化，虽会改变我们的知识，却毫不影响祂的知识，\BibleQuote{在祂内没有变化，也没有转动的阴影。}\BibleRef{雅 1:17}在祂的意念中，也没有从一个意念到另一个意念的增益，因为在祂的灵性洞见中，祂所知晓的一切是同时被涵盖的。因为正如祂推动一切时间性的事物而不带任何时间可测量的运动一样，祂也以时间不能测量的知识而知晓一切时间。因此，当祂看见所造之物是好的，也就是看见创造它是好的。而当祂看见它被造成时，并不因此就有了双重的知识或以任何方式增加的知识；仿佛在造成祂所看见的之前，祂的知识较少似的。因为，如果祂的知识不是那样完美，以致不因竣工的作品而有任何增益，那祂确实不会像现在这样是一位完美的创造者。因此，如果只是为了通知我们是谁创造了光，只说\Quote{天主创造了光}就够了；而如果还要告知我们它是通过什么方法造成的，那么说\BibleQuote{天主说：\InnerQuote{要有光}，就有了光}就已足够，因为这样我们不仅能知道天主创造了世界，也能知道祂是借着\OriginalTerm{圣言}{Word}创造的。然而，因为理应向我们提示关于受造物的三个主要真理，即：谁造的，以何种方式，以及为何而造，于是经上记着：\BibleQuote{天主说：\InnerQuote{要有光}，就有了光。天主看那光，认为好。}那么，若我们问是谁造的，就是\Quote{天主}；若问以何种方式，是祂说了\Quote{要有}，它便有了；若问为何而造，就是\Quote{它是好的}。也没有任何创作者比天主更卓越，没有任何技艺比天主的圣言更有效，没有任何理由比善的天主创造善的受造物更美好。这也正是\Person{柏拉图}所指出的创造世界最充足的理由：即由善的天主创造出善的作品；不论他是读过这段经文，还是或许从读过的人那里得知这些事，或是凭借他敏锐的天才，通过受造之物洞悉了灵性而不可见的事物，或是从已识透这些事的人那里领受了教导。\par

% structure: p0048 source=h2 target=subsection reason=relative-heading-rank; base=h2

\subsection{第二十二章 论那些不赞成此善的创造者所造的善的受造物中的某些部分，并以为有某种自然之恶的人}

然而，这美好受造的原因，即天主的善——我说，这个原因如此公义合宜，若虔诚而仔细地加以衡量，便能终结一切探究世界本原之人的争辩——却未被某些异端者所承认，因为，确实有许多事物，如火、霜、野兽等等，并不适合，反而伤害我们这目前正受公义惩罚的、血液稀薄且脆弱的肉身的必死性。他们不曾思量，这些事物各在其位是何等可赞叹，在其本性上是何等卓越，与受造界的其余部分配合得是何等美妙，它们通过各自的贡献，如同为共同体那样，给宇宙增添了多少恩宠；以及，我们若以认识其适宜用途的方式使用它们，它们对我们自己也是何等有益——甚至毒物，若使用不当则致人死命，若依照其性质和目的而使用，便成为有益和治病的；正如另一方面，那些令我们喜悦的事物，如食物、饮料和阳光，若过度或不当地使用，也会发现是有害的。这样，天主眷顾便劝诫我们，不要愚蠢地诋毁事物，而要仔细探究它们的用途；并且，在我们心智能力不足或软弱时，要相信其中必有用途，尽管是隐藏的，正如我们曾经历过，有其他事物我们几乎未能发现。因为这事物用途的隐藏，本身或是为操练我们的谦卑，或是为贬抑我们的骄傲；因为没有任何本性是恶的，这名称无非是指善的缺失。但从地上之物到天上之物，从可见之物到不可见之物，有些比另一些更好；为此它们是不平等的，为的是使一切都能存在。如今，天主在大事上是如此伟大的工匠，在小事上祂也同样伟大——因为这些小事不可按其自身的大小（其实并不存在）来衡量，而应按其设计者的智慧来衡量；如同在人可见的外貌上，若剃去一道眉毛，从身体上取走了几乎微不足道的东西，但从美上却取走了多少！——因为美不在于体积，而在于肢体的比例和安排。但我们并不太惊讶，那些以为有某种恶的本性，是由一种与它相配的对立本原产生并繁衍而来的人，拒绝承认受造的原因乃是良善的天主创造了美好的受造界。因为他们相信，祂是被迫于击退那与之争战的恶的急切需要，才进行这创造的工程，并且祂将自己善的本性混合于恶，为要抑制和战胜它；而这祂自己的本性，既被如此可耻地玷污，又被最残酷地压制和俘掳，祂便努力要清洗和解放它，虽尽了一切劳苦，却不能完全成功；但其中无法从那污秽中被洁净的部分，便要作为那被战胜且被囚禁的仇敌的监牢和锁链。摩尼教徒若是相信天主的本性，如其本然，是不变且绝对不朽的，不会遭受任何伤害；而且，若他们还以基督徒的清醒持守，那灵魂——这灵魂已显示出能因自己的意志而在恶中变坏，能被罪所败坏，并因此而被剥夺永恒真理的光照——我说，这灵魂并非天主的一部分，也非与天主同性同体，而是由祂所造，且与其造主截然不同。\par

% structure: p0050 source=h2 target=subsection reason=relative-heading-rank; base=h2

\subsection{奥力振的学说所包含的错误。}

但更令人惊奇的是，即便在我们当中，有些人相信万物只有一个本原，且凡非属神的本性若非由那位造主所产生便不能存在，却仍不肯以善良而单纯的信仰，接受这如此善良而单纯的世界受造之原因，即良善的天主创造了它，使之成为美好的；并且，受造之物既与天主不同，便低于祂，但仍是美好的，因为它们是由非祂莫属的那位所造。但他们却说，灵魂虽非天主的部分，却为祂所造，因离弃天主而犯了罪；按照各自不同的罪，它们配受从天到地的不同程度贬抑，以及作为监牢的各种身体；而这便是世界，这便是它受造的原因，不是为了产生美善的事物，而是为了抑制邪恶。奥力振因持有此意见而受到公正的责备。因为他题名为\WorkTitle{\OriginalTerm{论原理}{\GreekText{περὶ} \GreekText{αρχῶν}}}的书中，这就是他的观点，他的言论。而我无法充分表达我的惊奇，如此博学且精通教会文献的人，竟没有首先观察到，这与那具有权威的圣经的意图是何等相悖，圣经在记述天主的一切工作时，常加上一句：\BibleQuote{天主看了认为好}；且在一切完成后，加上这些话：\BibleQuote{天主看见祂所造的一切，看，一切都非常好}。\BibleRef{创 1:31} 难道这不是明显意味着，除了良善的天主应造出美好的受造物之外，世界受造别无其他原因吗？在这受造界中，若无人犯罪，世界本会充满并装饰着无一例外的美好本性；尽管有罪，万物却并不因此便被罪所充满，因为天界居民的大多数保全了其本性的完整。而那有罪的意志，虽违犯了其自身本性的秩序，却并未因此逃避天主的法律，天主公正地安排万物，使之趋向良善。因为如同一幅画的美，因巧妙处理的阴影而增加，同样，对那有辨识力的眼睛而言，宇宙甚至因罪人而更显美好，尽管就其本身而言，他们的丑恶是可悲的瑕疵。\par

其次，\Person{奥力振}及所有与他持相同看法的人本当看到：如果这真是正确的意见，即世界被造是为了让灵魂因自己的罪过而配置以身体，在其中如同被关在惩戒所里，较轻的小罪者得到更轻、更轻妙的身体，而较粗重、更严重的罪人得到更粗重、更为卑下的身体；那么随之而来的便是，那些处在最深邪恶中的魔鬼，应当比恶人更有理由拥有属土的身体，因为那是万物中最粗重、最不轻妙的。然而事实上，为让我们明白灵魂的功过不可按身体的品质来衡量，最邪恶的魔鬼却拥有一个轻妙的身体，而人——确实是恶的，但与魔鬼的恶相比又微不足道且属于小罪——甚至在犯罪之前就得到了一个泥土的身体。还有什么比下述断言更愚蠢的呢：天主用我们这个太阳，并不是有意要造福物质受造界，或为它的美丽增光添彩，因此才为这独一无二的世界创造了一个唯一的太阳；相反，这不过是恰好有一个灵魂犯了罪，且罪大恶极到堪当被封闭在这样一个身体里而已。按照这种原则，若非一个，而是两个，甚至十个、一百个灵魂同样犯了罪，且罪责程度相同，那么这个世界就会有一百个太阳。而事实并非如此，其原因不归因于创造者对物质事物的安全与美观所作的深思远虑，而应归因于只有一个灵魂达到了如此精准的犯罪品质，以至于唯有它堪当得到这样一个身体。显然，持守这等看法的人，应当致力约束的不是那些他们不知其所言的灵魂，而是他们自己，免得他们——确实理所当然地——远远偏离真理。至于我先前在论及任何受造物时，所推荐的那三个问题：谁造的？用什么方式？为什么？而其答案应为：天主，藉圣言，因为好——对于这三个答案，它是否如此奥秘地指示了\OriginalTerm{天主圣三}{Trinity}本身，即圣父、圣子及圣神，还是说在这段圣经经文中这样接受有着其他合理的理由——我且说，这是可存疑的，我们也不能期望在一卷书中解释一切。\par

% structure: p0053 source=h2 target=subsection reason=relative-heading-rank; base=h2

\subsection{第24章 论\OriginalTerm{天主圣三}{Divine Trinity}，及其临在散见于其化工之中的标识。}

我们相信、我们持守、我们忠信地宣讲：圣父生了圣言——即智慧，万物都是藉祂造成的——唯一圣子，与圣父同是一，与圣父同是永恒，且与圣父同等至善；而圣神是圣父和圣子共有的神，祂与二者同体且共永恒；这整个是因着位格的独立个性而为圣三，又因着不可分割的天主本体而为唯一的天主，也因着不可分割的全能而为唯一的全能者；然而，当我们分别询问每一位时，说每一位是天主、是全能者；当我们论及三位一起时，却不说有三位天主、三位全能者，而说一位全能的天主；这三者的不可分割的统一性就是如此之大，要求必须这样陈述。但是，圣父和圣子都是善的，圣神作为圣父和圣子的神，是否因为祂为二者所共有，便可恰当地称为二者的善，对此我不敢贸然断定。不过，我更少有犹豫地说，祂是二者的圣德，这并非好像祂仅仅是一种天主的属性，而是祂本身就是天主的本体，是圣三中的第三位。我之所以更胆敢如此陈述，是因为尽管圣父是神，圣子是神，圣父是圣的，圣子是圣的，然而第三位却特别地被称为圣神，就好像祂是与另两位同体的本体的圣德。但倘若天主的善无非就是天主的圣德，那么，确实，怀着合理的勤勉探究之心，而非冒昧侵扰，去询问：当圣经记载谁造了每一受造物、以什么方式造、以及为何而造时，是否以谜语般的言说方式暗示了这同一圣三，从而激发我们的探究——这并非不当。因为是圣言的父说了：\Quote{要有！}而当祂说话时所造成的，确实是通过圣言而造成的。而借着\BibleQuote{天主看了认为好}这句话，足以指明，天主造成所造的，不是出于任何必然，也不是为了弥补任何匮乏，而唯独是出于祂自己的善，\Emph{即}因为那是好的。而且这句话是在创造完成之后说的，为要使所造之物满足那促使它被造的善，这一点无可置疑。如果我们这样理解是正确的：这善就是圣神，那么整个圣三就在创造中向我们启示了出来。同样，在至高之处，在圣天使之中的那座圣城的起源、光照和真福也在于此。因为若我们问它的来源，是天主创造了它；或它的智慧来源，是天主光照了它；或它的真福来源，是天主就是它的福乐。它因存立于祂内而得其形式，因瞻仰祂而得其光照，因安居于祂内而得其喜乐。它存在；它观看；它爱。在天主的永恒中有它的生命；在天主的真理中有它的光明；在天主的善中有它的喜乐。\par

% structure: p0055 source=h2 target=subsection reason=relative-heading-rank; base=h2

\subsection{第25章 论哲学分为三部分。}

根据所能做出的判断，哲学家们正是出于同样的理由而寻求将科学划分为三个分支，或者说，他们得以认识到存在一个三分法（因为这不是他们发明的，而只是发现的），其中一个部分称为\OriginalTerm{物理的}{physical}，另一个称为\OriginalTerm{逻辑的}{logical}，第三个称为\OriginalTerm{伦理的}{ethical}。这些名称的拉丁对应词现已在许多作家的著作中通用，因此这些分科被称为\OriginalTerm{自然的}{natural}、\OriginalTerm{理性的}{rational}和\OriginalTerm{道德的}{moral}，我在第八卷中已稍作触及。我并不是要据此断定，这些哲学家在这个三分法中想到了在天主内的某种\OriginalTerm{三一性}{a trinity}，尽管据说柏拉图是最早发现并宣扬这一划分的人，并且他看到唯有天主才是自然的创造者、理智的赋予者，以及使生活变得良善和有福的爱的点燃者。但可以肯定的是，尽管哲学家们在事物的本性、探究真理的方式、以及我们一切行动所应趋向的善等问题上意见分歧，然而在这三大普遍问题上，他们耗尽了全部理智精力。尽管意见纷纭，令人困惑，每个人对于这些问题都力求确立自己的看法，但他们中没有人怀疑：自然有其原因，科学有其方法，生活有其目的和目标。此外，任何一个工匠若想有所成就，必须具备三样东西：\OriginalTerm{天性}{nature}、\OriginalTerm{教育}{education}和\OriginalTerm{实践}{practice}。天性由才能判断，教育由知识判断，实践由其成果判断。我明白，严格地说，成果是我们所\Emph{享受}的，而运用（实践）是我们所\Emph{使用}的。两者之间的区别似乎在于：我们说\Emph{享受}那些因其本身、而不考虑其他目的便使我们喜悦的事物；而\Emph{使用}我们为了追求某种更高目的而寻求的事物。因此，应当使用现世的事物，而非享受它们，这样我们才能堪当享受永恒的事物；不要像那些乖戾的人，他们乐意享受金钱，却使用天主——不是为天主的缘故花费金钱，而是为金钱的缘故崇拜天主。然而，在通常的说话方式中，我们既使用果实也享受使用。因为我们正确地谈论\Quote{田野的果实}，而我们在现世生活中确实都在使用它们。正是按照这一用法，我说在人身上应注意三样东西：天性、教育和实践。正如我说过的，哲学家们由此发展出了那门能使人获得有福生活的科学的三重划分：自然的部分关乎天性，理性的部分关乎教育，道德的部分关乎实践。那么，如果我们自己是本性的创造者，我们就会在自己内产生知识，而不需要经由教育，\Emph{即}从他人学习而获得知识。我们的爱，若是从我们自己出发又回归自身，便足以使我们的生命有福，而不需要任何外在的享受。但现在，既然我们的本性以天主作为其必需的创造者，那么可以肯定，我们必须以天主为师，才能成为智者；也必须由天主倾注给我们属灵的甘饴，我们才能有福。\par

% structure: p0057 source=h2 target=subsection reason=relative-heading-rank; base=h2

\subsection{第26章——论我们甚至在目前状态下的人性中亦发现\OriginalTerm{至高圣三}{the Supreme Trinity}的某种\OriginalTerm{肖像}{Image}}

我们确实在自己内认出\OriginalTerm{天主的肖像}{the image of God}，即\OriginalTerm{至高圣三}{the supreme Trinity}的肖像；这一肖像虽然并非与天主同等，更确切地说，虽然与祂相距甚远——既非\OriginalTerm{同永远}{co-eternal}，一言以蔽之，也非与祂\OriginalTerm{同性同体}{consubstantial}——但在本性上却比天主的任何其他工程都更接近祂，并且注定要得到恢复，以便能带有更进一步的相似性。因为我们既存在，也知道自己存在，并且喜悦于自己的\OriginalTerm{存在}{being}及对此存在的\OriginalTerm{知识}{knowledge}。此外，在这三件事上，没有任何看似真实的幻觉会扰乱我们；因为我们不是通过某种身体感官来接触这些的，就像我们感知外在事物那样——\Emph{例如}，颜色通过视觉，声音通过听觉，气味通过嗅觉，味道通过味觉，坚硬和柔软的物体通过触觉——对于所有这些可感对象，我们在心灵中察觉并保存在记忆中的，只是它们的相似肖像，而非它们本身，这些肖像激发我们去渴望那些对象。然而，无需任何欺骗性的肖像或幻象的呈现，我最为确定：我存在，我对此有知识并喜悦于此。对于这些真理，我毫不惧怕\OriginalTerm{学园派}{Academicians}学者的论证，他们说：你若受了欺骗，又当如何？因为如果我受了欺骗，则我存在。因为那不存在者，不可能受欺骗；如果我受了欺骗，凭这同一事实我便存在。既然我若受欺骗便存在，我如何会在相信\Quote{我存在}这件事上受欺骗呢？因为，如果我受欺骗，我肯定存在。所以，我——那受欺骗者——即使受欺骗，也必须存在，那么在我知道自己存在的这一知识上，我肯定没有受欺骗。因此，同样地，在知道自己知道这一点上，我也没有受欺骗。因为一如我知道自己存在，我也知道自己知道。当我爱这两者时，我便给它们加上了第三个东西，即我的\OriginalTerm{爱}{love}，同样具有重大价值。因为在这一件事上——即我爱——我也没有受欺骗，既然在我所爱的事物上我没有受欺骗；纵使它们都是虚假的，我\Emph{爱}那虚假的事物这一事实仍然是真实的。因为，假如我爱它们这件事是虚假的，我如何能被公正地责备并禁止去爱虚假的事物呢？然而，既然它们是真实而实在的，那么当它们被爱时，谁会怀疑那爱本身也是真实而实在的呢？再者，正如没有人不希望幸福，也没有人不希望存在。因为，如果他什么都不是，又如何能幸福呢？\par

% structure: p0059 source=h2 target=subsection reason=relative-heading-rank; base=h2

\subsection{第27章——论\OriginalTerm{存在}{Existence}、存在的\OriginalTerm{知识}{Knowledge}及对两者的\OriginalTerm{爱}{Love}}

确实，存在这一事实本身由于某种自然的魅力而如此令人愉悦，以至于甚至悲惨之人也仅因此而不愿灭亡；当他们感到自己悲惨时，并不希望自己被消灭，而是希望那悲惨继续如此。且看那些无论是自我评价还是事实上都极其悲惨的人，不仅智者因其愚昧视之为悲惨，那些自以为有福的人也因他们贫穷匮乏而视之为悲惨——如果有人赐给这些人一种\OriginalTerm{不死}{immortality}，在这种不死中他们的悲惨不会死亡，并提供另一种选择：如果他们不愿永远存在于同样的悲惨中，他们可以被消灭，无处存在，不处于任何状态，那么他们立刻会欣喜，不，狂喜地选择永远存在，即使是在那种状态下，而非根本不存在。这类人的众所周知的感觉为此作证。因为我们看到他们惧怕死亡，宁愿在那种不幸中活着也不愿以死亡结束它，这难道还不够明显，本性是如何畏惧\OriginalTerm{消灭}{annihilation}的吗？相应地，当他们知道自己必死时，他们寻求一种巨大的恩惠，即对他们显示这样的怜悯，让他们在同样的悲惨中多活一小会儿，延迟以死亡结束它。这样他们无可怀疑地证明，他们将以何等喜悦的急切接受不死，即使那不死带给他们的只是无尽的毁灭。什么！难道不是所有无理性的动物，它们不知这类盘算，从巨蟒到最小的蠕虫，不全都证明它们希望存在，并因此以它们能力内的各种运动来躲避死亡吗？不仅如此，那些没有我们所见的运动来躲避毁灭的生命的草木，难道不都以各自的方式寻求保存自己的存在，将根越扎越深入土中，以便吸收营养，并向天空伸出健康的枝条吗？最后，甚至无生命物体，它们不仅没有感觉而且没有种子生命，却也或升向高处，或沉入深处，或平衡在中间位置，以便在那种最符合其本性的处境下保护自己的存在。\par

而人类本性何等热爱对自己存在的知识，又何等畏惧被欺骗，这从此一事实可充分理解：每个人宁愿在清醒的心智中悲伤，也不愿在疯狂中快乐。这一宏大而奇妙的禀赋只属于一切动物中的人类；因为，虽然有些动物对此世之光具有比我们更敏锐的视觉，它们却不能抵达那\OriginalTerm{精神之光}{spiritual light}，我们心智以某种方式受到那光的照耀，因而能对万事作出正确的判断。我们判断的能力与我们接受这光的程度成比例。然而，无理性的动物虽没有知识，却确实具有某种类似知识的东西；而其他物质事物被称为可感的，不是因为它们有感觉，而是因为它们是我们的感官对象。然而在植物那里，它们的滋养与繁殖与可感生命有某种相似。不过，这些以及所有物质事物，其起因都深藏于本性之中；但它们的外在形式，赋予这有形世界结构以美，被我们的感官所感知，以致它们似乎想要通过为我们提供知识，来补偿自身知识的匮乏。但我们用身体感官感知它们的方式，并不通过这些感官去判断它们。因为我们拥有另一个、且远为优越的感觉，属于\OriginalTerm{内心的人}{inner man}，藉此我们感知何为义、何谓不义——义凭着一个\OriginalTerm{可理解的理念}{intelligible idea}，不义则由于缺乏这理念。这一\OriginalTerm{内在感觉}{inner sense}的运作，既不靠目力，也不靠耳孔，不靠鼻孔的气孔，不靠上颚的味觉，也不靠任何身体的触觉。凭着这感觉，我确知我存在，并且确知我知道这事实；这两者我皆爱，并以同一方式确知我爱它们。\par

% structure: p0062 source=h2 target=subsection reason=relative-heading-rank; base=h2

\subsection{第二十八章——我们是否应该爱那爱我们存在及其知识的爱本身，好使我们越加肖似\OriginalTerm{天主圣三}{Divine Trinity}的肖像}

关于这两件事，即我们的存在及对存在的认知，以及我们如何爱这两者，并且即使在低等造物中也某种程度地发现了与这些事物的相似之处，却又有所不同，我们已根据本书的范围作了足够论述。我们还需讨论那爱——我们以之爱这两者的爱，以确定这爱本身是否也被爱。无疑它是被爱的；这就是明证。因为在那些本当被爱的人身上，所爱的更是爱本身；因为一个人若知道善却不去爱它，便不能合理地被称为善人。那么，我们难道不是显然在我们自己内爱着那爱本身——我们以之爱一切我们所爱之善的爱吗？因为还有一种爱，我们以之爱那不该爱的事物；而这爱被那爱着应当被爱之事物的人所憎恨。因为两者完全可以共存于一人之内。而这种共存对人是有益的，为要使那有助于我们善度生活的爱得以增长，而那引我们走向邪恶的爱得以减少，直至我们整个生命完全治愈并转化为善。因为若我们是野兽，我们便会爱那肉欲和感官的生活，这便成了我们足够的美善；当我们于此安好时，便别无所求。同样，若我们是树木，严格说来，我们确实无法爱任何东西；然而我们似乎会渴望那能使我们更加丰盛地结出果实的东西。若我们是石头、波浪、风、火焰，或任何这类东西，我们确实既无感觉也无生命，却仍拥有一种趋向于自身适当位置和自然秩序的吸引。因为物体的比重，就如同它们的爱，无论它们因重量而下沉，或因轻浮而上升。因为物体被其重量所带动，如同精神被爱所带动，无论去往何方。但我们是人，是依照我们造物主的肖像所造的，祂的永恒是真实的，祂的真理是永恒的，祂的爱是永恒而真实的，祂自己就是永恒、真实且应受钦崇的天主圣三，无混淆，无分离；因此，当我们纵览祂所建立的一切工程时，我们甚至可以在那些低于我们的事物中，或多或少地辨认出祂的足迹，因为这些事物若没有那至高存在、至善、至智的祂所创造，便不可能存在，或呈现任何形态，或遵循任何规律；然而在我们自身内瞻仰祂的肖像时，让我们如同福音中那个小儿子，回归自我，起身归向祂——我们因罪背离了祂。在那里，我们的存在将没有死亡，我们的认知将没有错误，我们的爱将没有不幸。但现在，虽然我们确信自己拥有这三者，不是依赖他人的见证，而是通过我们自己意识到它们的存在，并通过我们最真实的内心视觉看到它们，但我们自己无法知道它们将持续多久，是否会永不止息，以及对它们善用或恶用将导致何种结局，我们便寻求他人——如果我们尚未发现的话——能告知我们这些事。关于这些见证者的可信度，以后有机会再讨论，而非现在。但在本书中，让我们如已开始的那样，依靠天主的助佑，继续论述\WorkTitle{天主之城}，不是处于旅居和必朽状态中的，而是那永远不朽地存在于天上的——即论述那些忠于天主、从未也永不背弃的圣天使，天主在起初，如我们已经说过的，已将他们与那些离弃永恒光明而成为黑暗者分开。\par

% structure: p0064 source=h2 target=subsection reason=relative-heading-rank; base=h2

\subsection{第29章 论圣天使如何在天主本体中认识天主的知识，以及他们如何在创造者的技艺中看到祂工程的起因，先于在艺术家的作品中看到它们}

那些圣天使认知天主，不是借着可听见的言语，而是借着不变真理、\Emph{即}天主的独生子圣言临在于他们的灵魂中；他们认识这圣言本身，也认识圣父和他们的圣神，并且知道这圣三是不可分的，其三位是同一实体，不是三位天主，而是一位天主；他们对此的认知，比我们对自己的认知更为透彻。同样，他们也认知受造物，不是在其自身内，而是以更好的方式，在天主的智慧中，如同在创造它的技艺中；因此，他们在天主内比在自身内更好地认识自己，尽管他们也有后一种认识。因为他们是受造的，与他们的造物主不同。所以在祂内，他们拥有一种如同正午的知识；在自身内，则是一种黄昏的知识，正如我们先前所解释的。因为，知道一物在设计中的样子——它是按这设计造成的——与知道它本身的样子，这中间有很大的差别：\Emph{例如}，线条的笔直和图形的准确，在心里构思时是一种认知方式，画在纸上则是另一种；正义在不变的真理中是一种认知方式，在义人的精神中又是另一种。其他一切事物也是如此——例如，上下水之间的穹苍，称为天；天下的水聚在一处，露出旱地，植物与树木的生长；日月星辰的创造；以及水中动物、飞鸟、鱼类、深海巨物的产生；和地上一切行走爬行的生物，以及超越地上万物的人自己——所有这些，天使在天主的圣言中以一种方式认识，他们在其中看到万有赖以造成的永远常存的根源与理则；而在万物自身内，则以另一种方式认识：前者是更清晰的知识，后者是较模糊的知识，与其说是设计的知识，不如说是对成品本身的认识。然而，当这些工程被指向对造物主本身的赞颂与钦崇时，便如同黎明在那些静观它们的人心中破晓。\par

% structure: p0066 source=h2 target=subsection reason=relative-heading-rank; base=h2

\subsection{第30章 论数字六的完美，它是第一个由其真因数组成的数字}

这些工程记于经上，在六天内完成（同一天重复了六次），因为六是一个完全数，——这并不是因为天主需要延长时间，似乎祂不能立即创造万物，并使万物以其固有的运动标记时间的进程，而是因为藉着数字六，表明了工程的完美。因为数字六是第一个由其各部分合成的数，\Emph{亦即}，由其六分之一、三分之一和一半合成，它们分别是一、二和三，总计为六。按照这种看待数字的方式，那些整除它的数被称为它的部分，例如一半、三分之一、四分之一，或任何分母的分数；\Emph{比如}，四是九的一部分，但并不是其整除部分；而一是，因为它是九分之一；三是，因为它是三分之一。然而这两个部分，九分之一和三分之一，或一和三，远未构成九的整体和。同样，在数字十中，四是其一部分，但并不整除它；而一是其整除部分，因为是十分之一；它有一个五分之一，即二；还有一个一半，即五。但十分之一、五分之一和一半，或一、二、五，加在一起不是十，而是八。再者，数字十二，其部分加在一起超过了整体；因为它有十二分之一，即一；六分之一，即二；四分之一，即三；三分之一，即四；以及一半，即六。但一、二、三、四、六加在一起，不是十二，而是十六。我认为这些叙述足以说明六这个数字的完美，它如我所言，是第一个恰好由其本身的部分相加而成的数；天主就在这六天内完成了祂的工程。因此，我们不可轻视数的学问，在圣经的许多章节中，这学问对细心解经者极有裨益。在天主的赞颂中，也并非无故地提到：\BibleQuote{你处置一切，原有一定的尺度、数目和衡量。} \BibleRef{智 11:21}\par

% structure: p0068 source=h2 target=subsection reason=relative-heading-rank; base=h2

\subsection{第31章。——论第七日，在第七日庆祝完满与安息。}

但是，在第七日（即同一日重复了七次，七这个数字也是一个完全数，尽管另有原因），宣示了天主的安息，而且我们首先听到了这日被定为圣日。因此，天主不愿以祂的工程，而是以祂的安息来圣化这一日；这安息没有暮色，因为它并非受造物；所以，一方面在圣言中以一种方式被认识，另一方面在自身中以另一种方式被认识，就构成了双重的知识，即昼与暮（日与晚）。关于七这个数字的完美，还有很多可说的，但本书已经太长了，我担心会显得我在炫耀那点肤浅的学问，这比博人一笑更为稚气而不够稳重。因此，我必须有所节制并保持庄重，免得因过于追求\Quote{数目}，而被指责忘记了\Quote{重量}和\Quote{尺度}。在此只需指出，三是最小的奇数整数，四是最小的偶数整数，而七由这两个数字组成。因此，它常用来代表全体，例如：\BibleQuote{义人虽七次跌倒，仍然要起来}\BibleRef{箴 24:16}——就是说，不论他跌倒多少次，也不致丧亡（这里应理解为并非指罪恶，而是指使人谦抑的磨难）。又如：\BibleQuote{我一日七次赞美你}，而在另一处则这样表达：\BibleQuote{我要\Emph{时时}赞美上主}。在圣经权威著作中，这类例子很多，其中数字七如我所言，常用来表示整体或事物的完备。因此，主所说的那位\BibleQuote{他要把你们引入一切真理}\BibleRef{若 16:13}的圣神，也以此数字作为象征。在这一日中有天主的安息，这安息是祂的子民在祂内所寻获的。因为安息在于整体，即在于完美的完整性中，而在部分中则有辛劳。因此，我们只要部分地认识，就必劳苦；\BibleQuote{及至那完全的一来到，局部的就必要消逝。} 我们甚至要辛苦地查考圣经本身。但是那些圣天使，我们在劳苦的旅途中叹息向往他们的团体与集会，他们既已永居于永远的家乡，便享有最从容的知识与安息的幸福。他们帮助我们，也毫不费力；因为他们纯洁而自由的神体活动，无需付出任何辛劳。\par

% structure: p0070 source=h2 target=subsection reason=relative-heading-rank; base=h2

\subsection{第32章。——论天使在创世之前受造的意见。}

但如果有人反对我们的意见，说当\WorkTitle{圣经}说\BibleQuote{要有光，就有了光}时，并不是指圣天使；如果他推测或教导，这指的是那时首先受造的一种物质的光，并认为天使的受造，不仅在那分开诸水并被称为\Quote{天}的穹苍之前，也在\BibleQuote{在起初天主创造了天地}这句话所指的时间之前；如果他断言，\BibleQuote{在起初}这个短语，并非意味着在此之前什么也没有受造（因为天使已经受造了），而是意味着天主藉着自己的智慧或圣言创造了万物，而这圣言在\WorkTitle{圣经}中被称为\BibleQuote{元始}，正如在福音中，祂自己对那些问祂是谁的犹太人回答说，祂就是元始；——我不会争辩这一点，尤其因为我极感欣慰地发现，在\BibleBook{创世}{纪}这卷书的开端，就有对天主圣三的颂扬。因为\WorkTitle{圣经}说\BibleQuote{在起初天主创造了天地}，意思是圣父在圣子内创造了它们（正如\BibleBook{}{圣咏集}所证明的，那里说：\BibleQuote{主，你的化工是何其浩繁！你以智慧造成了它们}），随后不久，又合宜地提到了圣神。因为，在我们得知天主起初造了怎样的地，或者说天主以\Quote{天地}之名预备了怎样的混沌或质料以构成世界，如后续经文所说：\BibleQuote{大地还是混沌空虚，深渊上还是一团黑暗}之后，便为使提及天主圣三完整，立即加上：\BibleQuote{天主的神在水面上运行}。那么，就任凭各人照自己的意见去理解吧；因为这段经文是如此深奥，足以提供许多见解，以锻炼读者的分辨力，而这些见解无一远违信德准则。同时，谁也不可怀疑，圣天使在他们天上的居所，虽不与天主同为永恒，却安稳而确定地享有永恒和真实的福乐。主教导我们，祂的小子们属于他们的团体；祂不仅说：\BibleQuote{他们将如同天主的天使一样}\BibleRef{玛 22:30}，而且还指示出天使们所享有的荣福的观瞻，说：\BibleQuote{你们小心，不要轻视这些小子中的一个，因为我告诉你们：他们的天使在天上，常见我在天之父的面}\BibleRef{玛 18:10}。\par

% structure: p0072 source=h2 target=subsection reason=relative-heading-rank; base=h2

\subsection{第三十三章 论天使的两个不同且不相似的团体，以光明与黑暗之名指称并无不当。}

宗徒\Person{伯多禄}极清晰地声明，有些天使犯了罪，被推落到这个世界的最低处，他们在那里如同被囚禁，直到审判之日的最终定罪；因为他说：\BibleQuote{天主既然没有宽免犯罪的天使，把他们投入了地狱，囚在幽暗的深坑，拘留到审判之时。} \BibleRef{伯后 2:4} 那么，谁能怀疑天主或是在预知中，或是在行动中，就已将他们与其他天使分开？又有谁会争辩说，其余的天使不应恰当地被称为\Quote{光明}呢？因为就连我们这些仍凭信德生活、只怀着希望、尚未享有与他们同等的地位的人，已被宗徒称为\Quote{光明}：\BibleQuote{从前你们原是黑暗，但现在你们在主内却是光明。} \BibleRef{弗 5:8} 至于那些背命的天使，凡是理解或相信他们比不信的人更为恶劣的人，都清楚他们被称为\Quote{黑暗}。因此，虽然在\WorkTitle{创世记}的那些章节中，光与暗应按字面意义理解，其中记载说：\BibleQuote{天主说：\InnerQuote{有光！}就有了光。}以及\BibleQuote{天主将光与黑暗分开。}然而，在我们看来，我们认为这两个天使的团体——一个享受天主，另一个自高自大；一个经上说：\BibleQuote{他的众天使，请赞美他，}另一个，其君王却说：\BibleQuote{这一切，我都要给你，只要你俯伏朝拜我。} \BibleRef{玛 4:9} 一个燃烧着天主的圣爱，另一个散发着自我抬举的不洁欲望。并且，正如经上所载：\BibleQuote{天主拒绝骄傲人，却赏赐恩宠于谦逊人，} \BibleRef{雅 4:6} 我们可以说，一个居住在诸天之天，另一个被驱逐出来，并在空气中的较低区域肆虐；一个在虔诚的光辉中宁静安详，另一个被昏暗的欲望所摇撼；一个依天主的旨意，温柔相助，公义惩罚——另一个则被自己的骄傲推动，沸腾着征服与伤害的欲望；一个是天主美善的使者，极其乐意地施善；另一个则被天主的大能制止，不得逞其害人之意；前者嘲笑后者，因为后者虽通过迫害行善，却并非出于自愿；后者嫉妒前者，因为前者收聚其朝圣者。那么，这两个天使的团体，彼此相异且相反：一个在本性上良善，意志也正直；另一个在本性上也良善，但意志却败坏了。正如在其他更明确的圣经章节中所展示的，我也认为它们在\WorkTitle{创世记}这本书中，是以光明和黑暗的名称被论及的。即使作者或许另有他意，我们讨论这隐晦的语言也不算浪费时间；因为，尽管我们未能发现他的本意，我们却持守了信德的准则，这准则是信友们从其他同等权威的圣经章节中，可以充分确定的。因为，虽然这里讲论的是天主的物质化工，但它们确实与精神的事理有相似之处，因此\Person{保禄}可以说：\BibleQuote{你们众人都是光明之子，白昼之子；我们不属于黑夜，也不属于黑暗。} \BibleRef{得前 5:5} 另一方面，如果\WorkTitle{创世记}的作者在那话语中看到了我们所看到的，那么我们的讨论便达到了一个更为满意的结论：这位如此卓越、具有神性智慧的天主的人，或者更确切地说，那借着他记述天主在第六日完工之化工的天主圣神，我们可以假设并未忽略对天使的一切提及——或是将它们包含在\Quote{在起初}这个词组内，因为天主首先创造了他们；或是，这似乎更为可能，因为天主在独生圣言内创造了他们。而在这\Quote{天}与\Quote{地}的称谓下，标示了整个受造界，或是将它划分为精神和物质的两部分——这似乎更为可能，或是将它分为世界的两大部分，其中包含一切受造之物，如此一来，首先概述了整个创造，然后按照那奥妙的日数，一一列举其各部分。\par

% structure: p0074 source=h2 target=subsection reason=relative-heading-rank; base=h2

\subsection{关于在论及穹苍分开诸水之处意指天使的看法，以及另一种认为诸水并未受造的看法。}

不过，有些人却认为天使的军旅以某种方式在水的名下被提及，以下这句话所指的正是此意：\BibleQuote{在水中间要有穹苍，将水分开！}\BibleRef{创 1:6}他们以为，上面的水应被理解为天使，下面的水则或指可见之水，或指众多邪恶天使，或指地上的万族。若果真如此，那么此处并未显示天使是何时受造的，而是显示他们何时被分开。然而，并不缺少如此愚蠢邪恶的人，竟因圣经中从未记载\Quote{天主说：要有水}，便否认水是天主所造。他们以同等的愚妄，也可能对大地说出同样的话，因为我们确实从未读到\Quote{天主说：要有地}。但他们会说，经上记载著：\BibleQuote{在起初天主创造了天地。}的确，这里水也被包含在内，因为二者都囊括在这一句话之中。因为如圣咏所说：\BibleQuote{海洋属祂，是祂所造，陆地也是祂手造成。}然而，那些愿意将天上的水理解为天使的人，却因元素的比重而感到困难，他们担心水因其流动性和重量，无法被安置在世界的上部。因此，如果他们按照自己的原则构造一个人，他们就不会把任何湿润的体液，即希腊人所称的\OriginalTerm{黏液}{phlegm}，即在我们身体元素中扮演水之角色的东西，安置在人的头脑里。然而，在天主所造的化工中，头部却是黏液的居所，且无疑极为合宜；然而，根据他们的设想，这又如此荒谬，以至于如果我们不是明知这一事实，并由这同一记载得知天主曾将一种湿润、寒冷因而沉重的体液放置在人体最上端的部分，这些权衡世界的人就会拒绝相信。而如果他们被圣经的权威所反驳，他们就会坚持说这些话必定另有所指。然而，假如我们要探究并发现这部神圣经典中关于世界创造的所有细节，我们将有太多话要说，并且会严重偏离本书预定的目标。因此，既然我们已就这两类相异而对立的天使团体（人类的两个团体也在其中找到根源，我们不久就要谈论到它们）似乎有必要说的话都说过了，就让我们立即将本书也作一个结束吧。\par

\SourceRecord{Translated by Marcus Dods.；Nicene and Post-Nicene Fathers, First Series；Vol. 2.；Edited by Philip Schaff.；Buffalo, NY: Christian Literature Publishing Co.,；1887.；<http://www.newadvent.org/fathers/120111.htm>.}

% structure: p0001 source=h1 target=section reason=page-title

\section{\WorkTitle{天主之城（卷十二）}}

\Person{圣奥斯定}首先就天使提出两个问题：天使中何以有的具有善意，有的具有恶意？又，善天使何以享真福，恶天使何以受永苦？然后他论及人的受造，教导说人并非从永恒就有，而是受造的，且惟独由\Person{天主}所造。\par

% structure: p0003 source=h2 target=subsection reason=relative-heading-rank; base=h2

\subsection{第一章：天使的本性，无论善恶，皆属同一}

前卷已阐明，两座城是如何在天使中起源的。在论及人的受造，并指出这两座城在有理性有死的人类中如何兴起之前，我看出，我必须先尽力引证一些能证明以下说法并非不合宜、不恰当的论述：即天使与人共同组成一个社会；故此，并非有四座城或社会——即天使两座，人类两座——而是总共只有两座，一座由善者组成，另一座由恶者组成，不分天使或人。\par

善恶天使之间相反的倾向，并非源于本性与起源的差异，因为\Person{天主}——万有的美善作者与创造者——创造了二者，而是源于意志与愿望的不同，这是无可置疑的。有些天使恒心坚守于万善所共有的善，即\Person{天主}自身，以及\Person{天主}的永恒、真理与爱；另有些天使却迷恋自己的能力，仿佛自己能成为自身的善，便从万善所共有的、更高的福乐之善，堕落到这属于私己的善。他们以永恒的高贵尊严换取骄傲的膨胀，以最确定的真理换取虚幻的狡诈，以合一的爱换取结党纷争，于是变得骄傲、受骗、嫉妒。因此，善天使享真福的原因，在于依附\Person{天主}。同样，恶天使受苦的原因，恰恰相反，在于他们不依附\Person{天主}。因此，若有人问善天使为何享真福，正确地回答是：因为他们依附\Person{天主}；若问恶天使为何受苦，正确地回答是：因为他们不依附\Person{天主}——那么，对于有理性的或理智的受造物而言，除了\Person{天主}以外，别无其他美好。这样，虽然并非一切受造物都能享真福（因为禽兽、树木、石头之类没有此能力），但那有此能力的受造物却不能凭自己享真福，因为它是从虚无中受造的，只能由那创造它的那一位获得真福。因为它拥有的，使之享福；失去的，使之受苦。而那不凭他者、惟凭自己而享真福者，不可能受苦，因为他不能失落自己。\par

因此，我们说，除了唯一、真实、享真福的\Person{天主}以外，没有不变的美善；\Person{天主}所造的万物，固然是好的，因为源于\Person{天主}，但却是可变的，因为不是从\Person{天主}本体所造，而是从虚无中造成。所以，它们虽然不是至善——因为\Person{天主}是更大的美善——但那些能依附不变之善、从而享真福的可变之物，也是非常美好的。因为\Person{天主}就是他们的美好，以致于若没有\Person{天主}，他们便不能避免不幸。宇宙间其他受造物，虽不能受苦，却并不因此就更优越。因为没有人会说，身体的其他肢体比眼睛优越，只因眼睛可能失明。然而，正如有感觉的本性，即使受苦时，也胜过无感觉的石块；同样，理性的本性，即使处于不幸中，也比那缺乏理性或感觉、因而无法经历不幸的本性更卓越。事情既是如此，那么在这被造得如此卓越的本性中——这本性虽然自身可变，却能通过依附不变的美善、至高的\Person{天主}来确保自己的真福；又因为它除非完全享福便不能满足，而若非在\Person{天主}内便不能如此享福——在这本性中，我说，不依附\Person{天主}，显然是一种缺失。但凡缺失，都伤害本性，因而与本性格格不入。因此，那依附\Person{天主}的受造物，与那不依附的受造物，其区别不在本性，而在缺失；然而，正是这缺失，反而证明了本性极其高贵与可敬。因为某种本性，其缺失受到公义的谴责，这本性便确实受到称扬。我们公义地谴责缺失，正是因为它损害了那值得称扬的本性。正如我们说失明是眼睛的缺失，就证明了视觉属于眼睛的本性；说耳聋是耳朵的缺失，就证明了听觉属于耳朵的本性——同样，我们说天使受造物不依附\Person{天主}是一种缺失，便以此最清楚地表明，依附\Person{天主}本属它的本性。谁能恰当地设想或言明，依附\Person{天主}，为\Person{天主}而生活，从\Person{天主}汲取智慧，在\Person{天主}内喜乐，并享受如此宏大的美好而无死亡、无错误、无忧伤，这是多么大的荣耀？因此，既然每一罪恶都是对本性的伤害，恶天使的罪恶，即他们背离\Person{天主}，便足以证明，\Person{天主}创造了他们如此良善的本性，以至于不与\Person{天主}同在，对于这本性就是一种伤害。\par

% structure: p0007 source=h2 target=subsection reason=relative-heading-rank; base=h2

\subsection{第二章：没有任何实体与\Person{天主}对立，因为\Quote{虚无}似乎是与那位至上且永恒的存在者全然相反者}

这样或许足以阻止任何人以为，当我们谈论\OriginalTerm{背命的天使}{apostate angels}时，它们是出于某种不同的来源，而非来自天主。我们将更加容易且清晰地摆脱这种极其不敬的错误，如果我们更清楚地理解，当天主派遣梅瑟到以色列子民那里时，祂藉着天使所说的话：\BibleQuote{我是自有者。}\BibleRef{出 3:14}因为天主是\OriginalTerm{至高存在}{supreme existence}，也就是说，祂是至高者，因而恒常不变；祂赋予受造之物以存在，但不是使它们像祂一样地至高存在。有的赋予较丰富的存在，有的则赋予较有限的存在，从而将万物的本性安排为不同的等级。因为正如\OriginalTerm{sapere}{sapere}产生\OriginalTerm{sapientia}{sapientia}，同样\OriginalTerm{esse}{esse}产生\OriginalTerm{essentia}{essentia}——这确实是一个新词，古老的拉丁作家未曾使用，但在我们这时代已通行，以便我们的语言不致缺少相当于希腊文\OriginalTerm{\GreekText{οὐσία}}{\GreekText{οὐσία}}的词。因为用\OriginalTerm{essentia}{essentia}即能逐字对译。因此，对于那至尊的、创造一切存在者之性体，唯有那不存在者才与之相反。因为\OriginalTerm{非存在}{nonentity}是与存在相反的。这样，没有任何存在相反于天主，那至高的存在和一切存在的创造者。\par

% structure: p0009 source=h2 target=subsection reason=relative-heading-rank; base=h2

\subsection{第三章——天主之敌所以为敌，非由本性，乃由意志；这意志在伤害自身时，也伤害了善的本性；因为恶习若不造成伤害，就不是恶习。}

在\WorkTitle{圣经}中，那些反对天主统治的，被称为天主的仇敌；他们并非由于本性，而是由于\OriginalTerm{恶习}{vice}；他们无力伤害天主，只能伤害自己。因为他们之为天主的仇敌，并不在于有伤害的能力，而在于有反对的意志。因为天主是恒常不变的，不受任何伤害。因此，那使被称为仇敌的人反抗天主的恶习，其恶并不在于天主，而在于他们自身。而这为他们是恶，只因为它败坏了他们本性的善。所以，相反的并不是本性，而是恶习。因为恶与善相反。而谁会否认天主是至善呢？因此，恶习相反天主，正如恶相反善。再者，它所败坏的乃是善的本性，所以它也与这本性的善相反。但是，它相反天主，只是作为恶相反善；它相反那受败坏的本性，则既是作为恶，又是作为伤害。因为对于天主，没有恶能造成伤害；恶只能伤害那些可变的、可朽坏的本性，尽管，根据恶习自身的证据，这些本性原先是善的。因为若非善的，恶习便不能伤害它们。因为除了剥夺它们的完整、美丽、安好、德行，简言之，一切本性之善——即恶习惯于减损或毁灭的——之外，恶习怎能伤害它们呢？但若没有善可以剥夺，就不会有伤害，因而也不会有恶习。因为不可能存在无害的恶习。由此我们可以推断，恶习虽然不能伤害那不变的善，却只能伤害善；因为它不存在时不造成伤害。这样，可以这样表述：恶习不能存在于至善之中，却只能存在于某些善之中。因此，纯粹是善的事物可在某种情况下存在；纯粹是恶的事物则永不可能；因为即使那些因邪恶意志而败坏的本性，就其受败坏而言是恶的，但就其仍是本性而言，仍是善的。并且，当败坏的本性受惩罚时，除了它作为本性所具有的善之外，它还具有了这不免受罚的性质。因为这是公义的，而凡是公义的，无疑都是善的。因为没有人是因本性的恶习，而是因自愿的恶习受惩罚。因为即使是那因习惯势力长期持续而成第二本性的恶习，也起源于意志。因为我们现在所讨论的，乃是那具有理智能力，能够接受那分辨义与不义的光照之本性的恶习。\par

% structure: p0011 source=h2 target=subsection reason=relative-heading-rank; base=h2

\subsection{第四章——论非理性与无生命受造物的本性：它们各从其类，各安其序，无损宇宙之美。}

但责备野兽、树木以及其他类似的、没有理智、感官或生命的、必死且可变的事物之缺陷，即便这些缺陷会摧毁其可朽坏的本性，也是荒谬的；因为按照造物主的旨意，这些受造物获得了一种适合它们的存在，通过消逝并为其他物让位，来确保那最低形式的美，即季节之美，这在它自己的位置上，是这个世界必需的一部分。因为地上的事物既不应被造得与天上的事物平等，它们虽较低劣，也不应完全从宇宙中略去。因此，在那些这类事物适宜存在的地方，有些消亡以便为取代它们而生的其他事物让路，较低的屈服于较高的，被胜过的事物转化为居于主宰地位的事物的性质，这便是被指定给短暂事物的秩序。这秩序的美并不使我们惊异，因为由于我们必死的脆弱，我们被深深卷入其中的一部分，以致不能觉察整体；而在整体中，那些冒犯我们的片段是以最精确的适宜与美和谐地联结在一起的。因此，在我们不能如此清楚地察觉造物主的智慧之处，我们被极为恰当地命令去相信它，以免以人类鲁莽的虚妄，竟敢对如此伟大的工匠的工程寻隙挑剔。同时，如果我们仔细思考地上事物中那些既非自愿也非惩罚性的缺陷，它们似乎反倒显明了这些本性的卓越，而万有都源于并受造于天主；因为我们所不悦于见被缺陷夺去的，正是在这本性中令我们喜悦的——除非这些本性本身也使人不悦，正如当它们对人变得有害时常常发生的那样，那时人们便不按其本性，而按其是否有用来评价它们；例如那些成群的动物击打埃及人的骄傲一事。但按照这种评价方式，他们可能会对太阳本身也吹毛求疵；因为某些罪犯或债务人被法官判处曝晒于太阳之下。因此，受造物并非相对于我们的方便或不便，而是相对于它们自身的本性，而彰显其工匠的荣耀。这样，甚至永火的本性，虽对被判决的罪人而言是惩罚性的，也毫无疑问地值得赞美。因为有什么比燃烧、闪耀、发光的火更美丽呢？有什么比用以取暖、复原、烹饪的火更有用呢？尽管没有什么比焚烧、吞噬的火更具毁灭性。因此，同一物，以一种方式使用时，是毁灭性的；但恰当地使用时，则极为有益。因为有谁能用言语诉说它在全世界的用途呢？因此，我们不应听从那些赞美火光却挑剔其热度的人，他们不按其本性，而是按自己的便利或不便来评判。因为他们想看，却不愿被烧。但他们忘了，正是这令他们如此愉快的亮光，却与弱眼不合，并伤害它们；而在那令他们不快的热度中，某些动物却找到最适宜于健康生活的条件。\par

% structure: p0013 source=h2 target=subsection reason=relative-heading-rank; base=h2

\subsection{在各等级各类别的万有本性中，天主都受光荣。}

因此，万有本性，就它们存在而言，并因此各具有其自身的等级和种类，以及一种内在的和谐，肯定是善的。当它们处于本性秩序所分配给它们的位置时，它们就保存所受的存在。而那些未领受永恒存在的，则被改变，或向好或向坏，以适应造物主之律法使其服役的那些事物的需求和运动；这样，它们便在天主的眷顾中，趋向那包含在宇宙治理总体规划中的目标。因此，虽然短暂易逝之物的朽坏导致其彻底毁灭，但这并不妨碍它们产生那被预定为其结果的事物。既然这样，那至高存在的天主，祂因此创造了一切的、并非最高存在的存在（因为从虚无中造成的，不可能与祂同等，若非祂造它，它甚至根本不能存在），便不因受造物的缺陷而可受指责，反倒因祂所造的万有本性而应受赞美。\par

% structure: p0015 source=h2 target=subsection reason=relative-heading-rank; base=h2

\subsection{善天使的真福之因是什么，及恶者的悲惨之因是什么。}

这样，善天使的真福原因便在于，他们依附于那至高存在者。如果我们询问恶者的悲惨之因，则会想到，并且不无道理，他们悲惨是因为他们离弃了那至高存在者，转而转向自己，而他们自身并不具有这样的本质。这种恶，除了骄傲之外，还能叫什么呢？因为\BibleQuote{骄傲是罪的开端。} \BibleRef{训 10:13}因此，他们不愿意为天主保持自己的力量；依附天主是他们享有更广大存在的条件，他们却因偏爱自己而减损了它。这是他们本性的第一个缺陷，第一个贫困，第一个瑕疵，这本性受造时，诚然并非至高存在，但其幸福在于享有至高存在者；而一旦离弃天主，它将成为，并非全无本性，而是一种存在较不广大、因而可怜的本性。\par

如果再进一步问：他们\OriginalTerm{邪恶意志}{evil will}的\OriginalTerm{动力因}{efficient cause}是什么？没有。因为，既然使行动成为恶的是意志本身，那么又是什么使意志成为恶的呢？因此，恶的意志是恶的行动的原因，但恶的意志却没有动力因。因为如果任何东西是原因，这东西要么有意志，要么没有意志。如果有意志，那意志要么是善的，要么是恶的。如果是善的，谁竟会糊涂到说一个善的意志使一个意志变恶？因为那样的话，善的意志就成了罪的原因；这是最荒谬的假设。另一方面，如果那假设之物有一个恶的意志，我想知道是什么使它如此；而且为了避免无限追溯，我立刻问：是什么使\Emph{第一个}邪恶意志变恶？因为第一个不是指自身被另一个邪恶意志败坏的意志，而是指不是被其他意志变恶的那个意志。如果它之前有使它变恶的意志，那么那个使它变恶的意志才是第一个。但如果有人回答说：\Quote{没有东西使它变恶；它从来就是恶的，}我要问它是否存在于某个本性中。因为如果不存在，那么它根本不曾存在；如果它确实存在于某个本性中，那么它便败坏、腐化了那本性，损害了它，并因此剥夺了它的善。因此，邪恶意志不可能存在于一个恶的本性中，而只能存在于一个既善又可变的、能被这恶习损害的本性中。因为如果它不造成损害，它就不是恶习；从而它所居的意志也就不能称为邪恶。但如果它造成损害，它就是通过夺走或减少善而造成的。因此，如所暗示的，在一个原先具有本性之善、而邪恶意志能藉腐化而减损它的东西中，不可能从永恒就存在一个邪恶意志。那么，如果它不是从永恒就有的，我问，是谁造成的？唯一可提出的回答是：某个本身没有意志的东西使意志变恶。那么我问，这东西是高于、低于还是等同于它？如果高于，那就更好。那么，它岂能没有意志，而不是拥有一个善的意志呢？如果等同于，同样的推理适用；因为只要两个东西同样拥有善的意志，一个就不可能使另一个产生邪恶意志。那么就剩下一种假设：那腐化了首先犯罪的\OriginalTerm{天使本性}{angelic nature}的意志的，是一个本身没有意志的低等东西。但是那个东西，哪怕属于最低下、最世俗的种类，也无疑是善的，因为它是一个本性、一个存在，在其自身的种类和秩序中具有自己的形式和等级。那么，一个善的东西怎能成为邪恶意志的动力因呢？我要说，善怎能成为恶的原因呢？因为当意志离弃在自己之上的，转向低下的时，它就变恶——不是因为所转向的东西是恶的，而是因为那转向本身是邪恶的。因此，不是低等东西使意志变恶，而是意志本身因邪恶而无序地贪求低等东西才变恶的。因为如果有两个人，体质和性情完全相同，看到同样的肉体之美，其中一人被这个景象激发而贪求非法享乐，另一人则坚定地以意志保持贞洁自持，我们会认为是什么造成一个有邪恶意志而另一个没有？在具有邪恶意志的那个人身上，是什么产生了它？不是肉体之美，因为它同样呈现在两人眼前，却没有在两人身上都产生邪恶意志。是那人的肉体在他观看时引起了欲望吗？但为什么另一人的肉体没有？或者是他的性情？但为什么不是两人的性情都有？因为我们假设两人身心禀赋相同。那么，我们是否必须说那人是受了恶灵的隐秘诱惑？仿佛他不是凭自己的意志同意那诱惑或任何引诱似的！那么，这种同意，这种对邪恶的劝诱所生出的恶意志——我们问：它的原因是什么？为了不在这难题上耽搁，如果两人同样受诱惑，一人屈服并同意诱惑，另一人却不为所动，那么除了说一人愿意、另一人不愿意离弃贞洁之外，我们还能给出什么解释呢？至少在像我们假设的、禀赋完全相同的情形下，除了他们各自的意志，还有什么造成了这种差异？同一美貌同样地显现在两人眼前；同样的隐秘诱惑以同等的力量压迫两人。因此，无论我们多么细致地审查，我们也看不出是什么使一人的意志变恶。因为如果我们说那人自己使他的意志变恶，那么在意志变恶之前，他本身岂不就是天主——不变的善——所造的一个善的本性吗？这里有两个人，在诱惑之前身心相同，其中一人屈服于说服他的诱惑者，而另一人却不能被说服去贪求那同样呈现于两人眼前的可爱身体。我们是否要说那被诱惑成功的人败坏了自己的意志，因为他在意志变恶之前无疑是善的？那么，他为什么这样做？是因为他的意志是一个本性，还是因为它是从虚无中造成的？我们会发现后一种情况才是事实。因为如果一个本性是邪恶意志的原因，我们除了说恶生于善，或者说善是恶的原因之外，还能说什么呢？而一个本性，虽是善而可变的，又怎能产生任何恶——也就是说，使意志本身变恶呢？\par

% structure: p0018 source=h2 target=subsection reason=relative-heading-rank; base=h2

\subsection{我们不应期望为邪恶意志寻找任何动力因}

因此，谁也不要寻求恶的意志的\OriginalTerm{产生因}{efficient cause}；因为它不是产生的，而是\OriginalTerm{缺失}{deficient}的，正如意志本身不是产生某物，而是一个缺失。因为从那至高存在者\OriginalTerm{背离}{defection}，转向那存在较少者——这就是开始有一个恶的意志。如今，试图发现这些背离的原因——如我所说，不是产生的，而是缺失的原因——就好像有人试图看见黑暗，或听见静寂。然而，这两者我们都知觉到，前者只凭眼目，后者只凭耳朵；但不是凭着它们实际上的积极存在，而是凭着它们的缺乏。那么，谁也不要向我求知我所知道而其实我不知道的事；除非他或许愿意学习，对于那我们知道其不可知的事保持无知。因为那些不是凭其实际存在，而是凭其缺乏而被知觉的事物，——如果可以说且能理解的话——是通过不认识它们而被知觉的，好借着认识它们而一无所知。因为，当视力观察触及感官的对象时，它只在开始看不见的地方才看到黑暗。同样，除了耳朵，没有别的感官能感知静寂，而它只是通过听不见而被感知。如此，我们的心灵通过理解而感知可理解的形式；但当它们缺失时，心灵就通过不理解它们而认知它们；因为\Quote{谁能理解缺失呢？}\par

% structure: p0020 source=h2 target=subsection reason=relative-heading-rank; base=h2

\subsection{第八章．论错乱的爱，它使意志从不变的美善堕落至可变的美善。}

这我确实知道：天主的本性绝不可能在任何时候、任何地方、以任何方式有缺失；而从虚无中造成的本性则能够有缺失。然而，这些受造物，它们拥有的存在越多，所做的善事越多（因为那时它们做的是积极之事），它们就拥有越多的\OriginalTerm{产生因}{efficient causes}；但就它们在存在上有缺失，从而做恶事而言（因为那时它们的作为岂不是虚无？），它们就有\OriginalTerm{缺失因}{deficient causes}。我同样知道，意志若非自己愿意，便不会成为恶的；因此它的失误应受公义的惩罚，因为它们不是必然的，而是自愿的。因为它的背离不是趋向恶的事物，而是其本身就是恶；也就是说，不是趋向那些本性上或自身上为恶的事物，而是意志的背离本身是恶，因为它相反本性的秩序，舍弃至高存在者而趋向存在较少者。因为贪婪并不是黄金固有的缺点，而是人过分贪爱黄金的缺点，以至于损害了正义，而正义本应被视为比黄金无可比拟地更贵重。淫逸也不是可爱迷人的对象的缺点，而是心灵过分贪爱感官享乐的缺点，以至于忽略了节制，而节制使我们依恋那些在灵性上更可爱、因其不朽性而更令人喜悦的对象。夸耀也不是人的赞美本身的缺点，而是灵魂过分喜爱人们的喝彩，而轻视良心之声的缺点。骄傲也不是授予权力者的缺点，也不是权力本身的缺点，而是灵魂过分爱慕自己权力的缺点，并藐视更高权威的更正义的统治。因此，谁若过分地爱任何本性所拥有的善，即使他得到了它，他自身也在那善中成为恶的，并因被剥夺了更大的善而成为可怜的。\par

% structure: p0022 source=h2 target=subsection reason=relative-heading-rank; base=h2

\subsection{第九章．天使除了从天主领受其本性外，是否也借着圣神赋予爱德而从祂领受了美善意志。}

所以，恶的意志既没有本性的\OriginalTerm{产生因}{natural efficient cause}，或者，如果容许我这样说，也没有\OriginalTerm{本质因}{essential cause}，因为它在可变的精神体中是恶的根源，借此，它们本性的善被削减和败坏；而意志成为恶的，除了从\OriginalTerm{背离天主}{defection from God}之外，别无他因——这一背离的原因也必定是缺失的。至于美善意志，如果我们说它没有产生因，我们就必须小心，不要散布这样一种意见：善天使的美善意志不是受造的，而是与天主同永恒的。因为如果天使自己是受造的，我们怎能说他们的美善意志是永恒的呢？但如果它是受造的，它是与他们一同受造的呢，还是他们曾一度没有它而存在呢？如果是一同受造的，那么无疑地，它是被创造他们的那位所创造的，并且他们一受造，就以祂在他们内创造的爱德，依附于创造他们的那位。他们与其余天使团体分离，因为他们持守了同一美善意志；而其他天使则堕入另一种意志，即邪恶的意志，这堕落本身就是从善的堕落；并且，我们可以补充说，如果他们不愿意堕落，他们本不会堕落。但如果善天使曾一度没有美善意志，而后在天主不干预的情况下，自己产生了它，那么接下来就会是：他们使自己变得比天主创造他们时更好。绝不可有这种想法！因为没有美善意志，他们除了是恶的，还能是什么呢？或者，如果他们不是恶的，因为他们既没有恶的意志，也没有善的意志（因为他们尚未背离他们尚未开始享用的那一位），那么他们肯定不同于后来获得美善意志时的状态，不如那样善。或者，如果他们不能使自己变得比那在工程上无人能胜过的天主所创造的更好，那么当然，如果没有祂的救助工程，他们不可能获得那使他们变得更好的美善意志。虽然他们的美善意志使得他们不转向自己——自己只有更有限的存在——而转向至高存在者，并且因与祂结合，他们自己的存在得以扩展，他们借着祂的通传度着智慧和真福的生活，但这证明了什么，岂不是证明了：无论意志如何美善，如果那位从虚无中创造了他们本性、且使之能享有祂的天主，没有先激发它渴望祂，然后用祂自己充满它，从而使之变得更美善，意志只会无助地仅仅渴望祂？\par

此外，还需探究的是：如果善的天使使自己的意志成为善的，他们这样做是有意志还是无意志？若无意志，便非他们所为。若有意志，那意志是善是恶？若是恶，恶的意志怎能产生善的？若是善，则他们已拥有善的意志。这已有的善的意志是谁造成的？正是那位以善的意志创造他们、或赋予他们那贞洁之爱——他们以此依恋于祂——的造物主；祂在同一行动中既创造了他们的本性，又赐予他们恩宠。因此，我们不得不相信，圣天使从未离开过善的意志或天主的爱。然而，那些虽受造为善、如今却成了恶的天使，是因自己的意志而变恶的。这意志之恶，并非由其善良的本性而来，除非因其自愿背离了善；因为善不是恶的原因，背离善才是。因此，这些天使要么比那些恒守于同一恩宠者领受了更少的天主圣爱之恩宠；要么，假若二者受造时同等的善，那么，当一者因自己的恶意堕落时，另一者则得到了更丰沛的助佑，达到了那样一种真福的顶峰，以至于确信自己永不会从那境地堕落——正如我们在前一卷书中已经阐明的那样。所以，我们应怀着对造物主的颂赞承认：不仅对于圣人，而且对于圣天使，都可这样说：\BibleQuote{天主的爱，藉着所赐与他们的圣神，已倾注在他们心中了。}\BibleRef{罗 5:5} 并且，不仅对于人，而且首先且主要地对于天使，经上所言也是真实的：\BibleQuote{亲近天主是多么美好。} 凡共同享有这一善的人，既与他们所亲近的那一位，又彼此之间，有着神圣的共融，并构成一个天主之城——祂的活的祭品，祂的活的殿宇。我看到，既然我已讲述了此城在天使中的兴起，现在是时候讲述那将来要与不死不灭的天使相结合、目前正从有死的人中聚集、或在地上旅居、或在那些已死亡者的无肉身灵魂的隐秘居所和寓所中安息的那一部分的起源了。因为整个人类都出自天主最初所造的那一个人，这是根据圣经的信仰——这圣经在全世界的万民中享有理应卓绝的权威；因为，在其别的真实陈述中，它以其神圣的预知预言了万民都要相信它。\par

% structure: p0025 source=h2 target=subsection reason=relative-heading-rank; base=h2

\subsection{第十章 论世界过去被赋予数千年之历史的虚妄性}

那么，我们且搁置那些不知其所言的人关于人类本性和起源的种种猜测。因为有些人关于人持有的意见，与他们关于世界本身的意见相同，即认为人类一直存在。例如，\Person{阿普列尤斯}在描述我们人类时这样说：\Quote{作为个体，他们是有死的；但作为集体，作为一个族类，他们却是不死的。} 若有人问他们：如果人类一直存在，他们如何为自身历史的真实性辩护，因为历史记述了谁是发明者、发明了什么、谁最先开创了人文修养和其他技艺、谁最先居住于此地或彼地、此岛或彼岛？他们回答说，大多数——即便不是全部——陆地都因火灾和洪水而间歇性地变成荒芜，以至于人类大为减少，而后再由这些幸存者繁衍至原先的数量；这样，每隔一段时间就有一个新的开端，而那些因严重毁坏而中断受阻的事物，虽仅是恢复，却看似在那时才开始；但人只能由人所生。然而，他们说的不过是自己的想法，而非他们所知道的事实。\par

他们也被那些极其虚假的文献所欺骗；这些文献声称给出了数千年之久的历史，然而，根据圣经的记载推算，我们发现还未满六千年。且不必多费言辞揭露这些文献的毫无根据——它们竟记载了如此多的千年——也不必证明其凭据全然不足，我只需引述\Person{亚历山大大帝}写给他母亲\Person{奥林匹娅丝}的一封信。信中他转述了从一位埃及祭司那里得来的、取自他们神圣档案的叙述，其中所述王国之事在希腊史家那里也有提及。在这封亚历山大的信中，\Place{亚述}王国被赋予了超过五千年的历史；而在希腊史书中，从\Person{贝洛斯}本人的统治算起——希腊人和埃及人都一致认为他是亚述的第一位国王——只算出了一千三百年。接着，对于波斯人和马其顿人的帝国，那位埃及人赋予了他们超过八千年，直算到他正在与之谈话的亚历山大时代；而希腊人只给马其顿人四百八十五年，算到亚历山大逝世，给波斯人二百三十三年，算到他的征服结束。因此，希腊人所给的年数比埃及人少得多；实际上，即便乘以三，希腊年表仍然更短。据说埃及人从前一年只算四个月；因此，按照现今他们以及我们中间所运用的更圆满、更真实的算法，一年将包含他们的三个旧年。但即便这样，如我所说，希腊历史在年代学上与埃及历史仍不相符。因此，应给予前者更大的信任，因为它并未超过由我们那些真正神圣的典籍所提供的关于世界存续时间的真实记述。此外，如果这封如此著名的亚历山大信件，在年代学问题上与有可能可信的记载有如此巨大的差异，那么，那些充斥着神话般和虚构的古老事物的文献，他们却妄想用来对抗我们众所周知且神圣的典籍的权威，我们又怎能相信它们呢？这些典籍预言了全世界都将相信它们，而全世界也确实相信了；它们也通过对未来事件的预言——这些预言已如此精准地应验——证明了它们对过去事件的叙述也是真实的。\par

% structure: p0028 source=h2 target=subsection reason=relative-heading-rank; base=h2

\subsection{第11章 论那些认为这个世界确实不是永恒的，而是要么有无数个世界，要么同一个世界在固定周期结束时不断分解为其元素并更新的人}

还有一些人，尽管不认为这个世界是永恒的，却主张要么这个世界不是唯一的，而是有无数个世界，要么它确实是唯一的，但它会在固定的时间间隔内死亡和重生，而且这种情况无数次发生；但他们必须承认，在还有其他人能生育他们之前，人类种族就已经存在了。因为他们不能假设，如果整个世界毁灭了，世界上还会留下一些活着的人，就像他们可能在洪水和烈火中幸存下来一样——其他那些思辨者认为这些灾难只是局部的，并因此能合理地推论，有少数人幸存，其后代将重新繁衍人口；但正如他们相信世界本身是由其自身的物质更新的一样，他们也必须相信，人类是由世界的元素产生的，然后，如同其他动物的后代是从父母繁衍出来的一样，人类的子孙也是从他们的父母繁衍出来的。\par

% structure: p0030 source=h2 target=subsection reason=relative-heading-rank; base=h2

\subsection{第12章 如何答复那些因人类受造日期较近而指责创造人类的人}

对于那些总是问，为什么在无穷无尽的过去的漫长世代里，人没有被创造出来，而是出现得如此之晚，按圣经的说法，自他开始存在以来还不到六千年，我要就人类的受造回答他们，正如我就世界的起源回答那些不愿相信世界不是永恒，而是有一个开端的人一样；这一点，就连\Person{柏拉图}本人也最明确地断言了，尽管有些人认为他的说法并不符合他真正的观点。如果令他们反感的是，按我们的权威记载，人类受造以来所经历的时间如此之短，他的年岁如此之少，那么让他们考虑这一点：任何有限的东西都不长，所有有限的时间年代，与无终的永恒相比，都是极其短暂的，甚至算不得什么。因此，假如人类受造以来所逝去的，我不说五六千年，就是六万年或六十万年，或六十倍于此，或六百倍或六十万倍于此，又或者这个数目相乘直至无法再用数字表示，同样的问题仍然可以提出：为什么他不在更早以前被造呢？因为天主没有创造人类的那段过去且无边的永恒，是如此浩瀚，无论你拿多么巨大、无法言说的年代来与之比较，只要这段时间有一个确定的终结，那就甚至不能比作你拿最小的一滴水与环绕大地四面流淌的海洋相比。因为这两者，一个确实很小，另一个无比巨大，但二者都是有限的；然而，那段从某个始点开始、被某个终点限定的时间，无论它有多长，若你将其与那没有开端的永恒相比，我真不知道该说应视它为最微末的存在，还是什么都不是。因为，拿这段有限的时间，从它的终点开始，一刻一刻地减去最短暂的刹那（就像你可以从一个人的生命中一天一天地减去，从他现在活着的这日，一直回溯到他出生的那天），尽管你在回溯中需要减去的刹那数目如此巨大，无以名状，但这种减法总有一天会带你回到开端。可是，如果你从一段没有开端的时间中减去——我不说一刻一刻短暂的刹那，也不是一小时、一天、一月、甚至成批的年，而是即使最精于计算的人也无法命名的巨大年限——减去像我们通过一刻刻减法所设想的那般巨大的年限，而且不是一次两次重复地减，而是不停地减，你又能达到什么效果，又能造成什么结果呢？因为你永远无法到达那个并不存在的开端。因此，我们现在五千年后才提出的问题，我们的后代在六十万年之后，或许也会怀着同样好奇提出，假设这些代代消亡又更新的世人持续如此长久地衰残和更新，假设后代仍像我们一样软弱无知。在我们之前生活的人，在人类还新鲜地生活在地上时，也可能问过同样的问题。第一个人类自己，甚至可能在他受造后的第二天，或者当天，就会问为什么他不更早受造。而且，无论他在更早或更晚的时候受造，关于这个世界历史开端的争论，都会遇到与现在完全相同的困难。\par

% structure: p0032 source=h2 target=subsection reason=relative-heading-rank; base=h2

\subsection{第13章 论时代的轮回——一些哲学家相信，经过一个固定的周期之后，万物会再次回到最初的秩序和形态}

一些哲学家认为，解决这一难题的唯一可行方法，就是引入\OriginalTerm{时间循环}{circuli temporum}；在循环中，宇宙的自然秩序不断更新和重复。因此，他们断言，这些循环将无休止地重现，一个循环逝去，另一个循环又来到，尽管他们并未达成一致，究竟是同一个永恒的世界会经历所有这些循环，还是世界会在固定的时间间隔内消亡，然后更新，从而展现出重复出现的相同现象——过去的事和未来的事正好重合。从这种幻想的变迁中，他们甚至也没有豁免那获得了智慧的不朽灵魂，反而将其交付于永不休止的\OriginalTerm{转生}{transmigratio}，徘徊在虚假的\OriginalTerm{幸福}{beatitudo}与真实的痛苦之间。因为，既然没有永恒存在的保证，那么，那些或处于对真理的无知，对即将临到的痛苦盲然不见，或虽然知道，却处于痛苦和恐惧之中的人，又怎能真正地被称为幸福呢？或者，如果灵魂过渡到永福，永远离开痛苦，那么在时间中就会发生一件时间不能终结的新事。那么，为什么世界不可以呢？为什么人也不可以呢？因此，沿着\OriginalTerm{健全的教义}{sana doctrina}的正道，我们便能逃脱那些骗人与被骗的哲人们所发现的、我不知为何物的曲折迷途。\par

有些人，在提倡这些使万物复归原状的循环时，为支持他们的假设，引用撒罗满在\BibleBook{训道}{篇}中所说的话：\BibleQuote{已有的事，后必再有；已行的事，后必再行。日光之下并无新事。岂有一件事人能指着说这是新的？哪知，在我们以前的世代早已有了。} 他这话或是指他刚才所谈之事——世代的更替、太阳的运行、河流的走向——或是指一切有生有死的受造物。因为在我们以前已有前人，在我们同时有今人，在我们以后也将有后人；一切生物和草木莫不如此。即使是怪异和不规则的产物，虽彼此互异，且有些据说是独一无二的个例，但就其奇异怪异而言，在某种意义上，它们也是已有过的，将来还会有，日光之下并无新鲜事。然而，有些人却想把这些话理解为：在\OriginalTerm{天主}{God}的\OriginalTerm{预定}{predestination}中，万物已然存在，因此日光之下没有新事。但无论如何，信友切不可认为，撒罗满这些话指的是那些哲学家们所说的循环，即同样的时期和事件重复出现；例如，哲学家\Person{柏拉图}曾在\Place{雅典}一所学园中教学，那么，在数不尽的世代之前，在漫长却固定的周期中，就曾有同样的\Person{柏拉图}、同样的学园、同样的门徒存在，并且在将来无数次的循环中，这一切还会重复——切不可，我说，我们切不可相信这等事。因为基督一次为我们的罪而死；从死者中复活后，祂就不再死；死亡再也不能统治祂了；\BibleRef{罗6:9}而我们自己在复活后，将\BibleQuote{永远与\OriginalTerm{主}{Lord}同在，}\BibleRef{得前4:16}正如我们现在按圣咏作者的话对祂说：\BibleQuote{上主啊，祢必保护我们，祢必由这一代永远护佑我们。} 我觉得接下来这句话也十分贴切：\BibleQuote{恶人\Emph{绕圈子}行走，} 这并非指他们的生命会按哲学家们所想像的那些循环而行，而是因为他们错误教义现在所走的路径是迂回的。\par

% structure: p0035 source=h2 target=subsection reason=relative-heading-rank; base=h2

\subsection{论人类在时间中的受造，以及此事如何在没有\OriginalTerm{天主}{God}任何新的计划或目的改变的情况下完成。}

难道这有什么奇怪吗？他们纠缠在这些循环中，既找不到入口，也找不到出路。因为他们不知道人类和我们这有死境况的起源，也不知道它将如何终结，因为他们无法参透\OriginalTerm{天主}{God}深不可测的智慧。因为，虽然\OriginalTerm{天主}{God}自己是永恒的、无始的，祂却使时间有了开端；人，祂先前并未曾造，却在时间中造了他，这不是出于新的、突然的决定，而是出于祂永恒不变的计划。谁能测透这计划的深不可测的底蕴，谁能探究这天主的上智，祂竟毫不改变自己的旨意，创造了未曾有过的人，使他存在于时间之中，并从一个人繁衍了人类？因为圣咏作者自己，当他首先说：\BibleQuote{上主啊，祢必保护我们，祢必由这一代永远护佑我们，} 随后便责备那些愚昧和不虔敬的道理不给人灵魂存留永远的解脱和真福，并立刻加了一句：\BibleQuote{恶人\Emph{绕圈子}行走。} 然后，仿佛有人对他说：\Quote{那么，你相信什么，感觉到什么，知道什么？我们是否该相信，\OriginalTerm{天主}{God}突然想到了要创造人，这个人在过去的永恒中从未造过？——对\OriginalTerm{天主}{God}来说，不会有任何新的事情发生，祂也没有任何变化？} 圣咏作者接着回答，好像对\OriginalTerm{天主}{God}本人说：\Quote{按祢智慧的高深，祢增多了世人的子孙。} 他似乎说：让人们随心所欲去幻想吧，让他们按自己的喜好去推测和争论吧，但祢却是按照祢那无人能理解的智慧的高深，增多了世人的子孙。因为这确实是一件高深事：\OriginalTerm{天主}{God}永远存在，而人，祂先前从未造过，却愿意在时间中造他，而且是在不改变祂的计划和旨意的情况下造的。\par

% structure: p0037 source=h2 target=subsection reason=relative-heading-rank; base=h2

\subsection{我们是否应相信，\OriginalTerm{天主}{God}既是永恒的主宰，就必有受造物始终受祂统治；以及我们在何种意义上可说受造物始终存在，却不能说它与\OriginalTerm{天主}{God}同为永恒。}

对于我自己，确实，我不敢说曾有一时主天主不是主，同样，我不应怀疑人没有时间之前的存在，而是首先在时间内受造。但当我考虑若非总有某些受造物，天主能是何种主时，我想起自己的卑微，以及经上记载的，就畏缩不敢作任何断言：\Quote{\BibleQuote{有谁能知道天主的旨意？有谁能想象主的意愿？因为有死之人的思想胆怯，我们的计谋不可靠。因为这必腐朽的肉身，重压着灵魂；这地上的帐篷，迫使心神多虑。} \BibleRef{智 9:13}}在这地上的帐篷中，我确实思虑许多事，因为在许多事中那唯一真实的事，或超越许多事的那事，我无法找到。那么，如果在我思绪万千之中，我说对于那位始终是且永远是主的天主，常有一些受造物为祂所主宰，但这些受造物并非始终是同样的，而是彼此接续（因为我们不愿显得说任何受造物与造物主\OriginalTerm{同永远}{co-eternal}，这种断言既为信仰也为健全理性所谴责），我必须小心，免得陷入荒谬无知的错误，竟主张通过这些接续和变化，有死的受造物曾永远存在，而不朽的受造物——即天使——直到我们的世界受造之日才开始存在；至少，如果天使是由那首先受造的光，或更好说，是由经上所说的那个天所意指的：\Quote{\BibleQuote{在起初天主创造了天地。} \BibleRef{创 1:1}}天使至少在受造之前并不存在；因为如果我们说他们曾永远存在，我们似乎就要使他们与造物主同永远。再者，如果我说天使不是在时间内受造的，而是存在于一切时间之前，作为那些由永远为\OriginalTerm{君王}{Sovereign}的天主对其行使王权的对象，那么我就会被问：如果他们在一切时间之前受造，作为受造物，他们是否能永远存在？或许可以回答说，为何不能是\Emph{永远}，因为那存在于一切时间内的，很可以适当地说是\Quote{永远？}如今，这些天使存在于所有时间中，以致甚至时间存在之前他们已受造，这是如此真实；至少如果时间与诸天一同开始，而天使先于诸天存在。而如果在这些天体之前，时间已经存在，并非由时、日、月、年标记——因为这些通常且适当地称为时间的时段度量，明显是随着天体的运动而开始的，因此天主在指定它们时说：\Quote{\BibleQuote{作为规定时节和年月日的记号，} \BibleRef{创 1:14}}——如果，我说，在这些天体之前，时间已由某种变化的运动所是，其部分前后相继且不能同时存在，并且如果在天使中有某种这样的运动使得时间必然存在，并且他们自受造之初就服从于这些时间性的变化，那么他们就存在于所有时间中，因为时间与他们一同产生。而谁能说那存在于所有时间中的，不是永远呢？\par

但是，如果我这样回答，他们又会对我说：那么，如果造物主与它们一直都是存在的，它们怎能不与造物主\OriginalTerm{同为永恒的}{co-eternal}呢？如果我们认为它们一直存在，又如何能说它们是被创造的呢？对此我们该如何回答呢？难道我们要说这两种说法都是正确的吗？即它们一直存在，因为它们存在于所有时间中，它们是与时间一同被造，或者时间是与它们一同被造，而它们也是被造的？因为同样，我们不会否认时间本身也是被造的，尽管无人怀疑时间存在于所有时间中；因为如果时间并非存在于所有时间中，那就有一个没有时间的时间存在。但是最愚蠢的人也不能作出这样的断言。因为我们能够合理地说：有一个时间\Place{罗马}不存在；有一个时间\Place{耶路撒冷}不存在；有一个时间\Person{亚巴郎}不存在；有一个时间人不存在，等等。总之，如果世界不是在时间的起始处被造，而是在经过了一些时间之后被造，我们可以说，有一个时间世界不存在。但是，如果说有一个时间时间不存在，那就如同说在没有人时有一个人一样荒谬；或者说，这个世界存在时这个世界不存在。因为如果我们不是指同一个对象，就可以使用这种表达方式，比如，当这个人不存在时，有另一个人存在。这样我们就可以合理地说，当这个时间不存在时，有另一个时间存在；但即使是最简单的人也不能说有一个时间时间不存在。因此，正如我们说时间是被造的，尽管我们也说它一直存在，因为在所有时间中时间都存在，所以即使天使一直存在，也不意味着它们不是被造的。因为我们说它们一直存在，是因为它们在所有时间中都存在；而它们之所以在所有时间中都存在，是因为时间本身绝不能没有它们。因为没有受造物（其变化运动容许连续）的地方，就完全没有时间。所以，即使它们一直存在，它们也是被造的；而且，即使它们一直存在，它们也并不因此就与造物主同为永恒。因为造物主是一直存在于不变的永恒中；而它们是被造的，并被称为一直存在，是因为它们在所有时间中都存在，没有受造物就不可能有时间。但时间通过其易变性流逝，不能与不变的永恒同为永恒。因此，虽然天使的不死不灭并不随时间流逝——并不变成过去仿佛现在已不存在，也不拥有未来仿佛尚未到来——但它们的变化运动，作为时间的基础，确实从未来流向过去；所以它们不能与造物主同为永恒，因为在造物主的变化中，我们不能说曾经有过现在已不存在的事物，或将要成为尚未存在的事物。因此，如果天主一直是主，祂就一直有受造物在祂的统治之下——然而这些受造物不是祂所生的，而是祂从虚无中创造的；它们也不是与祂同为永恒的，因为虽然祂在任何时候都不缺少它们，但祂在它们之前，不是因为时间的流逝，而是因为祂常存的永恒。但是，如果我这样回答那些诘问天主如何一直是创造者、一直是主——如果没有一直存在受造界——或如果受造界一直存在，它如何是被造的而不是与创造者同为永恒的人，我恐怕会被指责为轻率地肯定我所不知道的，而不是教授我所知道的。因此，我回到我们的造物主认为应该让我们知道的事情上；而那些祂容许更有能力的人在此生认识，或保留给成全的圣人在来生认识的事情，我承认超出了我的能力。但我认为讨论这些问题而不作肯定的断言是对的，好让读者得到警告，要避免冒险的问题，不要以为自己什么事都能做。相反，让他们努力遵循宗徒有益的命令，他说：\BibleQuote{因为我凭着所赐给我的恩宠，对你们各人说：不可自视过高，超过了应估计的；但应按照天主所分与各人的信德尺度，估计得恰如其分。}\BibleRef{罗 12:3} 因为如果一个婴孩得到的滋养适合他的力量，他随着成长就能接受更多；但如果他的力量和容量负担过重，他就会衰退而不是成长。\par

% structure: p0040 source=h2 target=subsection reason=relative-heading-rank; base=h2

\subsection{我们当如何理解天主在\Quote{永恒的时间}之前所许诺的永生}

我承认，我不知道在人类被创造之前经过了哪些时代，但我确信没有任何受造物与造物主同为永恒。但是，即使是宗徒也把时间说成是永恒的，而且不是指未来（这更令人惊奇），而是指过去。他说：\BibleQuote{希望在永生中，那是不能说谎的天主，在永恒的时间之前所预许的，并在适当的时期，藉着宣讲显示出了祂的圣言。}\BibleRef{铎 1:2-3} 你看他说过去曾有永恒的时间，但这些时间并不与天主同为永恒。既然天主在这些永恒的时间之前不但存在，而且还\Quote{预许}了永生——祂在适当的时间（即适当的时候）显示了这永生——这岂不是祂的圣言吗？因为这永生就是圣言。但是，祂是如何预许的呢？因为这预许是向人作的，而人在永恒的时间之前并不存在。这难道不是说，在祂自己的永恒内，并在祂的、同为永恒的圣言内，那要在适当的时候实现的事，早已被预定和确定了？\par

% structure: p0042 source=h2 target=subsection reason=relative-heading-rank; base=h2

\subsection{健全的信德如何为天主不变的旨意和意愿辩护，以反驳那些主张天主的工程在循环周期中永恒重复，使万物恢复原状之人的论证}

关于这件事，我也毫不怀疑：在第一个人受造之前，从未有过任何人，既不是这同一个人本身通过某种我所不知的周期循环往复，不知经历了多少循环，也不是任何其他类似本性的人。我不为哲学论证所吓倒，其中被认为最锐利的一个论证，是基于这样一种主张：无限不可能被任何认知方式所把握。因此，他们论证说，天主在自己心灵中对祂所造的一切有限事物都具有有限的概念。然而，我们不能设想祂的美善曾无所事事；因为倘若如此，就应归因于祂在时间中从以往永恒的无所作为中觉醒过来，开始活动，仿佛祂为自己那无始的闲散而懊悔，于是着手开始工作。既然如此，他们说，必然就是同样的事物永远重复，并且它们消逝之后，注定还要返回，不论在这一切变化中，世界始终如一——这个世界虽然永恒存在，却也是受造的——，还是世界在这些循环中不断消亡并更新；否则，如果我们指出天主的工作开始于某个时刻，人们就会认为祂视自己以往永恒的闲暇为懈怠和懒惰，并因此将其定罪并改变，因为那使祂不悦。现在，如果假设天主确实一直在创造时间性的事物，但这些事物彼此不同，且一个接一个，如此祂终于造了人，而人祂以前从未造过，那么似乎祂不是凭知识造的人（因为他们认为没有任何知识能把握受造物的无限序列），而是随时辰之指令，恰在那一刻触动祂，伴随着突然而偶然的心意改变。另一方面，他们说，如果承认那些周期，并假设同样的时间性事物重复出现，而世界要么在所有这些轮转中保持相同，要么消亡并更新，那么天主既不会被归因于过去永恒的怠惰闲逸，也不会被归因于轻率而预见不到的创造。而如果同样的事物不这样在周期中重复，那么它们无穷的多样性也就不能被任何知识或预知所把握。即使理智不能驳斥，信德也会嘲笑这些论证，无信仰者力图用它们将我们单纯的虔诚引离正道，好让我们与他们一同\Quote{在圈中绕行}。但是，靠着我们的主天主的助佑，连理智也不难击碎这些臆想所构造的旋转之圈。因为特别使他们误入歧途，将他们自己的圈套比照真理之正道的，乃是他们以自己那属人、易变而狭隘的理智，去度量那绝对不变、无限宽广、且无思维相继的天主心灵，这心灵无需数目而数算万有。因此，宗徒的那句话在他们身上应验了：\BibleQuote{他们以自己衡量自己，自然不明白。}因为他们每当有新的主意，便依此新主意去做面临的新事（他们的心灵是易变的），就断定天主也是如此行事；因此他们不是拿天主——因为他们无法构想天主，所以想到祂时，想到的是一个像他们自己一样的存在——不是拿天主，而是拿自己相比，且不是与天主相比，而是与自己相比。至于我们，我们不敢相信天主在工作时和安息时受着不同的感受影响。事实上，说祂受任何感受影响，都是语言的误用，因为那意味着祂的本性中出现了先前没有的东西。因为受感者是被动承受者，而凡被动承受者都是可变的。因此，祂的闲暇并非怠惰、懒散、无所事事；正如祂的工作并非辛劳、费力、勤勉。祂能在安息中行动，在行动中安息。祂能以（并非新的，而是）永恒的计画开始一项新的工作；祂先前未曾造的，如今祂并非因为懊悔先前的安息而开始造它。但是，当人们谈论祂先前的安息和后来的运作时（我不知道人们如何能理解这些事），\Quote{先前}和\Quote{后来}仅适用于受造物，它们先前不曾存在，而后得以存在。但在天主内，先前的定旨并不被后来的、不同的定旨所改变和抹除，相反，祂以同一永恒不变的意志，就祂所造的万物实现了这一点：先前它们不在时，它们就应该不在；后来当它们开始存在时，它们就应该存在。这样，或许祂会以一种非常引人注目的方式，向那些对此有眼力的人显示，祂是何等独立于祂所造的一切，以及祂创造是由于祂自己白白赐予的美善，因为自永恒以来，祂不偕受造物而寓居，其真福丝毫不减。\par

% structure: p0044 source=h2 target=subsection reason=relative-heading-rank; base=h2

\subsection{第十八章 驳斥那些主张无限事物不能为天主的认知所把握的人}

至于他们另一个主张，即天主的知识无法理解无限的事物，那么，为了让他们探测自己不虔敬的深渊，他们只须断言天主不知道所有的数字即可。因为数字无疑是无限的；因为，无论你假定一个数字的终点在哪里，这个数字不仅能够——我不用说——加一，而且不管它有多么巨大，不管它所理性而科学地表达的那个群体有多么庞大，它依然不仅可以加倍，甚至可以相乘。此外，每一数字都由其本身特性所界定，以至于没有两个数字是相等的。因此，它们彼此既不相等又各不相同；尽管作为个体是有限的，但总体而言则是无限的。那么，天主是否因着这种无限性而不知道数字呢？他的知识是否只延伸到数字的某一定高度，而对其他数字则一无所知呢？谁会愚昧到如此地步，竟敢这样说呢？然而，他们很难将数字排除在外，或主张数字与天主的知识无关；因为他们最崇敬的权威\Person{柏拉图}，就描绘天主按照数字原理构造世界。在我们圣经中也这样对天主说：\BibleQuote{你处置一切，是按数目、衡量和分量。}\BibleRef{智 11:20} 先知也说：\BibleQuote{他按数目领出他们的万象。}\BibleRef{依 40:26} 救主在福音中说：\BibleQuote{就是你们的头发，也都一一数过了。}\BibleRef{玛 10:30} 因此，我们绝不可怀疑，一切数字都为那\BibleQuote{他的智慧无限}者所知，正如圣咏作者所说的。数字的无限性，尽管无法数算无限的数字，但对那智慧无限者而言，却并非不可领悟。这样，如果一切被领悟的事物，都因认识它者的领悟而被限定或成为有限，那么，一切无限性都以一种不可言说的方式，为天主成为有限，因为他的知识能够领悟它。因此，如果数字的无限性对天主的知识——亦即他所领悟的——不能是无限的，那么我们这些可怜的人又算得什么，竟敢擅自为他的知识设限，说除非相同的时间性事物在同样的周期性循环中重复，否则天主既不能预知他要造的事物以造它们，也不能在造了它们之后认识它们呢？天主的知识纯一而多重，于变化中保持统一，他以一种如此不可理解的领悟，领悟一切不可理解之事，以至于即使他愿意让后续的工程新颖而不同于先前，他也不能无序无预知地产生它们，或瞬间构思它们，而是凭他永恒的预知。\par

% structure: p0046 source=h2 target=subsection reason=relative-heading-rank; base=h2

\subsection{第十九章 论无穷的世界，或称万世之世}

我不敢确定是否如此：这些被称为\Quote{\OriginalTerm{万世之世}{ages of ages}}的时间，是否连接成一个连续的系列，以有规的多样性彼此相继，惟有那些摆脱了苦难，并在永恒的洪福中永无止息安住的，才免于这种更替；或者，这些之所以被称为\Quote{\OriginalTerm{万世之世}{ages of ages}}，是要我们领悟到，那些世代在天主不变的智慧中保持永恒不变，并仿佛是那些在时间中流逝的世代的\OriginalTerm{动力因}{efficient causes}。也许\Quote{\OriginalTerm{万世}{ages}}被用作\Quote{\OriginalTerm{世}{age}}，因此\Quote{\OriginalTerm{万世之世}{ages of ages}}的含义无异于\Quote{\OriginalTerm{世之世}{age of age}}，正如\Quote{\OriginalTerm{天外之天}{heavens of heavens}}的含义无异于\Quote{\OriginalTerm{天外天}{heaven of heaven}}。因为天主称那在水之上的空间为\Quote{天}，而圣咏却说：\BibleQuote{愿那在\Emph{诸天}之上的水，赞美上主的名！}那么，我们对于\Quote{\OriginalTerm{万世之世}{ages of ages}}应赋予这两种含义中的哪一种，或者是否还有其他更好含义，这是一个极其深奥的问题；而我们目前所探讨的主题，并不妨碍我们暂时推迟讨论它，无论我们是能对此做出某种定论，还是只能因进一步探讨而更加谨慎，从而在面对如此隐晦之事时不敢妄下论断。因为目前我们正在驳斥那种主张存在着周期性循环，相同事物每隔一段时间便重复出现的观点。那么，无论关于\Quote{\OriginalTerm{万世之世}{ages of ages}}的这些假设中，哪一个为真，对于证实那些循环都毫无帮助；因为无论\OriginalTerm{万世之世}{ages of ages}并非同一世界的重复，而是不同的世界以有序的连接彼此相续，被赎的灵魂安住于确凿的洪福中，不再遭受任何苦难，还是\OriginalTerm{万世之世}{ages of ages}是那掌管时间中将要发生和现在存在之事的永恒原因，同样得出，那些使得相同事物循环出现的周期并不存在；而再也没有什么比圣人们的永生这一事实，更能彻底破除这些循环了。\par

% structure: p0048 source=h2 target=subsection reason=relative-heading-rank; base=h2

\subsection{第二十章 论那些主张享受真正圆满幸福之灵魂，仍须在这些周期性循环中屡次返回劳苦与痛苦中者的不虔敬}

虔诚的耳朵怎能忍受听到：在经历了充满无数且严重痛苦的一生之后（如果那还能称为生命，其实这更像是一种死亡，如此彻底的死亡，以至于我们因爱这现世的死亡而害怕那使我们脱离它的死亡），在藉着真宗教和智慧的助佑，终于补赎并结束了如此灾难性的邪恶和各种苦难之后，当我们因此得以见天主，并藉着默观灵性之光与分享祂永不改变的不朽而进入真福——这正是我们热烈渴望达到的——我们却必须在某个时候失去这一切，而那些失去这一切的人，从那永恒、真理和幸福中被抛下，落入地狱的死亡和可耻的愚妄，卷入可咒的灾祸中，在其中失落了天主，憎恨真理，并在不义的污秽中寻求幸福？而且这将无休止地一遍又一遍发生，按固定的间隔，并定期循环地重演？而且这种确定周期的永恒不息循环，依次除去并恢复真实的痛苦和虚假的幸福，其设计竟是为了让天主能够认识祂的工程，因为一方面祂不能停止创造，另一方面，如果祂一直创造受造物，就无法认识无限数量的受造物？我说，谁能够听这种事呢？谁能够接受或容许它们被说出来呢？如果它们是真的，那么不仅对此保持沉默更为明智，甚至（尽我所能表达）连不知道它们也是一种智慧。因为如果在来世我们将不记起这些事，并因这遗忘而蒙受祝福，那么我们为何现在要以这些知识增加我们本已够沉重的痛苦呢？另一方面，如果这些知识将来要强加给我们，那么至少现在让我们保持无知吧，这样在现时的期望中，我们可以享受一种未来的现实所不会赐予的幸福；因为在此生，我们期待获得永生，但在来世却会发现它虽是幸福的，却不是永恒的。\par

如果他们主张，除非在此生已经被教导那些循环——在其中真福和痛苦相互交替——否则无人能达到来世的真福，那么，他们如何宣称人越爱天主，就越容易达到真福呢？他们教导的恰是使爱本身瘫痪的道理。因为谁不会更加懈怠和冷淡地爱一个人，如果他认为将来必须抛弃这个人，并且将来会憎恨这个人的真理和智慧？而且这还是在他已经尽他所能达到了对天主最极致、最幸福的认知之后？即使是对人世的朋友，如果知道他注定要成为自己的敌人，谁还能保持爱的忠信呢？天主不容许这种意见有任何真实性，这种意见以永不休止却经常且无休止地被虚假幸福的间歇打断的真实的痛苦来威胁我们。因为，还有比这更虚假、更欺骗的幸福吗？在这种幸福真理的照耀下，我们仍然不知道我们将会不幸，或者在其最安全的堡垒中，我们仍然害怕变得不幸。因为，一方面，如果我们不知道将要来临的灾难，那么我们当前的痛苦就不是那么短视，因为这种痛苦确信未来的幸福。另一方面，如果在来世灾难并不向我们隐藏，那么最后将被交换为幸福状态的痛苦时间，比起将要回到痛苦结束的幸福时间，灵魂度过得更幸福。这样，我们对不幸的期待是幸福的，而对幸福的期待是不幸的。因此，既然我们在此生遭受当前的痛苦，又在来世害怕即将到来的痛苦，那么更真实的说法是，我们将永远不幸而不是有时幸福。\par

然而，宗教与真理的大声见证宣告这些事情是虚妄的；因为宗教真实地应许了真正的真福，我们将永远确知它，且任何灾祸都无法中断它。因此，让我们坚守那笔直的道路，即基督；并以祂为我们的向导与救主，从内心和思想上远离无神者那虚幻而无意义的\OriginalTerm{循环}{cycles}。\Person{波斐留斯}虽是柏拉图主义者，却摒弃了其学派的意见，即灵魂在这些循环中不断地消逝与复归；他或是被这观念的荒诞所震撼，或是因对基督宗教的认识而清醒。正如我在第十卷中所提到的，他宁愿说：灵魂被遣入世界，是为认识邪恶，并从中得到净化与解救，一旦返回\OriginalTerm{圣父}{the Father}那里之后，便永不再经历这样的遭遇。如果他尚且摒弃了其学派的教义，我们基督徒又当如何憎恶并避开一种如此没有根据且敌视我们信仰的见解呢？然而，我们既已处置了这些循环并从中脱身，就毫无必要设想人类在时间上没有起点，其理由是：自然界中并无任何新事物，是按其所谓循环，在先前某个时期不曾存在、今后也不再存在的。因为，如果灵魂一旦得解救——这是前所未有的——便永不再返回悲苦之中，那么它的经历中就发生了一件前所未有的事；而且这确实是一件意义至大之事，即安稳地进入\OriginalTerm{永福}{eternal felicity}。如果在一个不死不灭的\OriginalTerm{本性}{nature}中，能够出现一件新的、从未有过、也永不通过任何循环再现的事，为何对有死的本性可能发生同样的事却要争论呢？如果他们主张真福对灵魂来说并非新经历，而只是回归那永恒处于的状态，那么至少它从悲苦中得解救是新的，因为，按他们自己的说法，它所被解救的悲苦本身也是一种新经历。如果这一新经历是偶然发生的，并未被纳入\OriginalTerm{天主眷顾}{Divine Providence}所安排的万物秩序之中，那么，那些确定而有度量的循环在哪里呢？在那些循环里，并无新事发生，万物皆如其旧地重现。然而，如果这一新经历被纳入了那眷顾的\OriginalTerm{自然秩序}{order of nature}中（无论灵魂是为了管教之故而遭受此世之恶，还是因罪恶陷于其中），那么新事便有可能发生，它们前所未有，却又不与自然秩序相违。如果灵魂因其自身的愚昧，能够给自己造成一种新的悲苦——这一悲苦并非未被天主眷顾预见，而是在自然秩序中与对此的解救一同被预备好了——那么，即使怀着人类虚妄的胆大，我们又怎能妄自否认天主能创造新事物——对世界是新的，对祂却非如此——这些事祂先前从未创造过，却自永恒就预见了呢？如果他们声称，得赎的灵魂确实不再返回悲苦，但即便如此，也无新事发生，因为过去、现在、将来都持续有得赎的灵魂继出，那么他们至少必然承认，在此情况下，确有一些新的灵魂，对它们而言，这悲苦以及从中被解救都是新的。因为如果他们主张，那些每日由此构成新人的灵魂（如果这些人活得智慧，他们的灵魂便如此从身体中获得解救，自此永不复归悲苦）并非新的，而是亙古就存在的，那么他们逻辑上必得承认这些灵魂是无限的。因为，无论有限的灵魂数目有多大，都不足以自永恒以来持续产生新的世人——这些人灵魂将永恒地脱离这有死的状态，且永不再归于此。而我们的哲学家们将难以解释，在他们所要求的、必须是有限以便为天主所知的自然秩序中，如何可能存在无限数量的灵魂。\par

如今我们既已摒弃了那些被认为会按固定周期将灵魂带回同样悲苦之中的循环，那么有什么比相信天主既能创造前所未造的新事物，又能在如此施行时保持其意志不变，更合乎天主的理性呢？至于永恒得救的灵魂数目是否能持续增加，就让那些在决定无法容纳无限时如此敏锐的哲学家们自己去定夺吧。就我们而言，我们的论证在两种情形下皆能成立。因为如果灵魂数目可以无限增加，那有什么理由否认前所未造者能被创造呢？既然得赎之灵魂的数目先前从未存在过，而它不仅一次性被造出，且将永不断有新者进入存在。另一方面，如果更合适的是，永恒得救的灵魂数目是确定的，且此数目永不会增加，那么，无论这数目是多少，它确然先前从未存在过；若没有一个开端，它便无法增加而达到其所指的数目；而这个开端先前也从未存在过。因此，为使这开端得以存在，\OriginalTerm{第一个人}{first man}便被造了。\par

% structure: p0053 source=h2 target=subsection reason=relative-heading-rank; base=h2

\subsection{第21章——起初只造了一个个体，人类在他内受造。}

如今，我们已尽量解决了这个极其困难的问题——关于永恒的天主创造新事物而意志却无任何改变——便容易看出，天主乐于由祂所造的那一个人繁衍出人类，比由多人开始人类要好得多。至于其他动物，祂造了一些独居的，自然而然地寻求偏僻之处——如鹰、鸢、狮子、狼及诸如此类；另一些则是群居的，喜爱群聚，宁愿结队生活——如鸽子、椋鸟、雄鹿、小黇鹿之类。然而天主并没有让这两类动物由个别的个体繁殖，而是一次造成多只。另一方面，人的本性介于天使与野兽之间，天主这样创造人，若是他服从其创造者、他的合法之主，并虔诚地遵守祂的诫命，便得以进入天使的行列，无需经历死亡，获得有福的无尽不朽；但若是他因骄傲与悖逆地运用\OriginalTerm{自由意志}{free will}而冒犯了上主他的天主，就会屈服于死亡，像野兽一般生活——成为欲望的奴隶，并在死后被判受永罚。因此，天主只造了一个人，当然不是为了使他孤独无伴、失去一切社会，而是藉由此方式，更有效地向他推崇社会的统一与和谐的纽带，因为人不仅因本性相似，更因家庭亲情而彼此联结。而且事实上，祂甚至没有以造男人的方式创造那将要赐给他作妻子的女人，而是从男人身上造了她，好使全人类都出自一人。\par

% structure: p0055 source=h2 target=subsection reason=relative-heading-rank; base=h2

\subsection{天主预知第一个人会犯罪，并同时预见那众多虔敬者，将藉着祂的恩宠被接入天使的团体中}

而且天主并非不知人会犯罪，且一旦人屈服于死亡，就会繁衍注定死亡的子孙，这些必死者将陷入如此重大的罪恶，以致连那些没有理性意志、由水中和大地上大量受造的野兽，相比那本为推崇和谐而由一人繁殖的人类，都能更加安全和平地与同族生活。因为即使是狮子或龙，也从未像人类相互之间那样，与自己的同类发动如此战争。但天主也预见藉着祂的\OriginalTerm{恩宠}{grace}，有一个民族将被召叫为嗣子，他们因罪赦而得\OriginalTerm{成义}{justification}，将由\OriginalTerm{圣神}{Holy Ghost}联合于圣天使，在最后仇敌——死亡——被毁灭后，享受永恒的和平；祂知道这个民族将从这一事实获益：即天主使全人类出自一人，为要显明祂是多么珍视众多中的合一。\par

% structure: p0057 source=h2 target=subsection reason=relative-heading-rank; base=h2

\subsection{论按天主的肖像受造的人类灵魂的本性}

那么，天主按自己的肖像造了人。因为祂为人造了一个具有理性和智能的\OriginalTerm{灵魂}{soul}，使他能超越地上、空中和海里的所有受造物，它们没有被赋予如此恩赐。当祂用地上的尘土形成人，并愿意他的灵魂像我所说的那样——无论祂是早已造了灵魂，现在通过吹气赋予人，还是通过吹气而造成灵魂，因此天主吹气所形成的那个气息（因为\Quote{吹气}岂不就是制造气息吗？）就是灵魂——祂还为他造了一个妻子，以帮助他生养后裔，并以神圣的方式用人肋骨中取出的一根骨头造了她。因为我们不可按属肉体的方式设想这项工作，好像天主像我们通常看到的工匠那样工作，他们使用双手和所提供的材料，以他们的技能塑造某个物体。天主的手即天主的能力；祂以无形的方式工作，产生有形的结果。但这在人们看来，更像是神话而非真实，因为他们用习以为常的日常作品来衡量天主的能力与智慧，而天主无需种子就能明了并产生出种子本身；他们因无法理解起初受造的事物，便对其持怀疑态度——仿佛他们所知道的有关人类繁衍、受孕和出生的事情，若告诉毫无经验的人，就不会显得那么难以置信了；尽管许多人也将这些事归因于物理和自然的原因，而非天主心智的作为。\par

% structure: p0059 source=h2 target=subsection reason=relative-heading-rank; base=h2

\subsection{天使是否能被称为任何受造物的创造者，即使是最微小的受造物}

但在本书中，我们不与那些不相信天主心智创造并眷顾此世的人讨论。至于那些相信他们的\Person{柏拉图}的人，认为所有有死的动物——其中人居于卓越地位，并接近于诸神本身——并非由那位创造世界的至高天主所造，而是由至高者所造的其他较低级神明，在祂的掌控下行使受托的权力所造——只要这些人能从迷信中解脱出来，不再为敬拜并祭祀这些神明如他们的创造者寻找看似合理的理由，他们也就很容易摆脱这一错误。因为相信或说出天主以外的任何其他者是创造者，哪怕是最微小、有死的物的本性，也是亵渎之举（甚至在能够理解之前）。至于那些天使，柏拉图主义者更愿称之为诸神，虽然他们按所受许可与委派，确实协助产生我们周围的事物，但我们却不能因此称他们为创造者，正如我们不称园丁为果实和树木的创造者一样。\par

% structure: p0061 source=h2 target=subsection reason=relative-heading-rank; base=h2

\subsection{唯天主是一切种类受造物的创造者，无论其本性或形式如何}

因为存在一种形式，是从外面赋予一切有形之物的——例如陶匠、铁匠所构造的那类形式，以及画家和塑形艺术家所描摹的、类似动物身体的形式——但另有内在的形式，它自身并非被构造而成，而是作为\OriginalTerm{动力因}{efficient cause}，不仅产生自然的有形形式，甚至赋予生物生命本身；这种形式源自有理智和生命的本体之隐秘而深邃的抉择。前者可归于任何工匠，后者则唯独归于\Emph{天主}，那位\Emph{造物主与创始者}，祂创造了世界本身与天使，无需世界或天使的帮助。因为那同一的、可以说是创造性的神圣能力，本身非受造，却能创造，它赋予天地以圆浑之形——正是这神圣、有效而创造的能力，赋予了眼睛与苹果各自的圆浑形状；我们四处所见的其他自然事物，其形式也并非来自外部，而是源自造物主隐秘而深邃的大能，那位曾说：\BibleQuote{我不充满天地吗？}\BibleRef{耶 23:24}；祂的智慧正是\BibleQuote{她施展威力，从地极直达地极，从容治理万物}\BibleRef{智 8:1}。因此，我不知道那些首先受造的天使，在造物主创造其他事物时给予了何种助佑。我不能将他们或许做不到的事归于他们，也不应否认他们拥有的能力。然而，在不冒犯他们的情况下，我将那赋予万有本性的创造与肇始之工归于天主，他们自己也满怀感恩地将自身的存在归于祂。我们不称园丁为果实的创造者，因为经上说：\BibleQuote{栽种的不算什么，浇灌的也不算什麼，只在那使它生长的天主。}\BibleRef{格前 3:7} 甚至大地本身，我们也不称之为创造者，尽管她看似万象众生丰饶的母亲，协助种子萌芽破土，并将其植根于自己的胸怀；因为经上又说：\BibleQuote{可是天主随自己的心意给它一个形体，使各样种子各有自己的形体。}\BibleRef{格前 15:38} 我们甚至不应称女人为其子嗣的创造者；因为那创造者更是那位曾对祂的仆人说的：\BibleQuote{我还没有在母腹内形成你以前，我已认识了你。}\BibleRef{耶 1:5} 而且，尽管怀孕妇女种种的心理情绪的确会在她子宫的果实中产生相似的品质——正如\Person{雅各伯}以剥了皮的枝条使有斑纹的羊生出——但母亲创造自己后裔的程度，并不比创造她自身更多。所以，无论借助何种身体或种子的原因来产生万物，无论是通过天使、人或其他低等动物的合作，或经由两性生殖；无论母亲的欲望与心理情绪在多柔和可塑的胎儿身上，产生相应的容貌与色泽有多大的能力；但如此受到各种影响的本性本身，其产生只来自至高天主，而非任何其他。是祂那隐秘的能力渗透万物，临在于一切而不受玷污，赋予一切存在者以存在，并调节限定其存在；因此离了祂，事物便不能成为此或彼，甚或根本无从存在。那么，就工匠的手施加于其作品的外在形式而言，我们不说\Place{罗马}与\Place{亚历山大里亚}是由石匠与建筑师建造的，而是由那些凭其意志、规划和资源建造它们的君王所建，因而前者以\Person{罗慕洛}为奠基者，后者以\Person{亚历山大}为奠基者；那么，我们更应何等坚定地说，唯有天主是万有本性的创造者，因为祂的工程既未曾使用非祂所造的材料，也未曾使用非祂所造的工匠，而且，如果祂，可以这么说，从受造物收回祂的创造能力，它们便将立刻退回到受造之前的虚无之中？我指的是\Quote{前}，乃就永恒而言，非就时间而言。因为除了那创造其运动构成时间的万物的那一位，时间还能有别的创造者吗？\par

% structure: p0063 source=h2 target=subsection reason=relative-heading-rank; base=h2

\subsection{二十六、论柏拉图主义者的一种意见，即天使本身确由天主所造，但此后他们创造了人的身体}

很明显，当他把其他动物的创造归于那些由至高者所造的较低神明时，他的意思是，\OriginalTerm{不死部分}{immortal part}取自天主自身，而这些次要的创造者添加了\OriginalTerm{有死部分}{mortal part}；也就是说，他认为他们是我们身体的创造者，而不是我们灵魂的创造者。然而，由于\Person{波菲利}主张，如果灵魂要得到净化，就必须摆脱与身体的一切纠缠；同时，他同意\Person{柏拉图}和柏拉图主义者的看法，认为那些没有度过节制而尊贵生活的人，会回到有死的身体中作为惩罚（柏拉图认为回到兽类的身体，波菲利认为回到人类的身体）；由此推论，那些他们想让我们当作父母和创造者来崇拜、并因此可以貌似合理地称之为神明的，终究只是制造我们桎梏与锁链的工匠——不是我们的创造者，而是把我们关进最痛苦、最悲惨的感化院的狱吏和看守。那么，让柏拉图主义者们要么停止用我们的身体作为灵魂的惩罚来威胁我们，要么停止宣扬我们应该崇拜那些他们极力劝我们避而远之的、在我们身上做工的神明。然而，这两种意见都是完全错误的。灵魂再次回到此生受罚是假的；天地间有任何事物的创造者不是那位创造天地的造物主，也是假的。因为如果我们活在身体中只是为了补赎罪过，那么\Person{柏拉图}如何在另一处说，倘若世界没有被各种有死和不朽的受造物充满，它就不可能成为最美好、最良善的呢？但是，如果我们即便作为有死者的受造本身也是神圣的恩惠，那么回到身体中，即回到一种神圣的恩惠，又怎么是惩罚呢？而且，如果天主如\Person{柏拉图}所不断主张的那样，在祂的永恒理智中包含了宇宙和一切动物的\OriginalTerm{理念}{ideas}，那么祂为何不亲手创造这一切呢？难道祂会不愿意成为这些作品的构造者吗？而这些作品的理念和计划，正需要祂那不可言喻、且应受不可言喻赞颂的理智？\par

% structure: p0065 source=h2 target=subsection reason=relative-heading-rank; base=h2

\subsection{第27章 人类种族的整个满全都被包含在第一人内，天主在他内预知了其中将要受尊荣和赏报的那部分，以及将要被定罪和惩罚的那部分。}

因此，真宗教有充分的理由承认并宣扬：那位创造了整个宇宙的天主，也创造了一切动物，包括灵魂和身体。在陆生动物中，人是祂按自己的\OriginalTerm{肖像}{image}造成的，并且基于我已给出的理由，被造成一个个体，尽管他并非被遗弃在孤独中。因为没有什么种族像人类这样，在受造时如此倾向于社会，而在败坏后又如此反社会。对于防止纷争，或治愈已有的纷争，人性中最恰当的方法莫过于记起我们众人那\OriginalTerm{第一祖先}{first parent}：天主乐意单独创造了他，使一切人由一人而衍生，以此劝谕他们，要在整个群众中保持合一。而女人从第一人的肋旁为他受造这事，分明意味着我们当学到，夫妻之间的结合纽带应当何等亲密。天主的这些作为确实显得非同寻常，因为它们是起初的作为。不信这些事的人，也不应当相信任何\OriginalTerm{奇事}{prodigies}；因为这些奇事若不越出自然的常轨，就不会被称为奇事。然而，在\OriginalTerm{天主眷顾}{divine providence}如此宏大的治理之下，即使其缘由如何隐密，又岂会有任何徒然之事发生呢？一位圣咏作者说：\BibleQuote{来，观看上主的作为，祂在地上所行的奇事。}天主为何要从男人的肋旁造出女人，以及这第一件奇事预表了什么，我将在别处，靠天主的帮助，予以说明。但目前，由于本卷必须作一结束，我们只需说：在这起初受造的第一人内，就人性而论，已经奠定了这\OriginalTerm{两座城}{two cities}或两个社会的根基——虽然并非显然可见，却已在天主的\OriginalTerm{预知}{foreknowledge}之中。因为一切人都要从那人衍生：他们中有的将与善天使一同受赏，有的则与恶者一同受罚；这一切都由天主隐秘而正义的审判所安排。因为经上记载：\BibleQuote{上主的一切行径都是仁慈与真理，}祂的恩宠不会不义，祂的正义也不残酷。\par

\SourceRecord{Translated by Marcus Dods.；Nicene and Post-Nicene Fathers, First Series；Vol. 2.；Edited by Philip Schaff.；Buffalo, NY: Christian Literature Publishing Co.,；1887.；<http://www.newadvent.org/fathers/120112.htm>.}

% structure: p0001 source=h1 target=section reason=page-title

\section{\WorkTitle{天主之城（卷十三）}}

本书教导：死亡是惩罚性的，起源于亚当的罪。\par

% structure: p0003 source=h2 target=subsection reason=relative-heading-rank; base=h2

\subsection{第一章 论原祖的堕落，由此而承受了必死性。}

处理完关于世界起源和人类开端的极其困难的问题后，自然的顺序要求我们现在讨论原祖的堕落（我们可以说是原祖们）以及人类死亡的起源与传播。因为天主并没有将人造成像天使那样，即使他们犯罪，也绝不能死亡。祂造他们如此，如果履行服从的义务，就能在无需死亡干预下获得天使般的不朽和永福；但如果他们不服从，死亡就会以公义的判决临到他们——这也是在前一卷中已经谈到的。\par

% structure: p0005 source=h2 target=subsection reason=relative-heading-rank; base=h2

\subsection{第二章 论能影响不朽灵魂的死亡，以及身体所服从的死亡。}

但我觉得必须更仔细地谈论死亡的性质。虽然人的灵魂确实被认为是不朽的，但它也有其自身的某种死亡。它之所以被称为不朽，是因为在某种意义上它不会停止生活和感知；而身体被称为可朽的，是因为它可以被所有生命抛弃，并且自身根本无法生活。因此，灵魂的死亡发生在天主离弃它时，如同身体的死亡发生在灵魂离弃它时。因此，两者的死亡——即整个人——发生在灵魂被天主离弃时，灵魂离弃了身体。因为在这种情况下，天主不再是灵魂的生命，灵魂也不再是身体的生命。而全人的这种死亡之后，就是那凭神圣启示的权威我们称为第二次死亡的死亡。救主在说\BibleQuote{你们要害怕那能在地狱里毁灭灵魂和身体的}时所指的就是这个。\BibleRef{玛 10:28}既然这种死亡不会发生在灵魂与身体紧密结合以致它们根本无法分离之前，那么身体如何能被那种死亡杀死呢？在那种死亡中，它并未被灵魂抛弃，反而被灵魂赋予生命和感觉而受折磨。因为在那种惩罚性的永久刑罚中（我们将在适当的地方更详细地谈论），灵魂被公义地称为死亡，因为它不与天主相连；但我们怎么能说身体死了呢？既然它因灵魂而活着？否则它不会感受到复活后身体将遭受的痛苦。是不是因为各种生命都是好的，而痛苦是恶的，所以我们不愿意说那身体活着，其中灵魂不是生命的原因，而是痛苦的原因呢？因此，当灵魂生活良好时，它藉天主而生活，因为除非天主在其中运作善事，它无法生活得好；而身体藉灵魂而生活，当灵魂在身体中生活时，无论灵魂自身是否藉天主而生活。因为恶人在身体中的生活不是灵魂的生活，而是身体的生活，即使是已死的灵魂——即被天主离弃的灵魂——也能赋予身体，无论它们保留了多少自身不朽的适当生命。但在最后的惩罚中，虽然人不会停止感觉，但由于他的感觉既没有快乐甜蜜，也没有安息有益，而是痛苦惩罚，因此不无理由地被称为死亡而不是生命。它被称为第二次死亡，因为它跟随第一次死亡，第一次死亡分离了两个结合的本质，无论是天主和灵魂，还是灵魂和身体。因此，关于第一次和身体的死亡，我们可以说：对善人是好的，对恶人是坏的。但毫无疑问，第二次死亡，既然它不会临到任何善人，所以对任何人都不可能是好的。\par

% structure: p0007 source=h2 target=subsection reason=relative-heading-rank; base=h2

\subsection{第三章 论因我们原祖的罪过而传遍所有人的死亡，即使是善人，也作为罪的惩罚吗？}

但是，有一个不可回避的问题：使灵魂与身体分离的死亡，对善人来说是否确实是有益的？如果是，那么死亡这样的事怎能成为罪的惩罚呢？因为原祖若没有犯罪，就不会遭受死亡。如此，那只能发生在恶人身上的事，怎能对善人有益呢？再者，如果死亡只能发生在恶人身上，那么对善人而言，它应当不是有益，而是根本不存在。因为既然没有要惩罚的对象，为何要有惩罚呢？因此我们必须说：原祖确实被创造得如此，即如果他们未曾犯罪，就不会经历任何形式的死亡；但自从他们成为罪人后，就被死亡如此惩罚，以致凡从他们血统所生的，也都要受同样的死亡惩罚。因为除了他们自己已经成为的之外，从他们不能再生出别的。他们的本性依照其罪之判决的严重程度而败坏，以至于在首先犯罪者那里作为惩罚而存在的东西，在他们的后代身上成了自然的结果。因为人不是像从尘土造人那样由人生人。尘土是人被造的材料；人则是生人的父母。因此，土地与血肉并非同一回事，虽然血肉是由土地而成。但父母如何，子女也如何。所以，在原祖身上已经包含了整个人性，这人性要藉女人传递给后代——当那婚姻的结合受到天主对其自身的判决时；而人通过犯罪和受罚所变成的样子，而不是被造时的样子，正是他以其罪与死的根源所传播的。因为无论是犯罪还是受罚，都没有使他自身降低到我们在儿童身上看到的那种婴孩的软弱无助。因为天主安排，婴儿开始在世界上生活时，如同兽类的幼崽一样，因为他们的父母在生活方式和死亡方式上已经堕落到兽类的水平；正如经上所载：\BibleQuote{人虽居于尊位，竟不觉悟，反成了无知的畜类。} 不仅如此，我们看到，婴儿在肢体的使用和活动上甚至比最娇嫩的其它动物幼崽更为软弱，在择取和拒绝上也更为孱弱；仿佛居住在人性中的力量注定要超越所有其它生物，而其能量被抑制得越久，发挥的时间越延迟，箭就射得越高，就像弓拉得越满，箭飞得越高。原祖并没有因自己的放肆妄为和公正判决而陷入这种婴儿的孱弱；但人性在他身上被败坏和改变到了如此程度，以至于他在肢体中感受了悖逆的情欲的争斗，并屈从于死亡的必然性。而他自己因罪与罚所变成的，正是他生育的后代成为的——即屈从于罪与死。如果婴儿因救赎者的恩宠而从罪的奴役中被释放，他们只会遭受这种使灵魂与身体分离的死亡；但由于他们从罪债中被赎回，他们不会进入那第二次的、无止尽的、惩罚性的死亡。\par

% structure: p0009 source=h2 target=subsection reason=relative-heading-rank; base=h2

\subsection{第四章——为何作为罪罚的死亡，不向那些藉重生之恩而获免罪债的人免除？}

此外，若有人对此感到疑虑：既然死亡是罪的真正惩罚，为何那些藉恩宠而消除罪责的人仍要受死？这一难题已在我们另一部论婴儿洗礼的著作中处理并解决了。其中说，灵魂与身体的分离仍被保留，尽管它与罪的关联被除去，理由如下：如果身体的永生立即随着重生的圣事而来，那么信德本身就会因此削弱。因为只有在望德中等待那尚未在实体中见到的事物时，才算是信德。而正是藉着信德的活力与奋斗，至少在以往的年代里，对死亡的恐惧被克服了。这一点在圣殉道者身上尤为显著。如果在重生之洗之后，圣人们不能遭受身体的死亡，那么他们就不可能取得胜利，不可能获得荣耀，甚至不会有任何战斗。那样，谁不会带着受洗的婴儿一起奔往基督的恩宠，以便自己不离开身体呢？这样，信德就不会因未见之赏报而受考验，也就根本不成为信德，因为它在寻求并获得其行为的立即报酬。但现在，藉着救主更伟大、更令人钦佩的恩宠，罪的惩罚被转化为正义的服役。因为那时向人宣告说：\Quote{如果你犯罪，就必死；}如今向殉道者说：\Quote{死，以免你犯罪。}那时说：\Quote{如果你违犯诫命，就必死；}如今说：\Quote{如果你逃避死亡，你就违犯诫命。}从前被定为恐怖的对象，使人不致犯罪的事，如今如果我们要不犯罪，就得忍受它。这样，藉着天主不可言喻的慈爱，甚至邪恶的惩罚本身也成了德行的甲胄，罪人的惩罚成了义人的赏报。因为那时因犯罪而招致死亡，如今因死亡而成就正义。在圣殉道者身上正是如此：因为迫害者向他们提出选择：背教或死亡。因为义人宁愿藉着信仰而遭受那些首先背命者因不信而遭受的苦难。因为除非他们犯罪，他们不会死；但殉道者若不死，就是犯罪。前者因犯罪而死，后者因死而不犯罪。因前者的罪债招致惩罚；因后者的惩罚避免了罪债。并不是说死亡，那以前是恶的，变成善的；而是天主赐予信德这种恩宠：死亡，那公认与生命对立的事，竟成为达到生命的工具。\par

% structure: p0011 source=h2 target=subsection reason=relative-heading-rank; base=h2

\subsection{第五章——正如恶人滥用了本是善的法律，同样，善人善用了本是恶的死亡。}

宗徒愿意显示罪恶是多么有害的东西，当恩宠不帮助我们时，他不迟疑地说，罪的力量正是那禁止罪的法律。\BibleQuote{死亡的刺就是罪，而罪的力量就是法律。}\BibleRef{格前 15:56}这是千真万确的；因为禁令增加了对不法行为的欲望，如果义德不被爱到足以战胜罪的欲望。但除非神圣的恩宠帮助我们，我们不能爱慕真义德，也不能在其中喜乐。但为了避免法律被认为是恶的，因为它被称为罪的力量，宗徒在另一处处理类似问题时说：\BibleQuote{法律确实是圣的，诫命也是圣的、公义的、善的。那么，那圣的竟然成了我的死亡吗？绝不然！但罪，为了显出是罪，藉着那善的在我内造成了死亡，好使罪因着诫命成为极其罪恶的。}\BibleRef{罗 7:12}他所说的\Quote{\Emph{极其}}，是因为当经由罪的增加的情欲，法律本身也被藐视时，过犯就更加可恶了。我们为什么认为值得一提呢？因为，正如法律在增加犯罪者的情欲时不是恶的，同样，死亡在增加忍受它者的光荣时也不是善的，因为前者被邪恶地离弃，使人成为违犯者，后者为了真理而被拥抱，使人成为殉道者。因此，法律的确是善的，因为它是禁止罪恶的，死亡是恶的，因为它是罪的工价；但正如恶人不只对恶的事，也对善的事作恶用，同样，义人不只对善的事，也对恶的事作善用。由此导致，恶人滥用那本是善的法律，而善人善终，尽管死亡是恶的。\par

% structure: p0013 source=h2 target=subsection reason=relative-heading-rank; base=h2

\subsection{第六章——论死亡之恶的一般概念，视之为灵魂与身体的分离。}

因此，就肉身的死亡而言，即灵魂与身体的分离，在被我们称为临终的那些人承受时，对任何人都不是善的。因为身体与灵魂被强行扯开的那种暴力，在活着时它们是紧密结合、互相交织的，带来了严酷的体验，只要持续存在，就可怕地冲击着本性，直到完全失去感觉，而这感觉原本来自于精神与肉体的相互渗透。而所有这些痛苦有时被身体的突然一击或灵魂的急速飞离所预防，其迅速使得痛苦不被感觉到。但无论临终时那以剧烈痛苦的感觉剥夺所有感觉的是什么，然而，当它被虔诚而忠信地忍受时，它增加了忍耐的功绩，但并未使惩罚之名不适用。死亡，通过来自第一人的普通生殖而传递，是所有由他而生者的惩罚，然而，如果为了义德而忍受，它就变成那些重生者的光荣；并且尽管死亡是罪所应得的报应，它有时却确保没有任何应得报应归于罪。\par

% structure: p0015 source=h2 target=subsection reason=relative-heading-rank; base=h2

\subsection{第七章——论未受洗者为宣认基督而遭受的死亡。}

对于凡是未受洗而死于宣认基督的人，这一宣认对于赦免他们的罪具有同等效力，就如他们在圣洗圣事的圣泉中受洗一样。因为那一位曾说过：\BibleQuote{人除非由水和圣神而生，不能进天主的国。}\BibleRef{若 3:5}也为他们开了一个例外，就在祂同样绝对说的另一句话中：\BibleQuote{凡在人前承认我的，我在天上的父前也必承认他；}\BibleRef{玛 10:32}在另一处又说：\BibleQuote{凡为我的缘故丧失性命的，必要获得性命。}\BibleRef{玛 16:25}这就解释了这句话：\Quote{上主视为珍贵的，是祂圣人们的死亡。}因为有什么比一种死亡更珍贵，藉此一个人的所有罪都被赦免，他的功绩增加百倍？因为那些在无法逃脱死亡时才受洗，并带着所有罪过被消除而离世的人，他们的功绩并不等于那些没有推迟死亡——尽管他们有能力这样做——却宁愿以宣认基督结束生命，而不愿以否认祂来确保受洗机会的人。而且，即使他们在死亡的恐惧压力下否认了祂，这也会在那圣洗中得到赦免，因为在圣洗中连那些杀害基督的人的巨大邪恶也被赦免了。但在这些人身上，圣神的恩宠必定是何等丰沛，祂随意吹拂，因为他们如此热切爱慕基督，以至于在如此严峻的紧急关头不能否认祂，并且怀着如此确切的宽恕希望！因此，圣人的死亡是珍贵的，基督的恩宠以如此恩惠的效用施予他们，以至于他们毫不迟疑自己去面对死亡，好能与祂相遇。它也是珍贵的，因为它证明了那原本是为惩罚罪人而设立的，已被用于产生更丰硕的义德收获。但我们不应因此就把死亡视为善事，因为它被转用于如此有益的目的，并非由于它自身的任何德能，而是由于天主的干预。死亡最初被提出作为恐惧的对象，免得犯罪；现在它必须被经历，以免犯罪，或如果犯了罪，得蒙赦免，并将义德的赏报赐给那以胜利赢得它的人。\par

% structure: p0017 source=h2 target=subsection reason=relative-heading-rank; base=h2

\subsection{第八章——论圣人因真理而忍受第一次死亡，得以摆脱第二次死亡。}

因为如果我们稍加仔细地思考这事，就会看到，即使一个人为了真理而忠信地、可赞美地死去，他仍是在躲避死亡。因为他接受一部分死亡，正是为了躲避全部的死亡，以及外加的\OriginalTerm{第二次且永久的死亡}{second and eternal death}。他接受灵魂与肉身的分离，免得灵魂既与天主也与肉身分离，从而使\OriginalTerm{第一次死亡}{first death}完全实现，而第二次死亡永远地接纳他。因此，正如我所说，死亡在正被实际承受时，在它正使垂死者屈服于其权下时，对任何人都不是好的；但它为了保持或获得那\Emph{是}好的而受之有功德。至于死后所发生的事，说死亡对好人是好的，对恶人是坏的，并非荒谬。因为义人的脱体灵魂安息；但恶人的灵魂受罚，直到他们的肉身复活——义人复活进入永生，其他人复活进入永死，即所谓的第二次死亡。\par

% structure: p0019 source=h2 target=subsection reason=relative-heading-rank; base=h2

\subsection{第九章。——我们是否应该说死亡的那一刻（感觉停止之时）是发生在垂死者还是死者的经验中？}

善人和恶人的灵魂与肉身分离的时刻，我们是否应该说它是在死后还是正在死亡中？如果是在死后，那么死亡本身就不是善或恶，因为死亡已成为过去；而是灵魂现在进入的生命。死亡在它临在时是恶的，也就是说，当它被垂死者承受时；因为死亡给他们带来了严厉而痛苦的经历，而善人善用了它。但当死亡过去后，那不再存在的东西怎么能是善或恶呢？更进一步，如果我们更仔细地考察这事，我们将会看到，即使是垂死者所经历的那种剧烈而痛苦的感觉也不是死亡本身。因为只要他们有任何感觉，他们就肯定还活着；如果还活着，则更应说他们处于死亡之前的状态，而非死亡状态。因为当死亡实际来临时，它剥夺了我们所有肉身的感觉，而死亡在临近时是痛苦的。因此很难解释，我们如何称呼那些尚未死去，但在最后的、致命的痛苦中挣扎的人是在死亡的关头。然而，我们还能称他们为什么，只能说是垂死的人？因为当迫在眉睫的死亡实际来临时，我们不能再称他们为垂死者，而是死者。所以没有人是在活着的同时又是垂死的；因为即使是那些处于生命最后一刻、如我们所说正在断气的人，也还活着。因此同一个人同时是垂死和活着，但正接近死亡，离弃生命；然而活着，因为他的灵魂仍然停留在肉身内；尚未死亡，因为他的灵魂还没有离开肉身。但是，如果当灵魂离开时，那人甚至不在死亡中，而是在死后，那么谁会说他在死亡中呢？一方面，如果一个人不能同时既垂死又活着，那么没有人可以被称作垂死的；只要灵魂在肉身内，我们就不能否认他活着。另一方面，如果那接近死亡的人更应被称为垂死的，那么我不知道谁是活着的。\par

% structure: p0021 source=h2 target=subsection reason=relative-heading-rank; base=h2

\subsection{第十章。——论必死者的生命，它与其被称为生命不如被称为死亡。}

因为我们刚一在这必死的身体中开始生活，就开始不停地走向死亡。因为在生命的整个历程中（如果我们必须称它为生命的话），它的可变性趋向死亡。当然，没有人不是今年比去年更接近死亡，明天比今天更接近，今天比昨天更接近，一会儿之后比现在更接近，现在比一会儿之前更接近。因为我们活着的任何时间都是从我们的整个寿数中扣除的，剩余的部分日渐减少；因此我们整个生命无非是一场奔向死亡的竞赛，在其中没有人被允许稍作停留或走得更慢一些，而是所有人都被一种平等的运动以相等的速度向前驱赶。因为寿命短的人度过一天并不比寿命长的人更快。但是，当相等的时刻被平等地从两人身上夺走时，一个更近，另一个更远，以同样的速度达到目标。走更长的路程是一回事，走得更慢是另一回事。因此，在通往死亡的路上花更长时间的人，并不是以更悠闲的步伐前进，而是走过更多的里程。此外，如果每个人一当死亡开始在他身上显现（即通过夺取生命）就开始死去，也就是说，处于死亡中——因为当生命全部被夺走时，那人将不是处于死亡中，而是处于死后——那么他一当开始生活就开始死去。因为在他所有的日子、小时和时刻中，直到这缓慢作用的死亡完全完成，还有什么在进行呢？然后，死后的时刻到来，取代那生命正在被撤走的状态，我们称之为在死亡中。那么，从人居住在这必死而非活的身体中的那一刻起，他就从未处于生命中——至少如果他不能同时处于生命和死亡中。或者，我们是否应该说，他同时处于两者中？即处于生命中，因为他活着直到一切耗尽；但也处于死亡中，因为他随着生命的消耗而死去？因为如果他不在生命中，那被耗尽直至消失的是什么？如果他不在死亡中，那消耗本身又是什么？因为当整个生命耗尽后，\Quote{死后}这个说法将毫无意义，如果那消耗不是死亡的话。而且，如果当一切耗尽后，人不在死亡中而是在死后，那么他何时处于死亡中，除非当生命正在被消耗时？\par

% structure: p0023 source=h2 target=subsection reason=relative-heading-rank; base=h2

\subsection{第十一章。——一个人能否同时既是活着的又是死去的？}

但若说一个人在抵达死亡之前已在死亡之中，这是荒谬的（因为当他度过一生时，他正在奔向何处？若他已身在死亡中？），并且若说一个人同时活着又死了，这与常说同时睡着又醒着一样违反常理。那么，就不得不问：一个人何时正在死去？因为，在死亡来临之前，他不是正在死，而是活着；当死亡来临之后，他不是正在死，而是死了。一个是死亡之前，另一个是死亡之后。那么，他何时处在死亡之中，以致我们能说他正在死去？因为正如有三个时间：死亡之前、死亡之中、死亡之后，同样也有三种相应状态：活着、正在死、死了。而很难界定一个人何时在死亡之中或正在死去，因为此时他既非活着（那是死亡之前），也非死了（那是死亡之后），而是正在死（那是死亡之中）。因为只要灵魂在身体内，尤其是尚有意识时，人显然活着；因为身体和灵魂构成人。因此，死亡之前不能说他在死亡之中；而另一方面，当灵魂离去、一切身体感觉消失后，死亡已成过去，人便死了。在这两种状态之间，找不到正在死的状况；因为若人尚活着，死亡尚未到达；若他已停止活着，死亡已成过去。因此，他从未正在死去，即从未处在死亡状态中。同样，在时间的流逝中——你试图抓住当下，却找不到它，因为当下不占据空间，只是时间从未来到过去的过渡。那么，我们是否必须得出结论：根本没有身体死亡这回事？因为如果有，它在何处？既然它不在任何人身上，也没有人能在它里面？确实，如果生命仍在，死亡尚未到来；因为这是死亡之前的状态，而非死亡之中；如果生命已经停止，死亡并不存在；因为这是死亡之后的状态，而非死亡之中。另一方面，如果没有死亡之前或之后，那么我们说\Quote{死亡之后}或\Quote{死亡之前}是什么意思？如果不存在死亡，这种说法是愚蠢的。但愿我们在乐园里生活得如此之好，以至于真的根本没有死亡！但如今死亡不仅存在，而且是如此痛苦的事，以至于没有任何技巧足以解释或逃避它。\par

让我们按习惯的说法——谁也不该另样说——把死亡来临之前的时间称为\Quote{死亡之前}；正如经上所载：\BibleQuote{未死之前，不可赞美人。} \BibleRef{德 11:28} 当死亡发生后，让我们说\Quote{死亡之后}发生了这样那样的事。至于当下，我们尽量这样说，例如，我们说：\Quote{他临死时立了遗嘱，把这样那样的东西留给了某某人}——当然，除非他活着，他不能这样做，而这与其说是在死亡中，不如说是死亡之前。让我们使用圣经所用的同样说法；因为圣经毫不迟疑地说死人不是在死亡之后，而是在死亡之中。正如那节经文：\Quote{因为在死亡中，无人纪念你。} 因为在复活之前，人们恰当地被视为在死亡之中；正如一个人在醒来之前被视为在睡眠之中。然而，尽管我们可以说睡眠中的人正在睡觉，但我们不能这样对死者说，说他们正在死亡。因为，就我们现在所谈论的身体死亡而言，不能说自己经与身体分离的人还在不断死亡。但你看，这正是我所说的——没有任何言语能解释，如何可以说正在死的人还活着，或如何可以说死者即使在死亡之后仍处于死亡之中。因为如果他们处于死亡之中，怎么会在死亡之后呢？尤其当我们甚至不称他们为\Quote{正在死}，就像我们称睡眠中的人为\Quote{正在睡}，称困倦的人为\Quote{正在困倦}，称悲伤的人为\Quote{正在悲伤}，称活着的人为\Quote{正在活着}？然而，死者直到复活之前，被称为在死亡之中，但不能称为正在死。\par

因此我认为，虽然并非出于人的意图，但或许出于天主的旨意，发生了一件不无合适、不无恰当的事：拉丁词 \OriginalTerm{死}{moritur} 不能按照类似词的规则由语法学家进行变位。因为 \OriginalTerm{升起}{oritur} 的完成时形式是 \OriginalTerm{已升起}{ortus est}；所有类似动词都从它们的完成分词构成这个时态。但若问 \OriginalTerm{死}{moritur} 的完成时，我们得到常规的回答 \OriginalTerm{已死}{mortuus est}，带有双 u。因为 \OriginalTerm{死的}{mortuus} 的发音就像 \OriginalTerm{愚拙的}{fatuus}、\OriginalTerm{高耸的}{arduus}、\OriginalTerm{显眼的}{conspicuus} 以及类似词，这些词不是完成分词而是形容词，且不考虑时态而变格。但 \OriginalTerm{死的}{mortuus} 虽然形式上是一个形容词，却被用作完成分词，仿佛那本身不能变位的东西可以被变位；因此，恰好发生了这样的事：正如事情本身实际上不能变位，同样表示该行为的词也不能变位。然而，藉我们救赎者恩宠的帮助，我们至少可以设法变位第二种死亡。因为第二种更加痛苦，实际上是所有邪恶中最糟糕的，因为它不在于灵魂与身体的分离，而在于两者在永恒死亡中的结合。在那里，与我们现今的状况形成鲜明对比，人将不会在死亡之前或之后，而永远在死亡之中；因此永不活着，永不死去，而是无尽地死去。而当死亡本身成为不死时，没有比这更惨烈地处于死亡之中了。\par

% structure: p0027 source=h2 target=subsection reason=relative-heading-rank; base=h2

\subsection{第十二章——天主以死亡警告我们原祖父母，若他们违犯祂的命令，天主所意指的是哪种死亡。}

因此，当有人问天主用以威胁我们原祖父母的是什么死亡——如果他们违犯所领受的诫命，未能保持服从——是灵魂的死亡，还是身体的死亡，还是整个人（的死亡），或是那被称为第二次死亡的死亡？我们必须回答：是全部。因为第一次死亡由两种组成；第二次死亡是完整的死亡，由所有死亡组成。因为正如整个大地由许多土地组成，普世教会由许多教会组成，完整的死亡也由所有死亡组成。第一次死亡由两种组成：一种是身体的，另一种是灵魂的。因此第一次死亡是整个人（的死亡），因为灵魂没有天主、没有身体，暂时受罚；但第二次死亡是当灵魂没有天主、但有身体时，永远受罚。因此，当天主对那第一位被安置在乐园的人提到禁果时说：\BibleQuote{在你们吃的那一天，你们必定死。}\BibleRef{创 2:17}那威胁不仅包括第一次死亡的第一部分——灵魂被剥夺天主；也不仅包括第一次死亡的后续部分——身体被剥夺灵魂；也不仅包括第一次死亡的整体——灵魂在与天主和身体分离时受罚；而是包括了所有死亡，直到那被称为第二次的最终死亡，之后不再有死亡。\par

% structure: p0029 source=h2 target=subsection reason=relative-heading-rank; base=h2

\subsection{第十三章——我们原祖父母违命的首要惩罚是什么}

因为，我们原祖父母一违犯了诫命，天主的恩宠就离开了他们，他们因自己的邪恶而羞愧；于是他们拿无花果树叶（可能是他们忧虑心情中首先碰到的东西），遮盖了自己的羞耻；因为他们的肢体虽然依旧，他们现在却感到羞耻，而先前没有。他们体验到自己肉体的一种新冲动，这肉体已不再服从他们，作为他们自己不服从天主的严格报应。因为灵魂放纵自己的自由，不屑于服侍天主，它自身被剥夺了原来对身体的统御。又因为它故意离弃了它的至高主宰，它就不再管束它的低等仆役；它也不能使肉体服从，这原是它若自身服从天主就永远能做到的。于是\BibleQuote{肉体开始与圣神相争}\BibleRef{迦 5:17}，我们生在这争斗之中，从第一次违命继承了死亡的种子，并在我们的肢体和败坏的本性中带着肉体的争斗甚至胜利。\par

% structure: p0031 source=h2 target=subsection reason=relative-heading-rank; base=h2

\subsection{第十四章——人由天主造于何种状态，又因自己意志的选择堕入何种境况}

因为天主，本性的创造者而非罪恶的创造者，创造了正直的人；但人，因自己的意志败坏，又公正地被定罪，生出了败坏和定罪的子孙。因为我们都在那一个人里，因为我们都是那一个人，他因那在犯罪前由他而造的女人而陷入罪恶。因为那时还没有创造和分配给我们的个别形式，我们作为个体要活在其中，但已经有了我们由以繁衍的种子本性；这本性因罪受损，被死亡锁链束缚，又公正地被定罪，人不能从人而生于别的状态。这样，从自由意志的滥用，产生了整个罪恶的链条，它带着接连不断的苦难，把人类从其败坏的本源，如同从腐败的根，引向第二次死亡的毁灭，这毁灭没有终结，只有那些被天主的恩宠所释放的人除外。\par

% structure: p0033 source=h2 target=subsection reason=relative-heading-rank; base=h2

\subsection{第十五章——亚当在犯罪时离开天主先于天主离开他，而他离开天主正是灵魂的第一次死亡}

也许有人会认为，因为天主说：\BibleQuote{你们必定要死}\BibleRef{创 2:17}，而不是\Quote{各种死亡}，我们应当只理解为当灵魂被它的生命天主所抛弃时发生的那种死亡；因为灵魂不是被天主抛弃而离开祂，而是它离开祂，因而被祂抛弃。因为它自己的意志是它罪恶的起源，正如天主是它趋向善的运动的起源，既在它不存在时创造了它，又在它堕落灭亡时重造了它。但即使我们认为天主只指这种死亡，并且将\BibleQuote{在你们吃的那一天，你们必定要死}理解为此意——即\Quote{在你们因不服从而离开我的那一天，我将公义地离开你们}——然而，在这种死亡中，其他死亡也确实受到了威胁，因为它们是它的必然结果。因为在不顺从灵魂的肉体中第一次感到的悖逆冲动（这使我们原祖父母遮盖自己的羞耻）中，确实体验了一种死亡，即天主离开灵魂时发生的死亡。这由天主的话暗示：当那人因恐惧躲藏时，天主说：\BibleQuote{亚当，你在哪里？}\BibleRef{创 3:9}——祂说这话并非出于无知而询问，而是警告他省察自己在哪里，因为天主不与他同在。但当灵魂因年老腐败而离开身体时，就体验了另一种死亡，即天主在宣判人刑罚时所说的：\BibleQuote{你原是土，仍要归于土。}\BibleRef{创 3:19}这两种死亡构成了那整个人（的）第一次死亡。这第一次死亡最终被第二次死亡接续，除非人因恩宠得到释放。因为身体不会归于其所出的土，只有通过其自身的死亡，即当灵魂——它的生命——离开它时。因此，所有真诚持守公教信仰的基督徒都同意：我们受制于身体的死亡，并非出于自然律——天主原没有为人规定死亡——而是出于祂因罪的公义惩罚；因为天主为罪施行惩罚时，对当时我们众人都在其中的那人说：\BibleQuote{你原是土，仍要归于土。}\par

% structure: p0035 source=h2 target=subsection reason=relative-heading-rank; base=h2

\subsection{第十六章——论那些认为灵魂与身体分离并非惩罚的哲学家，尽管柏拉图描述至高的神向次等的神承诺他们永不从身体中被遣散}

但是那些反对我们为\WorkTitle{天主之城}（即祂的教会）辩护的哲学家，似乎有很好的理由嘲笑我们，因为我们说灵魂与身体的分离应当被视为对人的惩罚的一部分。因为他们认为，只有当灵魂完全脱离身体，以纯洁、简单、如同赤裸的灵魂回归天主时，其福乐才得以完全。关于这一点，如果我在他们自己的文献中找不到任何反驳这种观点的东西，我将不得不费力地证明，不是身体，而是身体的腐朽性，才是灵魂的负担。因此，我们在前一部书中引用的圣经句子说：\Quote{\BibleQuote{腐朽的身体压垮了灵魂。}} \BibleRef{智 9:15} \Quote{腐朽}一词被加上以表明，灵魂的负担不是来自任何身体，而是来自因罪而变得如此的身体。即使没有加上这个词，我们也无法理解别的。但当\Person{柏拉图}最明确地宣称由至高者制造的神灵具有不朽的身体，并且当祂介绍它们的制造者本人，应许它们一件伟大的恩赐，即它们将永远存留在它们的身体中，永不因任何死亡而脱离时，为什么我们的这些对手，为了扰乱基督信仰，假装不知道他们很清楚的事，甚至宁愿自相矛盾也不愿失去反驳我们的机会？以下是\Person{柏拉图}的话，按\Person{西塞罗}的翻译，其中他介绍至高者对祂所造的神灵说：\Quote{你们这些出于神圣族类的，要思考我是哪些作品的亲生父母和作者。这些（你们的身体）只要我愿意，就是不可毁坏的；尽管一切组合而成的都可以被毁坏。但是，解散理智所结合的东西是邪恶的。但是，既然你们已经出生，你们确实不能是不朽和不可毁坏的；然而你们决不会被毁坏，也没有任何命运将你们交付给死亡，并证明比我的意志更强，而我的意志是对你们永续性的更强保证，比你们出生时与之结合的身体更强。} 你看，\Person{柏拉图}说，神灵因身体与灵魂的结合而是可朽的，但因创造者的意志和命令而成为不朽的。因此，如果与任何身体结合是对灵魂的惩罚，为什么天主对它们说话，好像它们害怕死亡，即灵魂与身体的分离？为什么祂寻求安抚它们，应许它们不朽，不是基于它们的本性（本性是组合的而非单纯的），而是基于祂不可战胜的意志，借此祂能够使出生的东西不死，组合的东西不被解散，而是永久保存？\par

\Person{柏拉图}关于星体的这种观点是否正确，是另一个问题。因为我们不能立即同意他，这些发光的形体或球体，白天和黑夜以其物质实体的光照耀大地，也拥有理智和福乐的灵魂，每个灵魂赋予其身体生命，正如他自信地断言宇宙本身是一个巨大的动物，其中包含所有其他动物。但正如我说过的，这是另一个问题，我们目前不打算讨论。我只认为有必要提出这一点，以反对那些因自己是或被称为柏拉图主义者而如此骄傲，以至于以作基督徒为耻，并且不能忍受被称为普通人也有的名字，以免他们使哲学家的圈子庸俗化，而哲学家的圈子因其排他性而自豪。这些人寻找基督教教义的弱点，选择攻击身体的永恒性，好像这是一种矛盾：既主张灵魂的福乐，又希望它永远居于身体之中，仿佛被一条可悲的锁链束缚着；然而他们的创始人兼大师\Person{柏拉图}却断言，至高者作为一种恩赐赐给祂所造的神灵，使它们不会死，即不会与祂所连接它们的身体分离。\par

% structure: p0038 source=h2 target=subsection reason=relative-heading-rank; base=h2

\subsection{第17章——反对那些断言尘世身体不能变得不腐朽和永恒的人。}

这些同样的哲学家进一步争辩说，地上的身体不能是永恒的，尽管他们毫不怀疑整个大地，即他们神祇的中心部分——诚然，不是最伟大的神，而是一位伟大的神，即整个世界的中心部分——是永恒的。既然至高者为他们造了另一位神，即这个世界，高于祂之下的众神；既然他们设想这位神是一个有生命的个体，拥有，如他们所称，一个理性或理智的灵魂包藏在其庞大躯体的巨大体积中，并拥有，作为其躯体适当安置和调整的成员，四大元素，其结合他们希望是不可分解和永恒的，以免他们这位伟大的神某天或许会毁灭；有什么理由能使大地，即一个更伟大造物的躯体中的中心成员，是永恒的，而其他地上造物的身体却不能是永恒的，如果天主愿意如此的话？但是，他们说，泥土必须归于泥土，地上动物的身体是从泥土中取出的。因为，他们说，这就是它们死亡和解体的必然性的原因，也是它们回归到所来自的坚固永恒的大地的方式。但如果有人对火说同样的话，认为从火中产生以构成天上存在物的身体必须回归到宇宙之火，那么柏拉图认为这些神从至高者那里接受的永生，难道不在这场争论的热度中蒸发消失吗？或者，这些天上存在物不发生这种情况，是因为天主（其意志，如柏拉图所说，压倒一切力量）不乐意如此？那么，什么能阻止天主对地上的身体也这样规定呢？既然柏拉图承认，天主能够阻止出生之物死亡，阻止结合之物分离，阻止组成之物分解，并能规定灵魂一旦分配给它们的身体就永不离开，而与它们一同享有永生和永远的福乐，为什么祂不能也使地上的身体不死呢？难道天主无力成就基督徒信仰所特有的一切，而有力成就柏拉图主义者所欲求的一切吗？那些哲学家，当真，得以认识天主旨意和权能，而先知却被拒绝认识！事实是，天主的圣神教导祂的先知们祂的旨意，祂认为合适启示多少，就启示多少；而哲学家们在努力发现它时，却被人的猜测所欺骗。\par

但他们本不应如此被引入歧途——我不说由于他们的无知，而是由于他们的固执——以至于如此频繁地自相矛盾；因为他们全力主张，为灵魂的幸福，它必须放弃的不仅是地上的身体，而且是任何种类的身体。然而他们又坚持认为，那些灵魂最为有福的神祇，是被束缚于永恒的身体的：天上的神祇被束缚于火的身体，而朱庇特本人（或世界，如他们要我们相信的）的灵魂则被束缚于构成这整个从地到天的巨大质料的所有物理元素。因为柏拉图相信这个灵魂是由音乐数字扩展和散布的，从大地的内部中心（几何学家称之为圆心）向外通过其所有部分直到天空的最高处和极限；因此，这个世界是一个非常大且幸福的不朽动物，其灵魂既有完美的智慧幸福，又永不离开自己的身体，而其身体因灵魂而有永生，并且绝不拖累或阻碍灵魂，尽管身体本身不是单纯的，而是由如此多且巨大的材料组成的。既然他们容许自己的猜测如此之多，为什么他们拒绝相信，藉着天主的意志和能力，永生可以被赋予地上的身体，在其中灵魂既不被身体的负担所压迫，也不因任何死亡而与之分离，而是永恒而幸福地生活？难道他们不主张他们自己的神祇如此生活在火的身体中，而朱庇特本人，他们的王，如此生活在物理元素中？如果为了灵魂的幸福，它必须离开任何种类的身体，那就让他们的神祇从星界飞走，让朱庇特从地飞到天；或者，如果他们不能这样做，就让他们被宣布为不幸的。但这两种选择这些人都不会采取。因为，一方面，他们不敢将自己的神祇归于离开身体，以免他们似乎崇拜必死之物；另一方面，他们不敢否认他们的幸福，以免他们承认恶棍为神祇。因此，为获得幸福，我们不需要离开任何种类的身体，而只需要离开那腐朽的、累赘的、痛苦的、必死的身体——不是天主的美善为第一人所设计的那样的身体，而只是人的罪所导致的那样的身体。\par

% structure: p0041 source=h2 target=subsection reason=relative-heading-rank; base=h2

\subsection{第十八章——论地上的身体，哲学家们断言其不能在属天之处，因为凡是属于地的都被其自然重量吸引到地。}

但是，他们说，不可避免地，地上物体的自然重量要么使它们停留在地上，要么把它们拉向地上；因此它们不能在天上。我们的原祖父母确实在地上，在一个林木茂盛、果实丰饶的地方，被称为\Place{乐园}。但是，让我们的对手更仔细地思考一下这个关于地上重量的问题，因为它对基督身体的升天以及圣人们的复活身体都有重要的意义。如果人类的技艺可以通过某种设计，用那些一放到水上就会下沉的金属制造出浮动的器皿，那么，天主通过某种隐秘的运作方式，更确定地使这些土质的团块摆脱其重量的向下压力，岂不更加可信？这对于那全能的天主来说，并非不可能；按照\Person{柏拉图}的说法，依祂的旨意，所生之物不会消亡，所合成之物不会分解；尤其是，灵魂与身体本质的结合，比身体与其他物质实体的调整要更加奇妙，那么这一点就更有可能了。我们难道不能轻易相信，变得完全幸福的灵魂将被赋予能力，能以几乎自发的运动随心所欲地移动其土质但不可朽坏的身体，并以最敏捷的动作将它们放置在任何地方？如果天使们能随心所欲地从任何地方运送任何地上的受造物，并把它们带到任何地方，难道我们会相信他们这样做时没有劳苦和负担感吗？那么，我们为什么不能相信，由天主的恩宠变得完美而幸福的圣人们的灵魂，能将自己的身体带到任何地方，并将其安置在任何他们想要的地方呢？因为，尽管我们习惯于注意，在负重时，土质物体越大，重量越重，重量越重，负担越重；然而，灵魂携带自己肉体的肢体，在健康健壮时比在疾病消瘦时更不费力。而且，虽然健康强壮的人被他人抬着时比瘦弱病重的人更重，但这个人自己移动和携带自己的身体时，在健康强壮、体重更大的情况下，比因饥饿或疾病而消瘦到最小程度时，负担感更小。在估计土质物体的重量时，即使它们仍是可朽坏和会死的，所考虑的并不是死重量，而是各部分健康的平衡状态，这是多么重要啊。有什么言语能够描述我们现在所谓的健康与未来的不朽之间的区别呢？那么，让哲学家们不要用有关重量的论据来推翻我们的信仰；因为我不愿探究为什么他们不相信一个土质的身体能在天上，而整个地球却悬在虚无之上。也许世界保持其中心位置，正是依靠吸引所有重物到其中心的同一法则。但是，我要说，如果次级神明——\Person{柏拉图}曾将人类和其他地上受造物的创造委托给他们——如他所断言的那样，能够从火中撤去其燃烧的特性，而保留其照明的特性，使它通过眼睛发光；并且，如果\Person{柏拉图}也承认至高的天主有能力保存所生之物免于死亡，保存由如此不同的部分——身体和灵魂——所合成之物免于分解——那么，我们难道还要犹豫，不承认同一天主有能力对祂所赐予不朽之人的肉身施行作为，撤去其腐败而保留其本性，移除其沉重的重量而保留其美观的形式和肢体吗？至于我们对死者复活的信仰，以及关于他们不朽身体的信仰，若天主愿意，我们将在\WorkTitle{天主之城}的结尾更详细地讨论。\par

% structure: p0043 source=h2 target=subsection reason=relative-heading-rank; base=h2

\subsection{第19章 反对那些不相信原始人类若没有犯罪便会长生不死的人的意见}

目前，让我们继续按已开始的方式，对我们的原祖父母的身体做一些解释。我说，那么，除了作为罪的正当后果外，他们本不会遭受这种死亡，这种死亡对善人是善的——这种死亡，并非只有少数人认识并相信，而是众所周知的，即灵魂与身体分离，一具刚才还活着的动物的身体，现在明显死了。因为，尽管毫无疑义，正义和圣善的死者之灵魂活在平安的安息中，然而，对他们来说，活在健康良好的身体中要好得多，以至于那些认为摆脱任何身体都是最幸福的人，也不由自主地谴责这种意见。因为没有人敢将有智慧的人——无论即将死去或已经死去，换句话说，无论已经摆脱身体或即将摆脱——置于不朽的诸神之上，在柏拉图那里，至高者作为慷慨的恩赐，向诸神承诺了不可分解的生命，或与他们的身体永恒结合。但这位柏拉图认为，对人而言，没有比虔诚而正义地度过一生，在脱离身体后被接纳到永不抛弃身体的诸神怀抱中更好的事了；\Quote{他们忘却过去，可以重访上界，并渴望再次回到身体中。}维吉尔因从柏拉图体系中借用此说而受赞扬。确实，柏拉图认为，凡人的灵魂不能永远留在身体中，必须借由死亡离开；另一方面，他认为没有身体，灵魂不能永远存在，而是通过不断的交替，从生到死，再从死到生。然而，他在这点上区分了智者和其余的人：死后他们被带到星辰，每人可在他适合的星辰上休息一段时间，当他忘却了先前的痛苦并充满对拥有身体的渴望时，便从那里返回到凡人的劳苦和悲惨中。那些生活愚蠢的人则转入适合他们的身体，无论是人的还是野兽的。因此，他给善良和有智慧的灵魂安排了一个非常艰难的命运，因为他们没有得到可以永远甚至不朽地居住的身体，而是既不能永久保留也不能保持永恒纯洁的身体。关于柏拉图的这一观点，我们在前书中已经说过，波菲里在这些基督徒时代感到羞愧，以致他不仅将人类的灵魂从野兽身体的命运中解放出来，还主张智慧的灵魂应从一切身体束缚中解放，这样，摆脱所有血肉，他们作为赤裸而幸福的灵魂，就能与圣父永远同住。而且，为了不显得被基督对祂的圣人永生应许所超越，他还将净化的灵魂安置在无尽的幸福中，不再返回先前的苦难；但是，为了反驳基督，他否认不朽身体的复活，并坚持这些灵魂将永远活着，不仅没有尘世的身体，而且没有任何身体。然而，无论他通过这种教导意指什么，他至少没有教导这些灵魂不应向居住在身体中的诸神献上宗教敬拜。他为什么不？除非因为他并不认为灵魂即使与身体分离也优于那些诸神。因此，如果这些哲学家不敢（我想他们不会）将人类的灵魂置于最幸福却永远与身体结合的诸神之上，那么他们为何认为基督教信仰所宣讲的荒谬呢？即我们的原祖父母被如此创造，如果他们不曾犯罪，就不会因任何死亡而离开他们的身体，而是会因服从而获得不朽作为奖赏，并与他们的身体永远活在一起；此外，圣人们将在复活中拥有他们在此劳苦过的那些身体，但如此这般，既不会有任何朽坏或笨拙附着于他们的肉身，也不会有任何悲伤或烦恼遮蔽他们的幸福？\par

% structure: p0045 source=h2 target=subsection reason=relative-heading-rank; base=h2

\subsection{现今安息于平安的肉身将被提升到一种我们的原祖父母之肉身所不曾享有的完善。}

如此，\OriginalTerm{已亡圣者的灵魂}{souls of departed saints}并不受那令他们脱离身体的死亡所影响，因为他们的肉身在希望中安息，不论它在感觉消失后受到何等侮辱。因为他们并不希望自己的身体被遗忘，如同\Person{柏拉图}所认为的那样；相反，他们记念那位不欺骗人的、连他们头上的头发也保证保全的那位所应许的，他们怀着渴望的耐心，在希望中等待他们身体的\OriginalTerm{复活}{resurrection}，他们曾在这身体中忍受了许多艰苦，现在将永不再受苦。因为如果他们不曾\Quote{憎恨自己的肉身，}当它以其本性的软弱对抗他们的意志，必须受属神的法律约束时，那么当它自身也成为属神的时候，他们岂不更要爱它？因为正如当精神服事肉身时，它适当地被称为\OriginalTerm{属血肉的}{carnal}，同样，当肉身服事精神时，它将被公正地称为\OriginalTerm{属神的}{spiritual}。这并不是说它转化为精神，如同有些人从这话\BibleQuote{它种下的是可朽的，复活起来的是不朽的} \BibleRef{格前 15:42}所幻想的那样，而是因为它以完全而奇妙的乐意顺服于精神，并在一切事上响应那已进入不朽的意志——一切抗拒、一切朽坏、一切迟缓都被除去。因为身体不仅会比今世最佳健康状态更好，而且将超越我们原祖父母犯罪之前的身体。因为，虽然他们除非犯罪不会死亡，但他们仍如现在的人一样食用食物，他们的身体尚非属神的，而只是\OriginalTerm{生物的}{animal}。虽然他们不因年岁而衰败，也不接近死亡——这是天主奇妙的恩宠藉着\OriginalTerm{生命树}{tree of life}（与\OriginalTerm{禁树}{forbidden tree}一同生长在乐园中央）所赐予他们的状态——但他们仍食用其他果实，只是不从那棵禁树取用；那树被禁止，并非因为它本身不好，而是为了推崇一种纯洁而单纯的\OriginalTerm{服从}{obedience}，这是理性的受造物在造物主——他的主——之下所拥有的伟大德性。因为，若没有触碰任何恶事，但若触碰了被禁止的事，这\OriginalTerm{违命}{disobedience}本身就是罪。那么，他们由其他果实滋养，以使其生物性的身体不遭受饥渴的不适；但他们品尝了生命树，好使死亡不能从任何方面悄然临近，使他们不至于因年岁耗尽而衰败。其他果实，可说是他们的滋养，但这生命树却是他们的\OriginalTerm{圣事}{sacrament}。因此，在地堂中，生命树似乎就如天主的智慧在属神领域中所是的那一位，正如经上所载：\BibleQuote{她为把握她的人是一棵生命树} \BibleRef{箴 3:18}。\par

% structure: p0047 source=h2 target=subsection reason=relative-heading-rank; base=h2

\subsection{第二十一章 论乐园：可在属神的意义上理解，而不牺牲关于实际地点的叙述的历史真实性。}

为此，有些人将乐园本身的一切都作\OriginalTerm{寓意解经}{allegorize}，原祖——人类的第一对父母——根据圣经的真理，被记载曾在其中。他们把乐园里所有的树木和结果实的植物理解为德行和生活习惯，好像它们在外部世界并不存在，只是为了属灵的意义而如此叙述或记载。仿佛不可能有一个真实的尘世乐园！仿佛从未存在过这两个女人，\Person{撒辣}和\Person{哈加尔}，以及\Person{亚巴郎}所生的两个儿子，一个出于婢女，一个出于自由妇，因为宗徒说，在她们身上预表了两个盟约；又仿佛水从未从\Person{梅瑟}击打的磐石中流出，因为在那里\Person{基督}可以以预像被看见，正如同一宗徒所说：\BibleQuote{那磐石就是基督} \BibleRef{格前 10:4}。因此，无人否认乐园可以象征真福的生活；它的四条河流，\OriginalTerm{四枢德}{four cardinal virtues}：智德、勇德、节德、义德；它的树木，一切有用的知识；它的果实，虔敬者的习惯；它的\OriginalTerm{生命树}{tree of life}，\OriginalTerm{智慧}{wisdom}本身，万善之母；\OriginalTerm{知善恶树}{tree of the knowledge of good and evil}，\OriginalTerm{违反诫命的经验}{experience of a broken commandment}。天主所规定的惩罚本身就是正义的，因此是善的；但人的经验却不是善的。\par

这些事也可以更有效地理解为关于教会，从而成为未来之事的预言性预兆。因此，乐园就是教会，正如在\BibleBook{}{雅歌}中所称的；\BibleRef{歌 4:13}乐园的四条河流就是\OriginalTerm{四福音}{four gospels}；结果实的树是圣人，他们的果实是他们的善工；生命树是\OriginalTerm{至圣所}{holy of holies}，\Person{基督}；知善恶树是\OriginalTerm{意志的自由选择}{will's free choice}。因为人若轻视天主的旨意，只能毁灭自己；如此，他学到将自己奉献于公共善与沉溺于自我之间的区别。因为那爱自己的人被抛弃给自己，为使他被恐惧和忧伤压倒时，如果他的灵魂尚能感受他的痛苦，就可以用圣咏的话呼号：\Quote{我的灵魂在我内颓丧。}而在受管教后，可以说：\Quote{因他的力量，我将等候你。}这些以及类似的寓意诠释可以恰当地套用于乐园，而不冒犯任何人，同时我们相信历史的严格真实性，并由其具体的事实叙述所证实。\par

% structure: p0050 source=h2 target=subsection reason=relative-heading-rank; base=h2

\subsection{第二十二章 论圣人的身体在复活后将是属灵的，然而肉身并不变为灵。}

这样，义人的身体——他们复活后的身体——既不需要任何果实来防止疾病或衰老的腐朽，也不需要任何其他物质营养来缓解饥渴；因为他们将披戴如此确实和各方面不可侵犯的不朽，以至于他们只在愿意时才进食，并没有进食的必要，同时享有进食的能力。天使向人显现并让人看见和触摸时也是如此，并非因为他们不能以别的方式，而是因为他们能够并愿意藉着某种人形的服务来适应人。我们也不应当认为，当人们款待天使时，天使只是表面上进食，尽管对于那些不知道他们是天使的人，他们可能看起来像我们一样出于必要而进食。因此，在\WorkTitle{多俾亚传}中记载的话：\Quote{你看见我吃，但你看见的不过是幻象；} \BibleRef{多 12:19} 也就是说，你以为我像你一样为了恢复体力而进食。但如果关于天使，另一种观点似乎更可辩护，那么我们的信德毫无疑问地确信主自己，甚至在他复活之后，在属神但仍是真实的身体中，他同门徒们一起吃喝；因为从这些身体中被除去的是吃喝的必然性，而非能力。这样，它们将是属神的，不是因为它们不再是身体，而是因为它们将因赋予生命的圣神而存活。\par

% structure: p0052 source=h2 target=subsection reason=relative-heading-rank; base=h2

\subsection{第二十三章——关于属生灵的身体和属神的身体，或关于那些在亚当内死去和在基督内活着的人，我们该如何理解？}

因为我们那些有灵魂的身体（虽然还不是赋予生命的圣神）被称为有魂的身体，然而它们不是灵魂而是身体；同样，那些身体被称为属神的——但求不要让我们因此认为它们是神灵而不是身体——它们被圣神所赋予生命，具有肉体的实质，但没有肉体的笨重和腐朽。那时人将不再是属土的，而是属天的——不是因为身体不是那最初由土所造的身体，而是因为藉着天的禀赋，它将成为天堂的适宜居住者，这不是由于失去本性，而是由于改变性质。第一个人，出于地，属于土，被造为有生命的灵魂，而不是赋予生命的圣神——这等级是保留给他作为服从的赏报的。因此，他的身体需要饮食以满足饥渴，没有绝对和不可朽坏的不朽，而是藉着生命树避免死亡的必然性，并这样保持在青春的花季——我说，这个身体无疑不是属神的，而是属生灵的；然而除非它得罪天主而招致所警告的惩罚，它本不会死。虽然即使在乐园外，也不拒绝供应他食物，但一旦禁止他靠近生命树，他就被交托给时间的消耗——至少就那种生命而言，若他没有犯罪，他本可以在乐园中永久地保持它，尽管只是在属生灵的身体中，直到因服从而变得属神的那时候。\par

因此，虽然我们明白这种明显的死亡（即灵魂与身体的分离）也被天主在他说\BibleQuote{你吃的日子必定死} \BibleRef{创 2:17} 时所预示，但我们不应因此认为他们在吃了那被禁止的、带来死亡的果子当天就离开了身体。因为就在那一天，他们的本性被恶化、被败坏，并且由于他们被极公正地从生命树驱逐，他们被卷入甚至身体死亡的必然性中，我们正是在这种必然性中出生的。因此，宗徒不是说\Quote{身体因罪注定要死}，而是说\Quote{身体因罪是死的}。然后他加上：\BibleQuote{但若那使耶稣从死者中复活者的圣神住在你们内，那使基督从死者中复活的，也必要藉着那住在你们内的圣神，使你们必死的身体活过来。} \BibleRef{罗 8:10} 那时，身体将成为赋予生命的圣神，而现在它是有生命的灵魂；然而宗徒称它为\Quote{死的}，因为它已经处于必死的必然性之下。但在乐园中，它被造为有生命的灵魂，虽然不是赋予生命的圣神，以致不能正确地称为死的，因为除非犯罪，它不会陷入死亡的权势。现在，天主以\Quote{亚当，你在哪里？}的话指示了灵魂的死亡（那是当他遗弃它时发生的），并以\BibleQuote{你既是土，还要归于土} \BibleRef{创 3:19}的话指示了身体的死亡（那是当灵魂离开它时发生的）。因此，我们被引导相信他对第二次死亡什么也没说，愿意将它隐藏，为新约时代保留，在那里它最清楚地被启示。他这样做，为的是首先可能明显，这种第一次死亡（对众人是普遍的）是那在一个人身上成为所有人的罪的后果。但第二次死亡并非对所有人普遍，那些\BibleQuote{按天主旨意蒙召的人。因为他所预知的人，也预定他们与他自己儿子的肖像相同，使他可以在许多弟兄中作长子。} \BibleRef{罗 8:28}除外。天主的恩宠已藉着一位中保，将这些人从第二次死亡中解救出来。\par

因此，宗徒指明第一个人是在属生灵的身体中被造的。因为，他想要区分现在存在的属生灵的身体与将来在复活中要有的属神的身体，便说：\BibleQuote{它播种在可朽坏中，复活在不朽坏中；它播种在羞辱中，复活在光荣中；它播种在软弱中，复活在能力中；它播种一个属生灵的身体，复活一个属神的身体。}然后，为了证明这一点，他接着说：\BibleQuote{有一个属生灵的身体，也有一个属神的身体。}为表明属生灵的身体是什么，他说：\BibleQuote{正如经上所记：第一个人亚当成了有生命的灵魂，最后的亚当成了赋予生命的灵。}\BibleRef{格前 15:42}他这样是要表明属生灵的身体是什么，尽管圣经论及第一个人亚当，当他的灵魂由天主的气息造成时，并没有说\Quote{人是在属生灵的身体中造成的}，而是说\BibleQuote{人成了有生命的灵魂。}\BibleRef{创 2:7}因此，通过\BibleQuote{第一个人成了有生命的灵魂}这些话，宗徒愿意我们的属生灵的身体被理解。但他如何愿意属神的身体被理解，他在添加\BibleQuote{但最后的亚当成了赋予生命的灵}时表明了，这显然是指基督，祂已从死者中复活，再不能死。然后他接着说：\BibleQuote{但那属神的不是在先，而是属生灵的；然后才是属神的。}这里他更清楚地断言，当他说第一个人成了有生命的灵魂时，他指的是属生灵的身体；而当他说最后的亚当成了赋予生命的灵时，他指的是属神的身体。属生灵的身体是第一的，正如第一亚当所拥有的，如果他未曾犯罪，它就不会死；也正如我们现在所拥有的，其本性因罪而改变并败坏，以致我们不得不死；甚至基督也首先屈尊接受了这样一个身体，并非出于必然，而是出于选择；但以后到来的是属神的身体，它已经由基督作为我们的元首预先穿上了，并将在死者的复活中由祂的肢体穿戴上。\par

然后宗徒补充了这两个人之间的显著区别，说：\BibleQuote{第一个人出于地，属于土；第二个人是主，出自天。那属于土的怎样，凡属于土的也怎样；那属于天的怎样，凡属于天的也怎样。我们怎样带了那属于土的肖像，也要怎样带那属于天的肖像。}\BibleRef{格前 15:47}所以他在别处说：\BibleQuote{凡是你们受洗归于基督的，就是穿上了基督；}\BibleRef{迦 3:27}但这件事要真正实现，是在我们的复活中，那由生育而来的属于生灵的部分转变为属于神的时候。因为，再用他的话，我们是因望德而得救。\BibleRef{罗 8:24}现在我们因着罪与死的传递而带着属土的人的肖像，这传递藉着普通的生育临于我们；但我们因着宽恕与永生的恩宠而带着属天的人的肖像，这恩宠是由重生藉着天主与人之间的中保，那人基督耶稣赐予我们的。祂是保禄提到的属天之人，因为祂从天上来，穿上属土可朽坏的身体，为的是要给它穿上属天的不朽。祂称其他人为属天的，因为藉着恩宠他们成为祂的肢体，好使祂与他们一同成为一个基督，即头和身体。在同一封书信中，他更清楚地说明：\BibleQuote{因为死亡既因一人而来，死者的复活也因一人而来。正如在亚当内众人都死了，照样在基督内众人也都要复活，}\BibleRef{格前 15:21}——就是说，在属神的身体中，它将成为赋予生命的灵。这并不是说所有在亚当内死了的人都将是基督的肢体——因为绝大多数将在永死中受罚——而是他在两个子句中使用了\Quote{一切}这个词，因为正如没有人除非在亚当内而是在属生灵的身体中死去，同样也没有人除非在基督内而是在属神的身体中获享生命。因此，我们绝不能以为在复活中我们将拥有第一个人犯罪前那样的身体，也不能以为\BibleQuote{那属于土的怎样，凡属于土的也怎样}这句话是指由罪所产生的结果。因为我们不能认为亚当在堕落前曾拥有一个属神的身体，而为了惩罚他的罪，它被改变成了属生灵的身体。如果这样想，那就没有重视这样一位伟大导师的话，他说：\BibleQuote{有一个属生灵的身体，也有一个属神的身体；正如经上所记：第一个人亚当成了有生命的灵魂。}他是在犯罪之后才被造成这样的吗？或者这不是人的原始状态，而那位蒙福的宗徒正是从这种状态中选取见证，以表明属生灵的身体是什么？\par

% structure: p0057 source=h2 target=subsection reason=relative-heading-rank; base=h2

\subsection{第24章 我们必须如何理解那使\BibleQuote{第一个人成了有生命的灵魂}的天主的嘘气，以及主在说\BibleQuote{领受圣神}时将自己的圣神赐给门徒的那一次。}

有些人从\BibleQuote{天主将\OriginalTerm{生气}{breath of life}吹在亚当的鼻孔里，他就成了\OriginalTerm{活灵}{living soul}，}\BibleRef{创 2:7}这些话语仓促地推断，灵魂不是在那时才首次赐给人，而是那已赐给的灵魂由圣神所赋予生命。他们所以如此推测，是因主耶稣复活后向门徒们嘘了一口气，并说：\BibleQuote{你们领受圣神吧。}\BibleRef{若 20:22}由此他们推断，在两种情况下发生了同样的事，好像福音作者接着要说：他们就成为活灵了。但如果他添加了这一句，我们应明白：圣神在某种意义上是灵魂的生命，没有祂，有理性的灵魂必须算作死的，尽管其肉体在我们眼前似乎活着。然而，当人被造时发生的并非如此，叙事本身的词语充分显示：\BibleQuote{天主用地上的尘土造了人；}有些人认为用以下词语表达得更清楚：\Quote{天主用地上的泥土造了人。}因为在此之前已说过：\BibleQuote{有雾气从地上腾起，滋润整个地面，}\BibleRef{创 2:6}为使人理解由水气和尘土形成的泥土的所指。因为紧接这节经文之后就是宣告：\BibleQuote{天主用尘土造了人；}那些被译为拉丁文的希腊手稿就是这样记载的。但不论人喜欢读作\BibleQuote{\Emph{造}}或\BibleQuote{\Emph{形成}}，而希腊文读作\OriginalTerm{形成}{\GreekText{ἔπλασεν}}，这都不太重要；不过\BibleQuote{\Emph{形成}}是更好的译法。但那些偏好\BibleQuote{造}的人认为他们以此避免了由以下事实产生的歧义：在拉丁语中，习惯上，那些虚构和伪造事物的人被称为\Quote{形成}某物。因此，这个由地上的尘土、或湿润的尘土或泥土所造的人——即这\Quote{地上的尘土}（容我使用圣经明确的词语）——在获得灵魂时，如宗徒所教导的，成了一个有生命的身体。此人，他说，\BibleQuote{成了活灵；}也就是说，这被塑造的尘土成了活灵。\par

他们说道：他早已有灵魂，否则他就不能被称为人；因为人不仅是身体，也不仅是灵魂，而是由两者所构成的存有。这确实是真的：灵魂不是整个人，而是人的更好部分；身体不是整个，而是人的低劣部分；当两者结合时，他们获得人的名称，然而当我们单独谈论他们时，他们也并未分别失去这个名称。因为谁在日常用语中被禁止说\Quote{那个人死了，现在安息或受苦，}尽管这只能就灵魂而言；或\Quote{他被埋葬在某某地方，}尽管这仅指身体？他们会说圣经不遵循这样的用法吗？恰恰相反，圣经如此彻底地采用它，以至于即使当一个人活着，身体和灵魂结合时，圣经也单独称呼它们每一个为\Quote{\Emph{人}}，称灵魂为\BibleQuote{内在的人}，称身体为\BibleQuote{外在的人}，\BibleRef{格后 4:16}好像有两个人，虽然两者在一起确实只是一个。但我们必须理解人在何种意义上被称为天主的肖像，却又仍是尘土，并归于尘土。前者是指理性的灵魂，天主藉祂的嘘气，或更恰当地说，藉祂的吹气，将其传给人，即传给人的身体；后者则指身体，是天主用尘土所造的，并被赋予了一个灵魂，使它成为一个活的身体，也就是说，人成为一个活灵。\par

因此，当我们的主向祂的门徒嘘了一口气，并说：\Quote{你们领受圣神吧}，祂肯定希望我们明白，圣神不仅是圣父的灵，也是独生子自己的灵。因为同一位灵确实是圣父和圣子的灵，与他们形成圣父、圣子、圣神的天主圣三，不是受造物，而是造物主。因为那从祂肉身的口中发出的物质的呼吸，并非圣神本身的实体和本性，而只是表明（如我所说）圣神为圣父和圣子所共有；因为他们并非各自有单独的灵，而是两者拥有同一、相同的灵。如今，在圣经中，这灵总是用希腊词\OriginalTerm{灵}{\GreekText{πνεῦμα}}来称呼，正如主在所引的章节中称呼祂，当祂将祂赐予门徒时，用嘴唇的呼气暗示了这恩赐；而且我想不起在全部圣经中有任何地方用别的名称称呼祂。但在那记载着\BibleQuote{主天主用地上的尘土形成了人，并将生命的气息吹在他的脸上}的段落中，希腊文用的不是通常指圣神的\OriginalTerm{灵}{\GreekText{πνεῦμα}}，而是\OriginalTerm{气息}{\GreekText{πνοή}}，这个词更常用于受造物而非造物主；因此有些拉丁文译者宁可将它译为\Quote{气息}而非\Quote{灵}。因为这个词在希腊文依撒意亚第七章第十六节中也出现过，那里天主说：\BibleQuote{我造成了万灵}\BibleRef{依 7:16}，无疑是指所有的灵魂。因此，这个词\OriginalTerm{气息}{\GreekText{πνοή}}有时被译为\Quote{气息}，有时译为\Quote{灵}，有时译为\Quote{感动}，有时译为\Quote{灵感}，有时译为\Quote{灵魂}，即使当它用于天主时也是如此。另一方面，\OriginalTerm{灵}{\GreekText{Πνεῦμα}}则一致地被译为\Quote{灵}，无论是用于人（关于人，宗徒说：\BibleQuote{除了在人内的人灵，有谁知道人的事？}\BibleRef{格前 2:11}），还是用于兽（如在撒罗满的书中：\BibleQuote{谁知道人的灵往上走，兽的灵往下入地？}\BibleRef{训 3:21}），还是用于被称为风的物理之灵（因为圣咏作者这样称呼它：\BibleQuote{火与冰雹，雪与蒸气，暴风}），还是用于非受造的造物主之灵（主在福音中论及祂说：\BibleQuote{你们领受圣神吧}，用祂口中的呼气暗示了恩赐；当祂说：\BibleQuote{你们要去使万民成为门徒，因圣父、圣子、圣神之名给他们授洗}\BibleRef{玛 28:19}，这些话非常明确而卓越地赞颂了天主圣三；以及当它说：\BibleQuote{天主是灵}\BibleRef{若 4:24}，以及圣经中的许多其他处）。在所有上述圣经引文中，我们在希腊文里没有发现使用\OriginalTerm{气息}{\GreekText{πνοή}}，而是使用\OriginalTerm{灵}{\GreekText{πνεῦμα}}，在拉丁文里也不是\OriginalTerm{气息}{flatus}，而是\OriginalTerm{灵}{spiritus}。因此，再回到那处经文，它写着\BibleQuote{祂在他脸上吹了一口生命的气息}，即使希腊文没有用\OriginalTerm{气息}{\GreekText{πνοή}}（实际上用了）而是用了\OriginalTerm{灵}{\GreekText{πνεῦμα}}，也不能因此必然得出结论说那是指造物主之灵，即在天主圣三中特别被称为圣神的那一位，因为如前所述，\OriginalTerm{灵}{\GreekText{πνεῦμα}}不仅用于造物主，也用于受造物。\par

但他们会说：当圣经使用\Quote{灵}这个词时，它不会加上\Quote{生命}二字，除非是要我们理解为圣神；同样，当它说\Quote{人成了一个灵魂}时，也不会插入\Quote{活}字，除非是指那从上面由天主的恩赐赋予灵魂的生命。因为，既然灵魂本身有其自身固有的生命，他们问，何需加上\Quote{活}呢？除非是为了表示那由圣神赐予的生命？这岂不是在拼命为自己的猜测辩护，而漫不经心地忽略了圣经的教导？如果他们稍加留意，就在这本创世记的上一页就能找到这句话：\BibleQuote{让大地生出活的灵魂}\BibleRef{创 1:24}，那时所有地上的动物都被造了。然后，在同一本书中，稍隔几行，他们难道不可能注意到这节经文：\BibleQuote{凡在旱地上有生命气息的，都死了}？这表示所有在地上的动物都死于洪水。那么，如果我们发现圣经即使在论及野兽时也惯常使用\Quote{活的灵魂}和\Quote{生命的气息}；并且在这里，当它说\Quote{凡有生命气息的}时，希腊文用的是\OriginalTerm{气息}{\GreekText{πνοή}}而不是\OriginalTerm{灵}{\GreekText{πνεῦμα}}，那么我们为什么不能说：既然灵魂没有生命就不可能存在，何需加上\Quote{活}字？或者，在\Quote{灵}之后何需加上\Quote{生命}一词？但我们明白，当圣经谈论动物（即有灵魂的身体，其中灵魂是感觉的居所）时，它使用这些表达是其平常风格；但当论及人时，我们却忘记了圣经的通常和既定的用法，即它表明人获得了一个理性的灵魂，这个灵魂不像其他生物那样从水和土中产生，而是由天主的嘘气创造的。然而，这一创造被安排为人的灵魂应在动物性的身体中生活，就像那些其他动物一样，关于那些动物圣经说\BibleQuote{让大地生出每一个活的灵魂}，并且再次说它们里面有生命的气息，那里希腊文用的是\OriginalTerm{气息}{\GreekText{πνοή}}而不是\OriginalTerm{灵}{\GreekText{πνεῦμα}}，而且那里当然不是指圣神，而是指它们的灵。\par

但是，他们又反驳说，\OriginalTerm{气息}{breath}被理解为从\Person{天主}口中呼出；如果我们相信那就是\OriginalTerm{灵魂}{soul}，那么我们必须承认它与那智慧\OriginalTerm{同体}{consubstantial}且相等，那智慧说：\BibleQuote{我出自至高者的口。}\BibleRef{德 24:3}的确，智慧并没有说是从\Person{天主}口中所呼出的，而是从其中出来的。但正如我们呼吸时，能产生气息，不是出于我们人性的本质，而是出于我们吸入和呼出的周围空气；同样，全能的\Person{天主}也能产生气息，不是出于祂的本性，也不是出于祂以下的受造物，而是出于虚无；当祂将这气息赋予人的身体时，最恰当地称之为呼出或吹入——\OriginalTerm{非物质者}{Immaterial}吹入非物质者，但\OriginalTerm{不变者}{Immutable}并非也是不变的；因为它是受造的，祂是\OriginalTerm{非受造者}{uncreated}。然而，那些急于引用圣经却不了解其语言用法的人，要知道不仅与\Person{天主}同等同体的事物被称为从祂口中出来，请他们听或读\Person{天主}所说的话：\BibleQuote{所以，因为你既不冷也不热，而是温的，我必将你从我口中吐出。}\BibleRef{默 3:16}\par

那么，当我们\Person{宗徒}如此明确地区分\OriginalTerm{属魂的身体}{animal body}和\OriginalTerm{属神的身体}{spiritual body}——也就是说，我们现在所拥有的身体和我们将要拥有的身体时，我们就没有反对的理由了。他说：\BibleQuote{播种的是\OriginalTerm{自然身体}{natural body}，复活的是\OriginalTerm{属神身体}{spiritual body}；既有自然身体，也有属神身体。经上也这样记载：\InnerQuote{第一个人\Person{亚当}成了\OriginalTerm{有灵的灵魂}{living soul}，最后的\Person{亚当}成了\OriginalTerm{赋予生命的灵}{quickening spirit}。}但是，属神的不是在先，而是属自然的在先，然后是属神的。第一个人出于地，属于土；第二个人是\OriginalTerm{主}{Lord}，出于天。那属于土的怎样，凡属于土的也怎样；那属于天的怎样，凡属于天的也怎样。我们既然带了那属于土的肖像，也要带那属于天的肖像。}\BibleRef{格前 15:44}关于他所说的一切话，我们之前已经谈过了。因此，\Person{宗徒}所说的第一个人\Person{亚当}被造成的那种属魂的身体，并不是被造得完全不能死，而是除非他犯罪，否则就不会死。那将要被赋予生命的灵做成属神而\OriginalTerm{不朽}{immortal}的身体，是绝不能死的；正如灵魂被造为不朽的，所以虽然因罪可以说它死了，并且失去了它自己某种生命，即\Person{天主}的\Person{圣神}——灵魂靠祂才能明智而幸福地生活——但它并没有停止活某种生命，尽管是悲惨的，因为它因受造而是不朽的。同样，叛逆的天使，虽然因犯罪在某种意义上死了，因为他们离弃了生命的泉源——\Person{天主}，当他们饮用时，他们能够明智而良善地生活，但他们不能死到完全停止生活和感觉，因为他们因受造而是不朽的。因此，在最后审判之后，他们将被投入\OriginalTerm{第二次死亡}{second death}，在那里也不被剥夺生命或感觉，而是将受苦刑。但是那些被\Person{天主}的恩宠所拥抱，并成为恒常福乐中的圣天使的同胞的人，将永不再犯罪或死亡，因为他们拥有属神的身体；然而，穿上天使所享有的不朽，即使因犯罪也不能被剥夺，他们的肉身的性质将保持不变，但一切肉欲的败坏和笨重将被除去。\par

还有一个问题需要讨论，并借着\Person{真理之主天主}的帮助来解决：如果在我们原祖不顺服的肢体中\OriginalTerm{私欲偏情}{concupiscence}的运动源于他们的罪，并且仅当\OriginalTerm{天主的恩宠}{divine grace}离弃他们时才发生；如果正是在那时他们的眼睛开了，看到，更确切地说，注意到他们的赤身露体，并且他们遮盖了自己的羞耻，因为他们肢体无耻的运动不受他们意志的控制——那么，如果他们像受造时那样保持无罪，他们将如何生育子女呢？但由于本书必须结束，而如此大的一个问题不能草率处理，我们可以将其留到下一卷书中，在那里将更方便地讨论。\par

\SourceRecord{Translated by Marcus Dods.；Nicene and Post-Nicene Fathers, First Series；Vol. 2.；Edited by Philip Schaff.；Buffalo, NY: Christian Literature Publishing Co.,；1887.；<http://www.newadvent.org/fathers/120113.htm>.}

% structure: p0001 source=h1 target=section reason=page-title

\section{天主之城（第十四卷）}

奥斯定再度论及第一人的罪，并教导它是肉性生命及人邪恶情感的原因。他特别证明伴随情欲的羞耻正是那不顺从的公正惩罚，并探讨人若未犯罪，将如何能无淫欲地繁衍后代。\par

% structure: p0003 source=h2 target=subsection reason=relative-heading-rank; base=h2

\subsection{第一章．——第一人的不顺从原本会将全人类投入第二次死亡的无穷痛苦中，若非天主的恩宠拯救了许多人。}

我们在前几卷已经说明，天主不仅愿意人类因其本性的相似而能彼此交往，也愿他们因亲缘关系而和谐共处、缔结和平，因此乐意使全人类出自一人，并创造了具有如此本性的人，使本族的成员本不会死亡，若非第一对（其中一人由无中创造，另一人由他而出）因不顺从而招致此后果；因为他们犯下了如此大罪，以致人性因此而变得更糟，并传给了他们的后裔，使其易于犯罪并必死。死亡的王国如此统治人类，以致罪的应得惩罚本会将所有人毫无例外地投入那无尽头的第二次死亡，若非天主不应得的恩宠拯救了一些人。于是，虽然地上有极多且强大的邦国，其礼仪、习俗、言语、武器和服饰都有明显差异，但人类社会无非两种，我们可按圣经的用语恰当地称之为两座城。一座由那些愿意随从肉性生活的人组成，另一座由那些愿意随从灵性生活的人组成；当他们各自达成所愿时，便按各自的方式生活在平安中。\par

% structure: p0005 source=h2 target=subsection reason=relative-heading-rank; base=h2

\subsection{第二章．——论肉性生命，不仅指生活在肉体放纵中，也指生活在内里之人的恶习中。}

首先，我们必须看清何为随从肉性生活，何为随从灵性生活。任何人若不记得或未充分思考圣经的用语，初听我们所说的话，或许会以为伊壁鸠鲁派哲学家随从肉性生活，因为他们将人的至高幸福置于肉体快乐中；而另一些认为以某种形式肉体之善是人的至高幸福的人也是如此；以及大众也是如此，他们虽未对此进行教条化或哲学化思考，却如此倾向于情欲，以致除了从身体感官获得的快乐外，无法以任何其他快乐为乐；他或许会以为斯多亚派随从灵性生活，因为他们将人的至高幸福置于灵魂中；因为人的灵魂若不是灵，又是什么？但按圣经的意义，两者都被证明是随从肉性的。因为按圣经，\InnerQuote{肉}不仅指属地且必死动物的身体，如经上说：\BibleQuote{并非所有的肉都是同样的肉，但有一种是人的肉，另一种是走兽的肉，另一种是鱼的肉，另一种是鸟的肉。} \BibleRef{格前 15:39} 但它也在许多其他意义上使用这个词；在这些不同用法中，常见的一种是用\InnerQuote{肉}指人本身，以人的本性部分代整体，如经上说：\BibleQuote{凡有血肉的，不能因行法律而称为义人。} \BibleRef{罗 3:20} 他在这里说\InnerQuote{没有血肉的}不就是\InnerQuote{没有人}吗？事实上，他后来更清楚地说：\BibleQuote{没有人能因法律而成义。} \BibleRef{迦 3:11} 并且在致迦拉达人书中说：\BibleQuote{知道人不是因法律的作为而成义。} \BibleRef{迦 3:11} 因此我们理解这句话：\BibleQuote{圣言成了肉。} \BibleRef{若 1:14} 即成了人，有些人未能正确理解其意义，便以为基督没有人的灵魂。因为在福音中，玛利亚玛达肋纳的话以整体代部分：\BibleQuote{他们把我的主搬走了，我不知道他们把他放在哪里。} \BibleRef{若 20:13} 她此话仅指基督的肉，她以为这肉从其埋葬的坟墓中被搬走了。同样，以部分代整体，提到肉而指人，如上文引用的例子那样。\par

既然圣经以多种方式使用\Quote{肉躯}一词，我们无暇一一收集和研究；若想查明\Quote{活于肉躯}究竟是什么（这当然是邪恶的，尽管肉躯的本性并非邪恶），我们就必须仔细查考保禄宗徒写给迦拉达人的那封书信中的一段话，他在其中说：\BibleQuote{论肉躯的作为，是显而易见的，即：奸淫、邪淫、污秽、放荡、崇拜偶像、行邪术、仇恨、竞争、嫉妒、忿怒、争吵、分裂、异端、妒忌、凶杀、醉酒、宴乐，以及类似的事。关于这些，我预先告诉你们，如同我以前说过的那样：做这样事的人，不能继承天主的国。}\BibleRef{迦 5:19} 仔细考虑宗徒书信的这段全文，只要涉及当前问题，就足以回答\Quote{活于肉躯}是什么。因为在他列举并谴责的那些显而易见的肉躯作为中，我们不仅发现那些关乎肉躯享乐的，如奸淫、污秽、放荡、醉酒、宴乐，而且也发现那些虽远离肉躯享乐，却揭示灵魂恶习的。因为，谁看不出崇拜偶像、行邪术、仇恨、竞争、嫉妒、忿怒、争吵、异端、妒忌，与其说是肉躯的恶习，不如说是灵魂的恶习呢？因为人完全有可能为了崇拜偶像或某种异端谬误而戒绝肉躯享乐；然而，即使他这样做，根据这段宗徒的权威证明，他仍然是活于肉躯；他在戒绝肉躯享乐的同时，却在实践应当受罚的肉躯作为。有仇恨的人，仇恨不就在他的灵魂中吗？谁会向他的仇敌，或他认为的仇敌说：\Quote{你对我有坏的肉躯}，而不是说：\Quote{你对我有坏的心思}呢？最后，如果有人听到我所说的\OriginalTerm{肉性之事}{carnalities}，他必会将其归于人的肉性部分；因此，没有人会怀疑\OriginalTerm{仇视}{animosities}属于人的灵魂。那么，外邦人在信德和真理中的导师为什么称所有这些以及类似的事情为肉躯的作为呢？除非是因为，借着以部分指代整体的表达方式，他希望通过\Quote{肉躯}一词来理解人本身？\par

% structure: p0008 source=h2 target=subsection reason=relative-heading-rank; base=h2

\subsection{第3章——罪不是由肉躯造成，而是由灵魂造成，且从罪所染的腐朽不是罪，而是罪的惩罚。}

但是，若有人说肉躯是一切恶行和不良行为的根源，因为灵魂之所以邪恶地生活，只是因为它被肉躯所推动，那么他肯定没有仔细考虑人的全部本性。因为\BibleQuote{那可朽坏的身体，确实使灵魂沉重。}\BibleRef{智 9:15} 因此，宗徒论及这可朽坏的身体——他不久前提到的\BibleQuote{我们外在的人虽然日渐朽坏}\BibleRef{格后 4:16}——说：\BibleQuote{我们知道，如果我们地上帐棚式的寓所拆毁了，我们就由天主获得一所房舍，一所非人手所造，永远在天上的寓所。我们在这帐棚里叹息，切望穿上我们那来自天上的住所：但愿穿上之后，不至於赤身裸体。我们在这帐棚里的人，负担沉重，叹息不已：这并不是我们愿意脱去这帐棚，而是愿意穿上那另一件，好使有死的生命被生命所吞没。}\BibleRef{格后 5:1} 因此，我们被这可朽坏的身体所重压；但知道这重压的原因不是身体的本性和实体，而是它的腐朽，我们并不希望脱去身体，而是希望穿上它的不朽。因为那时，也将有一个身体，但它不再是重负，因为它不再是可朽坏的。目前，确实\BibleQuote{可朽坏的身体重压着灵魂，这地上的帐棚使充满挂虑的心智沉重}\BibleRef{智 9:15}，然而，那些以为灵魂的一切邪恶都来自身体的人，是错了。\par

维吉尔确实似乎用优美的诗句表达了柏拉图的思想，他说：\par

一股炽热的力量激励着他们的生命，\newline{}一种源自天堂的本质，\newline{}虽部分被泥土的肢体所阻碍\newline{}和那呆滞的\Quote{腐朽的衣袍}；\newline{}但是，尽管他接着提到四种最常见的心理情感——欲望、恐惧、喜悦、悲伤——意图表明身体是一切罪和恶的根源，说：\par

由此产生狂野的欲望和卑下的恐惧，\newline{}以及人的欢笑，人的眼泪，\newline{}被禁锢在似地牢的黑夜中，\newline{}他们向外张望，却看不见光明，\par

但我们的信念完全相反。因为身体的腐朽，它重压灵魂，并非原罪的成因，而是原罪的惩罚；并非可腐朽的肉身使灵魂有罪，而是有罪的灵魂使肉身成为可腐朽的。虽然从肉身的腐朽产生了某些趋向恶的诱惑，甚至邪恶的欲望，但我们不可将邪恶生活中的一切恶行都归于肉身，以免借此洗脱魔鬼的一切罪行，因为他没有肉身。因为虽然我们不能称魔鬼为奸淫者或酒徒，或归给他任何肉欲的放纵（虽然他是在这些方面犯罪之人的隐秘煽动者和唆使者），但他却极其骄傲和嫉妒。这种恶毒如此占据了他，以致他因此被拘禁在黑暗的锁链中，直至永罚。现在，宗徒将这些支配魔鬼的恶行归于肉身，而魔鬼显然没有肉身。因为他说\Quote{仇恨、纷争、忌恨、恼怒、嫉妒}是肉身的作为；而在所有这些邪恶中，骄傲是根源和首脑，它在魔鬼身上掌权，尽管他没有肉身。因为谁对圣人们显示更多的仇恨？谁与他们更加纷争？谁更嫉妒、苦毒、忌恨？既然他表现出这一切作为，虽然他并无肉身，它们又怎能是肉身的作为呢？除非因为它们是人——正如我所说，人被称为肉身——的作为？因为人成为像魔鬼，并非由于拥有肉身（魔鬼没有肉身），而是由于按自己而生活——即按人而生活。因为魔鬼也企图像自己而生活，当他不持守真理时；以致他说谎，不是出于天主，而是出于他自己，他不仅是说谎者，还是谎言之父，他是第一个说谎的，也是说谎与罪的创始者。\par

% structure: p0014 source=h2 target=subsection reason=relative-heading-rank; base=h2

\subsection{第四章——按人而生活与按天主而生活是什么意思}

因此，当人按人而生活，而非按天主而生活时，他就如同魔鬼。因为即使天使也不可按照天使而生活，而只能按照天主而生活，如果他愿意持守真理，并讲论天主的真理而不是自己的谎言。关于人，同一位宗徒在另一处也说：\BibleQuote{如果天主的真理因我的谎言更加丰盛}\BibleRef{罗 3:7}——他说\BibleQuote{我的谎言}和\BibleQuote{天主的真理}。当一个人按真理而生活时，他不是按自己而生活，而是按天主而生活；因为说\BibleQuote{我是真理}的那一位就是天主\BibleRef{若 14:6}。因此，当人按自己而生活——即按人，而非按天主——时，他确实是按谎言而生活；这并不是说人本身是个谎言，因为天主是他的创造者和作者，而天主当然不是谎言的作者和创造者，而是因为人被造为正直，原意是叫他不要按自己而生活，而要按创造他的那一位而生活——换句话说，叫他承行祂的旨意而非他自己的；而不按他被造时所应当的方式去生活，这就是谎言。因为他确实渴望蒙福，甚至不愿那样生活以便能够蒙福。如果这不算是谎言，那什么才是？因此，说\Quote{一切罪都是谎言}并非毫无意义。因为犯任何罪都是由于那种欲望或意志，我们渴望一切顺利，并惧怕不顺。因此，我们为了自己顺利而做的事，若反而使我们比先前更悲惨，那就是谎言。这是为什么呢？不是因为人的幸福之源只在于天主，他犯罪时离弃了天主，而不在于他自己——按自己而生活就是犯罪？\par

我们在陈述这一论点时，说有些人按照肉体生活，有些人按照精神生活，因而产生了两个互相对立、彼此冲突的城；我们本来也可以这样说：\Quote{有些人按照人生活，有些人按照天主生活。}因为保禄非常明确地对格林多人说：\Quote{你们中间有嫉妒和纷争，岂不是属血肉、按照人行事吗？}\BibleRef{格前 3:3}所以，按照人行事和属血肉是同一回事；因为人是借用了\Quote{肉体}这个词，即人的一部分。因为在此之前，他把那些后来称为属血肉的人称为属血气的人，他说：\Quote{除了人内里的心神，谁能知道人的事呢？同样，除了天主的神，也没有人知道天主的事。我们没有领受这世界的神，而是领受了出于天主的神，为使我们能知道天主白白赐予我们的事。我们也讲论这些事，不是用人的智慧所教导的言辞，而是用圣神所教导的言辞，将属神的事与属神的事相比较。然而，属血气的人不接受天主圣神的事，因为在他看这些事是愚妄的。}\BibleRef{格前 2:11}所以，对这种人，即属血气的人，他不久后说：\Quote{弟兄们，我从前不能把你们当作属神的人，只能当作属血肉的人。}\BibleRef{格前 3:1}这要用同样的用法来解释，即以部分代表整体。因为灵魂和肉体，作为人的组成部分，都可以用来指代整个人；因此属血气的人和属血肉的人不是两样不同的东西，而是同一样东西，即按照人生活的人。同样，\Quote{凡有\Emph{血肉}的，不得因行法律成义}\BibleRef{罗 3:20}这句话中的\Quote{血肉}无非是指\Quote{人}；而\Quote{与雅各伯一同下到埃及的，共有\Emph{七十五个灵魂}}\BibleRef{创 46:27}这句话中，\Quote{七十五个灵魂}是指七十五个人。并且，\Quote{不是用人的智慧所教导的言辞}这句话同样可以说是\Quote{不是用属血肉的智慧所教导的言辞}；而\Quote{你们按照人行事}可以说成\Quote{按照肉体行事}。这在随后的话中更明显：\Quote{因为有人说：我是属保禄的；有人说：我是属阿波罗的；你们岂不是人吗？}他先前所说的\Quote{你们是属血气的}、\Quote{你们是属血肉的}，现在用\Quote{你们是人}来表达；即你们按照人生活，不按照天主生活，因为如果你们按照天主生活，你们就成了神。\par

% structure: p0017 source=h2 target=subsection reason=relative-heading-rank; base=h2

\subsection{第五章——论柏拉图派关于灵魂与肉体本性的见解虽不如摩尼教那样应受责备，但仍有可驳之处，因为它将恶习的根源归于肉体的本性。}

因此，在我们的罪与恶习中，我们无需指责肉体的本性而冒犯造物主，因为肉体在其自身种类和层次上本是好的；但离弃那造物主——至善——而按照受造之善生活，这本身不是善，无论一个人选择按照肉体、按照灵魂、还是按照由灵魂和肉体组成的整个人性来生活。因为凡称赞灵魂的本性为至善、指责肉体的本性为邪恶的人，无疑是在灵魂的爱和肉体的恨中属于肉体；因为这些情感源于人的幻想，而非出自神圣的真理。柏拉图派并不像摩尼教那样愚昧，将我们现世的肉体憎恶为邪恶的本性；因为他们将构成这个可见和可触世界的一切元素及其属性都归于造物主天主。然而，他们相信灵魂受到身体那感染死亡的肢体和属土的构造的影响，以致由此产生欲望、恐惧、喜悦和悲伤这些疾病，按照西塞罗的说法是四种扰动，或按照大多数人的说法（用希腊人更常用的名称）是四种情欲，人的全部邪恶生活都包含在其中。但如果真是这样，为何维吉尔笔下的埃涅阿斯在冥府听父亲说灵魂要返回身体时，对这一宣告表示惊讶并喊道：\par

父亲呀！岂能想象，\newline{}快乐的灵魂会离开这境地，\newline{}寻求上面的天空，\newline{}与迟钝的泥土重新结合？\newline{}这令人恐惧的光明渴望，\newline{}从何而来，请说，为何？\par

那么，这种令人恐惧的渴望，即使在那自诩纯洁的离体灵魂中仍然存在，并且仍然来自感染死亡的肢体和属土的四肢吗？难道他不是断言，当他们开始渴望返回身体时，已经摆脱了所有所谓的身体瘟疫吗？由此我们推断，即使这种往复的净化与玷污——灵魂离开和返回——像其绝对虚假的那样真实，也不能断言所有有罪的邪恶灵魂运动都源于属土的身体；因为，按照他们那位高贵阐释者的话，这种\Quote{令人恐惧的渴望}如此远离身体，以至于它感动那已涤除一切身体玷污、且独立于任何身体而存在的灵魂，并且感动它再次进入身体。因此，连他们自己也承认，灵魂不仅因肉体而激动产生欲望、恐惧、喜悦、悲伤，而且它也能凭自身被这些情绪所扰动。\par

% structure: p0021 source=h2 target=subsection reason=relative-heading-rank; base=h2

\subsection{第六章——论人的意志的特性，它使灵魂的情感成为正当或邪恶。}

人的意志特質是重要的；因為，如果意志是錯的，這些靈魂的動情就會是錯的；但如果意志是對的，它們就不僅無可指責，甚至值得讚揚。因為意志在它們所有之中；是的，它們之中沒有一個不是意志。因為欲望和喜樂若不是我們對所願之物的同意意願，又是什麼呢？恐懼和悲傷若不是我們對所不願之物的厭離意願，又是什麼呢？但當同意採取追求擁有我們所願之物的形式時，稱之為欲望；當同意採取享受我們所願之物的形式時，稱之為喜樂。同樣地，當我們厭離那不願發生的事物時，此意願稱為恐懼；當我們背離那違背我們意願而發生的事物時，此意志行動稱為悲傷。一般而言，對於我們所追求或迴避的一切，正如人的意志被吸引或被排斥，它便改變並轉化為這些不同的情感。因此，按天主而活、不按人而活的人，應當是喜愛善者，因而也是憎恨惡者。既然沒有人本性是惡的，凡是惡的都是由於惡習而惡，那麼按天主而活的人應當對惡人懷有完全的憎恨，以便既不因人的惡習而恨其人，也不因其人而愛其惡習，而要恨惡習、愛其人。因為惡習被咒詛後，所有當愛的事物將存留，無一當恨。\par

% structure: p0023 source=h2 target=subsection reason=relative-heading-rank; base=h2

\subsection{第七章——論\OriginalTerm{愛}{Amor}與\OriginalTerm{摯愛}{Dilectio}在聖經中無區別地用於善與惡的情感。}

凡決意愛天主，並愛近人如己，不按人而按天主而愛的人，因這愛被稱為有善意；這在聖經中更常稱為愛德，但即使在同一些書卷中，也被稱為愛。因為宗徒說，被選為人民首領的人必須是喜愛善者。當主親自問伯多祿：\Quote{你\OriginalTerm{愛}{diligis}我超過這些人嗎？}伯多祿回答說：\Quote{主，祢知道我\OriginalTerm{愛}{amo}祢。}第二次主又問，不是問伯多祿是否\OriginalTerm{愛}{amaret}祂，而是問他是否\OriginalTerm{愛}{diligeret}祂，他再次回答：\Quote{主，祢知道我\OriginalTerm{愛}{amo}祢。}但第三次詢問時，主自己不再說：\Quote{你\OriginalTerm{愛}{diligis}我嗎？}而是說：\Quote{你\OriginalTerm{愛}{amas}我嗎？}然後聖史補充說，伯多祿因主第三次對他說\Quote{你\OriginalTerm{愛}{amas}我嗎？}而憂愁，雖然主不是說了三次，而是只說了一次\Quote{你\OriginalTerm{愛}{amas}我嗎？}和兩次\Quote{\OriginalTerm{愛}{Diligis}我嗎？}由此我們得知，即使主說\Quote{\OriginalTerm{愛}{diligis}}時，祂用的詞與\Quote{\OriginalTerm{愛}{amas}}等同。伯多祿也始終用同一個詞表達同一件事，第三次也回答說：\Quote{主，祢知道一切，祢知道我\OriginalTerm{愛}{amo}祢。}\par

我認為提說此事是合宜的，因為有些人認為\Quote{愛德}或\Quote{摯愛}（\OriginalTerm{摯愛}{dilectio}）是一回事，\Quote{愛}（\OriginalTerm{愛}{amor}）是另一回事。他們說，\OriginalTerm{摯愛}{dilectio}用於善良的情感，\OriginalTerm{愛}{amor}用於邪惡的愛。但非常確定的是，即使是世俗文學也不知道這種區別。然而，這留待哲學家去決定它們是否不同以及如何不同，儘管他們自己的著作充分證明他們高度重視置於善物之上甚至天主本身之上的\OriginalTerm{愛}{amor}。但我們想要表明，我們這宗教的聖經——其權威我們優於一切著作——在\OriginalTerm{愛}{amor}、\OriginalTerm{摯愛}{dilectio}與\OriginalTerm{愛德}{caritas}之間不作區分；我們已經表明\OriginalTerm{愛}{amor}用於良好的上下文。若有人幻想\OriginalTerm{愛}{amor}無疑既用於善愛也用於惡愛，但\OriginalTerm{摯愛}{dilectio}僅保留給善愛，就請他回想聖詠上所說：\Quote{\OriginalTerm{喜愛}{diligit}不義的，憎恨自己的靈魂}；以及宗徒若望的話：\Quote{\BibleQuote{若有人\OriginalTerm{愛}{diligere}世界，對聖父的\OriginalTerm{摯愛}{dilectio}就不在他內。} \BibleRef{若一 2:15}}這裡同一段落中，\OriginalTerm{摯愛}{dilectio}既用於善義也用於惡義。若有人要求\OriginalTerm{愛}{amor}用於惡義的例子（因為我們已表明它用於善義），就請他讀這些話：\Quote{\BibleQuote{因為那時人將成為\OriginalTerm{愛己者}{amantes}、\OriginalTerm{愛財者}{amatores}。} \BibleRef{弟後 3:2}}\par

因此，正确的意志是导向良善的爱，错误的意志是导向邪恶的爱。爱，渴望拥有所爱之物，便是欲求；拥有并享受所爱之物，便是喜乐；逃避与所爱相悖之物，便是恐惧；感受所爱相悖之物，当它降临时，便是忧伤。这些情感，若爱是邪恶的，则为邪恶；若爱是良善的，则为良善。让我们从圣经来证明我们所说的。宗徒\Quote{渴望离世，与基督同在}。\BibleRef{斐 1:23} 又说：\Quote{我的灵魂渴望切慕祢的典章}；或更恰当地说：\Quote{我的灵魂切慕渴望祢的典章}。又说：\Quote{对智慧的渴望引人进入王国}。\BibleRef{智 6:20} 然而，历来有一种用法，若对象不加以限定，则欲求和贪欲被理解为恶的意义。但喜乐用于善义：\Quote{义人哪，你们应当因主欢喜快乐。} 又说：\Quote{祢使我心中喜乐。} 又说：\Quote{祢以祢的容颜充满我喜乐。} 恐惧也被宗徒用于善义，他说：\Quote{当用恐惧战兢，作成你们救恩的工夫}。\BibleRef{斐 2:12} 又说：\Quote{不要心高气傲，倒要惧怕}。\BibleRef{罗 11:20} 又说：\Quote{我怕你们的心或偏于邪，失去那向基督所存纯一清洁的心，就像蛇用诡诈诱惑了夏娃一样}。\BibleRef{格后 11:3} 至于忧伤——西塞罗更愿意称之为\OriginalTerm{疾病}{œgritudo}，维吉尔称之为\OriginalTerm{痛苦}{dolor}（如他说\Quote{\Emph{Dolent gaudentque}}），而我更愿意称之为悲伤，因为疾病和痛苦更常用于指身体的痛苦——关于这种情感，我说，它能否用于善义，则是一个更困难的问题。\par

% structure: p0027 source=h2 target=subsection reason=relative-heading-rank; base=h2

\subsection{第八章 论斯多亚派承认智者灵魂中有三种扰动而排除忧伤或悲哀，刚毅之心不应经历之。}

那些希腊人称为\OriginalTerm{善情}{\GreekText{εὐπαθείαι}}、西塞罗称为\Emph{constantiae}的情感，斯多亚派将其限定为三种；他们以意志代替欲求，以满足代替喜乐，以谨慎代替恐惧，以此来取代智者灵魂中的三种\Quote{扰动}；至于疾病或痛苦（我们为避免含混，更愿称之为悲伤），他们否认它能存在于智者的心中。他们说，意志追求善，智者正是如此；满足的对象是已拥有的善，智者不断拥有它；谨慎躲避恶，智者应当躲避它。而悲伤产生于已经发生的恶；既然他们假设智者不可能遭遇任何恶，那么他的心中就不可能有悲伤的对应物。因此，按照他们，只有智者才能有意志、满足和谨慎；而愚人只能欲求、喜乐、恐惧和悲伤。前三种情绪西塞罗称为\Emph{constantiae}，后四种称为\Emph{perturbationes}。然而许多人称后者为\OriginalTerm{情感}{passions}；如我所说，希腊人称前者为\OriginalTerm{善情}{\GreekText{εὐπαθείαι}}，后者为\OriginalTerm{情感}{\GreekText{πάθη}}。当我仔细查考圣经，看这术语是否得到认可时，我见到了先知这句话：\Quote{恶人没有平安，}上主说\BibleRef{依 57:21}；好像恶人对于恶事更可欢喜而非满足，因为满足是良善和虔诚之人的特性。我也在福音中读到这句话：\Quote{凡你们愿意人给你们做的，你们也要照样给人做。}\BibleRef{玛 7:12}这似乎暗示邪恶或可耻之事可能是欲求的对象，却不是意志的对象。确实，有些译者添加了\Quote{善事}，以使表达更符合惯用说法，将其意思定为：\Quote{凡你们愿意人给你们做的善事，你们也要照样给人做。}因为他们认为这能防止任何人希望他人为他提供不雅甚至可耻的满足——例如奢侈的宴席——假设他以同样方式回报他人，就是在履行这条诫命。然而，拉丁文所译的希腊福音中，没有出现\Quote{善事}，只有：\Quote{凡你们愿意人给你们做的，你们也要照样给人做，}我相信，这是因为\Quote{愿意}这个词已包含了\Quote{善}；因为祂没有说\Quote{欲求}。\par

然而，虽然我们有时可以使用语言这些精确的特性，却不可始终为它们所束缚；当我们阅读那些权威不容置疑的作者时，我们必须接受上述意义，即在那些只有通过这种诠释才能得出正确含义的段落中，正如我们引用的部分来自先知、部分来自福音的例子。因为谁不知道恶人因欢乐而得意呢？然而主说：\Quote{恶人毫无\Emph{\OriginalTerm{满足}{contentment}}。}何以如此呢？岂非因为满足一词，若按其本义与特定意义使用，与欢乐有所不同吗？同样，谁会否认以下教训是错误的：要人将那些他们愿意他人为他们所做的事，也照样为他人去做，以免他们以可耻与不法的快乐彼此取悦？然而这诫命：\Quote{凡你们\Emph{愿意}人给你们做的，你们也要照样给他们做}，却是十分有益和公义的。这又是为何？岂非因为在此处，\OriginalTerm{意志}{will}一词是严格使用的，指向那种不可能以恶为对象的意志？然而日常用语不会允许\BibleQuote{不可存心捏造任何谎言}\BibleRef{德 7:13}，除非也有一种邪恶的意志，其邪恶与天使所赞颂的\BibleQuote{天主受享光荣于高天，主爱的人在世享平安}\BibleRef{路 2:14}相分离。因为如果除了善意志外别无其他意志，\Quote{善}字便是多余。宗徒为何在颂扬爱德时把它作为大事提及，就是\Quote{不以不义为乐}，除非因为邪恶也以此为乐？因为即使是世俗作家，这些词语也被不加区别地使用。西塞罗，这位最富辩才的演说家，说：\Quote{元老们，我渴望仁慈。}谁会如此迂腐，认为他应该说\Quote{我愿意}而非\Quote{我渴望}，只因为该词用于正面的语境？又如泰伦提乌斯笔下，沉溺于狂放情欲的青年说：\Quote{我所愿的唯有菲卢墨娜。}他的老仆人的回答充分表明此\Quote{愿}为情欲：\Quote{与其说这些话徒然燃起你的情欲，不如试着将那份爱从心中驱除，岂不更好！}至于满足一词被世俗作家用于贬义，维吉尔的诗句可以证明，他极其简洁地概括了这四种扰动——\par

\Quote{因此他们恐惧与渴望，悲伤与满足。}\par

同一作者也曾使用\Quote{心灵之邪恶的满足}这一表达。因此善人与恶人同样地意愿、谨慎和满足；换言之，善人与恶人同样地渴望、恐惧、喜悦，只是前者以善的方式，后者以恶的方式，取决于\OriginalTerm{意志}{will}是正确还是错误。悲伤本身，连斯多葛学派也不允许出现在智者心灵中，却被用于正面的意义，尤其是在我们的著作中。因为宗徒称赞格林多人，因为他们有了按照天主的悲伤。但也许有人会说，宗徒祝贺他们是因为他们忏悔的悲伤，而这种悲伤只存在于犯罪的人中。因为他的话是这样的：\BibleQuote{因为我察觉那封信曾使你们悲伤，不过只是暂时的。如今我欢喜，不是因你们悲伤，而是因你们悲伤以致悔改，因为你们是按照天主的意思悲伤，免得你们在我们手中受到什么损害。因为按照天主的意思悲伤，能产生引人悔改以致得救的悔改，而不必后悔；但世俗的悲伤却带来死亡。看，你们按照天主的意思悲伤，这本身产生了何等热忱！}\BibleRef{格后 7:8}因此，斯多葛学派可以辩护说，悲伤确实有益于悔罪，但这在智者的心灵中不可能存在，因为智者没有罪可悲伤悔改，也没有其他恶的可忍受或经历而使其悲伤。因为他们说，\Person{阿尔西比亚德}（如果我没有记错的话），他自认幸福，但当\Person{苏格拉底}与他辩论，证明他因愚蠢而不幸时，他流下了眼泪。在他的情况下，愚蠢便是这种有用而可欲的悲伤之因，人因哀悼自己本不当是的样子而悲伤。但斯多葛学派主张的不是愚人，而是智者不能悲伤。\par

% structure: p0032 source=h2 target=subsection reason=relative-heading-rank; base=h2

\subsection{第九章 论心灵之扰动，在义人生活中表现为正确情感。}

但就这些心灵扰动的问题而言，我们在本著第九卷中已经回答了这些哲学家，表明这更多的是言辞之争而非事实之争，他们寻求的是争辩而非真理。在我们中间，根据圣经和纯正的教义，天主之城——那在此生朝圣旅途中按天主生活——的公民，既有恐惧也有渴望，既有悲伤也有喜乐。因为他们的爱被正确安置，所有这些情感都是正当的。他们害怕永罚，渴望永生；他们悲伤，因为他们自己内心叹息，等待成为义子，等待身体的救赎\BibleRef{罗 8:23}；他们在希望中喜乐，因为\BibleQuote{必成就那经上所写的话：死亡被胜利吞灭。}\BibleRef{格前 15:54}同样，他们害怕犯罪，渴望持守；他们在罪中悲伤，在善工中喜乐。他们害怕犯罪，因为他们听到\BibleQuote{由于罪恶增多，许多人的爱必会冷淡。}\BibleRef{玛 24:12}他们渴望持守，因为他们听到经上写着：\BibleQuote{那坚持到底的，必得救。}\BibleRef{玛 10:22}他们为罪悲伤，听到\BibleQuote{如果我们说我们没有罪，就是欺骗自己，真理也不在我们内。}\BibleRef{若一 1:8}他们在善工中喜乐，因为他们听到\BibleQuote{主喜爱乐意捐献的人。}\BibleRef{格后 9:7}同样，根据他们是刚强还是软弱，他们害怕或渴望受试探，在试探中悲伤或喜乐。他们害怕受试探，因为他们听到训诫：\BibleQuote{如果一个人陷于过犯，你们属神的人要用温和的心神扶助他，但要小心，免得你也受试探。}\BibleRef{迦 6:1}他们渴望受试探，因为他们听到天主之城的一位英雄说：\Quote{上主，求你察验我，考验我，煅炼我的肺腑和心思。}他们在试探中悲伤，因为他们看到伯多禄哭泣\BibleRef{玛 26:75}；他们在试探中喜乐，因为他们听到雅各伯说：\BibleQuote{我的弟兄们，当你们落在各种试探中，要认为是大喜乐。}\BibleRef{雅 1:2}\par

而且，他们不仅为自己的缘故经历这些情绪，也为那些他们所渴望得救、所害怕丧亡的人，以及那些损失或得救使他们悲伤或喜乐的人。因为如果我们这些从外邦人中来到教会的人，可以适宜地举出那位高尚而强大的英雄——他以软弱为荣，是万民在信德和真理上的\OriginalTerm{导师}{doctor}，他比所有同僚宗徒更加劳苦，用他的书信教导天主子民的各部族，这些书信不仅造就了当时的人，也造就了所有将要被聚集的人——我说，这位英雄，基督的运动员，由基督教导，由祂的圣神傅油，与祂同钉十字架，在祂内得荣耀，合法地在世界这个舞台上进行了一场伟大的竞赛，成为天使和世人观看的对象\BibleRef{格前 4:9}，向着高天召叫的奖赏直奔\BibleRef{斐 3:14}——我们非常喜乐地用信德之眼看见他与喜乐的人同喜乐，与哭泣的人同哭泣\BibleRef{罗 12:15}；虽然内外交困，有争战在外，有恐惧在内\BibleRef{格后 7:5}；渴望离世与基督同在\BibleRef{斐 1:23}；切望见到罗马人，好使他们中间也如在外邦人中一样结些果子\BibleRef{罗 1:11}；对格林多人心怀嫉妒，且因这嫉妒而害怕他们的心思受到腐蚀，失去在基督内的纯洁\BibleRef{格前 11:1}；为以色列人常有忧愁、心中不断伤痛\BibleRef{罗 9:2}，因为他们不认识天主的正义，企图建立自己的正义，不服从天主的正义\BibleRef{罗 10:3}；不仅表达他的悲伤，而且对某些从前犯罪而没有悔改他们污秽和淫乱的人痛心哀恸\BibleRef{格后 12:21}\par

如果这些情绪和情感，源自对善的爱和圣洁的爱德，被称作恶习，那么让我们允许那些真正是恶习的情感以美德之名流传吧。但是，既然这些情感在适当地运用时，遵循正确理性的指导，谁敢说它们是疾病或邪恶的情欲呢？因此，连主自己，当他屈尊就卑，取了奴仆的形体，过人的生活时，也完全没有罪，却在他认为应该运用这些情感的时候运用了它们。因为在他内，有真实的人的身体和真实的人的灵魂，所以也有真实的人的情感。因此，当我们在福音中读到，犹太人的心硬使他忧伤愤慨\BibleRef{谷 3:5}，他说\BibleQuote{我为了你们而欢喜，为使你们相信}\BibleRef{若 11:15}，当他要复活拉匝禄时甚至流了泪\BibleRef{若 11:35}，他切望与门徒们一起吃逾越节晚餐\BibleRef{路 22:15}，当他的受难临近时，他的心灵悲伤\BibleRef{玛 26:38}，这些情感确实不是虚假地归给他。但是，正如他在所乐意的时候成为人，同样，在他特定旨意的恩宠中，在所乐意的时候，他在自己的人灵魂中经历了这些情感。\par

但是我们必须进一步承认，即使这些情感是调节得当并合乎天主旨意的，它们也属于现世生活，而非我们所期待的未来生活，而且我们常常违背自己的意志而向它们屈服。因此，有时我们不由自主地哭泣，身不由己，这并非出于有罪的欲望，而是出于可赞扬的爱德。所以，在我们身上，这些情感源于人性的软弱；但在主耶稣身上并非如此，因为连祂的软弱也是祂权能的结果。然而，只要我们仍带着此生的软弱，如果我们完全没有这些情绪，我们就会变得更坏，而非更好。因为宗徒谴责并憎恶某些人，如他所说，是\BibleQuote{没有亲情}\BibleRef{罗 1:31}。神圣的圣咏作者也指责那些人，他论到他们说：\Quote{我期待有人同我一起哀哭，却没有一个。}因为当我们处在这悲惨之地时，要完全免于痛苦，正如世上一位文人学者所察觉和评论的，是必须以心智和身体感觉的迟钝为代价才能获得的。因此，希腊人称之为\OriginalTerm{不动心}{\GreekText{ἀπαθεια}}，而拉丁人如果可以，会称之为\Quote{\OriginalTerm{无感受性}{impassibilitas}}，如果这个词是指灵魂而非身体的无感受性，换言之，是指摆脱那些违反理性、扰乱心灵的种种情绪，那么它显然是一种美好且最可取的品质，但在现世生活中却是无法获得的。因为宗徒的话——不是普通人的宣认，而是极其虔诚、正直和圣善之人的宣认——说：\BibleQuote{如果我们说我们没有罪，就是欺骗自己，真理也不在我们内}\BibleRef{若一 1:8}。当人身上不再有罪时，那时才会有这种\OriginalTerm{不动心}{\GreekText{απάθεια}}。目前，如果我们生活清白，就够了；而谁若认为自己没有罪，他撇开的不是罪，而是宽恕。如果那种心灵不受任何情绪影响的状态被称为不动心，那么谁不会认为这种麻木比一切恶习更糟呢？的确，有理由认为，我们所希望的真福将没有恐惧或悲伤的任何刺痛；但谁若不与真理完全隔绝，会说在那里既没有爱也没有喜乐呢？但是，如果不感动心被理解为一种没有恐惧令人战栗、没有痛苦令人烦扰的状态，那么如果我们愿意按照天主的旨意生活，就必须在此生摒弃这种状态，而可以希望在应许给我们作为永恒境况的那种真福中享有它。\par

因为宗徒若望所说的那种恐惧：\BibleQuote{在爱内没有恐惧，相反，圆满的爱把恐惧驱逐于外，因为恐惧内含着惩罚；那恐惧的，在爱内还未圆满}\BibleRef{若一 4:18}——那种恐惧与宗徒保禄所感受的、担心格林多人被蛇的诡计所诱惑的恐惧不是同一类；因为爱能感受这种恐惧，是的，唯独爱才能感受它。但是，不在爱内的恐惧，正是保禄本人所说的那一种：\BibleQuote{因为你们所领受的不是奴隶的精神，以致再度恐惧}\BibleRef{罗 8:15}。至于那\Quote{纯洁的敬畏，永远常存}，如果它要在来世存在（否则如何能说它永远常存呢？），它并非一种阻止我们遭遇可能发生的邪恶的恐惧，而是保存我们于那不能失落的善。因为对于已获得的善的爱若是不可改变的，那么，那避免邪恶的恐惧当然可以说是无忧无虑的。因为大卫以\Quote{纯洁的敬畏}之名表示那意志，我们凭借这意志必然回避罪并防范罪，不是带着软弱的焦虑（担心我们可能强烈地犯罪），而是带着圆满之爱的安宁。或者，如果在那种最不可动摇的、永恒福乐的安宁中根本不存在任何恐惧，那么\Quote{对主的敬畏是纯洁的，永远常存}这句话就必须与另一句话\Quote{贫民的忍耐不会永远丧亡}作同样的理解。因为忍耐只有在需要忍受苦难的地方才必要，它本身不会永恒，但忍耐所引导我们达到的将是永恒的。所以，也许这\Quote{纯洁的敬畏}之所以被称为永远常存，是因为敬畏所导向的将永远长存。\par

由于如此——由于我们必须过善生以达到福生，善生具有所有这些正确的情感，恶生则具有错误的情感。但在永恒的福生中，将会有爱和喜乐，不仅是正确的，而且也是确定的；但恐惧和忧愁则完全没有。由此已经可以大致看出，天国之城的公民在此世的旅途中，过的是按照圣神而非按照肉体的生活——也就是说，按照天主而非按照人——并且可以看出，他们在旅途中要到达的那个不朽之境中，将是怎样的人。而那恶者的城市或群体，他们不按照天主而按照人生活，并且在敬拜虚假的神祇和藐视真神时接受人或魔鬼的教义，被那些邪恶的情感所震撼，如同被疾病和骚乱所动摇。如果其中有某些公民似乎抑制了并多少调节了那些激情，他们因不敬神的骄傲而如此得意，以致他们的病患越大，痛苦却越小。如果有人，其虚荣心因罕见而变得畸形，他们迷恋自己，因为他们不受任何情绪的刺激或鼓动，不被任何情感所感动或动摇，这样的人与其说是获得了真正的安宁，不如说是失去了全部人性。因为一件事并不因为不可弯曲而就是正确的，也不因为没有感觉而就是健康的。\par

% structure: p0039 source=h2 target=subsection reason=relative-heading-rank; base=h2

\subsection{第十章 是否应当相信我们原祖父母在乐园中、在犯罪之前，免于一切骚动。}

但这是一个合理的问题：我们的原祖或原祖父母（因为那是两个人的婚姻），在他们犯罪之前，在他们属兽性的身体中，是否经历了我们将在罪被清除并最终消灭后不会再在属神身体中经历的那种情感？如果他们确实经历了，那么他们又怎能在那被称为乐园的福乐之地被称为有福呢？因为被恐惧或悲伤影响的人，谁能被称为绝对有福呢？在如此丰盛的福泽中，那些人又能恐惧或遭受什么？在那里，既不必害怕死亡或疾病，也没有任何善意所渴望的事物是缺乏的，也没有任何事物存在以打断人身心之享受。他们对天主的爱是无瑕的，彼此之间的感情是忠实而真诚的婚姻之爱；从这爱中涌流出奇妙的喜悦，因为他们始终享受所爱之物。他们避免犯罪是平静的；只要保持这种状态，就没有任何其他邪恶能侵入并带来悲伤。或者，他们或许渴望触摸并吃那禁果，却又害怕死亡？于是恐惧和欲望，即使在那乐园中，也已经开始折磨这些最初的人类？断乎不可这样想，在无罪之处哪里会有这种事！事实上，渴望天主法律所禁止的事物，并因害怕惩罚而非因喜爱正义而克制，这本身已经是罪。我说，断乎不可这样想：在任何罪之前，就已经对那果子犯了我们的主警告我们不要对妇女犯的那种罪：\BibleQuote{凡注视妇女而贪恋她的，已经在心里与她犯了奸淫。}\BibleRef{玛 5:28} 因此，我们的原祖原父母当时是多么有福啊，他们不受任何心神动荡的搅扰，也不受任何身体不适的困扰。整个人类本该也如此有福，倘若他们没有引入那邪恶并传于后代，且没有后代犯下应受惩罚的罪；但这原始的福乐持续着，直到因着那祝福（即\BibleQuote{你们要生育繁殖，}\BibleRef{创 1:28}）所预定圣人的数目得以满全，然后那更崇高的福乐，即至福天使所享的福乐，就会被赐予——那福乐中必有确然的保障，使无人犯罪，也无人死亡；这样，圣人们将在未尝劳苦、痛苦、死亡之后生活，正如他们如今在复活中，在经历了一切之后将要生活一样。\par

% structure: p0041 source=h2 target=subsection reason=relative-heading-rank; base=h2

\subsection{第十一章——论原祖的堕落，在其身上本性受造为善，且只能由其创造者恢复。}

但天主预知万事，因此并非不知道人也将堕落；我们应当结合天主所预知和所预定的来思考这座圣城，而不应根据我们自己的、不包含天主预定的观念。因为人的罪既不能扰乱天主的计划，也不能迫使天主改变祂已决定的事；因为天主的预知已预先知道两者——即祂所造为善的人将变得何等邪恶，以及祂自身甚至要从这恶中汲取何等善。因为尽管天主被说成会改变祂的旨意（所以圣经拟人化地说甚至天主后悔了），但这是相对于人的期望或自然因果顺序而言的，而非针对全能者预知祂将要行的事。因此，正如经上所载，天主造了正直的人，\BibleRef{训 7:29} 因而人具有善的意志。因为若没有善的意志，就不能正直。所以，善的意志是天主的工作；因为天主造人时赋予了它。但先于人类一切恶行的第一个恶的意志，与其说是某种积极的工作，不如说是从天主的工作堕落到其自身工作的一种堕落。因此，由此产生的行为是邪恶的，因为它们不以天主为终点，而以意志本身为终点；所以意志或人本身，就其意志是邪恶而言，就好比结坏果子的坏树。此外，恶的意志虽然与本性不相和谐，反而对立，因为它是恶习或污点，但就所有恶习而言，它只能存在于一个本性之中，并且只存在于从无中受造的本性中，而非存在于创造者从自身所生的本性中，就像祂生了圣言，万物藉祂而受造。因为虽然天主用地上的尘土造了人，但尘土本身和一切地上的质料都是绝对从无中受造的；人的灵魂也是天主从无中创造并接合于身体的。但是，恶被善如此彻底地战胜了，以至于尽管它们被允许存在，以彰显天主最公义的眷顾甚至能善用它们，然而善可以在没有恶的情况下存在，如同在真实至高的天主自身，以及在每一个存在于这黑暗大气之上的无形和有形的天上受造物中；但恶不能脱离善而存在，因为恶所存在于其中的本性，就其是本性而言，是善的。而恶的被消除，不是通过移除任何由恶引入的本性或本性的部分，而是通过治愈和纠正那被败坏和腐化的东西。因此，当意志不为恶习和罪恶所奴役时，它才是真正自由的。天主起初赐予我们的正是这样的意志；它因其自身的过错而失去后，只能由起初能够赐予它的那一位来恢复。因此，真理说：\BibleQuote{如果圣子使你们自由，你们就真正自由了；}\BibleRef{若 8:36} 这等于说：如果圣子拯救你们，你们就真正得救了。因为祂是我们的解放者，因为祂是我们的救主。\par

那时，人生活在天主的管理之下，在一个既是物质又是精神的乐园里。因为那乐园并非仅为身体的益处而只是物质的，也非仅为心灵的好处而只是精神的；显然，它对于两者而言都是兼具的。但在那骄傲因而嫉妒的天使（我已在本书第十一和十二卷中尽可能多地论述了他的堕落，以及他的同伙，他们从天主的天使变成了他的天使）宁愿以某种帝国式的排场统治，也不愿服从他人，从精神乐园堕落后，试图把他那具有说服力的诡计灌输给人的心灵；人的未堕落状态引起了他的嫉妒，因为自己已经堕落；他选择了蛇作为他的口舌，在那有形的乐园中，蛇和所有其他地上的动物都与那两个人——男人和他的妻子——生活在一起，服从他们，且无害；他选择蛇，因为蛇滑溜，行动曲折蜿蜒，适合他的目的。这动物因他天使性本质的临在和更高力量而被制服于他的邪恶目的，他滥用它作为工具，首先在女人身上尝试欺骗，攻击那人类结合体中较软弱的部分，以便逐渐得到全人类；他估计男人不会轻易听从他或被欺骗，但可能会屈服于女人的错误。正如亚郎并非被诱导去同意百姓盲目地要他制造偶像，而是屈服于强求；正如撒罗满不可能盲目到以为偶像应当被崇拜，而是被女人的谄媚引入这样的亵渎；所以我们不能相信亚当受了欺骗，以为魔鬼的话是真的，而是因同类的吸引而屈服于女人，丈夫屈服于妻子，这人屈服于那唯一的人。因为宗徒并非无意义地说：\Quote{亚当没有被欺骗，而是女人受了欺骗，陷于违命；} \BibleRef{弟前 2:14} 他之所以这样说，是因为女人相信了蛇告诉她的，而男人无法忍受与他的唯一伴侣分离，即使这涉及罪的合伙。他并不因此罪责较轻，而是明知故犯。所以宗徒不说：\Quote{他没有犯罪，} 而是说：\Quote{他没有受欺骗。} 因为他表明他犯了罪，当他说：\Quote{藉着一人，罪恶进入了世界，} \BibleRef{罗 5:12} 紧接着更明确地说：\Quote{同亚当的过犯相似。} 但宗徒的意思是，那些不判断自己所行是罪的人就是受了欺骗；而亚当知道。否则怎么会有\Quote{亚当没有受欺骗？} 但由于尚未体验天主的严厉，他可能错误地认为自己的罪是小罪。因此，他不像女人那样受了欺骗，但他对于自己辩护将受到的判断受了欺骗：\Quote{你所赐给我、与我同居的女人，她给了我，我就吃了。} \BibleRef{创 3:12} 还有什么可说的？虽然他们并非都因轻信而受骗，但两人都落入了魔鬼的陷阱，被罪擒获。\par

% structure: p0044 source=h2 target=subsection reason=relative-heading-rank; base=h2

\subsection{第十二章——论人类原罪的本质。}

如果有人难以理解，为什么其他罪不像最初那些人的过犯那样改变人性，以致人性因这过犯而遭受我们所感受和看见的巨大腐化，并导致死亡，又被如此多狂暴冲突的情感所扰乱和折磨，且肯定与犯罪前大不相同——即使那时它仍寄居于动物性的身体中——我再说，如果有人为此触动，他不应认为那罪是微小轻微的，因为它关乎食物，且那种食物并非不好或有害，只因被禁止；因为在那无比幸福之处，天主不能创造和种植任何邪恶之物。但通过他颁布的诫命，天主赞赏了服从，这在某种程度上是理性受造物一切德行之母与守护者；理性受造物受造如此，服从对他有益，而实现自己意志置于创造者意志之上则是毁灭。既然这诫命在大量其他食物中禁止一种，遵守起来极其容易——记忆负担如此之轻——且尤其没有遇到私欲对遵守的抵抗（私欲后来才作为罪的惩罚后果出现），那么违反它的不义，与遵守它的容易程度成正比地更大。\par

% structure: p0046 source=h2 target=subsection reason=relative-heading-rank; base=h2

\subsection{第十三章——论在亚当的罪中，恶意志先于恶行。}

我们的原祖父母之所以陷入公开的悖逆，是因为他们早已在暗中败坏了；因为若没有邪恶的意志在先，恶行就永远不会发生。而邪恶意志的根源是什么，岂不就是骄傲？因为\Quote{骄傲是罪恶的开端。}\BibleRef{德 10:13} 而骄傲又是什么，岂不就是对不当高举的渴求？这种不当高举，就是灵魂离弃那应当作为其终向而紧抓住的天主，而自身成为某种意义上的终向。当灵魂以自身为满足时，便发生这种情况。而当它离弃那不变之善——这善本应比它自身更令它满足——的时候，它便如此行。这种离弃是自发的；因为如果意志曾坚定地持守那更高且不变之善（意志由此被光照而获得理智，被点燃而发出爱），它就不会转向自身以求满足，从而变得冰冷而昏暗；女人就不会相信蛇说的是真话，男人也不会宁愿顺从妻子的请求而违背天主的命令，也不会认为即使在罪恶的合伙中与伴侣结合是一种小罪。因此，恶行——即吃禁果的违命行为——是由已经邪恶的人犯下的。那\Quote{坏果子}\BibleRef{玛 7:18} 只能由\Quote{坏树}结出。但树之坏并非出于本性；因为显然它只能因意志的瑕疵而变坏，而瑕疵是违反本性的。然而，本性若非从无中生有，就不会因瑕疵而堕落。因此，它之所以是本性，是因为它由天主所造；但它之所以离弃天主，是因为它从无中生有。但人并非如此堕落以至于完全化为乌有；而是转向自身后，他的存在变得比当初紧抓住那至高存在者时更为局促。因此，存在于自身之中，即在天主离弃后以自身为满足，并非完全成为虚无，而是接近虚无。因此，圣经用另一个名称称呼骄傲的人，即\Quote{自悦者}。因为心被高举是好的，但不是高举向自己——因为这是骄傲——而是高举向主——因为这是服从，且只有谦卑的人才能如此。因此，在谦卑中有某种奇妙地高举心灵的东西，而在骄傲中则有某种贬低心灵的东西。这看起来确实矛盾：高升反而贬低，卑下反而高升。但虔诚的谦卑使我们服从那高于我们的；没有什么比天主更高于我们；因此，谦卑通过使我们服从天主而高举我们。但骄傲是本性的瑕疵，由于拒绝服从并背叛那至高者这一行为本身，它便堕入低微的状态；于是就发生了经上所记的话：\Quote{他们自高自傲时，祢就把他们推倒。}因为经上不是说\Quote{当他们已被高举之后}好像他们先被高举，然后才被推倒；而是说\Quote{当他们自高自傲时}就在那时他们被推倒——也就是说，这高举本身就已经是一场堕落。因此，谦卑被特别推荐给在这世上作客旅的天主之城，并在天主之城及其君王基督身上特别彰显；而相反的骄傲之恶，根据圣经的见证，则特别统治着它的敌人——魔鬼。而这一点无疑正是我们所谈论的两座城之间的巨大区别：一座是虔诚之人的团体，另一座是不虔之人的团体，各自与支持己方的天使联合，一座被自爱引导和塑造，另一座被天主之爱引导和塑造。\par

因此，魔鬼之所以没有在人公然明显地犯下天主所禁止的罪之前就诱骗他，是因为人还未开始为自己而活。正是这一点使他会欣然地听从\Quote{你们将如同天主一样}\BibleRef{创 3:5} 这话；如果他们服从地紧抓住自己至高而真正的终向，而不是骄傲地为自己而活，他们本会更容易达到这一目的。因为受造的神之所以是神，并非凭自身之故，而是由于分享真实的天主。人因贪求更多而变得更少；因渴望自足而离弃了那真正使之满足的天主。因此，这种邪恶的欲望——它促使人仿佛自身就是光明那样悦乐自己，从而使人离弃那光明（若他追随它，自己就会成为光明）——我所说的这种邪恶欲望早已隐秘地存在于他里面，而公开的罪行只是它的结果。因为经上所记的是真实的：\Quote{骄傲在败坏以先，谦卑在荣耀之前}\BibleRef{箴 18:12}；也就是说，隐秘的败坏先于公开的败坏，而前者却不被算作败坏。因为谁会把高举视为败坏呢，尽管至高者一被离弃，堕落就开始了？但谁不承认，当显然且无可置疑地违犯诫命时，那就是败坏呢？因此，天主的禁令所针对的正是这样的行为，一旦犯下，便无法以任何行义的借口来辩护。我敢说，骄傲的人堕入公开且无可争辩的罪行，并因此而不悦自己，这对他们是有益的，因为他们早已因悦乐自己而堕落了。因为伯多禄在他哭泣并对自己不满时，比在他大胆自恃并满足自己时更为健康。这正是圣咏作者所宣称的：\Quote{上主，求祢使他们满面羞惭，好让他们寻求祢的名}；也就是说，那些在追求自己荣耀时悦乐自己的人，愿他们在寻求祢的荣耀时，因祢而喜悦并满足。\par

% structure: p0049 source=h2 target=subsection reason=relative-heading-rank; base=h2

\subsection{第十四章——论罪中之骄傲，其恶甚于罪本身。}

然而，有一种更恶劣、更该受罚的骄傲，即使在明显的罪过中也要寻找托辞，正如我们原祖所做的那样。\Quote{女人说：\InnerQuote{那蛇哄骗了我，我就吃了}；男人说：\InnerQuote{你赐给我与我同居的女人，她把树上的果子给了我，我就吃了。}}\BibleRef{创 3:12} 这里没有求赦的话，没有求医治的恳求。因为他们虽然不像加音那样否认自己做了这事，但他们的骄傲却试图将罪过推给别人——女人的骄傲推向蛇，男人的骄傲推向女人。但在明知违反天主诫命的情况下，这样做与其说是原谅自己，不如说是控告自己。因为女人因蛇的诱骗而犯罪，男人因女人的给予而犯罪，这并没有减轻违命的事实，仿佛有谁比天主更值得我们信赖或顺从似的。\par

% structure: p0051 source=h2 target=subsection reason=relative-heading-rank; base=h2

\subsection{第十五章——论我们原祖因不顺从所受惩罚的正义}

因此，因为罪过是轻视天主的权威——祂创造了人，按自己的肖像造了人，将人置于其他动物之上，将人安置在乐园中，使人享有各种丰富与安全，给人定的诫命不多、不重、也不难，但为了使顺服的操练易于实行，只给了一个极简短、极轻省的规诫，借以提醒那受造物，其服务应当是自由的，而祂是主——所以随之而来的惩罚是公义的，并且是这样的惩罚：人若遵守诫命，本可在肉体上也成为属神的，却变成了在精神上也属肉体的；并且，正如他在骄傲中寻求自我满足，天主按其公义舍弃了他，使他不能在他所妄求的绝对独立中生活，反而失去了他所渴望的自由，活在对自身不满的、艰难悲惨的奴役中，屈服于他因犯罪而顺从的那一位，注定要违背自己的意愿而肉身死去，正如他心甘情愿地在精神上死亡，甚至被判处永死（若非天主的恩宠拯救他），因为他离弃了永生。凡认为这样的惩罚过分或不公的人，都显示出他无法衡量在本来极易避免的罪中犯罪的严重不义。正如亚巴郎的服从被公认为伟大，因为所命之事——杀他的儿子——极其困难，同样，在乐园中，不顺从的罪更为重大，因为所命之事的难度微乎其微。正如第二人（基督）的服从更值得赞美，因为祂顺从\BibleQuote{至死}\BibleRef{斐 2:8}，同样，第一人（亚当）的不顺从更可憎，因为他至死都不顺从。因为当不顺从所附带的惩罚巨大，而造物主所命之事容易时，谁能充分估量，在如此容易的事上，不顺从如此伟大权威的命令，即或是那权威以如此可怕的惩罚来威慑，这罪孽是何等深重呢？\par

简言之，一言以蔽之，那罪中的不服從的惩罚岂不就是不服從吗？因为人的悲惨若不是他自己不服從自己，又是什么呢？其结果是，由于他不愿意做他能做的事，如今他却愿意做他所不能做的事。因为在他犯罪之前，他在乐园里虽不能做所有事，但他只愿意做他能做的事，因此他能做所有他愿意做的事。但现在，正如我们在他的后代中认出的，以及圣经所见证的：\Quote{人好似虚空。} 因为谁能数算他愿意却不能做多少事，只要他仍不服從自己，也就是说，只要他的心智和肉体不服从他的意志？因为不由自主地，他的心智常常被扰乱，他的肉体受苦、衰老并死亡；而且不由自主地，我们遭受我们所遭受的其他一切，如果我们的本性及其所有部分绝对服从我们的意志，我们就不会遭受这些。但阻碍肉体服务的不正是肉体的软弱吗？然而，它的服务\Emph{如何}受到阻碍又有什么要紧呢？只要事实仍然如此：由于至高天主的公义报应——我们拒绝服从和事奉的那位——原服于我们的肉体，现在以不服从来折磨我们，尽管我们的不服從给我们自己带来了麻烦，而不是给天主。因为祂并不像我们需用身体那般需用我们的事奉；因此我们所做的对祂不是惩罚，但我们所受的却是对我们的惩罚。所谓身体的痛苦是灵魂在身体内并经由身体所受的痛苦。因为肉体离开灵魂能独自感到什么痛苦或欲望呢？但当说肉体欲望或受苦时，如我们已解释，是指人如此行，或灵魂的某部分受到肉体感觉的影响，无论是引起痛苦的粗糙感觉，还是引起快乐的柔和感觉。但肉体的痛苦不过是灵魂因肉体而产生的不适，以及对其受苦的一种退缩，正如被称为忧伤的灵魂痛苦是对那些违反我们意志而发生之事的退缩。但忧伤常由恐惧先行，恐惧本身在灵魂中，不在肉体里；而身体的痛苦并不先有任何肉体的恐惧，这种恐惧在痛苦之前可以在肉体中感觉到。但快乐之前有一种在肉体中感觉到的欲望，如同渴求，例如饥渴和那生育的食欲，最常被冠以贪欲之名，尽管这是所有欲望的总称。因为愤怒本身被古人定义为不过是报复的贪欲；尽管有时人甚至对不能感受其报复的无生命物体发怒，如当人折断一支笔，或压碎一支书写不佳的羽毛笔时。然而即使如此，虽较不合理，在它自己的方式上也是一种报复的贪欲，并可说是报应的一种奥妙阴影——报应法则——即作恶者应受恶报。因此有报复的贪欲，称为愤怒；有金钱的贪欲，名为贪婪；有征服的贪欲，不论以何手段，称为好胜心；有喝采的贪欲，名为自夸。有许多不同的贪欲，其中一些有自己的名称，而另一些则没有。因为谁能轻易给统治的贪欲命名呢？它在暴君的灵魂中有强大影响，如内战所证明。\par

% structure: p0054 source=h2 target=subsection reason=relative-heading-rank; base=h2

\subsection{第十六章 —— 论贪欲之恶——一个词虽可适用于许多恶习，却特别指性不洁。}

因此，尽管贪欲有许多对象，但未指定对象时，贪欲一词通常使人联想到生殖器官的情欲激动。这种贪欲不仅占据整个身体和外在肢体，还在内部自我感受，并以一种心智情感与身体欲望混合的激情激动整个人，以至于所产生的快乐是所有身体快乐中最大的。这快乐如此具有占有性，以至于在它完成的那一刻，一切心智活动都暂停。智慧与神圣喜乐的朋友，身为已婚者，却知道如宗徒所说：\BibleQuote{如何以圣化和尊荣持守自己的身体，不在欲望的疾病中，如同不认识天主的外邦人，}\BibleRef{得前 4:4}难道不会更喜欢，如果可能的话，没有这种贪欲而生子女，以便在这个生育后代的功能中，为此目的而创造的身体器官不被贪欲的热火刺激，而是由他的意志所驱动，如同其他器官为各自的目的服务他一样？但即使是那些以这快乐为乐的人，也不能凭自己的意志被驱动，无论他们限于合法的快乐还是逾越到不合法的快乐；但有时这贪欲在他们不情愿时强求他们，有时在他们想要感受它时却失灵，以致虽然贪欲在心中沸腾，却不在身体中激动。因此，奇怪的是，这种情绪不仅不服从生育后代的合法欲望，也拒绝服务淫乱的贪欲；而且虽然它常常以其全部联合力量对抗抵抗它的灵魂，有时它也自相矛盾，当它激动灵魂时，身体却不动。\par

% structure: p0056 source=h2 target=subsection reason=relative-heading-rank; base=h2

\subsection{第十七章 —— 论我们原祖父母在他们卑劣可耻的罪之后所见的裸体。}

这番欲望与羞耻之感的联系尤为恰当；同样，这些肢体——它们不被我们的意志所驱动或约束，而是仿佛拥有某种独立的自主权——被恰如其分地称为\Quote{羞耻的}。它们在犯罪前的状态是不同的。因为正如经上所载：\BibleQuote{他们赤身露体，并不羞耻}\BibleRef{创 2:25}——并非因为他们不知道自己的赤身，而是因为赤身尚未成为可耻的，因为那时欲望尚未违背意志的同意而驱动这些肢体；肉体尚未因其不顺从而证实人的悖逆。因为他们并非如无知庸众所想象的那样生来盲目；因为亚当能看见他为之命名的动物，关于厄娃，我们读到：\BibleQuote{女人见那棵树的果子好作食物，也悦人的眼目}\BibleRef{创 3:6}。因此，他们的眼睛是睁开的，但尚未睁开到这一点，即尚未注意到并认识到恩宠的外衣赐予了他们什么，因为他们并未意识到自己的肢体与意志相争。但当他们被剥夺了这恩宠，为使其悖逆受到适当的惩罚时，他们身体肢体的运动中开始出现一种无耻的新奇，使赤身变得不雅；这立刻使他们注意到并感到羞耻。因此，在他们公然违犯天主的命令之后，经上记载：\BibleQuote{他们二人的眼睛就开了，才知道自己是赤身露体，便拿无花果树的叶子，为自己编作裙子}\BibleRef{创 3:7}。\BibleQuote{他们二人的眼睛就开了}，并非指他们此前不能看见，因为他们早已能看见，而是指他们分辨出所失去的善和所堕入的恶。因此，那棵禁止他们触摸的树本身也由此被称为知善恶树，因为若他们吃了它，它就会赋予他们这种知识。因为疾病的不适显露出健康的快乐。因此，\BibleQuote{他们知道}自己\BibleQuote{是赤身露体的}——失去了那使他们不为身体的赤身而羞耻的恩宠，那时罪的律法尚未与他们的心神对抗。于是他们获得了一种知识，如果他们信靠地服从天主，拒绝犯下那使他们陷入不忠和悖逆的有害后果的罪过，他们本会有福地对此无知。因此，他们为自己肉体的悖逆感到羞耻——这肉体的悖逆惩罚并见证了他们的悖逆——\BibleQuote{便拿无花果树的叶子，为自己编作裙子}，即遮羞布；因为有些译者将这个词译为\OriginalTerm{遮羞布}{succinctoria}。\OriginalTerm{运动短裤}{Campestria}确实是一个拉丁词，但用于指那些在广场上脱衣锻炼的年轻人所穿的类似用途的短裤或围裙；因此，那些这样束腰的人通常被称为\OriginalTerm{穿运动短裤的人}{campestrati}。羞耻感谦逊地遮盖了那些被欲望不服从地驱动、与意志相悖的东西，而这意志因其自身的不服从而受到惩罚。因此，所有民族，都从那一个始祖繁衍而来，有如此强烈的本能去遮盖这些羞耻的部位，以至于有些野蛮人甚至在沐浴时也不露出它们，而是穿着短裤洗澡。在印度的幽暗隐修处，虽然有些哲学家赤身露体，因而被称为\OriginalTerm{裸体智者}{gymnosophists}，但他们也对这些部位例外，加以遮盖。\par

% structure: p0058 source=h2 target=subsection reason=relative-heading-rank; base=h2

\subsection{十八、论与所有性行为相伴的羞耻。}

情欲需要黑暗和隐秘才能达到满足；这不仅在追求非法交合时如此，甚至在地上之城所合法化的那种淫乱也是如此。在没有刑罚恐惧之处，这些被允许的快乐仍回避公众视线。即使为这情欲作了准备，也同时提供了隐秘；而情欲发现很容易就废除法律的禁止，无耻却无法除去遮蔽的帷幕。因为就连无耻的人也把这事称为羞耻；尽管他们喜爱快乐，却不敢显露出来。怎么！难道连为生养子女而由法律认可的夫妻房事，尽管合法而尊贵，不也寻求避人耳目吗？新郎抚爱新娘之前，不也要摒退侍从，甚至伴郎伴娘以及因亲密关系而被允许进入洞房的朋友吗？罗马最伟大的雄辩家说，一切正当的行为都希望被置于光明之中，即希望被人知晓。然而，这个正当的行为虽然如此渴望被人知晓，却羞于被人看见。谁不知道夫妻之间为了生儿育女所行的事呢？不正是为此目的，妻子才以如此隆重的仪式出嫁吗？然而，当这一众所周知的行为旨在生育子女时，连他们可能已经出生的子女都不被容许作见证。这个正当的行为寻求光明，就其渴望被人知晓而言，但它却害怕被人看见。为什么如此，岂不是因为那本性上适宜和得体的事，竟然伴随着一种因罪而产生的羞耻刑罚吗？\par

% structure: p0060 source=h2 target=subsection reason=relative-heading-rank; base=h2

\subsection{十九、论如今正如人犯罪前一样，有必要用智慧的约束力来驾驭愤怒和情欲。}

因此，就连那些接近真理的哲学家也承认，\OriginalTerm{愤怒}{anger}和\OriginalTerm{私欲偏情}{lust}是邪恶的心灵情感，因为即使当它们指向智慧并未禁止的对象时，它们也是以一种不受控制、不合规矩的方式被激动，因此需要心智和理性的调节。他们断言，心灵的这第三部分仿佛驻扎在某种堡垒中，以统治其他部分，这样，当它统治而它们服侍时，人的正义便得以完整保存。那么，这些被他们承认即使在智慧而节制的人身上也是邪恶的部分，以致心灵必须以其镇定和约束的影响力，将它们从不法地趋向的对象上勒住并召回，并允许它们接近智慧法则所认可的对象——即愤怒，\Emph{例如}，可以被允许用于执行正当权威，私欲偏情可以被允许用于繁衍后代的责任——我说，这些部分在犯罪前的乐园里并非邪恶，因为它们从未违反神圣意志而趋向任何需要以理性的约束缰绳加以遏制的对象。因为虽然现在它们以这种方式被激动，并受到一种勒住和约束的力量的调节——那些度节制的、正义的、虔敬生活的人有时轻易、有时更困难地运用这种力量——但这并不是本性的健全健康，而是由罪而来的软弱。那么，为何羞耻并不掩盖由愤怒或其他情感所支配的行为和言语，就如它掩盖私欲偏情的冲动一样，岂不是因为我们用于完成这些行为和言语的身体肢体，并非由情感本身所驱动，而是由同意的意志的权威所驱动？因为一个在愤怒中责骂甚至打击别人的人，若不是他的舌头和手被意志的权威所驱动，他就不可能这样做，正如在没有愤怒时它们也被驱动一样。但是生殖器官如此受制于私欲偏情的统治，以至于它们的冲动除了私欲偏情所传递的之外别无其他。我们所羞耻的正是这个；正是这个在旁观者的眼前羞红地隐藏起来。一个人宁愿在一群见证者面前不义地向某人发泄愤怒，也不愿在一人眼前无辜地与妻子交合。\par

% structure: p0062 source=h2 target=subsection reason=relative-heading-rank; base=h2

\subsection{第二十章——论\OriginalTerm{犬儒派}{Cynics}的愚昧兽性。}

正是这一点被那些犬儒派或犬儒哲学家所忽视，他们违背了人类谦逊的本能，夸耀地宣扬其不洁而无耻的意见——实在配得上狗——即，既然婚姻行为是合法的，没有人应该羞于在街上或任何公共场所公开进行。本能的羞耻战胜了这一狂想。因为虽然据说\Person{第欧根尼}曾一度敢于将他的意见付诸实践，以为如果他那异乎寻常的无耻行为深深铭刻在人类的记忆中，他的学派就会更加出名，但此后这一榜样并未被效仿。羞耻对他们的影响，使他们能在人前脸红，胜过错误使他们模仿狗的愿望。而且可能，即使在\Person{第欧根尼}和那些模仿他的人的例子中，也仅是性交的外观和假装，而非真实。甚至在今天，仍可见到犬儒派哲学家；因为这些人是不满足于仅穿\OriginalTerm{外袍}{pallium}，还持一根棍棒的犬儒派；然而他们中没有一个人敢于做我们所说的这件事。如果他们做了，他们会被人群唾弃，甚至被石头砸死。那么，人性无疑对这等私欲偏情感到羞耻；而且这是正义的，因为这些身体肢体的不服从和对意志的蔑视，是人类第一次犯罪的惩罚的清晰见证。并且，这些部分特别显出这种惩罚是适当的，因为正是通过这些部分，那被第一次和伟大的罪所败坏的本性得以繁衍——那个罪，无人能逃脱其邪恶的后果，除非天主的恩宠在他个人身上赎偿了那在全体共有的、当所有人都包含在一人之中时所犯下、并由天主的公义所惩罚的罪。\par

% structure: p0064 source=h2 target=subsection reason=relative-heading-rank; base=h2

\subsection{第二十一章——人的过犯并未取消人在犯罪前所领受的多产祝福，却使之染上私欲偏情的疾病。}

因此，我们万万不可设想我们的原祖父母在乐园中有那种后来使他们羞愧并遮盖裸体的情欲，或是藉着情欲来履行天主的祝福：\BibleQuote{你们要生育繁殖，充满大地} \BibleRef{创 1:28}，因为情欲是在罪之后开始的。罪之后，我们的本性失去了对全身的控制，却没有失去所有羞耻，它察觉、注意、羞惭并遮盖了它。但那关于婚姻的祝福，鼓励他们生育繁殖充满大地，虽然在他们犯罪后仍然持续，却是犯罪前赐下的，为使人知道生育子女是婚姻的光荣，而非罪的惩罚。但现在，人们因不知乐园的幸福，便以为子女在那里只能以他们所知的现在的方式——即藉着连体面的婚姻也羞于启齿的情欲——被孕育；有些人不仅否认，还以怀疑的态度嘲笑神圣的圣经，其中我们读到我们的原祖父母犯罪后，羞于自己的裸体并遮盖它；另一些人虽然接受并尊敬圣经，却以为\Quote{你们要生育繁殖}这句并不指肉体的生育，因为类似的话用于灵魂，如\Quote{你必以力量在我的灵魂中加增我}；同样，在创世记随后的话\Quote{你们要充满大地，治理大地}中，他们以\Quote{大地}指肉体，灵魂以其临在充满肉体，并在力量加增时统治之。他们主张，子女的孕育无论当时还是现在都不能没有情欲，而情欲是在犯罪后被点燃、被察觉、被羞惭并被遮盖的；甚至子女本来不会被生在乐园，而只会生在乐园之外，正如实际发生的那样。因为他们在被逐出乐园后才结合生子。\par

% structure: p0066 source=h2 target=subsection reason=relative-heading-rank; base=h2

\subsection{第二十二章——论婚姻结合最初由天主建立并祝福}

但我们自己毫不怀疑，因天主祝福而生育繁殖充满大地，是天主起初在人类犯罪前设立婚姻时的恩赐，那时他造了他们男和女——即两个明显不同的性别。正是天主的工作蒙受了祂的祝福。因为圣经甫一说\BibleQuote{祂造了他们男和女} \BibleRef{创 1:27}，便立刻接着说：\BibleQuote{天主祝福了他们，天主对他们说：你们要生育繁殖，充满大地，治理大地}等等。虽然这一切都可以合宜地作属灵的解释，但\Quote{男和女}不能理解为一个人里面的两样事物，仿佛在他里面有一个统治的，另一个被统治的；但很明显，他们被造为男和女，具有不同性别的身体，正是为了生育后代，从而增加、繁殖并充满大地；否认这样明显的事实是极大的愚蠢。我们的主在回答是否可以因任何缘故休妻的问题时（因为摩西因以色列人心硬，准许写休书），所说的话，不是关于命令的灵和服从的身体，也不是关于理性的灵魂统治和非理性的欲望被统治，也不是关于至高的默观德性和从属的实践，也不是关于心智的理解和身体的感官，而是关于两性互相结合的婚姻结合。祂回答说：\BibleQuote{你们没有念过吗？那从一开始创造他们的，造了他们一男一女，并且说：为此，人要离开父母，与妻子结合，两人成为一体。这样，他们不再是两个，而是一体。所以，天主所结合的，人不可拆散} \BibleRef{玛 19:4}。所以，确实，人从一开始就被造为两性，即男和女，正如我们现在所看到和知道的，他们被称为一体，或由于婚姻结合，或由于女人的起源——女人是从男人的肋骨被造的。正是根据天主亲自设立的这个原始典范，宗徒劝告所有丈夫要特别爱自己的妻子。\BibleRef{厄 5:25}\par

% structure: p0068 source=h2 target=subsection reason=relative-heading-rank; base=h2

\subsection{第二十三章——若人不犯罪，乐园中是否仍应有生育，或那里是否应有贞洁与情欲之争}

但谁说若没有罪，就不应有交合或生育，实际上就是说人的罪对于完成圣人的数目是必需的。因为如果这两人不犯罪就该继续独居，因为据推测，如果他们不犯罪就不能生育子女，那么罪就必然是必需的，为的是不仅有两个义人，而是有许多义人。如果这一点不能荒谬地维持，我们就必须相信，即使没有人犯罪，那足以完成这最蒙福的城邑的圣人数目，也会像现今一样多，如今天主的恩宠从大群罪人中聚集其国民，只要这世界的儿女们生养并被生养。\BibleRef{路 20:34}\par

因此，那配享乐园幸福之婚姻，若无罪过，本应产生可欲之果实而无情欲之羞耻。但此事如何可能，如今已无实例可教我们。然而，这不应显得不可思议，因为既然如今有许多肢体服从意志，当时某一肢体亦能在无情欲下服从意志。我们现在不是随心所欲地移动手足，以完成我们藉这些肢体所欲行之事吗？我们在它们身上毫无抵抗，而感知它们甘为意志之仆，无论在我们自己或他人身上，尤其是在从事机械操作的工匠身上，因着勤勉的练习，天性的软弱与笨拙变得惊人地灵巧。我们岂能不信，正如所有那些肢体顺从地服务意志，同样地，那些肢体也应当履行生育的功能，尽管情欲——违命之罚——付之阙如？\Person{西塞罗}在其\WorkTitle{共和论}中讨论政体差异时，不是引用了一个来自人性的比喻，说我们命令身体肢体如同子女，它们如此顺从；但灵魂的恶习部分则须像奴隶般对待，以更严厉的权威加以压制吗？无疑，按自然秩序，灵魂优于身体；但灵魂命令身体比命令自身更容易。然而，我们此刻所谈论的情欲，因此更显可耻，因为灵魂在其中既不能自主以至于毫无情欲，也不能主身以致肢体服从意志；因为若如此被管治，便不会有羞耻。但如今，灵魂羞耻于那本性低劣且隶属于它的身体竟反抗其权威。因为在其他情绪中灵魂所遇的抵抗，羞耻较少，因为抵抗来自它自身，因此当它被自身征服时，它自身即是征服者，尽管这征服是无序邪恶的，因为它由灵魂中本应服从理性的部分完成，但既然由它自身的部分和能力完成，这征服，如我所说，是它自己的。因为当灵魂征服自身达到应有的服从，使其不合理之动受理性控制，同时它又服从于天主，这便是德性与可赞的征服。然而，当灵魂受到自身恶性部分的抵抗，相较于受到身体——虽与它不同且低劣，且依赖它获得生命本身——抵抗其意志和秩序，羞耻较少。\par

但是，只要意志将其他肢体保留在它的权威之下，没有这些肢体，为情欲所激动、反抗意志的肢体就无法达成它们所求的，贞洁就被保全，而罪中之乐就被舍弃。确实，若非有罪过的违命被惩罚性的违命所报应，乐园的婚姻就不会知道这种斗争和反叛，即意志与情欲之间的争吵，以使意志满足而情欲受抑；相反，那些肢体，像其余所有肢体一样，应当服从意志。生育的田地应当由为此目的而造的器官撒种，如同大地由手撒种一样。而如今，当我们试图更精确地探讨这一主题时，羞耻心阻碍我们，迫使我们请求贞洁之耳原谅；但若没有罪，就没有理由这样做，而我们可以自由地、不必担心似乎淫秽地谈论所有那些一旦思考此主题就会浮现的问题。甚至不会有任何可称为淫秽的词语；凡可能关于这些肢体所说的一切，都将如关于身体其他部分所说的那样纯洁。因此，凡以不洁之心来阅读这些篇章的人，让他责难自己的性情，而非他的本性；让他谴责他自己不洁的行为，而非我们迫于需要所用的言辞——每一个纯洁虔敬的读者或听者都会乐意原谅我，当我揭露那种仅仅依据自身经验、对任何超越之事毫无信心的怀疑论的愚昧。他不会因宗徒谴责妇女的可怕邪恶而震惊：她们\Quote{改变了本性的用途，而作了违反本性的事}\BibleRef{罗 1:26}，他将毫无震惊地阅读这一切，尤其因为我们不像\Person{保禄}那样引述并谴责一种该死的污秽，而是在尽可能解释人类的生育，同时我们与\Person{保禄}一样避免任何淫秽的语言。\par

% structure: p0072 source=h2 target=subsection reason=relative-heading-rank; base=h2

\subsection{如果人仍保持无罪和顺命于乐园，生育器官应如其他肢体一样服从意志。}

那么，男人将按需撒种，女人将接受它，生育器官由意志驱动，而非为情欲所激动。因为我们不仅随意移动那些由坚硬骨骼关节组成的肢体，如手、脚、手指，而且也随意移动那些由松弛柔软的神经组成的肢体：我们可以使它们运动、伸展、弯曲、扭转、收缩或僵硬，正如我们对口部和面部的肌肉所做的那样。肺脏，除了大脑之外最娇嫩的内脏，因此小心地藏在胸腔中，然而为了吸气和呼气的所有目的，以及发出和调节声音，当我们呼吸、呼气、说话、喊叫或歌唱时，它们服从意志，就像风箱服从铁匠或管风琴师一样。我不再强调有些动物天生就有能力移动全身覆盖的皮肤上的一个斑点，如果它们感觉到了想要赶走的东西——这种能力如此之大，通过这种皮肤的颤抖抖动，它们不仅能抖落落在身上的苍蝇，甚至能抖掉刺入肉中的长矛。人类确实没有这种能力；但这能成为理由，认为天主不能将这种能力赐给祂愿意拥有的受造物吗？因此，人本身也很可能曾对自己的肢体享有绝对权力，若不是因他的违命而丧失了这一权力；因为天主塑造他使他身体中如今只有情欲才能移动的部分只由意志移动，这并非难事。\par

我们也知道，有些人的身体构造与他人不同，拥有某种罕见而卓越的能力，能做他人无论如何努力也做不到的事，甚至听到别人这样做时也几乎难以置信。有人能移动耳朵，或单耳或双耳齐动。有人能不移动头部，将头发拉到前额，并随意前后移动整个头皮。有人轻轻按压腹部，便能吐出一大堆各式各样吞下的东西，完好无损，仿佛从袋中取出。有人能如此逼真地模仿鸟兽和其他人的声音，若非亲眼看见，根本无法分辨。有人对肠子有如此控制力，能随意连续放屁，产生歌唱的效果。我本人认识一个人，他能随己意出汗。众所周知，有人能随意哭泣，泪如泉涌。但更不可思议的是，我们的一些弟兄最近亲眼所见之事。在卡拉马教会教区有一位名叫雷斯提图特的司铎，每当他愿意时（那些想看如此非凡现象的人会请求他这样做），有人模仿哀悼者的哭声，他便变得毫无知觉，躺卧如死，不但他人拧他、刺他时毫无感觉，甚至用火烧他、灼伤他时，他也感觉不到疼痛，只是事后因伤口而痛。他的身体保持不动，并非出于自制，而是因为无知觉，这由他呼吸如同死人一般得到证明；然而他说，当有人说话特别清晰时，他能听到声音，但仿佛从远处传来。既然在此必死而悲惨的人生中，身体以许多超越自然常轨的显著动作和状态为某些人服务，那么，有什么理由怀疑，在人因罪而陷入这软弱可朽的状况之前，他的肢体可以无情欲地服从其意志以繁衍后代呢？人因抛弃天主而被交于自己，因为他寻求自我满足；不服从天主，他便连自己也不能服从。因此，他陷入显而易见的悲惨境地：无法按自己的意愿生活。因为如果他按自己的意愿生活，他会自以为有福；但如果他邪恶地生活，则不可能真有福。\par

% structure: p0075 source=h2 target=subsection reason=relative-heading-rank; base=h2

\subsection{第二十五章　论真正的幸福，此世生命无法享有。}

然而，若我们更仔细地审视，便会发现没有人能按自己的意愿生活，除非是有福之人，而只有义人才是有福的。但连义人本身也不能按自己的意愿生活，直到他抵达那不会死亡、不受欺骗、不受伤害的境界，并确信这将成为他永恒的状态。因为本性要求如此；本性在获得其所寻求之物以前，不能算是完全而圆满的幸福。但如今有谁能按自己的意愿生活呢？他连活着都不能自主：他愿意活着，却被迫要死。那么，一个不能活得如己所愿的人，怎能按自己的意愿生活呢？或者，若他愿意死，他怎能按自己的意愿生活，既然他甚至不愿意活着？或者，若他愿意死，并非因为厌恶生命，而是为了死后活得更好，那么他仍未按自己的意愿生活，只是期望通过死亡达到他所愿望的境地后，才能如此生活。但即使承认他按自己的意愿生活，因为他强迫自己，使自己不愿求得不到的东西，只愿求能得到的东西（如泰伦提乌斯所说：\Quote{既然你不能做你所愿做的事，就愿你所能做的事吧}），那么，他是否因忍耐悲惨而成为有福的呢？因为幸福的生活只属于爱它的人。如果它被爱并拥有，它必定比一切其他事物更被热切地爱；因为任何其他被爱之物，都必须为了幸福生活之故而爱。如果它被以应有的方式爱——而一个不以应有方式爱幸福生活的人，也不是有福的——那么，如此爱它的人，必不能不愿它成为永恒的。因此，只有当它成为永恒时，它才是有福的。\par

% structure: p0077 source=h2 target=subsection reason=relative-heading-rank; base=h2

\subsection{第二十六章　论我们当相信，在乐园中，我们的原祖父母生子而无羞惭。}

在乐园里，人只要渴望天主所命令的，就照他所愿的生活。他生活在享见天主的喜乐中，并因天主的良善而成为良善；他毫无匮乏，并且有能力永远这样生活下去。他有食物以免饥饿，有饮品以免口渴，有生命树以免衰老侵袭。他的身体没有腐朽，也没有腐朽的种子能在体内引起任何不适的感觉。他不惧怕内在的疾病，也不惧怕外在的意外。最健康的体质祝福他的身体，绝对的宁静祝福他的灵魂。正如乐园中没有过度的炎热或寒冷，其居民也免于恐惧和欲望的变迁。那里没有任何忧愁，也没有任何愚昧的喜乐；真正的喜悦从天主面前不断涌出，因为他以\Quote{纯洁的心、良好的良心和真诚无伪的信德}被爱。\BibleRef{弟前 1:5}夫妻之间真诚的爱情确保了和谐。身体与精神和谐共处，诫命得以不劳而守。没有倦怠使他们的休闲变得厌倦，没有睡意打断他们工作的愿望。\Emph{在如此顺利的环境和人的幸福中，我们切不可猜测：生育不能没有情欲的病症；而是：那些肢体，与其他肢体一样，是出于意志的示意而运动，并且在无欲火的诱惑刺激下，丈夫以身心平静、不损害完整，将精液注入妻子的子宫。因为尽管经验无法证明，也不应当不信；那时，那些身体部位不是由沸腾的情欲驱动，而是由自主能力在必要时加以运用；因此，那时能够保持女性生殖器的完整，将男性精液射入妻子的子宫，正如现在能够保持完整从童贞女的子宫流出月经血一样。同样的通道，既能射出，也能注入。因为为了生产，不是痛苦的呻吟，而是成熟的冲动使女性内脏放松；同样，为了} \Emph{受孕和怀胎，不是情欲的冲动，而是自愿的运用将两性结合。}我们谈论的是如今可耻的事，尽管我们尽力想象它们尚未变得可耻时的样子，但情势迫使我们宁愿将讨论限制在谦逊所设定的界限内，而不愿随着我们有限的言说能力加以扩展。因为我所谈论的事，即使那些可能经历过的人——即我们的原祖父母——也未曾体验（因为罪及其应得的逐出乐园，赶在他们无欲生育之前），如今谈论性交时，人们想到的不是原祖父母那种可能的意志的平静顺从，而是他们自己所经历的情欲的猛烈发作。因此，谦逊封住了我的口，尽管我的头脑清楚地理解这事。但是全能的天主，万有的至高无上、至善的创造者，祂帮助并赏报善的意志，抛弃并惩罚恶的意志，同时驾驭二者，祂并不缺乏一个计划，即从被定罪的人类中——不再按功绩区别，因为全群如同从腐败的根源被定罪——而是凭恩宠，选择预定数量的公民来充满祂的城邑，并表明不仅在被救赎者身上，也在未被解救者身上，祂赐予了多少恩宠。因为每个人都承认，自己从邪恶中被救，不是凭应得的，而是凭白白的良善，当他从那些本可公正地分担共同惩罚的人群中被选出，得以无伤地离开时。那么，天主为何不应创造祂预见会犯罪的那些人呢？因为祂能够在他们身上并通过他们，既彰显他们的罪过所应得的，也彰显祂的恩宠所赐予的；并且，在祂创造和治理的手中，甚至恶人的乖张无序也不能扰乱事物的正当秩序。\par

% structure: p0079 source=h2 target=subsection reason=relative-heading-rank; base=h2

\subsection{第二十七章——论犯罪的天使与人，以及他们的邪恶并未扰乱天主上智安排的秩序。}

人之罪与天使之罪，皆不能阻挠\Quote{主的伟大工程，祂的工程成就祂的旨意。}因为祂以祂的眷顾与全能，将各自的命运分赐予各人，不仅能善用善者，也能善用恶者。如此，祂善用了那恶天使——他因自己最初的邪恶意志\OriginalTerm{意志}{volition}而受罚，被注定顽固不化\OriginalTerm{顽固不化}{obduracy}，以致如今不能愿意任何善——那么，天主为何不应让他诱惑那最初被造为正直、即具有善意的人呢？因为人被造的状态是：若他仰望天主求助，人的良善将击败天使的邪恶；但若他因骄傲的自满而离弃他的创造者与维护者天主，他就会被击败。若他的意志因依靠天主的帮助而保持正直，他应受赏报；若因离弃天主而变为邪恶，他应受惩罚。然而，正是这信赖天主之助，若无天主的帮助也无法成就，尽管人有能力因自满而放弃天主恩宠的益处。因为正如我们没有能力在不以食物维持生命的情况下活于此世，但我们可以拒绝这滋养而停止生命，如同自杀者所做的那样；同样，即使是在乐园中，人也没有能力在没有天主帮助的情况下按应有的方式生活，但他有能力邪恶地生活，尽管这样会缩短他的幸福并招致非常公正的惩罚。因此，既然天主并非不知人将会堕落，为什么祂不应允许他被一个仇恨并嫉妒他的天使所诱惑呢？的确，并非祂不知道人会被击败，而是因为祂预见到，藉着那人的后裔，在神圣恩宠的助佑下，这同一个魔鬼本身将被击败，以增进圣人们的光荣。一切皆以如此方式实现，既没有让任何未来事件逃脱天主的预知，也没有让祂的预知强迫任何人犯罪，并且如此在理性受造物——人类与天使——的经验中，显明受造物的私自妄断\OriginalTerm{私自妄断}{presumption}与造物主的护佑之间有多么巨大的差异。因为谁敢相信或说，天主没有能力阻止天使和人类犯罪？但天主更愿意将这能力留在他们手中，以便既显示他们的骄傲能造成何等邪恶，也显示祂的恩宠能成就何等良善。\par

% structure: p0081 source=h2 target=subsection reason=relative-heading-rank; base=h2

\subsection{第28章——论两座城，属地之城与属天之城的性质。}

因此，两座城由两种爱所形成：属地之城由自爱而形成，甚至达到藐视天主；属天之城由爱天主而形成，甚至达到藐视自己。简言之，前者以自身为荣，后者以主为荣。因为前者从人寻求光荣；但后者的最大光荣是天主，良心的见证。前者在自己的光荣中昂首；后者向它的天主说：\Quote{祢是我的光荣，是我抬起头来的天主。}在前者中，其管辖的君王与民族受统治欲的支配；在后者中，君王与臣民在爱中彼此服侍，臣民服从，而君王关怀众人。前者以自己的力量为乐，这力量体现在统治者的身上；后者向它的天主说：\Quote{主，我的力量，我要爱慕祢。}因此，那座城中的智者，依人而生活，追求自身身体或灵魂或两者的利益；而那些认识天主的人，\Quote{没有当作天主光荣祂，也不感谢祂，反而在思念中成为虚妄，他们愚昧的心就昏暗了；他们自称聪明，}——即夸耀自己的智慧，被骄傲所占据——\Quote{成了愚拙，将不朽坏天主的荣耀变为偶像，仿佛必朽坏的人、飞禽、走兽、昆虫的样式。}因为他们或是领导或跟随人民崇拜偶像，\Quote{敬拜事奉受造之物，而不敬拜事奉造物主，祂是当受赞美的，直到永远。}\BibleRef{罗 1:21}但在另一座城中，没有人的智慧，只有虔诚，向真实的天主献上应有的敬拜，并在圣人的团体——包括圣天使和圣人——中指望奖赏，\Quote{使天主成为万物之中的万有。}\BibleRef{格前 15:28}\par

\SourceRecord{Translated by Marcus Dods.；Nicene and Post-Nicene Fathers, First Series；Vol. 2.；Edited by Philip Schaff.；Buffalo, NY: Christian Literature Publishing Co.,；1887.；<http://www.newadvent.org/fathers/120114.htm>.}

% structure: p0001 source=h1 target=section reason=page-title

\section{天主之城（第十五卷）}

在前四卷中论述了两个城市（地上的和天上的）的起源之后，奥斯定在接下来的四卷中解释它们的成长与进步；为此，他解释了与此主题相关的神圣历史中的主要段落。在第十五卷中，他通过解释\BibleBook{}{创世记}中从加音和亚伯尔时代到洪水的事件，开启了他这部分的著作。\par

% structure: p0003 source=h2 target=subsection reason=relative-heading-rank; base=h2

\subsection{第一章　论人类自始至终分为两系}

关于乐园的幸福、乐园本身、我们原祖在那里的生活、以及他们的罪与罚，许多人思考很多，谈论很多，写作很多。我们自己在之前的卷中也谈论了这些事情，并且写下了我们或是在圣经中读到、或是能合理推断出的内容。如果我们要更详细地探究这些问题，数不尽的无穷疑问将会出现，这将使我们陷入比当前时机所允许的更庞大的工作。我们不能指望有空回答每一个由无所事事和挑剔之人提出的问题，他们总是更愿意提问而不是能够理解答案。然而，我相信我们已经对这些关于世界、灵魂或人类本身起源的伟大而困难的问题做了适当的处理。我们把这个种族分为两部分：一部分是依照人而生活的人，另一部分是依照天主而生活的人。我们也奥秘地称它们为两个城市，或两种人类团体，其中一个被预定为与天主永远为王，另一个则与魔鬼一起承受永远的惩罚。然而，这是它们的结局，我们将在以后谈论。目前，既然我们已经充分讲述了它们的起源——无论是在我们不知数量的天使中，还是在最初的两个人类中——现在似乎适合尝试叙述它们的历程，从我们的两个原祖开始繁衍种族直到所有人类生育停止之时。因为这整个时间或世界时代，在其中死去的人让位，出生的人接续，正是我们讨论的这两个城市的历程。\par

在这人类的两个原祖中，加音是长子，他属于地上之城；在他之后出生了亚伯尔，他属于天主之城。因为在个人身上，可以辨明宗徒这句话的真理：\BibleQuote{那不是属神的在先，而是属生灵的，以后才是属神的，}\BibleRef{格前 15:46}由此产生的情况是：每个人出自被定罪的根源，首先生为亚当的恶的、属血肉的，只在后来通过重生被接枝到基督内，才成为善的、属神的。整个人类也是如此。当这两个城市通过一连串的死亡与出生开始其历程时，这世界之城的公民首先出生，随后是这世界中的旅客，即天主的城之公民，由恩宠预选、由恩宠拣选、由恩宠成为下界的旅客，并由恩宠成为上界的公民。由恩宠——因为就他自身而论，他是出自同一团材料，这材料全部在其起源中被定罪；但天主如同陶工（宗徒巧妙地、且非无深思地引入这个比较），\BibleQuote{从同一个泥团中制造一个器皿作贵重的，一个作卑贱的。}\BibleRef{罗 9:21}但首先制造的是卑贱的器皿，然后才是贵重的器皿。因为正像我已经说过的，在每一个人身上，首先有的是那被弃绝的，我们必须从它开始，但不必永远停留在其中；之后才是那被赞许的，我们可以通过进步达致，并且一旦达到，就可以在其中停留。确实，并非每个恶人都将成为善人，但没有人能不先成为恶人而成为善人；人越早成为善人，就越快获得这个称号，并在新名字中废弃旧名字。因此，记载上说加音建造了一座城，\BibleRef{创 4:17}但亚伯尔身为旅客，没有建造城。因为圣人的城在天上，虽然它在此下界产生公民，并在他们中旅居，直到它统治的时刻到来；那时要在复活之日聚集所有公民，然后将应许的王国赐给他们，使他们与他们的元首、万世之王一同为王，永无穷尽。\par

% structure: p0006 source=h2 target=subsection reason=relative-heading-rank; base=h2

\subsection{第二章　论血肉的儿女与恩许的儿女}

的确，在世上，只要仍有需要，便有这座城的一个预象和黑影，其目的是提醒人这座城将要来临，而非使之成为现实；这个预象本身也被称为圣城，作为未来之城的预象，但它本身并不是实体。关于这座作为预象的城，以及它所预表的自由之城，保禄在致迦拉达人书中这样写道：\Quote{你们愿意属于法律的，请告诉我：你们没有听过法律吗？法律上记载：\Person{亚巴郎}有两个儿子：一个生于婢女，一个生于自由妇。那生于婢女的，是按常例而生的；但那生于自由妇的，却是藉着恩许而生的。这些都含有\OriginalTerm{预象}{allegory}：这两个妇女代表两个盟约；一个是出于西乃山，生子为\OriginalTerm{奴役}{bondage}，那就是\Person{哈加尔}。这\Person{哈加尔}是指西乃山，在\Place{阿剌伯}，她对应着现今的\Place{耶路撒冷}，因为\Place{耶路撒冷}与她的子女同为奴隶。然而天上的\Place{耶路撒冷}是自由的，她是我们的母亲。因为经上记载：\InnerQuote{不生育的，喜乐罢！不生产的，欢呼罢！因为被弃者的子女比有夫者的子女还多。}弟兄们！你们像\Person{依撒格}一样，是恩许的子女。但是，当时那按常例而生的，迫害了那按神恩而生的，现今也是这样。尽管如此，经上说了什么？\InnerQuote{你将婢女和她的儿子赶出去，因为婢女的儿子不能与自由妇的儿子一同继承产业。}我们弟兄们，不是婢女的子女，而是自由妇的子女，藉着基督使我们自由。} \BibleRef{迦 4:21} 这段经文，借着宗徒权威传授给我们，指示我们应当如何理解旧约和新约两盟约的圣经。地上之城的一部分成了天上之城的预象，它并非自有意义，而是指向另一座城，因此它服役，或曰\Quote{处于奴役之中}。因为它被建立，并非为自身，而是为了预表另一座城；这座预象之城本身也被前一个预象所预表。因为\Person{撒辣}的婢女\Person{哈加尔}和她的儿子，正是这预象的预象。当全光来临，黑影必当消逝时，\Person{撒辣}，这位自由之妇（她预表自由之城，而那座预象之城\Place{耶路撒冷}也以另一种方式预表了它），便说：\Quote{你将婢女和她的儿子赶出去，因为婢女的儿子不能与我的儿子\Person{依撒格}一同继承产业}，或者如宗徒所说：\Quote{与自由妇的儿子。}于是，在地上之城中，我们发现两件事——它自身明显的存在，以及它对天上之城的象征性呈现。地上之城的公民是由被罪恶败坏的性生出的，而天上之城的公民则是由恩宠从罪恶中解放的性生出的；因此前者被称为\Quote{烈怒之器}，后者被称为\Quote{怜悯之器}。\BibleRef{罗 9:22} 这在\Person{亚巴郎}的两个儿子中得到了预表：\Person{依市玛耳}（\Person{哈加尔}婢女之子）是按肉身生的，而\Person{依撒格}则由自由妇\Person{撒辣}按恩许而生。两者固然都是\Person{亚巴郎}的后裔，但一个是由自然律而生，另一个是由恩许的恩赐而来。在一个生育中，人的行动得以显明；在另一个中，天主的美善得以彰显。\par

% structure: p0008 source=h2 target=subsection reason=relative-heading-rank; base=h2

\subsection{第三章　撒辣的荒胎因天主的恩宠而成为多产。}

事实上，\Person{撒辣}是荒胎的；她绝望于生育，并决心至少通过她的婢女获得她自知无法亲身获得的祝福，于是她将婢女给了丈夫，因为她自己已不能为他生孩子。她要求丈夫履行婚姻责任，在别人的子宫中行使她自己的权利。因此，\Person{依市玛耳}是按人类生育的普通律法，通过两性结合而生的。所以，经文说他\Quote{按肉身}而生，——并非因为这样的生育不是天主的恩赐，或不是天主的化工（其创造性的智慧，正如经上所载，\Quote{施展威力，从地极直达地极，从容治理万物}，\BibleRef{智 8:1}），而是因为，在要彰显天主的恩赐（并非人应得，而是恩宠的慷慨赠予）的情况下，必须以任何自然努力都无法企及的方式赐予一个儿子。大自然拒绝将孩子赐予\Person{亚巴郎}和\Person{撒辣}现在已达到的年龄；此外，\Person{撒辣}即使在壮年时也是荒胎的。这种无法期待后裔的自然状况，象征着被罪恶败坏并因此受公义定罪的人类性，它不配享有未来的福乐。因此，恩许之子\Person{依撒格}恰当地预表了恩宠的子女，即自由之城的公民，他们在永远的和平中共同生活，其中没有自私和自傲的位置，而是有服务的爱，以众人的共同喜乐为喜乐，使众多心灵合而为一，也就是说，达到完全的和谐。\par

% structure: p0010 source=h2 target=subsection reason=relative-heading-rank; base=h2

\subsection{第四章　论地上之城的争斗与和平。}

但地上之城既非永存（因为当它被处以极刑时，将不再是一座城），它在此世有它的福乐，并以此类事物所能提供的喜乐而喜乐。然而，这福乐并不能解除其追随者的一切困苦，因此这城常因诉讼、战争、争吵以及那些导致死亡或短暂的胜利而自相分裂。因为它的每一部分武装起来反对另一部分，试图在屈服于恶习的同时战胜列国。如果它征服后骄傲自大，则其胜利是致死的；如果它转而思考我们必死命运的常见变故，并对可能降临的灾难忧心忡忡，而非因已取得的成功而得意，则这种胜利虽较高级，仍只是短暂的；因为它不能永久地统治那些它所战胜的臣服者。但这城所渴望的事物不能公正地说是邪恶的，因为它本身，在其种类中，优于一切其他人类之善。因为它渴望地上的和平，为享受地上的福乐，并为获得这种和平而发动战争；如果它征服了，再无人抵抗，它就享有曾经因有对立党派争夺那些不足以同时满足双方的事物时所没有的和平。这种和平是通过艰苦的战争购得的；是通过他们称之为光荣的胜利而获得的。当胜利属于有更正义理由的一方时，谁不乐于祝贺胜利者，并称之为可取的和平呢？所以，这些事物都是善的，无疑都是天主的恩赐。但如果他们忽略了天上之城的更好事物——这些由永恒的胜利和永无止境的和平所确保——并如此过分贪求这些眼前的善，以至于相信它们是唯一可欲的事物，或者爱它们胜过那些被认为是更好的事物——如果是这样，那么苦难必然随之而来，且日益增加。\par

% structure: p0012 source=h2 target=subsection reason=relative-heading-rank; base=h2

\subsection{第五章。——论地上之城的建立者的杀害兄弟行为，以及罗马建立者的相应罪行。}

因此，地上之城的建立者是一个杀害兄弟者。他被嫉妒所胜，杀死了自己的兄弟，一位永恒之城的公民，地上的旅居者。所以我们不应对此感到惊讶：这第一个罪恶的范例，或如希腊人所说的\OriginalTerm{原型}{archetype}，在很久以后，在那注定要统治如此多民族并成为我们所说的地上之城之首的城市建立时，找到了相应的罪行。因为那座城市，如其一位诗人所提及的，\Quote{最初的城墙沾染了兄弟的鲜血}，或如罗马史所载，\Person{雷穆斯}被其兄弟\Person{罗慕路斯}所杀。因此，这座城的建立与地上之城的建立并无区别，除非认为\Person{罗慕路斯}和\Person{雷穆斯}都是地上之城的公民。二人都渴望获得建立罗马共和国的荣耀，但若只有一人独享，则两人都不能获得同样多的荣耀；因为想要统治荣耀的人，如果他的权力与活着的同僚分享，他统治的荣耀当然会更少。因此，为了使整个荣耀由一人独享，他的同僚被除去；通过这一罪行，帝国确实变得更大，但更劣；否则它会较小，但更好。然而，这两兄弟，\Person{加音}和\Person{亚伯尔}，并非被同样的地上欲望所驱动，凶手也不是因害怕两人共治会削弱自己的统治而嫉妒对方——因为\Person{亚伯尔}并不关心在他兄弟所建的城中统治——他是被那种恶人看待善人的魔鬼般的嫉妒仇恨所驱动，没有别的理由，只为他们是善而自己为恶。因为善的拥有绝不会因与一个永久的或临时的伙伴分享而减少；相反，善的拥有会随着分享者的和谐与爱德而成比例地增加。简言之，不愿分享此拥有者不能拥有它；而最愿意让别人分享的人，自己将拥有最大的丰裕。因此，\Person{罗慕路斯}与\Person{雷穆斯}之间的争斗显示了地上之城如何自相分裂；而\Person{加音}与\Person{亚伯尔}之间发生的事则表明了天主之城与人之城之间存有的仇恨。恶人与恶人相争；善人也与恶人相争。但善人与善人，或至少完全善良的人，不能相争；不过，在趋向完美的过程中，他们也会相争到这种程度：每个善人在他抵抗自己的那些事上抵抗他人。在每个个体中，\BibleQuote{肉情和圣神相争，圣神和肉情相争} \BibleRef{迦 5:17}。因此，这种属神的贪求可以与另一个人的肉情贪求相争；或者肉情贪求可以与另一个人的属神渴望相争，如同善人与恶人相争；或者更确切地说，两个善良但尚未完美的人的肉情贪求相互争斗，正如恶人与恶人相争，直到那些在恩宠治疗之下的人获得最终的胜利。\par

% structure: p0014 source=h2 target=subsection reason=relative-heading-rank; base=h2

\subsection{第六章。——论天主之城的公民在此世旅途中因惩罚罪恶而承受的软弱，以及他们如何因天主的眷顾而得痊愈。}

这种病患——即我们在第十四卷中所说的那悖逆——是原初悖逆的惩罚。因此它不是本性，而是恶习；所以对那些在恩宠中成长、凭信德在此朝圣之旅中生活的善人说：\Quote{你们要彼此担代负担，这样就满全了基督的法律。}\BibleRef{迦 6:2} 同样，别处也说：\Quote{你们要劝戒不守规矩的人，安慰胆怯的人，扶持软弱的人，对众人要忍耐。要小心，不要以恶报恶。}\BibleRef{得前 5:14} 又在另一处说：\Quote{如果一个人陷于过犯，你们属神的人，要用柔和的心将他挽回；要自己小心，免得也被诱惑。}\BibleRef{迦 6:1} 又在别处说：\Quote{不可让太阳在你们的愤怒中下山。}\BibleRef{弗 4:26} 又在福音中说：\Quote{如果你的弟兄得罪了你，去，要在你和他独处的时候，指出他的错。}\BibleRef{玛 18:15} 同样，宗徒论到可能引人跌倒的罪说：\Quote{犯罪的人，要在众人面前斥责，使其余的人也害怕。}\BibleRef{弟前 5:20} 为此，为要持守那平安，没有它，没有人能看见主，\BibleRef{希 12:14} 许多规诫被赐下，它们仔细地教导彼此宽恕；其中我们可以数算那句可怕的话，仆人要偿还先前被豁免的一万塔冷通的债务，因为他没有宽恕他同仆所欠的二百银钱。主耶稣为这比喻加上话说：\Quote{如果你们不从心里宽恕自己的弟兄，我的天父也要这样对待你们。}\BibleRef{玛 18:35} 就这样，天主之城的公民在地上旅居、渴慕天上家乡的平安时，他们得到了医治。圣神也在内工作，使得外部应用的药物能产生某种好结果。否则，即使天主亲自使用服从于祂的受造物，并以某种人形向我们的人性感官说话，无论我们在睡眠中或在某种外在显现中接受这些印象，如果祂不通过自己内在的恩宠影响并作用于心灵，任何真理的宣讲都无济于事。但天主凭藉自己极其隐密却又极其公义的上智安排，区分忿怒的器皿和仁慈的器皿，祂就是这样做的。当祂亲自以自己隐藏而奇妙的方式帮助灵魂，那住在我们肢体中的罪（如宗徒所教导的，它更是罪的惩罚）不再在我们可死的身体中作王，使我们顺从它的私欲，当我们不再将我们的肢体献为不义的武器时，\BibleRef{罗 6:12} 那时，灵魂就从自己邪恶自私的欲望中归向，并且天主占有它，它就在平安中占有自己，甚至在此生，然后，带着完美的健康并被赋予不朽，将在永久的平安中无罪地为王。\par

% structure: p0016 source=h2 target=subsection reason=relative-heading-rank; base=h2

\subsection{第七章——论加音罪恶的原因及其执拗，连天主圣言也无法制伏。}

不过天主虽然使用了我们一直在努力解释的这种说话方式，并以那种祂惯于迁就我们的原祖父母并与他们如同朋友交谈的形式对\Person{加音}说话，但对\Person{加音}有什么好的影响呢？难道他听到天主的警告后没有实现杀死兄弟的邪恶意图吗？因为天主区分了他们的祭献，忽视了\Person{加音}的，而看重了\Person{亚伯尔}的——这无疑是通过某种可见的迹象表明的；天主这样做是因为一人的行为是邪恶的，而他兄弟的行为是正义的；\Person{加音}便十分愤怒，脸色也沉了下来。因为经上这样记载说：\BibleQuote{上主对加音说：你为什么愤怒，为什么脸色沉了下来？如果你献得正当，却分得不正当，你不是犯了罪吗？不要烦恼，他要依恋你，你要管辖他。}\BibleRef{创 4:6} 在这项由天主对\Person{加音}发出的警告中，那一句\BibleQuote{如果你献得正当，却分得不正当，你不是犯了罪吗？}确实晦涩，因为不清楚它出于什么理由或目的而被说，许多学者提出了多种含义，各人按信德准则试图解释它。事实上，祭献要\BibleQuote{献得正当}，是指向真天主献上，因为我们只能向祂献祭。而\BibleQuote{分得不正当}，是指我们没有正确区分献祭的地点、时间、材料，或献祭的人，或祭品所献给的对象，或献祭后分食的人。\Emph{分}在这里指辨别——无论是祭品献在不该献的地方，或材料不应在那里献而应献在别处；或在错误的时间献祭，或材料不适合当时而适合别时；或献上在任何地方、任何时间都不该献的祭品；或人将同类中更好的留给自己而献给天主较次的；或他自己或任何无权领受的人亵渎地吃了祭品。在这些情况中，\Person{加音}在哪一点上得罪了天主，很难确定。但若望宗徒在谈到这两兄弟时说：\BibleQuote{不要像加音那样，他是属于那恶者，杀了他的兄弟。为什么杀了他呢？因为自己的行为是邪恶的，而他兄弟的行为是正义的。}\BibleRef{若一 3:12} 他以此让我们明白，天主没有悦纳他的祭献，是因为他没有\BibleQuote{分得正当}，在于他把一些自己的东西给了天主，却把自己保留给自己。因为所有不随从天主的旨意、而随从自己的私意，不是以正直的心、而是以歪曲的心生活的人，都这样做；他们献给天主这样的礼物，以为能藉此获得天主的帮助，不是藉着医治，而是藉着满足他们的邪恶欲望。这是世俗之城的特征：它敬拜天主或众神，希望它们帮助自己在世上得胜、和平地统治，不是出于爱德行善，而是出于统治的贪欲。善人利用世界，以便享受天主；恶人则相反，为了享受世界而想利用天主——至少那些已经相信天主存在并眷顾人事的恶人是这样。因为连这种信仰都没有达到的人，还处于更低的层次。所以，\Person{加音}看到天主悦纳了他兄弟的祭献，却没有悦纳自己的，本应谦卑地选择他的好兄弟为榜样，而不应骄傲地视他为对手。但他却愤怒了，脸色沉了下来。这种对他人——甚至是他兄弟——的善行生气后悔的举动，被天主视为大罪。祂在质问中指责他：\Quote{你为什么愤怒，为什么脸色沉了下来？}因为天主看出他嫉妒自己的兄弟，并为此指责他。因为对人来说，人隐藏的心意是可疑且不确定的，那忧愁究竟是哀悼自己的邪恶，即他已知的冒犯天主之事，还是哀悼他兄弟的善行，后者蒙天主悦纳，使他的祭献获得了天主的恩宠？但天主在解释为何拒绝\Person{加音}的祭献、以及\Person{加音}本应自我懊恼而非憎恨兄弟的原因时，向他显示：尽管他因\BibleQuote{不分得正当}，即生活不正，不配其祭献被悦纳，是不义的，但他无故憎恨他正义的兄弟，则更为不义。\par

然而，祂并未不给他忠告就打发他走，这忠告是神圣、正义而良善的。\BibleQuote{不要烦恼，}祂说，\BibleQuote{因为罪要归于你，你要管辖它。}是指他的兄弟吗？当然不是。那么是指什么呢？是指罪。因为祂曾说过：\BibleQuote{你犯了罪，}然后祂补充说：\BibleQuote{不要烦恼，因为罪要归于你，你要管辖它。}罪的\OriginalTerm{转向}{turning}人，可以理解为：罪咎不能归咎于别人，只能归咎于自己。这是补赎的良药，是求得宽恕的恰当祈求；因此，当它说\BibleQuote{罪要归于你}时，我们不应加上\Quote{将要}，而应读作\BibleQuote{愿罪归于你}，当作命令而非预言。当人不偏爱罪、不辩护罪，而是以悔改降服它时，他就能管辖罪；否则，若成为罪的保护者，必成罪的囚徒。但若我们理解这罪是指那\OriginalTerm{肉欲}{carnal concupiscence}，即宗徒所说的\BibleQuote{肉性欲相反灵性}\BibleRef{迦 5:17}，在这情欲的果实中他提到嫉妒，加音无疑是被此刺伤并激动要杀害他的兄弟，那么我们可以正当地加上\Quote{将要}一词，读作\BibleQuote{罪要归于你，你要管辖它}。因为当那肉性部分——宗徒称之为罪，他在那处说\BibleQuote{不是我做的，而是住在我内的罪}\BibleRef{罗 7:17}——这部分哲学家也称之为恶习，它不应引导心神，而心神应以理性管辖和约束它，防止非法冲动——当这部分被激动去行任何邪恶时，若被抑制且听从宗徒的话\BibleQuote{不要将你们的肢体献给罪作不义的武器}\BibleRef{罗 6:13}，它就转向心神，被心神征服和战胜，如此理性就如主人般统辖它。这正是天主对那被嫉妒之火点燃而欲除灭自己兄弟的人的诫命，他本应以兄弟为榜样。祂说：不要烦恼，或镇定下来，管住你的手不犯罪；不要让罪在你必死的身体上为王，顺从它的私欲，也不要将你们的肢体献给罪作不义的武器。\BibleQuote{因为罪要归于你，}只要你不用放纵来助长它，而是用扑灭其火焰来抑制它。\BibleQuote{你要管辖它；}因为当它不被允许任何外在行为时，它就顺服于统治心神和正义意志的管束，甚至内在的冲动也止息了。在同一部圣书中，对女人也有类似的说法，当天主在他们犯罪后审问并判决他们时，对全体宣告了判决——魔鬼以蛇的形象，女人和她的丈夫以自身。因为当祂对她说：\BibleQuote{我要多多增加你的忧愁和怀胎；你必在忧愁中生子，}然后祂加上：\BibleQuote{你要归于你的丈夫，他要管辖你。}\BibleRef{创 3:16}论及加音之罪（或他肉体的邪恶贪欲）的话，在此针对犯了罪的女人说了；我们应理解丈夫管辖妻子，如同灵魂管辖肉体。因此，宗徒说：\BibleQuote{那爱自己妻子的，就是爱自己，因为从来没有人恨恶自己的肉身}\BibleRef{弗 5:28}。这肉身是要被医治的，因为它是属于我们自己；不可像对待异于我们本性的东西一样弃之不顾。但加音以不愿改过的心接受了天主的劝诫。实际上，嫉妒的恶习在他内生长；他诱骗了他的兄弟并杀了他。这就是地上之城的创建者。他也是杀害基督（人群之牧者）的犹太人的预像，基督由牧羊的亚伯尔所预表；但这是寓意和预言的事，我现在不加以解释；此外，我记得在\WorkTitle{驳斥摩尼教徒富斯图斯}一文中我已有所评论。\par

% structure: p0019 source=h2 target=subsection reason=relative-heading-rank; base=h2

\subsection{加音在人类历史早期建城的理由是什么？}

目前我所要辩护的是这段历史，免得人以为圣经所记载的不可信，因为它说有一个时期地上似乎只有四个人，或者更确切地说只有三个人（在兄弟杀死另一兄弟之后）——即第一个人（万民之父）、加音本人、和他的儿子哈诺客，城市就是以哈诺客的名字命名的。但那些被这种想法所困扰的人忘记了，圣史的作者不一定提到那时可能活着的人，只提到他作品范围要求他点名的人。那位作者（在这件事上是圣神的工具）的计划是，通过从一人传下来的确定世系下到亚巴郎，然后从亚巴郎的后裔转到天主子民身上；在这子民中，他们与外邦人分离，预表和预言了一切与那统治永恒的城市及其君王和奠基者基督有关的事，这些事在圣神中已预见到必将到来；然而，这一目的的实现并非完全不提另一个人类社会，即我们所谓的地上之城，而是提及其必要之处，以便通过与对立面的对比来增强天上之城的荣耀。因此，当圣经在提到那些人的年岁时，对其所讲的每个人作结时，以这样的话结束：\Quote{他生了儿女，他的一生如此如此，他便死了}，我们是否要理解，因为它没有提到那些儿女的名字，所以在那些早期一人生命延续的漫长岁月里，就没有很多人生出来，由他们的联合人数不仅一座城市，而且几座城市都可能被建造呢？但天主的上智安排（由祂的默感而写成这些历史）乐于从一开始就安排并区分这两个社会的各代世系——也就是说，一方面是按人生活的人（即按人而活的人）的世系，另一方面是按天主生活的天主子女（即按天主而活的人）的世系，两者可以一起而又分开地被追溯，直到洪水时代；在那时，它们的分离和联合被展示出来：它们的分离，在于两系世系分别记载在单独的谱表中，一系从杀人者加音传下，另一系从舍特（舍特是代替被兄弟所杀的亚伯尔而生于亚当的）传下；它们的联合，在于善人如此堕落，以致整个种族变得如此败坏，以致被洪水冲走，除了一个义人诺厄（及其妻子、三个儿子和三个儿媳）之外，这八个人被认为唯独配得从那场毁灭全人类的浩劫中得救。\par

因此，虽然记载着：\Quote{加音认识了自己的妻子，她怀孕生了哈诺客，他建造了一座城，按他儿子的名字给那城起名叫哈诺客}\BibleRef{创 4:17}，但不可因此便相信这是他的长子；因为我们不能以为这是由\Quote{他认识了自己的妻子}这句话所证明的，好像那时他第一次与她同房。因为在万民之父亚当的情形中，这句话不仅用于加音（他似乎是他头胎所生）的受孕，而且之后同一部圣经说：\Quote{亚当认识了自己的妻子厄娃，她怀孕生了一个儿子，给他起名叫舍特}\BibleRef{创 4:25}。由此可见，圣经使用这句话既非总在记载生育时，也非仅在提到头胎生育时。也不必因哈诺客以他的名字命名城市而以为他是加音的长子。因为虽然他有其他儿子，但父亲因某种原因更爱他，这是完全可能的。犹大不是长子，但犹太和犹太人却以他的名字命名。但即使哈诺客是城市创始者的长子，也没有理由认为父亲在他一出生就以他的名字命名城市；因为那时他只是一个孤独的人，不可能建立城邦——城邦无非是由某种联系纽带结合在一起的人群。但当他的家庭繁衍到相当多的人口时，他就有可能建造一座城市，并在建成后以儿子的名字命名。因为那些洪水前的人寿命如此之长，以至于圣经中所记年岁的人中活得最短的也达到了753岁。虽然没有人活到一千岁，但有几个超过了九百岁。那么谁能怀疑，在一个人的一生中，人类可能繁衍到足以建造和居住不只一座、而是几座城市的人口呢？这一点很容易从以下事实推测：从一个人亚巴郎起，在不过四百多年内，希伯来民族的人数如此增长，以致该民族出离埃及时记载有六十万能作战的男子\BibleRef{出 12:37}，此外还有厄东人——他们虽不算以色列的后裔，却出自亚巴郎的兄弟（也是亚巴郎的孙子）——以及出于亚巴郎同一血统的其他民族（虽然不是通过撒辣），即他通过哈加尔和刻突辣所生的后裔：依市玛耳人、米德杨人等。\par

% structure: p0022 source=h2 target=subsection reason=relative-heading-rank; base=h2

\subsection{第九章——论洪水前之人的长寿与更伟岸的身材}

因此，凡慎重权衡事实者，无不相信\Person{加音}可能建造了一座城，而且是一座大城，尤其是考虑到当时人的寿命何等漫长；除非或许有怀疑论者对圣经作者所记述的洪水前之人的高龄提出异议，否认其可信。同样，他们也不相信那时人的体格比现在更大，尽管他们自己最受推崇的诗人\Person{维吉尔}也持此说，当他提及那块曾被用作地标的大石时，说那古代的一位壮士在战斗中将其抓起，奔跑，投掷，抛了出去——\par

\Quote{后世十二壮士，\newline{}难将此重负于肩。}\par

从而表明他（\Person{维吉尔}）认为当时的大地产生了更强壮的人。若在后世尚且如此，那么在举世闻名的大洪水之前的时代，岂不更甚？原始人体之高大，常因坟墓的暴露而向不信者证明，或因岁月磨损、或因洪流冲击、或因某种意外，其中发现或滚落出难以置信的巨骨。我本人曾与其他人在\Place{乌提卡}海岸看到一颗人的臼齿，其大小若切磨成我们这样的牙齿，我想足以做成一百颗。但我想那属于某个巨人。因为虽然当时普通人的身体比我们高大，但巨人更在身高上超越一切。而且无论是我们这时代还是任何其他时代，都并非完全没有巨人身高的实例，尽管可能很少。\Person{小普林尼}，一位博学之士，主张世界越老，人的身体就越小。他还提到\Person{荷马}在其诗篇中常哀叹同样的退化；对此他并没有当作诗意的虚构而嘲笑，而是以记录自然奇观的身份，将其作为历史事实予以接受。但是，正如我所说，间或发现的骨骼证明了古人体格之高大，并且也将向后世证明，因为它们不易腐朽。然而，洪水前之人的寿命长度，现在无法通过此类纪念碑般的证据来证明。但我们不能因此就不相信圣经史书，我们若不信其所记载的过去事实，就更不可原谅，因为我们看到其关于未来预言的准确性。就连那位同一\Person{小普林尼}也告诉我们，现今仍有一个民族的人活到二百岁。那么，如果在我们所不知的地方，人们相信有远超我们经验的寿数，为何我们就不相信在离我们遥远的时代也有呢？或者，我们相信别处有此处没有的事物，却不愿相信过去有不同于现今的事物吗？\par

% structure: p0026 source=h2 target=subsection reason=relative-heading-rank; base=h2

\subsection{第十章——论希伯来文抄本与我们抄本对洪水前之人寿命的不同计算}

首先，第一人\Person{亚当}，在他生儿子\Person{舍特}之前，在我们的抄本中记载为活了二百三十年，而在希伯来文抄本中则为一百三十年。但他生了\Person{舍特}之后，我们的抄本说他活了七百年，而希伯来文抄本则为八百年。因此，将两段时期相加，总数相符。在接下来的各代中，父亲生子前的时期在希伯来文抄本中总是少一百年，而生子后的时期在希伯来文抄本中则比我们的抄本多一百年。因此，将两段时期合并，结果相同。第六代则完全没有差异。第七代，即\Person{哈诺客}，因蒙主喜悦而被记载未经死亡即被提升，其差异与首五代相同，我们的抄本在他生子前多给了一百年。但结果仍一致，因为根据两种文献，他在被提升前都活了三百六十五年。第八代的差异较小，且性质不同。\Person{默突舍拉}，即\Person{哈诺客}所生，按照希伯来文读本，在他生继承人之前，并非少一百年，而是多一百年；而在我们的抄本中，这些年份又被加到了他生子后的时期；因此，在此例中总数依然相同。仅在第九代，即\Person{拉默客}（\Person{默突舍拉}之子、\Person{诺厄}之父）的寿命中，总数存在差异，甚至在此例中也十分微小。因为希伯来文抄本记载他比我们的抄本多活了二十四年。在他生儿子（称为\Person{诺厄}）之前，希伯来文抄本比我们的抄本少给六年；但在生此子之后，则多给三十年；因此，扣除前六年，正如我们所说，剩余二十四年的盈余。\par

% structure: p0028 source=h2 target=subsection reason=relative-heading-rank; base=h2

\subsection{第十一章——论默突舍拉的年龄，似乎超出了洪水十四年}

从希伯来文典籍与我们的版本之间的这一差异，产生了一个众所周知的问题，即关于默突舍拉的年龄；因为据计算，他在洪水之后还活了十四年，尽管 \WorkTitle{圣经} 记载当时地上所有生灵中，只有方舟内的八个人逃脱了洪水的毁灭，而默突舍拉并不在其中。因为，根据我们的典籍，默突舍拉在生他的儿子拉默客之前，活了一百六十七年；然后拉默客本人在他的儿子诺厄出生之前，活了一百八十八年，合计三百五十五年。再加上诺厄在洪水发生时的年龄六百岁，从默突舍拉出生到洪水之年总共便是九百五十五年。默突舍拉一生的总年数计算为九百六十九年；因为他活了一百六十七年生下拉默客的儿子后，又活了八百零二年，这样总计，正如我们所说，是九百六十九年。从中减去从默突舍拉出生到洪水时的九百五十五年，还剩十四年，据推测他在洪水之后还活了这些年。因此有些人推测，虽然他没有在地上（地上公认所有不能自然在水中生活的生物都灭亡了），但他有一段时期与已被提升的父亲在一起，并在那里一直住到洪水过去。他们采纳这一假设，是为了不给教会置于崇高权威地位的译本的可信度蒙上阴影，并且因为他们认为犹太人的手稿有误，而非我们的译本有误。因为他们不承认这是译者的错误，而坚持认为原本中存在虚假的陈述，而 \WorkTitle{圣经} 通过希腊文从该原本翻译成了我们的语言。他们说，七十贤士译者同时并一致地产生了一个译文，他们不可能犯错，或者在一个与他们自身利益无关的案例中伪造译文；而是犹太人嫉妒我们翻译了他们的法律和先知书，在他们的文本中作了改动，以削弱我们的权威。这个意见或猜测，各人可按自己的判断采纳。确实，默突舍拉没有活过洪水，而是在洪水发生的那一年去世，如果希伯来文手稿中的数字是真的。关于七十贤士译者的看法，我将在适当的时候，按照这部作品所要求的顺序，下到他们进行翻译的那个时期时，更仔细地陈述，依靠天主的助佑。就当前的问题而言，根据我们的版本，那个时代的人寿命如此之长，以至于在当时仅存的两位始祖的长子有生之年，人类繁衍得足够形成一个社群，这完全是可能的。\par

% structure: p0030 source=h2 target=subsection reason=relative-heading-rank; base=h2

\subsection{第十二章——论那些不相信在太古时期人能活得如此之久的人的意见。}

因为那些认为在那些时代年份的计算方式不同，并且如此之短，以致我们的一年可能相当于他们的十年的人，是绝不可听从的。所以他们说，当我们读到或听到某人活了几百岁时，我们应该理解为九十岁，因为他们的十年才等于我们的一年，而我们的十年等于他们的一百年。因此，按照他们的设想，亚当生舍特时是二十三岁，舍特本人生厄诺士时是二十岁零六个月，尽管 \WorkTitle{圣经} 称这些月份为二百零五年。因为，根据我们正在解释其意见的那些人的假设，他们习惯于将我们的一年分成十部分，并把每一部分称为一年。而每一部分由六天的平方组成；因为天主在六天内完成了祂的工程，以便在第七天安息。关于这一点，我在第十一卷中已尽己所能地讨论过。六的平方，即六乘以六，得三十六天；这乘以十便是三百六十天，或十二个太阴月。至于完成太阳年所需的剩余五天，以及每四年或闰年需加一天的四分之一天，古人加入了这些日子，罗马人称之为\Quote{插闰日}，以凑足年数。所以，舍特的儿子厄诺士生刻南时是十九岁，尽管 \WorkTitle{圣经} 称这些年为一百九十年。同样，在所有记载洪水前人物年龄的世代中，我们发现我们的版本里几乎没有人是在一百岁或以下，甚至一百二十岁左右生子的；而最年轻的父亲记载中都是一百六十岁以上。他们解释说，这是因为没有人能在十岁时生育，而那些人所说的一百岁正是十岁；但十六岁是青春期，此时有生育能力，而他们称之为一百六十岁的正是这个年龄。而且，为了使这些时代的年份计算与我们不同这一点不显得难以置信，他们引用了多位历史学家的记载：埃及人一年有四个月，阿卡纳尼亚人有六个月，拉维尼人有十三个月。小普林尼在提到一些作者记载某人活了一百五十二年、另一人多活了十年、其他人二百、三百，有些人甚至达到五百和六百，少数人活到八百岁时，表示他的看法是这一切都应归因于错误的计算。因为，他说，有些人把每个冬夏算作一年；另一些人把每个季节算作一年，例如阿卡迪亚人，他说他们的年只有三个月。他还补充说，埃及人的小年我们已经说过有四个月，他们有时在每个月的望日结束时结束一年；因此在他们那里就产生了千年的寿命。\par

有些人被这些似是而非的论点说服，并不想削弱这部神圣历史的可信度，反而想通过消除如此难以置信的长寿难题来促进对其的信仰，他们自认明智地说服他人：在那些日子，年岁如此短暂，他们的十年等于我们的一年，而我们的十年等于他们的一百年。但有最清楚的证据表明这完全是错误的。然而，在提出这个证据之前，似乎应该先提一个更为合理的猜测。从希伯来文手稿中，我们可以立即驳斥这种武断的说法；因为据希伯来文记载，\Person{亚当}在生他的第三个儿子之前，活了不是二百三十年，而是一百三十年。那么，如果这相当于我们通常计算的十三年，那么他生第一个儿子时大概只有十二岁左右。按照自然规律，谁能在这种年龄生儿育女呢？但暂且不提他——因为他可能被造后不久就能生育，（他不像我们的婴儿那么弱小，这不可信）——暂且不提他，他的儿子在生\Person{厄诺士}时，据我们的译本，不是二百零五岁，而是一百零五岁，因此按照这种观点，他还不到十一岁。至于他的儿子\Person{刻南}，据我们的译本是一百七十岁，而希伯来文是七十岁生了\Person{玛拉肋耳}，这又怎么说呢？如果那时的七十年只相当于我们的七年，那么七岁的孩子怎么能生儿育女呢？\par

% structure: p0033 source=h2 target=subsection reason=relative-heading-rank; base=h2

\subsection{计算年岁时，应否跟随希伯来文抑或\WorkTitle{七十贤士译本}}

但若我这样说，即刻会有人回答：\Quote{这是犹太人捏造的谎言。}然而，我们以上已处理过此问题，表明像七十贤士这样声名卓著的人不可能窜改他们的译本。不过，若我问他们，两者中哪一个更可信：是那个流散四方的犹太民族能一致同谋伪造这个谎言，从而因嫉妒他人而剥夺了自身圣经的真实性呢？还是那七十位因埃及王\Person{托勒密}召集而聚于一处的犹太人（他们自身也是犹太人），竟会嫉妒外邦人得享同样的真理，并一致同意插入这些错误？谁能看不出哪种解释更为自然、更易于令人相信呢？但任何明智者都绝不会相信，犹太人——无论他们多么恶毒固执——会篡改如此众多且广泛流传的手稿；也不会相信那些著名的七十位贤士会有共同意图去吝惜真理、不让外邦人得知。因此，更合理的说法是：当他们的译作最初从\Person{托勒密}图书馆的抄本被誊写时，某种类似的错误可能渗入了第一份抄本，并由此流传开来；这可能并非出于欺诈，而只是抄写员的笔误。这对于有关\Person{玛突舍拉}寿命的难题，以及另一个相差二十四年总数的案例，是一个足够合理的解释。但在那些有规律性篡改的案例中——即一个译本在生子继承前的年数一致多出百年，而在之后的年数一致少百年，\Emph{\OriginalTerm{反之}{vice versâ}}，以使总数相符——这种情况适用于第一、二、三、四、五及第七代——在这些案例中，错误似乎（如果我们可这样说）具有某种恒常性，不似偶然，而似有意为之。\par

因此，使希伯来文圣经抄本有别于希腊文和拉丁文抄本的那种数字差异——即在接连几个世代中，每个世代都统一在寿数上增加和减除一百年——既不应归咎于犹太人的恶意，也不应归咎于像七十贤士那样勤勉谨慎的人，而应归咎于那个最初获准从上述国王的图书馆抄录手抄本的抄写员的错误。因为即使在现在，当数字无助于更容易理解或更圆满地认识任何事物时，它们都既被粗心地抄录，又被更粗心地校订。因为谁肯费心去了解以色列各个支派各有多少千人呢？他看出这种知识并无益处。或者有多少人意识到这种知识所隐藏的巨大好处呢？但在这一情况下，在如此多的连续世代中，在一份手抄本中增加了一百年，而在另一份手抄本中却没有计算这一百年，然后在儿子和继承人出生后，再将缺少的年份补上，很明显，设计这种安排的抄写员意在暗示：洪水前的人之所以活了过多年数，只是因为每一年都过短；他试图通过陈述他们能够生育后代的青春期年龄来引起人们对这一事实的注意。因为，为了避免怀疑者因如此长久的寿命而绊倒，他暗示他们的一百年仅等于我们的十年；他通过每当发现年龄低于一百六十年左右时便增加一百年，然后又从儿子出生后的时期中减去这些年份，以使总数协调，来表达这一暗示。通过这种方式，他意图将生育后代归因于合适的年龄，而不减少归功于个人一生的总年数。而恰恰是在第六世代，他背离了这一统一做法，这一事实反倒使我们更倾向于相信：当我们所指的情况需要他修改时，他便修改；而既然没有这种情形，他便不做修改。因为在同一世代，他在希伯来文手抄本中发现，\Person{耶勒得}在生\Person{哈诺客}之前活了一百六十二年，按照短年计算法，这相当于十六年零不到两个月，是能够生育的年龄；因此没有必要再增加一百短年，从而将年龄变为通常的二十六年；当然也没有必要在儿子出生后减去他之前并未增加的年份。因此，在这一例中，两份手抄本之间没有差异。\par

这一点还由以下事实进一步证实：在第八世代，希伯来文经卷记载\Person{默突舍拉}在\Person{拉默客}出生前活了一百八十二年，而我们的经卷则记载他少了二十年，尽管通常在这一时期会增加一百年；然后，在\Person{拉默客}出生后，这二十年又得以恢复，以使两书的年数总和相等。因为如果他的意图是要将这一百七十年理解为十七年，以适合青春期，那么既然他无需增加任何东西，也无需减去任何东西；因为在这一情况下，他找到了一个适合生育后代的年龄，而正是为了这个年龄，他习惯在发现年龄不够充分的情况时增加那一百年。这二十年的差异，我们或许本可认为它是偶然发生的，若不是他随后小心地恢复这些年份，正如他之前从该时期中扣除它们一样，以使总数没有不足。或者，我们是否或许要设想，有一个更加巧妙的设计，旨在掩盖故意且统一地在前期增加一百年并在后期扣减这一事实——他是否意图通过做类似的事情来掩盖这一点，即不是增加和扣减一个世纪，而是增加和扣减一些年，甚至是在他无需这样做的情况下？但不论对此作何想法，不论是否相信他这样做了，终究不论是否如此，我绝不怀疑：当经卷中发现任何差异时，既然两者都不能符合事实，我们最好优先相信译者们据以译成另一种语言的那种原文。因为有三份希腊文手抄本、一份拉丁文手抄本和一份叙利亚文手抄本彼此一致，在这些手抄本中，\Person{默突舍拉}据说在洪水前六年去世。\par

% structure: p0037 source=h2 target=subsection reason=relative-heading-rank; base=h2

\subsection{第十四章——论上古时代的年份与我们现在的年份同样长短}

现在让我们来看，如何能清楚地证明，那些人的寿命极其漫长，其中的年份并不短，以至于他们的十年只等于我们的一年，而是和我们的一样长，由太阳的运行来衡量。这可由圣经记载洪水发生在诺厄生命的第六百年这一点来证明。但为什么同一个地方又写道：\BibleQuote{洪水泛滥在地上，是在诺厄生命的第六百年，二月二十七日}，如果那极短的年份（需要十个才等于我们的一年）只有三十六天，那又当如何？因为如此短的年份，如果古人尊称它为年，要么没有月份，要么这个月必须是三天，这样才可能有十二个月。那么，这里如何能说\BibleQuote{第六百年，二月，二十七日}呢？除非当时的月份和现在一样长。否则，怎么能说洪水在二月二十七日开始呢？然后后来，在洪水结束时，这样写道：\BibleQuote{七月二十七日，方舟停在阿辣辣特山上。水又渐消，直到十一月；十一月初一日，山顶都现出来了。}\BibleRef{创 8:4}但如果月份像我们的一样，那么年份也一样。而如果每个月只有三天，就不可能有一个二十七日。或者，如果每个时间度量都按比例缩小，那么三天中的三十分之一被称为一天，那么那场大洪水，记载持续了四十昼夜，实际上在我们不到四天内就结束了。谁能忍受这样的愚蠢和荒谬？愿我们远离这种错误——这种错误试图用虚假的猜测来建立我们对圣经的信心，却要在另一点上拆毁我们的信心。显然，那时的日子就和现在一样，是二十四小时，由昼夜交替界定；那时的月份和现在一样，由一个月的升起和完成来定义；那时的年份也和现在一样，由十二个阴历月完成，再加上五天和四分之一天来与太阳运行调整。正是这样长度的年份被算作诺厄生命的第六百年，在二月二十七日，洪水开始了——据记载，这场洪水是由持续四十天的大雨引起的，这些日子不仅有两个多小时，而是有二十四小时，构成一个昼夜。因此，那些洪水前的人活了九百多年，这些年份和后来亚巴郎活了一百七十五年、他的儿子依撒格一百八十年、他的儿子雅各伯将近一百五十年、后来梅瑟一百二十年、以及现在的人七十或八十年、或者稍长一些的年份一样长，关于这些年份，经上说：\BibleQuote{他们的力量是劳苦和愁烦}。\par

但是，在我们自己的文本与希伯来文本之间发现的那种数字上的差异，并不涉及古人的长寿；如果存在某种差异大到两个版本不能都真实，那么我们必须从我们的译本所依据的那个文本中获取对真实情况的理解。然而，尽管任何愿意的人都有权修改这个译本，但值得注意的是，在\WorkTitle{七十贤士译本}与希伯来文本似乎不一致的许多地方，没有人敢从希伯来文本修正\WorkTitle{七十贤士译本}。因为这种差异没有被视为篡改；就我个人而言，我相信它也不应被视为篡改。但是，当差异不仅仅是抄写员的错误，并且含义符合真理且阐明真理时，我们必须相信是圣神促使他们给出不同的译文，不是作为译者，而是在预言的自由中。因此，我们发现宗徒们公正地认可\WorkTitle{七十贤士译本}，当他们从圣经中引证时，既引用它，也引用希伯来文本。但正如我承诺过的，如果天主帮助，我将在更合适的地方更仔细地处理这个问题，现在我将继续处理手头的事情。因为毫无疑问，人的生命如此漫长，第一个人的长子能够建造一座城——然而，是属地的城，而不是那被称为\WorkTitle{天主之城}的城，描述这城是我们着手进行这项伟大工作的目的。\par

% structure: p0040 source=h2 target=subsection reason=relative-heading-rank; base=h2

\subsection{第十五章——原始时代的人是否可能直到记载他们生育子女的年龄才禁绝性交？}

那么，或许有人会说：是否可以相信一个有意生儿育女、并无意于节欲的人，竟能禁欲一百多年，或者按照希伯来文版本，仅仅是少一些，比如八十年、七十年或六十年？或者，如果他没有禁欲，却未能生育后代？这个问题有两种解答。要么青春期大大推迟，正如整个寿命更长；要么——在我看来更可能的是——这里提到的并非长子，而是那些名字被用来填补直到诺厄的世系所需的人，从诺厄我们又看到传承延续到亚巴郎，以及此后直到那个时间点，必须通过家谱来标志那座最荣耀之城的历程，这城在此世作为异乡旅居，寻求天上的国度。不可否认的是，加音是第一个由男女所生的人。因为如果他不是第一个在未出生的二人之外通过出生增添的人，亚当就不可能说出记载的话：\Quote{我藉主获得了一个人} \BibleRef{创 4:1}。之后是亚伯尔，被长兄杀害，他第一个通过某种预象，向旅居的天主之城预示了这座城将遭受恶人——即那些属土的人——的不义迫害，因为他们爱慕自己的尘世起源，并沉醉于尘世之城的地上幸福。但亚当生这些儿子时多大年纪，并未显示。此后，世系分岔，一支源于加音，另一支源于亚当替代被兄杀害的亚伯尔所生的儿子，他给他起名叫舍特，正如经上记载：\Quote{天主又赐给了我一个儿子，代替亚伯尔，因为加音杀了他} \BibleRef{创 4:25}。因此，这两条世系——加音的和舍特的——代表了这两座城的不同等级：一座是旅居于地上的天上之城，另一座是贪恋世俗之乐并沉溺其中、仿佛那是唯一喜乐的尘世之城。尽管洪水之前登记了八代（包括亚当），但在加音的后代中，没有记录生继承者的年龄。因为天主圣神拒绝在尘世之城的世代中标记洪水前的时间，而宁愿在天上世系中这样做，仿佛它更值得纪念。此外，当舍特出生时，记载了他父亲的年龄；但那时他已经生了其他儿子，谁能断言加音和亚伯尔是他先前所生的仅有的儿子呢？因为亚当只生了他们，并不等于只有他们被命名以延续那些需要提及的世代序列。尽管所有其他名字都被埋没在沉默中，但圣经说亚当生了儿女；谁若谨慎避免冒失，敢断言他的后代有多少呢？很可能亚当在舍特出生后得到天主的启示而说：\Quote{天主又赐给了我一个儿子，代替亚伯尔}，因为这个儿子能够代表亚伯尔的圣洁，而非因为他在时间上首先生于亚伯尔之后。接着，因为经上记载：\Quote{舍特活了二百零五岁}，或按希伯来文读法\Quote{一百零五岁，生了厄诺士} \BibleRef{创 5:6}，除了轻率的人，谁能断言这是他的长子呢？谁会这样引人惊奇，并使我们探究为何这么多年他虽无持久节欲之意却仍免于房事，或者若他没有禁欲却仍无子女呢？谁能这样，当经上论到他记载：\Quote{他生了儿女，舍特共活了九百一十二岁，然后死了} \BibleRef{创 5:8}？同样，对于后来记载了他们的年岁的人，也并不隐瞒他们生了儿女。\par

因此，完全看不出被提名者为儿子的人是否就是长子。不，因为那些先祖要么如此晚才达到青春期，要么不能娶妻，要么不能使妻子受孕，这是难以置信的；因此那些儿子是他们的长子也同样难以置信。但是，由于圣史的作者设计通过一系列代际、以明确的间隔追溯到诺厄的出生和生活（在诺厄时代发生了洪水），他没有提及那些最先出生的儿子，而是那些传承得以传递的儿子。\par

让我通过在此插入一个例子来使这一点更清楚，关于这个例子，没有人能怀疑我所断言的真实性。福音作者玛窦，当他打算通过一系列父母将主的肉身的世代铭记在我们记忆中时，从亚巴郎开始并意图到达达味，他说：\BibleQuote{亚巴郎生依撒格} \EditorialNote{玛窦福音第一章}；为什么他没有说依市玛耳，那是他生的第一个？然后\BibleQuote{依撒格生雅各伯}；为什么他没有说厄撒乌，那是长子？简单地说，因为这些儿子无助于他到达达味。接着是：\BibleQuote{雅各伯生犹大和他的兄弟们}：犹大是长子吗？\BibleQuote{犹大}，他说，\BibleQuote{生培勒兹和则辣黑}；但这双胞胎也不是犹大的长子，在他之前他已经生了另外三个儿子。因此，在世代顺序中，他保留了那些能使他到达达味的人，以便朝着他所设定的目标前进。由此我们可以理解，所提到的洪水前的人不是长子，而是那些能使后代顺序延续到族长诺厄的人。因此，我们不必为探讨他们青春期晚到这一不必要且晦涩的问题而自寻烦恼。\par

% structure: p0044 source=h2 target=subsection reason=relative-heading-rank; base=h2

\subsection{第十六章——论血亲之间的婚姻，关于此事，现代法律不能约束最早时期的人们。}

因此，在由尘土所造的男人和他的妻子（由他的肋骨所造）首次婚姻之后，人类需要男女结合才能繁衍，而且除了从这两人所生的人之外没有其他人，所以男人娶自己的姐妹为妻——这行为在古代确实为必然所迫，后来却被宗教禁令所谴责。因为非常合理而且正当的是，人们之间和谐是荣耀且有益的，应当通过各种关系彼此联结；一个人不应自己承担许多关系，而应将各种关系分散给不同的人，从而将最多的人联结在共同的社会利益中。\Quote{父亲}和\Quote{岳父}是两种关系的名称。因此，当一个人有一个人作他的父亲，另一个人作他的岳父时，友谊就扩展到更大的数量。但\Person{亚当}独自一人必须对他的儿子和女儿同时保持这两种关系，因为兄弟和姐妹通过婚姻结合。同样，他的妻子\Person{厄娃}对她的男女孩子既是母亲又是岳母；而如果有两个女人，一个是母亲，另一个是岳母，那么家庭情感就会有更广阔的领域。然后姐妹自己成为妻子，就独自承担了两个关系，如果这些关系分配给不同的人，一个为姐妹，另一个为妻子，那么家庭纽带就会包含更多的人数。但当时没有实现这一点的材料，因为除了从这两对最初父母所生的兄弟和姐妹之外，没有其他人类。因此，当人口充足使得可能时，人们应当选择不是自己姐妹的女人为妻；因为那时不仅没有必要娶姐妹，而且如果这样做，会是极其可憎的。因为如果第一对夫妇的孙子孙女，现在能够选择他们的表亲为妻，却娶了他们的姐妹，那么一个人就不再只是承担两种关系，而是三种关系，而每一种关系本应由不同的人承担，以便通过家庭情感联结更广的人数。因为一个人在这种情况下对他的孩子（兄弟和姐妹现在成为夫妻）既是父亲，又是岳父，又是伯父；他的妻子对他们既是母亲，又是姑妈，又是岳母；他们自己就不仅是兄弟和姐妹，而且是夫妻，还是表亲，因为是兄弟和姐妹所生的孩子。现在，所有这些关系，将三个男人合并为一，如果每种关系由一个人承担，就会包含九个人，这样一个人有一个姐妹，另一个妻子，另一个表亲，另一个父亲，另一个伯父，另一个岳父，另一个母亲，另一个姑妈，另一个岳母；这样社会纽带就不会紧缩以束缚少数人，而是宽松地包含更多数目的亲戚。\par

我们看到，自从人类增长繁衍以来，即使在崇拜众多假神的外邦人中，这也严格遵循，尽管他们的法律乖谬地允许兄弟娶姐妹，但风俗有更精细的道德，宁愿放弃这种许可；虽然在人类早期时代娶自己的姐妹是完全允许的，但现在却被视为任何情况都无法辩解的可憎之事。因为风俗对于吸引或冲击人类情感有非常强大的力量。在这件事上，当它在适当范围内克制情欲时，忽视和违背它的人被公正地谴责为可憎。因为如果为了贪婪的利益耕越自己的边界是不义的，那么因性欲而逾越公认的道德边界岂不更是不义？至于次一级血缘的婚姻，即表亲之间的婚姻，我们观察到在我们这个时代，习惯的道德已经阻止了这种婚姻频繁发生，尽管法律允许。它没有被神律禁止，也尚未被人律禁止；尽管如此，尽管合法，人们仍回避它，因为它与不合法如此接近，娶表亲几乎就像娶姐妹——因为表亲关系如此接近，他们被称为兄弟和姐妹，几乎就是真的。但古代先祖们担心近亲关系可能在世世代代中逐渐疏远，变为远亲关系，或完全不再是亲属关系，他们虔诚地试图在它变得疏远之前通过婚姻的纽带限制它，从而在它逃离时将其召回。因此，即使当世界充满人口时，虽然他们不选择他们的姐妹或同父异母的姐妹为妻，但他们仍然倾向于从同宗族中选择。但谁怀疑现代对甚至表亲婚姻的禁止是更恰当的规范——不仅因为我们一直主张的理由，即增加关系，这样一个人就不会吸收两个可以分配给两个人的关系，从而增加作为家庭联结的人数，而且因为人性中有一种我不知道是什么自然的和值得称赞的羞怯，它阻止我们渴望那种关系——尽管是为了繁衍，却仍是淫欲的，甚至婚姻的端庄也对此脸红——与任何血缘关系要求我们尊敬的人发生？\par

男女之间的性行为，对凡人来说，可以说是城邦的一种苗床。但地上之城为了人口只需要生育，而天上之城还需要重生，以除去生育的玷污。在洪水之前是否有任何有形或可见的重生标记，如同后来亚巴郎受割损时被吩咐的那样，或者那是一种怎样的标记，神圣历史没有告诉我们。但它确实告诉我们，甚至这些最早的人类也向天主献祭，这在最初的两兄弟身上也显明了。此外，诺厄也被说在洪水之后从方舟出来时向天主献了祭。关于这个主题，我们已在前几卷中说过了：魔鬼们僭取神性，并要求祭献以便被尊为神，它们喜爱这些尊荣，没有别的原因，只是因为它们知道真正的祭献应当归于真天主。\par

% structure: p0048 source=h2 target=subsection reason=relative-heading-rank; base=h2

\subsection{第十七章 论出自同一始祖的两位父亲和领袖。}

既然亚当是两支血脉的父亲——也就是说，他既是属于地上之城的那支血脉的父亲，也是属于天上之城的那支血脉的父亲——当亚伯尔被杀，并藉其死亡显示了一个奇妙的奥迹时，从此就有两支血脉分别出自两位父亲，加音和舍特，并在他们那些应当被记载的儿子们身上，这两座城的标记开始更清晰地显现出来。因为加音生了哈诺客，并以他的名建造了一座城，一座地上之城，这城并非以这个世界为家，而是满足于现世的和平与幸福。加音这个名也意为\Quote{获得}；因此在他出生时，他的父亲或母亲曾说：\Quote{我藉着天主获得了这个人}。哈诺客则意为\Quote{奉献}；因为这地上之城在这世界中被奉献，它正是在这个世界中找到了它所追求和向往的终结。此外，舍特意为\Quote{复活}；他的儿子厄诺士意为\Quote{人}，并非像亚当那样——亚当也意为人，但希伯来语中不分男女地使用，正如经上所载：\BibleQuote{祂造了他们一男一女，且祝福了他们，并称他们的名字为亚当}\BibleRef{创 5:2}。这无可置疑地表明，尽管女人被特别称为厄娃，但亚当这一名字，意为\Quote{人}，为两者所共有。但厄诺士在更严格的意义上意为人，希伯来语学者告诉我们，它不能用于女人：它等同于\Quote{复活之子}，\BibleQuote{那时他们不娶也不嫁}\BibleRef{路 20:35}。因为在重生将我们带至的那个地方，将不再有生育。因此，我认为注意到这一点并非无关紧要：在那些从被称为舍特的那一位传下来的世代中，虽说生了女儿也生了儿子，却没有一个女子被明确地记录名字；但在那些从加音传下来的世代中，在那支血脉的末端，最后被记录为所生的是一个女子。因为我们读到：\BibleQuote{默突沙耳生了拉默客。拉默客娶了两个妻子：一个名叫阿达，另一个名叫齐拉。阿达生了雅巴耳：他是居住在帐棚的牧人的始祖。他的兄弟名叫犹巴耳：他是所有弹琴吹箫者的始祖。齐拉也生了突巴耳加音，他是所有铜铁匠的师傅；突巴耳加音的姐妹是纳阿玛}\BibleRef{创 4:18}。加音的所有世代到此为止，共八代，包括亚当——即从亚当到拉默客共七代，他娶了两个妻子，他的子女（其中也列出了一个女子）构成了第八代。藉此优雅地表明，地上之城直到其终结都将有出于男女交合的肉体世代。因此，加音这一支最后一位被提及的父亲，他的妻子们各自以她们的名字被记录下来——这在洪水之前，除了厄娃之外，没有其他地方遵循过这种做法。现在，既然加音（意为\Quote{获得}），地上之城的建立者，以及他的儿子哈诺客（意为\Quote{奉献}），并以他的名建立此城，表明这座城在其开端和终结都是属地的——在这城中，人们所期望的不过是这个世界中可见的事物——那么，舍特（意为\Quote{复活}），作为与其他世代分开记录的世代之父，我们必须考虑这段神圣历史关于他儿子的记载。\par

% structure: p0050 source=h2 target=subsection reason=relative-heading-rank; base=h2

\subsection{第十八章 亚伯尔、舍特、厄诺士对基督及其身体教会的意义。}

\BibleQuote{给舍特也生了一个儿子，他给他起名叫厄诺士：他开始呼号上主天主的名。}\BibleRef{创 4:26} 这里有一个对真理的响亮的见证。人，这位复活之子，生活在希望中：只要天主之城（这城是因着对复活的信心而生的）在世上作客旅，他就生活在希望中。因为在这两个人身上——\Person{亚伯尔}，意为\Quote{哀伤}，和他的兄弟\Person{舍特}，意为\Quote{复活}——基督的死和祂从死者中的复活被预表了。藉着对这二者的信心，世上生出了天主之城，也就是说，那盼望呼求主名的人。\BibleQuote{因为我们得救，是在乎盼望；只是所见到的盼望不是盼望；有谁还盼望他所见到的呢？但我们若盼望那未看见的，就必须坚忍等待。}\BibleRef{罗 8:24} 谁能不将这联系到一个深奥的奥秘呢？难道\Person{亚伯尔}不也曾盼望呼号上主天主之名吗？他的祭品在圣经中被记载为天主所悦纳。难道\Person{舍特}自己不曾盼望呼号上主天主之名吗？关于他，经上说：\Quote{天主给了我另一个儿子代替\Person{亚伯尔}。}那么，为什么这件所有虔诚人共有的特质，特别归给\Person{厄诺士}呢？除非是因为，在\Person{厄诺士}身上——他被记载为那些被分别出来、归于天上之城的更好部分的世代的父亲的长子——理应有一个人或人群的预表，他们不按人的方式满足于地上的幸福，而是按天主的方式盼望着永恒的幸福。圣经没有说：\Quote{他盼望于上主天主}，也没有说\Quote{他呼号上主天主的名}，而是说\Quote{他盼望呼号上主天主的名}。这\Quote{盼望呼号}是什么意思呢？除非是预言，将来有一个百姓要兴起，按照恩宠的选择，呼号上主天主的名。这正是另一位先知所说的，而宗徒将此解释为属于天主恩宠的百姓：\BibleQuote{凡呼号主名的，必然得救。}\BibleRef{罗 10:13} 因为这两句话：\Quote{他给他起名叫\Person{厄诺士}，意思是人}和\Quote{他盼望呼号上主天主的名}，足以证明人不应该将希望寄托于自身；正如经上另处所写：\BibleQuote{凡信赖世人的人，是可咒骂的。}\BibleRef{耶 17:5} 所以，没有人应当信靠自己，以为他能成为那另一座城的公民，那城不是以\Person{加音}之子的名号为这名号奉献于今世——即在这个必朽世界的流逝过程中——而是奉献于永福的永生。\par

% structure: p0052 source=h2 target=subsection reason=relative-heading-rank; base=h2

\subsection{第十九章——\Person{哈诺客}被提升天的意义}

因为从\Person{舍特}来的那一条谱系，在从\Person{亚当}算起的第七代，也有\Quote{奉献}之称。从\Person{亚当}算起，第七代是\Person{哈诺客}，即\OriginalTerm{奉献}{Dedication}。这就是那位因取悦天主而被提升天的人，在世代序列中占据一个突出的位置，他是\Person{亚当}的第七代，这个数字因安息日的祝圣而显要。但是，从两条谱系的分叉点，即从\Person{舍特}算起，他是第六代。在第六天，天主造了人并完成了祂的工作。\Person{哈诺客}被提升天预表了我们那被延迟的奉献；因为虽然这在我们的元首基督内已经完成，祂复活了，不再死亡，且祂自己也被提升天，但整个房屋——基督自己是其根基——还有另一个奉献，这奉献被延迟到末时，那时所有人都将复活，不再死亡。无论被称为天主的房屋、天主的圣殿，还是天主的城，说到奉献，都是一回事，也符合拉丁语的习惯。因为\Person{维吉尔}本人称那最广阔帝国的城市为\Quote{阿萨拉科斯之家}，意指罗马人，他们通过特洛伊人出自\Person{阿萨拉科斯}。他也称他们为\Quote{埃涅阿斯之家}，因为罗马是由那些在\Person{埃涅阿斯}率领下来到意大利的特洛伊人所建。因为那位诗人模仿了圣经，在圣经中，希伯来民族虽然人数众多，却被称为\Person{雅各伯}之家。\par

% structure: p0054 source=h2 target=subsection reason=relative-heading-rank; base=h2

\subsection{第二十章——为何\Person{加音}的谱系在第八代终结，而\Person{诺厄}虽与\Person{加音}同出一父\Person{亚当}，却为第十代。}

有人会说：如果这历史的作者在列举从\Person{亚当}通过其子\Person{舍特}往下传的世代时，目的是要经由他们下传到\Person{诺厄}（在\Person{诺厄}时代发生了洪水），再从\Person{诺厄}继续追溯相连的世代直到\Person{亚巴郎}（\Person{玛窦}就从\Person{亚巴郎}开始列举基督——天主之城永恒之王的谱系），那么他为何又要列举从\Person{加音}而出的世代？他打算终止于何处？我们的回答是：止于洪水，因为地上之城的所有苗裔被洪水毁灭，但由\Person{诺厄}的儿子们得以重建。因为按肉体生活的地上之城和人类社会，直到今世的终结不会消灭，正如我们的主所说的：\BibleQuote{今世之子娶嫁生育，也被生育。}\BibleRef{路 20:34}但天主之城在此世作客旅，通过重生被引领到来世，来世的儿女既不娶也不嫁。在这世上，生育为两城所共有；尽管即使现在，天主之城已有成千上万的公民戒绝生育行为；然而另一城也有某些公民仿效他们，但却是错误地。因为那些从信仰中偏离并引入各种异端的人也属于那一城；因为他们按人而不按天主生活。还有印度的\OriginalTerm{裸体哲学家}{gymnosophists}，据说他们裸体在印度旷野中从事哲学思考，也是它的公民；他们戒绝婚姻。因为节制并非善事，除非是出于对至善即天主的信德而实行的。然而在洪水之前，没有人被发现实行了节制；因为就连\Person{亚当}的第七代子孙\Person{哈诺客}（据说他没有死亡而被提升天），在被提升天之前也生儿育女，其中就有\Person{默突舍拉}，通过他维持了所记载的世代传承。\par

那么，如果应当追溯\Person{加音}的世代直到洪水，而且并没有青春期延迟到一百岁或更多年才使人有望生育后代的情况，为何所登记的\Person{加音}世代数目如此之少？因为如果本书的作者心中并没有一个固定目标，要像他在那些从\Person{舍特}苗裔出发下传到\Person{诺厄}、再从\Person{诺厄}按严格秩序重新开始的世代中所设计的那样，严格追溯世系，那么他何必为了下传到\Person{拉默客}（即从\Person{亚当}算第八代，从\Person{加音}算第七代）而省略头生子？就好像从此处他想要转向另一条世系，通过它他可以到达以色列民族（地上的耶路撒冷在他们中预示了天上之城的形象），或是到达耶稣基督——\BibleQuote{按肉体说，祂是在万有之上，永远受赞美的天主}\BibleRef{罗 9:5}，天上之城的创造者和统治者。我说，既然\Person{加音}的所有后裔都在洪水中毁灭了，这还有什么必要？由此可见，在这谱系中登记的是头生子。那么，为何他们如此之少？如果青春期与他们的长寿不成比例，并且他们在百岁之前就有了孩子，那么在洪水之前这段时间里，他们的数目必然更大。因为假设他们开始在三十岁时生育孩子，那么包括\Person{亚当}和\Person{拉默客}的子女在内，共有八代，8乘以30得240年；难道在洪水之前剩余的所有时间里，他们没有再生孩子吗？那么，写这记录的人为何对后续世代只字不提？因为根据我们的圣经抄本，从\Person{亚当}到洪水共2262年，而根据希伯来文本是1656年。那么，假设较小的数字是正确的，从1656年中减去240年，难道在洪水之前余下的1400多年里，\Person{加音}的后裔就没有生育孩子吗？\par

但任何被这段话触动的人应当记得，我曾讨论过这样一个问题：那些原始人如何能够多年克制自己不生育子女？当时提出了两种解决办法：要么他们的青春期与长寿成正比地延迟，要么族谱中登记的儿子并非长子，而是作者为了达到其目的才纳入谱系的人，正如他打算通过\Person{舍特}的后代到达\Person{诺厄}一样。因此，如果在\Person{加音}的后代中，没有出现作者可以通过省略长子、插入那些适合此目的的人而抵达的目标，那么我们就必须求助于青春期延迟的假设，并说他们要到一百岁以后才能生育子女，这样谱系顺序便通过长子延续，甚至填满了洪水前那段漫长的时期。然而，也有可能出于某种我无法看透的更隐秘的原因，我们称之为地上之城的这座城，其全部谱系一直展示到\Person{拉默客}和他的儿子们，然后作者便不再记载洪水前可能存在的其余部分。而且，即使不假设这些人的青春期如此之晚，也有可能存在另一种理由，即谱系是通过非长子的儿子来追溯的：也就是说，\Person{加音}所建造并以他儿子\Person{哈诺客}命名的城，可能拥有一个极其广阔的统治区域和许多国王，他们不是同时而是相继统治，在位的国王总是生育他的继承人。\Person{加音}本人将是这些国王中的第一位；他的儿子\Person{哈诺客}，以其名字命名了他统治的那座城，将是第二位；第三位是\Person{哈诺客}所生的\Person{依辣得}；第四位是\Person{依辣得}所生的\Person{默胡雅耳}；第五位是\Person{默胡雅耳}所生的\Person{默突舍拉}；第六位是\Person{默突舍拉}所生的\Person{拉默客}，他是从\Person{亚当}经\Person{加音}而来的第七代。但是，并不一定是长子继承父亲的王位，而是那些因拥有对地上之城有益的某种德行而被推举的人，或因抽签选中的人，或最受父亲喜爱的儿子，会以一种世袭权利继承王位。洪水可能发生在\Person{拉默客}在世和统治期间，并将他与所有其他人一同毁灭，除了那些在方舟中的人。因为，从\Person{亚当}到洪水这么长的时期中，且各人的年龄各不相同，我们不能惊讶于两条谱系的代数并不相等：\Person{加音}系有七代，\Person{舍特}系有十代。正如我已经说过的，\Person{拉默客}是从\Person{亚当}算起的第七代，\Person{诺厄}是第十代；在\Person{拉默客}的谱系中，不像之前那样只登记了一个儿子，而是登记了多个，因为如果他死后与洪水之间还有任何统治的时间，那么他们中哪一个会继位是不确定的。\par

但是，无论\Person{加音}的后代是以怎样的方式向下追溯的（或是通过长子，或是通过王位继承人），在我看来，我绝不能忽略指出：当\Person{拉默客}被确定为从\Person{亚当}算起的第七代时，另外还有他的子女被列出来，使这个数字达到十一，即象征罪恶的数字；因为又加上了三个儿子和一个女儿。\Person{拉默客}的妻子们则有另一种含义，与我此时正要强调的不同。因为我目前讨论的是子女，而不是生育子女的人。既然法律由数字十所象征——由此有那令人难忘的十诫——那么毫无疑问，超过十的数字十一，象征着对法律的违犯，因而就是罪恶。因为这个缘故，在会幕中曾命令悬挂十一件山羊皮的幔子，这会幕在天主子民流浪期间充当移动的圣殿。在那苦衣中，有对罪的提醒，因为山羊要被放在审判者的左边；因此，当我们告明自己的罪时，我们身披苦衣，仿佛在诉说圣咏上所写的：\BibleQuote{我的罪常在我面前}。那么，由杀人者\Person{加音}而传的\Person{亚当}后代，在数字十一中完成了，这个数字象征罪恶；这个数字本身由一个女人组成，正如由同一性别开始了使我们所有人都死亡的罪恶。而且，它被犯下，以便肉体的快乐能够随之而来，这快乐反抗精神；因此，\Person{拉默客}的女儿\Person{纳阿玛}的意思是\Quote{快乐}。但是，从\Person{亚当}到\Person{诺厄}，在\Person{舍特}的谱系中，有十代。而\Person{诺厄}又添了三个儿子，其中一人陷入罪恶，两人受到父亲的祝福；因此，如果你减去那被弃绝的，并将蒙恩的儿子加入数目中，你便得到十二——这个数目在族长和宗徒的事例中显得突出，并且由数字七的部分相乘而成——因为三乘四或四乘三得十二。既然如此，我看到我必须思考并提及这两条谱系如何——它们通过各自的族谱描绘了那两座城市，一座是地上出生的，另一座是重生者的——后来变得如此混淆和混乱，以致全人类除了八个人之外，都该当在洪水中灭亡。\par

% structure: p0059 source=h2 target=subsection reason=relative-heading-rank; base=h2

\subsection{第21章——为何在提及\Person{加音}的儿子\Person{哈诺客}后，谱系便立即延续到洪水，而在提到\Person{舍特}的儿子\Person{厄诺士}后，叙事却又回到人的创造？}

我们必须首先看到，为何在列举\Person{加音}的后裔时，在提及建造该城所在的\Person{哈诺客}之后，其余的后裔立刻就被列举到我所说过的那终点，而那一族和整个谱系都在洪水中被消灭；然而，在提及\Person{舍特}之子\Person{厄诺士}之后，其余的后裔并没有立刻被列举到洪水，而是插入了一段话，内容如下：\Quote{这是亚当后裔的族谱。当\OriginalTerm{天主}{God}创造人的日子，是按\OriginalTerm{天主}{God}的样式造了他；造了一男一女；且祝福了他们，在他们被造的日子，称他们的名字为\Person{亚当}。} \BibleRef{创 5:1} 在我看来，插入这段话的目的，是使年代的推算再次从\Person{亚当}本身开始——作者在谈论地上之城时并无此意图，仿佛天主提及了它，却没有记下它的存续时间。但为何他在提及舍特之子、那希望呼求上主天主之名的人之后，又回到这一重述呢？除非是为了这样呈现这两座城：一座始于凶手，终于凶手（因为\Person{拉默客}也向他的两个妻子承认他杀了人），另一座则是由那希望呼求上主天主之名的人所建立。因为\OriginalTerm{天主之城}{city of God}，这世上的旅客，其最高且完全的地上职责，在那由预表被杀的\Person{亚伯尔}复活之人所生的人身上得到了体现。那一个人就是整个天上之城的合一，尚未完全，但将要完成，正如这预表的形象所预示的。因此，\Person{加音}之子，即占有之子（若非地上的占有，又是何物？），可以在以他名字建造的地上之城中拥有一个名字。诗人正是论到这样的人说：\Quote{他们以自己的名字称呼他们的土地。} 因此他们应验了另一篇诗篇所记载的：\Quote{主，祢在祢的城里必藐视他们的偶像。} 至于\Person{舍特}之子，复活之子，让他希望呼求上主天主之名吧。因为他预表了那会说\Quote{但我好像天主殿中的青橄榄树，我信赖天主的仁慈}的人群。但让他不要寻求世上盛名的空虚荣誉，因为\Quote{以主名为依靠的人是有福的，他不注意虚无及欺骗的愚妄。} 在呈现这两座城之后——一座建于今世的物质利益，另一座建于对天主的希望，但两者都从\Person{亚当}所打开的同一扇门进入这必死的状态，并且都奔向他们各自应得的终点——圣经开始推算年代，并在这一推算中包括其他世代，从\Person{亚当}开始重述。从天主所定罪的苗裔中，如同从一团该受惩罚的泥土中，天主造了一些忿怒之器用于羞辱，另一些仁慈之器用于尊荣；对前者施加惩罚以还其所应得，对后者施以恩宠赐予其所不应得；为的是借与忿怒之器的对比，那在地上作客旅的天上之城可以学会不倚靠自身意志的自由，而希望呼求上主天主之名。因为意志，作为一种由善的天主所造的性体，但因由无中生有，又由不变者所造，故为可变，既可以因行善而转向恶，这发生在它自由选择之时，也可以脱离恶而行善，这唯有赖于天主的协助。\par

% structure: p0061 source=h2 target=subsection reason=relative-heading-rank; base=h2

\subsection{第二十二章 论天主子因受人的女儿诱惑而堕落，为此众人除八人外都该当在洪水中灭亡。}

当人类在这一自由意志下增长和前进时，两座城因共同参与不义而产生了混杂和混乱。这一灾祸如同第一次一样，也是由女人引发的，尽管方式不同；因为那些女人自己并未被欺骗，也没有诱使男人犯罪，而是她们从起初就属于地上之城和地上社会，德行败坏，且因身体的美丽被天主子——即那另一座在世上作客旅的城的公民——所爱慕。美丽确实是天主的美好恩赐；但为了使善人不以为它是一种伟大的善，天主也把它赐予恶人。因此，当善人放弃了那伟大且专属于善人的善时，他们便堕入了一种微不足道的善，这善并非专属于善人，而是善人与恶人所共有；当他们被人的女儿所迷时，他们为了赢得她们为妻而采纳了地上之城的习俗，并抛弃了他们在自己神圣社会中曾遵循的敬虔之道。因此，美丽固然是天主的化工，但只是一种暂时的、肉体的、较低级的善，不宜被爱在永恒、属灵、不变的天主之先。当守财奴偏爱他的金子胜过正义时，这不是金子的过错，而是人的过错；每个受造物也是如此。因为它虽是善，却可以以恶爱，也可以以善爱：它被正当地爱，是当它被有秩序地爱；邪恶地爱，是当它被无秩序地爱。这正是某人在赞美造物主的诗句中简明扼要地说的：\Quote{这些都属于祢，它们是善的，因为创造了它们的祢是善的。它们之中没有一样属于我们，除非我们犯罪，当我们忘记事物的秩序，不去爱祢而爱祢所造之物时。}\par

但如果真正爱慕造物主，也就是说，爱慕祂本身而非任何替代祂的事物，那么这种爱不可能是邪恶的；因为爱本身应当以有秩序的方式去爱，我们爱那使我们生活良善而有德性的事物，这是正确的。因此，在我看来，对德性简短而真实的定义就是：爱慕的秩序；为此，在\BibleBook{}{雅歌}中，基督的新娘、天主的城唱道：\Quote{请在我内安排爱情。}\BibleRef{歌 2:4} 正是这种爱慕的秩序，即这种爱德或依恋，被天主的儿子们打乱了，因为他们离弃了天主，迷恋上了人的女儿们。借着这两个名称（天主的儿子们和人的女儿们），两座城得以充分区分。因为前者虽然按本性是人的子女，却因恩宠获得了另一个名号。因为在这同一部圣经中，当说天主的儿子们爱上了人的女儿们时，他们也被称为天主的天使；因此许多人认为他们不是人，而是天使。\par

% structure: p0064 source=h2 target=subsection reason=relative-heading-rank; base=h2

\subsection{第二十三章——我们是否应当相信，那些属于精神实体的天使，会爱上妇女的美貌，并寻求与她们结合，并且从这种结合中诞生了巨人？}

在本作品第三卷（第五章）中，我们曾简短提及这个问题，但并未断定天使作为精神体是否能够与妇女有肉体的结合。因为经上记载：\Quote{祂使祂的天使成为风}，也就是说，祂使那些本性是精神体的存在，通过委派他们传递信息而成为祂的天使。因为希腊文\OriginalTerm{天使}{\GreekText{ἄγγελος}}，在拉丁文中表现为\Quote{angelus}，意思是使者。但圣咏作者在接下来所说的\Quote{祂的仆役成为火焰}是否指他们的身体，或者是说天主的仆役应当如同属神的火焰般以爱火燃烧，这并不确定。然而，同样可靠的圣经证实，天使曾以不仅可见而且可触摸的身体向人显现。此外，有一种非常普遍的传闻，许多人凭自身经历证实，或者可信赖的人听他人证实，即森林神与农牧神（通常被称为\OriginalTerm{梦淫妖}{incubi}）常常对妇女进行邪恶的侵犯，并满足他们的欲望；而且某些被高卢人称为\OriginalTerm{杜斯}{Duses}的魔鬼不断企图并实现这种污秽之事，这一点被普遍断言，以至否认它将是厚颜无耻的。从这些断言中，我不敢断定是否有一些精神体寓于气态实体中（因为这种元素即使被扇子搅动，身体也能感觉到），并且能够产生情欲并与妇女有身体上的结合；但可以肯定的是，我绝不会相信天主的圣天使在那个时候会如此堕落，我也不会认为伯多禄宗徒指的是他们，他说：\Quote{因为天主既然没有宽恕犯罪的天使，反而把他们投入地狱，交在黑暗的锁链中，留待审判。}\BibleRef{伯后 2:4} 我认为他指的是那些最初与他们的首领魔鬼一起背弃天主的天使，魔鬼以蛇形嫉妒地欺骗了第一个人。但同一部圣经提供了充分的证据，证明即使虔诚的人也曾被称为天使；因为关于若翰，经上记载：\Quote{看，我派遣我的使者（天使）在你面前，他要在你前面预备你的道路。}\BibleRef{谷 1:2} 而先知玛拉基亚，因特别赐予他的特殊恩宠，也被称为天使。\BibleRef{拉 2:7}\par

但是，有些人由于我们读到那些被称为天主的使者的人与所爱的女子结合，所生的后代并非像我们同类的人，而是巨人，而感到困扰；仿佛在我们自己的时代（如上所述）就没有出生过比普通身材高大得多的人。几年前，当如今由哥特人完成的罗马毁灭临近时，不是有一个女子，与她的父母一起，因其巨大的身材远超众人吗？从各处涌来的好奇人群来观看她，最令他们惊讶的是，她的父母都未达到普通人的最高身材。因此，即使在天主之子（亦称天主的使者）与人的女儿，即那些按人生活的人结合之前，也就是说，在\Person{舍特}之子与\Person{加音}之女结合之前，巨人也很可能出生。因为连正典圣经本身，其中记载了此事，也这样说道：\BibleQuote{当人开始在地上繁衍，生养女儿时，天主的儿子看见人的女儿美貌[佳]，就随意选取，作为妻子。上主天主说：我的神不永远与人相争，因为人也是血肉；但他的寿数可到一百二十年。那时地上有巨人；以后，天主的儿子来到人的女儿那里，她们为他们生了孩子，这些人就成为巨人，即古代的有名望之人。}神圣圣经的这些话充分表明，在那些日子，当天主的儿子因爱她们的美貌而娶人的女儿时，地上已有巨人。因为这部圣经惯于把外貌美丽的人称为\Quote{佳}。但在此结合之后，也有巨人生出。因为圣经的话是：\BibleQuote{那时地上有巨人，\Emph{以后也有}，当天主的儿子来到人的女儿那里时。}因此，既有\Quote{在那时}的巨人，也有\Quote{以后}的巨人。而\Quote{她们为他们生了孩子}这句话清楚地表明，在天主的儿子以这种方式堕落之前，他们为天主生子，而非为自己——也就是说，并非出于性欲的冲动，而是履行传宗接代的职责，意图不是生养一脉以骄矜，而是产生公民以充实天主之城；并且作为天主的使者，他们要将这消息传达给他们：应寄望于天主，如同那从\Person{舍特}所生、成为复活之子并盼望呼号上主天主之名的人一样，在其中他们和他们的后裔将是永恒福乐的共同继承者，以及以天主为父的大家庭中的兄弟。\par

但是，关于那些使者并非不具人性的使者（如有些人所认为），圣经本身已作决断，它明确宣布他们是人。因为当初先是记载\BibleQuote{天主的使者看见人的女儿美貌，就随意选取，作为妻子}，紧接着又记载\BibleQuote{上主天主说：我的神不永远与这些人相争，因为他们是血肉。}因为藉着天主的神，他们曾成为天主的使者和天主的儿子；但他们堕向低级事物，就称为人，这是本性的名称，而非恩宠的名称；他们被称为血肉，作为圣神的离弃者，并因离弃而被（圣神）离弃。的确，七十贤士译本既称他们为天主的使者，又称他们为天主的儿子，尽管并非所有抄本都如此，有的只用了天主的儿子之名。而犹太人更偏爱的\Person{阿奎拉}译者，既未译为天主的使者，也未译为天主的儿子，而是译为诸神的儿子。但两者都正确。因为他们既是天主的儿子，因此也是自己父亲们的弟兄，这些父亲同样是同一天主的儿女；他们又是诸神的儿子，因为由诸神所生，他们自身也同诸神一起是神，正如圣咏所说：\BibleQuote{我曾说：你们是神，你们都是至高者的儿子。}因为七十贤士译者被公认为领受了预言的神恩；因此，即使他们在神恩权能下做了某些改动，而未坚持严格翻译，我们也不能怀疑这是出于天主的默示。然而，希伯来文词可能被认为有歧义，并可接受两种译法，即\Quote{天主的儿子}或\Quote{诸神的儿子}。\par

我们要略过那些被称为\OriginalTerm{伪经}{apocryphal scriptures}的作品中的神话，因为它们的模糊起源不为教父们所知，而真正的圣经权威正是通过一个最确实、最明确的传承从教父们那里传给了我们。因为这些\OriginalTerm{伪经}{apocryphal scriptures}作品虽有一些真理，却包含许多虚假陈述，因而没有\OriginalTerm{正典权威}{canonical authority}。我们不能否认，从亚当算起的第七代哈诺客留下了一些神圣著作，因为宗徒犹达在他的正典书信中这样断言。但这些著作没有列入那个由历任司祭勤勉保存在希伯来民族圣殿中的圣经正典，这并非没有缘故；因为它们的古老性引起了怀疑，无法确定它们是否真是哈诺客的作品，而谨慎保存了正典书籍并通过连续传承传递下来的人并未将它们作为真作提出。所以，那些以哈诺客之名流传、包含关于\OriginalTerm{巨人}{giants}的神话（说他们的父亲不是人）的著作，被明智之人正当地判定为不真实；正如异端者以其他先知之名、以及近来以宗徒之名制作了许多著作一样，所有这些著作经过仔细审查，皆被排除在正典权威之外，称为伪经。因此，根据希伯来和基督信仰的正典圣经，洪水之前确实有许多\OriginalTerm{巨人}{giants}，他们是地上人类社会的成员，而按肉身是\OriginalTerm{舍特的儿子}{sons of Seth}的\OriginalTerm{天主的儿子}{sons of God}，在他们背弃正义时沉沦于这个群体。我们也不奇怪，即或从他们也会生出\OriginalTerm{巨人}{giants}。因为他们的孩子并非全是巨人；但当时比洪水之后的其他时期更多。造物主乐意生出他们，为要向智者显明，美貌、身材和力量都没有多大意义，智者的福乐在于属灵和不朽的祝福，在于更美好、更持久的恩赐，在于专属于善人而不被善恶共享的美善事物。另一先知证实了这一点，他说：\newline{}\Quote{这些是巨人，从起初就著名，身材高大，善于作战。主不拣选他们，也未将知识之道赐给他们；他们因没有智慧而遭毁灭，因自己的愚昧而灭亡。}\par

% structure: p0069 source=h2 target=subsection reason=relative-heading-rank; base=h2

\subsection{第二十四章——如何理解主对将要在洪水中灭亡的人所说的话：\Quote{他们的日子要有一百二十年。}}

但天主所说\Quote{他们的日子要有一百二十年}，不应理解为预言从此以后人活不过一百二十年——因为洪水之后我们发现有人活过五百多年——而应理解为天主说这话时，诺厄已接近他生命的第五个世纪，即已活了四百八十年，圣经常用整体或大部分来代表全体，故称之为五百年。洪水是在诺厄生命第六百年、二月来临的；因此，这些被判罪之人剩余的年数被预言为一百二十年，这些年份一过，他们就要被洪水毁灭。合理地相信，洪水之所以如此来临，是因为地上再也找不到不值得共享如此明显审判性死亡的人——并非说善人（终有一死）在经历这种死亡后会变得更糟。然而，在洪水中死去的人中没有一个是圣经所记载的舍特的后裔。以下是神性记载的洪水原因：\newline{}\BibleQuote{上主天主见人在地上的罪恶很大，人心中思想的每一个意念都终日是恶。上主后悔造人在地上，心中忧伤。上主说：我要从地面上消灭我所造的人，以及走兽、爬虫和飞鸟：因为我后悔造了他们。}\BibleRef{创 6:5}\par

% structure: p0071 source=h2 target=subsection reason=relative-heading-rank; base=h2

\subsection{第二十五章——论天主的义怒，它不激动他的心灵，也不扰乱他不变的宁静。}

天主的义怒不是他心中激动的情绪，而是对罪施以惩罚的审判。他的思想和重新考虑也是改变事物的不变理性；因为他不像人一样对自己所做的任何事后悔，因为他在一切事上的决定与其预知一样确定不移。但如果圣经不使用上述这类表达，它就不能熟悉地渗入所有阶层的人心，它为了他们的益处而寻找机会接近他们，以警醒骄傲者，唤醒粗心者，锻炼好奇者，满足明智者；它若不先俯就他们，以某种方式降下到他们所在之处，就不能做到这一点。它宣告地上和空中所有动物都将死亡，则是表明即将来临的灾难之巨大；并非威胁要毁灭无理性的动物，仿佛它们也因罪而招致死亡一样。\par

% structure: p0073 source=h2 target=subsection reason=relative-heading-rank; base=h2

\subsection{第二十六章——诺厄奉命建造的方舟在每方面都是基督和教会的预象。}

此外，既然天主命令诺厄——一个义人，并且正如真实的圣经所说，在他的世代是一个完人（当然不是具有天主之城居民在永生境界中与天使相等的完美，而是就他们在这世间的旅居中所能达到的完美而言）——既然天主命令他造一只方舟，使他与其家人，即他的妻子、儿子和儿媳，以及那些听从天主命令来到方舟的动物，得以从洪水的毁灭中得救：这无疑是那在世间旅居的天主之城的预象，也就是说，教会的预象，教会是由那悬挂着天主与人之间的中保、那人基督耶稣的木头上得救的。\BibleRef{弟前 2:5}因为就连方舟的长、宽、高尺寸，也象征了他来临时所取的人体，正如早已预言的。因为人体的长度，从头顶到脚底，是其宽度（左右）的六倍，是其深度或厚度（从前到后）的十倍；也就是说，如果你测量一个仰卧或俯卧的人，他从头到脚的长度是其左右宽度的六倍，是其离地高度的十倍。因此，方舟的长是三百肘，宽是五十肘，高是三十肘。方舟的侧面开了一道门，这无疑预示了被钉者的肋旁被长矛刺透时所受的创伤；因为凡投奔他的人，都由此进入；因为从那里流出了圣事，信徒藉以受到入门。而方舟被命令用方木建造，这象征了圣人们坚定不移的稳定生活；因为无论你如何旋转一个立方体，它仍然挺立。方舟构造的其他特征也都是教会特征的预象。\par

但我们此刻没有时间深入探讨此问题；事实上，我们在反对摩尼教徒福斯特所写的著作中已详述过，他否认希伯来文书籍中有任何关于基督的预言。或许一人的解释优于另一人，而我们的解释并非最佳；但所有解释都必须指向我们所谈论的这座天主之城，它在这罪恶的世界中如同在洪水中旅居，至少解经者不应严重偏离作者的原意。例如，我在反对福斯特的著作中对那句话的解释——\Quote{你要用下层、第二层和第三层来造它}——我认为，因为教会是从万民中聚集的，所以它被说成有两层，代表两类人，即受割礼的和未受割礼的，或者如宗徒所称的，犹太人和外邦人；又因万民是由诺厄的三个儿子繁衍的，所以有三层。现在，任何人都可以反对这个解释，并提出另一个符合信仰准则的解释。因为方舟不仅有下层房间，还有上层房间，被称为\Quote{第三层}，以便从底层起第三层有可居住的空间，有人可能将这些解释为宗徒所赞扬的三项恩宠——信、望、爱德。或者更适当地，它们可能被设想为代表福音中的三种收成：三十倍、六十倍、一百倍——贞洁的婚姻居住在最底层，贞洁的寡居在中间层，贞洁的童贞在最顶层。也可以给出任何更好的解释，只要不偏离与这座城的关联。对于这段经文所需解释的所有其他细节，我也要作同样的声明，即：尽管有不同的解释，但它们都必须与唯一的、和谐的公教信仰相一致。\par

% structure: p0076 source=h2 target=subsection reason=relative-heading-rank; base=h2

\subsection{第二十七章 论方舟与洪水，以及我们既不能同意那些只接受字面历史而拒绝寓意解释的人，也不能同意那些主张寓意而否认历史意义的人。}

然而，没有人应当认为这些事情是毫无目的而写的，或者我们只应研究历史真理而不顾任何寓意；相反，也不应认为它们只是寓意，而根本没有这样的事实，或者无论是否如此，这里都没有对教会的预言。因为哪个有理智的人会争辩说，这些数千年来被虔诚保存、并经过如此有序传承的书籍，竟然是毫无目的而写的，或者我们阅读时只应考虑纯粹的历史事实？因为，且不提其他例子，如果动物的数量要求建造一只巨大尺寸的方舟，那么为什么需要把每种动物中的两个不洁和七个洁净的送进方舟，而它们本来可以保存相同数量的？难道那位命令它们得以保存以繁衍后代的天主，不能像祂创造它们那样重新恢复它们吗？\par

但那些主张这些事从未发生，而只是象征其他事物的人，首先就设想不可能有这样大的洪水，使水超过最高的山脉十五肘高，因为据说云层不能升到\Place{奥林匹斯山}顶之上，因为山顶直达天空，那里没有较稠密的大气层——风、云、雨都源自该大气层。他们没有认识到，最稠密的元素——土，可以在那里存在；或者他们可能否认山顶是土。那么，这些元素测量者和称量者凭什么争论说，土可以被提升到那些高空，而水却不能，尽管他们承认水比土更轻、更易上升？他们提出什么理由，说明为什么较重的、较低的元素土，能够如此长久地上升到宁静的苍穹，而较轻的、更易上升的水，却连短暂的时间都不能同样上升？\par

他们又说，那只方舟的面积容不下如此多类的动物，包括不洁的每种一对和洁净的每种七对。但在我看来，他们只计算了长三百肘、宽五十肘的一层面积，却忘了上面还有同样的一层，再上面还有同样大的一层；因而总面积为九百肘长、一百五十肘宽。如果我们接受\OriginalTerm{奥利振}{Origen}颇为恰当的建议，即天主的仆人\OriginalTerm{梅瑟}{Moses}，正如经上所载，\BibleQuote{精通埃及人的一切智慧}\BibleRef{宗 7:22}，喜爱几何学，可能指的是几何学意义上的肘，据说一几何肘等于我们通常的六肘，那么谁看不出这些尺寸赋予方舟多大的容量呢？至于他们反对说如此大的方舟无法建造，那是一种非常愚蠢的诽谤；因为他们知道巨大的城市被建造起来，也应当记得方舟的建造历时一百年。或者，也许石头与石头之间仅用石灰粘合就能建造几英里长的城墙，但木板却不能通过榫眼、螺栓、钉子以及沥青胶合在一起，建造一艘不是用弯曲的肋骨而是用直木料制成的方舟；这方舟不是由建造者下水，而是当水涨到时被水的自然浮力托起，并且它在漂浮时主要靠天主的看顾而非人的技艺得以避免沉没。\par

关于那些吹毛求疵者常问的微小生物，不仅包括老鼠和蜥蜴，还有蝗虫、甲虫、苍蝇、跳蚤等，是否方舟中存有的数目比天主要在祂命令中指定的更多，那些被这个问题困扰的人应当被告知：\Quote{地上的一切爬虫}这句话仅表明，那些能在水中生活的动物——无论是潜游在水中的鱼类，还是浮在水面的海鸟——无需保存在方舟里。而当说\Quote{公母一对}时，无疑是指种族的繁衍，因此那些从无生命之物或腐败之物中无需两性结合而生的生物，就不必在方舟中；或者，如果它们在方舟中，可以像平常在房屋中那样，没有确定的数量；或者，如果必须有一定数量的、所有不能自然生活在水中的动物，以便那正在进行的至圣奥秘得以在现实事物中体现并完美预表，那么这也不是\OriginalTerm{诺厄}{Noah}或其儿子的责任，而是天主的。因为诺厄并没有捕捉动物放入方舟，而是它们自己前来寻求进入时就让他进入。这正是\BibleQuote{它们要到你们这里来}\BibleRef{创 6:19}这句话的意味——也就是说，不是靠人的努力，而是靠天主的旨意。但我们当然不必相信那些没有性别的动物也来了；因为明确而肯定地说：\BibleQuote{它们要公母一对。}因为有些动物是从腐败中出生的，但后来它们自己交配并产生后代，如苍蝇；而另一些则没有性别，如蜜蜂。至于那些有性别但无法繁殖后代的动物，如公骡和母骡，它们很可能不在方舟中；保存它们的父母——即马和驴——就已足够；这也适用于所有杂种。然而，如果为了奥秘的完整性而有必要，它们也在那里；因为这类动物也有\Quote{公母一对}。\par

另一个常见的问题是有关肉食动物的食物：是否在不违反保存数量命令的前提下，方舟中必然还有别的动物作为它们的食物？或者，更可能的是，存在某种不是肉类的食物，却适合所有动物。因为我们知道，许多以肉为食的动物也吃植物产品和果实，特别是无花果和栗子。因此，如果那位智慧而正义的人蒙天主指示什么适合每种动物，从而不靠肉类而预备并储存适合每种物种的食物，那有什么奇怪呢？况且，有什么是饥饿不能使动物吃的呢？或者，有什么不能由天主使之变得甜美而有益呢？天主以神圣的轻易本可使它们完全不需要食物，但若不是为了如此伟大奥秘的完整需要它们被喂养，祂就会那样做。然而，只有好争论的人才会认为，在如此多样且详尽的情节中没有预表教会。因为万国已经如此充满教会，并被包括在其合一的框架内，洁与不洁的并存，直到指定的结局；这一非常明显的应验使我们毫不怀疑，我们也应当如何解释那些较为隐晦、不易辨别的部分。既然是这样，即使最放肆的人也不敢断言这些事是毫无目的地写下的，或者虽然事件真正发生却毫无意义，或者它们并非真正发生而只是寓意，或者它们根本没有对教会的预表性指涉；既然已经证明，我们反而必须相信，它们被记入记忆和文字是有智慧的目的，它们确实发生了，有意义，并且这意义对教会有预表性指涉，那么，这本书达到这个目的后，现在可以结束了，以便我们继续追溯洪水之后的历史中两个城的历程——属地之城，按人生活；属天之城，按天主生活。\par

\SourceRecord{Translated by Marcus Dods.；Nicene and Post-Nicene Fathers, First Series；Vol. 2.；Edited by Philip Schaff.；Buffalo, NY: Christian Literature Publishing Co.,；1887.；<http://www.newadvent.org/fathers/120115.htm>.}

% structure: p0001 source=h1 target=section reason=page-title

\section{天主之城（第十六卷）}

本书前半部分，从第一章到第十二章，依圣经展示了从诺亚到亚伯拉罕期间两座城，即属地之城与属天之城的发展；后半部分则专门论述从亚伯拉罕到以色列君王期间属天之城的发展。\par

% structure: p0003 source=h2 target=subsection reason=relative-heading-rank; base=h2

\subsection{第一章——洪水之后，从诺亚到亚伯拉罕，能否找到按天主生活的家族？}

从圣经中难以发现，洪水之后，圣城的踪迹是否持续不断，或是在无神时期的间隔中完全中断，以至于在人类中竟找不到一位唯一真神的崇拜者；因为从诺亚——他与其妻、三子及三位儿媳在方舟中从洪水的毁灭中得救——直到亚伯拉罕，我们在正经书中没有发现任何人的虔诚曾得到天主明确的证言，除非是诺亚，他以预言的祝福称赞他的两个儿子闪和雅弗，同时他看到并预见了很久以后将要发生的事。也是出于这预言的精神，当他那个中子——即比长子年幼、比幼子年长的儿子——冒犯了他时，他并没有直接诅咒他本人，而是诅咒他的儿子（即他自己的孙子），说：\BibleQuote{客纳罕童仆应被诅咒，他必作他兄弟的奴仆。}\BibleRef{创 9:25} 同样，他也接着祝福了他的其他两个儿子，即长子和幼子，说：\BibleQuote{闪的上主天主应受赞美；客纳罕应作他的奴仆。愿天主使雅弗扩大，使他住在闪的帐幕中。}\BibleRef{创 9:26} 同样，诺亚栽种葡萄园，因果实而醉酒，睡觉时赤身露体，以及当时所发生并记载的其他事情，都充满了预言意义，隐藏在奥秘之中。\par

% structure: p0005 source=h2 target=subsection reason=relative-heading-rank; base=h2

\subsection{第二章——诺亚的儿子们预表了什么}

那些当时隐藏的事情，现在通过随后发生的实际事件已充分揭示。因为谁能仔细而明智地思考这些事情，而看不出它们是在基督身上应验的呢？闪，基督按肉身是由他降生的，意为\Quote{有名}。还有什么名字比基督更伟大？祂名字的芬芳如今处处被感知，以致预言预先歌颂它，在\WorkTitle{雅歌}\BibleRef{歌 1:3}中将其比作倾注的香膏。万民的\Quote{扩张}不也是住在基督的房屋里，即在教会里吗？因为雅弗意为\Quote{扩张}。而含（\Emph{即}\Quote{热}），是诺亚的中子，可以说他使自己与两者分离，留在他们之间，既不属以色列的初果，也不属外邦人的圆满，他除了代表异端者的支派，还有什么意义呢？这些人热心，不是忍耐之灵，而是不忍耐之灵，异端者的心胸常以此燃烧，并以此扰乱圣徒的平安。但即使异端者也给长进的人带来益处，按宗徒所说：\BibleQuote{在你们中间也免不了有异端，好使那些被认可的人得以显明出来。}\BibleRef{格前 11:19} 因此，别处也说：\Quote{受教的儿子必成为智慧，他利用愚昧人作他的仆人。}因为当异端者火热的躁动激起关于许多天主教信德问题的争论时，捍卫它们的必要性迫使我们更精准地研析它们，更清晰地理解它们，更热切地宣讲它们；而由对手提出的问题便成了教导的契机。然而，不仅那些公开与教会分离的人，而且所有以基督徒之名夸耀却同时过着堕落生活的人，似乎都可以合理地由诺亚的中子所预表：因为那人的赤裸所预表的基督的苦难，既被他们的信仰宣认，又被他们的邪恶行为所玷污。因此，关于这样的人，经上说：\BibleQuote{凭他们的果实你们可以认出他们来。}\BibleRef{玛 7:20} 所以含在他的儿子身上受诅咒，那儿子仿佛是他的果实。同样，他的这个儿子客纳罕被适当地解释为\Quote{他们的行动}，这不是别的，正是他们的行为。但闪和雅弗，即受割礼与未受割礼，或者如宗徒所称的犹太人与希腊人，但他们蒙召并被成义，他们似乎发现了父亲的赤裸（这象征救主的苦难），就拿了一件衣服搭在肩上，倒退着进去，遮盖了父亲的赤裸，却没有看到他们所敬畏所隐藏的。因为我们既尊崇为我们完成的基督的苦难，也憎恨那些钉死祂的犹太人的罪行。衣服象征着圣事，他们的背象征着对过去之事的纪念：因为教会庆祝基督的苦难作为已经完成的事，不再期待，如今雅弗已住在闪的帐幕中，而他们邪恶的兄弟介于他们之间。\par

但邪恶的弟兄，凭其儿子（即他的作品），成为他善弟兄们的仆人或奴隶，当善人巧妙地利用恶人时，或为操练他们的忍耐，或为增进他们的智慧。因为宗徒作证说，有些人传基督并非出于纯正的动机；他说：\BibleQuote{无论出于假意，或是出于诚心，基督总被传扬；为此我欢喜，而且还要欢喜。}\BibleRef{斐 1:18} 因为正是基督亲自栽种了这葡萄树，先知论它说：\BibleQuote{万军上主的葡萄园就是以色列家；}\BibleRef{依 5:7} 祂喝这酒，无论我们这样理解那杯，就是祂所说的：\BibleQuote{你们能饮我将要饮的杯吗？}\BibleRef{玛 20:22} 以及\BibleQuote{父啊！如果可能，就让这杯离开我吧！}\BibleRef{玛 26:39} 祂显然是指祂的苦难。或者，由于酒是葡萄树的果实，我们也可以愿意如此理解：从这葡萄树，即从以色列民族，祂取了血肉，为能受苦；\BibleQuote{他喝醉了，}即祂受苦；\BibleQuote{且赤身露体，}即祂的软弱在受苦中显露出来，正如宗徒所说：\BibleQuote{祂虽因软弱而被钉在十字架上}\BibleRef{格后 13:4}。因此同一宗徒说：\BibleQuote{天主的软弱比人更强壮，天主的愚昧比人更智慧。}\BibleRef{格前 1:25} 并且当圣经在\BibleQuote{他赤身露体}的表述上加上\BibleQuote{在他的帐棚内}时，它巧妙暗示耶稣将要受十字架和死亡，是出于自己的家人，自己的亲属和同胞，即犹太人。基督的这苦难，只有被弃绝者才在外表上和言语上承认，因为他们所承认的，他们并不了解。但蒙选者却在内心持守这如此伟大的奥迹，并在心中内在尊崇天主的这软弱和愚昧。关于这一点，有一个预象：\Person{含}出去宣扬他父亲的赤身；而\Person{闪}和\Person{耶斐特}，为了遮盖或尊崇它，反而进去，也就是说，从内心去做。\par

天主的这些圣经奥秘，我们尽力探究。并非所有人都以同样的信心接受我们的解释，但所有人都确信，这些事的发生和记载，无一不是对未来事件的某种预象，而且它们只应指向基督和祂的教会，即天主之城，从人类历史之初就藉着预象被宣告，这些预象如今我们处处看到应验。从\Person{诺厄}两个儿子的祝福和中间儿子的诅咒，直到\Person{亚巴郎}，或超过一千年，如我所说，没有提到任何敬拜天主的义人。我并非因此断定没有义人；但一一提及未免冗长，而且会显示历史准确性而非预言远见。这些圣书的作者，更确切地说，他内里的天主圣神，其目的不仅是记录过去，而且是描绘未来，只要涉及天主之城；因为凡是对非此城公民的记述，若不是为训导她，就是作为陪衬以增加她的荣耀。然而我们不应认为凡所记载都有含义；但那些本身没有含义的事物，是为了有意义的事物而被交织在一起。只有犁铧能破土；但为此，犁的其他部分是必需的。只有竖琴和其他乐器的弦能发出悦耳的声音；但为使它们如此，乐器还有其他部分，这些部分虽不被歌唱者弹拨，却与弹拨的弦相连，并产生音符。同样，在这预言史中，有些叙述没有含义，但可以说是框架，有意义的事物附着其上。\par

% structure: p0009 source=h2 target=subsection reason=relative-heading-rank; base=h2

\subsection{第三章——论\Person{诺厄}三子的后代。}

因此，我们必须在本著中加入对诺厄三个儿子的族系的解释，以免它有助于阐明这两座城在时间进程中的发展。圣经首先提到幼子，他名叫\Person{耶斐特}：他有八个儿子，并藉着其中两个儿子有了七个孙子，一个生了三个，另一个生了四个；总共十五个后裔。诺厄的次子\Person{含}，有四个儿子，并藉着其中一个有了五个孙子，又藉着其中一个有了两个曾孙；总共十一个。列举这些之后，圣经又回到长子，并说：\Quote{\Person{雇士}生了\Person{尼默洛得}；他开始成为地上的巨人。他是在上主天主面前猎取的巨人：因此人们说，像\Person{尼默洛得}一样，在上主面前猎取的巨人。他王国的开始是\Place{巴比伦}、\Place{厄勒客}、\Place{阿加得}和\Place{加蓝}，都在\Place{史纳尔}地。从那地出去\Person{亚述}，建造了\Place{尼尼微}和\Place{勒曷波特}城，以及\Place{加拉}，并\Place{尼尼微}与\Place{加拉}之间的\Place{勒森}：这是一座大城。}现在，这个\Person{雇士}，巨人\Person{尼默洛得}的父亲，是\Person{含}的儿子中第一个被命名的，被归于他有五个儿子和两个孙子。但他要么是在他的孙子出生之后才生了这个巨人，要么——更为可信的——圣经因着他的杰出而单独提到他；因为也提到他的王国，始于那座壮丽的\Place{巴比伦}城，以及与之一起被提及的其他地方，无论城市或区域。但关于属于\Person{尼默洛得}王国的\Place{史纳尔}地所记载的，即\Person{亚述}从那地出去并建造了\Place{尼尼微}和与其他一起被提及的城市，是在很久之后发生的；但（圣经作者）藉此机会在这里提到它，是为了亚述王国的伟大，该王国由\Person{贝鲁斯}的儿子\Person{尼努斯}奇妙地扩展，他建造了以他命名的伟大城市\Place{尼尼微}。但\Person{亚述}，亚述人的祖先，不是\Person{诺厄}的儿子\Person{含}的儿子，而是列于\Person{闪}的儿子中。由此看来，在\Person{闪}的后裔中出现了人，他们后来占领了那个巨人的王国，并从那里出发建立了其他城市，第一座是从\Person{尼努斯}得名的\Place{尼尼微}。从那里，圣经又回到\Person{含}的另一个儿子\Person{米兹辣因}；他的儿子们不是作为七个人被列举，而是作为七个民族。从第六个，好像从第六个儿子，据说\Person{培肋舍特人}这个种族产生了；所以总共有八个。然后它又回到\Person{迦南}，\Person{含}因他而受到诅咒；并且列出了他的十一个儿子。然后列出了他们所占据的领土以及一些城市。因此，如果我们计算儿子和孙子，\Person{含}的后裔登记了三十一个。\par

还剩下要提及诺厄的长子闪的儿子们；因为这份族谱叙述是从幼子逐渐上升到他的。但在记录闪的儿子的开始时，有一个含糊之处需要解释，因为它与我们的研究对象密切相关。因为我们读到：\Quote{闪也生了孩子——他是厄贝尔的所有子孙的祖先，耶斐特的长兄。}\BibleRef{创 10:21} 这是这些话的顺序：闪生了厄贝尔，甚至对他自己，即闪自己生了厄贝尔，而闪是他所有孩子的父亲。这是要我们明白，闪是所有将被提及的后裔的始祖，无论是儿子、孙子、曾孙，还是任何远裔。因为闪并未生下厄贝尔，厄贝尔实际上是他的第五代后裔。因为闪在儿子中生了阿帕沙得；阿帕沙得生了卡南，卡南生了舍拉，舍拉生了厄贝尔。他被称为闪的后裔中的首位，甚至在他的儿子之前，尽管只是第五代的孙子，这是有充分理由的；因为据传统，希伯来人从厄贝尔得名，虽然将名字源于亚巴郎（似乎为\Emph{Abrahews}）的词源也可能正确。但几乎可以肯定前者是正确的词源，即他们按照厄贝尔被称为\Emph{Heberews}，然后去掉一个字母，成为希伯来人；而他们的语言被称为希伯来语，只有以色列子民讲这种语言，在以色列子民中有一座天主之城，神秘地预兆于全体子民中，并真实地存在于圣徒中。然后首先列出了闪的六个儿子，然后其中一个儿子生了四个孙子；然后又提到闪的另一个儿子，他生了一个孙子；而他的儿子，即闪的曾孙，是厄贝尔。厄贝尔生了两个儿子，给其中一个取名培肋格，意思是\Quote{分开}；圣经附加了这个名字的理由，说：\Quote{因为在他的时候大地就被分开了。}此事的含义将在后面显现。厄贝尔的另一个儿子生了十二个儿子；因此闪的所有后裔是二十七人。诺厄三个儿子的后裔总数是七十三人：耶斐特十五人，含三十一人，闪二十七人。然后圣经补充说：\Quote{以上是闪的儿子，按照他们的家族、按照他们的语言、在各自的土地上、按照他们的民族。}然后全数总结：\Quote{以上是诺厄儿子的家族，按照他们的世系、在他们各自的国家内；洪水以后，这些家族的子孙分散在各个岛屿上。}由此可见，这七十三人（或者更确切地说，正如我稍后将要指出的，七十二人）不是个人，而是民族。因为在前面一段，当列举耶斐特的儿子时，结束时说：\Quote{以上这些人是各个岛屿上的民族，按照他们的语言、部族和民族分散在他们的土地上。}\par

但是，正如我前面所指出的，在含的儿子们中明确提到了邦国。\Quote{米兹辣因生了那些称为路丁人的人；}以及其他七个民族也是如此。在列举了他们全部之后，经文总结说：\Quote{这些是含的儿子们，按照他们的家族、语言、地域和邦国。}那么，其中几个人的子孙没有被提及的原因，是由于他们按出生属于其他邦国，而他们自身并没有成为邦国。否则，为什么耶斐特虽算有八个儿子，却只提到其中两个的儿子；含算有四个儿子，却只提到三个有后代；闪算有六个儿子，却只追溯了两个的后裔？难道其余的人没有子女吗？我们不能这样认为；但他们并没有产生足够大的邦国以致被提及，而是被他们出生所属的邦国吞没了。\par

% structure: p0013 source=h2 target=subsection reason=relative-heading-rank; base=h2

\subsection{第四章——论语言的多样性与巴贝耳塔的建立。}

尽管这些民族据说是按照他们的语言分散的，但记述者回溯到所有人只有一种语言的时代，并解释了语言的多样性是如何产生的。他说：\BibleQuote{全地都是一种口音，言语也都一样。他们从东方迁移的时候，在史纳尔地方发现一块平原，就住在那里。他们彼此说：来，我们做砖，用火烧烤。他们就用砖当石，用沥青当灰泥。他们又说：来，我们建造一座城和一座塔，塔顶摩天，好使我们扬名，免得我们分散在全地面上。上主下来察看世人所建造的城和塔。上主天主说：看，他们都是一个民族，都说一样的语言。他们现在开始做这事；以后他们所想做的，就没有不成功的。来，我们下去，在那里混乱他们的语言，使他们彼此语言不通。于是上主将他们分散到全地面上；他们便停止建造那城。因此那城名叫混乱，因为上主在那里混乱了全地的语言，并从那里将他们分散到全地面上。}\BibleRef{创 11:1} 这座被称为混乱的城就是巴比伦，外邦人的历史也记载了它奇妙的建造。因为巴比伦就是混乱的意思。由此我们推断，巨人尼默洛德是它的建造者，正如稍前圣经在提到他时暗示的，说他的王国起于巴比伦，也就是说，巴比伦作为京城和王室驻地，凌驾于其他城市之上；尽管它并未达到其傲慢邪恶的建造者所设计的宏大尺寸。计划是建得极高直达天空，无论这是指他们打算建造的一座比其他塔更高的塔，还是指所有的塔（以单数表示，就如我们说士兵指军队，或当摩西用灾祸打击埃及时，提及青蛙或蝗虫指成群的青蛙或蝗虫一样）。\EditorialNote{出谷纪第十章} 但这些虚妄狂妄的人想要做什么呢？当他们把塔建得高于地上所有的山和云层时，他们怎能指望这高耸的建筑反对天主呢？任何精神或物质上的高升又能对天主造成什么伤害呢？通往天堂的安全正途是谦卑，它把心举向上主，而不是反对祂；正如这巨人被称为\Quote{反对上主的猎人}。有些人因希腊词\OriginalTerm{反对}{\GreekText{ἐναντίον}}的歧义而误解了这一点，把它译成\Quote{在上主面前}，而不是\Quote{反对上主}；因为该词兼具\Quote{在……面前}和\Quote{反对}两种含义。在圣咏中，这个词被译作：\Quote{让我们在上主我们的造主面前哭泣。}同样的词出现在约伯传中，那里写着：\Quote{你曾对天主发泄怒气。}\BibleRef{约 15:13} 所以，这巨人应被认作\Quote{反对上主的猎人}。而\Quote{猎人}一词除了指欺骗者、压迫者、地上动物的毁灭者之外，还有什么意思呢？因此，他和他的民众建造这座塔是为了反对上主，以此表达他们不虔敬的骄傲；他们的邪恶意图即使没有成功，也受到了天主的公正惩罚。但惩罚的性质是什么呢？既然舌头是统治的工具，骄傲就在舌头上受了惩罚；因此，不肯理解天主命令的人，在自己发令时也不被人理解。于是那个阴谋瓦解了，每个人都远离了他听不懂的人，而与能听懂他言语的人交往；民族按照语言被分开，分散到地上各处，一如天主所愿，祂以我们无法看见和理解的方式完成了这一切。\par

% structure: p0015 source=h2 target=subsection reason=relative-heading-rank; base=h2

\subsection{第五章——论天主降临以混乱建城者的语言。}

我们读到：\Quote{主下来要看看世人所建造的城和塔。}那不是天主的子女，而是那按纯人性方式生活的社会，即我们称为地上之城的。天主，祂总是完整地无所不在，并不作地方性的移动；但祂被称为\Quote{下来}，当祂在地上做出超乎寻常的事情时，这仿佛使祂的临在被人感觉到。同样，祂并不借着\Quote{观看}而学到什么新事物，因为祂不可能对任何事物无知；但祂被称为\Quote{看见并认识}，在时间中，祂使别人看见并认识祂所显示的事。因此，那座城先前并没有被看见，直到天主使它被看见，当祂显示那城在祂眼中是多么可憎时。我们确实可以将天主下降到那城解释为祂所居住的天使们的下降；因此后面的话：\Quote{上主天主说：看，他们都是一个民族，都说一样的语言，}以及随之而来的：\Quote{来，我们下去，混乱他们的语言，}是对前面提到的\Quote{主的下降}如何实现的复述。因为如果祂已经下去了，为什么祂又说：\Quote{来，我们下去，混乱他们？}——这些话似乎是向天使说的，并暗示那在天使中的祂藉着天使们的下降而下降。这些话最恰当的是，不是\Quote{你们下去混乱，}而是\Quote{我们混乱他们的语言；}表明祂这样借着祂的仆人们工作，以至于他们自己也是与天主同工的，正如宗徒所说：\BibleQuote{我们原是与天主同工的。} \BibleRef{格前 3:9}\par

% structure: p0017 source=h2 target=subsection reason=relative-heading-rank; base=h2

\subsection{第六章——我们应如何理解天主对天使说话}

我们读到在造人时说：\Quote{我们造人}而不是\Quote{我造人}，是向天使说的，如果祂没有加\Quote{照我们的肖像}的话；但既然我们不能相信人是照天使的肖像造的，或天主的肖像与天使的肖像相同，那么将这话归于天主圣三的复数就是适当的。然而，这位圣三，虽是一位天主，在说了\Quote{我们造}之后，接着说：\Quote{天主就照自己的肖像造了人}\BibleRef{创 1:26}，而不是\Quote{众神造}或\Quote{照他们的肖像}。若将\Quote{来，我们下去混乱他们的语言}这话应用于天使有任何困难，我们可将复数归于圣三，好似圣父对圣子和圣神说话；但更恰当的是，天使以神圣的运动，即虔诚的思想，接近天主，并借此利用那在天庭中作为永恒法律统治的不变真理。因为他们自己不是真理；但分享创造性的真理，他们被引向它，如同生命之源，好使他们自己所没有的，能在它那里获得。而他们的运动是稳定的，因为他们从不退离已达成的。对这些天使，天主并不像我们彼此说话，或对天主、或对天使、或天使对我们、或天主通过他们对我们那样说话：祂以一种不可言喻的自己的方式对他们说话，而祂所说的以适合我们能力的方式传达给我们。因为那先于并超越祂所有工作的天主的说话，是祂工作的不变理性：它没有喧嚣和流逝的声音，而是永恒持久的，在时间中产生效果。因此祂对圣天使说话；但对远离的我们，祂以别的方式说话。然而，当我们用内心的耳朵听到天主言语的部分时，我们就近似于天使。但在这部作品中，我不必费力解释天主说话的方式。因为不变真理要么以某种无法描述的方式直接对理性受造物的心灵说话，要么通过可变的受造物说话，或向我们的灵呈现精神图像，或向我们的肉体感觉呈现肉体声音。\par

\Quote{凡他们想做的事，没有做不成的}\BibleRef{创 11:6}这些话肯定不是肯定句，而是反问句，如同威胁者所使用的那样，例如，当狄多喊道：\par

\Quote{他们不会拿起武器追击吗？}\par

我们应当理解这些话，好像是在说：凡他们想做的事，没有做不成的吗？从这三个人，即诺厄的三个儿子，我们意指73个，或如谱系所示，72个民族和同样多的语言分散在地上，随着他们的增多，甚至充满了岛屿。但民族的增多远多于语言的增多。因为即使在非洲，我们知道有几个野蛮民族只有一种语言；谁又能怀疑，随着人类的增加，人们设法乘船到达了岛屿呢？\par

% structure: p0022 source=h2 target=subsection reason=relative-heading-rank; base=h2

\subsection{第七章——即使最遥远的岛屿，其动物群是否来自通过洪水在方舟中保存的动物}

关于所有那些未被驯化、也不像青蛙那样从土中产生、而是由雌雄双亲繁殖的野兽——例如狼和此类动物——有一个问题被提出：洪水之后，所有未进入方舟的动物都灭绝了，那么这些动物如何能在岛屿上被发现呢？除非繁殖是从方舟中成对保存的那些动物恢复的。的确，可以说它们是通过游泳到达岛屿的，但这只适用于那些离大陆很近的岛屿；然而有些岛屿如此遥远，以至于我们想象没有任何动物能游过去。但若人捕捉它们并携带它们过海，从而在新居地繁殖这些种类，这并不会意味着对狩猎的难以置信的喜爱。同时，不可否认，借着天使的干预，它们可能由天主的命令或许可而被转移。然而，如果它们如起初受造时那样从土中产生，那时天主说：\BibleQuote{让地生出活物}，\BibleRef{创 1:24}，那么这就更明显地表明，所有种类的动物被保存在方舟中，与其说是为了更新物种，不如说是为了预表将在教会中得救的各个民族；如果地在许多动物无法横渡的岛屿上产生出许多动物，那么这一点就更为明显。\par

% structure: p0024 source=h2 target=subsection reason=relative-heading-rank; base=h2

\subsection{第8章——某些畸形人种是否源自亚当或诺厄子孙的血统。}

关于某些在世俗史中提到的畸形人种，是否源自诺厄的儿子们，或者我更应该说，源自那同一人（他们自身就是从那人传下的），这也是一个被追问的问题。因为据说有些人额头正中只有一只眼睛；有些人脚从脚跟向后翻；有些人具有双性，右胸像男人，左胸像女人，并且他们交替生育和分娩；有些人据说没有嘴，只通过鼻孔呼吸；有些人只有一肘高，因此被希腊人称为\OriginalTerm{矮人}{Pigmies}：他们说在某些地方，女性在五岁时受孕，并且活不过八岁。同样，他们讲述有一族人，有两只脚但只有一条腿，奔跑速度惊人，尽管他们不弯曲膝盖：他们被称为\OriginalTerm{影子脚人}{Skiopodes}，因为在天热时，他们仰面躺下，用脚遮荫。另一些人据说没有头，眼睛在肩膀上；还有其他人类或类似人类的种族，被描绘在\Place{迦太基}港口的广场的镶嵌画中，依据奇珍异闻的记载。关于\OriginalTerm{狗头人}{Cynocephali}，他们的狗头和吠叫声表明他们是野兽而非人类，我还能说什么呢？但我们不必相信所有我们听到的关于这些畸形物的说法。但是，无论何处出生的一个人，即一个理性的、有死的动物，无论他在颜色、动作、声音上呈现多么不寻常的外貌，也无论他在其本性的某一种能力、部分或性质上多么奇特，没有一个基督徒可以怀疑他来自那同一位原祖。我们能够区分共同的人性与那因其独特而奇异的特性。\par

对于个别情况中的怪胎，可给出与奇异种族相同的解释。因为创造万有的天主知道每件事物应当在何时、何地存在或曾被创造，因祂看到能够促进整体美的相似与差异。但看不见整体的人，因局部的畸形而感到不快，因为他看不见平衡它的部分以及它所属的整体。我们知道，有些人出生时手上或脚上长出超过四个手指或脚趾：这是一个较小的问题；但远离我们这种愚蠢，即以为创造主数错了人的手指数目，尽管我们无法解释这种差异。同样，在偏离常规更大的情况下，祂的作为无人能公正地指责，祂知道祂所做的。在\Place{希波-迪亚里图斯}有一个人，他的手呈新月形，每只手只有两根手指，他的脚也如此形成。如果有一个像他这样的种族，它就会被添加到奇妙历史中。难道我们因此就要否认这个人是出自那第一个被创造的人吗？至于所谓的\OriginalTerm{阴阳人}{Androgyni}或\OriginalTerm{雌雄同体者}{Hermaphrodites}，尽管罕见，但偶尔会出现性别如此可疑的人，以至无法确定他们从哪一性别得名；尽管习惯给他们一个男性化的名字，因为更有价值。因为从未有人称他们为雌雄同体者们。几年前，就在我记忆中，东方出生了一个人，上半身双倍，下半身单一——有两个头、两个胸腔、四只手，但一个身体和两只脚，像一个普通人；他活了很久，许多人有机会看到他。但谁能列举所有与确知父母明显不同的人类出生案例？因此，既然无人会否认这些都是从那一人而来的后裔，那么，所有那些被报道在身体外貌上偏离了自然通常或几乎普遍保持的常规的种族，如果它们被包含在\Quote{理性且可朽的动物}这一定义中，无疑都将其谱系追溯到那第一位万民之父。我们假设这些关于彼此不同且与我们相异的种族的传闻是真实的；但可能它们并非如此。因为如果我们不知道猿、猴和斯芬克司并不是人，而是兽类，那些历史学家可能会将它们描述为人的种族，并且毫无顾忌地炫耀他们虚假而虚荣的发现。但假设这些奇事所记载的是人，如果天主有意这样创造一些种族，使我们不至于认为出现于我们中间的怪胎是祂塑造人性的智慧的失败，如同我们谈论一个不那么完美的工匠的失败，那又如何？因此，我们不应觉得荒谬，正如在个别种族中有怪胎，在整个人类种族中也有奇异的种族。因此，谨慎而小心地结束这个问题：要么这些关于某些种族的说法根本不存在；要么如果存在，它们不是人类种族；要么如果它们是人类，它们是从亚当而来的。\par

% structure: p0027 source=h2 target=subsection reason=relative-heading-rank; base=h2

\subsection{第九章——我们是否应当相信对跖人？}

但关于存在\OriginalTerm{对跖人}{Antipodes}的传说，即在地球的另一侧，太阳在我们这里落下时，在他们那里升起，他们的双脚与我们的相对而行，这毫无可信的依据。事实上，并非根据历史知识，而是根据科学推测，肯定了这一说法，理由是地球悬浮于天空的穹窿之中，它的一侧与另一侧有同样多的空间；因此他们说，下面的部分也一定有人居住。但他们没有注意到，即使假设或科学证明世界是圆形球体的形式，也不意味着地球的另一侧没有水；即使没有水，也不立即意味着有人居住。因为圣经通过预言的应验证明了其历史陈述的真实性，它不会给出虚假信息；说有些人可能乘船穿越整个广阔海洋，从世界的这一侧到另一侧，从而那个遥远地区的居民也是从那个第一人而来的后裔，这太荒谬了。因此，让我们看看是否能在那些被记载为分为七十二个民族和同样多种语言的人类种族中找到那旅居于地上的天主之城。因为天主之城一直延续到洪水与方舟时期，并且通过诺亚儿子的祝福，证明仍然存在于他们中间，尤其是在长子闪那里；因为耶斐特得到这个祝福：他应该住在闪的帐棚里。\par

% structure: p0029 source=h2 target=subsection reason=relative-heading-rank; base=h2

\subsection{第十章——论闪的家谱，天主之城在其世系中得以保存，直至亚巴郎的时代。}

因此，为了展示洪水后的天主之城，正如洪水前通过由舍特降下的世代系列所展示的那样，有必要保存由闪降下的世代系列。正因如此，圣经在展示了作为巴比伦或\Quote{混乱}的属地之城后，转回到族长闪，并从他开始重述直到亚巴郎的世代，此外还注明每个父亲在哪个年岁生下了属于这谱系的儿子，以及他活了多少年。无疑，这应验了我所作的承诺，即应当显明为何论到厄贝尔的儿子时说：\BibleQuote{一个名叫培肋格，因为在他的时代大地分裂了。}\BibleRef{创 10:25} 因为若不是语言的多样性，我们还能从大地的分裂中理解什么呢？因此，圣经省略了闪的其他与此无关的儿子，给出了谱系，使这条线延续到亚巴郎，正如洪水前给出了由舍特到诺厄的谱系。因此，这一系列世代这样开始：\BibleQuote{以下是闪的后代：闪一百岁生阿帕革沙得，是在洪水后两年。闪生阿帕革沙得后，又活了五百年，生了儿女。} 同样，它记录了其余的，注明每个父亲在其年岁中生下了那属于延续到亚巴郎的谱系的儿子。它也写明此后他又活了多久，生儿育女，免得我们幼稚地认为所记的人就是全部的人，而能理解人口如何增长，以及闪的后代如何能繁衍如此广大的地区和王国；特别是亚述王国，尼努斯以辉煌的繁荣征服了周围的国家，并留给他的后代一个庞大而完全巩固的帝国，这个帝国持续了许多世纪。\newline{}\par

但为避免不必要的冗长，我们将不提及这谱系中每个成员活了多少岁，而只提及他们生继承者的年岁，以便我们计算从洪水到亚巴郎的年数，并同时留出余地，简要而粗略地讨论一些对我们的论证必要的其他事项。因此，洪水后第二年，闪一百岁时生阿帕革沙得；阿帕革沙得一百三十五岁时生卡南；卡南一百三十岁时生舍拉；舍拉同样年龄时生厄贝尔；厄贝尔活了一百三十四岁，生培肋格，在他的时代大地分裂了；培肋格活了一百三十岁，生勒伍；勒伍活了一百三十二岁，生色鲁格；色鲁格一百三十岁，生纳曷尔；纳曷尔七十九岁，生特辣黑；特辣黑七十岁，生亚巴郎，天主后来将他的名字改为\Emph{亚巴辣罕}。因此，根据\WorkTitle{武加大}或\WorkTitle{七十贤士译本}，从洪水到亚巴郎共有一千零七十二年。在希伯来文抄本中给出的年数少得多；对此，或者没有理由，或者给出的理由不太可信。\newline{}\par

因此，当我们在这些七十二个民族中寻找天主之城时，我们不能断言，当人类只有同一嘴唇，即同一种语言时，人类已经背离了真天主的敬拜，而真正的虔诚只保留在那些从闪经由阿帕革沙得直到亚巴郎的世代中；但从他们骄傲地建造一座通天塔（这是不虔敬高抬自己的象征）时起，恶人的城邦或团体就显现出来了。它以前是隐藏着，还是不存在；洪水后是否两个城都留存——虔诚的城在诺厄那两个受祝福的儿子及其后裔中，不虔诚的城在受诅咒的儿子及其后裔中（那位对抗主的英勇猎人即出自此系）——这不容易确定。因为可能——而且这更可信——在巴比伦建立之前，两个儿子的后裔中已有轻视天主的人，而含的后裔中也有敬拜天主的人。当然，没有一个种族曾从地球上被除灭。因为在两篇圣咏中也说：\BibleQuote{他们都偏离了正道，一同变为污秽；没有行善的，连一个也没有，} 我们进一步读到：\BibleQuote{所有作孽的人，难道都不知吗？他们吞吃我的百姓，如同吃饭，并不呼求主。} 那时已经有天主的子民。因此，\BibleQuote{没有行善的，连一个也没有} 这句话是针对人子说的，而不是针对天主之子说的。因为之前已经说过：\BibleQuote{天主从天上俯看人子，要看看有谁领悟，有谁寻求天主。} 接着的这些话表明，所有属于那按人而不按天主的城的人，都是被弃绝的。\newline{}\par

% structure: p0033 source=h2 target=subsection reason=relative-heading-rank; base=h2

\subsection{第11章——论人类最初使用的语言就是后来被称为希伯来语的那一种，来自厄贝尔，当语言混乱发生时，这种语言保留在他的家族中。\newline{}}

因此，正如全人类使用同一种语言并不能保证罪恶感染的人不存于人间——因为在洪水之前，言语也是统一的，但除了义人\Person{诺厄}一家之外，所有人都被洪水所灭——同样，当外邦人因更傲慢的无神而招致分散和语言混乱的惩罚，那无神者的城被称为混乱或\Place{巴比伦}时，仍然有\Person{厄贝尔}的家存在，其中保留了人类原始的言语。因此，正如我已经提及的，当列举\Person{闪}的儿子们——他们各自建立了一个民族——时，\Person{厄贝尔}被首先提及，尽管他是\Person{闪}的第五代后裔。并且因为，当其他种族被各自特有的语言分开时，他的家族保存了那种语言——它被合理地认为是人类共同的语言——正因如此，这语言此后被称为\OriginalTerm{希伯来语}{Hebrew}。因为那时有必要用一个专名来区分这语言与其他语言；然而，当只有一种语言时，它没有别的名称，只叫作人的语言或人类言语，因为全人类都只说这一种语言。有人会说：如果土地是在\Person{厄贝尔}的儿子\Person{培肋格}的日子里因语言分割的，那么那先前为所有人共有的语言理应以\Person{培肋格}命名。但我们应当理解，\Person{厄贝尔}亲自给他儿子取名为\Person{培肋格}，意为\OriginalTerm{分割}{Division}；因为他在土地分割之时——即正当分割时——出生，并且这正是那句话的意思：\Quote{在他的日子里，土地被分割了。}\BibleRef{创 10:25}因为除非\Person{厄贝尔}在语言增多时仍然活着，否则保存在他家族中的语言就不会以他命名。我们被引导相信，这便是原始的共同语言，因为语言的增多和变化是被作为惩罚引入的，且将豁免此惩罚归于天主子民是合宜的。\Person{亚巴郎}保留了这语言，而他不能将它传给所有后裔，只传给\Person{雅各伯}一系的人——他们独特而卓越地构成了天主子民，领受了祂的盟约，并按照肉身是\Person{基督}的先祖——这一点并非没有意义。同样，\Person{厄贝尔}本人也没有将那语言传给他的所有后裔，只传给了\Person{亚巴郎}所出的那一支。因此，虽然并未明确说明当恶人建造\Place{巴比伦}时还有虔诚的苗裔存留，这种模糊性是为了激发探究而非回避。因为当我们看到最初有一种共同语言，并且\Person{厄贝尔}在所有\Person{闪}的儿子中被首先提及——尽管他是\Person{闪}的第五代后裔——而且圣祖和先知们不仅在谈话中，而且在圣经的权威语言中所使用的语言被称为希伯来语；当我们被问及语言混乱之后，那种原始共同语言保存在何处时，既然毫无疑问，保存了这语言的人免除了它所代表的惩罚，我们还能提出什么其他建议呢？只能是：它幸存于以其命名的那个人的家族中，并且这家族的公义有一个不小的证明，即降在其他家族身上的惩罚没有落在它身上。\par

但还有另一个问题被提出：赫贝尔和他的儿子培肋格，如果他们只有同一语言，如何各自建立一个民族？毫无疑问，由赫贝尔经由亚巴郎而繁衍、并藉着他成为一大民族的希伯来民族，是一个民族。那么，如果赫贝尔和培肋格没有这样做，挪亚三个支系的所有儿子又如何被逐一列举为各自建立一个民族呢？极有可能的是，巨人尼默洛得也建立了他的民族，而圣经因其帝国和他身体的特殊庞大而单独提到他，因此七十二个民族之数保持不变。但培肋格被提及，不是因为他建立了一个民族（因为他的种族和语言是希伯来语），而是因为他出生在关键时期，那时全地正在分裂。我们也不该诧异巨人尼默洛得一直活到巴比伦建立和语言混乱发生的时期，以及随之而来的大地分裂。因为虽然赫贝尔是挪亚的第六代，而尼默洛得是第四代，但这并不妨碍他们可能同时活着。因为当世系较少时，他们活得更久且出生较晚；但当世系较多时，他们活得更短且更早来到世界。我们要明白，当大地分裂时，那些被记录为民族创始人的挪亚后裔不仅已经出生，而且已经到了拥有庞大家族的年龄，足以被称为部落或民族。因此我们绝不能认为他们是按照被记录的顺序出生的；否则，赫贝尔的另一个儿子、培肋格的兄弟约刻堂的十二个儿子，怎能已经建立了民族呢？因为约刻堂是照记录在他兄弟培肋格之后出生的，而大地是在培肋格出生时分裂的。因此我们要明白，虽然培肋格首先被命名，但他比约刻堂出生得晚很多，约刻堂的十二个儿子已经拥有庞大的家族，以至于他们可以因不同语言而分裂。最小的后出生者被首先命名并不稀奇：在挪亚的儿子中，雅弗的后代先被命名；然后是含的儿子，含是第二个儿子；最后是闪的儿子，闪是第一和长子。这些民族的名字，部分留存至今，以至于我们今天可以看出他们从谁而来，如亚述人来自亚述，希伯来人来自赫贝尔；但部分名字在时间流逝中已经改变，以至于最有学问的人，通过深入研究古代文献，也几乎无法发现这些民族的起源，我并非指所有，而是其中一些。例如，埃及人这个名字中毫无显示他们来自含的儿子米兹辣殷，埃塞俄比亚人这个名字也显不出与雇士的联系，尽管据说这些民族的起源如此。如果我们对整个名字作一番考察，就会发现改变的比保留的更多。\par

% structure: p0036 source=h2 target=subsection reason=relative-heading-rank; base=h2

\subsection{第十二章——论亚巴郎生命中的时代，从此神圣传承的新时期开始。}

现在让我们从圣祖亚巴郎的时代起审视天主之城的进展，从他的时代起，它开始变得更加显著，并且那些现在在基督内应验的神圣许诺被更充分地揭示。我们从圣经的暗示得知，亚巴郎出生于加色丁地，属于亚述帝国。那时，不敬神的迷信在加色丁人中如同在其他民族中一样盛行。亚巴郎所属的特辣黑家族是唯一保留了对真天主崇拜的家族，也是我们猜测的唯一保留希伯来语的家族；尽管农的儿子若苏厄告诉我们，这个家族在美索不达米亚也曾事奉其他神祇。\BibleRef{苏 24:2}赫贝尔的其他后裔逐渐被其他种族和语言所吸收。因此，正如挪亚一家在洪水中被保存以更新人类一样，在淹没全世界的迷信洪流中，只剩下特辣黑一家，其中保存了天主之城的种子。如同圣经在列出挪亚之前的世代及其年龄、并在天主开始对挪亚谈论建造方舟之前说明了洪水的原因时，说：\Quote{以下是挪亚的族谱；}同样，现在在列出从挪亚的儿子闪直到亚巴郎的世代之后，它又以一个时代为标志说：\Quote{以下是特辣黑的族谱：特辣黑生亚巴郎、纳曷尔和哈郎；哈郎生罗特。哈郎在他的出生地，加色丁人的乌尔，死在他父亲特辣黑面前。亚巴郎和纳曷尔都娶了妻子：亚巴郎的妻子名叫撒辣依；纳曷尔的妻子名叫米耳加，她是哈郎的女儿；哈郎是米耳加的父亲，也是依色加的父亲。}\BibleRef{创 11:27}这个依色加被认为就是亚巴郎的妻子撒辣。\par

% structure: p0038 source=h2 target=subsection reason=relative-heading-rank; base=h2

\subsection{第十三章——为何在特辣黑离开加色丁人迁往美索不达米亚的记述中，没有提到他的儿子纳曷尔。}

接下来叙述\Person{特辣黑}如何与他的家属离开了\Place{加色丁人}的地区，来到\Place{美索不达米亚}，居住在\Place{哈兰}。但没有提到他一个名叫\Person{纳曷尔}的儿子，似乎他没有带他同去。因为叙述这样写道：\Quote{\Person{特辣黑}带着他的儿子\Person{亚巴郎}，以及他孙子、\Person{哈兰}的儿子\Person{罗特}，和他的儿媳、他儿子\Person{亚巴郎}的妻子\Person{撒辣}，领他们出了\Place{加色丁人}的地区，要往\Place{客纳罕}地去；他到了\Place{哈兰}，就住在那里。}\BibleRef{创 11:31} \Person{纳曷尔}和他的妻子\Person{米耳加}在这里都没有提到。但后来，当\Person{亚巴郎}派他的仆人为他的儿子\Person{依撒格}娶妻时，我们发现这样写着：\Quote{那仆人从他主人的骆驼中取了十匹骆驼，并带了他主人所有的各样财物，起身往\Place{美索不达米亚}去，到了\Person{纳曷尔}的城。}\BibleRef{创 24:10} 这一点以及其他神圣历史的见证表明，\Person{亚巴郎}的兄弟\Person{纳曷尔}也已离开\Place{加色丁人}的地区，定居在\Place{美索不达米亚}，\Person{亚巴郎}曾与他的父亲住在那里。那么，为什么\OriginalTerm{圣经}{Scripture}没有提到他，当\Person{特辣黑}带着他的家属从\Place{加色丁}国出来住在\Place{哈兰}时，却提到他带了不仅他的儿子\Person{亚巴郎}，还有他的儿媳\Person{撒辣}和他的孙子\Person{罗特}呢？我们唯一能想到的理由是：或许他背离了他父亲和兄弟的虔诚，执著于\Place{加色丁人}的迷信，后来或因悔改，或因被当作可疑人物受迫害，才从那里迁移。因为在名为\WorkTitle{友弟德传}的书中，当以色列人的敌人\Person{何乐弗尼}询问那是怎样的民族，是否应该攻打他们时，\Place{阿孟人}的首领\Person{阿希约尔}这样回答他：\Quote{请我主现在听你仆人一句话，我要向你声明这真理，关于那住在你附近这山地的人民，你仆人嘴里决不说谎。因为这人民是由\Place{加色丁人}后裔，他们从前住在\Place{美索不达米亚}，因为他们不愿随从他们在\Place{加色丁}地荣耀的祖先的神，他们离开了祖先的道路，朝拜他们所认识的天上的天主；他们就把他们从诸神面前赶出，他们逃到\Place{美索不达米亚}，在那里住了多日。他们的天主对他们说：你们要离开你们的住所，往\Place{客纳罕}地去；他们就住了，}等等，正如\Place{阿孟人}\Person{阿希约尔}所叙述的。由此可见，\Person{特辣黑}一家因他们敬拜独一真天主的真正虔诚而遭受了\Place{加色丁人}的迫害。\par

% structure: p0040 source=h2 target=subsection reason=relative-heading-rank; base=h2

\subsection{第十四章——论\Person{特辣黑}在\Place{哈兰}完成他一生的年岁。}

论到\Person{特辣黑}在\Place{美索不达米亚}去世，据说他在那里活了二百零五年，天主向\Person{亚巴郎}所作的许诺现在开始显明；因为经上这样记载：\Quote{\Person{特辣黑}在\Place{哈兰}的日子是二百零五年，他在\Place{哈兰}死了。}\BibleRef{创 11:32} 这不应被理解为他在那里度过了他所有的日子，而是说他在那里完成了他一生的岁月，即二百零五年；否则就无法知道\Person{特辣黑}活了多久，因为没有说他是在生命的哪一年来到\Place{哈兰}的；而且，假定在这一系列世系中（其中仔细记录每人活了多少年），唯独他的年龄没有被记载，这是荒谬的。因为虽然有些被同一\OriginalTerm{圣经}{Scripture}提及的人没有记录年龄，但他们不在这一系列中，这一系列的时间计算是通过父母的去世和子女的继承不断显示的。因为这一系列从\Person{亚当}到\Person{诺厄}，从他再到\Person{亚巴郎}，是按顺序记载的，没有一个人没有他年岁的数目。\par

% structure: p0042 source=h2 target=subsection reason=relative-heading-rank; base=h2

\subsection{第十五章——论\Person{亚巴郎}迁移的时间，他依照天主的命令，离开了\Place{哈兰}。}

在记载了\Person{亚巴郎}的父亲\Person{特辣黑}去世之后，我们接着读到：\BibleQuote{主对\Person{亚巴郎}说：\InnerQuote{离开你的故乡、你的家族和你的父家}}\BibleRef{创 12:1}等，我们不应因为这在叙事顺序中是接着的，就认为它也接着事件的时间顺序。因为如果真是这样，就会出现一个无法解决的难题。因为在主对\Person{亚巴郎}说了这些话之后，圣经记载：\BibleQuote{\Person{亚巴郎}就照主所吩咐的出发了，\Person{罗特}也与他同行。\Person{亚巴郎}离开\Place{哈兰}时已七十五岁。}\BibleRef{创 12:4}如果他是在父亲死后才离开\Place{哈兰}，这怎么可能呢？因为如上文所暗示，\Person{特辣黑}七十岁时生了\Person{亚巴郎}；如果我们把这个数字加上\Person{亚巴郎}离开\Place{哈兰}时的七十五岁，就得到一百四十五年。因此，\Person{亚巴郎}离开那座美索不达米亚的城市时，\Person{特辣黑}的年龄就是那么大；因为\Person{亚巴郎}到了七十五岁，而他的父亲（在七十岁时生了他）如前所述，已到一百四十五岁。因此，他并非在父亲死后（即父亲活到二百零五岁之后）离开那里；而是他离开那地方的那一年——即他七十五岁那年——无疑是他父亲的（他在七十岁时生了他）一百四十五岁。因此应当理解，圣经按照其惯常做法，回到了叙事已经越过的时间；正如上文，在提到\Person{诺厄}的孙子们时，说他们各有自己的国家和语言；但后来，好像这也按时间顺序接着，它又说：\BibleQuote{全世界只有一种口音，一种语言。}\BibleRef{创 11:1}那么，如果全世界只有一种语言，他们又怎能说各有自己的国家和语言呢？除非是因为叙事回到前面，把越过的事情重新拾起。同样，在说了\Quote{\Person{特辣黑}在\Place{哈兰}的日子共二百零五年，\Person{特辣黑}死在\Place{哈兰}}之后，圣经回到前面所越过的地方，以完成关于\Person{特辣黑}已开始的内容，说：\BibleQuote{主对\Person{亚巴郎}说：\InnerQuote{离开你的故乡}}\BibleRef{创 12:1}等。在这些主的话之后，又加上：\BibleQuote{\Person{亚巴郎}就照主所吩咐的出发了，\Person{罗特}也与他同行。\Person{亚巴郎}离开\Place{哈兰}时已七十五岁。}因此，这事发生在他父亲一百四十五岁时；因为那时正是他自己的七十五岁。但这个问题也可以用另一种方式解决：\Person{亚巴郎}离开\Place{哈兰}时的七十五岁，是从他从加色丁人的火中被救出的那一年算起，而不是从他出生那年算起，好像他当时才出生似的。\par

有福的\Person{斯德望}在《宗徒大事录》中叙述这些事时说：\BibleQuote{光荣的天主在\Place{美索不达米亚}，在\Person{亚巴郎}住在\Place{哈兰}之前，显现给我们的祖先\Person{亚巴郎}，对他说：\InnerQuote{离开你的故乡、你的家族和你的父家，往我将要指示你的地方去。}}\BibleRef{宗 7:2}根据\Person{斯德望}的这些话，天主对\Person{亚巴郎}说话，并非在他父亲死后——他父亲确实死在\Place{哈兰}，他的儿子也与他同住那里——而是在他住在那个城市之前，尽管他当时已身在\Place{美索不达米亚}。因此，他已经离开了加色丁人。因此，当\Person{斯德望}补充说：\BibleQuote{那时\Person{亚巴郎}离开了加色丁人的地方，住在\Place{哈兰}}\BibleRef{宗 7:4}，这并非指天主对他说话之后发生的事（因为他并非在天主说了这些话之后才离开加色丁人的地方，因为他说天主是在\Place{美索不达米亚}对他说话的），而是他用的\Quote{那时}一词，指的是从他离开加色丁人的地方到住在\Place{哈兰}的整个时期。同样在接下来的话中：\Quote{从此以后，他父亲死后，天主便将他安置在这地方，就是你们如今所住、你们祖先所住的地方}，他并没有说，他父亲死后他从\Place{哈兰}出去；而是说，他父亲死后，他从此被安置在这里。因此应当理解，天主在\Place{美索不达米亚}对\Person{亚巴郎}说话时，他尚未住在\Place{哈兰}；但他记着天主的诫命，与父亲一同来到\Place{哈兰}，并在自己七十五岁（即他父亲一百四十五岁）时离开那里。但他说，他定居在\Place{客纳罕}地（而非从\Place{哈兰}出发）发生在他父亲死后；因为当他购买那块地并亲自取得所有权时，他的父亲已经去世。但是，当他已定居在\Place{美索不达米亚}，即已离开加色丁人的地方时，天主说：\BibleQuote{离开你的故乡、你的家族和你的父家}\BibleRef{创 12:1}，这并不意味着他应当从那里拔身（因为他已经这样做了），而是他应当割舍自己的心。因为他若仍被返回的希望和渴望所束缚，他就还没有从那里在思想上离开——这种希望和渴望，藉着天主的命令和协助，以及他自己的服从，将被断绝。确有一种并非不可信的设想，即后来当\Person{纳曷尔}跟随他的父亲时，\Person{亚巴郎}才履行了主的诫命，带着他的妻子\Person{撒辣}和侄儿\Person{罗特}离开\Place{哈兰}。\par

% structure: p0045 source=h2 target=subsection reason=relative-heading-rank; base=h2

\subsection{第十六章——论天主对\Person{亚巴郎}所许诸诺言的次序与性质。}

现在要考虑天主向亚巴郎所作的应许；因为在这些应许中，我们天主的神谕，即真天主的神谕，关于那虔诚的民族——先知权威所预言过的——开始更公开地显现。第一个应许记载如下：\BibleQuote{主对亚巴郎说：\Quote{离开你的故乡、你的家族和你的父家，到我将指示你的地方去。我要使你成为一个伟大的民族，我要祝福你，使你出名；你将成为祝福。我要祝福那祝福你的人，咒骂那咒骂你的人；地上万民都要因你获得祝福。}}\BibleRef{创 12:1}现在应注意，有两个东西应许给了亚巴郎：一个是他的后裔将占有客纳罕地，这话暗示着：\BibleQuote{到我将指示你的地方去，我要使你成为一个伟大的民族}；但另一个更卓越的应许，不是关于\OriginalTerm{肉身的后裔}{carnal seed}，而是关于\OriginalTerm{属神的后裔}{spiritual seed}，藉此他不仅是唯一的以色列民族的父亲，而是所有追随他信德足迹的万民的父亲，这应许最初是以下列话语作出的：\BibleQuote{地上万民都要因你获得祝福}。欧瑟伯认为这个应许是在亚巴郎七十五岁时作出的，好像应许之后不久亚巴郎就离开了哈兰，因为圣经不能自相矛盾，其中我们读到：\BibleQuote{亚巴郎离开哈兰时已七十五岁}。但如果这应许是在那一年作出的，那么亚巴郎当然是和他父亲一起住在哈兰；因为他若不先住在那里，就无法离开。那么，这是否与斯德望所说的相矛盾呢？\BibleQuote{荣耀的天主曾显现给我们的父亲亚巴郎，当时他在美索不达米亚，住在哈兰之前。}\BibleRef{宗 7:2}但应理解为这一切都发生在同一年——既包括天主在亚巴郎住在哈兰之前所作的应许，也包括他住在哈兰，以及他从那里离开——这不仅因为欧瑟伯在\WorkTitle{编年史}中从这应许的年份算起，并指出430年后出了埃及，那时法律被赐予，而且因为宗徒保禄也提到了这一点。\par

% structure: p0047 source=h2 target=subsection reason=relative-heading-rank; base=h2

\subsection{第17章 论万国中三个最著名的王国，其中之一，即亚述王国，在亚巴郎出生时已非常显赫。}

在同一时期，有三个著名的万国王国，其中\OriginalTerm{属地之城}{city of the earth-born}，即按人而生活的人之社会，在堕落天使的统治下，特别兴盛，即\Place{西基昂}、\Place{埃及}和\Place{亚述}三个王国。其中，亚述最为强大和崇高；因为那位\Person{贝鲁斯}之子\Person{尼努斯}王，已征服了除\Place{印度}外的所有亚洲民族。现在我所说的亚洲，不是指作为这更大亚洲之一省的那部分，而是被称为全亚洲的地区，有些人将它定为世界的一半，但多数人定为世界的三分之一——即亚洲、欧洲和非洲，由此造成不等的划分。因为被称为亚洲的部分从南经东直到北；欧洲从北直到西；非洲从西直到南。由此可见，欧洲和非洲这两个包含世界的一半，而亚洲独自包含另一半。而这两部分是由这样的情形形成的：从大洋进入它们之间所有的地中海海水，形成了我们这个大海。因此，如果你将世界分为两部分，东和西，亚洲将在一部分，欧洲和非洲在另一部分。因此，在那时著名的三个王国中，一个，即西基昂，并不在亚述人统治之下，因为它位于欧洲；至于埃及，它怎能不臣服于那个统治除印度外整个亚洲的帝国呢？因此，在亚述，不虔敬之城的统治占据优势。它的首都是\Place{巴比伦}——一个属地之城，命名极为恰当，因为它的意思是混乱。在那里，\Person{尼努斯}在他父亲\Person{贝鲁斯}死后统治，贝鲁斯最初在那里统治了六十五年。他的儿子\Person{尼努斯}，在父亲死后继承王国，统治了五十二年，并在亚巴郎出生时已为王四十三年，那大约是\Place{罗马}建城前1200年，仿佛西方的另一个巴比伦。\par

% structure: p0049 source=h2 target=subsection reason=relative-heading-rank; base=h2

\subsection{第18章 论天主对亚巴郎的再次说话，其中祂应许将\Place{客纳罕}地赐给他和他的后裔。}

亚巴郎于是从他七十五岁、他父亲一百四十五岁时离开\Place{哈兰}，带着他兄弟的儿子\Person{罗特}和他的妻子\Person{撒辣}，进入\Place{客纳罕}地，甚至来到\Place{舍根}，在那里他再次接受了神圣的神谕，记载如下：\BibleQuote{主显现给亚巴郎，对他说：我要将这地赐给你的后裔。}\BibleRef{创 12:7}这里没有应许关于那使他成为\OriginalTerm{万民之父}{father of all nations}的后裔，而只是关于那使他成为唯一的以色列民族之父的后裔；因为藉着这后裔，那地才被占有。\par

% structure: p0051 source=h2 target=subsection reason=relative-heading-rank; base=h2

\subsection{第19章 论天主在\Place{埃及}保全\Person{撒辣}的贞洁，当时亚巴郎称她为妹妹而非妻子。}

在那里筑了一座祭台，呼求了天主之后，亚巴郎从那里前行，住在旷野，并因饥荒的逼迫不得不进入\Place{埃及}。在那里他称他的妻子为妹妹，没有说谎。因为她确实也是妹妹，由于血缘相近；正如\Person{罗特}，由于同样的亲近关系，是他兄弟的儿子，被称为他的兄弟。他并没有否认她是他的妻子，只是保持沉默，将他妻子贞洁的护卫交托给天主，并作为人防备人的诡计；因为他若没有尽所能地防备危险，就是在试探天主而非信赖祂。关于此事，针对摩尼教徒\Person{福斯图斯}的诽谤，我们已经说得足够多了。最后，亚巴郎所期望主做的事发生了。因为埃及王\Person{法郎}曾将她娶为妻子，但在遭受严重打击后，将她归还给她的丈夫。我们远不能相信她因与他人同寝而受玷污；因为更可信的是，由于这些巨大的苦难，法郎没有被允许这样做。\par

% structure: p0053 source=h2 target=subsection reason=relative-heading-rank; base=h2

\subsection{第二十章——论罗特与亚巴郎分离，并未破坏爱德。}

亚巴郎从埃及返回他先前所在之地时，他兄弟的儿子罗特与他分离，往索多玛地方去了，并未破坏爱德。因为他们都已富有，开始有许多牧放牲畜的人；当这些牧人彼此争吵时，他们就这样避免了家族内的争斗不和。的确，按人事常情，这种起因甚至可能导致他们彼此之间的冲突。因此，亚巴郎预防这恶事，对罗特说了这些话：\Quote{你我之间不可有争端，我的牧人和你的牧人之间也不可；因为我们是弟兄。看，全地不都在你面前吗？请你离开我：你若向左，我便向右；你若向右，我便向左。} \BibleRef{创 13:8} 或许由此便在人中产生了一种和平的习俗：在分割地上财物时，年长者应作划分，年幼者应作选择。\par

% structure: p0055 source=h2 target=subsection reason=relative-heading-rank; base=h2

\subsection{第二十一章——论天主的第三个诺言，借此祂允诺将客纳罕地永赐给亚巴郎和他的后裔。}

当亚巴郎与罗特分离，为了维持家计而彼此分居，并非出于卑劣的不和时，亚巴郎住在客纳罕地，罗特住在索多玛。主在第三个神谕中对亚巴郎说：\Quote{请你举起眼来，从你现在所在的地方，向北方、向非洲、向东方、向海观看；因为你看见的一切地，我都要赐给你和你的后裔，直到永远。我要使你的后裔如同地上的尘土一样；如果谁能够数算地上的尘土，你的后裔也能数算。起来，纵横走遍这地，因为我必赐给你。} \BibleRef{创 13:14} 这诺言是否也包括使他成为万民之父，并不明确。因为\Quote{我要使你的后裔如同地上的尘土一样}这句话，似乎是指这一点，这句话用的是希腊人所谓的\OriginalTerm{夸张法}{hyperbole}，这确实是比喻说法，而非字面意思。但明白的人不会怀疑圣经如何使用此及其他的修辞。因为那种修辞（即说话方式）用于所说的大于所意指的；谁不知道，地上的尘土数量比从亚当直到世界末日所有的人数要大得无可比拟？那么它比亚巴郎的后裔大多少呢？这不仅指属于以色列民族的后裔，也指按照信德的榜样存在于全世界万民中现今和将来的后裔。因为那个后裔与众多恶人相比确实微小，尽管这些少数人本身也构成了不可胜数的群体，通过夸张法被比作地上的尘土。诚然，应许给亚巴郎的那群众在天主前并非不可数，对人而言则不可数；但甚至地上的尘土在天主面前也不是不可数的。此外，这里的诺言不仅可以理解为关于以色列民族，也可理解为关于亚巴郎的整个后裔，这后裔因数量众多而恰当地比作尘土，因为有关它也有多子多孙的诺言，不是按肉躯，而是按精神。但我们之所以说这不明确，是因为甚至那一个民族（即按肉躯从亚巴郎经其孙雅各伯所生的民族）也已经增长到几乎充满世界各地。因此，就连它本身也可以通过夸张法因数量之多被比作尘土，因为仅它自己也是人无法计数的。当然，没有人怀疑这里所指的地只是称为客纳罕的那地。但是\Quote{我要赐给你和你的后裔直到永远}这句话，可能会使一些人心生疑虑，如果他们理解\Quote{直到永远}是\Quote{直到永恒}。但如果在这里他们将\Quote{直到永远}理解为像我们坚定相信的那样，即来世的开端要从今世的终结开始安排，那么就没有困难了，因为虽然以色列子民被逐出耶路撒冷，但他们仍留在客纳罕地的其他城市，并且要留到最后；当那整个土地被基督徒居住时，他们也是亚巴郎的真正的后裔。\par

% structure: p0057 source=h2 target=subsection reason=relative-heading-rank; base=h2

\subsection{第二十二章——论亚巴郎战胜索多玛的敌人，救出罗特并受到默基瑟德司祭的祝福。}

领受了这应许的神谕后，亚巴郎迁移，住在那地的另一处，即玛默勒橡树旁，即赫贝龙。后来索多玛被侵犯，五王与四王交战，罗特与战败的索多玛人一起被掳，亚巴郎率领三百一十八名家生的仆人与敌作战，救出了罗特，为索多玛王赢得了胜利，但当索多玛王将战利品给他时，他分文不取。那时，他公开受到默基瑟德的祝福，他是至高天主的司祭，关于他，在题为《希伯来书》的书信中记载了许多重要的内容，多数人认为是宗徒保禄所写，但也有人否认。因为那时首次出现了如今基督徒在全世界向天主奉献的祭献，并且实现了此后许久先知对那将要在肉躯中来临的基督所说的话：\Quote{你按照默基瑟德的品位，永为司祭。} \BibleRef{咏 110:4}——即不是按照亚郎的品位，因为那品位在阴影所预示的事物显现时将被废除。\par

% structure: p0059 source=h2 target=subsection reason=relative-heading-rank; base=h2

\subsection{第二十三章——论天主对亚巴郎的话语，应许他的后裔将如繁星之多；亚巴郎相信这话，尚在未受割损之时就被称为成义。}

主的话也藉异象传给了亚巴郎。因为当天主应许赐给他保护并极大的赏报时，他关切后裔，就说有一个大马士革的\Person{厄里厄则尔}，生在他家中，将作他的继承人。他立即被应许得到后裔，不是那家中生的仆人，而是要出自亚巴郎本身的后裔；并且后裔不可胜数，不像地上的尘土，而像天上的星辰——在我看来，这似乎是应许后裔在属天的福乐中被高举。因为就数量而言，天上的星辰与地上的尘土相比又如何？除非有人说这比喻在于星辰也无法数算？因为并不是相信所有星辰都能被看见。人观察得越敏锐，看到的就越多。因此，可以认为有些星辰即使最敏锐的观察者也看不见，更不用说那些据称在地球最遥远另一端升起并落下的星辰了。最后，这本书的权威谴责了那些像\Person{阿拉托斯}或\Person{欧多克索斯}的人，以及其他夸耀自己发现并记录了全部星辰数目的人。的确，此处记载了宗徒为了赞扬天主的恩宠而引用的那句话：\BibleQuote{亚巴郎相信了天主，天主就以此算为他的正义}，免得受割损者自夸，不愿接纳未受割损的外邦人进入基督的信德。因为当他相信，且他的信德被算为正义时，亚巴郎还未受割损。\par

% structure: p0061 source=h2 target=subsection reason=relative-heading-rank; base=h2

\subsection{第二十四章．论\Person{亚巴郎}奉命献祭的意义，当他恳求教导关于他所相信之事时。}

同样在异象中，天主对他说话说：\BibleQuote{我是领你出\Place{加色丁}地区的天主，要将这地赐给你为产业。}\BibleRef{创 15:7} 当亚巴郎问他如何知道必得这地为产业时，天主对他说：\BibleQuote{你取一头三岁的母牛，一只三岁的母山羊，一只三岁的公绵羊，一只斑鸠，一只雏鸽。亚巴郎就把这些都取了来，从中间剖开，将每块对着另一块摆列；只是没有剖开飞鸟。有鸷鸟落下来，}正如所记载的，\BibleQuote{落在肉块上，亚巴郎就坐在旁边。太阳快要西沉时，亚巴郎深感恐惧，看，有大黑暗落在他身上。上主对亚巴郎说：你当知道，你的后裔必在异邦寄居，他们要奴役他们，压迫他们四百年。但我必惩罚他们所事奉的民族；之后他们必带着大量财物出来。至于你，你要平安归到你的祖先那里，安享高年。到了第四代，他们必回到这里，因为阿摩黎人的罪恶还没有满盈。日落之时，有火焰，有冒烟的火炉，有火把从那些肉块间经过。在那一天，上主与亚巴郎立约，说：我要将\Place{埃及河}直到\Place{幼发拉底大河}之地赐给你的后裔：就是刻尼人、刻纳次人、卡德摩尼人、赫特人、培黎齐人、勒法因人、阿摩黎人、客纳罕人、希威人、基尔加斯人和耶步斯人。}\BibleRef{创 15:9}\par

所有这些事都是在来自天主的异象中说和做的；但若要详细准确地处理它们，将会花费很长时间，并且超出本书的范围。我们只需知道，在据说亚巴郎相信了天主，天主就以此算为他的正义之后，他并没有因说\Quote{主天主，我凭什么知道我将继承它呢？}而欠缺信德；因为那地的继承是应许给他的。他并非说\Quote{我如何知道}，好像他还不相信；而是说\Quote{我凭什么知道}，意思是可能给出一个标志，使他能够知道他所相信的那些事的方式，正如童贞玛利亚并非因缺乏信德而说\BibleQuote{这事怎能成就？因为我不认识男人}\BibleRef{路 1:34}；因为她询问的是她确信必将发生之事的方式。当她这样问时，她被告知\BibleQuote{圣神要临于你，至高者的能力要庇荫你}\BibleRef{路 1:35}。此处最终也给出了一个象征，包括三只动物：一头母牛、一只母山羊和一只公绵羊，以及两只鸟：一只斑鸠和一只鸽子，使他能够知道他所不怀疑必将发生的事将要按照这个象征成就。因此，无论是母牛是人民将被置于法律之下的标记，母山羊是同一人民将成为有罪的标记，公绵羊是他们将统治的标记（这些动物被称为三岁，理由是有三个显著的时期划分：从亚当到诺厄，从诺厄到亚巴郎，从亚巴郎到达味；在撒乌耳被弃绝后，达味按主的旨意首先被立为以色列民族之王：在这第三个时期，即从亚巴郎到达味，这个民族成长起来，仿佛经历了生命的第三阶段），还是它们有其他更合适的含义，我仍然毫不怀疑，属灵的事物由它们以及斑鸠和鸽子所预表。经上说\Quote{但鸟他没有劈开}，因为属血肉的人在他们自己中间是分裂的，但属灵的人却毫不分裂，无论是像斑鸠那样远离人群的繁忙交谈，还是像鸽子那样住在他们中间；因为这两种鸟都是单纯无害的，表明即使在将要获得那土地的以色列人民中，也会有作为应许的子女和将要留在永恒幸福中的国度继承者。但是，飞禽降在被分开的肉块上并不代表好事，而是代表这空中的灵，在属血肉者的分裂中为自己寻求某种食物。但亚巴郎坐在它们中间，表明即使在属血肉者的分裂中，真正的信徒也将坚持到底。至于日落时，大恐惧和深沉的黑暗降临在亚巴郎身上，表明在世界末期，信徒将处于极大的扰乱和磨难中，关于这磨难，主在福音中说\BibleQuote{那时必有大灾难，从起初直到如今没有这样的}\BibleRef{玛 24:21}。\par

但是对亚巴郎所说的\BibleQuote{你确实要知道，你的后裔要在不是他们自己的地方作客，并且他们要奴役他们，折磨他们四百年}\BibleRef{创 15:13}清楚地是关于以色列人民将要在埃及受奴役的预言。并不是说这个人民要在埃及人的压迫下受奴役四百年，而是预言这件事将发生在那些四百年期间。因为正如论到亚巴郎的父亲特辣黑所记载的\BibleQuote{特辣黑在哈兰的日子是二百零五年}\BibleRef{创 11:32}，并不是因为他全部时间都在那里度过，而是因为在那里完成，所以这里也说\Quote{他们要奴役他们，并折磨他们四百年}，是因为那个数字完成了，而不是因为全部时间都在那种折磨中度过。这四百年是约整数，虽然它们稍微多一点——无论你是从此时，当这些事应许给亚巴郎时起算，还是从依撒格的出生，即亚巴郎的后裔，这些事所预言的，起算。因为，正如我们上面已经说过的，从亚巴郎七十五岁，当第一个诺言赐给他时，直到以色列从埃及出离，总共算四百三十年，宗徒这样提到：\BibleQuote{我是说，天主所确认的盟约，那四百三十年后制定的法律，不能废除，以致应许失效}\BibleRef{迦 3:17}。所以这四百三十年可以称为四百年，因为它们相差不大，尤其因为当这些事在异象中显示并告诉亚巴郎时，或者当依撒格在他父亲一百岁时出生，即第一次应许后二十五年，那时这四百三十年还剩四百零五年，天主乐意称之为四百年。没有人会怀疑随后在先知性天主话语中的其他事情属于以色列人民。\par

当又加上\BibleQuote{日落的时候，有火焰，看，一个冒烟的炉和烧着的火把，从那些肉块之间经过}\BibleRef{创 15:17}，这象征着在世界末日，属血肉的人将受火的审判。正如天主之城以前从未有过的磨难，预计在假基督之下发生，由亚巴郎在日落时的深沉的黑暗所象征，即当世界末日临近时——同样，在日落时，即在世界的最末日，那火象征审判的日子，它将那些要经由火得救的属血肉者与那些要在火中被定罪的分开。然后与亚巴郎所立的盟约特别陈述了客纳罕地，并从中命名了十一个支派，从埃及河直到幼发拉底河。这不从那大埃及河，即尼罗河，而是从一条分隔埃及和巴勒斯坦的小河，那里有\Place{Rhinocorura}城。\par

% structure: p0066 source=h2 target=subsection reason=relative-heading-rank; base=h2

\subsection{第25章——论撒辣的使女哈加耳，撒辣愿意她作亚巴郎的妾。}

以下是亚巴郎众子的时期，一个是婢女哈加尔所生，另一个是自由妇人撒辣所生。关于此事，我们已在上一卷中谈论过。针对此事，亚巴郎绝不能被指责为对这婢女有罪，因为他与她同房是为了生育后代，不是为了满足情欲；他并非要侮辱妻子，而是要听从妻子，因妻子以为若能利用婢女多产的子宫来弥补自己本性的缺陷，便能缓解自己的不育；并且按照宗徒所说的那条法则：\BibleQuote{同样，丈夫也没有权力掌管自己的身体，而是妻子。}\BibleRef{格前 7:4} 他身为丈夫，在自己不能生育的情况下，可以借助他人为妻子生子。这里没有放纵的情欲，没有污秽的淫邪。婢女是由妻子交给丈夫以生育后代，丈夫为生育后代而接收她，双方追求的都不是有罪的纵欲，而是自然的果实。当怀孕的婢女轻视她不生育的主母时，撒辣出于女性的嫉妒将过错归咎于丈夫；即便如此，亚巴郎仍显出他不是奴性的情人，而是自由的生育者；他使用哈加尔时，保全了妻子撒辣的贞洁，顺从了她的意愿而非自己的意愿——他接受她并非出于寻求，亲近她并非出于留恋，使她怀孕并非出于爱她——因为他说：\BibleQuote{看，你的婢女在你手中；你随意待她吧！}\BibleRef{创 16:6} 他是一个能以男人应有的方式对待女人的人——对妻子有节制，对婢女顺从，不纵情于任何人。\par

% structure: p0068 source=h2 target=subsection reason=relative-heading-rank; base=h2

\subsection{第二十六章．论天主对亚巴郎的证明，借此向他保证，当他已年老时，将由不生育的撒辣得一个儿子，并立他为万民之父，又借着割损圣事，印下他对恩许的信德。}

这些事以后，依市玛耳由哈加尔出生。亚巴郎可能以为天主所应许他的话，就在这儿子身上应验了；当初他想立家中仆人为嗣时，天主曾说：\BibleQuote{这人不会是你的继承人；那要由你亲生的，才是你的继承人。}\BibleRef{创 15:4} 因此，为免他以为恩许已在婢女之子身上应验，\BibleQuote{当亚巴郎九十九岁时，天主显现给他，对他说：我是全能的天主；你当在我的面前行走，作个齐全的人，我要与你立约，并要大大增加你的后裔。}\par

这里有更清晰的关于因依撒格召叫万民的恩许，依撒格即是恩许之子，借此表明恩宠而非本性；因为儿子是由老人和不生育的老妇人应许的。因为虽然天主甚至促成生育的自然过程，但在自然力衰败或失效之处，天主的作为明显可见，恩宠就更加清楚地被认出。而此事要成就，不是借着生育，而是借着再生，所以在撒辣应许得子时，割损礼就被命令了。并且命令所有儿子，不仅是亲生的，也包括家生和购买的仆人，都要受割损，以此证明这恩宠属于众人。因为割损礼若非除旧更新，又意味着什么？第八日若非指基督，他在一周完结、即安息日之后复活，又意味着什么？父母的名字都改变了：一切都宣告新意，新约已在旧约中预表。因为旧约一词除了隐藏新约，还意味着什么？而新约一词除了揭示旧约，又意味着什么？亚巴郎的笑，是欢跃的笑，不是不信的嘲笑。他心中所说的话：\BibleQuote{百岁的人还能生子吗？撒辣已九十岁，还能生育吗？}不是怀疑的话，而是惊奇的话。当经上说：\BibleQuote{我要将你侨居之地，即客纳罕全地，赐给你和你的后裔，作为永久的产业。}如果有人疑惑这话是已经应验，还是仍要期待，因为任何民族的尘世产业都不可能永久，那么，要知道，我们的作者翻译为\Quote{永久}的那个词，就是希腊人所说的\OriginalTerm{永恒}{\GreekText{αἰών}}。\par

\SourceRecord{Translated by Marcus Dods.；Nicene and Post-Nicene Fathers, First Series；Vol. 2.；Edited by Philip Schaff.；Buffalo, NY: Christian Literature Publishing Co.,；1887.；<http://www.newadvent.org/fathers/120116.htm>.}

% structure: p0001 source=h1 target=section reason=page-title

\section{天主之城（卷十七）}

本书追溯天主之城的历史，从撒慕尔经达味直至基督的君王与先知时期；记载于\WorkTitle{列王纪}、\WorkTitle{圣咏集}及撒罗满的著作中的预言，皆被解释为关于基督与教会。\par

% structure: p0003 source=h2 target=subsection reason=relative-heading-rank; base=h2

\subsection{第一章　论先知时代}

藉天主的恩宠，我们已详细论述了祂对亚巴郎所作的许诺：即按血统的以色列民族和按信德的所有民族，都将是他的后裔；而天主之城按时间顺序前进，将指出这些许诺如何应验。因此，在前一卷中我们已讲到达味王朝，现在将从同一王朝开始，按本书宗旨所需，论述余下部分。从圣撒慕尔开始预言的时候起，直至以色列人被掳往巴比伦，并依照圣耶肋米亚的预言，在七十年后从那里回归，重建天主圣殿，这整个时期乃是先知时代。因为，虽然诺厄本人（其时日洪水毁灭了全地）以及他前后直至天主之民开始有君王时的其他人，因他们所预言或预表了关于天主之城和天国的某些事，并非不当被称为先知，尤其我们读到其中有些人如亚巴郎和梅瑟，曾被明确称为先知；然而，从撒慕尔开始预言的时候起，才最显着地被称为先知时期。撒慕尔奉天主之命首先傅油立撒乌耳为王，其后当撒乌耳被弃时，又立达味为王，而达味的后裔应相继继位，直到适当的时候。因此，若我想重述先知们关于基督所预言的一切，而同时天主之城以其成员不断生死相继，在那段时期中运行，这部著作将漫无边际。首先因为圣经本身，即使在按次序论述君王及其事迹和统治时，似乎以历史的详尽来叙述所发生的事，但若藉天主圣神的光照来考量它所处理的内容，就会发现它旨在预言未来，其程度不少于甚至多于叙述过去。凡是稍加思考的人，岂不知道这要何等费力和冗长，需要多少卷帙才能彻底查考并论证清楚？其次，因为那些无可争议地属于预言的内容，关于基督和天国（即天主之城）的预言如此之多，要解释它们所需讨论的篇幅，将超出本著作的适当比例。因此，我将尽己所能约束自己，在完成本著作时，赖天主助佑，既不赘言，也不遗漏必要之事。\par

% structure: p0005 source=h2 target=subsection reason=relative-heading-rank; base=h2

\subsection{第二章　天主的关于客纳罕地的许诺何时实现，连属血肉的以色列也得到了那地}

在前一卷中我们说过，天主对亚巴郎的许诺从起初便包含两件事：其一，他的后裔将拥有客纳罕地，这由以下的话暗示：\Quote{你要进到我要指示给你的地方，我要使你成为一个大民族；}\BibleRef{创 12:1} 但另一件更为卓越，关乎那属神的而非属血肉的后裔，亚巴郎藉此成为万民之父，不仅是以色列一个民族，而是所有追随他信德足迹的人；这许诺开始由以下的话应许：\Quote{地上万民都要因你获得祝福。}\BibleRef{创 12:3} 此后我们又用许多其他证据证明这两件事是应许的。因此，亚巴郎按血统的后裔，即以色列民族，早已在应许之地；在那里，不仅藉占领并拥有敌人的城市，而且藉拥有君王，已经开始统治；关于那民族的天主许诺，大部分已经应验：不仅是对那三位先祖亚巴郎、依撒格、雅各伯及其时代所作的许诺，也包括通过梅瑟本人所作的许诺，梅瑟曾使同一民族摆脱埃及的奴役，并在带领百姓穿越旷野时，在他的时代揭示了所有过往的事。但既不是由杰出的领袖农的儿子若苏厄——他带领百姓进入应许之地，驱逐异民，按天主的命令将地分给十二支派，然后去世——也不是在他之后整个民长时期，天主关于客纳罕地的许诺得到应验，即该地从埃及的某条河直到幼发拉底大河；这许诺仍未作为即将到来的事被预言，而是期待其应验。它藉达味和他的儿子撒罗满得以应验，他们的王国扩展到整个应许的地域；因为他们征服了所有那些民族，使他们进贡。因此，在这些君王治下，亚巴郎按血统的后裔，即在客纳罕地，建立在应许之地；以致关于那属地许诺的完全应验，只剩余一点：即就现世福乐而言，希伯来民族应藉后裔继承，在那个地上稳固地存留，直到这有死之世的终结，若它遵守上主它的天主的律法。但天主知道它不会这样做，于是祂也使用现世的惩罚来训练其中的少数忠信者，并给日后要在万民中的人应有的警告；在万民中，那在新约中启示的另一个许诺，将要藉基督的降生成人而应验。\par

% structure: p0007 source=h2 target=subsection reason=relative-heading-rank; base=h2

\subsection{第三章　论预言的三重意义，有时指向地上的耶路撒冷，有时指向天上的耶路撒冷，有时又指向两者}

因此，正如那向亚巴郎、依撒格和雅各伯发出的神圣神谕，以及在早期圣经中出现的所有其他先知标记或言论，同样，从君王时代开始的这些预言也部分涉及亚巴郎按血肉的后裔，部分涉及他的后裔——即那藉着新约、在基督内同为承继者、获得永生和天国的万民所蒙受祝福的后裔。因此，它们部分涉及那生子为奴的婢女，即地上的耶路撒冷，她与她的儿女一同为奴；但部分涉及天主的自由之城，即天上永恒的耶路撒冷，她的儿女是那些在地上按天主生活的人。然而，其中有些事被视为同时涉及二者：对婢女而言是直接的，对自由妇人而言则是预表性的。\BibleRef{迦 4:22}\par

因此，先知性的言论可分为三类：有些涉及地上的耶路撒冷，有些涉及天上的耶路撒冷，有些则同时涉及二者。我认为以例子证明我所说的颇为恰当。纳堂先知被派遣去谴责达味王的重罪，并向他预告由此而来的未来灾祸。谁能质疑这一点及类似的事属于地上的城呢？无论是公开的——即关乎民众的安全或援助，还是私下的——即为了个人现世生活的益处而发出预示未来的神圣言论。然而，当我们读到：\Quote{看，日子将到，上主说，我必要与以色列家和犹大家订立新约，不像我握住他们祖先的手，领他们出离埃及地的那天与他们订立的约；因为他们没有遵守我的约，我也就不理会他们了，上主说。因为这是我要与以色列家订立的约：那些日子之后，上主说，我要将我的法律放在他们的脑海里，写在他们的心坎上；我要眷顾他们；我要做他们的天主，他们要做我的子民。}\BibleRef{希 8:8}——无疑这是对天上的耶路撒冷预言的，她的赏报就是天主自己，她首要和全部的善就是拥有祂并属于祂。但以下情况则同时涉及二者：天主的城被称为耶路撒冷，并且预言天主的殿要建在其中；这预言似乎当撒罗满王建造那座最辉煌的圣殿时得以应验。因为这些事既按历史记载发生在属地的耶路撒冷，又是天上耶路撒冷的预像。而这种预言，好似由其他两类混合而成，记载于古代正典书籍中，包含历史叙述，具有非常重要的意义，并极大地操练了那些探究圣经者的才智。例如，我们读到的历史记载，在亚巴郎按血肉的后裔中预言并应验的事，我们也必须探究其寓意，因为它在亚巴郎按信德的后裔中将应验。而且情况如此，以致有些人认为这些书中没有任何事——无论是预言并实现的，还是虽未预言却实现的——不暗示着某种藉预表性意义指向高天的天主之城及其在世上作客旅的儿女。如果真是这样，那么先知们的言论，或者更确切地说，所有归入旧约标题下的圣经，将不是三类，而是两类不同的。因为如果其中所说并应验于她或关于她的每一件事，都表示某种藉预表性预像指向天上的耶路撒冷，那么那里就没有任何事只涉及地上的耶路撒冷；而只有两类：一类涉及自由的耶路撒冷，另一类涉及二者。但正如我认为那些认为这类著作中的历史记载除了实际发生之外别无他意的人犯了大错，同样我也认为那些主张其全部内容都在于寓意的人过于大胆。因此，我说它们是三类，而不是两类。然而，持此看法时，我并不责备那些能从其中每一件事中提取灵性意义的人，只要他们首先保留历史真实性。此外，哪个信徒会怀疑那些既不符合人类事务也不符合神圣事务的事——无论是说已经做的还是将要做的——都是徒然言说的呢？谁不会如果可能就将这些事归入灵性的理解，或承认凡有能力的人都应将它们归入呢？\par

% structure: p0010 source=h2 target=subsection reason=relative-heading-rank; base=h2

\subsection{第四章——论以色列王国和司祭职的预像性转变，以及撒慕尔母亲亚纳代表教会所预言的种种}

因此，天主之城的进展，当它到达列王时代时，提供了一个预象：在\Person{撒乌耳}被弃绝后，\Person{达味}首次获得王位，其基础是此后他的后裔将在地上\Place{耶路撒冷}连续统治；因为事情的进程所象征并预言的（这不可缄默不言）是关于未来事物的变化，涉及旧约和新约——其中司祭职和王国被一位既是司祭又是君王、新而永久的，即\Person{耶稣基督}所改变。因为\Person{厄里}被弃绝为司祭后，由\Person{撒慕尔}取代，他同时执行司祭和民长的职务；以及\Person{撒乌耳}被弃绝后，\Person{达味}被立为王，这两件事都预表了我所说的。至于\Person{撒慕尔}的母亲\Person{亚纳}本人，她从前不孕，后来因生育而喜乐，当她欢欣地向上主献上感恩，将同一个男孩（她以同样的虔诚怀孕并断奶）奉献给天主时，似乎没有预言别的事。因为她说：\BibleQuote{我的心因上主而坚强，我的头角因我的天主而高举；我的口张开胜过我的仇敌；我因祢的救恩而欢欣。因为没有谁像上主那样圣洁；没有谁像我们的天主那样公义：除祢以外没有圣者。不要如此骄傲地夸耀，不要说高傲的事，也不让夸口的话从你们口中出来；因为上主是识别的天主，是预备祂奇妙计划的天主。勇士的弓已被祂折断，软弱的人却束上了力量。曾饱足的，今已减少；饥饿的，已超越大地：因为不妊的生了七个，多子的反而衰弱。上主使人死，也使人活；使人降入阴府，也使人从阴府上来。上主使人穷，也使人富；使人卑微，也使人尊贵。祂从尘土中抬举穷乏人，从粪堆中高举乞丐，为使他们与祂子民中的大能者同坐，并使他们继承荣耀的宝座；祂赐予许愿者所许的，并祝福义人的年岁：因为人不是靠力量强盛。上主必使祂的敌人软弱：上主是圣的。智者不要夸耀自己的智慧，勇士不要夸耀自己的勇力，富人不要夸耀自己的财富：但夸耀的要夸耀这一点：明白并认识上主，并在大地上施行审判和正义。上主升上诸天，发出雷霆：祂要审判地极，因为祂是公义的：祂赐力量给我们的君王，高举祂的受傅者的角。}\par

你们会说，这是一个软弱妇人为生子感恩的话吗？人的心智岂能如此抗拒真理之光，以致看不出这妇人倾吐的话语超过她的度量？再者，那适当关注这些事的人（这些事甚至在此尘世旅途中已开始应验），他岂不应用心智，看出并承认，透过这妇人——她的名字\Person{亚纳}意为\Quote{祂的恩宠}——那基督信仰本身，那天主之城本身（其君王和建立者是基督），简言之，天主的恩宠本身，藉着先知之神如此发言，由此骄傲的被剪除而跌倒，谦卑的被充满而升高，这正是那首颂歌所主要赞美的？除非或许有人会说这妇人没有预言什么，只是因她曾祈祷而得了儿子，便以欢欣的赞美颂扬天主。那么，当她说：\BibleQuote{勇士的弓已被祂折断，软弱的人却束上了力量；曾饱足的，今已减少；饥饿的，已超越大地；因为不妊的生了七个，多子的反而衰弱。}她是什么意思呢？她自己虽曾不妊，却生了七个吗？她说这话时只有一个儿子；后来她也没有生七个，也没有生六个（\Person{撒慕尔}本身可能是第七个），而是生了三个男儿和两个女儿。再者，那时在那一民族中尚无君王，如果她没有预言，她最后说\BibleQuote{祂赐力量给我们的君王，高举祂的受傅者的角}又是什么意思呢？\par

因此，让基督的教会，伟大君王的城邑，满被恩宠，后裔繁盛，让她说出那预言如此之久以前通过这位虔诚母亲的唇舌所宣认的：\Quote{我的心因主而坚强，我的角因我的天主被高举。}她的心真正坚强，她的角真正被高举，因为不是在她自己内，而是在主、她的天主内。\Quote{我的口向我的仇敌张开；}因为即使在紧迫困境中，天主的圣言也不受束缚，即使在受束缚的宣讲者中也不受束缚。\Quote{我因你的救恩而欢欣，}她说。这就是耶稣基督自己，正如我们在福音中读到，年迈的西默盎将它拥抱如婴孩，却认出他是伟大的，说道：\BibleQuote{主，现在让你的仆人平安离去，因为我亲眼看见了你的救恩。}\BibleRef{路 2:25}因此，教会可以说：\Quote{我因你的救恩而欢欣。因为除主以外，没有圣者，除我们的天主以外，没有义者；}祂是圣洁的，也是圣化者，是公义的，也是成义者。\Quote{除你以外，没有圣者；}因为没有人能成为圣者，除非因着你。接着又说：\Quote{不要如此骄傲地自夸，不要说高傲的话，也不要让夸口之言出自你的口。因为主是洞悉万有的天主。}祂认识你，即使无人认识；因为\Quote{人若认为自己是什么，而其实什么都不是，就是欺骗自己。}\BibleRef{迦 6:3}这些话是对天主之城的仇敌说的，他们属于巴比伦，依恃自己的力量，以自己为荣，而不在主内；其中也包括属血肉的以色列人，地上耶路撒冷的土生居民，正如宗徒所说：\Quote{因为不认识天主的义，}\BibleRef{罗 10:3}即唯独公义且成义的天主赐给人的义，\Quote{而想建立自己的义，}即仿佛由自己获得、而非由祂赐予的义，\Quote{就不服从天主的义，}正因为他们是骄傲的，认为能以自己所有的、而非来自天主的一切来取悦天主，而祂是洞悉万有的天主，因此也监督良心，在那里洞察人的思念是虚妄的，若是出于人而不是出于祂。\Quote{祂准备，}她说，\Quote{祂奇妙的计划。}我们认为这些奇妙的计划是什么呢？无非是骄傲者必要跌倒，谦卑者必要兴起。她述说这些奇妙的计划，说：\Quote{强者的弓被折断，弱者以力量束腰。}弓被折断，即那些认为自己如此强大之人的意图，以为没有天主的恩赐和帮助，单凭人的自足就能遵守神命；而那些以力量束腰的人，他们心中的呼声是：\Quote{上主，求你怜悯我，因为我软弱。}\par

\Quote{那些饱食的，}她说，\Quote{反而减少，饥饿的却遍及大地。}谁是那些饱食的呢？不就是那些如同强者的以色列人吗？天主的圣言交给了他们。\BibleRef{罗 3:2}但在那个民族中，婢女的子女减少了——用\Emph{minus}这个词，虽然是拉丁文，却很好地表达了从更大变为更少的意思——因为，即使在那食物本身，即天主的圣言中，以色列人在万民中独受，他们却品味属世的事物。而那些没有领受法律的民族，借着新约来到这些圣言面前，因着极大的饥渴而遍及大地，因为他们在其中品味了不是属世而是属天的事物。这样做的原因似乎被寻求：\Quote{因为不生育的，}她说，\Quote{生了七个，而那多子的反而衰弱。}这里，所有预言都已向那些理解数目\Quote{七}的人显明，七象征普世教会的完美。为此，宗徒若望向七教会写信，\BibleRef{默 1:4}以此表明他写信给唯一教会的整体；并且在撒罗满的箴言中，早已预言此事说：\Quote{智慧建造了她的房屋，她立了七根柱子。}\BibleRef{箴 9:1}因为天主的城在万民中曾是不生育的，直到我们看到的那位子嗣兴起。我们也看到，那属地的耶路撒冷，曾有许多子女，现在却衰弱了。因为，凡在其中的自由妇人的儿子，曾是她的力量；但现在，既然那里只有文字，而没有精神，失去了力量，她就衰弱了。\par

\Quote{主殺人，也使人活：}祂殺了那多有兒女的，卻使這不育的活過來，以致她生了七個。雖然可以更合適地理解為，祂使那些祂所殺的人活了過來。因為她，彷彿，重複這一點，加上說：\Quote{祂使人下地獄，也使人上來。}宗徒實在對他們說：\Quote{你們既然與基督同死，就當追求上面的事，那裏基督坐在天主的右邊。}\BibleRef{哥 3:1}因此，他們被主以一種有益的方式殺死，以致他又加上：\Quote{你們要思念上面的事，不要思念地上的事；}這樣，他們就是那些因饑餓而超越了塵世的人。\Quote{因為你們已經死了，}他說：看，天主是如何拯救性地殺人！然後接著說：\Quote{你們的生命與基督一同藏在天主內：}看，天主是如何使同一人活過來！但祂使他們下地獄又再上來嗎？在信徒中間毫無爭議的是，我們最好看到這工作的兩部分都在祂身上實現，即我們的首領，宗徒說我們的生命與祂一同藏在天主內。\Quote{因為祂沒有憐惜自己的兒子，反而為我們眾人把祂交出來，}\BibleRef{羅 8:32}這樣，當然，祂殺死了祂。又因為祂使祂從死人中復活，祂使祂活了過來。既然在預言中承認祂的聲音：\Quote{祢不會將我的靈魂遺留在地獄，}祂使祂下地獄又使祂上來。因祂的貧窮，我們成了富足的；\BibleRef{格後 8:9}因為\Quote{主使人貧窮，也使人富足。}但為使我們知道這是什麼，讓我們聽聽接著說的是什麼：\Quote{祂降低，也升高；}實在，祂謙抑驕傲的人，高舉謙卑的人。我們也在別處讀到：\Quote{天主阻擋驕傲的人，卻賜恩寵給謙卑的人。}這就是那位名字被解釋為\Quote{祂的恩寵}的婦人整首詩歌的負擔。\par

再者，接著說：\Quote{祂從塵埃中提拔窮人，}我認為沒有比祂更合適的理解，正如不久之前所說：\Quote{祂本是富有的，為了我們卻成了貧窮的，好使我們因祂的貧窮而成為富有的。}因為祂如此迅速地將祂從地上舉起，以致祂的肉身未見腐朽。我也不將接著說的話從祂身上偏離：\Quote{並從糞堆中高舉窮人。}的確，那位窮人也就是乞丐。但從其中被高舉的糞堆，我們很有理由理解為迫害的猶太人，宗徒在講述他屬於他們時曾迫害教會，說：\Quote{凡從前對我是贏利的，我為了基督都看作是損失。我不僅把它們看作損失，而且看作糞土，為贏得基督。}\BibleRef{斐 3:7}因此，那窮人從塵埃中被高舉，超過一切富人；那乞丐從糞堆中被高舉，超過一切財主，\Quote{使他坐在人民首領之中，}祂對他們說：\Quote{你們要坐在十二個寶座上，}\BibleRef{瑪 19:27}\Quote{並使他們繼承榮耀的寶座。}因為這些首領曾說：\Quote{看，我們捨棄了一切，跟隨了祢。}他們曾極其有力地表白了這個誓言。\par

但他们从何处领受这些，除非从祂，就是此处直接所说的那一位：\Quote{赐予发愿者以所愿？}否则，他们就会是那些弓已折断的强者之列。\Quote{赐予}，她说，\Quote{发愿者以所愿。}因为没有人能向天主发愿蒙悦纳的事，除非他从祂领受了他所能发愿的。接着是：\Quote{祂祝福了义人的年岁，}即，使他能与那一位永远生活，对祂说：\Quote{祢的年岁没有穷尽。}因为那里的年岁长存，但这里的年岁消逝，甚至朽坏：它们未到之前不存在，到了之后也不存在，因为它们带着自身的终结。关于这两点，即\Quote{赐予发愿者以所愿}和\Quote{祂祝福了义人的年岁}，前者是我们所做的，后者是我们所领受的。但这后者并非从慷慨的赐予者天主领受，除非那帮助者祂亲自使我们有能力行前者；\Quote{因为人不能以力逞强。}\Quote{主要使他的仇敌软弱，}即那嫉妒发愿之人并抵挡他，免得他完成所愿的。由于希腊文的含混，也可以理解为\Quote{祂自己的仇敌。}因为当天主开始占有我们时，那曾是我们的仇敌的，立刻成为祂的，并被我们征服；但不是靠我们自己的力量，\Quote{因为人不能以力逞强。}因此\Quote{主要使祂自己的仇敌软弱；主是圣的，}使他被圣人征服，这些圣人是主，至圣者，所圣化的。为此，\Quote{智者不要夸耀自己的智慧，勇士不要夸耀自己的勇力，富人不要夸耀自己的财富；夸耀的，只应夸耀这一点：明瞭和认识主，并在大地中央施行审判和正义。}谁若明瞭并认识，甚至这一点——他能明瞭并认识主——是主赐给他的，这人便不小地明瞭并认识主。\Quote{因为你有什么，}宗徒说，\Quote{不是领受的呢？既然领受了，为什么还夸耀，好像不是领受的呢？}\BibleRef{格前 4:7}即是，好像你自己有什么可夸耀的。现在，谁正直生活，谁就行审判和正义。但谁听从天主的命令，谁就正直生活。\Quote{诫命的终极，}即诫命所指向的，\Quote{是出自纯洁的心、良善的良心和真诚的信德的爱德。}此外，这\Quote{爱德，}正如宗徒若望所证，\Quote{是由天主来的。}\BibleRef{若一 4:7}因此，行正义和审判是由天主来的。但什么是\Quote{在大地中央？}难道住在地极的人不应该行审判和正义吗？谁能这样说？那么，为何加上\Quote{在大地中央？}若没有加上这句，只说\Quote{行审判和正义，}这条诫命就同时适用于两种人——住内陆的和住海边的。但恐怕有人认为，在肉身生活结束后，还有时间行他在肉身时未行的审判和正义，从而可以逃避天主的审判，所以\Quote{在大地中央}在我看来是指每个人生活在肉身中的时间；因为在这生命中，每个人都带着自己的泥土，人死后，共同的泥土取回它，在复活时必然归还给他。因此\Quote{在大地中央}，即当我们的灵魂被禁锢在这属地的身体中时，必须行审判和正义，这将在我们死后有益，那时\Quote{每人要按他肉身所行的，或善或恶，领取报应。}\BibleRef{格后 5:10}因为宗徒在那里说\Quote{在肉身内}，意思是他在肉身中生活的那段时间。然而，若有人以恶毒的心和不敬的念头亵渎，没有动用身体的任何肢体，他并不因此无罪，因为他是在肉身中度过的那个时刻做的。同样，我们可以恰当地理解在圣咏中所读的：\Quote{但是天主，我们在万世以前的君王，在大地中央施行了救恩；}如此，主耶稣可以被理解为我们万世以前的天主，因为万有是藉着祂创造的，祂在大地中央施行了我们的救恩，因为圣言成了血肉，居住在属地的身体中。\par

然后，在\Person{亚纳}用这些话预言之后，即那夸耀的应该全在主内夸耀，而非在自己，她就因那即将来临的审判之日的报应而说：\Quote{主已升上高天，并发出雷声：祂要审判地极，因为祂是正义的。}她始终持守着基督徒信经的次序：因为主基督已升天，并要从那里来审判生者死者。\BibleRef{宗 10:42}因为，正如宗徒所说：\Quote{那上升的，不就是那曾下降到地下最深处的那位吗？那下降的，正是那上升到诸天之上，为充满万有的那一位。}\BibleRef{弗 4:9}因此，祂藉着祂的云彩发出雷声，这些云彩是祂上升时以祂的圣神所充满的。关于这，在\Person{依撒意亚}先知书中，那婢女耶路撒冷——即那不结果实的葡萄园——被警告说，将没有雨水降在她身上。但\Quote{祂要审判地极}这句话，就好像说\Quote{甚至地的最远端}。这并不是说祂不审判地的其余部分，祂无疑要审判所有的人。但最好将地极理解为人的终极状态，因为在中间时期发生变化的事物（变好或变坏）不会被审判，而是那被审判者被发现时的终结。为此缘故，经上说：\Quote{那坚持到底的，必得救。}\BibleRef{玛 24:13}因此，凡在地上坚持不懈地施行审判与正义的人，当地极受审判时，必不被定罪。她又说：\Quote{赐力量给我们的君王，}使祂在审判时不至于定他们的罪。祂赐予他们力量，使他们作为君王统治血肉，并借着那为他们倾流鲜血的祂征服世界。\Quote{并高举祂的受傅者的角。}基督如何高举祂的受傅者的角？因为上面所说的\Quote{主已升上高天}，指的是主基督自己，正如这里所说的，祂\Quote{要高举祂的受傅者的角}。那么，谁的受傅者是祂的受傅者？这是否意味着祂要高举每个信祂的子民的角，正如她在这首赞美诗开头所说：\Quote{我的角因我的天主而高举}？因为我们能正当地称所有受祂傅油的人为受傅者，因为整个身体与它的头是一个基督。\BibleRef{格前 12:12}这些就是\Person{亚纳}，即圣者备受赞美的\Person{撒慕尔}的母亲，所预言的事，其中古代司祭职的变更当时已被预示，如今得以应验，因为那多子的已衰败，使那不能生育却生了七个的，得以在基督内拥有新司祭职。\par

% structure: p0019 source=h2 target=subsection reason=relative-heading-rank; base=h2

\subsection{第五章  论一位天主的人藉圣神向司祭\Person{厄里}所说的事，预示那按\Person{亚郎}设立的司祭职位将被除去。}

但这一点更明确地由一位天主的人亲自向司祭\Person{厄里}所说的话表达出来，这位天主的人的名字虽未提及，但他的职务与职事无疑表明他是一位先知。因为经上这样记载：\Quote{有一位天主的人来到\Person{厄里}面前，对他说：上主这样说：当你们祖先在埃及法郎家中为奴时，我明确地向你们祖先的家显现；我从以色列所有的支派中拣选了你们祖先的家，使他们给我充任司祭的职务，上我的祭台，焚香并佩带\OriginalTerm{厄弗得}{ephod}；我将以色列子民的一切火祭食物都赐给你们祖先的家。你们为何以无耻之眼注视我的馨香和我的祭品，并尊崇你们的儿子超过我，在以色列中在我之先祝福每一祭品的初果呢？因此上主以色列的天主这样说：我曾说，你的家和你的祖先家要永远在我面前行走；但现在上主说：这事远离我吧！因为尊重我的，我必尊重；轻视我的，必受轻视。看，日子将到，我要剪除你的后裔和你祖先家的后裔，在你家中将永远没有一个老人。我要从我祭台上除掉你的人，使他的眼目昏花，他的心灵消融；你家中余下的每一个人都必丧身刀下。这给你两个儿子\Person{曷弗尼}和\Person{丕乃哈斯}必临之事的标记：他们二人必同日而死。我要为自己兴起一位忠信的司祭，他要照我心中和我灵魂所想的一切去行；我要为他建立一坚固的家，他要在我的受傅者面前永远行走。将来你家中剩下的人要前来向他下拜，手里拿着一块银币说：请将我安置在你司祭职的一部分中，使我得吃饼。}\par

我们不能说这个预言——其中如此清晰地预言了古代司铎职的改变——是在撒慕尔身上应验的；因为虽然撒慕尔并非来自另一个支派，而是来自那被天主指定事奉祭台的支派，但他也不是亚郎的子孙，而亚郎的后裔曾被分别出来，以从中选取司铎。因此，藉着那件事，也预表了那将要藉着耶稣基督而来的同一改变；而预言本身，在事实上（而非言语上），严格地说属于旧约，但预表性地属于新约，以事实来象征那藉着先知对厄里司铎所说的话。因为在达味王时代及以后，曾有属亚郎族系的司铎，如匝多克和阿彼雅塔尔等相继担任，直到那长久以来预言关于司铎职改变的事，必须藉着基督应验的时候来到。然而，如今谁以信德之眼观看这些事，却不看见它们已经应验了呢？因为，事实上，犹太人已不再有会幕、圣殿、祭台、祭献，因此也没有了司铎——按天主的法律，司铎应由亚郎的后裔中任命；这一点先知在此也提到了，他说：\Quote{主以色列的天主这样说：我曾说过，你的家和你父的家要永远在我面前往来；但现在主说：这事绝不可有我；因为尊重我的，我必尊重他；藐视我的，他必被轻视。}因为他在提到\Quote{你父的家}时，所指的不是他生身之父的家，而是亚郎的家——亚郎是最初被立为司铎的，其后裔继任——这从前面的话可以看出来，他说：\Quote{我启示于你父的家，当他们还在埃及地，在法老家中为奴的时候；我从以色列各支派中拣选了你父的家，为我充任司铎之职。}在那埃及为奴的时期，\Quote{你父}不是指亚郎吗？当以色列人被释放时，亚郎被选为司铎。因此，他在这里说，出自亚郎血统的人将不再为司铎；这一点我们已经看见应验了。若信德警醒，这些事就在我们眼前：它们被辨识，被把握，甚至强加于不情愿之人的眼目，以致他们看见：\Quote{看，日子将到，他说，我要砍断你的后裔和你父家的后裔，使你家中再没有老人。我要从我的祭台上除掉你家中的人，使他的眼睛失明，他的心消融。}看，那预言的日子已经来到。没有按亚郎次序的司铎了；凡属他血统的人，当他看见基督徒的祭献遍行全世界，而他自己那极大的尊荣被夺去时，他的眼目失明，他的灵魂因忧愁而消融。\par

但接下来的一段话严格说是对厄里家的，这些话正是对他说的：\Quote{你家中剩下的人，都必倒在人的刀下。这事给你作个兆头，必临到你这两个儿子曷弗尼和非尼哈身上；他们两人必在同一日内死亡。}因此，这事成了司铎职从这家改变的标记，藉此表明亚郎家的司铎职将要改变。因为这人之子的死所象征的不是这些人的死，而是亚郎子孙的司铎职本身的死。但接下来那段话涉及那位司铎，就是撒慕尔在继承这家时所预表的。因此，下面的话是关于耶稣基督——新约的真司铎：\Quote{我要为自己兴起一位忠信的司铎，他要照我的心意和我的灵魂行事；我要为他建立一座稳固的家。}这家就是天上的永远耶路撒冷。\Quote{他要在我的基督面前，永远往来。}\Quote{往来}即\Quote{生活交往}，正如他先前对亚郎家所说的：\Quote{我曾说过，你的家和你父的家要永远在我面前往来。}但他所说的\Quote{他要在我的基督面前往来，}应完全理解为指那家本身，而非指那司铎——基督本人，即中保和救主。因此，他的家要在祂面前往来。\Quote{往来}也可理解为从死亡到生命，在整个这必死的生命经过期间，直到世界终结。至于天主说\Quote{他要照我的心意和我的灵魂行事}，我们不可认为天主有灵魂，因为祂是灵魂的创造者；这话是以比喻的方式论及天主，而非字面意义，正如祂被说有手、脚和其他肢体一样。而且，免得有人从这样的说法中以为人是以这肉身的形状按天主的肖像受造的，翅膀也被归于祂，而人根本没有翅膀；又对天主说：\Quote{请将我隐藏于祢羽翼的荫下，}好使人明白，论及那不可言喻的本性，这些话不是以字面意义，而是以比喻方式说的。\par

但所附加的：\Quote{凡你家中所留下来的人，必来叩拜他}，这话不是指着这位厄里的家说的，而是指着那位亚郎的家说的，亚郎家中的人一直存留到耶稣基督来临，直到如今这一族中仍不乏其人。因为关于厄里的家，上面已经说过：\Quote{你家中凡留下来的人，都必倒在人的刀下}。既然如此，这里怎能正确地说：\Quote{凡留下来的，都必来叩拜他}呢？如果那一说（无人能逃脱复仇之剑）是真实的，那么除非他把这话理解为属于按照亚郎品位整个司祭的后裔。因此，如果这是关于预定遗民的，关于他们另一位先知说过：\BibleQuote{遗民必将得救}\BibleRef{依 10:21}；因此宗徒也说：\BibleQuote{同样，现在也有按恩宠拣选的遗民得救了}\BibleRef{罗 11:5}；既然容易理解这样的遗民被说成\Quote{你家中所留下来的}，他肯定相信基督；正如在宗徒时代，该国许多人都信了；现在也不乏其人，虽然很少，但仍有相信的，而在他们身上应验了这位天主的人在此紧接着所说的话：\Quote{他必来叩拜他，带着一块银钱}；叩拜谁呢？岂不是那位大司祭，祂也是天主吗？因为在按照亚郎品位的司祭中，人来到天主的圣殿或祭台前，并不是为了叩拜司祭。但祂说\Quote{带着一块银钱}是什么意思呢？岂不是信德的简短言语，关于这个宗徒引用了那句：\BibleQuote{主要在地上完成并缩短的话}？但圣咏证明了\Quote{银钱}代表言语，其中咏唱：\Quote{上主的话是纯洁的话，如同在火中锤炼过的银钱}。\par

那么，他来叩拜天主的司祭，即那位是天主的司祭，他说了什么？\Quote{请把我安置在祢司祭职的一分上，好得饼吃}。我不愿被立在我祖先的荣誉上，那是虚无的；请将我安置在祢司祭职的一分上。因为我宁愿住在祢的殿宇中；我愿意成为祢司祭职的一个肢体，无论什么，或多么微小。他在这里以司祭职指百姓本身，这百姓的司祭就是天主与人之间的中保，即基督耶稣\BibleRef{弟前 2:5}。宗徒伯多禄称这百姓为\BibleQuote{圣洁的子民，王家的司祭}\BibleRef{伯前 2:9}。但有些人翻译成\Quote{祢的祭品}，而非\Quote{祢的司祭职}，这同样指同一个基督徒百姓。因此宗徒保禄说：\BibleQuote{我们虽多，却是一个饼，一个身体}\BibleRef{格前 10:17}。[他又说：\BibleQuote{将你们的身体献上，当作生活的祭献}\BibleRef{罗 12:1}。] 因此，他所补充的\Quote{吃饼}，也优雅地表达了那种祭献本身，关于这祭献，司祭亲自说：\BibleQuote{我所要赐给的饼，就是我的肉，是为世界的生命而赐给的}\BibleRef{若 6:51}。这同一个祭献不是按照亚郎的品位，而是按照默基瑟德的品位：凡读懂的，请明白\BibleRef{玛 24:15}。因此，这简短而有益于得救的谦卑宣认，即所说的\Quote{请把我安置在祢司祭职的一分上，好得饼吃}，本身就是那块银钱，因为它既简短，又是居住在信者心中的天主圣言。因为上文曾说，祂已将旧约的祭物赐给亚郎的家族作为食物，祂说：\Quote{我已将以色列子民以火所献的一切，赐给你父家作食物}，这些确实就是犹太人的祭献；因此这里祂说：\Quote{吃饼}，这在新约中是基督徒的祭献。\par

% structure: p0025 source=h2 target=subsection reason=relative-heading-rank; base=h2

\subsection{论犹太人的司祭职与王国：虽然应许永远建立，却未能延续；因此应理解为其他事物，其永恒性确凿无误。}

因此，虽然这些事现今清楚显明，正如它们曾被崇高预言，但仍有人可能并非徒然被激起询问：如果这句由天主所说的话——\Quote{你的家和你的父家，必永远在我面前行走}——未能实现，我们如何能确信这些书中所预言的一切将来都要发生呢？因为我们看到司祭职已经改变；而且，既然在被拒绝和改变之后所继承的倒被预言为永恒，就不能指望应许给那个家的事有一天会实现。说这话的人尚未理解或没有记起，这正是按照亚郎品位的司祭职被设立为未来永恒司祭职的预象；因此，当永恒被应许给它时，应许不是给那单纯的预象和影像，而是给它所预示和预表的事物。但唯恐有人认为预象本身要存留，所以它的变化也必须被预言。\par

同样，撒乌耳自己的王国——他确实被摒弃和弃绝了——也是一个将来要永远存留的国度的预象。因为他受傅所用的油，从这圣油起他被称作基督，这应被理解为奥秘的意义，一个伟大的奥迹；达味本人对他如此敬畏，以至于当他在一个黑暗的洞穴里（撒乌耳因生理需要也进了这洞），达味悄悄从后面割下他袍子的一小块时，他内心战栗，为证明他本可杀他却饶了他，从而消除撒乌耳心中因激烈迫害圣达味（视他为仇敌）而产生的猜疑。因此达味非常害怕，担心自己因触碰撒乌耳的衣物而冒犯了如此伟大的奥迹。经上如此记载：\Quote{达味因割去撒乌耳外袍的衣边，心中击打自己。}但对那些劝他杀死这个被交到他手中的撒乌耳的同伙，他说：\Quote{上主决禁止我对我的主、上主的受傅者行这事，下手加害他，因为他是上主的受傅者。}因此他对这未来事物的预象表现出极大的尊敬，不是为它本身，而是为了它所预表的那位。同样，撒慕尔对撒乌耳所说的话：\Quote{因为你没有遵守上主吩咐你的诫命，上主本应使你的王权在以色列永远竖立，但如今你的王权决不会长久；上主将要寻找一个合他心意的人，并命他作他人民的首领，因为你没有遵守上主所吩咐的。}不应理解为天主曾决定撒乌耳本人永远为王，后来因他的罪而不遵守这应许；天主并非不知他会犯罪，而是建立了他的王国作为永恒国度的预象。因此他接着说：\Quote{但如今你的王权决不会为你而持久。}因此它所预表的事物已经站立并要站立；但它不会为这个人而站立，因为他本人不会永远统治，他的后裔也不能；这样一来，至少\Quote{永远}这个词似乎要通过他的后代代代相传来实现。他说：\Quote{上主将要寻找一个人，}指的要么是达味，要么是新约的中保\BibleRef{希 9:15}，他由那傅油预表，达味及其后裔也因这傅油而受傅。但这并非天主不知他在哪里而寻找一个人，而是通过人说话，像人一样说话，在此意义上寻找我们。因为不仅对天主圣父，也对那来寻找失丧者\BibleRef{路 19:10}的独生子，我们早已被认识，甚至在世界创立之前在基督里被拣选\BibleRef{弗 1:4}。因此\Quote{他将寻找他}的意思是，他将拥有自己的人（就如他说：他已认识的属己者，他将向他人显明为他的朋友）。因此在拉丁文中，这个词\OriginalTerm{寻找}{quaerit}加上一个介词变成了\OriginalTerm{获得}{acquirit}，其意义足够清楚；尽管即使不加介词，\OriginalTerm{寻找}{quaerere}也被理解为\OriginalTerm{获得}{acquirere}，由此收益被称为\OriginalTerm{利润}{quaestus}。\par

% structure: p0028 source=h2 target=subsection reason=relative-heading-rank; base=h2

\subsection{第七章——论以色列王国的分裂，借此预表了属神的以色列与属血肉的以色列之间的永久分离。}

撒乌耳又因不顺服而犯罪，撒慕尔再次奉上主的话对他说：\Quote{因为你轻视了上主的话，上主也轻视你，使你不作以色列王。}\BibleRef{撒上 15:23} 又因同样的罪，撒乌耳认罪并祈求宽恕，恳求撒慕尔与他同去平息上主的怒气，撒慕尔说：\Quote{我不与你同去；因为你轻视了上主的话，上主也必轻视你，使你不作以色列王。}撒慕尔转身要走，撒乌耳抓住他外袍的衣边，撕了下来。撒慕尔对他说：\Quote{上主今日将以色列王国从你手中撕裂，赐给你的邻人，一个比你更好的人，并将以色列分为二。上主不会改变，也不会后悔；因他不是人，不后悔；他威胁却不坚持。}那听到\Quote{上主轻视你，使你不作以色列王}和\Quote{上主今日将以色列王国从你手中撕裂}的人，却统治了以色列四十年——和达味本人一样长——然而这些话是在他统治初期听到的，为使我们明白，这话是说因为他家族中没有一人要作王，并使我们仰望达味的家族，按肉身\BibleRef{罗 1:3}从那家族中出生了天人之际的中保、基督耶稣\BibleRef{弟前 2:5}。\par

但是，圣经并没有大多数拉丁文抄本所读的：\Quote{主今日已将以色列的王国从你手中撕裂了，}而是正如我们所记录的那样，在希腊文抄本中发现的是：\Quote{主已将王国从以色列中从你手中撕裂了；}这样，\Quote{从你手中}这个词可以理解为\Quote{从以色列。}因此，这个人象征性地代表了以色列人民，他们将要失去王国，那时我们的主耶稣基督将要统治，不是按血肉，而是按神。当论及祂时说：\Quote{并将它给你的邻人，}这是指肉身的亲属关系，因为基督按肉身是出于以色列，撒乌耳也是从他而出。但所添加的\Quote{善良在你之上，}确实可以被理解为\Quote{比你好，}而且有些人确实这样翻译；但最好这样理解：\Quote{善良在你之上，}意思是，因为祂是良善的，所以祂必须在你之上，正如另一个先知预言所说的：\Quote{直到我把祢的所有仇敌放在祢的脚下。}其中也包括以色列，基督从他那里，作为他的迫害者，夺去了王国；虽然那没有诡诈的以色列也可能在那里，好比糠秕中的一些麦粒。因为的确，宗徒们是从那里来的，那么多殉道者是从那里来的，其中斯德望是首位，那么多教会是从那里来的，宗徒保禄在使他们归化时赞美天主，并提名这些教会。\par

关于这事，我毫不怀疑下面的话应当这样理解：\Quote{并且要将以色列一分为二，}即分为属于婢女的以色列和属于自由的以色列。因为这两种人起初是共处的，就如亚巴郎仍依恋着婢女，直到那不生育的，因天主的恩宠而成为多产的，喊道：\Quote{把婢女和她的儿子赶出去！}\BibleRef{创 21:10}确实，我们知道，由于撒罗满的罪，在他儿子勒哈贝罕统治时期，以色列分裂为二，并一直持续，各自有各自的君王，直到那整个民族在一次大毁灭中被倾覆，并被加色丁人掳去。但这与撒乌耳何干呢？因为如果这样的事受到警告，那威胁是针对达味本人，他的儿子是撒罗满。最后，希伯来民族现在不是内部分裂，而是毫无分别地分散在世界各地，同属一个错误的团体。但是，天主以撒乌耳（他代表他们）来威胁王国和人民的那个分裂，被添上的这句话表明是永恒不变的：\Quote{祂不会改变，也不会后悔：因为祂不像人，会后悔；人威胁却不坚持}——也就是说，人威胁却不坚持，但天主不是这样，祂不像人那样后悔。因为当我们读到祂后悔时，意思是情况发生变化，这源于天主不变的预知。因此，当天主被说成不后悔时，应当理解祂不变化。\par

我们看到，关于以色列人民分裂的这个判决，由天主用这些话神圣地宣告，是完全不可挽回且永久不变的。因为凡是从那里转向基督的人，无论是已经转向、正在转向还是将要转向，都是按照天主的预知，而不是按照人类种族同一的本性。当然，没有一个依附基督并持续在祂内的以色列人，会永远留在那些直到生命终结仍坚持作祂仇敌的以色列人中，而将永远停留在此处预言的分离中。因为出自西乃山、生子为奴的旧约\BibleRef{迦 4:25}毫无益处，除非因为它为新约作证。否则，无论梅瑟被阅读多久，帕子总是蒙在他们心上；但一旦有人从那里转向基督，帕子就被除去了\BibleRef{格后 3:15}。因为那些转向者的渴望本身从旧转变为新，以至于每个人不再渴望获得属血肉的福乐，而是属神的福乐。因此，那位大先知撒慕尔本人，在傅油撒乌耳之前，当他为以色列向主呼求，主垂听了他，当他献上全燔祭时，外邦人前来与天主的子民作战，主在他们上方打雷，他们惊惶失措，倒毙在以色列面前并被击败；于是，他拿起一块石头，立在旧的\OriginalTerm{米兹帕}{Massephat}与新的\OriginalTerm{米兹帕}{Massephat}之间，给那块石头起名叫\OriginalTerm{厄本厄则尔}{Ebenezer}，意即\Quote{帮助之石}，并说：\Quote{直到如今主都帮助了我们。}\OriginalTerm{米兹帕}{Massephat}被解释为\Quote{渴望。}那块帮助之石就是救主的媒介，藉此我们从旧的米兹帕走向新的米兹帕——也就是说，从在血肉的王国中期待血肉福乐的渴望，走向在天国中期待最真实的属神福乐的渴望；既然没有比这更好的，主就一直在帮助我们。\par

% structure: p0033 source=h2 target=subsection reason=relative-heading-rank; base=h2

\subsection{第8章——关于对达味在他儿子中所作的应许，这些应许绝不是在撒罗满身上，而是在基督身上完全应验。}

现在，我认为必须指出，天主就我所论述的主题，向\Person{达味}本人应许了什么。\Person{达味}继\Person{撒乌耳}为王，他的变迁预表了那最终的变迁，一切神圣之言和一切书写都是为此而发。当\Person{达味}王诸事亨通时，他想要为天主建造一座殿宇，即那座后来由他儿子\Person{撒罗满}王建造的、名声卓著的圣殿。他正思忖此事时，主的话临到\Person{纳堂}先知，\Person{纳堂}将这话传给王。天主先说，\Person{达味}本人不可为祂建造殿宇，且在那么长的时期里，祂从未命令祂的任何子民为祂建造香柏木的殿宇；然后祂说：\BibleQuote{如今你要这样对我的仆人\Person{达味}说：全能的天主这样说：我将你从羊圈中取来，使你作我人民\Place{以色列}的领袖；我无论往哪里去，都与你同在，从你面前剪除了你的一切仇敌，使你得到了大名，如同世上伟人的大名一样。我要为我的百姓\Place{以色列}安置一个地方，栽培他们，使他们安居在自己的处所，不再受人惊扰；恶人也不再像先前那样欺压他们，自从我任命士师治理我百姓\Place{以色列}的日子以来。我要使你安享太平，不被一切仇敌扰乱；主也告诉你，祂必为你建立家室。及至你的日子满足，与你祖先同眠时，我必在你以后兴起你的后裔，即你所生的儿子，并巩固他的王国。他要为我的名建造殿宇；我要坚定他的王位，直到永远。我要作他的父亲，他要作我的儿子。他若犯过，我必用人的杖，用世人的鞭责打他；但我的慈爱必不离开他，如同我离开在你面前所废弃的人一样。你的家室和王权，在我面前必永远坚定不移；你的王位也必永远坚定。}\BibleRef{撒下 7:8}\par

凡以为这伟大的应许在\Person{撒罗满}身上应验的人，就大错特错了；因为他只留意\Quote{他要为我建造殿宇}这句话，却不留意\Quote{他的家室在我面前将永远忠诚，他的王国永远坚立}这句话。让他留神看看\Person{撒罗满}的家室充满了崇拜偶像的外邦女子，而这位从前智慧的王也被她们引诱，陷入同样的偶像崇拜；他就不敢认为天主要么虚假地应许了这事，要么不能预先知道\Person{撒罗满}和他的家室会变成那样。但我们在这里不应怀疑，也不应看到这些事的应验只在我们主基督身上，祂按肉身是\Person{达味}的后裔\BibleRef{罗 1:3}，免得我们像属肉体的犹太人一样，徒劳无益地在此寻找别的应验。因为连他们也明白，那地方所应许给\Person{达味}的儿子不是\Person{撒罗满}；他们对那应许的、如今已如此显明地宣告的基督却盲目无知，以致说他们盼望另一位。的确，在\Person{撒罗满}身上也出现了未来事件的某种影像，因他建造了圣殿，按照他的名字享有和平（因为\Person{撒罗满}意为\Quote{和平}），并在登基之初令人惊异地值得称赞；但作为那要来者的预像，他预示了我们的主基督，但他本人并不相似基督。因此，关于他所写的一些事，好像是为他本人预言的，而圣经甚至通过事件预言，以某种方式在他身上描绘了未来之事的形象。因为除了记载他统治的神圣历史书之外，\BibleBook{圣咏}{第七十二篇}的标题也以他的名字题写，其中许多话完全不能应用于他，却如此明显地适用于主基督，以致显然，在一者中预像以某种方式被影射，在另一者中真理本身被呈现。因为\Person{撒罗满}的王国范围是有限的；然而在那篇圣咏中，且不说别的，我们读到：\BibleQuote{他将统治从这海到那海，从大河直到地极}，这我们在基督身上应验了。确实，祂从\Person{若翰}施洗的那条河开始祂的统治；因为当\Person{若翰}指出祂时，祂开始被门徒承认，门徒不仅称祂为\Quote{老师}，也称祂为\Quote{主}。\par

\Person{撒罗满}在他父亲\Person{达味}还在世时就开始统治，这件事没有发生在他们其他任何一位君王身上，这并非没有其他原因，而是为叫这也清楚显明，那对他父亲说的预言所指的并不是他自己。预言说：\BibleQuote{及至你的日子满足，与你祖先同眠时，我必在你以后兴起你的后裔，即你所生的，并巩固他的王国。}那么，怎能因后面说\Quote{他要为我建造殿宇}，就认为这预言是\Person{撒罗满}呢？难道不应因前面说\Quote{及至你的日子满足，与你祖先同眠时，我必在你以后兴起你的后裔}，而理解这是应许另一位和平者，祂被预言将要在\Person{达味}死后才兴起，不像\Person{撒罗满}那样在\Person{达味}生前兴起吗？因为无论耶稣基督来临之前的时间间隔有多长，毫无疑问，祂是在\Person{达味}王死后才应当来临的，\Person{达味}正是接受了这一应许。这位基督要建造天主的殿，不是用木石，而是用人，我们欣然看到祂正在建造这样的殿。因为关于这殿，即关于信徒，宗徒说：\BibleQuote{天主的殿是神圣的，而这殿就是你们。}\BibleRef{格前 3:17}\par

% structure: p0037 source=h2 target=subsection reason=relative-heading-rank; base=h2

\subsection{第九章——论\BibleBook{圣咏}{第八十九篇}关于基督的预言与\BibleBook{撒慕尔}{纪}中\Person{纳堂}预言所应许之事如何相似。}

因此，在第八十九篇圣咏中——其标题为\Quote{厄堂以色列人的训诲}——提及了天主向\Person{达味}王所作的许诺，并加入了一些与《撒慕尔纪》相似的词句，例如：\BibleQuote{我向我的仆役达味起誓，我要永远为他预备后裔。}又说：\BibleQuote{那时你曾在异象中向你的儿子们发言说：我已将扶助加给那有能者，从我百姓中高举了那被选者。我寻获了我的仆役达味，用我的圣油傅油了他。因为我的手必扶助他，我的臂膀必坚强他。仇敌不能胜过他，邪恶之子不再伤害他。我必在他面前击倒他的敌人，恨他的人我必使他们逃跑。我的真理和我的慈爱必与他同在，因我的名，他的角必被高举。我要将他的手安置在海上，他的右手安置在河上。他必向我呼号说：\InnerQuote{你是我的父，我的天主，我救恩的担保者。}我也要立他为我的首生者，为地上诸王中至高者。我必为他永远保留我的慈爱，我的盟约必与他坚定。我也必使他的后裔存到永远，他的宝座如天之寿数。}这些话语，若正确理解，都是关于主耶稣基督的，以\Person{达味}之名，因为同一中保从\Person{达味}后裔的童贞女那里取了仆人的形状\BibleRef{斐 2:7}。因为紧接着，便论及他子女的罪恶，正如《撒慕尔纪》所载，更容易被认为是指\Person{撒罗满}。因为在那里，即在《撒慕尔纪》中，祂说：\BibleQuote{他若作恶，我必用人的棍杖管教他，用世人的鞭伤责罚他；但我的慈爱必不从他身上收回}\BibleRef{撒下 7:14}，\Quote{鞭伤}意指纠正的打击。因此有那句话：\Quote{不要伤害我的受傅者。}因为这难道不就是\Quote{不要伤害他们}吗？但在圣咏中，当祂似乎以\Person{达味}发言时，祂也在那里说了类似的话。祂说：\BibleQuote{倘若他的子孙离弃我的法律，不遵行我的典章；亵渎我的正义，不守我的诫命；我必用棍杖惩罚他们的过犯，用鞭伤责罚他们的罪孽；但我的慈爱必不从他身上废去。}祂没有说\Quote{从他们身上}，尽管祂说的是他的子孙，而非他本人；但祂说\Quote{从他身上}，若正确理解，意思相同。因为对于基督本身——教会的头——不可能找到任何需要用人的惩罚来约束、而慈爱仍继续的罪；但这些罪存在于祂的身体和肢体，即祂的子民中。因此，在《撒慕尔纪》中说的是\Quote{他的过犯}，而在圣咏中说的是\Quote{他子孙的过犯}，为使我们明白，凡关于祂身体所说的，在某方面也是关于祂自己说的。因此，当\Person{撒乌耳}迫害祂的身体，即祂的信徒时，祂亲自从天上下来说：\BibleQuote{扫禄，扫禄，你为什么迫害我？}\BibleRef{宗 9:4}。接着在圣咏的后续话中，祂说：\BibleQuote{我必不违背我的真理，也不亵渎我的盟约；我嘴唇所出的我必不收回。我一次指着我的圣洁起誓：我决不向达味说谎}——即我绝不向达味说谎；因为圣经常这样表达。至于祂在什么事上不说谎，祂补充说：\BibleQuote{他的后裔必存到永远，他的宝座在我面前如太阳，又如永远完美的月亮，在天上有忠信的证人。}\par

% structure: p0039 source=h2 target=subsection reason=relative-heading-rank; base=h2

\subsection{第十章——地上\Place{耶路撒冷}王国中的行为与天主所应许的何其不同，以致应诺的真理应被理解为关乎另一位王及其王国的荣耀。}

为避免有人以为如此强烈表达并确认的应许是在\Person{撒罗满}身上应验的，仿佛他曾期望却未得着，祂说：\BibleQuote{但祢却弃绝了，并使之归于乌有，上主。}这确实发生在\Person{撒罗满}后代中间的王国，直至地上\Place{耶路撒冷}本身——王国所在地——的倾覆，尤其是\Person{撒罗满}所建的圣殿的毁灭。但为避免因此以为天主违背了祂的应许，祂紧接着说：\BibleQuote{祢延迟了祢的基督。}所以祂不是\Person{撒罗满}，也不是\Person{达味}本人，如果主的基督被延迟了。因为所有的王都被称为祂的受傅者，他们是用那神秘的圣油祝圣的，不但是自\Person{达味}王以下，甚至自那最初受傅为该民族之王的\Person{撒乌耳}——\Person{达味}本人也称他为上主的受傅者——然而有一位真正的基督，他们藉着先知性的傅油预表祂，按人的意见，人们以为祂在\Person{达味}或\Person{撒罗满}身上已经来到，但实际上祂被长期延迟，但按天主的安排，祂要在自己的时期来到。这篇圣咏的后续部分叙述了在祂被延迟期间，地上\Place{耶路撒冷}的王国——人们希望祂必在那里为王——将遭遇什么：\BibleQuote{祢推翻了祢仆人的盟约；祢亵渎了他在地上的圣所。祢拆毁了他所有的墙壁；祢使他的堡垒陷入恐惧。一切过路的人都劫掠他；他成了邻人的羞辱。祢高举了他敌人的右手；祢使他的众敌都欢乐。祢使他的刀剑退缩，在战争中未帮助他。祢消灭了他，使他不得洁净；祢把他的宝座摔在地上。祢缩短了他宝座的日子；祢将耻辱倾倒在他身上。}这一切都临到了婢女\Place{耶路撒冷}身上，在那里也有一些自由妇的儿女作王，他们以暂时的管治执掌那王国，但以真信德持守天上\Place{耶路撒冷}的王国——他们是其儿女——并寄望于真正的基督。但这些事如何临到那王国，历史记载了，若有人去读便知。\par

% structure: p0041 source=h2 target=subsection reason=relative-heading-rank; base=h2

\subsection{第十一章——论天主子民的本质，这本质因祂取得肉身而在基督内，惟有基督有能力将自己的灵魂从地狱中救出。}

但在预言了这些事之后，先知便转向祈祷天主；然而，这祈祷本身也是预言：\Quote{主，祢要转离到何时为止？}\Quote{祢的面容}是被理解的，正如别处所说：\Quote{祢转离你的面容不顾我要到何时？}因此，有些抄本在此处不是\InnerQuote{祢转离}，而是\InnerQuote{祢要转离}；虽然也可理解为：\Quote{祢转离了祢应许给达味的仁慈。}但祂说\Quote{到何时为止}，这除了指直到终结，又是什么意思呢？这终结应理解为末时，那时连那个民族也要相信基督耶稣；在此之前，祂刚才悲伤哀叹的事必将发生。正因如此，此处又加上：\Quote{祢的忿怒如火焚烧。求祢记念我的本质是什么。}这最好理解为耶稣自己，祂是祂子民的本质，祂的肉体即是其本性。\Quote{因为祢并非徒然}，祂说，\Quote{创造了所有的人子。}因为除非那独一的人子是以色列的本质，藉着这人子许多的人子得以获得释放，所有的人子就全然是徒然受造了。但现在，事实上，全人类因着第一人的堕落而从真理落入了虚妄；为此，另一篇圣咏说：\Quote{人相似于虚妄：他的时日如影消逝；}然而天主并没有徒然创造所有的人子，因为祂藉着中保耶稣将许多人从虚妄中解救出来；而那些祂预知不会被解救的人，祂在整全理性受造物最美丽、最公义的秩序中，并非全然徒然地造了他们，乃是为那些将被解救者的益处，以及为使两座城在相互对比中得以比较。此后接着是：\Quote{有哪一个人能活着而不见死亡？他能将自己的灵魂从地狱的手中抢救出来吗？}这不就是出于达味后裔的以色列的本质、基督耶稣吗？关于祂，宗徒说：\Quote{祂从死者中复活，不再死，死亡再不能统治祂了。}\BibleRef{罗 6:9}因为祂要如此活着而不见死亡，以致祂曾经死过；但祂却将自己的灵魂从地狱的手中解救出来，祂曾下到那里，为释放一些被地狱锁链束缚的人；祂是藉着那在福音中所说的权能解救的：\Quote{我有权舍掉我的性命，也有权再取回它。}\BibleRef{若 10:18}\par

% structure: p0043 source=h2 target=subsection reason=relative-heading-rank; base=h2

\subsection{第十二章——当他在圣咏中说\Quote{主，祢往日的怜悯在哪里？}等等时，祈求应许的话应理解为属于谁。}

但这篇圣咏的其余部分如下：\BibleQuote{主，祢的古老怜悯何在？祢曾凭祢的真理向达味起誓。主，求祢记念祢仆人所受的凌辱，就是我在我的怀中由许多民族所承当的；主，祢的敌人怎样凌辱了，他们怎样凌辱了祢的基督的转变。}如今很有理由问：这是否出于那些希望达味所获的许诺能向他们应验的以色列人之口？还是出于基督徒，他们是按圣神而不是按血气作以色列人的？\BibleRef{罗 3:28}这确实是在厄堂的时代所说出或写下的，这篇圣咏的标题即取自他之名，那与达味统治的时代相同；因此，若不是先知扮演了那些很久以后要来的人，对他们而言，这些事被应许给达味的时间已是古代，就不会说：\BibleQuote{主，祢的古老怜悯何在？祢曾凭祢的真理向达味起誓？}但也可以这样理解：许多民族在迫害基督徒时，以基督的受难来凌辱他们，圣经称之为祂的转变，因为祂借着死亡而成为不朽。根据这段经文，基督的转变也可以被理解为以色列人所凌辱，因为当他们希望祂是属于他们的时，祂却成了万民的救主；许多借着新约而相信祂的民族，如今以此凌辱那些仍留在旧约里的人；因此说：\BibleQuote{主，求祢记念祢仆人所受的凌辱。}因为主没有忘记，反而怜悯他们，连他们在这凌辱之后也要相信。但我首先提出的似乎是最合适的含义。因为对于基督的敌人，即那些因基督离开了他们转而向外邦人而被凌辱的人\BibleRef{宗 13:46}，把这番话\BibleQuote{主，求祢记念祢仆人所受的凌辱}加给他们是不合宜的，因为这样的犹太人不得称为天主的仆人；但这话却适合那些因基督之名受迫害而遭遇极大羞辱的人，他们能回想曾应许给达味后裔的崇高国度，并渴求它，不是绝望地，而是借着祈求、寻找、叩门\BibleRef{玛 7:7}说：\BibleQuote{主，祢的古老怜悯何在？祢曾凭祢的真理向达味起誓？主，求祢记念祢仆人所受的凌辱，就是我在我的怀中由许多民族所承当的。}也就是说，在我内心深处耐心忍受。\BibleQuote{主，祢的敌人怎样凌辱了，他们怎样凌辱了祢的基督的转变。}不认为那是转变，而是毁灭。但\BibleQuote{主，求祢记念}是什么意思，岂不就是愿祢怜悯，并因我耐心忍受的羞辱，以祢凭祢的真理向达味所起的尊荣来赏报我吗？但若我们把这话归于犹太人，那些在耶稣基督按人的样式诞生前，在世俗的耶路撒冷被攻陷时被掳为囚的天主仆人，他们能说这样的话，理解基督的转变，因为借着祂，确实应当期待的不是属世和属肉体的幸福，如同在撒罗满王那几年所显现的，而是属天和属神的幸福；当那些民族因不信而对此无知，并为着被掳而得意洋洋地凌辱天主的子民时，除了无知地凌辱那些理解基督转变的人所理解的基督转变，还能是什么呢？因此，当这篇圣咏结束时，接下来的话\BibleQuote{愿主的祝福永远常存，阿们，阿们}完全适合属于天上耶路撒冷的天主全体子民，无论是对新约启示前隐藏在旧约中的那些事，还是对新约中现已启示、清楚可见属于基督的那些事。因为主在达味后裔中的祝福不属于任何特定时期，如撒罗满的日子所显现的，而是永远可期待的；在这最确定的希望中说\BibleQuote{阿们，阿们}，因为这词的重复就是那希望的确认。因此达味明白这事，在\BibleBook{列王}{纪下}我们岔开去谈这篇圣咏的那段中说：\BibleQuote{祢也为祢仆人的家说了将来长远的事。}\BibleRef{撒下 7:19}因此稍后他又说：\BibleQuote{现在开始，并永远祝福你仆人的家}等等，因为那时即将生下一个儿子，他的后裔将延续到基督，借着祂，他的家将永远长存，并也将成为天主的家。因为它被称为达味的家，是出于达味的血统；但同一个家被称为天主的家，是由于天主的圣殿，由人而非石头所建，在那里，子民将永远与他们的天主同在并住在天主内，天主也与祂的子民同在并住在祂的子民内，以致天主充满祂的子民，子民也被他们的天主充满，同时天主将成为万有中的万有，祂自己在和平中是他们的赏报，在战争中是他们的力量。因此，在纳堂的话中说：\BibleQuote{主要告诉你，你要为祂建造怎样的家。}\BibleRef{撒下 7:8}之后在达味的话中说：\BibleQuote{因为祢，主全能者，以色列的天主，开启了祢仆人的耳朵，说：我要为你建造一个家。}\BibleRef{撒下 7:2}因为这家既由我们借着好好生活而建造，也由天主借着帮助我们好好生活而建造；因为\BibleQuote{若不是主建造房屋，建造的人就枉然劳力。}当这家的最终奉献举行时，天主在这里借着纳堂所说的话就要应验：\BibleQuote{我要为我的子民以色列指定一个地方，栽植他，他要独居，不再受扰；邪恶之子不得再像起初，从我设立民长治理我的子民以色列的日子以来那样欺压他。}\BibleRef{撒下 7:10}\par

% structure: p0045 source=h2 target=subsection reason=relative-heading-rank; base=h2

\subsection{第十三章—所许诺的和平之真理是否能归于撒罗满治下的那些已逝时光。}

凡在今世、在这地上希望获得如此巨大福乐的人，他的智慧不过是愚妄。难道有人能以为这在撒罗满统治的和平中实现了吗？圣经确实以卓越的赞美将那和平作为将来之事的预像。但必须警醒地反对这种见解，因为经上先说：\BibleQuote{不义之子再不能谦抑他}，紧接着又加上：\BibleQuote{如起初，自从我任命民长管理我人民以色列的日子以来}。\BibleRef{撒下7:10}因为自从他们承受了应许之地以后，在未有君王之前，就有民长被任命管理那人民。的确，不义之子，即外邦的敌人，曾屡次谦抑他们，那时我们读到和平与战争交替；并且那段时期中，有比撒罗满更长的和平时期，撒罗满只统治了四十年。因为在那位称为厄胡得的民长之下，曾有八十年的和平。\BibleRef{民3:30}因此，我们切不可相信这应许所预言的是撒罗满的时代，更不用说任何其他君王了。因为别的君王没有一个像他那样统治在如此巨大的和平中；那民族也从未拥有过这样的王国，以至于不必担心被敌人征服：因为在人事极其多变的境况中，没有任何民族能享有如此大的安全，以致不惧怕对今生有害的侵略。因此，这所应许的和平安居之所是永恒的，且理所当然地永恒属于自由之母耶路撒冷，真正的以色列人民将居住在那里：因为这个名称被解释为\Quote{看见天主}；为了渴求这赏报，应在这悲惨的旅途中，借着信德度虔敬的生活。\par

% structure: p0047 source=h2 target=subsection reason=relative-heading-rank; base=h2

\subsection{第十四章——论达味在写作圣咏中的参与}

因此，当天主之城历经各世代进展时，达味首先在属地的耶路撒冷作王，作为将来之事的预像。达味是一位擅长歌咏的人，他深切喜爱音乐的和谐，并非出于庸俗的乐趣，而是怀着信德的态度，并借此事奉他的天主，即真天主，因着一件伟大之事的奥秘象征。因为不同声响在和谐的多样中的合理而有秩序的协调，暗示了秩序井然之城的紧凑统一。他的预言几乎全部在圣咏中，其中一百五十篇包含在我们称为 \WorkTitle{圣咏集} 的书中；有些人认为只有那些题有\OriginalTerm{属于达味}{Of David}的，才是达味所作。但也有人认为，除了那些标有\Quote{属于达味}的之外，其余都不是他写的；而那些标题有\OriginalTerm{为达味}{For David}的，则是别人假托他的名义所作。这一见解被救主在福音中的亲口之言所驳斥，因为他说达味自己藉着圣神称基督为主；因为第一百一十篇圣咏以此开始：\BibleQuote{上主对我主说：你坐在我右边，等我使你的仇敌作你的脚凳。}而那一篇圣咏，如同许多别的，标题是\Quote{为达味}，而不是\Quote{属于达味}。但在我看来，那些认为达味是全部一百五十篇圣咏的作者，并相信他在某些篇前冠以其他人的名字（这些名字预表了与内容有关的事），而对于其余则不愿在标题中出现任何人的名字的人，所持的见解更可信；正如天主在这多样性（虽晦暗，却非无意义）的安排中感动了他。也不应因在书中某些圣咏的题词中读到一些比达味王时代晚得多的先知的名字，以及其中所说的话仿佛是藉着他们说的，就因而妨碍我们相信这一点。因为先知之神并非不能启示达味王（当他预言时）这些未来先知的名字，以便他能预言性地唱出适合他们身份的歌；正如曾启示某一位先知，约史雅王将在三百多年后兴起并作王，同时预言了他未来的事迹连同他的名字。\par

% structure: p0049 source=h2 target=subsection reason=relative-heading-rank; base=h2

\subsection{第十五章——圣咏中关于基督及其教会的一切预言是否都应纳入本文}

现在，我看出有人可能期望我在本书的这一部分阐明达味在圣咏中关于主耶稣基督或他的教会预言了什么。然而，尽管我已在一个实例中这样做了，但我之所以不能按那期望所要求的那样去做，与其说是材料匮乏，毋宁说是材料太多。因为避免冗长的必要性阻止我录下全部；然而我担心，若我选择一些，那些熟悉这些事的人会认为我遗漏了更必需的部分。此外，所援引的证据应当由整篇圣咏的上下文来支持，以便至少当并非所有内容都支持它时，也没有任何事物反对它；免得我们效法集句诗的方式，为所欲之事搜集，仿佛是出自一首伟大诗篇的诗句，而它们可能被证明不是论及那事，而是论及其他大不相同的事。但在每篇圣咏中能够指出这一点之前，必须解释整篇；而这是一项多么浩大的工程，别人以及我们自己的那些已完成的卷册已充分表明。因此，凡愿意或能够读这些卷册的人，就会发现达味，既是君王又是先知，在圣咏中预言了多少伟大的事，关于基督和他的教会，即关于君王和他所建造的城。\par

% structure: p0051 source=h2 target=subsection reason=relative-heading-rank; base=h2

\subsection{第十六章——论第四十五篇圣咏中明确或比喻性地提及基督与教会的事}

因为无论关于什么事，凡有直接而明显的预言性陈述，那些比喻性的陈述都必须与它们混合；这些比喻性的陈述，主要是由于那些悟性较慢的人，便将费力阐明和解释它们的劳苦任务加于更有学识的人身上。的确，其中有些陈述，在刚一说出时，便立即展现出基督和教会，尽管其中有些不太明白的东西仍有待于从容解释。我们在同一部\WorkTitle{圣咏集}中有一个例子：\BibleQuote{我的心涌出美辞：我向君王朗诵我的诗句。我的舌头是书记的笔，书写迅速。你在世人中最为美丽：恩宠倾注在你的唇上，因此天主永远祝福了你。大能者啊，佩剑在腰间！以你的荣耀和威仪，勇往直前，为真理、温和和正义而统治；你的右手必为你施行奇事。你的箭锋利无比，射中君王仇敌的心；万民都仆倒在你脚下。天主，你的宝座永远常存；你王国的权杖是正直的权杖。你喜爱正义，憎恨邪恶；因此天主，你的天主，以喜乐之油膏你胜过你的同伴。没药、香料和肉桂芬芳你的衣袍，发自象牙宫中，在其中君王们使你喜欢于你的尊荣。} 有谁，无论多么迟钝，若听到祂是天主，祂的宝座永远常存，祂被天主傅油——如同天主确实傅油，不是用可见的，而是用属神和可理解的圣油——而不在此认出我们所宣讲、所信仰的基督呢？因为谁在此宗教中如此无知，或对其广泛传播的声名如此聋聩，以至于不知道基督之名源自这圣油，即这傅抹呢？但一旦承认这位君王就是基督，每个已臣服于这因真理、温和和正义而统治的君王的人，便可以悠闲地追究这些在此以比喻方式说出的其他事情：祂的容貌如何在世人中最为美丽，具有一种越非形体的便越可爱可赞的美；以及祂的剑、箭和此类其他事物可能是什么，这些都是不按字面、而是按比喻记载的。\par

那么，让他观看祂的教会，与她如此伟大的新郎在属灵的婚姻和神圣的爱情中结合，关于此事在以下的话中这样说：\Quote{王后身穿金绣衣裳，佩带各样珠宝，侍立在祢右边。女儿，请听，请看，请侧耳细听：忘却妳的民族和妳父家！因为君王恋慕妳的美丽；祂是妳的主，妳应崇拜祂。提洛的女儿们要来奉献礼品，民间的富有人士也要寻求祢的容颜。王女的荣耀全在宫内，她的衣裳以金丝绣成。她身穿锦衣华服，被引到君王面前；陪伴她的童女也随从她，被带到祢面前。她们在欢乐喜乐中被引导，进入君王的宫殿。祢的子孙要继承祢的祖先，祢要立他们为普天下的首领。他们必将纪念祢的名号于万代；因此列国要永远称谢祢，直到永远。}我不认为有人如此愚蠢，竟相信此处所赞扬和描述的是一个可怜的女子，作为祂的配偶，即对祂所说：\Quote{祢的御座，天主，永存不朽；祢的国权是正义的权杖。祢喜爱正义，憎恨邪恶；因此天主，祢的天主，给祢傅了喜乐的油，胜过祢的同伴。}那就是明明白白的：\Person{基督}在基督徒之上。因为这些人都是祂的同伴，从他们的合一与和谐中，在万民中形成了那位王后，正如在另一圣咏中论到她所说：\Quote{伟大君王的城。}这同一位在属灵上是\Place{熙雍}，其拉丁文名称被解释为\OriginalTerm{发现}{speculatio}；因为她眺望来世的伟大美善，因为她的目光转向那里。同样，她在属灵上也是\Place{耶路撒冷}，关于我们已说了许多。她的仇敌是魔鬼的城\Place{巴比伦}，被解释为\Quote{混乱}。然而，从这巴比伦中，这位王后藉着新生而在万民中获得自由，并从最坏的君王——即魔鬼——归于最好的君王\Person{基督}。因此对她说：\Quote{忘却妳的民族和妳父家。}这邪恶之城也包括那些仅在肉体上是以色列人、而非因信德的人，他们也是这位伟大君王及其王后的仇敌。因为\Person{基督}来到他们那里，被他们杀害，却更成为其他在肉体上未曾见过祂的人的王。因此我们的君王亲自藉某一圣咏的预言说：\Quote{祢救我脱离人民的违抗，祢立我为万国的首领；我未曾认识的民族事奉了我；他们一听到耳朵所听的，就服从了我。}因此这外邦民族，\Person{基督}在肉身内并不认识，却相信了所传报的基督；所以有理由对它说\Quote{他们一听到耳朵所听的，就服从了我}，因为\Quote{信德出自听闻}。\BibleRef{罗 10:5}我说，这民族，再加上那些在肉体和信德上都是真以色列人的，就是天主的城，她按肉体生了\Person{基督}自己，因为\Person{基督}只在这些以色列人中。因为\Person{童贞玛利亚}从那里出来，\Person{基督}在她内取了肉体，为能成为人。论到这座城，另一圣咏说：\Quote{母亲熙雍，人要说，这人在她内出生，至高者自己奠定了她。}这至高者是谁，除了天主？因此，\Person{基督}，祂是天主，在祂藉\Person{玛利亚}在那城里降生成人以前，祂自己藉着圣祖和先知奠定了她。因此，正如如此久远以前在预言中对这位王后——天主的城——所说的，我们如今已能看到应验：\Quote{祢的子孙要继承祢的祖先，祢要立他们为普天下的首领。}同样，从她的儿子中，也确实在普天下设立了她的父亲（首领），当万民聚集到她那里，以永远的赞颂永远称谢她时。毫无疑问，无论对这里以比喻言语稍显隐晦表达的内容作何解释，都应与此等最明显的事实相符。\par

% structure: p0054 source=h2 target=subsection reason=relative-heading-rank; base=h2

\subsection{第十七章．论\WorkTitle{圣咏第一百一十篇}中关于\Person{基督}的司祭职的事，以及\WorkTitle{圣咏第二十二篇}中关于祂的受难的事。}

正如在那篇圣咏中，基督被最公开地宣告为司铎，如同祂在此处被宣告为君王一样，\BibleQuote{主对我主说：你坐在我右边，等我使你的仇敌作你的脚凳。}基督坐在天主圣父的右边，是信德之事，非目之所见；祂的仇敌被置于祂脚下，目前尚未显现；此事正在成就，因此终将显现：是的，现在相信，日后必得看见。但紧接着的经文：\BibleQuote{主从\Place{熙雍}伸出你能力的权杖，你要在你仇敌中统治他们。}是如此明确，以至于否认它不仅是无信与错误，更是明目张胆的无耻。甚至连仇敌也必须承认，从\Place{熙雍}传出了基督的法律，即我们所称的福音，并承认它是祂能力的权杖。但祂在仇敌中统治，这些在祂统治下的仇敌自身也为此作证，咬牙切齿，消耗殆尽，却无力反对祂。然后稍后所说的话：\BibleQuote{主发了誓，决不反悔。}祂以这些话暗示祂所添加的是不变的：\BibleQuote{你按照\Person{默基瑟德}的品位，永为司铎。}谁还允许怀疑这些话说的是谁？眼看如今已无处有按照\Person{亚郎}品位的司铎和祭献，而处处有人在基督这位司铎之下奉献，这正是\Person{默基瑟德}祝福\Person{亚巴郎}时所预表的。因此，当正确理解时，同一圣咏中那些记载得稍为隐晦的话，都应参照这些明显的事。我们在通俗讲道中已揭示了如何正确理解这些事。同样，在那篇圣咏中，基督借预言说出的祂受难的屈辱：\BibleQuote{他们刺穿了我的手和脚，我数尽了我所有的骨骼。他们又观看，注视了我。}祂无疑是用这些话指祂的身体被挂在十字架上，手和脚被钉透穿孔，祂就这样使自己成为观看和注视者眼中的景象。祂又加上：\BibleQuote{他们分了我的衣裳，为我的外衣拈阄。}这段预言如何应验，福音史已有叙述。诚然，其他说得不太明显的话，若与那些如此清晰的事物相符，也得以正确理解；尤其因为那些我们不视为过去、而是视作现在的事物，正被全世界看见，如今正如在这篇圣咏中很早以前所预言的那样展现出来。因为稍后那里说：\BibleQuote{大地四极都要记起，并归向主；万民的宗族都要在祂面前敬拜；因为王国是主的，祂要统治万邦。}\par

% structure: p0056 source=h2 target=subsection reason=relative-heading-rank; base=h2

\subsection{第十八章——论\BibleBook{}{圣咏集}第三、第四十一、第十五和第六十八篇，其中预言了主的死亡与复活。}

关于祂的复活，\BibleBook{圣咏}{集}的神谕也绝非缄默。因为在第三篇圣咏中，以祂的身份所唱的，除了这篇以外还有什么呢？\BibleQuote{我躺下安睡，我醒了，因主扶持了我。}难道会有如此愚钝的人，竟以为先知选择向我们指出祂睡了又醒来是一件大事，除非那睡眠是死亡，那醒来是复活，而这正是必须如此预言关于基督的吗？因为在第四十一篇圣咏中也显示得更清楚，那里以中保的身份，照例把预言为未来之事叙述为已成过去，因为这些未来之事在天主的预定和预知中，如同已经成就的一样，因为它们是确定的。祂说：\BibleQuote{我的仇敌恶言辱骂我：他何时死，他的名字几时消灭？若有人来探望我，口说假话，心怀恶意；他一出去，立即宣扬。恨我的人聚在一起窃窃私议，他们合计谋想害我的性命。他们对我计划了不义的事。那睡觉的岂不也再起来吗？}这些话确实如此记载，使人理解为他无异于说：那死了的岂不重新得生命吗？前文清楚表明，祂的敌人图谋并计划了祂的死，而这死是由那进来探望又出去出卖的人所执行的。但在这里，谁不想到犹达斯呢？他本为祂的门徒，却成了出卖祂的人。因此，因为他们将要实行他们所阴谋的——即要杀害祂——祂为向他们表明，他们怀着徒然的恶意要杀害那将要复活的一位，就加上这一节，仿佛说：你们在做什么虚妄的事？你们的罪行将成为我的睡眠。\BibleQuote{那睡觉的岂不也再起来吗？}然而他在随后的几节中表明，他们不会无罪地犯下如此大不敬，说：\BibleQuote{就连那与我相安、我信赖的人，那吃我饭的人，也举脚踢我。}即践踏我。\BibleQuote{但是祢，}他说，\BibleQuote{主，求祢怜悯我，使我起来，好报复他们。}如今看见犹太人，在基督受难和复活之后，被战争屠杀和毁灭从他们的住处连根拔起，谁能否认这一点呢？因为祂虽被他们杀害，却已复活，并以暂时的惩罚报复了他们，只是为那些不悔改的人，祂保留着直到祂审判活人死人的时候。因为主耶稣基督亲自向宗徒们指出那人便是出卖祂的人时，引用了这篇圣咏的同一节，说它已应验在祂身上：\BibleQuote{那吃我饭的人，举脚踢我。}但祂所说\BibleQuote{我所信赖的}这一句，并不适用于头，而是适用于身体。因为救主自己并非不知道那一位，祂早已说过：\BibleQuote{你们中有一个是魔鬼。}\BibleRef{若 6:70}但祂惯于承担祂肢体的位格，并将他们应说的话归于自己，因为头和身体是一个基督\BibleRef{格前 12:12}；因此福音中有那句话：\BibleQuote{我饿了，你们给了我吃的。}\BibleRef{玛 25:35}解释这句话时，祂说：\BibleQuote{你们为我这些最小兄弟中的一个所做的，就是为我做的。}\BibleRef{玛 25:40}因此祂说祂曾信赖，因为那时祂的门徒们关于犹达斯曾信赖他；因为他曾列于宗徒之中\BibleRef{宗 1:17}。\par

但是犹太人并不期望他们所期待的基督会死；因此他们不认为我们这一位是法律和先知所宣布的那一位，而是自己想象我不知道他们中的哪一位，免于死亡的痛苦。因此，以奇妙的空虚和盲目，他们争辩说我们所引用的那些话所指的，不是死亡和复活，而是睡眠和再醒。但是第十六篇圣咏也向他们呼喊：\BibleQuote{因此我心欢乐，我的舌欢跃，我的肉身也要安息于希望中；因为祢不会将我的灵魂遗弃在地狱，也不会让祢的圣者见到腐朽。}除了第三天复活的那一位，谁可以说他的肉身曾安息于这希望中；他的灵魂没有被遗留在地方，而是速速返回肉身，使之复活，不致像尸体通常那样腐朽——这一点是先知和君王达味绝不能说的。第六十八篇圣咏也呼喊：\BibleQuote{我们的天主是施救的天主：主脱离生命是由于死亡。}还有什么比这说得更明白的呢？因为施救的天主就是主耶稣，这名字被解释为救主，或医治者。为此，在祂由童贞女诞生之前，曾说过：\BibleQuote{你将生一个儿子，要给他起名叫耶稣，因为他要把自己的子民从他们的罪恶中拯救出来。}\BibleRef{玛 1:21}因为祂的血为赦免他们的罪而倾流，所以祂脱离此生除死亡外别无出路。因此，当说了\BibleQuote{我们的天主是施救的天主}之后，立刻加上\BibleQuote{主脱离生命是由于死亡}，为表明我们将因祂的死而得救。但这句话是奇妙的：\BibleQuote{主脱离生命}，仿佛是说：凡人的生命就是如此，连主自己也不得不借死亡才能脱离它。\par

% structure: p0059 source=h2 target=subsection reason=relative-heading-rank; base=h2

\subsection{第19章——论第六十九篇圣咏，其中宣告了犹太人顽梗的不信。}

但是，当\Person{犹太人}丝毫不肯屈服于这如此明显、而且已由事件实现到如此清晰而确定的预言证据时，在他们身上无疑应验了接下来那篇圣咏中所记载的话。因为当那里也以\Person{基督}的位格预言性地谈及祂苦难的种种时，就提到了福音中所揭示的：\BibleQuote{他们拿苦胆给我作食物；我口渴时，他们拿醋给我喝。}而且，仿佛在这样将宴席和美味给予祂之后，祂立即引出了这些话：\BibleQuote{愿他们的筵席在他们面前成为网罗、成为报酬、成为绊脚石；愿他们的眼目昏花，不得看见；愿他们的腰时常弯曲}等等。这些话并不是作为愿望说的，而是以预言的愿望形式所预测的。那么，那些眼目昏花看不见的人，看不见这些明显的事，有什么稀奇呢？那些腰总是弯曲、以便在地上爬行的人，不仰望天上的事，有什么稀奇呢？因为这些从身体借来的词语，指的是心灵的瑕疵。关于圣咏，即关于达味王的预言，说了这些就已足够，使我们得以保持一定的限度。但请那些已经知晓这一切的读者原谅我们；也请他们不要抱怨那些他们知道或认为我忽略了的、或许更有力的证据。\par

% structure: p0061 source=h2 target=subsection reason=relative-heading-rank; base=h2

\subsection{第二十章——论\Person{达味}的统治与功绩；及其子\Person{撒罗满}，以及关于\Person{基督}的预言，这预言见于与他所写的书卷相联的那些书中，或见于无疑属于他的那些书中。}

因此，\Person{达味}在地上\Place{耶路撒冷}作王，他是天上耶路撒冷之子，受到天主见证的大力赞扬；因为甚至连他的过错也被伟大的虔诚所克服，通过他最有益的悔改谦卑，使他完全成为他自己所说的那样的人：\Quote{罪恶蒙赦免、过犯得遮盖的人，是有福的。}在他之后，他的儿子\Person{撒罗满}统治了同一整个百姓，如前所述，他在父亲还活着的时候就开始作王。这个人，在好的开端之后，结局很坏。因为确实，\Quote{顺境磨损智者的心智}这句话对他造成的伤害，超过了那智慧带给他的益处——那智慧至今仍闻名于世，将来也会如此，当时更是受到远近的赞扬。他还被发现曾在自己的书中预言，其中三卷被接受为正典权威，即\WorkTitle{箴言}、\WorkTitle{训道篇}和\WorkTitle{雅歌}。但还有另外两卷习惯上被归于\Person{撒罗满}，一卷称为\WorkTitle{智慧篇}，另一卷称为\WorkTitle{德训篇}，因为风格有些相似——但更有学识的人毫不怀疑它们不是他的作品；然而古时教会，尤其是西方教会，将它们接纳为权威——在其中一卷，即称为\WorkTitle{智慧篇}的书中，\Person{基督}的苦难被最公开地预言。因为确实引用了祂不虔敬的凶手们的话，说：\BibleQuote{我们要窥伺义人，因为他令我们厌烦，反对我们的行为；他指责我们违犯法律，把我们的教养的过失归为我们可耻的事。他自称拥有天主的智慧，并自称为天主的儿子。他生来就是为纠正我们的思想。我们连见到他也感到痛苦，因为他的生活与众不同，他的道路与别人不同。我们被他视为伪君子；他像避污秽一样避开我们的道路。他夸耀义人的结局，并以天主为父而自豪。我们且看他的话是否真实；让我们试验他将有什么结局，并知道他的结局是什么。因为如果义人是天主的儿子，天主必会保护他，将他从反对他的人手中救出。我们要以侮辱和酷刑来拷问他，好知道他的温良，考验他的忍耐。我们要判他受最可耻的死刑；因为按他自己的话，他必被尊重。}他们这样思想，就错了；因为自己的恶意完全蒙蔽了他们。\BibleRef{智 2:12} 而在\WorkTitle{德训篇}中，以这样的方式预言了万民未来的信仰：\BibleQuote{万有的主宰天主！求祢怜悯我们，并向万国投下祢的敬畏；举起祢的手打击异邦，使他们见到祢的权能。正如祢在他们面前在我们身上显示为圣，求祢也在我们面前在他们身上显示为圣，使他们认识祢，如同我们也认识祢一样；因为除祢以外，没有别的天主，上主啊。}\BibleRef{德 36:1} 我们看到这个以愿望和祈祷形式出现的预言，通过\Person{耶稣基督}实现了。但是，那些没有写在犹太人正典中的事物，不能以同样大的效力被引用来反驳他们的矛盾。\par

但至于那三卷书——显然是属于\Person{撒罗满}，并为犹太人所承认为正典的——要指出其中与基督和教会相关的这类内容，需要一番劳苦的讨论；如果现在着手，就会不当延长这部著作。然而，我们在\WorkTitle{箴言}中读到恶人说：\BibleQuote{让我们不义地将义人藏在地下；是的，让我们活活吞下他如阴间，并从地上除去他的纪念；让我们夺取他宝贵的财产。}\BibleRef{箴 1:11}这并非如此隐晦，以至无需费力的解释就无法理解为是指基督和祂的产业——教会。的确，福音中关于恶园户的比喻表明主耶稣亲自说了类似的话：\BibleQuote{这是继承人；来，我们杀了他，产业就成为我们的。}\BibleRef{玛 21:38}同样，同卷书中我们之前谈论不生子的妇女生了七个孩子时所触及的那段经文，在说出之后不久，必被那些知道基督就是天主的智慧的人理解为只关乎基督和教会。\BibleQuote{智慧建造了房屋，凿了七根柱子；宰杀了牲畜，调和了酒，也摆设了筵席。她打发使女出去，在城中高处呼喊说：谁是无知的，可以转到我这里来。她又对那缺乏心智的说：你们来，吃我的饼，喝我调和的酒。}在这里，我们确实感知到天主的智慧，即与圣父同永的圣言，为自己建造了房屋，即在童贞母腹中的人性身体，并将教会作为肢体连接于元首，宰杀了殉道者作为祭品，摆设了有酒和饼的筵席，在此也出现了按照默基瑟德品位的大司祭，并召叫了无知者和缺乏心智者，因为正如宗徒所说：\BibleQuote{天主选择了世上愚拙的，为叫那有智慧的羞愧。}\BibleRef{格前 1:27}但对这些软弱的人，智慧说：\BibleQuote{你们舍弃愚拙，就能生存；寻求明智，就能获得生命。}\BibleRef{箴 9:6}但分享这宴席本身就是开始拥有生命。因为当他在另一卷称为\WorkTitle{训道篇}的书中说：\BibleQuote{人除了吃喝，并无幸福}，还有什么比这更可信地理解为属于参与这宴席的事呢？——这宴席是新约的中保、按照默基瑟德品位的大司祭，以祂自己的体血所摆设的。因为那祭献已取代了旧约的所有祭献——它们作为那将要来临者的影子而被宰杀；因此我们也认出第四十篇\WorkTitle{圣咏}中那位同一中保借预言说话的声音：\BibleQuote{祭品和素祭祢不喜爱，却为我预备了身体。}因为，代替所有这些祭献和供物，祂的身体被奉献，并分给领受者。因为这位\WorkTitle{训道篇}的作者在关于吃喝的这句话上——他屡次重复，极力称赞——并不沾染肉欲的享乐，这点在他这样说时已足够显明：\BibleQuote{往丧家去，胜过往宴家去。}\BibleRef{训 7:2}稍后他又说：\BibleQuote{智慧人的心在哀悼之家；愚顽人的心在欢乐之家。}\BibleRef{训 7:4}但我认为从这卷书中更值得引用的，是关于两座城——一座属魔鬼，另一座属基督——以及它们的君王——魔鬼和基督——的话：他说：\BibleQuote{邦国啊，你当有祸！当你的君王是孩童，你的诸侯清晨宴饮！邦国啊，你是有福的！当你的君王是贵胄之子，你的诸侯按时宴饮，在刚毅中，不在混乱中！}\BibleRef{训 10:16}他称魔鬼为孩童，因为愚昧、骄傲、鲁莽、不受约束以及其他在这一年龄常见的恶习；但基督是贵胄之子，即圣祖们、属于自由之城的人们的儿子，祂按肉身是由他们而生。那城及其他城的诸侯是清晨宴饮者，即在合适时辰之前，因为他们不期待在来世中那真正的、合乎时宜的幸福，渴望立即以今世的声望而快乐；但基督之城的诸侯耐心等候一种不虚假的幸福的时刻。这由\BibleQuote{在刚毅中，不在混乱中}这些话表达，因为希望不欺骗他们；宗徒论此说：\BibleQuote{望德不叫人蒙羞。}\BibleRef{罗 5:5}一首\WorkTitle{圣咏}也说：\BibleQuote{因为仰望祢的，必不蒙羞。}但\WorkTitle{雅歌}是圣洁心灵的一种神性愉悦，关乎那位君王与女王之城——即基督与教会的婚姻。但这种愉悦包裹在比喻的帷幕中，为的是使新郎更被热切渴望，更喜悦地被揭示，得以显现；在这同一首歌中，对祂说：\BibleQuote{正直的人爱慕祢。}\BibleRef{歌 1:4}而新娘在那里听到：\BibleQuote{爱情在你的欢愉中。}\BibleRef{歌 7:6}我们静默地略过许多事，渴望完成这部著作。\par

% structure: p0064 source=h2 target=subsection reason=relative-heading-rank; base=h2

\subsection{第二十一章——论撒罗满之后犹大和以色列的君王。}

撒罗满之后的希伯来其他君王，无论是犹大还是以色列，几乎都未能以他们隐晦的言语或行为预言到有关基督和教会的事；因为从撒罗满犯罪之后，他的儿子勒哈贝罕继承王位之时起，天主作为惩罚将那个民族分成了两部分。事实上，撒罗满的仆人雅洛贝罕得到了十支派，在撒玛利亚被立为王，这些支派特别称为以色列，尽管这原是全体人民的名号；但犹大和本雅明这两个支派——为了达味的缘故，免得王权完全从他的家族被夺走——仍然隶属于耶路撒冷城，被称为犹大，因为达味是从那个支派出来的。本雅明，另一个如前所述属于同一王国的支派，是撒乌耳的支派，撒乌耳在达味之前为王。但这两个支派合在一起，如上所述，被称为犹大，并以此名与以色列区分，以色列是十支派在自己君王下的专属称号。至于肋未支派，因为它属于司祭，侍奉天主而非君王，被视为第十三个支派。因为若瑟，以色列十二子之一，并不像其他儿子那样形成一个支派，而是形成了两个支派：厄弗辣因和默纳协。然而，肋未支派也更属于耶路撒冷王国，那里有它侍奉的天主的圣殿。因此，在人民分裂之后，撒罗满的儿子勒哈贝罕作为犹大的第一位王在耶路撒冷统治，而撒罗满的仆人雅洛贝罕在撒玛利亚作以色列王。当勒哈贝罕想以暴君姿态用武力攻击那个分离的部分时，天主阻止人民与他们的兄弟作战，藉一位先知告诉他们这是天主所做的；由此看出，在这件事上以色列王或人民都没有罪，而是天主作为报复者的旨意得以成就。当这事被知晓后，双方和平地安定下来，因为所做的分裂不是宗教性的，而是政治性的。\par

% structure: p0066 source=h2 target=subsection reason=relative-heading-rank; base=h2

\subsection{第二十二章——论雅洛贝罕，他以偶像崇拜的不敬亵渎了托付给他的人民，然而天主并未停止感动先知，并保守许多人免于偶像崇拜之罪。}

但以色列王雅洛贝罕心存邪僻，不信天主——他曾证明天主在应许并赐予他王位上是信实的——他担心人民若来到耶路撒冷的天主圣殿（按天主的法律，整个民族都要来那里献祭），会从他那里被引诱，回归到达味的王族血统，于是在他的王国设立偶像崇拜，并以可怕的不敬迷惑人民，诱使他们与自己一同敬拜偶像。然而天主并未完全停止藉先知责备，不仅责备那王，也责备他的后继者和效法他不敬的人，以及人民。因为那里兴起了伟大而卓越的先知厄里亚和他的门徒厄里叟，他们也行了许多奇事。甚至在那里，当厄里亚说：\Quote{主啊，他们杀了你的先知，拆毁了你的祭坛，只剩下我一个人，他们还要寻索我的命。}时，有回答说在那里有七千人未曾向巴耳屈膝。\par

% structure: p0068 source=h2 target=subsection reason=relative-heading-rank; base=h2

\subsection{第二十三章——论两个希伯来王国的多变状况，直到两国人民先后被掳，犹大后来被召回自己的王国，最终落入罗马人之手。}

同样，在属于耶路撒冷的犹大王国，历代君王时期也不乏先知，正如天主乐意派遣他们，或为预言所需，或为矫正罪恶和训导正义\BibleRef{弟后 3:16}。因为在那里，尽管远少于以色列，也出现了因不敬而严重冒犯天主的君王，他们和与他们相似的人民一起，受到适度的鞭打。那里确实赞扬了虔诚君主的不少功德。但在以色列，我们读到君王有恶有更恶，但都是邪恶的。因此，每一部分，按天主眷顾所命令或所容许的，都因各种繁荣而被高举，或因各种逆境而被压垮；不仅受外敌之苦，也因彼此的内战而受苦，为的是藉某些存在的原因显明天主的仁慈或愤怒；直到因天主的日益愤怒，整个民族被战胜的加色丁人不仅在其居所被推翻，而且大部分被迁移到亚述人的土地——先是那十三个支派中称为以色列的部分，后来犹大也被掳，那时耶路撒冷和那最尊贵的圣殿被拆毁——在那地，他们被掳七十年。此后，他们被释放出来，重建了被推倒的圣殿。虽然有许多人留在外邦地，但王国不再有两个分开的部分，各有不同的君王，而是在耶路撒冷有一位首领治理他们；在某些时候，从四面八方，无论他们在哪里，无论他们能从何处来，他们都来到那里的天主的圣殿。然而即使那时，他们也没有外敌和征服者；是的，基督发现他们是罗马人的进贡者。\par

% structure: p0070 source=h2 target=subsection reason=relative-heading-rank; base=h2

\subsection{第二十四章——论那些或是犹太人中的最后一批先知，或是福音史所记载的关于基督诞生时期的先知。}

但在他们从巴比伦回来之后的那整个时期，从玛拉基亚、哈盖、匝加利亚（那时他们还在说预言）以及厄斯德拉以后，直到救主降临之时，没有别的先知，只有另一位匝加利亚，即若翰的父亲，和他的妻子依撒伯尔——当时基督的诞生已经近在咫尺；当基督已诞生后，有年迈的西默盎和已十分老迈的寡妇亚纳；最后，若翰本人，他虽年轻，却并非预言当时已年轻的基督将要来到，而是以先知性的认识指出这位尚不为人所知的基督；为此主自己说：\BibleQuote{法律及先知到若翰为止}\BibleRef{玛 11:13} 这五位先知的预言是在福音中为我们所知的，在那里，吾主的童贞母亲自己也被发现在若翰之前说了预言。但恶毒的犹太人不接受他们的预言；然而有无数的人从他们中间相信了福音，接受了这些预言。因为那时以色列确实一分为二，正如撒慕尔先知向撒乌耳王预言为不可变更的分裂。但就连被弃绝的犹太人也承认玛拉基亚、哈盖、匝加利亚和厄斯德拉为最后收入正典权威的先知。因为他们也像其他人一样，在先知的庞大群体中只有极少数人写了那些获得正典权威的书卷；他们的预言中，有些关乎基督和祂的教会，我想在这部著作中放入一些；靠着主的助佑，这将在下一卷中更为方便地完成，以免我们使本已过长的这一卷更加冗长。\par

\SourceRecord{Translated by Marcus Dods.；Nicene and Post-Nicene Fathers, First Series；Vol. 2.；Edited by Philip Schaff.；Buffalo, NY: Christian Literature Publishing Co.,；1887.；<http://www.newadvent.org/fathers/120117.htm>.}

% structure: p0001 source=h1 target=section reason=page-title

\section{天主之城（卷十八）}

奥斯定追溯了从亚巴郎到世界末了，这天主之城和地上之城的平行进程；并提及了关于基督的神谕，既有西比尔所发出的，也有在罗马建立之后写作的神圣先知欧瑟亚、亚毛斯、依撒意亚、米该亚及其后继者所发出的神谕。\par

% structure: p0003 source=h2 target=subsection reason=relative-heading-rank; base=h2

\subsection{第一章——论及在前面十七卷中已经讨论过的、直到救主时代的事情。}

我曾应许要写作两座城的起源、进展和预定的终结：一座是天主之城，另一座是此世之城，在这世中，就人类而言，前者如今是客旅。但首先我着手，在祂的恩宠许可下，驳斥天主之城的敌人，他们宁愿选择自己的神祇而不选基督这城的建立者，并以最恶毒的仇恨猛烈地仇恨基督徒。这我已在前面十卷中完成了。接着，就我刚才提及的三重应许而言，我在第十卷之后的四卷中清楚地论述了两座城的起源。此后，我从第一个人一直写到洪水，在本书的第十五卷中。再从那里一直写到亚巴郎，我们的作品按编年顺序进行。从族长亚巴郎直到以色列诸王的时代，我们在第十六卷结束时止；从那里直到基督亲自在肉身中的来临，第十七卷一直延伸到那个时期。天主之城在我的写作方式上似乎是单独行进的，然而它在现世并非单独行进，因为两座城在人类中行进的历程，确实从一开始就一起经历了变化无常的时期。但我这样做是为了：首先，从天主的应许开始变得更为清晰之时，直到那从起初所应许之事在祂身上应验的那位的童贞诞生，那座属于天主的城的历程可以更清楚地显明，而不被另一座城的历史中的外来内容所掺杂；尽管直到新盟约的启示，它是在阴影中而非光明中行进的。因此，现在我认为适合做我先前所略过的事，并尽可能必要地展示另一座城如何从亚巴郎的时代开始行进，以便细心的读者可以比较两者。\par

% structure: p0005 source=h2 target=subsection reason=relative-heading-rank; base=h2

\subsection{第二章——论地上之城的诸王与时代，它们与圣人们的时代是同时的，自亚巴郎兴起算起。}

凡人的社会散布到全地各处，虽然因我们共同本性的某种联系而结合在一起，却大多分裂自相为敌，最强者压迫他人，因为所有人都追逐自己的利益和欲望，而所贪求的要么对无人足用，要么对所有人不够，因为它并非那真实之物。因为被征服者屈服于得胜者，宁愿要任何一种和平与安全胜于自由本身；因此那些宁死不屈的人曾备受惊叹。因为在几乎所有民族中，那自然之声仿佛宣告：那些碰巧被征服的人应该选择臣服于征服者，而不是被各种战争毁灭所杀戮。这一切绝不离天主的眷顾而发生，在祂的权力之下，任何人在战争中要么征服要么被征服；有些人被授予王国，另一些人被置于国王的统治之下。然而，在这地上许多王国中——社会因世俗利益或欲望而分裂为这些王国（我们统称之为此世之城）——我们看到有两个王国，在时间和地点上彼此固定且区分开来，变得远比其余王国著名：首先是亚述王国，然后是罗马王国。首先出现一个，然后另一个。前者在东方兴起，紧接着它结束时，后者在西方兴起。我可以把其他王国和其他君主说成是它们的附属。\par

那时，继承其父\Person{贝鲁斯}（第一位亚述王）的\Person{尼努斯}，当\Person{亚巴郎}在\Place{加色丁}之地出生时，已是该国的第二位君主。当时还有一个很小的\Place{西基雍}王国，那位极其博学的\Person{马库斯·瓦罗}在撰写罗马民族史时，从远古开始，便是从这个王国起笔。因为他从\Place{西基雍}的这些国王写到雅典人，从雅典人写到拉丁人，再从拉丁人写到罗马人。然而，在罗马建城之前，关于这些王国的记载，与\Place{亚述}相比都显得微不足道。因为即使是罗马史学家\Person{撒路斯提乌斯}也承认，雅典人在希腊非常著名，但他认为其名声大于实际。他在谈到雅典人时说：\Quote{我认为，雅典人的事迹确实非常伟大辉煌，但仍略逊于传闻。但由于他们中间出现了天才的作家，雅典人的事迹便作为最伟大的事迹传遍天下。因此，那些行事者的德能，就被认为是与旷世天才凭借赞美之词所能表达的同样伟大。}这座城也从文学和哲学中获得了不小的荣耀，这些学问在那里最为兴盛。但就帝国而言，早期没有哪个王国比\Place{亚述}更强大或更广袤。因为当\Person{贝鲁斯}之子\Person{尼努斯}为王时，据说他征服了整个亚洲，甚至到了\Place{利比亚}的边界——\Place{利比亚}以数量计被称为第三部分，但以面积计则相当于整个世界的一半。只有东方的\Place{印度}人是他未能统治的；但在他死后，其妻\Person{塞米拉米斯}曾向\Place{印度}人开战。因此，那些地区的所有民族和国王都臣服于亚述人的王国和权威，凡事听命。现在，\Person{亚巴郎}就是在这个王国中的\Place{加色丁}之地，在\Person{尼努斯}时期出生的。但由于希腊的事情比亚述的事情更为我们所知，而那些致力于考察罗马民族起源的古代学者，都是通过希腊追溯到拉丁人，再从拉丁人追溯到罗马人——罗马人本身也是拉丁人——因此，在必要时，我们应该提及\Place{亚述}诸王，以使\Place{巴比伦}——如同第一个罗马——与在这世上为客旅的天主之城并行其历程。然而，在本作品中对这两座城，即地上之城与天上之城进行比较时，所应插入的内容，大部分应从希腊和拉丁王国中选取，因为罗马本身就像第二个\Place{巴比伦}。\par

在\Person{亚巴郎}出生时，\Place{亚述}和\Place{西基雍}的君主分别是第二位国王\Person{尼努斯}和\Person{欧罗普斯}，第一位国王分别是\Person{贝鲁斯}和\Person{埃吉阿留斯}。但当\Person{亚巴郎}离开\Place{巴比伦尼亚}时，天主应许他要成为一个大民族，并且万民都要因他的后裔而受祝福；那时亚述人已有其第七位国王，\Place{西基雍}人则有第五位。因为\Person{尼努斯}的儿子在其母\Person{塞米拉米斯}之后统治了他们——据说\Person{塞米拉米斯}因企图与儿子乱伦而被其子处死。有些人认为她建造了\Place{巴比伦}，或许她的确重建了它。但我们在第十六卷中已讲述了它是在何时、由何人建造的。\Person{尼努斯}和\Person{塞米拉米斯}的儿子，即继承其母王位的那位，有些人也称之为\Person{尼努斯}，另一些人则称之为\Person{尼尼阿斯}，这是一个父名化的词。随后，\Person{特莱克西翁}统治了\Place{西基雍}人。在他的统治时期，如此平静喜乐，以至于他死后，人们以献祭和举行竞技来敬拜他为神，据说竞技便是由此首次设立。\par

% structure: p0009 source=h2 target=subsection reason=relative-heading-rank; base=h2

\subsection{第三章——据应许，\Person{亚巴郎}百岁时\Person{依撒格}出生，以及\Person{依撒格}六十岁时\Person{黎贝加}为他生下孪生子\Person{厄撒乌}和\Person{雅各伯}时，\Place{亚述}和\Place{西基雍}有哪些君王在位。}

在他的时代，因天主的恩许，亚巴郎之子\Person{依撒格}在他父亲一百岁时，由他的妻子\Person{撒辣}所生，撒辣已年老不孕，早已对生育失去希望。那时\Person{亚拉略}是亚述人的第五位君王。\Person{依撒格}自己六十岁时，孪生子\Person{厄撒乌}和\Person{雅各}出生，是他的妻子\Person{黎贝加}所生，他们的祖父\Person{亚巴郎}当时仍在世，已活到一百七十岁去世，当时正值他一百六十岁之年。那时，第七位君王——在亚述人中有更古老的\Person{泽尔克塞斯}，也称为\Person{巴勒乌斯}；在西基雍人中有\Person{图里亚库斯}，或按某些人的写法，\Person{图里马库斯}。\Place{阿耳戈斯}王国，最初由\Person{伊那科斯}统治，兴起于\Person{亚巴郎}孙子的时代。我也不可省略\Person{瓦罗}所记述的，西基雍人还习惯于在他们的第七位君王\Person{图里亚库斯}的墓前献祭。在亚述的\Person{阿尔马米特雷斯}和西基雍的\Person{留喀波斯}的第八王位，以及阿耳戈斯的\Person{伊那科斯}初为王时，天主对\Person{依撒格}说话，并将与他父亲相同的那两件事应许给他——即\Place{客纳罕}地给他的后裔，以及万民因他的后裔蒙受祝福。同样的事也向他的儿子，\Person{亚巴郎}的孙子应许了，他起初称为\Person{雅各}，后来称为\Person{以色列}，当时\Person{贝洛库斯}是亚述的第九位君王，\Person{伊那科斯}之子\Person{福洛纽斯}以阿耳戈斯第二位君王统治，而\Person{留喀波斯}仍为西基雍的君王。在那些时代，在阿耳戈斯王\Person{福洛纽斯}治下，\Place{希腊}因某些法律和审判官的设立而更加闻名。\Person{福洛纽斯}死后，他的弟弟\Person{斐古斯}在他的墓旁建造了一座庙宇，在其中他被当作神明崇拜，并向他献祭公牛。我相信他们认为他配得如此大的尊荣，因为在他的那部分王国（他们的父亲曾将领土分给他们，他生前各自统治），他建立了敬拜诸神的小圣堂，并教导他们以月与年来度量时间，从而对事件进行计数和推算。尚未开化的人们因这些新奇事物而崇拜他，要么幻想他死后是神，要么决定使他成为神。据说\Person{伊俄}也是\Person{伊那科斯}的女儿，后来称为\Person{伊西斯}，在\Place{埃及}被当作大神崇拜；虽然其他人记载她是从\Place{埃塞俄比亚}来的女王，因为她统治广袤而公正，并为她的臣民创立了文字和许多有用的事物，死后在那里得到了如此神圣的尊荣，以至于若有人说她曾是人，就会被处以死罪。\par

% structure: p0011 source=h2 target=subsection reason=relative-heading-rank; base=h2

\subsection{第四章——论雅各及其子若瑟的时代。}

在亚述的第九位君王\Person{巴勒乌斯}和西基雍的第八位君王\Person{梅萨普斯}（有些人说他也叫\Person{刻菲索斯}——如果真是同一人有这两个名字，而那些写下另一个名字的人并没有将其与另一个人混淆的话）统治期间，\Place{阿耳戈斯}的第三位君王\Person{阿匹斯}在位时，\Person{依撒格}去世，享年一百八十岁，留下他的孪生子，当时已一百二十岁。这两个儿子中，较年幼的\Person{雅各}属于我们正在论述的天主之城（年长的完全被弃绝），他有十二个儿子，其中一个叫\Person{若瑟}，被他的兄弟们卖给下\Place{埃及}的商人，当时他的祖父\Person{依撒格}仍在世。但当\Person{若瑟}三十岁时，他站在\Person{法郎}面前，从他所受的屈辱中被高举，因为他神圣地解释了王的梦，预言将有七个丰年，极其丰盛的收成将被随后而来的七个荒年所耗尽。为此，王立他为全\Place{埃及}的统治者，将他从监牢中释放出来——他因保持自己的贞洁不受玷污而被投入监牢；他勇敢地保护自己免受女主人（她邪恶地爱他，并向轻信的主人说谎）的诱惑，没有同意与她通奸，而是从她手中逃离，留下了他的外衣在她手中，因为当她抓住他时。在七个荒年的第二年，\Person{雅各}带着他所有的一切下到\Place{埃及}到他儿子那里，那时他一百三十岁，正如他自己回答王的问题时所说。\Person{若瑟}当时三十九岁——如果我们将七个丰年和两个荒年加到他受王尊荣时所算的三十岁上。\par

% structure: p0013 source=h2 target=subsection reason=relative-heading-rank; base=h2

\subsection{第五章——论阿耳戈斯王阿匹斯，埃及人称他为塞拉匹斯，并以神圣尊荣敬拜。}

在这些时代，\Place{阿耳戈斯}王\Person{阿匹斯}乘船渡到\Place{埃及}，死在那里后，被立为\Person{塞拉匹斯}，所有埃及人的主神。\Person{瓦罗}给出了一个非常现成的理由，说明他死后为什么被称为\Person{塞拉匹斯}，而不是\Person{阿匹斯}。他死后被放置的棺木，现在人人都称为石棺，当时在希腊文中称为\OriginalTerm{棺材}{\GreekText{σορὸς}}，他们在神庙建成之前就开始崇拜被埋葬在其中的他；从\OriginalTerm{Soros}{\GreekText{σορὸς}}和\Person{阿匹斯}，他最初被称为\OriginalTerm{Sorosapis}{Sorosapis}或\OriginalTerm{Sorapis}{Sorapis}，然后因一个字母的变换而很容易地变成\Person{塞拉匹斯}。还曾颁布关于他的法令，凡说他曾是人的人将被处死。而且，在每一座敬拜\Person{伊西斯}和\Person{塞拉匹斯}的庙宇中，都有一尊雕像，以手指按在唇上，似乎在警告人保持沉默，\Person{瓦罗}认为这表示他们曾是人的事实应被保密。但是那头公牛——埃及人以惊人的愚昧用丰盛的美食喂养它来尊崇他——并不称为塞拉匹斯，而是阿匹斯，因为他们崇拜的是活的公牛，没有石棺。那头公牛死后，当他们寻找并发现了一头同样颜色——即有类似白色斑点的标记——的小牛时，他们相信这是一件神迹，是神圣地为他们预备的。然而，对于魔鬼来说，为了欺骗他们，在母牛受孕怀孕时向她显示这样一头公牛的形像（这形像只有她能看见），并借此吸引母牛的繁殖激情，以便在她的幼崽身上以有形体的形状显现，这并非难事，正如\Person{雅各}曾用有斑点的枝条如此操作，使绵羊和山羊生出有斑点的幼崽一样。因为人能用真实的颜色和物质做的事，魔鬼很容易通过向繁殖动物显示不真实的形像来做到。\par

% structure: p0015 source=h2 target=subsection reason=relative-heading-rank; base=h2

\subsection{第六章——当雅各在埃及去世时，谁是阿耳戈斯和亚述的君王。}

当时在埃及去世的\Person{Apis}并非埃及王，而是\Place{阿尔戈斯}王。其子\Person{Argus}继位，该地因他得名\Place{阿尔戈斯}，民众称为\OriginalTerm{阿尔戈斯人}{Argives}；因为在此前的诸王统治下，该地及其民族尚未有此名称。当他在\Place{阿尔戈斯}为王，\Person{Eratus}在\Place{西锡安}为王，\Person{Balæus}仍为\Place{亚述}王时，\Person{雅各伯}在\Place{埃及}去世，享年一百四十七岁。他临终时祝福了儿子们和\Person{若瑟}的孙辈，最清晰地预言了\Person{基督}，在祝福\Person{犹大}时说：\Quote{权杖不离开犹大，权柄不离他腿间，直到那应许给他的事物来临；他是万民的期待。}\BibleRef{创 49:10}在\Person{Argus}统治期间，\Place{希腊}开始使用果实，田地中有了谷物的收成，种子从别国引进。\Person{Argus}死后也被视为神，受献祭并被立庙。这一荣耀在他统治期间先赐予了一个私人，只因他是第一个用牛耕地的人。此人名\Person{Homogyrus}，后被雷电击毙。\par

% structure: p0017 source=h2 target=subsection reason=relative-heading-rank; base=h2

\subsection{第七章——当\Person{若瑟}在\Place{埃及}去世时，有哪些君王在位。}

在\Place{亚述}第十二王\Person{Mamitus}和\Place{西锡安}第十一王\Person{Plemnæus}统治期间，\Person{Argus}仍为\Place{阿尔戈斯}王时，\Person{若瑟}在\Place{埃及}去世，享年一百一十岁。此后，天主子民奇妙地繁衍，在\Place{埃及}又停留了一百四十五年，起初平安无事，直到认识\Person{若瑟}的人都死了。后来，因他们人数增多而遭嫉妒，并因可能终获自由的猜疑，他们受迫害和无法忍受的奴役之苦；然而在此期间，他们凭天主赐予的丰饶仍然不断增长。此时期内，\Place{亚述}和\Place{希腊}的王国继续存在。\par

% structure: p0019 source=h2 target=subsection reason=relative-heading-rank; base=h2

\subsection{第八章——当\Person{梅瑟}出生时，有哪些君王在位，当时开始崇拜哪些神祇。}

当\Place{亚述}第十四王\Person{Saphrus}、\Place{西锡安}第十二王\Person{Orthopolis}和\Place{阿尔戈斯}第五王\Person{Criasus}统治时，\Person{梅瑟}在\Place{埃及}出生。藉着他，天主子民得以摆脱\Place{埃及}的奴役——他们正需如此受试炼，好渴望造物主的援助。有人认为\Person{普罗米修斯}生活于上述诸王统治时期。据传他用泥土造人，因他被尊为智慧的最高导师；但当时有哪些智者并不清楚。其弟\Person{阿特拉斯}据说是一位伟大的占星家，这便引发了传说说他擎天；不过，关于他擎天的凡俗说法更像是源于一座以他命名的高山。的确，从那时起，\Place{希腊}便开始编造诸多其他神话；然而，直到\Person{刻克洛普斯}成为\Place{雅典}王——该城在其统治期间得名，并且天主藉\Person{梅瑟}将祂的子民带出\Place{埃及}，据希腊人的虚妄迷信，只有少数死者被奉为神。其中包括\Person{Melantomice}——\Person{Criasus}王的妻子，其子\Person{Phorbas}作为\Place{阿尔戈斯}第六王继位，\Person{Triopas}之子\Person{Iasus}为第七王，以及第九王\Person{Sthenelas}（或\Person{Stheneleus}，或\Person{Sthenelus}，不同作者对其称呼各异）。在那时期，据书籍中普遍记载，\Person{墨丘利}也生活过，他是\Person{阿特拉斯}的外孙，母亲是\Person{Maya}。他精通多种技艺，并将之教给人类，因此人们决定使他死后成神，甚至相信他理应如此。\Person{赫拉克勒斯}出现得较晚，但仍属于同一时期；尽管有些人（我认为他们错了）将他置于\Person{墨丘利}之前。但无论他们生于何时，庄重的史家（他们将此事写入史书）都一致认为，二人皆是凡人，因给凡人带来诸多恩惠使此世生活更为愉悦而获得了神圣尊荣。\Person{弥涅耳瓦}远早于他们；据传她以处女之身出现在\Person{Ogyges}时期的\Place{特里同湖}畔，故又称\OriginalTerm{特里托尼亚}{Tritonia}。她确实是许多技艺的发明者，因其起源鲜为人知，更易被信为女神。至于关于她从头戴\Person{朱庇特}头颅中跃出的歌谣，则属于诗歌与神话领域，而非历史与真实事实。史家对\Person{Ogyges}活跃的时间莫衷一是，其时期曾发生一场大洪水——并非那场只有进入方舟者才能逃脱的至大洪水（因为希腊和拉丁史家均不知晓），但此洪水大于后来\Person{丢卡利翁}时代的洪水。\Person{瓦罗}在其前述著作中即以此时为起点，从\Person{Ogyges}洪水（即\Person{Ogyges}时期发生的洪水）出发追溯\Place{罗马}史事，不再以更古老之事为源头。而我们的编年史家——首先是\Person{优西比乌斯}，其后是\Person{热罗尼莫}——他们完全遵循较早史家的意见，认为\Person{Ogyges}洪水发生在\Place{阿尔戈斯}第二王\Person{Phoroneus}统治之后三百余年。但无论他生活在何时，\Person{弥涅耳瓦}在\Person{刻克洛普斯}统治\Place{雅典}时已被作为女神崇拜，而\Place{雅典}城正是在其统治期间据称得以重建或建立。\par

% structure: p0021 source=h2 target=subsection reason=relative-heading-rank; base=h2

\subsection{第九章——\Place{雅典}城何时建立，\Person{瓦罗}为其名称提出了何种理由。}

雅典无疑得名于\Person{弥涅耳瓦}，后者在希腊文中被称为\OriginalTerm{雅典娜}{\GreekText{Ἀθηνη}}，\Person{瓦罗}指出了这样命名的原由如下：当一棵橄榄树突然出现在那里，而另一处有泉水涌出时，这些异象促使国王派人前往\Place{德尔斐}的\Person{阿波罗}，询问这些迹象的含义以及他应当如何应对。\Person{阿波罗}回答：橄榄树象征\Person{弥涅耳瓦}，泉水象征\Person{尼普顿}，市民有权根据这两位神祇——他们的迹象就是这些——来选择城市的名称。\Person{刻克洛普斯}接到神谕后，召集所有男女公民前来投票，因为当时那一带习俗允许妇女也参加公共议事。众人商议时，男子投票支持\Person{尼普顿}，女子支持\Person{弥涅耳瓦}；由于女子多了一票，\Person{弥涅耳瓦}获胜。\Person{尼普顿}因此大怒，掀起海浪淹没了雅典人的土地；因为魔鬼要更广泛地散布水流并不困难。同一权威说，为平息他的怒气，雅典人应对妇女施以三重惩罚：她们不得再拥有投票权；她们的孩子不得随母姓；任何人不得称她们为雅典人。如此，那座城市——它是自由学说的母亲和保姆，孕育了如此众多、如此伟大的哲学家，希腊无有比它们更著名、更高贵者——由于魔鬼嘲弄男神与女神的争斗，以及通过妇女取得的女性胜利，便得名雅典；然而，在被战败的神祇损害后，它被迫惩罚胜利者的胜利本身，因为它害怕\Person{尼普顿}的水胜过\Person{弥涅耳瓦}的武器。因为在那些受罚的妇女身上，获胜的\Person{弥涅耳瓦}也失败了，甚至无法帮助她的投票者到这个地步：虽然投票权从此失去，母亲们不能以自己的名字给孩子命名，她们至少可以被称为雅典人，并配得上那位她们投票使之战胜男神的女神之名。关于此事，若非我们急于讨论其他问题，可以说出多少话来，是显而易见的。\par

% structure: p0023 source=h2 target=subsection reason=relative-heading-rank; base=h2

\subsection{第十章——瓦罗关于亚略巴古一词及杜卡利翁洪水的叙述。}

\Person{马尔库斯·瓦罗}不愿相信那些诽谤诸神的虚假故事，以免发现任何有损其威严的事；因此他不承认\Place{亚略巴古}——宗徒\Person{保禄}曾在那里与雅典人辩论——之所以得名，是因为\Person{马尔斯}（在希腊文中被称为\OriginalTerm{阿瑞斯}{\GreekText{ἌΑρης}}）曾被控杀人罪，并在那田野上受十二位神祇审判，因六票对六票而被判无罪；因为当时习俗，票数相等时，宁释勿判。对于这个广为流传的意见，他试图从一些隐晦的书籍的记载中为这个名称找出另一个理由，以免雅典人被以为是从\Person{马尔斯}和\Quote{田野}这两个词称它为亚略巴古，仿佛它是马尔斯之野，而这在他看来有损诸神的尊严，因为他认为诉讼和审判与神祇相去甚远。他还断言，关于\Person{马尔斯}所说的这一说法，其虚假程度不亚于关于三位女神——即\Person{尤诺}、\Person{弥涅耳瓦}和\Person{维纳斯}——的说法，她们为获得金苹果而在\Person{帕里斯}面前进行的美貌竞赛，不仅被讲述，而且还在剧院中通过歌舞受到赞扬，在戏剧中取悦那些以自身这些罪行（无论真实还是虚构）为乐的诸神。\Person{瓦罗}不相信这些事，因为它们与神祇和道德的本性不符；然而，在给出雅典之名的并非神话的而是历史的缘由时，他在自己的书中记载了\Person{尼普顿}与\Person{弥涅耳瓦}关于该城应以谁的名字命名的争执，这一争执如此重大，以至于在他们通过显示异象争斗时，就连被咨询的\Person{阿波罗}也不敢在他们之间判决；但为了结束诸神的纷争，正如\Person{尤皮特}将那三位女神送到\Person{帕里斯}面前一样，他也将她们送到人们面前，结果\Person{弥涅耳瓦}通过投票获胜，却又因对支持她的妇女的惩罚而遭失败，因为她无法将\InnerQuote{雅典人}的称号授予那些支持她的妇女，尽管她可以将它强加于反对她的男子。在这些时代，当\Person{克拉纳奥斯}作为\Person{刻克洛普斯}的继承者在雅典为王时（如\Person{瓦罗}所述，但根据我们的\Person{优西比乌}和\Person{热罗尼莫}，\Person{刻克洛普斯}本人仍然在位），发生了被称为\Person{杜卡利翁}的洪水，因为它主要发生在他统治的那些地区。但这洪水根本没有波及\Place{埃及}及其附近。\par

% structure: p0025 source=h2 target=subsection reason=relative-heading-rank; base=h2

\subsection{第十一章——\Person{梅瑟}率领百姓出\Place{埃及}之时；以及当他的继任者\Person{农}的儿子\Person{若苏厄}去世时，哪些人在为王。}

\Person{梅瑟}带领百姓出\Place{埃及}，是在\Place{雅典}王\Person{刻克洛普斯}末年，当时\Place{亚述}王是\Person{阿斯卡塔德斯}，\Place{西基昂}王是\Person{马拉图斯}，\Place{阿尔戈斯}王是\Person{特里俄帕斯}；带领百姓出来后，他在\Place{西乃山}将那从天主接受的律法赐给他们，这律法称为旧约，因为它含有地上的应许，并且因为藉着\Person{耶稣基督}，将有一新约，应许天上的国度。因为在这事上应遵守的次序，与每个在天主内兴旺的人所遵守的次序相同，正如宗徒所说的：\Quote{那不是属神的在先，而是属血气的在先}，因为他确实说：\Quote{第一个人是属土的，属于土；第二个人是属天的，属于天。}\BibleRef{格前 15:46}。现在\Person{梅瑟}在旷野治理百姓四十年，一百二十岁去世，之前在会幕、司祭职、祭献及许多其他奥秘仪式中，藉着肉体礼仪的预象预言了基督。\Person{农}的儿子\Person{若苏厄}继承\Person{梅瑟}，将所带领的百姓安置在应许之地，藉着天主的权能征服了原先占据那地的民族。他也去世了，在\Person{梅瑟}死后治理百姓二十七年，当时\Place{亚述}王是第十八位王\Person{阿明塔斯}，\Place{西基昂}王是第十六位王\Person{科拉科斯}，\Place{阿尔戈斯}王是第十位王\Person{达那俄斯}，\Place{雅典}王是第四位王\Person{埃里克托尼乌斯}。\par

% structure: p0027 source=h2 target=subsection reason=relative-heading-rank; base=h2

\subsection{第十二章——论希腊诸王在以色列人出埃及至\Person{农}的儿子\Person{若苏厄}去世期间所设立的伪神礼仪。}

在此期间，即从以色列人出埃及至\Person{农}的儿子\Person{若苏厄}去世（藉着他，那百姓获得应许之地），希腊诸王设立了伪神的礼仪，这些礼仪通过定期的庆祝，唤起对洪水以及人类从洪水中获救的记忆，以及当时他们在高地与平原之间迁徙的困苦生活。因为就连\OriginalTerm{卢佩尔克斯祭司}{Luperci}，当他们沿圣路上行下行时，据说代表那些因洪水泛滥而寻找山顶、在洪水退去后返回低地的人。在那些时代，\Person{狄俄倪索斯}（也称为\Person{利伯尔}神），死后被尊为神，据说他在\Place{阿提卡}向他的主人展示了葡萄藤。于是，为了平息\Place{德尔斐}的\Person{阿波罗}的愤怒，设立了音乐竞技，他们认为因为未能保护其神庙，\Place{希腊}各地区才遭受荒芜之苦；这神庙是\Person{达那俄斯}入侵那些土地时烧毁的。因为神谕警告他们，要设立这些竞技。但\Person{埃里克托尼乌斯}王首先在\Place{阿提卡}为他设立竞技，并且不仅为他，也为\Person{弥涅耳瓦}设立；在这些竞技中，橄榄枝作为胜利者的奖品，因为他们讲述说，\Person{弥涅耳瓦}是那果实的发现者，正如\Person{利伯尔}发现了葡萄。在那些年间，据说\Person{欧罗巴}被\Place{克里特}王\Person{克珊托斯}（我们发现有些人给他另一个名字）拐走，并为他生了\Person{拉达曼提斯}、\Person{萨尔珀冬}和\Person{米诺斯}；更常见的说法是，他们是\Person{朱庇特}与同一女子所生。那些崇拜这些神祇的人，将我们关于\Place{克里特}王\Person{克珊托斯}的叙述视为真实历史；但关于\Person{朱庇特}的这些事——诗人们歌唱、剧场喝彩、民众庆祝——则被视为空洞的神话，是为了给竞技提供理由以平息神祇的愤怒，甚至将罪行错误地归于他们。在那些时代，\Person{赫丘利}在\Place{提洛}受到尊敬，但这不是我们上面所说的那位\Person{赫丘利}。在更隐秘的历史中，据说有几位被称为\Person{利伯尔}神和\Person{赫丘利}。这位\Person{赫丘利}，其伟大事迹被记为十二件（不包括杀死非洲的\Person{安泰俄斯}，因为那件事属于另一位\Person{赫丘利}），据他们的记载，他在\Place{俄塔山}上自焚，因为他无法用征服怪物的力量来忍受使他憔悴的疾病。那时，国王（更确切地说是暴君）\Person{布西里斯}，据称是\Person{尼普顿}与\Person{利比亚}（\Person{厄帕福斯}之女）所生，据说他将宾客祭献给神祇。现在不可相信\Person{尼普顿}犯了这通奸罪，以免神祇被指控；但诗人与剧场却必须将此类事归于他们，以使神祇喜悦。\Person{伏尔甘}和\Person{弥涅耳瓦}据说是\Place{雅典}王\Person{埃里克托尼乌斯}的父母，在他统治末年，\Person{农}的儿子\Person{若苏厄}去世了。但因他们坚持\Person{弥涅耳瓦}是贞女，遂说\Person{伏尔甘}在他们争斗中受到搅扰，将精液洒在地上，因此从地所生的人得了那个名字；因为在希腊语中，\OriginalTerm{争斗}{\GreekText{ἔρις}}意为\Quote{争斗}，\OriginalTerm{地}{\GreekText{χθὼν}}意为\Quote{地}，这两个词复合而成\Person{埃里克托尼乌斯}。然而必须承认，更有学识的人否认并弃绝关于其神祇的此类故事，并宣称这神话信仰起源于：在\Place{雅典}的\Person{伏尔甘}与\Person{弥涅耳瓦}共有的神庙中，发现一个被遗弃的男孩缠在龙卷中，这预示他将成为伟人；由于父母不详，他被称作\Person{伏尔甘}与\Person{弥涅耳瓦}之子，因为他们共有那神庙。但那个神话比这段历史更能解释他名字的由来。但这与我们何干？让说实话的书籍里的记载熏陶宗教人士，而让说谎的神话中的记载取悦不洁的魔鬼。然而这些宗教人士却敬拜它们为神。即便如此，当他们否认关于神祇的这些事时，仍不能免除一切罪行，因为依他们的要求，他们表演戏剧，在其中那些他们明智地否认的事却被卑劣地演出，而神祇被这些虚假和卑劣的事所平息。现在，即使戏剧颂扬了神祇的不实之罪，但以不实之罪为乐本身即是真实的罪。\par

% structure: p0029 source=h2 target=subsection reason=relative-heading-rank; base=h2

\subsection{第十三章——在民长开始治理希伯来人时期编造了哪些神话？}

嫩之子若苏厄死后，天主的子民受民长治理；在他们时期，因着他们的罪恶，人民屡次遭受磨难而谦抑，又因天主的怜悯而得昌盛而安慰。在那时期，发明了关于\Person{特里普托勒摩斯}的传说——他奉\Person{刻瑞斯}之命，乘坐有翼的蛇，飞过贫瘠之地时赐予谷种；关于那被关在\Place{拉比林特斯}中的怪物\Person{弥诺陶洛斯}的传说，凡进入那无法解开的迷宫的人便不得出路；关于半人马（其形为半马半人）的传说；关于地狱的三头狗\Person{刻耳柏洛斯}的传说；关于\Person{佛里克索斯}及其妹\Person{赫勒}乘有翼的公羊飞逃的传说；关于\Person{戈尔贡}的传说，她的头发由蛇组成，凡注视她的人皆化为石头；关于\Person{贝勒罗丰}乘有翼马\Person{珀伽索斯}的传说；关于\Person{安菲翁}以竖琴之甜美迷引石头的传说；关于工匠\Person{代达罗斯}及其子\Person{伊卡洛斯}装上羽翼飞翔的传说；关于\Person{俄狄浦斯}的传说，他解开了通常被视为无解的人面四足怪物\Person{斯芬克斯}所出的谜语，迫使她投崖自毁；关于\Person{安泰俄斯}的传说，他是大地之子，因此倒地时反而更强，被\Person{赫拉克勒斯}所杀；也许还有其他我已忘记的传说。这些传说，在包含事件真实叙述的历史中很容易找到，把我们带到特洛亚战争，\Person{马库斯·瓦罗}在他的\WorkTitle{论罗马人种族}第二卷中以此作结；这些传说如此巧妙地由人编造，以致不牵涉对诸神的丑行。但无论谁冒充说\Person{朱庇特}强夺美少年\Person{伽倪墨得斯}是\Person{坦塔罗斯}国王所犯之罪，传说却归之于\Person{朱庇特}；或者说\Person{朱庇特}化作金雨使\Person{达那厄}怀孕，意指那女子的贞操被黄金玷污：这些事究竟是真实发生还是只是传说，或是他人所做而被妄加于\Person{朱庇特}，都无法说明人心中被视若当然的邪恶有多大，以致他们被认为能耐心听这样的谎言。然而他们却欣然接受它们，事实上，他们越虔诚崇拜\Person{朱庇特}，就越应严厉惩罚那些敢于说他这样事的人。但他们不仅不恼怒于那些编造这些事的人，反而害怕如果不在剧场中演出这些虚构故事，诸神会向他们发怒。在那时期，\Person{拉托娜}生了\Person{阿波罗}，不是我们上面多次提及的那位发布神谕的\Person{阿波罗}，而是那位据说与\Person{赫拉克勒斯}一同放牧\Person{阿德墨托斯}国王羊群的\Person{阿波罗}；然而他仍被信为神，以至许许多多、几乎所有人都认为他就是那同一位\Person{阿波罗}。那时父亲\Person{利柏耳}在印度作战，率领一群被称为巴克坎忒斯的妇女，她们与其说是勇敢，不如说是狂怒著称。有些人记载说，这位\Person{利柏耳}曾被战胜并捆绑；有些说他战死在波斯，甚至说出他被葬之处；然而以他的名义，不洁的魔鬼设立了神圣的，或者更确切地说亵圣的，酒神祭；其令人发指的污秽，经过多年后，元老院感到如此羞耻以至于在罗马城禁止它们。人们相信在那时期，\Person{珀耳修斯}和他的妻子\Person{安德洛墨达}死后被擢升到天上，所以他们不羞耻也不害怕用星座来标出他们的形象，并以它们的名字称呼它们。\par

% structure: p0031 source=h2 target=subsection reason=relative-heading-rank; base=h2

\subsection{第十四章　论神学诗人。}

在同一时期内兴起了诗人，他们也被称为\Emph{神学家}，因为他们创作关于神的颂歌；然而所论的这些神，虽是伟大人物，却不过是人，或是真天主所造的这个世界元素，或是按照造物主的旨意和自身功绩被定为统治者和掌权者的受造物。如果在许多虚妄和虚假的事中，他们歌颂了独一真天主的任何事，然而，由于同时崇拜祂和其他非神者，并向他们献上唯独应归于祂的敬拜，他们就根本没有正确事奉祂；甚至连\Person{奥耳甫斯}、\Person{穆赛乌斯}、\Person{利诺斯}这样的诗人，也不能免于用传说羞辱他们的神明。但这些神学家崇拜诸神，而不被作为神崇拜，虽然不虔敬者的城邦，我不知道为何，惯于将\Person{奥耳甫斯}置于神圣的，或更确切地说，亵圣的地狱仪式之上。\Person{阿塔玛斯}国王的妻子，名叫\Person{伊诺}，和她的儿子\Person{墨利刻耳忒斯}，投海而死，按照大众信仰，被列入诸神之中，如同同时期其他人一样，其中有\Person{卡斯托耳}和\Person{波鲁克斯}。希腊人确实称\Person{墨利刻耳忒斯}的母亲为\Person{琉科忒亚}，拉丁人称她为\Person{玛图塔}；但两者都认为她是一位女神。\par

% structure: p0033 source=h2 target=subsection reason=relative-heading-rank; base=h2

\subsection{第十五章　论\Place{阿尔戈斯}王国的覆灭，彼时\Person{萨图尔努斯}之子\Person{皮库斯}初获其父的\Place{劳伦图姆}王国。}

在那时期，\Place{阿尔戈斯}的王国终结了；转移到\Place{迈锡尼}，\Person{阿伽门农}就是从那里来的；\Place{劳伦图姆}王国兴起，其首任国王是\Person{萨图尔努斯}之子\Person{皮库斯}，那时妇女\Person{德波辣}审判希伯来人；但那是天主的圣神使用她作为祂的代理人，因为她也是一位女先知，虽然她的预言如此晦涩，以至于不经长篇讨论我们不能证明它是关于基督而说的。现在\Person{劳伦特人}已经在\Place{意大利}统治，罗马民族的起源显然继希腊人之后来自他们；但\Place{亚述}王国仍然持续，当\Person{皮库斯}开始在\Place{劳伦图姆}统治时，\Person{兰帕瑞斯}是第二十三位王。崇拜这类神明的人可以看到，他们应该如何看待\Person{皮库斯}的父亲\Person{萨图尔努斯}，那些人否认他是一个人；有些人甚至记载说，\Person{萨图尔努斯}本人曾在其子\Person{皮库斯}之前在\Place{意大利}统治；\Person{维吉尔}在他的著名著作中说：\par

\Quote{那桀骜不驯的种族，散居于高山之间，\newline{}他使之定居，赋予法律，\newline{}并称其国土为\Place{拉丁姆}，因为\newline{}隐蔽在他们的海岸之内，他安居无恙。\newline{}传统说那纯净的黄金时代\newline{}始于他作王的时候。}\par

但他们将这些视为诗意的幻想，并断言\Person{皮库斯}的父亲实为\Person{斯特尔刻斯}，并讲述道，因其为最熟练的农夫，他发现了田地可用动物粪便施肥，而粪便即以其名字被称为\OriginalTerm{粪}{stercus}。有人称他为\Person{斯特尔库提乌斯}。但无论他们因何理由选择称他为\Person{萨图尔努斯}，确然的是，他们使这位斯特尔刻斯或斯特尔库提乌斯因其农业功绩而成为神；并且他们也将他儿子\Person{皮库斯}纳入这些神祇之列，宣称他为著名的占卜者和战士。\Person{皮库斯}生了\Person{法乌诺斯}，\Place{劳伦图姆}的第二任国王；而他也被他们视为神。他们在特洛伊战争之前就将这些神圣尊荣给予死者。\par

% structure: p0037 source=h2 target=subsection reason=relative-heading-rank; base=h2

\subsection{第十六章——论\Place{特洛伊}毁灭后被置于诸神之中的\Person{狄俄墨得斯}，而其同伴据称变成了飞鸟。}

\Place{特洛伊}被推翻，其毁灭被到处传唱并广为人知，甚至孩童皆知；因其本身的伟大以及文笔卓越的作者之笔，此事被显著宣扬和传播。这事发生在\Person{法乌诺斯}之子\Person{拉丁努斯}统治时期，从那时起，王国开始被称为\Place{拉丁姆}而非\Place{劳伦图姆}。获胜的希腊人离开被毁的\Place{特洛伊}返回本国时，被各种可怕的灾祸撕裂和击垮。然而，即使从他们中间，他们也增加了神祇的数量，因为他们使\Person{狄俄墨得斯}成为神。他们断言，他的归家被神圣施加的惩罚所阻止，并且他们证明——不是通过神话和诗意的谎言，而是通过历史的佐证——他的同伴们变成了鸟。然而他们认为，尽管他成为了神，他既不能凭自己的力量使他们恢复人形，也无法从众神之王\Person{朱庇特}那里获得这一恩宠，作为对新天神的恩赐。他们还说他神庙在\Place{狄俄墨得亚}岛上，靠近\Place{阿普利亚}的\Place{加尔加诺山}，这些鸟在这座神庙周围飞翔，并以如此奇妙的顺从在庙中崇拜，以至于它们用喙装满水并洒水；如果希腊人或希腊血统的人来到那里，它们不仅安静，还飞上去迎接；但如果是外邦人，它们会飞起来攻击他们的头部，并以如此猛烈的打击伤他们甚至杀死他们。因为它们据称装备了坚硬而巨大的喙，足以进行这些战斗。\par

% structure: p0039 source=h2 target=subsection reason=relative-heading-rank; base=h2

\subsection{第十七章——\Person{瓦罗}论令人难以置信的人类变形。}

为支持这一故事，\Person{瓦罗}讲述了其他同样令人难以置信的故事，关于那位最著名的女巫\Person{喀耳刻}——她将\Person{尤利西斯}的同伴变成了野兽，以及关于\Person{阿卡迪亚人}——他们按抽签游过某个池塘，在那里变成狼，并与像他们一样的野兽生活在那片地区的沙漠中。但若他们九年不食人肉，则再游过同一池塘时恢复人形。最后，他明确点名一位\Person{德墨涅托斯}，在尝过\Person{阿卡迪亚人}按照习俗献给他们的神\Person{吕凯奥斯}的一名男孩祭品后，变成了一只狼，并在第十年恢复原形后，训练自己成为拳击手，并在奥林匹克运动会上获胜。这位历史学家还认为，\Place{阿卡迪亚}的\Person{潘}和\Person{朱庇特}被称为\Quote{吕凯奥斯}的别号，其缘由无非是人类变成狼的变形，因为人们认为除非藉神圣力量，否则不可能发生。因为狼在希腊语中称为\OriginalTerm{狼}{\GreekText{λυκὸς}}，吕凯奥斯这个名称似乎来源于此。他还说，罗马的\Person{卢佩尔奇}据说似乎是从这些奥秘的种子中产生的。\par

% structure: p0041 source=h2 target=subsection reason=relative-heading-rank; base=h2

\subsection{第十八章——关于我们应如何看待那些似乎通过魔鬼的伎俩发生在人身上的变形。}

或许我们的读者期望我们就这由魔鬼造成的大欺骗说些什么；而我们能说什么呢，除非人必须逃离巴比伦？\BibleRef{依 48:20} 因为这条先知性训诫应在此意义上属灵地理解，即藉着在永生天主内前进，凭着那藉着爱德而运作的信德之步伐，我们必须逃离这世界的城，它完全是邪恶天使和人类的团体。是的，我们看见魔鬼在这些深渊中的能力越大，我们就越发必须紧抓那位中保，藉着祂我们从最低之处上升到最高之处。因为若我们说这些事不可信，即使现在也不乏有人会断言，他们要么亲耳听到最可靠的权威所讲，要么甚至亲身经历过类似的事。事实上，我们在意大利时，曾听说关于某地区的事：那里的旅店女店主，浸染于这些邪恶技艺，据说惯于向她们选择或能应付的旅客提供一块奶酪中的某物，使其当场变成驮畜，搬运所需之物，待工作完成后又恢复原形。然而他们的心智并未变兽性，而是保持理性和人性，正如\Person{阿普列乌斯}在其题为\WorkTitle{金驴记}的著作中所述（或虚构的）关于他自身之事：服毒后变成了一头驴，同时保留着人的心智。\par

这些事要么是假的，要么是如此离奇以致有充分理由不被相信。但我们必须坚定地相信，全能的天主能随其意愿成就一切，无论是惩罚还是施恩；而魔鬼凭其本性力量（因为它们的受造本质本身是天使性的，虽因自身过错变为邪恶）什么事也做不成，除非得到祂的许可——祂的审判虽常隐藏，却从不不义。事实上，魔鬼若真做出我们讨论的这类事，它们并非创造真实实体，而只是改变了由真天主所造之物的表象，使之看似并非其本相。因此，我不能相信，凭魔鬼的任何理由、技艺或能力，能使身体（更不用说心智）真正变成兽类的形状和面目；但一个人的幻象，即使在思想或梦境中经历无数变化，当此人的感官被抑制或压倒时，可能以某种我无法描述的不可言喻的方式，以有形体的形式呈现在他人感官前，以致人的身体本身可能躺在某处，虽然活着，但感官被比睡眠更沉重、更紧实地锁住，而那个幻象，仿佛以某种动物的形态具身化，可能出现在他人感官前，甚至可能让此人自己感觉似乎发生了变化，正如他在睡眠中可能感觉自己变成了那样，并背负着重物；而这些重物，如果是真实实体，则是由魔鬼背负，为的是让人在同时看到重物的真实实体和驮兽的模拟身体时受到欺骗。有一位名叫\Person{普雷斯坦提乌斯}的人常讲述，他曾在家中的父亲身上发生过这样的事：他父亲吃了一块奶酪中的毒药，躺在床上如同睡着，却无论如何无法被唤醒。但他说，几天后他仿佛醒来，讲述了所经历的事如同梦境一般，即他曾变成一匹驮马，与其他驮兽一起为所谓的勒提亚军团的士兵运送补给，因为该军团被派往\Place{勒提亚}。而这一切被发现正如他所说那样发生了，然而在他看来那是他自己的梦。另一个人宣称，在他自己家中，夜晚睡前，他看见一位他很熟悉的哲学家来到他面前，向他解释一些柏拉图哲学中的内容，而这些内容他先前被问及时曾拒绝解释。当他问这位哲学家为什么在自己家里做了他在家中拒绝做的事时，他说：\Quote{我没有做这事，而是梦见自己做了。}这样，一个人睡眠中所见的，就通过一个幻象显示给了另一个人醒着的时候。\par

这些事并非来自那些我们可能认为不可信的人，而是来自我们无法认为会欺骗我们的提供消息的人。因此，人们所说的并已记录下来关于\Person{阿卡迪亚人}常被\Place{阿卡迪亚}神祇（或更确切地说，魔鬼）变成狼的事，以及歌曲中关于\Person{喀耳刻}把\Person{乌利西斯}的同伴变形的事，如果它们真的发生过，依我之见，可能是以我所说的方式发生的。至于\Person{狄奥墨得斯}的鸟，由于据说它们的种族通过不断繁衍而延续，我认为它们并非由人变形而来，而是在人被移走时被狡猾地替换掉了，正如一头母鹿被替换了\Person{阿伽门农}王的女儿\Person{伊菲革涅亚}一样。因为这种把戏，如果得到天主审判的许可，对魔鬼来说并不困难；而且既然那位少女后来被发现还活着，很容易看出是一头母鹿被狡猾地替换了她。但是，因为\Person{狄奥墨得斯}的同伴突然不见了，之后无处可寻，被邪恶的复仇天使消灭了，人们便相信他们变成了那些鸟，这些鸟是从其他地方有这种鸟的地方秘密带来，并突然以欺诈方式替换了他们。至于它们用喙衔水洒在\Person{狄奥墨得斯}的神庙上，以及它们对希腊种族的人表示亲昵而迫害异邦人，这并不奇怪，是由魔鬼的内在影响所致，魔鬼的利益在于说服人们相信\Person{狄奥墨得斯}成了神，从而引诱他们崇拜许多假神，大大羞辱了真天主；并以庙宇、祭台、祭献和祭司来侍奉死人——这些人在世时也并未真正活着——而这一切，若合于正理，只应归于唯一永生的真天主。\par

% structure: p0045 source=h2 target=subsection reason=relative-heading-rank; base=h2

\subsection{第十九章——\Person{埃涅阿斯}来到\Place{意大利}时希伯来人由民长\Person{阿波东}治理。}

特洛伊被攻陷并毁灭后，\Person{埃涅阿斯}率领二十艘满载特洛伊遗物的船只来到\Place{意大利}，当时\Person{拉丁努斯}在那里为王，\Person{墨涅斯透斯}在\Place{雅典}，\Person{波吕菲多斯}在\Place{西基安}，\Person{陶塔诺斯}在\Place{亚述}，\Person{阿布顿}是希伯来人的民长。\Person{拉丁努斯}死后，\Person{埃涅阿斯}在位三年，上述各地仍由原王统治，只是\Person{珀拉斯戈斯}现在成了\Place{西基安}王，\Person{三松}是希伯来人的民长，有人认为他就是\Person{赫丘利}，因他有神奇的力量。当时拉丁人将\Person{埃涅阿斯}列为他们的一位神，因为他死时无处可寻。萨宾人也把他们第一位国王\Person{桑库斯}（\Person{桑古斯}或\Person{桑克图斯}）列在众神之中。那时\Place{雅典}王\Person{科德鲁斯}乔装打扮\Emph{incognito}，让自己被该城的\Place{伯罗奔尼撒}敌人所杀，就这样被杀了。他们说，他以此拯救了自己的国家。因为\Place{伯罗奔尼撒}人曾从神谕得到回应，只有不杀他们的国王才能战胜\Place{雅典}人。所以他穿上穷人的衣服，诱使他们争吵并将其杀死。\Person{维吉尔}因此说：\Quote{或科德鲁斯的争吵。}\Place{雅典}人以祭献的尊荣敬拜此人为神。拉丁人的第四位国王是\Person{埃涅阿斯}之子\Person{西尔维乌斯}，他不是\Person{克瑞乌萨}所生（\Person{克瑞乌萨}生了第三王\Person{阿斯卡尼乌斯}），而是\Person{拉丁努斯}之女\Person{拉维尼亚}所生，据说他是遗腹子。\Person{奥纽斯}是\Place{亚述}第二十九位王，\Person{美兰修斯}是\Place{雅典}第十六位王，司祭\Person{厄里}是希伯来人的民长；\Place{西基安}王国就此终结，据说延续了九百五十九年。\par

% structure: p0047 source=h2 target=subsection reason=relative-heading-rank; base=h2

\subsection{第二十章。——民长时期之后以色列君王世系。}

当上述各地有这些王统治时，民长时期结束，以色列王国始于\Person{撒乌耳}王，当时\Person{撒慕尔}先知在世。在那一时期，开始有被称为西尔维的拉丁王，他们除了本名外，还从他们的前任（那位被称为\Person{西尔维乌斯}的\Person{埃涅阿斯}之子）那里继承了这一姓氏；正如很久以后，\Person{奥古斯都·凯撒}的继承者被称为凯撒。\Person{撒乌耳}被弃绝，他的子孙不得为王，在他死后，\Person{达味}接替他作王，\Person{达味}在位四十年。然后，\Person{科德鲁斯}死后，\Place{雅典}人不再有王，开始由执政官统治共和国。\Person{达味}也作王四十年之后，他的儿子\Person{撒罗满}成为\Place{以色列}王，他在\Place{耶路撒冷}建造了那座极其辉煌的天主圣殿。在他的时代，\Place{拉丁}人中间建立了\Place{阿尔巴}，此后王不再称为拉丁王，而称为阿尔巴王，尽管仍在同一\Place{拉丁}地区。\Person{撒罗满}由他的儿子\Person{勒哈贝罕}继承，该民族在他统治下分裂为两个王国，各自部分开始有各自的王。\par

% structure: p0049 source=h2 target=subsection reason=relative-heading-rank; base=h2

\subsection{第二十一章。——关于\Place{拉丁地区}的王，其中第一位和第十二位，即\Person{埃涅阿斯}和\Person{阿文提努斯}，被奉为神。}

继被奉为神的\Person{埃涅阿斯}之后，\Place{拉丁地区}有十一位王，无人被奉为神。但\Person{阿文提努斯}，即\Person{埃涅阿斯}之后的第十二位，战死并被埋葬在那座仍以他命名的山丘，被列入他们为自己所造的神灵之列。有些人确实不愿写他被杀于战场，而是说他无处可寻，并且说那座山不因其名，而因鸟的降临而得名为\Person{阿文提努斯}。此后，除\Place{罗马}的奠基人\Person{罗慕路斯}外，\Place{拉丁地区}再无神祇被造。但在二者之间有两位王，我打算用\Person{维吉尔}诗句描述第一位：\par

\Quote{接着来了\Person{普洛卡斯}，特洛伊种族的荣耀。}\par

在那个帝国中，最伟大的\Place{亚述}王国在他的时代走向终结，而\Place{罗马}的诞生临近了。因为\Place{亚述}帝国在将近一千三百零五年后转归\Place{米底亚}人，如果我们计入\Person{贝鲁斯}的统治，他生了\Person{尼努斯}，并满足于一个小王国，是那里的第一位王。\Person{普洛卡斯}在位早于\Person{阿穆略}。\Person{阿穆略}曾使他兄弟\Person{努弥托耳}的女儿——名叫\Person{瑞亚}，也称为\Person{伊利亚}——成为灶神贞女，她因\Person{玛尔斯}（如他们所述）怀了双胞胎儿子，他们以此方式美化或掩饰她的通奸，并加上证据说一只母狼哺育了被遗弃的婴儿。因为他们认为这种野兽属于\Person{玛尔斯}，所以认为母狼给婴儿哺乳是因为她知道他们是她主人\Person{玛尔斯}的儿子；尽管也有不少人说，当哭泣的婴儿被遗弃时，起初被我不知道什么样的妓女捡起并先吸了她的奶（妓女被称为\OriginalTerm{母狼}{lupae}，她们肮脏的居所至今仍被称为\OriginalTerm{妓院}{lupanaria}），然后才落到牧羊人\Person{福斯图卢斯}手中，由他的妻子\Person{阿卡}抚养。然而，如果天主为了责罚那个残忍下令将他们扔进水里的王，并愿意从水中神奇地救出他们之后，由一头野兽哺乳来扶助这些婴儿，而他们将要建立一个如此伟大的城市，这有什么奇怪的呢？\Person{阿穆略}在\Place{拉丁}王位上由他的兄弟\Person{努弥托耳}继承，他是\Person{罗慕路斯}的外祖父；\Place{罗马}建于\Person{努弥托耳}在位的元年，从此他与外孙\Person{罗慕路斯}共同统治。\par

% structure: p0053 source=h2 target=subsection reason=relative-heading-rank; base=h2

\subsection{第二十二章。——\Place{罗马}建于\Place{亚述}王国灭亡之时，同时\Person{希则克雅}在\Place{犹大}作王。}

简而言之，罗马城被建立，如同另一个巴比伦，且可谓前巴比伦之女；天主乐意藉此城征服全世界，并将其广纳入同一政府与法律的团体，使之降服。因为当时已有强盛勇敢的民族和国家，习于征战，不易屈服，征服他们必然伴随着巨大的危险与破坏，以及极大而可怖的劳苦。当亚述王国几乎征服了整个亚洲时，虽然是通过战斗完成的，但战争还不至于非常激烈或艰难，因为那些民族尚未养成抵抗的习惯，而且数量和规模也不如后来那样大；因为在那次最大而真正普世的洪水之后，只有八人从诺厄方舟中幸存，不到一千余年，尼努斯就征服了除印度以外的整个亚洲。但罗马并没有以同样的速度和轻易程度完全征服那些我们所见被纳入罗马帝国的东西方的所有民族，因为在它的逐步扩张中，无论向哪个方向延伸，都遇到强盛而好战的民族。当罗马建立时，以色列子民已在应许之地居住了七百一十八年。其中二十七年属于农的儿子若苏厄，之后的三百二十九年属于民长时期。但从以色列开始有君王统治时起，已过了三百六十二年。那时在犹大有一位君王名叫阿哈兹，或按其他人的计算，是他的继任者希则克雅，一位最优秀、最虔诚的君王，公认在罗慕路斯时代在位。在希伯来民族中称为以色列的那一部分，曷舍亚已开始为王。\par

% structure: p0055 source=h2 target=subsection reason=relative-heading-rank; base=h2

\subsection{第23章——论厄立特里亚女先知，她比别人更清楚地咏唱了许多关于基督的事。}

有人说厄立特里亚女先知在这时发预言。瓦罗声称有许多女先知，并非仅有一位。这位厄立特里亚女先知确实写下了一些关于基督的相当明显的话，我们最初是以拉丁文读到的，诗句拉丁文拙劣且不合韵律，这是由于我们后来得知的、某位我不认识的无能译者的缘故。因为弗拉奇亚努斯，一位非常著名的人物，也曾任总督，口才极佳且学识渊博，当我们谈论基督时，他拿出一份希腊文手稿，说是厄立特里亚女先知的预言，在其中他指出某一段落，各行首字母如此排列，可以读出这些词：\OriginalTerm{耶稣基督、天主子、救主}{\GreekText{Ἰησοῦς} \GreekText{Χριστὸς} \GreekText{Θεοῦ} \GreekText{υἱὸς} \GreekText{σωτήρ}}。而这几行诗，其首字母拼出那个意义，内容如下，由某人用拉丁文以好韵律译出：\par

\OriginalTerm{}{\GreekText{Ι}} 审判将用其旗帜的汗水湿润大地，\newline{} \OriginalTerm{}{\GreekText{Η}} 永存的，看哪，君王将穿越时代而来，\newline{} \OriginalTerm{}{\GreekText{Σ}} 被派遣在此降生成人，并在世界末日审判。\newline{} \OriginalTerm{}{\GreekText{Ο}} 天主，信者与不信者将一同看见你，\newline{} \OriginalTerm{}{\GreekText{Υ}} 与圣徒一同被高举，当时代终结之时。\newline{} \OriginalTerm{}{\GreekText{Σ}} 坐在他面前的，是带着肉身受审的灵魂。\par

\OriginalTerm{}{\GreekText{Χ}} 隐藏在浓密的云雾中，大地荒凉。\newline{} \OriginalTerm{}{\GreekText{Ρ}} 偶像和长久隐藏的宝物被人抛弃；\newline{} \OriginalTerm{}{\GreekText{Ε}} 大地被火吞噬，它搜寻海洋和天空；\newline{} \OriginalTerm{}{\GreekText{Ι}} 它涌出，摧毁地狱的可怕门户。\newline{} \OriginalTerm{}{\GreekText{Σ}} 圣徒在身体和灵魂中将继承自由与光明；\newline{} \OriginalTerm{}{\GreekText{Τ}} 有罪者将在火与硫磺中永远焚烧。\newline{} \OriginalTerm{}{\GreekText{Ο}} 隐秘的行为被揭示，每人将公布自己的秘密；\newline{} \OriginalTerm{}{\GreekText{Σ}} 每人内心的秘密，天主将在光明中揭示。\par

\OriginalTerm{}{\GreekText{Θ}} 那时将有哀哭和嚎啕，是的，还有切齿；\newline{} \OriginalTerm{}{\GreekText{Ε}} 太阳昏暗，星辰沉默于它们的合唱。\newline{} \OriginalTerm{}{\GreekText{Ο}} 月光的光辉逝去，诸天融化，\newline{} \OriginalTerm{}{\GreekText{Υ}} 山谷被他高举，山岳被他推倒。\par

\OriginalTerm{}{\GreekText{Υ}} 在人中，高贵与低贱的区分完全消失。\newline{} \OriginalTerm{}{\GreekText{Ι}} 山丘涌入平原，天空与海洋混合。\newline{} \OriginalTerm{}{\GreekText{Ο}} 噢，万物的终结！大地破碎，将要消亡；\newline{} \OriginalTerm{}{\GreekText{Σ}} ……水和火焰同时涌流成河。\par

\OriginalTerm{}{\GreekText{Σ}} 总领天使的号角将从天上轰鸣，\newline{} \OriginalTerm{}{\GreekText{Ω}} 降临于那在罪孽与众多忧愁中呻吟的恶人。\newline{} \OriginalTerm{}{\GreekText{Τ}} 大地战栗，将会裂开，显露混沌与地狱。\newline{} \OriginalTerm{}{\GreekText{Η}} 每个君王在那日都要站在天主面前受审。\newline{} \OriginalTerm{}{\GreekText{Ρ}} 火与硫磺的河流将从天上降下。\par

在这些拉丁文诗句中，希腊文的意思被正确传达了，尽管并非完全按照与首字母相连接的行序；因为在其中三行，即第五、第十八和第十九行，出现希腊字母\OriginalTerm{}{\GreekText{Υ}}，在拉丁文中找不到以相应字母开头且意义合适的词。因此，如果我们将拉丁译文全部行的首字母放在一起（除了那三行我们保留\OriginalTerm{}{\GreekText{Υ}}在适当位置之外），它们将以五个希腊词表达这个意思：\Quote{耶稣基督、天主子、救主}。而诗句共二十七行，这是三的立方。因为三乘三得九；九本身若再乘三，从平面正方形上升到立方体，便得二十七。但如果你将这五个希腊词的首字母连起来：\OriginalTerm{耶稣基督、天主子、救主}{\GreekText{Ἰησοῦς} \GreekText{Χριστὸς} \GreekText{Θεοῦ} \GreekText{υἱὸς} \GreekText{σωτήρ}}，它们将组成\OriginalTerm{鱼}{\GreekText{ἰχθύς}}这个词，基督在其中被神秘地理解，因为他能够在这必朽的深渊如同在水深之处生活，即无罪恶地存在。\par

但这女预言家，无论她是厄立特里亚的，或如某些人所更相信的库迈的，在她的整首诗（这只是其中极小部分）中，不仅没有任何与假神或虚构之神的崇拜相关的内容，反而如此地反对他们和他们的崇拜者，以至于我们甚至可以认为她应被算作属于天主之城的人。拉克坦提乌斯也在他的著作中引用了某位女预言家关于基督的预言，但他没有说明是哪一位。但我觉得适宜将他以许多简短引文所记载的内容合并成一个可能较长的摘录。她说：\Quote{此后祂将落入不信者的恶手中，他们将以亵渎的手掌击打天主，用不洁的口吐出毒涎；但祂将谦逊地献出祂圣洁的背脊受鞭打。当被拳击时，祂将缄默，以免有人发现祂要说什么、从何处来、到地狱去说话；祂将戴上一顶荆棘冠冕。他们给祂苦胆作食物，给祂醋解渴；他们摆上这无情的筵席。因为你自己，愚昧的人，没有认出你的天主，迷惑了凡人的心意，反而用荆棘冠冕祂，为祂调了苦胆。但圣殿的幔子将裂开；正午时，黑暗将胜过夜晚三小时。祂将死去，安眠三日；然后从地狱返回，首先来到光明中，向被召者显示复活的开始。}拉克坦提乌斯使用了这些女预言家的见证，在讨论过程中，根据他所要证明的事的需要，逐点引入，而我们则将其连续地排列在一起，不夹杂评论，只注意用大写字母标明，但愿后来的抄写者不忽略保留它们。有些作者确实说厄立特里亚的女预言家不是在罗慕洛时代，而是在特洛伊战争时代。\par

% structure: p0064 source=h2 target=subsection reason=relative-heading-rank; base=h2

\subsection{第二十四章——七贤者在罗慕洛统治时期兴盛，当时称为以色列的十支派被加色丁人掳去，罗慕洛死后被赋予神圣尊荣。}

当罗慕洛统治时，据说米利都的泰勒斯在世，他是七贤者之一，他们继神学诗人（以奥菲士最为著名）而起，被称为\OriginalTerm{贤者}{\GreekText{Σοφοί}}，即贤者。在那期间，因人民分裂而称为以色列的十支派被加色丁人征服，掳到他们的土地上，而称为犹大的两支派，其王国所在为耶路撒冷，仍留在犹太地。当罗慕洛死后无处可寻时，罗马人（这是众所周知的）将他置于诸神之中——这事在那时已经不再做了，并且此后直到凯撒时代才再做，而那时并非由于错误，而是出于谄媚；因此西塞罗大为赞颂罗慕洛，因为他不是在那粗野未开化、人们容易受骗的时代，而是在已经文雅开化的时代，虽然哲学家伶俐机智的雄辩尚未达到巅峰，却赢得了这样的尊荣。但虽然后来时代不再神化死人，他们仍然继续尊崇和敬拜那些早先被神化的神祇；不仅如此，通过古人从未有过的雕像，他们甚至增加了虚妄不敬的迷信的诱惑，不洁的恶魔在他们心中促成此事，并以虚假的神谕欺骗他们，以致就连诸神的神话罪行（这在更文明的年代甚至未曾想象过）也被卑劣地在戏剧中演出，以尊崇这些相同的假神。努马继罗慕洛之后为王；虽然他以为罗马的神越多就越能得到保护，但他自己死后却不被认为配得诸神中的一席之地，仿佛他们认为他已使天堂如此拥挤，以致找不到地方给他。据说萨摩斯的女预言家在他统治罗马时在世，而默纳舍开始统治希伯来人——一个不虔敬的国王，据说先知依撒意亚被他所杀。\par

% structure: p0066 source=h2 target=subsection reason=relative-heading-rank; base=h2

\subsection{第二十五章——当塔尔奎尼乌斯·普利斯库斯统治罗马人、漆德克雅统治希伯来人、耶路撒冷被攻取、圣殿被推翻时，哪些哲学家闻名于世？}

当漆德克雅统治希伯来人，而安库斯·马尔提乌斯的继承者塔尔奎尼乌斯·普利斯库斯统治罗马人时，犹太人民被掳到巴比伦，耶路撒冷和所罗门所建的圣殿被推翻。因为先知们责备他们的邪恶和不虔敬，预言了这些事要发生，尤其是耶肋米亚，他甚至说出了年数。据说米蒂利尼的彼塔库斯，另一位贤者，生活在那个时代。尤西比乌斯写道，当天主的子民被囚在巴比伦时，其他五位贤者生活着，他们必须与我们上面提到的泰勒斯和彼塔库斯加起来，才凑足七位。他们是：雅典的索伦、拉刻代蒙的喀隆、科林斯的佩里安德、林都斯的克莱奥布洛斯、普里耶涅的彼亚斯。他们兴盛于神学诗人之后，被称为贤者，因为他们在某种可称颂的生活方式上超越其他人，并以警句格言总结了一些道德教训。但他们没有给后人留下文学著作，除了据说索伦给了雅典人某些法律，而泰勒斯是一位自然哲学家，并以简短的箴言留下了他的学说著作。在那犹太人被掳的时代，自然哲学家阿那克西曼德、阿那克西美尼和色诺芬尼兴盛起来。毕达哥拉斯也生活在那时，并且在这个时代，哲学家这个名称首次被使用。\par

% structure: p0068 source=h2 target=subsection reason=relative-heading-rank; base=h2

\subsection{第二十六章——当犹太人的被掳在七十年满时结束，罗马人同时也摆脱了王权统治。}

\Person{居鲁士}是波斯王，当时也统治\Place{加色丁}和\Place{亚述}。他稍为放宽了犹太人的掳期，使他们中的五万人返回，以重建圣殿。他们只立了最初的根基，并建了祭台；但因敌军入侵，无法继续，工程延迟到\Person{达理阿}王时代。在同一时期，也发生了\WorkTitle{友弟德传}中所记载的事，据说犹太人并未将其收入圣经正典。后来，在波斯王\Person{达理阿}治下，\Person{耶肋米亚}先知所预言的七十年期满，犹太人的掳期结束，他们重获自由。那时，\Person{塔克文}作为罗马人的第七位王在位。他被放逐后，罗马人也开始脱离王权的统治。直到此时，以色列人仍有先知；但虽然先知众多，只有少数几位先知的著作被犹太人并也被我们保存为正典。在前一卷书结尾时，我曾答应在本卷中记述一些关于他们的事，我现在就来做。\par

% structure: p0070 source=h2 target=subsection reason=relative-heading-rank; base=h2

\subsection{第二十七章 论其预言载于书中、在罗马王国开始及亚述帝国终结时歌唱许多关于外邦人蒙召之事的先知时代}

为使我们能够考虑这些时代，让我们稍为回溯更早的时代。在欧瑟亚先知书（他被列为十二先知之首）的开端写着：\BibleQuote{上主的话在乌齐雅、约堂、阿哈次和希则克雅为犹大王时，传给欧瑟亚。}\BibleRef{欧 1:1} 阿摩司也记载他在乌齐雅的日子说预言，并提及同时代的以色列王雅洛贝罕的名字。\BibleRef{亚 1:1} 依撒意亚——阿摩司的儿子（可能是上述那位先知，或更明确地说，是另一位并非先知但同名的人）——也在他的书卷开头列出了欧瑟亚所提的这四位君王，在前言中说他在他们的日子里说预言。米该亚也把乌齐雅之后的日子列为他说预言的时期，因为他列出了欧瑟亚所提的三位君王——约堂、阿哈次和希则克雅。\BibleRef{米 1:1} 我们从他们自己的著作中发现这些人同时说预言。此外，在乌齐雅治下还有约纳，在继乌齐雅位的约堂治下还有岳厄尔。但这两位的年代我们是根据编年史，而非他们的著作，因为他们自己并未提及。现在这些时代涵盖从拉丁人的普罗卡王或其前任阿文提努斯，直到罗马人的罗慕路王，甚至到他的继任者努马·庞皮留斯统治之初。犹大王希则克雅肯定统治到那时为止。因此，正如我所称呼的，这些预言之泉在亚述帝国衰亡、罗马帝国兴起之时同时涌流；因此，正如在亚述帝国初期，亚巴郎出现，他得到最明确的许诺：万民要因他的后裔蒙受祝福，同样，在西巴比伦的初期，在其政权的时代基督将要降临，这些许诺将在祂内实现，先知的预言不仅以口传，也以文字形式赐下，作为如此伟大的事将要发生的见证。因为虽然以色列人从开始立王以来几乎从未缺少先知，但这些先知只是为他们自己，而非为万民。但当那更明显是为万民而立的预言性圣经开始形成时，它当然应在那座将要统治万民的城市建立时开始。\par

% structure: p0072 source=h2 target=subsection reason=relative-heading-rank; base=h2

\subsection{第二十八章 论\Person{欧瑟亚}与\Person{阿摩司}所预言关于基督福音之事}

\Person{欧瑟亚}先知发言非常深奥，令人费解才能领悟其含义。但按着约定，我们必须从他的书中摘录一些内容。他说：\BibleQuote{将来在给他们说\InnerQuote{你们不是我的子民}的地方，他们将被称为永生天主的子女。} \BibleRef{欧 1:10} 连宗徒们也将此理解为外邦人蒙召的预言性见证，他们从前不属于天主；又因为外邦人本身在灵性上也属于亚巴郎的子孙，因此有理由称为以色列，所以他接着说：\BibleQuote{犹大的子女和以色列的子女要聚集在一起，为自己立一个首领，并从地上兴起。} \BibleRef{欧 1:11} 我们若试图阐释这预言，只会减弱其韵味。但愿读者记住那角石和两堵隔墙，一是犹太人，一是外邦人，\BibleRef{迦 2:14} 就能认出他们：一方称为犹大的子孙，另一方称为以色列的子孙，都靠同一个首领支持，从地上兴起。但那些肉体上的以色列人现不愿信基督，后来却会相信，即他们的子女（因为他们自己当然会死去归入自己的地方），同一位先知作证说：\BibleQuote{以色列的子女将多日没有君王，没有首领，没有祭献，没有祭台，没有司祭，没有神示。} \BibleRef{欧 3:4} 谁不见犹太人现今正是如此？但我们且听他接着说：\BibleQuote{后来以色列的子女必回转，寻求主、他们的天主，和达味他们的君王；并在末日对主和他的恩惠感到惊奇。} \BibleRef{欧 3:5} 这预言再清楚不过了，其中以达味之名称，意指基督为王，他\Quote{成了}，如宗徒所说，\BibleQuote{按肉身是达味的后裔。} \BibleRef{罗 1:3} 这位先知也预言了基督第三日的复活，按预言的高峰理当如此预言，他说：\BibleQuote{两天后他必使我们复生，第三日我们必兴起。} \BibleRef{欧 6:2} 与此一致，宗徒对我们说：\BibleQuote{你们既然与基督一同复活了，就该追求天上的事。} \BibleRef{哥 3:1} 亚毛斯也就此事如此预言：\BibleQuote{以色列！你应为迎接你的天主做准备；因为他正在形成雷声，创造风，并向人宣示他们的基督。} \BibleRef{亚 4:12} 在另一处他说：\BibleQuote{在那天，我必重建达味已倒塌的帐幕，修补其中的破口，重建他废墟，再建如古时一样；好使剩余的人寻求我，并所有属于我名下的民族，这是行此事的主说的。}\par

% structure: p0074 source=h2 target=subsection reason=relative-heading-rank; base=h2

\subsection{第29章——\Person{依撒意亚}关于基督和教会预言了什么}

依撒意亚的预言不在十二先知书中，十二先知因著作简短而被称为小先知，相比之下，那些出版较大卷册的被称为大先知。依撒意亚属于后者，我却将他与上述两位连在一起，因为他同时代预言。于是依撒意亚，连同他对邪恶的斥责、对正义的训诫、对灾祸的预言，还预言了比其他人更多关于基督和教会的事，即关于那位君王和他所建立的城；所以有人说他应被称为福音传道者而非先知。但为完成这部作品，我只在此处从众多中引述一段。以圣父的口吻说：\BibleQuote{看，我的仆人必要明白，且被高举，极受尊荣。许多人将因你而惊讶。} 这是关于基督的。\par

但现在让我们听听关于教会接着说了什么。他说：\BibleQuote{不生育的，不生产的，你要欢乐！未受产痛的，你要欢呼，并高声喊叫！因为被弃者的子女比有夫之妇的子女更多。} \BibleRef{依 54:1} 但这些话也就够了；其中有些需要解释；但我认为这些部分已经足够明确，连敌人也必被逼违背己意去理解。\par

% structure: p0077 source=h2 target=subsection reason=relative-heading-rank; base=h2

\subsection{第30章——\Person{米该亚}、\Person{约纳}和\Person{岳厄尔}按新约预言了什么}

\Person{米该亚}先知用一座大山的形象代表基督，说：\BibleQuote{到末日，主的山岭必要矗立在群山之上，超乎一切山丘；万民都要涌向它。许多民族要前往，说：来，我们攀登上主的山，到雅各伯天主的殿里；他必指示我们他的道路，我们循他的路径行走；因为法律将出自熙雍，主的话将出自耶路撒冷。他将在许多民族中施行审判，斥责远方的强国。} \BibleRef{米 4:1} 这位先知预言了基督诞生的确切地点，说：\BibleQuote{厄弗辣大的白冷啊！你在犹大的千村中虽是最小的，但将由你为我生出一位统治以色列的领袖；他的根源来自亘古，来自永远之日。因此他将舍弃他们，直到那分娩的妇女生产；他其余兄弟必归向以色列子民。他立足牧放，以主的力量，以主他天主之名的尊严，牧放他的羊群；那时他必受尊崇，直到地极。} \BibleRef{米 5:2}\par

先知\Person{约纳}，与其说是藉言语，不如说是藉他自己痛苦的经历，更清楚地预言了\Person{基督}的死亡与复活，比他亲口宣告这些事还要清楚。因为何以他被吞入鲸腹，第三日又被复原，岂不正是为了作一个标记，预示\Person{基督}在第三日要从地狱的深处归来？\par

我本应用很多话解释\Person{岳厄尔}所预言的一切，以阐明那些关乎\Person{基督}和教会的事。但有一处经文我不应略过，就是当圣神按照\Person{基督}的应许从高天降临在聚集的信徒身上时，宗徒们也引用的那一段。他说：\BibleQuote{此后，我要将我的神倾注在一切有血肉的人身上；你们的儿子们和女儿们要说预言，你们的老年人要做梦，你们的青年人要见异像；在那些日子里，连我的仆人和婢女，我也要倾注我的神。}\BibleRef{岳 2:28}\par

% structure: p0081 source=h2 target=subsection reason=relative-heading-rank; base=h2

\subsection{第31章——论及在基督内世界的救恩的预言，见于\Person{亚北底亚}、\Person{纳鸿}和\Person{哈巴谷}。}

三位小先知——\Person{亚北底亚}、\Person{纳鸿}和\Person{哈巴谷}——的日期，既未被他们自己提及，也未在\Person{欧瑟比}和\Person{热罗尼莫}的\WorkTitle{编年史}中给出。因为虽然他们把\Person{亚北底亚}与\Person{米该亚}列在一起，但\Person{米该亚}何时说预言，从他们著作中注明日期的部分无法得知。我认为，这是由于他们在粗心抄写他人著作时出现了错误。但我们未能在我们所拥有的编年史抄本中找到现在提及的另外两位；然而，因他们被列入正典，我们不应忽略他们。\par

\Person{亚北底亚}，就其著作而言，是众先知中最短的一篇，他发言反对\Place{厄东}，即\Person{厄撒乌}的民族，\Person{厄撒乌}是\Person{依撒格}孪生子中的被弃的长子，\Person{亚巴郎}的孙子。现在，如果我们采取以部分代表全体的表达方式，把\Place{厄东}视为代表万民，那么我们可以理解他在其他话语中所说的关于\Person{基督}的话：\BibleQuote{但在\Place{熙雍山}上必有救援，且有一位圣者。}\BibleRef{北 1:17}随后不久，在同一预言的结尾，他说：\BibleQuote{那些再次得救的人要从\Place{熙雍山}上来，保卫\Place{厄撒乌山}，王国将归于上主。}\BibleRef{北 1:21}很明显，当那些从\Place{熙雍山}再次得救的人——即来自\Place{犹太}的基督信徒，其中主要为宗徒们——上去保卫\Place{厄撒乌山}时，这就应验了。他们如何保卫它呢？除非是藉着宣讲福音，使那些相信的人得以安全，好使他们\BibleQuote{脱离黑暗的权势，被迁到天主的国里？}\BibleRef{哥 1:13}他以推论的方式表达这一点，补充说：\Quote{王国将归于上主。}因为\Place{熙雍山}象征\Place{犹太}，那里预言必有救援和一位圣者，即\Person{耶稣基督}。而\Place{厄撒乌山}则是\Place{厄东}，象征外邦人的教会，正如我所解释的，那些从\Place{熙雍}再次得救的人保卫了它，使它成为归于上主的王国。这事在发生之前是隐晦的；但如今已然成就，哪个信徒不能发现呢？\par

至于先知\Person{纳鸿}，天主通过他说：\BibleQuote{我必除灭雕刻的和铸造的偶像，我必使你下葬。看啊，那传报佳音、宣布和平者的脚在山上何等快捷！犹大啊，庆祝你的庆节，还你的誓愿吧！因为从今以后，他们不再前行以致过时。它已成就，它已销毁，它已被夺去。那在你脸上呼吸、拯救你脱离苦难者上去了。}愿那记起福音的人回想，是谁从地狱上升，在犹大——即犹太门徒——的脸上吹了圣神；因为他们属于新约，其庆节如此属灵地更新，以致不会过时。此外，我们已经看到雕刻的和铸造的偶像——即假神的偶像——藉着福音被除灭，被遗忘如同进入坟墓，我们知道这预言就在这件事上应验了。\par

除了那将要来的\Person{基督}的来临，还能理解为\Person{哈巴谷}所说的是什么？\BibleQuote{上主回答我说：你该将神视清楚地写在黄杨木板上，使阅读的人能明白。因为这神视为指定的时期而存留，它将在终末出现，绝不会落空：如果它迟延，你要等候它；因为它必定来临，不会延迟吗？}\BibleRef{哈 2:2}\par

% structure: p0086 source=h2 target=subsection reason=relative-heading-rank; base=h2

\subsection{第32章——论及包含在\Person{哈巴谷}的祷词和诗歌中的预言。}

在他的祈祷中，伴着一首诗歌，他不是向主基督说：\BibleQuote{主啊，我听到了你的听闻，就害怕：主啊，我默想了你的作为，便极其战栗吗？} \BibleRef{哈 3:2} 这不是对那预知的、崭新的、突然而来的人类救恩的不可言喻的赞叹吗？\BibleQuote{在两个活物中间，你将被人认识。} 这不是指着两个约之间，或是两个强盗之间，或是梅瑟和厄里亚在山上与主谈话吗？\BibleQuote{当岁月临近时，你将被认识；在时代来临时，你将显现，} 甚至不需要解释。\BibleQuote{当我的灵魂因祂而烦乱时，在忿怒中，你必记念你的仁慈。} 这不是指祂为犹太人作了代祷，祂本是出自犹太民族，他们满腔愤怒，将基督钉在十字架上；而祂却记念仁慈，说道：\BibleQuote{父啊，宽恕他们吧，因为他们不知道他们所做的是什么。} \BibleRef{路 23:34} \BibleQuote{天主将从特曼而来，圣者从荫蔽而密实的山岭而来。} \BibleRef{哈 3:3} 这里所说的：\BibleQuote{祂将从特曼而来，} 有些人解释为\Quote{从南方来的，} 或\Quote{从西南方来的，} 因此象征正午，即爱德的火热和真理的光辉。\BibleQuote{荫蔽而密实的山岭} 可以有多种理解，但我更倾向于将其理解为天主圣经的深奥，其中预言了基督：因为在圣经中有许多荫蔽且密实之处，锻炼读者的心智；当有智慧的人在圣经中找到基督时，基督便从中而来。\BibleQuote{祂的能力遮蔽诸天，大地充满了祂的赞美。} 这不就是圣咏中所说的：\BibleQuote{天主，请显于诸天之上；愿你的光荣遍及大地！} \BibleQuote{祂的光辉如同光明。} 这岂不是指祂的名声将光照信徒吗？\Quote{角在祂的手中。} 这不是十字架的胜利吗？\BibleQuote{祂已安置了祂力量的坚固爱德} \BibleRef{哈 3:4} 无需解释。\BibleQuote{话语将行在祂面前，并随祂脚步进入田野。} 这不是指在祂降来之前和祂回返之后，都应宣布祂吗？\BibleQuote{祂站立，大地震动。} 这不是指\BibleQuote{祂站立} 是为了救助，\BibleQuote{大地震动} 是为了相信吗？\BibleQuote{祂注视，万民消融；} 即祂动了怜悯的心，使百姓悔改。\BibleQuote{山岭被猛烈地击碎；} 即藉着行奇迹者的能力，高傲者的骄傲被打碎。\BibleQuote{永久的山丘崩塌；} 即它们暂时被谦卑，以便能永远被高举。\BibleQuote{我看见祂的步履因祂的劳苦而成永恒；} 即我看见祂的爱之劳苦未得永生的赏报。\BibleQuote{厄提约丕亚的帐幕将极其害怕，米德杨地的帐幕亦然；} 即那些不在罗马权柄下的民族，被祢奇妙作为的消息突然惊吓，将成为基督徒百姓。\BibleQuote{主啊，你向河流发怒吗？你的怒气向河流发作吗？你的忿怒向海发作吗？} 这是说祂现在不是来审判世界，而是叫世界藉着祂得救。\BibleRef{若 3:17} \BibleQuote{因为你将骑上你的马，你的骑行将是救恩；} 即你的传福音者将承载你，因为他们受你引导，你的福音给信你的人带来救恩。\BibleQuote{主说：你将弯弓向权杖；} 即你也用你的审判威胁地上的君王。\BibleQuote{大地将被河流裂开；} 即藉着宣讲你的人的讲论涌入人心，人们的心将敞开以作宣认，对他们说：\BibleQuote{撕裂你们的心，不要撕裂你们的衣服。} \BibleRef{岳 2:13} \BibleQuote{百姓将看见你而忧伤} 除了指他们哀恸却蒙祝福，还指什么？\BibleRef{玛 5:4} \BibleQuote{行走中散布众水} 不就是指你行走在那些四处宣扬你的人中，把你的道理之流到处播散吗？\BibleQuote{深渊发出声音？} 岂不是人心的深处表达了它觉察到的？\Quote{它的幻象的深处} 这句话是对前一节的解释，因为深处就是深渊；而\Quote{发出声音} 是紧随其后的理解，即，如我们所说，它表达了它所觉察到的。这幻象就是异象，它未持有或隐藏，而是在宣认中涌流出来。\BibleQuote{太阳升起，月亮停留在它的轨道中；} 即基督升天，教会在她的君王之下建立。\BibleQuote{你的箭将在光明中发出；} 即你的言语不会暗中发出，而是公开的。因为祂曾对祂的门徒说：\BibleQuote{我在暗中告诉你们的，你们要在光明中说出来。} \BibleRef{玛 10:27} \BibleQuote{藉着威吓，你使大地缩小；} 即藉着那威吓你使人类谦卑。\BibleQuote{在忿怒中你将推翻万民；} 因为在惩罚那些自高自大的人时，你使他们彼此撞击。\BibleQuote{你出来为拯救你的百姓，以拯救你的基督；你将死亡降在恶人的头上。} 这些话都不需要解释。\BibleQuote{你举起了锁链，直到颈项。} 这也可以理解为智慧的善锁，使脚入其镣，颈入其枷。\Quote{你以心灵的惊异击碎了锁链} 必须理解为：祂举起善的而击碎恶的，关于此，有人对祂说：\Quote{你折断了我的锁链，} 而那是\Quote{在心灵的惊异中，} 即奇妙地。\BibleQuote{强者的头将在此中动摇；} 即在那惊异之中。\BibleQuote{他们张开了牙齿，如同穷人暗中进食。} 因为犹太人中有些强者将来到主前，钦佩祂的作为和言语，并因怕犹太人而暗中贪婪地吃祂道理的饼，正如福音所显示的。\BibleQuote{你将你的马放入海中，搅动了许多水，} 这些水无非就是许多人；因为如果所有人都没有被搅动，有些人就不会因恐惧而悔改，另一些人也不会因愤怒而被追逐。\BibleQuote{我留心，我的肚腹因我嘴唇祈祷的声音而战栗；战栗进入我的骨头，我的体态在我下面烦乱。} 他留心于他所说的，他自身也被自己先知般倾吐的祈祷所惊吓，他在其中辨明未来之事。因为当许多人被搅动时，他看见了教会遭受的威胁性磨难，并立刻承认自己是它的一员，说道：\BibleQuote{我将在磨难之日安息，} 正像那些在望德中喜乐、在磨难中忍耐的人。\BibleRef{罗 12:12} \Quote{使我登上，他说，在我寄居的百姓中，} 完全离开他肉亲的恶民，他们不是这地上的朝圣者，也不寻求上面的家乡。\BibleQuote{虽然无花果树，} 他说，\BibleQuote{不发旺，葡萄树不结果实；橄榄树的劳作落空，田地不出粮食；羊从肉中剪除，棚中没有牛。} 他看见那将要杀害基督的百姓即将失去精神供给的丰盈，他以先知的方式用地上的丰裕来描绘。因那百姓要遭受如此的天主忿怒，因为他们不认识天主的正义，而想建立自己的正义，\BibleRef{罗 10:3} 他立刻说：\BibleQuote{但我仍要因主喜乐；我要因我救恩的天主欢乐。主天主是我的力量，他必使我的脚步完全；他必把我置于高处，使我能在他的诗歌中得胜，} 那就是在那首诗歌中，圣咏中有类似的话：\BibleQuote{他使我脚立磐石，稳定我的脚步，使我口中唱新歌，赞美我们的天主。} 因此，凡在主之诗歌中得胜的，是在赞美他，而非自己的赞美；因为\BibleQuote{夸耀的，当在主内夸耀。} 但有些抄本写作：\Quote{我要因天主我的耶稣而欢乐，} 这在我看来比那些试图译成拉丁文却未写下那个名字——那个对我们而言更亲切、更甘甜的名字——的版本更好。\par

% structure: p0088 source=h2 target=subsection reason=relative-heading-rank; base=h2

\subsection{第三十三章——耶肋米亚与索福尼亚借着先知之神，预先论及基督与万民蒙召}

\Person{耶肋米亚}，如同\Person{依撒意亚}，是大先知之一，而非小先知，就像我刚刚引用其著作的其他小先知。他在\Person{约史雅}作王于\Place{耶路撒冷}，\Person{安古斯·马提乌斯}作王于\Place{罗马}时预言，当时犹太人的被掳已近在眼前；他继续预言直到被掳的第五个月，从著作可知。\Person{索福尼亚}，一位小先知，与他一起被提及，因为他说自己在\Person{约史雅}的日子预言；但未说预言到何时。因此，\Person{耶肋米亚}不仅预言了\Person{安古斯·马提乌斯}时期，也预言了\Person{塔奎尼乌斯·普里斯库斯}时期，这人是罗马人的第五位王。因为在那次被掳发生时，他已经开始统治。\Person{耶肋米亚}在预言基督时说：\BibleQuote{我们口唇的气息，主基督，因我们的罪被夺去}\BibleRef{哀 4:20}，简略地显示了基督是我们的主，以及祂为我们受难。也在另一处说：\BibleQuote{这是我的天主，没有别的能与祂相比；祂寻得一切明智之路，赐给祂的仆人\Person{雅各伯}，祂的爱人\Person{以色列}：此后祂在地上显现，与人往来。}有人不将这证词归于\Person{耶肋米亚}，而是归于他的秘书，名叫\Person{巴路克}；但更普遍地归于\Person{耶肋米亚}。同一位先知又论到祂说：\BibleQuote{看，日子将到——上主说——我必给\Person{达味}兴起一支正义的苗芽，一位君王将统治，必明智，在地上施行公道和正义。在那日子里，犹大必得救，以色列必安居；他们称祂的名字为：我们正义的主。}\BibleRef{耶 23:5}。关于将要发生的万民蒙召，我们现在已经看到应验，他这样说：\BibleQuote{上主，我的天主，我在患难之日的避难所，万民要从地极来到祢面前，说：我们的祖先实在敬拜了虚假的偶像，其中毫无益处。}\BibleRef{耶 16:19}。至于犹太人，祂甚至要被他们所杀，却不承认祂，这位先知暗示说：\BibleQuote{人心比万物都诡诈；祂是人，谁能认识祂？}\BibleRef{耶 17:9}。我在第十七卷中引用过的关于新约的那段经文也是他的，基督是新约的中保。因为\Person{耶肋米亚}自己说：\BibleQuote{看，日子将到——上主说——我必与\Person{雅各伯}家订立新约，}其余内容可在那里读到。\par

目前我将记录\Person{索福尼亚}先知关于基督的预言，他与\Person{耶肋米亚}一同说预言。\BibleQuote{你们等候我——上主说——直到我复活的日子，将来；因为我决意聚集万民，列国。}\BibleRef{索 3:8}。他又说：\BibleQuote{上主必向他们显现可畏，消灭地上一切众神；各人必在自己地方敬拜祂，甚至万民的海岛。}\BibleRef{索 2:11}。稍后他又说：\BibleQuote{那时我要转变万民的语言，使他们同声呼号上主的名，同心合意侍奉祂。从\Place{埃塞俄比亚}河流的边界，他们必给我献上祭品。在那日，你必不因你所行的一切叛逆而羞愧；因为那时我要从你中间除去你骄傲的罪行；你不再在圣山上自高。我必在你们中间留下谦卑贫弱的百姓，\Person{以色列}的遗民必敬畏上主的名。}这些就是\Person{保禄}所引用的别处预言的遗民：\BibleQuote{以色列子民的数目虽如海沙，但只有遗民得救。}这就是那民族中相信基督的遗民。\par

% structure: p0091 source=h2 target=subsection reason=relative-heading-rank; base=h2

\subsection{第三十四章——论\Person{达尼尔}与\Person{厄则克耳}的预言，另外两大先知}

\Person{达尼尔}和\Person{厄则克耳}，另外两大先知，也在\Place{巴比伦}被掳之初开始说预言。\Person{达尼尔}甚至精确地确定了基督降临和受难的时间。要详细计算显示出来会太冗长，而前人也已多次做过。但关于祂的权力和荣耀，他这样说：\BibleQuote{我在夜间的异象中观看，见有一位像人子的，乘着天上的云彩而来，直到万古常存者面前，被引到祂跟前。祂便赐给祂统治权、尊荣和国度；各民族、各邦国、各语言都侍奉祂。祂的统治是永远的统治，永不消逝，祂的国度也不会被毁灭。}\BibleRef{达 7:13}\par

\Person{厄则克耳}也以天主圣父的身份预言了基督，以预言的方式称祂为\Person{达味}，因为祂取了\Person{达味}的后裔的血肉，并取了仆人的形态成为人，祂是天主子，也被称为天主的仆人。他说：\BibleQuote{我必立一牧人牧放我的羊，就是我的仆人\Person{达味}；祂必牧养他们，作他们的牧人。我上主必作他们的天主，我的仆人\Person{达味}必在他们中间作领袖。我上主说了。}\BibleRef{厄 34:23}。在另一处说：\BibleQuote{一个君王将统治他们全体；他们不再成为两个民族，不再分为两个王国；他们不再因偶像、可憎之物和一切罪恶玷污自己。我要从他们犯罪居住的各地拯救他们，洁净他们；他们要作我的子民，我要作他们的天主。我的仆人\Person{达味}要作他们的君王，众人只有一个牧人。}\BibleRef{厄 37:22}\par

% structure: p0094 source=h2 target=subsection reason=relative-heading-rank; base=h2

\subsection{第三十五章——论三位先知\Person{哈盖}、\Person{匝加利亚}与\Person{玛拉基亚}的预言}

尚余三位小先知：\Person{哈盖}、\Person{匝加利亚}、\Person{玛拉基亚}，他们在被掳之末预言。其中\Person{哈盖}更明白地预言基督与教会，其言简短如下：\BibleQuote{万军的上主这样说：再过片刻，我将震动天、地、海和旱地；我将摇动万民，万民所渴望的必将来临。}\BibleRef{盖 2:6}这预言的应验，部分已见，部分留待末期待望。因为当基督降生成人时，祂藉天使和星辰的见证震动了天；祂藉由童贞玛利亚诞生的大奇迹震动了地；当基督在岛屿和普世被宣扬时，祂震动了海和旱地。如此，我们看见万民被震动而归于信德；而后续经文\Quote{万民所渴望的必将来临}的应验，则期待于祂末次的来临。因为人必须先信祂、爱祂，才能渴望并等候祂。\par

匝加利亚论及基督与教会说：\BibleQuote{熙雍女子，你应尽情欢乐！耶路撒冷女子，你应欢呼！看，你的君王来到你这里，正义的，拯救者，谦卑地骑着一匹驴，一匹驴驹。他的权柄从这海到那海，从大河直到地极。}\BibleRef{匝 9:9}主基督在行旅中使用此种役畜如何成就此事，我们在福音中读到，其中亦引用了此段预言中为上下文所需的部分。在另一处，他以预言之神对基督论及藉祂的血赦免罪过说：\BibleQuote{你藉你盟约的血，释放了你的囚徒脱离无水之坑。}\BibleRef{匝 9:11}关于这坑指什么，持正确信仰者可有不同见解。然而在我看来，最合适的解释是人的悲惨深渊，它仿佛干涸不毛，那里没有正义之流，只有罪孽的泥淖。正如圣咏论及它说：\BibleQuote{祂领我脱离悲惨的深渊，从污泥的沼泽中。}\par

\Person{玛拉基亚}预告我们现在所见藉基督传扬的教会，最公开地对犹太人讲话，以天主的口气说：\BibleQuote{我不喜悦你们，也不从你们手中收纳礼物。因为从日出到日落，我的名在万民中是伟大的；在各处要献祭，并要向我名献洁净的供物：因为我的名在万民中是伟大的，主说。}\BibleRef{拉 1:10} 既然我们现在已经看到这祭献从日出到日落，在各处藉基督按照默基瑟德品位的司祭职献给天主，而犹太人——曾对他们说过\BibleQuote{我不喜悦你们，也不从你们手中收纳礼物}——不能否认他们的祭献已经终止，为什么他们读到这预言并看见它实现（这只能藉祂实现）时，仍然寻找另一位基督呢？稍后，他论及祂，以天主的口气说：\BibleQuote{我的盟约是与祂所立的生命与平安的盟约：我将敬畏赐给祂，使祂敬畏我，在我名面前战栗。真理的律法在祂口中：祂以平安引导与我同行，并使许多人转离不义。因为司铎的嘴唇应守护知识，人应从祂口中寻求律法：因为祂是全能上主的天使。}\BibleRef{拉 2:5} 这也无怪乎基督耶稣被称为全能天主的天使。因为正如祂因取了奴仆的形状来到人间而被称为奴仆，同样祂也因所宣告的福音而被称为天使。如果我们解释这些希腊词，\OriginalTerm{福音}{evangel}就是\Quote{好消息}，而\OriginalTerm{天使}{angel}就是\Quote{使者}。他又论及祂说：\BibleQuote{看，我要派遣我的天使，祂要在我面前预备道路：你们所寻求的主，必忽然进入祂的殿，就是你们所渴望的盟约的天使。看，祂来了，全能的上主说，祂来临的日子谁能忍受？祂显现时谁能站立得住？}\BibleRef{拉 3:1} 在此处，他预告了基督的第一次和第二次来临：第一次来临，即祂所说的\BibleQuote{祂必忽然进入祂的殿}，即指祂的肉身，关于此，祂在福音中说：\BibleQuote{拆毁这座圣殿，三天内我要把它重建起来。}\BibleRef{若 2:19} 至于第二次来临，祂说：\BibleQuote{看，祂来了，全能的上主说，祂来临的日子谁能忍受？祂显现时谁能站立得住？} 但祂所说的\BibleQuote{你们所寻求的主，和你们所渴望的盟约的天使}的意思正是：连犹太人，按他们所读的经文，也将寻求和渴望基督。但许多人并未认出他们所寻求和渴望的那位已经来了，因为他们内心被蒙蔽，念念不忘自己的功德。那么，他在这里称为盟约的，无论是在上面他说\BibleQuote{我的盟约与祂同在}，还是在这里他称祂为盟约的天使，毫无疑问都应理解为新盟约，其中所许诺的是永恒的，而非旧盟约，其中所许诺的只是暂时的。然而许多软弱的人在看见恶人在这些暂时的事物上富足时感到困扰，因为他们看重大这些事物，并服事真天主以求得到这些赏报。为此，为区别新盟约中只赐给善人的永恒福乐与旧盟约中大多连恶人也得到的世俗幸福，同一位先知说：\BibleQuote{你们以言语使我不堪重负：你们却说：我们说了什么反对祢的话？你们说：凡服事天主的都是愚昧；我们遵守祂的规诫，在上主全能者面前如哀求者行走，有什么益处？现在我们称外邦人为有福；是的，凡行恶的都重新建立；是的，他们反对天主却得救。敬畏上主的人彼此对邻舍这样说：上主垂听且听见了；祂在祂面前写了一本纪念册，为那些敬畏上主和尊重祂名的人。}\BibleRef{拉 3:13} 那本书指的是新盟约。最后，让我们听祂接着说：\BibleQuote{他们必成为我的产业，全能的上主说，在我所定的日子；我必拣选他们，如同人拣选服事他的儿子。你们必回转，辨别义人与恶人，服事天主与不服事天主的人。因为，看，那日子来临，炽热如炉，它必焚烧他们；一切外邦人和行恶的必如碎秸；那将要来的日子必烧尽他们，全能的上主说，不留下根或枝。但为你们敬畏我名的人，义德的太阳必升起，其翼下有医治；你们必出来跳跃，如出栏的牛犊。你们必践踏恶人，他们在你们脚下成为灰烬，在我行这事的日子，全能的上主说。} 这日子就是审判的日子，如果天主愿意，我们将在适当的地方更详细地讨论它。\par

% structure: p0098 source=h2 target=subsection reason=relative-heading-rank; base=h2

\subsection{第36章——关于\BibleBook{厄斯德拉}{书}及\BibleBook{玛加伯}{书}}

继这三位先知哈盖、匝加利亚、玛拉基亚之后，在人民从巴比伦奴役中获得解放的同一时期，厄斯德拉也写了书，他的作品更偏重历史而非预言；同样，称为\WorkTitle{艾斯德尔}的书，也记述了为赞美天主而发生的近乎当时的事迹，除非我们应将厄斯德拉在以下段落中理解为预言基督：当时有几个青年争论什么是最强大的东西，一人说是君王，另一人说是酒，第三人说是妇女（妇女大多统治君王），但同是这第三个青年证明真理战胜一切。因为我们查阅福音得知基督就是真理。从圣殿重建之时起，直到亚里士多布鲁时期，犹太人没有君王，只有首领；他们的年代记不在称为正典的圣经中，而是在其他著作中，其中包括\WorkTitle{玛加伯}书。这些书虽不被犹太人视为正典，却被教会视为正典，因为其中记载了一些殉道者极其奇妙而惊人的苦难，他们在基督降生成人之前，就为天主的法律奋战至死，忍受了最可怕、最可怖的灾祸。\par

% structure: p0100 source=h2 target=subsection reason=relative-heading-rank; base=h2

\subsection{第三十七章——论先知记录比外邦哲学的任何泉源都更古老。}

因此，在我们的先知时代——他们的著作已为几乎所有民族所知——外邦的哲学家尚未出现，至少那些被称为哲学家的人尚未出现；这个名称源自撒摩斯岛的毕达哥拉斯，他在犹太囚禁结束时开始闻名。那么，其他哲学家更是在先知之后了。因为连雅典的苏格拉底——当时所有最著名学者的宗师，在被称为道德或实践哲学的领域中享有卓越地位——在年代记中也被发现晚于厄斯德拉。柏拉图出生也不迟很多，他远远超过了苏格拉底的其他弟子。此外，如果我们再考虑他们之前的那些人，当时尚未被称为哲学家，即七贤，以及继泰勒斯之后研究事物本性的自然哲学家，如阿那克西曼德、阿那克西美尼、阿那克萨戈拉等人，在毕达哥拉斯首次自称哲学家之前，这些人也并未在年代上早于我们所有的先知；因为泰勒斯——其他人继承了他——据说活跃于罗慕路斯统治时期，那时预言之流已从以色列的泉源喷涌而出，以那些传遍全世界的著作为表现。因此，只有那些神学诗人，如俄耳甫斯、利诺斯、穆赛俄斯，可能还有希腊人中的其他一些人，在年代上早于我们奉为权威的希伯来先知的著作。但即使是他们，在时间上也没有早于我们真正的神视者梅瑟，他权威地宣讲了唯一真天主，他的著作是正典中的首批；因此，希腊人——他们的语言主导了这一时代的文学——没有理由夸耀他们的智慧，因为我们的宗教（其中包含着真智慧）至少在年代上显然更古老，如果不是更优越的话。然而必须承认，在梅瑟之前，虽不在希腊人中，但在野蛮民族中，例如在埃及，已经存在某种可称为其智慧的学说；否则圣经中就不会记载梅瑟精通埃及人的一切智慧，\BibleRef{宗 7:22} 正如他那样，因为他生于埃及，被法郎的女儿收养并抚养，接受了全面的教育。然而，即使是埃及人的智慧，在时间上也不能早于我们先知的智慧，因为连亚巴郎也是一位先知。在伊西斯——她死后被尊为女神——赐予埃及人文字之前，埃及能有什么智慧呢？据说伊西斯是伊那科斯的女儿，而伊那科斯在阿尔戈斯开始为王时，亚巴郎的孙子们已知已经出生了。\par

% structure: p0102 source=h2 target=subsection reason=relative-heading-rank; base=h2

\subsection{第三十八章——论教会正典未收录某些著作，因其过于古老，以免通过它们将虚假之事掺入真理。}

若我回顾更古老的时代，我们的始祖诺厄确实在那场大洪水之前就已存在，我称他为先知并非不当，因为他所建造的方舟，他与家人借此获救，其本身便是对我们时代的预言。至于七世祖厄诺客，宗徒犹达的公函不是宣布了他曾预言吗？\BibleRef{犹 1:14}但这些人的著作，无论对犹太人或对我们，都不能视为权威，因其过于古老，以致必须以怀疑的眼光看待，以免将虚假之事当作真理。有些据称是他们所作的著作被那些凭自己喜好随意相信的人引述。但正典的纯洁性并未接纳这些著作，不是因为拒绝那些中悦天主的人的权威，而是因为不相信这些著作确实出自他们之手。若有人对如此古老著作的真实性存疑，这也不足为怪，因为即使在为记载犹大和以色列诸王事迹的史书中——我们相信这些史书属于圣经正典——也提及许多未予解释的事，并说明它们见于先知所著的另一些书中，有时甚至提及这些先知的名字，却未出现在天主之民所接受的正典中。我承认这原因对我隐藏；不过我认为，那些得蒙圣神启示，应当作为宗教权威来接受之事的人，可能会凭借历史性的勤奋书写某些事，又因天主默感而预言另一些事；这些事如此分明，可以判断前一类应归于他们自身，后一类应归于藉他们说话的天主。因此，前一类属于知识的丰富，后一类属于宗教的权威。在这权威中，正典得以保存。所以，如今若有人以古代先知之名提出正典以外的著作，它们甚至不能作为知识的辅助，因为不确定其真伪；因此它们不被信任，特别是那些其中包含与正典书籍真理相悖之事的著作，足以显明它们不属于正典。\par

% structure: p0104 source=h2 target=subsection reason=relative-heading-rank; base=h2

\subsection{第三十九章　论希伯来文字始终为该语言所拥有。}

现在，我们不可相信：希伯尔之名衍生出\Quote{希伯来}一词，他只将希伯来语言作为口语传给亚巴郎，而希伯来字母始于藉梅瑟颁布法律；相反，这语言连同其字母是藉祖先的传承而保存的。的确，梅瑟在天主子民中指派了一些人，在他们了解神圣法律之任何字母之前就教授文字。圣经称这些人为\OriginalTerm{文字导师}{\GreekText{γραμματεισαγωγεῖς}}，拉丁文可称为\OriginalTerm{引导者}{inductores\/introductores}，因为他们仿佛将字母引入学习者心中，或更确切地说，引导他们所教之人进入文字。因此，没有一个民族能因古代智慧而以其邪恶的虚荣心高过我们的族长和先知；因为即使是埃及——它惯于虚假而徒劳地夸耀其学说的古老——也未能发现其智慧（无论何等智慧）在时间上早于我们族长的智慧。谁也不敢说，在了解文字（即伊西斯来教他们之前）之前，他们就在奇妙科学上极为精通。此外，他们那值得纪念的所谓智慧学说，大部分难道不就是天文学，以及可能某些类似的科学，这些通常更能锻炼人的心智，而非以真智慧启迪人心吗？至于哲学，它自称教授使人幸福之事，这类学问在墨丘利（被称为特里斯墨吉斯忒斯）的时代在那些地方兴盛，远在希腊智者与哲学家之前，但仍晚于亚巴郎、依撒格、雅各伯和若瑟，甚至晚于梅瑟本人。事实上，当梅瑟诞生时，我们发现那位伟大的天文学家阿特拉斯尚在世，是普罗米修斯的兄弟，老墨丘利的外孙，而那位特里斯墨吉斯忒斯墨丘利则是其孙。\par

% structure: p0106 source=h2 target=subsection reason=relative-heading-rank; base=h2

\subsection{第四十章　论埃及人最虚假的虚荣，他们将科学归诸于十万年的古老。}

因此，某些人以最空洞的狂妄喋喋不休，说埃及理解星象已有超过十万年之久，这完全是徒劳。因为那位撰写这些数字的人，其书写文字是伊西斯女主人所授，时间不过两千多年前；而瓦罗在历史中宣布了这一事实，他并非微不足道的权威，且与神圣书籍的真理并无冲突。因为自从第一个人（被称为亚当）以来，尚未满六千年，那些试图说服我们相信如此不同且与已知真理相背的时间跨度的人，岂不应该被嘲笑而非被反驳吗？我们应当信靠哪一位历史学家关于过去的记载，甚于那位也预言了我们如今看见应验之事的人？而且，历史学家彼此之间的分歧本身提供了一个好理由，说明我们更应相信那位与我们持守的神圣历史并不矛盾的人。另一方面，不虔敬之城的人民散布世界各地，当他们阅读最博学的作家（其中似乎没有一位是等闲之辈的权威）时，发现他们在回忆我们时代最为遥远的事实时彼此矛盾，便无法决定该信靠谁。但我们，因我们宗教的历史所承托的天主权威而坚定，相信凡与此相悖者必为最虚假，至于世俗书籍中的其他事，无论真假，对于我们正确地生活并获得幸福，都无关紧要。\par

% structure: p0108 source=h2 target=subsection reason=relative-heading-rank; base=h2

\subsection{第四十一章　论哲学意见之纷争与教会视为正典的圣经之和谐。}

但让我们略过对历史的进一步考察，回到我们因此而离题的哲学家们。他们似乎辛苦钻研，只为求得一种适于把握\OriginalTerm{真福}{blessedness}的生活方式。那么，为何门徒与师傅意见分歧，同门之间也彼此不合，除非是因为他们作为人，凭人的感觉和人的推理来寻求这些事？如今，虽然他们中间可能有一种求名的欲望，以致每个人都希望被认为比他人更聪明、更敏锐，绝不受他人判断的影响，而是自己教义和意见的创始人，然而我可以承认，他们中间有一些人，甚至非常多的人，其爱真理之心使他们离开了自己的老师或同门，以便为他们认为是真理的东西而努力，无论它是否真是真理。然而，若没有\OriginalTerm{神圣权威}{divine authority}引导，人类的悲惨又能做什么，它如何、在哪里能够达到真福呢？最后，愿我们的作者们——在他们之中，\OriginalTerm{圣书正典}{canon of the sacred books}是确定而有界限的——在任何方面都绝不彼此相悖。因此，并非没有充分的理由，不仅少数在学园和体育馆里以挑剔的辩论夸夸其谈的人，而是许多伟大的人，无论有学问的还是没有学问的，在各国和各城市中，都相信天主向他们或通过他们，\Emph{即} \OriginalTerm{正典作者}{canonical writers}说话，当他们写这些书时。确实，这样的人应当为数不多，免得因数量众多而使本该以宗教虔诚尊重的东西变得廉价；但也不应太少，以致他们的和谐一致不令人惊奇。因为在众多哲学家中，那些在著作中留下了他们教义纪念碑的人，没有人能轻易找到在任何观点上都一致的人。但要证明这一点对于这部著作来说太长了。\par

但是，在这座崇拜魔鬼的城市中，有哪一个教派的创始人如此受认可，以至于其他在意见上与他不同或反对他的人都受到不认可？\OriginalTerm{伊壁鸠鲁派}{Epicureans}声称人类事务不受\OriginalTerm{众神的眷顾}{providence of the gods}；而\OriginalTerm{斯多葛派}{Stoics}持相反意见，同意它们由仁慈和护佑的神明所统治和保护。然而，这两个派别不都在雅典人当中赫赫有名吗？那么，我奇怪，为何\Person{阿那克萨戈拉}因说太阳是一块燃烧的石头，并否认它完全是神而被告犯罪；而在同一座城市中，\Person{伊壁鸠鲁}却辉煌地繁荣并安全地生活，尽管他不仅不相信太阳或任何星辰是神，而且坚持认为朱庇特或任何神都根本不住在世界中，以致人的祈祷和恳求无法到达他们那里！难道\Person{亚里斯提卜}和\Person{安提斯泰尼}不都在那里，两位高贵的哲学家，且都是\Person{苏格拉底}的门徒吗？然而，他们将人生的终极目标置于如此多样且矛盾的范围之内：前者把\OriginalTerm{身体的快乐}{delight of the body}当作\OriginalTerm{至善}{chief good}，而后者则断言人主要因\OriginalTerm{心灵的德行}{virtue of the mind}而幸福。一位还说智者应逃离共和国；另一位说他应管理共和国的事务。然而，不是每个人都聚集了门徒来追随自己的派别吗？确实，在显眼而著名的柱廊、体育馆、花园、公共场所和私密之地，他们成群结队地公开为自己的意见而争辩：有些人断言只有一个世界，其他人则有无数个；有些人认为这个世界有开端，其他人则认为没有；有些人认为它会毁灭，其他人则认为它永远存在；有些人认为它由神圣的心智统治，其他人则认为是机遇和偶然；有些人认为\OriginalTerm{灵魂不朽}{souls are immortal}，其他人则认为灵魂会死——而在那些主张灵魂不朽的人中，有些人说灵魂会\OriginalTerm{转生为兽}{transmigrate through beasts}，其他人则说绝非如此；而在那些主张灵魂会死的人中，有些人说灵魂在身体死后立即消亡，其他人则说它们会存活一小段时间或更长时间，但并非永远；有些人在身体中确定\OriginalTerm{至善}{supreme good}，有些人在心智中，有些人则在两者中；其他人则在心智和身体之外再加上外在的善；有些人认为身体的感官应始终被信任，有些人认为并非始终，另一些人则认为从不信任。现在，这座不虔敬之城的哪个民族、元老院、权力机构或公共尊严曾费心在所有这些以及几乎无数的哲学家分歧之间进行判断，认可并接受一些，而不认可并拒绝另一些呢？它不是毫无判断地、混乱地将其怀抱中这么多不合之人的争论随意包容吗？这些争论不是关于田地、房屋或任何金钱性质的东西，而是关于那些使生活变得悲惨或幸福的事物。即使其中有些真实的东西被说出来，谎言也同样自由地被说出；以致这样一座城市没有错误地获得了\OriginalTerm{神秘的巴比伦}{mystic Babylon}的称号。因为巴比伦意为混乱，正如我们记得我们已经解释过的。而且，对它的君王魔鬼来说，他们如何在相互矛盾的错误中争吵并不重要，因为由于他们巨大而多样的不虔敬，他们都同样理应属于他。\par

但那个民族、那个百姓、那座城、那个共和国，这些以色列人，天主的谕言交托给他们，他们绝没有以类似的放纵将假先知与真先知混淆；而是彼此一致，毫无分歧，承认并维护他们圣书的正统作者。这些就是他们的哲学家，这些就是他们的智者、神学家、先知、以及正直和虔诚的导师。凡有智慧并按照他们生活的人，便是按照天主生活，不是按照人生活，因为天主藉着他们说话。若那里禁止亵渎，天主便禁止了它。若它说：\Quote{应孝敬你的父亲和你的母亲，}\BibleRef{出 20:12} 天主便命令了它。若它说：\Quote{不可奸淫，不可杀人，不可偷盗，} 以及其他类似的诫命，那不是人的口，而是神圣的谕言宣告的。某些哲学家在他们的错误见解中能够看到的任何真理，并努力通过艰难的讨论来说服人们——例如，天主创造了这个世界，并自己以极大的眷顾治理它；或关于德性的高贵、对祖国的爱、友谊的忠诚、善行以及一切与美德有关的事——尽管他们不知道这些事要归于什么目的和什么准则——所有这些，都通过先知的话语，即神圣的话语，虽然是藉着人说的，都被推荐给那座城里的百姓，而不是通过辩论中的争论来灌输，所以知道这些事的人会害怕藐视，不是人的才智，而是天主的谕言。\par

% structure: p0112 source=h2 target=subsection reason=relative-heading-rank; base=h2

\subsection{因天主眷顾的何种安排，旧约的圣经从希伯来文译成希腊文，以便为万民所知。}

埃及的一位托勒密王，渴望知道并拥有这些圣书。因马其顿的亚历山大（亦被称为大帝）曾以他神奇但毫不持久的武力，征服了整个亚洲，甚至几乎整个世界——部分凭武力，部分凭恐怖——并且在进入并占据犹太（也是东方诸王国之一）之后，他死时，他的将领们并未和平地瓜分那最广阔的王国作为产业，反而因战争消耗了一切，把它弄得四分五裂。于是埃及开始有托勒密诸王作她的君主。他们中的第一个，拉各斯之子，曾从犹太携许多俘虏到埃及。但另一位继位的托勒密，称为斐拉德尔弗斯，却允许所有被他征服的人自由返回；不仅如此，他还向天主的圣殿赠送了王室的礼物，并请求大司祭厄肋阿匝尔将圣经给他。他风闻那圣经确实是神圣的，因此非常渴望将它收入他所建造的那最著名的图书馆中。当大司祭将希伯来文的圣经送给他后，他又向大司祭请求翻译者。于是给了他七十二人，从十二支派中各选出六人，精通两种语言，即希伯来文和希腊文。他们的译本如今习惯上称为\OriginalTerm{七十贤士译本}{Septuagint}。据说，他们的文字如此奇妙、惊人、且显然神圣地一致，以至于当他们各自单独（因为托勒密愿意这样考验他们的忠诚）进行这项工作的时候，在任何一个具有相同意义和力量的词上，甚至在词序上，彼此并无差异；仿佛译者是一个人，所以所有人翻译的都相同，因为实际上只有一个圣神在他们众人内。他们接受了天主如此奇妙的恩赐，是为了使这些圣经的权威被推崇为神圣的而非人为的——事实确是如此——以造福那些终将相信的万民，正如我们现在所看到的。\par

% structure: p0114 source=h2 target=subsection reason=relative-heading-rank; base=h2

\subsection{论七十贤士译本的权威——在不损害希伯来原文的尊崇之下，应优先于一切译本。}

因为，虽然有其他译者将这些神圣神谕从希伯来语译成希腊语，例如\Person{阿奎拉}、\Person{辛马库}和\Person{特奥多提翁}，此外还有译本（其作者的名字不详）被引用为第五版，但教会接受这\WorkTitle{七十贤士译本}，仿佛它是唯一的一本；并且希腊基督徒民众使用它，他们中的大多数人不知道还有其他译本。从这一译本，也产生了一个拉丁语译本，供拉丁教会使用。然而，我们的时代有幸得益于司铎\Person{热罗尼莫}，一位博学者，精通三种语言，他将同一圣经从希伯来语而非希腊语译成拉丁语。但是，尽管犹太人承认他这一非常博学的工作是忠实的，他们却主张七十贤士译者在许多地方有误，然而基督的教会评判：没有人应优先于由当时的大司祭\Person{厄肋阿匝尔}选来承担这一极其伟大工作的众多人的权威；因为即使在他们身上没有显现出同一个无疑神圣的神，而七十位有学问的人按人的方式比较了他们译文的措辞，以便他们所有人都同意的内容得以成立，也不应偏爱任何单个译者；但由于在他们身上出现了如此伟大的神圣标记，当然，如果有人将他们的圣经从希伯来语译成其他语言是忠实的，那么他就与这七十位译者一致；如果发现他与他们不一致，那么我们就应当相信先知之恩在他们那里。因为同一圣神，当先知们说这些事时，也在他们里面，也在那七十位译者里面，当他们翻译它们时；因此，他们确实也可以说些别的东西，就好像先知本人说了两者，因为说两者的会是同一圣神；并且可以用不同的方式说同样的事情，这样，虽然措辞不同，但同一意义对善解人意者闪耀出来；并且可以省略或添加某些东西，以便甚至借此表明，在那工作中不是人的束缚（译者所欠于措辞的），而是神圣的能力，它充满并支配译者的心智。有些人却认为，七十贤士译本的希腊抄本应从希伯来抄本修正；但他们不敢去掉希伯来文所缺而七十贤士译本所有的东西，只添加了希伯来抄本中有而七十贤士译本所缺的东西，并通过在每节开头放置某些星形标记（他们称之为星号）来标注它们。而那些希伯来抄本没有、但七十贤士译本有的东西，他们也类似地在每节开头用水平线形标记（像我们表示盎司的标记）来标注；许多带有这些标记的抄本甚至在拉丁语中流通。但是，不检查两种抄本，我们就无法发现那些既未省略也未添加而是不同表述的东西，无论它们产生另一个本身并非不合适的含义，还是可以被显示为以另一种方式解释同一含义。因此，如果我们，如同应当的，在圣经中只看到圣神通过人所说的话，那么，如果希伯来抄本中有某些东西而七十贤士译本中没有，圣神就不曾选择通过他们说它，而只通过先知说它。但是，凡存在于七十贤士译本中而不存在于希伯来抄本中的，圣神反而选择通过前者说它，从而表明两者都是先知。因为祂按照祂的选择说话，有些通过\Person{依撒意亚}，有些通过\Person{耶肋米亚}，有些通过几位先知，或者同一件事通过这位先知和那位先知。此外，凡在两个版本中都找到的，同一位圣神意愿通过两者说它，但前者在预言中领先，后者在预言的诠释中跟随；因为，正如和平的同一圣神在前者中，当他们说真实一致的话语时，所以同一圣神也出现在后者中，当他们没有互相商议时，却仿佛异口同声地解释了一切。\par

% structure: p0116 source=h2 target=subsection reason=relative-heading-rank; base=h2

\subsection{第四十四章：如何理解对尼尼微人毁灭的威胁——在希伯来文中原为四十天，而在\WorkTitle{七十贤士译本}中缩短为三天。}

但有人可能说：\Quote{我怎么知道先知约纳对尼尼微人说的究竟是\InnerQuote{再过\Emph{三}天，尼尼微就要被倾覆}，还是\Emph{四十}天呢？} \BibleRef{纳 3:4} 因为谁看不出，先知被派遣来以毁灭迫在眉睫的威胁恐吓那城时，不可能说两种话？如果它的毁灭在第三天发生，当然就不可能在第四十天；但如果在第四十天，就肯定不在第三天。那么，如果有人问我约纳说的究竟是哪一句，我宁可认为希伯来文所读的是：\Quote{再过四十天，尼尼微就要被倾覆。} 但七十贤士译员在很久以后翻译时，可以说出不同但与主题相关的话，并在同一意义上相符，尽管措辞不同。这可以提醒读者不要轻视任何一方的权威，而要超越历史，去探寻历史本身所要阐明的那些事。这些事确实发生在尼尼微城，但也预示着远超过那城的某种大事；正如先知自己在鲸鱼腹中度过了三天，这也预示着万有的主——所有先知的主——将在地狱的深处度过三天。因此，若合理地认为那城预言性地代表外邦人的教会，即因补赎而不再是从前的样子（这事是由基督在外邦人的教会中完成的，尼尼微预表了它），那么基督本身就既由四十天、也由三天所预表：由四十天，因为祂复活后与门徒共度了那些日子，然后升天；由三天，因为祂在第三日复活了。所以，如果读者只希望把握事件的历史，他就可以被七十贤士译员以及先知们从沉睡中唤醒，去探究预言的深意，好像他们说了：在四十天中寻求祂，你也必能找到那三天——其一在祂的升天中，另一在祂的复活中。因为那最适宜由两个数字所表示的事，其中一个由约纳先知所用，另一个由七十贤士译本的预言所用，是由同一位圣神所说的。我担心冗长，因此不愿用许多例子来证明，即使七十贤士译员似乎与希伯来文相异，但若正确理解，就会发现它们相符。为此，我也按照自己的能力，追随宗徒们的足迹——他们本人曾引用过两种圣经的预言见证，即希伯来文和七十贤士译本——认为两者都应用作权威，因为两者都是一体的、神圣的。但是，现在让我们尽可能继续讨论剩下的问题。\par

% structure: p0118 source=h2 target=subsection reason=relative-heading-rank; base=h2

\subsection{第45章——犹太人自从圣殿重建后不再有先知，从那时起直到基督诞生，不断遭受患难，为证明另一座圣殿的建造已由预言之声所应许。}

犹太民族无疑在失去先知之后变得更糟，恰好在巴比伦被掳后、圣殿重建之时，他们曾希望变得更好。因为那属肉体的百姓正是这样理解哈盖先知所预言的：\Quote{这后一座殿的光荣将大于前一座。} \BibleRef{盖 2:9} 现在，这说的是新约，先知稍早一点就显示了，他显然在应许基督时说：\Quote{我要震动万民，万民的渴望者将要来到。} \BibleRef{盖 2:7} 在这段经文中，七十贤士译员给出了另一个更切合身体而非元首，即更切合教会而非基督的解释，他们凭先知权威说：\Quote{万民中上主所拣选的人将要来到。} 即\Emph{那些人}，耶稣在福音中论到他们时说：\Quote{许多人蒙召，但少数人被拣选。} \BibleRef{玛 22:14} 因为正是通过新约，用那些被拣选的万民，以活石建成了一座远比所罗门王所建造、被掳后所重建的圣殿更光荣的天主居所。为此，从那以后，那民族不再有先知，而是被外族君王和罗马人本身用许多灾祸所打击，免得他们以为哈盖的这预言是由那次圣殿重建所应验的。\par

因为不久之后，随着亚历山大到来，它就被征服了；当时虽然没有抢掠，因为他们不敢抵抗他，于是很容易就被征服，和平地接待了他，然而那殿的光荣并不如在他们自己君王自由统治下那样大。亚历山大确实在天主的圣殿中献了祭，但不是作为皈依祂的敬拜而怀着真虔敬，而是以不敬的愚妄认为祂应与假神一同受敬拜。然后，我已经提到过的拉古斯之子托勒密，在亚历山大死后把他们掳到埃及。他的继任者托勒密·非拉德尔弗最仁慈地释放了他们；正是他促成了我们稍早叙述的，使我们有了七十贤士译本圣经。后来他们被玛加伯书中所记载的战争所压垮。之后，他们又被号称厄丕法乃的亚历山大城王托勒密所掳。然后，叙利亚王安提约古用许多最严重的灾难强迫他们崇拜偶像，并以外邦人的亵渎迷信充满了圣殿本身。然而他们最英勇的领袖犹大，又称玛加伯，在击败了安提约古的将领之后，清除了偶像崇拜的一切污秽。\par

但不久之后，一位名叫阿尔基摩的人，虽非司祭支族，却因野心被立为大司祭，这是一件亵渎之事。将近五十年期间，他们从未享有和平，虽在某些事务上兴盛，但亚里士托布鲁首先在他们中间戴上王冠，成为君王兼大司祭。在此之前，确实自从他们从巴比伦囚掳中回归并重建圣殿以来，他们没有君王，只有将军或\OriginalTerm{首领}{principes}。虽然君王本身可以因治理的元首地位而被称为元首，因率领军队而被称为领袖，但这并不意味着所有元首和领袖都可被称为君王，如同那位亚里士托布鲁一样。他的继任者是亚历山大，也是君王兼大司祭，据报告他残酷地统治了他们。他之后，他的妻子亚历山德拉成为犹太人的女王，从她那时起，更严重的灾祸追随着他们；因为这位亚历山德拉的儿子们，亚里士托布鲁和依尔卡诺，在彼此争夺王位时，召请罗马军队来对抗以色列民族。因为依尔卡诺向他们请求援助来对抗他的兄弟。那时，罗马已经征服了阿非利加和希腊，并在世界其他地区也广泛统治，然而，仿佛无法承受自身的重量，在某种程度上因自身庞大而破裂了。事实上，她已陷入严重的内部叛乱，进而发展为同盟战争，随后又发生内战，以至于衰弱损耗，共和国的政体即将转变，由君王统治的时刻已经迫近。于是，罗马人中最显赫的元首庞培率领军队进入犹太，攻占城池，打开圣殿，并非以恳求者的虔诚，而是以征服者的权威，他进入至圣所，不是恭敬地，而是亵渎地，那里除大司祭外无人可进。他立依尔卡诺为大司祭，并派安提帕特治理被征服的民族，担任当时所谓的监护人即总督，然后他捆绑着亚里士托布鲁将他带走。从那时起，犹太人也开始成为罗马的进贡者。此后，卡西乌斯洗劫了圣殿。又过了几年，他们活该有一个外族出生的君王黑落德，在他统治时基督诞生了。因为那时已经来到那由预言之神通过族长雅各伯的口所预示的时刻，他说：\BibleQuote{犹大支派中必不缺元首，从他腰间必不出教师，直到那为他所保留者来到；祂就是万民的期待。}\BibleRef{创 49:10}因此，直到黑落德为止，犹太人从不缺少犹太人的元首；黑落德是他们所接纳的第一个外族君王。因此，现在是那为他所保留者应来到的时候，在新约中所应许的，祂应成为万民的期待。但是，万民不可能期待祂来到（正如我们看见他们确实期待了）在权能的光辉中施行审判，除非他们先相信祂，当祂在忍耐的谦卑中前来接受审判时。\par

% structure: p0122 source=h2 target=subsection reason=relative-heading-rank; base=h2

\subsection{第四十六章 论我们救主的诞生，圣言成了血肉；以及犹太人散居万民之中，正如所预言的}

当黑落德在犹太统治，奥古斯都凯撒在罗马作皇帝时，共和国政体已经改变，天下被他平定，基督在犹太的白冷诞生，从童贞女而来显现为可见的人，但从天主圣父而来隐藏为不可见的天主。先知曾预言说：\Quote{看，童贞女要怀孕生子，人要称他为\OriginalTerm{厄玛奴耳}{Immanuel}，意思是：天主与我们同在。}他行了众多奇迹，为在自身内彰显天主；其中有些足以宣扬他自己，载于福音书。第一个奇迹是他如此奇妙地诞生，最后一个则是他带着从死者中复活的身体升天。但那些杀害他、不肯相信他的犹太人，因为他必须死而复活，后来被罗马人更加悲惨地摧毁，从自己的王国中被连根拔除——外邦人早已统治他们——且分散到各地，以致没有一处没有他们；他们就这样用自己的圣经作我们的见证，证明我们并没有伪造关于基督的预言。他们中有许多人，考虑到这一点——尤其是在他受难前，但主要是他复活后——信了他；关于这些人曾预言说：\Quote{以色列子民的数量虽如海沙，但仅剩的遗民将得救。}其余的人则被蒙蔽；关于他们曾预言说：\Quote{愿他们的眼睛昏暗，不得看见；使他们的脊背常常弯曲。}因此，当不信我们的圣经时，他们自己所盲目诵读的圣经就在他们身上应验了，免得有人以为基督徒伪造了那些以西彼拉或其他人（若有的话，不属于犹太民族）的名义引用的关于基督的预言。对我们而言，那些从仇敌之书引用的预言已足够；我们向他们承认，因着他们拥有这些书卷而无意中提供的见证，他们自己被分散到所有民族中间，凡基督的教会所及之处。关于这件事的诗篇中早有预言，他们也诵读，其中写道：\Quote{我的天主，他的仁慈必先迎接我。我的天主，关于我的仇敌，你指示了我：你不可杀他们，免得他们最后忘记你的法律；要用你的能力分散他们。}因此，天主在敌人犹太人身上向教会显示了他怜悯的恩宠，因为正如宗徒所说：\Quote{他们的过犯成了外邦人的救恩。}\BibleRef{罗 11:11}因此他没有杀死他们，即没有让他们丧失自己身为犹太人的意识——尽管他们被罗马人征服——免得他们忘记天主的法律，免得他们的见证在我们所讨论的事情上失效。但仅仅说\Quote{不可杀他们，免得他们最后忘记你的法律}是不够的，除非他又加上\Quote{分散他们}；因为如果他们只在自己的土地上带着那圣经的见证，而不在各地，那么普世的教会就不能在所有民族中拥有他们作为关于基督预言的见证人了。\par

% structure: p0124 source=h2 target=subsection reason=relative-heading-rank; base=h2

\subsection{第四十七章 —— 在基督时代之前，是否有任何不属于以色列民族的人属于天上之城的团体？}

因此，如果我们读到任何外邦人——即既非生于以色列、也未被那民族接纳进入神圣经卷正典的人——曾预言过关于基督的事，只要它曾或将要为我们所知，我们便可以额外引用它；这并不是必需的，即使缺少也无妨，但相信即使在别的民族中也可能有人蒙启示得知这奥迹、并被催逼去宣扬它，也是合理的，无论他们是分享了同一恩宠，还是未曾经验它，而是由恶天使所教导，我们知道，连那些天使也承认了眼前的基督，而犹太人却不承认。我也不认为犹太人自己敢断言除了以色列人之外无人属于天主，因为以色列的壮大始于其长兄的被弃。事实上，确实没有别的民族被特称为天主的子民；但他们不能否认，即使在别的民族中，也有人——不凭属地而凭属天的团契——属于真正的以色列人，即那在上之国的公民。因为，如果他们否认这一点，他们很容易被圣洁非凡的约伯的例子所驳倒，他既非本地人也非皈依者，即不是加入以色列民族的外邦人，而是出生于厄东族，在那里出生也在那里去世，且受天主的启示如此称赞，以致在他那个时代，没有人在公义和虔敬上能与他相提并论。虽然我们在编年史中找不到他的年代，但从他那本因功德而被以色列人接纳为权威正典的书中，我们推断他生活在以色列之后第三代。我深信这是天主的安排，好让我们从这个例子知道，在别的民族中也可能有属于属灵的耶路撒冷的人，他们按照天主生活并中悦了他。也不应认为这恩赐曾赐予任何人，除非那唯一的天主与人之间的中保、降生成人的基督耶稣，\BibleRef{弟前 2:5}曾以神圣方式启示给他；祂从前曾向古圣先贤预告，说要来临于肉身中，正如祂现已来临而被宣告给我们一样，好使同一的信德藉着祂引导所有预定成为天主之城、天主之家、天主之殿的人归向天主。但是，凡引用来论及藉基督耶稣而来的天主恩宠的预言，都可能被以为是基督徒伪造的。因此，在驳斥各类外邦人（若他们就此争辩）并支持我们的朋友（若他们真有智慧）时，最有力的莫过于引用那些关于基督的神圣预言，它们写在犹太人的书中，而犹太人已从其本土被拔除、分散到全世界以作此见证，从而使基督的教会处处扩展。\par

% structure: p0126 source=h2 target=subsection reason=relative-heading-rank; base=h2

\subsection{第四十八章——哈盖的预言，他说天主殿宇的荣耀将大于先前的，这预言实际上并非在圣殿的重建中、而是在基督的教会中应验。}

天主的这座殿宇比那第一座由木头、石头、金属及其他贵重材料建成的殿宇更为荣耀。因此，哈盖的预言并未在那圣殿的重建中应验。因为那座圣殿重建后，绝无法证明它享有像撒罗满时代那样的荣耀；相反，那殿宇的荣耀表现为减少：首先预言停止了，然后民族本身遭受如此大的灾祸，甚至最终被罗马人毁灭，正如上述之事所证明的。但属于新约的这座殿宇，其荣耀程度正比于构成它的活石——即信德、被更新的更好的人。但之所以由那圣殿的重建来预表它，是因为那建筑本身的重建在预言性的神谕中预表了另一个称为\Quote{新}的盟约。因此，当天主藉上述先知说：\BibleQuote{我要在这地方赐予平安}\BibleRef{盖 2:9}时，我们应当理解祂是指那被那象征性地方所预表的那位；因为既然那被重建的地方预表了将由基督建成的教会，那么\BibleQuote{我要在这地方赐予平安}这句话，除了解释为\Quote{我要在那地方（即此地所象征的地方）赐予平安}之外，不可能有其他含义。因为一切象征性事物在某种程度上都仿佛代表它们所象征的那些事物，正如宗徒所说：\BibleQuote{那磐石就是基督。}因此，这新约殿宇的荣耀大于旧约殿宇的荣耀；当它被祝圣时，它将显得更大。因为那时\BibleQuote{万民所渴望的必来临}\BibleRef{盖 2:7}，正如我们读到的希伯来文本。因为在祂来临之前，祂尚未被万民所渴望。因为他们不认识祂，既然未信祂，就不应渴望。此外，按照\WorkTitle{七十贤士译本}的解释（因为它也含有预言意义）：\BibleQuote{那些被上主从万民中拣选的人必来临。}因为那时确实只有那些被拣选的人会来，关于他们，宗徒说：\BibleQuote{就如祂在创世以前，在祂内拣选了我们。}\BibleRef{弗 1:4}因为那位建筑大师说：\BibleQuote{许多人是被召，但少数人是被选}\BibleRef{玛 22:11}，祂不是指那些被召而进入宴席却被赶出的人，而是指那座由被选者建成的房屋，它从此将不再惧怕任何毁灭。然而，由于教会中也充满了那些将在扬谷时被分开的人，如同在禾场上一样，这殿宇的荣耀如今并不那么明显，而当其中每一个人都永远在那里时，它才会彰显。\par

% structure: p0128 source=h2 target=subsection reason=relative-heading-rank; base=h2

\subsection{第四十九章——论教会不加区别地增长，其中许多被弃绝者在此世与蒙选者混杂。}

在这个邪恶的世界，在这些罪恶的日子里，教会以她现世的谦卑衡量她未来的崇高，被刺痛的恐惧、折磨的悲伤、不安的劳苦和危险的诱惑所锻炼，她清醒地喜乐，只在望德中喜乐，有许多被弃绝的人与善人混杂在一起，两者都被福音像网一样聚集起来；\BibleRef{玛 13:47} 在这个世界上，如同在海中，两者毫无区别地被网围在一起游泳，直到网被拉上岸，那时恶人必须从善人中分离出来，以便在善人中，如同在祂的圣殿里，天主成为万有中的万有。我们确实承认，那在圣咏中发言并说：\Quote{我已宣布并发言；他们多得数不胜数。}的祂的话如今应验了。这事如今发生了，因为祂已经发言，首先通过祂的前驱\Person{若翰}的口，然后通过祂自己的口，说：\Quote{你们悔改吧！因为天国临近了。}祂拣选了门徒，祂也称他们为宗徒，\BibleRef{路 6:13}他们出身卑微、不受尊敬、没有学问，这样，无论他们可能成为或做什么伟大的事，祂可以在他们内成为并做这些事。祂在他们中间有一个，其邪恶祂能善加利用，以完成祂预定的苦难，并为祂的教会提供一个容忍恶人的榜样。祂以肉身临在撒下了神圣福音的种子，照所应做的，祂受苦、死亡并复活了，以祂的苦难显示我们应为真理忍受什么，以祂的复活显示我们在逆境中应盼望什么；始终保留圣事的奥迹，祂的血为了赦免罪过而倾流。祂在地上与门徒交谈了四十天，并在他们眼前升天，十天后派遣了所应许的圣神。这赐给了那些信者，作为祂来临的首要且最必要的标记：他们每个人都讲万国的语言；从而预示天主教会将包含万国，并且同样将讲万国的语言。\par

% structure: p0130 source=h2 target=subsection reason=relative-heading-rank; base=h2

\subsection{第50章——论福音的宣讲，由于宣讲者的苦难而变得更加著名和有力。}

那时，那预言应验了：\Quote{法律将出自熙雍，上主的话将出自耶路撒冷；}\BibleRef{依 2:3} 以及主基督自己的预言，当祂复活后，\Quote{开启了}祂惊愕的门徒的\Quote{理智，使他们理解圣经，并对他们说：经上这样记载：基督必须受苦，第三日从死人中复活，并且因祂的名，从耶路撒冷开始，向万国宣讲悔改和罪过的赦免。}\BibleRef{路 24:45} 又当祂回答他们关于祂最后来临的日子的提问时，祂说：\Quote{父以自己的权柄所定的时候和日期，不是你们应当知道的；但当圣神降临于你们身上时，你们将接受能力，并要在耶路撒冷、全犹太和撒玛黎雅，直到地极，为我作证人。}\BibleRef{宗 1:7} 首先，教会从耶路撒冷向外扩展；当犹太和撒玛黎雅有许多人相信后，她也通过那些宣讲福音的人前往其他国家，这些人如同光明，祂曾以自己的言语预备他们，并以圣神点燃他们。因为祂曾对他们说：\Quote{不要害怕那些杀害身体而不能杀害灵魂的。}\BibleRef{玛 10:28} 并且，为了不因恐惧而冻结，他们以爱德之火燃烧。最后，基督的福音传遍了全世界，不仅由那些在祂受难前和复活后见过、听过祂的人宣讲，而且在这些人去世后，由他们的继承者在可怕的迫害、各样的酷刑和殉道者的死亡中宣讲，天主也以征兆和奇事、各种奇迹和圣神的恩赐为他们作证，\BibleRef{希 2:4} 以便万民相信那为他们得救赎而被钉十字架者，能以基督徒的爱敬奉那些被魔鬼般狂怒倾流的殉道者的血，并且那些曾以法律使教会荒芜的君王们，能有益地服从他们曾残忍地试图从地上抹去的那个名字，并开始迫害那些因之真天主的敬拜者曾受迫害的假神。\par

% structure: p0132 source=h2 target=subsection reason=relative-heading-rank; base=h2

\subsection{第51章——论天主教的信德甚至可由异端者的分歧而得到证实。}

但魔鬼看见邪魔的庙宇荒凉，人类纷纷奔向那解救者的中保名下，便煽动异端者以基督徒的名义来反抗基督的教义，仿佛他们可以不加任何纠正就被容忍留在天主的城里，就如混乱之城容忍那些持有相反意见的哲学家一样。因此，在基督的教会里，凡是品味病态和堕落事物的人，受到纠正以品味健全和正确事物时，却顽固抵抗，不肯修改他们那有害而致命的教义，反而坚持捍卫它们，他们就成了异端者，并离开教会，应被视为服务于教会纪律的敌人。因为即使如此，他们以自己的邪恶也为基督真正的公教会肢体带来益处，因为天主连恶人也善加利用，一切事都使爱祂的人获得益处。\BibleRef{罗 8:28} 因为教会的一切仇敌，无论他们是被何种错误蒙蔽或被何种恶意败坏，如果他们得到能力在肉体上折磨教会，她们就操练教会的忍耐；如果他们只是以邪恶的思想反对，她们就操练教会的智慧；但同时，如果这些仇敌被爱，她们就操练教会的仁爱甚至慈善，无论教会是以劝导性的教义还是以严厉的纪律对待他们。因此，魔鬼，这恶城的首领，当他煽动自己的器皿反对那在世上旅居的天主之城时，被允许对她毫无伤害。因为毫无疑问，天主的眷顾为她既通过顺境带来安慰，使她不被逆境压垮，又通过逆境带来考验，使她不被顺境腐蚀；这样两者互相调和，正如我们在圣咏中听到的那声音，它毫无其他缘由：\Quote{依我心中的诸多忧苦，你的安慰使我欢乐。} 因此也有宗徒的那句话：\Quote{在希望中喜乐，在困苦中忍耐。} \BibleRef{罗 12:12}\par

因为不应认为那位同一导师所说的任何时候会落空：\Quote{凡愿意在基督内虔敬生活的，都必受迫害。} \BibleRef{弟后 3:12} 因为即使当外面的人并不发怒，因此似乎有了并确实有了平安，这给软弱者带来极大的安慰，但并非没有，甚至有许多里面的人以他们放荡的品行折磨那些虔敬生活者的心，因为通过他们，基督徒和公教会的名声被亵渎；而那名声对于那些愿意在基督内虔敬生活的人越是宝贵，他们就越因着那些在里面有地位的恶人而使它变得不如虔敬心灵所期望的那样被爱而忧伤。异端者们自己，既然他们被认为拥有基督徒的名号、圣事、圣经和信仰宣认，也给虔敬者的心中带来巨大的忧伤，因为许多愿意成为基督徒的人因他们的分歧而被迫犹豫，许多恶言者也从他们那里找到亵渎基督徒名号的素材，因为他们无论如何也被称为基督徒。由于这些以及类似的堕落品行和人的错误，那些愿意在基督内虔敬生活的人遭受迫害，即使没有人骚扰或折磨他们的身体；因为他们所受的这迫害，不在他们的身体里，而在他们的心里。因此有那句话：\Quote{依我心中的诸多忧苦}，因为他没说在我身体里。然而，另一方面，他们中没有一个人能灭亡，因为那不变的天主应许被思念。并且因为宗徒说：\Quote{主认识属于祂的人；\BibleRef{弟后 2:19} 因为祂所预知的人，也预定他们肖似祂儿子的肖像，\BibleRef{罗 8:29}} 他们中没有一人能灭亡；因此在那一篇圣咏中接着有：\Quote{你的安慰使我欢乐。} 但那在虔敬者心中因坏基督徒或假基督徒的品行而被逼迫而生的忧伤，对受痛苦者是有益的，因为它出自那爱德，他们不愿这些人灭亡，也不愿他们阻碍他人的得救。最后，从对他们的惩罚中产生了极大的安慰，这些安慰以同样多的痛苦充满虔敬者的灵魂，正如他们曾为这些人的丧亡而焦虑。因此，在这世界，在这邪恶的日子，不仅从基督及祂的宗徒的肉身临在之时，甚至从亚伯尔，他因是义人而被他的邪恶弟兄首先杀害，\BibleRef{若一 3:12} 从此直到这世界的终结，教会一直在世界的迫害和天主的安慰中前行、朝圣。\par

% structure: p0135 source=h2 target=subsection reason=relative-heading-rank; base=h2

\subsection{第五十二章\newline{}——我们是否应相信某些人的想法，即正如已经完成的十次迫害，除了第十一次——它必须在反基督者实际的时代发生——之外，再也没有别的了。}

诚然，我不认为某些人已经或可能认为的轻率说法或信仰是正确的，即直到假基督的时代，基督的教会除了已经遭受的迫害之外，将不再遭受任何迫害——也就是说，\Emph{十}次——而第十一次，即最后一次，将由假基督施加。他们算定第一次是尼禄所为，第二次是多米提安，第三次是图拉真，第四次是安东尼，第五次是塞维鲁，第六次是马克西米努斯，第七次是德西乌斯，第八次是瓦勒里安，第九次是奥勒良，第十次是戴克里先和马克西米安。因为正如天主子民开始出离之前在埃及有十灾一样，他们认为这应当被理解为表明假基督的最后迫害必须如同第十一灾，在那一灾中，埃及人敌意追赶希伯来人，而天主子民从干地经过红海时，埃及人却灭亡于海中。然而，我不认为埃及所发生的事是预言性地象征了迫害，无论那些持此观点的人似乎多么巧妙而机智地将二者作了详细比较——他们不是凭着先知之神，而是凭着人类心智的猜测，这种猜测有时触及真理，有时却被欺骗。但那些持此观点的人对主自身被钉十字架的那次迫害又有什么可说呢？他们将把它算在第几次？如果他们认为计算时应将这一次排除，仿佛只应计算那些属于身体的，而不应包括头自身被攻击并被杀的那一次，他们对基督升天之后在耶路撒冷发生的那次迫害又作何解释呢？那时，真福斯德望被石头砸死；若望的兄弟雅各伯被刀剑杀死；宗徒伯多禄被囚禁待杀，却被天使释放；弟兄们被驱逐并分散离开耶路撒冷；后来成为宗徒保禄的扫禄摧残了教会；而他自己在传播他曾经迫害过的信仰的喜讯时，无论是在犹太人还是在外邦人中，都遭受了他曾施加的苦难，他在各处最热切地宣讲基督。那么，他们为何认为适合从尼禄开始呢？当时教会正在成长，已经到了尼禄的时代，其间充满了最残酷的迫害；关于这些，说来话长，难以尽述。但如果他们认为只应计算由君王施行的迫害，那么黑落德王在主升天之后也曾施加了一次最严重的迫害。他们对尤里安又作何说明呢？他没有把他们算在十次之中。他不是迫害了教会，禁止基督徒教授或学习自由科目吗？在他治下，年长的瓦楞提尼亚努斯——即他之后的第三位皇帝——作为基督信仰的宣信者挺身而出，并被免去了军职。我对他（尤里安）在安提约基雅所作的事不拟多言，只提一点：他对一位极其忠实而坚毅的青年的自由和喜乐感到惊讶。当许多人被捉拿受刑时，这位青年受了一整天的酷刑，并在刑具下歌唱，直到皇帝担心他会在持续的酷刑下屈服，最终使他蒙羞，这使皇帝畏惧和害怕自己会被其余的人更加不光彩地羞辱。最后，在我们记忆所及，前述瓦楞提尼亚努斯的兄弟，亚略派瓦伦斯，不是在整个东方以大规模迫害摧残了天主教会吗？但不考虑这一点是多么不合理啊：教会，在全世界结果并成长，可能在有些国家遭受君王的迫害，即使在别国没有遭受迫害！然而，或许哥特人的国王在哥提亚本地以惊人的残酷迫害基督徒，而那里只有天主教徒，其中许多人被殉道者的荣冠所加冕——我们曾从一些当时在那里做孩童的弟兄那里听说，他们毫不迟疑地回忆起曾目睹这些事——难道这不算迫害吗？最近在波斯又发生了什么？对基督徒的迫害不是如此炽热（即便现在已有所缓和）吗？以至于有些逃难者甚至逃到了罗马城镇。当我想到这些以及类似的事情时，在我看来，教会将要经受的迫害次数无法明确确定。但另一方面，断言除了那最后一次（所有基督徒对此都毫不怀疑）之外，还会有一些由君王施行的迫害，同样是鲁莽的。因此，我们对此不作定论，既不支持也不反驳这一问题的任何一方，只是阻止人们冒昧地断言其中任何一方。\par

% structure: p0137 source=h2 target=subsection reason=relative-heading-rank; base=h2

\subsection{第五十三章 论最后迫害的隐秘时期}

诚然，耶稣自己要藉祂的来临消灭假基督所要施加的最后迫害。因为经上记载：\BibleQuote{祂要用口中的气杀死他，并以祂来临的光辉消灭他。}人们习惯于问：那将在何时？但这很不合理。因为如果知道这事对我们有益，谁能比天主自己、那位老师更好地告诉我们呢？那时门徒们曾问祂。他们与祂在一起时并没有沉默，而是询问祂说：\BibleQuote{主，祢要在这时候将王国呈现给以色列吗？或者是在什么时候？} \BibleRef{宗 1:6} 但祂说：\BibleQuote{你们不需要知道时间，父已将这些放在祂自己的权柄内。}当他们得到这个回答时，他们根本未曾问祂关于时辰、日子或年岁，而是关于时期。因此，我们试图精确计算这个世界可能剩余的年岁是徒然的，因为我们从真理的口中听到，我们不需要知道这事。然而，有些人说，从主升天到祂最后的来临，可能要经过四百年，有些人说五百年，另一些人说一千年。但指出每个人如何支持自己的观点将花费太多时间，且没有必要；因为他们确实使用人类的猜测，并从正典圣经的权威中提不出任何确定的东西。但关于这个问题，祂搁置了计算者的数字，并命令静默，祂说：\BibleQuote{你们不需要知道时间，父已将这些放在祂自己的权柄内。}\par

然而，因为这句经文在福音中，所以毫不奇怪，崇拜众多假神的人仍然没有被制止假装：通过他们奉为神明的魔鬼的回应，已经确定了基督宗教将持续多久。因为当他们看到它不能被如此众多和重大的迫害所消灭，反而从中获得了奇妙的增长，他们就编造了不知什么希腊诗句，仿佛由某个神谕涌出，给某个求问者一样；在这些诗句中，他们确实使基督对这件亵渎的罪行无罪，但补充说，\Person{伯多禄}通过魔法使得\Person{基督}的名被崇拜三百六十五年，并且，在那个年数完成后，立刻结束。哦，有学识之人的心！哦，有学识的才智，你们适于相信\Emph{关于}\Person{基督}的这样的事，却不愿相信\Emph{在}\Person{基督}内的这样的事，即：祂的门徒\Person{伯多禄}没有从祂那里学到魔法，然而，虽然祂是无辜的，祂的门徒却是个巫师，并且选择通过他的魔法、他的巨大劳苦和危险，甚至最终通过流血，使祂的名而非自己的名被崇拜！如果巫师\Person{伯多禄}使世界如此爱\Person{基督}，那么无辜的\Person{基督}做了什么使\Person{伯多禄}如此爱祂？那么让他们自己回答吧，并且，如果他们能做到，让他们明白：世界为了永生，被那同样的至高恩宠所造去爱\Person{基督}，这恩宠也使\Person{伯多禄}为了从祂那里领受永生而爱祂，甚至甘愿为祂忍受暂时的死亡。那么，这些是什么样的神呢？他们能够预言这些事，却不能避免它们，就这样屈服于一个巫师和邪恶魔法师（据说他杀了一个一岁的男孩，把他撕成碎片，用邪恶的礼仪埋葬了他），以致他们允许敌视自己的教派在如此长的时间内获得力量，并且通过忍受而非抵抗，克服了如此众多大迫害的可怕残酷，并促成了他们自己的塑像、庙宇、礼仪和神谕的覆灭？最后，是什么神——当然不是我们的神，而是他们自己的神——被如此巨大的邪恶所诱惑或强迫，去行这些事？因为那些诗句说，\Person{伯多禄}捆绑了，不是任何魔鬼，而是一个神去行这些事。那些没有\Person{基督}的人，就有这样的神。\par

% structure: p0140 source=h2 target=subsection reason=relative-heading-rank; base=h2

\subsection{第五十四章——论异教徒极其愚蠢的谎言，他们假装基督宗教不会超过三百六十五年就结束。}

我本可收集这些以及许多类似的论据，倘若那一年尚未过去，虚假的占卜曾预言，而被欺骗的虚荣竟相信了那一年。但既然几年前，自从基督之名被敬拜的确立，藉着祂在肉身的临在，以及藉着宗徒，已经满了三百六十五年，我们还需要寻求什么别的证据来驳斥那谬误呢？因为，不将此时期的起点置于基督的诞生，因为祂作为婴儿和孩童还没有门徒，然而，当祂开始有门徒时，基督的教义和宗教无疑是通过祂肉身显现而为世人所知，即祂在\Place{约旦河}中，藉着\Person{若翰}的职务受了洗之后。因为这缘故，关于祂早有预言说：\Quote{祂将统治从这海到那海，从河流直到地极。}但既然在祂受苦并从死复活之前，信德尚未对所有人确定，而是在基督复活时确定的（因为宗徒\Person{保禄}这样对雅典人说：\BibleQuote{但现在祂宣告给众人，叫他们各处都要悔改，因为祂已定了一日，要藉着祂所设立的人按公义审判天下，叫他从死人中复活，给众人确定了信德。}\BibleRef{宗 17:30}），那么，在处理这个问题时，我们最好从那一点开始，尤其是因为圣神当时被赐下，正如祂应在基督复活后在那座城里被赐下一样，第二部法律，即\WorkTitle{新约}，应从那里开始。因为第一部，称为\WorkTitle{旧约}，是从\Place{西乃山}藉着\Person{梅瑟}赐下的。但关于由基督所要赐下的这部，曾有预言说：\BibleQuote{法律将出自\Place{熙雍}，上主的话将出自\Place{耶路撒冷}；}\BibleRef{依 2:3}由此祂自己说，因祂的名而悔改，应传给万民，但要从\Place{耶路撒冷}开始。\BibleRef{路 24:47}因此，在那里，这名的敬拜开始兴起，即相信耶稣，祂曾死而复活。在那里，这信德以如此崇高的开端燃烧起来，以致数千人带着奇妙的热情归向基督之名，变卖财物分给穷人，从而藉着神圣的决心和最热烈的爱德，自愿处于贫穷，并准备了自身，在那些狂怒并渴求他们血的犹太人中间，为真理争战至死，不是靠武力，而是靠更强大的忍耐。如果这并非藉着魔法成就，他们为何犹豫相信，藉着成就此事同一天主之能，另一事也能在全世界成就？但假设\Person{伯多禄}行那法术，以致在\Place{耶路撒冷}如此众多的人被点燃来敬拜基督之名，这些人曾抓住祂并钉祂在十字架上，或在祂被钉时辱骂祂，我们仍必须探究那三百六十五年应从那年算起满期。基督死于\Person{Gemini}担任执政官时，在四月朔日前第八日。祂第三日复活，正如宗徒们以自己的感官所见为证。随后四十天后，祂升天。十天后，即祂复活后第五十日，祂派遣了圣神；于是宗徒们宣讲祂时，三千人相信了。因此，那时，那名的敬拜兴起了，按我们相信的，并按照真实情况，是靠圣神的效能，但按不虔的虚荣所虚构或认为的，是靠\Person{伯多禄}的魔法。不久后，又在一个奇异的征兆发生时，当因\Person{伯多禄}的话，一个乞丐，从母胎中跛脚，被他人抬来放在圣殿门口，在那里乞讨，在耶稣基督之名内痊愈，跳起来，五千人相信了，从此教会因信徒的不断增加而成长。由此我们推断那年开始的那一天，即圣神被派遣的那一天，也就是五月望日。并且，数算执政官，在\Person{Honorius}和\Person{Eutychianus}担任执政官时，同一望日满了三百六十五年。现在，在下一年，在\Person{Mallius Theodorus}担任执政官时，按照那魔鬼的神谕或人的虚构，基督教已然不复存在，便没有必要询问或许在世界其他部分发生了什么事。但是，据我们所知，在\Place{阿非利加}最著名卓越的城市\Place{迦太基}，\Person{Honorius}皇帝的官员\Person{Gaudentius}和\Person{Jovius}，在四月朔日前第十四日，推倒了庙宇，打碎了假神的偶像。从那时到现在，近三十年间，谁看不出对基督之名的敬拜增加了多少？尤其许多曾被那占卜所阻，认为它将应验而未信的人，在见到那同样年数满期后，发现它空洞可笑，便成了基督徒。因此，我们这些被称为基督徒，并确实是基督徒的人，并不信\Person{伯多禄}，而是信\Person{伯多禄}所信的那位——我们因\Person{伯多禄}关于基督的讲道而受建树，并非因他的咒语而中毒；我们未被他的法术欺骗，而是因他的善行而蒙助。基督自己，是\Person{伯多禄}在引领永生之道上的师傅，也是我们的师傅。\par

但现在，我们终于应当结束这一卷了，因为在之前我们已经足够详细地讨论并展示了天上之城与地上之城从起初直到终结的混合在一起的尘世历程。其中，地上之城从任何其他地方，甚至从人类中，为自己制造了她所愿意的假神，以祭献事奉它们；但天上之城（她在此世是旅居者）并不制造假神，而是由真天主所创造，她自己必须成为真天主的真实祭献。然而，两者同样都有享现世的美物，或遭受现世的灾祸，但怀着不同的信德、不同的望德、不同的爱德，直到她们被最后审判分开，各自承受她那没有终结的结局。关于这两座城的结局，我们接下来必须加以讨论。\par

\SourceRecord{Translated by Marcus Dods.；Nicene and Post-Nicene Fathers, First Series；Vol. 2.；Edited by Philip Schaff.；Buffalo, NY: Christian Literature Publishing Co.,；1887.；<http://www.newadvent.org/fathers/120118.htm>.}

% structure: p0001 source=h1 target=section reason=page-title

\section{天主之城（第十九卷）}

在本卷中，论述了属地之城与属天之城这两个城的终结。\Person{奥古斯丁}回顾了哲学家们关于至高之善的意见，以及他们在这尘世生命中为自己制造幸福的徒劳努力；在他驳斥这些意见的同时，他借机指出属天之城或基督的子民如今和将来所拥有的平安与幸福是何等性质。\par

% structure: p0003 source=h2 target=subsection reason=relative-heading-rank; base=h2

\subsection{第一章——\Person{瓦罗}指出，关于至高之善的各种意见可形成二百八十八种不同的哲学派别。}

既然我仍需讨论这两个城——属地之城与属天之城——的合适应得的结局，我必须首先，在这部著作允许的范围内，解释人们曾试图在这不幸的尘世生命中为自己制造幸福的各种推理，以便不仅从神圣权威，而且从可向不信者提出的理由中，显明哲学家们空洞的梦想与天主赐予我们的希望，以及祂将作为我们的福乐而赐予我们的那实在的满全，是何等不同。关于善与恶的终极，哲学家们已表达了大量各不相同的意见，他们热心探讨这一问题，以求尽可能发现什么使人幸福。因为我们之善的终极，是为了其自身而应欲求的，而其他事物则是为了它而被欲求；恶的终极，则是因其自身而被回避的，而其他事物则是为了它而被回避。因此，我们目前所谓的\Emph{善之终极}，并非指那使善消灭、不再存在的终结，而是指那使之完成、成为圆满的终结；而\Emph{恶之终极}，也并非指那使之消灭的终结，而是指那使之发展完全的终结。这两个终极，因而就是至高之善与至高之恶；如我所说，那些在这虚妄生命中自称追求智慧的人，曾竭力发现这些终极，并在今生获得至高之善、避免至高之恶。尽管他们以各种方式误入歧途，但自然洞察力尚未使他们偏离真理到如此地步，以致未将至高之善与至高之恶置于灵魂中、身体中，或两者兼有。基于这种对哲学派别的三分法，\Person{瓦罗} 在其著作 \WorkTitle{De Philosophia} 中，引出了如此大量的不同意见，以致通过细微详尽的分析，他毫不费力地数出了二百八十八个派别——并非这些派别实际存在过，而是可能存在的派别。\par

为了简要说明他的意思，我必须从他上述著作中的绪论开始：他说有四件事物，人几乎是出于本性、无需老师、无需任何教导、无需勤劳或所谓德行的生活技艺（而德行无疑是学得的）所渴望的：或是快乐，即身体感觉的一种愉悦激动；或是休息，即排除一切身体不适；或是二者兼有，\Person{伊壁鸠鲁}称之为同一个名称：快乐；或是自然的首要事物，它包括已提及的事物以及其他事物，或是身体方面的，如健康、安全、肢体完整，或是精神方面的，如人身上所有的大小精神天赋。现在这四件事——快乐、休息、二者结合、以及自然的首要事物——存在于我们之内，其方式是我们必须要么为了它们而渴望德行，要么为了德行而渴望它们，要么两者都因其自身而被渴望；因此，从这一区分产生了十二派别，因为每个项目由此被三重化。我将用一个例子来说明，一旦明了，就不难理解其他。那么，按照身体快乐是被置于德行之下、优先于德行、或与德行结合，就有三派。当快乐被选择为德行的侍从时，它就是被置于德行之下。因此，为祖国而活，并为其生育子女，是德行的责任，而这两者没有身体快乐是无法做到的；因为饮食有快乐，性交也有快乐。但当快乐优先于德行时，它就是因其自身而被渴望，德行仅仅为了快乐而被选择，且除了获得或保持身体快乐之外别无他用。而这确实使生活变得丑陋；因为当德行成为快乐的奴隶时，它就不再配称为德行。然而，这种可耻的歪曲却也找到了一些哲学家来庇护和辩护它。然后，当德行与快乐结合时，两者都不是为了对方而被渴望，而是都因其自身。因此，正如快乐按照被置于德行之下、优先于德行或与德行结合而产生三派，同样地，休息、快乐与休息结合、以及自然的首要事物，也各自产生三派。因为人们的意见分歧，这四样东西有时被置于德行之下，有时优先于德行，有时与德行结合，这样就产生了十二派别。但这个数目又因添加一个差别而加倍，即社会生活；因为任何依附于这些派别之一的人，要么仅为自己，要么也为一个同伴，他应希望同伴得到他自己所渴望的。于是，将有十二个派别认为这些意见之一应为自己而持守，另有十二个派别决定他们应遵循这种或那种哲学，不仅为了自己，也为了他人，他们渴望他人获得如自己一般的善。这二十四个派别再次加倍，变成四十八个，因为加入了源自新学园派的差别。因为每一个这样的二十四个派别都可以将其意见当作确定的来持有和辩护，如同斯多亚派辩护人的至善仅在于德行；或者它们可以被当作或然的而非确定的来持有，如同新学园派所做的那样。因此，有二十四个派别将其哲学当作确定真确的，另有二十四个派别将其意见当作或然的而非确定的。再者，由于每个依附于这些派别的人可以采纳要么犬儒学派要么其他哲学家的生活方式，这一区分将使数目加倍，从而变成九十六派别。最后，每一个这些派别可能被两种人所遵循：要么喜爱悠闲生活的人，例如那些出于选择或必要而致力于学术的人；要么喜爱忙碌生活的人，例如那些在从事哲学的同时还大量参与国务和公共事务的人；要么选择混合生活的人，仿效那些将时间分配给有学问的闲暇和必要工作的人。通过这些差别，派别的数目被三倍化，变成二百八十八个。\par

我依我所能，以最简洁清晰的方式，用我自己的话陈述了\Person{瓦罗}在其著作中表达的意见。然而，他如何驳斥其余所有学派，并选择其一——由\Person{柏拉图}创立、延续至该哲学学派第四任教师\Person{波勒莫}的旧学院派，认为其体系是确定的；他如何以此为基础将其与始于\Person{波勒莫}的继承者\Person{阿尔克西劳}的新学院派区分开来，后者认为一切皆不确定；以及他如何试图证明旧学院派既无错误亦无怀疑——所有这一切，我说，要详细论述未免过于冗长，然而我又不能完全置若罔闻。\Person{瓦罗}首先摒弃了所有那些导致学派数目增多的差异，其根据是这些差异并非关于至善的差异。他主张，在哲学中，一个学派之所以成立，仅因其在终极目的问题上持有与其他学派不同的见解。因为人从事哲学活动别无他因，唯求幸福；而使人幸福者，即是至善本身。换言之，至善是哲学活动的理由；因此，若不追寻通向至善的独特途径，便不能称之为哲学学派。因此，当有人问及智者是否将采纳社群生活，并如同关切自身至善一样关切并致力于朋友的至善，抑或相反，所做一切仅为自身利益时，这里并非关于至善的问题，而仅是关于是否应当让朋友参与分享之适当性问题：智者这样做是否不是为己，而是为友，并因其友之善如己之善而喜乐。同样，当有人问及哲学所涉一切事物是否应如新学院派所认为的那样视为不确定，抑或如其他哲学家所主张的那样确定时，这里的问题并非追求何种目的，而是我们是否应当相信该目的具有实质存在；或者更明确地说，追求至善者是否必须坚持其为真善，抑或仅认为其看起来为真，尽管可能为幻象——二者追求的是同一善。此外，基于犬儒学派的服饰与习俗的区分，亦不触及至善问题，而仅涉及追求那自认为真善者是否应如犬儒学派那样生活。事实上，有些人虽然追求不同的至善——有人选择快乐，有人选择德行——却采纳了使犬儒学派得名的生活方式。因此，任何将犬儒学派与其他哲学家区分开来的东西，都与构成幸福之善的选择与追求无关。因为若其有关，则相同的生活方式必导致追求相同的至善，而不同的生活方式则必导致追求不同的目的。\par

% structure: p0007 source=h2 target=subsection reason=relative-heading-rank; base=h2

\subsection{第二章——\Person{瓦罗}如何通过消除所有不构成学派、仅属次要问题的差异，得出我们必须择一的三种至善定义。}

关于三种生活，即勤勉的闲暇与追求真理的生活、轻松从事公务的生活、以及两者混合的生活，也可作同样说明。当问及其中哪一种应当被采纳时，这不涉及关于至善的争论，而是探究这三种生活中哪一种使人处于最佳位置，以发现并持守至善。因为一旦人发现这善，它便使人幸福；但博学的闲暇、公共事务或这两者的交替，并不必然构成幸福。事实上，许多人可能采纳这些生活方式中的一种或另一种，却未能获得使人幸福之物。因此，关于至善和至恶的问题（这正是区分哲学派别的依据）是一个问题；而那些关于社会生活、学园派的怀疑、犬儒派的服饰与饮食、三种生活模式（即行动、默观与混合）的问题，则是不同的问题，至善的问题并不进入其中。因此，正如\Person{马尔库斯·瓦罗}通过引入源于社会生活的这四个差别（即社会生活、新学园派、犬儒派、以及三重生活形式）将派别数量增至二百八十八（或任何他选择的更大数字），同样地，他将这些与至善无关因而不构成可恰当称为派别的差别除去，从而回到那十二个学派（这些学派关心探究那一使人幸福的善是什么），并指出其中一个是真的，其余是假的。换言之，他舍弃基于三重生活模式的区分，从而使总数减少三分之二，将派别降至九十六个。接着，去掉犬儒派的特性，数量减半，成为四十八个。再除掉因新学园派的犹疑而产生的区分，数量再次减半，降至二十四个。以类似方式处理因考虑社会生活而引入的多样性后，只剩下十二个（这一差别曾使其加倍至二十四个）。对于这十二个，没有任何理由可以解释为何它们不应被称为派别。因为其中唯一探究的就是至善和终极之恶——也就是说，关于至善，因为此善一旦被发现，与之相对的恶也就被发现了。如今，为构成这十二个派别，他将四者——即快乐、安息、快乐与安息之结合、以及\Person{瓦罗}称之为\OriginalTerm{原初之物}{primigenia}的自然的原始对象——乘以三。因为这四者有时从属于德行，因而似乎不是为其本身而被欲求，而是为了德行之故；有时则被优先于德行，因而德行似乎不是因其自身而必要，而是为了获得这些东西；有时则与德行结合，因而这些东西和德行都为自身而被欲求——我们必须将四乘以三，从而得到十二个派别。但从那四者中，\Person{瓦罗}排除了三个——快乐、安息、快乐与安息之结合——并非因为他认为这些不值得赋予它们的地位，而是因为它们已被包含在自然的原始对象之中。而且，无论如何，有何必要从这两个目的（快乐和安息）中做出三重划分——先分别取之，再结合取之——既然它们以及许多其他事物都已包含在自然的原始对象之中？剩下的三个派别中，哪一个应当被选取？这是\Person{瓦罗}再三思考的问题。因为无论选取这三个中的一个还是其他，理性都禁止有多个为真。这一点我们将在以后看到；但与此同时，让我们尽可能简明清晰地解释\Person{瓦罗}如何从这三个中做出选择，即从那些分别坚持以下看法的派别：自然的原始对象应为德行之故而欲求、德行应为这些对象之故而欲求、以及德行与这些对象应各自为其本身之故而欲求。\par

% structure: p0009 source=h2 target=subsection reason=relative-heading-rank; base=h2

\subsection{第三章——据\Person{瓦罗}（他追随\Person{安提约古}与旧学园派），关于至善的三种主要见解中哪一种应被偏好。}

这些三者中哪一个为真且当采纳，他试图以下列方式表明。既然哲学所寻求的至善，不是一棵树、一头野兽或一个神的至善，而是人的至善，他认为首先必须界定人。他意见是人性有两部分，即身体与灵魂，并且毫不怀疑在这两者中灵魂是更优越、远为更可贵的部分。但是否只有灵魂才是人，以至于身体与它的关系如同马与骑马者的关系，这一点他认为有待确定。骑马者并非马加人，而仅仅是个人，却因他与马有某种关系而被称为骑马者。再者，是否只有身体才是人，而与灵魂的关系如同杯与饮的关系？因为被称为杯的不是杯与其所盛的饮，而只是杯，但它之所以如此称呼，是因为被造来盛饮。抑或最后，既非单独的灵魂，也非单独的身体，而是两者结合才为人，身体与灵魂各为部分，而完整的人是两者结合，如同我们称两匹并辔的马为一对，其中左边和右边的马各为一部分，但我们不把任何一匹，无论它与另一匹如何相连，称为一对，而只有两匹一起才成一对？在这三种可能中，\Person{瓦罗}选择了第三种，即人既非单独的身体，也非单独的灵魂，而是两者结合。因此，幸福之所在至善，是由两种善——身体之善与灵性之善——共同构成的。从而他认为，自然原初物应为自身而追求，而德行——即生活艺术，且可藉教导传授——乃是灵性之善中最卓越的。于是，这种德行，或称调节生活的艺术，当它接受了那些独立于它且先于任何教导而存在的自然原初物时，它自身也为自身而寻求它们全体，并寻求它自身；它使用它们，也使用自身，以便从所有这一切中获得利益与享受，根据它们本身的大小而或多或少；并且当它享受它们全体时，它在必要时轻视较小的，以便获得或保持较大的。如今，在所有善中，无论灵性或身体，没有一样能与德行相比。因为德行善用自身以及构成人之幸福的一切其他善物；若德行缺席，则无论一个人拥有多少善物，那些善物对他都不利，因而它们不应被称为善物，因为它们属于一个因滥用而使它们无益的人。因此，当一个人的生命享受德行以及其他这些为德行所不可或缺的灵性善物和身体善物时，这生命被称为幸福的。若它享受某些或许多并非德行所必需的其他善物，则称为更幸福的；而若它不缺任何属于身体与灵魂的善物，则称为最幸福的。因为生命与德行并非同一回事；不是每一个生命，而是明智地调节的生命，才是德行；然而，没有德行可以有某种生命，但绝不能没有生命而有德行。这一点我可应用于记忆与理性，以及诸如此类的精神官能；因为这些先于教导而存在，没有它们便不可能有任何教导，因而也不可能有什么德行，因为德行是学得的。但身体上的优势，如脚快、美貌或力量，并非德行所必需，德行也非它们所必需，然而它们仍是善物；并且，按照我们的哲学家，甚至这些优势也为德行自身所渴求，并以适宜的方式为德行所使用和享受。\par

他们说这幸福生活也是社会性的，它爱朋友的益处如同自己的益处，并为他们盼望它为自己所愿望的，无论这些朋友是生活在同一家庭，如妻子、儿女、仆役；或是在一个人的家乡所在地，如同城公民；或是在普世范围，如因共同人类兄弟情谊而结合的诸民族；或是在宇宙本身，包容天地万物，如同他们称为神祇的那些，为智者提供为朋友，而我们更习惯称之为天使。此外，他们说，关于至善与至恶，不容怀疑，因此他们在这方面区别于新学园，并且他们并不关心一位哲学家是以犬儒学派的服饰和生活方式还是以其他方式追求他们认为真实的目的。最后，关于三种生活方式——默观、行动与混合的生活，他们声明赞同第三种。这些是老学园的见解和教训，\Person{瓦罗}以\Person{安提约古}为权威而断言，\Person{安提约古}是\Person{西塞罗}和他自己的老师，尽管\Person{西塞罗}把他描写得较常符合斯多葛学派而非老学园。但这于我们何干？我们应当按其本身的功过来判断事情，而非准确理解不同的人对此想了些什么。\par

% structure: p0012 source=h2 target=subsection reason=relative-heading-rank; base=h2

\subsection{第四章——基督徒所信的至善与至恶，与哲学家之主张相对，哲学家以为至善在于自身。}

若有人问我们，天主之城对这些问题有何说法，首先，对其关于至善与至恶的意见，它将回答说：永生是至善，永死是至恶，为获得前者而避免后者，我们必须正当地生活。因此经上记载：\Quote{义人因信德而生活}\BibleRef{哈 2:4}，因为我们尚未看见我们的善，所以必须因信德而生活；我们也没有能力凭自己正当地生活，惟有那赐给我们信德以相信祂帮助的主，在我们相信并祈祷时帮助我们，才能如此。至于那些认为至善与至恶存在于今世，并将其置于灵魂或肉体，或二者之中，或更明确地说，置于快乐或德行，或二者之中；置于安宁或德行，或二者之中；置于快乐与安宁，或德行，或三者之结合；置于自然之初因，或德行，或二者之中——所有这些人都以惊人的肤浅，试图在今世和自身中找到幸福。真理已对这些想法倾注了轻蔑，先知说：\Quote{上主知道人的思念}（或如宗徒保禄引用的经文，\Quote{上主知道\Emph{智者}的思念}）\Quote{尽都是虚幻}。\par

因为有什么雄辩的洪流足以详述今生的悲惨呢？西塞罗在悼念其女的\WorkTitle{慰藉}中，竭尽其才能哀叹；但连他的才能在这里也显得多么不足啊！因为何时、何地、如何在这今生中，这些自然之初因能如此被拥有，以致不受意外之灾的侵袭呢？智者的肉体能免除任何驱散快乐的痛苦，任何消除安宁的烦扰吗？肢体的截断或腐烂毁其完整，畸形损其美丽，虚弱损其健康，倦怠损其活力，困倦或迟钝损其行动——这些之中哪一个不会侵袭智者的肉身呢？身体优雅合宜的姿态和运动被视为自然的基本福分；但若某种疾病使肢体颤抖呢？若一个人脊椎弯曲得如此厉害，以致手能触地，像四足动物一样爬行呢？这不就破坏了身体的一切美和优雅，无论静止或运动？我该说什么关于灵魂的基本福分——感官和理智——其中一者被赋予感知真理，另一者被赋予理解真理？然而，当人变得又聋又瞎时，还剩什么样的感官呢？理智何在，当疾病使人谵妄时？我们几乎——甚至完全——无法忍住眼泪，当我们想到或看到这些疯狂者的言行，并考虑他们当前的行为与他们清醒的判断和平常的举止有多么不同，甚至相反。我该说什么关于那些受魔鬼附身的人呢？在他们自己的理智被恶灵按其意志使用他们的身体和灵魂时，被隐藏和埋葬了？谁能完全确定，在今生中不会发生这样的事于智者身上？至于对真理的感知，即使在身体内，我们又能对此抱什么希望呢？正如我们在真实的智慧书中读到的：\Quote{可朽坏的身体重压着灵魂，地上的帐幕压迫着多思虑的心灵？}\BibleRef{智 9:15}而热忱或行动的欲望，如果这是希腊文\OriginalTerm{冲动}{\GreekText{ὁρμη}}的正确含义，也被列为自然的基本优点；然而，难道不是它产生了那些可怜的精神错乱者的动作，以及那些我们见了战栗的行为，当感官被欺骗、理性被搅乱时？\par

最后，德行本身，虽非自然之初因，而是随学习而来，虽在人类美善中占有最高地位，它的职责除了与恶习——不是我们外在的，而是内在的；不是别人的，而是我们自己的——不断作战，还能是什么？这场战争特别由那希腊人称为\OriginalTerm{节制}{\GreekText{σωφροσυνη}}、我们称为节制的德行所进行，它约束肉欲，阻止它们赢得精神对邪恶行为的同意。因为我们切勿以为我们里面没有恶习，正如宗徒所说：\Quote{肉情相反圣神；}\BibleRef{迦 5:17}因为这恶习有相反的德行，正如同一作者所说：\Quote{圣神相反肉情。}\Quote{因为这两者，}他说，\Quote{彼此相反，致使你们不能做你们所愿意的事。}但我们寻求达到至善时，我们希望做什么呢？除非是肉情不再相反圣神，我们里面没有恶习以致圣神与之相反。既然我们在今生无法达到这一点，无论多么热切渴望，至少让我们藉天主的助佑完成这一点：保守灵魂不向相反它的肉情屈服和投降，拒绝同意犯下罪恶。因此，当我们还处于这场内战之中时，切勿幻想我们已通过胜利找到了所追求的福乐。有谁如此智慧，以致与自己的恶习毫无冲突呢？\par

我要如何谈论那被称为\OriginalTerm{智德}{prudence}的美德？难道它所有的警惕不都用于辨别善恶，以免我们于所欲所避上犯错吗？因此它本身便证明我们正处于邪恶之中，或邪恶在我们之内；因为它教导我们：同意犯罪是邪恶，拒绝同意是良善。然而，这邪恶——智德教导我们、\OriginalTerm{节德}{temperance}使我们能够不同意的邪恶——却既非智德亦非节德能从今生移除的。\OriginalTerm{义德}{justice}——其职责是按各人应得而给予，由此在人自身内产生一种合乎自然的正当秩序，使灵魂服从天主，肉体服从灵魂，进而灵魂与肉体都服从天主——这美德岂不证明它仍在努力达至其目标，而非安息于已完成的工作吗？因为灵魂越少思念天主，便越少服从天主；肉体越强烈地悖逆精神，便越少服从精神。因此，只要我们还被这软弱、这瘟疫、这疾病所困扰，我们怎敢说我们是安全的？若不安全，我们又怎能已享有终焉的\OriginalTerm{真福}{beatitude}？那名为\OriginalTerm{勇德}{fortitude}的美德，更是人生苦难最明显的证据，因为它被迫忍耐的正是这些苦难。而且，这一点即便与最成熟的智慧共存，也仍然成立。我无法理解\OriginalTerm{斯多葛派哲学家}{Stoic philosophers}怎敢断言这些并非苦难，尽管他们同时允许智者——若苦难变得如此沉重，使他不愿或不当忍受时——自杀以离开今生。但这些人愚蠢的骄傲竟幻想\OriginalTerm{至善}{supreme good}可在今生寻得，且能凭自身资源而获得幸福，以至于他们的智者——或至少他们虚构的如此之人——总是幸福的，即使他变得盲目、聋哑、残缺、遭受剧痛，或承受任何可能迫他自尽的灾难；他们竟不羞于将那充满这些邪恶的生命称为幸福。哦，幸福的生命，竟求助于死亡来终结它？若它是幸福的，就让智者留在其中吧；但若这些苦难将他驱逐出去，它又怎能算是幸福？或者，他们怎能说那些战胜了勇德、迫使它不只屈服而且如此狂乱——以致它一面高呼生命幸福，一面又劝人放弃它——的邪恶不是邪恶？因为，谁有如此盲目，竟看不出若生命幸福，便无人逃离它？若他们声称我们应因环绕的软弱而逃离生命，为何不低下他们的骄矜，承认它是悲惨的呢？请问，促使\Person{加图}自杀的是勇德还是软弱？因为他若非太软弱而不能忍受\Person{恺撒}的胜利，便不会这样做。那么，他的勇德何在？它已屈服，已败退，已完全被征服，以致它抛弃、离弃、逃离这幸福的生命。或者，生命已不再幸福？那么，它就是悲惨的。那么，这些使生命悲惨、成为可逃避之事的，又怎能不是邪恶呢？\par

因此，那些承认这些是恶的人，如同逍遥学派和老学园派（即瓦罗所拥护的派别）那样，提出了一种更易于理解的学说；但他们也犯了一个惊人的错误，因为他们主张，被这些恶所包围的生活就是幸福的生活，即使这些恶如此严重，以至于忍受它们的人应当自杀以逃避它们。\Quote{身体的痛苦和折磨，}瓦罗说，\Quote{是恶，而且其严重程度越甚，恶也越甚；为了逃避它们，你必须离开此生。}请问：是什么生活呢？他说，就是被这样的恶所压迫的生活。那么，难道它正是因为这些恶而幸福吗？而你说我们必须因这些恶而离开此生？或者你称它为幸福，是因为你可以自由地通过死亡逃避这些恶？那么，如果因天主的某种隐秘审判，你被牢牢束缚，不许死亡，也不许没有这些恶而活着，又会怎样呢？至少在那时，你会说这样的生活是悲惨的。它很快就被放弃了，毫无疑问，但这并不使它不悲惨；因为如果它是永恒的，你自己也会宣称它是悲惨的。因此，它的短暂并不能清除其悲惨；也不应当因为它是一段短暂的悲惨而被称为幸福。的确，这些恶有一种强大的力量，迫使一个人——甚至按照他们的说法，一个智者——为了逃避它们而不再做人，尽管他们说（而且说得对），人爱护自己，因此自然而然地逃避死亡，应当成为自己的朋友，渴望并强烈地追求作为一个活物继续存在，并在灵魂与身体的结合中存续，这可以说是自然的首要和最强的要求。这些恶有一种强大的力量，能够克服这种使人们千方百计、竭尽全力避免死亡的自然本能，并如此彻底地克服它，以至于原本避免的东西被渴望、被追求，如果无法以其他方式获得，就由人自己施加于自身。这些恶有一种强大的力量，使刚毅成为杀人犯——如果那被称为刚毅的东西被这些恶如此彻底地征服，以至于它不仅不能通过忍耐来保全它所承担的管理和保卫的人，反而被迫杀死他。我承认，智者应当忍耐地承受死亡，但那是他人所施加的死亡。那么，如果如这些人所主张的，他必须自己施加死亡，那就必须承认，迫使他这样做的灾祸不仅是恶，而且是无法忍受的恶。因此，那种或受偶然事件支配，或被如此重大而痛苦的恶所包围的生活，永远不可能被称为幸福，除非那些赋予它此名的人肯于向真理屈服，并在探究幸福生活时被有效的论证所征服，正如他们屈服于不幸并被压倒性的恶所征服而自杀一样，并且除非他们幻想至善可以在今生找到；因为今生的德行本身——这无疑是其最好、最有用的财产——在其有助于对抗危险、辛劳和苦难的暴力时，反而更加证明了它的悲惨。因为如果这些是真正的德性——而且它们只能存在于拥有真正虔诚的人身上——它们并不会宣称能够使拥有它们的人免除一切悲惨；真正的德性不会撒这样的谎，而是宣称借着对来世的希望，这被今世许多大恶所悲惨地缠绕的生活，既是幸福的，也是安全的。因为如果还不安全，它怎能幸福呢？因此，宗徒保禄，不是谈论没有明智、节欲、刚毅和正义的人，而是谈论那些生活受真正虔诚规范、因而其德性也是真正的人，他说：\BibleQuote{因为我们得救，是在乎盼望；只是所见的盼望不是盼望，谁还盼望他所见的呢？但我们若盼望那所不见的，就必须忍耐等候。}\BibleRef{罗 8:24}因此，我们如何因盼望而得救，我们也如何因盼望而成为幸福的。正如我们目前尚未拥有现在，而是盼望未来的救恩，我们的幸福也是如此，而这\Quote{要忍耐等候}；因为我们被恶所包围，应当忍耐地承受，直到我们来到那不可言喻的纯善的享受中；因为那时将不再有任何需要忍受的东西。救恩，如同在来世将要有的那样，本身将是我们最终的幸福。这些哲学家拒绝相信这种幸福，因为他们看不见它，并试图为自己编织一种基于德性的今世幸福，而这种德性既是骗人的，也是骄傲的。\par

% structure: p0018 source=h2 target=subsection reason=relative-heading-rank; base=h2

\subsection{第五章  论社会生活在令人极为向往的同时，却常被许多痛苦所扰乱。}

我们更加完全地赞同他们所认为的，智者的生活必须是社会性的。因为天主之城（关于它，我们正在这部著作中写不少于第十九卷）如果没有圣徒们的社会性生活，又怎能开始、发展并达到它应有的归宿呢？但是，谁能数尽人类社会中充满的、在这个可朽坏状态下的一切巨大苦难？谁能衡量它们？听一听他们的一个喜剧作家如何让他的一个角色表达在这方面所有人的共同情感：\Quote{我结婚了；这是一个痛苦。孩子出生了；它们是额外的烦恼。} 关于爱情的痛苦，特伦提乌斯也叙述——\Quote{轻视、猜疑、争吵、今日战争、明日和平？} 难道人类生活不是充满这样的事吗？它们甚至在光荣的友谊中不也常常发生吗？我们处处经历这些轻视、猜疑、争吵、战争，所有这些都是确定无疑的恶；而另一方面，和平却是一个不确定的善，因为我们不知道朋友的心，即使我们今天知道了，我们也不明白它明天会如何。谁应该比生活在同一家庭中的人更友好？又有谁能信赖这种友谊，既然秘密的背叛常常破坏它，并产生与友爱同样苦涩的敌意，或者通过最完美的伪装显得甜蜜？正因如此，西塞罗的话感动每个人的心，引起一声叹息：\Quote{没有比伪装成职责或亲属之名的陷阱更危险的了。因为你可以轻易地用防备来挫败公开的敌人；但隐藏的、内部的、家庭的危险不仅存在，而且在你预见和审视之前就压倒了你们。} 同样，神圣的言语也提到了这个，\BibleQuote{人的仇敌就是自己的家人，} \BibleRef{玛 10:36} ——这些话没有人能不痛苦地听到；因为即使一个人有足够的坚毅去平静地忍受，有足够的智慧去挫败假装朋友的恶意，但如果他自己是个好人，发现恶人的背信弃义时，他不能不感到极大的痛苦，无论他们是始终邪恶而只是假装善良，还是从更好的状态堕落到恶意。如果家庭这个逃避人生苦难的天然避难所本身都不安全，那么对于城市，它规模更大，充满了民事和刑事诉讼，并且从不免于令人不安和血腥的起义和内战的恐惧（即使不是实际爆发），我们又能说什么呢？\par

% structure: p0020 source=h2 target=subsection reason=relative-heading-rank; base=h2

\subsection{第六章——当真理被隐藏时人类判断的错误。}

关于人们在他人身上所作出的、在社团（无论它们享有怎样的外在和平）中所必需的判断，我该说什么呢？它们是令人忧伤和可悲的判断，因为审判者是些不能辨别受审者良心的人，因此常常被迫对无辜的证人施以酷刑，以查明关于他人罪行的真相。对被告本人施加的酷刑，我又该说什么呢？他被施以酷刑是为了查明他是否有罪，因此，尽管他是无辜的，他却为仍然存疑的罪行遭受了最确实的惩罚，不是因为证明他犯了罪，而是因为没有查清他没有犯罪。因此，审判者的无知常常使无辜的人遭受痛苦。而更令人难以忍受的是——这确实是一件值得哀悼、如果可能的话，用泪泉来浇灌的事——即当审判者讯问被告，以免无意中错杀无辜时，这种可悲无知的结果是，正是这个他为了如果不成立时免于定罪而施以酷刑的人，既被施以酷刑，又是无辜的，却被判了死刑。因为如果他服从了哲学给智者的教导，宁愿离开此生也不愿继续忍受这样的酷刑，他就会宣称自己犯了实际上没有犯的罪。当他被定罪处死后，审判者仍然不知道他处死的是无辜者还是有罪者，尽管他对被告施以酷刑正是为了不判处无辜者死刑；因此，他既为了查明无辜而施以酷刑，又未查明真相就处死了他。如果这样的黑暗笼罩着社会生活，一位明智的审判者还会坐上审判席吗？毫无疑问，他会。因为人类社会（他认为离开它是邪恶的）约束并迫使他履行此职责。而他认为以下不算邪恶：无辜的证人因他人被指控的罪行而受酷刑；被告被施以酷刑，以致常常被痛苦压倒，尽管无辜，却作出关于自己的虚假供认，并受惩罚；即使他们不被判死刑，也常常在酷刑中或因此死亡；有时，控告者（也许是出于通过将罪犯绳之以法来造福社会的愿望）因审判者的无知而被定罪，因为他们无法证明其指控的真实性（尽管是真实的），因为证人撒谎，而被告在酷刑下不为所动，不肯招供。这些众多而重大的恶行，他不认为是罪；因为明智的审判者做这些事，并非出于任何害人的意图，而是因为他的无知迫使他，并且人类社会要求他担任审判者。但是，虽然我们因此宣告审判者没有恶意，我们却不能不能不谴责人类生活的悲惨。而如果他因职务和无知的束缚而被迫对无辜者施以酷刑和惩罚，他既是无罪的，又是否幸福呢？诚然，如果他能认识到这些必要之事的悲惨，并回避自己卷入其中，那便是更深刻的体察和更细腻的情感；如果他有些敬虔，他会向天主呼求：\Quote{从我的需要中拯救我。}\par

% structure: p0022 source=h2 target=subsection reason=relative-heading-rank; base=h2

\subsection{第七章 语言的多样性，它如何阻碍人与人之间的交往；以及战争，甚至所谓正义战争，其悲惨境况。}

在家庭或城之后，便是世界，人类社会的第三个圈——第一个是家庭，第二个是城。世界越大，危险也越多，正如大海越大越危险。首先，语言的不同将人与人隔离开来。因为如果两个人彼此不懂对方的语言，他们相遇时并不一定需要分开，反而必须在一起，那么即使是不同种类的哑巴动物，也会比他们更容易交流，尽管他们都是人。因为当他们因语言不同而不能传达彼此的思想时，他们共同的人性对友谊毫无帮助；所以一个人与自己的狗交流比与一个外国人交流更容易。但是帝国之城曾尽力将它的枷锁和语言强加于被征服的民族，作为和平的纽带，因此翻译人员不仅不匮乏，而且不计其数。这是真的；但是，为了实现这种统一，发生了多少大战，多少屠杀和流血！虽然这些已成过去，但苦难的终结尚未到来。因为帝国之外从不缺少敌对民族，现在也不缺少，战争一直在进行；但是，即使没有这样的民族，帝国本身的规模也产生了更令人憎恶的战争——社会战争和内战——整个人类都因此而动荡不安，无论是实际冲突还是害怕再次爆发。如果我想恰当地描述这些多样的灾难、这些严峻而持久的必然，尽管我完全无力胜任，但我又能设定什么界限呢？但是，他们说，智者会发动正义的战争。好像他不会更痛惜正义战争的必要性，如果他记得自己是一个人；因为如果战争不是正义的，他就不会发动，因此他将免除一切战争。因为敌对方的恶行迫使智者发动正义战争；而这种恶行，即使没有引起战争，也仍会是人感到悲伤的事，因为它是人的恶行。因此，让每一个带着痛苦想到所有这些巨大邪恶的人，这些如此可怕、如此残忍的邪恶，都承认这是苦难。如果有人忍受或思考它们时没有精神痛苦，那是一种更悲惨的状况，因为他认为自己快乐，因为他失去了人类的情感。\par

% structure: p0024 source=h2 target=subsection reason=relative-heading-rank; base=h2

\subsection{第八章 善人的友谊不能安稳依靠，只要今生的危险迫使我们忧虑。}

在我们目前悲惨的境况中，我们常常把朋友当成敌人，把敌人当成朋友。如果我们避免了这种可怜的盲目，那么，在充满误解和灾难的人类社会中，真诚可信赖的真正善友的相互信任和爱，难道不是我们唯一的慰藉吗？然而，朋友越多，分布越广，我们就越担心生活中的巨大灾难可能会降临到他们身上。因为我们不仅担心他们遭受饥荒、战争、疾病、被俘或难以想象的奴役恐怖，而且更痛苦地害怕他们的友谊可能变成背信弃义、恶意和不义。当这些情况实际发生时——朋友越多、分布越广，发生的频率就越高——当我们得知这些时，除了亲身经历的人，谁能说出内心有多么痛苦？事实上，我们宁愿听到他们去世的消息，尽管听到这个我们也不会不痛苦。因为如果他们的生活以友谊的魅力安慰了我们，那么他们的去世难道不会使我们悲伤吗？不愿有这种悲伤的人，如果可能，就不要有任何友好的交往。让他禁止或熄灭友爱的情感；让他以冷酷无情的麻木破裂所有人际关系的纽带；或者让他设法利用这些关系，使任何甜蜜都不渗入他的心灵。但如果这完全不可能，我们怎能设法在那些生命对我们甜蜜的人去世时不感到痛苦呢？因此产生了那种像伤口或瘀伤一样影响柔嫩心灵的悲伤，而这种悲伤通过和善的安慰得以治愈。因为虽然灵魂状况越好，治愈就越容易、越快，但我们不能因此认为根本没有什么需要医治的。因此，尽管我们现世的生活有时因亲密之人的去世而受折磨，有时较轻，有时更痛苦，尤其是对公众有用的人，但我们仍然宁愿听到他们去世，也不愿听到或察觉他们已背弃信仰或美德——换言之，他们在精神上已死亡。大地充满了这巨大的苦难材料，所以经上记载：\BibleQuote{人的生命在世上不就是考验吗？}\BibleRef{约 7:1} 同样，主说：\BibleQuote{祸哉，世界因了恶事！}\BibleRef{玛 17:7} 又说：\BibleQuote{因为罪恶增多，许多人的爱德冷淡了。}\BibleRef{玛 24:12} 因此，当我们的好朋友去世时，我们感到一些安慰；因为虽然他们的去世使我们悲伤，但我们有安慰的保证，他们已超越了今生连最好的人也会被压垮或腐化的种种邪恶，或同时面临两种危险。\par

% structure: p0026 source=h2 target=subsection reason=relative-heading-rank; base=h2

\subsection{第九章 神圣天使的友谊，在今生人不能确信，因为魔鬼的欺骗束缚了崇拜多神的人。}

那些希望我们与神明为友的哲学家，将圣天使的友谊列在社会之第四圈，从地上的三个社会圈进阶至宇宙，并将天堂本身包括在内。在此友谊中，我们确实无需恐惧天使会因死亡或堕落而使我们忧伤。但因我们无法像与人一样亲密地与他们交往（这本身便是此生的一项痛苦），又因撒殚，如我们读到，\BibleQuote{有时装作光明的天使\BibleRef{格后 11:14}}，以诱惑那些需受管教或仅是受欺骗的人，我们极需天主的仁慈，以免在我们自以为拥有善天使为友时，却与伪装成天使的魔鬼结交；因为这些恶灵的狡诈与欺骗不亚于其危害。这岂不是人类生命的一大不幸，即我们陷于如此无知，若非天主的仁慈，便成为这些魔鬼的猎物？而且，无神之城的哲学家们，他们声称神明是他们的朋友，却确实已成了统治那座城的恶毒的魔鬼的猎物，且要与之共享永远的惩罚。这些魔鬼的本性，已从其崇拜中所行的神圣或更确切说是亵圣的仪式，以及那场他们罪行被庆祝的污秽游戏，得到了充分证明，这些游戏是魔鬼们自己发起并强求其崇拜者作为合宜的赎罪。\par

% structure: p0028 source=h2 target=subsection reason=relative-heading-rank; base=h2

\subsection{第十章 圣徒在经历此生考验后所预备的赏报}

但即使是圣徒和唯一至高真天主的忠实崇拜者，也不能免于魔鬼的各种诱惑和欺骗。因为在这软弱的居所和邪恶的日子里，这种焦虑的状态也有其用处，激励我们以更热切的渴望寻求那完全的、不可侵犯的平安。在那里，我们将享有本性的恩赐，即万有的创造者天主赐予我们本性的所有恩赐——不仅是好的，而且是永恒的——不仅是已因智慧治愈的精神部分，还有在复活中更新的身体部分。在那里，德行将不再与任何邪恶或罪恶争战，而是享受胜利的赏报，即没有仇敌能扰乱的永恒平安。这是最终的福乐，这是终极的圆满，是无尽的终结。在此，当我们拥有在善生中所能享受的那种平安时，我们可被称为有福；但这种福乐相比于那最终的幸福，不过是痛苦而已。当我们有死之人拥有此生所能提供的平安时，如果生活正直，德行便正确运用这平安境况的益处；若我们无此平安，德行则善用人类所遭受的恶事。但这是真正的德行，即它将其善用的一切益处、其在善用善与恶之事时所作的一切、以及自身，都指向那个我们将享有最佳且最大可能的平安的终向。\par

% structure: p0030 source=h2 target=subsection reason=relative-heading-rank; base=h2

\subsection{第十一章 论永恒平安之福乐，此乃圣徒之终向或真成全}

因此，关于平安，我们可以如同论及永生一样说，它是我们善的终向；尤其因为圣咏作者关于天主之城（这部辛劳之作的主题）说：\BibleQuote{耶路撒冷，请赞美主；熙雍，请赞美你的天主：因祂坚固了你的城门，祝福了在你内的子女；使你的边界成为平安。} 因为当她城门的门闩被坚固时，无人能出入；因此我们应当理解她边界的平安就是我们想要宣告的最终平安。连这城本身的奥秘名称，即\Emph{耶路撒冷}，如我已说过的，意为 \Quote{平安的异象}。但由于平安一词也与今世的事物相关联，而永生当然不在其中，我们宁愿称此城的终向或至善为永生而非平安。关于此终向，宗徒说：\BibleQuote{但现今你们既从罪恶中得了释放，作了天主的奴仆，就有你们的果实成圣，结局就是永生。\BibleRef{罗 6:22}} 但另一方面，那些对圣经不熟悉的人可能以为恶人的生命也是永生，或因为灵魂的不朽（某些哲学家也承认这一点），或因为恶人无尽的惩罚（这是我们信仰的一部分，且似乎除非恶人永远活着，否则不可能），因此，为了使每个人都能易于理解我们所说的，也许宜说此城的终向或至善或是永生中的平安，或是在平安中的永生。因为平安是如此大的善，即使在这必死的尘世生命中，也没有任何词如此悦耳、没有任何事物如此热切地被渴望或更令人完全满足。因此，如果我们在这主题上多停留片刻，我认为不会使读者厌烦，读者将既为了理解我们所说的此城的终向，也为了所有人珍视的平安的甘甜而倾听。\par

% structure: p0032 source=h2 target=subsection reason=relative-heading-rank; base=h2

\subsection{第十二章 连战争的残酷和人类的一切骚动都趋向这每一个本性所渴望的唯一平安终向}

但凡稍稍留意人事与共同本性者，必将承认：若无人不愿喜乐，亦无人不愿享有和平。盖凡作战者，所望无非胜利——即望以荣耀达成和平。胜利若非征服反抗者，又为何物？此既成，则和平存焉。故人即使以好战之性为乐，发号施令、驰骋沙场，亦因渴求和平而战。由此显然，和平乃战争所求之终点。盖人人借战争以求和平，却无人借和平以求战争。即使故意破坏所居和平者，亦非憎恶和平，唯愿将其转为更合己意之和平耳。彼等并非欲无和平，唯求更随心之所欲。至于叛乱之际，人虽离群独行，然若不与其同谋者保持某种和平，则所愿亦不得遂。故盗贼亦留意与同伙保持和平，以期更有效、更安全地侵扰他人之和平。设若一人勇力无双，嫉视结伴，不信赖任何同伴，独行其谋，独自劫掠杀戮，则彼仍与那些无法杀害且欲隐藏其行者保持一丝和平之影。在其家中，亦以与妻儿及家中其他成员和睦为务；盖彼等对其眼神之迅速服从无疑令其愉悅。若不如此，则怒而责罚；即以此风暴，随需而保其家庭之宁静和平。盖彼见和平不得维持，除非同一家室所有成员皆服从于一首，即如彼在家中自身之所是。故设若某城或某国愿归顺于他，照其治家之方式侍奉他，则彼将不再藏匿于匪窟，而白日昂首为王，虽其贪欲与邪恶依旧如故。如此，人人皆愿与自己所辖之圈子和平共处，依己意治理之。甚至彼等与之作战者，亦欲使其归为己有，强加其和平之法。\par

然让我们设想一个如诗歌与神话所言之人——如此孤僻野蛮，与其称之为人，毋宁称之为半人。虽其王国乃阴森洞穴之孤寂，而其心性之邪僻至极，竟被名为\OriginalTerm{坏}{\GreekText{Κακός}}——希腊语中\Emph{坏}之意；虽无娇妻以温言相慰，无子女嬉戏，无子嗣奉命，无友交谈以悦之，甚至无其父\Person{武尔坎}（虽有一事较其父为幸：未生如己之怪胎）；虽不与人以物，但随心所欲，遇人则夺，无所不能夺之时；然在那孤寂之洞中——如\Person{维吉尔}所言，其地面常浸最新杀戮之血——所寻求者无非和平，即无人侵扰、无人攻击或惊吓之和平。彼愿与自身躯体和平相处，其满足程度唯视此和平之多寡而定。盖其统御肢体，肢体服从之；为平息其必朽本性之反叛（当其有所需时）并平息饥馑之骚动（此骚动几欲逐灵魂离躯体），彼乃劫掠、杀戮、吞食，然其行为之残暴凶悍，仅用于保全自身生命之和平。如此，若彼愿与他人建立如同在洞穴中与自己建立的和平，则彼既不被称作\Emph{坏}，亦不被视作怪物或半人。若因其外貌与口中喷烟之火使人生畏而不敢与之交往，则其凶暴之行或非出于为恶之欲，而系生存之必然。然此人或许未曾存在，或至少不如诗人虚构那般——彼辈为颂扬\Person{赫拉克勒斯}，不惜贬损\Person{卡库斯}。故宁可相信此类人或半人从未存在，此不过与诗人诸多幻想一般，纯属虚构。盖最残暴之兽类（据说彼几近野兽）亦以保护性之和平环绕其同类。彼等交配、生育、产仔、哺乳、养育幼崽——虽其中许多并非群居，而为独居——非如绵羊、鹿、鸽、椋鸟、蜜蜂，而如狮、狐、鹰、蝙蝠。试问哪只母虎不向幼仔柔声低吟，不敛其凶猛以抚爱之？哪只鸢——虽盘旋猎物时孤身独处——不寻求配偶、筑巢、孵卵、养育雏鸟，并尽己所能与妻子维持和平之家？人之本性法则岂非更强烈地促使其与人交往、与人保持和平，尽其所能？盖恶人亦为维持其圈内和平而战，且希望——若可能——所有人与事物皆归其所有，唯服于一首，或因爱或因惧，而向彼屈服于和平之中！骄傲以其乖张模仿天主：憎恶与他人同在天主之下的平等，非以天主之统治，而欲将其自身之统治强加于平等者。亦即，憎恶天主之正义和平，而爱自身之不义和平；然无论如何，彼不能不爱和平。盖无任何恶习如此违逆本性，以致抹煞本性最微末之痕迹。\par

因此，那偏爱正义而非不义、偏爱秩序而非混乱的人，看到不义之人的平安与义人的平安相比，简直不配称为平安。然而，即便是混乱的事物，也必然与事物的秩序和谐一致、依赖并属于其中一部分，否则它根本不可能存在。假设一个人头朝下悬挂着，这当然是身体的一种反常姿态和肢体的错误安排；因为本性要求在上方的部分却到了下方，\Emph{反之亦然}。这种反常扰乱了身体的平安，因此是痛苦的。然而，精神仍与身体保持平安，并努力维护它，由此产生了痛苦；但如果痛苦使精神脱离身体，那么只要身体的框架还保持在一起，残余部分中就有一种肢体的平安，因此身体仍然悬挂着。而且，由于尘世的身体趋向大地，并依靠悬挂它的绳索而休息，它趋向于它自然的平安，它自身重量的呼声要求一个安息之处；尽管现在已无生命、无感觉，它并未脱离其在受造界中位置所固有的平安，无论它已经拥有这种平安，还是正在趋向它。因为如果你施用防腐香料防止身体框架腐烂分解，那么一种平安仍然将各部分联结在一起，并将整个身体保持在大地上一个合适的位置——也就是说，在与身体平安的位置上。另一方面，如果身体没有得到这样的照料，而是顺其自然，它就会因不相和谐且刺激我们感官的臭气而受到扰乱；因为在腐烂过程中，直到它融入世界的元素，并逐渐与它们进入平安，我们所感知的正是这一点。然而在整个过程中，至高创造者和治理者的法律被严格地遵守，因为宇宙的平安是由祂管理的。虽然从较大动物的尸体中会产生微小的生物，但所有这些小原子，按照同一创造者的法律，平安地服务于它们所属的动物。尽管死动物的肉被其他动物吃掉，无论它被带到哪里，与什么接触，或被转化和改变成什么，它仍然受着贯穿万物的同一法律的支配，这些法律是为了保存每一个有死的种类，并使彼此适合的事物和谐一致。\par

% structure: p0036 source=h2 target=subsection reason=relative-heading-rank; base=h2

\subsection{第十三章——论自然律在一切混乱中保存的普遍平安，以及每个人如何按照公义审判者的规定获得他应得的归宿。}

身体的平安在于其各部分比例的适当安排。非理性灵魂的平安是欲望的和谐宁静，理性灵魂的平安则是知识与行动的和谐。身体与灵魂的平安是有生命者的有序和谐的生命与健康。人与天主之间的平安是信德对永恒法律的有序服从。人与人之间的平安是有序的和谐。家庭的平安是家中统治与服从者之间的有序和谐。公民的平安是公民之间类似的和谐。天上之城的平安是在天主内以及彼此在天主内的完全有序和谐的享受。万物的平安是秩序的安宁。秩序是分配使平等与不平等各得其所的安排。因此，尽管悲惨者就其悲惨而言确实不享平安，而是与那种没有纷扰的秩序安宁分离，然而，由于他们理应且公正地悲惨，他们的悲惨本身使他们与秩序相连。他们确实不与有福者结合，而是因秩序的法律与他们分离。虽然他们不安，但他们的处境仍与之相适，因此他们拥有某种秩序的安宁，从而也有某种平安。但他们是悲惨的，因为虽然并非完全悲惨，但他们不在那不可能有任何悲惨混杂的地方。然而，如果他们不曾拥有那种来自与事物自然秩序相和谐的平安，他们会更加悲惨。当他们受苦时，他们的平安在某种程度上被扰乱；但他们的平安继续存在，只要他们不受苦，只要他们的本性继续存在。因此，正如可能有生命而无痛苦，却不可能有痛苦而无某种生命，同样，可能有平安而无战争，却不可能有战争而无某种平安，因为战争预设了进行战争的本性的存在，而这些本性若无某种平安就不可能存在。\par

因此，有一种本性，其中恶不存在，甚至不可能存在；但不可能有一种本性中没有任何善。所以，连魔鬼自己的本性也不是恶的，就其为本性而言，但因被扭曲而成为恶的。因此，他没有持守真理\BibleRef{若 8:44}，却不能逃脱真理的审判；他没有持守秩序中的安宁，但也没有因此逃脱安排者的大能。天主赋予他本性的善，并没有遮蔽天主的正义，这正义在惩罚中保持了秩序；天主也没有惩罚他所创造的善，而是惩罚魔鬼所犯的恶。天主并没有收回他赋予本性的全部，而是取走了一些，留下了一些，好留下足够的东西，使他对被取走的东西感到失落。而这种痛苦的感觉正是被取走的善和留下的善的证据。因为，如果没有留下任何善，就不会有因失去善而产生的痛苦。因为犯罪的人若因失掉义德而欢喜，则更糟。但受苦的人，即使从中得不到益处，至少为失去健康而哀伤。义德和健康都是善的事物，失去任何善的事物都是悲伤的事，而不是欢喜的——至少如果没有补偿，如同属灵的义德可以补偿身体的健康——那么，恶人受惩罚时悲伤，总比他犯罪时欢喜更为合适。因此，正如一个放弃善的罪人的欢喜是恶意的证据，同样，他受惩罚时对失去的善的悲伤，是善的本性的证据。因为那哀叹自己本性失去平安的人，是被一些平安的残余所驱动，这些残余使他的本性与自己友好。在最终的惩罚中，邪恶和不敬虔的人痛苦地哀叹失去他们曾享有的自然益处，并意识到这些益处正是被他们所藐视的天主那慈祥的慷慨所极其公正地夺走的，这是非常公正的。因此，天主，最智慧的创造者和最公正的万有安排者，把人类放在地上作为其最伟大的装饰，赐予人类一些适合今生的善物，即今世的平安，如我们今生从健康、安全和人类共融中能够享受到的，以及为保存和恢复这种平安所需的一切，例如适合我们外在感官的事物——光、夜、空气、合适的水，以及身体所需的一切来维持、遮蔽、医治或美化它——而且这一切都是在最公平的条件下：凡善用这些适合此生平安之优势的人，将获得更广大的和更好的祝福，即不朽的平安，伴随着在无尽的生命中的荣耀和尊荣，这生命适于享受天主并在天主内彼此享受；但那恶用现世祝福的人，既失去它们，也得不到其他的。\par

% structure: p0039 source=h2 target=subsection reason=relative-heading-rank; base=h2

\subsection{论天地间的秩序和法则，由此人类社会为掌权者所服事。}

一切属世之物的全部用处，都指向这尘世社群中属世和平的结果，而在天主之城，它则与永恒的和平相连。因此，倘若我们是无理性的动物，我们便不会渴望别的，只求身体各部分的适当安排和欲望的满足——也就是说，只求身体的舒适和感官快乐的丰盈，使身体之和平有助于灵魂之和平。因为若缺少身体之和平，甚至无理性的灵魂之和平也受到阻碍，因它无法满足其欲望。这两者共同促成灵魂与身体相互的和平，即和谐生活与健康的和平。正如动物通过躲避痛苦表明它们喜爱身体之和平，又通过追求快乐以满足欲望表明它们喜爱灵魂之和平，同样，它们对死亡的退避也充分显示它们对那紧密联结灵魂与身体之和平的强烈热爱。但是，由于人拥有理性的灵魂，他将这些与动物共有的一切都从属于其理性灵魂之和平，以便其理智得以自由运作并规范其行为，从而享有认识与行动的有序和谐——这正如我们所说，构成了理性灵魂之和平。为此，他必须愿意既不为痛苦所困扰，也不为欲望所搅扰，更不为死亡所灭绝，以便获得某种有益的知识，据此规范他的生活与行为。但因人心易于犯错，这种对知识的追求本身可能成为他的陷阱，除非他有一位神圣的导师，他能够无疑虑地服从他，同时导师也给他帮助以保全其自由。并且，由于他尚在这可朽的身体中，是远离天主的旅客，他凭着信德行走，并非凭着目睹；因此他将一切和平——无论是身体的、灵魂的，或两者的——都指向那必朽之人与不朽天主之间的和平，从而展现出对永恒律法有秩序的信德服从。但这位神圣导师教导两条诫命——爱天主与爱近人——而在这两条诫命中，人找到三样他必须爱的事物：天主、自己及近人——并且那爱天主的人亦因此爱自己。由此推出，他必须努力使他的近人爱天主，因为他被命令要爱人如己。他应当为其妻子、儿女、家人及所有在他能力范围内的人作此努力，正如他若需要时也愿他的近人为他如此行；因此，他将尽可能与众人和平相处，即有秩序地和谐一致。这种和谐的秩序是：首先，不伤害任何人；其次，向一切可能的人行善。因此，他的首要责任是他的家人，因为自然律与社会律使他更容易接近他们，并有更多机会服务他们。因此宗徒说：\Quote{谁若不照顾自己的亲属，尤其自己的家人，便是背弃了信德，比不信者更坏。} \BibleRef{弟前 5:8} 这便是家庭和平的起源，即家中统治者和服从者之间的有序和谐。因为照顾他人者治理——丈夫治理妻子，父母治理子女，主人治理仆人；而受照顾者服从——妻子服从丈夫，子女服从父母，仆人服从主人。但在那凭信德生活、尚在朝圣途中奔赴天上之城的义人的家庭中，即使是那些治理者，也服务那些他们似乎管辖的人；因为他们治理不是出于对权力的爱，而是出于对他人所尽职责的意识——不是因为骄傲于权威，而是因为喜爱仁慈。\par

% structure: p0041 source=h2 target=subsection reason=relative-heading-rank; base=h2

\subsection{第15章——论人性本有的自由及因罪引入的奴役——在此奴役中，意志邪恶之人是他自己私欲的奴隶，尽管就他人而言他是自由的。}

这是本性的秩序所规定的：天主就是这样创造了人。因为祂说：\BibleQuote{让他们，}祂说，\BibleQuote{管理海中的鱼、空中的鸟，以及在地上爬行的一切爬虫。} \BibleRef{创 1:26}祂并不愿意按祂肖像所造的理性受造物统治任何事物，除非是无理性的受造物——不是人统治人，而是人统治牲畜。因此，远古时代的义人被立为羊群的牧人，而不是人的君王，天主以此教导我们受造物的相对位置以及罪的后果；因为我们相信，奴隶的境况是罪的结果，这是公义的。这就是为什么我们在圣经任何部分都找不到\Quote{奴隶}一词，直到正义的\Person{诺厄}用这个名字标记他儿子的罪。因此，这个名字是因罪而来，而非本性。拉丁文表示奴隶的词，其起源据推测是这样的：那些按战争法本应被杀的人，有时被胜利者保留下来，从而被称为\Quote{仆人}。这种情况若非因罪，绝不可能发生。因为即使我们进行正义的战争，我们的对手也必定在犯罪；而每一场胜利，即使是由恶人取得的，都是天主第一次审判的结果，祂使被征服者谦卑，无论是为了消除他们的罪，还是惩罚他们的罪。请看那位属天主的人\Person{达尼尔}，他在被掳时向天主承认自己和自己人民的罪，并以虔诚的悲痛宣告这些是掳掠的原因。\EditorialNote{达尼尔书第九章}因此，奴隶制的首要原因是罪，它使人受制于同类的统治——这只有通过天主的审判才会发生，在祂那里没有不义，祂知道如何对每一类过犯给予适当的惩罚。但我们在天堂的导师说：\BibleQuote{凡犯罪的，就是罪的奴隶。} \BibleRef{若 8:34}因此，有许多邪恶的主人拥有虔敬的人作他们的奴隶，而他们自己却仍处于束缚之中；\BibleQuote{因为人被谁制伏，就是谁的奴隶。} \BibleRef{伯后 2:19}毫无疑问，作人的奴隶比作情欲的奴隶更为幸福；因为就连这种统治他人的欲望——更不用说其他欲望——也以最残酷的统治蹂躏人心。此外，当人们在和平秩序中彼此服从时，卑微的地位对仆人有好处，正如骄傲的地位对主人有伤害一样。但就本性而言，正如天主最初创造我们时，没有人是人的奴隶，也没有人是罪的奴隶。然而，这种奴役是惩罚性的，是由那命令维持自然秩序并禁止扰乱它的法律所规定的；因为如果没有违反那法律的事发生，就没有必要用惩罚性的奴役来约束。因此，宗徒劝诫奴隶要服从他们的主人，要真心诚意地服事他们，这样，如果他们不能被主人释放，他们也可以使自己的奴隶以某种方式变得自由，不是以诡诈的恐惧，而是以忠信的爱来服事，直到一切不义都过去，一切权柄和一切人的能力都被消灭，天主成为万物中的万有。\par

% structure: p0043 source=h2 target=subsection reason=relative-heading-rank; base=h2

\subsection{第十六章 论公平的统治}

因此，虽然我们正义的祖先拥有奴隶，并管理他们的家务，在此生福乐方面区分奴隶的地位和儿子的继承权，但在关于天主的崇拜上——我们盼望从祂获得永福——他们对所有家庭成员同等关爱地加以监督。这非常符合自然秩序，以至于一家之长被称为\OriginalTerm{家长}{paterfamilias}；这个名称已被普遍接受，以致那些统治不义的人也乐于将其用于自身。但那些是真正家长的人，愿意并努力使所有家庭成员，如同自己的儿女一样，崇拜天主并赢得天主，并应达到那个天上家园，在那里统治人的责任不再需要，因为照顾他们永恒幸福的职责也已终止；但是，在到达那家园之前，主人应感到他们的权威地位比仆人的服务更重的负担。如果任何家庭成员因不顺从而扰乱家庭和平，他就会被用言语或责打，或某种社会所允许的公正合法的惩罚加以纠正，使他自己从中得益，并重新融入他被隔断的家庭和谐。因为正如以牺牲某人可能获得的更大利益来帮助他并非仁慈，同样，冒着使其陷入更重罪过的危险而宽恕他也并非无罪。要成为无罪，我们不仅不伤害任何人，还必须阻止他犯罪或惩罚他的罪，以使受惩罚者本人能从其经历中获益，或使他人引以为戒。因此，家庭应当成为城邦的开端或元素，而每一个开端都指向其种类的某个终点，每一个元素都指向其所从属的整体之完整，因此显然，家庭和平与城邦和平有关联——换言之，家庭服从与家庭统治的有序和谐，与公民服从与公民统治的有序和谐有关联。由此进一步得出，家长应当按照城邦的法律来制定其家庭规则，以使家庭与城邦秩序相和谐。\par

% structure: p0045 source=h2 target=subsection reason=relative-heading-rank; base=h2

\subsection{第十七章 论什么在天上之城和地上之城之间产生和平与不和}

但是，那些不凭信德而生活的家庭，在此生中寻求尘世的利益；而凭信德生活的家庭，则期盼那被预许的永恒福乐，并像旅居者一样，使用那些既不诱惑他们、也不使他们偏离天主的暂时的和尘世的利益，反而帮助他们更轻松地忍受那可朽坏的身体加于灵魂的重担。因此，这两种人和家庭同样使用此可朽生命所需之物，但各自在使用时有自己特殊且截然不同的目的。地上之城，不凭信德而生活，寻求地上的和平，它在公民服从与统治的有序和谐中所提出的目标，是结合人的意志以获得对此生有益之物。天上之城，或者说它在世上旅居并凭信德生活的那一部分，只是出于必要才使用这种和平，直到需要这种和平的可朽状态消逝为止。因此，当它像俘虏和异乡人一样生活在世上之城中时，虽然已领受了救赎的预许和圣神作为凭据，它仍毫无顾忌地遵守地上之城的法律，因为这些法律管理着维持此可朽生命所需之物；因此，既然此生为两城所共有，那么在与此生相关的事物上，它们之间便存在着和谐。然而，地上之城曾有一些哲学家，他们的学说受到神圣教诲的谴责，他们被自己的猜测或魔鬼所欺骗，认为必须邀请许多神介入人类事务，并为每尊神分配不同的职能和部门——一位管身体，另一位管灵魂；在身体本身，一位管头，另一位管脖子，其他肢体各归一位神；同样，在灵魂中，将自然能力分配给一位神，教育给另一位，愤怒给另一位，欲望给另一位；如此，生活的各种事务也被分配——牲畜归一位神，谷物归另一位，酒归另一位，油归另一位，森林归另一位，金钱归另一位，航海归另一位，战争与胜利归另一位，婚姻归另一位，生育与丰饶归另一位，其他事物归其他神；另一方面，天上之城知道只有一位天主应当受敬拜，惟有祂才配受那希腊人称为\OriginalTerm{敬礼}{\GreekText{λατρεία}}的服侍，而这种服侍只能献给一位神。因此，这两座城在宗教的法律上不可能达成一致，天上之城被迫在这件事上持异议，并招致那些持不同思想者的反感，承受他们的愤怒、仇恨和迫害，除非敌人的心被基督徒的众多人数所震慑，并被天主赐予他们的明显保护所平息。因此，天上之城在世上旅居时，从所有民族中召叫公民，并聚集一个由各种语言组成的旅居者团体，对于用以确保和维护地上和平的风俗、法律和制度上的差异毫不介意，反而认识到，无论这些差异多么不同，它们都趋向地上和平这同一个目的。因此，它远非废除或取消这些差异，反而保存并采纳它们，只要不因此引入对唯一至高真天主的敬拜的阻碍。因此，即使天上之城在旅居状态中，也利用地上的和平，并在不损害信德与虔诚的范围内，渴望并维护人们之间关于获取生活必需品的共同协议，并使这地上的和平指向天上的和平；因为只有这和平才能被真正称为并视为理性受造物的和平，它在于完全有序且和谐地享见天主，并因天主而彼此共享。当我们达到那和平之后，这必死的生命将让位于永恒的生命，我们的身体将不再是这因腐败而沉重压着灵魂的动物身体，而是毫无匮乏的、所有肢体都服从意志的神性身体。天上之城在旅居状态中凭信德拥有这和平；并凭这信德正义地生活，将其对天主和人的每一善行都指向那和平的获得；因为城的生活是一种社会生活。\par

% structure: p0047 source=h2 target=subsection reason=relative-heading-rank; base=h2

\subsection{第十八章——新学园派的不确定性如何不同于基督徒信德的确定性。}

关于瓦罗断言为新学园派区别性特征的对一切事物的不确定性，天主之城彻底憎恶这种怀疑，视之为疯狂。对于它凭理智和理性所认识的事物，它具有最绝对的确定性，尽管由于可朽坏的身体压着心灵，它的知识是有限的，因为正如宗徒所说：\BibleQuote{我们知道，只是局部的。}\BibleRef{格前 13:9}。它相信心智借助身体所运用的感官的证据；因为如果有时信赖感官的人受了欺骗，那么幻想永远不应信赖它们的人则被更可怜地欺骗了。它相信新旧圣经，我们称之为正典，它们是义人凭信德而生活的信德之源\BibleRef{哈 2:4}，也是我们因不在主那里而凭此毫无疑惑地行走的信德之源\BibleRef{格后 5:6}。只要这信德保持不受侵犯且坚定，我们就可以毫无责备地怀疑某些事物，这些事物我们既未凭感官也未凭理性察觉，也未被正典圣经启示给我们，也未通过我们荒谬地不能不相信的证人所知。\par

% structure: p0049 source=h2 target=subsection reason=relative-heading-rank; base=h2

\subsection{第十九章——关于基督徒人民的衣着和习惯。}

在天主之城内，一个人接受引人归向天主的信德，无论他采取何种服饰和生活方式，只要他按天主的诫命生活，就无关紧要。因此，当哲学家们自己成为基督徒时，他们确实被迫放弃错误的学说，但不放弃不妨碍宗教的服饰和生活方式。所以我们不重视\Person{瓦罗}在提到犬儒学派时所提出的那种教派区分，只要不保留任何不雅或放纵之事。至于这三种生活方式，即默观的、行动的与混合的，虽然只要人的信德得以保持，他可以任选一种而不损害其永恒利益，但他绝不可忽视真理与义务的要求。没有人有权过一种默观的生活，以至于在自己安逸中忘记了对邻人应尽的义务；也没有人有权如此沉浸于行动的生活，以至于忽略了默观天主。闲暇的魅力绝不能是懒散的空虚，而是探求或发现真理，以便每个人都能取得扎实的成就，而不嫉妒他人同样如此。在行动的生活中，我们不应贪图今世的荣誉或权力，因为太阳底下的一切都是虚幻，而应以我们已有的地位和影响力（如果这些是以可敬的方式获得的）来造福那些在我们之下的人，如我们已说明的那样。宗徒指着这一点说：\BibleQuote{谁渴望主教职，就是渴望一件善工。}\BibleRef{弟前 3:1} 他想表明，主教职是工作的名称，而非荣誉的名称。这是一个希腊词，意思是说，治理者监督或照料他所治理的人：因为\OriginalTerm{在上}{\GreekText{ἐπί}}意为\Emph{在上}，\OriginalTerm{看}{\GreekText{σκοπεῖν}}意为\Emph{看}；所以\OriginalTerm{监督}{\GreekText{ἐπισκοπεῖν}}意为\Quote{监督}。因此，喜爱治理胜过行善的人不是主教。因此，没有人被禁止探求真理，因为闲暇可以用来做最可赞誉的事；但贪图治理人民所需的高位却是不合宜的，即使那个位置被持有，且治理以合宜的方式进行。因此，圣善的闲暇因爱真理而被渴望；但爱德的需要驱使我们承担必要的事务。若无人强加此担子于我们，我们便可自由地筛选和默观真理；但若担子被加于我们，我们便因爱德不得不承担它。然而，即使在这种情况下，我们也毋须完全放弃默观的甘甜；因为若这些甘甜被撤去，担子可能变得不堪承受。\par

% structure: p0051 source=h2 target=subsection reason=relative-heading-rank; base=h2

\subsection{第二十章——圣人们在今生因望德而蒙福}

既然天主之城的至善是完美而永恒的和平，不是凡人经由出生与死亡而进出的那种和平，而是脱离一切邪恶的和平，不朽者永远安息其中；谁能否认那未来的生命是最有福的，或者与此相比，我们现在所过的今生是最悲惨的，即使它充满身体、灵魂和外在的一切福乐？然而，如果有人以此生参照他所热切爱与坚信的那另一生命，那么他甚至现在也可被称为有福的，尽管不是在实际上而是在望德中。但实际拥有今生的幸福而没有对来世的希望，只是一种虚假的幸福和深重的悲惨。因为灵魂的真正福乐现在并不享受；因为那并非真正的智慧，它不将其所有谨慎的观察、刚毅的行动、有德的自制和公正的安排导向那个终向，即天主在永恒安全与完美和平中成为万有中的万有。\par

% structure: p0053 source=h2 target=subsection reason=relative-heading-rank; base=h2

\subsection{第二十一章——是否曾有符合\Person{西塞罗}对话录中\Person{斯基皮奥}定义的罗马共和国}

因此，这正是我应当履行在本作品第二卷中所作承诺的地方，要尽可能简明清楚地解释，如果我们接受\Person{西塞罗}\WorkTitle{论共和国}中\Person{斯基皮奥}所下的定义，那么从来就不存在什么罗马共和国；因为他简要地将共和国定义为人民的福祉。如果这定义是正确的，那么从来就不存在罗马共和国，因为在罗马人中从未实现过人民的福祉。因为按他的定义，人民是由共同承认权利和共同利益而联合起来的群体。而他所指的\Quote{共同承认权利}他详细解释，表明没有正义就无法管理共和国。因此，何处没有真正的正义，何处就没有权利。因为出于权利而行的事是正义地行出的，而不正义地行出的事不能出于权利而行。因为人的不义发明既不应被视为也不应被称为权利；就连他们自己也说，权利是源于正义之泉的，并否认那些误解此事的人通常给出的定义，即正义是对强者有利的东西。因此，何处没有真正的正义，何处就没有由共同承认权利而联合起来的人的群体，因此也就没有\Person{斯基皮奥}或\Person{西塞罗}所定义的人民；如果没有人民，就没有人民的福祉，而只是一群不配人民之名的乌合之众。因此，如果共和国是人民的福祉，而没有由共同承认权利而联合起来的人民就没有人民，而没有正义就没有权利，那么毫无疑问，没有正义就没有共和国。此外，正义是那种给予每个人应得之物的德性。那么，当人离弃真天主而屈服于不洁的魔鬼时，他的正义在哪里呢？这是给予每个人应得之物吗？或者，一个从购买者那里扣留一块土地并把它给一个无权拥有的人是不义的，而一个把自己从创造他的天主那里扣留下来并服事邪恶神灵的人却是正义的吗？\par

这本\WorkTitle{De Republica}以极大的力量和锐气为正义辩护，反对不正义。首先听到了为不正义辩护的言论，断言没有不正义，共和国既不能扩展，甚至不能存续，因为一个绝对不容置疑的立场是：一些人统治、一些人服侍是不正义的；然而这个共和国所属的帝国城市，若不诉诸这种不正义，就无法统治其行省。为正义辩护的回答是：统治行省是正义的，因为奴役对行省居民可能是有益的，而且当正确管理时——也就是当无法无天的人被阻止作恶时——它确实如此。此外，他们在自由状态下变得越坏，受制于人就会改善。为了确认这一推理，加入了一个来自自然的显著例子：\Quote{为什么}有人问，\Quote{天主统治人，灵魂统治身体，理性统治激情和灵魂的其他邪恶部分吗？}这个例子毫不怀疑，对一些人来说，奴役是有用的；而且，事实上，服侍天主对所有人都是有用的。当灵魂服侍天主时，它才能正确地控制身体；在灵魂本身中，理性必须服从天主，才能适当地统治激情和其他恶习。因此，当一个人不服侍天主时，我们能赋予他什么正义呢？因为在这种情况下，他的灵魂不能公正地控制身体，他的理性也不能控制他的恶习。如果个人没有正义，那么由这样的人组成的团体当然也不可能有正义。因此，这里没有那种共同承认的正义，这种正义使一群人成为一个我们称之为共和国的人民。我又何必谈论功利——根据定义，使人民成为人民的共同参与？因为，如果你仔细考虑，你会发现对不虔敬生活的人没有任何益处，因为每一个不服侍天主而服侍魔鬼的人——这些魔鬼的邪恶你可以从他们渴望接受人的敬拜来度量，尽管他们是最不洁的邪灵——然而，我所说的关于共同承认的正义足以表明，根据上述定义，没有正义，就没有人民，因此也没有共和国。因为如果他们断言，在他们的共和国中，罗马人服侍的不是不洁的邪灵，而是善良圣洁的神，那么我们必须再次回应这种遁词，尽管我们已经说得足够多，甚至过多来揭露它？一个特别愚蠢或无耻好辩的人，读完前面这几卷到这里，还能怀疑罗马人服侍的是邪恶不洁的魔鬼吗？但是，不说他们的品性，真天主的法律中写着：\BibleQuote{凡祭祀其他神祇，只祭祀上主一神的，应完全毁灭。}\BibleRef{出 22:20}因此，说出如此威胁性诫命的那位，命定不应向善神或恶神祭献。\par

% structure: p0056 source=h2 target=subsection reason=relative-heading-rank; base=h2

\subsection{第二十二章——基督徒所事奉的天主是否是那惟独应向他祭献的真天主？}

但可能会有人反问：这位天主是谁，有什么证据表明唯独祂配受罗马人的祭献？一个人必须非常盲目才会仍在问这位天主是谁。祂就是那位先知们所预言、而我们看到已经应验之事的天主。祂就是那位亚巴郎从祂那里得到保证的天主：\BibleQuote{地上万民都要因你的后裔蒙受祝福。}\BibleRef{创 22:18}这一点在按肉身来自那后裔的基督身上实现了，即使那些继续敌对这名号的人，无论他们是否愿意，都承认这一点。祂就是那位天主，其圣神借着我们在前面几卷中引用其预言的那些人说话，而这些预言在遍及世界的教会中已经应验。这就是那位罗马最有学问的\Person{瓦罗}以为是朱庇特的天主，尽管他不知道自己在说什么；然而我认为有必要指出，这样一个有学问的人无法设想这位天主不存在或可鄙，而是相信祂与至高的天主是同一的。总之，祂就是那位最有学问的哲学家、即使是最敌视基督徒的\Person{波斐利}，也根据他所尊崇之神明的谕旨，承认是伟大的天主的那个天主。\par

% structure: p0058 source=h2 target=subsection reason=relative-heading-rank; base=h2

\subsection{第二十三章——\Person{波斐利}关于神明谕旨论及基督的回答的陈述。}

因为在他题为\OriginalTerm{神谕哲学}{\GreekText{ἐκ} \GreekText{λογίων} \GreekText{φιλοσοφίας}}的书中，他收集并评论了他假托神明关于神圣事物所发出的神谕，他说——我用从希腊文翻译过来的他的话：\Quote{有人询问他应当讨好哪一位神，以使他的妻子远离基督教，阿波罗用下面的诗句回答。}接着，以下的话被当作阿波罗的话：\Quote{你或许会发现，在水上写永恒的字符，或像鸟一样轻快地飞过空中，比恢复你不敬神的妻子的正当情感更容易，一旦她玷污了自己。让她随她所愿留在愚昧的欺骗中，为她死去的天主唱虚假的哀歌，这位天主被明智的法官定罪，并因暴力死亡而可耻地灭亡。}然后，在这些阿波罗的诗句之后（我们给出拉丁文译本，未保留格律形式），他接着说：\Quote{在这些诗句中，阿波罗揭露了基督徒不可救药的腐败，说犹太人比基督徒更认识天主。}看他是如何歪曲基督，在认识天主方面偏袒犹太人而贬低基督徒。这就是他对阿波罗诗句的解释，其中他说基督是被明智或公正的法官处死的——换言之，他该死。我把这个关于基督的神谕的责任留给阿波罗的撒谎解释者，或者留给这位哲学家，他相信了它，或者可能自己编造了它；至于它是否与波菲利的意见或其他神谕一致，我们稍后会说。但在这一段中，他说犹太人作为天主的解释者，公正地判定基督应受最可耻的死亡。那么，他应该听从这位他为之作证的犹太人的天主，因为这位天主说：\Quote{凡祭献给其他神，惟独不祭献给主者，应被彻底消灭。}但让我们来看更明白的表述，听听波菲利认为犹太人的天主是多么伟大的神。他说，当阿波罗被问及\InnerQuote{话语}（即\Emph{即}理性或法律）哪一样更好时，用下面的诗句回答。然后他给出阿波罗的诗句，我从其中挑选以下这些作为足够：\Quote{神，万有之先的创造者和君王，在他面前，天、地、海和阴间的隐秘处都战栗，众神明自己也恐惧，因为他们的法律就是圣洁的希伯来人所尊崇的父。}在他的神阿波罗的这个神谕中，波菲利承认希伯来人的天主是如此伟大，以至于众神明自己在祂面前都恐惧。因此，我感到惊讶，当天主说\InnerQuote{祭献给其他神的，应被彻底消灭}时，波菲利自己竟然不害怕他会因祭献给其他神而被消灭。\par

然而，这位哲学家对基督也有好评，好像忘记了他刚才所说的那种侮辱；或者好像他的神只在睡着时才说基督的坏话，而醒来时却认出他是善的，并给予他应得的赞美。因为，好像要宣布某种难以置信的奇妙事情，他说：\Quote{我们即将说的话肯定会让一些人吃惊。因为众神已宣告基督非常虔诚，并且成为不朽，他们珍视对他的记忆；然而，基督徒是污秽的、不洁的，并陷入错误。}他说：\Quote{众神还说了许多其他反对基督徒的话。}然后，他列举了众神控告他们的例子，接着又说：\Quote{但有些人问赫卡忒基督是否是神，她回答说：你知道脱离肉体的不朽灵魂的境况，如果它与智慧分离，它就总是犯错。你提到的灵魂是一个极其虔诚之人的灵魂：他们敬拜它，因为他们误解了真理。}对于这个所谓的神谕回应，他加上了自己的话：\Quote{赫卡忒说，这个非常虔诚的人，他的灵魂像其他好人的灵魂一样，死后获得了不朽，而基督徒由于无知敬拜它。对那些问为什么他被判死刑的人，女神的答复是：身体确实总是受折磨，但虔诚之人的灵魂安居在天上。你所询问的灵魂已成为其他灵魂致命错误的原因，这些灵魂的命运不是接受众神的恩赐，也不是认识不朽的朱庇特。这样的灵魂因此被众神憎恨；因为那些命运注定不接受神恩、不认识神的人，命运注定要因你所说的那个人而陷入错误。然而，他本身是善的，天堂已像对其他好人一样向他敞开。所以，你们不可说他的坏话，而要怜悯人的愚昧；通过他，人的危险迫在眉睫。}\par

谁如此愚蠢，竟看不出这些神谕要么是某个对基督徒怀有强烈敌意的聪明人杜撰的，要么是不洁的魔鬼为类似目的——即为了使他们对\Person{基督}的赞美赢得人们对他们诋毁基督徒的信任——而发出的回答，以便他们借此尽可能堵住永生救恩之路，而这永生救恩正等同于基督教？因为他们相信，只要人们对基督徒的诽谤也被接受，他们就绝不会因促进对\Person{基督}的信仰而阻碍自己的有害诡计；因为他们这样确保了，即使一个人对\Person{基督}有好感，也拒绝成为基督徒，因此就不会被他所赞美的那位\Person{基督}从他们自己的统治下解救出来。此外，他们对\Person{基督}的赞美是如此设计的，以至于任何相信如此被描绘的\Person{基督}的人都不会是一个真正的基督徒，而是一个\OriginalTerm{斐利济恩异端者}{Photinian heretic}，只承认\Person{基督}的人性而不承认祂的神性，从而被排除在救恩之外，无法摆脱这些魔鬼谎言的罗网。就我们而言，\Person{赫卡忒}对\Person{基督}的赞美并不比\Person{阿波罗}对\Person{基督}的诋毁更令我们满意。\Person{阿波罗}说\Person{基督}是\Quote{被公正的审判官处死的}，暗示祂是不义的。\Person{赫卡忒}说祂是\Quote{一个最虔诚的人，但也仅此而已}。两者的意图是相同的：阻止人们成为基督徒，因为如果这一点得到确保，人们就永远无法从他们的权势下被解救出来。但是，我们的哲学家——或者更确切地说，那些相信这些针对基督徒的虚假神谕的人——有责任首先，如果可能的话，使\Person{阿波罗}和\Person{赫卡忒}对\Person{基督}的看法一致，以便要么两者都谴责\Person{基督}，要么两者都赞美祂。即使他们成功了，我们这边也无论如何会拒绝魔鬼的见证，无论这见证对\Person{基督}有利还是不利。但当我们的反对者发现他们自己的一个神和一个女神对\Person{基督}意见不一致——一个赞美祂，另一个诋毁祂——时，他们如果有任何判断力，肯定不能相信那些亵渎基督徒的凡人了。\par

当波斐利或赫卡忒赞美基督，并补充说祂将自己作为一份致命的礼物赐给了基督徒，好使他们陷入错误时，他（按他自己的看法）揭示了这错误的原因。但在引用他此意的言论之前，我要问：如果基督如此将自己赐给基督徒，使他们陷入错误，祂是甘心如此，还是违背自己的意愿？若是甘心，祂如何是正义的？若是违背意愿，祂如何是有福的？不过，让我们听听这错误的原因吧。他说：\Quote{有一些，}在某个地方非常小的属地的神体，隶属于邪恶魔鬼的权下。希伯来的贤者——其中就有这位耶稣，正如你们从上文引用的阿波罗神谕所听到的——使宗教人士远离这些极其邪恶的魔鬼和较低的神体，教导他们反而要崇拜天上的诸神，尤其是要敬拜天主圣父。他说：\Quote{这是诸神所吩咐的；并且我们已经表明，他们如何劝告灵魂转向天主，并命令灵魂崇拜祂。但是那些无知和不敬神的人，他们不配领受来自诸神的恩惠，也不认识不朽的朱庇特，不听诸神和他们的讯息，便远离了所有神，不仅不憎恨，反而尊敬那些被禁止的魔鬼。虽然声称崇拜天主，却拒绝做那些唯独能崇拜天主的事。因为天主，作为万有之父，实在一无所缺；但对我们来说，藉着正义、贞洁和其他德行来敬拜祂，并藉着探究和效法祂的本性，使生命本身成为向祂的祈祷，这是好的。因为探究}他说，\InnerQuote{净化我们，而效法使我们神化，藉此将我们移近祂。}他在宣告天主圣父以及我们应当崇拜祂的行为方面是正确的。希伯来的先知书充满了这样的训诫，它们在称赞或责备圣人的生活时也如此。但在谈论基督徒时，他犯了错误，并按照他所视为神的魔鬼所愿的程度毁谤他们，仿佛任何人要回想在剧场和神庙中为取悦神祇而进行的那些可耻和羞愧的行为，并将这些与我们在教会中所听到的、向真天主所奉献的相比较，从而推论出在哪里品格得到建造、在哪里被毁坏，是困难的事。但除了魔鬼般的精神之外，还有谁向他述说或暗示了如此明显而虚妄的谎言，即基督徒尊敬而非憎恨魔鬼，而希伯来人却禁止崇拜魔鬼？但是，希伯来贤者所崇拜的那个天主，禁止甚至在祭献中将祭品献给天上的圣天使和神圣的权能，而我们在这一旅途中，将这些人尊为最有福的同胞并热爱他们。因为在天主赐给祂的希伯来人民的法律中，祂以雷霆般的声音发出了这样的威胁：\BibleQuote{凡祭献给任何神祇的，除非只祭献给主，他应被彻底消灭。}\BibleRef{出 22:20}并且，为了使没有人以为这禁令只针对那些非常邪恶的魔鬼和属地的神体——这位哲学家称之为非常小而低下的东西——因为甚至在圣经中，这些也被称为神，不过不是希伯来人的神，而是外邦人的神，正如\WorkTitle{七十贤士译本}在圣咏中所表明的，经上说：\BibleQuote{因为万邦的神都是魔鬼}——我说，为了使没有人以为只是祭献这些魔鬼被禁止，而祭献给所有或部分天上的神祇却可以，紧接着又加上：\BibleQuote{只祭献给主。}因此，希伯来人的天主——这位著名的哲学家对祂作了如此显著的见证——将一部用希伯来语写成的法律赐给了祂的希伯来人民，这法律并非隐晦不明，而是如今已公布于万邦，在这法律中记载着：\BibleQuote{凡祭献给任何神祇的，除非只祭献给主，他应被彻底消灭。}何必在法律或先知书中寻找同一事的更多证据呢？不必说\Emph{寻求}，因为经文既不少也不难找到；但是，何必收集并应用那些密集而明显、并且藉之如白昼般清楚表明只能向至高而真天主奉献祭献的论据呢？这里有一句简短而果断、甚至带有威胁而确实真实的言语，出自那位我们对手中最有智慧的人所高度赞扬的天主。愿这句话被听从、畏惧、遵行，使没有不顺从的灵魂被剪除。\BibleQuote{凡祭献的，}祂说，不是因为祂需要什么，而是因为我们应该成为祂的产业。因此，希伯来圣经中的圣咏作者歌唱说：\BibleQuote{我曾对主说：祢是我的天主，因为祢不需要我的美物。}因为我们自己，作为祂的圣城，乃是祂最高贵、最相称的祭献，而这正是我们在我们的祭献中所庆祝的奥迹——这些祭献为信友们所熟知，正如我们在前几卷书中已经解释过的。因为藉着先知，天主的神谕宣告了犹太人作为将来之事的影象所奉献的祭献将要终止，并且万邦从日出之地到日落之地，将奉献一种祭献。从这些我们现已看到应验的神谕中，我们选取了那些似乎适合本书宗旨的内容。因此，凡没有这种义德的地方——即独一至高天主按其恩宠统治顺服的圣城，使它只向祂奉献祭献；并且在这顺服的圣城的所有公民中，灵魂按正当秩序统治身体，理性统治私欲偏情，以致如同每个义人一样，义人的团体和人民也因以爱德运作的信德而生活，即人爱天主当如何被爱，并爱近人如己——我说，凡没有这种义德的地方，就没有因共同承认正义和共同利益而结合的团体。但如果没有这个，就没有人民，如果我们的定义是正确的话，因此也就没有共和国；因为没有人民的地方，就不可能有共和国。\par

% structure: p0063 source=h2 target=subsection reason=relative-heading-rank; base=h2

\subsection{第二十四章——必须给予人民和共和国的定义，以辩护罗马人和其他王国采用这些称号。}

但如果我们摒弃这个关于民族的定义，而另取一个定义，说一群有理性的存在者因对所爱之物的共同认同而结合在一起则称为民族，那么，要发现任何民族的特征，我们只需观察他们所爱之物。然而，无论他们爱什么，只要是一群有理性的存在者，而非野兽，并且因对爱之物的认同而结合在一起，就有理由称之为一个民族；并且，根据他们因更高利益结合的程度，民族就更优越；因较低利益结合，则更卑微。按照我们这一定义，罗马人民是一个民族，其善无疑便是公共福祉或共和国。但罗马人民在其早期及后来的品味如何，以及它如何堕落为血腥的叛乱，继而导致社会战争和内战，从而使构成民族健康的和谐之链断裂或腐烂，历史已经表明，并且我在前面的书中已详细叙述。然而，我并不因此就说它不是一个民族，或它的治理不是一个共和国，只要仍有这样一群有理性的存在者因对爱之物的共同认同而结合在一起。但我对这一民族和这一共和国所说的话，必须理解为同样适用于雅典人或任何希腊城邦、埃及人、早期的亚述巴比伦，以及任何其他大小民族，只要它们曾有过公共政府。因为，一般而言，那不服从天主的命令（即只向祂献祭）的不虔之城，因此不能使灵魂对身体有适当的统治，也不能使理性对恶行有正当的权威，它缺乏真正义。\par

% structure: p0065 source=h2 target=subsection reason=relative-heading-rank; base=h2

\subsection{第二十五章 哪里没有真宗教，哪里就没有真德性。}

因为即使灵魂似乎极好地统治身体，理性似乎统治恶行，但如果灵魂和理性本身不服从天主，正如天主命令他们事奉祂那样，它们就对身体和恶行没有适当的权威。因为那不认识真天主，并且非但不服从祂的权威，反而献身于最邪恶的魔鬼的腐化影响的思想，能是什么样的身体和恶行的主母呢？正因如此，那些它自认为拥有、并借此约束身体和恶行以获得并保持其所欲之事的德性，只要在这些事上没有指向天主，就宁可说是恶行而非德性。因为，虽然有些人认为那些仅以自身为指向、且仅为其本身而被欲求的德性，仍是真实而真正的德性，但事实是，即使那时它们也充满骄傲，因此应被视为恶行而非德性。因为正如赋予肉体生命者并非来自肉体，而是高于肉体，同样，赋予人幸福生命者并非来自人，而是高于人的某物；我对人所说的，对一切天上的能力和德性亦然。\par

% structure: p0067 source=h2 target=subsection reason=relative-heading-rank; base=h2

\subsection{第二十六章 论远离天主之民族所享有的和平，以及天主之民在朝圣时期对此和平的运用}

因此，正如肉体的生命是灵魂，同样，人的幸福生命是天主，关于祂，希伯来人的圣经说：\BibleQuote{以主为天主的人民是有福的。}因此，那远离天主的人民是不幸的。然而，这人民也有它自己的和平，这和平不应被轻视，尽管它最终不会享有这和平，因为它在这和平终结之前没有善用它。但我们的利益是它在此期间在此生享有这和平；因为只要这两座城混杂在一起，我们也享有巴比伦的和平。因为天主之民从巴比伦被释放出来，以至于暂时寄居在它的陪伴中。因此，宗徒也劝勉教会为君王和掌权者祈祷，提出的理由是：\BibleQuote{我们可以在一切虔诚和爱德中，安静而平稳地生活。}而先知耶肋米亚在预言将要降临在古代天主之民身上的被掳，并给予他们顺从前往巴比伦从而事奉他们的天主的命令时，也劝告他们为巴比伦祈祷，说：\BibleQuote{在它的和平中你们将有和平，}\BibleRef{耶 29:7}——这现世和平是善人和恶人共同享有的。\par

% structure: p0069 source=h2 target=subsection reason=relative-heading-rank; base=h2

\subsection{第二十七章 论事奉天主者的和平在此生不能达致圆满}

然而，那为我们所独有的平安，如今我们借着信德与天主一同享有，将来还要借着亲眼看见与他永远享有。但我们在今世所享有的平安，无论是众人共有的还是我们独有的，与其说是幸福的实在享受，不如说是我们痛苦的慰藉。我们的义德本身，虽然就它指向真正的善而言是真的，但在今世却仍属于那样的性质：它更在于罪的赦免，而不在于德行的成全。请看整个天主之城在朝圣状态中的祈祷，因为它借着所有成员的嘴向天主呼求：\BibleQuote{免我们的债，如同我们免了人的债。}\BibleRef{玛 6:12}这祈祷并非为那些信德\BibleQuote{没有行为，是死的}\BibleRef{雅 2:17}的人有效，而是为那些信德\BibleQuote{借着爱德行事}\BibleRef{迦 5:6}的人有效。因为理性虽然服从天主，却仍\BibleQuote{被可朽坏的身体压住，}\BibleRef{智 9:15}只要它处于这必死的状态，它就不能完全掌管恶习，因此义人需要这祈祷。因为虽然它行使主权，恶习却不会毫无反抗地屈服。因为无论一个人如何善守战斗，无论他如何彻底地制伏了这些仇敌，总会有某种邪恶的东西潜入，如果它没有在行动中立时表现出来，就会从嘴唇滑出，或潜入思想中；因此，只要他还在与恶习争战，他的平安就不是完全的。因为他与那些抵抗者进行着一场胜负未卜的战斗，而他对于被击败者的胜利并不稳固，而是充满焦虑和努力。因此，在这些诱惑之中——关于这一切，在神圣的训谕中概括地说：\BibleQuote{人在世上岂不是一个诱惑吗？}\BibleRef{约 7:1}——除了骄傲的人，谁敢自认为他生活得如此之好，以至于不需要向天主说：\BibleQuote{免我们的债？}这样的人并不伟大，而是被虚荣所肿胀和膨胀，并正当地被那丰盛地赐恩宠给谦卑者的那一位所抵抗。因此经上说：\Quote{天主拒绝骄傲人，却赐恩宠给谦卑人。}义人的义德就在于：他使自己服从天主，使身体服从灵魂，使恶习——即使它们反叛——服从理性，理性要么击败它们，至少也抵抗它们；并且他也祈求天主赐予恩宠以尽本分，以及赦免他的罪，并为所领受的一切祝福感谢天主。但是，在最终的平安中——我们所有的义德都指向它，并为了它而持守——由于我们的本性将享有健全的不朽和不腐，将不再有恶习，我们将不再经历来自自身或他人的抵抗，理性就不再需要统治不再存在的恶习，而是天主统治人，灵魂统治身体，带着与脱离束缚的生命之福乐相宜的甜蜜和轻便。并且这状态将是永恒的，我们将确知它的永恒；因此，这福乐的平安和这平安的福乐将是至善。\par

% structure: p0071 source=h2 target=subsection reason=relative-heading-rank; base=h2

\subsection{第二十八章 恶人的结局。}

但是，另一方面，那些不属于这天主之城的人将承受永苦，这永苦也称为第二次死亡，因为那时灵魂将与它的生命——天主——分离，因此不能说是活着，而身体将遭受永久的痛苦。因此，这第二次死亡将更为严重，因为没有死亡能终止它。但既然战争与和平相对，如同痛苦与幸福、生命与死亡相对，那么，有理由问：在恶人的结局中，能够找到怎样的战争，来与被称为义人结局的平安相对应呢？提出这问题的人只需观察战争中什么是有害和毁灭性的，他就会发现，那无非是事物之间的相互对立和冲突。他还能构想出比这更痛苦、更惨烈的战争吗？在这战争中，意志与情欲如此对立，情欲与意志也如此对立，以至于它们的敌意永远不能因任何一方的胜利而终止；痛苦的暴力与身体的本性如此冲突，以至于谁也不向谁让步。因为在此生中，当这种冲突发生时，要么痛苦获胜，死亡驱散痛苦的感觉；要么本性获胜，健康驱散痛苦。但在来世，痛苦持续着以便折磨，本性持续着以便感受痛苦；两者都不会消失，以免惩罚也停止。既然人是通过最后审判而达到这些结局——善人达到至善，恶人达到至恶——我将在下一卷中讨论这审判。\par

\SourceRecord{Translated by Marcus Dods.；Nicene and Post-Nicene Fathers, First Series；Vol. 2.；Edited by Philip Schaff.；Buffalo, NY: Christian Literature Publishing Co.,；1887.；<http://www.newadvent.org/fathers/120119.htm>.}

% structure: p0001 source=h1 target=section reason=page-title

\section{天主之城（第二十卷）}

论末次审判，以及旧约与新约中关于审判的宣告。\par

% structure: p0003 source=h2 target=subsection reason=relative-heading-rank; base=h2

\subsection{第一章——尽管天主常常审判，但将本书注意力集中于祂末次审判仍是合理的。}

我倚赖天主的恩宠，打算谈论祂末次审判的日子，并以此反驳不虔敬和不信的人；首先，我们必须如同奠定建筑物的根基一般，将神圣的宣告放置其中。那些不信这些宣告的人，尽其所能用自己的虚假和虚幻的诡辩反对它们，或声称从圣经引用的内容另有含义，或全然否认那是天主的言语。因为我认为，任何理解所写内容并相信它是由至高而真实的天主通过圣人们传达的人，都不会拒绝屈服并同意这些宣告，无论他是在口头上承认同意，还是因为某种邪恶影响而羞于或害怕这样做；甚至，以一种近乎疯狂的固执，极力捍卫他所知道和相信是虚假的东西，反对他所知道和相信是真实的东西。\par

因此，整个真实天主的教会所持守并宣认为信经的，即基督将从天降来审判生者死者，我们称之为天主审判的末日本身，或末次时间。因为我们不知道这审判会持续多少天；但凡阅读圣经的人，即使疏忽大意，也无需被告知，其中\Quote{日}通常用以表示\Quote{时间}。当我们谈及天主审判的日子时，加上\Quote{末次}或\Quote{最后}一词，原因在于：即使现在，天主也在审判，并且从人类历史之初就已在审判——将犯下如此大罪的最初之人逐出乐园，并排除在生命树之外。是的，祂当然也在施行审判，当祂没有宽恕犯罪的天使时，他们的首领因嫉妒而跌倒，并在自己跌倒后诱惑了人类。魔鬼和人类的生活——一个在空中，一个在地上——充满了悲惨、灾难和错误，这并非没有天主深奥而公正的审判。即使没有人犯罪，也只有藉着天主良善而公正的审判，整个理性受造物才能通过坚忍地依附其主而在永恒的福乐中得以保持。祂不仅整体地审判，因那些族类的原罪而判处魔鬼族类及人类族类遭受悲惨，祂也审判个人自愿和自主的行为。因为连魔鬼也祈求不要受折磨，\BibleRef{玛 8:29}，这表明若无不当，它们本可依其功过或得宽免或受折磨。人类因自己的罪受天主惩罚，常常是可见的，总是隐密的，或在此生，或死后，尽管没有人能藉助神圣帮助之外的行为而行义；也没有人或魔鬼能在神圣而最公正之审判的许可之外行不义。因为正如宗徒所说：\Quote{天主没有不义；}\BibleRef{罗 9:14}又如在别处说：\Quote{祂的审判深不可测，祂的道路难以寻踪。}\BibleRef{罗 11:33}因此，在本卷中，我将照天主许可，不谈那些最初的审判，也不谈这些中间的审判，而只论末次审判，即基督要从天降来审判生者死者之时。因为那日子才正当地被称为审判之日，因为在那时，不再有空间留给愚昧的疑问——为何这恶人幸福而那义人痛苦。在那一天，真正而圆满的幸福必只属于善人，而应得的最高悲惨必是恶人且唯独恶人的份。\par

% structure: p0006 source=h2 target=subsection reason=relative-heading-rank; base=h2

\subsection{第二章——在人类事务交错的网络里，天主的审判虽存在，却无法辨识。}

在此現世，我們學習以平靜的心承受連善人也會遭遇的災禍，並輕視連惡人也享有的福樂。因此，即使在生活中天主的正義並不顯明，祂的教導仍是有益的。因為我們不知道天主憑什麼審判，使這個善人貧窮而那個惡人富有；為什麼那個根據我們的意見，理應為其放蕩生活而痛苦的人卻享樂，而悲傷卻緊隨著那個值得讚賞的生活使我們以為他應該快樂的人；為什麼無辜者在法庭上不僅未得昭雪，甚至被定罪，或因法官的不義而受冤屈，或因虛假證據而被壓倒，而有罪的對手卻不僅免於刑罰，甚至其訴求得到承認；為什麼不虔誠的人身體健康，而虔誠的人卻在病中憔悴；為什麼強壯的惡人體格最好，而那些連一句話都不能傷害任何人的人卻從嬰兒時期就受複雜疾病的折磨；為什麼對社會有用的人過早死亡，而那些似乎本不該出生的人卻壽命異常長久；為什麼罪惡累累的人被榮譽加冕，而無可指責的人卻被埋沒在忽視的黑暗中。但誰能收集或列舉所有這類的對比呢？如果在這生命中——如同聖詠作者所說：\BibleQuote{人好像虛幻，他的年日如影飛逝}——這種反常情況是統一的，只有惡人獲得地上短暫的繁榮，而只有善人遭受其災禍，那麼這可以歸結於天主公正甚至仁慈的審判。我們可以設想，那些註定得不到構成人類幸福之永恆福樂的人，要么被短暫的福樂所迷惑，作為他們邪惡的正當報酬，要么因天主的仁慈而得到安慰；而那些註定不會遭受永恆痛苦的人，則因其罪孽而受暫時的懲罰，或被激勵去追求更高的德行。但如今，既然我們不僅看到善人遭遇生活中的災禍而惡人享受生活中的福樂——這似乎不公——而且邪惡也常降臨惡人，福樂也意外臨到善人，正因如此，天主的審判是不可測量的，祂的道路是難以追尋的。因此，雖然我們不知道憑什麼審判，這些事由天主——祂擁有最高德性、最高智慧、最高正義，沒有軟弱、沒有魯莽、沒有不義——做或允許做的，但對我們有益的教訓是：要輕視那些善人和惡人都同樣會有的好壞事物，要渴望只屬於善人的善事，逃避只屬於惡人的惡事。但當我們到達那審判——它的日期被特別稱為審判日，有時也稱為主的日子——我們將認識到所有天主審判的正義，不僅是那時將宣判的，而且包括從起初就生效或在那時之前可能生效的一切審判。在那一天，我們也將認識到，現世中那麼多——幾乎所有——天主正義的審判，如何躲避人類感官或洞察力的審查，儘管在這件事上，虔誠的心靈並非不明白隱藏的事是正義的。\par

% structure: p0008 source=h2 target=subsection reason=relative-heading-rank; base=h2

\subsection{第三章——\Person{所罗门}在\BibleBook{}{训道篇}中关于善人和恶人所遭遇相同之事所说的话。}

\Person{撒罗满}，以色列最智慧的君王，在\Place{耶路撒冷}作王，开始了一本称为\BibleBook{}{训道篇}的书，犹太人将其列于圣经正典中：\Quote{虚而又虚，训道者说：虚而又虚，万事皆虚。人在太阳之下所操劳的一切，对他有何益处？}接着，他以这句话为主题，列举了今生的灾难和迷惑，以及现世的易变本质，其中没有什么是实在的、持久的。他在太阳之下的种种虚幻中，还哀叹这一点：虽然智慧胜于愚昧，如同光明胜于黑暗，智者的眼睛在头上，愚者却在黑暗中行走，见：\BibleRef{训 2:13}，但众人遭遇相同，即在此生阳光之下，无疑是指我们所见善恶之人同样遭遇的那些灾祸。他还说，善人如同作恶者般遭受今生的祸患，恶人却如同善人般享受今生的福乐。\Quote{在世上有一件虚幻的事，就是义人所遭遇的，反照恶人所行的；又有恶人所遭遇的，反照义人所行的。我说，这也是虚幻。}见：\BibleRef{训 8:14}这位最智慧的人用了整卷书来充分揭露这种虚幻，显然目的只在于使我们渴望那种生命，其中在太阳之下没有虚幻，而是在创造太阳的那一位之下有真理。那么，在这种虚幻中，人变得如同虚幻，岂不是按天主公义公正的审判而注定要消逝吗？但在这些虚幻的日子里，他是抗拒还是顺服真理，是缺乏真虔诚还是参与其中，这有很大的不同——不同之处不在于获得这短暂虚幻生命的福乐或躲避灾祸，而是关系到未来的审判，那审判将使善人得到永久的、不可剥夺的善事，恶人得到恶事。最后，这位智慧人在他的书结束时说：\Quote{应敬畏天主，遵守祂的诫命：因为这是人的一切。因为天主必审问一切的作为，无论善恶，连同每一个被轻视的人。}见：\BibleRef{训 12:13}有什么比这更真实、更简洁、更有益的宣告呢？\Quote{敬畏天主，遵守祂的诫命：因为这是人的一切。}因为凡是真实存在的人，就是遵守天主诫命的人；不这样做的人，就是虚无。因为只要他仍处于虚幻的相似之中，他就没有在真理的肖像中更新。\Quote{因为天主必审问一切的作为}——即人在今生所做的一切——\Quote{无论善恶，连同每一个被轻视的人}——即每一个在此世看似可鄙、因而被忽视的人；因为天主看见他，并不轻视他，也不在审判中忽略他。\par

% structure: p0010 source=h2 target=subsection reason=relative-heading-rank; base=h2

\subsection{第四章——论证明末次审判的证据，将先取自新约，再取自旧约。}

那么，我提出证明天主这末次审判的经文，应首先取自新约，然后取自旧约。因为虽然旧约在时间上优先，但新约在内在价值上有优先权；旧约扮演新约的先驱角色。因此我们将首先引用新约的章句，并以旧约的引文证实之。旧约包含法律和先知，新约包含福音和宗徒书信。如今宗徒说：\Quote{因为法律只能使人认识罪过。但是，天主的正义，在法律之外已显示出来，有法律和先知作证；天主的正义，因对耶稣基督的信德，赐给凡信的人。}见：\BibleRef{罗 3:20}这属于新约的天主的正义，在旧约的书中，即法律和先知中，有证据。那么，我首先陈述案情，然后传唤证人。这个次序是耶稣基督亲自指示我们要遵守的，祂说：\Quote{凡成为天国门徒的经师，就像一个家主，从他的宝库里拿出新的和旧的。}见：\BibleRef{玛 13:52}祂没有说\Quote{旧的和新的}，如果祂不打算遵循价值的次序而非时间的次序，祂肯定会那样说。\par

% structure: p0012 source=h2 target=subsection reason=relative-heading-rank; base=h2

\subsection{第五章——救主声明在世界末了将有神圣审判的章节。}

救主亲自责备那些祂行过许多异能却未信的城市，并将它们与外邦城市作不利对比时说：\Quote{但是我告诉你们：在审判的日子，\Place{提洛}和\Place{漆冬}所受的惩罚，比你们所受的还要容易忍受。}见：\BibleRef{玛 11:22}稍后祂又说：\Quote{我实在告诉你们：在审判的日子，\Place{索多玛}地方所受的惩罚，比你们所受的还要容易忍受。}见：\BibleRef{玛 11:24}这里祂最明确地预言审判的日子将要来临。在另一处祂说：\Quote{\Place{尼尼微}人在审判时，将同这一代人起来，定他们的罪，因为他们因\Person{约纳}的宣讲而悔改了；看，这里有一位大于\Person{约纳}的！南方的女王在审判时，将同这一代人起来，定他们的罪，因为她从地极而来，听\Person{撒罗满}的智慧；看，这里有一位大于\Person{撒罗满}的！}见：\BibleRef{玛 12:41}从这段经文我们学到两件事：审判将要发生，且是在死人复活时发生。因为当祂提到尼尼微人和南方的女王时，祂显然是指死去的人，然而祂说他们将在审判的日子起来。祂没有说\Quote{他们要定罪}，好像他们自己将是审判者，而是因为与他们相比，其他人将被公正地定罪。\par

同样，在另一段经文中，祂论及善人与恶人在现世中的混杂及未来的分离——这分离将在公审判之日完成——祂引用了一个撒在田里的麦子和莠子的比喻，并向祂的门徒这样解释：\BibleQuote{那撒好种子的就是人子，}等等。的确，在这里祂没有明说审判或公审判之日，但通过描述情节更清楚地表明了它，并预言它将在世界终末发生。\par

同样，祂对门徒说：\BibleQuote{我实在告诉你们，你们这些跟从我的人，在重生的时代，当人子坐在祂光荣的宝座上时，你们也要坐在十二个宝座上，审判以色列十二支派。}\BibleRef{玛 19:28}由此我们得知耶稣将与祂的门徒一同审判。因此，祂在别处对犹太人说：\BibleQuote{我若仗贝耳则步驱魔，你们的子弟仗谁驱魔？因此，他们将是你们的审判者。}\BibleRef{玛 12:27}我们也不应以为只有十二个人将与祂一同审判，尽管祂说他们将坐在十二个宝座上；因为数字十二象征了那些将要审判者的全数。数字七（通常象征完整）的两部分，即四和三，相乘得十二。因为四乘三或三乘四是十二。数字十二还有其他含义。若不是对十二宝座的这种正确解释，那么既然我们读到玛弟亚被立为宗徒取代叛徒犹大，宗徒保禄尽管比他们众人更劳苦\BibleRef{格前 15:10}，却不应有审判的宝座；但他明确地把自己包括在审判者之中，当他说：\BibleQuote{你们不知道我们要审判天使吗？}\BibleRef{格前 6:3}在将数字十二应用于那些将要受审判者时，也要遵守同样的规则。因为虽然经上说\BibleQuote{审判以色列十二支派，}但作为第十三支派的肋未支派并不因此免于审判，审判也不仅针对以色列而不针对其他民族。而\BibleQuote{在重生的时代}这句话，祂无疑是指死者的复活；因为我们的肉躯将因不朽而重生，如同我们的灵魂因信德而重生一样。\par

我略过许多段落，因为尽管它们似乎指向最后的审判，但仔细考察后发现它们含糊不清，或者暗示的是其他事件——或是救主在祂教会中持续来临，即在祂的肢体中，逐渐地、一块一块地来临，因为整个教会是祂的身体，或是地上耶路撒冷的毁灭。因为当祂谈到这件事时，祂也常常使用适用于世界终末和那最后大审判之日的语言，以至于若不将三位福音作者——玛窦、马尔谷和路加——关于这主题的所有相应段落加以比较，就无法区分这两件事——因为有些事一位福音作者说得较为隐晦，另一位说得较为清楚——这样就能看出哪些事是指同一事件。我曾在写给可敬的萨洛纳主教赫西基乌斯的一封信中费力地做了这件事，那封信题为\WorkTitle{论世界终末}。\par

我现在要从玛窦福音中引用一段经文，它谈到基督最有效且最终的审判将善人与恶人分开：\BibleQuote{当人子在自己的光荣中来临……}祂说，\BibleQuote{那时祂也要对左边的人说：\InnerQuote{可咒骂的，离开我，到那给魔鬼及其使者预备的永火里去。}}然后祂同样向恶人列举他们没有做的事，但祂曾说过右边的人做了这些事。当他们问何时见过祂需要这些事物时，祂回答说，既然他们没有给祂这些兄弟中最小的一位做，就是没有给祂做，并以这些话结束祂的演讲：\BibleQuote{这些人要进入永罚，而义人却要进入永生。}而且，福音作者若望最清楚地表明祂预言了审判将在死者复活时发生。因为祂说过：\BibleQuote{父不审判任何人，但将一切审判交给了子，为叫众人尊敬子如同尊敬父；不尊敬子的，就是不尊敬派遣祂的父；}祂紧接着又说：\BibleQuote{我实实在在告诉你们：那听我的话，并相信派遣我来的那一位，就有永生，不会遭受审判；但已从死亡进入生命。}\BibleRef{若 5:22}这里祂说信祂的人不会进入审判。那么，他们如何通过审判与恶人分离，并站在祂右边呢？除非此处的审判是指定罪？因为那些听祂的话并相信派遣祂的那一位的人，不会进入审判（在此意义上）。\par

% structure: p0018 source=h2 target=subsection reason=relative-heading-rank; base=h2

\subsection{第六章　第一次复活是什么，第二次复活是什么。}

此后祂又说：\BibleQuote{我实实在在告诉你们：时候要来，且现在就是，死者要听见天主子的声音；听见的人就必活着。正如圣父在自己内有生命，照样祂也赐给圣子在自己内有生命。}\BibleRef{若 5:25} 祂还未论及第二次复活，即身体的复活，那是在末日的复活，而是论及第一次复活，即现在的复活。祂为了区分这一点，说：\BibleQuote{时候要来，且现在就是。}这复活不是关乎身体，而是关乎灵魂。因为灵魂也有自己的死亡，就是邪恶与罪恶，这些灵魂的死者正是祂所说的：\BibleQuote{让死人埋葬他们的死人}\BibleRef{玛 8:22}——即让那些灵魂死了的人埋葬那些身体死了的人。正是论及这些死者——即不虔敬与邪恶中的死者——祂说：\BibleQuote{时候要来，且现在就是，死者要听见天主子的声音；听见的人就必活着。}\BibleQuote{听见的人}即那些听从、相信并坚持到底的人。这里没有区分善人与恶人。因为对所有人来说，听见祂的声音而活着，从邪恶之死进入虔敬之生，都是好的。关于这种死亡，圣保禄宗徒说：\BibleQuote{所以众人都死了，祂为众人死了，为使活着的人不再为自己活，而为主死了又复活的那位活。}\BibleRef{格后 5:14} 因此，所有人毫无例外地都死在罪中，无论是原罪或本罪，无知之罪或明知故犯之罪；为了所有这些死者，唯独那一位活着的人死了——即那位毫无罪愆者——好让那些因罪过得赦免而活着的人不再为自己活，而为那替众人死的主活，祂为我们的罪死了，并为我们的成义而复活，使我们信靠那称不虔敬者为义的天主，从不虔敬中被称义，或从死亡中被苏醒，得以达到那现存的第一次复活。因为在这第一次复活中，只有那些将永远蒙福的人有份；但在第二次复活中——祂随即要论及——正如我们将会知道的，所有人都有份，无论有福者还是受祸者。第一次是仁慈的复活，第二次是审判的复活。因此圣咏上记载：\BibleQuote{我要歌颂仁慈与审判；上主，我要向祢歌颂。}\par

祂接着论及这审判时说：\BibleQuote{并且赐给祂行审判的权柄，因为祂是人子。}这里祂表明祂要在那曾受审判的肉身中来临审判。因为祂说\BibleQuote{因为祂是人子}正是为表明这一点。接着是我们所关注的话：\BibleQuote{你们不要惊奇这事；时候要到，凡在坟墓里的，都要听见祂的声音，而出来：行过善的，复活进入生命；行过恶的，复活进入审判。}\BibleRef{若 5:28} 此处祂使用的\Quote{审判}一词与前面略早时所使用的意义相同，那时祂说：\BibleQuote{那听见我的话，而相信派遣我来者的，就有永生，并且不受\Emph{审判}，而已出死入生}——即由于在第一次复活中有份，由此在现世完成了从死亡到生命的过渡，他就不进入定罪，祂称定罪为审判，正如祂在另一处说：\BibleQuote{行过恶的，复活进入审判}，即进入定罪。因此，谁若不愿在第二次复活中被定罪，就让他先在第一次复活中复活吧。因为\BibleQuote{时候要来，且现在就是，死者要听见天主子的声音；听见的人就必活着}，即他们不会进入称为第二次死亡的定罪；那些在第一次（属灵）复活中没有复活的人，在第二次（身体的）复活之后，将被投入第二次死亡。因为\BibleQuote{时候要到}（但这里祂没有说\BibleQuote{且现在就是}，因为那将在世界末日的最后最大审判中到来）\BibleQuote{凡在坟墓里的，都要听见祂的声音，而出来。}祂不像在第一次复活中那样说：\BibleQuote{听见的人就必活着。}因为并非所有人都将活着——至少是那唯一应被称为生命的生命，因为唯独它是幸福的。他们必须具有某种生命，才能听见并从坟墓中带着复活的躯体出来。至于为何不是所有人都活着，祂以随后的话教导说：\BibleQuote{行过善的，复活进入生命}——这些人是将活着的人；\BibleQuote{行过恶的，复活进入审判}——这些人是不会活着的人，因为他们将在第二次死亡中死去。他们行恶是因为他们的生命曾是邪恶的；他们的生命邪恶是因为它没有在现存的第一次（属灵）复活中得到更新，或者因为他们没有在更新后的生命中坚持到底。因此，正如有两种再生——我已提及过——一种是在现世通过圣洗圣事凭信德发生的；另一种是在身体方面，将通过伟大的最后审判在不朽与不灭中成就——同样也有两种复活：一种是第一次的属灵复活，在现世进行，保护我们免于进入第二次死亡；另一种是第二次复活，不发生在现在，而是在世界末日，是身体的复活而非灵魂的复活，并通过最后审判将一些人送入第二次死亡，将另一些人送入那没有死亡的生命。\par

% structure: p0021 source=h2 target=subsection reason=relative-heading-rank; base=h2

\subsection{关于\WorkTitle{若望默示录}中所记载的两个复活与一千年，以及关于这些事可合理持守的观点。}

福音作者\Person{若望}在名为\WorkTitle{默示录}的书中论述了这两种复活，但他的论述方式使一些基督徒不理解其中的第一次复活，从而将这段经文解释成荒谬的幻想。因为宗徒\Person{若望}在上述书中说：\BibleQuote{我看见一位天使从天降下……那参与第一次复活的是有福的、圣洁的：第二次的死亡对他没有权柄；他们要做天主和基督的司祭，并要与祂一同为王\OriginalTerm{一千年}{a thousand years}。}那些依据这段经文怀疑第一次复活是未来和身体的复活的人，特别受到千年之数的影响，认为圣人在这段时间内应当享受一种安息日般的休息，即从人类被创造以来六千年劳苦之后的圣洁闲暇（人类因其大罪从乐园的幸福中被逐入今世生活的苦难），因此，正如经上所载：\BibleQuote{主看一日如千年，千年如一日，}\BibleRef{伯后 3:8}在六千年之后，如同六天之后，应在接下来的千年中有一种第七日的安息日；圣人复活正是为此目的，即庆祝这个安息日。如果人们相信圣人在那个安息日中的喜乐是属灵的，并来自天主的同在，那么这个意见并非不可接受；因为我本人也曾持有这种看法。但是，他们声称那时复活的人将享受无节制的肉欲宴席，提供大量的饮食，不仅震惊节制者的感受，甚至超越了可信度的尺度，这样的断言只有肉性的人才能相信。这些相信者被属灵的人称为\OriginalTerm{千禧年主义者}{Chiliasts}，我们可以直译为\OriginalTerm{千禧年派}{Millenarians}。要逐点驳斥这些观点是一件繁琐的事；我们更倾向于说明应该如何理解这段经文。\par

主耶稣基督亲自说：\BibleQuote{没有人能进入壮士的家，抢夺他的家具，除非先捆绑那壮士}——以壮士指魔鬼，因为他有能力掳掠人类；以他要抢夺的家具指那些曾被魔鬼困在各种罪恶中，但将要相信祂的人。正是为了捆绑这强者，宗徒在\WorkTitle{默示录}中看见\BibleQuote{一位天使从天降下，手里拿着无底坑的钥匙和一条链子。他捉住}，他说，\BibleQuote{那龙，就是古蛇，被称为魔鬼和撒殚的，把他捆绑\OriginalTerm{一千年}{a thousand years}}——即抑制并约束了他的能力，使他不能诱惑和占有那些将要被释放的人。以我之见，对于这一千年可以有两种理解：要么是因为这些事发生在第六千年（即第六个千年期，其后期正在过去），如同在第六天，之后将有一个没有黄昏的安息日，即圣人无尽的安息，因此他以整体之名称呼部分，将千年期的最后部分（即世界结束前尚要经过的部分）称为一千年；要么他用一千年作为今世全部时段的等价物，以完美的数目来标志时间的圆满。因为一千是十的立方。十个十是一百，即平面上的平方。但要使这个平面具有高度而成为立方体，再将一百乘以十，就得出一千。此外，如果一百有时用来表示整体，如主应许那舍弃一切跟随祂的人时所说的：\BibleQuote{今世要得百倍；}\BibleRef{玛 19:29}关于此，宗徒仿佛作了解释，他说：\BibleQuote{似乎一无所有，却拥有万有；}\BibleRef{格后 6:10}——因为自古就曾说\Quote{整个世界都是信徒的财富}——那么，一千比一百更有理由用来表示整体，因为一千是立方，而一百只是平方。同理，我们解释圣咏的话：\BibleQuote{祂永远记念祂的盟约，祂所吩咐的话直到千代，}最好理解为\InnerQuote{直到万代}。\par

\BibleQuote{祂把他投入深渊}——即\Emph{即}把魔鬼投入深渊。所谓\Quote{深渊}，是指无数恶人的众多，他们心中对天主的教会怀着深不可测的恶意；并非魔鬼原先不在那里，而是说他被投入那里，因为当他无法伤害信徒时，他就更完全地占据不敬虔的人。因为那不但远离天主，而且无故仇恨事奉天主的人，被魔鬼占据得更深。\BibleQuote{祂把他关闭起来，并且加上封印，使他在一千年完成以前，不再迷惑万民。}\BibleQuote{祂把他关闭起来}——\Emph{即}，禁止他出去，禁止他做被禁止的事。再加上\BibleQuote{加上封印}，在我看来意味着要保密哪些属于魔鬼的党羽，哪些不属于。因为在这个世界上这是秘密，因为我们不知道那看似站立的人是否会跌倒，也不知道那看似跌倒的人是否会重新站起。但借着这禁令的锁链和监牢，魔鬼被禁止和约束，不能诱惑那些属于基督的万民——这些民他从前曾诱惑或使之臣服。因为在创世以前，天主就选择将这些民从黑暗权势中拯救出来，并迁到祂爱子的国度里，如宗徒所说。\BibleRef{哥 1:13}因为哪个基督徒不知道，他甚至在现在也迷惑万民，把他们与自己一起拖向永罚，但不是那些预定得永生的人呢？也不要因魔鬼常常迷惑那些在基督内重生、开始行走在天主道路上的人而惊慌。因为\BibleQuote{主认识属于祂的人}\BibleRef{弟后 2:19}，魔鬼不能诱惑他们中的任何人到永罚中去。因为主认识他们，是作为天主，没有什么未来的事能向祂隐藏；不像人，人只能看见当前的人（如果可以称作看见那心他看不见的人），甚至不能看见自己到了能知道自己将来会成为什么样的人的程度。因此，魔鬼被捆绑并关闭在深渊中，使他不能诱惑那些教会从中聚集的万民——这些民在教会存在之前他曾诱惑。因为经上不是说\BibleQuote{使他不再诱惑任何人}，而是说\BibleQuote{使他不再迷惑万民}——无疑是指那些教会所在的万民——\BibleQuote{直到一千年完成}——\Emph{即}，要么指由一千年组成的第六日的剩余部分，要么指直到世界末了所要经过的所有年数。\par

\BibleQuote{使他在一千年完成以前不再迷惑万民}这句话不应理解为，此后他只迷惑那些从预定教会组成的万民，而他现在被锁链和监牢约束而不得迷惑他们；而是按照圣经中经常使用的用法，正如圣咏所说：\BibleQuote{我们的眼睛仰望主我们的天主，直到祂怜悯我们}——并非指祂的仆人们的眼睛在祂怜悯他们之后就不再仰望主他们的天主。或者，毫无疑问，词序是这样：\BibleQuote{祂把他关闭起来，并且加上封印，直到一千年完成}；而插入的短语\BibleQuote{使他不再迷惑万民}不应按照其所在语序理解，而应独立地、仿佛是后来加上去的，这样整个句子可以读作：\BibleQuote{祂把他关闭起来，并且加上封印，直到一千年完成，使他不再迷惑万民}——即他之所以被关闭直到一千年完成，是为了使他不再欺骗万民。\par

% structure: p0026 source=h2 target=subsection reason=relative-heading-rank; base=h2

\subsection{第八章 论魔鬼的被缚与释放}

\BibleQuote{此后,} \Person{若望}说，\BibleQuote{他必被释放一个短时期。} 如果束缚和关闭魔鬼意味着使他不能诱惑教会，那么他的释放是否意味着恢复这种能力？绝对不是。因为教会是在世界创立以前被预定和拣选的，论到教会曾说：\BibleQuote{主认识那些属于祂的人}，她的信德永远不会被他诱惑。然而，即使当魔鬼被释放时，世上仍会有教会，正如从起初以来一直有，并且将永远有，死去的人的位置由新信徒填补。因为稍后\Person{若望}说，魔鬼被释放后，将引诱他在全世界所诱惑的列国，与教会作战，这些敌人的数目如海沙。\BibleQuote{他们上到地极，包围了圣人的营和蒙爱的城；就有火从天主那里从天降下，吞灭了他们。那诱惑他们的魔鬼被扔进硫磺火湖里，那里有兽和假先知，他们必昼夜受痛苦，直到永永远远。}\BibleRef{默 20:9} 这是关于最后审判的，但我认为现在提及是合适的，以免有人以为在魔鬼被释放的那短时期内，世上没有教会，无论是由于魔鬼找不到教会，还是用许多迫害摧毁了它。因此，魔鬼在本书所涵盖的整个时期——即从基督首次来临到世界末了、祂第二次来临——并没有被束缚，意思不是说在这称为一千年的间隔期间，他不会诱惑教会，因为即使他被释放，他也不会诱惑她。因为如果他被束缚意味着他不能或不被允许诱惑教会，那么他被释放岂不意味着他能够或被允许这样做？但願天主不允許這樣！但魔鬼的束缚是阻止他运用全部力量诱惑人，无论是用暴力强迫还是用诡计欺骗他们与他为伍。如果他在这么长时期内被允许攻击人的软弱，许多人——就是天主不愿让他们受如此试探的人——的信德就会被打倒，或者他们就不会相信；而为了避免这种情况发生，他被束缚了。\par

但当短时期来临，他必被释放。因为他要带着他自己和其天使的全部力量发怒三年六个月；与他作战的人将有能力抵抗他所有的暴力和诡计。如果他从未被释放，他的恶意力量就不会那么明显，圣城的坚固勇气也得不到更多的证明：简言之，全能者如何善用他的大恶，就不会那么显明。因为全能者并非完全使圣人远离他的诱惑，而只保护他们的内在之人——信德所在之处，使他们通过外在的诱惑在恩宠中成长。祂束缚他，使他不能自由且热切地行使恶意，阻碍或摧毁那些无数软弱之人的信德，这些人已经相信或将要相信，教会必须由他们增长并得以完成；最终祂会释放他，使天主的城看到它战胜了多么强大的敌人，以彰显其救赎者、帮助者、解救者的伟大光荣。与那时将要存在的信徒和圣人相比，我们算什么？他们将因敌人的释放而受考验，而我们在魔鬼被束缚时就与他作战，尚且极其危险。尽管也确定，即使在现今这个期间，也曾有并仍有一些基督的士兵如此智慧而坚强，如果他们能在魔鬼被释放时仍活在这必死的状态下，他们必能最智慧地防范并最耐心地忍受他所有的网罗和攻击。\par

魔鬼这样被束缚，不仅是在教会开始从犹太地以外扩展到万民中之时，而且现在一直要到世界末日他被释放为止，都一直是被束缚的。因为甚至现在，人们也无疑直到世界末日，都会从他束缚他们的不信中归信。这个壮士在每一次被抢夺其财物时就被束缚；而他被关闭的深渊，当最初被关闭时那些活着的人死去时并未终结，而是由后来出生、同样仇恨基督徒的人接续，并且直到世界末日都将接续；在他被关闭于他们盲目的心灵深处如同深渊。但是问题在于：当他被释放、用尽全力肆虐的那三年半期间，是否会有先前不信的人归信？因为那样的话，这话如何成立：\BibleQuote{谁能进入壮士的家，抢夺他的财物，除非先把那壮士捆住？} 因此这节经文似乎迫使我们相信，在那短短的时间内，将没有人加入基督徒团体，而是魔鬼要与先前已成为基督徒的人争战；并且虽然其中有些人可能被征服而背叛魔鬼，但他们不属于天主子女的预定数目。因为记载这\WorkTitle{默示录}的同一位宗徒\Person{若望}，在他的\WorkTitle{若望一书}中论到某些人，并非没有理由地说：\BibleQuote{他们从我们中间出去了，但他们并不属于我们；因为如果他们属于我们，也许就会留在我们中间。}\BibleRef{若一 2:19} 但那些小孩子将如何？因为在这些日子里，竟然没有一些出生但尚未受洗的基督徒儿童，甚至在那时期也没有出生的，这完全难以置信；如果有这样的人，我们不能相信他们的父母会找不到办法带他们去接受重生之洗。但若是如此，当魔鬼被释放时，这些财物怎能从他手中夺去呢？因为除非先捆住他，否则无人能进入壮士的家抢夺他的财物。相反，我们更应相信，在那日子里，无论背叛教会的人还是加入教会的人都将不缺；但父母为他们的孩子寻求圣洗的决心，以及那些届时初次相信的人，将如此坚定，以至于他们能战胜那甚至未被捆绑的壮士——也就是说，他们会警醒地识破他并耐心地抵挡他，尽管他施展前所未有的诡计和力量；这样，即使他未被捆绑，他们也会从他手中被夺去。然而福音的话不会不真实：\BibleQuote{谁能进入壮士的家，抢夺他的财物，除非先把那壮士捆住？} 因为根据这句真话，顺序是：先捆住壮士，然后抢夺他的财物；因为教会因远近万民的弱者和强者增长，以致它能以对神圣预言与成就之事最坚固的信德，甚至能抢夺未被捆绑的魔鬼的财物。因为我们必须承认，\BibleQuote{由于罪恶的增加，许多人的爱情冷淡了，}\BibleRef{玛 24:12} 并且那些名字没有写在生命册上的人，将大批屈服于现在被释放的魔鬼的严厉空前迫害和诡计；所以我们不能不想：不仅那时在信德上健全的人，而且一些到那时仍未信的人，都会在他们先前拒绝的信德上变得坚定，并且强大到足以战胜即使未被捆绑的魔鬼，天主的恩宠帮助他们理解圣经——圣经中，除了其他事，也预言了那他们自己看到正在来临的结局。如果这样，那么就必须说他的捆绑先于对被捆绑和被释放的他的掠夺；因为经上说：\BibleQuote{谁能进入壮士的家，抢夺他的财物，除非先把那壮士捆住？}\par

% structure: p0030 source=h2 target=subsection reason=relative-heading-rank; base=h2

\subsection{第九章 何谓圣人与基督一同为王一千年，以及它如何与永恒国度不同}

但是，当魔鬼被捆绑时，圣人们与基督一同为王，同样的一千年，以同样方式理解，即指祂第一次来临的时期。因为，暂且不计祂在末日将要说的那个王国（\BibleQuote{凡我父所祝福的，你们来承受为你们预备的王国吧！}\BibleRef{玛 25:34}），除非祂的圣人们即使现在也与祂一同为王，尽管是另一种且十分不同的方式，否则教会现在就不能被称为祂的王国或天国。因为祂对祂的圣人们说：\BibleQuote{看！我同你们天天在一起，直到今世的终结。}\BibleRef{玛 28:20} 确实是在这现世，那位通晓天国、我们已经提过的经师，从他的宝库里取出新和旧的东西。并且，那些收割者要从教会里收集祂容许与麦子一同生长直到收割的莠子，正如祂用这些话解释的：\BibleQuote{收割就是今世的终结；收割者是天使。所以，正如将莠子收集起来用火焚烧，在今世终结时也将如此。人子要派遣祂的天使，他们将从祂的王国里收集一切绊脚石。}\BibleRef{玛 13:39} 祂能指那个没有绊脚石的王国吗？那么，他们必须是从祂现世的王国，即教会中，被收集出去。所以祂说：\BibleQuote{谁若废除这些诫命中最小的一条，也这样教训人，在天国里他将称为最小的；但谁若实行并这样教训人，这人在天国里将称为大的。}\BibleRef{玛 5:19} 祂说到两者都在天国里，既包括那人不实行祂所教训的诫命——因为违背意思就是不遵守、不实行——也包括那人实行并像祂那样教训人；但祂称前者为最小的，后者为大的。祂紧接着又说：\BibleQuote{我告诉你们：除非你们的义德超过经师和法利塞人的义德}——即那些违背自己所教训的人之义德；因为关于经师和法利塞人，祂在别处说：\BibleQuote{因为他们只说而不做}\BibleRef{玛 23:3}——除非因此，你们的义德超过他们，即你们不违背，反而实行你们所教训的，\BibleQuote{你们决不能进入天国。}\BibleRef{玛 5:20} 我们必须在一个意义上理解天国，其中既有违背自己所教训的人，也有实行它的人，一个是最小的，另一个是大的；在另一个意义上理解天国，只有实行他所教训的人才能进入。因此，凡这两种人都存在的地方，就是教会现在的情况，但只有一类人存在的地方，就是教会未来的情况，那时教会内不再有恶人。所以，教会现在就是基督的王国，也是天国。相应地，甚至现在祂的圣人们就与祂一同为王，尽管不同于他们将来为王的方式；然而，尽管莠子在教会中与麦子一同生长，他们却不与祂一同为王。因为与祂一同为王的，是那些遵行宗徒所说的人：\BibleQuote{你们既然与基督一同复活了，就该追求天上的事，在那里基督坐在天主的右边。你们要思念上面的事，不要思念地上的事。}\BibleRef{哥 3:1} 关于这样的人，他也说他们的家乡在天上。\BibleRef{斐 3:20} 总之，与祂一同为王的，是那些在祂的王国里以至于他们自己就是祂的王国的人。但是，那些——姑且不说别的——虽然在祂的王国里直到今世终结时一切绊脚石被收集出去，却在此寻求自己的事，而非基督的事的人，在哪一种意义上他们是基督的王国呢？\BibleRef{斐 2:21}\par

这是关于这战斗的国度，在其中与敌人的冲突仍在进行，战争与争战的私欲进行，或当它们屈服时施加统治，直到我们进入那最和平的国度，在那里我们将没有敌人而统治；正是关于这现世生命中的第一次复活，默示录在以刚引用的言语中说话。因为，在说了魔鬼被捆绑一千年，之后被释放一小段时间之后，它继续描述教会在那时期做什么或教会中发生什么，用以下言语：\BibleQuote{我看见座位，和坐在其上的，并给予了审判。} 不应认为这指的是最后审判，而是指统治者的座位和统治者本身，教会现在由他们治理。对于\BibleQuote{审判给予了}没有比我们在以下言语中所有的更好的解释：\BibleQuote{你们在地上捆绑的，在天上也要被捆绑；你们在地上释放的，在天上也要被释放。} \BibleRef{玛 18:18} 因此宗徒说：\BibleQuote{我与判断外人有什么关系？你们不是判断内部的人吗？} \BibleRef{格前 5:12} 若望说：\BibleQuote{那些为了耶稣的见证和天主的圣言被杀害者的灵魂，}——理解他后面所说的\BibleQuote{与基督一同为王的一千年} \BibleRef{默 20:4} ——即那些尚未回归身体的殉道者的灵魂。因为虔诚亡者的灵魂并未与教会分离，教会如今就是基督的国度；否则，在参与基督身体的祭台上就不会有对他们的纪念，在危险中赶快领受祂的圣洗，以免我们未领受而离开此生，也无益处；同样，若因补赎或不良良心而与祂的身体分离，和好也无益。因为如果不因信友即使死了也是祂的肢体，这些事为何实行呢？因此，当这一千年流逝时，他们的灵魂与祂一同为王，虽尚未与身体结合。因此在同一书的另一部分我们读到：\BibleQuote{从今而后，在主内死去的死者是有福的，圣神说，使他们安息于劳苦，因为他们的功业随着他们。} \BibleRef{默 14:13} 因此，教会现在在活人和死人中开始与基督一同为王。因为，如宗徒所说：\BibleQuote{基督死了，为成为活人和死人的主。} \BibleRef{罗 14:9} 但他只提到了殉道者的灵魂，因为他们为真理争战至死的人，他们自己在死后主要为王；但是，以部分代表全体，我们理解这些话包括所有属于教会——基督的国度——的其他人。\par

至于随后的话：\Quote{若有人没有朝拜那兽和他的像，也没有在额上或手上接受他的标记}，我们必须理解为包括活人和死人。至于这兽是什么，虽然需要更仔细的研究，但将其理解为不虔敬的城市本身，以及反对信众和天主之城的无信者的团体，并不违背真信仰。\Quote{他的像}在我看来是指他的伪装，即那些自称相信却像无信者一样生活的人。因为他们假装成他们不是的样子，被称为基督徒，不是由于真正的相似，而是由于欺骗性的像。因为属于这兽的不仅有基督圣名及其最荣耀之城公然的仇敌，也有在世界终结时将从祂的国度——教会——中收集的莠子。谁不是朝拜那兽和他的像呢？如果不是那些遵行宗徒所说\BibleQuote{不要与无信者同负一轭}的人？\BibleRef{格后 6:14} 因为这样的人不朝拜，\Emph{即}不同意、不服从；他们也不接受那罪孽的印记，在额上通过信仰宣认，在手上通过实践。因此，那些摆脱这些玷污的人，无论他们仍活在这有死的肉身中，还是已死，现在都在这由一千年标示的整个期间，以适合此时的方式与基督一同为王。\par

他说：\BibleQuote{其余的人没有活过来。} 因为现在是死人要听见天主圣子声音的时候，听见的人要活过来；其余的人不会活过来。所加的\BibleQuote{直到一千年完毕}这句话，是指他们没有在应该借由从死亡过渡到生命而活过来的时间内活过来。因此，当肉身复活的日子来到时，他们将走出坟墓，不是为生命，而是为审判，即受罚，称为第二次死亡。因为凡没有活到一千年完毕的，\Emph{即}在这第一次复活进行的整个时期内——凡没有听见天主圣子的声音，并从死亡过渡到生命的人——那人必然在第二次复活，即肉身的复活中，与他的肉身一同进入第二次死亡。因为他接着说：\BibleQuote{这是第一次复活。有分于第一次复活的是有福而圣洁的，} 或体验它的人。现在体验它的人不仅从罪之死亡中复生，而且继续在这种更新了的生命中。\BibleQuote{第二次死亡在这些身上没有权柄。} 因此它在其余的人身上有权柄，就是他在上面所说\BibleQuote{其余的人直到一千年完毕还没有活过来}的人；因为在这整个被称为一千年的中间时期，无论他们怎样精力旺盛地活在身体里，他们并没有从那恶毒困锁他们的死亡中被苏醒进入生命，以便通过这种复活的生命，他们成为第一次复活的分享者，从而使第二次死亡在他们身上没有权柄。\par

% structure: p0035 source=h2 target=subsection reason=relative-heading-rank; base=h2

\subsection{第十章——应如何回答那些认为复活只关乎身体而不关乎灵魂的人。}

有些人认为复活只能就身体而言，因此他们主张（默示录中的）第一次复活是身体的复活。因为他们说，\Quote{复活}只能用于会跌倒的事物。身体在死亡中跌倒。所以，不可能有灵魂的复活，只有身体的复活。但他们如何回答那谈及灵魂复活的宗徒呢？的确，那些人对宗徒说\BibleQuote{你们既然与基督一同复活了，就该追求天上的事。}\BibleRef{哥 3:1}同样的意思他在别处用其他话语表达说：\BibleQuote{基督藉着圣父的光荣从死者中复活了，这样我们也能在新生活中行走。}\BibleRef{罗 6:4}同样，\BibleQuote{你这沉睡的，醒来，从死者中起来，基督必要光照你。}\BibleRef{弗 5:14}至于他们所说，只有跌倒的才能复活，因而他们断定复活只属于身体而不属于灵魂，因为身体跌倒；他们为何不顾这些经文：\BibleQuote{你们敬畏主的，等候祂的仁慈；不要偏离，以免跌倒；}\BibleRef{德 2:7}以及\BibleQuote{他或站立或跌倒，自有他的主人；}\BibleRef{罗 14:4}还有\BibleQuote{那自以为站立的，要小心，免得跌倒？}\BibleRef{格前 10:12}因为我想这要小心的跌倒乃是灵魂的跌倒，而非身体的。那么，如果复活属于跌倒的事物，而灵魂会跌倒，就必须承认灵魂也复活。在\Quote{第二次死亡在他们身上无权}这句话之后，又加上了\Quote{他们将成为天主和基督的司铎，并要与祂一同为王一千年；}这不仅指主教和长老——他们在教会中特别被称为司铎——而是因着那神秘的\OriginalTerm{傅油}{\GreekText{χρίσμα}}，我们称所有信徒为基督徒，同样我们也称所有人为司铎，因为他们是那位唯一大司祭的肢体。宗徒\Person{伯多禄}论他们说：\BibleQuote{圣洁的子民，王家的司铎。}\BibleRef{伯前 2:9}诚然，他用一种附带和间接的方式暗示基督是天主，说\Quote{天主和基督的司铎}，即圣父和圣子的司铎，尽管基督是在仆人的形象中、作为人子，按\Person{默基瑟德}的品位永远为司铎。但这些我们已经多次解释过了。\par

% structure: p0037 source=h2 target=subsection reason=relative-heading-rank; base=h2

\subsection{第十一章——关于歌革和玛各，当魔鬼在世界末了被释放时，他们将被他唤醒来逼迫教会。}

\BibleQuote{那一千年完了以后，撒殚将从狱中被释放，出来迷惑地上四方的万民，就是歌革和玛各，并召集他们作战，他们的数目多如海沙。}这就是他迷惑他们的目的，即召集他们作战。因为在此以前，他也常用他能持续的各种迷惑。而\Quote{他出来}的意思是，他从潜伏的仇恨中爆发为公开的迫害。因为这场迫害，在最后审判临近时，将是圣教会普世所忍受的最后一场，基督的整个城市被魔鬼的整个城市攻击，如同各自在地球上存在一样。因为他所命名的这些民族，歌革和玛各，不应理解为地球上某个地方的野蛮民族，如有些人从起首字母推论的革特人和玛撒革特人，或是其他不在罗马政府统治下的外族。因为\Person{若望}指出他们遍及全地，当他说\BibleQuote{地上四方的万民}时，他又补充说这就是歌革和玛各。我们发现这些名字的含义是：\OriginalTerm{歌革}{a roof}，\OriginalTerm{玛各}{from a roof}——好比房子和从房子里出来的人。因此，他们就是我们在其中发现魔鬼像关在\OriginalTerm{深渊}{abyss}里一样的那些民族，而魔鬼自己从他们中间出来并走出去，所以他们就是屋顶，而他来自屋顶。或者，如果我们把两个词都指民族，而不是一个指他们一个指魔鬼，那么他们两者都是屋顶，因为老仇敌目前被关在他们里面，如同被屋顶覆盖；他们将是从屋顶，因为他们从隐藏的仇恨爆发为公开的仇恨。经文说：\BibleQuote{他们在大地的广阔区域上走上来，包围了圣徒的营地和蒙爱的城，}这并不是说他们已来到或将要来到一个地方，好像圣徒的营地和蒙爱的城在某个地方；因为这营地不是别的，正是遍布全世界的基督的教会。因此，无论教会在哪里——她将在所有民族中，如\Quote{大地的广阔区域}所象征的——那里也是圣徒的营地和蒙爱的城，她将被所有敌人的猛烈迫害所包围；因为他们也将在所有民族中与她共存——也就是说，她将受到束缚、受压榨、被围困在磨难中，但不会放弃她的军事职责，这由\Quote{营地}一词所象征。\par

% structure: p0039 source=h2 target=subsection reason=relative-heading-rank; base=h2

\subsection{第十二章——从天上降下的火吞灭了他们，是否指的是对恶人的最后惩罚。}

那些话：\Quote{有火从天降下，吞灭了他们}，不应理解为那最后的惩罚，即当祂说：\Quote{可咒骂的，离开我，到永火里去}\BibleRef{玛 25:41}时所要施加的；因为那时他们将被投入火中，而非火从天降在他们身上。在此处，\Quote{从天降下的火}宜理解为圣人的坚毅，他们借此拒绝屈服于那些向他们发怒的人。因为穹苍就是\Quote{天}，这些攻击者将因穹苍的坚固而受灼热之恨的痛苦，因为他们无力将圣人拉向敌基督的阵营。这火就是要吞灭他们的火，这火是\Quote{来自天主的}；因为藉着天主的恩宠，圣人成为不可战胜的，从而折磨他们的仇敌。因为正如在善的意义上所说的：\Quote{我为祢的殿宇所消耗的热忱}；同样，在恶的意义上所说的：\Quote{热忱抓住了未受教的人民，现在火要吞灭仇敌}\BibleRef{依 26:11}。\Quote{现在}，即指不是那最后审判的火。或者，\Person{若望}藉这从天降下并吞灭他们的火，意指基督来临时要击打那些祂在地上所发现的教会迫害者，即祂用口中的气杀死敌基督的那一击\BibleRef{得后 2:8}，那么这也不是恶人最后的审判；最后的审判乃是他们要在肉身复活后所受的。\par

% structure: p0041 source=h2 target=subsection reason=relative-heading-rank; base=h2

\subsection{论迫害或敌基督的时期是否应算在千年内}

敌基督的最后迫害将持续三年半，正如我们已经说过的，并在\BibleBook{若望}{默示录}和\BibleBook{达尼尔}{先知书}中得以肯定。这段时间虽短，但有人有理由质疑：它是包含在魔鬼被束缚、圣人与基督一同为王的千年之内，还是这小段时间应加在这些年之上？因为如果我们说它们包含在千年之内，那么圣人与基督一同为王的时期比魔鬼被束缚的时期更长。因为他们要同他们的王和得胜者强有力地一同为王，即使在那最顶峰的迫害中，此时魔鬼已被释放，尽其全力向他们发怒。那么，如果魔鬼的束缚在圣人与基督为王的三年半之前就停止了，圣经如何将魔鬼的束缚和圣人的为王都定为同一千年呢？另一方面，如果我们说这短暂的迫害时期不应算作千年的一部分，而应作为一个额外的时期，那么我们确实可以解释这些话：\Quote{天主和基督的司祭要同祂一起为王一千年；及至一千年满了，撒殚就要从监牢里被释放}；因为这样它们表明圣人的为王和魔鬼的束缚同时结束，因此我们所说的迫害时期既不与圣人的为王同时，也不与撒殚的被囚同时，而应作为额外增加的时间段来计算。但这样一来，我们就不得不承认圣人在那次迫害中不与基督一同为王。然而，谁敢说祂的肢体在那一时刻，当他们最坚强、最大胆地依附于祂，并且抵抗的荣耀和殉道的冠冕随战况的激烈而越发显赫时，不与祂一同为王呢？或者，如果有人认为他们因所受的苦难而不能称为为王，那么所有先前在千年期间受过苦难的圣人都不能说在他们受苦期间与基督一同为王，因此，连这书作者\Person{若望}说他看见的那些为耶稣的见证和天主的话而被杀的灵魂，当他们遭受迫害时也没有与基督一同为王，他们自己也不是基督的国度，尽管基督那时已卓越地拥有他们。这实在是完全荒谬、应受唾弃的。但是，光荣的殉道者的得胜灵魂，既已克服并结束了所有悲痛和劳苦，并且放下了他们可朽的肢体，就一直在与基督一同为王，直到千年结束，以便以后当他们领受不朽的身体时，能与祂一同为王。因此，在这三年半中，那些为祂的见证而被杀的灵魂，无论是先前离开身体的，还是在那最后迫害中要离开的，都要与祂一同为王，直到这有死的世界终结，并进入那没有死亡的国度。这样，圣人与基督一同为王的时期比魔鬼的捆绑和监禁更长久，因为他们要在魔鬼不再被捆绑的这三年半中，与他们的王、天主子一同为王。因此，当我们读到\Quote{天主和基督的司祭要同祂一起为王一千年；及至一千年满了，魔鬼就要从监牢里被释放}时，我们应当理解：要么圣人为王的千年并不终止，尽管魔鬼的监禁终止了——这样双方都有各自的千年，即他们完整的时间，但各有不同的实际长度，圣人的国度更长，魔鬼的监禁更短——要么至少，由于三年半是很短的时间，它既不被视为从撒殚监禁的总时间中扣除，也不被视为加在圣人为王的总时间上，正如我们在\WorkTitle{第十六卷}中关于四百年这个整数所说明的，虽然实际稍多，却写作四百年；在圣经中若留意，也常有类似的表述。\par

% structure: p0043 source=h2 target=subsection reason=relative-heading-rank; base=h2

\subsection{第十四章　论魔鬼及其党羽的受罚；以及论全体死者的肉身复活和最后报应审判的概述。}

在提及这末次迫害之后，他扼要地指出了魔鬼及其所统治的城在最后审判中将要遭受的一切。因为他说：\Quote{那迷惑他们的魔鬼被投入那烧着硫磺的火湖里，就是那兽和假先知所在的地方；他们必昼夜受痛苦，直到永永远远。}我们已经说过，那兽可理解为那邪恶的城。它的假先知，或是假基督，或是我们在同一处所论及的那像或幻象。此后，他简短叙述了最后审判本身，即将在第二次或肉身复活时发生的事，正如\Person{若望}所启示的：\Quote{我看见一个白色的大宝座，和坐在上面的那位；天地都从祂面前逃避，再也找不到它们的地方了。}他没有说：\Quote{我看见一个白色的大宝座，和坐在上面的那位，天地都从祂面前逃避，}因为那时尚未发生，\Emph{即}在审判活人与死人之前；但他说他看见那位坐在宝座上，天地从祂面前逃避，但那是后来。因为当审判结束时，这天地将不复存在，将有新天新地。因为这个世界将因转化而消逝，并非彻底毁灭。因此宗徒说：\Quote{因为这世界的局面正在逝去。我愿你们无所挂虑。}\BibleRef{格前 7:31}所以消逝的是局面，而非本性。\Person{若望}在说他看见那位坐在宝座上，天地从祂面前逃避（虽然那是后来）之后，又说：\Quote{我又看见死了的人，无论大小，都站在宝座前；案卷展开了，并另有一卷展开，就是生命册；死了的人都凭着这些案卷所记载的，照他们所行的受审判。}他说案卷展开了，又有一卷；但他并未说明这卷是什么，只是说\Quote{就是生命册}。因此，他首先提及的那些案卷，我们应理解为旧约与新约的圣经，以表明天主吩咐了哪些诫命；而那生命册则表明各人遵行或未遵行哪些诫命。若将这卷书按物质理解，谁能计算它的尺寸或长度，或阅读记录每一个人全部生活的书所需的时间呢？难道会有与人数相等的天使，每人由负责他的天使宣读他的生活吗？那样就不止一卷包含所有生活的书，而是每人一卷。但我们的经文要求我们只想到一卷。它说：\Quote{另有一卷展开了。}所以我们必须将之理解为某种神圣的能力，藉此将促使每一个人回忆起自己所有的事，无论善恶，并以惊人的速度在心中审视它们，以致这知识或控告或辩护良心，从而所有人和每一个人同时受审。这神圣的能力被称为书，因为我们将如同在其中读到它使我们想起的一切。为表明谁是受审的死者（无论大小），他回到他之前省略或更确切地说是推迟的内容，说：\Quote{于是海交出其中的死者；死亡和地狱也交出其中的死者。}这当然发生在死者受审之前，但在此却被提及。所以我说，他返回到他省略的内容。但现在他遵循事件的顺序，为了展示它，他在适当的位置重复了他已经说过的关于受审的死者的话。因为他说了\Quote{于是海交出其中的死者；死亡和地狱也交出其中的死者}之后，立即加上他已经说过的话：\Quote{他们都照各人所行的受审判。}这正是他先前所说的：\Quote{死了的人都凭着这些案卷所记载的，照他们所行的受审判。}\par

% structure: p0045 source=h2 target=subsection reason=relative-heading-rank; base=h2

\subsection{第十五章　论由海、死亡和地狱交出受审的死者是谁。}

但海中的死者是谁，海又\Emph{交出}了谁？因为我们不能认为那些死在海中的人不在地狱，也不能认为他们的身体保存在海中；更荒谬的是，海保留了善人，而地狱接受了恶人。谁能相信这点呢？但有些人很明智地推测，此处\Quote{海}代表这个世界。当若望要表示基督在肉身中仍要找到的人将同复活的人一起受审时，他称他们为\Quote{死者}，既指善人——经上对他们说：\BibleQuote{因为你们死了，你们的生命与基督一同藏在天主内} \BibleRef{哥 3:3}，也指恶人——经上说：\BibleQuote{任凭死人埋葬他们的死人} \BibleRef{玛 8:22}。他们也可以被称为死者，因为他们带有必死的身体，正如宗徒所说：\BibleQuote{身体固然因罪而死亡，但神魂却因正义而生活} \BibleRef{罗 8:10}，证明在活着的肉身之人中，既有死亡的身体，也有生活的神魂。然而他并没有说身体是必死的，而是说它是死的，尽管随后他用更常见的说法提到必死的身体。这些就是海中的死者，海将它们交出来，即那些仍在这个世界上、尚未死去的人，世界将他们交出来受审。他说：\BibleQuote{死亡和地狱也交出了其中的死者}。海\Emph{交出}了他们，因为他们只需被发现所在之处；但死亡和地狱\Emph{交出}或\Emph{归还}了他们，因为他们召回了他们已经离开的生命。也许并非没有理由，死亡和地狱两者都被提及——死亡表示善人，他们只经历死亡而非地狱；地狱表示恶人，他们还遭受地狱的惩罚。因为如果我们相信，那些相信基督及其未来降临的古圣们被保存在远离恶人折磨的地方，但仍处于地狱，直到基督的血和祂降至这些地方释放了他们，那么被那已付出的宝贵代价赎回的好基督徒，在等待复活和接受赏报期间，就完全与地狱无涉了。在说了\BibleQuote{他们都按照各自的行为受了审判}之后，他简短地补充了审判是什么：\BibleQuote{死亡和地狱也被投入火湖}，用这些名称指代魔鬼及其所有使者，因为他是死亡和地狱痛苦的制造者。因为这是他先前用更明确的语言预先说过的：\BibleQuote{那欺骗他们的魔鬼被投入了硫磺火湖}。他在\BibleQuote{那兽和假先知也在其中}这句话中含糊的补充，在此解释为：\BibleQuote{凡没有记载在生命册上的，就被投入火湖}。这书不是为了提醒天主，好像祂会因遗忘而漏掉什么，而是象征祂对那将得永生之人的预定。因为并非天主无知，要读这书来了解自己，而是祂无误的预知就是生命册，在其中他们被记载，也就是说，被预知。\par

% structure: p0047 source=h2 target=subsection reason=relative-heading-rank; base=h2

\subsection{第十六章——论新天新地}

在完成了关于恶人的审判预言之后，还需论及善人。在简要解释了主的话\BibleQuote{这些人要进入永罚}之后，还需解释与之相连的\BibleQuote{义人却要进入永生} \BibleRef{玛 25:46}。他说：\BibleQuote{我看见一个新天新地，因为先前的天和先前的地已经过去了，海也不再有了} \BibleRef{默 21:1}。这将按照他预先宣告的顺序发生：\BibleQuote{我看见一位坐在宝座上的，从祂面前天地都逃跑了}。因为一旦那些没有记载在生命册上的人受审并被投入永火——这火的性质，或它在世界或宇宙中的位置，我想没人知道，除非圣神启示某人——然后这个世界的形态将在普世的烈火中消逝，如同先前世界在普世的洪水中被淹没一样。通过这场普世的烈火，适合我们可朽身体的腐朽元素的品质将彻底消灭，我们的实体将获得这样的品质，通过奇妙的转变，与我们不朽的身体相协调，以致世界本身更新为更美好的事物，也适宜地适应那些在肉身中更新为更美好事物的人。至于\BibleQuote{海也不再有了}这句话，我不敢轻率地说它是因那过度的热而干涸，还是本身也变成了某种更美好的事物。因为我们读到有新天新地，但我记得没有在任何地方读到新海，除非在这同一卷书中找到的：\BibleQuote{好像一个玻璃海，如同水晶} \BibleRef{默 15:2}。但他那时并非在说世界的终结，他似乎也不是在说真正的海，而是\BibleQuote{好像一个海}。可能，因为预言文体喜欢混合比喻和实际的语言，从而在一定程度上遮蔽意义，所以\BibleQuote{海也不再有了}这句话也可与前面\BibleQuote{海交出了其中的死者}取同一意义。因为那时不再有这个世界的动荡和人类生活的不安，而这就是海所象征的。\par

% structure: p0049 source=h2 target=subsection reason=relative-heading-rank; base=h2

\subsection{第十七章——论教会无尽的荣耀}

\Quote{我看见了}，他说，\Quote{一座伟大的城，新\Place{耶路撒冷}，从天上由天主那里降下，好像一位装饰好了的新娘，等候她的丈夫。我又听见从宝座发出一个宏亮的声音说：看哪，天主的帐幕在人间，祂要与人同住，他们要作祂的子民，天主亲自与他们同在。天主要从他们的眼中擦去一切眼泪；不再有死亡，也不再有悲哀、哭号、疼痛，因为先前的事都过去了。那位坐在宝座上的说：看哪，我更新一切。}\BibleRef{默 21:2}这座城被说成从天而降，因为天主用来建造它的恩宠来自天上。因此，祂通过\Person{依撒意亚}对它说：\Quote{我是上主，那塑造你的。}\BibleRef{依 45:8}的确，它从起初就从天而降，因为它的市民在这世界的进程中，借着从天上降下的圣神所施行的重生之洗，从天主那里获得成长的恩宠。然而，在天主最后的审判中，那将由祂的儿子\Person{耶稣基督}施行，因天主的恩宠，要彰显出一种如此弥漫且崭新的光荣，以致旧的一切痕迹都不复存在；因为连我们的身体也要从旧有的腐朽和必死性，转变为新的不朽和不死。因为将这一应许应用于现在这个时代，就是圣人们与他们的君王一同作王一千年的时候，在我看来是极其无耻的，因为那里明明白白地说：\Quote{天主要从他们的眼中擦去一切眼泪；不再有死亡，也不再有悲哀、哭号，疼痛也不再有。}谁还会如此荒谬，被好争论的固执蒙蔽，胆敢断言，在这必死生命的灾难中，天主的子民，甚至是一位圣人，曾活着、或曾活过、或将要活着，没有眼泪或疼痛？事实是，一个人越圣洁，越充满神圣的渴望，他的恳求就越是泪流满面。难道这些不是天上\Place{耶路撒冷}的一位市民的呼喊吗：\Quote{我的眼泪昼夜成为我的食物}，以及\Quote{我每夜使我的床铺漂流，用我的眼泪浸透我的床榻}，以及\Quote{我的呻吟不向你隐藏}，以及\Quote{我的痛苦重新发动？}难道那些叹息、负重、不是想脱去衣服，而是想穿上衣服，好让必死的被生命吞没的，不正是天主的儿女吗？\BibleRef{格后 5:4}就连那些拥有圣神初果的，不也是在自己内叹息，等候着义子的名分，即身体的救赎吗？\BibleRef{罗 8:23}难道宗徒\Person{保禄}本人不也是天上\Place{耶路撒冷}的一位市民吗？当他为他的以色列同胞心中怀着极大的忧愁和不断的痛苦时，他不是更加如此吗？\BibleRef{罗 9:2}但在那座城里，什么时候才会不再有死亡呢？除非当它被说：\Quote{死亡啊，你的争战在哪里？死亡啊，你的毒刺在哪里？死亡的毒刺就是罪。}显然，当能被说\Quote{在哪里}的时候，就不再有罪了。但就目前而言，不是这座城市的某个软弱的市民，而是这位\Person{若望}宗徒自己说：\Quote{如果我们说我们没有罪，便是自欺，真理不在我们内。}\BibleRef{若一 1:8}无疑，虽然这本书被称为\WorkTitle{默示录}，其中有许多深奥的段落来锻炼读者的心智，而且即使我们费劲，也几乎没有几段如此浅显能帮助我们解释其他段落；而且由于相同的事以不同的形式重复，这种困难增加了，以至于所指的事似乎不同，尽管实际上只是不同方式陈述。但在\Quote{天主要从他们的眼中擦去一切眼泪；不再有死亡，也不再有悲哀、哭号，疼痛也不再有}这些话中，如此明显地指向未来的世界和圣人的不朽与永生——因为只有那时和那里才能实现这样的状况——如果我们认为这是晦涩的，那么我们就不必期望在圣经的任何部分找到任何明白的东西。\par

% structure: p0051 source=h2 target=subsection reason=relative-heading-rank; base=h2

\subsection{第十八章。—— 宗徒\Person{伯多禄}关于最后审判的预言。}

现在让我们来看宗徒伯多禄关于这审判的预言。\Quote{将有人来，}他说，\Quote{在末世的好讥诮者……但我们按照祂的恩许，等候新天新地，其中有正义常住。}这里关于死人复活什么也没有说，但关于这个世界的毁灭肯定说得足够。他提及洪水，似乎是在暗示我们，应在多大的程度上相信世界的毁灭会延伸至世界末日。因为他说，那时的世界消灭了，不仅大地本身，而且连诸天——我们理解为空气，其位置被水所占据。因此，整个或几乎整个多风的大气层（他称之为天，或更确切地说诸天，意指地球的大气层，而不是放置太阳、月亮和星辰的上层空气）变成了水汽，并这样连同大地一同消灭了，因洪水而摧毁了大地原先的样子。\Quote{但现在的诸天和大地，凭同一的圣言被保存存留，直到审判及不虔敬之人灭亡的日子被火焚烧。}因此，那从水中被保存以取代在洪水中消灭的世界的诸天和大地，本身也最终在审判及不虔敬之人灭亡的日子被留给了火。他毫不迟疑地肯定，在这巨大的变化中，人也将会灭亡；然而，他们的本性仍将存续，尽管是在永罚中。也许有人会问：如果审判宣告后世界本身要燃烧，那么在火灾期间，以及在它被新天新地取代之前，圣人们会在哪里呢？因为他们在某处必须存在，因为他们有物质的身体。我们可以回答：他们将在上层区域，那场火灾的火焰不会升到那里，就像洪水的水也没有达到一样；因为他们将拥有那样的身体，使他们可以随意去任何地方。此外，当他们成为不朽的和不可朽坏的，他们将不会非常害怕那场火灾的火焰，正如那三个人的可朽坏和必死的身体能够在烈火窑中安然无恙地活着一样。\par

% structure: p0053 source=h2 target=subsection reason=relative-heading-rank; base=h2

\subsection{第十九章——宗徒保禄关于那将先于主的日子显现的敌基督所写与得撒洛尼人的话。}

我看到我必须省略许多关于这最后审判的福音书和书信中的论述，以使本卷不至于过分冗长；但我绝不能省略宗徒保禄在写给得撒洛尼人的书信中所说的话：\Quote{弟兄们，我们因我们的主耶稣基督的来临，恳求你们，}等等。\par

没有人能怀疑他（保禄）所写的是关于敌基督和审判的日子——他在这里称为主的日子——也没有人能怀疑他宣布这日子不会来到，除非那称为背教者的先来——即背弃主天主的背教者。如果这话可以公正地适用于所有不虔敬的人，那么对于他（敌基督）更该怎样呢？但不确定他将坐在什么圣殿里，是在撒罗满所建圣殿的废墟中，还是在教会中；因为这位宗徒不会将任何偶像或魔鬼的殿称为天主的殿。因此，有些人认为这段经文中的敌基督不仅指那元首本身，也指他的整个身体，即附从他的广大群众（连同他们的首领）；他们也认为，如果我们将希腊文更准确地译为\Quote{在天主殿内}，而读作\Quote{为天主的殿}或\Quote{作为天主的殿}，好像他自己就是天主的殿——教会，那么译法就更精确。至于\Quote{如今你们也知道那阻拦的是什么，}\Emph{即}，你们知道有什么阻碍或迟延的原因，\Quote{使他能在他自己的时候显露出来；}这表明他不愿意明说，因为他声称他们知道。因此，我们这些没有他们知识的人，即使费尽心力也还是不能明白这位宗徒指的是什么，尤其是因为他添加的话使他的意思更加模糊。因为他所说的\Quote{因为这不法的奥义已经运行，只是那现今阻止的，让他阻止，直到他被除去；那时那不法的人将要显露出来}，到底是什么意思呢？我坦白承认我不知道他指的是什么。然而，我将提到我所听到或读到的一些猜测。\par

有些人认为宗徒保禄指的是罗马帝国，他不愿使用更直白的语言，以免招致诽谤，说他希望这个被认为永恒的帝国遭殃；因此，当他说\Quote{不义的奥秘已经运行}时，他暗指\Person{尼禄}，尼禄的所作所为已经像\OriginalTerm{假基督}{Antichrist}的行径。于是有些人推测尼禄将复活，成为假基督。另有些人则推测他根本没有死，而是被隐藏起来，让人以为他被杀了；他现在以他被认为死去时的同样年龄活着，隐藏起来，直到他适时显现，恢复他的王国。但我惊讶于人们竟敢如此大胆地猜测。然而，相信宗徒的话\Quote{只有那抑制者现在抑制着，直到他被除去}指的是罗马帝国，并非荒谬；就好像是说\Quote{只有那现在统治的，让他统治，直到他被除去}。\Quote{然后那恶者将要显露出来}：无人怀疑这指的是假基督。但另一些人认为\Quote{你们知道那阻止的是什么}和\Quote{不义的奥秘已在运行}这些话只指教会中的恶人和伪善者，直到他们的人数多到能为假基督提供一大群民众，而这就是不义的\Emph{奥秘}，因为它似乎是隐藏的；宗徒劝勉信友坚忍地持守他们所持的信仰，当他说\Quote{只有那抑制者现在抑制着，直到他被除去}时，意思是直到现在隐藏的\OriginalTerm{不义的奥秘}{mystery of iniquity}离开教会。因为他们认为，\Person{若望}在书信中提到同一奥秘时说：\BibleQuote{孩子们，现在是末时了；你们既然听说过假基督要来，如今已经有许多假基督出来了，由此我们知道现在是末时了。他们从我们中间出去，却并不是属于我们的；因为如果他们是属于我们的，就必会留在我们中间。}\BibleRef{若一 2:18}。因此，正如在末时之前有许多异端者从教会出去，\Person{若望}称之为\Quote{许多假基督}，他也称之为\Quote{末时}，同样，在末时，那些不属于基督、而属于最后的假基督的人将要出去，然后假基督将被显露。\par

因此，对于宗徒隐晦的话，就有各种猜测性的解释。他所说的一点是毫无疑问的：基督不会来审判活人死人，除非他的仇敌假基督先来诱惑那些灵魂已死的人；虽然他们的诱惑是出于天主已经作出的秘密审判的结果。因为，正如经上所说\Quote{他的来临，是出于撒殚的作为，具有各种德能、征兆和欺诈的奇迹，以及各种不义的欺骗，在那些丧亡的人身上。}因为那时撒殚将被释放，并藉着那假基督，以一切德能作工，虽是欺诈却奇异。通常有人质疑这些作为为何被称为\Quote{征兆和欺诈的奇迹}：是因为他要以虚假的显现欺骗人的感官，还是因为他所做的事，虽然是真正的奇事，但对于那些认为只有天主才能行这样的事、不知魔鬼的能力、尤其不知他那时将首次施展的空前能力的人而言，成了谎言。因为当他像火一样从天上坠落，一举扫荡了圣约伯众多的仆人和大群的牲畜，然后如旋风冲击房屋，杀死他的儿女时，这些并非欺骗性的显现，却是撒殚的作为，而天主给了他这权能。为何它们被称为征兆和欺诈的奇迹，到那时我们更有可能知道。但无论名称的原因是什么，它们将是这样的征兆和奇迹，足以诱惑那些该受诱惑的人，\Quote{因为他们没有接受爱真理的心，好得救。}宗徒也毫不迟疑地接着说：\Quote{为此，天主就给他们派遣了一种错误的力量，使他们相信谎言。}因为天主将要\Emph{派遣}，是因为天主将准许魔鬼做这些事，这准许是由于他公义的审判，虽然魔鬼的作为是出于他不义和恶毒的意图，\Quote{为使一切不信真理、反而喜爱不义的人，都被定罪。}因此，他们被定罪，就受诱惑；受诱惑，就被定罪。但藉着天主那些秘密而公正、公正而秘密的审判——自从理性受造物首次犯罪以来，天主从未停止审判——他们被定罪，就受诱惑；受诱惑，就在那最后而明显的审判中被定罪，这审判由耶稣基督执行，他自己曾是最不公正地被审判，而他将最公正地审判。\par

% structure: p0058 source=h2 target=subsection reason=relative-heading-rank; base=h2

\subsection{第二十章——同一宗徒在\BibleBook{得撒洛尼}{前书}关于死人复活所教导的。}

但宗徒在这里关于死人复活什么也没有说；但在他的\BibleBook{得撒洛尼}{前书}中，他说：\BibleQuote{弟兄们，关于安眠的人，我们不愿你们不明白}\BibleRef{得前 4:13}等等。宗徒的这些话最清楚地宣告了将来死人的复活，那时主基督要来审判活人死人。\par

但常有人问：吾主在世上将发现的那些活着的人（在此段中由宗徒及与他同活的人所代表），是否根本不会死，还是会在他们与那些复活者一同被提到空中迎接主的同一瞬间，以不可思议的速度穿越死亡而达于不朽？因为我们不能说，当他们被高举在空中时，不可能既死又复生。因为这话：\Quote{我们就这样常与主同在}，不可理解为他在说我们将永远与主留在空中；因为主自己不会留在那里，只是在来临时经过而已。因为我们将去迎接那来临的主，不是留在主所在之处；但\Quote{我们就这样与主同在}，意思是，无论我们在何处与他同在，我们都将拥有不朽的身体与他同在。我们似乎必须这样理解这话，并假定：主在世上将发现的那些活着的人，将在那短暂期间既经历死亡又获得不朽；因为同一位宗徒说：\Quote{\BibleQuote{在基督内众人都要复活}} \BibleRef{格前 15:22}；而他在别处论到同一身体的复活时又说：\Quote{\BibleQuote{你所播种的，若不先死，决不会活过来}} \BibleRef{格前 15:36}。那么，基督在世上将发现的那些活着的人，若他们不死，又怎能在他内获得永生呢？因为正是为此才说：\Quote{\BibleQuote{你所播种的，若不先死，决不会活过来}}？或者，如果我们不能恰当地说人的身体是被播种的，除非就因死亡而多少回归大地而言，正如天主对犯罪的人类始祖所宣告的判词说：\Quote{\BibleQuote{你原是土，仍要归于土}} \BibleRef{创 3:19}。那么我们必须承认，基督来临时仍将在肉身中发现的那些人，不在宗徒的这些话和创世纪的话之内；因为他们被提到云中，既不是被播种，也没有进入或回归大地，无论他们完全不死，还是在空中暂时死去。\par

但另一方面，同一位宗徒在给格林多人论到身体复活时的话也出现在我们面前：\Quote{\BibleQuote{我们都要复活}}，或如其他抄本所读：\Quote{\BibleQuote{我们都要安眠}} \BibleRef{格前 15:51}。既然没有死亡就没有复活，既然我们在此处只能将安眠理解为死亡，那么，如果许多基督在肉身中要发现的人既不睡眠也不复活，又怎能说\Emph{我们都要}安眠或复活呢？因此，如果我们相信基督来临时将被发现活着并被提到迎接他的圣人们，会在同一升腾中从可朽的身体转为不朽的身体，那么宗徒的话就不会让我们为难，无论他说\Quote{\BibleQuote{你所播种的，若不先死，决不会活过来}}，还是说\Quote{\BibleQuote{我们都要复活}}或\Quote{\BibleQuote{都要安眠}}，因为连圣人们若不先死（哪怕短暂），也不会被赋予不朽的生命；因此，他们也不会免除由安眠所预表的复活，无论这安眠多么短暂。而且，为什么我们觉得难以置信：那一大群身体仿佛播在空中，又立即在空中复活成为不朽、不腐的？因为我们凭同一宗徒的见证相信，复活要在眨眼之间发生，早已死去的身体的尘土要以难以理解的速度和敏捷地回归到那些现在要永远活着的肢体上。我们也不认为在这些圣人身上，\Quote{\BibleQuote{你原是土，仍要归于土}}的判决无效，尽管他们的身体在死时并未落在地上，而是在被提到空中时既死又立即复活。因为\Quote{\BibleQuote{你要归于土}}的意思是，你在死时要归回生命开始前的状态。你失去灵魂时，将成为你被赋予灵魂之前的状态。因为天主在造人成为有灵的生物时，向地上的面孔吹了生命之气；这就好像是说，你是有灵魂的土，你原本不是如此；你将成为没有灵魂的土，就像你原本那样。这就是所有死了的身体在腐烂前的状态；也是那些圣人如果在死时（无论何处死）的身体状态，他们一旦放弃那生命，就立即重新获得。因此，他们就这样归回或进入土，因为从活人变成了土，就像烧成灰的东西说是归灰，腐烂的东西说是归腐，以及类似的一百多种说法。但这种情形如何发生，我们现在只能微弱地猜测，只有到了那时才会明白。因为基督来审判生者死者时，必有身体的复活，这是我们若为基督徒就必须相信的。但如果我们无法完全理解其情形，我们的信仰也不会因此落空。但现在，我们应当如先前所许诺的，尽必要地展示古代先知书关于天主这最后审判的预言；我想，如果读者细心利用我们已经提供的帮助，我们无需花太多时间讨论和解释这些预言。\par

% structure: p0062 source=h2 target=subsection reason=relative-heading-rank; base=h2

\subsection{第21章——先知依撒意亚关于死者复活和报应审判的言论。}

先知\Person{依撒意亚}说：\BibleQuote{死人要复活，凡在坟墓里的都要复活；凡在地上的都要欢乐，因为从你而来的露水是他们的健康；恶人的地却要倒塌。} \BibleRef{依 26:19} 前半部分都与义人的复活有关；但\BibleQuote{恶人的地却要倒塌}这话正确理解为：恶人的身体要坠入永罚的毁灭。如果我们更精确仔细地审视关于义人复活的话，我们可以将\BibleQuote{死人要复活}这话应用于首次复活，而将\BibleQuote{凡在坟墓里的都要复活}这话应用于第二次复活。如果我们问及那些主来临时将在地上仍活着的圣人，下面的话可以合适地应用于他们：\BibleQuote{凡在地上的都要欢乐，因为从你而来的露水是他们的健康。}此处的\Quote{健康}最好理解为不朽。因为那是最完美的健康，无需藉营养如同每日的补救来修复。同样，同一位先知为义人提供希望，为恶人显明审判之日的恐怖，说：\BibleQuote{主这样说：看，我要像和平的河流一样涌向他们，像急流一样涌向外邦人的荣耀；他们的儿子必被扛在肩上，在膝上得安慰。就如母亲安慰人，我也要这样安慰你们；你们要在耶路撒冷得安慰。你们必看见，你们的心必欢乐，你们的骨头必如草发芽；主的手必被祂的敬拜者认识，祂必威胁悖逆的人。因为，看，主必像火一样来临，祂的战车像旋风，要施行报应，发怒并用火焰毁灭。因为主要用火审判全地，用祂的刀审判凡有血肉的；许多人必被主击伤。} 在祂对义人的应许中，祂说祂要像和平的河流一样涌下，就是说，以最丰盛的和平。我们终将因这和平而得更新；但我们在前一卷书中已充分论述此事。正是这条河流，祂说祂要涌向那些祂应许如此大福乐的人，好使我们明白，在那幸福之境——即天上——万物都从这河得满足。但因为从那里，甚至要流到地上的身体，即不朽与不死的和平，因此祂说祂要像这河一样涌下，好使祂既从上倾注于下，使人同天使平等。同样，\BibleQuote{耶路撒冷}，我们不应理解为那与儿女一同为奴的，而要按宗徒的话，理解为那是我们自由的母亲，永恒在天上。\BibleRef{迦 4:26} 在她内，我们将得安慰，如同历经尘世忧患与灾祸的疲惫者，并如她的儿女被抱在膝上和肩上。我们对这等抚慰既生疏又新鲜，将被接纳入不寻常的福乐中。在那里我们将看见，我们的心将欢乐。祂没有说我们将看见什么；但除了天主，还有什么呢？好使福音中的应许在我们身上应验：\BibleQuote{心里洁净的人是有福的，因为他们要看见天主？} \BibleRef{玛 5:8} 我们将看见什么？岂不是那些我们现在看不见、但相信的一切吗？我们按自己微弱的能力所形成的观念，远不及现实。他说：\BibleQuote{你们必看见，你们的心必欢乐。}这里你们相信，那里你们将看见。\par

但因为他说：\Quote{你们的心要欢乐}，免得我们以为那耶路撒冷的福乐仅是属灵的，他接着又说：\Quote{你们的骸骨要像青草一样兴起}，意指身体的复活，仿佛补充了他之前遗漏的内容。因为这事不会在我们看见时发生，而是在它发生之后我们才看见。因为他早已论及新天新地，多次用许多比喻讲述应许给圣人们的那些事，说：\BibleQuote{将有新天新地；先前的事不再被记念，也不再涌上心头；但他们要在其中找到喜乐和欢欣。看，我要使耶路撒冷成为欢欣，使我的百姓成为喜乐。我要因耶路撒冷而欢欣，因我的百姓而喜乐；其中不再听见哭泣的声音；}\BibleRef{依 65:17}还有其他一些许诺，有人试图将它们解释为千年期间的肉欲享乐。因为，按预言的方式，\OriginalTerm{比喻性表达和字面表达}{figurative and literal expressions}是混合在一起的，以便认真的人通过有益而救恩的努力，可以达到\OriginalTerm{灵性意义}{spiritual sense}；但\OriginalTerm{肉欲的迟钝}{carnal sluggishness}，或未受教育、未受训练之心灵的迟缓，停留在表面的字句，以为下面没有什么可寻求的。但对于刚才引用的那些预言表达的风格，就说到这里。现在，回到它们的解释。当他说：\Quote{你们的骸骨要像青草一样兴起}时，为了表明他所指的是善人的复活，虽然是身体的复活，他又接着说：\Quote{主的手必被敬拜祂的人所认识。}这岂不就是那位将敬拜祂的人与藐视祂的人分别出来的手吗？关于这些人，上下文紧接着说：\Quote{祂要威胁悖逆的人}，或如另一位译者所译的：\Quote{不信的人。}祂那时并不是实际威胁，而是现在所说的威胁那时将实际应验。他说：\BibleQuote{因为看，主要像火一样来临，祂的战车像旋风，要发怒执行报复，并用火焰施行毁灭。因为全地都要用主的火来审判，凡有血肉的都要用祂的剑来审判：许多人要被主击伤。}藉着\Emph{火}、\Emph{旋风}、\Emph{剑}，他意指天主的审判惩罚。因为他说主自己要像火一样来临，即对那些祂的来临将是惩罚的人。藉着祂的\Emph{战车}（因为该词是复数），我们恰当地理解为天使的服役。当他说凡有血肉的和全地都要用祂的火和剑来审判时，我们并不理解为包括灵性的和圣洁的人，而是指属地和属血肉的人，关于这些人，经上说他们\Quote{思念地上的事}\BibleRef{斐 3:19}，\Quote{随从肉性的人，就是死亡}\BibleRef{罗 8:6}，并且主在说\Quote{我的神不永远住在这些人里面，因为他们是血肉}\BibleRef{创 6:3}时，简单地称他们为血肉。至于\Quote{许多人要被主击伤}这句话，这伤将导致\OriginalTerm{第二次死亡}{second death}。事实上，也有可能从好的意义上理解\Emph{火}、\Emph{剑}和\Emph{伤}。因为主曾说祂愿意把火投在地上\BibleRef{路 12:49}。当圣神降临时，有似火焰的舌头显现给他们\BibleRef{宗 2:3}。我们的主说：\BibleQuote{我来不是为把平安带到地上，而是刀剑}\BibleRef{玛 10:34}。圣经说天主的话是双刃的利剑\BibleRef{希 4:12}，因为有两刃，即两约。在\WorkTitle{雅歌}中，圣教会说她因爱而受伤\BibleRef{歌 2:5}，仿佛被爱之箭刺透。但在这里，当我们读到或听到主要来执行报复时，显然这些表达该按何种意义理解。\par

在简短地提及那些将在这次审判中被消灭的人之后——他用旧法律所禁止的肉类来比喻邪恶者和罪人，因为他们没有禁戒这些——他概括地讲述了新约的恩宠，从救主的首次来临直到我们此刻所论的最后审判；并以此结束他的预言。因为他叙述说，主宣布祂要来聚集万民，使他们来并看见祂的荣耀\BibleRef{依 66:18}。因为，如宗徒所说：\BibleQuote{所有人都犯了罪，缺少天主的荣耀}\BibleRef{罗 3:23}。他又说祂要在他们中间施行奇事，使他们惊奇而信从祂；并且要从他们中派遣得救的人到各国，到那未曾听过祂名、未曾见过祂荣耀的遥远海岛；他们要列国中宣扬祂的荣耀，并要\Emph{带来}先知所说话之人的弟兄们，\Emph{也就是说}，将选民以色列人的弟兄们带到天父天主面前的信仰中；他们要从万国中用驮畜和车辆（这应理解为天主以天使或人事的服役所提供的帮助）带来供物献给主，送到圣城耶路撒冷，这城目前散布在全地，存在于忠信的圣人中。因为哪里赐予了天主的助佑，哪里的人就相信；哪里的人相信，他们就前来。主又用比喻将他们比作以色列的子民在祂的殿中弹琴献祭，这已由教会在各处实行；祂又应许要从他们中为自己拣选司铎和肋未人，这也已看到实现。因为我们现在看到，司铎和肋未人的拣选，不再是按原来的规矩，从某一家族和血统中，如同亚郎品位的司祭职，而是符合新约，在默基瑟德品位之下基督是大司祭，按每人由天主的恩宠所赐予的功德。这些司铎不应仅凭其名衔来评判——这名衔常由不配的人所拥有——而应凭那善人与恶人所不共有的圣德。\par

在这样论及天主对教会现今所施予的仁慈——这是显而易见并为我们所熟悉的——之后，他又预言了当最后审判将善人与恶人分开时，人们将面临的结局。先知如此说，或是先知本人代替天主发言说：\Quote{因为新天新地怎样在我面前存留，主说，你们的后裔和你们的名字也必照样存留。他们将有月复一月，安息日复安息日。凡有血肉的，都必来到耶路撒冷，在我面前朝拜，主说。他们要出去，看见那些得罪我之人的\OriginalTerm{肢体}{members of the men}：他们的虫不死，火也不灭；他们必成为一切有血肉之人的景象。} \BibleRef{依 66:22}先知在这里结束了他的书卷，正如世界在此处终结一样。有些译者确实将\Quote{肢体}译为\Quote{\OriginalTerm{尸体}{carcasses}}，用\Emph{尸体}指代身体明显的惩罚，尽管\Quote{尸体}通常只用于死去的肉体，而此处所论的身体是有生命的，否则它们不会感到任何痛苦；但或许可以毫不荒谬地将它们称为尸体，因为它们是那些将要落入\OriginalTerm{第二次死亡}{second death}之人的身体。出于同一原因，同一先知也说过，正如我已引用的：\Quote{恶人的地必坠落。}显然，那些对人采用不同译词的译者，并非只指男性，因为没有人会说犯罪的妇女不会在那审判中出现；但男性更为尊贵，且女人由男人而出，因此以男性涵盖两性。然而，与我们主题特别相关的是：既然\Quote{凡有血肉的都必来到}这句话适用于善人，因为天主的子民将由各种族的人组成——并非所有人都会到场，因为大部分将在惩罚中——但如我所说，既然\Emph{血肉}用于善人，而\Emph{肢体}或\Emph{尸体}用于恶人，那么毫无疑问，那区分善恶归宿的审判，必在身体复活之后发生；我们对此的信德，因这些词的使用而得以彻底建立。\par

% structure: p0067 source=h2 target=subsection reason=relative-heading-rank; base=h2

\subsection{第二十二章——善人出去观看恶人的惩罚有何含义。}

但善人将以何种方式出去观看恶人的惩罚呢？他们是否要通过身体的移动，离开他们福乐的居所，前往惩罚之地，从而亲眼目睹恶人的折磨？当然不是；他们乃是借着知识出去。因为\Quote{出去}这个词，表示那些受惩罚的人是在外面。因此，主也称这些地方为\Quote{\OriginalTerm{外黑暗}{outer darkness}} \BibleRef{玛 25:30}，与之相对的是那入口，关于它曾对良善的仆人说：\Quote{进入你主的福乐罢}，免得让人以为恶人能进入那里并被认识，而是善人借知识出去到他们那里，因为善人要知道那在外面的事。因为那些在痛苦中的人，不会知道在主的福乐里面发生的事；但那些进入那福乐的人，将会知道在外黑暗中外面发生的事。因此经上说：\Quote{他们要出去}，因为他们将知道那些在外面的人所做的事。如果先知们，借着那时住在他们心中（虽然有限）的天主，能够知道尚未发生的事，那么当天主成为万物之中的万有时 \BibleRef{格前 15:28}，不朽的圣人们岂不能知道已经发生的事吗？因此，圣人们的后裔和名字将在那福乐中长存——后裔，即\Person{若望}所说的：\Quote{他的后裔存留在他内} \BibleRef{若一 3:9}；名字，即通过\Person{依撒意亚}本人所说的：\Quote{我要赐给他们一个永远的名} \BibleRef{依 56:5}。\Quote{他们将有月复一月，安息日复安息日}，如同说月复一月，安息复安息，两者都将是他们自己，当他们从时间的旧影子进入永恒的新光明时。\OriginalTerm{虫}{worm}不死，\OriginalTerm{火}{fire}不灭，构成恶人的惩罚，对此有不同的解释。有些人将两者都归于身体，有些人将两者都归于灵魂；又有些人将火字面地归于身体，将虫象征地归于灵魂，这似乎更可信。但现在不是讨论这种分歧的时候，因为我们承担了本书论述最后审判的任务，其中善人与恶人被分开；他们的赏报和惩罚，我们将在别处更详细地讨论。\par

% structure: p0069 source=h2 target=subsection reason=relative-heading-rank; base=h2

\subsection{第二十三章——\Person{达尼尔}关于\OriginalTerm{假基督}{Antichrist}的迫害、\OriginalTerm{天主的审判}{Judgment of God}以及\OriginalTerm{圣者的王国}{Kingdom of the Saints}的预言。}

达尼尔预言最后审判的方式，是藉以指出假基督将先来，并将描述延续到圣人的永恒统治。因为当他在先知异象中看见四兽，象征四个王国，而第四兽被某一位王——他被认作假基督——征服，在此之后是人子（即基督）的永恒王国，他说：\Quote{我的神魂惊惶，我达尼尔在身躯之中，我头脑的异象使我烦乱}等。有些人解释这四个王国为亚述、波斯、马其顿和罗马。凡欲了解这解释之恰当性者，可读热罗尼莫所著\WorkTitle{论达尼尔书}，此书写得足够仔细且博学。但读此段者，即使半睡，也不能不见假基督的王国将在天主最后审判引入圣人的永恒统治之前，猛烈地（虽然为时短暂）攻击教会。因为从上下文明显可见，\Emph{一个时期、两个时期和半个时期}意指一年、两年和半年，即三年半。有时在圣经中同一件事以月份表示。因为虽然拉丁文中\Emph{时期}一词似乎在此处无定数地使用，那只是因为拉丁语没有双数，如同希腊语有双数，且据说希伯来语也有。因此，\Emph{时期}被用来表示两个时期。至于十位君王，假基督似乎在他来到时将在十个个体身上找到，我承认我担心我们可能在此受骗，而他可能在罗马世界中没有十位君王活着时意外地来到。因为如果这数字十象征所有在他来临之前君王的全部数目，正如整体常以千、百、七或其他数字来象征，而这些不必细数呢？\par

在另一处，同一达尼尔说：\Quote{那时必有大灾难，自从有国以来直到那时，未曾有过这样的灾难；那时，凡写在册子上的你的人民，都必得救。有许多睡在尘土中的人要醒起，有的得到永生，有的受到羞辱，永远蒙羞。明智者要发光，如同穹苍的光辉；许多义人要像星辰，直到永远。}\BibleRef{达 12:1}这段经文与我们从福音所引用的\BibleRef{若 5:28}非常相似，至少就死人复活而言。因为那里所说的\Quote{在坟墓里的}，这里说\Quote{睡在尘土中的}，或者如其他译本所说\Quote{在尘土中的}。那里说\Quote{他们要出来}，这里说\Quote{他们要醒起}。那里说\Quote{行善的，复活得生命；作恶的，复活受审判}；这里说\Quote{有的得到永生，有的受到羞辱，永远蒙羞}。也不应以为有差别，因为福音中说的是\Quote{所有在坟墓里的}，而先知不说\Quote{所有}，却说\Quote{许多睡在尘土中的}。因为\Emph{许多}在圣经中有时用作\Emph{所有}。例如，对亚巴郎说：\Quote{我已立你为许多民族之父}，而在另一处对他说：\Quote{地上万民都要因你的后裔蒙福}。关于这样的复活，稍后对先知本人说：\Quote{你且去等待，因为还有日子直到完成末期；你要安息，在末日你的份中你要起来。}\BibleRef{达 12:13}\par

% structure: p0072 source=h2 target=subsection reason=relative-heading-rank; base=h2

\subsection{第二十四章——从\WorkTitle{达味圣咏}中预言世界终结和最后审判的经文。}

\BibleQuote{在起初您奠定了大地的根基，上主，苍天是您双手的化工。它们要消灭，而您常存；是的，一切都要像衣服一样陈旧；您要改变它们像更换衣裳，它们就被改变。但您永是同一，您的岁月没有穷尽。}为什么\Person{波斐利}一方面赞扬希伯来人的虔诚，他们崇拜一位伟大真实、令众神本身也畏惧的天主，另一方面却追随这些神明的神谕，指责基督徒极端愚蠢，因为他们说这个世界要消灭？因为我们在希伯来人的圣书中，正好看到向这位伟大的哲学家承认令众神都畏惧的天主所说的话：\BibleQuote{苍天是您双手的化工；它们要消灭。}如果苍天，这个世界较高且较稳固的部分要消灭，世界本身还能保全吗？如果这个观念不为\Person{尤庇特}所喜爱——这位哲学家引用尤庇特的神谕作为无可置疑的权威来斥责基督徒的轻信——为什么他不同样斥责希伯来人的智慧为愚蠢，既然这预言出现在他们至圣的书中？但如果\Person{波斐利}如此迷恋希伯来人的智慧，乃至通过他自己神明的言词来赞扬它，而这位智慧宣告苍天要消灭，他怎么又如此昏聩，部分地——即使不是主要地——因为基督徒相信世界要消灭而憎恨他们的信仰？——虽然如果世界不消灭，苍天如何消灭并不容易明白。事实上，在我们独有的圣书中——即福音书和宗徒著作——而不是希伯来人与我们共通的书中，用了以下说法：\BibleQuote{这世界的景象正在逝去；}\BibleRef{格前 7:31}\BibleQuote{世界正在逝去；}\BibleRef{若一 2:17}\BibleQuote{天地将要逝去，}\BibleRef{玛 24:35}——这些说法，我想，比起\BibleQuote{它们要\Emph{消灭}}要温和一些。在宗徒\Person{伯多禄}的书信中，也说到那当时的世界被水淹没而消灭了，由此足以看出整个世界所指的世界部分是哪里，以及\Emph{消灭}一词应如何理解，哪些天被保留、存留于火，直到审判和不信的人灭亡的日子\BibleRef{伯后 3:6}。当稍后他说：主的日子要像贼一样来临；在那日子，天要哗然消逝，元素要因焚烧而熔化，大地及其中的工程都要被烧尽，然后加上：\BibleQuote{既然这一切都要如此消散，你们应该成为何等样的人？}\BibleRef{伯后 3:10}——这些将要消灭的天，可以理解为与前述被保留存留于火的天是相同的；而那些将要被烧的元素，则是这世界最底层充满风暴和骚乱的元素，他说这些天被保留在其中；因为较高天穹上的星辰是安全的，保持完整。因为圣经说\BibleQuote{星辰要从天上坠落}\BibleRef{玛 24:29}，且不提另一种更可取的解释，这反而显示天本身要存留，如果星辰从它们上面坠落。那么，这一说法或是比喻性的——这更可信——或是这一现象将发生在最底层之天，如维吉尔所提及的那样——\par

\Quote{一颗流星拖着光尾，\newline{}划过天空，耀眼明亮，\newline{}然后在\Place{伊达}山林中消失。}\par

但是我所引用的\WorkTitle{圣咏}中的那段话，似乎没有把任何天从毁灭的命运中排除出去；因为他说：\BibleQuote{苍天是您双手的化工：它们要消灭；}因此，既然没有一个天被排除在天主化工的范畴之外，也就没有一个天被排除于毁灭之外。因为我们的对手不愿屈尊俯就，用宗徒\Person{伯多禄}的话来为希伯来人的虔诚辩护——这种虔诚赢得了他们神明的赞许——而他们又强烈憎恨伯多禄；他们也不会论证说，正如宗徒在他的书信中谈到全世界在洪水中消灭时，所指的只是世界最低的部分和相应的天被毁灭，同样，在圣咏中，整体被用来指部分，说\BibleQuote{它们要消灭}，而实际上只有最低的天要消灭。但是，正如我所说，他们不愿如此推理，以免显得认同\Person{伯多禄}的意思，或者像我们看待洪水那样重视最终的大火，而他们却坚持认为，没有水或火焰能够毁灭全人类，这样一来，他们只能坚持认为，他们的神明赞扬希伯来人的智慧，是因为他们没有读过这篇\WorkTitle{圣咏}。\par

也正是天主的最后审判，正如\BibleBook{}{圣咏}第五十篇所言：\Quote{天主必要来临，我们的天主，决不缄默；烈火在祂面前吞噬，暴风在祂四周旋转。祂呼唤高天和大地，为审判自己的子民：你们应聚集我的圣徒，就是那些用祭品与我立约的人。}我们明白这指的是我们的主耶稣基督，就是我们期待祂从天降来审判生者死者的那一位。因为祂曾隐秘地来临，为不义者所不义地审判；如今祂将显明地来临，以正义审判义人与不义者。我说，祂必显明地来临，决不缄默，就是藉着审判的声音使自己为人所知。而祂以前隐秘地来临时，在审判者面前缄默不语，正如我们读到\BibleBook{依撒意亚}{先知书}关于祂的预言：\BibleQuote{祂被牵去待宰，像羔羊在剪毛者前缄默，不开口}\BibleRef{依 53:7}，又看到这在\BibleBook{}{福音}中应验了\BibleRef{玛 26:63}。至于\WorkTitle{圣经}所说的\Emph{火}和\Emph{暴风}，我们在解释\BibleBook{依撒意亚}{先知书}类似段落时已经说明应如何理解。至于\Quote{祂呼唤高天}这句话，既然圣徒和义人正当地被称为\Emph{天}，无疑这正是指宗徒所说的：\InnerQuote{我们将被提到云彩中，在空中与主相遇}\BibleRef{得前 4:17}。因为如果按字面意义理解，如何能\Quote{呼唤高天}呢？难道天还能在别处而不在高处吗？接下来的话：\Quote{和大地，为审判自己的子民}，如果我们只补充\Quote{祂呼唤}这个词，就是说\Quote{祂也呼唤大地}，而不补充\Quote{高处}，那么似乎给我们一种合乎正确教义的理解：天象征那些要与基督一同审判的人，地象征那些将要受审判的人。因此，\Quote{祂呼唤高天}这句话的意思不是\Quote{祂要将人提到空中}，而是\Quote{祂要提升他们到审判的宝座}。也可能，\Quote{祂呼唤高天}的意思是，祂呼唤在高天之上的天使，好同他们一起降来进行审判；而\Quote{祂也呼唤大地}则指祂呼唤地上的人来受审判。但如果我们在\Quote{和大地}这句话中不仅理解为\Quote{祂呼唤}，还理解为\Quote{高处}，从而使完整的意思成为\Quote{祂呼唤高天}，也\Quote{呼唤高地}，那么我认为这最好理解为那些要被提到空中迎接基督的人，他们被称为\Emph{天}是指灵魂，称为\Emph{地}是指身体。那么，\Quote{为审判自己的子民}若不是指藉审判将善人与恶人分开，正如将绵羊从山羊中分别出来，又是什么呢？然后，圣咏转向对天使说话：\Quote{你们应聚集祂的圣徒到祂面前。}因为如此重要的事，当然必须藉天使的服务来完成。如果我们问，那些由天使聚集到祂面前的圣徒是谁，我们被告知：\Quote{那些用祭品与祂立约的人。}这就是圣徒的全部生活：藉祭品与天主立约。因为\Quote{藉祭品}或者指慈悲的工作，在天主眼中这些胜于祭品，因为祂说：\Quote{我喜欢仁爱胜过祭献}\BibleRef{欧 6:6}；或者如果\Quote{藉祭品}意味着在祭品中，那么这些慈悲的工作本身就是天主所喜悦的祭品，正如我记得在本著作第十卷中所论述的；圣徒在这些工作中与天主立约，因为他们为了祂的新约或盟约中所包含的应许而做这些工作。因此，当祂的圣徒被聚集到祂面前，在最后审判中站在祂右边时，基督会说：\Quote{你们蒙我父降福的，来吧！承受自创世以来为你们预备了的国度吧！因为我饿了，你们给了我吃的}\BibleRef{玛 25:34}等等，列举了善人的善行，以及由审判者的最后判决所指定的永恒赏报。\par

% structure: p0077 source=h2 target=subsection reason=relative-heading-rank; base=h2

\subsection{第二十五章——\Person{玛拉基亚}的预言：关于最后审判，以及某些人藉净化性惩罚所受的炼净。}

先知\Person{玛拉基亚}，也称为天使，有些人（因为\Person{热罗尼莫}告诉我们这是希伯来人的意见）认为他就是司铎\Person{厄斯德拉}，后者的其他著作已被收入正典，预言了最后审判，说：\Quote{看，他来了，全能的上主说；谁又能忍受他来临的日子？……因为我，上主，你们的天主，是不改变的。} 从这些话更明显地看出，有些人在最后审判中将遭受某种\OriginalTerm{炼净性的惩罚}{purgatorial punishments}；因为还能怎样理解这话：\Quote{谁又能忍受他来临的日子，或谁能注视他？因为他进入如同熔炉之火，又如同漂布者的碱草；他要坐下熔炼和净化，如同对待金银；他要净化肋未的子孙，倾倒他们如同金银？} 同样，依撒意亚说：\Quote{上主将以审判的神和焚烧的神，洗净熙雍子女的污秽，涤除他们中间的血。} \BibleRef{依 4:4} 除非我们或许应该说，当他们因惩罚性的审判而与恶人分离时，他们从污秽中被洁净，在某种程度上被澄清，以致一方的消灭和定罪成为另一方的净化，因为他们从此将摆脱这类人的污染而生活。但当他说：\Quote{他要净化肋未的子孙，倾倒他们如同金银，他们就要以公义向上主奉献祭献；犹大和耶路撒冷的祭献将中悦上主，} 他宣告那些被净化的人将以公义的祭献中悦上主，因而他们自己将从自己的不义中被净化，这曾使他们不中悦天主。现在，他们自己，当被净化后，将成为完全和完美公义的祭献；因为这样的人能向上主奉献什么比他们自身更悦纳的祭品呢？但关于炼净性的惩罚这一问题，我们必须推迟到另一时间，以作更充分的处理。借着肋未、犹大和\Place{耶路撒冷}的子孙，我们应当理解教会本身，不仅从希伯来人中，也从其他民族中聚集；也不是像她现在这样，那时\Quote{我们若说自己没有罪，便是自欺，真理不在我们心里了，} \BibleRef{若一 1:8} 而是像她那时将要有的样子，被最后审判净化，如同被簸扬的风所净化的禾场，她的那些需要被火洁净的成员被净化，以致绝对没有一个人为自己的罪奉献祭献。因为所有奉献这种祭献的人，确实仍在他们的罪中，他们奉献祭献是为了罪的赦免，以便在向上主奉献了悦纳的祭献后，得以被赦免。\par

% structure: p0079 source=h2 target=subsection reason=relative-heading-rank; base=h2

\subsection{论圣人们向上主所奉献的、中悦于他的祭献，如同在起初的日子和往年一样。}

正是为了表明他的城市那时将不再遵守这一习俗，天主说肋未的子孙应当以公义奉献祭献——因此不是处于罪中，因而也不是为罪而献。因此我们看到，犹太人多么徒劳地根据以下的话，自己应许回归旧约律法下的旧时祭献，他们引证的话是：\Quote{犹大和耶路撒冷的祭献将中悦上主，如同在起初的日子，如同在往年。} 因为在律法时代，他们奉献祭献不是出于公义而是出于罪，特别是首先为罪而奉献，以至于甚至那位司铎本人——我们应假定他是他们中最正义的人——也按天主的诫命习惯先为自己的罪，然后为人民的罪奉献祭献。因此，我们必须解释如何理解这些话：\Quote{如同在起初的日子，如同在往年；} 因为他或许暗指我们原祖父母在乐园的时代。那时，他们确实完整纯洁，毫无罪的污点和瑕疵，将自己作为最纯洁的祭献奉献给天主。但由于他们因过犯而被逐出那里，人性在他们身上受到了谴责，除了一位中保以及那些已受洗且尚为婴儿的人之外，\Quote{没有一个人是无污点的，即使是在地上只活了一天的婴孩。} \BibleRef{约 14:4} 但如果有人回答说，那些凭信德奉献的人可以说是以公义奉献，因为义人因信德而生活，\BibleRef{罗 1:17}——然而，若他说自己无罪，便是自欺，因此他不这样说，因为他因信德而生活——有谁会说这个信德的时代能与那\OriginalTerm{圆满终结}{consummation}相提并论，那时那些以公义奉献祭献的人将因最后审判的火而得以净化呢？因此，既然必须相信在如此净化之后，义人将不再保有罪恶，那么毫无疑问，就无罪的状态而言，那个时代不能与任何其他时期相比，除非与我们原祖父母在犯罪之前生活在乐园中那最纯洁的幸福时期相比。所以，当经上说\Quote{如同在起初的日子，如同在往年}时，应当正确理解为这个时期。因为依撒意亚也说过，在应许了新天新地之后，在圣人们的福乐中（那里以比喻和象征描绘了其他元素，但因欲避免冗长，我不作充分解释），经上说：\Quote{我人民的岁月，要像树木的岁月。} \BibleRef{依 65:22} 又哪个阅读圣经的人，不知道天主在哪里栽种了生命树，当我们的原祖父母因自己的邪恶被驱逐出乐园时，上主禁止他们吃树上的果子，并且树的周围设置了可怕的火篱笆？\par

但若有人争辩说，依撒意亚先知所提到的那些生命树的日子，是指基督教会现今的时代，而基督自己先知性地被称为生命树，因为祂是智慧，而关于智慧，撒罗满说：\Quote{智慧是生命树，给拥抱它的人}；\BibleRef{箴 3:18}并且若他们主张，我们原祖父母在乐园中并未度过\Emph{年岁}，而是很快被逐出，以致没有子女在那里诞生，因此那段时间不能被指称为\Quote{如在原始的日子，如在昔日的岁月}，我避免进入这个问题，以免因讨论一切而变得冗长，使整个主题悬而不决。因为我看到另一种含义，应阻止我们相信先知曾应许给我们恢复原始日子和旧律法祭献的岁月，作为极大的恩赐。因为在旧律下被选作祭品的动物必须是无玷的，没有任何瑕疵，并象征无罪的圣人，唯一具有这种品格的实例是在基督身上。因此，在审判之后，那些堪当如此净化的人，甚至将用火净化，成为完全无罪，并将在正义中将自身献给天主，确实是完全无玷、没有任何瑕疵的祭品，他们那时必将\Quote{如在原始的日子，如在昔日的岁月}，当时最纯洁的祭品被献上，正是这未来实体的预像。因为在圣人们的肉身和灵魂中，那时将有由这些祭品之身体所象征的纯洁。\par

然后，论到那些堪当不是净化而是受罚的人，祂说：\Quote{我要走近你们施行审判，我要迅速作证，攻击作恶的人和犯奸淫的人；}在列举其他该受罚的罪行后，祂接着说：\Quote{因为我是上主你们的天主，我不改变。}这好像祂说：虽然你们的过犯使你们变得更糟，而我的恩宠使你们变得更好，但我没有改变。祂说祂自己将是一位证人，因为在祂的审判中祂不需要证人；并且祂将是\Quote{迅速的}，或是因祂要突然来临，那似乎迟延的审判将因祂的出人意料而非常迅速，或是因祂将直接地、不须任何冗长说教就说服人的良心。因为经上记载：\Quote{在恶人的思念中，祂的查考必被进行。}\BibleRef{智 1:9}宗徒也说：\Quote{他们的思念互相控告或辩护，就在天主藉耶稣基督按我的福音审判人隐秘事的日子。}\BibleRef{罗 2:15}因此，主将是一位迅速的证人，当祂突然将那些将说服并惩罚良心的回忆带回来的时候。\par

% structure: p0083 source=h2 target=subsection reason=relative-heading-rank; base=h2

\subsection{第二十七章——论善人与恶人的分离，这宣告了最后审判的区分效力。}

此外，我以前从这位先知引用的另一处经文也指最后审判，他说：\Quote{他们必属于我，万军的上主说，在我收集我的财宝的日子里}等等。当这种区分义人与恶人的赏罚之差异，在那正义的太阳下，在永生光明中显现时——这种差异在此照耀着今世虚妄的太阳下是无法辨别的——那时将有一次前所未有的审判。\par

% structure: p0085 source=h2 target=subsection reason=relative-heading-rank; base=h2

\subsection{第二十八章——论梅瑟律法必须被属灵地理解，以避免按肉体解释的该受诅咒的抱怨。}

紧接着的话：\Quote{你要记住我仆人梅瑟的法律，即我在曷勒布山颁赐给他，为全以色列的法令和典章}\BibleRef{拉 4:4}，先知适当地提到了诫命和律例，在宣告了将来遵行法律者和藐视法律者之间的重要区别之后。他也意图使他们学习属灵地解释法律，并在其中找到基督，因祂的审判而将善人与恶人分开。因为主自己并非无缘故地对犹太人说：\Quote{如果你们信梅瑟，也必信我；因为他曾经论及我。}\BibleRef{若 5:46}因为他们按肉体接受法律，没有看出它地上的许诺是属灵事物的预像，便陷入了这样的怨言，竟敢说：\Quote{事奉天主是徒然的；遵守祂的诫命，在万军的上主面前行走哀伤，有什么益处呢？现在我们称外邦人为有福的；是的，作恶的人反而被高举。}\BibleRef{拉 3:14}正是他们这些话在某种程度上迫使先知宣布最后的审判，在那审判中，恶人将甚至表面上也不算有福，而明显极其悲惨；善人将不受甚至暂时的苦难压迫，而享受纯洁永恒的福乐。因为他先前曾引用了那些人的类似说法，他们说：\Quote{凡作恶的，在上主眼里看为善，并且祂喜悦他们。}\BibleRef{拉 2:17}我再说，正是因为他们按肉体理解梅瑟法律，他们才这样向天主抱怨。因此，第七十三篇圣咏的作者说，他的脚几乎失闪，他的脚步险些滑倒，因为他嫉妒罪人，观看他们的兴旺，以致他这样说，天主怎会知道？至高者有知识吗？又，我徒然圣化我心，在清白中洗手吗？他接着说，他试图解决这个最难的问题——当善人似乎受苦而恶人享福时——他的努力是徒然的，直到他进入天主的圣所，明白了最后的事。因为在最后审判中，情况不再如此；而是在义人明显的福乐与恶人明显的悲惨中，将显出完全不同的光景。\par

% structure: p0087 source=h2 target=subsection reason=relative-heading-rank; base=h2

\subsection{第二十九章——论厄里亚在审判前来临，为使犹太人因他的讲道和圣经解释而归向基督。}

在劝诫他们留心梅瑟法律之后，因为他预见他们很长时间内不会属灵地、正确地理解它，他接着说：\Quote{看，在上主伟大而可畏的日子来临以前，我派遣提市贝人厄里亚到你们这里来；他将使父亲的心转向儿子，使人的心转向他的近人，免得我来打击那地，完全毁灭。} \BibleRef{拉 4:5} 这是信友们常常谈论和心中所想的事：在审判之前的最后日子，犹太人将藉着这位伟大可敬的先知厄里亚而相信真基督，即我们的基督，因为他将向他们解释法律。我们并非没有理由盼望厄里亚在我们审判者与救主来临之前先来，因为我们有充分理由相信他现在还活着；正如圣经最清楚地告诉我们\BibleRef{列下 2:11}，他被火马车从这生命提去了。因此，当他来到时，他将对那法律做出属灵的解释——犹太人目前按肉体来理解法律——从而\Quote{使父亲的心转向儿子}，即父亲们的心转向他们的子女；因为七十贤士译本经常用单数代替复数。意思是，儿子们，即犹太人，将像父亲们，即先知们，其中包括梅瑟自己，那样理解法律。因为当子女们像父亲们那样理解法律时，父亲的心就转向了子女；当子女们持有与父亲们相同的见解时，子女的心就转向了父亲。七十贤士译本用了\Quote{使人的心转向他的近人}，因为父亲和子女彼此是最亲近的。在七十贤士译者的词语中可以找到另一种更可取的含义，他们以预言的眼光翻译了圣经，即厄里亚将使天主圣父的心转向圣子——当然不是说他促成圣父对圣子的这一爱，而是说他要使这一爱被知晓，并且那些先前曾憎恨圣子的犹太人，那时将爱圣子，即我们的基督。因为就犹太人而言，天主的心背离了我们的基督，这是他们对天主和基督的观念。但在犹太人方面，当他们自己回转心意，并认识到圣父对圣子的爱时，天主的心就转向了圣子。接着的话：\Quote{使人的心转向他的近人}——即厄里亚也将使人的心转向他的近人——我们如何能比作使人转向那人基督更好地理解呢？因为虽然按天主的形体他是我们的天主，但取了奴仆的形体，他屈尊成为我们的近人。因此，这就是厄里亚将要做的，\Quote{免得，}他说，\Quote{我来打击那地，完全毁灭。}因为那些思念地上事物的人就是地。至今肉体性的犹太人正是如此；因此他们向天主发怨言说：\Quote{恶人使他喜悦，}以及\Quote{事奉天主是徒然的。}\par

% structure: p0089 source=h2 target=subsection reason=relative-heading-rank; base=h2

\subsection{第三十章——在旧约经书中，当说到天主将审判世界时，基督的位格并未明确指明，但从某些上主天主说话的段落中清楚看出指的是基督。}

还有许多其他圣经章节论及天主的最后审判——事实上，若全部引用，本书将膨胀到不可原谅的篇幅。足以证明旧约和新约都宣告了审判。但在旧约中，不像新约那样明确宣告审判将由基督施行，即基督将从天降下为审判者；因为当旧约中由上主天主或祂的先知说上主天主将来到时，我们不一定理解为基督。因为圣父、圣子和圣神都是上主天主。然而，我们不能不提供证据。因此，我们必须首先说明耶稣基督如何在先知书中以上主天主的头衔说话，而同时毫无疑问是耶稣基督在说话；这样，在其他段落中，当这不是立即明显，而同时却说上主天主将来到那最后审判时，我们可以理解为是指耶稣基督。在依撒意亚先知书中有一处经文说明我的意思。因为天主藉先知说：\Quote{雅各伯和我所召叫的以色列，请听我说：我是元始，我是永在的；我的手奠定了大地，我的右手铺设了诸天。我要呼唤他们，他们便一同站立，并聚集聆听。谁向他们宣告了这些事？出于对你的爱，我向巴比伦行了你所喜悦的事，为除去加色丁人的子孙。我说了，也召叫了；我带领了他，使他的道路顺利。你们走近我，听这话：我从起初没有在隐秘中说话；当他们被造时，我在那里。现在，上主天主和祂的圣神派遣了我。} \BibleRef{依 48:12} 这是祂自己以上主天主的身份说话；然而，若祂没有加上\Quote{现在，上主天主和祂的圣神派遣了我}，我们就不会明白那是耶稣基督。因为祂说这话是指奴仆的形体，把将来的事说成过去的事，正如在同一位先知书中我们读到：\Quote{祂被牵去像一只待宰的羔羊，} \BibleRef{依 53:7} 而不是\Quote{祂将被牵去；}但用过去时态表达将来。预言经常这样说话。\par

在\WorkTitle{匝加利亞先知书}中还有另一段经文，它明确宣告全能者派遣了全能者；除了天主圣父和天主圣子，这些话还能理解为何人呢？因为经上记载：\BibleQuote{全能的上主这样说：在荣耀之后，祂派遣了我到那掠夺你们的万邦；因为触摸你们的，就是触摸祂眼中的瞳人。看，我要伸手攻击他们，他们必成为他们仆人的掠物；那时你们便知道全能的上主派遣了我。} \BibleRef{匝 2:8} 请注意，全能的上主说全能的上主派遣了祂。除了\Person{基督}，谁还能理解这些话呢？祂正向以色列家迷失的羊说话。因为祂在福音中说：\BibleQuote{我受派遣，无非是为以色列家迷失的羊。} \BibleRef{玛 15:24} 祂在这里将以色列家比作天主眼中的瞳人，以表明极深的爱。宗徒们自己也属于这类羊。但在祂复活——即荣耀——之后（因为在复活之前，福音作者说\BibleQuote{\Person{耶稣}还没有受光荣} \BibleRef{若 7:39}），祂藉着宗徒们被派遣到万邦中；如此便应验了圣咏的话：\BibleQuote{祢救我脱离人民的反对，立我作万民的首领。} 这样，那些掠夺以色列人、并使以色列人在被征服时服事他们的民族，自己不再以同样方式被掠夺，反而自身成为以色列人的掠物。因为主曾应许宗徒们说：\BibleQuote{我要使你们成为捕人的渔夫。} \BibleRef{玛 4:19} 又对其中一位说：\BibleQuote{从今以后，你要捕人。} \BibleRef{路 5:10} 于是他们成为掠物，但这是好的意义上的掠物，就像那壮士被更强的一位捆绑后从他手中被夺回的人一样。\BibleRef{玛 12:29}\par

同样，主藉着同一位先知说：\BibleQuote{在那一天，我要设法毁灭凡前来攻击耶路撒冷的万邦。我要把恩宠与慈悲之神倾注在\Person{达味}家和\Place{耶路撒冷}的居民身上；他们必仰望我，因为他们曾侮辱了我；他们必为祂哀悼，如哀悼一位最亲爱的人，为祂悲痛，如悲痛独生子。} \BibleRef{匝 12:9} 毁灭所有敌对圣城\Place{耶路撒冷}并\Quote{前来攻击它}——即反对它，或如某些译本所谓\Quote{前来压服它}，仿佛要将它踏在脚下——的万邦，以及将恩宠与慈悲之神倾注在\Person{达味}家和\Place{耶路撒冷}的居民身上，这岂不是惟独属于天主的作为？这无疑属于天主，先知也将这些话归于天主；然而\Person{基督}却表明自己就是施行这些伟大神圣事工的天主，因为祂接着说：\BibleQuote{他们必仰望我，因为他们曾侮辱了我；他们必为祂哀悼，如哀悼一位最亲爱的人（或爱子），为祂悲痛，如悲痛独生子。} 因为在那一天，犹太人——至少那些接受恩宠与慈悲之神的犹太人——当他们看见祂以威严降来时，并认出祂就是他们（藉其祖先）在祂初次谦卑来临时所侮辱的那位，他们将懊悔在祂受难时侮辱了祂。而他们的祖先，即犯下这极大不敬之罪的元凶，当他们复活时将看见祂；但这对他们只是惩罚，而非改正。以下的话并不是指他们：\BibleQuote{我要把恩宠与慈悲之神倾注在\Person{达味}家和\Place{耶路撒冷}的居民身上，他们必仰望我，因为他们曾侮辱了我；} 而是指他们的后裔，即那时要藉着\Person{厄里亚}而相信的人。正如我们对犹太人说：你们杀害了基督，虽然真正动手的是他们的祖先，同样，这些人也要因他们（在某种意义上）做了他们祖先所做的事而哀悼。因此，那些接受慈悲与恩宠之神而相信的人，虽然不会与他们的不虔敬的祖先一同被定罪，但他们仍要哀悼，仿佛他们自己做了祖先所做的事。他们的悲伤不是出于罪疚，而是出于孝爱之情。确实，\WorkTitle{七十贤士译本}所译的\BibleQuote{他们必仰望我，因为他们曾侮辱了我，} 在希伯来文中是\BibleQuote{他们必仰望我，就是他们所刺透的。} 这个字更清楚地指出了\Person{基督}的十字架受难。但\WorkTitle{七十贤士译本}的译者更倾向提及祂整个苦难中所包含的侮辱。因为事实上，他们在祂被捕时、被捆绑时、被审判时、被穿上紫红袍并屈膝朝拜时、被戴上茨冠并被苇秆击打头部时、在祂背十字架时，以及最后被悬在树上时，都侮辱了祂。因此，当我们不局限于一种解释，而是将两者结合起来，既读\Quote{侮辱}又读\Quote{刺透}时，我们就能更全面地认识主的苦难。\par

因此，当我们读到先知书上说天主末时要来执行审判时，仅凭\Quote{审判}一词，即使别无其他确定其含义的内容，我们也必须推知指的是基督；因为虽然圣父将审判，但祂将藉圣子的来临而审判。因为祂自己，藉其显现的临在，\Quote{不审判任何人，却将一切审判交给了圣子；}\BibleRef{若 5:22} 因为正如圣子曾作为人受审判，祂也要以人的形体审判。因为天主藉依撒意亚以雅各伯和以色列之名所说者，无非就是祂；基督从那后裔取了肉身，正如经上所载：\Quote{雅各伯是我的仆人，我必扶持祂；以色列是我所拣选的，我的灵已覆庇祂：我已将我的灵赐予祂；祂必将审判带给外邦人。祂不呼喊，不喧嚷，也不使自己的声音被听见于外。压伤的芦苇，祂不折断；将熄的灯芯，祂不吹灭；祂要凭真理施行审判。祂要照耀，不被折断，直到祂在地上设立审判；海岛都仰望祂的名。}\BibleRef{依 42:1} 希伯来原文并没有\Quote{雅各伯}和\Quote{以色列}；但\WorkTitle{七十子译本}的译者，为了表明\Quote{我的仆人}这一表述的意义，并指明它是指那至高者自己谦卑所取的仆人之形像，便插入了那位的人名，基督从其后裔取了仆人的形像。圣神赐予了祂，并如福音所记载，以鸽子的形状显现。\BibleRef{若 1:32} 祂将审判带给外邦人，因为祂预言了外邦人所不知道的事。祂以柔和之气，不呼喊，也不停止宣讲真理。但祂的声音未曾被听见，也未被听见于外，因为那些在祂身体之外的人不听从祂。连那迫害祂的犹太人，祂也没有折断；他们虽如压伤的芦苇失去了完整，如将熄的灯芯其光已灭，但祂宽恕了他们，因为祂来是为受审判，而非施行审判。祂凭真理施行了审判，宣告他们若坚持作恶，将受惩罚。祂的面容在山上发光，\BibleRef{玛 17:1}；祂的名声传遍世界。祂未被折断，也未被克服，因为在祂自身或在祂的教会中，迫害未能摧毁祂。因此，敌人过去所说或现在所说的：\Quote{祂几时死，祂的名何时消灭？} 这事未曾发生，也绝不会发生。\Quote{直到祂在地上设立审判。} 看，我们所寻求的隐密事情被发现了。因为这就是最后的审判，当祂从天降临时，祂将在地上设立它。而且，在祂身上，我们已经看到预言的最后一句应验了：\Quote{列国都仰望祂的名。} 由于这无人能否认的应验，人们被鼓励去相信那最无耻地被否认的事。因为，谁能期望那连尚未相信基督的人现在也看到在我们当中应验了的事呢？这事如此确凿，令他们只能咬牙切齿、日渐憔悴。我说，谁能期望列国会仰望基督的名呢？当祂被逮捕、捆绑、鞭打、戏弄、钉死时，连门徒们自己也失去了起初对祂怀有的希望。那时只被十字架上的一个强盗所怀有的希望，如今被地上列国的人们所珍视，他们都被印记了祂在其上死亡、使他们得以永生的十字架的记号。\par

那么，最后的审判将由耶稣基督按圣经所预言的方式施行，这是无人否认或怀疑的，除非那些因某种难以置信的仇恨或盲目，拒绝相信这些经书的人——尽管其真实性已向全世界证实。在那一审判之时或与之相关，将发生以下事件，正如我们所知：提市贝特人厄里亚要来；犹太人将相信；假基督将迫害；基督将审判；死人将复活；善人和恶人将被分开；世界将被焚烧并更新。我们相信这一切都将发生；但其方式或顺序，人的理智无法完全教导我们，只有事件本身的经历才能。然而，我的意见是，它们将按我叙述的顺序发生。\par

我还有两本书待写，以靠天主的助佑完成我所应许的。其中之一将解释恶人的惩罚，另一本则论义人的福乐；在这些书中，我将特别费心，藉天主的恩宠，驳斥那些可怜之人自认为可以削弱天主应许与警告的论据，并将最有益于信德滋养的陈述嘲笑为空言。但那些通达神圣之事的人，视天主的真理和全能作为维护这些事情最有力的论据；这些事尽管在人看来难以置信，却包含在圣经中，而圣经的真理已从多方面得到证明；因为他们确信，天主绝无谎言，并且祂能做不信者认为不可能的事。\par

\SourceRecord{Translated by Marcus Dods.；Nicene and Post-Nicene Fathers, First Series；Vol. 2.；Edited by Philip Schaff.；Buffalo, NY: Christian Literature Publishing Co.,；1887.；<http://www.newadvent.org/fathers/120120.htm>.}

% structure: p0001 source=h1 target=section reason=page-title

\section{\WorkTitle{天主之城（第二十一卷）}}

论为魔鬼之城所保留的结局，即恶人的永罚；以及不信者对此提出的异议。\par

% structure: p0003 source=h2 target=subsection reason=relative-heading-rank; base=h2

\subsection{第一章——论讨论的顺序，要求我们先论丧亡者与魔鬼同受永罚，然后论圣人的永福。}

我愿凭天主所赐的能力，在这本书中更详尽地讨论那要加给魔鬼及其所有随从的惩罚的性质，当时那两座城——一座属于天主，另一座属于魔鬼——将藉着审判生者死者的主耶稣基督达到它们各自的结局。我采取这个顺序，并宁愿先论魔鬼的惩罚，后论圣人的福乐，因为身体分享这两种命运；身体在永久的折磨中忍受，比起身体在永福中毫无痛苦地持续存在，似乎更令人难以置信。因此，当我证明那惩罚不该是不可信的之后，这将大大有助于我证明那更可信的事，即圣人的身体脱离一切痛苦而不朽。这个顺序也与神圣的著作不矛盾，其中有时先论善人的福乐，如经上说：\Quote{行善的，复活得生；作恶的，复活定罪}\BibleRef{若 5:29}；但有时也后论，如：\Quote{人子要差遣祂的天使，把祂国里一切使人跌倒的事和作恶的人收集起来，丢在火窑里；那里要有哀号和切齿。那时，义人要在他们父的国里发光如同太阳}\BibleRef{玛 13:41}；以及\Quote{这些人要进入永罚，而义人却要进入永生}\BibleRef{玛 25:46}。虽然我们无暇列举实例，但任何研读先知书的人都会发现，他们时而采用这一顺序，时而采用另一顺序。我采用后一种顺序的原因已如上述。\par

% structure: p0005 source=h2 target=subsection reason=relative-heading-rank; base=h2

\subsection{第二章——身体能否在烈火中永久存留？}

那么，我能提出什么来说服那些拒绝相信人类有生命和灵魂的身体不仅能在死后存活，而且能在永火折磨中持续存在的人呢？他们不允许我们简单地将此归因于全能者的能力，而要求我们用某些实例说服他们。那么，如果我们回答他们说，有些动物固然是可朽坏的，因为它们会死，但却能在火焰中生活；同样，在热得无人能安全把手伸进去的温泉中，发现有一种蠕虫不仅在那里生活，而且不能在其他地方生活；他们要么拒绝相信这些事实，除非我们能指给他们看；要么，如果我们能够通过亲眼目睹或充分的证词证明这些事实，他们又以同样的怀疑态度争辩说，这些事实并非我们所寻求证明的例子，因为这些动物不会永远活着，而且它们在火焰中生活时没有痛苦，火元素与它们的本性相宜，使之兴旺而非受苦——就好像在这样的环境中，兴旺比受苦更不可思议似的。奇怪的是，任何东西在火中受苦却仍活着，但更奇怪的是，它在火中活着却不受苦。那么，如果后者可信，为什么前者不可信呢？\par

% structure: p0007 source=h2 target=subsection reason=relative-heading-rank; base=h2

\subsection{第三章——身体的痛苦是否必然导致肉身的毁灭？}

但那些人说，没有身体能够受苦而不能死。我们怎么知道呢？因为当魔鬼承认自己备受折磨时，谁能断言魔鬼不在他们的身体中受苦呢？如果回答说，没有地上的身体——即没有坚固可感知的身体，或一言以蔽之，没有肉身——能够受苦而不能死，这不只是告诉我们人们从经验和身体感官中收集到的东西吗？因为他们除了那会死的肉身外，确实不认识任何肉身；而他们的全部论据就是，他们认为自己所没有经验过的事是完全不可能的。因为我们不能把以痛苦作为死亡的前提称为推理，而事实上痛苦反倒是生命的标志。因为虽然受苦者是否能永远活着是一个问题，但可以肯定的是，每一个感到痛苦的事物都活着，而且痛苦只能存在于一个活的主体中。因此，受苦者必然是活着的，而不是痛苦必然杀死他；因为即使是我们这些注定会死的可朽身体，也并非每一种痛苦都会杀死它们。而某种痛苦杀死它们，是由于灵魂与身体的结合方式使灵魂屈服于巨大的痛苦而离去；因为我们肢体的结构以及生命器官如此脆弱，无法承受那引起巨大或极度的剧痛。但在来世，灵魂与身体的结合是如此，它既不会因时间的流逝而解散，也不会因任何痛苦而破裂。因此，虽然在这个世界上确实没有能受苦却不能死的肉身，但在来世，却会有如今所没有的肉身，正如也会有如今所没有的死亡。因为死亡不会被废除，而会是永久的，因为灵魂既不能享受天主而活，也不能死去而逃脱身体的痛苦。第一次死亡违背灵魂的意愿把灵魂从身体中驱出；第二次死亡违背灵魂的意愿把灵魂禁锢在身体中。这两者有共同之处：灵魂违背自己的意愿遭受自己的身体所施加的痛苦。\par

我们的对手也强调这一点，即在这个世界上，没有能受苦却不能死的肉体；但他们却忽略了有比身体更伟大的东西。因为灵魂，其临在使身体有生命并支配身体，既能受苦，却不能死。因此，有一种东西虽能感到痛苦，却是不朽的。这种能力，我们现在在所有人的灵魂中看到，将来也会在受罚者的身体中。此外，如果我们更仔细地思考这个问题，就会发现所谓的身体痛苦更应该归于灵魂。因为感到痛苦的是灵魂，而不是身体，即使痛苦源于身体——灵魂在身体受伤的地方感到痛苦。正如我们说身体有感觉、有生命，尽管身体的感觉和生命来自灵魂，同样我们也说身体受苦，尽管没有灵魂，身体不能承受任何痛苦。因此，灵魂与身体一起受苦，在身体受到伤害的部分；而灵魂独自受苦，即使它在身体内，当某种不可见的原因使灵魂痛苦而身体安然无恙时。即使不与身体结合，灵魂也会受苦；因为那富翁在地狱里确实受苦，当他说：\BibleQuote{我在这火焰里受折磨。} \BibleRef{路 16:24} 但至于身体，当它没有灵魂时，它不会受苦；即使有灵魂，也只能通过灵魂的受苦来受苦。因此，如果我们能公平地从痛苦的存在推断死亡的存在，并得出结论说，哪里能感到痛苦，哪里就能发生死亡，那么死亡更应该是灵魂的属性，因为痛苦更特别属于灵魂。但是，既然那最受苦的东西不能死，那么凭什么认为那些身体因为注定要受苦，就注定要死呢？柏拉图派确实主张，这些尘世的身体和必死的肢体引起了灵魂的恐惧、欲望、悲伤和快乐。\Quote{因此，} 维吉尔说（\Emph{即}，从这些尘世的身体和必死的肢体中），\par

\Quote{因此狂野的欲望和卑下的恐惧，\newline{}以及人的欢笑，人的眼泪。}\par

但在本书第十四卷中，我们已经证明，按照柏拉图派自己的理论，即使灵魂洗除了身体的一切污染，仍会怀着一种怪异的欲望要回到身体中。但哪里有可能存在欲望，也必然存在痛苦；因为欲望受挫，无论是因未达到目标还是失去已获得之物，都会转为痛苦。因此，如果灵魂——它是唯一的或主要的受苦者——本身仍拥有某种不朽性，那么就无法推论说，因为受罚者的身体将要受苦，所以它们将会死亡。总之，如果身体导致灵魂受苦，为什么身体不能同样导致死亡，而只能导致痛苦呢？除非是因为，导致痛苦的并不一定也导致死亡。那么，为什么不能相信，这些火既能引起我们所说的那些身体的痛苦，却不能引起死亡，正如身体本身引起灵魂的痛苦，却不一定引起死亡？因此，痛苦并不必然预示着死亡。\par

% structure: p0012 source=h2 target=subsection reason=relative-heading-rank; base=h2

\subsection{第四章——来自自然的例证，证明身体可以在火中不被焚毁而继续活着。}

因此，如果火蜥蜴在火中生活，如博物学家所记载的那样，并且如果西西里的某些著名山脉从远古至今一直燃烧，却仍然完整，这些就是足够令人信服的例证，证明并非一切燃烧的东西都会被烧毁。同样，灵魂也证明，并非一切能受苦的东西都能死。那么，他们为什么还要要求我们提供真实的例证，来证明被判永久惩罚之人的身体在火中仍然保留灵魂、燃烧而不被烧毁、受苦而不消亡并非不可信呢？因为适合的特性将被赋予肉身的实体，由那位赋予我们所见的万物如此奇妙且多样的特性的主，其数量之多反而使我们不再惊奇。因为除了万物的造物主天主，谁给了孔雀肉防腐的特性？当我最初听到这种特性时，觉得难以置信；但碰巧在迦太基，有人烹饪并端给我一只这种鸟，我从其胸部取了一小块合适的肉，命令保存起来；当它被保存了足够使其他任何肉变臭的天数之后，被取出放在我面前，并没有发出难闻的气味。而且，在放置了三十多天后，它的状态仍然相同；一年之后，仍然如此，只是稍微更干瘪一些。谁给糠这么强大的冷冻能力，使它能够保存埋在它下面的雪，又有这么强大的加热能力，使它能够使未熟的水果成熟？\par

但是谁能解释火本身的奇特属性呢？火本身明亮，却将所烧的一切熏黑；火本身颜色极为绚烂，却几乎污染它所接触并吞噬的一切，将熊熊燃烧的燃料变为灰黑的灰烬？不过这并非一条绝对统一的规律；因为相反，在炽热的火中烧制的石头本身也会发光；尽管火略带红色，而石头呈白色，但白色与光明相合，黑色与黑暗相合。因此，虽然火在烧制石头时燃烧了木材，但这些相反的效果并非由材料的对立性造成。因为木材和石头虽然不同，但它们并非如同黑白那样互相对立；同一火的作用在石头中产生一种颜色，在木材中产生另一种颜色：火将自身的光明给予石头，却将木材熏黑；而且如果没有木材作燃料，火就无法对石头产生任何作用。那么，在木炭中我们又发现何等奇妙的神奇特性呢？木炭如此脆弱，轻轻一击便碎裂，稍加压力便化为粉末；然而它却又如此坚固，湿气不能使它腐烂，时间也不能使它朽坏。它如此耐久，以至于人们在立地界标时习惯在下面放置木炭；这样，即使过了最长的间隔，如果有人提起诉讼并辩称没有界石，他就可以被下面的木炭所证明。那么，究竟是什么使得木炭在潮湿的土壤中（原来的木材在其中腐烂）久久不烂呢？岂非这同一场吞噬万物的火吗？\par

再让我们考虑石灰的奇迹。除了它在使其他东西变黑的火中变白（这一点我已经说得够多了）之外，它还有一种神秘的属性，即在它内部孕育着火。它本身摸起来是冷的，却潜藏着火，这火并不立即被我们的感官察觉，但经验告诉我们，它仿佛沉睡在其中，即使看不见。正因如此，它被称为\Quote{生石灰}，仿佛火是那无形的灵魂，使这可见的实体或身体具有生命。但奇妙的是，这火是在它被熄灭时点燃的。因为为了释放隐藏的火，石灰被浸湿或用水浸泡，然后，虽然它之前是冷的，却由于那种通常冷却热物的处理而变热。仿佛火正在离开石灰并呼出它的最后一息，它不再隐藏，而是显现出来；然后石灰处于死亡的冰冷之中，无法再复活，我们以前称之为\Quote{生}，现在称之为\Quote{熟}。还有什么比这更奇怪的呢？然而还有一个更大的奇迹。因为如果你不用水，而是用油（油是火的燃料）来处理石灰，那么无论多少油都不能使它变热。现在，如果这个奇迹是关于某种我们无法实验的印度矿物的，我们可能会立即认为它是假的，或者肯定会感到非常惊讶。但是，那些每天呈现在我们眼前的事物，我们却轻视它们，不是因为它们实际上不那么奇妙，而是因为它们平常；以至于甚至印度本地产的一些东西，尽管离我们很远，一旦我们能够悠闲地欣赏它们，它们也就不再引起我们的赞叹了。\par

钻石是一种许多人拥有的石头，尤其是珠宝商和宝石匠；这种石头非常坚硬，既不能用铁也不能用火加工，而且据说除了山羊血之外，没有任何东西能加工它。但是，您认为那些拥有它并熟悉其特性的人，和那些第一次见到它的人，对它一样赞叹吗？没有见过它的人或许不相信关于它的说法，或者如果相信，也会像遇到超出经验的事物那样惊叹；而如果他们碰巧看见了，他们仍然会因不习惯而惊叹，但逐渐熟悉它的经验会减弱他们的赞叹。我们知道磁石有吸引铁的神奇力量。当我第一次看到时，我感到震惊，因为我看到一个铁环被磁石吸引并悬挂着；然后，仿佛磁石将它自身的特性传递给了它所吸引的铁，并使它变成了与自己类似的物质，这个铁环靠近另一个铁环，并将其提起；正如第一个铁环附着在磁石上，第二个铁环也附着在第一个上。第三个和第四个同样被添加，因此从磁石上悬挂着一串铁环，它们的环连接在一起，不是互相套接，而是通过外表面附着在一起。谁不会对这种石头的力量感到惊叹呢？这种力量不仅存在于自身，而且通过许多悬挂的铁环传递，用看不见的纽带将它们连接在一起。然而，我从我的主教同仁，米莱维主教塞维鲁那里听到的关于这种石头的事情更加惊人。他告诉我，曾任阿非利加伯爵的巴塔纳里乌斯，当主教与他共进晚餐时，拿出一块磁石，放在一个银盘下面，盘上放了一小块铁；然后他用手在盘下移动磁石，盘上的铁也随之移动。中间的银完全没有受到影响，但无论磁石在下面多快地来回移动，上面的铁都被吸引。我已经叙述了我亲眼所见的事；我也叙述了从那位我信任如同自己眼睛的人那里听到的事。让我进一步说说我关于此磁石读到的东西。当钻石靠近它时，它就不吸铁；或者如果它已经吸起了铁，一旦钻石靠近，它就会松开。这些石头来自印度。但是，如果我们因为熟悉而不再赞叹它们，那么那些非常容易获得并送到我们这里的人，又怎么会赞叹它们呢？也许他们视它们如我们视石灰一样便宜，因为石灰很普通，我们毫不在意，尽管它具有奇怪的特性：当水（通常灭火）浇在它上面时，它会燃烧；而与油（通常助火）混合时，它却保持冷却。\par

% structure: p0017 source=h2 target=subsection reason=relative-heading-rank; base=h2

\subsection{第五章 有许多事理性无法解释，却仍然是真实的。}

然而，当我们宣告天主已经施行或将要施行、且无法呈现在人眼前的奇迹时，怀疑者一直要求我们用理性解释这些奇事。因为我们无法做到——它们超出人类理解——他们就以为我们在说谎。因此，这些人自己应当说明我们能够或确实看见的一切奇事。如果他们认识到这是人所不能的，就应当承认：不能因为某事无法与理性相符，就断定它从前不曾存在或将来不会存在，因为现在就有许多同样无法与理性相符的事物存在。因此，我不打算详述书籍中记载的众多奇事——这些事并非指那些发生过就消逝的事物，而是在某些地方恒常存在，任何人若有兴趣和机会，都可证实其真实性——我只叙述少数几件。以下就是人们告诉我们的若干奇迹：——\Place{西西里}的\Place{阿格里真托}的盐，投入火中便像在水中一样化为液体，但在水中却像在火中一样噼啪作响。\Place{加拉曼泰}人有一泉水，白昼冷得无人能饮，夜间热得无人能触。\Place{伊庇鲁斯}也有一泉水，与所有其他泉水一样能熄灭点燃的火炬，但又与众不同：它能点燃已熄灭的火炬。在\Place{阿卡迪亚}发现一种石头，称为\OriginalTerm{不燃石}{asbestos}，因为一旦点燃就无法熄灭。某种埃及无花果树的木材沉入水中，不像其他木材那样漂浮；更奇怪的是，沉入水底一段时间后，它又浮上水面，尽管自然规律要求木材浸水后应变得更重。还有\Place{索多玛}的苹果，表面看似成熟，但当用手或牙齿触碰时，果皮裂开，化为尘土和灰烬。\OriginalTerm{波斯黄铁矿}{Persian stone pyrites}紧握在手中时会烧手，因此得名于火。在\Place{波斯}还发现另一种石头，称为\OriginalTerm{月长石}{selenite}，因其内部光泽随月相盈亏而增减。在\Place{卡帕多西亚}，母马受风而孕，马驹仅活三年。\Place{提隆}，一个印度岛屿，优于所有其他土地：其上生长的树木永不落叶。\par

这些记载在历史中的、并非关于过去事件而是关于恒久地方的无数奇事，我没有时间详述而偏离我的主要目的；但让那些拒绝相信圣经的怀疑者，如果他们能够，给我一个对这些奇事的理性说明吧。因为他们不信圣经的唯一理由，就是其中含有不可信的事，恰如我刚才叙述的那些。因为，他们说，理性不能承认肉体会被焚烧而不化为灰烬，会受苦而不死亡。真是了不起的理性主义者，竟能给出一切存在奇迹的理性解释！那么，让他们给我们解释一下我们引述的这几件事吧。如果它们不知道这些事是存在的，而只是被我们告知它们将在未来某个时候发生，那么他们相信它们的可能性，会比现在因我们的话语而拒绝相信的那些事还要低。因为他们中间有谁会相信我们呢？如果我们不说将来人的肉身将如此：能忍受永久的痛苦和火而不死亡，反而说来世将有盐在火中变为液体就像在水中一样，在水中噼啪作响就像在火中一样；或者有一泉水，其水在夜间寒风中热得不能触摸，而在日间炎热中冷得不能饮用；或者有一石头紧握时因其自身的热而灼伤手，或者有一石头一旦点燃任何部分就无法熄灭；或者我引述的任何奇迹——我还没提数不清的其他奇迹呢？如果我们说这些事将在来世出现，而我们的怀疑者回答说：\Quote{如果你们要我们相信这些事，那就请用理性为我们一一解释清楚}，我们就应当承认我们无法做到，因为人类软弱的理解力无法掌握天主工作中的这些以及类似的奇迹；然而，我们的理性完全确信全能者所做的一切无不有其理由，尽管人脆弱的头脑无法解释这理由；并且，虽然我们在许多情况下不确定祂的意图，但始终最确定的是：祂所意图的没有一件对祂是不可能的；而且，当祂宣告祂的心意时，我们相信祂，因为我们不能相信祂是软弱或虚假的。然而，这些对信德吹毛求疵、要求理性解释的人，是如何处理那些无法给出理性解释、却实际存在、且显然与事物本性相矛盾的事物呢？如果我们宣告这些事将要发生，这些怀疑者会向我们要求它们的理性解释，就像他们对待我们宣告为注定要发生的事那样。因此，正如这些现存的奇事并非不存在，尽管人的理性和言语在天主的这些工程中迷失，同样，我们所说的事物也不会因为无法解释就是不可能的；因为在这一方面，它们与地上的奇事处于相同的境地。\par

% structure: p0020 source=h2 target=subsection reason=relative-heading-rank; base=h2

\subsection{所有奇迹并非都出于自然，有些是出于人的技巧，有些是出于魔鬼的计谋。}

此时他们或许会回答说：\Quote{这些事并不存在；我们一件也不相信；它们都是旅行者的传闻和虚构的故事；}并且他们可能加上看似有理的论据，说道：\Quote{如果你相信诸如此类的事，那么也相信同一书籍中所记载的，在\Person{维纳斯}庙中有一个露天放置的\OriginalTerm{烛台}{candelabrum}，上面有一盏灯，燃烧得如此强烈，以至于没有风暴或雨水能熄灭它，因此它像上面提到的石头一样被称为\OriginalTerm{石棉}{asbestos}或不灭灯。}他们说这话也许是为了使我们陷入两难：若我们说这事不可信，那么我们就会否定其他记录中的奇迹的真实性；若反之我们承认此事可信，那么我们就会拥护外教的神祇。但是，正如我在本作品第十八卷中已经说过的，我们并不认为有必要相信一切世俗历史所包含的内容，因为正如\Person{瓦罗}所说，连史学家本身在许多点上都不一致，使人以为他们有意且费力地如此做；但我们愿意相信那些与这些书籍不相矛盾的内容，我们毫不迟疑地说我们必须相信这些书籍。至于那些持久的自然界奇迹，我们想以此说服怀疑者关于未来世界的奇迹，那些我们自身能观察到或不难找到可靠见证的奇迹已足够用于我们的目的。此外，那座\Person{维纳斯}庙及其不灭灯非但没有使我们陷入困境，反而为我们的论点开辟了有利的领域。因为这盏不灭灯，我们可以加上许多由人类或魔法——即由受魔鬼影响的人，或由魔鬼直接——造成的奇迹，因为我们不能否定这些奇迹而不否定我们所相信的圣经的真理。因此，那盏灯要么是通过某种机械和人工装置装上了\OriginalTerm{石棉}{asbestos}，要么是通过魔法技艺安排以使敬拜者惊讶，要么是某个以\Person{维纳斯}之名的魔鬼如此显著地显现自己，以致这一奇事既开始又持久。魔鬼们通过向它们提供适合它们各种口味的受造物（天主的受造物，不是它们的）而被吸引居住在某些庙宇中。它们不像动物那样被食物吸引，而是像神灵一样，被适合它们口味的象征——各种石头、木材、植物、动物、歌曲、礼仪——所吸引。为了使人们提供这些吸引物，魔鬼首先狡猾地引诱他们，要么用秘密的毒药渗透他们的心，要么以友善的面貌显现自己，从而使少数人成为它们的门徒，这些人再成为大众的教师。因为除非它们首先教导人，否则不可能知道每个魔鬼想要什么、害怕什么、应以什么名字被呼召或约束而出现。因此产生了魔法和巫师。但最重要的是，它们拥有人的心，并尤其以在变成光明天使时拥有此权而骄傲。因此，许多事情是它们做的；这些事我们越认为惊人，就越应当小心避开。然而，这些事本身却能推进我现在的论点。因为如果这些奇迹是由不洁的魔鬼完成的，那么圣天使的能力该有多大啊！那位使天使本身能行奇迹的天主，又能做什么呢！\par

因此，如果许多效果可以由人的技艺创造出来，其惊奇程度使未入门者以为它们是神迹，例如，在某座庙中，两个磁铁被调整好，一个在屋顶，另一个在地面，从而使一个铁像悬在两者之间的空中，若不知上下有磁铁，有人会以为这是神的力量；或者，如我们已提到的那盏\Person{维纳斯}灯，它是对\OriginalTerm{石棉}{asbestos}的巧妙运用；再者，如果借助那些圣经称之为巫术师和行邪术者的巫师，魔鬼能获得如此大的能力，以至于高贵的诗人\Person{维吉尔}认为自己有理由用以下诗句描述一位极强大的巫师：\par

她的魅力能治愈她所愿的灵魂，\newline{}夺去其他心灵的康宁，\newline{}使河流倒流回源头，\newline{}并使星辰忘记轨道，\newline{}并从黑夜召唤鬼魂：\newline{}大地将在你脚下轰鸣：\newline{}山梣将离开原位，\newline{}并从高处移下；——\newline{}\par

如果是这样，那么天主更能做那些对怀疑者来说不可思议，但对祂的能力却很容易的事，因为正是祂赐予了石头和一切事物它们的力量，并赐予人巧妙运用它们的能力；祂赐予了天使一种比地球上一切生物更强大的本性；祂的能力超越一切奇迹，祂在作为、规定和允许中的智慧，在治理万物时与在创造万物时一样令人惊叹！\par

% structure: p0025 source=h2 target=subsection reason=relative-heading-rank; base=h2

\subsection{第七章——相信奇迹的根本原因是造物主的全能。}

那麼，天主——祂所創造的世界充滿了天空、大地、空氣和水中無數的奇蹟，而這個世界本身就是一個無疑比它所充滿的一切奇蹟更偉大、更令人讚嘆的奇蹟——為什麼不能既使死者的身體復活，又使受罰者的身體在永不熄滅的火中受折磨呢？但那些與我們爭論或反對我們的人，他們既相信有一位創造世界的天主，又相信祂所創造的諸神作為祂的代治者管理世界的法則；我們的對手，我說，他們絕不否認，反而斷言世界上有能產生驚人效果的力量（無論是自發的，還是因某種儀式或祈禱而被召喚的，或是通過某種魔法的方式）；當我們向他們陳述其他事物的奇妙特性時，這些事物既不是理性的動物，也不是理性的靈體，而是像我們剛才引用的那些物質對象；他們習慣回答說：\Quote{這是它們的自然屬性，它們的本性；這些是它們天然擁有的力量。}因此，阿格里真托的鹽在火中溶解、在水中噼啪作響的全部理由就是：\Quote{這是它的本性。}然而，這似乎有悖於本性，因為本性賦予水而非火融化鹽的能力，賦予火而非水燒焦鹽的能力。但他們說，這是\Emph{這種}鹽的天然屬性，表現出與之相反的效果。出於同樣的理由，他們解釋了加拉曼提亞的泉水，同一條小溪白天冰冷，夜間沸騰，以至於兩種極端都無法觸摸。還有另一處泉水，雖然摸起來是冷的，像其他泉水一樣熄滅點燃的火炬，但不像其他泉水，卻以驚人的方式點燃熄滅了的火炬。還有石綿石，雖然本身沒有熱量，但被火點燃後卻無法熄滅。其餘的我不勝枚舉，在這些事物中，雖然似乎有一種與本性相反的特殊屬性，但除了說\Quote{這是它們的本性}之外，沒有給出其他理由——這確實是一個簡短的理由，我承認，是一個令人滿意的回答。但是，既然天主是一切本性的創造者，為什麼我們的對手在拒絕相信我們所肯定的（因為他們認為這不可能）時，不願意接受我們比他們自己更好的解釋，即這是全能天主的旨意——因為祂被稱為全能，正是因為祂有能力做一切祂願意做的事——祂能夠創造如此多的奇蹟，不僅是未知的，而且是眾所周知的，正如我一直展示的那樣；如果這些奇蹟不是我們親眼所見，或者不是由近期可信的證人報導的，它們肯定會被認為是不可能的？至於那些除了作者在書中記載之外沒有任何其他證明的奇蹟，而這些作者並非受天主啟示寫作，因此容易犯人的錯誤，我們不能公正地責備任何拒絕相信它們的人。\par

就我自己而言，我不希望我所引用的所有奇蹟都被輕率地接受，因為我自己也不是全部盲目相信，除非那些是我親眼所見的，或者任何人都容易驗證的，例如：石灰被水加熱、被油冷卻；磁鐵以其神秘而不可感知的吸力吸引鐵，但對稻草毫無作用；孔雀的肉戰勝腐敗，而\Person{柏拉圖}的肉也無法免於腐敗；穀殼如此寒冷，以至於阻止雪融化，又如此炎熱，以強迫蘋果成熟；熾熱的火，按照其熾熱的外觀，將它烘烤的石頭變白，而與其熾熱的外觀相反，它將它燃燒的大部分東西變黑（正如油，無論多麼純淨，也會留下污漬，而白色銀子畫出的線條是黑色的）；還有木炭，通過火的作用與原來如此徹底不同，以至於一塊紋理清晰的木頭變得難看，堅韌的變得脆弱，腐朽的變得不可腐敗。其中一些事物，我和許多其他人一樣知道；另一些，我和所有人一樣知道；還有許多其他的事物，我沒有空間在這本書中插入。但在那些我所引用的、我本人沒有親眼見過、只是從書中讀到的事物中，除了那處泉水——其中燃燒的火炬被熄滅，熄滅的火炬被點燃——以及\Place{索多瑪}的蘋果——外表成熟，但裡面充滿灰塵——我未能找到可信的證人來確定它們是否屬實。事實上，我沒有遇到任何人說他們見過\Place{以彼羅}的泉水，但有些人知道在\Place{高盧}離\Place{格勒諾布爾}不遠有一處類似的泉水。然而，\Place{索多瑪}樹上的果子不僅在可靠的書籍中談到，而且許多人說他們見過，因此我無法懷疑這個事實。至於其他奇蹟，我既不肯定也不否定地接受；我引用它們是因為我在我們對手的著作中讀到它們，並且是為了證明：他們自己中有許多人相信這麼多事物，因為它們寫在他們自己文學家的作品中，儘管沒有給出合理的解釋；然而他們卻輕視相信我們斷言全能的天主將要做超越他們經驗和觀察的事情；即使我們為祂的工作給出理由，他們也是如此。因為對於這些事，有什麼比說全能者能夠實現它們，並且將會實現它們，正如祂在那些書中所預言的那樣——這些書中預言了許多其他已經實現的奇蹟——更好、更強的理由呢？那些被認為不可能的事，將根據那位預言並實現了使不信的民族相信難以置信的奇蹟的天主的話語和權能而成就。\par

% structure: p0028 source=h2 target=subsection reason=relative-heading-rank; base=h2

\subsection{在已知本性的物體中發現其已知的自然屬性的改變，並不違反本性。}

然而，如果他们反驳说，他们不相信我们所说的人的身体将永远焚烧却不死亡的理由，是因为已知的人的身体的性质完全不同于此；如果他们说，对于这一奇迹，我们无法给出在自然奇迹案例中有效的理由，即这是事物的自然属性——因为我们知道这不是人肉身的自然属性——那么我们在圣经中找到答案：这同一人的肉身在犯罪之前有一种构造，实际上被构造得不能死亡，而在犯罪之后有另一种构造，成为我们在今世悲惨的必朽状态中所见的，不能保持持久的生命。同样，在死者的复活中，它将不同于现今众所周知的状态而被构造。但由于他们不相信我们的这些著作，我们在其中读到人在乐园中的性质，以及他如何远离死亡的必然——实际上，如果他们相信这些，我们当然不必费力与他们辩论被罚者将来的惩罚——我们必须从他们自己最博学的权威的著作中举出一些例子，以表明一件事物可能变得不同于它以前已知的特征性质。\par

我从\Person{马库斯·瓦罗}所著\WorkTitle{论罗马人民的种族}一书中逐字引用以下实例：\Quote{发生了一个显著的天象异兆；因为\Person{卡斯托尔}记载，在明亮的金星——\Person{普劳图斯}称之为\OriginalTerm{晚星}{Vesperugo}，\Person{荷马}称之为美丽的\OriginalTerm{昏星}{Hesperus}——上，发生了如此奇异的怪事，竟改变了它的颜色、大小、形状和运行，这前所未有，此后也未再发生。\Place{库齐库斯}的\Person{阿德拉斯图斯}和\Place{那不勒斯}的\Person{狄翁}，著名的数学家，说此事发生在\Person{俄古革斯}王统治时期。}\par

\Quote{使江河倒流回源头，\newline{}令星辰忘却其轨道。}\par

因为在我们的圣经中，我们也读到这事曾发生：一条河\Quote{倒流回}，在上游停住而下游继续流淌，当百姓在上述领袖、\Person{农的儿子若苏厄}带领下渡过时；也当先知\Person{厄里亚}渡过时；后来，当他的门徒\Person{厄里叟}渡过时也是如此；我们刚才提及，在\Person{希则克雅}王的事例中，最大的星辰\Quote{忘却了它的轨道}。但根据瓦罗，发生在金星上的事他并未说是因任何人的祈祷而发生的。\par

但愿怀疑者不要因这种对事物本性的知识而自我蒙蔽，仿佛天主的大能不能在一件事物上成就任何超出他们经验所显示其本性之事。即使是最寻常的自然事物，若人们不习惯只惊叹稀有之事，也不会不那么奇妙，也不那么令所有看见的人感到惊奇。因为谁若细心思索那无数的人群，以及他们本性的相似，怎能不惊讶而赞叹地注意到每个人容貌的独特性：这暗示我们，除非人们彼此相似，否则他们不会与其他动物区别开来；而另一方面，除非他们不相似，否则他们不能彼此区分，以至我们声称相似的人，也发现他们是不相似的？而不相似是两者中更令人惊奇的考虑；因为共同的本性似乎更要求相似。然而，正因为事物的稀有性才使之令人惊奇，当我们遇到两个人如此相似，以至于我们总是或经常在试图区分他们时弄错，我们便感到更大的惊奇。\par

但可能，尽管\Person{瓦罗}是外邦历史学家，极博学之士，他们仍不相信我从他那里引用的例子确曾发生；或他们说这例子无效，因为那颗星并未长时间保持新的轨道，而是回到了通常的轨道。于是，现在另有一个现象呈现在他们眼前，我认为足以说服他们：尽管他们观察并确定了某种自然规律，但他们仍不应以此为理由给天主规定限制，好像祂不能改变它并将其转化为与他们所观察到的截然不同的东西。\Place{索多玛}之地并非总是现在这个样子；它一度有与其他地方相似的外观，享有同样甚至更丰饶的肥力；因为在天主的叙述中，它被比作天主的乐园。但自从它被天上的火触碰之后，正如外邦历史所见证，也为访问该地者今日所目睹，它变得异常可怖的焦黑；它的苹果，在欺骗性的成熟外表下，内藏灰烬。这是一件曾是一种样子、现在是另一种样子的事。你看，它的本性是如何被万有的创造者奇妙的转化改变成如此令人厌恶的变异的——这一改变是在如此漫长的岁月后才发生，并持续了如此长的时日。因此，正如天主不可能不能创造祂所乐意的本性，祂也不可能不能将祂自己所创造的本性变成祂所乐意的任何形态，从而散布大量被称作\Quote{怪物}、\Quote{异兆}、\Quote{奇事}、\Quote{现象}的奇迹，如果我有意列举并记录它们，这部著作何时才有个结束？他们说，它们被称为\Quote{怪物}，因为显示或象征某事；被称为\Quote{异兆}，因为预示某事；等等。但让他们的占卜者看看，他们要么是被欺骗，要么即使他们预言了真实之事，那也是因为他们被那些一心要使人的灵魂（诚然配得上如此命运）陷入有害好奇心的网罗的邪灵所感动，要么是因为他们做出众多预测而偶尔碰巧蒙对。然而，对于我们而言，这些相反本性发生、并被称为相反本性之事（正如宗徒依照人的说法所说，将野橄榄枝接在好橄榄树上、共享其肥汁，是相反本性的），且被称为怪物、现象、异兆、奇事的，应当显示、预示、预言：天主将实现祂关于人身体所预言的一切，没有任何困难能阻止祂，没有任何自然规律能为祂设定界限。祂如何预言了祂将要做的事，我认为在前一卷书中已经充分表明，我从新旧约圣经中搜集了——当然不是所有与此相关的段落，而是我认为足以完成本工作的那些。\par

% structure: p0035 source=h2 target=subsection reason=relative-heading-rank; base=h2

\subsection{第九章——论地狱及永罚的本质。}

因此，天主藉其先知论及受罚者永罚的预言必将实现——必然实现——\BibleQuote{他们的虫子总不死，他们的火总不灭。}\BibleRef{依 66:24}为了令我们最深切地铭记这一点，主耶稣基督亲自在命令我们砍掉我们的肢体时——以此意指那些被某人如身体最有用肢体般爱着的人——说：\BibleQuote{你残废进入生命，比有两只手而往地狱，进入永不熄灭的火里去更好，那里他们的虫子总不死，火也不灭。}同样，关于脚：\BibleQuote{你瘸腿进入生命，比有两只脚而被投入地狱，进入永不熄灭的火里去更好，那里他们的虫子总不死，火也不灭。}同样，关于眼：\BibleQuote{你一只眼进入天主的国，比有两只眼而被投入地狱火里去更好，那里他们的虫子总不死，火也不灭。}\BibleRef{谷 9:43}祂在一段经文中毫不迟疑地三次使用相同的话。谁不被这重复以及主亲自以如此强烈语气宣布的惩罚威胁所恐惧呢？\par

现今，那些主张火与虫皆归于灵魂而非身体的人坚称：恶人既被摒弃于天主的国之外，便要因灵魂徒然太晚悔改的苦痛而仿佛被焚烧；他们争辩说，用火来表达这种焚烧之苦并非不当，正如宗徒疾呼：\BibleQuote{谁软弱，我不软弱呢？谁跌倒，我不焦急呢？}\BibleRef{格后 11:29}他们也认为虫也应作同样理解。经上说：\BibleQuote{如蛀虫蛀蚀衣服，如蠹虫蛀蚀木头，同样，忧愁蛀蚀人的心。}\BibleRef{依 51:8}然而，那些毫不怀疑在将来的惩罚中身体与灵魂都将受苦的人断言：身体将被火焚烧，而灵魂则将被痛苦的虫啮咬。这种见解更为合理——因为假定身体或灵魂在将来的惩罚中能免于痛苦是荒谬的。但就我个人而言，更容易将两者都理解为指身体，而非认为两者都不指身体；我认为圣经对受罚者的灵性痛苦保持缄默，是因为虽然没有明言，但可以理解的是，在一个如此受折磨的身体中，灵魂也必因徒劳的悔改而痛苦。因为在古代经卷中我们读到：\BibleQuote{不虔者的肉，火与虫为其报应。}\BibleRef{德 7:17}本可说\Quote{不虔者的报应}，为何却说\Quote{不虔者的肉}？岂不是因为火与虫皆是肉身的惩罚？或者，作者说\Quote{肉身的报应}，是为了指明那些按肉性而活的人所受的惩罚（因为这导致第二次死亡，如宗徒所暗示：\BibleQuote{如果你们随从肉性生活，必要死亡。}\BibleRef{罗 8:13}）。但愿各人自择：或指火归于身体、虫归于灵魂——一为象征，一为真实；或指两者都真实地归于身体。因为我已经充分说明，动物可以活在火中，被焚烧却不被烧毁，受痛苦却不死亡，这是全能的造物主的奇迹；若有人不是无知于谁创造了自然界一切神奇之事，就不能否认这是可能的。正是天主自己，在这个世界上行了所有那些我提过的大大小小的奇迹，以及更多我未提及的奇迹，并将这些奇事封闭在这世界之中，而这世界本身便是最大的奇迹。因此，各人可自行选择：认为虫是真实的并属于身体，或认为属灵之事由身体表征并属于灵魂。但究竟孰真，待圣人们拥有那样的知识时——无需凭借亲身经验来认识这些刑罚的性质，而是凭其圆满与完善便足以知晓此事——便更容易从事实本身发现。因为\BibleQuote{我们现在所知道的，只是局部的，及至那圆满的一来到，局部的就过去了。}\BibleRef{格前 13:9}只是，关于那些将来的身体，我们相信：它们必定会被火所痛苦。\par

% structure: p0038 source=h2 target=subsection reason=relative-heading-rank; base=h2

\subsection{第十章——假如地狱之火是物质的火，它能否焚烧恶灵，即无形的魔鬼？}

这里产生一个问题：如果那火不是非物质的、类似于灵魂的痛苦，而是物质的、以接触来焚烧，以便身体在其中受折磨，那么恶灵又怎能在这火中受罚呢？因为无疑，用以惩罚人和魔鬼的是同一火，按基督的话：\BibleQuote{可咒骂的，离开我，到那给魔鬼和他的使者预备的永火里去！}\BibleRef{玛 25:41}——除非，如一些有学问的人所想的，魔鬼有一种由稠密潮湿的空气构成的身体，正如我们感觉起风时所体会到的那种物质。如果这类物质不受火的影响，那么它在浴池中就不能被烧热。因为要焚烧，它首先被烧热，然后以自身被影响的方式影响其他事物。但若有人主张魔鬼没有身体，这既无需费力探究，亦无需激烈辩论。因为，若人的灵魂（无疑也是非物质的）如今被包含在身体的有形肢体中，且在未来将不可分解地与自己的身体结合，那么即使魔鬼没有身体，它们的灵魂（即魔鬼本身）也将彻底接触物质的火而受其折磨；并非因为火本身因与这些灵魂结合而获得生命，并成为由身体与灵魂组成的生物，而是如我所说，这种结合将以奇妙而不可言喻的方式实现，使它们从火中感受痛苦，却不给予火生命。事实上，另一种结合方式——身体与灵魂结合成为生物——本身已极其奇妙，超乎人的理解，尽管人正是如此形成的。\par

我确实想说，这些魔鬼将没有任何自己的躯体而燃烧，正如那个财主在地狱中焚烧时喊道：\BibleQuote{我在这火焰里极其痛苦，} \BibleRef{路 16:24} 如果我不知道可以恰当地回答说，那火焰与他举目望\Person{拉匝禄}的眼睛性质相同，与他请求滴一点凉水到他的舌头上的舌头性质相同，也与\Person{拉匝禄}的手指性质相同——这一切都发生在没有躯体的灵魂所在之处。因此，那焚烧他的火焰和他所求的那滴水都是非物质的，犹如睡梦中或出神者的异象，在其中非物质的对象以有形的样式出现。因为处于这种状态的人，虽然只是在神魂中而非在身体中，却看到自己如此像自己的身体，以致无法辨别任何差异。但那个地狱，也称为火和硫磺的湖，\BibleRef{默 20:10} 将是物质之火，将折磨被定罪者的身体，无论是人还是魔鬼——人的实体身体，魔鬼的气态身体；或者，如果只有人既有身体又有灵魂，那么恶灵虽无身体，却将与物质之火如此联结，以致感受到痛苦却不赋予生命。同样的火必然是他们两者的命运，因为真理已经如此宣告。\par

% structure: p0041 source=h2 target=subsection reason=relative-heading-rank; base=h2

\subsection{第11章——罪罚持续长于罪恶本身是否公义？}

然而，在我们为之辩护的\WorkTitle{天主之城}的反对者中，有些人认为，任何人为罪恶遭受永恒的惩罚是不公义的，因为无论这些罪恶多么严重，都是在短时间内犯下的；仿佛任何法律都曾根据犯罪的时间长短来规定惩罚的持续时间！\Person{西塞罗}告诉我们，法律承认八种刑罚：损害赔偿、监禁、鞭笞、赔偿、耻辱、流放、死刑、奴役。这其中的任何一种能否压缩到与犯罪的迅速完成成比例的短暂时间，以便惩罚所花的时间不超过犯罪所花的时间？除非是赔偿？因为赔偿要求犯罪者承受他所做的，正如法律条款所说：\BibleQuote{以眼还眼，以牙还牙。} \BibleRef{出 21:24} 因为，确实，犯罪者因法律的严厉报复而失去一只眼睛，其所花的时间可能与他因自己的不法残酷而剥夺他人眼睛的时间一样短暂。但是，如果鞭笞是对亲吻他人之妻的合理惩罚，那么瞬间的过失是否要用长时间的补赎来偿还，瞬间的愉悦是否要以持久的痛苦来惩罚？我们该如何说监禁呢？罪犯是否只应被监禁与他所犯的罪行所花的时间一样长？或者，对那个用一句话激怒主人或迅速打了一拳的奴隶，难道不是处以多年的监禁吗？至于损害赔偿、耻辱、流放、奴役，这些刑罚通常施加时不允许任何减轻或赦免，它们不是类似于永恒的惩罚吗？只要这短暂的生命允许类似。它们不是永恒的，只是因为承受它们的生命不是永恒的；然而，以这些最持久的痛苦来惩罚的罪行，却是在非常短暂的时间内犯下的。没有人会认为惩罚的痛苦应与罪行的时间一样短；或者谋杀、奸淫、亵渎或任何其他罪行，不应根据伤害或邪恶的严重性来衡量，而应根据犯罪所花的时间长短来衡量。那么，对于任何重大罪行判处死刑，法律认为惩罚在于执行死刑的短暂瞬间呢，还是在于罪犯永远被逐出活人的社会？正如第一次死亡的惩罚将人从这现世的易朽之城剪除，同样，第二次死亡的惩罚将人从那未来的不朽之城剪除。因为正如这现世之城的法律没有规定被处决的罪犯可以返回，同样，被判处第二次死亡的人也不会被召回永生。但如果暂时的罪行要受永恒的惩罚，那么，他们如何说，你们的基督所说的这句话是真实的：\BibleQuote{你们用什么尺度量给人，也要用什么尺度量给你们？} \BibleRef{路 6:38} 他们没有注意到，\Quote{同样的尺度}指的并不是相等的时间，而是恶有恶报，换句话说，就是根据那使作恶者承受恶的律法。此外，这些话可以恰当地理解为涉及我们的主在使用它们时所说的事情，即审判和定罪。因此，如果一个人不公正地审判和定罪，他自己也被公正地审判和定罪，他领受的是\Quote{同样的尺度}，尽管他给的和他领受的不是同一件事。因为他给了审判，他领受的也是审判，尽管他给的审判是不义的，但他领受的审判是公义的。\par

% structure: p0043 source=h2 target=subsection reason=relative-heading-rank; base=h2

\subsection{第12章——论原初过犯的伟大，因这过犯，所有不在救主恩宠范围之内的人都应受永恒的惩罚。}

然而，永恒的惩罚在人的感知中显得严酷和不公，因为我们有死之境的微弱理智缺乏那种最高而最洁的智慧，使人能看清那第一次过犯犯下了何等巨大的邪恶。人越是在天主内找到喜乐，他离弃天主的邪恶就越大；那在自己内毁灭了本可永恒之善的人，成了永恒之恶的当受者。因此，全人类都被定罪；因为那首先引入罪恶的人，连同所有在他内如同在根内的后代，都受了惩罚，以致无人能免除这公正而应得的惩罚，除非藉着仁慈和白白的恩宠得救；人类如此分配，以致在一些人身上显示仁慈恩宠的功效，在其余人身上显示公正报应的功效。因为两者不能同时在所有人身上显示；如果所有人都留在公正定罪的惩罚下，就没有人身上显出救赎恩宠的仁慈。反之，如果所有人都从黑暗被迁移到光明，那么报应的严厉就没有在任何一人身上显明。但留在惩罚下的人远比得救的人多，为要显示众人都应得的惩罚。如果这惩罚加于所有人，没有人能公正地指责施行报应者的正义；然而，从这公正判决中得救的如此多人，却成了向那施救者的白白恩赐致以最真挚感谢的理由。\par

% structure: p0045 source=h2 target=subsection reason=relative-heading-rank; base=h2

\subsection{第十三章　驳斥那些认为恶人死后所受惩罚是炼狱性的意见}

确实，柏拉图主义者虽然主张没有不受惩罚的罪，但他们认为所有惩罚——不论是依人法或神法，今生或死后——都是为了医治；因为一个人在此世可能不受伤害，或虽受惩罚却仍不改过。故维吉尔有诗云，当他说到我们地上的身体和有死的肢体时，我们的灵魂源自——\par

\Quote{Hence wild desires and grovelling fears,\newline{}And human laughter, human tears;\newline{}Immured in dungeon-seeming night,\newline{}They look abroad, yet see no light,}\newline{}接着又说：\par

\Quote{Nay, when at last the life has fled,\newline{}And left the body cold and dead,\newline{}Ee'n then there passes not away\newline{}The painful heritage of clay;\newline{}Full many a long-contracted stain\newline{}Perforce must linger deep in grain.\newline{}So penal sufferings they endure\newline{}For ancient crime, to make them pure;\newline{}Some hang aloft in open view,\newline{}For winds to pierce them through and through,\newline{}While others purge their guilt deep-dyed\newline{}In burning fire or whelming tide.}\par

持此意见的人认为死后的所有惩罚都是炼狱性的；既然气、火、水元素高于土，它们中的某一个便可用作洗净和清除因土的染污而沾染的污点。故维吉尔以\Quote{Some hang aloft for winds to pierce;}暗示气，以\Quote{whelming tide}暗示水，以\Quote{in burning fire}暗示火。就我们而言，我们承认即使在今生，某些惩罚也是炼狱性的——不过不是针对那些生命并未因此变好，反而更糟的人，而是针对那些因这些惩罚而被约束改过自新的人。其余所有惩罚，无论暂时的或永恒的，由天主眷顾加于每个人，或是由于过去的罪，或是由于今生仍在犯的罪，或是为操练并显明人的恩宠。它们可藉恶人和善人、恶天使和善天使作为工具施行。即使有人因他人的恶意或错误而受伤害，确实那因无知或不公而作恶的人犯了罪；但天主，因其公正而隐藏的审判允许此事发生，并不犯罪。但暂时的惩罚，有的只在今生受，有的在死后受，有的今生和死后都受；但所有这一切都在那最后最严厉的审判之前。至于那些在死后受暂时惩罚的人，并非所有人都注定要受那审判之后的永恒痛苦；因为对有些人，如我们已说过的，今世未获赦免的会在来世获得赦免，即他们不会受来世永恒的惩罚。\par

% structure: p0050 source=h2 target=subsection reason=relative-heading-rank; base=h2

\subsection{第十四章　论此世人类处境所受的暂时惩罚}

那些今生不受惩罚、仅在日后受罚的人是相当例外的。然而，有些人活到风烛残年，甚至从未经历过丝毫疾病，且一直享受生命，我既闻于传闻，也亲眼目睹。然而，我们凡人所过的生命本身就是全部的惩罚，因为它完全是诱惑，正如圣经所宣告的，其中记载：\Quote{人的生命在地上岂不是一个诱惑？}\BibleRef{约 7:1} 因为无知本身就是一个不小的惩罚，或缺乏教化，人们正当地认为必须逃避它，以至于男孩们被迫在严厉惩罚的痛苦之下学习手艺或文字；而他们被惩罚驱赶去学习本身对他们来说就是如此大的惩罚，以至于他们有时宁愿选择驱赶他们的痛苦，也不愿忍受被驱赶去学习的痛苦。如果有人在两者之间——要么受死，要么重新成为婴儿——作选择，谁会不退缩，而选择死亡呢？实际上，我们的婴儿时期以眼泪而非欢笑引我们进入此生，似乎无意识地预兆我们将遭遇的灾祸。据说只有\Person{琐罗亚斯德}出生时笑了，而那反常的征兆对他并非吉兆。据称他是魔术的发明者，然而这些技艺甚至无法确保他享有今世那可怜的幸福，以抵御敌人的攻击。因为，他本人是\Place{巴克特里亚}人的国王，却被\Place{亚述}人的国王\Person{尼努斯}征服。总之，圣经的话：\Quote{沉重的轭加在亚当子孙身上，从他们出母胎之日直到他们回归万物之母之日}\BibleRef{德 40:1} ——这些话如此无误地应验，以至于那些幼童，虽藉着圣洗圣事从原罪的束缚中获得释放，却仍遭受许多灾祸，甚至在某些情况下还暴露于恶灵的攻击之下。但我们片刻不可认为这痛苦会损害他们未来的幸福，即使这痛苦已大到使灵魂与身体分离，并在那幼年就结束他们的生命。\par

% structure: p0052 source=h2 target=subsection reason=relative-heading-rank; base=h2

\subsection{第十五章——天主恩宠在拯救我们脱离我们所深陷的根深蒂固的邪恶中所做的一切，都属于未来的世界，在其中万物都成为新的。}

然而，在这\Quote{沉重的轭加在亚当子孙身上，从他们出母胎之日直到他们回归万物之母之日}中，我们发现了一个虽令人痛苦却令人钦佩的训诫，教导我们保持清醒，并说服我们：此生已成为惩罚性的，因为那在乐园中犯下的滔天罪恶；而且新约所邀请的一切都属于那在来世等待我们的未来产业，并作为定金提供给我们，使我们在适当的时候获得它所担保的实体。所以，现在让我们在望德中行走，并藉着圣神致死肉体的行为，从而日日进步。因为\Quote{主认识属于祂的人；}\BibleRef{弟后 2:19} 并且\Quote{凡被天主的圣神所引导的，都是天主的子女，}\BibleRef{罗 8:14} 但这是因恩宠，而非因本性。因为只有一位天主子因本性是子，祂出于怜悯为我们成为人子，好使我们这些因本性是人的子女的，能藉着祂因恩宠成为天主的子女。因为祂，不变地存在，取了我们的本性，为要借此将我们带到祂自己那里；并且，坚守祂自己的神性，祂成为我们软弱的分享者，好使我们，被改变成更好的事物，能因参与祂的义德与不死而失去我们自己的罪与死亡的属性，并保守祂曾植入我们本性中的一切善质，这本性如今因分享祂本性的善而得以成全。因为正如因一人的罪，我们跌入如此可悲的悲惨中，同样因一人——祂也是天主——的义德，我们将获得不可言喻的高超幸福。任何人也不应相信自己已经从那人转到这人，除非他到达了没有诱惑的地方，并进入了他在此战争的各种冲突中所寻求的平安，在这战争中\BibleQuote{肉情相反圣神，圣神相反肉情。}\BibleRef{迦 5:17} 如今，这样的战争本不会存在，如果人性在运用自由意志时，坚定持守于它被造时所具有的义德中。但现在，在它的悲惨中，它向自己开战，因为它在幸福时不肯与天主保持和平；而这一点，尽管是悲惨的灾祸，仍好于此生的早期阶段，那些阶段并不认识一场战争必须持续。因为与恶习斗争好过不加抵抗而被它们征服。我再说，带着永恆平安的战争，好过毫无得救之念的囚禁。我们确实渴望这战争的止息，并受天主之爱的火焰点燃，渴望进入那井然有序的平安，在其中凡较低者永远从属于较高者。但若（愿天主不许）没有如此幸福结局的希望，我们仍宁愿忍受这冲突的艰难，而不愿因不抵抗而向恶习的统治屈服。\par

% structure: p0054 source=h2 target=subsection reason=relative-heading-rank; base=h2

\subsection{第十六章——恩宠的规律，它扩展到重生者生命的所有时期。}

然而，天主对祂所预备要得荣耀的怜悯之器，其仁慈是如此伟大，以致人的最初年龄，即婴幼儿时期，它毫不抵抗地顺从肉体；以及第二个年龄，称为童年，它尚无足够理解力来参与这场战争，因此几乎屈服于每一种邪恶的享乐（因为尽管这个年龄已有言语能力，因而看似已过了婴幼儿期，但心智仍太软弱，无法理解诫命）；然而，若这两个年龄中的任何一个领受了中保的圣事，那么，即使现世生命立即终结，这个孩子既已从黑暗的权势迁到基督的国度，就不仅会得救免于永罚，而且死后也不会遭受炼狱之苦。因为属灵的再生本身足以阻止因肉身的生育所形成的与死亡的联结在死后产生任何恶果。但是，当我们达到那个能理解诫命并服从法律统治的年龄时，我们必须向罪恶宣战，并热切地进行这场战争，以免陷入该死的罪中。如果罪恶没有因习惯性的胜利而积聚力量，它们就更容易被克服和制伏；但如果它们已经习惯于征服和统治，那么要制服它们就只有通过艰难和劳苦了。事实上，这种胜利只有通过喜爱真正的正义才能真诚而真实地获得，而这正义是由对基督的信德所赐予的。因为如果法律带着它的命令存在，而圣神没有带着祂的帮助同在，那么禁令的存在只会增加犯罪的欲望，并增加违法的罪过。有时，明显的罪恶被其他隐藏的罪恶所克服，这些隐藏的罪恶被当作德行，尽管骄傲和一种毁灭性的自满自足是它们的形成原则。因此，只有当罪恶被对天主的爱所征服时，才能认为它们被克服了；而这爱唯独是天主自己赐予的，并且祂只通过天主与人之间的中保、降生成人的基督耶稣赐予，祂成为我们可朽坏性的分享者，为使我们成为祂神性的分享者。然而，真正幸运的人是很少的，他们能够在青年时期没有犯下任何该死的罪，无论是通过放荡或暴力的行为，还是追随某些不敬神且不合法的意见，而是凭借他们灵魂的伟大，制伏了所有能使他们成为肉体享乐的奴隶的东西。大多数人首先成为他们所领受的法律的违犯者，并让罪恶在他们身上占据上风，然后逃向恩宠寻求帮助，这样，通过比原本更痛苦的忏悔和更激烈的斗争，他们使灵魂顺服于天主，从而给予灵魂对肉体的合法权柄，成为胜利者。因此，凡渴望逃避永罚的人，不仅要受洗，还要在基督内成义，如此才能真正从魔鬼过渡到基督。并且，他不应幻想有任何炼狱的痛苦存在于那最终而可怖的审判之前。然而，我们不可否认，永恒的火也将按恶人的功过成比例，使一些人更痛苦，另一些人较不痛苦；这结果要么通过火本身的温度变化实现，按每人的功过分级，要么热度保持不变，但并非所有人都以同等的痛苦强度感受它。\par

% structure: p0056 source=h2 target=subsection reason=relative-heading-rank; base=h2

\subsection{第十七章 论那些幻想无人会受永罚的人。}

我看我现在必须进入与那些心软的基督徒的友好争论，他们不愿相信任何被绝对公义的审判者判定应受地狱之罚的人，或他们中所有的人，会永远受苦，而猜想他们会在一个固定的惩罚期限后获得释放，期限长短按各人罪的大小而定。关于此事，\Person{奥利振}甚至更为宽纵；因为他相信就连魔鬼本身及其天使，在遭受了他们罪过所应得的更严厉和更长久的痛苦之后，也会从折磨中释放出来，并与圣天使结合。但教会并非无理由地谴责了他这一错误以及其他错误，尤其是他关于幸福与苦难不断交替、以及在一定时期内有规律地从一种状态过渡到另一种状态的理论；因为在这一理论中，他甚至连仁慈的名声也失去了，因为他为圣人分派了真实的苦难来赎罪，以及虚假的幸福，这幸福不能带给他们真正而可靠的喜乐，即对永恒福乐的无畏确据。然而，我们现在所说的错误却大不相同，它是由那些心软的基督徒的温情所驱使，他们认为审判中被定罪者的苦难将是暂时的，而所有迟早获得释放者的福乐将是永恒的。这种观点，如果因为它仁慈而算为良善真实的，那么它越仁慈，就越良善真实。那么，让这仁慈的泉源扩展到失落的众天使，也让他们在经过看似合适的许多漫长年代后获得释放吧！为何这仁慈之流流向整个人类，一触及天使就干涸了呢？然而，他们不敢将怜悯进一步扩展，提议释放魔鬼本身。如果有人胆敢这样做，他确实使他们的爱德蒙羞，但自己却被更丑恶的错误和更悖谬的天主真理的曲解所定罪，其程度与他表面上的仁慈慷慨成正比。\par

% structure: p0058 source=h2 target=subsection reason=relative-heading-rank; base=h2

\subsection{第十八章 论那些幻想因圣人的转求，人在最后审判中不会被定罪的人。}

还有一些人，我与他们交谈中了解到他们的观点，他们虽然似乎敬重圣经，但生活却应受谴责，因此，为了自己的利益，他们将一种更大的慈悲归于天主，施予人类。因为他们承认，圣言中确实预言了恶人和不信者应受惩罚，但他们断言，当审判来临时，慈悲将得胜。他们说，因为天主怜悯他们，必将他们交托给圣人们的祈祷和转求。因为如果圣人们在遭受他们残酷的仇恨时曾为他们祈祷，那么当他们看到他们俯伏谦卑地哀求时，岂不是更会这样做吗？他们说，我们不能相信圣人们一旦达到最完美完全的圣德，就会失去怜悯的心肠；既然他们还是罪人时就曾为仇敌祈祷，如今他们已脱免罪恶，怎会不为哀求者转求呢？难道天主会拒绝聆听祂如此众多爱子的祈祷吗？尤其是他们的圣德已净化了他们的祈祷，使之毫无阻碍，能蒙祂垂允。而那些承认恶人和不信者将受长时间惩罚，但最终会从一切痛苦中获释的人所引用的圣咏经文，也被我们现在所谈论的人认为更有利于他们。那经文是：\BibleQuote{难道天主忘记施恩了吗？难道祂在愤怒中关闭了祂的怜悯吗？}他们说，祂的愤怒会判定一切不配得永生的人受无尽的惩罚。但如果祂允许他们受长时间惩罚，甚至有任何惩罚，那么祂岂不关闭了祂的怜悯？而圣咏作者暗示祂不会这样做。因为祂并没有说：\BibleQuote{难道祂在愤怒中长时间关闭了祂的怜悯吗？}而是暗示祂根本不会关闭。\par

他们否认，即使天主不定任何人的罪，天主审判的威胁就会被证明是虚假的，正如我们不能说祂摧毁尼尼微的威胁是虚假的，尽管绝对预言的毁灭并未实现。因为祂并没有说：\BibleQuote{若尼尼微人不悔改、改善他们的行为，尼尼微就要被倾覆，}而是没有任何条件地预言该城将被倾覆。他们坚持认为，这个预言是真实的，因为天主预言了他们应受的惩罚，尽管祂并未施行。因为虽然祂因他们悔改而赦免了他们，但祂当然知道他们会悔改，然而祂仍然绝对而明确地预言该城将被倾覆。他们说，这在严厉的真实性上是真实的，因为他们理应受罚；但从抑制祂愤怒的慈悲来看，祂饶恕了哀求者，免除了祂先前威胁悖逆者的惩罚，这就不真实了。那么，如果祂饶恕了那些连祂自己的圣先知都因祂的饶恕而恼怒的人，祂岂不更会饶恕那些更可怜的哀求者，祂所有的圣人都会为他们转求吗？而且他们认为，圣经并未暗示他们的这一猜想，是为了激励许多人因恐惧极其漫长或永恒的惩罚而改过迁善，并激励其他人为那些尚未改过的人祈祷。然而，他们认为天主的圣言并非对此完全沉默；因为他们问，\BibleQuote{祢为敬畏祢的人所积存的恩惠何其广大！}这句话若非为了教导我们，天主的慈悲那伟大而隐藏的甘饴被隐藏起来，好使人畏惧，那又是为何呢？他们认为，使徒也说了同样的话：\BibleQuote{因为天主把众人都禁锢在不信服之中，是要怜悯众人。}\BibleRef{罗 11:32}表明天主不应定任何人的罪。然而，持此观点的人并未将此延伸至魔鬼及其使者的开脱或释放。他们的人性同情只针对人，而且主要为自己辩护，通过这种天主对全人类的准慈悲，为他们堕落的生命提供有罪不罚的虚假希望。因此，那些甚至应许魔鬼及其爪牙也有罪不罚的人，更是充分展示了天主的仁慈。\par

% structure: p0061 source=h2 target=subsection reason=relative-heading-rank; base=h2

\subsection{第十九章——论那些藉参与基督的身体，而应许赦免一切罪过乃至异端者的人。}

同样，另有一些人应许从永罚中获得释放，虽非给所有人，而只给那些曾在基督的圣洗圣事中被洗净，并成为基督身体分享者的人，无论他们如何生活，或堕入何种异端或不敬。他们基于耶稣的话：\BibleQuote{这是从天上降下来的食粮，谁吃了，就不死。我是从天上降下来的活粮，谁若吃了这食粮，将永远活着。}\BibleRef{若 6:50}因此，他们说，这些人必从永死中获得解救，并在某一时刻被引入永生。\par

% structure: p0063 source=h2 target=subsection reason=relative-heading-rank; base=h2

\subsection{第二十章——论那些不将此恩惠应许给所有人，而只给那些曾受洗为天主教徒，但后来犯下诸多罪行和异端的人。}

还有另一些人，他们甚至不将此应许给予所有领受过基督的圣洗圣事和圣体圣事的人，而只给予天主教徒，无论他们生活多么败坏。因为他们不仅是在圣事上，而且是实际地吃了基督的身体，与祂的身体结合，正如宗徒所说：\BibleQuote{我们虽多，却是一个饼，一个身体。}\BibleRef{格前 10:17}因此，即使他们后来堕入某种异端，甚至外教和偶像崇拜，但因这一件事，即他们已在基督的身体内，也就是在天主教会内，领受了基督的圣洗圣事并吃了基督的身体，他们不会永远死亡，而终有一时会获得永生；他们的一切邪恶并不能使他们的惩罚成为永久的，而只是与时间长短和严厉程度相称。\par

% structure: p0065 source=h2 target=subsection reason=relative-heading-rank; base=h2

\subsection{第二十一章　论那些主张凡持守信仰的公教徒，即使因生活败坏应受地狱之火，仍将因他们信德的\Quote{根基}而得救的人。}

有些人也依据圣经的这句话：\Quote{凡坚忍到底的，必会得救。}\BibleRef{玛 24:13}他们只将救恩应许给那留在天主教会中的人；并且，他们说，即使这样的人生活败坏，仍将因那根基而如经火一样得救——宗徒所说的根基是：\Quote{除已奠定的根基外，任何人不能再奠立别的，而这根基就是耶稣基督。人若在这根基上建造，或用金、银、宝石，或木、草、禾秸；各人的工程必将显露出来，因为主的日子要把它揭露，原来它要被火揭示；各人的工程是何种性质，火也要试验出来。谁在这根基上建造的工程若存留得住，他就要获得赏报；但谁的工程若被焚毁，他就受了损失，他自己却要得救，可是这样得救，像从火中经过的一样。}\BibleRef{格前 3:11}因此，他们声称，无论公教徒的生活如何，都有基督作为他的根基，而任何与祂身体的合一分离的异端都不拥有这根基。所以，基于这根基，即使公教徒因生活的不恒定，如在根基上堆积木、草、禾秸，他们相信他仍要得救，像从火中经过一样——换言之，他在尝过那在最后审判中要惩罚恶人的火之苦后，将得释放。\par

% structure: p0067 source=h2 target=subsection reason=relative-heading-rank; base=h2

\subsection{第二十二章　论那些幻想与施舍相混的罪在审判之日不被记念的人。}

我也曾遇到一些人，他们以为只有那些忽略用施舍遮盖自己罪过的人，才受永火的惩罚；他们引用了雅各伯宗徒的话：\Quote{未施怜悯者，必受无怜悯之审判。}\BibleRef{雅 2:13}因此，他们说，那没有改善自己行为，却将放荡邪恶的行为与慈悲善工相混的人，将在审判中获得怜悯，以至于要么完全逃脱定罪，要么在或长或短的时间之后从刑罚中得到释放。他们以为，这就是审判生者和死者的主自己在赋予右边的人永生、左边的人永罚时，只提及行了或未行慈悲善工的原因。\BibleRef{玛 25:33}他们说，同样的道理也见于我们在主祷文中的每日祈求：\Quote{宽免我们的罪债，如同我们也宽免得罪我们的人。}\BibleRef{玛 6:12}因为，无疑，那宽恕得罪自己的人，行了慈善之事。而这点受到主亲自的高度赞扬，祂说：\Quote{因为你们若宽免别人的过犯，你们的天父也必宽免你们的；但你们若不宽免别人的，你们的父也必不宽免你们的过犯。}\BibleRef{玛 6:14}因此，雅各伯宗徒的话就是针对这类施舍说的：\Quote{未施怜悯者，必受无怜悯之审判。}他们又说，我们的主并未区分大罪小罪，而是说：\Quote{你们的天父要宽免你们的罪，如果你们宽免别人的罪。}于是他们得出结论：即使一个人直到生命的最后一天都过着放荡的生活，无论他的罪如何，只要他留意这一件事——即当那些伤害他的人请求宽恕时，他从心里宽恕他们——那么，藉着这每日的祈祷，所有的罪就都被赦免了。\par

在天主的助佑下，当我答复了所有这些错误后，我将结束本书（第二十一卷）。\par

% structure: p0070 source=h2 target=subsection reason=relative-heading-rank; base=h2

\subsection{第二十三章　驳斥那些认为魔鬼和恶人的惩罚并非永恒的人。}

首先，我们应当探究并认识，为何教会一直无法容忍那种认为魔鬼即使经历了最严厉、最持久的惩罚之后，也会得到洁净或宽恕的观点。因为如此众多圣洁的人，充满\WorkTitle{旧约}和\WorkTitle{新约}的精神，他们并不嫉妒任何等级或品性的天使在经过痛苦的净化后享受天上王国的福乐，而是他们意识到自己无法推翻或取消主所说、并在审判中宣布的神圣判决：\BibleQuote{你们可咒骂的，离开我，到那给魔鬼和他的使者预备的永火里去。}\BibleRef{玛 25:41}因为这里清楚显示魔鬼和他的使者将在永火中焚烧。还有\WorkTitle{默示录}中的宣告：\BibleQuote{那欺骗他们的魔鬼，被投入硫磺火湖里，那里也有兽和假先知。他们必昼夜受痛苦，直到永远。}\BibleRef{默 20:10}在前一段经文中用了\Quote{永远}，在后一段用了\Quote{永永远远}；圣经通常用这些词来指无限的时间。因此，没有其他理由，没有更明显、更正当的理由，可以使我们坚定地相信（作为最真实虔敬的不移信仰）魔鬼和他的使者绝不会再归向圣人的正义和生命，这理由就是：圣经——它从不欺骗人——说天主没有宽恕他们，而是预先定了他们的罪，将他们投入地狱的黑暗牢狱中，\BibleRef{伯后 2:4}等候最后审判的日子，那时永火将吞没他们，他们在其中受痛苦，永无穷尽。如果真是这样，怎能相信所有人，甚至有些人，在遭受一段时间的惩罚之后会被从痛苦中解救出来？如果不削弱我们对魔鬼永久惩罚的信仰，怎能相信这一点？因为如果所有或部分将听到\Quote{你们可咒骂的，离开我，到那给魔鬼和他的使者预备的永火里去}\BibleRef{玛 25:41}的人，不总是待在那火里，那么有什么理由相信魔鬼和他的使者将永远在那里？或许天主对恶人和天使都宣告的判决，对天使是真的，对人却是假的？显然，如果人的臆测比天主的话更重，就会如此。但因为这是荒谬的，那些希望免除永久惩罚的人，应当停止与天主争论，而是趁还有机会时，服从天主的命令。那么，认为永久惩罚意味着长时间的惩罚，而永生意味着无终的生命，这是多么愚蠢的幻想！因为在同一段经文中，基督在同一句话中用了相似的词语提到两者：\BibleQuote{这些人要进入永罚，义人却要进入永生。}\BibleRef{玛 25:46}如果两种结局都是\Quote{永远的}，那么我们必须要么将两者理解为持续很久但最终终止，要么两者都是无终的。因为它们是相关的——一方面是永罚，另一方面是永生。而在一句话中说永生是无终的，永罚却会终止，这是极其荒谬的。因此，正如圣人的永生是无终的，同样，那些被判永罚之人的永久惩罚也将没有终点。\par

% structure: p0072 source=h2 target=subsection reason=relative-heading-rank; base=h2

\subsection{第二十四章——驳斥那些幻想在天主的审判中，所有被告都会因圣人的祈祷而得以幸免的人}

这个推理同样驳斥了那些人，他们为了自身利益，却假借更大的温柔精神，试图否定天主的话，并断言这些话是真实的，不是因为人们要承受天主所威胁的那些事，而是因为他们理应承受这些事。他们说，天主必会应允祂的圣人们的祈祷，那时圣人们因更加成全的圣德，将更加热切地为他们的敌人祈祷，他们的祈祷也将更有效、更堪当天主的垂听，因为他们现在已从一切罪恶中净化了。那么，如果在那成全的圣德中，他们的祈祷如此纯洁且全能，为什么他们不为那些预备了永火的天使使用这些祈祷，好使天主减轻祂的判决，改变它，并使他们脱离那火呢？或许会有人大胆断言，即使圣天使也会与圣人（那时成为与天主天使同等的人）联合，为有罪的，无论是人还是天使，代祷，好使慈悲饶恕他们，避开真理已宣告他们应受的惩罚？但没有任何信仰纯正的人说过这样的话；将来也不会。否则，教会为何现在不为魔鬼及其天使祈祷，因为她的主天主命令她为仇敌祈祷。所以，阻止教会现在为邪恶的天使祈祷（她知道他们是她的敌人）的理由，正是那阻止她（无论圣德多么成全）在最后审判时为那些将受永火惩罚的人祈祷的同一个理由。目前她为人类中的敌人祈祷，因为他们还有机会作有效的悔改。她特别为他们祈求什么，不就是宗徒所说的\Quote{天主赐予他们悔改}，\Quote{使他们从魔鬼的网罗中醒悟过来，魔鬼按自己的意思俘虏了他们吗？}\BibleRef{弟后 2:25}但如果教会知道那些人是谁，虽然他们仍活在世上，却已预定与魔鬼进入永火，那么她就不能再为他们祈祷，如同不为魔鬼祈祷一样。但既然她对任何人都没有这种确定性，她就为所有仍在世上活着的敌人祈祷；然而她的祈祷并非都为所有人所垂听。她只在那些虽然反对教会，却预定通过她的代祷而成为她子女的人身上得到垂听。但如果有人怀着不肯悔改的心直到死亡，没有从敌人转变为子女，教会还继续为他们祈祷吗？为那些灵魂，\Emph{即}，已故之人的灵魂？为什么她停止为他们祈祷？除非因为此人尚在肉身时未被迁入基督的王国，如今被判定属于撒殚的随从吗？\par

我说，正是同一个理由阻止教会任何时候为邪恶天使祈祷，也正是这个理由阻止她以后为那些将受永火惩罚的人祈祷；这也是为什么，尽管教会甚至为恶人祈祷，只要他们还活着，但她甚至今世也不为已死的不信者和不敬神的人祈祷。因为对于某些死者，教会或虔诚个人的祈祷确实蒙垂听；但那是为那些人，他们已在基督内重生，生活并非邪恶到被认为不配得到这种怜悯，也并非良善到被认为不需要它。同样，在复活之后，将有一些死者，在经历了亡者灵魂应有的痛苦后，将获得慈悲，并从永火的惩罚中获释。因为如果没有一些人，他们的罪虽未在今世得赦，却要在来世得赦，那么就不能真实地说：\Quote{他们不得赦免，既不在今世，也不在来世。}\BibleRef{玛 12:32}但当审判生者死者的那一位说过：\Quote{我父所祝福的，你们来吧！承受自创世以来为你们预备的国度吧！}又对另一边的人说：\Quote{你们这些受诅咒的，离开我，到那为魔鬼及其天使预备的永火里去吧！}以及\Quote{这些人要进入永罚，而义人却要进入永生。}那么，若说天主所说将进入永罚的人中任何一人的惩罚不是永恒的，便是极其狂妄的，并因此给永生对应的应许带来绝望或怀疑。\par

那么，谁也不应当这样理解圣咏作者的话：\BibleQuote{难道天主忘记慈悲了吗？难道祂在怒气中关闭了祂的怜悯吗？}似乎天主的判决对于善人是真的，对于恶人是假的，或者对于善人和恶天使是真的，但对于恶人是假的。因为圣咏作者的话是指慈悲的器皿和恩许的子女，先知自己就是其中之一；因为当他说\BibleQuote{难道天主忘记慈悲了吗？难道祂在怒气中关闭了祂的怜悯吗？}之后，立刻加上\BibleQuote{我说：现在我开始了；这是至高者右手所行的改变。}他清楚地解释了他所说\BibleQuote{难道祂在怒气中关闭了祂的怜悯吗？}的意思。因为天主的怒气就是这必死的生命，人在其中变得如同虚空，他的日子如同影儿过去。然而在这怒气中，天主没有忘记慈悲，使祂的太阳升起，降雨给义人和不义的人；\BibleRef{玛 5:45}因此，祂并没有在怒气中减少祂的怜悯，特别是在圣咏作者所说的\BibleQuote{现在我开始了；这改变是从至高者的右手而来的。}因为祂为慈悲的器皿改变得更好，即使他们仍然处于这最悲惨的生命中，这就是天主的怒气，即使祂的怒气在这悲惨的腐朽中显明；因为\BibleQuote{在祂的怒气中，祂没有关闭祂的怜悯。}既然这首神圣诗歌的真理因这样的应用而完全满足，就没有必要把它引用到那些不属于天主之城的人在永火中受罚的地方。但如果有人坚持将其应用扩展到恶人的痛苦，至少让他们这样理解：天主的怒气，曾以永罚威胁恶人，将继续存在，但会与怜悯相混合，以减轻本可公正施加的痛苦；这样，恶人既不会完全逃脱，也不会仅仅暂时忍受这些威胁的痛苦，而是所受的痛苦会比他们应得的更轻、更可忍受。因此，天主的怒气将继续，同时祂不会在这怒气中关闭祂的怜悯。但我并不应当被认为肯定这一假设，因为我并不坚决反对它。\par

至于那些认为这些经文只是空洞的威胁而非真理的人：\BibleQuote{离开我，可咒骂的，进入永火里去！}和\BibleQuote{这些人要进入永罚里去；}\BibleRef{玛 25:41}以及\BibleQuote{他们必被折磨直到永永远远；}\BibleRef{默 20:10}和\BibleQuote{他们的虫子不死，他们的火也不灭，}\BibleRef{依 66:24}——我说，这样的人被极其有力地、充分地驳斥了，与其说是被我，不如说是被神圣的圣经本身。因为尼尼微人在今生悔改了，因此他们的悔改是富有成果的，因为他们在那块主愿意人们流泪播种、日后可以欢乐收割的田地里播种了。然而，谁能否认天主的预言在他们身上应验了呢？如果至少他观察到天主摧毁罪人不仅出于愤怒，也出于怜悯。因为罪人被摧毁有两种方式：要么像索多玛人那样，他们自身因罪受罚；要么像尼尼微人那样，他们的罪被悔改摧毁。因此，天主的预言应验了——邪恶的尼尼微被推翻了，而善良的尼尼微被建立起来。因为它的城墙和房屋仍然屹立；城市是在它堕落的风俗中被推翻的。因此，虽然先知因他所预言的毁灭没有发生而生气（居民因他的预言而恐惧），但天主预知所预言的确实发生了，因为预言毁灭的那一位知道它应如何以一种较轻的灾难方式应验。\par

然而，为使这些反常富有同情心的人明白这些话的含意：\BibleQuote{上主，祢的慈爱何其丰沛，祢为敬畏祢的人所珍藏的}，让他们阅读下文：\BibleQuote{祢为仰望祢的人所成全的}。因为\BibleQuote{祢为敬畏祢的人所珍藏的}和\BibleQuote{祢为仰望祢的人所成全的}是什么意思呢？岂不是说：那些因惧怕刑罚而企图借法律建立自己的正义的人，不认识天主的正义，因他们对此无知？他们没有尝到它。因为他们仰望自己，而非仰望祂；因此天主的丰沛慈爱向他们隐藏。他们确实敬畏天主，但那是奴性的恐惧，\BibleQuote{那不在爱内，因为完美的爱驱逐恐惧}。\BibleRef{若一 4:18}因此，祂为仰望祂的人成全祂的慈爱，以祂自己的爱激励他们，使他们怀着一种不被爱驱逐、反而永远长存的圣洁敬畏，当他们夸耀时，在主内夸耀。因为天主的正义就是基督，\BibleQuote{祂由天主成为我们的}，如宗徒所说，\BibleQuote{智慧、正义、圣化与救赎：正如经上记载：夸耀的，当在主内夸耀}。\BibleRef{格前 1:30}这天主的正义，即无功劳的恩宠恩赐，不被那些企图建立自己的正义、因此不服从天主的正义（即基督）的人所认识。\BibleRef{罗 10:3}但正是在这正义中，我们找到天主慈爱的丰沛，圣咏说：\BibleQuote{你们要尝尝，并看看上主是何等甘饴}。而在此朝圣之旅中，我们是尝到它，而非饱享它。我们现在渴慕它，为的是将来当我们在祂内看见祂时，得以满足，并应验所记载的：\BibleQuote{当我看见祢的光荣显现时，我就心满意足}。如此，基督为仰望祂的人成全祂丰沛的慈爱。但若天主向敬畏祂的人隐藏祂的慈爱，如这些反对者所幻想的，即人对天主对恶人的慈悲旨意的无知，可以引导他们敬畏祂而生活得更好，并且可以为那些生活不当的人祈祷，那么，祂如何成全祂对仰望祂的人的慈爱呢？因为，如果他们的梦想成真，正是这慈爱会阻止祂惩罚那些不仰望祂的人？那么，让我们寻求祂的慈爱，就是祂为仰望祂的人所成全的，而不是祂被认为为那些藐视和亵渎祂的人所成全的；因为在今生之后，人若忽略了在此生所当预备的，便徒然寻求。\par

然后，关于宗徒的那句话：\BibleQuote{因为天主把众人都禁锢在不信之中，为要怜悯众人}，\BibleRef{罗 11:32}这并不意味着祂不谴责任何人，但前文表明了其含意。宗徒写信给已经信主的外邦人；当他向他们谈及尚未信主的犹太人时，他说：\BibleQuote{正如你们从前不信天主，但如今因他们的不信而蒙受了怜悯；同样，这些人现在也不信，为的是因你们的怜悯，他们也可获得怜悯}。然后他加上了这些有争议的话，这些人是用来欺骗自己的：\BibleQuote{因为天主把众人都禁锢在不信之中，为要怜悯众人}。所有人？若不是指他所谈论的所有人，就好比他说：\Quote{你们和他们双方？}天主于是把那些祂所预知和预定要与祂圣子的肖像相似的犹太人和外邦人，都禁锢在不信之中，为的是使他们因不信的苦楚而羞愧，并悔改，以信心转向天主慈悲的甘饴，并应取用圣咏中的感慨：\BibleQuote{上主，祢的慈爱何其丰沛，祢为敬畏祢的人所珍藏的，却为仰望祢的人所成全的}，不是在自己内，而是\Quote{在祢内}。因此，祂怜悯所有怜悯之器。而\InnerQuote{所有人}是什么意思？就是那些祂所预定、召叫、成义、光荣的犹太人和外邦人：这些人中没有一个会被祂定罪；但我们不能说所有的人都无一例外。\par

% structure: p0079 source=h2 target=subsection reason=relative-heading-rank; base=h2

\subsection{第25章——那些接受异端圣洗、后来堕入邪恶生活的人；或那些接受天主教圣洗、后来转向异端和分裂的人；或那些留在他们领受圣洗的天主教会内、但继续过着不道德生活的人——能否藉圣事的德能希望获得永恒惩罚的赦免？}

但现在让我们回答那些人，他们许诺从永火中得救，不是给魔鬼和他的天使（正如我们刚才谈论的那些人也不许诺），甚至也不是给所有的人类，而只给那些被基督的圣洗洗涤、并成为祂体血分享者的人，无论他们如何生活，无论他们陷入何种异端或不虔敬。但他们被宗徒反驳，他说：\BibleQuote{肉身的行为是明显的，就是：淫乱、污秽、放荡、崇拜偶像、邪术、仇恨、争执、嫉妒、愤怒、纷争、异端、嫉妒、酗酒、狂欢，以及类似的事。我事先告诉你们，如同我以前说过的：行这样事的人，不能承受天主的国。}\BibleRef{迦 5:19}显然，如果这样的人经过一段时间之后得救，并承受天主的国，那么宗徒的这句话就是假的。但既然它不是假的，他们就一定永远不能承受天主的国。而且，如果他们永远不能进入那国，那么他们就将永远被拘留在永罚中；因为没有被接纳进入那国的人，没有一个中间之地可以使他不受惩罚地生活。\par

因此，我们可以合理地询问，应当如何理解主耶稣的这些话：\Quote{这是从天上降下来的食粮，人吃了就不死。我是从天上降下的生活的食粮；谁若吃了这食粮，就永远活着。}\BibleRef{若 6:50}事实上，我们现在所答复的人，他们对这段经文的解释，被我们将要简短答复的人所驳斥；后者并不将这种解脱应许给所有领受了圣洗圣事和主圣体的人，而只应许给天主教徒，无论他们生活多么邪恶；因为他们说，这些人不仅圣事性地，而且实际地吃了主的圣体，因为他们成了祂身体的肢体，关于这身体，宗徒说：\Quote{我们虽多，却是同一个饼，同一个身体。}\BibleRef{格前 10:17}因此，凡在基督身体的合一内（即在天主教徒的团体内）的人——这身体的圣事，信徒惯常在祭台上领受——此人真正被称为吃了基督的身体、喝了基督的血。因此，异端者和分裂者既与这身体的合一分离，他们能够领受同样的圣事，但对自己毫无益处——反而有害，以致他们更严厉地受审判，而不是在某些时候后获得解脱。因为他们不在那由这圣事所象征的和平联系中。\par

但是，那些充分理解不在基督身体内的人不能被称为吃基督身体的人，他们也同样错误：他们应许那些从这身体的合一堕入异端、甚至堕入外教迷信的人，会从永罚之火中解脱。首先，他们应当考虑这是多么不可容忍，多么与健全教义不符：竟认为许多、甚至几乎所有离开天主教会并创立不虔敬异端、成为异端首领的人，其命运要优于从未是天主教徒却落入这些人陷阱的人；也就是说，如果他们天主教圣洗和最初在基督真身体内领受基督身体之圣事的事实，足以使这些异端首领从永罚中解脱。因为那离弃信德、并从离弃者变为攻击者的人，比起从未持守所不信之信德的人更坏。其次，他们被宗徒驳斥，宗徒在列举血肉之工后，论及异端说：\Quote{做这样事的人，不能继承天主的国。}\par

因此，那些过着放肆和该受惩罚的生活的人，也不应对救恩有信心，即使他们在天主教会共融中坚持到底，并以此话安慰自己：\Quote{那坚持到底的，必得救。}他们因生活的邪恶，离弃了那在基督内所是正义的生活，无论是通过奸淫，或在身体上犯下宗徒连提都不愿提的其他不洁行为，或放荡奢侈，或行任何他所说\Quote{做这样事的人，不能继承天主的国}的事。因此，做这样事的人，除了在永罚中，不可能在任何地方存在，因为他们不能在神的国里。因为如果他们坚持这样的事直到生命终结，就不能说他们是在基督内坚持到底；因为留在祂内就是留在对基督的信德内。而这信德，按照宗徒的定义，\Quote{以爱德行作}。\BibleRef{迦 5:6}而\Quote{爱德}，他处说，\Quote{不作恶}。\BibleRef{罗 13:10}这些人也不能被称为吃了基督的身体，因为他们甚至不能被算作祂的肢体。且不说其他理由，他们不能同时是基督的肢体和娼妓的肢体。最后，祂自己说：\Quote{谁吃我的肉，并喝我的血，便住在我内，我也住在他内，}\BibleRef{若 6:56}表明实际地、而非圣事性地吃祂的肉、喝祂的血是什么；这就是住在基督内，祂也住在我们内。所以，就好像祂说：那不住在我内、我也不住在他内的人，不要说，也不要以为他吃了我的肉或喝了我的血。因此，不是基督肢体的人，不住在祂内。而那些使自己成为娼妓肢体的人，不是基督的肢体，除非他们悔改而放弃那恶，并回到这善中，与它和好。\par

% structure: p0084 source=h2 target=subsection reason=relative-heading-rank; base=h2

\subsection{第二十六章——以基督为根基是什么意思，以及那些应许如经火而得救的人是谁。}

但是，他们说，天主教基督徒以基督为根基，并且无论他们在这根基上建造多么堕落的生活，如木、草、禾秸，都没有脱离与祂的合一；因此，使基督成为他们根基的正确信德，足以使他们经过一段时间从那持续的火中得救，虽有损失——因为他们在上面建造的东西要被烧毁。让宗徒雅各伯简要地回答他们：\Quote{若有人说他有信德，却没有行为，这信德能救他吗？}\BibleRef{雅 2:14}那么，他们问，宗徒保禄所说的\Quote{他自己却要得救，可是像从火里经过一样}是指谁？让我们与他们一同询问；有一点是非常确定的，那指的不是雅各伯所说的人，否则我们就会使两位宗徒互相矛盾，如果一位说：\Quote{即便人的行为是邪恶的，他的信德仍会像从火里一样救他，}而另一位说：\Quote{他若没有善行，这信德能救他吗？}\par

我们将先确定谁能得救却要经过火，如果我们首先发现以基督为根基是什么意思。这我们从图像本身很容易学到。在建筑中，根基是首要的。因此，凡是以基督在心中，以至于没有任何属地或暂时的事物——即使是那些合法的、允许的——被置于祂之上，这人就是以基督为根基。但如果这些事物被置于优先，那么即使一个人似乎有信仰基督，基督却不是那人的根基；更何况如果他藐视健全的诫命，追求被禁止的享乐，他显然是被定罪为不把基督放在首位而是末位，因为他藐视祂作为他的主宰，并且宁愿满足自己邪恶的私欲，藐视基督的命令和许可。因此，如果有基督徒爱一个妓女，并依附于她，成为一体，他现在就没有以基督为根基。但如果有人爱自己的妻子，并按照基督所愿的方式爱她，谁能怀疑他以基督为根基呢？但如果他按世俗的方式爱她，属肉体的，如同情欲的病态所驱使，如同不认识天主的外邦人所爱，即使是这个，宗徒，或者更确切地说，基督借着宗徒，许可为小罪。因此，即使这样的人也可能以基督为根基。因为只要他不将这样的情感或快乐置于基督之上，基督就是他的根基，尽管他在其上建造木、草、禾秸；因此他将得救，却要经过火。因为苦难的火将烧尽此类奢侈的快乐和属地的爱慕，尽管它们并非可咒骂的，因为是在合法的婚姻中享受的。而这火的燃料是丧亲之痛，以及所有消耗这些喜乐的灾祸。因此，上层建筑对于建造它的人将是损失，因为他不能保有它，反而会因失去那些他曾从中找到快乐的事物而痛苦。但借着这火，他将因根基的德能而得救，因为即使一个迫害者问他是否愿保留基督或这些事物，他会选择基督。你愿意听宗徒自己的话，谁是在根基上建造金、银、宝石的呢？\Quote{没有结婚的，}他说，\Quote{挂虑主的事，怎样悦乐主。}\BibleRef{格前 7:32} 你愿意听谁建造木、草、禾秸呢？\Quote{但结了婚的，挂虑世俗的事，怎样悦乐妻子。}\BibleRef{格前 7:33} \Quote{各人的工程将来必显露：因为那日子会将它表明，}——那日子无疑是苦难的日子——\Quote{因为它要，}他说，\Quote{被火揭露。}\BibleRef{格前 3:13} 他称苦难为火，正如别处所说：\Quote{炉火考验陶匠的器皿，患难的考验考验义人。}\BibleRef{德 27:5} 并且\Quote{火要试验各人的工程是怎样的。若有人在这根基上建造的工程存留得住}——因为一个人挂虑主的事，怎样悦乐主，是存留得住的——\Quote{他就要得赏报}——也就是说，他将收获他挂虑的果实。\Quote{但若有人 的工程被烧毁，他就要受亏损}——因为他所爱的，他不能保留——\Quote{他自己却要得救}——因为没有任何苦难能把他从那稳固的根基上移动——\Quote{可是要像从火中经过一样；}\BibleRef{格前 3:14}\par

但是，如果这段（哥林多前书）经文是指那火，关于它主将要对祂左边的人说：\Quote{可咒骂的，离开我，到永火里去，}\BibleRef{玛 25:41} 以致我们在这些人中相信有那些在根基上建造木、草、禾秸的人，并且他们借着善的根基，过一段时间后将从作为他们恶行报应的火中被释放，那么我们对右边的人，就是那些将要听到\Quote{我父所祝福的，来吧！承受为你们预备的国度，}\BibleRef{玛 25:34} 的人，将作何想法？除非他们是那些在根基上建造金、银、宝石的人。但如果我们的主所说的火与宗徒所说的\Quote{可是要像从火中经过一样，}\BibleRef{格前 3:14} 那么两边——即右边和左边——都要被投入其中。因为那火要考验两者，正如经上所说：\Quote{主的日子将会表明，因为它要被火揭露；火要试验各人的工程是怎样的。}\BibleRef{格前 3:13} 因此，如果火要考验两者，以便若有人工程存留——\Emph{i.e.}，若上层建筑不被火烧毁——他可得赏报，而若他的工程被烧毁，他受亏损，那么那火肯定不是永火本身。因为只有左边的人才会被投入这后者之火，且带着最终永久的惩罚；但前者之火考验右边的人。其中一些人的工程它考验的方式是不烧毁和消耗他们在基督根基上建造的结构；而另一些人它却以另一种方式考验，烧毁他们所建造的，使他们受亏损，但他们自己得救，因为他们持守了基督——祂被放置作为他们稳固的根基——并且爱祂超越一切。但如果他们得救，那么他们当然要站在右边，并与其余的人一同听到判决：\Quote{我父所祝福的，来吧！承受为你们预备的国度；} 而不是在左边，那将是那些不得救的人所在之处，因而他们将听到判决：\Quote{可咒骂的，离开我，到永火里去。} 因为没有人能从那里得救，因为他们都要进入永罚，那里的虫子不死，火也不灭，他们在那里日夜受折磨，直到永远。\par

但若有人说，在这身体死后与那审判和报应的末日（即在复活之后）之间的时段里，死者的身体将遭受一种火，其性质是：它不伤害那些在此生中没有沉溺于那些如木、草、禾秸般将被烧尽的享乐和追求的人，却伤害那些携带了这类建构的人；若有人说，这种世俗性（作为小罪）将在磨炼之火中被烧尽，或仅在此世，或在此世与来世皆有，或在此世以便不在来世——对此我不反驳，因为可能是真的。因为也许连身体死亡本身也是这磨炼的一部分，它源自首次违命，以致死后那段时间，因各人建筑的质性而各具色彩。那些使殉道者获得荣冠的迫害，以及各类基督徒所受的迫害，像火一样考验两种建筑：有些建筑连同建筑者本身一起被烧毁，如果基督不是它们的基础；有些建筑则被烧毁而建筑者不被烧毁，因为找到了基督在他们内，他们便得救，但如同经过火一样；另有一些建筑则不被烧毁，因为在其中找到了永恒的材料。在世界的终结，在假基督的时代，将有前所未有的磨炼。那时将有多少建筑，或金或草，建于最好的基础基督耶稣之上，那火将考验它们，给一些人带来喜乐，给另一些人带来损失，但都不致毁灭，因为基础稳固！但无论何人，我不说他的妻子（他与之过肉欲生活），而是任何亲属，他们不提供这类快乐，而理应被爱——无论何人将这些置于基督之上，并以人性的、肉欲的方式爱他们，便没有基督为基础，因此不会因火而得救，实际上根本不得救；因为他不可能与救主同住，救主关于这事说得非常明确：\BibleQuote{那爱父亲或母亲超过我的，不配属于我；那爱儿子或女儿超过我的，不配属于我。}\BibleRef{玛 10:37}但若有人肉欲地爱他的亲属，却并不将他们置于基督之上，反而在受考验时宁可失去他们也不愿失去基督，他将因火而得救，因为失去这些亲属，他必须按他爱的程度受苦。而那按基督的方式爱父母、儿女的人，为帮助他们获得祂的王国并依附于祂，或因为他们基督的肢体而爱他们，天主禁止这爱如木、草、禾秸般被烧毁，反而应算作金、银、宝石的建筑。因为一个人若只为基督的缘故而爱那些人，又怎能爱他们超过基督呢？\par

% structure: p0089 source=h2 target=subsection reason=relative-heading-rank; base=h2

\subsection{第二十七章。——反对那些认为与施舍相伴的罪过不会伤害他们的人之信念。}

尚需答复那些主张只有忽视与自己罪过相称的施舍的人才会在永火中焚烧的人，他们以雅各伯宗徒的话为依据：\BibleQuote{不行怜悯的人，将受无怜悯的审判。}\BibleRef{雅 2:13}因此，他们说，那施行了怜悯的人，即使他没有改掉放荡的生活，反而在施舍丰富的同时仍邪恶不义地生活，他将受到怜悯的审判，以致他要么完全不被定罪，要么在时间之后从最终审判中被释放。同样，他们设想基督将区分右边的和左边的，将一些人送进祂的王国，另一些人送入永罚，唯一的标准是他们对爱德工作的注意或忽视。此外，他们还试图用主亲自教导的祈祷作为他们意见的证明和屏障，就是那些从未放弃的日常罪过可以通过施舍得赦免，不论它们多么冒犯或何种类。因为，他们说，如同没有一天基督徒不该用这祈祷，所以没有任何一种罪过（即使每天犯）不被赦免，当我们说\BibleQuote{宽免我们的罪债，}如果我们注意履行随后的话\BibleQuote{如同我们宽免我们的债务人。}\BibleRef{玛 6:12}因为他们接着说，主没有说：\Quote{如果你们宽免别人的过犯，你们的天父也必宽免你们的小罪；}而是说\Quote{必宽免你们的罪。}因此，他们推测，不论何种或多大罪的，即使每日犯且在此生从未放弃或制服，都可以通过施舍得赦免。\par

但他们正确地强调要施舍与过去的罪相称；因为若他们说任何施舍都能获得天主对每天犯下且习惯性严重的大罪的赦免，若他们说这样的罪能每天被赦免，他们就会看到自己的道理是荒谬可笑的。因为这样一来，他们就会被逼承认，一个非常富有的人可以通过每天施舍十枚小钱而买得谋杀、奸淫和一切邪恶的赦免。如果承认这一点是极其荒谬和疯狂的，并且我们仍然要问，那些合适的施舍是什么，连基督的前驱也说过：\Quote{你们要结出与悔改相称的果实}，\BibleRef{玛 3:8} 无疑会发现它们并不是那些每天甚至到末了都严重败坏自己生命的人所行的施舍。因为他们以为，用强取豪夺所得财富的一小部分施舍给穷人，就能平息基督，从而可以无罪地犯下最该受罚的罪，自以为已从祂那里买得了犯罪的许可，或者更确切地说，每天买得宽恕。即使他们为了一件罪行而将全部财产分给基督的有需要的肢体，这也不会对他们有益，除非他们停止所有类似的行为，并达到不作恶的爱德。因此，那行与罪相称的施舍的人，必须先开始从自身做起。因为一个人对邻人施行爱德，却不向自己施行，这是不合理的，因为他听到主说：\Quote{你当爱你的邻人如同你自己}，\BibleRef{玛 22:39} 又说：\Quote{你当怜恤你的灵魂，使天主喜悦}，\BibleRef{德 30:24} 那么，一个人若没有怜恤自己的灵魂以取悦天主，又怎能说他行了与罪相称的施舍呢？与此相同，经上写着：\Quote{对自己刻薄的人，对谁能好呢？} \BibleRef{德 21:1} 因此，我们应当行施舍，好使我们祈求赦免过去之罪的祈祷蒙垂听，而不是在我们继续犯罪时，以为可以通过施舍为自己提供犯罪的许可证。\par

因此，我们预言祂将把右边的人所行的施舍归给他们，并指控左边的人没有行施舍，理由是这样祂可以表明爱德对于消除过去之罪的效力，而不是允许他们持续犯罪而不受惩罚。而且，那些拒绝放弃邪恶的生活习惯、不转向更好道路的人，不能说他们行了爱德行为。因为这话的意思正是：\Quote{你们没有给这些最小中的一个做，就是没有给我做。} \BibleRef{玛 25:45} 祂向他们表明，即使他们自以为在行慈善，实际上并没有行。因为如果他们给饥饿的基督徒面包，是因为他是基督徒，那么他们肯定也不会拒绝将义德之粮，即基督本身，给予自己；因为天主看重的不是受赠的对象，而是行善的心意。因此，谁若在基督徒身上爱基督，他施行爱德给基督徒的心意，与他接近基督的心意相同，而不是那种若能免受惩罚就会离弃基督的心意。因为一个人越爱基督所不赞同的事，就越离弃基督。一个人受了洗，若不成义，对他有什么益处呢？难道那位说：\Quote{人除非由水和圣神而生，不能进天主的国}，\BibleRef{若 3:5} 的主，不也说：\Quote{除非你们的义德超过经师和法利塞人的义德，你们决进不了天国}，\BibleRef{玛 5:20} 吗？为什么许多人因害怕第一句话而跑去领洗，却很少有人因害怕第二句话而寻求成义？因此，正如一个人称他的兄弟为\Quote{愚人}时，如果他愤怒的不是兄弟的身份，而是犯罪者的罪过，那么他就不算对这兄弟说话（否则他就要受地狱的火）；同样，一个人向基督徒施行爱德，如果他在基督身上不爱基督，那他就不算是向基督徒施行爱德。谁若不愿在基督内成义，就不爱基督。或者，如果有人犯了称兄弟为愚人的罪，不正当地辱骂他，而没有去除他的罪的意愿，那么他的施舍对于弥补这一过失作用甚微，除非他再加上同一段经文所要求的和好的补救。因为那里说：\Quote{所以，你若在祭台前献你的礼物时，在那里想起你的兄弟有什么怨你的事，就把你的礼物留在祭台前，先去与你的兄弟和好，然后再来献你的礼物。} \BibleRef{玛 5:23} 同样，无论施舍有多大，只要犯罪者继续行罪，对于任何罪来说，施舍都是微不足道的。\par

至于主亲自教导的每日祈祷，因此称为主祷文，它确实消除当日的罪过——当我们每日念说：\BibleQuote{免我们的债}，并且我们不仅口念，也实行随后的话：\BibleQuote{如同我们免了人的债}；\BibleRef{玛 6:12}——但我们发出这个恳求，是因为罪已犯下，而不是要人去犯罪。因为藉此，我们的救主设计教导我们：无论我们在这一软弱与黑暗的生命中活得多正义，我们仍会犯罪，为此我们应当祈求赦免，同时我们必须宽恕得罪我们的人，好使我们自己也得蒙宽恕。主随后说出的话：\BibleQuote{如果你们宽恕别人的过犯，你们的天父也必宽恕你们的过犯}，\BibleRef{玛 6:14}并非要我们从此恳求中获得自信，以致可以日日安然犯罪——或自认为凌驾于人间法律之上，或狡猾地欺骗人关于我们的行为——而是为了叫我们由此学习不要以为我们没有罪，即使我们远离大恶；正如天主也为了同样的目的提醒旧约的司祭，关于他们的祭献，祂命令他们先为自己的罪献祭，再为百姓的罪献祭。因为连如此伟大的导师与主的言语，也需用心思考。祂没有说：如果你们宽恕人的罪，你们的天父也必宽恕你们的罪，无论是什么罪；祂却说：你们的罪。因为祂教导的是每日的祈祷，而祂说话的对象显然是已经成义的门徒。那么，祂说\BibleQuote{你们的罪}是什么意思呢？岂不就是那些连你们这些已经成义且被圣化的人也无法免除的罪吗？因此，那些藉此恳求寻求放纵于习惯性之罪的人认为主是指大罪，因为祂没有说：祂会宽恕你们的小罪，而是说\BibleQuote{你们的罪}；另一方面，我们考虑到祂所训诲之人的品性，无法将\BibleQuote{你们的罪}解释为其他，只能是小罪，因为这样的人已不再犯大罪。然而，即使是那些我们须以彻底改变生活来逃避的大罪本身，若祈祷者不遵守附带的诫命\BibleQuote{如同你们也宽恕你们的债主}的话，也不会得到赦免。因为如果连义人生活中所沾染的最小罪过，若不凭此条件，也不能获得赦免，那么那些犯了许多大罪的人，如果停止行恶，却仍毫不留情地拒绝宽恕别人对自己的任何亏欠，他们离获得宽恕岂不更远？因为主说：\BibleQuote{但如果你们不宽恕别人的过犯，你们的天父也必不宽恕你们的过犯}。\BibleRef{玛 6:15}因为这也是宗徒雅各伯所说的话的意思：\BibleQuote{那不怜悯人的，也要受无怜悯的审判}。\BibleRef{雅 2:13}因为我们应记住那个仆人，他的主人取消了他一万塔冷通的债，但后来却命令他还清，因为那仆人自己对他欠他一百德纳的同伴没有怜悯。\BibleRef{玛 18:23}宗徒雅各伯接着说的话：\BibleQuote{怜悯胜过审判}，\BibleRef{雅 2:13}应验在那些恩许的子女和慈悲的器皿身上。因为甚至那些以如此圣洁生活、以致也接纳那些以不义的钱财赢得他们友谊的人进入永远的帐幕的义人，\BibleRef{路 16:9}也是藉着那位使不虔敬的人成义、按恩宠而非按债赐予报酬的祂的慈悲拯救而成为这样的。因为在这数目中，有那位宗徒，他说：\BibleQuote{我蒙了怜悯，得以忠信}。\BibleRef{格前 7:25}\par

但必须承认，那些被接入永远的帐幕的人，并非具有那种品性，以致他们自己的生命足以拯救他们而无须圣人的帮助；因此，在他们身上，怜悯特别胜过审判。然而，我们不可因此认为，每个不悔改的放荡之徒，只要他曾经用不义的钱财帮助过圣人——也就是说，用非法获得的金钱或财富，或者即使合法获得，却并非真正的财富，而只是不义所认为的财富，因为它不认识那真正的财富，即那些甚至也接纳别人进入永远帐幕的人所富有的——就可以被接入永远的帐幕。因此，存在着一种生命，它既不是如此恶劣，以致那些采取这种生命的人，即使有任何慷慨的施舍来缓解圣人的需要，并结交能够接他们进入永远帐幕的朋友，也无法被帮助进入天国；也不是如此良善，以致它本身足以为他们赢得那巨大的幸福，如果他们不通过他们所结交之人的功德而获得怜悯的话。我常常惊讶，连维吉尔也应和了主的这句话，祂说：\BibleQuote{你们要用不义的钱财结交朋友，好使他们接你们进入永远的帐幕}；\BibleRef{路 16:9}以及这非常相似的话：\BibleQuote{那因先知的名接纳先知的，必得先知的赏报；那因义人的名接纳义人的，必得义人的赏报。}\BibleRef{玛 10:41}因为当那位诗人描述厄吕西翁原野——他们以为有福者的灵魂住在那里——时，他不仅在那里安置了那些凭自身功德能够到达那居所的人，还加上了——\par

\Quote{和那些因施惠于人而赢得\newline{}感恩怀念的人}；\par

即那些曾服侍他人，因而配得被他们纪念的人。就像他们使用了基督徒口中常说的表达，即某个谦卑的人将自己托付给一位圣人，并说：\Quote{纪念我}，并通过在他面前行善来确保如此行。但我们所谈论的这种生命是什么，以及哪些罪会阻止人凭自己赢得天主的国，却又允许他利用圣人的功德，这很难确定，界定起来也很危险。就我而言，尽管经过了各种研究，直到此刻仍未能发现。或许这是向我们隐藏的，以免我们疏忽于避免这些罪，从而停止进步。因为如果知道了哪些罪在持续存在且不因更高尚的生活而弃绝时，仍不妨碍我们寻求并期望圣人的代祷，那么人的懒惰就会放肆地裹挟在这些小罪之中，不会采取任何步骤通过任何德行的灵巧努力从中解脱，而只希望凭借他人的功德得救，这些人是用不义的钱财\OriginalTerm{不义的钱财}{mammona iniquitatis}赢得的友谊。但现在我们对自己所不知道的、即使坚持也是属于小罪的罪过\OriginalTerm{小罪}{venial}的确切性质无知，我们当然在祈祷和努力进步方面更加警醒，也更加谨慎地用不义的钱财在圣人中为自己结交朋友。\par

但这种救赎，无论是通过自己的祈祷还是圣人的代祷，均保证人不会被投入永火，但并非一旦被投入后，过一段时间就能被救出来。因为甚至那些以为好土地结出丰硕果实、有三十倍、六十倍、一百倍的\BibleRef{玛 13:8}说法是指圣人，以致按他们的功德，有的救三十人，有的救六十人，有的救一百人的人——甚至坚持这种观点的人也普遍倾向于认为这救赎发生在审判之日，而非其后。在这种印象下，有人察觉到人们根据所有人都会包含在这救赎方法中而许诺自己不受惩罚的愚蠢之举，便据说非常机智地评论说：我们应竭力生活得如此美好，以至我们全都应列在那些为别人得解放而代祷的人之中，免得这样的人数太少，以至于在一个人救了三十、另一个救了六十、另一个救了一百之后，仍有许多人无法通过他们的代祷而得救，而其中就有每一个虚妄而轻率地许诺别人劳动果实的人。但针对那些承认与我们相同的圣经权威，却因错误解释而按自己愿望而非圣经教导设想未来的人，已说得够多了。作了这样的答复后，我现在按照承诺结束本书。\par

\SourceRecord{Translated by Marcus Dods.；Nicene and Post-Nicene Fathers, First Series；Vol. 2.；Edited by Philip Schaff.；Buffalo, NY: Christian Literature Publishing Co.,；1887.；<http://www.newadvent.org/fathers/120121.htm>.}

% structure: p0001 source=h1 target=section reason=page-title

\section{天主之城（卷二十二）}

本书论及天主之城的终向，即圣人们的永恒真福；建立并阐明肉身复活的信德；最后以圣人穿上不朽的、属神的身体后将如何生活，作为全书的结尾。\par

% structure: p0003 source=h2 target=subsection reason=relative-heading-rank; base=h2

\subsection{第一章 论天使与人的受造。}

正如我们在前一卷中所许诺的，这最后一卷将讨论天主之城的永恒真福。这真福被称为永恒，并非因为它将持续许多世代而终有尽头，而是因为按照福音所言，\Quote{祂的王国没有终结}\BibleRef{路 1:33}。它也不是享受那种仅靠新的一代接替死去的一代来维持的表面的永久性，如同常青树，新叶不断长出代替枯萎凋落的叶子，而始终保持同样的青翠和茂密的外观；在那城里，所有市民都将不朽，此时人类才第一次享有神圣天使从未失去的福乐。这将由天主，那最全能的城的建立者来实现。因为祂已应许，且不能撒谎，祂已履行了许多应许，并已对祂现在要求相信祂也将成就此事的人，施予了许多未曾应许的恩惠。\par

因为正是祂在起初创造了充满一切有形与可理解存在物的世界，其中祂没有创造比那些受造为有理智、能默观并享受祂、并与我们结为社会的灵更好的受造物；这社会我们称为圣而天上的城，在其中他们赖以生存和得福的养料就是天主自己，犹如他们共同的食物和营养。正是祂赋予了这理智性体以自由意志，其性质是：若他愿意离弃天主，\Emph{即}他的真福，立即产生痛苦。正是祂，预知某些天使将因骄傲而想自足于自己的真福，并离弃他们的大善，却没有剥夺他们这种能力，认为从恶中引出善比阻止恶的发生更符合祂的能力和良善。的确，恶永远不会存在，要不是那可变的性体——虽然是善的，但可变的，由至高天主和不变的善所创造，祂创造了万善——因罪恶而将恶带给了自己。而这罪恶本身证明其性体原本是善的。因为若它当初不是非常好，尽管不及它的造物主，那么离弃天主作为它的光明就不会对它成为恶。因为正如瞎眼是眼睛的缺陷，而这事实本身就表明眼睛是为看见光明而被造的，因此缺陷本身表明眼睛比其它肢体更卓越，因为它能接受光明（否则，缺乏光明就不会是眼睛的缺陷）；同样，一度享有天主的性体，即使通过它的缺陷也教导我们它被造为最好的，因为现在它因不享受天主而痛苦。正是祂以最公正的刑罚，将自愿堕落的天使判罚永远的痛苦，并奖赏那些继续依附至高之善的天使，以无尽的稳定作为他们忠诚的奖赏。正是祂也造了人正直，有同样的自由意志——一种地上的动物，确实，但如果他忠于其造物主，就适合天上；如果他离弃祂，就注定要承受适合这种性体的痛苦。正是祂，预知人将因离弃天主并违反祂的律法而犯罪，却没有剥夺他自由意志的能力，因为祂同时预见了祂自己将从恶中引出什么善，以及祂将如何从这应受公正审判的必死种族中，藉恩宠聚集，正如祂现在所做的，如此众多的子民，以至于祂因此填补并修复因堕落天使而留下的空缺，并且那亲爱的天上的城不致因缺少其应有的市民而受亏，或许甚至因更丰盈的人口而欢欣。\par

% structure: p0006 source=h2 target=subsection reason=relative-heading-rank; base=h2

\subsection{第二章 论天主永恒不变的旨意。}

确实，恶人做了许多违反天主旨意的事；但天主的智慧和能力如此伟大，以致一切看似与祂目的相悖的事，仍都趋向于祂自己预知的那些正义而美善的结局和结果。因此，当人们说天主改变祂的旨意时——例如，祂对先前温和的人发怒——这与其说是祂改变，不如说是他们改变，并且他们发现祂改变了，只因他们从祂手中初次体验到痛苦，就像太阳对于受伤的眼睛，变得仿佛从温和变为猛烈，从令人愉悦变为有害，而它本身仍然保持不变。那也称为天主的旨意，是祂在那些遵守祂诫命的人心中所行的；关于这一点，宗徒说：\BibleQuote{是天主在你们内工作，使你们愿意。}\BibleRef{斐 2:13}正如天主的\Quote{义}不仅指祂自己成为义的那义，也指祂在祂所成义的人身上产生的义，同样，那被称为祂的法律的，虽由天主所赐，却更是人的法律。因为耶稣确实是对人说的：\BibleQuote{你们的法律上记载着……}\BibleRef{若 8:17}尽管在另一处我们读到：\Quote{天主的法律在他的心中。}根据这种天主在人心中工作的旨意，祂也被说成愿意那些祂自己并不愿意、却使祂的子民愿意的事；正如祂被称为知道那些祂原本无知的人因祂而知道的事。因为当宗徒说：\BibleQuote{但现今你们认识了天主，更好说为天主所认识}\BibleRef{迦 4:9}时，我们不能认为天主那时才初次认识那些在创世以前已被祂预知的人；而是说祂那时认识他们，是因为祂那时使他们认识。但我记得我在前几卷中讨论过这些表达方式。因此，按照这种旨意——我们说天主愿意那些祂使别人愿意的事，而这些人对未来的事是隐藏的——祂愿意许多祂并不实行的事。\par

因此，祂的圣人们，受祂神圣旨意的感动，渴望许多从未发生的事。例如，他们为某些人祈祷——他们虔诚而圣洁地祈祷——但他们所求的祂并未成就，尽管祂自己借着祂的圣神在他们心中产生了这祈祷的意愿。因此，当圣人们顺应天主的意愿，愿意并祈求众人得救时，我们可以用这种表达方式：天主愿意却未成就——意思是，祂使他们愿意这些事，祂自己也愿意这些事。但如果我们讲论祂那永恒如预知一样的旨意，那么祂确实早已成就了祂在天地间所愿意的一切——不仅过去和现在的事，甚至未来的事。但在祂预知并提前安排好的事发生的时间到来之前，我们说：\Quote{当天主愿意时，这事必发生。}但如果我们不仅不知道它发生的时间，甚至不知道它是否会发生，我们就说：\Quote{如果天主愿意，这事就会发生。}——这并不是因为天主那时会有一个祂先前没有的新意愿，而是因为那从永恒以来就在祂不变的旨意中预备好的事件，那时将要实现。\par

% structure: p0009 source=h2 target=subsection reason=relative-heading-rank; base=h2

\subsection{第三章　论对圣人们永生福乐的应许，以及对恶人永远的惩罚。}

因此，除了其他许多例证之外，正如我们现在在基督身上看到天主向亚巴郎所应许的应验了，当祂说：\BibleQuote{地上万民都要因你的后裔获得祝福}\BibleRef{创 22:18}，同样，祂向同一民族所应许的也必将应验，当祂借先知说：\BibleQuote{在坟墓里的人要复活}\BibleRef{依 26:19}，又说：\BibleQuote{将有新天新地；先前的不再被记念，不再被人思念；但你们要因我所造的永远喜乐欢欣，因为我创造耶路撒冷为人所喜，创造其中的居民为人所乐。我要因耶路撒冷而喜乐，因我的百姓而欢欣；其中不再听见哭泣的声音。}\BibleRef{依 65:17}又借着另一位先知宣告了同样的预言：\BibleQuote{那时，你的人民，凡记录在书上的，都必得救。睡在尘埃中的（或如某些人解释为\InnerQuote{在土丘中}）必有许多人苏醒，有的得永生，有的受羞辱，永远被憎恨。}\BibleRef{达 12:1}又借着同一位先知在另一处说：\BibleQuote{至高者的圣民要得国，永远享有这国，直到万世万代。}\BibleRef{达 7:18}稍后他说：\BibleQuote{祂的王国是永远的王国。}\BibleRef{达 7:27}有关同一主题的其他预言，我在第二十卷中已经提出，还有许多我尚未提出的也记载在同一圣经中；这些预言必将应验，正如那些不信之人曾以为会落空的预言也已应验。因为应许了这两者、并预言两者都将实现的是同一位天主——就是外教神祇在祂面前战栗的那位天主，连外教最杰出的哲学家\Person{波斐利}也为此作证。\par

% structure: p0011 source=h2 target=subsection reason=relative-heading-rank; base=h2

\subsection{第四章　驳斥世上的智者，他们幻想人的尘世身体不能转移到天上的居所。}

然而，那些运用学识和才智来抵抗那伟大权威的力量的人——这权威应验了许久以前所预言的，使所有民族都归向对其许诺的信德和望德——自以为是在巧妙地反对肉身复活，同时他们引用西塞罗在\WorkTitle{论共和国}第三卷中所提及的话。因为当他主张赫拉克勒斯和罗慕路斯的神化时，他说：\Quote{他们的身体并没有被提升到天上；因为大自然不允许一个属土的身体存在于地上以外的任何地方。} 这诚然是智者们深奥的推理，但天主知道他们的思念是虚妄的。因为如果我们只是灵魂，即没有身体的精神，并居住在天上，对地上的动物一无所知，有人告诉我们，我们将通过某种奇妙的结合纽带与地上的身体结合，并赋予它们生命，难道我们不会更加坚决地拒绝相信这一点，并坚持认为大自然不允许一个无形实体被有形束缚所束缚吗？然而大地充满了活的灵，属地的身体与它们相联，并以奇妙的方式与它们交织在一起。既然如此，那位创造这些存有的天主若也愿意如此，有什么能阻止属地的身体被提升为属天的身体呢？既然一个比所有身体都更优越，因此甚至比属天的身体更优越的灵，已经被系于一个属地的身体上？如果如此微小的一粒尘土都能与自身结合一个比属天身体更美好的东西，使其获得感觉和生命，难道天会不屑于接纳，或至少保留这一有感觉、有生命的微粒吗？特别是这微粒的生命和感觉来自一个比任何属天身体都更优越的实体。若此事现在尚未发生，那是因为那位已经行了比这些人不肯相信的事更奇妙得多的天主所定的时刻尚未到来。为什么我们不更强烈地感到惊奇：无形灵魂，其等级高于属天的身体，却被束缚于属地的身体；而不是相反，属地的身体被提升到一个虽然属天却仍是物质的住所？无非是因为我们习惯了看见前者，并且我们自己就是如此，而后者我们尚未见到，也从未见过。诚然，若我们诉诸清醒的理性，就会发现两者中更奇妙的天主工程，是以某种方式将有形之物与无形之物联结，而非将地上之物与天上之物相连，因为它们尽管不同，却同属有形。\par

% structure: p0013 source=h2 target=subsection reason=relative-heading-rank; base=h2

\subsection{第五章 论肉身的复活，有些人拒绝相信，尽管全世界都相信它。}

但即便这曾是不可信的，看，如今，世界已经相信基督的肉身身体被接升天。有学问的和无学问的都已经相信肉身的复活及其升天，只有极少数有学问或无学问的人仍为此动摇。如果所信之事是可信的，那么那些不信的人就应看到自己何等愚钝；但若它是不可信的，那么这也是一件不可信的事：那不可信的事竟然获得了如此的信任。因此我们便有两件不可信之事——即我们身体复活进入永生，以及世界相信如此不可信之事；而这两件不可信之事，同一位天主在它们任何一件发生之前就预言它们必会实现。我们已经看到两件中的一件已经实现，因为世界相信了那不可信的事；为何我们应当绝望于剩下的一件也会实现，即那世界所信的（尽管它不可信）本身会发生呢？因为那同样不可信的事已经实现了：世界相信了一件不可信的事。两者都曾是不可思议的：一件我们看见成就了，另一件我们相信必将成就；因为两者都在那使世界相信的同一部\WorkTitle{圣经}中被预言了。而世界获得信德的方式，若我们细想，甚至更为不可思议。没有受过任何文科教育的人，没有一丝异教学问的修养，不通文法，不精辩证，不饰修辞，而是平凡的渔夫，人数极少——这些人就是基督用信德的网撒向这世界的海，从而从各族中捕得如此之多的鱼，甚至连那些哲学家本身，\Quote{虽稀少却奇妙}。如果你们愿意，或者因为你们理当愿意，让我们在这两件之上加上这第三件不可信之事。现在我们有三件不可信之事，全都已经实现了。不可信：耶稣基督肉身复活并带着肉身升天；不可信：世界竟然相信如此不可信之事；不可信：极少数出身卑微、地位低贱、没有学问的人，竟能如此有效地说服世界，甚至其中博学之士，相信如此不可信之事。在这三件不可信之事中，与我们辩论的那些人不肯相信第一件；他们不能拒绝看见第二件，如果他们不信第三件，他们便无法解释第二件。无疑，基督的复活以及祂带着复活后的肉身升天，已在全世界被宣讲和相信。若它不可信，它如何又已在全世界得到信任？如果许多尊贵、崇高、有学问的人说他们曾目睹此事，并费力传播他们所见的，那么世界相信它就不足为奇，而不肯相信倒是极为固执；但若事实确如那样，世界相信了少数卑微、微不足道、没有学问的人，他们宣称并书写他们目睹了此事，那么一小撮执迷不悟的人反对全世界的信条、拒绝相信，岂非不合理？如果世界相信了少数出身卑微、地位低贱、没有学问的人，那是因为事物本身的神性在如此卑微的见证者身上显得更为明显。确实，使他们的信息具有说服力的口才，是奇妙的作为，而非言辞。因为他们没有亲眼见过基督肉身复活，也没有见过祂带着复活的身体升天，却相信了那些讲述他们如何看见这些事的人，那些人不仅以言语，而且以奇妙的征兆作证。因为他们所认识的这些人原本只懂一种语言，至多两种，却惊奇地听见他们用万国的方言说话。他们看见一个生来瘸腿的人，四十多年后，因这些人以基督之名的一句话而站起来痊愈；从他们身上取下的手帕能治愈病人；无数患有各种疾病的人被一排排放在他们将要经过的路上，以便他们走路时影子能落在病人身上，病人立即得到痊愈；还有许多其他惊人的奇迹以基督的名由他们施行；最后，他们甚至复活死人。如果承认这些事如所记载的那样发生了，那么我们就有了无数不可信之事加到那三件不可信之事上。为使耶稣基督复活升天这一件不可信之事得以被相信，我们积累了无数不可思议奇迹的见证，但即便如此，我们仍未能折服这些怀疑者可怕的顽固。但如果他们不信这些奇迹是由基督的宗徒们为使他们宣讲基督复活升天获得信任而行的，那么这一个伟大的奇迹就足够我们了：全世界没有奇迹却已相信。\par

% structure: p0015 source=h2 target=subsection reason=relative-heading-rank; base=h2

\subsection{第六章——罗马因爱其创立者\Person{罗慕路斯}而尊之为神；教会却因信基督为天主而爱祂。}

让我们在此引述\Person{图利乌斯}表达他对\Person{罗慕路斯}被奉为神竟然得人相信一事感到惊异的那段话。我把他的话照录如下：\Quote{最值得注意的，是关于\Person{罗慕路斯}的一点：其他据说成了神的人，生活在较为蒙昧的时代，那时更倾向于虚构故事，未受教育的人容易相信任何事物。但\Person{罗慕路斯}的时代距今不过六百年，而文学与科学已经驱散了与未开化时代相联的错误。}紧接着，他又说了几句关于同一\Person{罗慕路斯}的话，大意如下：\Quote{由此我们可以看出，\Person{荷马}远在\Person{罗慕路斯}之前就已兴盛，而到那时，个人学问已如此深厚，普遍的启蒙已如此广泛，几乎容不下任何传说。因为古代接纳传说，有时甚至是很拙劣的传说；但\Person{罗慕路斯}的时代已经足够开明，足以拒绝任何缺乏真实感的东西。}因此，一位最有学问、当然也是最雄辩的人，\Person{玛尔库斯·图利乌斯·西塞罗}，说\Person{罗慕路斯}的神性竟得人相信乃是一件惊奇的事，因为时代已经如此开明，不会接受虚构的传说。但除了当时尚且弱小、处于襁褓中的\Place{罗马}本身，还有谁相信\Person{罗慕路斯}是神呢？后来，后代子孙有必要保留祖先的传统；借着这种迷信，如同吮吸母亲的乳汁，国家得以成长并达到如此强大，以致可以从一个有利的位置，将这一信仰强加于其统治之下的所有民族。这些民族虽然可能并不相信\Person{罗慕路斯}是神，但至少如此说，以免因拒绝给予其建立者以\Place{罗马}所授予他的称号而冒犯了这位主权国家——\Place{罗马}之所以采纳这一信仰，并非出于对错误的爱，而是出于爱的错误。然而，\Person{基督}是天上永恒之城的建立者，这城并非因为由祂建立而相信祂是天主，而是相反，因着信仰而被祂建立。\Place{罗马}在建成并献城之后，在神庙中敬拜其建立者为神；而这座\Place{耶路撒冷}则将\Person{基督}——它的天主——作为基石，以便建筑和献城得以进行。前者爱其建立者，因此相信他是神；后者相信\Person{基督}是天主，因此爱祂。前者的爱有一个前因，并且相信它所爱之物具有一种虚假的尊严；同样，后者的信仰也有一个前因，并且爱那由正确的信仰（而非轻率的猜测）归于其对象的真正尊严。因为，且不提证明\Person{基督}是天主的大量惊人奇迹，还有神性的预言为祂作证，这些预言最值得相信，并且已经应验，我们不像教父们那样需要等待其应验。至于\Person{罗慕路斯}，关于他建立\Place{罗马}并在那里统治，我们读到或听到的是实际发生的事件的叙述，而非预先预言说这些事将要发生。就他进入神列而言，历史只记载这一信仰，并未将其陈述为事实；因为没有神迹奇事证实这一点的真实性。至于那头据说哺养了孪生兄弟的狼，它被视为一个伟大的奇事，但这如何证明他具有神性呢？即使这位乳母真的是一头狼，而非仅仅是一个妓女，她也哺养了兄弟二人，而\Person{雷穆斯}并不被算作神。此外，有什么能阻止任何人声称\Person{罗慕路斯}、\Person{赫丘利}或任何一个这样的人是神呢？又有谁宁愿选择死，也不愿承认他的神性呢？难道有哪一个民族将\Person{罗慕路斯}列入其众神之中，除非是迫于对\Place{罗马}之名的恐惧？但是，谁能数得清那些选择最残酷的死亡而不愿否认\Person{基督}神性的群众呢？因此，对\Place{罗马}人心中可能存在的（也许毫无根据的）轻微愤怒的畏惧，迫使一些臣服于\Place{罗马}的国家将\Person{罗慕路斯}当作神来敬拜；然而，对不仅是轻微的精神打击，而是严酷且多样的刑罚，乃至对最可怕的死亡本身的畏惧，却不能阻止全世界无数殉道者不仅敬拜\Person{基督}，而且宣认祂是天主。\Person{基督}之城，虽然在地上仍是异乡客，却有无数公民，它并没有为了暂时的安全而向它那无神的迫害者开战，而是宁愿藉着避免战争来赢得永生的救恩。他们被捆绑、被监禁、被鞭打、被折磨、被焚烧、被肢解、被残杀，然而他们越发增多。他们不能为了永恒的救恩而争斗，除非为了救主的缘故而轻视眼前的救恩。\par

我知道\Person{西塞罗}在他的\WorkTitle{论共和国}第三卷中（如果我没有记错）主张一流政权不会发动战争，除非是为了荣誉或安全。关于安全问题以及他所谓的安全是什么意思，他在另一处解释了，他说：\Quote{私人经常通过迅速死亡来逃避贫困、流放、镣铐、鞭笞以及其他即使最麻木的人也感到的痛苦。但对于国家而言，那似乎使个人免于一切惩罚的死亡本身却是一种惩罚；因为国家应当被构建为永恒的。因此，死亡对于一个共和国来说并不像对于一个人那样是自然的，对一个人说死亡不仅是必要的，而且常常是值得向往的。但是当一个国家被摧毁、抹去、消灭时，就好像（用大事比小事）整个世界都灭亡和崩塌了。}\Person{西塞罗}这样说是因为他和柏拉图主义者一样相信世界不会灭亡。因此，根据\Person{西塞罗}，国家应当为了那种安全而战，这种安全使国家尽管有公民更替，仍能永久存在；正如橄榄树或月桂树或任何这类树木的叶子是常绿的，旧叶被新叶取代。因为，正如他所说，死亡对于个人不是惩罚，反而使他们摆脱所有其他惩罚，但是对于国家却是一种惩罚。因此合理地问，\Person{萨贡托人}是否做得对，他们选择整个国家灭亡而不愿背弃对\Place{罗马}共和国的信义；因为他们的这一行为受到了地上共和国的公民的赞扬。但我看不出他们如何能遵循\Person{西塞罗}的建议，他告诉我们除非为了安全或荣誉决不开战；他也没有说明两者中哪个应该优先，如果出现不能保全一个而不失去另一个的情况。因为很明显，如果\Person{萨贡托人}选择安全，他们必须背弃信义；如果他们持守信义，他们必须放弃安全；结果也确实如此。但是天主之城的安全是这样的：它可以通过信德并伴随着信德而被保存，或者更确切地说获得；但如果信德被放弃，没有人能获得它。正是这种最坚定和最坚忍的精神思想造就了许多高贵的殉道者，而\Person{罗慕路斯}却连一个为他神性而死的人都没有，也不可能有。\par

% structure: p0018 source=h2 target=subsection reason=relative-heading-rank; base=h2

\subsection{第七章——世界对\Person{基督}的信仰是神圣力量的结果，而非人的说服。}

但是将\Person{罗慕路斯}的虚假神性拿来与\Person{基督}相比是极其可笑的。然而，如果\Person{罗慕路斯}大约生活在\Person{西塞罗}之前六百年，在那个已经相当开明、拒绝一切不可能之事的时代，那么在这位\Person{西塞罗}本人以及\Person{奥古斯都}和\Person{提庇留}皇帝的时代——这肯定是一个更加开明的时代（六百年后）——人的心智怎么会拒绝听信或相信\Person{基督}身体的复活和它的升入天堂，并把它当作不可能之事而嗤之以鼻呢？如果不是真理本身的神性（或神性的真理）以及确凿的奇迹证明了这事能够发生并且已经发生？凭着这些见证，尽管有如此多残酷迫害的反对和恐怖，肉身的复活与不朽——首先在\Person{基督}内，随后在新世界中的众人中——被相信，被无畏地宣讲，并播撒到全世界，并用殉道者的血浇灌。因为先知们关于事件的预言被阅读，被强有力的神迹所证实，真理被看作不仅不违反理性，而且只是与习惯观念不同，所以最终世界拥抱了它曾疯狂迫害的信仰。\par

% structure: p0020 source=h2 target=subsection reason=relative-heading-rank; base=h2

\subsection{第八章——论为使世界信\Person{基督}而行的奇迹，以及自世界相信以来这些奇迹未曾止息。}

他们问道：为何如今不再施行你们声称从前所行的那些奇迹呢？我确实可以这样回答：在世人相信之前，奇迹是必要的，为使他们得以相信。而如今，任何要求看到异事才肯相信的人，本身就是一个大异事，因为尽管全世界都已相信，他却不信。但他们提出这些反对意见的唯一目的，是想暗示从前那些奇迹也从未发生过。那么，为何基督的复活与升天被如此坚定地信仰，到处受人称颂？为何在这开明的时代，一切不可能之事都被拒绝，而世界却无需任何奇迹就相信了极为难以置信的事？或者他们会说，这些事本身可信，因而被相信了？那么他们自己为何不信？因此，我们的论证简明扼要：要么是那些未曾目睹的难以置信之事，使世界相信了那些曾经发生并被目睹的难以置信之事；要么是此事本身如此可信，以致无需奇迹来证明，从而证明这些不信者不可原谅的怀疑态度。我这样说，是为驳斥这些最无聊的反对者。但我们不能否认，为了证实基督带着他所复活的肉体升天这一伟大而赐予救恩的奇迹，曾施行了许多奇迹。因为我们这些最可信的典籍，在同一叙述中既包含了所行的奇迹，也包含了奇迹所证实的信经。奇迹被宣扬是为了产生信德，而信德一旦产生，又使奇迹更加显赫。因为它们在会众中被诵读，为的是让人相信，但若非已经相信，它们也不会被这样诵读。因为即使如今，因基督之名，借着祂的圣事或圣人们的祈祷与圣髑，奇迹仍在施行，但它们不如从前那些奇迹那样辉煌显赫，足以被如此荣耀地宣扬。因为神圣的经典正典需要固定下来，使得那些（从前）的奇迹在各地被传诵，并深入所有会众的记忆；而现今的奇迹，即便在其发生的地方，也几乎不为全体居民所知，至多局限于一处。因为它们往往只为极少数人所知，其余人一概不知，尤其若城市很大；而当它们被报告给其他地方的人时，也没有足够的权威使人立即坚定地相信，尽管是信徒向信徒报告的。\par

当我在\Place{米兰}时，那里发生了一个奇迹，一个盲人得以恢复视力；这个奇迹可为许多人知晓，因为不仅城市很大，而且当时皇帝也在那里，一大群人聚集到殉道者\Person{普罗塔西乌斯}和\Person{革尔瓦西乌斯}的遗体处，这些遗体长期隐藏不被知晓，后来在梦中启示给\Person{安波罗修}主教，由他发现。凭借这些遗体的效能，那盲人的黑暗被驱散，他重见天日。\par

然而，除了极少数人之外，有谁知晓因诺森提乌斯（前任副行政长官代理人）所蒙受的治愈之恩呢？这治愈发生在迦太基，就在我眼前，在我的亲眼目睹之下。那时，我和我的兄弟阿利比乌斯尚未担任神职人员，却已是天主的仆人，我们从外地回来，此人接待了我们，并让我们与他同住，因为他全家都虔心敬主。他因患瘘管而接受医生治疗，这些瘘管数量众多，且深藏在直肠中。他已接受过一次手术，外科医生们竭尽全力为他缓解病痛。在那次手术中，他忍受了漫长而剧烈的疼痛；然而，在肠道的众多皱褶中，有一条窦道完全逃过了手术刀，以致他们本应用刀将其切开，却从未触及它。因此，尽管所有已切开的瘘管都已痊愈，这一条却仍保持原状，使他们的努力全部落空。病人因由此造成的延误而生疑，又极度害怕第二次手术——这是他家里的一位医生告知他必须接受的，尽管此人甚至未被允许目睹第一次手术，且被赶出家门，费尽周折才得以回到盛怒的主人面前——病人终于对外科医生们爆发了，说道：\Quote{你们还要再给我动刀吗？你们终究要应验那个你们甚至不许在场之人的预言吗？}外科医生们嘲笑那位不熟练的医生，用好话和承诺安抚病人的恐惧。就这样过了几天，但他们尝试的一切都无济于事。不过，他们仍坚持说他们能用药物治好那条瘘管，无须手术。他们还请来了另一位在该领域享有盛誉的年老医生阿摩尼乌斯（他当时仍在世）；他检查了患处后，也向他们保证，凭借他们的照料和技艺，能取得同样的结果。凭着这位权威，病人变得有信心了，仿佛已然痊愈，便以诙谐的言辞发泄他的好心情，嘲笑那位预言要二次手术的家庭医生。长话短说，在徒劳地度过了若干天后，外科医生们感到厌倦和困惑，最终不得不承认他只能通过手术治愈。他极度恐惧，惊慌失措，脸色煞白；当他镇定下来能够说话时，他命令他们离开，永远不要再回来。他因哭泣而疲惫不堪，迫于无奈，他想到请一位亚历山大里亚的医生，此人当时被认为是一位技艺惊人的手术师，让他来执行手术——他的愤怒不容许那些外科医生做这件事。但当这位医生到来后，用专业的眼光检查了他们精细工作的痕迹，他做了件好事，说服他的病人允许同一双手来完成治愈，因为这双手以令他赞叹的技艺开始了手术，他补充说，毫无疑问，他治愈的唯一希望在于手术，但根据他的本性，他绝不愿意通过完成那一点点剩余的工作而赢得治愈的荣誉，从而剥夺那些人的报酬——他看见他们工作的痕迹时，不得不赞叹他们精湛的技艺、谨慎和勤奋。因此，他们再次获得宠爱；并商定在亚历山大里亚医生面前，他们要对那条瘘管进行手术，所有人都同意，现在只能通过手术刀来治愈。手术推迟到第二天。但他们离开后，家中因主人的极度沮丧而发出一片哀号，那情景在我们看来如同葬礼上的哀悼，我们几乎无法抑制。圣德之人惯常每日探望他：蒙福的撒杜尼诺（当时是乌扎利的主教）、司铎革洛苏斯，以及迦太基教会的执事们；其中还有主教奥勒留，他是他们中唯一尚存的人——我们当以应有的敬意提及此人——我常与他谈及此事，当我们一起谈论天主的奇妙作为时，我发现他清楚地记得我现在所叙述的。当这些人照常于当晚来探望他时，他含着可怜的泪水恳求他们，请他赏脸次日出席他所认为的葬礼，而非他的手术。因为先前的疼痛使他如此恐惧，以至于他毫不怀疑自己会死在外科医生手中。他们安慰他，劝勉他信赖天主，并鼓起勇气像个男子汉。然后我们去祈祷；但当我们像往常一样跪下并俯伏在地时，他俯身倒下，仿佛有人将他猛地摔在地上，开始祈祷；但他是以怎样的方式、怎样的热忱和情感、怎样的泪如泉涌、怎样的呻吟和呜咽（全身颤抖，几乎说不出话）来祈祷，谁能描述呢？其他人是否祈祷，是否没有被他的举动完全吸引注意力，我不知道。至于我自己，我根本不能祈祷。我在心里只简短地说：\Quote{主啊，祢若不垂听这些祈祷，祢还垂听祢子民的什么祈祷呢？}因为在我看来，这祈祷已无以复加，除非他在祈祷中咽气。我们站起身来，接受了主教的祝福，便离开了，病人恳求探望者次日早晨到场，他们则鼓励他保持信心。可怕的日子破晓了。天主的仆人们按所承诺的到场了；外科医生们来了；所需的一切都已就绪；可怕的手术器械拿了出来；所有人在惊奇和悬念中注视。当那些对病人最有影响力的人鼓舞他低落的精神时，他的四肢被安放在长椅上，以适合手术师的操作；绷带的结被解开；患处暴露；外科医生检查它，手执手术刀，急切地寻找要切开的窦道。他用眼睛搜寻，用手指触摸，用各种方法检查：他找到了一个完全坚固的疤痕！没有任何言语能描述从所有人口中涌出的喜乐、赞美和对慈悲全能天主的感恩，伴随着喜悦的泪水。与其描述，不如想象那场景！\par

在同城\Place{迦太基}住着一位极为虔敬、地位最高的贵妇\Person{依诺森提雅}。她一侧乳房患有癌症，照医师说，这是一种不治之症。因此，他们通常要么切除患病的肢体，使之与身体分离；要么放弃一切治疗，以便病人的生命能稍延长——虽然死亡即使推迟也是不可避免的——正如他们所说，遵循\Person{希波克拉底}的忠告。我们所谈的这位妇人曾得到一位熟悉她家事的熟练医师的同样建议，但她惟独藉祈祷投靠了天主。临近复活节时，她在梦中得到指示：等候第一位从洗礼池出来、受过洗的妇女，请她在自己的患处划\OriginalTerm{基督的标记}{the sign of Christ}。她照做了，立刻痊愈。那位曾劝告她若想多活些时日就不要治疗的医师，事后检查她时，发现她——先前他检查时还受此病折磨——现已完全痊愈，便急切地问她用了什么疗法，渴望（我们完全可以相信）找到一种能推翻\Person{希波克拉底}论断的药物。但当她告诉他发生了什么事时，据说他带着宗教性的礼貌，却以轻蔑的语气和表情回答说，这使她担心他会说出亵渎\Person{基督}的话：\Quote{我还以为你会告诉我什么重大发现呢。}她对他漠不关心的态度感到颤抖，立刻答道：\Quote{\Person{基督}治好癌症算什么大事，他不是曾叫一个死了四天的人复活了吗？}因此，当我听到这事，我极其愤慨：如此伟大的奇迹发生在那个著名的城市，而且发生在一位绝非无名的人身上，竟然不被传扬。我认为应当为此事对她说几句话，若非严责的话。而当她回答我说她并未对此事缄默时，我问那些与她最亲近的妇女们是否以前听说过这事。她们告诉我说一无所知。\Quote{你看，}我说，\Quote{你所谓的不缄默究竟是怎么回事，既然连那些与你如此亲近的人都不知道。}由于我只是简短地听了这故事，我便让她从头到尾讲述整个事情如何发生，其他妇女极其惊讶地听着，并光荣了天主。\par

同城的一位患痛风的医师，当他报名领受圣洗，而在领受圣洗的前一天，有一些黑发卷毛的男孩出现在他梦中，他理解为魔鬼，禁止他当年受洗；虽然他们践踏他的双脚，使他忍受了前所未有的剧痛，但他拒绝服从他们，反而战胜了他们，不愿推迟在\OriginalTerm{重生的洗涤}{laver of regeneration}中受洗；就在领受圣洗的行动中，他不仅解除了折磨他的异常疼痛，也解除了疾病本身，以致此后他活了很久，从未再患痛风；然而有谁知道这个奇迹？不过，我们知道它，邻近的一些兄弟，那些可能听到此事的少数人，也知道。\par

一位来自\Place{库鲁比斯}的老年喜剧演员在圣洗时不仅瘫痪被治愈，疝气也被治愈，而且，既已从两种痛苦中解脱，他从\OriginalTerm{重生的泉源}{font of regeneration}中上来，仿佛身体从未有任何不适。在\Place{库鲁比斯}之外，谁知道这事？或在别处听到的极少数人之外，还有谁知道？但我们听到此事后，便奉圣主教\Person{奥勒留}之命，使那人来到\Place{迦太基}，尽管我们已经根据一些我们不可怀疑其言辞的人的报告，确认了事实。\par

\Person{赫斯佩里乌斯}，出身护民官家族，是我们的一位邻居，在\Place{富萨利亚}地区有一个名为\Place{祖贝迪}的农场；他发觉他的家人、牲畜和仆役都因邪魔的恶毒而受折磨，便在我缺席期间，请求我们的司铎中有一人随他同去，藉祈祷驱赶那些邪魔。有一人去了，在那里献上了\OriginalTerm{基督身体的祭献}{the sacrifice of the body of Christ}，竭尽全力祈祷，使那搅扰得以止息。因天主的仁慈，搅扰立刻止息了。赫斯佩里乌斯曾从一位朋友那里得到一些从\Place{耶路撒冷}带来的圣土，那里是\Person{基督}被埋葬后第三日复活之处。他将这土挂在卧室中，以保护自己免受伤害。但当他的房屋被清除那魔鬼的入侵后，他开始思考这土应如何处理；因为他对此土的敬意使他不愿再将它留在卧室。恰巧我和\Place{西尼塔}的主教\Person{马克西米努斯}（当时是我的同僚）正在附近。赫斯佩里乌斯请我们拜访他，我们便去了。当他讲述了全部情况后，他恳求将这土埋在某处，并将那地点用作祈祷之所，以供基督徒聚集朝拜天主。我们没有反对：照他所愿而行。在那附近有一个患瘫痪的年轻乡下人，他听说此事后，恳求父母毫不延迟地带他去那圣地。当他被带到那里时，他祈祷，立刻完全痊愈地用自己的双脚离开了。\par

距\Place{希波勒吉乌斯}不到三十英里，有一处名为\Place{维多利亚纳}的庄园。那里有一座纪念\Place{米兰}殉道者\Person{普罗塔西乌斯}与\Person{革尔瓦西乌斯}的圣墓。某夏日正午，一青年在河中水潭饮马时，被魔鬼附身，被人抬到圣墓旁。他躺在那里，奄奄一息，几乎如同死人。庄园的女主人照常带领女仆和虔敬的侍从前来参加晚间祈祷和赞颂，他们开始唱圣诗。听到这声音，那青年仿佛触电般彻底惊醒，发出可怕的尖叫，抓住祭台，紧紧不放，仿佛不敢或不能松手，如同被钉在或绑在上面；附在他身上的魔鬼大声哀号，求饶命，并供认何时、何地、如何附上这青年。最后，魔鬼说要离开他，便逐一说出他身体各部分，扬言离开时要将其伤残；说完这话，他就离开了那人。但青年的眼睛掉出，挂在脸颊上，仅靠一根细血管如根般连着，整个原本黑色的瞳孔变成了白色。在场的人见此情景（另有一些人听到喊声也赶来，一同为他祈祷），虽然为他恢复神志而高兴，另一方面却为他的眼睛担忧，说应当求医。但带他来的妹夫说：\Quote{天主既然驱走了魔鬼，也必能因圣人们的祈祷复原他的眼睛。}于是他将掉出悬挂的眼睛放回原位，用手帕尽其所能地绑好，并嘱咐他七天之内不可解开绷带。他照做后，发现眼睛完全健康。那里也有其他人得到治愈，但说来话长，兹不赘述。\par

我知道\Place{希波}的一位少妇，用掺有曾为她祈祷的司铎眼泪的油傅抹自己后，立即被驱逐了魔鬼。我也知道一位主教曾为一个他从未见过的附魔青年祈祷，那青年当场痊愈。\par

我们在\Place{希波}有一位同乡，名叫\Person{弗洛伦提乌斯}，是一位虔敬的穷老人，靠裁缝为生。他丢失了外套，无钱再买，便向在本城拥有著名纪念圣堂的\Person{二十位殉道者}祈祷，大声恳求能得着衣服。几个嘲弄人的青年恰好听见，便在他离去时一路讥讽他，好像他竟向殉道者讨五十文钱买外套似的。但他默默前行，在海边看见一条大鱼，像是刚被冲上岸，喘息未定；他得到那些青年好心帮助，捉住了鱼，卖给一位名叫\Person{卡托苏斯}的厨师——一位好基督徒，告诉他得鱼的经过，得了三百文钱。他用这些钱买了羊毛，让妻子施展手艺，为他缝制外套。厨师剖鱼时，在鱼腹中发现了一枚金戒指；他立刻动了怜悯之心，也出于宗教的敬畏，将戒指还给了那人，说：\Quote{看，二十位殉道者如何给你穿了衣服。}\par

当\Person{普罗耶克图斯}主教将最光荣的殉道者\Person{斯德望}的圣髑送往\Place{提比利斯}的泉水时，大批群众在圣堂迎接他。那里有一位盲妇求人领她到携带圣髑的主教跟前。主教将手中拿着的花给了她。她接过花，敷在眼上，立刻看见了。在场的人惊愕不已，而她满心欢喜地走在前面，不再需要向导。\par

\Place{西尼塔}的主教\Person{路济路斯}（该地靠近殖民城镇\Place{希波}）曾捧着同一殉道者的圣髑列队游行，这些圣髑原供奉在\Place{西尼塔}的城堡中。他长期患有瘘管病，私人医生正伺机切开，而仅仅捧着那神圣的包裹，瘘管突然痊愈了——至少此后他身上再未留下任何痕迹。\par

一位住在\Place{卡拉马}的西班牙司铎\Person{欧加里乌斯}长期患结石症。\Person{波西迪乌斯}主教带来的同一殉道者的圣髑治愈了他。后来，同一司铎因另一种疾病倒下，已经死去，人们正在捆绑他的双手。同一殉道者出手相助：人们从小圣堂取来司铎的外衣，盖在尸体上，他便复活了。\par

那里有一位名叫\Person{马尔提亚}的老年贵族，非常厌恶基督信仰，但他的女儿是基督徒，女婿也在同年受了洗。他病重时，女儿和女婿含泪恳求他成为基督徒，他坚决拒绝，怒气冲冲地把他们赶走。女婿想到去\Person{斯德望}的小圣堂，在那里恳切祈祷，求天主赐予他岳父正确的心志，使他不再迟延信仰基督。他怀着极大的呻吟和眼泪，以及诚挚虔敬的热忱这样做了；离开时，他拿了那里的一些花，当时已是夜晚，将花放在岳父的头旁，让老人入睡。看，天还未亮，老人就喊人快去找主教；但那时主教恰好在\Place{希波}与我在一起。他听说主教不在家，便请司铎们来。司铎们来了。在众人的喜悦和惊讶中，他宣称自己相信了，并受了洗。此后他在世的日子里，口中常挂着一句话：\Quote{基督，收我的灵魂！}他并不知道这正是最蒙福的\Person{斯德望}被犹太人用石头砸死时所说的最后一句话\BibleRef{宗 7:59}。这也是他自己的临终遗言，因为不久之后他也离世了。\par

同样，在同一\OriginalTerm{殉道者}{martyr}之处，有两人——一位市民，一位外乡人——治愈了痛风；市民完全痊愈，而外乡人仅被告知当疼痛复发时应涂抹何种药物；他遵从此建议，疼痛顿时缓解。\par

奥杜鲁斯是一处庄园之名，那里有一座教堂，内设殉道者斯德望的纪念祠。有一日，一个小男孩在庭院中玩耍，拉车的牛偏离道路，车轮碾过他，他立时气息奄奄。母亲抱起他，放在祠前，他不仅复活，而且毫无损伤。\par

一位虔诚的妇女，住在邻近的庄园卡斯帕利乌姆，病重至绝望，人们将她的衣服带到此祠前，但衣服尚未带回，她已去世。然而，她的父母用那衣服包裹她的尸体，她便恢复了呼吸，完全痊愈。\par

在希波，一位名叫巴苏斯的叙利亚人正在同一殉道者的\OriginalTerm{圣髑}{relics}前为病危的女儿祈祷。他也将女儿的衣服带到祠前。祈祷时，他的仆人们从家中跑来告知他女儿已死。但他的朋友们拦住他们，禁止他们告知，以免他在众人前哀哭。他回家后，家中已哭声震天，他将所带的衣服盖在女儿身上，她便复活了。\par

同样，在那里，我们的一位税吏依勒乃的儿子患病而亡。尸体已无气息，人们正准备后事，众人哭泣哀悼时，一位安慰父亲的朋友建议用同一殉道者的油敷抹尸体。如此行，他便复活了。\par

同样，我们中一位护民官阶层的厄肋乌西努斯，将已死的幼子放在殉道者的祠前（此祠在他所住的郊区），他含泪祈祷后，抱起孩子，孩子已复活。\par

我该当如何？我受完成此著作的诺言所迫，无法记录我所知的一切奇迹；无疑，我的一些同道读到我所叙述的内容，会遗憾我遗漏了许多他们与我确知的事。此刻我恳请这些人原谅我，并请他们思量，若我叙述全部奇迹——我因完成手头著作之需而不得不省略——将耗费多少时日。假若我只字不提其他，只记录通过此殉道者（即最光荣的斯德望）在卡拉马和希波地区所行的治愈奇迹，也足以写满数卷；然而，即使这些奇迹也无法全部收录，仅能收录那些已写成公开宣读的叙述。因为当我亲眼目睹在我们时代，屡次出现类似古代的神圣大能征兆时，我渴望有人写成记述，认为民众不应对这些事一无所知。这批圣髑首次被迎至希波勒吉乌斯尚不满两年，尽管据我所知，许多借之而行的奇迹未被记录，但在撰写此刻，已公布的奇迹几近七十宗。但在卡拉马，圣髑已有更长时间，且为公众信息而宣读的奇迹更多，其数量远胜于此。\par

同样，在靠近乌提卡的殖民城乌札利，据我所知，同一殉道者也行了许多显著奇迹，其圣髑早在希波获得之前，便已由埃沃迪乌斯主教安排安放于此。但那里并无公开叙述的习惯，或说，过去没有——或许如今已开始。因我近日在那里时，一位贵妇佩特洛尼娅，久患重病，医药罔效，奇迹般痊愈；经该地上述主教同意，我鼓励她公开记述此事，以供民众宣读。她欣然从命，并在其记述中插入一事，我虽急于推进本工作所需的主题，却不可不提。她说，曾有一犹太人说服她，在贴身衣物下束一毛发腰带，带上戴一戒指，戒上非宝石，而是一枚牛肾中取出的石头。她佩此符咒，前往圣殉道者的门槛。但离开迦太基后，在巴格拉达河畔的自家庄园歇宿，正要起身继续行程时，她看见戒指掉在脚前。她大为惊讶，检查毛发腰带，见其仍如先前般牢牢系住，便推测戒指已磨损脱落；但发现戒指本身完好无损，她猜想，藉此大奇迹，她已以某种方式获得了痊愈的保证；于是她解下腰带连同戒指，一同投入河中。此奇事不为那些不信主耶稣基督从母腹中生出而未损其童贞、且关着门进入门徒中的人所信；但可让他们仔细查考此奇迹，若证实为真，便当相信其他奇迹。这位女士出身高贵，嫁与贵族，居于迦太基。城邑显赫，人物显赫，查询者必能满意。无疑，那殉道者本人，因她因之祈祷而痊愈，相信了那位童贞之母之子；相信了那位关着门进入门徒中者；总之——我们所述这一切皆指向——相信了那位带着复活后的肉身升天者；正因他为这信仰舍命，如此奇迹才藉他而施行。\par

因此，即使在现今，仍有许多奇迹发生，那位曾施行我们所读到之奇迹的天主，仍按祂的意愿，借着祂所选择的人继续施行奇迹；但这些奇迹并不为人所知，也没有因频繁的诵读而被像沙砾一样锤入记忆，以至于无法从脑海中消失。因为即使在我们中间，现今也注意让那些获得恩惠者的小册子被公开宣读，但那些在场的人只听一次叙述，而许多人不在场；因此，即使那些在场的人，几天后也会忘记他们所听到的，几乎找不到一个人会将他所听到的告诉他知道不在场的人。\par

在我们中间曾发生一个奇迹，虽然并不比我所提及的那些更大，却是如此显著和引人注目，以至于我认为希波没有一个居民没有看见或听说过它，也没有一个人可能忘记它。卡帕多细亚凯撒利亚的一个贵族家庭有七个兄弟和三个姐妹，他们的母亲是新寡，因他们对她所做的一些错事而怨恨他们并咒骂他们，于是他们受到了上天如此严厉的惩罚，以致他们所有人都四肢剧烈颤抖。由于带着这种可憎的模样无法忍受同胞的目光，他们几乎走遍了整个罗马世界，各人按自己的方向行走。其中两人来到希波，是一对兄妹，保禄和巴拉第亚，他们悲惨的遭遇在许多地方已经出名。他们是在复活节前大约十五天来到的，每天到教堂来，特别是到最光荣的斯德望的圣髑前，祈求天主现在息怒，恢复他们昔日的健康。在那里，以及他们所到之处，他们都引起了每个人的注意。有些曾在别处见过他们并知道他们颤抖原因的人，一有机会就告诉别人。复活节到了，在主日的早晨，当时已有大批人群在场，那位青年正扶着圣髑所在圣所的铁栏祈祷，突然他倒下了，躺卧在那里，完全像睡着了一样，不再像他平时甚至在睡眠中也颤抖。所有在场的人都惊讶不已。有些人惊慌，有些人感到怜悯；有人想扶他起来，另一些人阻止他们，说应该等待看看结果如何。看啊！他站了起来，不再颤抖，因为他已痊愈，完全健康地站着，打量着那些打量他的人。那时谁能克制自己而不赞美天主呢？整个教堂充满了那些欢呼和祝贺他的声音。然后他们跑来找我，当时我正坐着准备进入教堂。他们一个接一个地涌进来，最后来的人像新闻一样告诉我第一个已经告诉我的事；我心中欢喜，暗中感谢天主时，那青年本人也与许多人一起进来，跪在我的膝前，我扶他起来接受我的亲吻。我们进去参加集会：教堂满了，回荡着喜悦的呼声：\Quote{感谢天主！赞美天主！}每个人都一同从四面八方呼喊：\Quote{我治愈了百姓}，然后更加响亮地呼喊。终于安静下来后，按惯例诵读圣经的章节。轮到我讲道时，我作了一些适合当时欢乐愉快气氛的简短评论，不愿他们听我，而宁愿他们考虑天主在这神圣工程中的口才。那人同我们一起用餐，并仔细告诉我们他自己、他母亲和他家庭的灾祸。因此，第二天，在我讲道之后，我许诺次日将他的叙述读给百姓听。当我这样做时（复活节后的第三天），我叫那兄妹俩都站在我通常讲话的高台的台阶上；当他们站在那里时，他们的叙述被宣读。全体会众，无论男女，都看见那一个站着毫无异常动作，另一个四肢颤抖；这样，那些先前没有见过那男子本人的人，在他妹妹身上看到了天主慈悲从他身上所移除的东西。在他身上他们看到了值得庆贺的事，在她身上则看到了祈求的对象。与此同时，他们的叙述读完，我吩咐他们退离众人的视线；我开始更仔细地讨论整个事件，看啊！当我正说下去时，从殉道者的墓穴那里传来新的祝贺声。我的听众转过身，开始向墓穴跑去。那年轻女子从她站立的台阶下来，到圣髑前祈祷，她一触到铁栏，就像她哥哥一样，像睡着般倒下，然后起来痊愈了。当我们正在问发生了什么事，以及什么引起了这喜悦的喧嚷时，他们进入我们所在的巴西利卡，从殉道者的墓穴领着她，已完全健康。那时，男女一齐发出如此惊讶的呼喊，以至于那欢呼和眼泪似乎永远不会停止。她被领到她刚才颤抖着站立的地方。他们现在欢喜她像她哥哥一样，先前他们曾为她依然不像他而悲伤；他们还没有为她发出祈祷，就觉察到他们这样做的意向已被迅速垂听。他们无言地赞美天主，但声音如此响亮，我们的耳朵几乎承受不住。在这些欢欣鼓舞的人心中，除了基督的信德（斯德望曾为此流血）之外，还有什么呢？\par

% structure: p0045 source=h2 target=subsection reason=relative-heading-rank; base=h2

\subsection{第九章　论借着殉道者因基督之名所行的一切奇迹，都证明殉道者对基督所有的信德。}

这些奇迹为甚么目的作证，岂不正是为这信德作证，这信德宣讲基督肉身复活，并带着同一肉身升天？因为殉道者本身正是殉道者，即这信德的见证人，他们藉着自己的见证吸引了世界的仇恨，并征服了世界，不是透过抵抗，而是透过死亡。他们为这信德而死，如今能以他们被杀害的主的名义祈求这些恩惠。正是为了这信德，他们惊人的坚毅得到了操练，以致在这些奇迹中彰显了伟大的能力。因为倘若肉身的复活——那进入永生的复活——没有在基督身上发生，也不会如基督或预言基督来临的先知所预言的那样，在祂的子民身上实现，那么，那些为这宣讲复活的信德而被杀的殉道者，为何会有如此能力？无论天主自己是藉着那种奇妙的工作方式施行这些奇迹，即祂虽是永恒的，却在时间中产生效果；或是祂藉着仆人施行，若是如此，祂是否使用殉道者的灵魂，如同使用仍在肉身中的人，或者藉着天使成就这一切奇事，祂对天使行使着无形、不变、无体的统治，以致所谓由殉道者所做之事，并非出于他们的操作，而仅是由于他们的祈祷和请求；或是最终，有些事以一种方式成就，另一些以另一种方式，以至于人完全无法理解它们——无论如何，这些奇迹证实了这宣讲肉身复活、进入永生的信德。\par

% structure: p0047 source=h2 target=subsection reason=relative-heading-rank; base=h2

\subsection{第十章——论殉道者获得许多奇迹，使真天主受到敬拜，远较那些行奇事以图自称为神的魔鬼更值得尊敬。}

在此，我们的对手或许会说，他们的神也行过一些奇妙的事，如果他们现在开始将他们的神与我们的死人相比。或者，他们还会说，他们拥有一些从死人中选出的神，如\Person{赫丘利}、\Person{罗慕路斯}以及许多其他他们幻想被纳入诸神之列的人。然而，我们的殉道者并非我们的神；因为我们知道，殉道者与我们同有一位天主，而且是同一位。况且，他们声称藉着他们的神庙所行的奇迹，也根本无法与藉着我们殉道者的坟墓所行的奇迹相比。即便看似相似，他们的神已被我们的殉道者击败，如同\Person{法郎}的术士被\Person{梅瑟}击败一样。实际上，魔鬼行这些奇事，是出于那不洁的骄傲，企图成为万国的神；而殉道者行这些奇事，或者更确切地说，天主在她们祈祷和协助下行这些奇事，是为了激励信德，使我们相信他们并非我们的神，而是与我们同有一位天主。总之，他们为他们的神建造了神庙，设立了祭台，任命了司铎，规定了祭献；但对于我们的殉道者，我们建造的并非神庙，好像他们是神，而是纪念碑，如同为死人建造的，他们的灵魂与天主同在。我们也不在这些纪念碑前设立祭台，以向殉道者献祭，而是向殉道者与我们同有的唯一天主献祭；在这祭献中，他们因其本身的位置和等级被提及，作为因承认天主而征服世界的人，但他们并不被献祭的司铎所呼求。因为司铎是向天主献祭，而非向他们，尽管他在他们的纪念碑前献祭；因为他是天主的司铎，而非他们的司铎。再者，这祭献本身是基督的身体，这身体并非向他们献上，因为他们本身就是这身体。那么，谁更可能被认为是行奇迹的呢？是那些希望因所行奇迹而被人们视为神的人，还是那些行任何奇迹的唯一目的是促使人们相信天主，并相信基督也是天主的人？是那些甚至想将自己的罪行转为神圣礼仪的人，还是那些甚至不愿自己的赞美被圣化，而寻求将一切他们理应受赞美的事归于他们所赞美的那位的荣耀的人？因为他们的灵魂是在主内受到赞美。因此，让我们相信那些既说真话又行奇事的人。因为他们因说真话而受苦，从而赢得了行奇事的能力。他们宣称的首要真理是：基督从死者中复活，并首先在自己的肉身中显示了复活的永生不朽，这复活祂应许将赐予我们，或是在来世之初，或是在今世终结之时。\par

% structure: p0049 source=h2 target=subsection reason=relative-heading-rank; base=h2

\subsection{第十一章——驳斥柏拉图派，他们从元素的重力性质论证尘世身体不能居住在天堂。}

但是对于天主的这一伟大恩赐，这些推理者——\Quote{他们的思念原是虚幻的，主知道}——从元素的重量中提出论证；因为他们从他们的老师\Person{柏拉图}那里学到，世界的两大元素，相距最远的，由中间的两者——空气和水——连结并统一。因此他们说，既然地是元素系列中从底部开始的第一，水是第二，在土之上，空气第三，在水之上，天第四，在气之上，那么一个地上的身体就不能住在天上；因为每个元素由其自身的重量保持平衡，以维持其自己的位置和等级。看，人的软弱，为虚妄所占据，竟以这样的论证反驳天主的全能！那么，这么多地上的身体在空气中做什么呢？既然空气是从地起的第三元素？除非那位赐给鸟类的尘世身体羽毛和翅膀的轻便以在空中飞翔的，就不能赐给人那成为不朽的身体以停留在至高天上的能力？同样，那些不能飞的陆生动物，其中包括人，按此理论应该住在地下，就像水中的动物鱼住水下一样。为什么一个陆生动物不能生活在第二元素，即水中，却能生活在第三元素中？为什么它虽属于地，若被迫生活在紧邻地的第二元素中就立刻窒息，却能生活在第三元素中，且不能离开它而活？是元素次序有误，还是错误在于他们的推理，而不在于事物的本性？我不再重复我在第十三卷中所说的：许多地上的身体，虽沉重如铅，却因工匠之手获得一种形式而使它们能在水中浮游；然而竟有人否认那全能的工匠能赐给人一个性质，使其能升入天上并住在那里。\par

但是针对我先前所说的，他们无言以对，尽管他们提出并充分利用了他们所信赖的这种元素次序。因为若次序是：地第一，水第二，气第三，天第四，那么灵魂就在一切之上。因为\Person{亚里士多德}说灵魂是第五个身体，而\Person{柏拉图}则根本否认它是一个身体。若它是一个第五身体，那么它当然高于其他；若它根本不是身体，则它更超越一切。那么，它在一个地上的身体里做什么呢？这比一切都精细的灵魂，在如此一大团物质中做什么呢？最轻的实体在这种沉重中做什么呢？这最快捷的实体在如此迟钝中做什么呢？难道身体不会因如此卓越的本性而被提升到天上吗？若如今地上的身体能将灵魂拘于下界，难道灵魂终有一天不能将地上的身体提升到高处吗？\par

如果我们现在来看他们所反对我们殉道者的、由他们的神所行的奇迹，难道这些不也为我们作证，并有助于我们的论证吗？因为若他们的神的任何奇迹是伟大的，那么\Person{瓦罗}所提到的那个灶神贞女的事迹当然是伟大的：她受到不贞洁的虚假指控的威胁，用\Place{台伯河}的水装满了一个筛子，端到审判官面前，筛子没有漏出一点水。谁使水的重量保持在筛子里？谁阻止任何一滴水从那么多小孔中漏出？他们会回答：某个神或某个魔鬼。若是一个神，他比那创造世界的天主更大吗？若是一个魔鬼，他比那事奉创造世界之天主的天使更强吗？那么，若一个较小的神、天使或魔鬼都能如此支撑这液体元素的重量，使水似乎改变了它的本性，难道全能的天主，祂亲自创造了所有元素，就不能从地上的身体中除去它的沉重，使那复活的身体能住在赋予生命的灵所愿的任何元素中吗？\par

再者，既然他们给空气以火在上水在下的中间位置，为何我们经常发现它位于水与水之间，以及水与地之间？他们如何解释那些水汽的云彩？在它们和海洋之间，空气总是被夹在中间。我想知道，根据元素的什么重量和次序，非常猛烈而暴风雨般的洪流在顺着空气下方的地面奔涌之前，被悬在空中的云彩里？最后，为何在整个地球上，空气位于最高天与地之间，如果它的位置是在天与水之间，正如水的位置是在天与地之间？\par

最后，如果元素的秩序如此安排，按照\Person{柏拉图}的看法，两个极端——火与地——由空气和水这中间物联结，火占据天空的最高处，地占据最低处，如同世界的基础；因此，地不能在天空，那么火如何在地内？因为按照这个推理，这两个元素——地与火——应当被限制在它们自己的位置，最高与最低，使得最低者不能升至最高者的位置，最高者也不能降至最低者的位置。因此，正如他们认为没有一粒尘土现在或将来会在天空中，我们也不应在世上看见一粒火。但事实是，火不仅在世上，甚至在地上也存在到如此程度，以致山顶喷发它们；此外，我们看见它存在于世上供人使用，甚至从大地生出，因为它由木柴和石头点燃，而这些无疑都是地上的物体。但那种上界的火，他们说，是宁静、纯净、无害、永恒的；而这种地上的火则是混浊、多烟、可朽坏且使人朽坏的。然而，它并未朽坏那不断在其中肆虐的山脉和洞穴。但即使地上的火如此不同于那种火，以适应其地上的位置，那么他们为什么反对我们相信，某一天地上身体的本质将变为不朽坏并适合天空，正如现在火是腐朽的并适合大地？因此，他们从元素的重量和秩序中，无法证明全能的天主不能使我们的身体如此，以至于它们可以在天空中居住。\par

% structure: p0055 source=h2 target=subsection reason=relative-heading-rank; base=h2

\subsection{第十二章——驳不信者以嘲笑加于基督信仰关于肉身复活的毁谤。}

但他们惯于假装在探究此问题时心怀谨慎的焦虑，并嘲笑我们对身体复活的信仰，问道：流产的胎儿会复活吗？而正如主所说：\BibleQuote{我实在告诉你们：你们连一根头发也不会失落。}\BibleRef{路 21:18}那么，所有的身体都会有相等的身材和力量吗？还是会有大小之别？因为如果有平等，那些流产的胎儿，假若他们复活，从何处获得他们在此世没有的体型？或者，如果他们不复活，因为未出生而被抛弃，他们也会提出同样的问题：那些幼年夭折的孩子，我们从何处获得他们在此世没有的身材？因为我们不会说那些不仅出生了，而且重生了的，不会复活。然后，他们进一步问，这些相等的身体会是多大。因为如果所有人都如世上最高大者一样高大，他们问我们，如何不仅孩童，而且许多成年人，将要得到他们在此世所没有的，如果每个人都要得到他在这里所拥有的？而如果宗徒的话——我们都达到\BibleQuote{基督圆满年龄的尺度}\BibleRef{弗 4:13}，或另一句话——\BibleQuote{使与祂圣子的肖像相似}\BibleRef{罗 8:29}——被理解为意味着基督身体的身材和大小将是所有在祂王国中之人身体的尺度，那么，他们说，许多人的大小和高度必须缩小；如果身体本身有如此多部分丧失，那么\BibleQuote{你们连一根头发也不会失落？}这句话又如何成立？此外，关于头发本身，他们可能会问：理发师剃掉的所有头发都会恢复吗？如果恢复，谁不因这样的畸形而畏缩？因为对指甲剪掉的同样恢复，许多为了外表美观而剪掉的部分会重新放在身体上。那么，它的美丽何在？在那种不朽的状态中，美丽无疑应当比在这个腐坏的状态中更大。另一方面，如果这些东西不恢复，它们就必须消失；那么，他们怎么说不使一根头发失落？同样，他们也推论关于肥胖和消瘦；因为如果所有人都相等，那么当然不会有胖有瘦。因此，有些人会有所得，有些人会有所失。因此，不仅不会有对从前存在之物的简单恢复，反而一方面会添加从未存在之物，另一方面会失去先前存在之物。\par

关于死体的腐败和分解——有的化为尘土，有的化为空气；有的被野兽吞噬，有的被火焚烧，有的因船难或形形色色的溺水而消亡，以致身体化为液体——这些困难使他们极度惊慌，他们相信所有那些分解的元素不能重新聚集并重建为身体。他们还热切地利用所有因意外或出生而产生的畸形和瑕疵，于是，怀着恐惧和嘲笑，引用怪胎，并问在复活中每一种畸形是否都会被保存。因为如果我们说在人的身体中不会再现这样的东西，他们就以为可以通过引用我们在主基督复活的身体上发现的伤痕来驳倒我们。但在所有这些问题中，最难的是：那被另一个因饥饿所迫而吃下并同化的人肉，将归于谁的身体？因为它已转化为食用者的肉，并填补了饥荒造成的肉体损失。因此，为了嘲笑复活，他们问：这肉是归还给它最初所属的人，还是后来成为其肉的人？如此，他们也试图为人的灵魂提供真实的痛苦和虚假的幸福之交替，按照\Person{柏拉图}的理论；或者按照\Person{波斐利}的理论，即在多次转移到不同身体之后，灵魂结束其痛苦，永不再返回，然而不是通过获得不朽的身体，而是通过逃离每一种身体。\par

% structure: p0058 source=h2 target=subsection reason=relative-heading-rank; base=h2

\subsection{第13章——流产儿若算在死者中，是否也将分享复活？}

现在，我要回复以上详述的敌人们的异议，信赖天主慈悲助佑我的努力。至于流产儿——即便假定他们在母腹中活着，也死在那里——是否会复活，我既不敢肯定也不敢否认，尽管我看不出，如果他们不排除在死者之外，为何不能达到死人的复活。因为要么所有死者都不复活，那么将有一些灵魂永远没有身体——尽管他们曾在母腹中拥有过身体；要么，如果所有人类灵魂都将重新获得他们曾在任何地方生活时所拥有并离世时留下的身体，那么我看不出如何能说，连那些死在母腹中的人也得不到复活。但无论谁对这些流产儿采取哪种观点，如果他们复活，我们至少必须将我们对已出生婴儿所说的一切应用于他们。\par

% structure: p0060 source=h2 target=subsection reason=relative-heading-rank; base=h2

\subsection{第14章——婴儿是否将以他们长大后会有的身体复活？}

那么，我们该如何论及婴儿呢？难道不是：他们将不会以死时微小的身体复活，而是要借由天主神奇而迅速的作为，领受时间本来缓慢进程会赐予他们的身体？因为主的话中，祂说：\Quote{连你们的一根头发也不会失落，}这话肯定所拥有的都不会缺少；但并没有说未曾拥有的就不会赐予。对于死去的婴儿，他们缺少身体的完美身材；因为即使是完美的婴儿也缺乏体型上的完美，尚有生长潜力。这种完美身材，在某种意义上，所有人都在受孕和出生时拥有——即他们潜在地拥有它，虽然尚未在实际上拥有；正如身体的所有肢体潜在地存在于种子中，即使孩子出生后，某些肢体（例如牙齿）可能仍未长出。在每一样体的这种种子原理中，似乎有一切尚未存在，或更确切地说，尚未显现，但随着时间的推移将要存在或显明的事物之开端。因此，在这个原理中，将来高或矮的孩子已经高或矮了。同样，在身体的复活中，我们无需害怕任何身体的损失；因为即使所有人都应具有相同大小，达到巨人般的尺寸，恐怕今生最大的人会失去他们的体积而使其消失，与基督的话相悖（祂说他们的头发也不会失落），然而，那位奇妙的工者（祂是创造者，曾从无中创造了万物）为何会缺少增添身体的手段呢？\par

% structure: p0062 source=h2 target=subsection reason=relative-heading-rank; base=h2

\subsection{第15章——是否所有死者的身体都要以主身体的同一尺寸复活？}

确然，基督复活时带着祂死时同样体型的身躯，若说当普通复活来临时，祂的身体为了与最高者平等，会采取祂在门徒面前以他们熟悉的身形显现时所没有的比例，这是错误的。但如果我们说，甚至高大者的身体也要缩减至主身体的尺寸，那么许多身体将遭受巨大的损失，尽管祂曾应许连他们的一根头发也不会失落。因此，剩下的结论是：每个人将得到他在青年时期曾拥有的尺寸（尽管他是老人而死），或者如果他在壮年前去世，他将拥有的尺寸。至于宗徒所说的\Quote{基督圆满年龄的尺度}，我们要么理解他指别的事物，即：当基督教团体中的所有肢体都添加到元首时，基督的尺度将完成；要么如果我们将它应用于身体的复活，那么意思是所有人都将复活在既不超越也不低于青年，而是在我们所知基督已达到的活力与年龄中。因为甚至世间最智慧的人也把青春年华定在三十岁左右；一旦过了这个时期，人就开始衰退，走向有缺陷、更迟钝的老年时期。因此宗徒所说的不是身体的尺度，也不是身量的尺度，而是\Quote{基督圆满年龄的尺度}。\par

% structure: p0064 source=h2 target=subsection reason=relative-heading-rank; base=h2

\subsection{第16章——圣人符合圣子肖像意指什么？}

再者，这些话说：\Quote{预定符合圣子的肖像，}\BibleRef{罗 8:29}可以理解为就内心的人而言。同样，在另一处祂对我们说：\Quote{不要符合这世界，但要在心意的革新中变形。}\BibleRef{罗 12:2}那么，只要我们被变形，从而不符合世界，我们就符合圣子。它也可以如此理解：正如祂藉着取死亡本性而符合我们，我们将藉着不死而符合祂；这确实与身体的复活有关。但如果这些话也教导我们身体以何种形式复活——正如我们之前谈到的尺度——那么这种符合也不是指尺寸，而是指年龄。因此，所有人都将以他们已达到的，或者若活到壮年将会达到的身量复活，尽管即使身体形态是婴儿或老年，也没有大碍，因为任何衰弱都不会留在头脑或身体本身中。所以即使有人主张每个人都将以他死时的同一身体形态复活，我们无需花太多工夫与他争论。\par

% structure: p0066 source=h2 target=subsection reason=relative-heading-rank; base=h2

\subsection{第17章——妇女的身体在复活中是否将保留其性别？}

从\BibleQuote{直到我们众人都达到成全的人，达到基督圆满年龄的尺度}\BibleRef{弗 4:13}这句话，以及\BibleQuote{肖似天主子肖像}\BibleRef{罗 8:29}这句话，有些人就得出结论说，女人复活后将不再是女人，而是所有的人都将成为男人，因为天主只用地上的尘土造了男人，而女人是从男人出来的。就我而言，那些毫不怀疑两性都将复活的人，似乎更为明智。因为那里不再有情欲，这情欲是如今混乱的原因。因为在犯罪之前，男人和女人都是赤身露体的，并不羞耻。因此，从那些身体中，恶习将被除去，而本性将被保存。女人的性别不是恶习，而是本性。那时它确实将超越肉体交合和生育；然而女性的器官仍将保留，不是为旧日的用途，而是为一种新的美丽，这美丽非但不会激起如今已灭除的情欲，反而会激发对天主的智慧和仁慈的赞美，祂造了原本不存在的东西，并拯救祂所造之物脱离败坏。因为在人类之初，女人是从男人睡眠时从他身侧取出的肋骨造成的；因为那时看来，连在这一事件中预表基督和祂的教会也是合适的。因为男人的睡眠预表基督的死亡，当祂毫无生气地悬在十字架上时，祂的肋旁被长矛刺透，流出了血和水，我们知道这些就是圣事，教会就是藉此\Quote{建造}起来的。因为圣经使用了这个词，没有说\Quote{形成了}或\Quote{造了}，而是说\BibleQuote{将她建造成一个女人}\BibleRef{创 2:22}；因此宗徒也谈到基督身体的\Emph{建立}\BibleRef{弗 4:12}，这身体就是教会。所以，女人和男人一样都是天主的受造物；但藉着由男人受造这一点，合一受到推崇；而她受造的方式，如前所述，预表了基督和教会。因此，造了两性的那一位，也将恢复两性。耶稣自己也曾回答否认复活的撒杜塞人，他们问：若按照法律，兄弟相继娶妻为兄立嗣，那复活后这妇人应是七兄弟中哪一个的妻子？祂说：\BibleQuote{你们错了，既不明白圣经，也不明了天主的能力。}\BibleRef{玛 22:29}虽然祂本可以说你们所问的这位妇人，她自己将变成男人，不再是女人，但祂并没有这样说；反而说：\BibleQuote{在复活时，他们不娶也不嫁，而是像天上的天使一样。}\BibleRef{玛 22:30}他们在不死不灭和幸福方面将与天使同等，但不是肉身方面，也不是复活方面，因为天使不需要复活，他们不会死。因此，主否认了复活时会有婚姻，但没有否认女人；祂在此时如此否认，是因为如果祂真的预知女性不存在，那么提出的问题本可以更轻易、更迅速地通过否认女性存在来解决。但事实上，祂甚至肯定了女性将存在，因为祂说：\BibleQuote{她们不嫁}，这句话只适用于女性；\BibleQuote{他们不娶}，这句话适用于男性。因此，在这世上习惯娶妻和嫁人的那些人，在复活时不会再行这样的婚姻。\par

% structure: p0068 source=h2 target=subsection reason=relative-heading-rank; base=h2

\subsection{第十八章　论成全的人，即基督；以及论祂的身体，即教会，也就是祂的圆满。}

为了理解宗徒所说\Quote{我们众人都要到达一个完人}的意思，我们必须考虑整段经文的联系，经文如下：\BibleQuote{那下降的，也正是那上升到诸天之上的，为要充满万有。祂赐予一些人作宗徒，一些人作先知，一些人作传福音者，一些人作牧者和教师，为成全圣人，为服务的工作，为建立基督的身体，直到我们众人达到信仰的一致和对天主子的认识，达到一个完人，达到基督圆满年龄的尺度，使我们不再作小孩子，被各种教义之风吹来吹去，被人蒙骗的伎俩和狡猾的诡计所摇动，他们设伏欺骗；而要在爱德中持守真理，在各方面长进到祂里面，祂是元首，就是基督；整个身体靠祂紧密配合，藉着每个关节的供应，按每部分的有效运作，使身体增长，在爱德中建立自己。}\BibleRef{弗 4:10} 看，那完人是什么——是元首和身体，这身体由所有肢体组成，这些肢体将在适当的时候被成全。但每日都有新的成员加入这身体，就是教会在被建立时，对这身体说：\BibleQuote{你们是基督的身体和祂的肢体；}\BibleRef{格前 12:27} 又说：\BibleQuote{为了祂的身体，}他说，\BibleQuote{就是教会；}\BibleRef{哥 1:24} 又说：\BibleQuote{我们虽是众多，却是一个身体。}\BibleRef{格前 10:17} 正是论到这身体的建立，这里也说：\BibleQuote{为成全圣人，为服务的工作，为建立基督的身体；}然后加上了我们现在讨论的这段：\BibleQuote{直到我们众人达到信仰的一致和对天主子的认识，达到一个完人，达到基督圆满年龄的尺度，}等等。他又指出我们应当理解这尺度是针对哪个身体，当他说：\BibleQuote{使我们在各方面长进到祂里面，祂是元首，就是基督：整个身体靠祂紧密配合，藉着每个关节的供应，按每部分的有效运作。} 因此，正如每部分都有其尺度，同样，由所有部分组成的整个身体的圆满也有其尺度，正是论到这一尺度（\OriginalTerm{基督圆满年龄的尺度}{the measure of the age of the fullness of Christ}）说：\BibleQuote{达到基督圆满年龄的尺度。} 他在另一处也论到这种圆满，论到基督说：\BibleQuote{将祂赐给教会作万有之上的元首，教会是祂的身体，是那充满万有者所充满的。}\BibleRef{弗 1:22} 但即使这应当被理解为每个人复活时的形态，有什么能阻止我们将那明确论到男人的话应用于女人，将两性都包含在通用的术语\Quote{人}之下呢？因为确实，在\Quote{敬畏上主的人是有福的}这句话中，也包含了敬畏上主的妇女。\par

% structure: p0070 source=h2 target=subsection reason=relative-heading-rank; base=h2

\subsection{第十九章——所有在今世损毁人体优美的身体瑕疵，在复活时都将被除去，身体的本体物质保持不变，但质与量有所改变，以产生优美。}

我现在该怎么论头发和指甲呢？一旦明白了身体的任何部分都不会消失以致在身体中造成畸形，那么同时也明白了，那些由于过度比例会造成畸形的部分，将被加入到身体的总体量中，而不是加入到那些会破坏优美比例的部分。正如用粘土制作一个器皿后，想要用同一块粘土重新制作它，不必使原来构成把手的粘土再次构成新把手，也不必使原来构成底部的粘土再次构成底部，而只需整块粘土构成整个新器皿，且没有任何部分被浪费。因此，如果修剪的头发和剪下的指甲若回到原位会造成畸形，它们就不会复原；然而在复活时没有人会失去这些部分，因为它们将被改变为同样的血肉，其材质被如此改变以保持身体各部的比例。然而，主所说的：\BibleQuote{你们连一根头发也不会失落，}或许更适宜被理解为指头发的数量，而不是长度，正如祂在别处所说：\BibleQuote{你们的头发都已被数过了。}\BibleRef{路 12:7} 我这样说并不是因为我以为身体自然所属的任何部分会消失，而是说任何畸形——它原本服务于我们必死之人所处的受罚状态——将被如此恢复：在完全保存材质的同时，畸形将消失。因为如果连一个手工艺人，由于某种原因制作了一个畸形的雕像，也能重新铸造并使之非常优美，且丝毫不损失材质，只让畸形消失——例如，他可以移除某个不雅或不成比例的部分，不是通过切除和分离这部分，而是通过打破和混合整个材料来消除缺陷而不减少材料的量——我们难道不更应如此看待全能的工匠吗？祂难道不能移除并消除人体的一切畸形，无论是常见的还是罕见的怪胎？这些畸形虽然与今世悲惨的生活相称，但在未来的福乐中却不应被提及；祂难道不能如此移除它们，使自然但不雅观的缺陷终结，而自然材质却不受丝毫减损？\par

所以，过于肥胖或过于消瘦的人不必担心：他们在天堂中会长成他们即使在今世也不愿有的形状。因为一切身体之美都在于各部分的匀称，加上某种悦目的颜色。没有匀称，眼睛就会感到不适，或是因为有所欠缺、过于细小，或是因为过于粗大。因此，在那一切错误都被纠正、一切缺陷都由造物主所知的资源补足、一切多余都被去除而不损本质完整的状态中，不会有因比例失调而产生的畸形。至于那悦目的颜色，当\BibleQuote{义人在他们父的国里要发光如同太阳！}\BibleRef{玛 13:43}时，该是何等显著！我们宁可相信，这光辉在基督复活时并未缺失，而是对门徒的眼睛隐藏了。因为软弱的人眼无法承受，而他们必须看见祂，以便认出祂来。为此，祂也容许他们触摸祂伤口的痕迹，并且进食——不是因为祂需要营养，而是因为祂若愿意就能吃。现在，当一个物体虽然在场，却对那些看见其他在场之物的人不可见时——正如我们说那光辉在场，却被那些看见其他事物的人所不见——这在希腊文中称为\OriginalTerm{不可见}{\GreekText{ἀορασία}}；我们的拉丁文译者因缺乏更好的词，在\BibleBook{}{创世纪}中将其译为\Emph{\OriginalTerm{盲目}{cæcitas}}。\Place{索多玛}人寻找义人\Person{罗特}的门却找不着时，就遭受了这种盲目。但如果那是盲目，即他们什么也看不见，那么他们就不会问进入那屋子的门，而会求人领他们离开了。\par

但我们对于真福殉道者的爱，不知何故，使我们在天国渴望看见他们为基督之名所受的伤痕，也许我们真会看见。因为这不是畸形，而是光荣的标记，将为他们的容貌增添光辉，一种属灵的美，即使不是身体的美。然而，我们不必相信那些曾被告知\BibleQuote{你们连一根头发也不会失落}的人，在复活时会缺少他们在殉道中所失去的肢体。但如果在那新的国度里，某些伤痕的标记仍然可见于那不朽的肉身是合宜的，那么他们受伤或被截肢的地方将保留伤疤，而不会失去任何肢体。因此，虽然身体所受的任何瑕疵都不会在复活中出现，但我们不应将这些美德的标记视为瑕疵。\par

% structure: p0074 source=h2 target=subsection reason=relative-heading-rank; base=h2

\subsection{第二十章  论复活时我们的身体即使分解，其本质也将完全重新结合。}

我们切莫害怕，造物主的全能无法为了我们身体的复苏和重生而召回那些已被野兽或火吞噬，或已化为尘土、灰烬，或已分解为水，或蒸发到空气中的所有部分。我们切莫认为，在自然最隐秘的角落中任何逃过我们观察的事物会逃避万有造物主的认知或超越其能力。我们对手中的权威\Person{西塞罗}，在尽可能精确地定义天主时说：\Quote{天主是自由而独立的心智，没有物质性，感知并推动万物，本身具有永恒的运动。}这是他从最伟大的哲学家中找到的。那么，我以他们自己的语言请问：有什么东西能够对那感知万物的祂隐藏，或者能从那推动万物的祂手中不可挽回地逃脱呢？\par

这就引我去回答那个似乎最难的问题：在复活时，死者已变成活人之肉的肉将属于谁？假如有人因匮乏而饥饿，被迫以人肉为食——这种极端情况并非不为人知，正如古代史和我们这个时代的不幸经历所教导的那样——能否合理地争辩说，所有被吃的肉都已排出，没有任何部分被同化到食者的本质中？尽管先前存在的消瘦现已消失，这充分表明多少巨大的匮乏已被这食物填补？但我已作过一些评论，足以解决这个难题。因为所有被饥饿消耗的肉都通过蒸发进入了空气，从那里，正如我们所说，全能的天主能够召回它。因此，那肉应归还给最初使它成为人肉的那个人。因为它必须被视为被另一个人所借用，如同金钱贷款，必须归还给贷方。然而，他因饥荒失去的自己的肉，将由那位甚至能恢复蒸发之物者归还给他。即使它被完全消灭，以至于其本质的任何部分都没有留在自然的任何隐秘角落，全能者也能用祂认为合适的方式恢复它。因为真理所说的话：\BibleQuote{你们连一根头发也不会失落}，禁止我们以为，虽然人头上的一根头发都不能失落，但被饥饿者吃掉和消耗的大块肉却可以失落。\par

从我们如此思考并用我们所能支配的如此贫乏的能力所讨论的一切，我们得出这个结论：在肉身的复活中，身体将具有它在青春的花季中已经达到或应当达到的尺寸，并将享受因保持所有肢体的匀称与比例而产生的美丽。并且，有理由假定，为保持这种美丽，身体物质的任何部分，如果放在一个地方会产生畸形，都将分布到全身，这样既不会失去任何部分，也不会失去整体的匀称，而只是身体的整体高度因分布了那些在一个地方会难看的部分而略有增加。或者，如果有人主张每个人都将以他死时身体的同样身材复活，我们不会固执地争论这一点，只要没有畸形、没有软弱、没有倦怠、没有腐朽——任何不适合那个国度的事物都没有，在那个国度里，复活和恩许的子女将与天主的众天使同等，即使身体和年龄上不同，至少在幸福上相同。\par

% structure: p0078 source=h2 target=subsection reason=relative-heading-rank; base=h2

\subsection{第二十一章——论圣人肉身将要转化成的新的属神身体。}

因此，凡是从身体取走的，无论是在生前还是死后，都将归还给它，并与坟墓中残留的结合起来，再次复活，从属生灵身体的旧形态转变为属神身体的新形态，并穿上不朽坏与永生。即使身体因某种严重事故或敌人的残忍而被完全磨成粉末，并且被精心撒向风中或水中，以致毫无踪迹留下，它也不会超出造物主的全能——不，连一根头发也不会损失。肉身那时将是属神的，隶属于神魂，但仍是肉身，不是神魂，正如神魂自身，当隶属于肉身时，曾是属肉体的，但仍是神魂，不是肉身。关于这一点，我们在我们受罚状况的畸形中有经验的证明。因为那些人是属肉体的，不是以肉体方式，而是以属神方式，宗徒曾对他们说：\BibleQuote{我无法对你们讲话，如同对属神的人，而是如同对属肉体的人。} \BibleRef{格前 3:1} 一个人在今世如此属神，以至于他对自己的身体仍是属肉体的，并且看到在自己的肢体中有另一条法律，与他的心神的法律交战；但甚至在他的身体里，他也会成为属神的，当同样的肉身获得了那种复活，关于它这些话说：\BibleQuote{播种的是属生灵的身体，复活起来的是属神的身体。} \BibleRef{格前 15:44} 但这个属神的身体是什么，它的恩宠有多大，我担心鲁莽断言，因为我们尚未经验它。然而，既然我们希望的喜乐应当表达出来，从而彰显天主的赞美，并且既然圣咏作者出于最深刻的热烈而神圣的爱而呼喊：\BibleQuote{上主，我喜爱你住所的华美。} 我们可以，藉天主的帮助，谈论祂在此最悲惨的生命中倾注于善恶之人的恩赐，并尽力揣测那状态的伟大荣耀，我们因尚未经验而无法相称地述说。因为我不谈及天主造人正直的时候；我不谈及\Quote{男人和他的妻子}在富饶乐园中的幸福生活，因为它如此短暂，以致他们的子女中无人经历过它：我只谈我们所知的这个生命，我们现在所处的这个生命，我们在此生命中，无论取得多大进步，都无法逃避诱惑，因为它全是考验；我问，谁能描述甚至在此生命中也延伸到人类身上的天主仁慈的标记呢？\par

% structure: p0080 source=h2 target=subsection reason=relative-heading-rank; base=h2

\subsection{第二十二章——论人类因原罪而正义地遭受的苦难与灾祸，除非藉基督的恩宠无人能得救。}

整个人类在其最初根源上已被定罪，这生命本身（如果还能称为生命的话）通过所充满的无数残酷灾祸来作证。这不正是由那深重可怕的愚昧所证明的吗？这愚昧产生了一切包围亚当子孙的错误，除非藉劳苦、痛苦和恐惧，无人能从中得救。这不正是由他对那么多虚妄有害事物的爱所证明的吗？这爱产生折磨人的忧虑、不安、悲伤、恐惧、狂喜、争吵、诉讼、战争、背叛、愤怒、仇恨、欺骗、谄媚、欺诈、偷盗、抢劫、背信、骄傲、野心、嫉妒、谋杀、弑亲、残忍、凶猛、邪恶、奢侈、傲慢、无耻、不知羞耻、奸淫、通奸、乱伦，以及无数不洁和两性间违反本性的行为，这些行为甚至提及都是可耻的；还有亵渎圣物、异端、亵渎、伪证、欺压无辜、诽谤、阴谋、谎言、假见证、不义的审判、暴力行为、掠夺，以及任何类似的邪恶，它们进入了人类的生活，尽管不能进入纯洁心灵的概念？这些确实是恶人的罪行，然而它们出自那错误和错爱之根，这根在每个亚当的儿子出生时即与他同在。因为谁没有观察到，人带着何其深重的无知（甚至在婴儿期就表现出来）以及何其泛滥的愚蠢欲望（在童年开始显现）进入今世生命，以至于若任其随心所欲，为所欲为，他必将陷入我所提及的（以及未能提及的）所有罪过或不义之中？\par

但是，由于天主并未完全抛弃祂所定罪的人，也未曾恼怒地封闭祂的怜悯，人类受法律和教育的约束，这些法律和教育防范困扰我们的无知，抵挡恶习的袭击，但本身却充满劳苦和忧伤。那些用来遏制孩童愚昧的种种威胁意味着什么？学监、教师、桦条、皮带、藤鞭，以及圣经所说必须给予孩童的管教，\BibleQuote{击打他的腰，免得他变得倔强}，\BibleRef{德 30:12}，否则几乎不可能或根本不可能驯服他，这一切意味着什么？所有这些惩罚，若不是为了克服无知和约束邪恶的欲望——这些我们与生俱来的邪恶——那又是为什么？为什么我们记得困难，忘记却容易？学习困难，不学习却容易？勤勉困难，懒惰却容易？这岂不表明败坏的本性以其自身重量趋向和倾向何方，以及若想得救需要何种帮助？怠惰、散漫、懒惰、疏忽，都是逃避劳苦的恶习，因为劳苦虽有益处，本身却是一种惩罚。\par

然而，除了孩童时期的惩罚——没有这些惩罚，孩子们就无法学到父母所愿之事（而父母很少愿教授有用之事）——谁能描述、谁能想象折磨人类的惩罚之繁多与严酷？这些痛苦不仅是无神之徒邪恶的伴随，更是人类处境和共同苦难的一部分：丧亲与哀悼、损失与定罪、欺诈与谎言、虚假猜疑，以及他人所有的罪行与恶行，引发了何等的恐惧与悲伤？因为从他们手中，我们遭受抢劫、俘虏、锁链、监禁、流放、酷刑、截肢、失明、贞洁受辱以满足压迫者的淫欲，以及其他许多可怕的灾祸。无数意外从外部威胁我们的身体：酷热与严寒、风暴、洪水、泛滥、闪电、雷霆、冰雹、地震、房屋倒塌；或因马匹的失足、惊跳或恶癖；或因水果、水、空气、动物中无数的毒物；或因野兽疼痛甚至致命的咬伤；或因疯狗传播的狂犬病，以至在所有动物中对主人最温顺友好的动物，竟变得比狮子或龙更令人恐惧；而那人若偶然被这瘟疫传染，变得如此疯狂，以致他的父母、妻子、儿女都怕他胜过任何野兽！旅行者在陆地和海上遭受何等灾难！有谁出门时能不遭遇到处无法预见的意外？四体健全地回家，却在自家门阶上滑倒，摔断腿，永不康复。还有什么比人坐在椅子上更安全的呢？司祭厄里从椅子上跌倒，摔断了脖子。农夫们——或者不如说所有人——多么担心庄稼受到天气、土壤或破坏性动物的侵害？通常，当庄稼收割入库时，他们感到安全。然而，据我所知，突发的洪水驱散了劳工，将最好的收成从谷仓中席卷一空。天真无邪能充分抵御魔鬼的各种攻击吗？为使无人这样想，就连领受圣洗的婴儿——他们无疑在纯真上无可超越——有时也受如此折磨，以致允许此事的天主藉此教导我们哀叹此生的灾祸，并渴慕来世的幸福。至于身体疾病，它们如此繁多，连医书也未能尽录。而在许多或几乎所有的疾病中，治疗和药物本身就是折磨，以致人通过痛苦的疗法摆脱致命的痛苦。难道干渴的疯狂没有驱使人们喝人尿，甚至自己的尿？难道饥饿没有驱使人们吃人肉，且不是尸体的肉，而是特意杀害的人肉？难道饥荒的剧痛没有驱使母亲吃掉自己的孩子——这似乎残酷得难以置信？最后，睡眠本身，虽恰被称为休息，但有时被梦境和幻象扰乱，其中何曾有安宁？可怜的心灵被那些呈现出来、仿佛如此明显地站在感官面前以致我们无法将它们与现实区分开来的事物形象，是何等恐惧地压垮！在某些疾病中，虚假幻象如何悲惨地分散人的注意力？即使是健康的人，有时也被恶神以惊人多样的幻象所欺骗，这些恶神制造这些错觉以迷惑受害者的感官，如果它们不能成功地引诱他们到自己的阵营！\par

从这地上的地狱中，除非藉着救主\Person{耶稣}基督、我们的天主和主的恩宠，是无法逃脱的。\Person{耶稣}这名就表明这一点，因为它意思是救主；祂特别拯救我们脱离今世生命，使我们不至于进入一个更悲惨和永恒的状态，那状态与其说是生命，不如说是死亡。因为在此生中，虽然圣人圣洁的追求给予我们极大的安慰，但人类所渴望的祝福并不一定都赐给他们，免得宗教因这些暂时的利益而被培养，而它本应为了那排除一切邪恶的另一生命而被培养。因此，恩宠也帮助善人在当前的灾难中，使他们能够以与信德相称的恒心忍受它们。世上的哲人声称哲学对此有所贡献——这种哲学，根据\Person{西塞罗}，神祇只将它纯粹地赐给了少数人。他说，神祇从未给过，也永远不会给人类更大的恩赐。所以，即使那些与我们争论的人也不得不以某种方式承认，天主的恩宠对于获得——不是任何哲学，而是真正的哲学——是必要的。如果真正的哲学——这对此生苦难的唯一支持——是由上天只赐给了少数人，那就足以表明人类已被判定要承受这悲惨的惩罚。而且，根据他们的承认，天主从未赐下更大的恩赐，那么必须相信，它只能由那个他们自己承认比他们所敬拜的所有神祇都大的天主所赐。\par

% structure: p0085 source=h2 target=subsection reason=relative-heading-rank; base=h2

\subsection{第二十三章 论今世苦难中特别属于善人劳苦的那些，而不论那些善恶所共有的。}

但是，撇开在此生中善恶所共有的苦难不谈，义人还承受着他们特有的劳苦，因为他们与自己的恶习争战，并卷入这种争战的诱惑和危险之中。因为虽然有时更猛烈，有时更松弛，但肉情与心神相争，心神与肉情相争，永不间断，以致我们不能做我们想做的事，\BibleRef{迦 5:17}并且不能根除一切私欲，只能拒绝同意它，正如天主赐给我们能力，如此克制它，警惕地守望，免得似是而非的真理欺骗我们，免得微妙的言论蒙蔽我们，免得错误使我们陷入黑暗，免得我们把善当作恶或把恶当作善，免得恐惧阻碍我们做该做的事，或欲望驱使我们做不该做的事，免得太阳在我们含怒时落下，免得仇恨激起我们以恶报恶，免得不当或过度的悲伤消耗我们，免得忘恩负义的心使我们迟钝于认识所领受的恩惠，免得诽谤困扰我们的良心，免得我们轻率的猜疑欺骗我们对待朋友，或别人对我们的虚假猜疑使我们过于不安，免得罪在我们必死的身体中作王，顺从它的私欲，免得我们的肢体被用作不义的器具，免得眼睛追随私欲，免得报复的渴望带走我们，免得视线或思想久久停留在某些使我们愉悦的邪恶之事上，免得我们乐意听邪恶或下流的话，免得我们做愉快但非法的事，并且在这场充满如此多劳苦和危险的战争中，我们或希望靠自己的力量获得胜利，或将在获得胜利时归功于自己的力量，而不是归功于祂的恩宠，宗徒说：\BibleQuote{感谢天主，祂藉我们的主\Person{耶稣基督}赐给我们胜利。} \BibleRef{格前 15:57}在另一处他说：\BibleQuote{靠着爱我们的主，我们在这一切事上得胜有余。} \BibleRef{罗 8:37}然而我们要知道，无论我们多么英勇地抵抗我们的恶习，多么成功地克服它们，只要我们还在这身体内，我们总有理说：\BibleQuote{求祢宽恕我们的罪债。} \BibleRef{玛 6:12}但在那我们将永远居住、披上不朽身体的国度里，我们不再有争战或罪债——事实上，如果我们的本性保持受造时的正直，我们任何时候、任何情况下都不会有。因此，我们这场冲突——在其中我们暴露于危险，并希望藉着最终的胜利得释放——本身就是今生的苦难之一，而这生命被如此多严重邪恶的见证证明是受咒诅的生命。\par

% structure: p0087 source=h2 target=subsection reason=relative-heading-rank; base=h2

\subsection{第二十四章 论造物主在这虽受咒诅却仍充满的今世生命中所赐的祝福。}

但我们现在必须默想天主的良善——祂眷顾自己所造的一切——如何以丰富而数不胜数的恩赐充满人类这厄运本身，而此厄运正反映祂报应的公义。在堕落之前祂曾说的第一个祝福，即祂说：\BibleQuote{你们要滋生繁殖，充满大地} \BibleRef{创 1:28}，在人犯罪后祂并未禁止，而原初赐予的生殖力仍存留在被定罪的族类中；罪的恶习虽使我们不得不面对死亡，却并未剥夺我们那奇妙的种子能力，或更确切地说，那更为奇妙的能力，即产生种子的能力，它仿佛深植并交织于人体中。但是在这条我称之为人类的河流或激流中，两种元素一同被携带——既包括来自生养者的恶，也包括来自创造我们的天主所赐的善。在原初的恶中有两样：罪与惩罚；在原初的善中也有两样：\OriginalTerm{繁衍}{propagation}与\OriginalTerm{成形}{conformation}。关于恶，其中一种是罪，来自我们的妄为；另一种是惩罚，来自天主的审判，我们已经按照当前目的讲了足够多。我现在要谈论的是天主赐给或仍在赐给我们那被败坏又被定罪的本性的恩赐。因为在定罪它时，祂并未收回所有赐予它的，否则它早已消灭；祂在惩罚中将之置于魔鬼的权下，也并未使之脱离祂自己的权能；因为连魔鬼本身也不在天主的治理之外，因为魔鬼的本性只由那至高的创造者维系，祂给予一切以任何形式存在的事物以存在。\par

那么，关于这两种恩赐——我们说过它们如泉源从天主的美善流向我们那被罪败坏并被定罪为惩罚的本性——其中一种，\OriginalTerm{繁衍}{propagation}，是天主在最初的工作中——祂在第七天从这些工作中安息——借着祂的祝福赐予的。但另一种，\OriginalTerm{成形}{conformation}，则是在祂\BibleQuote{一直工作至今} \BibleRef{若 5:17}的那项工作中赐予的。因为如果祂把有效的力量从万物中收回，它们就无法继续并完成其有节律的运动所指定的时期，甚至无法保存它们被创造时所拥有的本性。所以，天主创造人时，赐予了我们所谓的生殖力，藉此他可以繁衍其他的人，赋予他们一种与生俱来的繁衍同类的才能，但并没有强加任何必然性。天主随意从个人夺走这种才能，使之不生育；但从整个族类，祂并没有收回那一次赐予的繁衍祝福。但这繁衍的能力虽因罪未被收回，却已不是若无罪时所当有的样子。因为\Quote{人虽然居尊位，却堕落，变得像畜生一样}，并且像它们一样生殖，尽管那一点理性之火——即天主在他身上的肖像——尚未完全熄灭。但若在繁衍之外没有成形，就不会有同类的再生产。因为即使没有交合之事，而天主愿意用人类居民充满大地，祂可以创造所有这些，如同创造一个人一样，无需借助人类的生育。而且，事实上，即使现在，那些交合者除非藉着天主的创造大能，什么也不能产生。所以，如同在属灵成长方面——人由此被塑造为虔敬和正义——宗徒说：\BibleQuote{可见，栽种的算不得什么，浇灌的算不得什么，只在那使它生长的天主} \BibleRef{格前 3:7}。同样必须说，生育的算不得什么，只在那赋予本质形式的天主；怀孕和哺乳的算不得什么，只在那使它生长的天主。因为唯独祂，借着那\BibleQuote{一直工作至今}的大能，使种子发育，并从某些隐秘不可见的褶皱中展开为我们所看见的可视的美形。唯独祂，以某种奇妙的方式将精神和物质的本性连接起来，一个命令，一个服从，使之成为一个活物。祂的这工作如此伟大奇妙，以至于不仅是人——理性的动物，因而比地上所有其他动物更优越——甚至是最微小的昆虫，都不能不带着惊叹和对造物主的赞颂而被仔细端详。\par

是祂赐予人的灵魂一个理智，在幼年时，理性与理解力仿佛沉睡其中，如同不存在一般；然而，随着年岁的增长，它们必将被唤醒并受到锻炼，从而能够获取知识与接受教导，适宜于领悟真理并爱慕美善。藉此能力，灵魂汲取智慧，并获赋予那些德行，藉着智德、勇德、节德与义德，与谬误及其他与生俱来的恶习作战，并通过将渴望只专注于至高不变之善而战胜它们。即或结果并不总是如此，然而，谁能恰如其分地言说，甚至谁能设想，全能者的这项伟大工程，以及祂赐予我们理性本性的不可言喻的恩惠——仅仅是赋予我们达致如此成就的能力？因为，除了那些被称为德行、教导我们如何善度此生、并获致无尽幸福的技艺——这些技艺是藉着唯独在基督内的天主恩宠，赐予恩许之子嗣与国度的——人类的才智岂非已经发明并应用了无数令人惊叹的技艺，部分出于需要，部分出于丰沛的创造？这心智的活力，不仅在发现多余之物，甚至是在发现危险与毁灭性事物时都如此活跃，不正表明那能发明、学习或运用此类技艺的本性具有无穷的财富吗？人类在纺织、建筑、农业与航海技艺上取得了何等奇妙——几乎令人难以置信——的进步！陶艺、绘画与雕塑的设计何其无穷无尽，制作又何其巧妙！剧场中展示的奇观何等令人惊叹，未亲眼目睹者简直无法置信！捕捉、杀害或驯服野兽的巧计何其精妙！为伤害他人，又发明了多少种类的毒药、武器、毁灭器械，而为保存或恢复健康，所应用的器具和药物更是无穷无尽！为刺激食欲、取悦味觉，调制了多少种佐料！为表达并传达思想，有多少种多样的符号，其中言语与书写居于首位！雄辩拥有何等华丽的辞藻来愉悦心灵！歌曲何等丰富以吸引耳朵！发明了多少乐器与和谐曲调！在度量与数字方面获得了何等的技巧！以何等的睿智发现了星辰的运动与联系！谁能尽述耗费在自然之上的思虑，即使他放弃详述，仅仅试图给出一个概观？最后，就连谬误与误解的辩护——这彰显了异端者与哲学家的才智——也难以充分言说。因为目前我们是在赞颂装点此生的人性心智，而非那导向永生的信德与真理之路。既然这伟大的本性确实由真实而至高的天主所创造，祂以绝对的权能与公义治理祂所创造的一切，若非那首先之人——其余众人皆由他而生——犯下了极其重大的罪过，这本性绝不会陷入这些悲惨境况，也不会由此进入永久的悲惨——唯有那些获得救赎者除外。\par

此外，即使在这身体中，尽管它像野兽的身体一样会死，且在许多方面比它们更软弱，天主的慈善、伟大造物主的眷顾何等显明！感官器官和其他肢体，它们的位置、整体的形象、形状和身材，岂不正是被塑造得表明它是为服务理性灵魂而造的吗？人没有被造得向地面俯伏，像无理性的动物；但他的身体形态，直立且仰望上天，提醒他思念上面的事。再者，赐予舌头和双手的奇妙敏捷，使它们适于说话、书写、执行如此多的职分、从事如此多的技艺，岂不证明那灵魂的高贵，为此就提供了这样的助手吗？甚至撇开它对其所需工作的适应不谈，其各部分之间存在着如此匀称，保持如此美丽的比例，以致人难以决定，在创造身体时，是更注重实用还是美丽。确实，没有哪一部分是为了实用而造，却不对其美丽也有所贡献。如果我们更精确地知道所有部分是如何相互连接和适配的，并且不限于观察表面显现的东西，这一点就会更加明显；至于被遮盖隐藏在我们视线之外的东西，静脉和神经的复杂网络，皮肤下所有东西的生机部分，无人能发现它们。因为，尽管有些被称为解剖学家的医生，带着残酷的科学热忱，解剖过死者的尸体，有时甚至是被他们刀下致死的病人，并且不人道地探查人体的秘密，以了解疾病的性质及其确切位置，以及如何治愈，然而我所说的这些关系，它们形成了整个身体内外的一致，或如希腊人所称的\OriginalTerm{和谐}{harmony}，如同某种乐器一样，至今无人能够发现，因为无人足够大胆去追寻它们。但如果这些能够被知晓，那么即使是那些似乎没有美丽的内在部分，也会以其精妙的适宜性如此愉悦我们，给心灵带来比满足眼睛的显见美丽更深沉的满足——眼睛不过是心灵的仆役。还有一些东西在身体中占据这样的位置，它们显然不服务于任何有用目的，而纯粹为了美丽，\Emph{例如}男子胸部的乳头，或他脸上的胡须；因为这是为了装饰而非保护，这一点由女人的光滑面孔所证明，而女人作为较弱的性别，反倒更应享有这样的防御。因此，如果在所有暴露于我们视野的肢体中，确实没有一个部分是为了实用而牺牲美丽，而有一些部分除了美丽之外别无用处，那么我认为可以容易地得出结论：在创造人体时，悦目比实用更受重视。确实，需要是暂时的；时候将要来到，我们将毫无淫欲地彼此享受美丽——这情形将特别归于对造物主的赞美，如圣咏所说，祂\BibleQuote{披上了赞美与美丽}。\par

我如何能讲述其余的受造物，它们的一切美丽与实用，都是天主的仁慈赐予人，以悦其目、以供其用——尽管人已受定罪，被投入这些劳苦与苦难中？\newline{}我是否该谈论天空、大地和海洋那多样而迷人的美景？谈论光明那丰盛的供应与奇妙的特性？谈论太阳、月亮和星辰？谈论树木的荫蔽？谈论花朵的色彩与芬芳？谈论众多鸟类，它们羽毛与歌声各有不同？谈论动物的多样性，其中体型最小的往往最为奇妙——蚂蚁和蜜蜂的工效比鲸鱼的庞大躯体更令我们惊叹？\newline{}我是否该谈论海洋本身，它是一幅何等宏伟的景象，当它仿佛披上各色衣裳，时而呈现各种碧绿，时而又变成紫色或蓝色？在风暴中观看它，并因意识到我们自己未曾被颠簸、未曾失事而感到那种振奋的满足，岂不令人愉悦？\newline{}关于无数种缓解饥饿的食物，以及大自然随处散布、并非依靠烹饪艺术而得的刺激食欲的各种调味品，我还能说什么？有多少天然的物品用于保持和恢复健康！日夜交替是多么令人感激！使空气凉爽的微风是多么宜人！树木和动物为我们提供的衣物是多么丰盛！谁能尽数我们所享受的一切祝福？\newline{}如果我只试图详述并展开这些我粗略提及的少数几样，这样的列举也足以写满一卷书。而所有这些都只是受痛苦、被定罪之人的慰藉，并非蒙福之人的赏报。那么，如果被定罪的状态尚有如此祝福，那赏报将会是何等样？祂既已预定得永生的人赐予这些，那么祂将赐予那些祂预定死亡的人什么呢？祂既已愿意让祂的独生子为这些人——即使在这苦难的状态中——忍受苦难直至死亡，那么祂将在永生中为那些预定得永生的人倾注何等恩赐呢？\newline{}因此，宗徒论到那些预定得那国度的人说：\Quote{祂没有怜惜自己的儿子，反而为我们众人把祂交出，岂不也把一切与祂一同赐给我们吗？} \BibleRef{罗 8:32}\newline{}当这应许实现时，我们将是什么？我们将在那国度领受什么恩赐，既然我们已经领受了基督的死作为它们的抵押？人的精神将处于何种状态，当它不再有任何恶习；当它既不向任何恶习屈服，也不受任何恶习束缚，也不必与任何恶习争战，而是成全的，并享有与自身的平安？难道它那时不会确知一切，毫无辛劳与错误，不受阻碍、喜乐地畅饮智慧的活泉吗？\newline{}身体将如何，当它在各方面都完全服从精神，从精神那里获得足够丰富的生命，以致无需其他营养？因为它将不再是属生灵的，而是属神的，虽然具有肉体的实质，却没有丝毫肉体的腐朽。\par

% structure: p0093 source=h2 target=subsection reason=relative-heading-rank; base=h2

\subsection{第二十五章　论那些质疑肉身复活之人的执拗——尽管正如所预言的，整个世界已相信此事}

最杰出的哲学家们在来世蒙福之人所享有的属神幸福上与我们意见一致；只有肉身复活是他们质疑并竭力否认的。\newline{}但众人，无论有学问的或无学问的，世上的智者或愚者，都已相信，使不信者孤立无援，并归向了基督——祂以自己的复活证明了那些在我们对手看来荒谬之事的事实。因为世界相信了天主所预言的事，正如也曾预言世界会相信——这预言并非出于伯多禄的巫术，因它在如此久远之前已发出。\newline{}预言这些事的那位（如我早已说过，并乐于重复），是连其他诸神都在祂面前战栗的天主——波菲里本人也承认这一点，并试图通过来自这些神祇神谕的见证来证明，甚至称祂为天主圣父与君王。\newline{}我们切不可像那些不信的人一样解释这些预言——他们并未与整个世界一同相信那预言说世界将相信的事。因为我们为何不按世界的方式理解它们——其相信已被预言——而让那一小撮不信者去空谈、固执、孤立地不信呢？\newline{}如果他们主张，他们以不同方式解释预言，只是为了避免指责圣经为愚妄，从而冒犯那位他们为之作如此显著见证的天主，那么他们这样做，岂不是更严重地冒犯了祂？当他们说祂的预言必须被理解为与世界所相信的相反——而祂自己却称赞、应许并成就了世界这方面的信仰——这岂非更大的冒犯？\newline{}为什么祂不能使身体复活并永远活着？或者，我们不应该相信祂会这样做，因为这件事不受欢迎，且不配于天主？关于祂的全能（成就如此多伟大奇迹），我们已说得够多了。如果他们想知道全能者不能做什么，我要告诉他们：祂不能撒谎。让我们因此相信祂所能做的，拒绝相信祂所不能做的。拒绝相信祂会撒谎，就让他们相信祂会实现祂所应许的；让他们像世界那样相信——祂预言了世界的信德、称赞了它、应许了它、如今又指明了它。\newline{}但他们如何证明复活是不受欢迎的事呢？那时将不再有腐朽，这是身体唯一有害的东西。我已就元素的秩序以及其他人们提出的幻想式反对意见说得足够多了；在第十三卷中，我依我自己的判断，已充分说明了不朽身体运动的便捷性——这可以从我们即使现在身体健康时所体验到的轻松与活力来推断。那些没有读过前几卷书，或希望重温的人，可以自己阅读。\par

% structure: p0095 source=h2 target=subsection reason=relative-heading-rank; base=h2

\subsection{第26章——论\Person{波斐利}的意见（灵魂为了享真福，必须脱离一切种类的身体）被\Person{柏拉图}所粉碎，\Person{柏拉图}说至高天主曾应允众神，他们永不会被逐出自己的身体。}

但是，他们说，\Person{波斐利}告诉我们，灵魂为了享真福，必须脱离与一切种类身体的联系。因此，如果说未来的身体将是不朽的，这是无用的，因为灵魂若不能摆脱一切身体，就不能得真福。不过，我在上述书籍中已经充分讨论过这个问题。我只重复这一点——让他们的主人\Person{柏拉图}订正他的著作，并说他们的众神为了享真福，必须离开他们的身体，或换句话说，必须死亡；因为他曾说，众神被关在天体之内，然而，那造他们的天主却应许他们永生——即永久享有这些同样的身体，这并非自然所供给他们的，而仅仅是藉着祂旨意的进一步干预，从而他们可以确保幸福。在此，他显然推翻了他们的主张，即：身体的复活不能相信，因为这是不可能的；因为按照他的说法，当那非受造的天主应许受造的众神永生时，祂明确地说祂要做不可能的事。因为\Person{柏拉图}告诉我们，祂说：\Quote{你们既有起源，便不能永生和不朽；然而你们不会腐朽，也不会有任何命运毁灭你们或比我意志更强，我的意志将你们更确定地绑在永生上，胜过你们本性的束缚。}如果听见这些话的人——我们不是说理解，而是有耳朵——就不能怀疑\Person{柏拉图}相信天主曾应许祂所造的众神，祂将做一件不可能的事。因为那说\Quote{你们不能永生，但藉着我的意志，你们将是永生}的，岂非说\Quote{我将使你们成为你们所不能是的}？因此，身体将被复活为不朽的、永生的、属神的，由那位根据\Person{柏拉图}应许做不可能之事的天主来做。那么，他们为何仍声称，天主所应许的、世界依据预言的应许所相信的，是不可能的呢？因为我们所说的是，那位即使按照\Person{柏拉图}也行不可能之事的天主将做这事。那么，灵魂的真福不必脱离任何种类的身体，而要得到一个不朽的身体。又有什么不朽的身体比那当他们陷入可朽时曾为之叹息的身体更适合令他们喜悦呢？因为他们将不会再感受到那种\Person{维吉尔}模仿\Person{柏拉图}所归给他们的可怕渴望，即希望回到他们的身体。我再说，他们不会感到这种回到身体的渴望，因为他们将拥有那些他们曾渴望回归的身体，并将如此完全地拥有它们，以致永远不会失去它们，哪怕一刻，也永不再在死亡中放下它们。\par

% structure: p0097 source=h2 target=subsection reason=relative-heading-rank; base=h2

\subsection{第27章——论\Person{柏拉图}与\Person{波斐利}表面冲突的意见；倘若他们能从彼此让步，两者都会被引向真理。}

\Person{柏拉图}与\Person{波斐利}各人曾提出一些见解，倘若他们能够看到共同持有这些见解的路，他们或许可能成为基督徒。\Person{柏拉图}说，灵魂不能没有身体而永远存在；因为正是基于这个原因，他说，甚至智者之魂也必须在一定时候返回它们的身体。\Person{波斐利}又说，净化了的灵魂，当它回到圣父那里时，永不再回到这世界的苦患。因此，如果\Person{柏拉图}将他所见到的真理（即灵魂即使完全净化，属于智者与义者，也必须返回人的身体）传授给\Person{波斐利}，而\Person{波斐利}又将他所见到的真理（即圣洁的灵魂永不再回到可朽身体的痛苦）传授给\Person{柏拉图}，使得他们不各坚持己见，而是二者共同持有两个真理——我想，他们将会看到，结论是：灵魂要返回它们自己的身体，而且这些身体将是如此，以致赋予它们一种有福的永生。因为按照\Person{柏拉图}，甚至圣洁的灵魂也要返回身体；按照\Person{波斐利}，圣洁的灵魂不会回复这世界的苦患。让\Person{波斐利}随同\Person{柏拉图}说：它们要返回身体；让\Person{柏拉图}随同\Person{波斐利}说：它们不会返回从前的痛苦；并且他们将会同意，它们返回的身体中不再受苦。而这正是天主所应许的——祂将赐予灵魂与它们自己身体结合的永恒幸福。因为，我想，两人都会欣然承认：如果圣人的灵魂要与身体重新结合，那必是与它们自己的身体结合，就是那些他们在此世忍受苦难的身体，并为了逃避这些苦难，他们曾以虔诚和忠诚事奉了天主。\par

% structure: p0099 source=h2 target=subsection reason=relative-heading-rank; base=h2

\subsection{第28章——\Person{柏拉图}、\Person{拉贝奥}或甚至\Person{瓦罗}，倘若他们采纳彼此的意见合成一体，可能对复活的真实信德作出何等贡献。}

有些基督徒因着\Person{柏拉图}文体之华丽及他偶尔吐露的真理而喜爱他，竟说他对死人复活也持有与我们相似的见解。然而，\Person{西塞罗}在他的\WorkTitle{共和国}中提及此点时，断言\Person{柏拉图}更愿意将其视为一种游戏般的幻想而非现实；因为他引入一个死而复生的人，叙述其经历以证实\Person{柏拉图}的学说。\Person{拉贝奥}也说，有两人在同一天死去，在一个十字路口相遇，后来奉命返回他们的身体，他们便约定终生为友，直至再次死亡。但这些作者所举的复活实例，类似于我们自己所知的一些复活之人，他们确实回到此生，却并非不再死亡。不过，\Person{马库斯·瓦罗}在他的\WorkTitle{论罗马人的起源}中记载了更奇特的事；我认为应当引用他本人的话。\Quote{某些占星家，}他说，\Quote{写道：人类将注定得到一种新生，希腊人称之为\OriginalTerm{重生}{\GreekText{παλιγγενεσία}}。此事将在四百四十年后发生；届时，从前曾在人身上结合过的同一个灵魂和同一个身体，将再次结合在一起。}这位\Person{瓦罗}，或者那些无名的占星家——因为他没有引述其言论之人的名字——所宣称的确实并非全然正确；因为灵魂一旦回到它们曾穿过的身体，就永不再离开它们。然而，他们的话推翻并摧毁了我们的对手关于复活不可能性的许多空谈。因为持有或曾经持有这种看法的人，并不认为那些已经分解为空气、尘土、灰烬、水，或进入吞食它们的野兽甚至人体中的身体，能够恢复为它们原先所是的样子。因此，如果\Person{柏拉图}和\Person{波菲利}，或者不如说他们现今的门徒，同意我们：圣洁的灵魂要返回身体（如\Person{柏拉图}所说），而且它们不再返回苦难（如\Person{波菲利}所主张），——如果他们接受基督教信仰所教导的这两个命题的推论，即它们将获得身体，在其中永远生活而不受任何苦难——那么，也让他们从\Person{瓦罗}那里接受这样的意见：它们将返回原先所在的身体，这样，关于身体永恒复活的整个问题就可从他们自己的口中得到解决。\par

% structure: p0101 source=h2 target=subsection reason=relative-heading-rank; base=h2

\subsection{第二十九章 论荣福直观}

现在，让我们按天主可能赐予的能力，思考当圣人们穿上不朽的、属神的身体，且血肉不再按属血肉的方式而按属神的方式生活时，他们将从事什么。事实上，老实说，我对于那种从事的性质——或者不如说安息与安逸——不得而知，因为它从未进入我身体感官的范围。若我谈论我的心灵或理智，我们的理智与那种卓越相比又算什么呢？因为那时将有\BibleQuote{天主的平安，}如宗徒所说，\BibleQuote{超越一切理智，}\BibleRef{斐 4:7}——即超越一切人的，或许也超越一切天使的理智，但当然不超越天主的。它超越我们的理智是无疑的；但若它也超越天使的理智——而说\BibleQuote{一切理智}的人似乎没有将天使除外——那么我们就必须理解他是指：我们和天使都不能像天主所理解的那样，理解天主自身所享有的平安。这无疑超越除祂自己外的一切理智。然而，因为我们将来有一天会按我们有限的能力，在自身内、与邻人、与我们的至善天主一起，分享祂的平安，就此而言，天使按他们的尺度理解天主的平安，而人类虽远逊于他们，但无论他们在灵性上有了多大进步，也能理解。因为我们必须记住，那位伟大的圣者曾说：\BibleQuote{我们现在所知道的，只是局部的；我们作先知所讲的，也只是局部的；及至那完全的来到，局部的就必归于无有了。}\BibleRef{格前 13:9}又说：\BibleQuote{现在我们对着镜子观看，模糊不清；到那时，就要面对面了。}\BibleRef{格前 13:12}圣洁的天使的观看也是如此，他们也被称为我们的天使，因为我们被救出黑暗的权势，领受了圣神的凭据\OriginalTerm{圣神的凭据}{\GreekText{ἀρραβών}}，被迁到基督的王国，并已经开始属于那些天使，我们将与他们一起享受那神圣而最令人喜悦的天主之城，对此我们已经写了许多。因此，天主的天使是我们的天使，正如基督是天主的，也是我们的。他们是天主的，因为他们没有离弃祂；他们是我们的，因为我们是他们的同乡。主耶稣也说：\BibleQuote{你们小心，不要轻视这些小子中的一个；我告诉你们：他们在天上的天使，时常看见我在天之父的面。}\BibleRef{玛 18:10}因此，他们怎样看见，我们也将要怎样看见；但我们现今还不这样看见。所以，宗徒用了刚才引用的那些话：\BibleQuote{现在我们对着镜子观看，模糊不清；到那时，就要面对面了。}这观看被保留作为我们信德的赏报；关于它，宗徒若望也说：\BibleQuote{祂若显现，我们必要相似祂，因为我们要看见祂实在怎样。}\BibleRef{若一 3:2}所谓天主的\Quote{面}，我们应理解为祂的显示，而不是像我们身体中称为面的那部分身体。\par

因此，当有人问我，圣人们在那属神身体中将如何行动时，我不是说自己看见的，而是说自己相信的，依照我在圣咏中所读的：\Quote{我信了，所以我所说的话。}我于是说：他们将在身体内看见天主；但他们是否藉着身体看见祂，如同现在我们看见日、月、星辰、海、地及其中一切那样，这是一个困难的问题。因为很难说圣人们那时会有如此的身体，以致不能随意闭上和睁开他们的眼睛；而更难说的是，每一个闭上眼睛的人都会失去天主的神视。因为如果先知\Person{厄里叟}虽在远处，却看见了他的仆人\Person{革哈齐}——那仆人以为自己的恶行能逃过主人的眼目，接受了\Place{叙利亚}人\Person{纳阿曼}的礼物，而先知曾医治了他的可怕癞病——那么，圣人们在属神身体中岂不更能看到一切事物，不仅当他们的眼睛闭上，甚至当他们自己在远处时也是如此？因为那时将成为\Quote{那成全的}，宗徒论此说：\Quote{我们现在所知道的，只是局部的；我们作先知，也是局部的；但当那成全的一来到，局部的，就必要消逝。}然后，为尽可能以比喻说明未来的生命比现今的生活——不仅普通人的生活，甚至圣人们中的杰出者——更加优越，他说：\BibleQuote{我做孩子的时候，说话像孩子，思想像孩子，判断像孩子；但一成了人，就把孩子的事丢掉了。我们现在是藉着镜子观看，模糊不清；但那时，就要面对面了；我现在所知道的，只是局部的；但那时，我就要完全知道，如同我被完全知道一样。}\BibleRef{格前 13:11}那么，如果在这种生命中——即使是在这生命中，杰出人物的先知能力与未来生命的神视相比，就像童年与成年的差别一样——\Person{厄里叟}虽远离他的仆人，却看见了他接受礼物；我们能否说：当那成全的来到，可朽坏的身体不再压迫灵魂，而是成为不朽的，不造成任何阻碍时，圣人们却还需要肉体的眼睛去看，而\Person{厄里叟}却不需要它们去看他的仆人？因为，按\WorkTitle{七十贤士译本}，先知的话是：\BibleQuote{当那人从战车上下来迎接你的时候，我的心不是同你一起去了吗？你拿了他的礼物。}\BibleRef{列下 5:26}或者，如同司铎\Person{热罗尼莫}从希伯来文所译的：\BibleQuote{当那人从战车上转身迎接你的时候，我的心不是在你那里吗？}先知说他是用心看到的，并藉着天主的奇妙帮助，这是无人能否认的。但当天主成为万有中的万有时，圣人们更将如何丰富地享有此恩赐呢？然而，肉体的眼睛也将有其职务和位置，并由灵魂藉着属神身体来使用。因为先知并没有放弃用他的眼睛去看眼前的事物，虽然他不需要眼睛去看他不在场的仆人，而且他本可以用精神、闭上眼睛去看这些眼前的对象，正如他看见远处自己不在场的事物一样。因此，我们绝不能断言，在将来的生命中，圣人们闭上眼睛时看不见天主，因为他们将常以精神看见祂。\par

但问题出现了：当他们的眼睛睁开时，是否以肉体的眼睛看见祂？如果属神身体的眼睛不比我们现在拥有的眼睛有更大的能力，那么显然，不能藉着它们看见天主。它们必须具有非常不同的能力，才能注视那非物质的、不被任何地方所包含、却在每一处都是全部的本性。因为尽管我们说天主在天上并在地上，如祂自己藉先知所说：\BibleQuote{我充满了天地。}\BibleRef{耶 23:24}我们并不是说天主有一部分在天上，另一部分在地上；而是祂全部在天上，全部在地上，不是轮流交替，而是同时如此，如同任何有形体的存有所不能做到的。那么，眼睛将拥有无比优越的能力——不是那种归诸蛇或鹰的锐利视力，因为这些动物无论看得多么敏锐，除了有形质的事物外什么也不能分辨——而是看见非物质的灵性事物的能力。或许正是这种伟大的视力，暂时赋予了\Person{圣约伯}的眼目，当他仍在可朽坏的身体中时，他向天主说：\BibleQuote{我以前只听见了有关祢的事，但现在我亲眼看见了祢；为此我憎恶我自己，在尘埃和灰烬中忏悔。}\BibleRef{约 42:5}虽然我们也没有理由不将之理解为心灵的眼目，宗徒论此说：\BibleQuote{照亮你们心灵的眼目。}\BibleRef{弗 1:18}但是，凡以信德接受我们的天主和师父所说的话的基督徒，无人怀疑天主将用这些眼睛被看见：\BibleQuote{心里洁净的人是有福的，因为他们要看见天主。}\BibleRef{玛 5:8}但是，在将来的生命中，天主是否也将以肉体的眼睛被看见，这正是我们现在的问题。\par

圣经的话说：\BibleQuote{凡有血肉的，都要看见天主的救恩。}\BibleRef{路 3:6} 可以毫无困难地理解为相当于说：\Quote{凡人都要看见天主的基督。} 祂当然曾在肉身内被看见，将来也要在肉身内被看见，当祂审判活人死人之时。而且基督就是天主的救恩，还有许多别的圣经章节为之作证，尤其是可敬的西默盎的话，当他将婴孩基督抱在手中时说：\BibleQuote{主啊！现在可照祢的话，让祢的仆人平安去了；因为我亲眼看见了祢的救恩。}\BibleRef{路 2:29} 至于前面提到的\Person{约伯}的话，正如在希伯来文抄本中发现的：\Quote{在我肉体内我要看见天主}，无疑这是肉身复活的预言；但他并没有说\Quote{藉着肉身}。的确，即使他这样说了，仍然有可能\Quote{天主}指的是基督；因为基督是要由肉身在肉身内被看见。但即使理解为天主，这也不过等于说：当我看见天主时，我将是在肉身内。然后宗徒的话：\BibleQuote{面对面}\BibleRef{格前 13:12}，并不迫使我们相信我们将用肉身的脸（其中有肉身的眼睛）看见天主，因为我们将在精神内不停地看见祂。如果宗徒不是指内在的人的脸，他就不会说：\BibleQuote{我们众人，以揭开的脸面反映主的光荣，逐渐被转化成同一肖像，光荣上加光荣，这乃是由主圣神所成就的。}\BibleRef{格后 3:18} 同样的意思，我们理解圣咏作者所唱的：\Quote{你们亲近祂，并要光照；你们的面容，决不致羞惭。} 因为我们藉着信德亲近天主，而信德是精神的行为，不是肉体的行为。但我们不知道灵性的身体要达到何等程度的完美——因为我们在此谈论的是我们所没有经验过的事，而且圣经的权威对此也没有明确宣告——因此，需要\WorkTitle{智慧篇}的话在我们身上实现：\BibleQuote{必朽之人的思念何等怯懦，我们的计谋多么不可靠！}\BibleRef{智 9:14}\par

因为如果哲学家的那种推理——他们试图证明，可理解的或精神的对象由心灵以这种方式看见，可感觉的或肉身的对象由身体以那种方式看见，以至于前者不能藉着身体被心灵辨识，后者也不能不藉着身体被心灵本身辨识——如果这种推理是可信的，那么必然得出结论：甚至连灵性的身体的眼睛也不能看见天主。但这一推理已被正确的理性和先知的权威所驳倒。因为谁与真理如此疏远，竟敢说天主对可感觉的对象毫无认知？那么祂是否有一个身体，其眼睛给祂这种知识呢？再者，我们刚才有关先知\Person{厄里叟}的叙述，岂不足够表明肉身的事物可以不靠身体的帮助而被精神所辨识吗？因为当那个仆人接受礼物时，这当然是肉身或物质的事，但先知不是藉着身体，而是藉着精神看见的。因此，既已同意身体由精神所看见，那么如果灵性的身体的力量如此之大，以致精神也由身体所看见，那又如何呢？因为天主是神。此外，每个人以内在感觉，而不是以肉体的眼睛，认识他自己的生命——即他现在身体内所活着的生命，它赋予这些尘世的肢体以生命并使之生长——但其他人的生命，虽然不可见，他却用肉体的眼睛看见。因为我们如何区分活的身体和死的身体，除非同时看见身体和生命，而这生命除了藉着眼睛我们无法看见？但一个没有身体的生命，我们不能这样看见。\par

因此，很可能而且完全可信的是，在将来的世界中，我们将以这样的方式看见新天新地的物质形体，以至于我们将最清晰地认出天主处处临在并治理一切事物，物质的和精神的，并且看见祂，不像现在藉着受造之物了解天主不可见的事物\BibleRef{罗 1:20}，模糊地如在镜中观看，而且只是部分地，更多地是藉着信德而非肉身的视觉对物质现象的认识，而是藉着我们将来要穿着的身体，并且无论我们转目何处，我们都会看见这身体。正如我们不凭信德，而是看见我们周围活着的人正在行使生命功能，我们知道他们是活着的，虽然我们不能没有他们的身体而看见他们的生命，却最清晰地藉着他们的身体看见它，同样，无论我们将来用我们身体的这些灵性的眼睛看向何处，我们那时也将藉着物质实体看见天主，虽然祂是神，却统治万有。因此，要么眼睛将具有某种类似心灵的性质，藉此能够辨别属灵的事物，包括天主——这一假设很难甚至根本不可能在圣经中找到支持——要么，更易于理解的是，天主将被我们如此认识，并且如此呈现在我们面前，以至于我们将在自己内、在彼此内、在祂自己内、在新天新地内、在一切那时存在的受造物内藉着精神看见祂；并且也藉着身体，在灵性身体的眼睛敏锐视力所及的每一个身体内看见祂。我们的思念也将为众人所见，因为那时宗徒的话将应验：\BibleQuote{时候未到，你们什么也不要判断，只等主来，祂要揭发黑暗中的隐情，显露人心中的意念；那时，各人才会从天主那里得到称赞。}\BibleRef{格前 4:5}\par

% structure: p0108 source=h2 target=subsection reason=relative-heading-rank; base=h2

\subsection{第三十章 论天主之城永恒的福乐与永久安息日。}

那幸福将是何等伟大，它不被任何邪恶所玷污，不缺少任何美善，并且提供闲暇来赞美天主，天主将是万有之中的万有！因为我不知道还有什么其他工作，在那里不会有疲倦松弛活动，也不会有任何缺乏刺激劳动。我也受到圣歌的提醒，在其中我读到或听到这些话：\BibleQuote{住在你殿中的，是有福的；他们仍要赞美你。}不朽身体的所有肢体和器官，我们现在看到它们适合各种必要的用途，都将有助于赞美天主；因为在那个生命中，必要性将没有位置，只有圆满、确定、安全、永恒的幸福。对于身体和谐的所有那些部分，它们分布在整个身体内外，我刚才说它们目前逃避我们的观察，那时将被辨别出来；并且，连同其他伟大而奇妙的发现，这些发现将点燃理性心灵去赞美那位伟大的工匠，那时还将有对理性美的享受。这些身体将拥有怎样的运动能力，我不敢鲁莽地定义，因为我没有能力构想。然而我要说，无论如何，在运动或静止中，它们都将是，如同在它们的外表上，得体的；因为没有什么不得体的东西会被允许进入那个状态。有一件事是确定的：身体将立刻到达精神所愿意的任何地方，而精神将不愿意任何对精神或身体不体面的事。真正的尊荣将在那里，因为它不会被拒绝给任何配得的人，也不会让给任何不配的人；也没有任何不配的人会去追求它，因为只有配得的人会在那里。真正的和平将在那里，没有人会遭受来自自己或他人的反对。天主自己，是德行的创造者，将在那里成为德行的赏报；因为没有什么比天主更大或更好，祂应许了祂自己。祂通过先知所说的话还有什么别的意思：\BibleQuote{我要作你们的天主，你们要作我的子民，}\BibleRef{肋 26:12}不就是：我将成为他们的满足，我将成为人们所光荣渴望的一切——生命、健康、滋养、丰足、荣耀、尊荣、和平，以及一切美善？这也是对宗徒所说\BibleQuote{叫天主成为万物之中的万有。}\BibleRef{格前 15:28}的正确解释。祂将是我们渴望的终结，我们将无尽头地看见祂，不厌腻地爱祂，不疲倦地赞美祂。这种情感的流露，这种工作，当然将像永生本身一样，为所有人共有。\par

但是谁能构想，更不用说描述，何种程度的尊荣和光荣将赏给不同等级的功德？然而不能怀疑那里会有等级之分。在那有福的城里，将有这一伟大的祝福：没有下属会嫉妒任何上级，就如现在总领天使不被天使嫉妒一样，因为没有人会愿意成为他未曾领受的，尽管与领受者处于最紧密的和谐之中；如同在身体中，手指不寻求成为眼睛，尽管两个肢体都被和谐地包括在身体的完整结构中。因此，伴随着他的恩赐，无论大小，每个人还将领受这一进一步的恩赐：满足于不渴望超出他所拥有的。\par

我们也不应以为，因为罪没有能力使他们快乐，自由意志就必须被收回。相反，它将更加真正地自由，因为从以犯罪为乐中解放出来，而转向以不犯罪为永不衰竭的快乐。因为当人受造正直时所接受的第一种自由意志，在于有能力不犯罪，也有能力犯罪；而最后这种自由意志将更卓越，因为它将不能犯罪。这确实不是一种自然能力，而是天主的恩赐。因为成为天主是一回事，分享天主是另一回事。天主按其本性不能犯罪，但分享天主的人从天主那里领受了这种不能犯罪的能力。在这神圣的恩赐中，可以看到这样的等级：人首先领受一种自由意志，使他能够不犯罪，最后领受一种自由意志，使他不能犯罪——前者适于获取功德，后者适于享受赏报。但如此构成的本性，在有能力犯罪时犯了罪，正是通过更丰富的恩宠，它才被拯救出来，达到那不能犯罪的自由。因为正如亚当因犯罪失去的第一种不死性在于他能够不死，而最后一种不死性将在于他不能死；同样，第一种自由意志在于他能够不犯罪，最后一种在于他不能犯罪。因此，虔诚和正义将像幸福一样不可剥夺。因为确实，我们因犯罪而失去了虔诚和幸福；但当我们失去了幸福，我们并没有失去对它的爱。我们岂能说天主自己因不能犯罪而不自由吗？那么，在那城里，将有自由意志，在全体公民中是同一的，在每个人中是不可分割的，从一切邪恶中解放出来，充满一切美善，不可剥夺地享受永恒喜乐的欢愉，忘却罪恶，忘却痛苦，然而并没有忘记它的救赎，以致对它的救赎者忘恩负义。\par

灵魂将拥有对其过去苦难的理智记忆；但就感性经验而言，它们将被完全遗忘。因为一个高明的医生，确实在专业上几乎知道所有疾病；但就经验而言，他对自己从未患过的许多疾病一无所知。因此，正如认识恶事有两种方式——一种是通过理智洞察，另一种是通过感性经验，因为藉受过熏陶的心灵的智慧理解一切恶习是一回事，藉堕落生活的愚昧理解它们则是另一回事——同样，遗忘恶事也有两种方式。一个受过良好教育且有学识的人以一种方式遗忘它们，一个曾因经验而受苦的人以另一种方式遗忘它们——前者是通过忽略他所学的，后者是通过逃脱他所受的苦。圣人将以后一种方式遗忘他们过去的苦难，因为他们将完全摆脱它们，以至于它们将完全从他们的经验中抹去。然而他们伟大的理智知识，不仅将使他们了解自己过去的苦难，还将使他们了解沉沦者永久的痛苦。因为他们若不知道自己是可怜的，又怎能像圣咏作者所说的那样永远歌颂天主的仁慈呢？那城确实没有比庆祝基督的恩宠、祂用祂的血救赎我们更大的喜乐了。那时要应验圣咏的话：\Quote{你们要静止，要知道我是天主。}那里要有伟大的安息日，没有夜晚，天主在祂最初的作为中庆祝了这安息日，正如经上所记：\BibleQuote{天主在第七天完成了祂所做的一切工作，就在第七天安息了。天主降福了第七天，祝圣了它，因为在那天祂安息了，停止了祂所做的一切工作。}\BibleRef{创 2:2}因为我们自己将成为第七天，当我们充满并饱享天主的祝福和圣化时。在那里我们将静止，知道祂是天主；祂是我们自己当背弃祂时曾渴望成为的，我们听从了诱惑者的声音：\BibleQuote{你们将如同神一样，}\BibleRef{创 3:5}于是离弃了天主，而天主本要使我们成为神，不是通过离弃祂，而是通过参与祂。因为离开祂，我们除了在祂的震怒中灭亡外，还能成就什么呢？但当我们被祂恢复，并以更大的恩宠得以成全时，我们将有永恒的闲暇来认识祂是天主，因为我们将在祂成为万有中的万有时充满祂。因为即使是我们的善工，当它们被理解为属于祂而非属于我们时，也被算作我们的，好使我们能享受这安息日的安息。因为我们若将它们归于自己，它们便成为奴役的；因为关于安息日经上说：\BibleQuote{你不可在那天作任何劳役的工作。}\BibleRef{申 5:14}因此厄则克耳先知也说：\BibleQuote{我将我的安息日赐给他们，作为我与他们之间的记号，好叫他们知道我是祝圣他们的主。}\BibleRef{则 20:12}这种知识将在我们完全安息并完全知道祂是天主时得以成全。\par

如果我们按照圣经所定义的时间段，把时代当作日子来计算，这个安息日就会显得更加清楚，因为那个时代将被发现是第七个。第一个时代，如同第一天，从亚当延续到洪水；第二个时代，从洪水到亚巴郎，与第一个时代在时间长度上不等，但在世代数目上相等，每个都有十代。从亚巴郎到基督的来临，按圣史玛窦的计算，有三个时期，每个时期有十四代——一个时期从亚巴郎到达味，第二个时期从达味到充军，第三个时期从充军到基督按其肉身诞生。这样共有五个时代。第六个时代正在过去，无法用任何世代数目来衡量，正如经上所说：\BibleQuote{父用自己的权柄所定的时间和时期，不是你们应当知道的。}\BibleRef{宗 1:7}在这时期之后，天主将如同第七天一样安息，那时祂将把我们（即第七天）安息在祂自己内。但现在没有篇幅来讨论这些时代；只需说第七天将是我们安息日，它的结束不是黄昏，而是主的日子，作为第八个也是永恒的一天，由基督的复活祝圣，并预表不仅灵魂而且身体的永恒安息。在那里我们将安息并看见，看见并爱慕，爱慕并赞美。这就是那将要在无终的终结中所成就的。因为我们给自己定什么别的目标，若不是达到那没有终点的王国呢？\par

我想，我现在已藉天主的帮助，完成了写作这部巨著的义务。让那些认为我说得太少的人，以及那些认为我说得太多的人，原谅我；并让那些认为我说得恰好的人，与我一同感谢天主。阿门。\par

\SourceRecord{Translated by Marcus Dods.；Nicene and Post-Nicene Fathers, First Series；Vol. 2.；Edited by Philip Schaff.；Buffalo, NY: Christian Literature Publishing Co.,；1887.；<http://www.newadvent.org/fathers/120122.htm>.}
